\documentclass[11pt]{article}
\usepackage[utf8]{inputenc}
\usepackage[T1]{fontenc}
\usepackage[a4paper,margin=1in]{geometry}
\usepackage{lmodern}
\usepackage{microtype}
\usepackage{amsmath,amssymb,mathtools}
\usepackage{mathrsfs}
\usepackage{tikz-cd}
\usepackage{enumitem}
\usepackage{longtable,booktabs,array}
\usepackage[hidelinks]{hyperref}
\usepackage{hyperxmp}
\newcommand{\sgaStableTarget}[1]{\hypertarget{#1}{}}
\newcommand{\sgaStableTargetInline}[1]{\pdfdest name {#1} xyz\relax}
\setlength{\parindent}{0pt}
\setlength{\parskip}{0.55em}
\emergencystretch=3em
\hbadness=10000
\hfuzz=20pt
\newcommand{\K}{K}
\newcommand{\Gr}{\operatorname{Gr}}
\newcommand{\Fil}{\operatorname{Fil}}
\newcommand{\ch}{\operatorname{ch}}
\newcommand{\Todd}{\operatorname{Todd}}
\newcommand{\cl}{\operatorname{cl}}
\newcommand{\Parf}{\operatorname{Parf}}
\newcommand{\Tor}{\operatorname{Tor}}
\newcommand{\Spec}{\operatorname{Spec}}
\newcommand{\RRR}{\mathrm{RRR}}
\providecommand{\codim}{\operatorname{codim}}
\providecommand{\Inv}{\operatorname{Inv}}
\providecommand{\Pic}{\operatorname{Pic}}
\providecommand{\Z}{\mathbb Z}
\providecommand{\Q}{\mathbb Q}
\providecommand{\R}{\mathbb R}
\providecommand{\C}{\mathbb C}
\providecommand{\Ahat}{\widehat A}
\providecommand{\Ahp}{\widehat A^{+}}
\providecommand{\That}{\widehat T}
\providecommand{\rad}{\operatorname{rad}}
\providecommand{\Gl}{\operatorname{Gl}}
\providecommand{\GL}{\operatorname{GL}}
\providecommand{\ch}{\operatorname{ch}}
\providecommand{\Todd}{\operatorname{Todd}}
\providecommand{\Gr}{\operatorname{Gr}}

\providecommand{\Atil}{\widetilde A}
\providecommand{\Ctil}{\widetilde C}
\providecommand{\K}{\mathrm K}
\providecommand{\cC}{\mathcal C}
\providecommand{\cF}{\mathscr F}
\providecommand{\eps}{\varepsilon}
\providecommand{\ellc}{\ell}
\providecommand{\Tor}{\operatorname{Tor}}
\providecommand{\Coker}{\operatorname{Coker}}
\providecommand{\Ker}{\operatorname{Ker}}
\providecommand{\Supp}{\operatorname{Supp}}
\providecommand{\codim}{\operatorname{codim}}
\providecommand{\Gr}{\operatorname{G}}
\long\def\SGAArchiveOnly#1{}
\providecommand{\rank}{\operatorname{rank}}
\providecommand{\rang}{\operatorname{rang}}
\providecommand{\cO}{\mathcal O}

\title{SGA 6: Intersection Theory and the Riemann--Roch Theorem}
\author{P. Berthelot, A. Grothendieck, L. Illusie}
\date{22 July 2026}

\hypersetup{
  hypertexnames=false,
  pdftitle={SGA 6: Intersection Theory and the Riemann--Roch Theorem -- English Translation},
  pdfauthor={P. Berthelot, A. Grothendieck, L. Illusie},
  pdfsubject={English translation of SGA 6},
  pdfkeywords={SGA 6, intersection theory, Riemann--Roch, English translation}
}
\providecommand{\N}{\mathbb N}
\providecommand{\Zell}{\mathbb Z_\ell}
\providecommand{\Qell}{\mathbb Q_\ell}
\providecommand{\C}{\mathbb C}
\providecommand{\Ch}{\operatorname{Ch}}
\providecommand{\Sym}{\operatorname{Sym}}
\providecommand{\Td}{\operatorname{Td}}

\begin{document}
%%%%%%%%%%  scan idx 0  |  folio cover  %%%%%%%%%%
\SGAPrintPage{cover}%
\begin{center}
{\Large Lecture Notes in Mathematics}\par
A collection of informal reports and seminars\par
Edited by A. Dold, Heidelberg and B. Eckmann, Zürich\par
\bigskip
225\par
\bigskip
Algebraic Geometry Seminar of Bois Marie 1966/67\par
SGA 6\par
Directed by\par
P. Berthelot, A. Grothendieck, L. Illusie\par
With the collaboration of\par
D. Ferrand, J.P. Jouanolou, O. Jussila,\par
S. Kleiman, M. Raynaud, J.P. Serre\par
\bigskip
{\Large Intersection Theory\par
and the Riemann--Roch Theorem}\par
\bigskip
Springer-Verlag\par
Berlin \textperiodcentered{} Heidelberg \textperiodcentered{} New York
\end{center}

%%%%%%%%%%  scan idx 1  |  folio cover-verso  %%%%%%%%%%
\SGAPrintPage{cover-verso}%
% The contemporary-template restitution is documented in the new edition preface and source records.

% The masked notice is restored from the contemporary primary template
% in Springer LNM 228 (1971), front-matter PDF page 1; see the separate
% ASTRA-FRONT-NOTICE-1971-RESTORATION evidence packet. No obscured stroke
% in the supplied SGA 6 image is claimed to have become directly readable.
This series aims to report new developments in mathematical research and teaching -- quickly, informally
and at a high level. The type of material considered for publication includes:
\begin{enumerate}
\item Preliminary drafts of original papers and monographs
\item Lectures on a new field, or presenting a new angle on a classical field
\item Seminar work-outs
\item Reports of meetings
\end{enumerate}
Texts which are out of print but still in demand may also be considered if they fall within these categories.

The timeliness of a manuscript is more important than its form, which may be unfinished or tentative.
Thus, in some instances, proofs may be merely outlined and results presented which have been or will
later be published elsewhere.

Manuscripts should comprise not less than 100 pages.

Publication of \textit{Lecture Notes} is intended as a service to the international mathematical community, in
that a commercial publisher, Springer-Verlag, can offer a wider distribution to documents which would
otherwise have a restricted readership. Once published and copyrighted, they can be documented in the
scientific literature. --

\subsection*{Manuscripts}
Manuscripts are reproduced by a photographic process; they must therefore be typed with extreme care.

Symbols not on the typewriter should be inserted by hand in indelible black ink.
Corrections to the typescript should be made by sticking the amended text over
the old one, or by obliterating errors with white correcting fluid. Authors
receive 75 free copies.

The typescript is reduced slightly in size during reproduction; best results
will not be obtained unless the text on any one page is kept within the overall
limit of $18\times26.5$ cm ($7\times10\frac12$ inches). The publishers will be
pleased to supply on request special stationery with the typing area outlined.

Manuscripts in English, German or French should be sent to Prof. Dr. A. Dold,
Mathematisches Institut der Universität Heidelberg, 69 Heidelberg/Germany,
Tiergartenstraße or Prof. Dr. B. Eckmann, Eidgenössische Technische Hochschule,
CH-8006 Zürich/Switzerland.

\subsection*{Lecture Notes in Physics}
\textbf{Bisher erschienen/Already published}

Vol. 1: J. C. Erdmann, Wärmeleitung in Kristallen, theoretische Grundlagen und
fortgeschrittene experimentelle Methoden. II, 283 Seiten. 1969. DM 20,--

Vol. 2: K. Hepp, Théorie de la renormalisation. III, 215 pages. 1969. DM 18,--

Vol. 3: A. Martin, Scattering Theory: Unitarity, Analytic and Crossing.
IV, 125 pages. 1969. DM 14,--

Vol. 4: G. Ludwig, Deutung des Begriffs physikalische Theorie und axiomatische
Grundlegung der Hilbertraumstruktur der Quantenmechanik durch Hauptsätze des
Messens. XI, 469 Seiten. 1970. DM 28,--

Vol. 5: M. Schaaf, The Reduction of the Product of Two Irreducible Unitary
Representations of the Proper Orthochronous Quantummechanical Poincaré Group.
IV, 120 pages. 1970. DM 14,--

Vol. 6: Group Representations in Mathematics and Physics. Edited by V. Bargmann.
V, 340 pages. 1970. DM 24,--

Vol. 7: R. Balescu, J. L. Lebowitz, I. Prigogine, P. Résibois, Z. W. Salsburg,
Lectures in Statistical Physics. V, 181 pages. 1971. DM 18,--

Vol. 8: Proceedings of the Second International Conference on Numerical Methods
in Fluid Dynamics. Edited by M. Holt. IX, 462 pages. 1971. DM 28,--

%%%%%%%%%%  scan idx 2  |  folio I  %%%%%%%%%%
\SGAPrintPage{I}%
\begin{center}
{\Large Lecture Notes in Mathematics}\par
A collection of informal reports and seminars\par
Edited by A. Dold, Heidelberg and B. Eckmann, Zürich\par
\bigskip
225\par
\bigskip
{\Large Intersection Theory\par
and the Riemann--Roch Theorem}\par
\bigskip
Springer-Verlag\par
Berlin \textperiodcentered{} Heidelberg \textperiodcentered{} New York 1971
\end{center}

%%%%%%%%%%  scan idx 3  |  folio II  %%%%%%%%%%
\SGAPrintPage{II}%
AMS Subject Classifications (1970): 14-02, 14C15, 14F15, 14K30, 18E30

ISBN 3-540-05647-5 Springer-Verlag Berlin \textperiodcentered{} Heidelberg \textperiodcentered{} New York\par
ISBN 0-387-05647-5 Springer-Verlag New York \textperiodcentered{} Heidelberg \textperiodcentered{} Berlin

This work is subject to copyright. All rights are reserved, whether the whole
or part of the material is concerned, specifically those of translation,
reprinting, re-use of illustrations, broadcasting, reproduction by photocopying
machine or similar means, and storage in data banks.

Under \S~54 of the German Copyright Law where copies are made for other than
private use, a fee is payable to the publisher, the amount of the fee to be
determined by agreement with the publisher.

\textcopyright{} by Springer-Verlag Berlin \textperiodcentered{} Heidelberg 1971.
Library of Congress Catalog Card Number 73-178287. Printed in Germany.

Offsetdruck: Julius Beltz, Hemsbach/Bergstr.


%%%%%%%%%%  scan idx 4  |  folio III  %%%%%%%%%%
\SGAPrintPage{III}%
\begin{center}
{\Large Algebraic Geometry Seminar of Bois Marie}\par
{\large 1966--67}\par
\vspace{1.5em}
{\Large Intersection Theory and the Riemann--Roch Theorem}\par
{\large (SGA 6)}\par
\vspace{1.5em}
a Seminar directed by\par
P. Berthelot, A. Grothendieck, L. Illusie\par
with the collaboration of\par
D. Ferrand, J.-P. Jouanolou, O. Jussila, S. Kleiman, M. Raynaud, J.-P. Serre
\end{center}

%%%%%%%%%%  scan idx 5  |  folio IV  %%%%%%%%%%
\SGAPrintPage{IV}%
%%%%%%%%%%  scan idx 6  |  folio V  %%%%%%%%%%
\SGAPrintPage{V}%
\sgaStableTarget{sga6.front.preface}\section*{\underline{Preface}}
The present edition of the Algebraic Geometry Seminar 6 reproduces, with very few changes, the first edition, distributed by the I.H.E.S. as mimeographed exposés. The corrections made concern essentially printing errors, except in Exposés I and III, where L. Illusie made slight modifications to some statements that were incorrect in the original version. In addition, footnotes have been added, in particular in Exposé XIV by A. Grothendieck, to indicate recent progress on questions that were still open at the time of the Seminar.

Finally, we point out that Exposé XI, by D. Ferrand, concerning the theory of the determinant for perfect complexes, has not been written up and will therefore not be published; to avoid errors in references, we have kept the numbering of the subsequent exposés unchanged.

\begin{flushright}June 1971\end{flushright}

%%%%%%%%%%  scan idx 7  |  folio VI  %%%%%%%%%%
\SGAPrintPage{VI}%
%%%%%%%%%%  scan idx 8  |  folio VII  %%%%%%%%%%
\SGAPrintPage{VII}%
\sgaStableTarget{sga6.front.introduction}\section*{\underline{Introduction}}
The aim of the present Seminar is to develop a \underline{global} intersection theory, and a Riemann--Roch formula, for arbitrary schemes. The reader will find a description of the Seminar's program in Exposé 0. The possibility, in principle, of proving a Riemann--Roch formula for general \underline{regular} schemes, along the lines of the well-known report of Borel--Serre in 1958, was clear from that time, at least for a projective morphism. The extension to the general case is less evident; the program of which it forms a part (and which is partially carried out in our seminar) dates from 1960. As was also the case for the duality theory of coherent sheaves (cf.\ R.~Hartshorne, Residues and Duality, Lecture Notes in Mathematics n\textsuperscript{o}~20, Springer), an essential tool for formulating a satisfactory theory is Verdier's theory of \underline{derived categories}, knowledge of which is indispensable for understanding the Seminar.

Apart from this theory, studying the Seminar requires scarcely more than a general knowledge of the foundations of Algebraic Geometry, as set out in EGA, Chapters I, II, III; in addition, we shall have only occasional references to make to EGA IV, for certain facts concerning flat or smooth morphisms, which for the most part also appear in the first exposés of SGA~1. Finally, to develop certain results with all the generality desirable for applications, we sometimes use the notions of a \underline{ringed site} and a \underline{ringed topos}, for which we refer to SGA~4, Exposés I to IV. A reader unfamiliar with the language of sites and topoi may everywhere replace these objects by ordinary topological spaces, the objects of the topos then being replaced by open subsets of these spaces; nevertheless, we advise the reader, preferably, to assimilate the language of topoi, which provides an extremely convenient principle of unification.

The fundamental notion for the theory presented here is that of a \underline{perfect complex of Modules}, and its various variants, developed in Exposés I, II, III. It seems clear that these notions, which are also important in other contexts in Algebraic Geometry (notably for formulas of Lefschetz--Verdier type in cohomology with discrete (SGA 5) or coherent coefficients), will also have their part to play outside Algebraic Geometry, notably for the formulation of a com%
%%%%%%%%%%  scan idx 9  |  folio VIII  %%%%%%%%%%
plex\SGAPrintPage{VIII}%
 analytic variant of the Riemann--Roch theorem presented here, or suitable variants of the Atiyah--Singer index theorem. Some indications in this direction will be given in Exp.\ II. Unfortunately, at the present time there is still no statement, even a heuristic one, that would encompass these two theorems (the first of which, for the moment, remains conjectural). A new idea is apparently lacking, as one is also lacking in Algebraic Geometry to obtain a proof of the Riemann--Roch theorem beyond projectivity hypotheses (cf.\ Exp.\ XIV, n\textsuperscript{o}~2).

We shall make no use in this seminar of the \underline{local theory} of intersections, set out in J.P. Serre, Algèbre Locale. Multiplicités (Lecture Notes in Mathematics n\textsuperscript{o}~11, Springer), and shall simply content ourselves with pointing out here that this course of Serre had an evident influence on the genesis of the theory in 1957. Nor shall we use the theory of the \underline{Chow ring} developed in the Chevalley Seminar 1958, Exposés 2 and 3 (École Normale Supérieure). This theory is essentially linked to nonsingularity hypotheses, whereas the aim of our seminar, on the contrary, is to develop an intersection theory on general schemes (and even on general locally ringed topoi) see on this subject the comments in Exp.\ XIV n\textsuperscript{o}~4 and n\textsuperscript{o}~8, giving the relations between our theory and that of the Chow ring, and raising the question of a generalization of the latter. We may say that the Seminar presents a coherent and self-contained intersection theory, but by no means exhausts the different points of view that are useful (or indeed indispensable) in intersection theory, and that it should consequently be supplemented by the cited exposés of Serre and Chevalley.

We have attached to the seminar (as an Appendix to Exp.\ 0) Grothendieck's 1957 report ``classes de Faisceaux et théorème de Riemann--Roch'', which served as the basis for the seminar of Borel and Serre at Princeton in the same year, as well as for their already cited article. This report outlines two proofs of the Riemann--Roch theorem for nonsingular quasi-projective varieties; the first, valid at present only in characteristic zero, but giving in return a slightly more precise result in the case of an immersion, does not yet appear in the literature (including the work of Borel--Serre and the present seminar). Reading this report of course presupposes none of the other exposés of the seminar, and may even serve as an introduction to its study on the same footing as Exposé 0, especially for a reader not yet familiar with Borel--Serre. Moreover, the proof to which we have just alluded applies essentially as it stands to the case
%%%%%%%%%%  scan idx 10  |  folio IX  %%%%%%%%%%
\SGAPrintPage{IX}%
of a projective morphism of complex analytic spaces, and will undoubtedly also apply in arbitrary characteristic once the problem of the ``exterior power'' operations in the derived category has been solved (cf.\ Exp.\ XIV, n\textsuperscript{os}~1, 2, 3). The absence of a study of this question undoubtedly constitutes the most serious gap in this Seminar, even from the point of view we have adopted.

One will notice the absence, from the table of contents of the present Seminar, of two of the participants mentioned on the cover page. One of them, J.P. Jouanolou, had taken an active part in the technical development of the first part of the seminar, but had been prevented from taking part in the oral exposés; to him we owe, in particular, the linear independence assertion contained in the important statement VI 1.1. giving the structure of the ring $K^{\bullet}$ of classes of locally free Modules on a projective bundle (which had previously been proved only under the assumption that the base scheme admits an ample invertible Module). The other, J.P. Serre, had given two oral exposés, logically independent of the rest of the Seminar, which he therefore preferred to publish as a separate article (cf.\ J.P. Serre, Groupes de Grothendieck des schémas en groupes réductifs déployés, to appear in Publications Mathématiques, n\textsuperscript{o}~34).

As in the other parts of the Algebraic Geometry Seminar of Bois Marie, the sigla SGA 1 to SGA 6 refer to the different parts of that Seminar, the Roman numeral following the siglum indicating the number of the exposé, and the Arabic numerals that follow it corresponding to the internal numbering of the exposé in question; for internal references within the present Seminar, the siglum SGA 6 is omitted, using references of the form I 4.9 (meaning: Exposé I, statement 4.9). The siglum EGA refers to the Éléments de Géométrie Algébrique of J. Dieudonné and A. Grothendieck.

%%%%%%%%%%  scan idx 11  |  folio X  %%%%%%%%%%
\SGAPrintPage{X}%
%%%%%%%%%%  scan idx 12  |  folio XI  %%%%%%%%%%
\SGAPrintPage{XI}%
\sgaStableTarget{sga6.front.table-of-contents}\section*{\underline{Table of contents}}

\noindent
\underline{EXPOSÉ 0}\\
Outline of a Program for an Intersection Theory on General Schemes\\
by A. Grothendieck \dotfill 1\\[8pt]
Classes of Sheaves and the Riemann--Roch Theorem\\
by A. Grothendieck \dotfill 20\\[8pt]
\underline{EXPOSÉ I}\\
Generalities on Finiteness Conditions in Derived Categories\\
by L. Illusie \dotfill 78\\[8pt]
\underline{EXPOSÉ II}\\
Existence of Global Resolutions\\
by L. Illusie \dotfill 160\\[8pt]
\underline{EXPOSÉ III}\\
Relative Finiteness Conditions\\
by L. Illusie \dotfill 222\\[8pt]
\underline{EXPOSÉ IV}\\
Grothendieck Groups of Ringed Topoi\\
by L. Illusie \dotfill 274\\[8pt]
\underline{EXPOSÉ V}\\
Generalities on $\lambda$-Rings\\
by P. Berthelot \dotfill 297\\[8pt]
\underline{EXPOSÉ VI}\\
The $K^{\bullet}$ of a Projective Bundle: Calculations and Consequences\\
by P. Berthelot \dotfill 365\\[8pt]
%%%%%%%%%%  scan idx 13  |  folio XII  %%%%%%%%%%
\SGAPrintPage{XII}%
\underline{EXPOSÉ VII}\\
Regular Immersions and Calculation of the $K^{\bullet}$ of a Blown-up Scheme\\
by P. Berthelot \dotfill 416\\[8pt]
\underline{EXPOSÉ VIII}\\
The Riemann--Roch Theorem\\
by P. Berthelot \dotfill 466\\[8pt]
\underline{EXPOSÉ IX}\\
Some Calculations of $K$ Groups.\\
by P. Berthelot \dotfill 498\\[8pt]
\underline{EXPOSÉ X}\\
Intersection Formalism on Proper Algebraic Schemes\\
by O. Jussila\\
With an Appendix by A. Grothendieck\\
Specialization in Intersection Theory \dotfill 519\\[8pt]
\underline{EXPOSÉ XI} -- Not written up\\[8pt]
\underline{EXPOSÉ XII}\\
A Relative Representability Theorem for the Picard Functor\\
by M. Raynaud (written up by S. Kleiman) \dotfill 595\\[8pt]
\underline{EXPOSÉ XIII}\\
Finiteness Theorems for the Picard Functor\\
by S. Kleiman \dotfill 616\\[8pt]
\underline{EXPOSÉ XIV}\\
Open Problems in Intersection Theory\\
by A. Grothendieck \dotfill 667\\[8pt]
Terminological Index \dotfill 691\\[4pt]
Index of Notations \dotfill 696\par

%%%%%%%%%%  scan idx 14  |  folio 1  %%%%%%%%%%
\SGAPrintPage{1}
\sgaStableTarget{sga6.e00.expose}\section*{\underline{Exposé 0}}
\begin{center}
\underline{\underline{OUTLINE OF A PROGRAMME FOR A THEORY OF INTERSECTIONS}}\\[3pt]
\underline{\underline{ON GENERAL SCHEMES}}\\[8pt]
\underline{by A. Grothendieck}
\end{center}

The present exposé is introductory in nature, and reading it is not logically indispensable for the study of the Seminar. It is addressed more particularly to readers familiar with the provisional version of the Riemann--Roch theorem, as it is presented in the report of Borel--Serre [2] or in the report of Grothendieck already mentioned in the Introduction (subsequently cited as [RRR]), which is reproduced as an appendix at the end of the present exposé.

\sgaStableTarget{sga6.e00.sec.1}\subsection*{\normalfont 1.}
Recall the Riemann--Roch formula for a proper morphism
\[
f:X\longrightarrow Y
\]
of quasi-projective smooth schemes over a field $k$, and a coherent sheaf $F$ on $X$, whose class in the group $K(X)$ of classes of coherent sheaves on $X$ is denoted by $\cl(F)$:
\sgaStableTarget{sga6.e00.eq.1-1}\begin{equation}\tag{1.1}
\Todd(T_Y)\,\ch_Y\bigl(f_*(\cl(F))\bigr)=f_*\bigl(\Todd(T_X)\,\ch_X(F)\bigr).
\end{equation}
This formula holds in $A(Y)\otimes_{\mathbb Z}\mathbb Q$, where $A(Y)$ is the Chow ring of $Y$, the $f_*$ on the right-hand side being obtained by tensoring with $\mathbb Q$ the homomorphism ``direct image of cycles''
\[
f_*:A(X)\longrightarrow A(Y),
\]
and that on the left-hand side being the Euler--Poincaré characteristic of $F$ relative to $f$:
\[
f_*(\cl(F))=\sum_i(-1)^i\cl\bigl(R^if_*(F)\bigr)~;
\]
$\ch_X$, $\ch_Y$ denote the Chern characters on $X$, respectively on $Y$, and $T_X$, $T_Y$ the tangent bundles of $X$, respectively of $Y$. As is known, $\Todd(-)$ and $\ch(-)$ are universal polynomials with coefficients in $\mathbb Q$ in the Chern classes of the argument. Since the constant term of $\Todd(-)$ is $1$, it is an invertible element for every value of the argument, so that, after multiplication by $\Todd(T_Y)^{-1}$, formula (1.1) takes the equivalent form, more convenient for our purpose,
\sgaStableTarget{sga6.e00.eq.1-2}\begin{equation}\tag{1.2}
\ch_Y\bigl(f_*(\cl(F))\bigr)=f_*\bigl(\Todd(T_f)\cdot\ch_X(F)\bigr),
\end{equation}
where we have put
\sgaStableTarget{sga6.e00.eq.1-3}\begin{equation}\tag{1.3}
T_f=T_X-f^*(T_Y)\in K(X),
\end{equation}
%%%%%%%%%%  scan idx 15  |  folio 2  %%%%%%%%%%
\SGAPrintPage{2}
so that $T_f$ plays the role of a \underline{virtual relative tangent bundle} of $X$ over $Y$. When the morphism $f$ is smooth (i.e. its tangent map is everywhere surjective), we simply have
\[
T_f=T_{X/Y}\qquad\text{(tangent bundle along the fibres)},
\]
whereas when $f:X\to Y$ is an immersion, we find
\[
T_f=-\check N_{X/Y},
\]
where $\check N_{X/Y}$ denotes the normal sheaf of $X$ in $Y$.

One of the main aims of the Seminar is to generalize (1.2) simultaneously in two directions:
\begin{enumerate}[label=\alph*)]
\item To remove the hypothesis of the existence of a base field $k$.
\item To replace the regularity hypotheses on $Y,X$ by a hypothesis of ``local regularity'' of $f$.
\end{enumerate}
Finally, along the way, we shall also consider the problem:
\begin{enumerate}[label=\alph*),start=3]
\item To eliminate the quasi-projectivity hypotheses which, in the absence of a base field, are expressed by the existence of ample invertible Modules on $X$ and on $Y$.
\end{enumerate}

\sgaStableTarget{sga6.e00.sec.2}\subsection*{\normalfont 2.}
Let us first examine generalization a), while retaining the hypotheses of regularity and of the existence of ample invertible Modules on $X$ and on $Y$.

The definition of $K(X)$, $K(Y)$ and of the homomorphism $f_*:K(X)\to K(Y)$ then presents no new problem, because $X$ and $Y$ are regular. The most natural way of giving a meaning to (1.2) therefore seems to consist in defining \underline{Chow rings} $A(X),A(Y)$ and a group homomorphism
\[
f_*:A(X)\longrightarrow A(Y),
\]
in establishing a theory of Chern classes, providing maps
\[
c_i:K(X)\longrightarrow A(X),
\]
and similarly for $Y$, and finally in describing a virtual relative tangent bundle element
\[
T_f\in K(X).
\]

%%%%%%%%%%  scan idx 16  |  folio 3  %%%%%%%%%%
\SGAPrintPage{3}
\sgaStableTarget{sga6.e00.sec.2.1}\subsubsection*{\normalfont 2.1}
As for the latter, a definition presents itself fairly evidently. We use the fact that, by the hypothesis of the existence of an ample invertible Module on $X$ and on $Y$, the morphism $f$ can be factored as
\sgaStableTarget{sga6.e00.eq.2-1}\begin{equation}\tag{2.1}
X\xrightarrow{i}X'\xrightarrow{f'}Y,
\end{equation}
where $i$ is a closed immersion and $f'$ is a smooth morphism; for example, we may take $X'=\mathbb P^r_Y$ for suitable $r$. We then put
\sgaStableTarget{sga6.e00.opening.eq.2-2}\begin{equation}\tag{2.2}
\check T_f=\cl({\underline\Omega}^1_{X'/Y})-\cl({\underline N}_{X/X'}),
\end{equation}
where ${\underline\Omega}^1_{X'/Y}$ is the locally free Module of relative $1$-differentials of $X'$ over $Y$ (or relative cotangent Module of $X'$ over $Y$), and ${\underline N}_{X/X'}$ is the conormal Module of $X$ in $X'$, equal by definition to ${\underline J}/{\underline J}^2$, where ${\underline J}$ is the Ideal on $X'$ defining the closed subscheme $X$; since $X$ and $X'$ are regular, we also conclude that ${\underline N}_{X/X'}$ is locally free. Finally, $\check T_f$ denotes the dual of $T_f$. We readily verify (Exp. VIII) that the element $T_f$ defined by (2.1) does not depend on the chosen factorization of $f$.

\sgaStableTarget{sga6.e00.sec.2.2}\subsubsection*{\normalfont 2.2.}
As for a definition, under the conditions considered, of a Chow ring and a corresponding theory of Chern classes, although it is plausible that such a theory should be developable on the model of the theory already known for smooth quasi-projective algebraic schemes, it seems that this work has not been done at present, and a fortiori it has not been undertaken in the Seminar. Instead, we introduce another ring which will play a role analogous to that of the Chow ring, and whose tensor product with $Q$ is indeed isomorphic to $A(X)\otimes_{\mathbb Z}\mathbb Q$ when $X$ is quasi-projective and smooth over a field $k$. This ring is the associated graded ring of $K(X)$ for a suitable filtration of that ring. A first natural filtration of $K(X)$, the one considered in [RRR], is the ``\underline{topological filtration}'' by the codimension of the support of coherent sheaves, $\operatorname{Filt}^{i}_{\mathrm{top}}K(X)$ being generated by the $\cl(F)$ with $F$ coherent and $\operatorname{codim}(\operatorname{supp}F,X)\ge i$. A first question is the compatibility of this filtration with the ring structure, a question resolved in [RRR], when $X$ is smooth over a field $k$, by means of Chow's ``moving lemma''. Granting this point, we therefore obtain an associated graded ring, which we shall denote here by $\Gr^{\bullet}_{\mathrm{top}}(X)$. In [RRR] a surjective ring homomorphism is defined:
\[
\phi:A(X)\longrightarrow \Gr^{\bullet}_{\mathrm{top}}(X),
\]
%%%%%%%%%%  scan idx 17  |  folio 4  %%%%%%%%%%
\SGAPrintPage{4}
(sending the class of a prime cycle $Z$ to $\cl(\underline O_Z)$,) whose kernel subsequently turns out (Exp. XIV 4.2) to be a torsion subgroup. Moreover, [RRR] shows how the images in $\Gr^{\bullet}_{\mathrm{top}}(X)$ of the Chern classes $c_i(F)$ can be calculated directly in terms of $\cl(F)\in K(X)$ and the exterior-power operations $\lambda^i$ in $K(X)$.

This shows that, for $X$ a noetherian regular scheme equipped with an ample invertible Module, the ring $\Gr^{\bullet}_{\mathrm{top}}(X)$ is a convenient substitute for the hypothetical Chow ring $A(X)$; it also has the advantage over the latter that certain calculations in [RRR] take place directly in this ring and lead to important formulas in it, formulas which, even for $X$ projective and smooth over a field, have not at present been established in $A(X)$ itself (Exp. XIV, 4.3 and 4.4). To justify this point of view completely, however, one would have to resolve the problem mentioned above of the compatibility of the topological filtration of $K(X)$ with its ring structure, a problem which seems essentially connected with Chow's ``moving lemma'', which is itself the essential technical ingredient of the theory of the Chow ring that we have claimed to be able to avoid.

\sgaStableTarget{sga6.e00.sec.2.3}\subsubsection*{\normalfont 2.3.}
To avoid this problem, we give $K(X)$ a second filtration, finer than the topological filtration, which can be described in purely algebraic terms from the augmented $\lambda$-ring structure of $K(X)$. The definition of this filtration is suggested by the expression given in [RRR] for the Chern classes of an $x\in K(X)$ in $\Gr^{\bullet}_{\mathrm{top}}(X)$,
\sgaStableTarget{sga6.e00.eq.2-3}\begin{equation}\tag{2.3}
c_i(x)=\gamma^i(x-\epsilon(x))\mod \operatorname{Filt}^{i+1}_{\mathrm{top}}K(X),
\end{equation}
where $\gamma^i$ denotes a certain linear combination of the $\lambda^j$ ($j\leq i$) made explicit in loc. cit., $\epsilon$ being the augmentation. We show that $\gamma^i(x-\epsilon(x))$ does have topological filtration $\ge i$, which gives a meaning to formula (2.3). We then define the $\lambda$-\underline{filtration} of $K(X)$ by the ideals $\operatorname{Filt}^i K(X)$, where $\operatorname{Filt}^i K(X)$ consists of sums of expressions of the form
\[
\gamma^{i_1}(x_1-\epsilon(x_1))\cdots \gamma^{i_k}(x_k-\epsilon(x_k)),\qquad i_1+\cdots+i_k\ge i.
\]
The associated graded ring of this filtration will be denoted by $\Gr^\bullet(X)$, and it is this ring that will replace the Chow ring of $X$ in the Seminar. We can moreover show that, in the case considered here, the canonical homomorphism
\sgaStableTarget{sga6.e00.eq.2-4}\begin{equation}\tag{2.4}
\Gr^\bullet(X)\longrightarrow \Gr^{\bullet}_{\mathrm{top}}(X)
\end{equation}
%%%%%%%%%%  scan idx 18  |  folio 5  %%%%%%%%%%
\SGAPrintPage{5}
is an isomorphism modulo torsion, which justifies the point of view that this ring tensored with $\mathbb Q$ can play in every respect the same role as the Chow ring tensored with $\mathbb Q$. On the other hand, we have done exactly what is needed for formula (2.3) (with $\operatorname{Filt}^i_{\mathrm{top}}$ replaced by $\operatorname{Filt}^i$) to define a theory of Chern classes with values in $\Gr^\bullet(X)$, from which the Chern classes in $\Gr^{\bullet}_{\mathrm{top}}(X)$ are deduced using the homomorphism (2.4). We show (Exp. V), using the fact (Exp. VI) that $K(X)$ is a \underline{special} $\lambda$-ring in the sense of [RRR] (i.e. a $\lambda$-ring in the sense of Exp. V), that this notion of Chern class satisfies all the usual formal properties reviewed in [RRR].

\sgaStableTarget{sga6.e00.sec.2.4}\subsubsection*{\normalfont 2.4.}
We note a very appreciable advantage of the definition of $\Gr^\bullet(X)$ over that of $A(X)$ and $\Gr^{\bullet}_{\mathrm{top}}(X)$: it remains meaningful for an absolutely arbitrary scheme $X$, and more generally for an arbitrary locally ringed topos $X$, taking for $K(X)$ the group of classes constructed from locally free Modules on $X$, which is still an augmented $\lambda$-ring, hence equipped with a corresponding filtration allowing an associated graded $\Gr^\bullet(X)$ to be defined. As the price of this generality, however, we must observe that, when we are unwilling to disregard torsion phenomena, i.e. to tensor with $\mathbb Q$, the properties of $\Gr^\bullet(X)$ are less satisfactory in various respects than those of its two competitors (Exp. XIV no.~4). Thus, supposing again that $f:X\to Y$ is a proper morphism of noetherian regular schemes equipped with ample invertible Modules, it is not true in general that the homomorphism $f_*:K(X)\longrightarrow K(Y)$ maps $\operatorname{Filt}^i K(X)$ into $\operatorname{Filt}^{i-d} K(Y)$, where $d$ is the relative dimension of $X$ over $Y$, and hence defines, by passage to the associated graded rings, a homomorphism of degree $-d$ from $\Gr^\bullet(X)$ to $\Gr^\bullet(Y)$. However, we can prove (Exp. VII) that this holds after tensoring with $\mathbb Q$, so that we obtain a homomorphism of degree $-d$
\sgaStableTarget{sga6.e00.eq.2-5}\begin{equation}\tag{2.5}
f_*:\Gr^\bullet(X)_{\mathbb Q}\longrightarrow \Gr^\bullet(Y)_{\mathbb Q},
\end{equation}
playing the role of the homomorphism ``direct image of cycles'' $A(X)\to A(Y)$.

With this understood, we see that we have all the structural elements needed to give a meaning to formula (1.2) without postulating the existence of a base field, $f:X\to Y$ being a proper morphism of noetherian regular schemes admitting ample invertible Modules.

%%%%%%%%%%  scan idx 19  |  folio 6  %%%%%%%%%%
\SGAPrintPage{6}
\sgaStableTarget{sga6.e00.sec.3}\subsection*{\normalfont 3.}
Let us now show how we can still give a meaning to formula (1.2) when we abandon the regularity hypothesis on $X,Y$, replacing it by a hypothesis of a local nature on the morphism $f$, which we shall specify.

\sgaStableTarget{sga6.e00.sec.3.1}\subsubsection*{\normalfont 3.1}
A first difficulty arises because the observation that served as the starting point of the theory [RRR], namely that the group $K(X)$ of classes of sheaves on $X$ is the same whether it is defined in terms of coherent sheaves or in terms of locally free sheaves, is essentially tied to the regularity hypothesis on $X$. When $X$ is a general scheme (say noetherian, for simplicity), we must distinguish these two groups, which we shall denote respectively by $K_\bullet(X)$ and $K^\bullet(X)$, the first being covariant in $X$ for proper morphisms, and the second contravariant in $X$ for arbitrary morphisms. We note that $K^\bullet(X)$ has a natural ring structure (and even an augmented $\lambda$-ring structure), whereas $K_\bullet(X)$ in general no longer has a ring structure, but only a module structure over $K^\bullet(X)$. The homomorphism $x\to x\,\cl_\bullet(\underline O_X)$ (where $\cl_\bullet(\underline F)$ denotes the class of $\underline F$ in $K_\bullet(X)$) from $K^\bullet(X)$ to $K_\bullet(X)$ is in general no longer an isomorphism. \uline{For our formulation of the generalized Riemann--Roch theorem, we shall work exclusively with $K^\bullet(X)$, $K^\bullet(Y)$, to the exclusion of $K_\bullet(X)$, $K_\bullet(Y)$}. As we noted in 2.3., these augmented $\lambda$-rings, equipped with their natural $\lambda$-filtration, give rise to graded rings $\Gr^\bullet(X)$, $\Gr^\bullet(Y)$, and ``Chern class'' homomorphisms
\[
c_i:K^\bullet(X)\longrightarrow \Gr^\bullet(X),
\]
and similarly for $Y$. The difficulty, however, arises in defining a direct-image homomorphism
\sgaStableTarget{sga6.e00.eq.3-1}\begin{equation}\tag{3.1}
f_*:K^\bullet(X)\longrightarrow K^\bullet(Y),
\end{equation}
which can no longer be defined by the naive formula
\[
f_*(\cl^\bullet(\underline F))=\sum_i(-1)^i\cl^\bullet\bigl(R^if_*(\underline F)\bigr),
\]
because, even when $\underline F$ is locally free, the Modules $R^if_*(\underline F)$ are in general no longer locally free, nor even of finite tor-dimension, and the preceding formula loses its meaning. A solution to this difficulty is provided by the point of view of \underline{derived categories}. Indeed, although the $R^if_*(\underline F)$ are in general ``bad'' sheaves, say of infinite tor-dimension,
%%%%%%%%%%  scan idx 20  |  folio 7  %%%%%%%%%%
\SGAPrintPage{7}
to which we cannot associate a class in $K^\bullet(Y)$, the complex $Rf_*(\underline F)$ on $Y$, the direct image of $\underline F$ in the sense of derived categories, whose cohomology sheaves are the $R^if_*(\underline F)$, nevertheless tends to be excellent. More precisely, suppose that $\underline F$ has finite tor-dimension over $\underline O_Y$, which is the case when $\underline O_X$ has finite tor-dimension over $\underline O_Y$ and, in addition, $F$ is locally free. We then show (Exp. III) that the complex $Rf_*(F)$ is ``\underline{perfect}'', by which we mean here: it has coherent cohomology sheaves (which we already knew in any case) \underline{and} finite global tor-dimension; or, equivalently: on every affine open $U$ of $Y$, this complex is isomorphic, in the derived category $D(U)$, to a complex bounded in degree and locally free of finite type in each degree. From the point of view of Homological Algebra, perfect complexes are ``as good'' as locally free Modules, of which they are the natural generalization. They moreover form a ``triangulated subcategory'' of the derived category $D(Y)$, i.e. a subcategory stable under cones (or ``mapping cylinders'') of morphisms. On the other hand, to every triangulated category $C$ we associate a group $K(C)$ of classes, a natural variant of the construction of the group of classes of an abelian category, by taking the free group generated by the objects of $C$, modulo the relations
\[
\cl(K)-\cl(K')-\cl(K'')
\]
for every distinguished triangle
\[
\begin{tikzcd}[column sep=large]
 & K'' \arrow[dl] & \\
K' \arrow[rr] & & K \arrow[ul] .
\end{tikzcd}
\]
What replaces the fundamental observation of [RRR] in the present case is the fact that, when the scheme $Y$ admits an ample invertible Module, every perfect complex $L_\bullet$ on $Y$ is globally (and not only locally) isomorphic in $D(Y)$ to a complex $L'_\bullet$ bounded in degree and locally free of finite type in each degree (Exp. II); it follows readily (Exp. IV) that the canonical homomorphism
\[
K^\bullet(Y)\longrightarrow K(\Parf(Y)),
\]
where $\Parf(Y)$ is the triangulated category of perfect complexes on $Y$, is an isomorphism. Thus to every perfect complex $L_\bullet$ on $Y$ we can associate a class
\[
\cl^\bullet(L_\bullet)\in K^\bullet(Y),
\]
%%%%%%%%%%  scan idx 21  |  folio 8  %%%%%%%%%%
\SGAPrintPage{8}
which is moreover simply
\[
\cl^\bullet(L_\bullet)=\sum_i(-1)^i\cl^\bullet(L'_i).
\]
where $L'_\bullet$ is as above, so that the $L'_i$ are indeed locally free Modules on $Y$ whose classes we can take in $K^\bullet(Y)$.

With this understood, we define a homomorphism (3.1) whenever $f:X\to Y$ is a \underline{proper morphism of finite tor-dimension} of noetherian schemes such that $Y$ admits an ample invertible Module, simply by putting
\sgaStableTarget{sga6.e00.eq.3-2}\begin{equation}\tag{3.2}
f_*(\cl^\bullet(\underline F))=\cl^\bullet(Rf_*(\underline F))
\end{equation}
for every locally free Module $\underline F$ on $X$ (so that $Rf_*(\underline F)$ is indeed a perfect complex on $Y$, and the right-hand side is defined). When $X$ also admits an ample invertible Module, $K^\bullet(X)$ can likewise be interpreted as $K(\Parf(X))$, and the homomorphism (3.1) is precisely the one associated, by functoriality, with the exact functor of triangulated categories induced by
\[
Rf_*:\Parf(X)\longrightarrow \Parf(Y).
\]

\sgaStableTarget{sga6.e00.sec.3.2}\subsubsection*{\normalfont 3.2.}
To give a meaning to (1.2), we must still impose on $f$ a hypothesis allowing a meaning to be given to the homomorphism (2.5)
\sgaStableTarget{sga6.e00.eq.3-3}\begin{equation}\tag{3.3}
f_*:\Gr^\bullet(X)_{\mathbb Q}\longrightarrow \Gr^\bullet(Y)_{\mathbb Q},
\end{equation}
i.e. ensuring that the homomorphism (3.1) just defined is (up to a shift and modulo torsion) compatible with the filtrations of these rings; and finally we must still define a ``virtual relative tangent bundle'' element
\[
T_f\in K^\bullet(X).
\]
For the latter, returning to the definition proposed in 2.1, we see that we should make a hypothesis ensuring that (with the notation of loc. cit.) the conormal Module $\underline N_{X/X'}$ is locally free. The most natural hypothesis for this purpose is that $i$ is a ``regular immersion'', i.e. identifies $X$ with a subscheme of $X'$ defined by an Ideal which locally can be defined by an $\underline O_{X'}$-regular sequence of sections of $\underline O_{X'}$. This hypothesis, of a purely local nature, does not in fact depend on the choice of the factorization of $f$ into an immersion followed by a smooth morphism (Exp. VIII), and, as above, the same holds for the class $T_f$ defined by formula (2.2). When the hypothesis just considered holds, we say that the morphism $f$ is a ``\underline{locally complete intersection}'' \underline{morphism}\SGAAuthorFootnote{raisedstar}{We shall also allow ourselves to omit the word ``locally'' in this terminology, which entails no risk of confusion.}. Such a morphism is moreover locally of finite tor-dimension.

%%%%%%%%%%  scan idx 22  |  folio 9  %%%%%%%%%%
\SGAPrintPage{9}
For a locally complete intersection morphism $f$, we introduce a natural notion of ``virtual relative dimension'' $d(f)$ by the formula
\[
d(f)(x)=\dim_x\bigl(f'^{-1}f(x)\bigr)-\codim_x(X,X').
\]
This is an integer-valued function, locally constant in the point $x\in X$, hence constant when $X$ is connected; it likewise does not depend on the choice of the factorization (2.1). When this virtual relative dimension $d$ is constant, $f$ is proper, and $X$ and $Y$ admit ample invertible Modules, we can show (Exp. VIII), although this is far from a trivial result, that
\sgaStableTarget{sga6.e00.eq.3-4}\begin{equation}\tag{3.4}
f_*\bigl(\operatorname{Filt}^iK^\bullet(X)_{\mathbb Q}\bigr)\subset \operatorname{Filt}^{i-d}K^\bullet(Y)_{\mathbb Q},
\end{equation}
hence a homomorphism (3.3) of degree $-d$.

\sgaStableTarget{sga6.e00.sec.3.3}\subsubsection*{\normalfont 3.3.}
We have thus assembled all the structural elements needed to give a meaning to formula (1.2) \uline{for a proper locally complete intersection morphism of noetherian schemes $X,Y$ both admitting ample invertible Modules}. With a little additional care, this formulation moreover remains meaningful independently of any noetherian hypothesis.

In this general form, this Riemann--Roch formula will be proved in Exp. VIII. The reader will see that the proof given in this Seminar combines elements of the two proofs in [RRR] which were discussed in the Introduction. The need to use the methods of the first of these proofs, involving more $\lambda$-formalism than the second (followed in Borel--Serre), is due to the need for a proof of inclusion (3.4), a question that did not arise in the more restricted setting of [RRR] and Borel--Serre [2], where we could work with the Chow ring or the ring $\Gr^\bullet_{\mathrm{top}}$. From the second proof, we retain the introduction of the blow-up of $Y$ along $X$ when $f$ is an immersion. This technical device will doubtless disappear, and the proof for an immersion will extend to the case where we no longer postulate the existence of an ample invertible Module on $Y$, once the question raised in Exp. XIV no.~1 concerning the introduction of exterior-power operations in the derived category $D^-(Y)$ has been resolved.

\sgaStableTarget{sga6.e00.sec.4}\subsection*{\normalfont 4.}
Finally, we must give some indications of how a meaning can still be given to (1.2) without a hypothesis of the existence
%%%%%%%%%%  scan idx 23  |  folio 10  %%%%%%%%%%
\SGAPrintPage{10}
of ample invertible Modules on $X$ and $Y$.

\sgaStableTarget{sga6.e00.sec.4.1}\subsubsection*{\normalfont 4.1.}
First, we note that \uline{for an arbitrary scheme $X$ (or, more generally, locally ringed topos), it seems beyond doubt that the ``right'' invariant to replace the ring $K(X)$ of [RRR]} is not \uline{the ring $K^\bullet(X)$ considered above, defined in terms of locally free Modules on $X$, but the group $K(\Parf(X))$, considered in 3.1., defined in terms of perfect complexes on $X$, a group that we shall denote here by} $K^\bullet(X)$, and which we shall distinguish from the ``naive'' group of 3.1, denoted by $K^\bullet_{\mathrm{naive}}(X)$. This point of view is forced by the need to define a direct-image homomorphism (3.1) for a proper morphism $f:X\to Y$ of finite tor-dimension between, say, noetherian schemes (the non-noetherian case requires somewhat stronger local finiteness properties of $f$, cf. Exp. III). With the definition of $K^\bullet(X)$ now adopted, this homomorphism $f_*$ is indeed defined trivially by formula (3.2), where $\underline F$ is now an arbitrary perfect complex on $X$, which implies that $Rf_*(\underline F)$ is also perfect.

\sgaStableTarget{sga6.e00.sec.4.2}\subsubsection*{\normalfont 4.2}
This new definition of $K^\bullet(X)$ nevertheless raises a new problem of homological algebra: that of defining a $\lambda$-structure on $K^\bullet(X)$ compatible with its natural ring structure arising from the tensor product in the derived category. (NB When $L_\bullet$ and $L'_\bullet$ are two perfect complexes, their total tensor product $L_\bullet\otimes^{\mathbf L}L'_\bullet$ is also perfect). More precisely, we should define ``exterior-power'' operations in the derived category $D^-(X)$, taking perfect complexes to perfect complexes. This question has not yet been clarified at present, and, as was stated in the Introduction, it constitutes the most serious gap in the Seminar. Once this question has been resolved, and proceeding as in 2.3, we obtain a natural filtration on $K^\bullet(X)$, allowing us to form the associated graded ring $\Gr^\bullet(X)$, which will play the role of the Chow ring, and to define a theory of Chern classes
\[
c_i:K^\bullet(X)\longrightarrow \Gr^\bullet(X).
\]

\sgaStableTarget{sga6.e00.sec.4.3}\subsubsection*{\normalfont 4.3}
To give a meaning to formula (1.2) for a proper locally complete intersection morphism $f:X\to Y$ of noetherian schemes, we must still include in the statement of this formula the nontrivial inclusions (3.4) (where $d$ is the virtual relative dimension) to allow (3.3) to be defined, and finally define the virtual relative tangent class
%%%%%%%%%%  scan idx 24  |  folio 11  %%%%%%%%%%
\SGAPrintPage{11}
\[
T_f\in K^\bullet(X).
\]
When $f$ admits a factorization (2.1) through an immersion into a smooth relative scheme, the definition of $T_f$ presents no new difficulty. We can moreover reformulate it in this case, in a manner more in keeping with the ``derived categories'' point of view, by considering the familiar canonical homomorphism
\sgaStableTarget{sga6.e00.eq.4-1}\begin{equation}\tag{4.1}
\underline N_{X/X'}\xrightarrow{d}\Omega^1_{X'/Y}\otimes_{\underline O_{X'}}\underline O_X
\end{equation}
induced by the exterior differential
\[
d_{X'/Y}:\underline J\subset\underline O_{X'}\longrightarrow\Omega^1_{X'/Y},
\]
and the chain complex
\sgaStableTarget{sga6.e00.opening.eq.4-2}\begin{equation}\tag{4.2}
L^{X/Y}_\bullet\in\operatorname{Ob}D(X)
\end{equation}
defined from (4.1) ``by extending by zeros'', a complex equal in degree $1$, respectively $0$, to the source, respectively the target, of the homomorphism (4.1). Formula (2.2) then takes the more pleasant form
\sgaStableTarget{sga6.e00.opening.eq.4-3}\begin{equation}\tag{4.3}
\check T_f=\cl^\bullet\bigl(L^{X/Y}_\bullet\bigr),
\end{equation}
a formula which makes sense in $K^\bullet(X)$ when $f$ is a locally complete intersection morphism, since in this case the complex $L^{X/Y}_\bullet$ is manifestly perfect. The assertion of invariance of $T_f$ relative to the chosen factorization of $f$ can then readily be made more precise as an assertion of invariance of the complex $L^{X/Y}_\bullet$: up to canonical isomorphism in the derived category $D(X)$, this complex is likewise independent of the chosen factorization (Exp. VIII).

\sgaStableTarget{sga6.e00.sec.4.4}\subsubsection*{\normalfont 4.4.}
When $f:X\to Y$ is an arbitrary morphism of schemes, we can still construct a canonical object $L^{X/Y}_\bullet$ of the derived category $D(X)$, defined up to unique isomorphism, which on each affine open of $X$ over an affine open of $Y$ is isomorphic to the complex constructed according to the principle explained in the preceding paragraph (where, however, in the general case ``smooth'' should be replaced by ``formally smooth'', so that every algebra $B$ over a ring $A$ occurs as a quotient of a formally smooth algebra, for example a polynomial algebra). For such a definition, we cannot proceed by simple ``gluing'' from the affine constructions, because the objects of the derived category $D(X)$ are by their nature
%%%%%%%%%%  scan idx 25  |  folio 12  %%%%%%%%%%
\SGAPrintPage{12}
essentially not glueable. Here is a general construction of $L^{X/Y}_\bullet$, valid more generally for every morphism $f:X\to Y$ of (commutatively) ringed topoi: let $B=\underline O_X$, $A=f^{-1}(\underline O_Y)$, so that $A$ is a Ring on the space $X$, and $B$ is an $A$-Algebra. Let $T\to B$ be a homomorphism of sheaves of \underline{sets} generating $B$ as an $A$-Algebra, i.e. such that the corresponding homomorphism of $A$-Algebras
\[
\phi:C=A[T]\longrightarrow B
\]
is an epimorphism of sheaves of sets, where $A[T]$ denotes the free commutative $A$-Algebra generated by $T$ (i.e. the sheaf associated with the presheaf $U\longmapsto A(U)[T(U)]$ = free commutative $A(U)$-algebra generated by $T(U)$).
Let $J$ be the kernel of $\phi$, and consider, as in 4.3, the chain complex of length $1$ defined by the natural homomorphism
\[
J/J^2\xrightarrow{d}\Omega^1_{C/A}\otimes_C B.
\]
It is not difficult to see that, up to canonical isomorphism in the category $D(X)$, the complex so constructed is independent of the choice of $T$ and of $T\to B$, so that it deserves the name \uline{relative cotangent complex $L^{X/Y}_\bullet$ of $X$ over $Y$}. We also verify that, on every affine open of $X$ over an affine open of $Y$, it is given (up to canonical isomorphism) by the procedure of 4.3.

When $f$ is locally complete intersection, we moreover see that we even obtain a \underline{perfect} complex, which then allows us again to define an element $T_f$ in $K^\bullet(X)$ by formula (4.3). We point out that here again, in general, it would not have been possible to define $T_f$ as an element of $K^\bullet_{\mathrm{naive}}(X)$, and nothing permits us to suppose that $T_f$ is \underline{globally} isomorphic in $D(X)$ to a bounded complex with locally free components of finite type. This therefore provides further justification for the point of view advanced in 4.1 in favour of the ``sophisticated'' $K^\bullet(X)$, in preference to $K^\bullet_{\mathrm{naive}}(X)$.

\sgaStableTarget{sga6.e00.sec.4.5}\subsubsection*{\normalfont 4.5.}
Subject to the question raised in 4.2, we have therefore again assembled all the structural elements needed to give a meaning to (1.2) for a proper locally complete intersection morphism of arbitrary schemes (the noetherian hypothesis implicitly made in the preceding sections can in fact be removed without essential difficulty). Unfortunately, this does not mean that, even subject to this reservation, the formula in question has been established, even when $Y$ is the spectrum of a field
%%%%%%%%%%  scan idx 26  |  folio 13  %%%%%%%%%%
\SGAPrintPage{13}
that is algebraically closed of characteristic $p>0$, and $X$ is, say, a proper smooth algebraic scheme of dimension $3$. See the comments in Exp. XIV no.~2.

\sgaStableTarget{sga6.e00.sec.5}\subsection*{\normalfont 5.}
In the form presented in no.~4, the Riemann--Roch formula (1.2) also has a meaning for a proper locally complete intersection morphism of complex analytic spaces (or of rigid analytic spaces in the sense of Tate). But of course the proof given in [RRR] applies to this case only when we also assume that the morphism $f$ is \underline{projective}, and, as in the algebraic case, it seems that an essentially new idea is needed to treat the general case. The formula we envisage here has values in an ``analytic Chow ring'' $\Gr^\bullet(X)$, or rather $\Gr^\bullet(X)_{\mathbb Q}$, constructed according to the principle explained in 4.2 (the difficulty encountered there moreover disappears because we are in characteristic zero, as noted in Exp. XIV no.~1).

When $X$ and $Y$ are nonsingular, another version of the Riemann--Roch formula can be given as follows. First let $X$ be an arbitrary complex analytic space, and $X^{\mathrm{top}}$ the same space equipped with the sheaf of continuous complex-valued functions. We thus have a canonical morphism of ringed spaces
\[
X^{\mathrm{top}}\longrightarrow X,
\]
which, by the evident contravariance of the functor $K^\bullet(X)$ in the variable ringed space $X$, defines a ring homomorphism
\sgaStableTarget{sga6.e00.eq.5-1}\begin{equation}\tag{5.1}
K^\bullet(X)\longrightarrow K^\bullet(X^{\mathrm{top}}).
\end{equation}
For simplicity, suppose that $X^{\mathrm{top}}$ is compact; then, as in the case (already mentioned in 3.1) of a scheme admitting an ample invertible module, we see (Exp. II) that every perfect complex on $X^{\mathrm{top}}$ is globally isomorphic to a bounded complex locally free in each degree, i.e. to a complex of complex vector bundles on $X^{\mathrm{top}}$, and consequently $K^\bullet(X^{\mathrm{top}})$ is canonically isomorphic to the topologists' familiar group $K(X^{\mathrm{top}})$ [1], defined in terms of complex vector bundles on the compact space $X^{\mathrm{top}}$. Thus we can interpret the homomorphism (5.1) as a homomorphism from $K^\bullet(X)$, defined in terms of the analytic structure, to the topologists' familiar ring $K(X^{\mathrm{top}})$. Now, for a morphism $f:X\longrightarrow Y$ of compact nonsingular analytic spaces, we know how to associate, by Atiyah--Hirzebruch
%%%%%%%%%%  scan idx 27  |  folio 14  %%%%%%%%%%
\SGAPrintPage{14}
[1], a group homomorphism
\sgaStableTarget{sga6.e00.eq.5-2}\begin{equation}\tag{5.2}
f^{\mathrm{top}}_*:K^\bullet(X^{\mathrm{top}})\longrightarrow K^\bullet(Y^{\mathrm{top}}).
\end{equation}
Using likewise the homomorphism $f_*$ of (3.1) defined by formula (3.2) (which here implicitly uses Grauert's finiteness theorem [3] for a proper morphism of analytic spaces), we obtain a square of natural homomorphisms involving only groups of type $K^\bullet$ (to the exclusion of rings of Chow-ring or cohomology type):
\sgaStableTarget{sga6.e00.eq.5-3}\begin{equation}\tag{5.3}
\begin{array}{ccc}
K^\bullet(X) & \longrightarrow & K^\bullet(X^{\mathrm{top}}) \\
f_*\downarrow && \downarrow f^{\mathrm{top}}_* \\
K^\bullet(Y) & \longrightarrow & K^\bullet(Y^{\mathrm{top}}).
\end{array}
\end{equation}
The Riemann--Roch type formula to which we alluded consists in asserting that this square commutes. We note the difference in nature between this assertion (which does not disregard torsion phenomena, and lies in a group $K^\bullet(Y^{\mathrm{top}})$ of an essentially discrete nature) and the Riemann--Roch formula of the type considered in no.~4, which lies in a group of ``Chow'' type that is no longer discrete in nature but must instead be tensored with $\mathbb Q$. We note that formula (5.3) is proved, in any case, when $X$ and $Y$ are nonsingular projective varieties [7], as is the formula of the type considered in no.~4, which is proved in this case in [RRR]. This is a plausible variant of the Atiyah--Singer index theorem (and of its generalizations due to Shih [6] and Atiyah), which plainly calls for a common generalization. Once the commutativity of (5.3) has been proved, the ``differentiable Riemann--Roch theorem'' of Atiyah--Hirzebruch [1] would immediately give a cohomological variant of the complex analytic Riemann--Roch formula, namely the commutativity of the square
\sgaStableTarget{sga6.e00.eq.5-4}\begin{equation}\tag{5.4}
\begin{array}{ccc}
K^\bullet(X) & \xrightarrow{\;x\longrightarrow \ch(x)\Todd(T_X)\;} & H^{2*}(X,\mathbb Q) \\
\downarrow && \downarrow \\
K^\bullet(Y) & \xrightarrow{\;y\longrightarrow \ch(y)\Todd(T_Y)\;} & H^{2*}(Y,\mathbb Q),
\end{array}
\end{equation}
where the horizontal arrows are obtained from the cohomological Chern classes, obtained from the homomorphism (5.1) and the ordinary Chern class homomorphism
\[
K^\bullet(X^{\mathrm{top}})\longrightarrow H^{2*}(X,\mathbb Z),
\]
%%%%%%%%%%  scan idx 28  |  folio 15  %%%%%%%%%%
\SGAPrintPage{15}
the first vertical arrow is the same as in (5.3), and the second is the Gysin homomorphism in rational cohomology.

\sgaStableTarget{sga6.e00.sec.6}\subsection*{\normalfont 6.}
Throughout this discussion, we have considered almost exclusively the contravariant invariant $K^\bullet(X)$, to the exclusion of the covariant invariant $K_\bullet(X)$ whose existence was mentioned in 3.1. Regarding $K^\bullet(X)$ as an analogue of the integral \underline{cohomology} ring of a topological space, we must regard $K_\bullet(X)$ as an analogue of integral \underline{homology}, the Module structure of $K_\bullet(X)$ over $K^\bullet(X)$ then corresponding to the cap-product operation\SGAAuthorFootnote{raisedstar}{This analogy is moreover not purely verbal: after passing to the associated graded objects $\Gr^\bullet(X)$, $\Gr_\bullet(X)$, it can be made precise by an explicit compatibility relation between the module structure on $\Gr_\bullet(X)$ and the ``cap'' operation.}. The fact that, when $X$ is regular, the natural homomorphism $K^\bullet(X)\to K_\bullet(X)$ is an isomorphism must be regarded as the analogue of the Poincaré duality theorem on oriented topological manifolds, asserting that cap product with the fundamental homology class defines an isomorphism from cohomology to homology. These analogies also suggest the desirability of defining, say for a finite polyhedron $X$, an invariant $K_\bullet(X)$ of a covariant nature, whose relations to integral homology (spectral sequence, Chern homomorphism\dots) would be analogous to those between $K^\bullet(X)$ and integral cohomology; nor does it seem excluded that there might exist, in terms of these groups $K_\bullet(X)$, variants of the differentiable Riemann--Roch theorem valid for spaces more general than compact differentiable manifolds.

Returning to an arbitrary noetherian scheme $X$, we should equip $K_\bullet(X)$ with an increasing filtration, $\operatorname{Filt}_iK_\bullet(X)$ being generated by classes $\cl_\bullet(F)$ with $F$ a coherent sheaf on $X$ such that $\dim \operatorname{supp}.F\le i$. We shall prove (Exp. X) the compatibility of this filtration with the module structure and the filtration of $K^\bullet(X)$, i.e. the formula
\sgaStableTarget{sga6.e00.eq.6-1}\begin{equation}\tag{6.1}
\operatorname{Filt}^iK^\bullet(X)\cdot \operatorname{Filt}_j K_\bullet(X)\subset \operatorname{Filt}_{j-i}K_\bullet(X),
\end{equation}
which equips the associated graded object of the filtered module $K_\bullet(X)$, denoted by
\[
\Gr_\bullet(X),
\]
with a module structure over $\Gr^\bullet(X)$. Its elements can moreover be interpreted as \uline{classes of algebraic cycles} on $X$, (for an equivalence relation less fine than rational equivalence, when $X$ is of finite type over a base field $K$). If
%%%%%%%%%%  scan idx 29  |  folio 16  %%%%%%%%%%
\SGAPrintPage{16}
\[
f:X\longrightarrow Y
\]
is a proper morphism of noetherian schemes, then
\sgaStableTarget{sga6.e00.eq.6-2}\begin{equation}\tag{6.2}
f_*:K_\bullet(X)\longrightarrow K_\bullet(Y)
\end{equation}
is compatible with the filtrations and defines a homomorphism of graded groups
\sgaStableTarget{sga6.e00.eq.6-3}\begin{equation}\tag{6.3}
f_*:\Gr_\bullet(X)\longrightarrow \Gr_\bullet(Y),
\end{equation}
giving a ``projection formula'', i.e. it is $\Gr^\bullet(Y)$-linear, a formula which is deduced by passage to the associated graded objects from the analogous projection formula for the homomorphism (6.2) (which is $K^\bullet(Y)$-linear).

\uline{From the point of view proposed in the Seminar}, \uline{a workable theory of intersections}, \uline{on noetherian schemes} $X$ \uline{that are not necessarily regular}, \uline{must consist in the combined study of the contravariant invariants} $K^\bullet(X),\Gr^\bullet(X)$ \uline{and the covariant invariants} $K_\bullet(X),\Gr_\bullet(X)$, \uline{which play roles of an essentially different nature}, \uline{and whose properties complement one another}. Thus the covariant groups $K_\bullet(X),\Gr_\bullet(X)$ can be calculated more readily for certain nonproper fibrations, by means of the ``homotopy exact sequence'' that can be established for them ([RRR] and Exp. IX); for example, we prove (Exp. IX) canonical isomorphisms such as
\[
K_\bullet(X[t])\simeq K_\bullet(X),\qquad \Gr_\bullet(X[t])\simeq \Gr_\bullet(X),
\]
which fail in general for $K^\bullet$ when $X$ is a nonregular scheme. On the other hand, $K^\bullet(X)$ and $\Gr^\bullet(X)$ lend themselves well to calculation because of their ring structures.

For schemes $X$ proper over a field, the projection
\[
f^X:X\longrightarrow(\operatorname{point})
\]
defines an ``\uline{Euler--Poincaré characteristic}'' homomorphism
\sgaStableTarget{sga6.e00.eq.6-4}\begin{equation}\tag{6.4}
\lambda_X=f^X_*:K_\bullet(X)\longrightarrow K_\bullet(\operatorname{point})=\mathbb Z,
\end{equation}
inducing on the subgroup $\Gr_0(X)=\operatorname{Filt}_0K_\bullet(X)$ the ``degree of $0$-cycles'' homomorphism
\sgaStableTarget{sga6.e00.eq.6-5}\begin{equation}\tag{6.5}
\deg_X:\Gr_0(X)\longrightarrow \mathbb Z.
\end{equation}
The multiplication homomorphism
\[
\Gr^i(X)\times\Gr_i(X)\longrightarrow \Gr_0(X),
\]
composed with the degree homomorphism (6.5), then gives a pairing
\sgaStableTarget{sga6.e00.eq.6-6}\begin{equation}\tag{6.6}
\Gr^i(X)\times \Gr_i(X)\longrightarrow \mathbb Z,
\end{equation}
which explains why, in certain respects, $\Gr^\bullet(X)$ and $\Gr_\bullet(X)$ play dual roles,
%%%%%%%%%%  scan idx 30  |  folio 17  %%%%%%%%%%
\SGAPrintPage{17}
just as cohomology and homology do. Thus, for a morphism $f:X\to Y$ of proper algebraic schemes, the two homomorphisms
\[
f^*:\Gr^\bullet(Y)\longrightarrow \Gr^\bullet(X)\qquad\text{and}\qquad f_*:\Gr_\bullet(X)\longrightarrow \Gr_\bullet(Y)
\]
are formally transposes of one another, with respect to the pairings (6.6).

The homomorphisms (6.4) and (6.5), and the pairing (6.6) obtained from them, can be regarded as the source of \uline{numerical relations} in intersection theory on proper algebraic schemes over a given field $k$. This point of view will be developed in some detail in Exp. X, which is largely independent of Exp. VI, VII, VIII (centred on the proof of the Riemann--Roch theorem) and of Exp. IX. Some notions concerning the \uline{behaviour of} $K^\bullet(X),\Gr^\bullet(X)$ \uline{under specialization} will also be given there.

\sgaStableTarget{sga6.e00.sec.7}\subsection*{\normalfont 7.}
As in every theory of intersections, it is important to specify the place of the \uline{Picard group} in the theory (or the group of \uline{divisor classes} in the classical terminology). Subject to a slight restriction, automatically satisfied for example when $X$ is quasi-compact, Exp. X establishes a canonical isomorphism
\sgaStableTarget{sga6.e00.eq.7-1}\begin{equation}\tag{7.1}
c_1:\Pic(X)\xrightarrow{\sim}\Gr^1(X),
\end{equation}
which shows the exact role of the Picard group in the theory we propose. This role justifies a more detailed study of this Picard group in Exposés XII, XIII, which are logically independent of the text of the Seminar. A study of \uline{numerical equivalence} and certain \uline{finiteness theorems for the Picard group} (announced without proof in [5]) is made by Kleiman in Exp. XIII, using techniques of a purely numerical nature (which do not use the theory of the Picard functor and its representability), due to himself and to Mumford, markedly simpler than Grothendieck's initial proofs. The only result refractory to Kleiman's methods (and which he is obliged to use) is [5, C-07 (i)], whose proof and certain refinements (including certain representability theorems for the Picard functor) occupy Exposé XII. The latter moreover has a spirit and a technique entirely foreign to the rest of the Seminar, being connected with [5]; this ex%
%%%%%%%%%%  scan idx 31  |  folio 18  %%%%%%%%%%
\SGAPrintPage{18}%
posé is logically independent of all the others.

Finally, in Exp. XI, of a distinctly more general nature than XII and XIII and closer to the general spirit of the Seminar, we study, on an arbitrary locally ringed topos, a remarkable functor called the ``\uline{determinant functor}'':
\sgaStableTarget{sga6.e00.eq.7-2}\begin{equation}\tag{7.2}
\det^*:\underline{\operatorname{Parfis}}(X)\longrightarrow \underline{\Inv}(X)
\end{equation}
from the category of perfect complexes on $X$, whose morphisms consist only of the isomorphisms in the sense of the derived category $D(X)$, to the category of invertible Modules on $X$. This functor is described almost completely by the requirement that it be of a ``local nature'' and that, for a bounded complex $L_\bullet$ with locally free components, it be given by the formula
\sgaStableTarget{sga6.e00.eq.7-3}\begin{equation}\tag{7.3}
\det^*(L_\bullet)\simeq\bigotimes_i \det(L_i)^{(-1)^i},
\end{equation}
where $\det(L_i)$ denotes the exterior power of maximal degree of the locally free Module $L_i$. Although this exposé too is logically independent of the rest of the Seminar except for Exposé I, it nevertheless forms an important part of the general ``yoga'' proposed in the Seminar and developed mainly in Exposés I to XI.

%%%%%%%%%%  scan idx 32  |  folio 19  %%%%%%%%%%
\SGAPrintPage{19}
\sgaStableTarget{sga6.e00.bibliography}\section*{\underline{BIBLIOGRAPHY}}
\begin{enumerate}[label={[\arabic*]},leftmargin=2.5em]
\item M. Atiyah and F. Hirzebruch, Vector bundles and homogeneous spaces, Symposia in Pure Mathematics, vol. 3, Differential Geometry p. 7--38, 1961.
\item A. Borel and J.P. Serre, Le théorème de Riemann--Roch, Bull. Soc. math. France, t. 86, p. 97--136 (1958).
\item H. Grauert, Ein Theorem der analytischen Garbentheorie, Pub. Math. no.~5, p. 5--64 (1960).
\item A. Grothendieck, Classes de faisceaux et théorème de Riemann--Roch, mimeographed notes, Princeton 1957, (cited as [RRR] and reproduced as an Appendix to Exp. 0 of the Seminar).
\item A. Grothendieck, Technique de descente et théorèmes d'existence en Géométrie Algébrique, in Fondements de la Géométrie Algébrique, Extracts from the Bourbaki Seminar 1957--1962, Secrétariat Mathématique, Paris.
\item W. Shih, Séminaire de Topologie à l'IHES, 1966/67.
\item M. Atiyah and F. Hirzebruch, The Riemann--Roch theorem for analytic embeddings, Topology, vol. 1, p. 151--166 (1962).
\end{enumerate}

\clearpage

%%%%%%%%%%  scan idx 33  |  folio 20  %%%%%%%%%%
\SGAPrintPage{20}
\sgaStableTarget{sga6.e00.appendix-to-expose-0-rrr}
\begin{center}
{\Large \underline{\underline{CLASSES OF SHEAVES AND THE RIEMANN--ROCH THEOREM}}}\par
\vspace{0.5em}
\underline{by A. GROTHENDIECK}\SGAAuthorFootnote{*}{This is the verbatim reproduction of the report cited in the Introduction to the present Seminar. The footnotes indicated by the signs (*), (**) were added in October 1967.}
\end{center}

\sgaStableTarget{sga6.e00.chapter-i-rings-formal-preliminaries}\subsection*{CHAP. I -- $\lambda$-\uline{Rings} (\uline{formal preliminaries})}
\begin{longtable}{@{}p{0.74\linewidth}r@{}}
1. Definitions & 1\\
2. Examples & 4\\
3. The $\lambda$-ring defined by a graded ring & 8\\
4. The operations $\lambda^p(N,x)$ & 13\\
\end{longtable}

\sgaStableTarget{sga6.e00.chapter-ii-classes-of-coherent-algebraic-sheaves}\subsection*{CHAP. II -- \uline{Classes of coherent algebraic sheaves and Chern classes}\hfill 18}
\begin{longtable}{@{}p{0.74\linewidth}r@{}}
1. Chow's theory & 18\\
2. Definition of the Chern classes of coherent algebraic sheaves & 22\\
3. Functorial generalities on $K(X)$ & 27\\
4. Some technical results & 34\\
5. Sheaf-theoretic definition of Chern classes. Application to the study of injection morphisms & 43\\
6. The Riemann--Roch theorem & 49
\end{longtable}

\begin{center}
Proof of the Riemann--Roch theorem (1 November 1957).
\end{center}

\clearpage

%%%%%%%%%%  scan idx 34  |  folio 21  %%%%%%%%%%
\SGAPrintPage{21}
\sgaStableTarget{sga6.e00.chapter-i-rings-formal-preliminaries.occ2}\section*{CHAPTER I}
\begin{center}
$\lambda$-Rings (Formal preliminaries)
\end{center}

\sgaStableTarget{sga6.e00.sec.1.occ2}\subsection*{\S 1. Definitions}

A $\lambda$\uline{-ring}\SGAAuthorFootnote{*}{In this seminar we instead say \uline{pre-}$\lambda$\uline{-ring} (V 2.1), reserving the name $\lambda$-ring for those satisfying the additional condition on page 3 (cf V 2.4).} is a commutative ring $K$ with unit, equipped with a family of maps
\sgaStableTarget{sga6.e00.eq.1-1.occ2}\begin{equation}\tag{1.1}
\lambda^i:K\longrightarrow K
\end{equation}
($i$ an integer $\geq 0$) satisfying the following conditions:
\sgaStableTarget{sga6.e00.eq.1-2.occ2}\begin{equation}\tag{1.2}
\lambda^0x=1,\qquad \lambda^1x=x,\qquad
\lambda^n(x+y)=\sum_{i=0}^n\lambda^ix\cdot\lambda^{n-i}y.
\end{equation}
For every $x\in K$, put
\sgaStableTarget{sga6.e00.eq.1-3.occ2}\begin{equation}\tag{1.3}
\lambda_t(x)=\sum_{n\geq 0}\lambda^n(x)t^n\in K[[t]] ;
\end{equation}
the preceding relations then express that $x\longmapsto\lambda_t(x)$ is a homomorphism from the additive group $K$ to the multiplicative group $1+K[[t]]^+$ of formal series over $K$ with augmentation $1$, lifting the canonical homomorphism $1+\sum_{i>0}x_it^i\longrightarrow x_1$.

Thus, on the ring $\Z$, there exists a unique $\lambda$-structure compatible with its ring structure and such that
%%%%%%%%%%  scan idx 35  |  folio 22  %%%%%%%%%%
\SGAPrintPage{22}
\sgaStableTarget{sga6.e00.eq.1-4}\begin{equation}\tag{1.4}
\lambda_t(1)=1+t.
\end{equation}
We then have
\sgaStableTarget{sga6.e00.eq.1-5}\begin{equation}\tag{1.5}
\lambda_t(n)=(1+t)^n,
\end{equation}
hence $\lambda^i(n)=\binom{n}{i}$.

The notion of a $\lambda$-homomorphism of $\lambda$-rings is clear; an \uline{augmented} $\lambda$\uline{-ring} is a $\lambda$-ring equipped with a $\lambda$-homomorphism to $\Z$.

To define the notion of a \uline{special} $\lambda$\uline{-ring}, the following considerations are useful. Let $k$ be a commutative ring with unit, and let $K=1+k[[t]]^+$ be the group of formal series with coefficients in $k$ and augmentation $1$; we shall introduce on it a $\lambda$-ring law whose additive structure is multiplication of formal series. The multiplication of this ring will be denoted by $f\circ g$. It is completely determined by the condition that the coefficients $c_n=c_n(f\circ g)$ of $f\circ g$ are given by \uline{universal} polynomials with integer coefficients in the coefficients $a_i=c_i(f)$ and $b_i=c_i(g)$ of $f$ and $g$, respectively:
\sgaStableTarget{sga6.e00.eq.1-6}\begin{equation}\tag{1.6}
c_n=P_n(a_1,\ldots,a_n;b_1,\ldots,b_n),
\end{equation}
that it is distributive with respect to the ``addition'' $fg$, and is given for linear factors by:
\sgaStableTarget{sga6.e00.eq.1-7}\begin{equation}\tag{1.7}
(1+at)\circ(1+bt)=1+abt.
\end{equation}
The $P_n$ are calculated by calculations of symmetric functions; $P_n$ is isobaric of weight $n$ in the $a_i$ and isobaric of weight $n$ in the $b_i$ ($a_i$ and $b_i$ having weight $i$), and is moreover symmetric in $a=(a_i)$ and $b=(b_i)$. Likewise, the $\lambda^i(f)$ are determined unambiguously by the condition that
%%%%%%%%%%  scan idx 36  |  folio 23  %%%%%%%%%%
\SGAPrintPage{23}
the coefficients $c_n(\lambda^i(f))=c_n\text{ of }\lambda^i(f)$ are calculated using \uline{universal} polynomials with integer coefficients in the coefficients $c_n(f)=a_n$ of $f$:
\sgaStableTarget{sga6.e00.eq.1-8}\begin{equation}\tag{1.8}
c_n=P_{i,n}(a_1,\ldots,a_{in}),
\end{equation}
that relations (1.2) hold, and that, for $i>1$, we have:
\sgaStableTarget{sga6.e00.eq.1-9}\begin{equation}\tag{1.9}
\lambda^i(1+at)=1
\end{equation}
($1$ is the zero element of the ring $K=1+k[[t]]^+$). The $P_{i,n}$ are again calculated by calculations of symmetric functions, and we find that $P_{i,n}$ is an isobaric polynomial of weight $in$ in the $a_i$. Equipped with the preceding operations, $K$ becomes a $\lambda$-ring; its zero element is $1$, and its unit element is $1+t$. We also note the relation:
\sgaStableTarget{sga6.e00.eq.1-7bis}\begin{equation}\tag{1.7 bis}
(1+at)\circ f=f(at)
\end{equation}

Now let $K$ be an arbitrary $\lambda$-ring. We say that $K$ is \uline{special} when the additive homomorphism $\lambda_t$ from $K$ to $1+K[[t]]^+$ is a $\lambda$-homomorphism. This therefore means that we have the formulas:
\sgaStableTarget{sga6.e00.eq.1-10}\begin{equation}\tag{1.10}
\lambda_t(1)=1+t,\quad
\lambda_t(xy)=\lambda_t(x)\circ\lambda_t(y),\quad
\lambda_t(\lambda^i(x))=\lambda^i(\lambda_t(x)),
\end{equation}
or, explicitly:
\sgaStableTarget{sga6.e00.eq.1-11}\begin{equation}\tag{1.11}
\begin{cases}
\lambda^i(1)=0 & \text{for }i>1\\
\lambda^n(xy)=P_n(\lambda^1x,\ldots,\lambda^n x;\lambda^1y,\ldots,\lambda^n y)\\
\lambda^n(\lambda^i(x))=P_{i,n}(\lambda^1x,\ldots,\lambda^{in} x),
\end{cases}
\end{equation}
where the $P_n$ and $P_{i,n}$ are the universal polynomials considered above. In particular, formulas (1.5) follow. (We can show that the $\lambda$-ring $1+k[[t]]^+$ constructed from a commutative ring $k$ is spe%
%%%%%%%%%%  scan idx 37  |  folio 24  %%%%%%%%%%
\SGAPrintPage{24}%
cial; but we shall not need this fact).

\sgaStableTarget{sga6.e00.sec.2.occ2}\subsection*{\S 2. Examples}

\sgaStableTarget{sga6.e00.ex.2-1}1) Let $G$ be a group and $k$ a commutative ring; we consider the group $\underline{K}(G)$, the quotient of the free group generated by the classes of representations of $G$ in $k$-modules of finite type by the subgroup generated by elements of the form $(\rho\oplus\rho')-(\rho)-(\rho')$ (we write $\oplus$ for direct sum). Tensor product and exterior powers define a $\lambda$-ring structure on $\underline{K}(G)$. We can also pass to the quotient by elements of the form $(\rho)-(\rho')-(\rho'')$, where $\rho$ is a representation that is an extension of $\rho'$ by $\rho''$, trivial as an extension of $k$-modules. We obtain the group $\underline{K}_r(G)$ of reduced classes of representations of $G$, and we find that the $\lambda$-ring law passes to the quotient.

When $k$ is an \uline{algebraically closed field of characteristic} $0$\SGAAuthorFootnote{*}{It is unnecessary for $k$ to be algebraically closed, cf. (VI 3.3) for a still more general result (including all those that will be used in the present report).}, we can show that $\underline{K}(G)$, and a fortiori $\underline{K}_r(G)$, is a \uline{special} $\lambda$-ring; this no longer holds when the characteristic is $\neq 0$, but nevertheless, when $k$ is algebraically closed\SGAAuthorFootMark{*}, $\underline{K}_r(G)$ is still a special $\lambda$-ring. To prove this, we reduce to the case where $G$ is a product $\Gl(m,k)\times\Gl(n,k)$ and only rational representations of $G$ are considered; our assertion follows from the \uline{explicit determination} of $\underline{K}_r(G)$ \uline{to be made below} (example 3); finally, when $k$ has characteristic $0$, complete reducibility of the rational representations of $G$ shows that $\underline{K}(G)=\underline{K}_r(G)$.

%%%%%%%%%%  scan idx 38  |  folio 25  %%%%%%%%%%
\SGAPrintPage{25}
Variants of the preceding example: $G$ is an algebraic group defined over $k$, and only the rational representations of $G$ defined over $k$ are considered; we can consider a topological group, or a Lie group, and continuous representations, etc\dots\SGAAuthorFootnote{1}{Other examples of $\lambda$-rings can be constructed, for example the set of equivalence classes of quadratic forms, the set of classes of representations of a Lie algebra, etc\dots We promise that we shall not need all these examples to treat the Riemann--Roch theorem!}

\sgaStableTarget{sga6.e00.ex.2-2}2) Let $X$ be an algebraic variety, irreducible or not, defined over the field $k$. We can consider the quotient of the free group generated by the classes of vector bundles on $X$ (defined over $k$) by the subgroup generated by elements of the form $(E\oplus E')-(E)-(E')$, on which the $\lambda$-ring operations will be defined by tensor product and exterior powers. When $X$ is complete, this is a free group generated by the indecomposable vector bundles on $X$ (by Krull--Remak--Schmidt); knowing this ring and its basis is equivalent to the complete classification of vector bundles on $X$ and knowing the elementary tensor operations on those bundles. When $k$ is algebraically closed of characteristic $0$, all the other tensor operations on classes of bundles are then known, using what follows in 3). We can make a more drastic passage to the quotient by the group generated by the elements $(E)-(E')-(E'')$, where $E$ is a vector bundle that is an extension of $E'$ by $E''$; we then obtain a $\lambda$-ring, denoted by $\underline{K}(\xi(X))$, where $\xi(X)$ denotes the additive category of vector bundles on $X$ defined over $k$. Using the results of example 1), we show that, when $k$ is algebraically closed\SGAAuthorFootnote{*}{Unnecessary condition, cf. the preceding footnote.}, the $\lambda$-ring $\underline{K}(\xi(X))$ is special, and the same holds for the
%%%%%%%%%%  scan idx 39  |  folio 26  %%%%%%%%%%
\SGAPrintPage{26}%
first $\lambda$-ring defined here, if the characteristic of $k$ is $0$, but not in the general case.

\sgaStableTarget{sga6.e00.ex.2-3}3) \uline{Determination of} $\underline{K}(G)$ \uline{and of} $\underline{K}_r(G)$ \uline{in certain important cases}: We suppose that the field $k$ is algebraically closed, that the group $G$ is algebraic, \uline{affine connected} and ``globally \uline{reductive}'' (the radical of $G$ is isomorphic to $k^{*s}$), and we consider only rational linear representations. In characteristic $0$, the complete reducibility of the rational linear representations of $G$ shows that $\underline{K}(G)=\underline{K}_r(G)$ is the free group generated by the classes of irreducible rational representations. Now, if $T$ is a maximal torus of $G$, we have $\underline{K}(T)=\underline{K}_r(T)=\Z(\That)$, the algebra of the group $\That$ of rational characters of $T$; we have $\That\simeq\Z^r$ if $r$ is the rank of $G$ (this is independent of the characteristic). The Weyl group $W$ operates on $T$, hence on $\underline{K}(T)$, in a way that depends only on the operations of $W$ on the free group $\That$. We have an evident homomorphism of $\lambda$-rings, deduced from the injection of $T$ into $G$:
\sgaStableTarget{sga6.e00.eq.1-13}\begin{equation}\tag{1.13}
\underline{K}_r(G)\longrightarrow \underline{K}_r(T)^W=\underline{K}(T)^W,
\end{equation}
where $\underline{K}(T)^W$ denotes the set of fixed points of $W$ in $\underline{K}(T)$. The theory of linear representations of $G$\SGAAuthorFootnote{*}{Cf.\ Séminaire Chevalley 1956/58 Groupes de Lie algébriques (École Normale Supérieure).} then proves the following theorem:

\sgaStableTarget{sga6.e00.thm.1-1}\textbf{Theorem 1.1.} \uline{The homomorphism} (1.13) \uline{is an isomorphism}.

We thus see that $\underline{K}_r(G)$, once the ``type'' of $G$ is known (characterized by $W$ operating on a group $\Z^r$), does not depend on the characteristic, and that it is a special $\lambda$-ring. We have a canonical basis of $\underline{K}_r(G)$ corresponding to the orbits of $W$ in $\That$; another basis is obtained as follows:
%%%%%%%%%%  scan idx 40  |  folio 27  %%%%%%%%%%
\SGAPrintPage{27}%
let $(\rho_i)_{1\leq i\leq r'}$ be the representation classes of $G$ corresponding to the \uline{fundamental} representations of $G/\rad G$, and let $(\sigma_i)_{1\leq i\leq r''}$ be a basis of the group of rational homomorphisms from $G$ to $k^*$; we thus have $r'+r''=r$ (the rank of $G$).

\sgaStableTarget{sga6.e00.cor.unnamed}\textbf{Corollary:} \uline{The ring} $\underline{K}_r(G)$ \uline{is identified with the subring}
\[
\Z[\rho_1,\ldots,\rho_{r'};\sigma_1,\ldots,\sigma_{r''};\sigma_1^{-1},\ldots,\sigma_{r''}^{-1}]
\]
\uline{of the field of rational functions in the} $\rho_i$ \uline{and the} $\sigma_j$ \uline{with coefficients in} $\Q$\SGAAuthorFootnote{*}{The derived group $G'$ of $G$ must be assumed simply connected for this statement to be correct.}.

In particular, $\underline{K}_r(G)$ has no divisors of $0$; if $G=\prod_i\Gl(n_i,k)$, the ring $\underline{K}_r(G)$ is made explicit as follows: let $\rho_i$ be the representation of $G$ deduced from the identity representation of $\Gl(n_i,k)$; if we then consider the field of rational functions over $\Q$ in indeterminates $\lambda^j(\rho_i)$ ($1\leq j\leq n_i$), then $\underline{K}_r(G)$ is identified with the subring generated by these indeterminates and the inverses of the $\lambda^{n_i}(\rho_i)$. This justifies our assertion that, at least in characteristic $0$, the operations $\lambda^i$ allow all tensor operations on vector bundles to be reconstructed; this is still true over an algebraically closed field of arbitrary characteristic, provided reduced classes are taken.

\sgaStableTarget{sga6.e00.rem.unnamed}\textbf{Remarks:} a) In example 1), the ring $\underline{K}(G)$ is augmented, the augmentation being defined by the degree of the representations, and, by passage to the quotient, this augmentation defines an augmentation on $\underline{K}_r(G)$. Likewise, the $\lambda$-rings considered in example 2) are augmented when $X$ is $k$-irreducible; in $\underline{K}(\xi(X))$, the augmentation is defined by the \uline{rank of a vector bundle}, that is, the dimension of its fibres.

%%%%%%%%%%  scan idx 41  |  folio 28  %%%%%%%%%%
\SGAPrintPage{28}%
b)
When the field $k$ is not algebraically closed, it is plausible that $\underline{K}(G)$ is a special $\lambda$-ring in characteristic $0$, and likewise for $\underline{K}_r(G)$ in all characteristics.

c) In nonzero characteristic, theorem 1.1 was kindly supplied to me by Chevalley.

\sgaStableTarget{sga6.e00.sec.3.occ2}\subsection*{\S 3. The $\lambda$-ring defined by a graded ring}

Let $A$ be a commutative graded ring with unit; for simplicity we shall assume that the degrees of $A$ are positive and that $A^0=\Z$. We denote by $\Ahat$ the product of the $A^i$ (it is a commutative ring with unit), and by $\Ahp$ the kernel of the evident augmentation $\Ahat\to A^0=\Z$. Then $1+\Ahp$ is a subgroup of the multiplicative group of invertible elements of $\Ahat$. Consider the abelian group
\sgaStableTarget{sga6.e00.eq.1-14}\begin{equation}\tag{1.14}
\widetilde A=\Z\times(1+\Ahp),
\end{equation}
whose elements will be denoted $[n,x]$ with $n\in\Z$ and $x=1+\sum_{i\geq 1}x^i\in1+\Ahp$ ($x^i\in A^i$). We shall write this group additively, thus:
\sgaStableTarget{sga6.e00.eq.1-15}\begin{equation}\tag{1.15}
[n,x]+[n',x']=[n+n',xx'].
\end{equation}
The zero element of this group is $[0,1]$. We are going to introduce on it an augmented $\lambda$-ring structure, compatible with its additive structure, the augmentation being $\varepsilon:[n,x]\longrightarrow n$. The product $[m,x][n,y]$ is completely characterized by being written $[mn,x^n y^m(x*y)]$, where $x*y$ is an element of $1+\Ahp$ whose components are expressed as universal polynomials with integer coefficients in the $x^i$ and the $y^i$:
\sgaStableTarget{sga6.e00.eq.1-16}\begin{equation}\tag{1.16}
(x*y)^i=Q_i(x^1,\ldots,x^i;y^1,\ldots,y^i)\qquad (i\geq1)
\end{equation}
%%%%%%%%%%  scan idx 42  |  folio 29  %%%%%%%%%%
\SGAPrintPage{29}%
the
polynomial $Q_i$ being isobaric of weight $i$. Moreover, $x*y$ must be distributive with respect to the ``addition'' $xy$ and, for ``linear'' factors, must reduce to the following:
\sgaStableTarget{sga6.e00.eq.1-17}\begin{equation}\tag{1.17}
[1,1+x^1][1,1+y^1]=[1,1+x^1+y^1].
\end{equation}
The polynomials $Q_i$ are evaluated by calculations with elementary symmetric functions (this is the calculation of the Chern classes of a tensor product of vector bundles by means of the Chern classes $x^i$ and $y^i$ of the factors and their respective dimensions $m$ and $n$). From (1.17) we deduce:
\sgaStableTarget{sga6.e00.eq.1-17bis}\begin{equation}\tag{1.17 bis}
(1+x^1)*(1+y^1)=(1+x^1+y^1)/(1+x^1)(1+y^1),
\end{equation}
and an arbitrary product $x*y$ is calculated term by term by decomposing ``formally'' $x$ into a product $\prod_{i=1}^n(1+\alpha_i)$ and $y$ into a product $\prod_{i=1}^n(1+\beta_i)$, where the $\alpha_i$ and the $\beta_i$ have degree $1$.

In this way, $\widetilde A$ becomes an augmented commutative ring, with a unit $[1,1]$, and it can be interpreted as the ring obtained by adjoining a unit to the ring $1+\Ahp$. Moreover, on $\widetilde A$, we define a \uline{filtration} by putting $\widetilde A^0=\widetilde A$ and denoting by $\widetilde A^n$ the set of the $[0,x]$ with $x^i=0$ for $1\leq i<n$; this filtration is compatible with the product, for we have
\sgaStableTarget{sga6.e00.eq.1-18}\begin{equation}\tag{1.18}
\left(1+\sum_{i\geq m}x^i\right)*\left(1+\sum_{j\geq n}y^j\right)
=1-\frac{(m+n-1)!}{(m-1)!(n-1)!}x^m y^n+\cdots,
\end{equation}
the unwritten terms having degree $>m+n$. This formula shows that the graded ring $G(\widetilde A)$ associated with $\widetilde A$ is identified with $A$ from the additive point of view, but not from the multiplicative point of view; we can only say that the map $\eta'$ from $A$ to $G(\widetilde A)$ defined by
%%%%%%%%%%  scan idx 43  |  folio 30  %%%%%%%%%%
\SGAPrintPage{30}%
\sgaStableTarget{sga6.e00.eq.1-19}\begin{equation}\tag{1.19}
\eta'(x^i)=(-1)^{i-1}(i-1)!\,x^i\qquad (x^i\in A^i)
\end{equation}
is a \uline{ring homomorphism}.

We introduce on $\widetilde A$ operations $\lambda^i$ satisfying the axioms of $\lambda$-rings (\S 1, formulas (1.2)), well determined by the condition that the components of $\lambda^i[0,x]$ be deduced from the components $x^i$ of $x$ by \uline{universal} polynomials with integer coefficients
\sgaStableTarget{sga6.e00.eq.1-20}\begin{equation}\tag{1.20}
\lambda^i[0,x]=[0,\lambda^i x]\qquad (\lambda^i x)^{(n)}=Q_{i,n}(x^1,\ldots,x^n)
\end{equation}
($Q_{i,n}$ is isobaric of weight $n$), and that we have
\sgaStableTarget{sga6.e00.eq.1-21}\begin{equation}\tag{1.21}
\lambda^i[1,1+x^1]=0\qquad\text{for }i>1.
\end{equation}
The polynomials $Q_{i,n}$ are evaluated by a calculation with elementary symmetric functions, which is the calculation of the Chern classes of the exterior powers of a vector bundle, by means of the Chern classes and the rank of the latter. We then find that $\widetilde A$ is a \uline{special} augmented $\lambda$-ring and that the $\widetilde A^i$ are stable under the operations $\lambda^j$.

The augmented $\lambda$-ring $\widetilde A$ may be considered as a functor in $A$; consequently, the principal automorphism
\[
x=\sum_{i\geq0}x^i\longrightarrow \check x=\sum_{i\geq0}(-1)^i x^i
\]
of $A$ defines an automorphism of $\widetilde A$, denoted by the same symbol $\check{\ }$.

\sgaStableTarget{sga6.e00.prop.1-2}\textbf{Proposition 1.2.} \uline{For} $x\in\widetilde A$ \uline{of the form} $[n,1+x^1+\cdots+x^n]$ (we therefore have $x^i=0$ for $i>n$), \uline{we have} $\lambda^i(x)=0$ \uline{for} $i>n$, \uline{we have} $\lambda^n(x)=[1,1+x^1]$, \uline{and finally}
\sgaStableTarget{sga6.e00.eq.1-22}\begin{equation}\tag{1.22}
\sum_{i=0}^n(-1)^i\lambda^i(x)=[0,1-(n-1)!x^n+\cdots]
\end{equation}
From this proposition and from the fact that the $\widetilde A^i$ are stable under the operations $\lambda^j$, we deduce that formula (1.22) is valid whenever $\varepsilon(x)=n$, whether or not we have $x^i=0$ for $i>n$.

%%%%%%%%%%  scan idx 44  |  folio 31  %%%%%%%%%%
\SGAPrintPage{31}%
Formula
(1.22) admits the following generalizations: let $K$ be any $\lambda$-ring, otherwise arbitrary; for $x\in K$, put
\sgaStableTarget{sga6.e00.eq.1-23}\begin{equation}\tag{1.23}
\gamma^n(x)=(-1)^n\sum_{i=0}^n(-1)^i\lambda^i(x+n)=\lambda^n(x+n-1)
\end{equation}
(the identity of the last two members follows from an immediate induction). In particular, we have $\gamma^0(x)=1$ and $\gamma^1(x)=x$; next put:
\sgaStableTarget{sga6.e00.eq.1-24}\begin{equation}\tag{1.24}
\gamma_t(x)=\sum_{n\geq0}\gamma^n(x)t^n\in K[[t]]
\end{equation}
whence we deduce the formula:
\sgaStableTarget{sga6.e00.eq.1-25}\begin{equation}\tag{1.25}
\gamma_t(x)=\lambda_{t/(1-t)}(x),
\end{equation}
which is equivalent to
\sgaStableTarget{sga6.e00.eq.1-25bis}\begin{equation}\tag{1.25 bis}
\lambda_s(x)=\gamma_{s/(1+s)}(x).
\end{equation}
This proves that the $\lambda^i$ are expressed by means of the $\gamma^i$ and that $\gamma_t(x+y)=\gamma_t(x)\gamma_t(y)$, and consequently that the maps $\gamma^i$ define a new $\lambda$-ring structure on $K$. Formula (1.22) then shows that, for $K=\widetilde A$ and every $x\in\widetilde A$, we have
\sgaStableTarget{sga6.e00.eq.1-22bis}\begin{equation}\tag{1.22 bis}
\gamma^n(x-\varepsilon(x))=[0,1+(-1)^{n-1}(n-1)!x^n+\cdots]
\end{equation}
and consequently, the \uline{component} $x^n$ \uline{is obtained, up to the factor} $(-1)^{n-1}(n-1)!$, \uline{by reducing modulo} $\widetilde A^{n+1}$ \uline{the element} $\gamma^n(x-\varepsilon(x))$ \uline{of} $\widetilde A^n$.

\sgaStableTarget{sga6.e00.the-chern-homomorphism}\paragraph{The Chern homomorphism:}
For every graded ring $A$ satisfying the preceding conditions, we are going to define a homomorphism of augmented rings:
\sgaStableTarget{sga6.e00.eq.1-26}\begin{equation}\tag{1.26}
\ch:\widetilde A\longrightarrow \widehat{A\otimes\Q}.
\end{equation}
This homomorphism is defined unambiguously by the conditions that it be an additive homomorphism, take $1$ to $1$, be such that the com%
%%%%%%%%%%  scan idx 45  |  folio 32  %%%%%%%%%%
\SGAPrintPage{32}%
ponents of $\ch(x)$ in degree $>0$ be given by universal isobaric polynomials with rational coefficients in the components of $x$, and that
\sgaStableTarget{sga6.e00.eq.1-27}\begin{equation}\tag{1.27}
\ch([1,1+x^1])=\exp x^1=\sum_{n\geq0}(x^1)^n/n!.
\end{equation}
We then immediately obtain the explicit expression for $\ch$: let $\eta$ be the additive endomorphism of $\widehat{A\otimes\Q}$ defined, for $x^i$ of degree $i$, by
\sgaStableTarget{sga6.e00.eq.1-28}\begin{equation}\tag{1.28}
\eta(x^i)=1/(-1)^{i-1}(i-1)!\,x^i
\end{equation}
we then have
\sgaStableTarget{sga6.e00.eq.1-29}\begin{equation}\tag{1.29}
\ch\left([n,1+\sum_{i\geq1}x^i]\right)=n+\eta\left(\log\left(1+\sum_{i\geq1}x^i\right)\right).
\end{equation}
As we can ``invert'' this formula, at least when $A$ has only finitely many homogeneous components, we see that, in this case, we have an \uline{isomorphism} $\ch$ \uline{from the ring} $\widetilde A\otimes\Q$ \uline{onto the ring} $\widehat{A\otimes\Q}$ (which completely elucidates the structure of $\widetilde A\otimes\Q$).

Recall that, according to Hirzebruch [1], a formal series $f\in\Q[[t]]$ defines, by universal polynomial formulas, an additive homomorphism $1+\Ahp\longrightarrow 1+(\widehat{A\otimes\Q})^+$, denoted $\mathcal T_f$; when $f(t)=t/(1-\exp(-t))$, we simply put $\mathcal T_f=\mathcal T$ ($\mathcal T$ is the initial of Todd). We extend the definition of $\mathcal T(x)$ to $x\in\widetilde A$. This said, we have the following result:

\sgaStableTarget{sga6.e00.prop.1-3}\textbf{Proposition 1.3.} \uline{Let} $N\in\widetilde A$ \uline{be of the form} $[q,1+N^1+\cdots+N^q]$ (we have $N^i=0$ for $i>q$). \uline{If we put}:
\[
\lambda_{-1}(N)=\sum_{i\geq0}(-1)^i\lambda^i(N)=\sum_{i=0}^{q}(-1)^i\lambda^i(N),
\]
\uline{we shall have}:
\sgaStableTarget{sga6.e00.eq.1-30}\begin{equation}\tag{1.30}
\ch(\lambda_{-1}(N))=(-1)^q N^q\mathcal T(\check N)^{-1},
\end{equation}
%%%%%%%%%%  scan idx 46  |  folio 33  %%%%%%%%%%
\SGAPrintPage{33}%
\uline{which is equivalent to}:
\sgaStableTarget{sga6.e00.eq.1-30bis}\begin{equation}\tag{1.30 bis}
\ch(\lambda_{-1}(\check N))=N^q\mathcal T(N)^{-1}.
\end{equation}

\sgaStableTarget{sga6.e00.sec.4.occ2}\subsection*{\S 4. The operations $\lambda^p(N,x)$}

Let $A$ be the ring of formal series, with coefficients in $\Z$, in two infinite sequences of indeterminates, denoted $\lambda^iN$ and $\lambda^i x$ ($i\geq1$). Put $\lambda^1N=N$ and $\lambda^1x=x$, and introduce on $A$ the unique special $\lambda$-ring structure for which we have $\lambda^i(N)=\lambda^iN$ and $\lambda^i(x)=\lambda^ix$ and for which the operations $\lambda^i$ are ``continuous'' (if we restricted ourselves to \uline{polynomials} in indeterminates $\lambda^iN$ and $\lambda^ix$, we would have the ``free $\lambda$-ring generated by $N$ and $x$''; our ring $A$ is a completion of it). We can then form
\sgaStableTarget{sga6.e00.eq.1-31}\begin{equation}\tag{1.31}
\lambda_{-1}(N)=\sum_{i\geq0}(-1)^i\lambda^iN,
\end{equation}
which is a formal series with constant term $1$, hence invertible; we can therefore define an element $\lambda^p(N,x)$ of $A$ unambiguously by putting
\sgaStableTarget{sga6.e00.eq.1-32}\begin{equation}\tag{1.32}
\lambda^p(N,x)\lambda_{-1}(N)=\lambda^p(x\lambda_{-1}(N))
\end{equation}

\sgaStableTarget{sga6.e00.thm.1-4}\textbf{Theorem 1.4.} \uline{Let} $J_q$ \uline{be the closed ideal of} $A$ \uline{generated by the} $\lambda^iN$ \uline{with} $i>q$; \uline{identify} $A/J_q$ \uline{with the ring of formal series in the} $\lambda^i x$ \uline{and the} $\lambda^jN$ ($1\leq j\leq q$). \uline{Then} $\lambda^p(N,x)$ \uline{reduced modulo} $J_q$ \uline{is a} polynomial, \uline{isobaric of weight} $p$ \uline{with respect to the variables} $\lambda^i x$.

\uline{Let} $N$ \uline{and} $F$ \uline{be two vector spaces of finite dimensions} $q,q'$ \uline{over an algebraically closed field} $k$ \uline{of characteristic} $0$; \uline{let} $N^{(p)}$ \uline{be the kernel of the map} $(a_1,\ldots,a_p)\longrightarrow a_1+\cdots+a_p$ from $N^p$ \uline{to} $N$.
We
%%%%%%%%%%  scan idx 47  |  folio 34  %%%%%%%%%%
\SGAPrintPage{34}%
\uline{let the symmetric group} $\mathfrak S_p$ \uline{operate on} $N^p$, \uline{hence on} $N^{(p)}$ \uline{hence on} $\Lambda(N^{(p)})$, \uline{then on} $\bigotimes^p F$ \uline{and on} $\Lambda(N^{(p)})\otimes\bigotimes^p F$. \uline{Consider this vector space as a graded module over} $\mathfrak S_p\times\Gl(N)\times\Gl(F)$ \uline{and take its ``alternating part'', which is a graded module over} $\Gl(N)\times\Gl(F)\eqqcolon G$. \uline{We then have}
\sgaStableTarget{sga6.e00.eq.1-33}\begin{equation}\tag{1.33}
\lambda^p(N,F)=\underline{E}_G\left(\left(\Lambda(N^{(p)})\otimes\bigotimes^p F\right)^{\mathrm{alt}}\right).
\end{equation}

N.B. If $M$ is a complex of $G$-modules, we denote by $\underline{E}_G(M)$ the alternating sum in $\underline{K}(G)$ of the classes of the components of $M$; this is the Euler--Poincaré characteristic of $M$; in formula (1.33), the first member $\lambda^p(N,F)$ denotes the element of the ring $\underline{K}(G)$ obtained by substituting for the variables $\lambda^iN$ and $\lambda^iF$ the classes in $\underline{K}(G)$ of the $G$-modules $\Lambda^i(N)$ and $\Lambda^i(F)$ respectively.

\textbf{Proof:} By expanding the second member of (1.32) by the second formula (1.11), we see that the first assertion is proved if we prove the analogous assertion obtained by putting $x=1$, i.e. if we prove that, in the $\lambda$-ring $\Z[\lambda^1N,\ldots,\lambda^qN]$ where $\lambda^iN=0$ for $i>q$, the element $\lambda^p(\lambda_{-1}(N))$ is divisible by $\lambda_{-1}(N)$. A priori, the quotient $\lambda^p(N,1)=\lambda^p(\lambda_{-1}(N))/\lambda_{-1}(N)$ is a formal series contained in the field of fractions $K$ of the ring of formal series in the $\lambda^iN$ for $1\leq i\leq q$. On the other hand, also denote by $N$ a vector space of dimension $q$ over an algebraically closed field $k$ of characteristic $0$, and consider the element $\ell^p(N)=\underline{E}_G((\Lambda(N^{(p)}))^{\mathrm{alt}})\in \underline{K}(G)$, where we now have $G=\Gl(N)$. We know (\S 2, theorem 1.1) that $\underline{K}(G)$ is identified
%%%%%%%%%%  scan idx 48  |  folio 35  %%%%%%%%%%
\SGAPrintPage{35}%
with the
subring of $K$ generated by the $\lambda^iN$ for $1\leq i\leq q$ and by the inverse of $\lambda^qN$. On the other hand, we readily verify in this ring that we have $\ell^p(N)\lambda_{-1}(N)=\lambda^p(\lambda_{-1}(N))$. As we are in a large field $K$ and $\lambda_{-1}(N)\neq0$, we conclude that $\ell^p(N)=\lambda^p(N,1)$; both members are in the intersection of $\underline{K}(G)$ and the ring of formal series in the $\lambda^iN$, and are therefore polynomials in the $\lambda^iN$.

To prove (1.33) in the general case, we verify that the second member satisfies in $\underline{K}(G)$ the formula analogous to (1.32), that is, its product by $\lambda_{-1}(N)$ is $\lambda^p(F\lambda_{-1}(N))$. Likewise, the product of $\lambda^p(N,F)\in \underline{K}(G)$ by $\lambda_{-1}(N)$ is $\lambda^p(F\lambda_{-1}(N))$, as follows from the fact that $\underline{K}(G)$ is \uline{special}. As $\underline{K}(G)$ is an integral domain (cf.\ \S 2) and $\lambda_{-1}(N)\neq0$ (same reference), we have indeed proved formula (1.33). q.e.d.

\sgaStableTarget{sga6.e00.rem.unnamed.occ2}\textbf{Remarks:} a) When the dimension $q'$ of $F$ is $\geq p$, then (1.33) characterizes the polynomial $\lambda^p(N,x)$ in the $\lambda^iN$ ($1\leq i\leq q$) and the $\lambda^jx$ ($1\leq j\leq p$).

b) It is plausible that (1.33) still holds when the field $k$ is not algebraically closed\SGAAuthorFootnote{*}{This is indeed proved by the same method.}. This question would have to be resolved to prove the Riemann--Roch theorem over a field of characteristic $0$ that is not algebraically closed by the method we are going to explain. Moreover, the passage from an $\mathfrak S_p$-module to its alternating part has a reasonable meaning only when the base field has characteristic $0$.

Now let $K$ be any \uline{special} $\lambda$-ring and let $N$ be an
ele%
%%%%%%%%%%  scan idx 49  |  folio 36  %%%%%%%%%%
\SGAPrintPage{36}%
ment of $K$ such that $\lambda^iN=0$ for sufficiently large $i$. We can then, for every $x\in K$ and every integer $p\geq1$, consider the elements $\lambda^p(N,x)$, which are expressed by universal polynomials with integer coefficients in the $\lambda^iN$ and the $\lambda^jx$. We have $\lambda^1(N,x)=x$ and we have formula (1.32) \uline{in} $K$. To express the properties of the operations $x\mapsto\lambda^p(N,x)$, it is convenient to introduce a new $\lambda$-ring denoted $K_N$. As a commutative ring, $K_N$ is obtained by adjoining a unit to the ring whose underlying additive group is $K$ and whose multiplication is given by
\sgaStableTarget{sga6.e00.eq.1-34}\begin{equation}\tag{1.34}
x_Ny=xy\mu\qquad (\mu=\lambda_{-1}(N)).
\end{equation}
The map $\lambda_t$ from $K_N$ to $K_N[[t]]$ is defined by:
\sgaStableTarget{sga6.e00.eq.1-35}\begin{equation}\tag{1.35}
\lambda_t(n,x)=(1+t)^n\sum_{p\geq0}\lambda^p(N,x)t^p,
\end{equation}
where we have put $\lambda^0(N,x)=(1,0)$ (the unit element of $K_N$). I say that, \uline{equipped with these operations}, $K_N$ \uline{is a special} $\lambda$\uline{-ring}. We therefore have to verify the formulas $\lambda_t(X+Y)=\lambda_t(X)\lambda_t(Y)$, $\lambda_t(XY)=\lambda_t(X).\lambda_t(Y)$, and $\lambda_t(\lambda^i(X))=\lambda^i(\lambda_t(X))$ for $X,Y\in K_N$. We can restrict ourselves to verifying these formulas when $X$ and $Y$ are in $K$, and we find that it suffices to verify them when $K$ is the free special $\lambda$-ring generated by $X,Y$ and $N$ subject to the relations $\lambda^iN=0$ for $i>q$. But the group homomorphism
\sgaStableTarget{sga6.e00.eq.1-36}\begin{equation}\tag{1.36}
K_N\longrightarrow K,
\end{equation}
which takes unit to unit and on $K$ reduces to the homomorphism $x\longrightarrow x\mu$, is a homomorphism compatible with the ring structures and the maps $\lambda^i$, as we verify immediately. In the present case,
%%%%%%%%%%  scan idx 50  |  folio 37  %%%%%%%%%%
\SGAPrintPage{37}%
its
restriction to $K$ is \uline{injective}; as $K$ is a special $\lambda$-ring, it follows readily that the above relations do hold for $X,Y\in K$, and this completes the proof of our assertion.

According to formula (1.23), it is natural, for $x\in K$ and $n\geq1$, to put
\sgaStableTarget{sga6.e00.eq.1-37}\begin{equation}\tag{1.37}
\gamma^n(N,x)=\sum_{i=0}^{n-1}\binom{n-1}{i}\lambda^{n-i}(N,x),
\end{equation}
in other words, we put $\gamma^n(N,x)=\gamma^n(x_N)$, where $x_N$ denotes $x$ considered as an element of the $\lambda$-ring $K_N$.

\sgaStableTarget{sga6.e00.prop.1-5}\textbf{Proposition 1.5.} \uline{Let} $A$ \uline{be a commutative graded ring with unit}, $N$ \uline{an element of} $\widetilde A$ \uline{of the form} $N=[q,1+N^1+\cdots+N^q]$, \uline{and} $x$ \uline{an arbitrary element of} $\widetilde A$: $x=[\nu,1+\sum_{i\geq1}x^i]$. \uline{Then, for every integer} $j\geq1$, \uline{we have} $\gamma^{q+j}(N,x)\in\widetilde A^j$, \uline{and the component of degree} $j$ \uline{of} $\gamma^{q+j}(N,x)$ \uline{is of the form}
\sgaStableTarget{sga6.e00.eq.1-38}\begin{equation}\tag{1.38}
\left(\gamma^{q+j}(N,x)\right)^{(j)}=(-1)^{j-1}(j-1)!G_{q,j}(\nu,x^1,\ldots,x^j;N^1,\ldots,N^j),
\end{equation}
\uline{where} $G_{q,j}$ \uline{is a universal polynomial with integer coefficients characterized by the condition}
\sgaStableTarget{sga6.e00.eq.1-39}\begin{equation}\tag{1.39}
G_{q,j}(\nu,x^1,\ldots,N^j)N^q=(-1)^q(x*\lambda_{-1}(N))^{(q+j)}.
\end{equation}

\textbf{Proof.} We can assume that $A$ is the ring of polynomials with integer coefficients in indeterminates $x^i$ and $N^j$ (with $1\leq j\leq q$). As we evidently have, combining (1.32) and (1.37),
\sgaStableTarget{sga6.e00.eq.1-40}\begin{equation}\tag{1.40}
\gamma^n(N,x)\lambda_{-1}(N)=\gamma^n(x\lambda_{-1}(N)),
\end{equation}
and, for $n=q+j$, the second member has filtration $\geq q+j$ (use formula (1.22 bis), taking into account that $x\lambda_{-1}(N)$ has
%%%%%%%%%%  scan idx 51  |  folio 38  %%%%%%%%%%
\SGAPrintPage{38}%
zero augmentation), whereas $\lambda_{-1}(N)$ has filtration $q$ (formula (1.22)), it follows that $\gamma^n(N,x)$ has filtration $j$, since, by (1.18), the associated graded ring of $A$ is here an \uline{integral domain}. Using the preceding formulas, we immediately obtain the indicated form for $(\gamma^{q+j}(N,x))^{(j)}$.
\begin{flushright}q.e.d.\end{flushright}

\sgaStableTarget{sga6.e00.chapter-ii-classes-of-coherent-algebraic-sheaves.occ2}\section*{CHAPTER II}
\begin{center}
Classes of coherent algebraic sheaves\par
and Chern classes
\end{center}

\sgaStableTarget{sga6.e00.sec.1.occ3}\subsection*{\S 1. Chow's theory}

Throughout this chapter, we assume a \uline{base field} $k$ fixed, which we shall assume algebraically closed, for simplicity. A large part of what follows, and perhaps all of it, is nevertheless valid without this restriction. The algebraic spaces and regular maps considered will be defined over $k$.

An algebraic space is called \uline{quasi-projective} if it is isomorphic to a locally closed subset of a projective space.

Let $X$ be a connected nonsingular quasi-projective variety of dimension $n$. We denote by $A(X)$, and call the \uline{Chow ring} of $X$, the graded ring of classes of cycles on $X$ modulo rational equivalence. ($A^i(X)$ consists of the classes of cycles of dimension $n-i$). The fact that $A(X)$
%%%%%%%%%%  scan idx 52  |  folio 39  %%%%%%%%%%
\SGAPrintPage{39}%
is indeed a graded ring follows from the following theorem of Chow:

\sgaStableTarget{sga6.e00.thm.2-1}\textbf{Theorem 2.1.}\SGAAuthorFootnote{*}{As a result of a misunderstanding on the author's part, the theorem states more than is currently known. In fact, only Chow's ``moving lemma'' is used, saying that two cycles can be moved within their classes so that their intersection has no excess; cf.\ Séminaire Chevalley 1958, Anneau de Chow et applications.} \uline{Every cycle on} $X$ \uline{is rationally equivalent to a nonsingular cycle. Given elements of} $A^i(X)$ \uline{and} $A^j(X)$, \uline{we can find nonsingular representatives of these elements, say} $Y^{n-i}$ \uline{and} $Z^{n-j}$, \uline{which meet transversely everywhere}.

(A cycle is called nonsingular if its components are nonsingular subvarieties. We say that two cycles meet transversely everywhere if every component of one meets every component of the other transversely everywhere, i.e. if the two intersecting varieties are nonsingular at their points of intersection and their tangent varieties there are in general position \dots)

Let $f$ be a morphism (i.e. a regular map) from $X$ to $Y$ ($X$ and $Y$ being quasi-projective and nonsingular); then $f$ defines (using the preceding theorem) a homomorphism of graded rings:
\sgaStableTarget{sga6.e00.eq.2-1.occ2}\begin{equation}\tag{2.1}
f^*:A(Y)\longrightarrow A(X)
\end{equation}
(\uline{inverse image of cycle classes}). Thus $A(X)$ becomes a contravariant functor in $X$.

Let $f$ be a \uline{proper} morphism from $X^n$ to $Y^m$ (i.e. in Weil's terminology, the graph of $f$ is complete over $Y$, cf.\ [3]). Then $f$ also defines an additive homomorphism
\sgaStableTarget{sga6.e00.eq.2-2}\begin{equation}\tag{2.2}
f_*:A(X^n)\longrightarrow A(Y^m)
\end{equation}
preserving the dimension of cycles, and therefore increasing the degree by $m-n$.

% --- Source pages 47--75 ---
%%%%%%%%%%  scan idx 53  |  folio 40  %%%%%%%%%%
\SGAPrintPage{40}%
Thus, with respect to proper morphisms, and disregarding the multiplicative structures and gradings, $A(X)$ is a covariant functor in $X$. Let us note the classical formula
\sgaStableTarget{sga6.e00.eq.2-3.occ2}\begin{equation}\tag{2.3}
 f_*(x.f^*(y))=f_*(x).y .
\end{equation}

We shall also use Chow's theory of Chern classes. It consists in attaching to every vector bundle $E$ on $X$ (still assumed quasi-projective and nonsingular) \uline{Chern classes} $C^i(E)\in A^i(X)$, ($i\geq 1$), in such a way as to satisfy the following conditions; put $C^0(E)=1$ and
\sgaStableTarget{sga6.e00.eq.2-4.occ2}\begin{equation}\tag{2.4}
 \Ctil(E)=\left[\rang E,\sum_{i\geq 0}C^i(E)\right]\in \Atil(X)
\end{equation}
(see chapter I, \S 3 for the definition of $\Atil(X)$). The conditions to be satisfied by a theory of Chern classes are then expressed as follows:
\sgaStableTarget{sga6.e00.eq.2-5.occ2}\begin{equation}\tag{2.5}
 \Ctil(E)=\Ctil(E')+\Ctil(E'')
\end{equation}
if the vector bundle $E$ is an extension of $E'$ by $E''$ (putting $C(E)=\sum_{i\geq0}C^i(E)$, this condition is expressed by $C(E)=C(E').C(E'')$). Formula (2.5) allows $\Ctil$ to be extended to an additive homomorphism from $\underline{K}(\xi(X))$ (cf.\ Chap.\ I, \S 2, example 2) to $\Atil(X)$:
\sgaStableTarget{sga6.e00.eq.2-6}\begin{equation}\tag{2.6}
 \Ctil:\underline{K}(\xi(X))\longrightarrow \Atil(X),
\end{equation}
and we want $\Ctil$ to be a $\lambda$\uline{-homomorphism of} $\lambda$\uline{-rings}; this last condition is also written
\sgaStableTarget{sga6.e00.eq.2-7}\begin{equation}\tag{2.7}
 \Ctil(E\otimes F)=\Ctil(E)\Ctil(F)\qquad
 \Ctil(\Lambda^iE)=\lambda^i\Ctil(E),
\end{equation}
for two vector bundles $E$ and $F$ (of course, $\Ctil$ is automatically compatible with the augmentations). Let $D$ be a divisor on $X$ and $L(D)$ the vector bundle it defines; we want
%%%%%%%%%%  scan idx 54  |  folio 41  %%%%%%%%%%
\SGAPrintPage{41}%
\sgaStableTarget{sga6.e00.eq.2-8}\begin{equation}\tag{2.8}
 C^1(L(D))=\ellc(D)\qquad C^i(L(D))=0\quad (i>1)
\end{equation}
(where,
for an arbitrary cycle $Z$, we denote by $\ellc(Z)$ its class in $A(X)$). Finally, we require the $C^i$ to be \uline{functorial}, i.e. that, when $f$ is a morphism from $X$ to $Y$ and $E$ is a vector bundle on $Y$, we have
\sgaStableTarget{sga6.e00.eq.2-9}\begin{equation}\tag{2.9}
 C^i(f^!(E))=f^*(C^i(E)),
\end{equation}
where $f^!(E)$ is the inverse image of $E$ on $X$. We can also write (2.9) in the following form: the notion of inverse image of a bundle defines a $\lambda$-homomorphism of augmented $\lambda$-rings
\sgaStableTarget{sga6.e00.eq.2-10}\begin{equation}\tag{2.10}
 f^!:\underline{K}(Y)\longrightarrow \underline{K}(X)\SGAAuthorFootnote{*}{The notation $f^!$ used here conflicts with that used in duality theory (cf.\ for example R. Hartshorne, Residues and Duality, Lecture Notes in Math. no.\ 20, Springer, 1966). This is why in this Seminar we replace it by the notation $f^*$.}
\end{equation}
(here we write $\underline{K}(X)$ instead of $\underline{K}(\xi(X))$), and (2.9) means that the diagram commutes:
\[
\begin{tikzcd}[column sep=large,row sep=large]
K(Y) \arrow[r,"f^!"] \arrow[d,"\Ctil_Y"'] & K(X) \arrow[d,"\Ctil_X"]\\
\Atil(Y) \arrow[r,"f^*"] & \Atil(X).
\end{tikzcd}
\]
(The space on which the symbol $\Ctil$ is considered has been put as a subscript.)

\sgaStableTarget{sga6.e00.thm.2-2}\textbf{Theorem 2.2.} \uline{There exists a theory of Chern classes satisfying conditions} (2.5), (2.7), (2.8) \uline{and} (2.9).

We shall also admit the following proposition:

\sgaStableTarget{sga6.e00.prop.2-3}\textbf{Proposition 2.3.} \uline{Let} $Y$ \uline{be a closed subset of} $X$ \uline{and} $U$ \uline{its complement; if} $x$ \uline{is an element of} $A(X)$ \uline{inducing} $0$ \uline{on} $U$, \uline{there exists a representative of} $x$ \uline{that is a cycle supported on} $Y$.

%%%%%%%%%%  scan idx 55  |  folio 42  %%%%%%%%%%
\SGAPrintPage{42}%
\sgaStableTarget{sga6.e00.cor.unnamed.occ2}\textbf{Corollary.} \uline{Let} $p=\dim X-\dim Y$. \uline{Then, if} $i<p$, \uline{the homomorphism} $A^i(X)\longrightarrow A^i(U)$ \uline{is bijective. Every element of the kernel of} $A^p(X)\longrightarrow A^p(U)$ \uline{is a linear combination of components of dimension} $n-p$ \uline{of} $Y$ (\uline{and is therefore proportional to} $\ellc(Y)$ \uline{if} $Y$ \uline{is irreducible})

\sgaStableTarget{sga6.e00.sec.2.occ3}\subsection*{\S 2. Definition of the Chern classes of coherent algebraic sheaves}

Let $\mathbb{C}$ be an abelian category, and let $\mathbb{C}_0$ be a subclass of $\mathbb{C}$. We denote by $K(\mathbb{C}_0)$ the quotient group of the free group generated by the classes (up to isomorphism) of elements of $\mathbb{C}_0$, by the subgroup generated by elements of the form $(E)-(E')-(E'')$, where $E$ is an \uline{extension} of $E'$ by $E''$ ($E,E',E''\in\mathbb{C}_0$). We must assume that the classes of elements of $\mathbb{C}_0$ for the isomorphism relation form a set. We denote by $\gamma_{\mathbb{C}_0}(E)$ the class of an $E\in\mathbb{C}_0$ in $K(\mathbb{C}_0)$, and omit the subscript $\mathbb{C}_0$ when there is no risk of confusion. The group $\underline{K}(\mathbb{C}_0)$ is the solution of a \uline{universal problem} concerning functions $E\longrightarrow f(E)$ defined on $\mathbb{C}_0$, with values in an arbitrary abelian group $A$, which are \uline{additive} in the sense that $f(E)=f(E')+f(E'')$ whenever $E$ is an extension of $E'$ by $E''$.

Let $F$ be an additive functor from $\mathbb{C}$ to $\mathbb{C}'$, and $\mathbb{C}'_0$ a subclass of $\mathbb{C}'$ satisfying the same condition as $\mathbb{C}_0$, such that, for every $A\in\mathbb{C}_0$, the $R^qF(A)$ are in $\mathbb{C}'_0$ and are zero for sufficiently large $q$ (we assume that injective resolutions exist in $\mathbb{C}$, so that the derived functors $R^qF$ of $F$ exist); then $F$ defines a homomorphism from $\underline{K}(\mathbb{C}_0)$ to $\underline{K}(\mathbb{C}'_0)$, still denoted $F$, by the formula
\sgaStableTarget{sga6.e00.eq.2-11}\begin{equation}\tag{2.11}
 F(\gamma_{\mathbb{C}_0}(E))=\sum_{i\geq 0}(-1)^i\gamma_{\mathbb{C}'_0}(R^iF(E)).
\end{equation}
More generally, when we have a cohomological functor $(F^i)$ from $\mathbb{C}$ to $\mathbb{C}'$
%%%%%%%%%%  scan idx 56  |  folio 43  %%%%%%%%%%
\SGAPrintPage{43}%
such that,
for every $E\in\mathbb{C}_0$, the $F^i(E)$ are in $\mathbb{C}'_0$ and are zero except for finitely many, we deduce a homomorphism from $\underline{K}(\mathbb{C}_0)$ to $\underline{K}(\mathbb{C}'_0)$. These remarks extend to multifunctors; moreover, the ``classical'' spectral sequences give transitivity properties that we shall not make explicit.

In particular, when $F$ is an exact functor, formula (2.11) simplifies:
\sgaStableTarget{sga6.e00.eq.2-11bis}\begin{equation}\tag{2.11 bis}
 F(\gamma_{\mathbb{C}_0}(E))=\gamma_{\mathbb{C}'_0}(F(E)).
\end{equation}
This applies in particular when $\mathbb{C}=\mathbb{C}'$, when $F$ is the identity functor, and when $\mathbb{C}_0\subset\mathbb{C}'_0$: we have a canonical homomorphism from $\underline{K}(\mathbb{C}_0)$ to $\underline{K}(\mathbb{C}'_0)$.

\sgaStableTarget{sga6.e00.thm.2-3}\textbf{Theorem 2.3.} \uline{Let} $\mathbb{C}$ \uline{be an abelian category and} $\mathbb{C}_0\subset\mathbb{C}_1$ \uline{two subclasses such that the classes} (up to isomorphism) \uline{of objects of} $\mathbb{C}_1$ \uline{form a set. We assume that the following conditions hold}
\begin{enumerate}[label=(\roman*)]
\item \uline{For every exact sequence} $0\to A'\to A\to A''\to0$ \uline{in} $\mathbb{C}$ \uline{such that} $A$ \uline{and} $A''$ \uline{are in} $\mathbb{C}_0$ (resp.\ $\mathbb{C}_1$), \uline{we have} $A'\in\mathbb{C}_0$ (resp.\ $A'\in\mathbb{C}_1$).
\item \uline{Let} $u:A\to B$ \uline{and} $v:C\to B$ \uline{with} $A,B,C\in\mathbb{C}_1$ \uline{and} $u$ \uline{surjective; then there exist} $L\in\mathbb{C}_0$ \uline{and homomorphisms} $u':L\to C$ \uline{and} $v':L\to A$ \uline{such that} $u'$ \uline{is surjective and} $vu'=uv'$.
\item \uline{For every} $A\in\mathbb{C}_1$, \uline{there exists an integer} $d=d(A)$ \uline{such that, for every left resolution} $\underline{L}$ \uline{of} $A$ \uline{by objects of} $\mathbb{C}_0$, \uline{we have} $Z_d(\underline{L})\in\mathbb{C}_0$ ($Z_d$ denotes the cycles of dimension $d$).
\end{enumerate}
\uline{Under these conditions, the canonical homomorphism from} $K(\mathbb{C}_0)$ \uline{to} $K(\mathbb{C}_1)$ \uline{is an isomorphism}.

\textbf{Principle of the proof:} let $A\in\mathbb{C}_1$; by (ii) and (iii) there exists a finite resolution $\underline{L}$ of $A$ by objects of $\mathbb{C}_0$; put
%%%%%%%%%%  scan idx 57  |  folio 44  %%%%%%%%%%
\SGAPrintPage{44}%
then $f(A)=\sum_{i\geq 0}(-1)^i\gamma_{\mathbb{C}_0}(L^i)$.
We show that the element of $\underline{K}(\mathbb{C}_0)$ thus constructed does not depend on the chosen resolution $\underline{L}$: given a second resolution $\underline{L}'$, we ``cover'' $\underline{L}$ and $\underline{L}'$ by a third resolution $\underline{L}''$, with \uline{surjective} homomorphisms from $\underline{L}''$ to $\underline{L}$ and $\underline{L}'$ (which is possible by (i) and (ii)); we are then reduced to proving $\underline{E}_{\mathbb{C}_0}(\underline{L})=\underline{E}_{\mathbb{C}_0}(\underline{L}'')$. Now the difference of the two is $\underline{E}_{\mathbb{C}_0}(\underline{N})$, where $\underline{N}$ is the kernel of the surjective homomorphism from $\underline{L}''$ onto $\underline{L}$, a kernel that is in $\mathbb{C}_0$ by (i). Now $\underline{N}$ is an acyclic complex in all dimensions, whence $\underline{E}_{\mathbb{C}_0}(\underline{N})=0$, whence the conclusion (N.B. $\underline{E}_{\mathbb{C}_0}(\underline{L})$, for a finite complex with values in $\mathbb{C}_0$, denotes the Euler--Poincaré characteristic $\sum_{i\geq0}(-1)^i\gamma_{\mathbb{C}_0}(L^i)$ of $\underline{L}$).

We prove analogously that the function $f:\mathbb{C}_1\longrightarrow\underline{K}(\mathbb{C}_0)$ thus constructed is \uline{additive}, and therefore defines a homomorphism from $\underline{K}(\mathbb{C}_1)$ to $\underline{K}(\mathbb{C}_0)$, and it is immediate that this homomorphism is inverse to the canonical homomorphism from $\underline{K}(\mathbb{C}_0)$ to $\underline{K}(\mathbb{C}_1)$. q.e.d.

\sgaStableTarget{sga6.e00.rem.unnamed.occ3}\textbf{Remark:} To verify (ii), it suffices to verify that $\mathbb{C}_1$ is contained in a subclass $\mathbb{C}_2$ of $\mathbb{C}$ satisfying the two conditions
\begin{enumerate}[label=(ii \alph*)]
\item Every $A\in\mathbb{C}_2$ is isomorphic to a quotient of an $L\in\mathbb{C}_0$.
\item If $u:A\to B$, with $A,B\in\mathbb{C}_2$, then $\Ker u\in\mathbb{C}_2$.
\end{enumerate}

\sgaStableTarget{sga6.e00.cor.unnamed.occ3}\textbf{Corollary:}
\uline{Let} $X$ \uline{be a quasi-projective algebraic space. Let} $\cF(X)$ \uline{be the category of coherent algebraic sheaves on} $X$, $\xi(X)$ \uline{the category of locally free coherent algebraic sheaves} (equivalent to the category of vector bundles on $X$), \uline{and} $\cF_0(X)$ \uline{the category of coherent algebraic sheaves} $F$ \uline{such that, for every} $x\in X$, $F_x$ \uline{is an} $\underline{O}_x$\uline{-module}
%%%%%%%%%%  scan idx 58  |  folio 45  %%%%%%%%%%
\SGAPrintPage{45}%
\uline{of finite cohomological dimension (we therefore have} $\xi(X)\subset\cF_0(X)\subset\cF(X)$). \uline{Then the canonical homomorphism}
\sgaStableTarget{sga6.e00.eq.2-12}\begin{equation}\tag{2.12}
 K(\xi(X))\longrightarrow K(\cF_0(X))
\end{equation}
\uline{is an isomorphism}.

We are going to verify the conditions of theorem 2.3. Conditions (i) and (iii) follow from the definition of cohomological dimension and from the fact that a projective module of finite type over the noetherian local ring $\underline{O}_x$ is free. To verify (ii), we shall show that $\mathbb{C}=\cF(X)$ satisfies conditions (ii a) and (ii b) of the remark. The second is trivial; it therefore remains to prove the following result:

\sgaStableTarget{sga6.e00.prop.2-4}\textbf{Proposition 2.4. (Serre).} \uline{Let} $X$ \uline{be a quasi-projective algebraic space. Then every coherent algebraic sheaf on} $X$ \uline{is isomorphic to the quotient of a locally free algebraic sheaf}.

This proposition holds when $X$ is a closed subset of projective space, by [2], since, for sufficiently large $n$, $F(n)$ is generated by its sections, which means that $F$ is a quotient of $\underline{O}(-n)^k$. In the general case, $X$ is an open subset of a closed subset of a projective space, and it suffices to use the following result (due to Cartier and Serre in particular cases):

\sgaStableTarget{sga6.e00.prop.2-5}\textbf{Proposition 2.5.} \uline{Let} $F$ \uline{be a coherent algebraic sheaf on an open subset} $U$ \uline{of an algebraic space} $X$. \uline{Then} $F$ \uline{is isomorphic to the restriction of a coherent algebraic sheaf on} $X$.

By applying Zorn's lemma to the open subsets containing $U$ to which $F$ can be extended, we are reduced to the case where $X$ is affine. Let $\overline F$ be the
alge%
%%%%%%%%%%  scan idx 59  |  folio 46  %%%%%%%%%%
\SGAPrintPage{46}%
braic sheaf on $X$ whose sections on an open subset $V$ are the sections of $F$ on $U\cap V$. We readily find that $\overline F$ is \uline{quasi-coherent} (i.e. defined on every affine open subset $V$ -- hence here on all of $X$ -- by a module, of finite type or otherwise, over the coordinate ring of the open subset in question). This is immediate when $U$ itself is affine, for $\overline F$ is then defined on the open subset $V$ by $\Gamma(U\cap V,F)$ considered as a module over $\Gamma(V,\underline{O}_X)$. In the general case, $U$ is a finite union of affine open subsets $U_i$, so $\overline F$ is the kernel of an evident homomorphism of quasi-coherent sheaves:
\[
 \sum_i \overline{F|U_i}\longrightarrow \sum_{i,j}\overline{F|U_i\cap U_j}
\]
and is therefore itself quasi-coherent. As $X$ is assumed affine, $\overline F$ is therefore generated by its sections. As $\overline F|U=F$ is coherent, there exist finitely many sections of $\overline F$ generating $\overline F|U$. These sections generate a coherent subsheaf of $\overline F$, whose restriction to $U$ is isomorphic to $F$.
\begin{flushright}q.e.d.\end{flushright}

Now assume $X$ quasi-projective and \uline{nonsingular}. Then, by the syzygy theorem, we have, with the notation of the corollary to theorem 2.3, $\cF_0(X)=\cF(X)$. In this case we therefore have:
\sgaStableTarget{sga6.e00.eq.2-13}\begin{equation}\tag{2.13}
 \underline{K}(\xi(X))=\underline{K}(\cF(X)).
\end{equation}
We can express this isomorphism as follows:

\sgaStableTarget{sga6.e00.cor.2}\textbf{Corollary 2 to theorem 2.3.} \uline{Let} $X$ \uline{be a nonsingular quasi-projective variety. For every abelian group} $A$, \uline{there is a one-to-one correspondence between additive functions} $f$ \uline{of vector bundles on} $X$, \uline{with values in} $A$, \uline{and additive functions} $g$ \uline{of coherent algebraic sheaves on} $X$, \uline{with values in} $A$. \uline{To every} $f$ \uline{corresponds the unique} $g$ \uline{such that we have}:
%%%%%%%%%%  scan idx 60  |  folio 47  %%%%%%%%%%
\SGAPrintPage{47}%
\sgaStableTarget{sga6.e00.eq.2-14}\begin{equation}\tag{2.14}
 f(E)=g(\underline{O}_X(E))
\end{equation}
\uline{for every vector bundle} $E$ \uline{on} $X$, \uline{denoting by} $\underline{O}_X(E)$ \uline{the sheaf of germs of regular sections of} $E$ \uline{on} $X$.

To calculate $g(F)$ for an arbitrary coherent algebraic sheaf, we consider a finite resolution of $F$ by sheaves of the form $\underline{O}_X(E_i)$, and we then have:
\sgaStableTarget{sga6.e00.eq.2-14bis}\begin{equation}\tag{2.14 bis}
 g(F)=\sum_{i\geq0}(-1)^if(E_i).
\end{equation}

In particular, when $X$ is connected, we considered in \S 1 an additive function $\Ctil(E)$ of the vector bundle $E$, with values in the group $\Atil=\Atil(X)$. We deduce from it an additive function of coherent sheaves, say $\Ctil(F)$, and, by passing to components, functions $F\longrightarrow C^i(F)\in A^i(X)$. The $C^i(F)$ are called the \uline{Chern classes of the coherent sheaf} $F$; they are theoretically calculated from the Chern classes of vector bundles by (2.14 bis) (where, however, we must pass to multiplicative notation in the second member). As for the component of $\Ctil(F)$ along $\Z$, it is the \uline{rank} of the sheaf $F$, which corresponds to the rank of a module over an integral domain: when $X$ is affine (a case to which we can always reduce), $F$ is defined by a module of finite type over the coordinate ring of $X$, and we take the rank of this module\dots.

\sgaStableTarget{sga6.e00.sec.3.occ3}\subsection*{\S 3. Functorial generalities on $\underline{K}(X)$}

When $X$ is an algebraic space, we put, for brevity, $\underline{K}(X)=\underline{K}(\cF(X))$, and, when $X$ is nonsingular and quasi-projective, we identify this group with $\underline{K}(\xi(X))$ by corollary 2 to theorem 2.3. As $\underline{K}(\xi(X))$ is
%%%%%%%%%%  scan idx 61  |  folio 48  %%%%%%%%%%
\SGAPrintPage{48}%
not
only an abelian group, but a $\lambda$-ring, it follows by transport of structure that $\underline{K}(X)$ is also a $\lambda$-ring, and is therefore equipped with a multiplicative structure and maps $\lambda^i$. It is important to define these structures directly, without passing to vector bundles; this will be done below.

If $E$ is a vector bundle on $X$, we denote by $\gamma(E)$ its class in $\underline{K}(X)$; if $F$ is a coherent algebraic sheaf on $X$, $\gamma(F)$ likewise denotes its class in $\underline{K}(X)$. If $Y$ is an irreducible subvariety of $X$, we put $\gamma(Y)=\gamma(\underline{O}_Y)$, where $\underline{O}_Y$ denotes the sheaf of local rings of $Y$, considered as a coherent algebraic sheaf on $X$. We define the symbol $\gamma(Y)$ by linearity when $Y$ is an arbitrary cycle on $X$. When there is a risk of confusion, we shall write $\gamma_X$ for $\gamma$.

Let $X$ be an arbitrary algebraic space. We introduce a \uline{filtration} on $\underline{K}(X)$, denoting by $\underline{K}(X)^i$ the subgroup of $\underline{K}(X)$ generated by the $\gamma(F)$ with dim.supp.\ $F\leq n-i$ (where $n$ is the dimension of $X$). From Serre's ``dévissage lemma'', in the form given in [3], follows the result:

\sgaStableTarget{sga6.e00.prop.2-6}\textbf{Proposition 2.6.} \uline{The group} $\underline{K}(X)^i$ \uline{is generated by the} $\gamma(Y)$, \uline{where} $Y$ \uline{runs through the set of subvarieties of} $X$ \uline{of dimension} $\leq n-i$. \uline{If} $X$ \uline{is irreducible, we have} $G^0(\underline{K}(X))=\Z$.

The latter isomorphism is given by the homomorphism from $\underline{K}(X)$ to $\Z$ defined by means of the notion of \uline{rank} of a coherent algebraic sheaf.

Now assume $X$ \uline{nonsingular}. Following the developments
%%%%%%%%%%  scan idx 62  |  folio 49  %%%%%%%%%%
\SGAPrintPage{49}%
at the beginning of \S 2, we
define a multiplication on $\underline{K}(X)$ by putting
\sgaStableTarget{sga6.e00.eq.2-15}\begin{equation}\tag{2.15}
 \gamma(F)\gamma(G)=\sum_{i\geq0}(-1)^i\gamma(\underline{\Tor}_i(F,G))
\end{equation}
(the $\underline{\Tor}_i$ being taken, of course, with respect to the local rings of $X$). The exact sequence of the $\underline{\Tor}_i$ shows that formula (2.15) does define a biadditive map from $\underline{K}(X)\times\underline{K}(X)$ to $\underline{K}(X)$; associativity follows from the spectral sequence of multiple $\Tor$. We also verify in this way that the product of several factors can be calculated by the formula:
\sgaStableTarget{sga6.e00.eq.2-16}\begin{equation}\tag{2.16}
 \gamma(F_1)....\gamma(F_p)=\sum_{i\geq0}(-1)^i\gamma(\underline{\Tor}_i(F_1,\ldots,F_p))
\end{equation}
where, in the second member, the $\underline{\Tor}_i$ are ``simultaneous $\underline{\Tor}$'', calculated by simultaneous projective resolution of the arguments $F_i$.

If, in (2.15), we assume $G$ \uline{locally free}, we find:
\sgaStableTarget{sga6.e00.eq.2-15bis}\begin{equation}\tag{2.15 bis}
 \gamma(F)\gamma(G)=\gamma(F\otimes G),
\end{equation}
which shows in particular that $\gamma(\underline{O}_X)$ is the unit element of $\underline{K}(X)$, and that the canonical homomorphism from $\underline{K}(\xi(X))$ to $\underline{K}(X)$ is a ring homomorphism.

Formula (2.15) is evidently suggested by Serre's ``$\Tor$ formula'' in intersection theory. This formula proves the second part of the following proposition, the first following immediately from formula (2.15) and a calculation of local $\underline{\Tor}_i$:

\sgaStableTarget{sga6.e00.prop.2-7}\textbf{Proposition 2.7.} \uline{Let} $X$ \uline{be a nonsingular variety, and} $Y$ \uline{and} $Z$ \uline{two cycles in} $X$. \uline{If} $Y$ \uline{and} $Z$ \uline{meet transversely everywhere, we have}:
\sgaStableTarget{sga6.e00.eq.2-17}\begin{equation}\tag{2.17}
 \gamma(Y)\gamma(Z)=\gamma(Y.Z).
\end{equation}
\uline{If} $Y$ \uline{and} $Z$ \uline{are homogeneous and have respective dimensions} $n-i$ \uline{and} $n-j$, \uline{and if they meet without an excess component, we then have}:
%%%%%%%%%%  scan idx 63  |  folio 50  %%%%%%%%%%
\SGAPrintPage{50}%
\sgaStableTarget{sga6.e00.eq.2-18}\begin{equation}\tag{2.18}
 \gamma(Y)\gamma(Z)=\gamma(Y.Z)\qquad\mathrm{mod.}\ \underline K(X)^{i+j+1}.
\end{equation}

We
must take care that equality does not hold in general in (2.18). We shall see in \S 4 that, when $X$ is \uline{nonsingular and quasi-projective, the ring structure of} $\underline K(X)$ \uline{is compatible with its filtration and that the associated graded ring} $G(\underline K(X))$ \uline{is identified with a quotient of} $A(X)$ \uline{by a torsion subgroup (perhaps always reduced to} $0$\SGAAuthorFootnote{*}{No; for a counterexample cf.\ this Seminar XIV 4.7.}).

I know how to define the $\lambda^i(F)$ directly in terms of sheaves only when the field $k$ has characteristic $0$ (this is the reason that has prevented us from proving the Riemann--Roch theorem when $k$ is not of characteristic $0$). Let $F$ be a coherent algebraic sheaf on $X$ on which the group $\mathfrak S_p$ of permutations of $p$ letters operates; we then denote by $(F)^{\mathrm{alt}}$ the set of the $x\in F$ such that $s.x=\varepsilon(s)x$ for every $s\in\mathfrak S_p$ ($\varepsilon(s)$ is the \uline{signature} of $s$). This is evidently a coherent algebraic subsheaf of $F$. This said, we have:

\sgaStableTarget{sga6.e00.thm.2-8}\textbf{Theorem 2.8.} \uline{Let} $X$ \uline{be a nonsingular quasi-projective variety over the field} $k$ of characteristic $0$. \uline{Let} $F$ \uline{be a coherent algebraic sheaf on} $X$; \uline{we then have}:
\sgaStableTarget{sga6.e00.eq.2-19}\begin{equation}\tag{2.19}
 \lambda^p(\gamma(F))=\sum_{i\geq0}(-1)^i\gamma\left((\underline{\Tor}_i(F,F,\ldots,F))^{\mathrm{alt}}\right)
\end{equation}

N.B. It is clear that the group $\mathfrak S_p$ acts on the sheaf $\underline{\Tor}_i(F,\ldots,F)$.

\uline{Principle of the proof}: we consider a finite resolution $\underline L$ of $F$ by locally free sheaves $L_i$, and use the following general lemma:

%%%%%%%%%%  scan idx 64  |  folio 51  %%%%%%%%%%
\SGAPrintPage{51}%
\sgaStableTarget{sga6.e00.lem.2-9}\textbf{Lemma 2.9.} \uline{Let} $\underline L=(L_i)$ \uline{be a graded locally free coherent algebraic sheaf, with finitely many nonzero components. In} $\underline K(X)$ \uline{we have the formula}:
\sgaStableTarget{sga6.e00.eq.2-20}\begin{equation}\tag{2.20}
 \underline E\left((\overset{p}{\otimes}\underline L)^{\mathrm{alt}}\right)=\wedge^p(\underline E(\underline L)).
\end{equation}
In this formula, we consider $\overset{p}{\otimes}\underline L$ as a sheaf on which the group $\mathfrak S_p$ operates, taking the gradings into account (which therefore introduce signs in the well-known way, cf.\ Cartan--Eilenberg); moreover, if $\underline M$ is a graded coherent algebraic sheaf on $X$, we put $\underline E(\underline M)=\sum_{i\geq0}(-1)^i\gamma(M_i)$.

To prove formula (2.20), we are reduced to proving the analogous formula in $\underline K(G)$, where $G$ is an algebraic group (a product of groups $\operatorname{Gl}(n,k)$, if desired) and $\underline L$ a $G$-module of finite dimension over $k$ (cf.\ Chap.\ I, \S 2, example 1, for the definition of $\underline K(G)$). In this case, this formula has already been used in theorem 1.4 to show that the second member of (1.33) satisfies the same functional equation as the first member; its verification is elementary.
\begin{flushright}q.e.d.\end{flushright}

Let $f$ be a morphism from $X$ to $Y$, where $Y$ is an algebraic space. The notion of inverse image of a vector bundle defines a homomorphism $f^!:\underline K(\xi(Y))\longrightarrow\underline K(\xi(X))$ of $\lambda$-rings. Identifying vector bundles with locally free algebraic sheaves, the notion of inverse image of a vector bundle corresponds to the notion of inverse image of an algebraic sheaf (already used in [3, no.\ 8]). When $Y$ is nonsingular, we define more generally a group homomorphism:
\sgaStableTarget{sga6.e00.eq.2-21}\begin{equation}\tag{2.21}
 f^!:\underline K(Y)\longrightarrow\underline K(X)
\end{equation}
%%%%%%%%%%  scan idx 65  |  folio 52  %%%%%%%%%%
\SGAPrintPage{52}%
by a
formula involving an alternating sum of $\Tor$, inspired by the developments at the beginning of \S 2, which the reader will make explicit. The following diagram then commutes:
\[
\begin{tikzcd}[column sep=large,row sep=large]
\underline K(\xi(Y)) \arrow[r,"f^!"] \arrow[d] & \underline K(\xi(X)) \arrow[d]\\
\underline K(Y) \arrow[r,"f^!"] & \underline K(X){,}
\end{tikzcd}
\]
which shows in particular, when $X$ and $Y$ are both nonsingular and quasi-projective, that the two definitions of $f^!$ coincide under the identification of $\underline K(X)$ with $\underline K(\xi(X))$ and of $\underline K(Y)$ with $\underline K(\xi(Y))$.

Another case in which a homomorphism (2.21) can be defined, without $Y$ necessarily being nonsingular, is that in which the functor $F\to f^!(F)$ from $\cF(Y)$ to $\cF(X)$ is exact, i.e.\ when the local rings of $X$ are \uline{flat} modules over the local rings of $Y$. It then suffices to put $f^!(\gamma_Y(F))=\gamma_X(f^!(F))$.
This case occurs in particular when $X$ is a locally trivial fibre space over $Y$ (to see this, we are reduced to the case of a product $X=T\times Y$). Note that we also verify that, in this case, we have $f^!(\underline{O}_Z)=\underline{O}_{f^{-1}(Z)}$ if $Z$ is an irreducible subvariety of $Y$, whence:
\sgaStableTarget{sga6.e00.eq.2-22}\begin{equation}\tag{2.22}
 f^!(\gamma_Y(Z))=\gamma_X(f^{-1}(Z))
\end{equation}
\uline{if} $X$ \uline{is ``flat'' over} $Y$. Without this last condition, formula (2.22) is in general incorrect, but we have the following analogue of proposition 2.7:

\sgaStableTarget{sga6.e00.prop.2-8}\textbf{Proposition 2.8.}
\uline{Let} $f:X\longrightarrow Y$ \uline{be a morphism of nonsingular varieties, and} $Z$ \uline{a cycle on} $Y$. \uline{If} $Z$ \uline{is everywhere transverse to} $f$, \uline{formula}
%%%%%%%%%%  scan idx 66  |  folio 53  %%%%%%%%%%
\SGAPrintPage{53}%
(2.22) \uline{is exact; if} $Z$ \uline{is homogeneous of codimension} $i$ \uline{and} $f^{-1}(Z)$ \uline{has no excess component, we have}:
\sgaStableTarget{sga6.e00.eq.2-23}\begin{equation}\tag{2.23}
 f^!(\gamma_Y(Z))\equiv\gamma_X(f^{-1}(Z)).\qquad\mathrm{mod.}\ \underline K(X)^{i+1}
\end{equation}
Proof analogous to that of proposition 2.7.

Finally, let us mention that we shall see in \S 4 that, \uline{when} $X$ \uline{and} $Y$ \uline{are quasi-projective and nonsingular}, $f^!$ \uline{is compatible with the filtrations of} $\underline K(X)$ \uline{and of} $\underline K(Y)$.

Again let $X$ and $Y$ be two arbitrary algebraic spaces, and $f$ a \uline{proper} morphism from $X$ to $Y$. For every coherent algebraic sheaf $F$ on $X$, the direct image sheaf $f_!(F)$ is coherent algebraic, and the same holds more generally for the sheaves $R^qf_!(F)$ (cf.\ [3]). Recall that, by definition, $R^qf_!(F)$ is the sheaf on $Y$ associated with the presheaf assigning to an open subset $U$ the group $H^q(f^{-1}(U),F)$; in particular, we have:
\sgaStableTarget{sga6.e00.eq.2-24}\begin{equation}\tag{2.24}
 R^0f_!(F)=f_!(F),\qquad R^qf_!(F)=0\quad(q>\dim X),
\end{equation}
moreover:
\sgaStableTarget{sga6.e00.eq.2-25}\begin{equation}\tag{2.25}
 \Gamma(U,R^qf_!(F))=H^q(f^{-1}(U),F)
\end{equation}
\uline{if the open subset} $U$ \uline{is affine}. By \S 2, we define a homomorphism from $\underline K(X)$ to $\underline K(Y)$, still denoted $f_!$, by the formula
\sgaStableTarget{sga6.e00.eq.2-26}\begin{equation}\tag{2.26}
 f_!(\gamma_X(F))=\sum_{q\geq0}(-1)^q\gamma_Y(R^qf_!(F)).
\end{equation}
Thus $\underline K(X)$ becomes a \uline{covariant functor in} $X$, with respect to \uline{proper} morphisms. The relation $(fg)_!=f_!g_!$ (just like the relation $(fg)^!=g^!f^!$ that we forgot to mention in its place) is a parti%
%%%%%%%%%%  scan idx 67  |  folio 54  %%%%%%%%%%
\SGAPrintPage{54}%
cular case of the transitivity relations to which we alluded in \S 2, which follow from the spectral sequences of composite functors.

Let us note the following formula, valid for $Y$ nonsingular (and particularly easy to verify when $Y$ is also quasi-projective, by taking for $y$ the class of a vector bundle on $Y$):
\sgaStableTarget{sga6.e00.eq.2-27}\begin{equation}\tag{2.27}
 f_!(x.f^!(y))=f_!(x).y\qquad(x\in\underline K(X),\ y\in\underline K(Y)).
\end{equation}
In the particular case where the inverse images under $f$ of the points of $Y$ are finite and contained in an affine open subset\SGAAuthorFootnote{*}{This last condition is evidently a consequence of the first (something the author still did not know in 1957!).}, we find that the functor $F\longrightarrow f_!(F)$ is exact and that $R^qf_!(F)=0$ for $q>0$; in this case, we therefore have $f_!(\gamma_X(F))=\gamma_Y(f_!(F))$. In particular, if $X$ is an algebraic subspace of $Y$ and $i$ is the canonical injection from $X$ to $Y$, we have:
\sgaStableTarget{sga6.e00.eq.2-28}\begin{equation}\tag{2.28}
 i_!(\gamma_X(F))=\gamma_Y(F),
\end{equation}
where, in the second member, $F$ is considered as a sheaf on $Y$ that is zero outside $X$.

\sgaStableTarget{sga6.e00.sec.4.occ3}\subsection*{\S 4. Some technical results}

Let $X$ be an algebraic space; we saw in \S 3 that the projection $p$ from $X\times k$ onto $X$ defines a homomorphism
\sgaStableTarget{sga6.e00.eq.2-29}\begin{equation}\tag{2.29}
 p^!:\underline K(X)\longrightarrow\underline K(X\times k).
\end{equation}

\sgaStableTarget{sga6.e00.prop.2-9}\textbf{Proposition 2.9.} \uline{The homomorphism} (2.29) \uline{is surjective}.

We argue by induction on $n=\dim X$; the assertion is trivial if $n<0$; assume $n\geq0$; by proposition 2.6, it suffices to
%%%%%%%%%%  scan idx 68  |  folio 55  %%%%%%%%%%
\SGAPrintPage{55}%
verify that, for every irreducible closed subset $Y$ of $X\times k$, $\gamma(Y)$ is in the image of $p^!$. It suffices to continue the reasoning assuming $X$ irreducible; if $p(Y)$ is not dense in $X$, our assertion follows from the induction hypothesis; otherwise, either $Y=X\times k$ and the assertion is trivial, or $Y$ is a hypersurface in $X\times k$. Now let $U$ be an affine open subset of $X$ consisting of nonsingular points; if $D$ is the divisor on $U\times k$ induced by $Y$, we know (commutative algebra) that it is linearly equivalent to a divisor of the form $D'\times k$, where $D'$ is a divisor on $U$. Then $\gamma_{U\times k}(D)-\gamma_{U\times k}(D'\times k)$ belongs to $\underline K(U\times k)^2$ (corollary of proposition 2.10), and, by what we have already proved, this element is of the form $p^!(x')$ with $x'\in\underline K(U)$. Thus $\gamma_{U\times k}(D)=p^!(y')$ with $y'=x'+\gamma_U(D')$; but $y'$ is the restriction of an element $y$ of $\underline K(X)$ (prop.\ 2.10), and consequently the restriction of $\gamma_{X\times k}(Y)-p^!(y)$ to $U\times k$ is zero; by proposition 2.10, $\gamma_{X\times k}(Y)-p^!(y)$ comes from an element of $\underline K(X\times k-U\times k)=\underline K(X'\times k)$ with $X'=X-U$. The proof is completed by applying the induction hypothesis again.
\begin{flushright}q.e.d.\end{flushright}

\sgaStableTarget{sga6.e00.cor.unnamed.occ4}\textbf{Corollary:} \uline{Let} $X$ \uline{be an algebraic space and} $p$ \uline{the projection from} $X\times k^n$ \uline{onto} $X$; \uline{then the homomorphism} $p^!:\underline K(X)\longrightarrow\underline K(X\times k^n)$ \uline{is surjective; it is even bijective when} $X$ \uline{is nonsingular}.

Induction on $n$ for the first assertion. For the second, we put $f(x)=(x,0)$, whence $pf=\underline{1}$ and $f^!p^!=1$, which proves that $p^!$ is injective.

\sgaStableTarget{sga6.e00.prop.2-10}\textbf{Proposition 2.10.} \uline{Let} $X$ \uline{be an algebraic space}, $U$ \uline{an open subset of} $X$, \uline{and} $Y$ \uline{its complement in} $X$. \uline{Then the sequence of natural homomorphisms}:
%%%%%%%%%%  scan idx 69  |  folio 56  %%%%%%%%%%
\SGAPrintPage{56}%
\sgaStableTarget{sga6.e00.eq.2-30}\begin{equation}\tag{2.30}
 \underline K(Y)\longrightarrow\underline K(X)\longrightarrow\underline K(U)\longrightarrow 0
\end{equation}
\uline{is exact}.

N.B. Compare with proposition 2.3.

The fact that $\underline K(X)\longrightarrow\underline K(U)$ is surjective follows immediately from proposition 2.5 (or also from proposition 2.6). It remains to prove that, if $x\in\underline K(X)$ has zero restriction to $\underline K(U)$, it is in the image of $\underline K(Y)$. For this, we return to the explicit definition of $\underline K(X)$ and $\underline K(U)$: we have two coherent algebraic sheaves $F$ and $G$ on $X$ with $x=\gamma_X(F)-\gamma_X(G)$ and $\gamma_U(F)-\gamma_U(G)=0$; this last relation means that we can find coherent sheaves $P'$ and $Q'$ on $U$, equipped with filtrations whose factors are pairwise isomorphic on $U$, and such that we have an isomorphism on $U$:
\[
 P=F+Q'\longrightarrow Q=G+P'.
\]
By proposition 2.5, the sheaves $P'$ and $Q'$ are restrictions to $U$ of coherent sheaves on $X$ (denoted by the same letters); likewise, their filtrations come from filtrations on the corresponding sheaves on $X$, by lemma 2.11 below. We then have:
\[
 \gamma_X(F)-\gamma_X(G)=(\gamma_X(P)-\gamma_X(Q))+(\gamma_X(P')-\gamma_X(Q')).
\]
Introducing the factors $P'_i$ and $Q'_j$ of the filtrations of $P'$ and $Q'$ on $X$, the last term is a sum of terms $\gamma_X(P'_i)-\gamma_X(Q'_i)$, where $P'_i|U\simeq Q'_i|U$. We are thus reduced to proving the following particular case: if $R$ and $S$ are coherent sheaves on $X$ such that $R|U\simeq S|U$, then $\gamma_X(R)-\gamma_X(S)$ is in the image of $\underline K(Y)$. Now let $T$ be the graph of the given isomorphism from $R|U$ onto $S|U$; it is a coherent subsheaf of $(R\times S)|U$, and therefore, by lemma 2.11, it is the restriction to $U$ of a coherent subsheaf,
%%%%%%%%%%  scan idx 70  |  folio 57  %%%%%%%%%%
\SGAPrintPage{57}%
again
denoted $T$, of $R\times S$. We then have $\gamma_X(R)-\gamma_X(S)=(\gamma_X(R)-\gamma_X(T))-(\gamma_X(S)-\gamma_X(T))$, and it suffices to prove, for example, that the first term of the second member is in the image of $\underline K(Y)$. Now the projection $u:T\longrightarrow R$ is an isomorphism over $U$, so $\gamma_X(R)-\gamma_X(T)=\gamma_X(\operatorname{Coker}u)-\gamma_X(\operatorname{Ker}u)$, and $\operatorname{Ker}u$ and $\operatorname{Coker}u$ have their support in $Y$.

It therefore remains to prove the following lemma:

\sgaStableTarget{sga6.e00.lem.2-11}\textbf{Lemma 2.11.} \uline{Let} $X$ \uline{be an algebraic space}, $F$ \uline{a coherent algebraic sheaf on} $X$, $U$ \uline{an open subset of} $X$, \uline{and} $G$ \uline{a coherent subsheaf of} $F|U$. \uline{Then} $G$ \uline{is the restriction to} $U$ \uline{of a coherent subsheaf of} $F$.

Let $\overline G$ be the subsheaf of $F$ whose sections on the open subset $V$ of $X$ are the sections of $F$ on $V$ whose restriction to $U\cap V$ is a section of $G$. Everything reduces to proving that this sheaf $\overline G$ is coherent. For this, we can assume $X$ and $U$ affine (for $U$ is a union of affine open subsets $U_i$, and $\overline G$ is then the intersection of the analogous sheaves $\overline{G|U_i}$). Then the coordinate ring of $U$ is a ring of fractions $A_S$ of the coordinate ring $A$ of $X$, and we verify immediately that, if $F$ is defined by the $A$-module $M$, then $F|U$ is defined by the $A_S$-module $M_S=A_S\otimes_A M$. The subsheaf $G$ of $F|U$ is then defined by a submodule $N'$ of $M_S$, and we see that $\overline G$ is precisely defined by the submodule $N$ of $M$ that is the inverse image of $N'$ under the canonical homomorphism from $M$ to $M_S$.
\begin{flushright}q.e.d.\end{flushright}

\sgaStableTarget{sga6.e00.cor.1}\textbf{Corollary 1 to proposition 2.10:} \uline{Let} $x\in\underline K(X)$. \uline{For} $x$ \uline{to have filtration} $\geq i$, \uline{it is necessary and sufficient that there exist a closed subset} $Y$ \uline{of} $X$, \uline{of dimension} $\leq n-i$, \uline{such that the restriction of} $x$ \uline{to} $X-Y$ \uline{is zero}.

%%%%%%%%%%  scan idx 71  |  folio 58  %%%%%%%%%%
\SGAPrintPage{58}%
\sgaStableTarget{sga6.e00.cor.2.occ2}\textbf{Corollary 2.} \uline{Assume} $X$ \uline{connected of dimension} $n$, \uline{without singularities in dimension} $n-1$. \uline{If} $D$ \uline{and} $D'$ \uline{are two linearly equivalent divisors on} $X$, \uline{we have} $\gamma(D)\equiv\gamma(D')\quad\mathrm{mod.}\ \underline K(X)^2$.

By corollary 1, we can assume $X$ nonsingular; the corollary then follows from the general formula:
\sgaStableTarget{sga6.e00.eq.2-31}\begin{equation}\tag{2.31}
 \gamma(D)=1-\gamma(L(-D))\qquad\mathrm{mod}\ \underline K(X)^2,
\end{equation}
where $L(-D)$ denotes the vector bundle of rank $1$ defined by the divisor $-D$. To verify (2.31), we write it in the form:
\sgaStableTarget{sga6.e00.eq.2-31bis}\begin{equation}\tag{2.31 bis}
 \gamma(L(-D))=1-\gamma(D)\qquad\mathrm{mod}\ \underline K(X)^2,
\end{equation}
and note that both members are homomorphisms from the group of divisors to the multiplicative group of invertible elements of the ring $\underline K(X)/\underline K(X)^2$ (note that $\underline K(X)^1/\underline K(X)^2$ has square zero, which is very easily verified from the explicit formula (2.15), using corollary 1). We are thus reduced to the case where $D$ is an irreducible hypersurface. More generally, if $D$ is a hypersurface whose components are disjoint, we have the equality:
\sgaStableTarget{sga6.e00.eq.2-32}\begin{equation}\tag{2.32}
 \gamma(D)=1-\gamma(L(-D))
\end{equation}
as follows from the exact sequence of sheaves
\[
0\longrightarrow\underline{O}_X(L(-D))\longrightarrow\underline{O}_X\longrightarrow\underline{O}_D\longrightarrow0.
\]

\sgaStableTarget{sga6.e00.thm.2-12}\textbf{Theorem 2.12.} \uline{Let} $X$ \uline{be a nonsingular quasi-projective variety, and} $Z$ \uline{and} $Z'$ \uline{two homogeneous cycles of codimension} $p$ \uline{on} $X$ \uline{that are linearly equivalent. We then have}:
\sgaStableTarget{sga6.e00.eq.2-33}\begin{equation}\tag{2.33}
 \gamma_X(Z)\equiv\gamma_X(Z')\qquad\mathrm{mod}\ \underline K(X)^{p+1}.
\end{equation}
By hypothesis, we can find a cycle $Y$ of dimension $n-p+1$ on
%%%%%%%%%%  scan idx 72  |  folio 59  %%%%%%%%%%
\SGAPrintPage{59}%
$X\times k$ ($n$ is the
dimension of $X$), and two points $a,a'\in k$ such that we have:
\[
 Z=f_a^{-1}(Y),\qquad Z'=f_{a'}^{-1}(Y),
\]
where we put $f_u(x)=(x,u)$. By formula (2.23), we then have:
\[
 \gamma_X(Z)\equiv f_a^!(\gamma_{X\times k}(Y))\qquad\mathrm{mod}\ \underline K(X)^{p+1}
\]
and an analogous formula for $\gamma_X(Z')$. It then suffices to prove that $f_a^!=f_{a'}^!$, which follows immediately from proposition 2.9.
\begin{flushright}q.e.d.\end{flushright}

\sgaStableTarget{sga6.e00.cor.1.occ2}\textbf{Corollary 1:} \uline{The filtration of} $\underline K(X)$ \uline{is compatible with its ring structure, i.e.\ we have} $\underline K(X)^i\underline K(X)^j\subset\underline K(X)^{i+j}$ \uline{for all integers} $i,j\geq0$. \uline{Moreover, we have a natural surjective homomorphism of graded rings}:
\sgaStableTarget{sga6.e00.eq.2-34}\begin{equation}\tag{2.34}
 \varphi:A(X)\longrightarrow G(\underline K(X))
\end{equation}
(where the second member denotes the graded ring associated with $\underline K(X)$ equipped with the filtration $(\underline K(X)^i)_{i\geq0}$) \uline{defined by the condition}:
\[
 \varphi(Z^{n-p})=\gamma(Z^{n-p})\qquad\mathrm{mod}\ \underline K(X)^{p+1}.
\]
This corollary follows immediately from theorem 2.12, propositions 2.6 and 2.7, and Chow's theorem 2.1, by induction on the integer $i+j$.

\sgaStableTarget{sga6.e00.cor.2.occ3}\textbf{Corollary 2:} \uline{Let} $X$ \uline{and} $Y$ \uline{be two nonsingular quasi-projective varieties and} $f$ \uline{a morphism from} $X$ \uline{to} $Y$. \uline{Then the homomorphism} $f^!$ \uline{from} $\underline K(Y)$ \uline{to} $\underline K(X)$ \uline{is compatible with the filtrations, and we have a commutative diagram}:
\[
\begin{tikzcd}[column sep=large,row sep=large]
G(\underline K(Y)) \arrow[r,"G(f^!)"] & G(\underline K(X))\\
A(Y) \arrow[r,"f^*"] \arrow[u,"\varphi_Y"'] & A(X) \arrow[u,"\varphi_X"'] .
\end{tikzcd}
\]
%%%%%%%%%%  scan idx 73  |  folio 60  %%%%%%%%%%
\SGAPrintPage{60}%
Using the graph of $f$, we decompose $f$ into the product of an \uline{immersion} of $X$ into $X\times Y$ and a \uline{projection} from $X\times Y$ onto $Y$. For the latter the corollary is trivial, and for an immersion we proceed as for corollary 1, using proposition 2.8 in place of proposition 2.7.

\sgaStableTarget{sga6.e00.prop.2-13}\textbf{Proposition 2.13 (``cellular decomposition''):} \uline{Let} $P$ \uline{be an algebraic space (irreducible or not), equipped with an increasing family} $P^0\subset P^1\subset\cdots\subset P^n=P$ \uline{of closed subsets such that, for every} $i$ \uline{with} $0\leq i\leq n$, \uline{the connected components of} $P^i-P^{i-1}$ \uline{are algebraic spaces} $V_{i,\alpha}$ \uline{isomorphic to} $k^i$ (\uline{put} $P^{-1}=\varnothing$). \uline{Then}:
\begin{enumerate}[label=\alph*)]
\item \uline{For every algebraic space} $X$, \uline{the natural map from} $\underline K(X)\otimes\underline K(P)$ \uline{to} $\underline K(X\times P)$ \uline{deduced from the tensor product of sheaves on} $X$ \uline{and} $P$ \uline{is surjective}.
\item \uline{The} $\gamma_P(V_{i,\alpha})$ \uline{form a system of generators of the group} $\underline K(P)$.
\end{enumerate}
The proof proceeds by induction on $n$, using the corollary of proposition 2.9 and proposition 2.10. (N.B. Unfortunately, I do not know whether the $\gamma_P(V_{i,\alpha})$ form a basis of $\underline K(P)$ and whether the homomorphism considered in a) is bijective. We shall not need to know this in what follows).

\sgaStableTarget{sga6.e00.thm.2-14}\textbf{Theorem 2.14.}
\uline{Let} $x_i$ ($1\leq i\leq q$) \uline{and} $E_j$ ($1\leq j\leq r$) \uline{be letters, and} $n_j$ ($1\leq j\leq r$) \uline{integers} $>0$; \uline{consider the symbols} $\lambda^k(x_i)$ \uline{with} $k\geq1$ \uline{and} $\lambda^k(E_j)$ \uline{with} $1\leq k\leq n_j$ \uline{as indeterminates, and let} $R$ \uline{be a polynomial with integer coefficients in these indeterminates. Assume that, for every graded ring} $A$, \uline{when in} $R$ \uline{we substitute for the} $x_i$ \uline{elements} $\xi_i$ \uline{and for the} $E_j$ \uline{elements} $\eta_j$ \uline{of}
%%%%%%%%%%  scan idx 74  |  folio 61  %%%%%%%%%%
\SGAPrintPage{61}%
$\Atil$ \uline{subject to the restrictions} $\varepsilon(\xi_i)=0$ \uline{and} $\eta_j=[n_j,1+\eta_j^{(1)}+\cdots+\eta_j^{(n_j)}]$ (thus $\eta_j^{(k)}=0$ for $k>n_j$), \uline{we have}:
\[
 R((\lambda^k(\xi_i)),(\lambda^k(\eta_j)))\in\Atil^p
\]
\uline{Now let} $X$ \uline{be an irreducible nonsingular quasi-projective variety. For arbitrary elements} $x'_i$ \uline{of} $\underline K(X)^1$ \uline{and vector bundles} $E'_j$ \uline{of dimension} $n_j$ \uline{on} $X$ ($1\leq i\leq q$, $1\leq j\leq r$), \uline{we have}
\sgaStableTarget{sga6.e00.eq.unnamed}\[
 R((\lambda^k(x'_i)),(\lambda^k(\gamma(E'_j))))\in\underline K(X)^p\qquad\text{\SGAAuthorFootMark{*}}\ .
\]
\SGAAuthorFootText{*}{We readily prove that the hypothesis made implies that, if $R$ is expressed in terms of the $\gamma^k(\xi_i)$, $\gamma^k(E_j-\varepsilon(E_j))$, then all terms of $R$ have weight $\geq p$ (where $\gamma^k$ has weight $k$). Then 2.14 is a particular case of X 1.3.2.}

\uline{Proof}: By \S 2, the $x'_i$ themselves are differences of elements of the form $\gamma(E)$, where $E$ is a vector bundle. On the other hand, assuming $X$ locally closed in a projective space $P$, it follows from \S 2 and proposition 2.5 that, for every $F\in\cF(X)$, the sheaf $F(n)$ is generated by its sections for sufficiently large $n$. If $F$ is defined by a vector bundle $E$, this proves that $E(n)$ is the inverse image of the canonical vector bundle on a Grassmannian variety $G$ under a morphism from $X$ to $G$. Considering the corresponding morphism from $X$ to $P\times G$, we see that $E$ is the inverse image of a vector bundle on $P\times G$. Thus, applying this to all the vector bundles involved, and using corollary 2 of theorem 2.12, we see that \uline{we are reduced to the case where} $X$ \uline{is a product of Grassmannians}. In this case, the conclusion will follow from the following lemma:

\sgaStableTarget{sga6.e00.lem.2-15}\textbf{Lemma 2.15.}
\uline{Assume that} $X$ \uline{is a product of Grassmannians. Then we can find a} $\lambda$\uline{-homomorphism from} $\underline K(X)$ \uline{to a} $\lambda$\uline{-ring without torsion of the form}
%%%%%%%%%%  scan idx 75  |  folio 62  %%%%%%%%%%
\SGAPrintPage{62}%
$\Atil$ (\uline{where} $A$ \uline{is a graded ring), compatible with the filtrations, and such that the corresponding homomorphism} $G(\underline K(X))\longrightarrow G(A)$ \uline{is injective}.

To prove the lemma, we shall use Chow's theory of Chern classes; a direct proof of the lemma, or of theorem 2.14, would yield a purely ``sheaf-theoretic'' theory of Chern classes (cf.\ \S 5).

We shall consider the homomorphism $\Ctil$ from $\underline K(X)$ to $\Atil(X)$ and put $A=A(X)$; we must prove that the homomorphism $\psi=G(\Ctil)$ from $G(\underline K(X))$ to $G(A(X))\simeq A(X)$ is injective. The fact that $\Ctil$ is compatible with the filtrations holds for every nonsingular quasi-projective variety and follows readily from the corollary to proposition 2.3. On the other hand, as $X$ is a product of Grassmannians $C_j$, we have on $X$ vector bundles $E_j$, inverse images of the canonical bundles on the factors. We shall use the following (known) facts:
\begin{enumerate}[label=\alph*)]
\item On $X$, linear equivalence of cycles is equivalent to numerical equivalence, and $A(X)$ is a group of finite rank;
\item The $C^i(E_j)$ generate the ring $A(X)$.
\end{enumerate}

From a) it follows that $\varphi:A(X)\longrightarrow G(\underline K(X))$ is injective (hence bijective), by the following lemma:

\sgaStableTarget{sga6.e00.lem.2-15.occ2}\textbf{Lemma 2.15.} \uline{Let} $X^n$ \uline{be a nonsingular projective variety, and} $Z^{n-p}$ \uline{a cycle on} $X$ \uline{such that} $\varphi(Z)\in G^p(\underline K(X))$ \uline{is zero; then} $Z^{n-p}$ \uline{is numerically equivalent to zero}.

To prove this lemma, it suffices to use the additive function $\chi(F)=\sum_{i\geq0}(-1)^i\dim H^i(X,F)$ and the fact that this function equals
%%%%%%%%%%  scan idx 76  |  folio 63  %%%%%%%%%%
\SGAPrintPage{63}%
$1$
for $F=\underline{O}_{(x)}$ for every $x\in X$.

Thus, in the case considered, $A(X)$ and $G(\underline K(X))$ have the same finite rank in every degree, and $G(\underline K(X))$ has no torsion, since this holds for $A(X)$. From assertion b) above it follows, in a ``formal'' way, that the map $\Ctil$ from $\underline K(X)$ to $\widetilde{A(X)}$ is surjective, modulo finite groups. It readily follows, since $G(\underline K(X))$ and $A(X)$ have the same finite rank in every degree, that the homomorphism $G(\Ctil)$ from $G(\underline K(X))$ to $G(\widetilde{A(X)})$ is bijective modulo finite groups in every degree (use induction on the degree). As $G(\underline K(X))$ has no torsion, this homomorphism $G(\Ctil)$ is therefore injective.
\begin{flushright}q.e.d.\end{flushright}

N.B. The homomorphism $\Ctil$ is not surjective in general, even when $X$ is a projective space of dimension $2$.

\sgaStableTarget{sga6.e00.sec.5.occ2}\subsection*{\S 5. Sheaf-theoretic definition of Chern classes. Application to the study of injection morphisms.}

\sgaStableTarget{sga6.e00.thm.2-16}\textbf{Theorem 2.16.} \uline{Let} $X$ \uline{be a nonsingular quasi-projective variety. Then, for every} $x\in\underline K(X)$ \uline{and every integer} $i\geq0$, \uline{the element} $\gamma^i(x-\varepsilon(x))$ \uline{has filtration} $\geq i$. \uline{Let} $c^i(x)\in G^i(\underline K(X))$ \uline{be the class of this element modulo} $\underline K(X)^{i+1}$; \uline{if we put}:
\[
 \widetilde c(x)=\left[\varepsilon(x),1+\sum_{i\geq1}c^i(x)\right]\in\widetilde{G(\underline K(X))},
\]
\uline{then} $\widetilde c$ \uline{satisfies the conditions analogous to those of theorem} 2.2.

The first assertion follows from theorem 2.14, as does the fact that $x\longrightarrow\widetilde c(x)$ is a homomorphism of $\lambda$-rings (formula (1.25) implies $\gamma_t(x+y)=\gamma_t(x)\gamma_t(y)$, which already shows that $\widetilde c$ is an additive homomorphism). The ``functoriality'' of $\widetilde c$ is verified immediately, and
%%%%%%%%%%  scan idx 77  |  folio 64  %%%%%%%%%%
\SGAPrintPage{64}%
(2.8) follows from (2.31).
\begin{flushright}q.e.d.\end{flushright}

\sgaStableTarget{sga6.e00.cor.1.occ3}\textbf{Corollary 1:} \uline{We have} $c^i(x)=\varphi(C^i(x))$.

This follows from a uniqueness theorem for Chern classes, proved as in Hirzebruch by passage to flag varieties. We need only show that, if $Y$ is a flag variety associated with a vector bundle on $X$, then $G(\underline K(X))\longrightarrow G(\underline K(Y))$ is injective, which is verified by explicitly determining $\underline K(Y)$ with all its structures from $\underline K(X)$ (I shall not give the details here).

\sgaStableTarget{sga6.e00.cor.2.occ4}\textbf{Corollary 2:} \uline{Assume that, for every element} $Z\in A(X)$, \uline{there exists a morphism} $f$ \uline{from} $X$ \uline{to a product} $X'$ \uline{of Grassmannians such that} $Z$ \uline{is in the image of} $f^*$. \uline{Then the canonical homomorphisms}
\[
 \varphi:A(X)\longrightarrow G(\underline K(X)),\qquad \psi:G(\underline K(X))\longrightarrow A(X),
\]
\uline{satisfy in dimension} $i$ \uline{the following relations}:
\sgaStableTarget{sga6.e00.eq.2-35}\begin{equation}\tag{2.35}
 \psi\varphi=(-1)^{i-1}(i-1)!\,\underline 1,\qquad
 \varphi\psi=(-1)^{i-1}(i-1)!\,\underline 1.
\end{equation}

As $\varphi$ is surjective, it evidently suffices to prove the first formula, and for this we can assume $X$ a product of Grassmannians. As $\varphi$ is then injective, it suffices to prove the second formula (2.35), which, by corollary 1, is written
\[
 \gamma^i(x-\varepsilon(x))\equiv(-1)^{i-1}(i-1)!\,x\qquad\mathrm{mod}\ \underline K(X)^{i+1}
\]
for every $x\in\underline K(X)^i$. As $\psi$ is \uline{injective}, it suffices to prove the analogous relation for $\Ctil(x)\in\Atil(X)^i$, and this is then a consequence of (1.22 bis).

\sgaStableTarget{sga6.e00.rem.unnamed.occ4}\textbf{Remark:} Corollary 2 makes the validity of (2.35) in all cases plausible; there should be a direct proof of it. In
%%%%%%%%%%  scan idx 78  |  folio 65  %%%%%%%%%%
\SGAPrintPage{65}%
any case, we readily prove the second formula (2.35) modulo the torsion subgroup of $G(\underline K(X))$, and, in characteristic $0$, this same formula is a consequence of the Riemann--Roch theorem. The interest of (2.35) (when it holds) is that it shows that $\varphi$ is injective modulo torsion groups (perhaps always zero!\SGAAuthorFootnote{*}{No, cf.\ the note on page 30.}), so that we ``do not lose much'' by defining the Chern classes as elements of $G(\underline K(X))$; after taking the tensor product with $\Q$, as we shall do later for the Riemann--Roch theorem, the two definitions of Chern classes would then be identical.

Let $X$ still be nonsingular and quasi-projective; let also $Y$ be a nonsingular subvariety of $X$, and $i$ the injection of $Y$ into $X$. For $i^!$ and $i_!$, we then have the following formulas:
\begin{enumerate}[label=(\roman*)]
\item $i^!$ is a homomorphism of $\lambda$-rings
\item $i_!(x.i^!(y))=i_!(x).y$, whence in particular $i_!i^!(y)=\xi y$ with $\xi=i_!i^!(1)=\gamma_X(Y)$,
\item $i^!(i_!(x))=x.\lambda_{-1}(N)$, whence in particular $i^!(\xi)=i^!i_!(1)=\lambda_{-1}(N)$,
\item $i_!(x)i_!(x')=i_!(xx'\lambda_{-1}(N))$,
\end{enumerate}
where, by definition, $N$ is the class in $\underline K(Y)$ of the conormal bundle of $Y$ in $X$. Formulas (i) and (ii) hold for arbitrary morphisms, not necessarily injections, and (iii) and (iv) are verified from definition (2.15) and the analogous formula for inverse images (taking for $x$ and $x'$ classes of vector bundles, and using the fact that
\sgaStableTarget{sga6.e00.eq.2-36}\begin{equation}\tag{2.36}
 \Tor^{\underline{O}_X}_i(\underline{O}_Y,\underline{O}_Y)=\underline{O}_Y(\wedge^i N).
\end{equation}
)
%%%%%%%%%%  scan idx 79  |  folio 66  %%%%%%%%%%
\SGAPrintPage{66}%
Formula (iv) expresses that $i_!$ is a ring homomorphism when we introduce on $\underline K(Y)$ a new multiplicative structure, by the definition $x_N x'=xx'\lambda_{-1}(N)$ (cf.\ Chapter I, \S 4). It is completed by the formula:
\begin{enumerate}[label=(\roman*)]
\setcounter{enumi}{4}
\item $i_!(\lambda^p(N,x))=\lambda^p(i_!(x))$ for $p\geq1$,
\end{enumerate}
which means that $i_!$ defines a $\lambda$-homomorphism from $\underline K(Y)_N$ to $\underline K(X)$. This formula (v), certainly true in all cases, has for the moment been proved only in characteristic $0$, since it follows from the conjunction of theorems 1.4 and 2.8.

We have analogous formulas for $i_*$ and $i^*$:
\begin{enumerate}[label=(\roman* bis)]
\item $i^*$ is a homomorphism of graded rings.
\item $i_*(x.i^*(y))=i_*(x).y$, whence in particular $i_*(i^*(y))=\xi' y$, with $\xi'=i_*(1)$,
\item $i^*(i_*(x))=x.(-1)^q N^q$, whence in particular $i^*(\xi')=(-1)^q N^q$,
\item $i_*(x)i_*(x')=i_*(xx'(-1)^q N^q)$,
\item $i_*$ increases degrees by $q$ units.
\end{enumerate}

In these formulas, $q$ is the codimension of $Y$ in $X$, and we put $N^i=C^i(N)$. Only formulas (iii bis) and (iv bis) require a proof; I do not know whether they are well known (or even whether they are true!!!), and merely note that they follow by reduction of formulas (iii) and (iv), \uline{provided that we work in} $G(\underline K(Y))$ \uline{and} $G(\underline K(X))$ \uline{and not in} $A(Y)$ \uline{and} $A(X)$. We then use the fact that $i_*$ and $i^*$ arise by passage to the quotient from $i_!$ and $i^!$, noting that $i_!$ shifts the filtration by $q$ units.

%%%%%%%%%%  scan idx 80  |  folio 67  %%%%%%%%%%
\SGAPrintPage{67}%
N.B. All the ``universal'' formulas relating $i_!$ and $i^!$ (resp.\ $i_*$ and $i^*$) are formal consequences of the preceding ones, as one can convince oneself without much difficulty.

From formula (v) follows the formula
\[
 i_!(\gamma^p(N,x))=\gamma^p(i_!(x)).
\]
The conjunction of theorem 2.14 and proposition 1.5 shows that $\gamma^p(N,x)$ has, for $p=q+j$, filtration $\geq j$; denoting by $G_p(N,x)$ its class in $G^j(\underline K(Y))$, theorem 2.16 shows that we have
\[
 i_*(G_p(N,x))=c^p(i_!(x)).
\]
Finally, proposition 1.5 together with theorem 2.14 allows us to express $G_p(N,x)$ as a polynomial with integer coefficients in the $c^i(N)$, the $c^i(x)$ and $\varepsilon(x)$; we deduce the theorem:

\sgaStableTarget{sga6.e00.thm.2-17}\textbf{Theorem 2.17.} \uline{Assume that the field} $k$ \uline{has characteristic} $0$ \uline{and let} $i:Y\longrightarrow X$ \uline{be an injection morphism of nonsingular quasi-projective varieties. We then have formulas (i) to (v), and, moreover, for every integer} $j\geq0$, \uline{the formula}:
\sgaStableTarget{sga6.e00.eq.2-37}\begin{equation}\tag{2.37}
 c^{q+j}(i_!(x))=i_*\left(\frac{\widetilde c(x)*\lambda_{-1}(\widetilde c(N))}{(-1)^qN^q}\right)^{(j)}.
\end{equation}

The right-hand side of formula (2.37) is an abbreviated way of denoting the polynomial in the components of $\widetilde c(x)$ and $\widetilde c(N)$ described in proposition 1.5.

\sgaStableTarget{sga6.e00.cor.unnamed.occ5}\textbf{Corollary:} \uline{Under the preceding conditions, we have}:
\sgaStableTarget{sga6.e00.eq.2-38}\begin{equation}\tag{2.38}
 \ch(\widetilde c(i_!(x)))=i_*\left(\ch(\widetilde c(x))\,\mathcal T(\check N)^{-1}\right)
\end{equation}
\uline{where} $\ch$ \uline{is the Chern homomorphism defined by formula} (1.26).

%%%%%%%%%%  scan idx 81  |  folio 68  %%%%%%%%%%
\SGAPrintPage{68}%
It suffices to use (2.37) written in the form:
\sgaStableTarget{sga6.e00.eq.2-37bis}\begin{equation}\tag{2.37 bis}
 \widetilde c(i_!(x))=1+i_*\left(\frac{\widetilde c(x)*\lambda_{-1}(\widetilde c(N))-1}{(-1)^qN^q}\right),
\end{equation}
the explicit formula (1.29), formula (iv bis), and finally (1.30).

We can write (2.38) in the form
\sgaStableTarget{sga6.e00.eq.2-38bis}\begin{equation}\tag{2.38 bis}
 \ch_X(F)=i_*\left(\ch_Y(F)\,\mathcal T(T_{X/Y})^{-1}\right),
\end{equation}
where $F$ is a coherent algebraic sheaf on $Y$ (considered on the left-hand side as a sheaf on $X$ zero outside $Y$), $\ch_X$ and $\ch_Y$ are the Chern homomorphisms taken respectively on $X$ and $Y$, and $T_{X/Y}$ is the normal bundle of $Y$ in $X$.

\sgaStableTarget{sga6.e00.rem.unnamed.occ5}\textbf{Remarks:} Unlike (2.38), formula (2.37) is a ``formula without denominators''. It has been proved (even in characteristic $0$) only by working in $G(\underline K(X))$ and not in $A(X)$, which may not be the same thing. It is certainly true in every characteristic and in $A(X)$. Let us note that we immediately verify that the images of the two sides under $i^*$ are the same, whence it follows, on applying $i_*$, that the product of their difference by $\xi'$ is zero. This suggests a completely different method of proof, by trying to obtain $Y$ as the transverse inverse image of a subvariety $Y'$ in a variety $X'$ such that, in dimension $\leq n=\dim X$, the class of $Y'$ in $A(X')$ is not a divisor of $0$\dots. Unfortunately, we cannot hope that $X'$ is a Grassmannian or something close to one.

We shall note that, if we disregard torsion elements, formula (2.38) is equivalent to (2.37).




% Additional macro definitions
\providecommand{\cS}{\mathcal S}
\providecommand{\cD}{\mathcal D}
\providecommand{\cE}{\mathscr E}
\providecommand{\cG}{\mathscr G}
\providecommand{\cL}{\mathscr L}
\providecommand{\cM}{\mathscr M}
\providecommand{\cP}{\mathscr P}
\providecommand{\cK}{\mathcal K}
\providecommand{\Db}{\mathrm D}
\providecommand{\Kb}{\mathrm K}
\providecommand{\Trc}{\operatorname{Tr}}
\providecommand{\ob}{\operatorname{ob}}
\providecommand{\Coh}{\operatorname{Coh}}
\providecommand{\Cone}{\operatorname{Cone}}
\providecommand{\im}{\operatorname{Im}}
\providecommand{\Id}{\operatorname{Id}}
\providecommand{\RR}{\mathrm{R.R.}}
\providecommand{\RRR}{\mathrm{R.R.R.}}
\providecommand{\GrR}{\mathrm{G}}
\providecommand{\pr}{\operatorname{pr}}


%%%%%%%%%%  scan idx 82  |  folio 69  %%%%%%%%%%
\SGAPrintPage{69}%
\sgaStableTarget{sga6.e00.sec.6.occ2}\section*{\S 6. The Riemann--Roch theorem}

Its statement is as follows: \uline{let} \(f\) \uline{be a} proper \uline{morphism from a nonsingular quasi-projective variety} \(Y\) \uline{to a variety} \(X\) \uline{of the same kind. Then we have}:
\sgaStableTarget{sga6.e00.eq.2-39}\begin{equation}\tag{2.39}
 f_*\bigl(\ch_Y(x)\,\mathcal T(Y)\bigr)
 =
 \ch_X\bigl(f_!(x)\bigr)\,\mathcal T(X).
\end{equation}
(Both sides are in \(A(X)\otimes\mathbb Q\) and \(\mathcal T(X)=\mathcal T(C(T_X))\).)

It allows us to calculate, up to torsion elements, the Chern classes of \(f_!(x)\) from the Chern classes of \(x\), those of \(T_X\) and \(T_Y\), and the homomorphism \(f_*\). Or, again, it allows us to determine the homomorphism
\(f_!\) from \(\underline K(Y)\otimes\mathbb Q\) to \(\underline K(X)\otimes\mathbb Q\), when we identify these groups by means of the homomorphisms \(\ch_Y\) and \(\ch_X\) respectively with \(A(Y)\otimes\mathbb Q\) and \(A(X)\otimes\mathbb Q\) (if formula (2.35) is true!), using \(f_*\) and the Chern classes of \(T_X\) and \(T_Y\).

I know how to prove formula (2.39) only in characteristic \(0\)\SGAAuthorFootnote{*}{See below for a proof valid in every characteristic, based on a different principle (p.\,51).}, and as an equality in \(G(\underline K(X))\otimes\mathbb Q\). For this, since \(Y\) is locally closed in a projective space \(P\), we decompose \(f\) into the product of an injection morphism from \(Y\) to \(P\times X\), and a projection morphism from \(P\times X\) onto \(X\), and prove (2.39) for each of these two morphisms.

For an injection morphism, we see immediately that formulas (2.38) and (2.39) are equivalent (the characteristic \(0\) hypothesis entered only through theorem 2.8, which allows us to express the \(\lambda^i(\gamma(F))\), for a coherent algebraic sheaf \(F\), by \uline{local} cohomological invariants of \(F\)).

For a projection morphism \(f:P\times X\longrightarrow X\), we can, by proposition 2.13, assume that \(x\) is of the form \(y\otimes z\), with \(y\in\underline K(P)\) and
%%%%%%%%%%  scan idx 83  |  folio 70  %%%%%%%%%%
\SGAPrintPage{70}%
\(z\in\underline K(X)\).
Then the K\"unneth formula gives \(f_!(x)=\chi(y)z\), where, for a coherent algebraic sheaf \(F\) on \(P\), we put:
\[
  \chi(F)=\sum_{i\geq 0}(-1)^i\dim H^i(P,F).
\]
We are then immediately reduced to proving the formula
\sgaStableTarget{sga6.e00.eq.2-40}\begin{equation}\tag{2.40}
  \chi(F)=\varkappa_r\bigl(\ch_P(F)\,\mathcal T(P)\bigr)
\end{equation}
where \(\varkappa_r\) denotes the component of degree \(r=\dim P\). By proposition 2.13, we can assume that \(F=\underline{O}_Q\), where \(Q\) is a linear subvariety of \(P\). Then \(\chi(F)=1\), since it is the arithmetic genus of \(Q\); it remains to see that the right-hand side of formula (2.40) equals \(1\).

Indeed, if \(V\) is a hyperplane of \(P\) and we put \(\xi=\gamma_P(\underline{O}_V)\) and \(L=\ell(V)\), we shall have:
\[
  1-\xi=\gamma_P(E(-L))
\]
(where \(E(-L)\) is the vector bundle defined by \(-L\)), whence:
\[
  C_P(1-\xi)=[1,1-L],
\]
and consequently
\[
  \ch_P(1-\xi)=e^{-L},\qquad\text{whence}\qquad \ch_P(\xi)=1-e^{-L},
\]
and therefore
\[
  \ch_P(\xi^p)=(1-e^{-L})^p.
\]
Now, if \(p\) is the codimension of \(Q\) in \(P\), we have \(\xi^p=\gamma_P(Q)\); we deduce
\[
  \ch_P(F)\,\mathcal T(P)
  =
  \bigl(L/(1-e^{-L})\bigr)^{r+1}(1-e^{-L})^p
\]
and we must verify that the coefficient of \(L^r\) in this formal series in \(L\) equals \(1\), which follows from the characteristic property of the formal series \(L/(1-e^{-L})\), namely that the coefficient of \(L^q\) in \(\bigl(L/(1-e^{-L})\bigr)^{q+1}\) equals \(1\) (put \(q=r-p\)).


%%%%%%%%%%  scan idx 84  |  folio 71  %%%%%%%%%%
\SGAPrintPage{71}%
\sgaStableTarget{sga6.e00.bibliographical-references}\subsection*{Bibliographical references:}
\begin{enumerate}[label={[\arabic*]}]
\item \uline{Hirzebruch}: Neue topologische Methoden in der algebraischen Geometrie.
\item \uline{Serre}: Faisceaux algébriques cohérents, Ann. of Maths.
\item \uline{Grothendieck}: Sur les faisceaux algébriques cohérents, in Séminaire Cartan 1956/57.
\end{enumerate}

\sgaStableTarget{sga6.e00.proof-of-the-riemann-roch-theorem}\section*{Proof of the Riemann--Roch theorem}
\begin{center}
(Grothendieck, 1 November 1957)
\end{center}

\sgaStableTarget{sga6.e00.reminder-of-notation}\subsection*{* Reminder of notation}
\SGAAuthorFootText{unmarked}{The passages between asterisks \textsuperscript{*}\dots\textsubscript{*} are interpolations by the copyist.}
[These are those of the ``report on Riemann--Roch'', cited as R.R.R. in what follows.]

All varieties considered are nonsingular and quasi-projective (i.e.\ isomorphic to a locally closed subvariety of a projective space). The base field is algebraically closed, of arbitrary characteristic.

For such a variety \(X\), we denote by \(K(X)\) the abelian group generated by the classes of coherent sheaves on \(X\) (modulo the identification of an extension with a sum); the alternating sum of Tor makes \(K(X)\) a ring. If \(\underline F\) is a sheaf (resp.\ a vector bundle \(E\), resp.\ a sub%
%%%%%%%%%%  scan idx 85  |  folio 72  %%%%%%%%%%
\SGAPrintPage{72}%
variety \(Y\) of \(X\)), we denote by \(\gamma_X(\underline F)\) (resp.\ \(\gamma_X(E)\), resp.\ \(\gamma_X(Y)\)) the element of \(K(X)\) corresponding to \(\underline F\) (resp.\ to \(\underline E\), resp.\ to \(\underline{O}_Y\)). (In fact, we often write \(\underline F\), \(E\) instead of \(\gamma_X(\underline F)\)\dots).

We denote by \(A(X)\) the graded ring of classes of cycles on \(X\) (for rational equivalence---cf.\ Chow); \(A^i(X)\) consists of cycles of codimension \(i\).

For a vector bundle \(E\), we denote by \(C(E)=\sum C^i(E)\), \(C^i(E)\in A^i(X)\), its Chern class (defined as the inverse image of Schubert cycles); by linearity we deduce the \(C^i(y)\), for every \(y\in K(X)\). If \(C(y)=\prod(1+a_i)\), we put, following Hirzebruch:
\[
  \ch(y)=\sum e^{a_i},\qquad
  T(y)=\prod\frac{a_i}{1-e^{-a_i}}.
\]
These are elements of \(A(X)\otimes\mathbb Q\).

When \(y\) is the tangent bundle of \(X\), we write \(T(X)\) instead of \(T(y)\).

For a vector bundle \(G\), we denote by \(\lambda^pG\) its \(p\)-th exterior power, and put \(\lambda_{-1}(G)=\sum(-1)^p\lambda^p(G)\); this is an element of \(K(X)\).

For a proper regular map \(f\) from a variety \(Y\) to a variety \(X\), we define (additive) homomorphisms:
\[
  f^*:A(X)\longrightarrow A(Y)\qquad\text{and}\qquad
  f_*:A(Y)\longrightarrow A(X).
\]

The first is homogeneous of degree \(0\) and multiplicative. The second is homogeneous of degree \(\dim.X-\dim.Y\).

We also define homomorphisms:
\[
  f^!:K(X)\longrightarrow K(Y)\qquad\text{and}\qquad
  f_!:K(Y)\longrightarrow K(X).
\]

For a coherent sheaf \(\underline F\) on \(X\), \(f^!(\underline F)\) equals, by definition, the alternating sum of the \(\Tor_p^{\underline{O}_X}(\underline{O}_Y,\underline F)\). For vector bundles, this is simply the inverse-image operation.

%%%%%%%%%%  scan idx 86  |  folio 73  %%%%%%%%%%
\SGAPrintPage{73}%
For a coherent sheaf \(G\) on \(Y\), \(f_*(\underline G)\) is coherent on \(X\). Denoting by \(R^pf_*\) the \(p\)-th derived functor of the functor \(f_*\) thus defined, we have, by definition, \(f_!(\underline G)=\sum(-1)^p R^pf_*(\underline G)\).\quad\textsubscript{*}

\sgaStableTarget{sga6.e00.riemann-roch-theorem}\textbf{RIEMANN--ROCH THEOREM.} \uline{Let} \(f:Y\longrightarrow X\) \uline{be a proper morphism} (\textsuperscript{*} regular map \textsubscript{*}) \uline{from a nonsingular quasi-projective variety} \(Y\) \uline{to another} \(X\), \uline{and let} \(y\in K(Y)\). \uline{We have}:
\sgaStableTarget{sga6.e00.eq.1}\begin{equation}\tag{1}
  f_*\bigl(\ch(y)T(Y)\bigr)=\ch(f_!(y))T(X)
  \quad\text{\uline{in} } A(X)\otimes\mathbb Q.
\end{equation}

\uline{Proof} (\textsuperscript{*} abridged \textsubscript{*}).

\sgaStableTarget{sga6.e00.lem.1}\textbf{LEMMA 1.} \uline{Consider proper morphisms of varieties} \(Z\xrightarrow{g}Y\xrightarrow{f}X\). \uline{If the theorem} \(\RR\) \uline{is true for} \(f\) \uline{and} \(g\), \uline{it is true for} \(f\circ g\). \uline{If it is true for} \(fg\) \uline{and} \(g\), \uline{it is true for} \(f\). \uline{If} \(f_*\) \uline{is injective, and the theorem} \(\RR\) \uline{is true for} \(fg\) \uline{and} \(f\), \uline{it is true for} \(g\).

Trivial. (\textsuperscript{*} In fact, when we assume \(\RR\) true for \(fg\) and \(g\), we must also assume that \(g_!\) is surjective to be able to assert that \(\RR\) is true for \(f\). \textsubscript{*})

As a proper morphism is the composite of an injection and a projection \(X\times P\longrightarrow X\) (where \(P\) is a projective space), we are reduced to treating these two cases separately. The second, as we have seen, reduces to a standard calculation, by proposition 2.13 (``cellular decomposition'') of R.R.R. We shall therefore assume \(Y\subset X\), with the injection morphism \(i\). Formula (1) then becomes:
\sgaStableTarget{sga6.e00.eq.2}\begin{equation}\tag{2}
  \ch(i_!(y))=i_*\bigl(\ch(y)T(E)^{-1}\bigr),
\end{equation}
where \(E=T_{X/Y}\) is the normal bundle of \(Y\) in \(X\).

%%%%%%%%%%  scan idx 87  |  folio 74  %%%%%%%%%%
\SGAPrintPage{74}%
\sgaStableTarget{sga6.e00.lem.2}\textbf{LEMMA 2.} \uline{Assume that there exists on} \(X\) \uline{a vector bundle} \(G\) \uline{of rank} \(p=\codim_XY\), \uline{such that}:
\sgaStableTarget{sga6.e00.eq.3}\begin{equation}\tag{3}
  \gamma_X(Y)=i_!(1)=\lambda_{-1}(G)\,;\qquad
  i^!(G)=\check E\,;\qquad
  (-1)^pC^p(G)=Y.
\end{equation}
\uline{If we then have}:
\sgaStableTarget{sga6.e00.eq.4}\begin{equation}\tag{4}
  y=i^!(x),\qquad \text{\uline{with} }x\in K(X),
\end{equation}
\uline{formula (2) is valid}.

The proof is immediate from the fact that \(\ch\) is a ring homomorphism, and that we have:
\sgaStableTarget{sga6.e00.eq.5}\begin{equation}\tag{5}
  \ch(\lambda_{-1}(G))=(-1)^pC^p(G)T(\check G)^{-1}.
\end{equation}

\sgaStableTarget{sga6.e00.cor.unnamed.occ6}\textbf{COROLLARY.} \(\RR\) \uline{is true if} \(Y\) \uline{is a hypersurface and} \(y\) \uline{is of the form} \(i^!(x)\), \uline{with} \(x\in K(X)\).

\sgaStableTarget{sga6.e00.lem.3}\textbf{LEMMA 3.} \uline{Assume that} \(X=Y\times P\), \uline{with} \(P\) \uline{a projective space, and that} \(i\) \uline{is the injection} \(u\longrightarrow (u,v_0)\) \(v_0\in P\). \uline{Then} \(i_*\) \uline{is injective and} \(\RR\) \uline{is true}.

This is pure standard calculation.

Assume again that \(X,Y,y\) are arbitrary (still with \(Y\subset X\)). Let \(X'\) be the variety obtained by blowing up \(X\) along \(Y\), and \(Y'\) the inverse-image hypersurface of \(Y\). We have a diagram of canonical maps
\[
\begin{tikzcd}
Y \arrow[r,"i"] & X\\
Y' \arrow[u,"g"] \arrow[r,"j"] & X' \arrow[u,"f"']
\end{tikzcd}
\]
(\(j\) being the canonical map). The inverse image of \(E\) on \(Y'\) has a subbundle of rank \(1\), whose dual is denoted by \(L\). Moreover, \(L^{-1}\) is also the restriction to \(Y'\) of the bundle \(L(Y')\) defined on \(X'\) by the divisor \(Y'\). We put \(F=E/L^{-1}\). Elementary calculations give:

%%%%%%%%%%  scan idx 88  |  folio 75  %%%%%%%%%%
\SGAPrintPage{75}%
\sgaStableTarget{sga6.e00.lem.4}\textbf{LEMMA 4.} \uline{We have the formulas (where} \(p=\codim_XY\)):
\sgaStableTarget{sga6.e00.eq.6}\sgaStableTarget{sga6.e00.eq.7}\sgaStableTarget{sga6.e00.eq.8}\sgaStableTarget{sga6.e00.eq.9}\begin{align}
  f^!i_!(y)&=j_!\bigl(\lambda_{-1}(\check F)y\bigr). \tag{6}\\
  \lambda_{-1}(\check F)&\equiv \gamma^{p-1}(E)\lambda^p(\check E)
    \quad\text{mod.}(1-L)\text{\SGAAuthorFootMark{*}}. \tag{7}\\
  f_*(1)&=1,\quad\text{\uline{whence} }f_*f^*=1. \tag{8}\\
  f^*i_*(\alpha)&\equiv j_*\bigl(\alpha C^{p-1}(F)\bigr)
    \quad\text{mod.}\Ker f_*\quad\text{\uline{for every} }\alpha\in A(Y). \tag{9}
\end{align}
\SGAAuthorFootText{*}{This unnecessarily complicated formula can be replaced, for example, by VII 4.1.}

[To simplify notation, we have made \(K(Y)\) operate on \(K(Y')\) and \(A(Y)\) on \(A(Y')\) by means of the homomorphisms \(g^!\) and \(g^*\). {*}In formula (7), it seems that \(\gamma^{p-1}(E)\) denotes \(\lambda^{p-1}(E-2)\), that is, \(\sum_{i+j=p-1}\binom{-p}{j}\lambda^i(E)\); the congruence holds in \(K(Y')\).\textsubscript{*}]

Formula (6) follows from the equality:
\sgaStableTarget{sga6.e00.eq.6bis}\begin{equation}\tag{6 bis}
  \Tor_i^{\underline{O}_X(v)}\bigl(\underline{O}_Y(v),\underline{O}_{Y'}(v')\bigr)
  =
  \wedge^i\check F(v')
  \quad (v'\in Y',\ v=g(v')).
\end{equation}

Formula (7) follows from the formulas
\sgaStableTarget{sga6.e00.eq.10}\begin{equation}\tag{10}
  g_!(L^{-i})=0\quad\text{for }0<i\leq p-1,\qquad g_!(1)=1,
\end{equation}
which themselves follow from K\"unneth and the cohomology of projective spaces.

Formula (8) is trivial; finally, (9) follows from \(g_*(C^{p-1}(F))=1\), which is deduced, for example, from the immediate formulas:
\sgaStableTarget{sga6.e00.eq.10bis}\begin{equation}\tag{10 bis}
  g_*(\xi^i)=0\quad\text{if }0\leq i<p-1,\qquad
  g_*(\xi^{p-1})=1,
\end{equation}
with \(\xi=C^1(L)\in A^1(Y')\).

[In fact, (9) is very probably an identity, and not merely a congruence.]


%%%%%%%%%%  scan idx 89  |  folio 76  %%%%%%%%%%
\SGAPrintPage{76}%
Formula (7) immediately gives \(y\lambda_{-1}(\check F)\in j^!(K(X'))\), provided that we have:
\sgaStableTarget{sga6.e00.eq.11}\begin{equation}\tag{11}
  y\,\gamma^{p-1}(E)\lambda^p(\check E)\in i^!(K(X)).
\end{equation}
Assume that this condition (11) is satisfied. By the corollary to lemma 2, \(\RR\) is true for \((Y',X',y\lambda_{-1}(\check F))\). On the other hand, to prove (2), it suffices, by formula (8), to prove that
\sgaStableTarget{sga6.e00.eq.12}\begin{equation}\tag{12}
  f^*\ch(i_!(y))
  \equiv
  f^*\bigl(i_*(\ch(y)T(E)^{-1})\bigr)
  \quad\text{mod.}\Ker f_*.
\end{equation}
Now, this last formula is verified immediately, using (6), (9), \(\RR\) for \((Y',X',y\lambda_{-1}(\check F))\) and the formula:
\sgaStableTarget{sga6.e00.eq.13}\begin{equation}\tag{13}
  \ch(\lambda_{-1}(\check F))=C^{p-1}(F)T(F)^{-1}\quad\text{(cf.\ formula (5)).}
\end{equation}

Thus, \(\RR\) is proved whenever \(y\,\gamma^{p-1}(E)=0\), in particular (since \(\gamma^{p-1}(E)\) has filtration \(\geq p-1\) \textsuperscript{*}in \(K(Y)\) filtered by the codimension of supports\textsubscript{*}) whenever \(p-1>\dim.Y\), i.e.\ whenever \(\dim.Y<\frac{\dim.X-1}{2}\). Now, we reduce to this case by combining lemma 1 and lemma 3 (taking for \(P\) a projective space of sufficiently large dimension \(r\), for example \(r>\dim.X+1\)). Q.E.D.

\sgaStableTarget{sga6.e00.rem.unnamed.occ6}N.B. We have used the fact that \(\gamma^{p-1}(E)\) has filtration \(\geq p-1\). We can prove this fact directly, without passing through Grassmannians as in \(\mathrm{R.R.R.}\) 2.16, by studying the bundle \(Y'\) of projective spaces associated with \(E\). This method gives, more generally, th.\ 2.14 of \(\mathrm{R.R.R.}\)

\sgaStableTarget{sga6.e00.extract-from-a-letter-of-grothendieck}\subsection*{Extract from a letter of Grothendieck}
\dots{} ``morally'', the new method rests on the determination of \(K(X)\) and \(A(X)\) when \(X\) is either a projective bundle associated with a vector bundle, or a variety obtained by blowing up a non%
%%%%%%%%%%  scan idx 90  |  folio 77  %%%%%%%%%%
\SGAPrintPage{77}%
singular subvariety. In this second case, I have not made the complete determination; a little something is missing. With the notation of the short paper, we must prove formula (9) (but as an \uline{equality}):
\sgaStableTarget{sga6.e00.eq.9.occ2}\begin{equation}\tag{9}
  f^*(i_*(\alpha))=j_*(\alpha C^{p-1}(F)),\qquad \alpha\in A(Y).
\end{equation}
We observe that the two sides have the same image under \(f_*\), and the same image under \(j^*\), so that the question is equivalent to:
\sgaStableTarget{sga6.e00.eq.9bis}\begin{equation}\tag{9 bis}
  \Ker f_*\cap\Ker j^*=0.
\end{equation}
In any case, (9) is true upon reduction to \(G(K(X'))\) (which destroys a little torsion, at most), since it then follows from (6); there is therefore hardly any doubt that it is true.

\dots I also point out to you that the \uline{equality} (9) follows very easily from (6) without passing through the troublesome (6 bis)\dots

\begin{flushleft}
p.c.c.J-P.S.
\end{flushleft}

\newpage
%%%%%%%%%%  scan idx 91  |  folio 78  %%%%%%%%%%
\SGAPrintPage{78}%
\sgaStableTarget{sga6.e01.expose}\section*{\underline{\underline{GENERALITIES ON FINITENESS CONDITIONS IN}}\\[3pt] \underline{\underline{DERIVED CATEGORIES}}}
\begin{center}
\underline{by L. ILLUSIE}
\end{center}

\sgaStableTarget{sga6.e01.sec.0}\subsection*{0. Introduction}
The purpose of Exposés I to IV is to develop, with the generality appropriate to this Seminar, the formalism of finiteness conditions in the derived categories of ringed topoi.

As was said in Exp.\ 0, it was the need to define, on arbitrary schemes, ``Grothendieck groups'' with good variance properties that required us to generalize and relax the notions of finiteness used until now.

A classical notion such as that of a coherent sheaf on a ringed space \((X,\underline{O}_X)\) becomes uninteresting as soon as \(\underline{O}_X\) ceases to be coherent. One might think of replacing it by the notion of finite presentation, but this has the drawback that the kernel of an epimorphism of Modules of finite presentation is in general no longer of finite presentation. We arrive at a satisfactory notion by observing that, when \(\underline{O}_X\) is coherent, a coherent sheaf \(F\) is not only of finite presentation but of \(n\)-finite presentation for every \(n\in\mathbb N\), ``\(n\)-finite presentation'' meaning that there locally exists an exact sequence
%%%%%%%%%%  scan idx 92  |  folio 79  %%%%%%%%%%
\SGAPrintPage{79}%
\[
  L_n\longrightarrow L_{n-1}\longrightarrow \ldots \longrightarrow L_0\longrightarrow F\longrightarrow 0,
\]
where the \(L_i\) are free of finite type. With the sheaf \(\underline{O}_X\) no longer assumed coherent, let us say that an \(\underline{O}_X\)-Module \(F\) is \uline{pseudo-coherent} when it satisfies the preceding condition, i.e.\ is of \(n\)-finite presentation for every \(n\). Then the notion of pseudo-coherence has, with respect to short exact sequences, the same stability property as coherence: when two of the terms of a short exact sequence of \(\underline{O}_X\)-Modules are pseudo-coherent, the third is too. We readily generalize the preceding notion to complexes of sheaves. A \uline{complex} \(F\) of \(\underline{O}_X\)-Modules is called \uline{pseudo-coherent} when it is \(n\)-pseudo-coherent for every \(n\in\mathbb Z\), ``\(n\)-pseudo-coherent'' meaning that there locally exists a quasi-isomorphism \(L\to F\), where \(L\) is a complex with degrees bounded above and components free of finite type in degree \(\geq n\). For a complex \(F\) concentrated in degree \(0\), it amounts to the same thing to say that \(F\) is pseudo-coherent as a complex or as a Module. The notion of a pseudo-coherent complex has excellent stability properties, described in Exp.\ I \S\S~1 and 2. First, it is stable under isomorphism in the derived category \(\Db(X)\) of the category of \(\underline{O}_X\)-Modules; moreover, when two of the objects of a distinguished triangle of \(\Db(X)\) are pseudo-coherent, the third is too: in other words, the full subcategory \(\Db(X)_{\mathrm{coh}}\) of \(\Db(X)\) formed by pseudo-coherent complexes is a triangulated subcategory. Furthermore, pseudo-coherence is preserved by derived inverse image and tensor product. When \(\underline{O}_X\) is coherent, saying that a complex \(F\) is pseudo-coherent simply means that \(F\) has locally bounded-above cohomology and that all the \(\underline H^i(F)\) are coherent.

When \(\underline{O}_X\) is a sheaf of regular local rings, every \(\underline{O}_X\)-Module
%%%%%%%%%%  scan idx 93  |  folio 80  %%%%%%%%%%
\SGAPrintPage{80}%
\(F\) that is coherent locally admits a \uline{finite} resolution to the left by free Modules of finite type, i.e.\ there locally exists an exact sequence
\[
  0\longrightarrow L_n\longrightarrow \ldots \longrightarrow L_0\longrightarrow F\longrightarrow 0,
\]
where the \(L_i\) are free of finite type. It is therefore natural, when \(\underline{O}_X\) is an arbitrary sheaf of local rings, to consider Modules with the preceding property. Such Modules are called \uline{perfect}. More generally, a \uline{complex} \(F\) of \(\underline{O}_X\)-Modules is called \uline{perfect} when there locally exists a quasi-isomorphism \(L\to F\), where \(L\) is a complex with bounded degrees and components free of finite type. It is shown in Exp.\ I \S~3 that the full subcategory \(\Db(X)_{\mathrm{parf}}\) of \(\Db(X)\) formed by perfect complexes is, like \(\Db(X)_{\mathrm{coh}}\), a triangulated subcategory, ``stable'' under derived inverse image and tensor product. We evidently have \(\Db(X)_{\mathrm{parf}}\subset\Db(X)_{\mathrm{coh}}\), but the inclusion is strict in general. Pseudo-coherence can be redefined ``by passage to the limit'' from perfection: a complex \(F\) is pseudo-coherent precisely when, for every \(n\in\mathbb Z\), \(F\) can be ``locally approximated to order \(n\) by a perfect complex'', by which we mean that there locally exists a distinguished triangle
\[
\begin{tikzcd}
& R \arrow[dl] &\\
L \arrow[rr] && F \arrow[ul] ,
\end{tikzcd}
\]
where \(L\) is perfect and \(R\) acyclic in degree \(\geq n\). Conversely, perfection is recovered from pseudo-coherence through an additional regularity condition: a complex is perfect precisely when it is pseudo-coherent and locally of finite tor-dimension (Exp.\ I \S~5). We thus recover the fact that, when the local rings are regular, every coherent sheaf is perfect, and, more generally, every pseudo-coherent complex with locally bounded cohomology is perfect.


%%%%%%%%%%  scan idx 94  |  folio 81  %%%%%%%%%%
\SGAPrintPage{81}%
Pseudo-coherence and perfection are the two fundamental notions of finiteness with which we shall work in this Seminar. Sections~1 to 5 of Exposé I are devoted to their definition and to establishing their elementary stability properties. For this, we use only two or three local properties of the category of free Modules of finite type with respect to the category of all Modules. It therefore proves convenient---and useful, particularly in Exp.\ II---to axiomatize the situation by introducing a notion of pseudo-coherence (resp.\ perfection) in a category fibered over a site relative to a fibered subcategory satisfying suitable conditions.

Sections~6 to 8 of Exp.\ I generalize to perfect complexes a number of familiar notions for locally free sheaves of finite type: rank, duality, trace of an endomorphism.

In Exp.\ II, we examine the problem of ``the existence of global resolutions'': under what conditions, on a ringed topos \((X,\underline{O}_X)\), is a pseudo-coherent (resp.\ perfect) complex \uline{globally} isomorphic in \(\Db(X)\) to a complex of locally free Modules of finite type? The importance of this question lies essentially in the fact that, at present, we do not know how to generalize to perfect complexes certain usual constructions on locally free Modules: exterior and symmetric powers, Chern classes (for more details on this subject, see Exp.\ XIV). This is why it is convenient to have workable sufficient criteria for the existence of global resolutions, which allow us to reduce certain questions about pseudo-coherent (resp.\ perfect) complexes to analogous questions about locally free sheaves of finite type. Such criteria
%%%%%%%%%%  scan idx 95  |  folio 82  %%%%%%%%%%
\SGAPrintPage{82}%
are given in Exp.\ II. In particular, it is shown that global resolutions exist in the following cases: \(X\) is the Zariski topos of an affine scheme, or more generally of a quasi-compact scheme with an ample invertible sheaf, or even an ample family of invertible sheaves in Kleiman's sense (for example a regular scheme); or again \(X\) is the topos of sheaves of sets on a compact topological space, \(\underline{O}_X\) the sheaf of continuous complex-valued functions. As an illustration---and to bring out the flexibility of the notions introduced---the appendix gives a ``purely sheaf-theoretic'' definition of the index of a family of elliptic operators.

Exposé III studies the stability of finiteness conditions under derived direct image. To obtain reasonable statements, we must ``relativize'' pseudo-coherence and perfection. We work here in the setting of ordinary schemes, which suffices for the Seminar, but there is no doubt that sooner or later an analogous theory will have to be developed for relative analytic schemes or spaces. Let \(p:X\to S\) be an \(S\)-scheme locally of finite type and \(F\) a complex of \(\underline{O}_X\)-Modules; we say that \(F\) is \uline{pseudo-coherent} (\uline{resp.\ perfect}) \uline{relative to} \(p\) when we can locally embed (by a closed immersion) \(X\) in a smooth \(S\)-scheme \(X'\) so that \(F\), extended by zero to \(X'\), is pseudo-coherent (resp.\ perfect). Pseudo-coherence relative to \(p\) is especially interesting when \(S\) is not locally noetherian, since in the opposite case it coincides with ordinary pseudo-coherence. Moreover, as one might expect, saying that \(F\) is perfect rel.\ to \(p\)
%%%%%%%%%%  scan idx 96  |  folio 83  %%%%%%%%%%
\SGAPrintPage{83}%
amounts to saying that \(F\) is pseudo-coherent rel.\ to \(p\) and locally of finite tor-dimension rel.\ to \(p\) (i.e.\ rel.\ to the sheaf of rings \(p^{-1}(\underline{O}_S)\)). The central theorem of Exp.\ III is the \uline{finiteness theorem}, which asserts (in a slightly more precise form) that, for a \uline{proper} morphism \(f:X\to Y\) of \(S\)-schemes locally of finite type, the functor \(\mathbb{R}f_*\) sends complexes pseudo-coherent rel.\ to \(S\) to complexes pseudo-coherent rel.\ to \(S\); this theorem is in fact a conjecture, but is nevertheless proved in the following two special cases: a)~\(S\) is locally noetherian b)~\(f\) is projective; unfortunately, it seems that extending it to the general case is of the same order of difficulty as the analogous theorem in Analytic Geometry (Grauert's ``theorem''). Combining the finiteness theorem with an essentially trivial formula called the \uline{projection formula}, we obtain workable criteria for the stability of relative perfection under direct image. As a corollary, we recover ``Grauert's continuity and semicontinuity theorems'' (EGA III 7.6).

Exposé IV translates the results of Exposés I to III into the language of ``Grothendieck groups''. On a ringed topos \((X,\underline{O}_X)\), we denote by \(K^\bullet(X)\) (resp.\ \(K_\bullet(X)\)) the ``Grothendieck group'' of the category of complexes that are perfect and of finite tor-dimension (resp.\ pseudo-coherent and with bounded cohomology). The group \(K^\bullet(X)\) is a ring, and \(K_\bullet(X)\) a module over \(K^\bullet(X)\). \(K^\bullet\) is a contravariant functor (as the raised dot suggests), while \(K_\bullet\) is a covariant functor for proper pseudo-coherent morphisms (of schemes or analytic spaces, and subject to having the finiteness theorem\dots). There are also unusual variances,
%%%%%%%%%%  scan idx 97  |  folio 84  %%%%%%%%%%
\SGAPrintPage{84}%
covariance of \(K^\bullet\), contravariance of \(K_\bullet\), for morphisms satisfying suitable regularity hypotheses. Finally, the globalization criteria of Exp.\ II allow us to relate these to the ``Grothendieck groups'' defined naively from coherent sheaves or locally free sheaves of finite type.

\sgaStableTarget{sga6.e01.sec.1}\section*{1. Preliminary definitions}

\sgaStableTarget{sga6.e01.sec.1.1}\subsection*{1.1. Fibered categories with additive (resp. abelian, resp. triangulated) fibres}

\sgaStableTarget{sga6.e01.para.1-1-1}1.1.1. Let \(\underline{S}\) be a category, and \(\underline{C}\) a fibered \(\underline{S}\)-category (SGA 1 VI 6.1) whose fibres are additive (resp. triangulated) categories. We shall say that \(\underline{C}\) is an \uline{additive} \(\underline{S}\)\uline{-category} (\uline{resp.\ triangulated}) (or that \(\underline{C}\) is additive (resp.\ triangulated) over \(\underline{S}\)) when, for every arrow \(f:X\to Y\) of \(\underline{S}\), the inverse-image functor (determined up to unique isomorphism) \(f^*:\underline{C}_Y\longrightarrow \underline{C}_X\) is additive (resp.\ exact). We shall say that \(\underline{C}\) is an \uline{abelian} \(\underline{S}\)-category when \(\underline{C}\) is an additive \(\underline{S}\)-category with abelian fibres, and that \(\underline{C}\) is a \uline{flat} abelian \(\underline{S}\)-category when, moreover, the inverse-image functors are exact.

\sgaStableTarget{sga6.e01.para.1-1-2}1.1.2. Let \(\underline{C}\) be a triangulated \(\underline{S}\)-category. We define, for example using a cleavage of \(\underline{C}\) (SGA 1 VI 7.1), a fibered \(\underline{S}\)-category \(\Trc(\underline{C})\) whose fibre at each \(X\in \ob(\underline{S})\) is the category \(\Trc(\underline{C}_X)\) of distinguished triangles of \(\underline{C}_X\).

\sgaStableTarget{sga6.e01.para.1-1-3}1.1.3. Let \(\underline{C}\) and \(\underline{C}'\) be additive (resp. triangulated) \(\underline{S}\)-categories and \(F:\underline{C}\to\underline{C}'\) a functor. We call \(F\) an \uline{additive} (resp.\ \uline{exact}) \(\underline{S}\)\uline{-functor} when \(F\) is a cartesian \(\underline{S}\)-functor that is additive (resp.\ exact) on each fibre.
% Additional macro definitions
\providecommand{\cC}{\mathcal C}
\providecommand{\cD}{\mathcal D}
\providecommand{\cK}{\mathcal K}
\providecommand{\cO}{\mathcal O}
\providecommand{\cS}{\mathcal S}
\providecommand{\Cpx}{\mathrm C}
\providecommand{\Kpx}{\mathrm K}
\providecommand{\Dpx}{\mathrm D}
\providecommand{\D}{\mathrm D}
\providecommand{\Ho}{\mathrm H}
\providecommand{\RHom}{\mathrm R\!\operatorname{Hom}}
\providecommand{\Ob}{\operatorname{ob}}
\providecommand{\Coker}{\operatorname{Coker}}
\providecommand{\Ker}{\operatorname{Ker}}
\renewcommand{\Im}{\operatorname{Im}}
\providecommand{\Hom}{\operatorname{Hom}}
\providecommand{\Cone}{\operatorname{Cone}}
\providecommand{\id}{\operatorname{Id}}
\providecommand{\qis}{\xrightarrow{\sim}}
%%%%%%%%%%  scan idx 98  |  folio 85  %%%%%%%%%%
\SGAPrintPage{85}%
\sgaStableTarget{sga6.e01.para.1-1-3.occ2}When \(\underline{C}\) and \(\underline{C}'\) are triangulated \(\underline{S}\)-categories and \(F\) is an exact \(\underline{S}\)-functor, \(F\) defines a cartesian \(\underline{S}\)-functor \(\Trc(\underline{C})\longrightarrow \Trc(\underline{C}')\).

\sgaStableTarget{sga6.e01.para.1-1-4}1.1.4. Let \(\underline{C}\) be a triangulated \(\underline{S}\)-category, and \(\underline{C}'\) a full fibered sub-\(\underline{S}\)-category of \(\underline{C}\). We call \(\underline{C}'\) a \uline{triangulated sub-}\(\underline{S}\)\uline{-category} of \(\underline{C}\) when, for every \(X\in\Ob \underline{S}\), \(\underline{C}'_X\) is a triangulated subcategory of \(\underline{C}_X\) ([V] chap.\ I \S~1 2.3).

\sgaStableTarget{sga6.e01.para.1-1-5}1.1.5. Let \(\underline{C}\) be an additive \(\underline{S}\)-category. We define (SGA 4 XVII 2.2, 2.4) an additive (resp.\ triangulated) \(\underline{S}\)-category \(\Cpx(\underline{C})\), (resp.\ \(\Kpx(\underline{C})\)), whose fibre at \(X\in\Ob \underline{S}\) is the category \(\Cpx(\underline{C}_X)\) (resp.\ \(\Kpx(\underline{C}_X)\)) of complexes\SGAAuthorFootnote{1}{with differential of degree \(+1\).} (resp.\ complexes ``up to homotopy'') of \(\underline{C}_X\). When \(\underline{C}\) is a flat abelian \(\underline{S}\)-category, we define (loc.\ cit.) a triangulated \(\underline{S}\)-category \(\Dpx(\underline{C})\), whose fibre at \(X\in\Ob \underline{S}\) is the derived category \(\Dpx(\underline{C}_X)\); for an arrow \(f:X\to Y\) of \(\underline{S}\), we shall denote by \(f^*:\Dpx(\underline{C}_Y)\longrightarrow \Dpx(\underline{C}_X)\) the inverse-image functor, obtained from \(f^*:\Kpx(\underline{C}_Y)\longrightarrow \Kpx(\underline{C}_X)\) by simple passage to the quotient\SGAAuthorFootnote{2}{for \(a\in\Ob \Dpx(\underline{C}_Y)\) (resp.\ \(\mathrm{Fl}\,\Dpx(\underline{C}_Y)\)), we shall often write, in abbreviated form, \(a_X\) instead of \(f^*(a)\).}. We define full triangulated sub-\(\underline{S}\)-categories \(\Kpx^*(\underline{C})\) (resp.\ \(\Dpx^*(\underline{C})\)) \((*=+,-,b)\) by the usual degree conditions.

\sgaStableTarget{sga6.e01.sec.1.2}1.2. \uline{Terminology.} Let \(\underline{S}\) be a site, \(\underline{C}\) a flat abelian \(\underline{S}\)-category (1.1), and \(\underline{C}_0\) a strictly full fibered sub-\(\underline{S}\)-category.

Let \(F\in\Ob\underline{C}_X\), for \(X\in\Ob \underline{S}\). The set of \(f:U\to X\) in \(\underline{S}_{/X}\) such that there exists an epimorphism \(E\longrightarrow f^*F\) with \(E\in\Ob\underline{C}_{0U}\) is a sieve on \(X\). When it is covering, we shall say that \(F\) is of \(\underline{C}_0\)-\uline{finite type} (or simply of finite type when there is no risk of confusion).


%%%%%%%%%%  scan idx 99  |  folio 86  %%%%%%%%%%
\SGAPrintPage{86}%
We shall say that \(\underline{C}_0\) is \uline{locally stable under kernels of epimorphisms} when the following condition is satisfied: for every \(X\in\Ob \underline{S}\) and every epimorphism \(u:E\longrightarrow F\) of \(\underline{C}_X\) with \(E\) and \(F\in\Ob\underline{C}_{0X}\), the kernel of \(u\) is locally in \(\underline{C}_0\) (i.e.\ the set of \(f:U\to X\) in \(\underline{S}_{/X}\) such that \(f^*\Ker(u)\in\Ob\underline{C}_{0U}\) is a refinement of \(X\)).

By the standard procedure of decomposing a long exact sequence into short exact sequences, we see that the preceding condition is equivalent to the following (apparently stronger) condition: for every \(X\in\Ob \underline{S}\) and every exact sequence of \(\underline{C}_X\)
\sgaStableTarget{sga6.e01.eq.unnamed}\begin{equation}\tag{*}
  0\longrightarrow E^0\longrightarrow E^1\longrightarrow\ldots\longrightarrow E^n\longrightarrow 0
\end{equation}
with \(E^i\in\Ob\underline{C}_{0X}\) for \(1\leq i\leq n\), \(E^0\) is locally in \(\underline{C}_0\).

We shall say that \(\underline{C}_0\) is \uline{locally quasi-stable under kernels of epimorphisms} when, for every exact sequence such as \((*)\), \(E^0\) is of \(\underline{C}_0\)-finite type.

On the other hand, we shall say that \(\underline{C}_0\) is \uline{locally liftable in} \(\underline{C}\) when, for every \(X\in\Ob \underline{S}\), every epimorphism \(u:E\longrightarrow F\) of \(\underline{C}_X\) with \(F\in\Ob\underline{C}_{0X}\) locally admits a section. For \(\underline{C}_0\) to be locally liftable in \(\underline{C}\), it is necessary and sufficient that, for every diagram of \(\underline{C}_X\) (\(X\in\Ob \underline{S}\))
\[
\begin{tikzcd}
& H \arrow[d]\\
E \arrow[r,"u"] & F
\end{tikzcd}
\]
where \(u\) is an epimorphism and \(H\in\Ob\underline{C}_0\), the set of \(U\) over \(X\) such that there exists a commutative diagram
\[
\begin{tikzcd}
& H_U \arrow[d] \arrow[dl]\\
E_U \arrow[r,"u_U"] & F_U
\end{tikzcd}
\]
is a refinement of \(X\).

%%%%%%%%%%  scan idx 100  |  folio 87  %%%%%%%%%%
\SGAPrintPage{87}%
We shall say that \(\underline{C}_0\) is \uline{locally quasi-liftable in} \(\underline{C}\) when the following weaker condition is satisfied: for every \(X\in\Ob \underline{S}\), and every epimorphism \(u:E\longrightarrow F\) of \(\underline{C}_X\) with \(F\in\Ob\underline{C}_{0X}\), the set of \(f:U\to X\) in \(\underline{S}_{/X}\) such that there exists \(v:F'\longrightarrow f^*E\) with \(F'\in\Ob\underline{C}_{0U}\) such that \(f^*(u)v\) is an epimorphism is a refinement of \(X\).

This condition is equivalent to the following: for every diagram of \(\underline{C}_X\)
\[
\begin{tikzcd}
& H \arrow[d]\\
E \arrow[r,"u"] & F ,
\end{tikzcd}
\]
where \(u\) is an epimorphism, and \(H\in\Ob\underline{C}_0\), the set of \(U\) over \(X\) such that there exists a commutative diagram
\[
\begin{tikzcd}
G \arrow[r,"v"] \arrow[d] & H_U \arrow[d]\\
E_U \arrow[r,"u_U"] & F_U ,
\end{tikzcd}
\]
with \(G\in\Ob\underline{C}_0\) and \(v\) an epimorphism is a refinement of \(X\).

When \(\underline{S}\) is a chaotic site, for example the punctual site (one object, one morphism), we shall omit the adverb ``locally'' in the preceding definitions.

\sgaStableTarget{sga6.e01.sec.1.3}\subsection*{1.3. Examples}

a) Let \(\underline{S}\) be a punctual site, \(A\) a ring, and \(\underline{C}\) the category of left \(A\)-modules. The full subcategory \(\underline{C}_0\) of \(\underline{C}\) formed by the projective (left) \(A\)-modules of finite type is stable under kernels of epimorphisms and liftable in \(\underline{C}\); an \(A\)-module is of \(\underline{C}_0\)-finite type if and only if it is of finite type in the
%%%%%%%%%%  scan idx 101  |  folio 88  %%%%%%%%%%
\SGAPrintPage{88}%
ordinary sense. The full subcategory \(\underline{C}_1\) of \(\underline{C}\) formed by the \(A\)-modules of finite type is quasi-liftable in \(\underline{C}\) (but not liftable in general); when \(A\) is noetherian, \(\underline{C}_1\) is stable under kernels of epimorphisms. The full subcategory \(\underline{C}'_0\) of \(\underline{C}\) formed by the free \(A\)-modules of finite type is liftable in \(\underline{C}\) and quasi-stable (not stable in general) under kernels of epimorphisms.

b) Let \((S,\underline{O}_S)\) be a space ringed by local rings, \(\underline{S}\) the site of open subsets of \(S\), and \(\underline{C}\) the fibered \(\underline{S}\)-category (even split (SGA 1 VI 9)) defined by the functor \(U\longmapsto\) (category of left \(\underline{O}_U\)-Modules); \(\underline{C}\) is a flat abelian \(\underline{S}\)-category (1.1). Let \(\underline{C}_0\) be the strictly full additive sub-\(\underline{S}\)-category of \(\underline{C}\) defined by \(U\longmapsto\) (category of free \(\underline{O}_U\)-Modules of finite type). An \(\underline{O}_S\)-Module \(M\) is locally of \(\underline{C}_0\)-finite type if and only if \(M\) is of finite type in the ordinary sense (EGA \(0_{\mathrm I}\) 5.2). It follows from (EGA \(0_{\mathrm I}\) 5.2.7 and 5.4.9) that \(\underline{C}_0\) is locally stable under kernels of epimorphisms and locally liftable in \(\underline{C}\).

c) Let \(G\) be a group, \(\underline{S}\) the topos of \(G\)-sets (SGA 4 V 3.11), \(k\) a field (regarded as a constant Ring of \(\underline{S}\)), and \(\underline{C}\) the split \(\underline{S}\)-category defined by the functor \(U\longmapsto\) (category of left \(k_U\)-Modules); \(\underline{C}\) is a flat abelian \(\underline{S}\)-category. Let \(\underline{C}_0\) be the strictly full additive sub-\(\underline{S}\)-category of \(\underline{C}\) defined by \(U\longmapsto\) (category of free \(k_U\)-Modules of finite type).
Since localization in \(\underline{S}\) is equivalent to forgetting the operations of \(G\), it is clear that a \(G\)-\(k\)-module \(M\) is locally of \(\underline{C}_0\)-finite type if and only if \(M\) is finite-dimensional as a vector space over \(k\), and that \(\underline{C}_0\) is locally stable under kernels of epimorphisms and locally liftable in \(\underline{C}\).

The reader is left to examine the ``Galois'' variant of c) in which \(G\) is a profinite group and \(\underline{S}\) the topos of sets on which \(G\) operates continuously.

%%%%%%%%%%  scan idx 102  |  folio 89  %%%%%%%%%%
\SGAPrintPage{89}%
In fact, examples b) and c) are special cases of the following example:

d) Let \((\underline{S},\underline{O}_{\underline{S}})\) be a ringed topos (cf.\ SGA 4 IV 2.10)\SGAAuthorFootnote{*}{Here \(\underline{O}_{\underline{S}}\) is a ring of \(\underline{S}\), with identity element, but not necessarily commutative.}, and \(\underline{C}\) the split \(\underline{S}\)-category defined by \(U\longmapsto\) (category of left \(\underline{O}_U\)-Modules); \(\underline{C}\) is a flat abelian \(\underline{S}\)-category. Let \(\underline{C}_0\) be the strictly full additive sub-\(\underline{S}\)-category of \(\underline{C}\) defined by \(U\longmapsto\) (category of free \(\underline{O}_U\)-Modules of finite type). An \(\underline{O}_{\underline{S}}\)-Module locally of \(\underline{C}_0\)-finite type will simply be said to be of \uline{finite type}. The category \(\underline{C}_0\) is locally quasi-stable under kernels of epimorphisms and locally liftable in \(\underline{C}\), as follows from

\sgaStableTarget{sga6.e01.lem.1-3-1}\textbf{Lemma 1.3.1.} \uline{Let} \(P\) \uline{be an} \(\underline{O}_{\underline{S}}\)\uline{-Module. Equivalent conditions}:
\begin{enumerate}[label=(\roman*)]
\item \(P\) \uline{is of finite type and the functor} \(\underline{\Hom}(P,.)\) \uline{is exact};
\item \(P\) \uline{is of finite type and every epimorphism} \(M\to P\) \uline{locally admits a section};
\item \(P\) \uline{is locally a direct factor of a free Module of finite type}.
\end{enumerate}
The proof is immediate and is left as an exercise to the reader.

\sgaStableTarget{sga6.e01.sec.1.4}\subsection*{1.4. An ``extension'' lemma (cf. EGA \(0_{\mathrm{III}}\) 11.9.1)}

\sgaStableTarget{sga6.e01.prop.1-4-1}\textbf{Proposition 1.4.1.} \uline{Let} \(\underline{S}\) \uline{be a site,} \(\underline{C}\) \uline{a flat abelian} \(\underline{S}\)\uline{-category (1.1),} \(\underline{C}_0\) \uline{a strictly full additive sub-}\(\underline{S}\)\uline{-category (1.1) of} \(\underline{C}\)\uline{; assume that} \(\underline{C}_0\) \uline{is locally quasi-liftable in} \(\underline{C}\) \uline{(1.2). Let} \(S\in\Ob \underline{S}\)\uline{,} \(u:E\longrightarrow F\) \uline{be a morphism (of degree} \(0\)\uline{) of complexes of} \(\underline{C}_S\) \uline{(differentials of degree} \(+1\)\uline{),} \(G\) \uline{the cone of} \(u\) \uline{(cf.\ [V] chap.\ I \S~1 2.1),} \(n\in Z\)\uline{; assume that} \(\Ho^n(G)\) \uline{is}
%%%%%%%%%%  scan idx 103  |  folio 90  %%%%%%%%%%
\SGAPrintPage{90}%
\uline{of} \(\underline{C}_0\)-%
\uline{finite type (1.2). Then the set of} \(X\in\Ob \underline{S}_{/S}\) \uline{satisfying property} \((P)\) \uline{below is a refinement of} \(S\).

\((P)\): there exist an object \(M\) of \(\underline{C}_{0X}\) and arrows of \(\underline{C}_X\)
\[
  r:M\longrightarrow Z^{n+1}(E_X)=
  \Ker\!\left(E_X^{n+1}\xrightarrow{d^{n+1}}E_X^{n+2}\right)
\]
\[
  \text{and}\quad s:M\longrightarrow F_X^n
\]
such that the following diagram of \(\underline{C}_X\)
\[
\begin{tikzcd}[column sep=tiny]
\ldots \arrow[r] & E_X^{n-1}
 \arrow[r,"{\binom{d_{E_X}^{n-1}}{0}}"]
 \arrow[d,"u_X^{n-1}"] &
E_X^n\oplus M
 \arrow[r,"{(d_{E_X}^n,r)}"]
 \arrow[d,"{(u_X^n,s)}"] &
E_X^{n+1} \arrow[r] \arrow[d,"u_X^{n+1}"] &
\ldots\\
\ldots \arrow[r] & F_X^{n-1} \arrow[r] &
F_X^n \arrow[r] &
F_X^{n+1} \arrow[r] &
\ldots
\end{tikzcd}
\]
is a morphism of complexes \(u':E'\to F_X\) (i.e.\ \(d_{F_X}^n s=u_X^{n+1}r\)), and that
\begin{enumerate}[label=(\roman*)]
\item \(\Ho^{n+1}(u')\) is a monomorphism, with \(\Im\Ho^{n+1}(u')=\Im\Ho^{n+1}(u_X)\),
\item \(\Ho^n(u')\) is an epimorphism.
\end{enumerate}

In pictorial terms, we may say that, by modifying \(E\) and \(u\) locally, \uline{in degree \(n\) only}, by adding to \(E^n\) an object of \(\underline{C}_0\), we can locally make \(\Ho^{n+1}(u)\) injective, without changing its image, and \(\Ho^n(u)\) surjective.

\uline{Proof}. The procedure is as follows: first we show that we can choose \(r\) and \(s\) locally so as to make \(\Ho^{n+1}(u)\) injective; we verify 
%%%%%%%%%%  scan idx 104  |  folio 91  %%%%%%%%%%
\SGAPrintPage{91}%
that the finiteness hypothesis on \(\Ho^n(G)\) is preserved; a new local correction in degree \(n\) then allows us to make \(\Ho^n(u)\) surjective.

The exact cohomology sequence
\[
  \longrightarrow \Ho^n(G)\longrightarrow
  \Ho^{n+1}(E)\longrightarrow
  \Ho^{n+1}(F)\longrightarrow
\]
shows that the kernel of \(\Ho^{n+1}(u)\) is of \(\underline{C}_0\)-finite type. In other words, the set of \(X\) over \(S\) such that there exists an exact sequence
\[
  M\longrightarrow \Ho^{n+1}(E_X)
  \xrightarrow{\Ho^{n+1}(u)}
  \Ho^{n+1}(F_X),
\]
with \(M\in\Ob\underline{C}_{0X}\), is a refinement of \(S\). Since \(\underline{C}_0\) is locally quasi-liftable in \(\underline{C}\), we can moreover require the arrow \(M\longrightarrow \Ho^{n+1}(E_X)\) to factor through \(Z^{n+1}(E_X)\) as
\[
  M\xrightarrow{r}Z^{n+1}(E_X)\longrightarrow \Ho^{n+1}(E_X).
\]
The composite arrow
\[
  M\xrightarrow{u_X^{n+1}r}Z^{n+1}(F_X)\longrightarrow \Ho^{n+1}(F_X)
\]
is zero by construction, so \(u_X^{n+1}r\) factors through \(B^{n+1}(F_X)=\operatorname{Im}\!\left(F_X^n\xrightarrow{d^n}F_X^{n+1}\right)\). Applying quasi-liftability again, we can require \(u_X^{n+1}r\) to factor through \(F_X^n\) as
\[
  M\xrightarrow{s}F_X^n\longrightarrow B^{n+1}(F_X)\longrightarrow Z^{n+1}(F_X).
\]
We then check trivially that \(u'\), defined as in (1.4.1), is a morphism of complexes satisfying (i).

Let \(G'\) be the cone of \(u':E'\to F_X\). Let us show that \(\Ho^n(G')\) is of \(\underline{C}_0\)-finite type. By construction, there is an exact sequence
\[
  0\longrightarrow E_X\xrightarrow{v}E'\longrightarrow M[-n]\longrightarrow 0,
\]
%%%%%%%%%%  scan idx 105  |  folio 92  %%%%%%%%%%
\SGAPrintPage{92}%
with \(u'v=u_X\). To compare \(G_X\) and \(G'\), we can, for example, form the ``octahedron'' of \(\Kpx(\underline{C}_X)\)
\[
\begin{tikzpicture}[>=stealth,baseline=(GX.base),node distance=1.0cm]
  \node (EX) at (0,0) {$E_X$};
  \node (Ep) at (2.6,0) {$E'$};
  \node (FX) at (5.2,0) {$F_X$};
  \node (M)  at (1.35,1.25) {$M[-n]$};
  \node (GX) at (2.65,2.55) {$G_X$};
  \node (Gp) at (3.95,1.25) {$G'$};
  \node at (3.22,1.72) {$(*)$};
  \draw[->] (EX) -- node[below] {$v$} (Ep);
  \draw[->] (Ep) -- node[below] {$u'$} (FX);
  \draw[->] (Ep) -- (M);
  \draw[->] (M) -- (GX);
  \draw[->] (GX) -- (Gp);
  \draw[->] (FX) -- (Gp);
  \draw[->] (M) to[bend left=38] node[left] {$+1$} (EX);
  \draw[->] (Gp) to[bend left=28] node[above] {$+1$} (M);
  \draw[->] (Gp) to[bend left=28] node[right] {$+1$} (Ep);
  \draw[->] (GX) to[bend left=55] node[left] {$+1$} (EX);
  \draw[->] (FX) to[bend right=42] (GX);
  \node[anchor=base] at (6.3,0 |- FX.base) {.};
\end{tikzpicture}
\]
The exact cohomology sequence of the distinguished triangle \((*)\) shows that \(\Ho^n(G_X)\longrightarrow \Ho^n(G')\) is an epimorphism, hence that \(\Ho^n(G')\) is indeed of \(\underline{C}_0\)-finite type.

The preceding correction therefore reduces us to the case where \(\Ho^{n+1}(u)\) is a monomorphism, which we shall now assume. The exact cohomology sequence of \(E\to F\to G\) then gives an exact sequence
\begin{equation}\tag{.}
  \Ho^n(E)\xrightarrow{\Ho^n(u)}
  \Ho^n(F)\longrightarrow \Ho^n(G)\longrightarrow 0.
\end{equation}
Since \(\Ho^n(G)\) is of \(\underline{C}_0\)-finite type, the set of \(X\) over \(S\) such that there exists an epimorphism \(M\longrightarrow \Ho^n(G_X)\) with \(M\in\Ob\underline{C}_{0X}\) is, by definition, a refinement of \(S\). Using the quasi-liftability of \(\underline{C}_0\) in \(\underline{C}\), we can require this epimorphism to factor through \(Z^n(F_X)\) as
\[
  M\xrightarrow{s}Z^n(F_X)\longrightarrow \Ho^n(F_X)\longrightarrow \Ho^n(G_X).
\]
It is clear that \(u'\), defined as in (1.4.1) with \((r=0)\), is a morphism of complexes and that \(\Ho^{n+1}(u')=\Ho^{n+1}(u_X)\). On the other hand, it follows immediately from the exact sequence (.) that
%%%%%%%%%%  scan idx 106  |  folio 93  %%%%%%%%%%
\SGAPrintPage{93}%
\[
  \Ho^n(u'):\Ho^n(E_X)\oplus M\longrightarrow \Ho^n(F_X)
\]
is an epimorphism. This completes the proof of (1.4.1).

\sgaStableTarget{sga6.e01.lem.1-4-2}\textbf{Lemma 1.4.2.} \uline{Let} \(\underline{C}\) \uline{be an abelian category,} \(u:E\to F\) \uline{a morphism of complexes of} \(\underline{C}\) \uline{(resp.\ an arrow of} \(\Dpx(\underline{C})\)\uline{),} \(G\) \uline{the cone of} \(u\)\uline{. Equivalent conditions for a given} \(n\in\mathbb Z\):
\begin{enumerate}[label=(\roman*)]
\item \(\Ho^i(G)=0\) \uline{for} \(i\geq n\);
\item \(\Ho^i(u)\) \uline{is an isomorphism for} \(i>n\) \uline{and an epimorphism for} \(i=n\).
\end{enumerate}
The proof follows immediately from the exact cohomology sequence.

\sgaStableTarget{sga6.e01.def.1-4-3}\textbf{Definition 1.4.3.} \uline{We shall say that} \(u\) \uline{is an} \(n\)-quasi-isomorphism (\uline{resp.\ an} \(n\)-isomorphism) \uline{when} \(u\) \uline{satisfies the equivalent conditions of (1.4.2).}

That being so, we deduce at once from (1.4.1) the

\sgaStableTarget{sga6.e01.cor.1-4-4}\textbf{Corollary 1.4.4.} \uline{Let} \(\underline{S},S,\underline{C},\underline{C}_0\) \uline{be as in (1.4.1). Let} \(u:E\to F\) \uline{be an} \((n+1)\)\uline{-quasi-isomorphism of complexes of} \(\underline{C}_S\) \uline{(}\(n\in\mathbb Z\) \uline{given) such that} \(\Ho^n(G)\) \uline{is of} \(\underline{C}_0\)\uline{-finite type. Then the set of} \(X\) \uline{over} \(S\) \uline{such that there exist} \((M,r,s)\) \uline{as in (1.4.1) such that} \(u'\) \uline{is an} \(n\)\uline{-quasi-isomorphism is a refinement of} \(S\)\uline{.}

In pictorial terms, using the finiteness condition on \(\Ho^n(G)\), we can locally ``gain one notch'', i.e.\ modify \(u\) locally in degree \(n\) by adding to \(E^n\) an object of \(\underline{C}_0\) so as to make \(u\) an \(n\)-quasi-isomorphism.

\sgaStableTarget{sga6.e01.cor.1-4-5}\textbf{Corollary 1.4.5.} \uline{Let} \(\underline{C}\) \uline{be an abelian category,} \(n\in\mathbb Z\)\uline{, and} \(u:E\to F\) \uline{an} \(n\)\uline{-quasi-isomorphism of complexes of} \(\underline{C}\)\uline{. There exists a commutative diagram of} \(\Cpx(\underline{C})\)
\[
\begin{tikzcd}
E \arrow[r,"v"] \arrow[d,"u"'] & E' \arrow[dl,"u'"]\\
F & ,
\end{tikzcd}
\]
%%%%%%%%%%  scan idx 107  |  folio 94  %%%%%%%%%%
\SGAPrintPage{94}%
\uline{where} \(u'\) \uline{is a quasi-isomorphism and} \(v\) \uline{a monomorphism inducing the identity in degree} \(\geq n\) (i.e.\ such that \(E'^i=E^i\) and \(v^i=\id(E^i)\) for \(i\geq n\)).

\uline{Proof.} Applying (1.4.4) in the case where \(\underline{S}\) is the punctual site and \(\underline{C}_0=\underline{C}\), we obtain a commutative triangle of \(\Cpx(\underline{C})\)
\[
\begin{tikzcd}
E \arrow[r,"v_{n-1}"] \arrow[d,"u"'] & E'_{\,n-1} \arrow[dl,"u'_{n-1}"]\\
F & ,
\end{tikzcd}
\]
where \(u'_{n-1}\) is an \((n-1)\)-quasi-isomorphism and \(v_{n-1}\) a monomorphism inducing the identity in degree \(\geq n\). Iterating the procedure, we find, for \(i\leq n-1\), a sequence of commutative triangles of \(\Cpx(\underline{C})\)
\[
\begin{tikzcd}
E'_i \arrow[r,"v_{i-1}"] \arrow[d,"u'_i"'] & E'_{\,i-1} \arrow[dl,"u'_{i-1}"]\\
F & ,
\end{tikzcd}
\]
where \(u'_i\) is an \(i\)-quasi-isomorphism and \(v_{i-1}\) a monomorphism inducing the identity in degree \(\geq i\). The corollary follows by passage to the limit. More precisely, denoting by \(d'_i\) the differential of \(E'_i\), we define \(E'\) by
\[
  E'^i=E^i,\quad d'^i=d^i\quad\text{for }i\geq n,
\]
and
\[
  E'^i=(E'_i)^i,\quad d'^i=(d'_i)^i\quad\text{for }i<n,
\]
then \(v\) and \(u'\) by
\[
  v^j=(v_i v_{i+1}\ldots v_{n-1})^j,\qquad u'^j=(u'_i)^j\qquad\text{for }i\leq j,
\]
and we check at once that \(E'\), \(v\), \(u'\) so defined answer the question.

%%%%%%%%%%  scan idx 108  |  folio 95  %%%%%%%%%%
\SGAPrintPage{95}%
In other words, briefly, we can extend an \(n\)-quasi-isomorphism to a quasi-isomorphism ``without changing anything given in degree \(\geq n\)''.

\sgaStableTarget{sga6.e01.sec.2}\section*{2. Pseudo-coherent complexes}

\sgaStableTarget{sga6.e01.para.2-0}2.0. In this section, \(\underline{S}\) denotes a site, \(S\) an object of \(\underline{S}\), \(\underline{C}\) a flat abelian \(\underline{S}\)-category (1.1), and \(\underline{C}_0\) a strictly full additive sub-\(\underline{S}\)-category of \(\underline{C}\). We assume that \(\underline{C}_0\) is locally quasi-stable under kernels of epimorphisms (1.2) and locally quasi-liftable in \(\underline{C}\) (1.2).

\sgaStableTarget{sga6.e01.def.2-1}\textbf{Definition 2.1.} \uline{Let} \(X\in\Ob \underline{S}\)\uline{,} \(E\) \uline{be a complex of} \(\underline{C}_X\)\uline{, and} \(n\in\mathbb Z\)\uline{. We shall say that} \(E\) \uline{is} strictly \(\underline{C}_0\)-\(n\)-pseudo-coherent (resp.\ \(\underline{C}_0\)-pseudo-coherent, resp.\ \(\underline{C}_0\)-perfect) \uline{when} \(E^i=0\) \uline{for} \(i\) \uline{sufficiently large and} \(E^i\in\Ob\underline{C}_{0X}\) \uline{for} \(i\geq n\) \uline{(resp.\ for every} \(i\)\uline{, resp.\ for every} \(i\) \uline{and} \(E^i=0\) \uline{for almost every} \(i\)\uline{).}

When there is no risk of confusion, we shall sometimes say ``strictly \(n\)-pseudo-coherent (resp\dots)'' instead of ``strictly \(\underline{C}_0\)-\(n\)-pseudo-coherent (resp\dots)''.

\sgaStableTarget{sga6.e01.lem.2-1-1}\textbf{Lemma 2.1.1.} \uline{Let} \(E\in\Ob\Cpx(\underline{C}_S)\) \uline{and} \(n\in\mathbb Z\)\uline{. When} \(E\) \uline{is strictly} \(n\)\uline{-pseudo-coherent and} \(\Ho^i(E)=0\) \uline{for} \(i>n\)\uline{, then} \(\Ho^n(E)\) \uline{is of finite type.}

\uline{Proof.} The sequence
\[
  0\longrightarrow Z^n(E)\longrightarrow E^n\longrightarrow E^{n+1}\longrightarrow\ldots\longrightarrow 0
\]
is exact and \(E^i\in\Ob\underline{C}_0\) for \(i\geq n\). Since \(\underline{C}_0\) is locally quasi-stable under kernels of epimorphisms, it follows that \(Z^n(E)\) is of finite type, hence so is \(\Ho^n(E)\), which is a quotient of it.

%%%%%%%%%%  scan idx 109  |  folio 96  %%%%%%%%%%
\SGAPrintPage{96}%
\sgaStableTarget{sga6.e01.prop.2-2}\textbf{Proposition 2.2.} \uline{Let} \(F\in\Ob\Dpx(\underline{C}_S)\) \uline{and} \(n\in\mathbb Z\)\uline{. Then}:

a) \uline{When} \(F\) \uline{is strictly} \(n\)\uline{-pseudo-coherent (relative to} \(\underline{C}_0\)\uline{) and} \(u:E\to F\) \uline{is a quasi-isomorphism, the set of} \(X\) \uline{over} \(S\) \uline{such that there exists a quasi-isomorphism} \(v:F'\longrightarrow E_X\) \uline{with} \(F'\) \uline{strictly} \(n\)\uline{-pseudo-coherent is a refinement of} \(S\)\uline{.}

b) \uline{The following conditions are equivalent}:
\begin{enumerate}[label=(\roman*)]
\item \uline{the set of} \(X\) \uline{over} \(S\) \uline{such that there exists a quasi-isomorphism} \(u:E\longrightarrow F_X\) \uline{of} \(\Cpx(\underline{C}_X)\)\uline{, with} \(E\) \uline{strictly} \(n\)\uline{-pseudo-coherent, is a refinement of} \(S\);
\item[\((\mathrm{i}')\)] \uline{the same condition as (i), but with} \(u\) \uline{an isomorphism of} \(\Dpx(\underline{C}_X)\);
\item \uline{the set of} \(X\) \uline{over} \(S\) \uline{such that there exists an} \(n\)\uline{-quasi-isomorphism (1.4.3)} \(u:E\longrightarrow F_X\) \uline{of} \(\Cpx(\underline{C}_X)\)\uline{, with} \(E\) \uline{strictly perfect, is a refinement of} \(S\);
\item[\((\mathrm{ii}')\)] \uline{the same condition as (ii), but with} \(u\) \uline{an} \(n\)\uline{-isomorphism of} \(\Dpx(\underline{C}_X)\).
\end{enumerate}

\uline{Proof.} a) For an integer \(i\) sufficiently large, there exists an \(i\)-quasi-isomorphism \(v_i:F'_i\longrightarrow E\), with \(F'_i\) strictly \(n\)-pseudo-coherent (take \(F'_i=0\)). Let us show that, for \(i>n\), there exists, after refining \(S\), an \((i-1)\)-quasi-isomorphism \(v_{i-1}:F'_{i-1}\longrightarrow E\), with \(F'_{i-1}\) strictly \(n\)-pseudo-coherent. By (1.4.4), it suffices to show that \(\Ho^{i-1}(M')\), where \(M'\) is the cone of \(v_i\), is of \(\underline{C}_0\)-finite type. Now, with \(M''\) denoting the cone of \(uv_i\), the hypothesis that \(u\) is a quasi-isomorphism and (1.4.2) imply that
\[
  \Ho^j(M')\qis \Ho^j(M'')\qquad\text{for every }j
\]
and \(\Ho^j(M')=0\) for \(j\geq i\).
%%%%%%%%%%  scan idx 110  |  folio 97  %%%%%%%%%%
\SGAPrintPage{97}%
The cone formula \(M''{}^j=F'_i{}^{\,j+1}\oplus F^j\) shows, on the other hand, that \(M''\) is strictly \((i-1)\)-pseudo-coherent. Thus, by (2.1.1), \(\Ho^{i-1}(M'')=\Ho^{i-1}(M')\) is of finite type, whence our assertion. By iteration, and after refining \(S\), we deduce an \(n\)-quasi-isomorphism \(v_n:F'_n\longrightarrow E\), where \(F'_n\) is strictly \(n\)-pseudo-coherent. But, by (1.4.5), we can extend \(v_n\) to a quasi-isomorphism \(v:F'\longrightarrow E\) without changing anything in degree \(\geq n\), hence in particular with \(F'\) strictly \(n\)-pseudo-coherent, which proves a).

b) The equivalence of conditions (i) and \((\mathrm{i}')\) follows at once from a) and the definition of the isomorphisms of \(\Dpx^-(\underline{C}_X)\). For \(E\in\Ob\Dpx^-(\underline{C}_S)\) strictly \(n\)-pseudo-coherent, there exists an exact sequence
\[
  0\longrightarrow E'\longrightarrow E\longrightarrow E''\longrightarrow 0,
\]
where \(E'\) is strictly perfect, and \(E''{}^i=0\) for \(i\geq n\); consequently \(E'\to E\) is an \(n\)-quasi-isomorphism, giving the implication \((\mathrm{i})\Rightarrow(\mathrm{ii})\). The implication \((\mathrm{ii})\Rightarrow(\mathrm{ii}')\) is trivial, so it remains to prove \((\mathrm{ii}')\Rightarrow(\mathrm{i})\). Let, then, \(u:E\longrightarrow F\) be an \(n\)-isomorphism of \(\Dpx^-(\underline{C}_S)\), with \(E\) strictly perfect; \(u\) is therefore defined by a diagram of \(\Cpx^-(\underline{C}_S)\)
\[
\begin{tikzcd}
& E' \arrow[dl,"s"'] \arrow[dr,"u'"]\\
E && F ,
\end{tikzcd}
\]
where \(s\) is a quasi-isomorphism and \(u'\) an \(n\)-quasi-isomorphism. By a), there exists, after refining \(S\), a quasi-isomorphism \(t:E''\longrightarrow E'\), with \(E''\) strictly \(n\)-pseudo-coherent. We then apply (1.4.5) to extend \(u't:E''\to F\) to a quasi-isomorphism \(E'''\longrightarrow F\) with \(E'''\) strictly \(n\)-pseudo-coherent, which proves \((\mathrm{ii}')\Rightarrow(\mathrm{i})\) and completes the proof of (2.2).

%%%%%%%%%%  scan idx 111  |  folio 98  %%%%%%%%%%
\SGAPrintPage{98}%
\sgaStableTarget{sga6.e01.def.2-3}\textbf{Definition 2.3.} \uline{Under the hypotheses of (2.2), we shall say that} \(F\in\Ob\Dpx(\underline{C}_S)\) \uline{is} \(\underline{n}\)-\(\underline{C}_0\)-pseudo-coherent \uline{when} \(F\) \uline{satisfies the equivalent conditions b) of (2.2). We shall say that} \(F\) \uline{is} \(\underline{C}_0\)-pseudo-coherent \uline{when} \(F\) \uline{is} \(n\)\uline{-}\(\underline{C}_0\)\uline{-pseudo-coherent for every} \(n\)\uline{.}

When this will cause no confusion, we shall sometimes allow ourselves to say \(n\)-pseudo-coherent, pseudo-coherent, instead of \(n\)-\(\underline{C}_0\)-pseudo-coherent, \(\underline{C}_0\)-pseudo-coherent.

\sgaStableTarget{sga6.e01.rem.2-3-1}\textbf{Remark 2.3.1.} Let \(F\in\Ob\Dpx(\underline{C}_S)\) be a pseudo-coherent complex. The set of \(X\) over \(S\) such that there exists a quasi-isomorphism \(E\longrightarrow F_X\) with \(E\) strictly pseudo-coherent (2.1) is in general not a refinement of \(S\). However, we shall see later (2.7. a) and Exp.II) important cases where it is.

The proof of (2.2 a)) gives the following result:

\sgaStableTarget{sga6.e01.prop.2-4}\textbf{Proposition 2.4.} \uline{Let} \(n,i\in\mathbb Z\)\uline{, and let} \(u:E\to F\) \uline{be an} \(i\)\uline{-quasi-isomorphism of} \(\Cpx(\underline{C}_S)\)\uline{, with} \(F\) \(n\)\uline{-pseudo-coherent and} \(E\) \uline{strictly} \(n\)\uline{-pseudo-coherent. Then, after refining} \(S\)\uline{, there exists a commutative triangle of} \(\Cpx(\underline{C}_S)\)
\[
\begin{tikzcd}
E \arrow[r,"v"] \arrow[d,"u"'] & E' \arrow[dl,"u'"]\\
F & ,
\end{tikzcd}
\]
\uline{where} \(E'\) \uline{is strictly} \(n\)\uline{-pseudo-coherent,} \(u'\) \uline{a quasi-isomorphism and} \(v\) \uline{a monomorphism inducing the identity in degree} \(\geq i\)\uline{.}

\sgaStableTarget{sga6.e01.prop.2-5}\textbf{Proposition 2.5.} a) \uline{When} \(F\in\Ob\Dpx(\underline{C}_S)\) \uline{is} \(n\)\uline{-pseudo-coherent (resp. pseudo-coherent),} \(F[m]\) \uline{is} \((n-m)\)\uline{-pseudo-coherent (resp. pseudo-coherent).}

%%%%%%%%%%  scan idx 112  |  folio 99  %%%%%%%%%%
\SGAPrintPage{99}%
b) \uline{Let}
\sgaStableTarget{sga6.e01.eq.unnamed.occ2}%
\[
\begin{tikzcd}
& G \arrow[dl,"+1"'] & \\
E \arrow[rr] && F \arrow[ul]
\end{tikzcd}
\tag{*}
\]
\uline{be a distinguished triangle of} \(\Dpx(\underline{C}_S)\)\uline{. When} \(E\) \uline{and} \(G\) \uline{are} \(n\)\uline{-pseudo-coherent (resp. pseudo-coherent), then} \(F\) \uline{is} \(n\)\uline{-pseudo-coherent (resp. pseudo-coherent).}

\uline{Proof.} a) Trivial from the definitions.

b) It of course suffices to prove the assertion without the ``resp.'' alternatives. The triangle \((*)\) can also be written
\[
\begin{tikzcd}
& F \arrow[dl,"+1"'] & \\
G[-1] \arrow[rr] && E \arrow[ul] & .
\end{tikzcd}
\]
By a), \(G[-1]\) is \((n+1)\)-pseudo-coherent. We are therefore reduced to proving that, when \(E\) is \((n+1)\)-pseudo-coherent and \(F\) \(n\)-pseudo-coherent, then \(G\) is \(n\)-pseudo-coherent. After localizing, we can assume that \(E\) (resp.\ \(F\)) is strictly \((n+1)\)-pseudo-coherent (resp.\ \(n\)-pseudo-coherent). The arrow \(E\xrightarrow{u}F\) of \(\Dpx^-(\underline{C}_S)\) is defined by a pair of arrows of \(\Cpx^-(\underline{C}_S)\)
\[
\begin{tikzcd}
& E' \arrow[dl,"s"'] \arrow[dr,"u'"]\\
E && F ,
\end{tikzcd}
\]
where \(s\) is a quasi-isomorphism. By (2.2 a)), there exists, after localizing, a quasi-isomorphism \(t:E''\to E'\), where \(E''\) is strictly \((n+1)\)-pseudo-coherent. We are therefore reduced to the case where \(u\) is defined by an
%%%%%%%%%%  scan idx 113  |  folio 100  %%%%%%%%%%
\SGAPrintPage{100}%
``actual'' arrow of \(\Cpx^-(\underline{C}_S)\). The assertion then follows from the definitions and the formula \(C(u)^i=E^{i+1}\oplus F^i\), where \(C(u)\) denotes the cone of \(u\).

\sgaStableTarget{sga6.e01.schol.2-6}\textbf{Scholium 2.6.} By ``rotation'' of the triangles, we see that pseudo-coherence properties of two of the objects of the triangle \((*)\) imply another such property of the third. The different cases can be summarized in the table below (the lowercase letters denote degrees of pseudo-coherence and each row corresponds to an implication whose conclusion is underlined):
\[
\begin{array}{ccc}
E&F&G\\[0.2em]
n+1&n&\underline n\\
n&\underline n&n\\
\underline n&n&n-1 .
\end{array}
\]

We shall denote by \(\Dpx(\underline{C})_{\underline{C}_0\text{-}n\text{-coh}}\) (resp.\ \(\Dpx(\underline{C})_{\underline{C}_0\text{-}\mathrm{coh}}\)) the strictly full fibered subcategory of \(\Dpx(\underline{C})\) whose fibre at each \(X\in\Ob \underline{S}\) consists of the objects of \(\Dpx(\underline{C}_X)\) that are \(n\)-pseudo-coherent (resp.\ pseudo-coherent) relative to \(\underline{C}_0\). By (2.5), the category \(\Dpx(\underline{C})_{\underline{C}_0\text{-}\mathrm{coh}}\) is a \uline{triangulated} sub-\(\underline{S}\)\uline{-category} of \(\Dpx(\underline{C})\) (1.1).

\sgaStableTarget{sga6.e01.prop.2-7}\textbf{Proposition 2.7.} \uline{Assume that} \(\underline{S}\) \uline{is a punctual site, so that} \(\underline{C}=\underline{C}_S\)\uline{,} \(\underline{C}_0=\underline{C}_{0S}\)\uline{. Then}:

a) \uline{When} \(F\in\Ob\Dpx(\underline{C})\) \uline{is pseudo-coherent, there exists a quasi-isomorphism} \(E\longrightarrow F\)\uline{, with} \(E\) \uline{strictly pseudo-coherent} (2.1).

b) \uline{Denote by} \(\Kpx^{-,\varnothing}(\underline{C}_0)\) \uline{the full triangulated subcategory of} \(\Kpx^-(\underline{C}_0)\) \uline{consisting of the acyclic complexes. The canonical functor}
\[
  \Kpx^-(\underline{C}_0)/\Kpx^{-,\varnothing}(\underline{C}_0)
  \longrightarrow \Dpx(\underline{C})
\]
%%%%%%%%%%  scan idx 114  |  folio 101  %%%%%%%%%%
\SGAPrintPage{101}%
\uline{is fully faithful and has essential image} \(\Dpx(\underline{C})_{\underline{C}_0\text{-}\mathrm{coh}}\).

\uline{Proof.} a) Let \(F\in\Ob\Dpx(\underline{C})_{\underline{C}_0\text{-}\mathrm{coh}}\), and suppose that an \(i\)-quasi-isomorphism (of \(\Cpx(\underline{C})\)) \(u_i:E_i\longrightarrow F\) has been constructed with \(E_i\) strictly pseudo-coherent (this is possible for \(i\) sufficiently large with \(E_i=0\)). Let \(G_i\) be the cone of \(u_i\). It follows from (1.4.2) and (2.1.1) that \(\Ho^{i-1}(G_i)\) is of \(\underline{C}_0\)-finite type. Thus, by (1.4.4), there exists a commutative triangle of \(\Cpx(\underline{C})\)
\[
\begin{tikzcd}
E_i \arrow[r,"v_{i-1}"] \arrow[d,"u_i"'] & E_{i-1} \arrow[dl,"u_{i-1}"]\\
F & ,
\end{tikzcd}
\]
where \(E_{i-1}\) is strictly pseudo-coherent, \(u_{i-1}\) an \((i-1)\)-quasi-isomorphism, and \(v_{i-1}\) a monomorphism inducing the identity in degree \(\geq i\). By passage to the limit (cf.\ proof of (1.4.5)), we deduce a quasi-isomorphism \(u:E\to F\) with \(E\) strictly pseudo-coherent.

b) Apply the criterion ([V] chap I \S\ 2 4.2 b)(i)) and a).

Observe that the proof of a) gives the following more precise result, which generalizes (1.4.5):

\sgaStableTarget{sga6.e01.prop.2-7-1}\textbf{Proposition 2.7.1.} \uline{Under the hypotheses of (2.7), let} \(u:E\to F\) \uline{be an} \(n\)\uline{-quasi-isomorphism of} \(\Cpx(\underline{C})\)\uline{, with} \(F\) \uline{pseudo-coherent and} \(E\) \uline{strictly pseudo-coherent. There then exists a commutative triangle of} \(\Cpx(\underline{C})\)
\[
\begin{tikzcd}
E \arrow[r,"v"] \arrow[d,"u"'] & E' \arrow[dl,"u'"]\\
F & ,
\end{tikzcd}
\]
\uline{where} \(E'\) \uline{is strictly pseudo-coherent,} \(u'\) \uline{a quasi-isomorphism, and} \(v\) \uline{a}
%%%%%%%%%%  scan idx 115  |  folio 102  %%%%%%%%%%
\SGAPrintPage{102}%
\uline{monomorphism inducing the identity in degree} \(\geq n\).

\sgaStableTarget{sga6.e01.def.2-8}\textbf{Definition 2.8.} \uline{Let} \(n\in\mathbb N\) \uline{and} \(F\in\Ob\underline{C}_S\)\uline{. We shall say that} \(F\) \uline{is of} \(\underline{n}\)-\(\underline{C}_0\)-finite presentation \uline{when the set of} \(X\) \uline{over} \(S\) \uline{such that there exists an exact sequence}
\[
  E^{-n}\longrightarrow\ldots\longrightarrow E^0\longrightarrow F_X\longrightarrow 0,
\]
\uline{with} \(E^i\in\Ob\underline{C}_{0X}\) \uline{for} \(-n\leq i\leq 0\)\uline{, is a refinement of} \(S\)\uline{. We shall say that} \(F\) \uline{is of} \(\infty\)-\(\underline{C}_0\)-finite presentation \uline{when} \(F\) \uline{is of} \(n\)\uline{-}\(\underline{C}_0\)\uline{-finite presentation for every} \(n\in\mathbb N\).

Of course, we shall sometimes simply say \(n\)-finite presentation when no ambiguity concerning \(\underline{C}_0\) is possible.

We shall say ``of finite presentation'' for ``of 1-finite presentation'', and ``of finite type'' for ``of 0-finite presentation'', the latter notion coinciding with the one already introduced in (1.2).

\sgaStableTarget{sga6.e01.prop.2-9}\textbf{Proposition 2.9.} \uline{Let} \(n\in\mathbb N\) \uline{and} \(F\in\Ob\underline{C}_S\)\uline{. Equivalent conditions}:

(i) \(F\) \uline{is of} \(n\)\uline{-finite presentation (resp.\ of} \(\infty\)\uline{-finite presentation)};

(ii) \(F\) \uline{is} \((-n)\)\uline{-pseudo-coherent (resp.\ pseudo-coherent)}.

\uline{In the case where} \(\underline{S}\) \uline{is a punctual site,} \(F\) \uline{is pseudo-coherent if and only if there exists an exact sequence}
\[
  \ldots\longrightarrow E^i\longrightarrow\ldots\longrightarrow E^0\longrightarrow F\longrightarrow 0,
\]
\uline{with} \(E^i\in\Ob\underline{C}_0\) \uline{for every} \(i\).

\uline{Proof.} For the equivalence \((\mathrm{i})\Leftrightarrow(\mathrm{ii})\), it clearly suffices to treat the case without the ``resp.'' alternatives. The implication \((\mathrm{i})\Rightarrow(\mathrm{ii})\) follows immediately from the definitions, and
%%%%%%%%%%  scan idx 116  |  folio 103  %%%%%%%%%%
\SGAPrintPage{103}%
the implication \((\mathrm{ii})\Rightarrow(\mathrm{i})\) follows at once from (2.4): the arrow \(0\to F\) is a 1-quasi-isomorphism; we deduce (after localizing) a quasi-isomorphism \(E\longrightarrow F\) with \(E\) strictly \((-n)\)-pseudo-coherent and \(E^i=0\) for \(i\geq 1\), i.e.\ an \(n\)-finite presentation of \(F\). The second assertion is proved in a similar way (use (2.7.1)).

\sgaStableTarget{sga6.e01.prop.2-10}\textbf{Proposition 2.10.} a) \uline{Let} \(m,n\in\mathbb Z\)\uline{, with} \(n\leq m\)\uline{. When} \(F\in\Ob\Dpx^-(\underline{C}_S)\) \uline{is} \(n\)\uline{-pseudo-coherent and acyclic in degree} \(\geq m\)\uline{, there locally exists a quasi-isomorphism} \(E\longrightarrow F\)\uline{, with} \(E\) \uline{strictly} \(n\)\uline{-pseudo-coherent and} \(E^i=0\) \uline{for} \(i\geq m\)\uline{. When} \(\underline{S}\) \uline{is punctual, and} \(F\) \uline{pseudo-coherent and acyclic in degree} \(\geq m\)\uline{, there exists a quasi-isomorphism} \(E\longrightarrow F\) \uline{with} \(E\) \uline{strictly pseudo-coherent and} \(E^i=0\) \uline{for} \(i\geq m\)\uline{.}

b) \uline{Let} \(F\in\Ob\Dpx^-(\underline{C}_S)\) \uline{be acyclic in degree} \(>n\)\uline{. For} \(F\) \uline{to be} \(n\)\uline{-pseudo-coherent, it is necessary and sufficient that} \(\Ho^n(F)\) \uline{be of finite type.}

\uline{Proof.} a) Since \(F\) is acyclic in degree \(\geq m\), the arrow \(0\to F\) is an \(m\)-quasi-isomorphism. The assertion follows from (2.4) and (2.7.1).

b) There exists a distinguished triangle of \(\Dpx^-(\underline{C}_S)\):
\[
\begin{tikzcd}
& F' \arrow[dl,"+1"'] & \\
F \arrow[rr] && \Ho^n(F)[-n] \arrow[ul] & ,
\end{tikzcd}
\]
(to see this, begin by truncating \(F\), replacing \(F^n\) by \(Z^n(F)\) and \(F^i\) by \(0\) for \(i>n\)). The exact cohomology sequence implies \(\Ho^i(F')=0\) for \(i\geq n-1\), so \(F'\) is \((n-1)\)-pseudo-coherent. We conclude using (2.5) and (2.6).

% Additional notation for Expose I continuation
\providecommand{\RHom}{\mathrm R\!\Hom}
\providecommand{\Ext}{\operatorname{Ext}}
\providecommand{\coh}{\operatorname{coh}}
\providecommand{\PreCoh}{\operatorname{precoh}}
\providecommand{\Tr}{\operatorname{Tr}}
\providecommand{\Qis}{\operatorname{Qis}}
\providecommand{\Rad}{\operatorname{rad}}
\providecommand{\Spec}{\operatorname{Spec}}
\providecommand{\End}{\operatorname{End}}
\providecommand{\Id}{\operatorname{Id}}
\providecommand{\varinjlim}{\mathop{\mathrm{lim}}\limits_{\longrightarrow}}
\providecommand{\varprojlim}{\mathop{\mathrm{lim}}\limits_{\longleftarrow}}
\providecommand{\qcoh}{\operatorname{qcoh}}
\providecommand{\amp}{\operatorname{amp}}

%%%%%%%%%%  scan idx 117  |  folio 104  %%%%%%%%%%
\SGAPrintPage{104}%
\sgaStableTarget{sga6.e01.expose.occ2}%
\sgaStableTarget{sga6.e01.prop.2-11}\textbf{Proposition 2.11.} \uline{Let} \(n\in\mathbb Z\)\uline{.}

a) \uline{Let} \(F\in\Ob\Cpx^-(\underline{C}_S)\)\uline{. When, for every} \(i\geq n\) \uline{(resp.\ for every} \(i\)\uline{),} \(F^i\) \uline{is of} \((i-n)\)\uline{-finite presentation (resp.\ of} \(\infty\)\uline{-finite presentation), then} \(F\) \uline{is} \(n\)\uline{-pseudo-coherent (resp.\ pseudo-coherent).}

b) \uline{Let} \(F\in\Ob\Dpx^-(\underline{C}_S)\)\uline{. When, for every} \(i\geq n\) \uline{(resp.\ for every} \(i\)\uline{),} \(\Ho^i(F)\) \uline{is of} \((i-n)\)\uline{-finite presentation (resp.\ of} \(\infty\)\uline{-finite presentation), then} \(F\) \uline{is} \(n\)\uline{-pseudo-coherent (resp.\ pseudo-coherent).}

\uline{Proof.} It suffices to treat the assertions without the ``resp.'' alternatives.

a) We argue by induction on \(N\) such that \(F^i=0\) for \(i>n+N\). For \(N=0\), the assertion follows at once from (2.10 b)). Suppose \(N\geq1\). There is an exact sequence
\[
  0\longrightarrow F^{n+N}[-n-N]\longrightarrow F\longrightarrow F'\longrightarrow 0,
\]
By the induction hypothesis, \(F'\) is \(n\)-pseudo-coherent; on the other hand, \(F^{n+N}[-n-N]\) is \(n\)-pseudo-coherent by hypothesis (2.9 and 2.5.a)); thus, by (2.5.b)), \(F\) is \(n\)-pseudo-coherent.

b) We argue by induction on \(N\) such that \(\Ho^i(F)=0\) for \(i>n+N\). For \(N=0\), the assertion follows from (2.10 b)). Suppose \(N\geq1\). There is a distinguished triangle of \(\Dpx^-(\underline{C}_S)\) (cf.\ proof of (2.10.b))):
\[
\begin{tikzcd}
& F' \arrow[dl,"+1"'] & \\
F \arrow[rr] && \Ho^{n+N}(F)[-n-N] \arrow[ul] & .
\end{tikzcd}
\]
The exact cohomology sequence implies
\[
  \Ho^i(F')=0\quad\text{for }i\geq n+N-1,\quad\text{i.e. }i>(n-1)+(N-1),
\]
and
\[
  \Ho^i(F')\qis\Ho^{i+1}(F)\quad\text{for }i<n+N-1 .
\]
%%%%%%%%%%  scan idx 118  |  folio 105  %%%%%%%%%%
\SGAPrintPage{105}%
By the induction hypothesis (applied with \(n-1\) in place of \(n\)), \(F'\) is therefore \((n-1)\)-pseudo-coherent. On the other hand, \(\Ho^{n+N}(F)[-n-N]\) is \(n\)-pseudo-coherent (2.9 and 2.5.a)). Thus, by (2.6), \(F\) is \(n\)-pseudo-coherent.

\sgaStableTarget{sga6.e01.rem.2-11-1}\textbf{Remark 2.11.1.} When \(F\) is pseudo-coherent, the \(F^i\) (resp.\ \(\Ho^i(F)\)) are in general not pseudo-coherent, nor even of finite type. However, in the next paragraph, in the sorites on notions of coherence, we shall see important examples where the pseudo-coherence of \(F\) implies that of the \(\Ho^i(F)\).

\sgaStableTarget{sga6.e01.prop.2-12}\textbf{Proposition 2.12.} \uline{Let} \(F',F''\in\Ob\Dpx(\underline{C}_S)\)\uline{, and} \(F=F'+F''\)\uline{. For} \(F\) \uline{to be} \(n\)\uline{-pseudo-coherent (resp.\ pseudo-coherent), it is necessary and sufficient that} \(F'\) \uline{and} \(F''\) \uline{be} \(n\)\uline{-pseudo-coherent (resp.\ pseudo-coherent).}

\uline{Proof.} It suffices to prove the assertion without the ``resp.'' alternatives. Sufficiency is clear; let us prove necessity. By induction, we may suppose we have already proved that \(F'\) and \(F''\) are \((n+1)\)-pseudo-coherent. After localizing, we can even suppose that \(F'\) and \(F''\) are strictly \((n+1)\)-pseudo-coherent and that \(F=F'+F''\) in \(\Cpx^-(\underline{C}_S)\). Denote by \(F'_{>n}\) (resp.\ \(F'_{\leq n}\)) the complex obtained from \(F'\) by replacing \(F'^i\) by \(0\) for \(i\leq n\) (resp.\ \(i>n\)), and similarly \(F''_{>n}\) (resp.\ \(F''_{\leq n}\))\dots{} There is an exact sequence
\[
  0\longrightarrow F'_{>n}+F''_{>n}\longrightarrow F
  \longrightarrow F'_{\leq n}+F''_{\leq n}\longrightarrow 0 .
\]
\(F'_{>n}+F''_{>n}\) is strictly perfect and \(F\) is \(n\)-pseudo-coherent, so (2.6) \(F'_{\leq n}+F''_{\leq n}\) is \(n\)-pseudo-coherent. By (2.10 b)), this means that \(\Ho^n(F'_{\leq n}+F''_{\leq n})\) is of finite type. But we have
\[
  \Ho^n(F'_{\leq n}+F''_{\leq n})
  =\Ho^n(F'_{\leq n})+\Ho^n(F''_{\leq n}),
\]
%%%%%%%%%%  scan idx 119  |  folio 106  %%%%%%%%%%
\SGAPrintPage{106}%
so \(\Ho^n(F'_{\leq n})\) (resp.\ \(\Ho^n(F''_{\leq n})\)) is of finite type, so again, by (2.10.b)), \(F'_{\leq n}\) (resp.\ \(F''_{\leq n}\)) is \(n\)-pseudo-coherent. Applying (2.5) to the exact sequence
\[
  0\longrightarrow F'_{>n}\longrightarrow F'\longrightarrow F'_{\leq n}\longrightarrow 0,
\]
(resp.\dots), we deduce that \(F'\) (resp.\ \(F''\)) is \(n\)-pseudo-coherent, QED.

\sgaStableTarget{sga6.e01.prop.2-13}\textbf{Proposition 2.13.} \uline{Let} \((\underline{S}',\underline{C}',\underline{C}'_0)\) \uline{be a triple satisfying the same hypotheses} (2.0) \uline{as} \((\underline{S},\underline{C},\underline{C}_0)\)\uline{. Let} \(f:\underline{S}'\to \underline{S}\) \uline{be a morphism of sites} ([G] I 1.4)\uline{, corresponding to an underlying functor} \(f^{-1}:\underline{S}\to \underline{S}'\)\uline{, and let}
\(u:\Dpx^*(\underline{C})\longrightarrow \Dpx^*(\underline{C}')\) \((*=+,-,\varnothing)\)
\uline{be a cartesian functor over} \(f\)\uline{,} i.e.\ \uline{a cartesian} \(\underline{S}\)\uline{-functor} \(\Dpx^*(\underline{C}')\times_{\underline{S}'}\underline{S}\longleftarrow\Dpx^*(\underline{C})\)\uline{. Assume that} \(u\) \uline{is exact on each fibre and that there exist} \(a,b\in\mathbb Z\) \uline{such that the following conditions hold for every} \(S\in\Ob \underline{S}\):

(i) \uline{when} \(E\in\Ob\underline{C}_{0S}\)\uline{,} \(u(E)\) \uline{is} \(a\)\uline{-}\(\underline{C}'_0\)\uline{-pseudo-coherent (resp.} \(\underline{C}'_0\)\uline{-pseudo-coherent)}

(ii) \uline{when} \(F\in\Ob\Dpx^*(\underline{C}_S)\) \uline{is acyclic in degree} \(>0\)\uline{,} \(u(F)\) \uline{is acyclic in degree} \(>b\)\uline{.}

\uline{Then, when} \(F\in\Ob\Dpx^*(\underline{C}_S)\) \uline{is} \(m\)\uline{-pseudo-coherent rel.\ to} \(\underline{C}_0\) \uline{and acyclic in degree} \(>n\)\uline{,} \(u(F)\in\Ob\Dpx^*(\underline{C}'_{S'})\) \uline{(where} \(S'=f^{-1}(S)\)\uline{) is} \(\sup(m+b,n+a)\)\uline{-pseudo-coherent (rel.\ to} \(\underline{C}'_0\)\uline{) (resp.} \((m+b)\)\uline{-pseudo-coherent) and acyclic in degree} \(>n+b\)\uline{.}

\uline{In particular, in the case of the ``resp.'' alternative, when} \(F\in\Ob\Dpx^*(\underline{C}_S)\) \uline{is pseudo-coherent} (rel.\ to \(\underline{C}_0\)), \(u(F)\) \uline{is pseudo-coherent (rel.\ to} \(\underline{C}'_0\)\uline{).}

\uline{Proof.} We can restrict ourselves to the case without the ``resp.'' alternative. Let us first show that, for \(F\in\Ob\Dpx^*(\underline{C}_S)\) strictly perfect (2.1) with \(F^i=0\) for \(i>n\), \(u(F)\) is \((n+a)\)-\(\underline{C}'_0\)-pseudo-coherent. We argue by induction on the number \(N\) of indices \(i\) such that \(F^i\neq0\). For \(N=1\),
%%%%%%%%%%  scan idx 120  |  folio 107  %%%%%%%%%%
\SGAPrintPage{107}%
we have \(F=F^p[-p]\) for some \(p\leq n\), so \(u(F)\) is \((p+a)\)-\(\underline{C}'_0\)-pseudo-coherent, hence a fortiori \((n+a)\)-pseudo-coherent. For \(N>1\), there is an exact sequence
\[
  0\longrightarrow F'\longrightarrow F\longrightarrow F''\longrightarrow 0,
\]
where \(F'\) and \(F''\) are strictly perfect of length \(\leq N-1\) and satisfy \(F'^i=F''^i=0\) for \(i>n\). Applying (2.5.b)) to the distinguished triangle
\[
\begin{tikzcd}
& u(F'') \arrow[dl,"+1"'] & \\
u(F') \arrow[rr] && u(F) \arrow[ul] & ,
\end{tikzcd}
\]
we deduce from the induction hypothesis that \(u(F)\) is \((n+a)\)-pseudo-coherent.

Let us now pass to the general case: \(F\) \(m\)-\(\underline{C}_0\)-pseudo-coherent and acyclic in degree \(>n\). After localizing, we can, by (2.10.a)), suppose that \(F\) is strictly \(m\)-pseudo-coherent and satisfies \(F^i=0\) for \(i>n\). There is then an exact sequence
\[
  0\longrightarrow F'\longrightarrow F\longrightarrow F''\longrightarrow 0,
\]
where \(F''\) is zero in degree \(\geq m\), and \(F'\) strictly perfect and zero in degree \(>n\). In the distinguished triangle
\[
\begin{tikzcd}
& u(F'') \arrow[dl,"+1"'] & \\
u(F') \arrow[rr] && u(F) \arrow[ul] & ,
\end{tikzcd}
\]
\(u(F'')\) is acyclic in degree \(\geq m+b\), hence \((m+b)\)-pseudo-coherent, and \(u(F')\) is \((n+a)\)-pseudo-coherent by what we have seen. We conclude by (2.5.b)).

\sgaStableTarget{sga6.e01.cor.2-14}\textbf{Corollary 2.14.} \uline{Let} \(\underline{C}_1\) \uline{be a strictly full additive sub-}\(\underline{S}\)\uline{-category of} \(\underline{C}\) \uline{such that} \(\underline{C}_0\subset\underline{C}_1\subset\underline{C}\)\uline{. Assume that every object of} \(\underline{C}_1\) \uline{is} \(\underline{C}_0\)\uline{-pseudo-coherent.}
%%%%%%%%%%  scan idx 121  |  folio 108  %%%%%%%%%%
\SGAPrintPage{108}%
\uline{Then}:

(i) \(\underline{C}_1\) \uline{is locally quasi-liftable in} \(\underline{C}\) (1.2) \uline{and locally quasi-stable under kernels of epimorphisms} (1.2)\uline{.}

(ii) \uline{Let} \(S\in\Ob \underline{S}\) \uline{and} \(F\in\Ob\Dpx(\underline{C}_S)\)\uline{. For} \(F\) \uline{to be} \(n\)\uline{-pseudo-coherent (resp.\ pseudo-coherent) relative to} \(\underline{C}_1\)\uline{, it is necessary and sufficient that it be so relative to} \(\underline{C}_0\)\uline{.}

\uline{Proof.} (i) Let \(E\to F\) be an epimorphism of \(\underline{C}_S\), with \(F\in\Ob\underline{C}_{1S}\). Since \(F\) is \(\underline{C}_0\)-pseudo-coherent (hence a fortiori of \(\underline{C}_0\)-finite type), there exists (after localizing) an epimorphism \(E'\to F\) with \(E'\in\Ob\underline{C}_{0S}\). Using the quasi-liftability of \(\underline{C}_0\) in \(\underline{C}\), we can suppose (after localizing again) that this factors through \(E\), whence the quasi-liftability of \(\underline{C}_1\). On the other hand, when \(E\in\Ob\underline{C}_S\) admits a finite right resolution by objects of \(\underline{C}_1\), \(E\) is \(\underline{C}_0\)-pseudo-coherent by (2.11.a)), hence a fortiori of \(\underline{C}_1\)-finite type.

(ii) It suffices to prove necessity. For this we can apply, as we choose, (2.11.a)) or (2.13) with \(\underline{S}'=\underline{S}\), \(\underline{C}'=\underline{C}\), \(f=\text{Identity}\), \(u=\text{Identity}\), replacing \(\underline{C}'_0\) and \(\underline{C}_0\) respectively by \(\underline{C}_0\) and \(\underline{C}_1\).

\sgaStableTarget{sga6.e01.schol.2-14-1}\textbf{Scholium 2.14.1.} There thus exists a largest strictly full additive sub-\(\underline{S}\)-category of \(\underline{C}\) giving the same notion of \(n\)-pseudo-coherence (resp.\ pseudo-coherence) as \(\underline{C}_0\), namely the category formed by the \(\underline{C}_0\)-pseudo-coherent objects. The latter is locally stable under kernels of epimorphisms.

\sgaStableTarget{sga6.e01.ex.2-15}\textbf{Example 2.15.} Let \((\underline{S},A)\) be a ringed topos and \(\underline{C},\underline{C}_0\) the \(\underline{S}\)-categories considered in (1.3.d)). The fibered category \(\Dpx(\underline{C})\) will often be denoted by \(\Dpx(\underline{S},A)\), or \(\Dpx({}_A \underline{S})\) (to
%%%%%%%%%%  scan idx 122  |  folio 109  %%%%%%%%%%
\SGAPrintPage{109}%
suggest that these are left \(A\)-Modules), or even \(\Dpx(\underline{S})\) when there is no ambiguity concerning \(A\). We shall denote by \(S\) the final object of \(\underline{S}\) and put \(\Dpx(\underline{C}_S)=\Dpx({}_A S)\) (or \(\Dpx(S)\)). An object of \(\Dpx(\underline{S})\) will be called \(\underline{n}\)\uline{-pseudo-coherent} (resp.\ \uline{pseudo-coherent}) when it is \(n\)-\(\underline{C}_0\)-pseudo-coherent (resp.\ \(\underline{C}_0\)-pseudo-coherent) (likewise for the strict notions). The fibered category \(\Dpx(\underline{C})_{n\text{-}\underline{C}_0\text{-}\coh}\) (resp.\ \(\Dpx(\underline{C})_{\underline{C}_0\text{-}\coh}\)) (2.6) will be denoted by \(\Dpx({}_A \underline{S})_{n\text{-}\coh}\) (resp.\ \(\Dpx({}_A \underline{S})_{\coh}\)). Observe that, by (2.14), we obtain the same notion of pseudo-coherence on replacing \(\underline{C}_0\) by the \(\underline{S}\)-category \(\underline{C}_1\) defined by \(U\longmapsto\) (category of left \(A_U\)-Modules locally direct factors of free Modules of finite type). When an \(A\)-Module \(M\) is a direct factor of a free \(A\)-Module of finite type, it does not follow, in general, that \(M\) is locally free of finite type. However, we have:

\sgaStableTarget{sga6.e01.lem.2-15-1}\textbf{Lemma 2.15.1.} \uline{Assume that} \(\underline{S}\) \uline{satisfies one of the following two hypotheses:}

(i) \(\underline{S}\) \uline{``has enough points''} (SGA 4 IV 6.4) \uline{and, for every point} \(s\) \uline{of} \(\underline{S}\)\uline{, the projective} \(A_s\)\uline{-modules} (on the left) \uline{of finite type are free, this last condition being satisfied for example} (BOURBAKI Alg. Comm. chap.\ II \S~3 n\textsuperscript{o}~2 cor.~2 of prop.\ 5) \uline{when} \(A_s/\underline{\operatorname{radical}}(A_s)\) \uline{is a field.}

(ii) \((\underline{S},A)\) \uline{is a ``locally ringed'' topos in the sense of} [1]\uline{, which means that} \(A\) \uline{is commutative and that, for every} \(X\in\Ob \underline{S}\) \uline{and every} \(f\in\Gamma(X,A)\)\uline{, we have} \(X=X_f\cup X_{1-f}\)\uline{, where} \(X_f\) \uline{(resp.} \(X_{1-f}\)\uline{) denotes the largest subobject of} \(X\) \uline{on which} \(f\) \uline{(resp.} \(1-f\)\uline{) is invertible.}

\uline{Then every direct factor of a free} \(A\)\uline{-module (on the left) of finite type is locally free of finite type.}

\uline{Proof.} Suppose (i). We first observe that, for two \(A\)-Modules \(E\) and \(F\), with \(E\) of finite presentation, the canonical arrow
%%%%%%%%%%  scan idx 123  |  folio 110  %%%%%%%%%%
\SGAPrintPage{110}%
\[
  \underline{\Hom}(E,F)_s\longrightarrow\Hom(E_s,F_s)
\]
is an isomorphism for every point \(s\) of \(\underline{S}\). It follows that, when \(E\) and \(F\) are both of finite presentation, and \(E_s\simeq F_s\) at a point \(s\) of \(\underline{S}\), there exists an \(s\)-pointed object \(U\) of \(\underline{S}\)\SGAAuthorFootnote{*}{i.e.\ such that \(s\) factors through \(\underline{S}_{/U}\).} such that \(E|U\simeq F|U\). The conclusion follows at once.

Suppose (ii). Let \(n\in\mathbb N\), and denote by \(\mathbb D^n\) the functor on the category of locally ringed topoi (belonging to a fixed universe\dots) defined by \(T\longmapsto\Gamma(T,\underline{O}_T)^n\). We know [1] that the functor \(\underline{\End}_{\mathbb D}(\mathbb D^n)\) is represented by \(M_n=\Spec\mathbb Z[T_{ij}]\ (1\leq i\leq n,\ 1\leq j\leq n)\), equipped with the endomorphism \(u_n\) of \((\underline{O}_{M_n})^n\) defined by the matrix \((T_{ij})\). Now let \(p:A^n\longrightarrow A^n\) be a projector; let us show that \(p(A^n)\) is locally free of finite type. By what we have just said, there exists a morphism \(s:\underline{S}\longrightarrow M_n\) of locally ringed topoi such that \(s^*(u_n)=p\). We can interpret \(s\) as a point of \(M_n\) with values in \(\Gamma(S,A)\), i.e.\ a matrix of order \(n\) with coefficients in \(\Gamma(S,A)\); since this matrix is a projector, \(s\) factors through the closed subscheme \(i:P_n\to M_n\) corresponding to projectors as \(\underline{S}\xrightarrow{s'}P_n\xrightarrow{i}M_n\). Thus \(p=s'^*i^*(u_n)\), and we are done, because by the assertion under hypothesis (i), the image of \(i^*(u_n)\) is a locally free sheaf of finite type.

\sgaStableTarget{sga6.e01.cor.2-16}\textbf{Corollary 2.16.} \uline{Let} \(f:(\underline{S}',A')\to(\underline{S},A)\) \uline{be a morphism of ringed topoi (defined by a functor} \(f^{-1}:\underline{S}\to \underline{S}'\) \uline{and a morphism of sheaves of rings} \(f^{-1}(A)\to A'\)\uline{). Denote by} \(\Dpx({}_{A'}\underline{S}'_A)\) \uline{the derived category of the category of} \(A'\)\uline{-}\(A\)\uline{-biModules (on the left over} \(A'\)\uline{, on the right over} \(A\)\uline{), and by}
\[
  \otimes_A^{\mathrm L}:\Dpx^-(_{A'}\underline{S}'_A)\times_{\underline{S}'}\Dpx^-(_A \underline{S})
  \longrightarrow \Dpx^-(_{A'}\underline{S}')
\]
%%%%%%%%%%  scan idx 124  |  folio 111  %%%%%%%%%%
\SGAPrintPage{111}%
\uline{the derived functor of the tensor product} \((E',E)\longmapsto E'\otimes_{f^{-1}(A)} f^{-1}(E)\)\uline{.} \uline{Let} \(X\in\Ob \underline{S}\)\uline{,} \(X'=f^{-1}(X)\)\uline{,} \(E'\in\Ob\Dpx^-(_{A'}X'_A)\)\uline{,} \(E\in\Ob\Dpx^-(_A X)\)\uline{. Then, when} \(E'\) \uline{is} \(m'\)\uline{-pseudo-coherent (as an object of} \(\Dpx^-(_{A'}\underline{S}')\)\uline{) and acyclic in degree} \(>n'\) \uline{and} \(E\) \uline{is} \(m\)\uline{-pseudo-coherent and acyclic in degree} \(>n\)\uline{,} \(E'\otimes_A^{\mathrm L}E\) \uline{is} \(\sup(m+n',m'+n)\)\uline{-pseudo-coherent and acyclic in degree} \(>n+n'\)\uline{.}

\uline{Proof.} Replacing \(\underline{S}\) by \(\underline{S}_{/X}\), we can suppose that \(X=S\), and we apply (2.13) with \(u=E'\otimes_A^{\mathrm L}:\Dpx^-(_A \underline{S})\to\Dpx^-(_{A'}\underline{S}')\). Let \(U\in\Ob \underline{S}\), \(U'=f^{-1}(U)\), and \(E\) be a complex of left \(A_U\)-Modules. When \(E\) is acyclic in degree \(>n\), \(E'\otimes_A^{\mathrm L}E\) is acyclic in degree \(>n+n'\) (definition of \(\otimes_A^{\mathrm L}\)); on the other hand, for \(E=(A_U)^p\), we have \(E'\otimes_A^{\mathrm L}E=(E'|U')^p\), which is \(m'\)-pseudo-coherent by hypothesis (and (2.12)). The conditions of (2.13) are therefore satisfied, with \(a=m'\), \(b=n'\), whence the conclusion.

In particular:

\sgaStableTarget{sga6.e01.cor.2-16-1}\textbf{Corollary 2.16.1.} \uline{Under the hypotheses of} (2.16), \uline{we have}:

a) \uline{The functor} \(\mathbb{L}f^*: \Dpx^-(_A \underline{S})\to \Dpx^-(_{A'}\underline{S}')\) \uline{(which is simply the functor} \(A'\otimes_A^{\mathrm L}\)\uline{) induces a functor}
\[
  \mathbb{L}f^*: \Dpx^-(_A \underline{S})_{\coh}\longrightarrow \Dpx^-(_{A'}\underline{S}')_{\coh}
\]
(\uline{notation of} (2.6)).

b) \uline{When} \(A\) \uline{is commutative, the functor} \(\otimes_A^{\mathrm L}:\Dpx^-(\underline{S})\times_{\underline{S}}\Dpx^-(\underline{S})\longrightarrow \Dpx^-(\underline{S})\) \uline{induces a functor}
\[
  \otimes_A^{\mathrm L}:\Dpx^-(\underline{S})_{\coh}\times_{\underline{S}}\Dpx^-(\underline{S})_{\coh}
  \longrightarrow \Dpx^-(\underline{S})_{\coh} .
\]

%%%%%%%%%%  scan idx 125  |  folio 112  %%%%%%%%%%
\SGAPrintPage{112}%
\sgaStableTarget{sga6.e01.exercise.2-16-2}\textbf{Exercise 2.16.2.} Let \(\underline{S}\) be a topos, and \(A\to A'\) a morphism of Rings of \(\underline{S}\). We make one of the following hypotheses:
\begin{enumerate}[label=(\roman*)]
\item \(A\to A'\) is surjective with kernel a nilpotent Ideal
\item the functor \(A'\otimes_A-\) on the category of left \(A\)-modules is exact and faithful.
\end{enumerate}
Show that \(E\in\Ob\Dpx^-(_A \underline{S})\) is \(n\)-pseudo-coherent (resp.\ pseudo-coherent) if and only if \(A'\otimes_A^{\mathrm L} E\in\Ob\Dpx^-(_{A'}\underline{S})\) is \(n\)-pseudo-coherent (resp.\ pseudo-coherent).

Hints: First show that a left \(A\)-Module \(L\) is of finite type if and only if \(A'\otimes_A L\) is (cf.\ Bourbaki, Alg. commutative, chap.\ I, \S\ 3, prop.\ 11 and chap.\ II, \S\ 3, prop.\ 4), then reduce to this case by truncation, using (2.5) and (2.10 b)).

\sgaStableTarget{sga6.e01.sec.3}\section*{3. Link with the classical notion of coherence}

\sgaStableTarget{sga6.e01.para.3-0}3.0. The hypotheses and notation are those of (2.0), except that \(\underline{C}_0\) is no longer assumed locally quasi-stable under kernels of epimorphisms.

\sgaStableTarget{sga6.e01.def.3-1}\textbf{Definition 3.1.} \uline{Let} \(F\in\Ob\underline{C}_S\)\uline{. We shall say that} \(F\) \uline{is} precoherent \uline{(relative to} \(\underline{C}_0\)\uline{)}\SGAAuthorFootnote{*}{We introduce this notion here as a technical intermediate step; the only notion to be retained is that of coherence.} \uline{when, for every} \(U\) \uline{over} \(S\) \uline{and every morphism} \(u:E\to F|U\) \uline{with} \(E\in\Ob\underline{C}_{0U}\)\uline{,} \(\Ker(u)\) \uline{is of} \(\underline{C}_0\)\uline{-finite type} (1.2)\uline{. We shall say that} \(F\) \uline{is} coherent \uline{(relative to} \(\underline{C}_0\)\uline{) when} \(F\) \uline{is precoherent and of finite type.}

\sgaStableTarget{sga6.e01.rem.3-2}\textbf{Remark 3.2.} Let \(F\in\Ob\underline{C}_S\). It is clear from the definitions that, when \(F\) is precoherent (resp.\ of finite type), every subobject (resp.\ quotient) of \(F\) is precoherent (resp.\ of finite type).

\sgaStableTarget{sga6.e01.lem.3-3}\textbf{Lemma 3.3.} a) \uline{Let}
\[
  0\longrightarrow F\longrightarrow G\longrightarrow H
\]
%%%%%%%%%%  scan idx 126  |  folio 113  %%%%%%%%%%
\SGAPrintPage{113}%
\uline{be an exact sequence of} \(\underline{C}_S\)\uline{. When} \(G\) \uline{is of finite type and} \(H\) \uline{precoherent,} \(F\) \uline{is of finite type.}

b) \uline{Let}
\[
  F\longrightarrow G\longrightarrow H\longrightarrow 0
\]
\uline{be an exact sequence of} \(\underline{C}_S\)\uline{. When} \(F\) \uline{is of finite type and} \(G\) \uline{precoherent,} \(H\) \uline{is precoherent.}

c) \uline{Let}
\[
  0\longrightarrow F\longrightarrow G\longrightarrow H\longrightarrow 0
\]
\uline{be an exact sequence of} \(\underline{C}_S\)\uline{. When} \(F\) \uline{and} \(H\) \uline{are precoherent (resp.\ of finite type),} \(G\) \uline{is too.}

\uline{Proof.} a) After localizing, we have an epimorphism \(G'\to G\) with \(G'\in\Ob\underline{C}_0\). Hence a commutative diagram with exact rows
\[
\begin{tikzcd}
0\arrow[r]&F'\arrow[r]\arrow[d]&G'\arrow[r]\arrow[d]&H\arrow[d,"\Id"]\\
0\arrow[r]&F\arrow[r]&G\arrow[r]&H .
\end{tikzcd}
\]
\(F'\) is of finite type because \(H\) is precoherent, but \(F'\to F\) is an epimorphism, so \(F\) is of finite type.

b) Let \(U\) be over \(S\), and \(H'\xrightarrow{f}H|U\) a morphism of \(\underline{C}_U\) with \(H'\in\Ob\underline{C}_{0U}\). Changing notation, we can suppose that \(U=S\). By the quasi-liftability of \(\underline{C}_0\) in \(\underline{C}\), there exists (after localizing) a commutative square
\[
\begin{tikzcd}
G' \arrow[r,"p"] \arrow[d] & H' \arrow[d,"f"]\\
G \arrow[r] & H,
\end{tikzcd}
\]
%%%%%%%%%%  scan idx 127  |  folio 114  %%%%%%%%%%
\SGAPrintPage{114}%
where the horizontal arrows are epimorphisms and \(G'\in\Ob\underline{C}_0\). Since \(\Ker(f)\) is a quotient of \(\Ker(fp)\), it suffices to prove that \(\Ker(fp)\) is of finite type, so we can suppose that \(H'=G'\). The arrows \(F\to G\) and \(H'\to G\) define \(g:F\oplus H'\to G\). Putting \(K=\Ker(G\to H)\), we therefore have a commutative diagram in which the rows are exact and \(q\) is an epimorphism:
\[
\begin{tikzcd}
0\arrow[r]&F\arrow[r]\arrow[d,"q"']&F\oplus H'\arrow[r]\arrow[d,"g"]&H'\arrow[r]\arrow[d,"f"]&0\\
0\arrow[r]&K\arrow[r]&G\arrow[r]&H\arrow[r]&0 .
\end{tikzcd}
\]
The snake diagram shows that \(\Ker(g)\to\Ker(f)\) is an epimorphism. But, since \(F\oplus H'\) is of finite type and \(G\) precoherent, \(\Ker(g)\) is of finite type by a), and we are done.

c) Let us first prove the assertion without the ``resp.'' alternative. Up to a change of notation, we must show that, for a morphism \(g:G'\to G\) of \(\underline{C}_S\) with \(G'\in\Ob\underline{C}_0\), \(\Ker(g)\) is of finite type. Denoting by \(H'\) the image of the composite arrow \(G'\to H\), we have a commutative diagram whose rows are exact
\[
\begin{tikzcd}
0\arrow[r]&F'\arrow[r]\arrow[d,"f"']&G'\arrow[r]\arrow[d,"g"]&H'\arrow[r]\arrow[d,"h"]&0\\
0\arrow[r]&F\arrow[r]&G\arrow[r]&H\arrow[r]&0 .
\end{tikzcd}
\]
Since \(h\) is a monomorphism, the snake diagram gives an isomorphism \(\Ker(f)\qis\Ker(g)\). But \(F'\) is of finite type by a), so (again by a)) \(\Ker(f)\) is of finite type, which proves the assertion without the ``resp.'' alternative. Let us pass to the ``resp.'' assertion (strictly speaking, we cannot invoke (2.5 b)) because we do not assume that \(\underline{C}_0\) is locally quasi-stable under kernels of epimorphisms). After localizing, there exist epimorphisms \(F'\xrightarrow{u}F\), \(H'\xrightarrow{v}H\), with
%%%%%%%%%%  scan idx 128  |  folio 115  %%%%%%%%%%
\SGAPrintPage{115}%
\(F'\in\Ob\underline{C}_0\), \(H'\in\Ob\underline{C}_0\). After localizing further, we can suppose that \(v\) factors through \(G\), whence a commutative diagram with exact rows
\[
\begin{tikzcd}
0\arrow[r]&F'\arrow[r]\arrow[d,"u"']&F'\oplus H'\arrow[r]\arrow[d,"w"]&H'\arrow[r]\arrow[d,dashed,"v"]&0\\
0\arrow[r]&F\arrow[r]&G\arrow[r]&H\arrow[r]&0 .
\end{tikzcd}
\]
Since \(u\) and \(v\) are epimorphisms, so is \(w\), and we are done. This completes the proof of (3.3).

We deduce immediately the

\sgaStableTarget{sga6.e01.cor.3-4}\textbf{Corollary 3.4.} \uline{The (strictly) full subcategory of} \(\underline{C}_S\) \uline{formed by the coherent objects is stable under finite projective (resp.\ inductive) limits. In particular, the kernel, cokernel and image of an arrow between two coherent objects are coherent. Let}
\[
  E^0\longrightarrow E^1\longrightarrow E^2\longrightarrow E^3\longrightarrow E^4
\]
\uline{be an exact sequence of} \(\underline{C}_S\)\uline{. When} \(E^0,E^1,E^3,E^4\) \uline{are coherent,} \(E^2\) \uline{is too.}

\sgaStableTarget{sga6.e01.cor.3-5}\textbf{Corollary 3.5.} \uline{Assume that the objects of} \(\underline{C}_0\) \uline{are} \(\underline{C}_0\)\uline{-coherent (which implies in particular that} \(\underline{C}_0\) \uline{is locally quasi-stable under kernels of epimorphisms,} i.e.\ \uline{that the conditions of} (2.0) \uline{hold}).

a) \uline{Equivalent conditions on} \(F\in\Ob\underline{C}_S\):
\begin{enumerate}[label=(\roman*)]
\item \(F\) \uline{is of finite presentation};
\item \(F\) \uline{is coherent};
\item \(F\) \uline{is of} \(\infty\)\uline{-finite presentation} (2.8).
\end{enumerate}

b) \uline{Let} \(F\in\Ob\Dpx^-(\underline{C}_S)\), \(n\in\mathbb Z\). \uline{We have}:
\[
\begin{split}
F\ n\text{-\uline{pseudo-coherent}}
\quad\Longleftrightarrow\quad&
\Ho^i(F)\text{ \uline{coherent for} }i\geq n+1,\\
&\text{\uline{and} }\Ho^n(F)\text{ \uline{of finite type}.}
\end{split}
\]
%%%%%%%%%%  scan idx 129  |  folio 116  %%%%%%%%%%
\SGAPrintPage{116}%
\uline{In particular, for} \(F\) \uline{to be pseudo-coherent, it is necessary and sufficient that} \(\Ho^n(F)\) \uline{be coherent for every} \(n\)\uline{.}

\uline{Proof.} Assertion a) and the implication \(\Rightarrow\) of b) follow trivially from (3.4). The implication \(\Leftarrow\) of b) follows at once from a) and (2.11.b)).

\sgaStableTarget{sga6.e01.ex.3-6}\textbf{Examples 3.6.} Let \((\underline{S},A)\) be a ringed topos, and \(\underline{C},\underline{C}_0\) the \(\underline{S}\)-categories of (1.3.d)) (cf.\ also (2.15)). For the objects of \(\underline{C}_0\) to be coherent, it is necessary and sufficient that \(A\) be coherent (i.e.\ precoherent). Here are examples where this holds:
\begin{itemize}[label=-]
\item \(A\) is constant with value a left noetherian ring (SGA 4 IX 2.1).
\item \(\underline{S}\) is the topos of sheaves of sets (for the Zariski topology) on a locally noetherian scheme \(S\)\SGAAuthorFootnote{*}{In accordance with the new conventions, we shall say ``scheme'' for ``prescheme'', and ``separated scheme'' for ``scheme''.}, and \(A=\underline{O}_S\).
\item \(\underline{S}\) is the topos of sheaves of sets on a complex analytic space \(S\), \(A=\underline{O}_S\).
\end{itemize}
On the other hand, the sheaf of continuous complex-valued functions on a topological space is in general not coherent.

Of course, when \((\underline{S},A)\) is defined by a ringed space, the notion of coherence considered here coincides with Serre's usual notion (EGA \(0_I\) 5.3.1).

\sgaStableTarget{sga6.e01.prop.3-7}\textbf{Proposition 3.7.} a) \uline{Let} \(f:(\underline{S}',A')\to(\underline{S},A)\) \uline{be a morphism of ringed topoi. Assume that} \(f\) \uline{is flat on the right,} i.e.\ \uline{that the functor} \(f^*=A'\otimes_A.\) \uline{is exact. For} \(F\in\Ob\Dpx(_A S)\)\uline{,} \(G\in\Ob\Dpx^+(_A S)\)\uline{, there is a functorial canonical arrow:}
\[
  f^*\mathbb{R}\underline{\Hom}(F,G)\longrightarrow\mathbb{R}\underline{\Hom}(f^*(F),f^*(G)).
\]
%%%%%%%%%%  scan idx 130  |  folio 117  %%%%%%%%%%
\SGAPrintPage{117}%
\uline{This is an isomorphism when} \(F\) \uline{is pseudo-coherent.}

b) \uline{Suppose that} \(A\) \uline{is commutative and coherent, and let} \(F\in\Ob\Dpx^-(S)\)\uline{,} \(G\in\Ob\Dpx^+(S)\)\uline{. If} \(\underline{\Ho}^i(F)\) and \(\underline{\Ho}^i(G)\) \uline{are coherent for every} \(i\)\uline{, the same is true of} \(\underline{\Ext}^i(F,G)\)\uline{.}

\uline{Proof.} a) To define the arrow (cf.\ SGA 4 XVII), we can suppose that \(G\) has injective components, and take the composite arrow
\[
  f^*\underline{\Hom}(F,G)\longrightarrow\underline{\Hom}(f^*F,f^*G)\longrightarrow\mathbb{R}\underline{\Hom}(f^*F,f^*G).
\]
Let us show that it is an isomorphism for \(F\) pseudo-coherent. For brevity, write \(T\) (resp.\ \(T'\)) for the functor \(f^*\mathbb{R}\underline{\Hom}(.,G)\) (resp.\ \(\mathbb{R}\underline{\Hom}(f^*.,f^*G)\)), and let \(a\in\mathbb Z\) be such that \(\Ho^i(G)=0\) for \(i<a\). It is clear that \(TA\qis T'A\), hence, by dévissage, \(TF\qis T'F\) for \(F\) strictly perfect (2.1). On the other hand, when \(F\) is acyclic in degree \(>n\), \(TF\) and \(T'F\) are acyclic in degree \(<a-n\). To verify that \(TF\qis T'F\) for \(F\) pseudo-coherent, we can, since the question is local, suppose that \(F\) is strictly \(n\)-pseudo-coherent, i.e.\ that there is a distinguished triangle
\[
\begin{tikzcd}[column sep=small]
& F'' \arrow[dl,"+1"'] & \\
F' \arrow[rr] & & F \arrow[ul]
\end{tikzcd}
\;,
\]
with \(F'\) strictly perfect and \(F''\) acyclic in degree \(>n-1\). We deduce a morphism of distinguished triangles
\[
\begin{tikzcd}[column sep=small]
& TF'' \arrow[dr] & \\
TF' \arrow[ur,"-1"'] & & TF \arrow[ll]
\end{tikzcd}
\qquad\longrightarrow\qquad
\begin{tikzcd}[column sep=small]
& T'F'' \arrow[dr] & \\
T'F' \arrow[ur,"-1"'] & & T'F \arrow[ll]
\end{tikzcd}
\;.
\]
Now \(TF'\qis T'F'\) since \(F'\) is strictly perfect, and \(TF''\) and \(T'F''\) are
%%%%%%%%%%  scan idx 131  |  folio 118  %%%%%%%%%%
\SGAPrintPage{118}%
acyclic in degree \(\leq a-n\); writing the induced morphism on the exact cohomology sequences of the preceding triangles, we deduce that \(TF\to T'F\) induces an isomorphism \(\Ho^i(TF)\qis\Ho^i(T'F)\) for \(i\leq a-n\). Since we can choose \(n\) arbitrarily negative, we are done.

b) The proof is left as an exercise to the reader: see, for example, ([H] II 3.3).

\sgaStableTarget{sga6.e01.rem.3-7-1}\textbf{Remark 3.7.1.} We shall see below (7.1.2) another case of ``compatibility'' of \(\mathbb{R}\underline{\Hom}\) with \(\mathbb{L}f^*\).

\sgaStableTarget{sga6.e01.cor.3-7-2}\textbf{Corollary 3.7.2.} (cf.\ EGA \(0_I\) 5.2.7). \uline{Let} \((\underline{S},A)\) \uline{be a ringed topos,} \(s\) \uline{a point of} \(\underline{S}\) (SGA 4 VIII 7.8)\uline{,} \(E,F\in\Ob\D^b(_A S)_{\coh}\)\uline{. If} \(E_s\simeq F_s\)\uline{, there exists an} \(s\)\uline{-pointed object} \(U\) \uline{of} \(\underline{S}\) \uline{such that} \(E|U\simeq F|U\)\uline{.}

\uline{Proof.} Let \(\alpha:E_s\to F_s\) and \(\beta:F_s\to E_s\) be reciprocal isomorphisms. By (3.7 a)), we have
\[
  \underline{\Hom}(E,F)_s\qis\Hom(E_s,F_s)
  \qquad
  (\text{resp. }\underline{\Hom}(F,E)_s\qis\Hom(F_s,E_s)).
\]
There are therefore a ``neighbourhood'' \(U\) of \(s\) and a section \(u\) (resp.\ \(v\)) of \(\underline{\Hom}(E,F)\) (resp.\ \(\underline{\Hom}(F,E)\)) over \(U\) such that \(u_s=\alpha\) (resp.\ \(v_s=\beta\)). The homomorphisms \(vu:E|U\to E|U\) and \(uv:F|U\to F|U\) induce the identity at \(s\), hence over a suitable \(V\) over \(U\), as required.

%%%%%%%%%%  scan idx 132  |  folio 119  %%%%%%%%%%
\SGAPrintPage{119}%
\sgaStableTarget{sga6.e01.sec.4}\section*{4. Perfect complexes}

\sgaStableTarget{sga6.e01.para.4-0}4.0. In this section, \(\underline{S}\) denotes a site, \(S\) an object of \(\underline{S}\), \(\underline{C}\) a flat abelian \(\underline{S}\)-category, and \(\underline{C}_0\) a strictly full additive sub-\(\underline{S}\)-category of \(\underline{C}\) (1.1). Unless otherwise mentioned, we assume that \(\underline{C}_0\) is locally liftable in \(\underline{C}\) and locally stable under kernels of epimorphisms (1.2).

\sgaStableTarget{sga6.e01.lem.4-1}\textbf{Lemma 4.1}\SGAAuthorFootnote{raisedstar}{This lemma was communicated to the editor by D. Ferrand.}. \uline{Let} \(u:E\to F\) \uline{be an arrow of} \(\Cpx(\underline{C}_S)\)\uline{. If} \(E\) \uline{is strictly} \(\underline{C}_0\)\uline{-perfect} (2.1) \uline{and} \(F\) \uline{is acyclic, then} \(u\) \uline{is locally homotopic to zero.}

\uline{Proof.} We must construct locally a homotopy \(k\) such that \(u=dk+kd\). Since \(E\) is strictly perfect, there exist \(a,b\in\mathbb Z\), with \(a\leq b\), such that \(E^i=0\) for \(i\notin[a,b]\). We can therefore take \(k^i=0\) for \(i\notin[a,b]\). We then construct \(k^i\) for \(i\geq a\) by descending induction on \(i\). Suppose that \(k^i:E^i\to F^{i-1}\) has been constructed so that \(u^i=k^{i+1}d_E^i+d_F^{i-1}k^i\) for \(i\geq n>a\), and let us show that, after localizing, there exists \(k^{n-1}:E^{n-1}\to F^{n-2}\) such that \(u^{n-1}=k^nd_E^{n-1}+d_F^{n-2}k^{n-1}\). By the induction hypothesis, the arrow
\[
  v^{n-1}=u^{n-1}-k^nd_E^{n-1}:E^{n-1}\to F^{n-1}
\]
factors through \(Z^{n-1}(F)\). But \(Z^{n-1}(F)=B^{n-1}(F)\) since \(F\) is acyclic. Since \(\underline{C}_0\) is locally liftable in \(\underline{C}\), it follows that, after localizing, there exists \(k^{n-1}:E^{n-1}\to F^{n-2}\) such that \(v^{n-1}=d_F^{n-2}k^{n-1}\), as required.

\sgaStableTarget{sga6.e01.cor.4-2}\textbf{Corollary 4.2.} \uline{Let}
\[
\begin{tikzcd}
& P \arrow[d,"v"] \\
E \arrow[r,"f"] & F
\end{tikzcd}
\]
%%%%%%%%%%  scan idx 133  |  folio 120  %%%%%%%%%%
\SGAPrintPage{120}%
\uline{be a diagram of} \(\Cpx(\underline{C}_S)\)\uline{, where} \(P\) \uline{is strictly} \(\underline{C}_0\)\uline{-perfect and} \(f\) \uline{a quasi-isomorphism. The set of} \(X\) \uline{over} \(S\) \uline{for which there exists an arrow} \(u:P_X\to E_X\) \uline{of} \(\Cpx(\underline{C}_X)\) \uline{such that the diagram}
\[
\begin{tikzcd}
& P_X \arrow[d,"v_X"] \arrow[dl,"u"']\\
E_X \arrow[r,"f_X"] & F_X
\end{tikzcd}
\]
\uline{is commutative in} \(\Kpx(\underline{C}_X)\) \uline{is a refinement of} \(S\)\uline{.}

\uline{Proof.} Let \(G\) be the cone of \(f\). We have the exact sequence
\[
  \longrightarrow \Hom_{\Kpx(\underline{C}_S)}(P,E)
  \longrightarrow \Hom_{\Kpx(\underline{C}_S)}(P,F)
  \longrightarrow \Hom_{\Kpx(\underline{C}_S)}(P,G)
  \longrightarrow .
\]
Since \(G\) is acyclic, the image of \(v\) in \(\Hom_{\Kpx(\underline{C}_S)}(P,G)\) is locally zero by (4.1). Thus, after localizing, there exists \(u:P\to E\) such that \(v=fu\) in \(\Kpx(\underline{C}_S)\), as required.

\sgaStableTarget{sga6.e01.cor.4-3}\textbf{Corollary 4.3.} \uline{Let} \(f:E\to F\) \uline{be a quasi-isomorphism of} \(\Cpx(\underline{C}_S)\)\uline{, with} \(F\) \uline{strictly perfect (relative to} \(\underline{C}_0\)\uline{). The set of} \(X\) \uline{over} \(S\) \uline{for which there exists a quasi-isomorphism} \(g:F_X\to E_X\) \uline{such that} \(f_Xg\) \uline{is homotopic to} \(\Id_{F_X}\) \uline{is a refinement of} \(S\)\uline{.}

\uline{Proof.} This is a special case of (4.2).

\sgaStableTarget{sga6.e01.schol.4-4}\textbf{Scholium 4.4.} Let \(F\in\Ob\Kpx(\underline{C}_S)\). We can define (for example using a cleavage of \(\underline{C}\) (SGA 1 VI 7.1)) a category \((\underline{\Qis/F})\) fibered over \(\underline{S}_{/S}\), whose fibre at \(X\) over \(S\) is the category \((\Qis/F_X)^\circ\) (opposite to the category of quasi-isomorphisms with target \(F_X\)). We can then express (4.3) by saying that,
%%%%%%%%%%  scan idx 134  |  folio 121  %%%%%%%%%%
\SGAPrintPage{121}%
when \(F\) is strictly perfect, the arrow \(\Id_F\) is locally cofinal in \((\underline{\Qis/F})\).

\sgaStableTarget{sga6.e01.notation.4-5}\textbf{Notation 4.5.} Let \(\underline{M}\) be a fibered \(\underline{S}\)-category. Let \(X\in\Ob \underline{S}\), and \(E,F\in\Ob\underline{M}_X\). We shall denote by \(\underline{\Hom}_{\underline{M}}(E,F)\) the sheaf on \(\underline{S}_{/X}\) associated to the presheaf of \(\underline{S}\)-morphisms from \(E\) to \(F\) in \(\underline{M}\) (for the definition of this presheaf, see [G] I 2.6; take care that, even when \(\underline{C}\) is a stack ([G]), the corresponding presheaf for \(\underline{M}=\Kpx(\underline{C})\) (or \(\Dpx(\underline{C})\)) is in general not a sheaf).

\sgaStableTarget{sga6.e01.cor.4-6}\textbf{Corollary 4.6.} \uline{Let} \(E,F\in\Ob\Cpx(\underline{C}_S)\)\uline{, with} \(E\) \uline{strictly perfect. The canonical arrow}
\[
  \underline{\Hom}_{\Kpx(\underline{C})}(E,F)
  \longrightarrow
  \underline{\Hom}_{\Dpx(\underline{C})}(E,F)
\]
\uline{is an isomorphism}.

\uline{Proof.} This is an immediate consequence of (4.4) and of the formula
\[
  \Hom_{\Dpx(\underline{C}_X)}(E_X,F_X)
  = \varinjlim_{(\Qis/E_X)^\circ}\Hom_{\Kpx(\underline{C}_X)}(\,.\,,F_X)
\]
for \(X\) over \(S\).

\sgaStableTarget{sga6.e01.def.4-7}\textbf{Definition 4.7.} \uline{Let} \(F\in\Ob\Dpx(\underline{C}_S)\) \uline{and let} \(a,b\in\mathbb Z\) \uline{be such that} \(a\leq b\)\uline{. We shall say that} \(F\) \uline{has} \(\underline{C}_0\)-perfect amplitude contained in the interval \([a,b]\), \uline{and we shall write}
\[
  \underline{C}_0\text{-}\mathrm{parf}\text{-}\amp(F)\subset[a,b],
\]
\uline{when the set of} \(X\) \uline{over} \(S\) \uline{for which there exists a quasi-isomorphism}
%%%%%%%%%%  scan idx 135  |  folio 122  %%%%%%%%%%
\SGAPrintPage{122}%
\(F'\to F\) \uline{of} \(\Cpx(\underline{C}_X)\) \uline{with} \(F'\) \uline{strictly} \(\underline{C}_0\)\uline{-perfect} (2.1)\uline{, and such that} \(F'^i=0\) \uline{for} \(i\notin[a,b]\)\uline{, is a refinement of} \(S\). (\uline{By} (4.3) \uline{this condition is equivalent to the one obtained by replacing the expression ``quasi-isomorphism of} \(\Cpx(\underline{C}_X)\)\uline{'' by ``isomorphism of} \(\Dpx(\underline{C}_X)\)\uline{''.}) \uline{We shall say that} \(F\) \uline{has} finite \(\underline{C}_0\)-perfect amplitude \uline{when there exist integers} \(a\) \uline{and} \(b\) \uline{such that} \(\underline{C}_0\)-\(\mathrm{parf}\)-\(\amp(F)\subset[a,b]\). \uline{We shall say that} \(F\) \uline{is} \(\underline{C}_0\)-perfect \uline{when} \(F\) \uline{locally has finite} \(\underline{C}_0\)\uline{-perfect amplitude}\SGAAuthorFootnote{raisedstar}{Of course, when no confusion can arise, we shall omit \(\underline{C}_0\) from the preceding expressions.}.

\sgaStableTarget{sga6.e01.ex.4-8}\textbf{Example 4.8.} Let \((\underline{S},A)\) be a ringed topos, \(\underline{C}\) the \(\underline{S}\)-category defined by \(U\longmapsto\) (category of left \(A_U\)-Modules), and \(\underline{C}_0\) the sub-\(\underline{S}\)-category defined by \(U\longmapsto\) (category of left \(A_U\)-Modules that are direct factors of free Modules of finite type). By (1.3.1), the pair \((\underline{C},\underline{C}_0)\) satisfies the hypotheses (4.0). An object \(F\) of \(\Dpx(_A \underline{S})\) (cf.\ (2.15)) will be said to have perfect amplitude contained in \([a,b]\) (resp.\ finite perfect amplitude, resp.\ to be perfect) when \(F\) has \(\underline{C}_0\)-perfect amplitude contained in \([a,b]\) (resp.\dots). When the direct factors of free Modules of finite type are locally free (cf.\ (2.15.1)), saying that \(F\) is perfect therefore means that \(F\) is locally isomorphic (in \(\Dpx(\underline{S})\)) to a complex with bounded degrees and components free of finite type.

\sgaStableTarget{sga6.e01.notation.4-9}\textbf{Notation 4.9.} We shall denote by \(\Dpx(\underline{C})_{\underline{C}_0\text{-}\mathrm{parf}}\) or \(\Parf(\underline{C},\underline{C}_0)\) the strictly full sub-\(\underline{S}\)-category of \(\Dpx(\underline{C})\) formed by the \(\underline{C}_0\)-perfect objects. In the case of Example (4.8), we shall write \(\Dpx(_A \underline{S})_{\mathrm{parf}}\), or \(\Parf(_A \underline{S})\), or simply \(\Parf(\underline{S})\) when there is no possible ambiguity about \(A\).

%%%%%%%%%%  scan idx 136  |  folio 123  %%%%%%%%%%
\SGAPrintPage{123}%
\sgaStableTarget{sga6.e01.prop.4-10}\textbf{Proposition 4.10.} \uline{The category} \(\underline{\Parf}(\underline{C},\underline{C}_0)\) \uline{is a triangulated sub-}\(\underline{S}\)\uline{-category of} \(\Dpx(\underline{C})\) (1.1.4)\uline{. The canonical functor}
\[
  \Kpx^b(\underline{C}_0)\longrightarrow\Parf(\underline{C},\underline{C}_0)
\]
(resp.\ \(\Tr(\Kpx^b(\underline{C}_0))\to\Tr(\Parf(\underline{C},\underline{C}_0))\) (1.1.3)) \uline{induces an} \(\underline{S}\)\uline{-equivalence on the associated stacks} ([G] II 2.2.4)\uline{, which means}:
\begin{enumerate}[label=(\roman*)]
\item \uline{for every} \(X\in\Ob \underline{S}\) \uline{and all} \(E,F\in\Ob\Kpx^b(\underline{C}_{0X})\) (resp.\ \(\Tr(\Kpx^b(\underline{C}_{0X}))\)), \uline{the canonical arrow}
\[
  \underline{\Hom}_{\Kpx^b(\underline{C}_0)}(E,F)
  \longrightarrow
  \underline{\Hom}_{\Dpx(\underline{C})}(E,F)
\]
(\uline{resp.}\ \(\underline{\Hom}_{\Tr(\Kpx^b(\underline{C}_0))}(E,F)\to\underline{\Hom}_{\Tr(\Dpx(\underline{C}))}(E,F)\)) \uline{is an isomorphism, and}
\item \uline{every object of} \(\underline{\Parf}(\underline{C},\underline{C}_0)\) (resp.\ \(\Tr(\Parf(\underline{C},\underline{C}_0))\)) \uline{is locally isomorphic in} \(\Dpx(\underline{C})\) \uline{(resp.}\ \(\Tr(\Dpx(\underline{C}))\)) \uline{to an object of} \(\Kpx^b(\underline{C}_0)\) \uline{(resp.}\ \(\Tr(\Kpx^b(\underline{C}_0))\)).
\end{enumerate}

\uline{Proof.} Assertion (i) without the ``resp.'' alternative is a special case of (4.6), and assertion (ii) without the ``resp.'' alternative is true by definition ((4.7) and (4.9)). It is also clear that \(\Parf(\underline{C},\underline{C}_0)\) is stable under translation of degrees. To prove that \(\Parf(\underline{C},\underline{C}_0)\) is a triangulated sub-\(\underline{S}\)-category of \(\Dpx(\underline{C})\), it therefore suffices to show that, when
\sgaStableTarget{sga6.e01.eq.unnamed.occ3}\[
\begin{tikzcd}
& G \arrow[dl,"+1"'] & \\
E \arrow[rr,"u"] && F \arrow[ul]
\end{tikzcd}
\tag{*}
\]
is a distinguished triangle of \(\Dpx(\underline{C})\) such that \(E\) and \(F\) are perfect, \(G\)
%%%%%%%%%%  scan idx 137  |  folio 124  %%%%%%%%%%
\SGAPrintPage{124}%
is perfect. After localizing, we can suppose that \(E\) and \(F\) are objects of \(\Kpx^b(\underline{C}_0)\). After localizing further, we can suppose, by assertion (i) without the ``resp.'' alternative, that \(u\) is an arrow of \(\Kpx^b(\underline{C}_0)\); then \(G\) is perfect because it is isomorphic to the cone of \(u\), which is strictly perfect. Thus \(\Parf(\underline{C},\underline{C}_0)\) is indeed a triangulated sub-\(\underline{S}\)-category of \(\Dpx(\underline{C})\). Moreover, the preceding argument also proves the ``resp.'' assertion (ii). It remains to prove the ``resp.'' assertion (i). Let \(X\in\Ob \underline{S}\), and let \(E=(E_1\xrightarrow{a_1}E_2\xrightarrow{a_2}E_3\xrightarrow[+1]{a_3}E_1)\), \(F=(F_1\xrightarrow{b_1}F_2\xrightarrow{b_2}F_3\xrightarrow[+1]{b_3}F_1)\) be distinguished triangles of \(\Kpx^b(\underline{C}_{0X})\). Let us show that the arrow
\sgaStableTarget{sga6.e01.eq.unnamed.occ4}\[
  \underline{\Hom}_{\Tr(\Kpx^b(\underline{C}_0))}(E,F)
  \longrightarrow
  \underline{\Hom}_{\Tr(\Dpx(\underline{C}))}(E,F)
  \tag{**}
\]
is an isomorphism. By assertion (i) without the ``resp.'' alternative, we already know that \((**)\) is injective. Let \((u_i):E\to F\) be a morphism of triangles in \(\Dpx(\underline{C})\). After localizing, we can suppose, by assertion (i) without the ``resp.'' alternative, that \(u_i\) (\(1\leq i\leq3\)) is an arrow of \(\Kpx^b(\underline{C}_0)\). The arrow \(u_{i+1}a_i-b_iu_i\), for \(1\leq i\leq3\) (with the convention \(u_4=u_1\)), is zero in \(\Dpx(\underline{C})\), hence locally zero in \(\Kpx^b(\underline{C}_0)\) by assertion (i) without the ``resp.'' alternative. In other words, \((u_i)\) locally defines a morphism of triangles in \(\Kpx^b(\underline{C}_0)\). This proves that \((**)\) is surjective and completes the proof of (4.10).

\sgaStableTarget{sga6.e01.rem.4-10-1}\textbf{Remark 4.10.1.} In (4.10), we can replace \(\Parf(\underline{C},\underline{C}_0)\) by \(\Dpx^*(\underline{C})_{\underline{C}_0\text{-}\mathrm{parf}}=\Dpx^*(\underline{C})\cap\Parf(\underline{C},\underline{C}_0)\) \((*=-\text{ or }b)\); nothing in the proof needs to be changed.

\sgaStableTarget{sga6.e01.complement.4-11}\textbf{Complement 4.11.} We can describe more precisely the triangulated structure of \(\Parf(\underline{C},\underline{C}_0)\) in terms of perfect amplitude (4.7). When \(E\in\Ob\Dpx(\underline{C})\) has perfect amplitude contained in \([a,b]\), \(E[n]\) has perfect amplitude contained in \([a-n,b-n]\).

%%%%%%%%%%  scan idx 138  |  folio 125  %%%%%%%%%%
\SGAPrintPage{125}%
Let
\[
\begin{tikzcd}
& G \arrow[dl,"+1"'] & \\
E \arrow[rr] && F \arrow[ul]
\end{tikzcd}
\]
be a distinguished triangle of \(\Dpx(\underline{C})\). We can draw up a table analogous to (2.6), in which the entries indicate perfect amplitudes and each row corresponds to an implication whose conclusion is underlined.
\[
\begin{array}{ccc}
E&F&G\\[0.3em]
{}[a+1,b+1]&{}[a,b]&\underline{[a,b]}\\
{}[a,b]&\underline{[a,b]}&{}[a,b]\\
\underline{[a,b]}&{}[a,b]&{}[a-1,b-1].
\end{array}
\]
From (4.10) we immediately deduce the

\sgaStableTarget{sga6.e01.cor.4-12}\textbf{Corollary 4.12.} \uline{Suppose that} \(\underline{S}\) \uline{is a punctual site. For} \(F\in\Ob\Dpx(\underline{C})\) \uline{to have perfect amplitude (rel.}\ to \(\underline{C}_0\)\uline{) contained in} \([a,b]\)\uline{, it is necessary and sufficient that there exist a quasi-isomorphism} \(F'\to F\)\uline{, with} \(F'\) \uline{strictly perfect and such that} \(F'^i=0\) \uline{for} \(i\notin[a,b]\)\uline{. The canonical functors}
\[
  \Kpx^b(\underline{C}_0)\longrightarrow\Parf(\underline{C},\underline{C}_0)
\]
\uline{and}
\[
  \Tr(\Kpx^b(\underline{C}_0))\longrightarrow\Tr(\Parf(\underline{C},\underline{C}_0))
\]
\uline{are equivalences of categories}.

\sgaStableTarget{sga6.e01.rem.4-12-1}\textbf{Remark 4.12.1.} When \(\underline{S}\) is a punctual site, the objects of \(\underline{C}_0\) are projective. When \(\underline{C}_0\) contains all projective objects of \(\underline{C}\), one says ``projective amplitude'' instead of ``\(\underline{C}_0\)-perfect amplitude''.

% Additional notation for Expose I continuation
\providecommand{\Cpx}{\mathrm C}
\providecommand{\Kpx}{\mathrm K}
\providecommand{\Dpx}{\mathrm D}
\providecommand{\Ho}{\mathrm H}
\providecommand{\RHom}{\mathrm R\!\operatorname{Hom}}
\renewcommand{\Ob}{\operatorname{ob}}
\providecommand{\qis}{\xrightarrow{\sim}}
\providecommand{\RHom}{\mathrm R\!\Hom}
\providecommand{\Ext}{\operatorname{Ext}}
\providecommand{\Hom}{\operatorname{Hom}}
\providecommand{\Tr}{\operatorname{Tr}}

\providecommand{\Add}{\operatorname{Add}}
\providecommand{\Cart}{\underline{\operatorname{Cart}}}
\providecommand{\perfamp}{\operatorname{parf.amp}}
\providecommand{\perfdim}{\operatorname{parf.dim}}
\providecommand{\toramp}{\operatorname{tor.amp}}
\providecommand{\torf}{\operatorname{torf}}
\providecommand{\injf}{\operatorname{injf}}
\providecommand{\rg}{\operatorname{rg}}
\providecommand{\is}{\mathrm{is}}
\providecommand{\Mod}{\operatorname{Mod}}
\sgaStableTarget{sga6.e01.expose.occ3}%

%%%%%%%%%%  scan idx 139  |  folio 126  %%%%%%%%%%
\SGAPrintPage{126}%
\sgaStableTarget{sga6.e01.rem.4-12-2}\textbf{Remark 4.12.2.} Let \(\underline{C}\) be an abelian category (which can be regarded as fibered over a punctual site), and let \(\underline{C}_0\) be a strictly full additive subcategory. Suppose that \(\underline{C}_0\) is stable under kernels of epimorphisms and \uline{quasi-liftable} in \(\underline{C}\). Then an acyclic object of \(\Kpx^b(\underline{C}_0)\) is in general not zero, so there is no hope that the functor
\(\Kpx^b(\underline{C}_0)\longrightarrow \Dpx(\underline{C})\)
is faithful. However, when \(\Kpx^{b,\varnothing}(\underline{C}_0)\) denotes the thick subcategory of \(\Kpx^b(\underline{C}_0)\) formed by the acyclic objects (in \(\underline{C}\)), the canonical functor
\[
  \Kpx^b(\underline{C}_0)/\Kpx^{b,\varnothing}(\underline{C}_0)
  \longrightarrow \Dpx(\underline{C})
\]
is fully faithful.

Hints: We already know (2.7) that the functor
\[
  \Kpx^-(\underline{C}_0)/\Kpx^{-,\varnothing}(\underline{C}_0)
  \longrightarrow \Dpx^-(\underline{C})
\]
is fully faithful. Show, for example using the criterion ([V] 4.2.b)(ii)'), that the functor
\[
  \Kpx^b(\underline{C}_0)/\Kpx^{b,\varnothing}(\underline{C}_0)
  \longrightarrow
  \Kpx^-(\underline{C}_0)/\Kpx^{-,\varnothing}(\underline{C}_0)
\]
is fully faithful.

\sgaStableTarget{sga6.e01.lem.4-13}\textbf{Lemma 4.13.} \uline{Let} \(F\in\Ob\Dpx(\underline{C}_S)\)\uline{,} \(a,b,c\in\mathbb Z\) \uline{be such that} \(a\leq b\) \uline{and} \(a\leq c\)\uline{. If} \(F\) \uline{has perfect amplitude contained in} \([a,c]\) (4.7) \uline{and} \(\Ho^i(F)=0\) \uline{for} \(i>b\)\uline{, then} \(F\) \uline{has perfect amplitude contained in} \([a,d]\)\uline{, where} \(d=\inf(b,c)\)\uline{.}

\uline{Proof.} We can suppose that \(b\leq c\), and that \(F\) is strictly perfect with \(F^i=0\) for \(i\notin[a,c]\). Then \(F\) is isomorphic to its truncation \(F'\) defined by
%%%%%%%%%%  scan idx 140  |  folio 127  %%%%%%%%%%
\SGAPrintPage{127}%
\(F'^i=F^i\) for \(i<b\), \(F'^b=Z^b(F)\), \(F'^i=0\) for \(i>b\). But \(Z^b(F)\) admits a finite right resolution by objects of \(\underline{C}_0\), so it is locally an object of \(\underline{C}_0\), whence the assertion.

\sgaStableTarget{sga6.e01.cor.4-14}\textbf{Corollary 4.14.} \uline{Let} \(F\in\Ob\underline{C}_S\) \uline{and} \(n\in\mathbb N\)\uline{. Equivalent conditions}:
\begin{enumerate}[label=(\roman*)]
\item \(\underline{\perfamp}(F)\subset[-n,0]\);
\item \(F\) \uline{locally admits a left resolution of length} \(n\) \uline{by objects of} \(\underline{C}_0\)\uline{, i.e.\ the set of} \(X\) \uline{over} \(S\) \uline{for which there exists an exact sequence}
\[
  0\longrightarrow L^{-n}\longrightarrow L^{-n+1}\longrightarrow \ldots
  \longrightarrow L^0\longrightarrow F_X\longrightarrow 0,
\]
\uline{with} \(L^i\in\Ob\underline{C}_{0X}\) \uline{for every} \(i\), \uline{is a refinement of} \(S\).
\end{enumerate}

\uline{Proof.} This is an immediate consequence of (4.7) and (4.13).

\sgaStableTarget{sga6.e01.def.4-14-1}\textbf{Definition 4.14.1.} \uline{We shall say that} \(F\) \uline{has perfect dimension} \(\leq n\)\uline{, and we shall write} \(\perfdim(F)\leq n\)\uline{, when} \(F\) \uline{satisfies the equivalent conditions of} (4.14)\uline{. We shall call the smallest integer} \(n\) \uline{such that} \(\underline{\perfdim}(F)\leq n\) \uline{the} perfect dimension \uline{of} \(F\) \uline{(when there is no such integer, we shall say that} \(F\) \uline{has infinite perfect dimension).}

When \(\underline{S}\) is the punctual site and \(\underline{C}_0\) is the full subcategory of \(\underline{C}\) formed by the projective objects, perfect dimension is the usual projective dimension.

\sgaStableTarget{sga6.e01.prop.4-15}\textbf{Proposition 4.15.} \uline{Let} \(a,b,n\in\mathbb Z\) \uline{(}\(a\leq b\)\uline{,} \(n\geq0\)\uline{), and let} \(F\in\Ob\Dpx(\underline{C})\) \uline{be such that} \(F^i=0\) \uline{(resp.}\ \(\Ho^i(F)=0\)\uline{) for} \(i\notin[a,b]\)\uline{. If, for every} \(i\in[a,b]\)\uline{,} \(F^i\) \uline{(resp.}\ \(\Ho^i(F)\)\uline{) has perfect dimension} \(\leq n+(i-a)\)\uline{, then} \(F\) \uline{has perfect amplitude contained in} \([a-n,b]\)\uline{. Consequently, if for every} \(i\in[a,b]\)\uline{,} \(F^i\) \uline{(resp.}\ \(\Ho^i(F)\)\uline{) is perfect,} \(F\) \uline{is too.}

%%%%%%%%%%  scan idx 141  |  folio 128  %%%%%%%%%%
\SGAPrintPage{128}%
\uline{Proof.} Proceed by induction on \(b\) (with \(a\) and \(n\) fixed), truncating \(F\) and applying (4.11) (cf.\ the proof of (2.11)).

\sgaStableTarget{sga6.e01.lem.4-16}\textbf{Lemma 4.16.} \uline{Let} \(a,b\in\mathbb Z\) \uline{(}\(a\leq b\)\uline{), and} \(F\in\Ob\Dpx(\underline{C})\) \uline{be such that} \(F^i=0\) \uline{for} \(i\notin[a,b]\)\uline{. Assume that} \(\underline{\perfamp}(F)\subset[a,b]\)\uline{, and that} \(F^i\in\Ob\underline{C}_0\) \uline{for} \(i>a\)\uline{. Then} \(F^a\) \uline{is locally an object of} \(\underline{C}_0\)\uline{.}

\uline{Proof.} We have an exact sequence (obtained by truncating \(F\))
\[
  0\longrightarrow F'\longrightarrow F\longrightarrow F^a[-a]\longrightarrow 0,
\]
with \(\perfamp(F')\subset[a+1,b+1]\). From (4.11) and (4.13) we deduce that \(\perfdim(F^a)=0\), i.e.\ that \(F^a\) is locally an object of \(\underline{C}_0\).

\sgaStableTarget{sga6.e01.prop.4-17}\textbf{Proposition 4.17.} \uline{Let} \(F',F''\in\Ob\Dpx(\underline{C}_S)\)\uline{, and} \(a,b\in\mathbb Z\)\uline{, with} \(a\leq b\)\uline{. Consider the conditions}:
\begin{enumerate}[label=(\roman*)]
\item \(\perfamp(F')\subset[a,b]\) and \(\perfamp(F'')\subset[a,b]\);
\item \(\perfamp(F'+F'')\subset[a,b]\).
\end{enumerate}
\uline{We have} (i) \(\Rightarrow\) (ii)\uline{, and, when} \(\underline{C}_0\) \uline{is locally stable under direct factors (}i.e.\ \uline{the relation} \(E'+E''\in\Ob\underline{C}_0\)\uline{, for} \(E',E''\in\Ob\underline{C}_X\) \uline{and} \(X\in\Ob \underline{S}\)\uline{, implies that locally} \(E'\) \uline{and} \(E''\) \uline{are objects of} \(\underline{C}_0\)\uline{), we have} (i) \(\Leftrightarrow\) (ii).

\uline{Proof.} The implication (i) \(\Rightarrow\) (ii) is a special case of (4.11). Let us prove the converse. Put \(F=F'+F''\). Since \(F\) is pseudo-coherent, \(F'\) and \(F''\) are too (2.12). On the other hand, the relation \(\Ho^i(F)=0\) for \(i\notin[a,b]\) implies \(\Ho^i(F')=\Ho^i(F'')=0\) for \(i\notin[a,b]\). By (2.10 a)), we can therefore suppose (after localizing) that \(F'^i\) (resp.\ \(F''^i\)) \(\in\Ob\underline{C}_0\) for \(i\geq a+1\) and \(F'^i=F''^i=0\) for \(i>b\), then, by truncating, that
%%%%%%%%%%  scan idx 142  |  folio 129  %%%%%%%%%%
\SGAPrintPage{129}%
\(F'^i=F''^i=0\) for \(i<a\). Since \(\perfamp(F)\subset[a,b]\), we deduce from (4.16) that \(F^a\) is locally an object of \(\underline{C}_0\). But \(F^a=F'^a+F''^a\), so \(F'^a\) and \(F''^a\) are locally objects of \(\underline{C}_0\), since \(\underline{C}_0\) is locally stable under direct factors.

\sgaStableTarget{sga6.e01.prop.4-18}\textbf{Proposition 4.18.} \uline{Let} \((\underline{S}',\underline{C}',\underline{C}'_0)\) \uline{be a triple satisfying the same hypotheses} (4.0) \uline{as} \((\underline{S},\underline{C},\underline{C}_0)\)\uline{. Let} \(f:\underline{S}'\to \underline{S}\) \uline{be a morphism of sites (defined by} \(f^{-1}:\underline{S}\to \underline{S}'\)\uline{), and let}
\(u:\Dpx^*(\underline{C})\longrightarrow \Dpx^*(\underline{C}')\qquad (*=-,+,\varnothing)\)
\uline{be a functor over} \(f\) (cf.\ (2.13)). \uline{If} \(u\) \uline{takes objects of} \(\underline{C}_0\) \uline{to} \(\underline{C}'_0\)\uline{-perfect complexes, then} \(u\) \uline{defines a functor}
\(\Dpx^*(\underline{C})_{\underline{C}_0\text{-}\mathrm{parf}}\longrightarrow\Dpx^*(\underline{C}')_{\underline{C}'_0\text{-}\mathrm{parf}}\).

\uline{Proof.} Let \(F\in\Ob\Dpx^*(\underline{C})_{\underline{C}_0\text{-}\mathrm{parf}}\); let us show that \(u(F)\in\Ob\Dpx^*(\underline{C}')_{\underline{C}'_0\text{-}\mathrm{parf}}\). Since the question is local, we can suppose that \(F\) is strictly perfect. We need only use dévissage, taking into account that \(\Dpx^*(\underline{C}')_{\underline{C}'_0\text{-}\mathrm{parf}}\) is a triangulated subcategory of \(\Dpx^*(\underline{C}')\) (4.10).

\sgaStableTarget{sga6.e01.rem.4-18-1}\textbf{Remark 4.18.1.} Suppose that there exist \(m,n\in\mathbb Z\) (\(m\leq n\)) such that \(F\in\Ob\underline{C}_0\) implies \(\underline{C}'_0\text{-}\perfamp(u(F))\subset[m,n]\). Then, for \(F\in\Ob\Dpx^*(\underline{C})\), we have
\[
  \underline{C}_0\text{-}\perfamp(F)
  \subset[a,b]
  \quad\Rightarrow\quad
  \underline{C}'_0\text{-}\perfamp(u(F))
  \subset[a+m,b+n].
\]
The proof is left as an exercise to the reader (use (4.11)).

\sgaStableTarget{sga6.e01.cor.4-19}\textbf{Corollary 4.19.} \uline{Under the hypotheses of} (2.16), \uline{if} \(E'\in\Ob\Dpx^-({}_{A'} \underline{S}'_A)\) \uline{is perfect as an object of} \(\Dpx^-({}_{A'} \underline{S}')\) \uline{and} \(E\in\Ob\Dpx^-({}_A \underline{S})\) \uline{is perfect, then}
\(E'\overset{\mathrm L}{\otimes}_A E\) \uline{is perfect}.

\uline{Proof.} We can suppose that \(E'\in\Ob\Dpx^-({}_{A'} S'_A)\) (\(S'\)= final object of \(\underline{S}'\)), and apply (4.18) to
\(u=E'\overset{\mathrm L}{\otimes}_A:\Dpx^-({}_A \underline{S})\longrightarrow \Dpx^-({}_{A'} \underline{S}')\).
In particular:

%%%%%%%%%%  scan idx 143  |  folio 130  %%%%%%%%%%
\SGAPrintPage{130}%
\sgaStableTarget{sga6.e01.cor.4-19-1}\textbf{Corollary 4.19.1.} a) \uline{The functor}
\(\mathbb L f^*:\Dpx^-({}_A \underline{S})\longrightarrow \Dpx^-({}_{A'} \underline{S}')\)
\uline{induces a functor}
\[
  \mathbb L f^*:\Dpx^-({}_A \underline{S})_{\mathrm{parf}}
  \longrightarrow
  \Dpx^-({}_{A'} \underline{S}')_{\mathrm{parf}}.
\]

b) \uline{If} \(A\) \uline{is commutative, the functor}
\(\overset{\mathrm L}{\otimes}_A:\Dpx^-(\underline{S})\times_{\underline{S}}\Dpx^-(\underline{S})\longrightarrow \Dpx^-(\underline{S})\)
\uline{induces a functor}
\[
  \overset{\mathrm L}{\otimes}_A:
  \Dpx^-(\underline{S})_{\mathrm{parf}}\times_{\underline{S}}\Dpx^-(\underline{S})_{\mathrm{parf}}
  \longrightarrow \Dpx^-(\underline{S})_{\mathrm{parf}}.
\]

\sgaStableTarget{sga6.e01.rem.4-19-2}\textbf{Remark 4.19.2.} We can make (4.19) more precise in terms of perfect amplitude:
\[
  \perfamp(E)\subset[a,b]\quad\text{and}\quad
  \perfamp(E')\subset[a',b']
  \quad\Rightarrow\quad
  \perfamp(E'\overset{\mathrm L}{\otimes}_A E)\subset[a+a',b+b'].
\]
In particular:
\[
  \perfamp(E)\subset[a,b]
  \quad\Rightarrow\quad
  \perfamp(\mathbb L f^*(E))\subset[a,b].
\]

\sgaStableTarget{sga6.e01.sec.5}\section*{5. Finite Tor-dimension and perfection}

\sgaStableTarget{sga6.e01.para.5-0}5.0. The reader will surely have asked: what does a pseudo-coherent complex \(E\) lack in order to be perfect? Is it enough, for example, that \(E\) have bounded cohomology? The answer is no. Indeed, take for \(\underline{C}\) the category of modules over a commutative noetherian ring \(A\), and for \(\underline{C}_0\) the full subcategory of projective \(A\)-modules of finite type. Any \(A\)-module of finite type \(E\) is pseudo-coherent and has bounded cohomology, but \(E\) is perfect only when \(E\) also has finite projective dimension (or, equivalently, finite tor-dimension). This example suggests that the difference between a pseudo-coherent complex and a perfect complex lies in a question of cohomological dimension. This is what we shall now examine
%%%%%%%%%%  scan idx 144  |  folio 131  %%%%%%%%%%
\SGAPrintPage{131}%
in the case of the category of Modules of a ringed topos (cf.\ (3.3)). We shall need some preliminary sorites on the notion of finite tor-dimension. We shall be brief, since this question already appears in the literature: cf.\ ([H] II 4) and (SGA 4 XVII 4.1.9.).

\sgaStableTarget{sga6.e01.prop.5-1}\textbf{Proposition 5.1.} \uline{Let} \((\underline{S},A)\) \uline{be a ringed topos,} \(X\in\Ob \underline{S}\)\uline{,} \(F\in\Ob\Dpx^-({}_A X)\)\uline{,} \(m,n\in\mathbb Z\) \uline{be such that} \(m\leq n\)\uline{. Equivalent conditions}:
\begin{enumerate}[label=(\roman*)]
\item \(F\) \uline{is isomorphic, in} \(\Dpx^-({}_A X)\)\uline{, to a complex} \(F'\) \uline{such that} \(F'^i\) \uline{is flat for every} \(i\) \uline{and zero for} \(i\notin[m,n]\);
\item \uline{for every right} \(A_X\)\uline{-Module} \(E\)\uline{, we have} \(\underline{\Ho}^i(E\overset{\mathrm L}{\otimes}_A F)=0\) \uline{for} \(i\notin[m,n]\).
\end{enumerate}

\uline{Proof.} The implication (i) \(\Rightarrow\) (ii) is trivial; the implication (ii) \(\Rightarrow\) (i) follows, by truncation, from the

\sgaStableTarget{sga6.e01.lem.5-1-1}\textbf{Lemma 5.1.1.} (cf.\ (4.16)). \uline{Let} \(F\in\Ob\Dpx^-({}_A X)\)\uline{. Assume that} \(F\) \uline{satisfies hypothesis} (ii) of (5.1)\uline{, that} \(F^i=0\) \uline{for} \(i\notin[m,n]\)\uline{, and that} \(F^i\) \uline{is flat for} \(i\ne m\)\uline{. Then} \(F^m\) \uline{is flat.}

\uline{Proof.} We have a distinguished triangle
\[
\begin{tikzcd}
& F^m[-m] \arrow[dl,"+1"'] & \\
F' \arrow[rr] && F \arrow[ul] & ,
\end{tikzcd}
\]
where \(F'\) has flat components that vanish outside the interval \([m+1,n]\) (suppose \(m<n\)). Apply the functor \(E\overset{\mathrm L}{\otimes}_A.\) for \(E\) an arbitrary right \(A_X\)-Module, and write the exact cohomology sequence.

%%%%%%%%%%  scan idx 145  |  folio 132  %%%%%%%%%%
\SGAPrintPage{132}%
\sgaStableTarget{sga6.e01.def.5-2}\textbf{Definition 5.2.} \uline{An object} \(F\) \uline{of} \(\Dpx^-({}_A \underline{S})\) \uline{satisfying the equivalent conditions of} (5.1) \uline{will be said to have} flat amplitude \uline{contained in the interval} \([m,n]\)\uline{, and we shall write}
\[
  \toramp(F)\subset[m,n].
\]
\uline{We shall say that} \(F\) \uline{has finite flat amplitude (or finite tor-dimension) when there exist integers} \(m\) \uline{and} \(n\) \uline{such that the preceding condition holds.}

\sgaStableTarget{sga6.e01.para.5-3}5.3. Let \(X\in\Ob \underline{S}\). We shall denote by \(\Dpx^-({}_A X)_{\torf}\) the full subcategory of \(\Dpx^-({}_A X)\) formed by the objects of finite tor-dimension. By criterion (ii) of (5.1), it is a triangulated subcategory of \(\Dpx^-({}_A X)\). As \(X\) varies, the categories \(\Dpx^-({}_A X)_{\torf}\) form a triangulated sub-\(\underline{S}\)-category (1.1.4) of \(\Dpx^-({}_A \underline{S})\), which we shall denote by \(\Dpx^-({}_A \underline{S})_{\torf}\). The functor \(\overset{\mathrm L}{\otimes}_A\) induces an exact \(\underline{S}\)-functor
\[
  \overset{\mathrm L}{\otimes}_A:
  \Dpx^b(\underline{S}_A)\times_{\underline{S}} \Dpx({}_A \underline{S})_{\torf}
  \longrightarrow \Dpx^b(\underline{S}_{\mathbb Z}).
\]

\sgaStableTarget{sga6.e01.para.5-4}5.4. For a left \(A\)-Module \(F\) to have flat amplitude contained in \([-n,0]\) (\(n\in\mathbb N\)), it is necessary and sufficient that \(F\) have tor-dimension \(\leq n\) in the usual sense, i.e.\ that \(F\) admit a left resolution of length \(n\) by flat \(A\)-Modules.

%%%%%%%%%%  scan idx 146  |  folio 133  %%%%%%%%%%
\SGAPrintPage{133}%
\sgaStableTarget{sga6.e01.para.5-5}5.5. Let \(X\in\Ob \underline{S}\), and \(m,n\in\mathbb Z\) (\(m\leq n\)). For \(F\in\Ob\Dpx^-({}_A X)\) to have flat amplitude contained in \([m,n]\), it is necessary and sufficient that there exist a covering family \((X_i\to X)\) such that \(F|X_i\) has flat amplitude contained in \([m,n]\). In particular, when \(\underline{S}_{/X}\) is of finite type (SGA 4 VI 1.7), and \(F\) is locally of finite tor-dimension, \(F\) has finite tor-dimension.

\sgaStableTarget{sga6.e01.prop.5-6}\textbf{Proposition 5.6.} \uline{Under the hypotheses of} (2.16)\uline{, if} \(E'\) \uline{has flat amplitude contained in} \([m',n']\) \uline{as an object of} \(\Dpx^-({}_{A'} \underline{S}')\) \uline{and} \(E\) \uline{has flat amplitude contained in} \([m,n]\)\uline{, then} \(E'\overset{\mathrm L}{\otimes}_A E\) \uline{has flat amplitude contained in} \([m+m',n+n']\)\uline{.}

\uline{Proof.} This is immediate: use form (5.1 (i)) of the notion of finite tor-dimension and the fact that, when \(E\) is a flat left \(A\)-Module, \(f^{-1}(E)\) is a flat left \(f^{-1}(A)\)-Module (SGA 4 IV App.\ II).

In particular:

\sgaStableTarget{sga6.e01.cor.5-6-1}\textbf{Corollary 5.6.1.} a) \uline{The functor} \(\mathbb L f^*:\Dpx^-({}_A \underline{S})\longrightarrow \Dpx^-({}_{A'} \underline{S}')\) \uline{induces a functor}
\[
  \mathbb L f^*:\Dpx({}_A \underline{S})_{\torf}\longrightarrow \Dpx({}_{A'} \underline{S}')_{\torf}.
\]

b) \uline{If} \(A\) \uline{is commutative, the functor} \(\overset{\mathrm L}{\otimes}_A:\Dpx^-(\underline{S})\times_{\underline{S}}\Dpx^-(\underline{S})\longrightarrow \Dpx^-(\underline{S})\) \uline{induces a functor}
\[
  \overset{\mathrm L}{\otimes}_A:
  \Dpx(\underline{S})_{\torf}\times_{\underline{S}}\Dpx(\underline{S})_{\torf}
  \longrightarrow \Dpx(\underline{S})_{\torf}.
\]

\sgaStableTarget{sga6.e01.rem.5-6-2}\textbf{Remark 5.6.2.} Let \(E\in\Ob\Dpx^-({}_A S)\) (\(S\)= final object of \(\underline{S}\)), and \(m,n\in\mathbb Z\) (\(m\leq n\)). When \(E\) has flat amplitude contained in \([m,n]\), so does \(E_s\), by (5.6), for every point \(s\) of \(\underline{S}\). The converse is true when \(\underline{S}\) has enough points (use form (5.1.(ii)) of the notion of flat
%%%%%%%%%%  scan idx 147  |  folio 134  %%%%%%%%%%
\SGAPrintPage{134}%
amplitude, and the fact that \(\overset{\mathrm L}{\otimes}_A\) commutes with passage to the fibres).

\sgaStableTarget{sga6.e01.exercise.5-7}\textbf{Exercises 5.7.} a) Suppose that \(A\) is commutative and \((\underline{S}',A')\) faithfully flat over \((\underline{S},A)\) (cf.\ (2.16.2)). For \(E\in\Ob\Dpx^-({}_A S)\)\SGAAuthorFootnote{1}{\(S\) = final object of \(\underline{S}\)} to have flat amplitude contained in \([m,n]\), it is necessary and sufficient that \(f^*(E)\) have that property.

b) Let \(F\in\Ob\Dpx^+({}_A S)\)\SGAAuthorFootMark{1}. Show the equivalence of the conditions:
\begin{enumerate}[label=(\roman*)]
\item \(F\) is isomorphic, in \(\Dpx^+({}_A S)\), to a complex \(F'\) such that \(F'^i\) is injective for every \(i\) and zero for \(i\notin[m,n]\);
\item for every left \(A_S\)-Module \(E\), we have \(\Ext^i(E,F)=0\) for \(i\notin[m,n]\).
\end{enumerate}
When \(F\) satisfies the preceding conditions, we say that \(F\) has injective amplitude contained in \([m,n]\). The full subcategory of \(\Dpx^+({}_A S)\) formed by the objects of finite injective amplitude is a triangulated subcategory, denoted by \(\Dpx^+({}_A S)_{\injf}\). Show that, when \(A\) is commutative, the functor \(\mathbb R\underline{\Hom}:\Dpx(S)^\circ\times\Dpx^+(S)\longrightarrow \Dpx(S)\) induces a functor
\[
  \mathbb R\underline{\Hom}:\Dpx(S)^\circ_{\torf}\times\Dpx(S)_{\injf}
  \longrightarrow \Dpx(S)_{\injf}.
\]

We can now answer the question of (5.0) and state the main result of this section.

\sgaStableTarget{sga6.e01.thm.5-8}\textbf{Theorem 5.8.} \uline{Let} \((\underline{S},A)\) \uline{be a ringed topos,} \(E\in\Ob\Dpx({}_A \underline{S})\)\uline{,} \(m,n\in\mathbb Z\)\uline{, with} \(m\leq n\)\uline{. Equivalent conditions}:
\begin{enumerate}[label=(\roman*)]
\item \(\perfamp(E)\subset[m,n]\) (4.8);
\item \(E\) \uline{is} \((m-1)\)\uline{-pseudo-coherent} (2.15) \uline{and we have} \(\toramp(E)\subset[m,n]\) (5.2).
\end{enumerate}
\uline{Since} ``\uline{perfect}'' \uline{obviously implies} ``\uline{pseudo-coherent}''\uline{, we deduce}:

%%%%%%%%%%  scan idx 148  |  folio 135  %%%%%%%%%%
\SGAPrintPage{135}%
\sgaStableTarget{sga6.e01.cor.5-8-1}\textbf{Corollary 5.8.1.} \uline{For} \(E\) \uline{to be perfect} (4.8)\uline{, it is necessary and sufficient that} \(E\) \uline{be pseudo-coherent} (2.15) \uline{and locally of finite tor-dimension} (5.2)\uline{.}

For the proof, we shall need some lemmas.

\sgaStableTarget{sga6.e01.lem.5-8-2}\textbf{Lemma 5.8.2.} \uline{Let} \(B\) \uline{be a ring of} \(\underline{S}\)\uline{. If} \(E\) \uline{is a left} \(A\)\uline{-Module,} \(G\) \uline{a right} \(B\)\uline{-Module, and} \(F\) \uline{an} \(A\)\uline{-}\(B\)\uline{-Bimodule (on the left over} \(A\)\uline{, on the right over} \(B\)\uline{), there is a functorial canonical homomorphism}
\sgaStableTarget{sga6.e01.eq.5-8-2-1}\begin{equation}\tag{5.8.2.1}
  \underline{\Hom}_B(F,G)\otimes_A E
  \longrightarrow
  \underline{\Hom}_B(\underline{\Hom}_A(E,F),G).
\end{equation}
\uline{This is an isomorphism for} \(G\) \uline{injective and} \(E\) \uline{of finite presentation}.

\uline{Proof.} Let us first define the arrow. By the adjunction formula dear to Cartan, it amounts to defining a homomorphism of \(B\)-Modules:
\sgaStableTarget{sga6.e01.eq.unnamed.occ5}\begin{equation}\tag{*}
  \underline{\Hom}_B(F,G)\otimes_A E\otimes_{\mathbb Z_S}\underline{\Hom}_A(E,F)
  \longrightarrow G\text{\SGAAuthorFootMark{1}}.
\end{equation}
\SGAAuthorFootText{1}{we write \(S\) for the final object of \(\underline{S}\) and \(\mathbb Z_S\) for the constant sheaf with value \(\mathbb Z\).}
Now we have a canonical homomorphism of \(A\)-\(B\)-biModules
\sgaStableTarget{sga6.e01.eq.unnamed.occ6}\begin{equation}\tag{**}
  E\otimes_{\mathbb Z_S}\underline{\Hom}_A(E,F)
  \longrightarrow F
\end{equation}
(coming by adjunction from the identity arrow of \(\underline{\Hom}_A(E,F)\)). Likewise, we have a canonical homomorphism of \(B\)-Modules
\sgaStableTarget{sga6.e01.eq.unnamed.occ7}\begin{equation}\tag{***}
  \underline{\Hom}_B(F,G)\otimes_A F\longrightarrow G.
\end{equation}
We then define (*) as the composite of (**) and (***). To see that the arrow (5.8.2.1) is an isomorphism when \(G\) is injective and \(E\) of finite presentation, it suffices to observe that, when \(G\) is injective, the source and target of the arrow are right-exact functors in \(E\), and that the arrow induces the identity for \(E=A\).

%%%%%%%%%%  scan idx 149  |  folio 136  %%%%%%%%%%
\SGAPrintPage{136}%
Here is now a special case of (5.8).

\sgaStableTarget{sga6.e01.lem.5-8-3}\textbf{Lemma 5.8.3.} \uline{Let} \(E\) \uline{be a left} \(A\)-\uline{Module}. \uline{Equivalent conditions}:
\begin{enumerate}[label=(\roman*)]
\item \(E\) \uline{is locally a direct factor of a free Module of finite type};
\item \(E\) \uline{is flat and of finite presentation}.
\end{enumerate}

\uline{Proof} (after BOURBAKI Alg. Comm. chap.\ I, \S\ 2 Exer 15). It is clear that (i) implies (ii). Let us prove (ii) \(\Rightarrow\) (i). By (1.3.1), we must show that the functor \(\underline{\Hom}_A(E,.)\) is exact. Let \(F\longrightarrow F''\longrightarrow 0\) be an exact sequence of \(A\)-Modules (on the left). We must show that the sequence
\[
  \underline{\Hom}_A(E,F)\longrightarrow \underline{\Hom}_A(E,F'')\longrightarrow 0
\]
is exact. For this it suffices to prove that, for every injective \(\mathbb Z_S\)-Module \(I\), the sequence obtained by applying the functor \(\underline{\Hom}_{\mathbb Z}(.,I)\) is exact. Now we have a commutative diagram, in which the vertical arrows are the canonical arrows of (5.8.2.1):
\[
\begin{array}{ccccccccc}
0&\to&\underline{\Hom}_{\mathbb Z_S}(F'',I)\otimes_A E&\to&\underline{\Hom}_{\mathbb Z_S}(F,I)\otimes_A E\\
&&\downarrow&&\downarrow\\
0&\to&\underline{\Hom}_{\mathbb Z_S}(\underline{\Hom}_A(E,F''),I)&\to&\underline{\Hom}_{\mathbb Z_S}(\underline{\Hom}_A(E,F),I).
\end{array}
\]
The vertical arrows are isomorphisms by (5.8.2). The top row is exact because \(E\) is flat. Thus the bottom row is exact, as required.

\uline{Proof of} (5.8). The implication (i) \(\Rightarrow\) (ii) follows trivially from the definitions. Let us prove (ii) \(\Rightarrow\) (i). Since the question is local, we can suppose that \(E\) is strictly \((m-1)\)-pseudo-coherent. Since, on the other hand, \(E\) is acyclic in degree \(<m\), the
%%%%%%%%%%  scan idx 150  |  folio 137  %%%%%%%%%%
\SGAPrintPage{137}%
truncation arrow
\[
\begin{array}{ccccccc}
&\to&E^{m-1}&\to&E^m&\to&E^{m+1}\to\ldots\\
&&\downarrow&&\downarrow&&\downarrow\\
&\to&0&\to&E^m/B^m&\to&E^{m+1}\to\ldots
\end{array}
\]
is a quasi-isomorphism. Denote the truncated complex by \(E'\). By construction, \(E'\) has flat amplitude \(\subset[m,n]\), and \(E'^i\) is free of finite type for \(i>m\). Thus, by (5.1.1), \(E'^m\) is flat. On the other hand, we have a distinguished triangle
\[
\begin{tikzcd}
& E'^m[-m] \arrow[dl,"+1"'] & \\
E'' \arrow[rr] && E' \arrow[ul] & ,
\end{tikzcd}
\]
where \(E''\) is strictly perfect. By (2.6), \(E'^m[-m]\) is \((m-1)\)-pseudo-coherent, i.e.\ (2.9) \(E'^m\) is of finite presentation. By (5.8.3), we conclude that \(E'^m\) is locally a direct factor of a free Module of finite type, hence that \(\operatorname{parf.amp}(E)\subset[m,n]\).

\sgaStableTarget{sga6.e01.rem.5-8-4}\textbf{Remarks 5.8.4.} a) Using (2.5), we recover the fact that perfect complexes form a triangulated subcategory of \(\Dpx({}_A \underline{S})\).

b) A perfect complex, even with bounded cohomology, does not necessarily have finite tor-dimension. It does, however, when \(\underline{S}\) is of finite type (5.5). Denote by \(\Dpx({}_A \underline{S})_{\underline{\mathrm{parf}}}\) the full sub-\(\underline{S}\)-category of \(\Dpx^b({}_A \underline{S})\) formed by the perfect complexes of finite tor-dimension. It follows at once from (2.16.1 b)) and (5.3) that, when \(A\) is commutative, the functor \(\overset{\mathrm L}{\otimes}_A\) induces a functor
\[
  \overset{\mathrm L}{\otimes}_A:
  \Dpx(\underline{S})_{\underline{\mathrm{parf}}}\times_{\underline{S}} \Dpx^b(\underline{S})_{\mathrm{coh}}
  \longrightarrow \Dpx^b(\underline{S})_{\mathrm{coh}}.
\]
%%%%%%%%%%  scan idx 151  |  folio 138  %%%%%%%%%%
\SGAPrintPage{138}%
This formula is the starting point of the theory of intersections on schemes. Observe, on the other hand, that the (derived) tensor product of two pseudo-coherent complexes with bounded cohomology does not in general have bounded cohomology.

\sgaStableTarget{sga6.e01.cor.5-9}\textbf{Corollary 5.9.} \uline{Suppose that} \(\underline{S}=(\mathrm{Ens})\) \uline{(so} \(A\) ``\uline{is}'' \uline{an ordinary ring) and that every left} \(A\)\uline{-module has finite tor-dimension (this holds, for example, when} \(A\) \uline{has finite global cohomological dimension, see} (MULT VI 2.6) \uline{and} (EGA \(0_{\mathrm{IV}}\) 17.2.8)). \uline{Then we have}
\[
\begin{aligned}
\text{(i)}\quad & \Dpx^b({}_A \underline{S})=\Dpx({}_A \underline{S})_{\torf},\quad\text{\uline{and}}\\
\text{(ii)}\quad & \Dpx^b({}_A \underline{S})_{\mathrm{coh}}=\Dpx({}_A \underline{S})_{\underline{\mathrm{parf}}}.
\end{aligned}
\]
\uline{In other words, for} \(E\in\Ob\Dpx({}_A \underline{S})\) \uline{to have finite tor-dimension (resp.\ finite tor-dimension and be perfect), it is necessary and sufficient that} \(E\) \uline{have bounded cohomology (resp.\ bounded cohomology and be pseudo-coherent)}.

\uline{Proof}. By (5.8), it suffices to prove (i). Let \(E\in\Ob\Dpx^b({}_A \underline{S})\). We can suppose (extending \(E\) by \(0\)) that \(E\in\Ob\Dpx^b({}_A S)\), where \(S\) is the final object of \(\underline{S}\). Proceeding by induction on the number of nonzero cohomology objects of \(E\), we reduce by dévissage to the case where \(E\) is an ordinary \(A\)-Module, and we are done.

More generally, we can state:

\sgaStableTarget{sga6.e01.cor.5-10}\textbf{Corollary 5.10} \uline{Let} \(X\in\Ob \underline{S}\)\uline{. Suppose that} \(\underline{S}_{/X}\) \uline{has enough points} (SGA 4 IV 6.4.1.) \uline{and that, for every point} \(s\) \uline{of} \(\underline{S}_{/X}\)\uline{, every left} \(A_s\)\uline{-module has finite tor-dimension. Then we have}
\[
  \Dpx^b({}_A X)_{\mathrm{coh}}=\Dpx^b({}_A X)_{\mathrm{parf}}.
\]
%%%%%%%%%%  scan idx 152  |  folio 139  %%%%%%%%%%
\SGAPrintPage{139}%
In other words, for \(E\in\Ob\Dpx^b({}_A X)\) to be pseudo-coherent, it is necessary and sufficient that \(E\) be perfect.

\uline{Proof}. Let \(E\in\Ob\Dpx^b({}_A X)_{\mathrm{coh}}\) and \(s\) be a point of \(\underline{S}_{/X}\). We must show that \(E\) is perfect ``in a neighbourhood of \(s\)'', i.e.\ that there exists an \(s\)-pointed object \(U\) of \(\underline{S}\) over \(X\) such that \(E|U\) is perfect. Since the functor \(s^*\) is exact, by (2.16.1 a)) we have \(E_s\in\Ob\Dpx^b(A_s)_{\mathrm{coh}}\). Thus, by (5.9), \(E_s\) is perfect. The assertion will therefore follow from (3.7.2) and the following lemma:

\sgaStableTarget{sga6.e01.lem.5-10-1}\textbf{Lemma 5.10.1.} \uline{Let} \(s\) \uline{be a point of} \(\underline{S}\) \uline{and let} \(F\in\Ob\Dpx(A_s)_{\mathrm{parf}}\)\uline{. There exist an} \(s\)\uline{-pointed object} \(U\) \uline{of} \(\underline{S}\) \uline{and} \(F'\in\Ob\Dpx({}_A U)_{\mathrm{parf}}\) \uline{such that} \(F'_s=F\)\uline{.}

\uline{Proof}. This is evident when \(F\) is a free \(A_s\)-module of finite type. Now suppose that \(F\) is projective of finite type, hence the image of a projector \(p:L\longrightarrow L\), where \(L\) is free of finite type. There then exist a ``neighbourhood'' \(U\) of \(s\), a free \(A_U\)-Module \(L'\), and a homomorphism \(p':L'\to L'\) such that \(p'_s=p\). Since \({p'_s}^2=p'_s\), we can suppose, after shrinking \(U\), that \(p'^2=p'\). Thus \(F'=\Im(p')\) is a direct factor of \(L'\) and satisfies \(F'_s=F\). This generalizes immediately to the case where \(F\) is strictly perfect, i.e.\ a complex with bounded degrees made up of projective modules of finite type.

\sgaStableTarget{sga6.e01.ex.5-11}\textbf{Examples 5.11.} Corollary (5.10) applies in particular when \((\underline{S},A)\) is defined by a ringed space with \uline{regular} local rings, for example a regular scheme, or a complex analytic space smooth over \(\mathbb C\).

%%%%%%%%%%  scan idx 153  |  folio 140  %%%%%%%%%%
\SGAPrintPage{140}%
\sgaStableTarget{sga6.e01.sec.6}\section*{6. Rank of a perfect complex}

\sgaStableTarget{sga6.e01.para.6-0}6.0. Let \((X,\underline{O}_X)\) be a ringed space with local rings. It is well known that to every locally free sheaf of finite type \(L\) on an open set \(U\) of \(X\) we can attach a function \(\rg_U(L)\), locally constant on \(U\) and with integer values, called the \uline{rank} of \(L\), the family of functions \(\rg_U(L)\) being characterized by the following properties:
\begin{enumerate}[label=(\roman*)]
\item for an inclusion \(V\subset U\) of open sets, we have \(\rg_V(L|V)=\rg_U(L)|V\);
\item for every exact sequence of locally free \(\underline{O}_U\)-Modules of finite type
\[
  0\longrightarrow L'\longrightarrow L\longrightarrow L''\longrightarrow 0,
\]
we have \(\rg_U(L)=\rg_U(L')+\rg_U(L'')\);
\item \(\rg_X(\underline{O}_X)=1\).
\end{enumerate}
In this section, we propose to extend the notion of rank to perfect complexes.

\sgaStableTarget{sga6.e01.def.6-1}\textbf{Definition 6.1} (cf.\ SGA 5 VIII 2). \uline{Let} \(C\) \uline{be an additive (resp.\ triangulated) category and} \(F\) \uline{an abelian group. A function} \(f:\Ob(C)\longrightarrow F\) \uline{is called an} additive function from \(C\) to \(F\) \uline{when for all} \(L,L',L''\in\Ob(C)\) \uline{such that} \(L\simeq L'+L''\) \uline{(resp.\ there is a distinguished triangle}
\[
\begin{tikzcd}
& L'' \arrow[dl, dashed] & \\
L' \arrow[rr] && L \arrow[ul]
\end{tikzcd}
\]
) \uline{we have} \(f(L)=f(L')+f(L'')\).

\uline{The additive functions from} \(C\) \uline{to} \(F\) \uline{form an abelian group, denoted by} \(\Add(C,F)\).

%%%%%%%%%%  scan idx 154  |  folio 141  %%%%%%%%%%
\SGAPrintPage{141}%
\sgaStableTarget{sga6.e01.para.6-2}6.2. Let \(C\) be a triangulated category and \(f\) an additive function from \(C\) to an abelian group \(F\). Then \(f\) induces an additive function on the additive category \(C^{\mathrm{ad}}\) underlying \(C\), hence an inclusion \(\Add(C,F)\subset\Add(C^{\mathrm{ad}},F)\). In particular, we have \(f(0)=0\) and \(f(L)=f(M)\) when \(L\simeq M\). On the other hand, we have \(f(L[n])=(-1)^n f(L)\) for every \(L\in\Ob(C)\) and \(n\in\mathbb Z\).

\sgaStableTarget{sga6.e01.para.6-3}6.3. Let \(C\) be an additive (resp.\ triangulated) category. The group \(\Add(C,F)\) depends functorially on \(F\). The functor \(\Add(C,.)\) is representable (cf.\ SGA 5 VIII 2) by an abelian group denoted by \(k(C)\)\SGAAuthorFootnote{raisedstar}{we shall not use here the usual notation \(\Kpx(C)\), since \(\Kpx(C)\) already denotes, when \(C\) is additive, the category of complexes of \(C\) ``up to homotopy''.} endowed with an additive map \(\cl_{C}:\Ob(C)\longrightarrow k(C)\) (when \(C\) is additive, we take for \(k(C)\) the symmetrization of the monoid of isomorphism classes of objects of \(C\); when \(C\) is triangulated, we take for \(k(C)\) the quotient of the free abelian group generated by \(\Ob(C)\) by the relations identifying \(L\) with \(L'+L''\) whenever there is a distinguished triangle \(L'\longrightarrow L\longrightarrow L''\longrightarrow L'[1]\)).

The group \(k(C)\) depends functorially on \(C\) (with respect to additive (resp.\ exact) functors). When \(u,u':C\to D\) are isomorphic additive (resp.\ exact) functors, we have \(k(u)=k(u')\).

When \(C\) is a triangulated category, we have a canonical epimorphism \(k(C^{\mathrm{ad}})\longrightarrow k(C)\) (cf.\ (6.2)).

\sgaStableTarget{sga6.e01.lem.6-4}\textbf{Lemma 6.4.} \uline{Let} \(C\) \uline{be an additive category}. \uline{The canonical functor} \(C\longrightarrow\Kpx^b(C)\) \uline{induces an isomorphism} (\uline{functorial in} \(C\)) \(k(C)\longrightarrow k(\Kpx^b(C))\).

\uline{Proof}. This is an immediate consequence of the following formula (which follows trivially from (6.2)):
\sgaStableTarget{sga6.e01.eq.6-4-1}\begin{equation}\tag{6.4.1}
  f(L)=\sum (-1)^i f(L^i)
\end{equation}
%%%%%%%%%%  scan idx 155  |  folio 142  %%%%%%%%%%
\SGAPrintPage{142}%
for \(L\in\Ob\Kpx^b(C)\) and \(f\) an additive function on \(\Kpx^b(C)\).

\sgaStableTarget{sga6.e01.def.6-5}\textbf{Definition 6.5.} \uline{Let} \(\underline{S}\) \uline{be a category,} \(\underline{C}\) \uline{an additive (resp}.\ \uline{triangulated)} \(\underline{S}\)\uline{-category} (1.1.1)\uline{, and} \(F\) \uline{an abelian presheaf on} \(\underline{S}\). \uline{A family of functions} \(f_X:\Ob(\underline{C}_X)\longrightarrow F(X)\) \uline{indexed by} \(X\in\Ob \underline{S}\) \uline{is called an additive} \(\underline{S}\)\uline{-function from} \(\underline{C}\) \uline{to} \(F\) \uline{when the following conditions are satisfied}:
\begin{enumerate}[label=(\roman*)]
\item \uline{for every} \(X\in\Ob \underline{S}\)\uline{,} \(f_X\) \uline{is an additive function} (6.1);
\item \uline{for every arrow} \(u:X\to Y\) \uline{of} \(\underline{S}\)\uline{, the following diagram is commutative}
\[
\begin{tikzcd}
\Ob\underline{C}_Y \arrow[r,"f_Y"] \arrow[d,"u^*"'] & F(Y) \arrow[d,"F(u)"]\\
\Ob\underline{C}_X \arrow[r,"f_X"] & F(X).
\end{tikzcd}
\]
\end{enumerate}
\uline{The additive} \(\underline{S}\)\uline{-functions from} \(\underline{C}\) \uline{to} \(F\) \uline{form an abelian group, denoted by} \(\Add_{\underline{S}}(\underline{C},F)\).

\sgaStableTarget{sga6.e01.rem.6-6}\textbf{Remarks 6.6.} a) Denote by \(\underline{C}^{\is}\) the fibered \(\underline{S}\)-category whose fiber over \(X\in\Ob \underline{S}\) is the category with the same objects as \(\underline{C}_X\), but whose arrows are the isomorphisms of \(\underline{C}_X\). On the other hand, denote by \(\underline F\) the fibered \(\underline{S}\)-category whose fiber over \(X\in\Ob \underline{S}\) is the discrete category (i.e.\ the only arrows are the identity arrows) \(F(X)\), the inverse-image functor by \(u\in\mathrm{Fl}(\underline{S})\) being \(F(u)\). Then an additive \(\underline{S}\)-function from \(\underline{C}\) to \(F\) is precisely a cartesian \(\underline{S}\)-functor from \(\underline{C}^{\is}\) to \(\underline F\), inducing an additive function in the sense of (6.1) on each fiber. It will be convenient for us to interpret additivity as follows. To fix ideas, suppose that \(\underline{C}\) is triangulated. The category \(\Cart_{\underline{S}}(\underline{C}^{\is},\underline F)\) (resp.\ \(\Cart_{\underline{S}}(\Tr(\underline{C})^{\is},\underline F)\) (1.1.2)) of cartesian \(\underline{S}\)-functors from \(\underline{C}^{\is}\) (resp.\ \(\Tr(\underline{C})^{\is}\)) to \(\underline F\) is the discrete category associated to an abelian
%%%%%%%%%%  scan idx 156  |  folio 143  %%%%%%%%%%
\SGAPrintPage{143}%
group, and we have a homomorphism
\[
  \phi:\Cart_{\underline{S}}(\underline{C}^{\is},\underline F)
  \longrightarrow
  \Cart_{\underline{S}}(\Tr(\underline{C})^{\is},\underline F),
\]
defined by \(\phi(f)(T)=f(L)-f(L')-f(L'')\) for \(f\in\Cart_{\underline{S}}(\underline{C}^{\is},\underline F)\) and \(T=(L'\to L\to L''\to L'[1])\in\Ob\Tr(\underline{C})\). To say that \(f\in\Cart_{\underline{S}}(\underline{C}^{\is},\underline F)\) is additive means that \(\phi(f)=0\); in other words, we have a (canonical functorial) exact sequence
\sgaStableTarget{sga6.e01.eq.6-6-1}\begin{equation}\tag{6.6.1}
  0\longrightarrow \Add_{\underline{S}}(\underline{C},F)
  \longrightarrow \Cart_{\underline{S}}(\underline{C}^{\is},\underline F)
  \xrightarrow{\phi}
  \Cart_{\underline{S}}(\Tr(\underline{C})^{\is},\underline F).
\end{equation}

b) Let \(u:X\to Y\) be an arrow of \(\underline{S}\). The ``inverse-image by \(u\)'' functors, being all isomorphic to one another, define (6.3) a unique homomorphism \(k(Y)\longrightarrow k(X)\). Consequently, \(X\mapsto k(X)\) defines an abelian presheaf on \(\underline{S}\), denoted by \(k(\underline{C})\). An additive \(\underline{S}\)-function from \(\underline{C}\) to \(F\) can then be interpreted simply as a homomorphism of presheaves \(k(\underline{C})\longrightarrow F\); in other words, we have
\[
  \Add_{\underline{S}}(\underline{C},F)=\Hom(k(\underline{C}),F).
\]
The presheaf \(k(\underline{C})\) of course depends functorially on \(\underline{C}\) with respect to additive (resp.\ exact) \(\underline{S}\)-functors.

\sgaStableTarget{sga6.e01.prop.6-7}\textbf{Proposition 6.7.} \uline{Let} \(\underline{S}\) \uline{be a site,} \(\underline{C}\) \uline{a flat abelian} \(\underline{S}\)\uline{-category, and} \(\underline{C}_0\) \uline{a strictly full additive sub-}\(\underline{S}\)\uline{-category of} \(\underline{C}\) \uline{satisfying the hypotheses} (4.0). \uline{Then the morphism of presheaves}
\[
  k(\underline{C}_0)\longrightarrow k(\Parf(\underline{C},\underline{C}_0))\qquad(\text{cf. }(4.9)),
\]
\uline{induced by the inclusion of} \(\underline{C}_0\) \uline{in} \(\underline{\Parf}(\underline{C},\underline{C}_0)\)\uline{, defines an isomorphism on the associated sheaves. In other words, for every abelian sheaf} \(F\) \uline{on} \(\underline{S}\)\uline{,}
%%%%%%%%%%  scan idx 157  |  folio 144  %%%%%%%%%%
\SGAPrintPage{144}%
\uline{the restriction homomorphism}
\[
  \Add_{\underline{S}}(\Parf(\underline{C},\underline{C}_0),F)
  \longrightarrow
  \Add_{\underline{S}}(\underline{C}_0,F)
\]
\uline{is an isomorphism}.

\uline{Proof}. Let \(F\) be an abelian sheaf on \(\underline{S}\); let us show that the arrow
\[
  \Add_{\underline{S}}(\Parf(\underline{C},\underline{C}_0),F)
  \longrightarrow
  \Add_{\underline{S}}(\underline{C}_0,F)
\]
is an isomorphism. It factors as
\[
\Add_{\underline{S}}(\Parf(\underline{C},\underline{C}_0),F)
 \xrightarrow{a}
\Add_{\underline{S}}(\Kpx^b(\underline{C}_0),F)
 \xrightarrow{b}
\Add_{\underline{S}}(\underline{C}_0,F).
\]
Now, by (6.4), the arrow \(b\) is an isomorphism. We are therefore reduced to proving that \(a\) is an isomorphism.

Consider the following commutative diagram, in which the rows are exact (6.6.1), and the vertical arrows are the ``restriction'' arrows:
\[
\begin{array}{ccccccccc}
0&\to&\Add_{\underline{S}}(\Parf(\underline{C}),F)&\to&\Cart_{\underline{S}}(\Parf(\underline{C})^{\is},\underline F)&\to&\Cart_{\underline{S}}(\Tr(\Parf(\underline{C}))^{\is},\underline F)\\
&&\downarrow a&&\downarrow&&\downarrow\\
0&\to&\Add_{\underline{S}}(\Kpx^b(\underline{C}_0),F)&\to&\Cart_{\underline{S}}(\Kpx^b(\underline{C}_0)^{\is},\underline F)&\to&\Cart_{\underline{S}}(\Tr(\Kpx^b(\underline{C}_0))^{\is},\underline F).
\end{array}
\]
Since \(F\) is a sheaf, \(\underline F\) is a stack (G II 2.2.4); thus, by (4.10), the two vertical arrows on the right are isomorphisms.

The same is therefore true of \(a\), q.e.d.

\sgaStableTarget{sga6.e01.rem.6-7-1}\textbf{Remark 6.7.1.} (cf.\ 4.10.1). In (6.7), we can replace \(\Parf(\underline{C},\underline{C}_0)\) by \(\Dpx^*(\underline{C})_{\underline{C}_0\text{-}\mathrm{parf}}\) (\(*=-\) or \(b\)).

%%%%%%%%%%  scan idx 158  |  folio 145  %%%%%%%%%%
\SGAPrintPage{145}%
\sgaStableTarget{sga6.e01.para.6-8}6.8. Let \((\underline{S},A)\) be a ringed topos and \(\underline{C}\) and \(\underline{C}_0\) the fibered \(\underline{S}\)-categories considered in (4.8). We have homomorphisms of abelian presheaves
\[
  k(\underline{C}_0)
  \xrightarrow{\alpha}
  k(\Dpx^-({}_A \underline{S})_{\mathrm{parf}})
  \xrightarrow{\beta}
  k(\Parf({}_A \underline{S})),
\]
which, by (6.7), induce isomorphisms on the associated sheaves. The homomorphism \(\alpha\) is compatible with inverse images, i.e.\ when \(u:(\underline{S}',A')\longrightarrow(\underline{S},A)\) is a morphism of ringed topoi, we have a commutative diagram, in which the vertical arrows are induced by the functor \(\mathbb L u^*\) (cf.\ (4.19.1)):
\[
\begin{tikzcd}
 k(\underline{C}_0) \arrow[r] \arrow[d] & k(\Dpx^-({}_A \underline{S})_{\mathrm{parf}}) \arrow[d]\\
 k(\underline{C}'_0) \arrow[r] & k(\Dpx^-({}_{A'} \underline{S}')_{\mathrm{parf}})
\end{tikzcd}
\]
(\(\underline{C}'_0\) being defined analogously to \(\underline{C}_0\)). When \(A\) is commutative, the tensor product \(\overset{\mathrm L}{\otimes}\) equips \(k(\underline{C}_0)\) and \(k(\Dpx^-({}_A \underline{S})_{\mathrm{parf}})\) with structures of presheaves of rings (cf.\ (4.19.1)), and \(\alpha\) is a homomorphism of presheaves of rings.

\sgaStableTarget{sga6.e01.para.6-9}6.9. In the situation of (6.8), suppose that \((\underline{S},A)\) is locally ringed (cf.\ (2.15.1)). Then (2.15.1), for every \(X\in\Ob \underline{S}\), \(\underline{C}_{0X}\) is the category of locally free \(A_X\)-Modules of finite type. Denote by \(\underline{C}_1\) the fibered \(\underline{S}\)-category defined by \(X\longmapsto\) (the category of free \(A_X\)-Modules of finite type). It is clear that \(\underline{C}_1\) satisfies the hypotheses (4.0). An object of \(\Dpx(\underline{S})\) is perfect if and only if it is perfect relative to \(\underline{C}_1\). We have homomorphisms of abelian presheaves
\[
  k(\underline{C}_1)\longrightarrow k(\underline{C}_0)\longrightarrow
  k(\Dpx^-(\underline{S})_{\mathrm{parf}})\longrightarrow k(\Parf(\underline{S}))
\]
%%%%%%%%%%  scan idx 159  |  folio 146  %%%%%%%%%%
\SGAPrintPage{146}%
(the two on the left being homomorphisms of presheaves of rings), which, by (6.7), induce isomorphisms on the associated sheaves. We shall show that the associated sheaves are in fact canonically isomorphic to the constant sheaf \(\mathbb Z_{\underline{S}}\).

\sgaStableTarget{sga6.e01.lem.6-10}\textbf{Lemma 6.10.} \uline{Let} \((\underline{S},A)\) \uline{be a locally ringed topos,} \(U\) \uline{an object of} \(\underline{S}\) \uline{distinct from the empty object, and} \(p,q\in\mathbb N\)\uline{. If we have} \(A_U^p\simeq A_U^q\)\uline{, then} \(p=q\)\uline{.}

\uline{Proof}. By an argument similar to the one given in the proof of (2.15.1), we reduce to the case where \((\underline{S},A)\) is the locally ringed topos defined by a prescheme, and in this case the assertion is trivial (take a point in \(U\)).

\sgaStableTarget{sga6.e01.def.6-11}\textbf{Definition 6.11.} \uline{We call} rank, \uline{and denote by} \(\rg\), \uline{the additive} \(\underline{S}\)\uline{-function from} \(\underline{C}_1\) \uline{to} \(\mathbb Z_{\underline{S}}\) \uline{defined by}
\[
  \rg(L)=p\quad\text{if } L\simeq A_X^p,
  \qquad X\ne\varnothing,
\]
\[
  \rg(L)=0\quad\text{if } L\in\Ob\underline{C}_{1X},
  \qquad X=\varnothing.
\]
By (6.7), the rank function extends uniquely to an additive \(\underline{S}\)-function from \(\Parf(\underline{S})\) to \(\mathbb Z_{\underline{S}}\) (which we shall again call rank and denote by \(\rg\)).

\sgaStableTarget{sga6.e01.prop.6-12}\textbf{Proposition 6.12.} a) \uline{For} \(L,M\in\Ob\Dpx^-(X)_{\mathrm{parf}}\) (\(X\in\Ob \underline{S}\)), \uline{we have} \(\rg(L\overset{\mathrm L}{\otimes} M)=\rg(L)\,\rg(M)\).

b) \uline{Let} \(u:(\underline{S}',A')\longrightarrow(\underline{S},A)\) \uline{be a morphism of locally ringed topoi. For} \(M\in\Ob\Dpx^-(\underline{S})_{\mathrm{parf}}\), we have \(\rg(\mathbb L u^*(M))=u^*(\rg(M))\). \uline{In other words, the rank homomorphism}: \(k(\Dpx^-(\underline{S})_{\mathrm{parf}})\longrightarrow \mathbb Z_{\underline{S}}\) \uline{is a homomorphism of presheaves of rings, compatible with inverse images}.

% Additional notation for Exposé I
\providecommand{\Cpx}{\mathrm C}
\providecommand{\Kpx}{\mathrm K}
\providecommand{\Dpx}{\mathrm D}
\providecommand{\Ho}{\mathrm H}
\providecommand{\RHom}{\mathrm R\!\operatorname{Hom}}
\renewcommand{\Ob}{\operatorname{ob}}
\providecommand{\Hom}{\operatorname{Hom}}
\providecommand{\RHom}{\mathrm R\!\operatorname{Hom}}
\providecommand{\Ext}{\operatorname{Ext}}
\providecommand{\Tor}{\operatorname{Tor}}
\providecommand{\Tr}{\operatorname{Tr}}
\providecommand{\Add}{\operatorname{Add}}
\providecommand{\Cart}{\underline{\operatorname{Cart}}}
\providecommand{\Parf}{\operatorname{Parf}}
\providecommand{\coh}{\operatorname{coh}}
\providecommand{\amp}{\operatorname{amp}}
\providecommand{\cl}{\operatorname{cl}}
\providecommand{\rg}{\operatorname{rg}}
\providecommand{\Id}{\operatorname{Id}}
\providecommand{\Ker}{\operatorname{Ker}}
\renewcommand{\Im}{\operatorname{Im}}
\providecommand{\cH}{\mathcal H}
\providecommand{\Zs}{\mathbb Z_S}
\providecommand{\A}{\mathcal A}
\providecommand{\cC}{\mathcal C}
\providecommand{\cCzero}{\mathcal C_0}
\providecommand{\otL}{\mathop{\otimes}^{\mathrm L}}
\providecommand{\tamp}{\operatorname{tor.amp}}
\providecommand{\pamp}{\operatorname{parf.amp}}
\providecommand{\pdim}{\operatorname{parf.dim}}
\providecommand{\bon}{\mathrm b}
\providecommand{\parf}{\operatorname{parf}}
\providecommand{\parff}{\operatorname{parf}}
\providecommand{\ac}{\mathrm{ac}}


% Additional notation for Exposé I, completion pages 154--166
\providecommand{\End}{\operatorname{End}}
\providecommand{\tr}{\operatorname{tr}}

%%%%%%%%%%  scan idx 160  |  folio 147  %%%%%%%%%%
\SGAPrintPage{147}%
\uline{Proof}. It suffices to show that the homomorphism induced by rank on the sheaf associated to \(k(\Dpx^-(\underline{S})_{\mathrm{parf}})\) is a homomorphism of sheaves of rings compatible with inverse images. By (6.9), we can replace \(k(\Dpx^-(\underline{S})_{\mathrm{parf}})\) by \(k(\underline{C}_1)\). We are therefore reduced to checking a) and b) for \(L\) and \(M\) free of finite type, which is trivial.

\sgaStableTarget{sga6.e01.para.6-13}6.13. With \((\underline{S},A)\) still denoting a locally ringed topos, let \(X\in\Ob \underline{S}\) and \(\underline n\in H^0(X,\mathbb Z_{\underline{S}})\). We define an object of \(\Parf(X)\), denoted by \(A_X^{\underline n}\), as follows. When \(\underline n\) is constant with value an integer \(n\geq0\) (resp.\ \(n<0\)), we put \(A_X^{\underline n}=A_X^n\) (resp.\ \(A_X^{\underline n}=A_X^{-n}[1]\)); in the general case, we decompose \(X\) into disjoint pieces \(X_i\) such that \(\underline n|X_i\) is constant, and \(A_X^{\underline n}\) is defined as the unique object of \(\Parf(X)\) inducing \(A_{X_i}^{\underline n|X_i}\) on each \(X_i\). (More generally, we can define \(L^{\underline n}\) for \(L\in\Ob\Parf(X)\), by putting \(L^{\underline n}=A_X^{\underline n}\otimes L\).).
The map \(i_X:H^0(X,\mathbb Z_{\underline{S}})\longrightarrow k(\Parf(X))\), defined by \(i_X(\underline n)=\cl(A_X^{\underline n})\), is a ring homomorphism (such that \(i_X(1)=1\)), and defines, as \(X\) varies, a homomorphism of presheaves of rings
\[
  i:\mathbb Z_{\underline{S}}\longrightarrow k(\Parf(\underline{S})).
\]

\sgaStableTarget{sga6.e01.prop.6-14}\textbf{Proposition 6.14.} \uline{Under the hypotheses of} (6.13), \uline{we have}
\[
  \rg\, i=\Id(\mathbb Z_{\underline{S}})
\]
(\uline{in other words,} \(k(\Parf(\underline{S}))\)\uline{, considered as a presheaf of} \(\mathbb Z_{\underline{S}}\)\uline{-algebras by means of} \(i\)\uline{, is augmented toward} \(\mathbb Z_{\underline{S}}\) \uline{by rank}). \uline{Moreover, rank}
\[
  \rg:k(\Parf(\underline{S}))\longrightarrow\mathbb Z_{\underline{S}}
\]
\uline{defines an isomorphism on the associated sheaves}.

%%%%%%%%%%  scan idx 161  |  folio 148  %%%%%%%%%%
\SGAPrintPage{148}%
\uline{Proof}. The first assertion follows trivially from the relations \(\cl(L[1])=-\cl(L)\), \(\cl(L^n)=n\cl(L)\) for \(L\in\Ob\Parf(\underline{S})\). Let us prove the second. First observe that \(i\) factors in the evident way through \(k(\underline{C}_0)\). Since \(\rg\, i\) is the identity, we are reduced (by (6.9)) to proving that the homomorphism \(i\, \rg:k(\underline{C}_0)\longrightarrow k(\underline{C}_0)\) induces an isomorphism on the associated sheaves, i.e.\ that, for \(L\in\Ob\underline{C}_0\), \(i\, \rg(L)\) is locally isomorphic to \(L\), which is immediate.

We deduce, by (6.6.b)),

\sgaStableTarget{sga6.e01.cor.6-15}\textbf{Corollary 6.15.} \uline{Under the hypotheses of} (6.13), \uline{for every abelian sheaf} \(F\) \uline{on} \(\underline{S}\)\uline{, the group of additive} \(\underline{S}\)\uline{-functions from} \(\underline{\Parf}(\underline{S})\) \uline{to} \(F\) \uline{is canonically identified with} \(H^0(S,F)\) (\(S\) = \uline{final object of} \(\underline{S}\)), \uline{the rank function corresponding to} \(F=\mathbb Z_{\underline{S}}\) \uline{and to the unit section of} \(\mathbb Z_{\underline{S}}\).

\sgaStableTarget{sga6.e01.sec.7}\section*{7. Duality of perfect complexes}

\sgaStableTarget{sga6.e01.para.7-0}7.0. In this section, \((\underline{S},A)\) denotes a topos ringed in commutative rings; we denote by \(S\) the final object of \(\underline{S}\). We propose to generalize to perfect complexes the notions of duality well known for locally free sheaves of finite type.

\sgaStableTarget{sga6.e01.prop.7-1}\textbf{Proposition 7.1.} \uline{Let} \(E\in\Ob\Dpx(S)_{\mathrm{parf}}\), \(F\in\Ob\Dpx^+(S)\).

a) \uline{When} \(F\) \uline{has locally bounded cohomology, the same holds for} \(\mathbb{R}\underline{\Hom}(E,F)\). \uline{When} \(F\in\Ob\Dpx(S)\) \uline{is pseudo-coherent (resp.\ locally of finite tor-dimension, resp.\ perfect), the same holds for} \(\mathbb{R}\underline{\Hom}(E,F)\).

%%%%%%%%%%  scan idx 162  |  folio 149  %%%%%%%%%%
\SGAPrintPage{149}%
b) Suppose that \(E\) is moreover of finite tor-dimension (a condition satisfied when \(\underline{S}\) is of finite type). Then, when \(F\) has bounded cohomology (resp.\ finite tor-dimension), the same holds for \(\mathbb{R}\underline{\Hom}(E,F)\). More precisely, let \(a,b,m,n\) be integers such that \(\tamp(E)\subset[a,b]\) (or, equivalently, by (5.8), \(\pamp(E)\subset[a,b]\)) and \(\Ho^i(F)=0\) for \(i\notin[m,n]\) (resp.\ \(\tamp(F)\subset[m,n]\)). Then we have \(\Ho^i(\mathbb{R}\underline{\Hom}(E,F))=0\) for \(i\notin[m-b,n-a]\) (resp.\ \(\tamp(\mathbb{R}\underline{\Hom}(E,F))\subset[m-b,n-a]\)).

\uline{Proof}. Assertion a) is local, so we can suppose that \(E\) is strictly perfect, and we conclude by dévissage on \(E\) (reducing to \(E=A\)). For assertion b), we can suppose that \(E\) is strictly perfect with \(E^i=0\) for \(i\notin[a,b]\), and \(F^i=0\) for \(i\notin[m,n]\) (resp.\ \(F^i\) flat for every \(i\) and \(F^i=0\) for \(i\notin[m,n]\)). It then suffices to apply the following lemma:

\sgaStableTarget{sga6.e01.lem.7-1-1}\textbf{Lemma 7.1.1.} \uline{Let} \(E\in\Ob\Dpx^-(S)\), \(F\in\Ob\Dpx^+(S)\). \uline{When} \(E\) \uline{is strictly pseudo-coherent} (i.e.\ (2.1) \uline{with components direct factors of free Modules of finite type}), \uline{the canonical arrow of} \(\Dpx^+(S)\)
\[
  \underline{\Hom}^\bullet(E,F)\longrightarrow\mathbb{R}\underline{\Hom}(E,F)
\]
\uline{is an isomorphism}.

\uline{Proof}. The spectral sequence
\[
  E_2^{pq}=\underline{\Ho}^p(\underline{\Ext}^q(E,F)^\bullet)\Longrightarrow \underline{\Ext}^*(E,F),
\]
where
\[
  \underline{\Ext}^q(E,F)^i=\coprod_{i_2-i_1=i}\underline{\Ext}^q(E^{i_1},F^{i_2}),
\]
degenerates since \(\underline{\Ext}^q(E,F)^\bullet=0\) for every \(q>0\) (1.3.1). We deduce isomorphisms
%%%%%%%%%%  scan idx 163  |  folio 150  %%%%%%%%%%
\SGAPrintPage{150}%
\[
  E_2^{p0}=\underline{\Ho}^p(\underline{\Hom}^\bullet(E,F))
  \xrightarrow{\sim}\underline{\Ext}^p(E,F)=\underline{\Ho}^p(\mathbb{R}\underline{\Hom}(E,F)),
\]
which proves the assertion.

\sgaStableTarget{sga6.e01.prop.7-1-2}\textbf{Proposition 7.1.2.} \uline{Let} \(f:(\underline{S}',A')\longrightarrow(\underline{S},A)\) \uline{be a morphism of ringed topoi. Let} \(E\in\Ob\Dpx^-(S)\) \uline{and} \(F\in\Ob\Dpx^b(S)\) \uline{be such that} \(\mathbb{R}\underline{\Hom}(E,F)\in\Ob\Dpx^b(S)\) \uline{and} \(\mathbb{L}f^*(F)\in\Ob\Dpx^b(S')\). \uline{There exists a canonical functorial arrow}
\[
  \mathbb{L}f^*\mathbb{R}\underline{\Hom}(E,F)
  \longrightarrow
  \mathbb{R}\underline{\Hom}(\mathbb{L}f^*E,\mathbb{L}f^*F).
\]
\uline{It is an isomorphism when} \(E\) \uline{is perfect}.

\uline{Proof}. By the adjunction of the functors \(\mathbb{L}f^*\) and \(\mathbb{R}f_*\), it suffices to define an arrow
\[
  \mathbb{R}\underline{\Hom}(E,F)
  \longrightarrow
  \mathbb{R}f_*\mathbb{R}\underline{\Hom}(\mathbb{L}f^*E,\mathbb{L}f^*F).
\]
For this we use the adjunction arrow \(F\to\mathbb{R}f_*\mathbb{L}f^*F\) and the duality isomorphism
\[
  \mathbb{R}f_*\mathbb{R}\underline{\Hom}(\mathbb{L}f^*E,\mathbb{L}f^*F)
  \simeq
  \mathbb{R}\underline{\Hom}(E,\mathbb{R}f_*\mathbb{L}f^*F).
\]
To see that the arrow of (7.1.2) is an isomorphism for \(E\) perfect, we reduce trivially, by localization and dévissage, to the case \(E=A\), in which it is a tautology.

\sgaStableTarget{sga6.e01.prop.7-2}\textbf{Proposition 7.2.} \uline{There exists} (cf.\ SGA 5 I 1.6), \uline{for} \(E\in\Ob\Dpx(S)\) \uline{and} \(F\in\Ob\Dpx^+(S)\)\uline{, a canonical functorial arrow}
\[
  E\longrightarrow\mathbb{R}\underline{\Hom}(\mathbb{R}\underline{\Hom}(E,F),F).
\]
\uline{It is an isomorphism when} \(E\in\Ob\Dpx(S)_{\mathrm{parf}}\) and \(F=A\).

%%%%%%%%%%  scan idx 164  |  folio 151  %%%%%%%%%%
\SGAPrintPage{151}%
\uline{Proof}. We define the arrow by resolving \(F\) injectively and using the canonical arrow of \(\Kpx(S)\)
\[
  E\longrightarrow\underline{\Hom}^\bullet(\underline{\Hom}^\bullet(E,F),F).
\]
To see that the arrow in question is an isomorphism when \(E\) is perfect and \(F=A\), we reduce, by localization and dévissage, to \(E=A\), in which case it is evident.

\sgaStableTarget{sga6.e01.def.7-3}\textbf{Definition 7.3.} \uline{Let} \(E\in\Ob\Dpx(S)\). \uline{We shall put} \(\mathbb{R}\underline{\Hom}(E,A)=E^\vee\), \uline{and say that} \(E^\vee\) \uline{is the} dual \uline{of} \(E\).

By (7.1) and (7.2), when \(E\in\Ob\Dpx(S)\) is perfect (resp.\ perfect and of finite tor-dimension), \(E^\vee\) is also, and the canonical arrow \(E\to E^{\vee\vee}\) is an isomorphism. In other words, the functor \(\mathbb{R}\underline{\Hom}(\quad,A)\) defines an anti-auto-equivalence of the category \(\Dpx(S)_{\mathrm{parf}}\) (resp.\ \(\Dpx(S)_{\underline{\mathrm{parf}}}\) (5.8.4.b))) with itself.

\sgaStableTarget{sga6.e01.lem.7-4}\textbf{Lemma 7.4} (CARTAN isomorphism) (cf.\ ([H] II 5.15)). \uline{For} \(E\in\Ob\Dpx(S)\), \(F\in\Ob\Dpx^-(S)\), \(G\in\Ob\Dpx^+(S)\), \uline{there exist canonical functorial isomorphisms}:
\[
\begin{aligned}
\mathrm R\underline{\Hom}(E\otL F,G)&\xrightarrow{\sim}\mathrm R\underline{\Hom}(E,\mathrm R\underline{\Hom}(F,G)),\\
\RHom(E\otL F,G)&\xrightarrow{\sim}\RHom(E,\mathrm R\underline{\Hom}(F,G)),\\
\Hom(E\otL F,G)&\xrightarrow{\sim}\Hom(E,\mathrm R\underline{\Hom}(F,G)).
\end{aligned}
\]

%%%%%%%%%%  scan idx 165  |  folio 152  %%%%%%%%%%
\SGAPrintPage{152}%
\uline{Proof}: This is standard: resolve \(G\) injectively and \(F\) flatly.

\sgaStableTarget{sga6.e01.lem.7-5}\textbf{Lemma 7.5.} \uline{For} \(E\in\Ob\Dpx^-(S)\), \(F\in\Ob\Dpx^+(S)\), \uline{there exists a canonical functorial homomorphism}
\[
  E\otL\mathrm R\underline{\Hom}(E,F)\longrightarrow F
\]

\sgaStableTarget{sga6.e01.prop.7-6}\textbf{Proposition 7.6} (cf.\ ([H] II 5.13)). \uline{Let} \(E\in\Ob\Dpx(S)\), \(F\in\Ob\Dpx^+(S)\), \(G\in\Ob\Dpx^+(S)\). \uline{Under one of the following hypotheses}
\begin{enumerate}[label=\roman*)]
\item \(G\) \uline{of finite tor-dimension},
\item \(E\in\Ob\Dpx^-(S)\), \(F\) \uline{of finite tor-dimension and} \(G\in\Ob\Dpx^b(S)\),
\end{enumerate}
\uline{there exists a canonical functorial homomorphism}
\[
  \mathrm R\underline{\Hom}(E,F)\otL G
  \longrightarrow
  \mathrm R\underline{\Hom}(E,F\otL G).
\]
\uline{It is an isomorphism when} \(E\) \uline{is perfect}.

\uline{Proof}. Let us define the arrow. Under hypothesis (i), it suffices to resolve \(F\) injectively and \(G\) flatly. Under hypothesis (ii), it is equivalent, by (7.4), to define an arrow
\[
  E\otL\mathbb{R}\underline{\Hom}(E,F)\otL G\longrightarrow F\otL G.
\]
We apply the functor \(\otL G\) to the arrow of (7.5).

To see that the arrow is an isomorphism when \(E\) is perfect, we use the standard method (cf.\ the proof of (7.2)): we reduce, by localization and dévissage, to the case \(E=A\).

%%%%%%%%%%  scan idx 166  |  folio 153  %%%%%%%%%%
\SGAPrintPage{153}%
\sgaStableTarget{sga6.e01.cor.7-7}\textbf{Corollary 7.7.} \uline{There exists, for} \(E\in\Ob\Dpx(S)_{\underline{\mathrm{parf}}}\) and \(F\in\Ob\Dpx^+(S)\), \uline{a functorial isomorphism}
\[
  E^\vee\otL F\xrightarrow{\sim}\mathbb{R}\underline{\Hom}(E,F),
\]
\uline{with the notation of} (7.3).

\uline{Proof}. Apply (7.6), replacing \(F\) and \(G\) respectively by \(A\) and \(F\), and use the fact (7.1) that \(E^\vee\in\Ob\Dpx(S)_{\underline{\mathrm{parf}}}\).

\sgaStableTarget{sga6.e01.cor.7-8}\textbf{Corollary 7.8.} \uline{Let} \(E,F\in\Ob\Dpx(S)_{\underline{\mathrm{parf}}}\). \uline{There exists a natural arrow}
\[
  \mathbb{R}\underline{\Hom}(E,F)\otL\mathbb{R}\underline{\Hom}(F,E)\longrightarrow A,
\]
\uline{which is a} ``\uline{perfect duality}'', i.e.\ \uline{identifies, modulo} (7.4), \uline{each of the complexes} \(\mathbb{R}\underline{\Hom}(E,F)\), \(\mathbb{R}\underline{\Hom}(F,E)\) \uline{with the dual of the other} (\uline{in the sense of} (7.3)).

\uline{Proof}. By (7.5), we have natural pairings \(E^\vee\otL E\to A\), \(F^\vee\otL F\to A\), hence, by tensor product, a natural pairing
\[
  E^\vee\otL F\otL F^\vee\otL E\longrightarrow A,
\]
which, in view of (7.7), gives the stated pairing on the \(\mathbb{R}\underline{\Hom}\). Let us show that it is a perfect duality. We have the canonical isomorphisms
\[
\begin{aligned}
\mathbb{R}\underline{\Hom}(F^\vee\otL E,A)&\simeq\mathbb{R}\underline{\Hom}(F^\vee,E^\vee) &&(7.4)\\
&\simeq F^{\vee\vee}\otL E^\vee &&(7.6)\\
&\simeq F\otL E^\vee &&(7.2).
\end{aligned}
\]
Thus, after checking compatibility, we are done. We could also localize and use dévissage to reduce to \(E=F=A\).

%%%%%%%%%%  scan idx 167  |  folio 154  %%%%%%%%%%
\SGAPrintPage{154}%
\sgaStableTarget{sga6.e01.exercise.7-9}\textbf{Exercise 7.9.} Let \(E,F\in\Ob\Dpx(S)_{\underline{\mathrm{parf}}}\). Show that there exists a canonical isomorphism
\[
  E^\vee\otL F^\vee \xrightarrow{\sim} (E\otL F)^\vee .
\]

\sgaStableTarget{sga6.e01.sec.8}\section*{8. Traces and cup-products}

\sgaStableTarget{sga6.e01.para.8-0}8.0. In this section, we generalize to perfect complexes the notion of the trace of an endomorphism of a projective module of finite type (see SGA 5 III for applications). The hypotheses and notation are those of (7.0).

\sgaStableTarget{sga6.e01.para.8-1}8.1. Let \(E\in\Ob\Dpx(S)_{\underline{\mathrm{parf}}}\) (5.8.4.b)). There exists a natural arrow
\[
  \tau:\mathbb{R}\underline{\Hom}(E,E)\longrightarrow A,
\]
composed of the inverse isomorphism of (7.7), \(\mathbb{R}\underline{\Hom}(E,E)\xrightarrow{\sim}E^\vee\otL E\), and the canonical arrow (7.5) \(E^\vee\otL E\to A\).

Let \(U\in\Ob \underline{S}\). Applying the functor \(H^0(U,.)\) to \(\tau\), and using the fact that \(H^0(U,\mathbb{R}\underline{\Hom}(E,E))\) is identified with \(\Hom_{\Dpx(U)}(E|U,E|U)\), we deduce a homomorphism of \(H^0(U,A)\)-modules,
\[
  \Tr_{E|U}:\Hom_{\Dpx(U)}(E|U,E|U)\longrightarrow H^0(U,A),
\]
which is called the \uline{trace homomorphism}\SGAAuthorFootnote{raisedstar}{The finite tor-dimension hypothesis on \(E\) is in fact superfluous for the definition of \(\Tr_{E|U}\), see (8.5).}.

Before going further, let us give some properties of the trace that follow immediately from the definition.

%%%%%%%%%%  scan idx 168  |  folio 155  %%%%%%%%%%
\SGAPrintPage{155}%
\sgaStableTarget{sga6.e01.para.8-1-1}8.1.1. When \(E=A\), \(\Tr_{A|U}:\Hom(A|U,A|U)\longrightarrow H^0(U,A)\) is precisely the well-known isomorphism. When \(E\) is a direct factor of a free Module of finite type, a homomorphism from \(E\) to \(E\) over \(U\) is locally defined by
\[
  e\longmapsto \sum\langle e,x_i^\vee\rangle x_i,
\]
where \(x_i\) (resp.\ \(x_i^\vee\)) is a local section of \(E\) (resp.\ \(E^\vee\)); the corresponding trace is \(\sum\langle x_i,x_i^\vee\rangle\). We easily deduce that, when \(E\) and \(F\) are direct factors of free Modules of finite type, and \(u:(E+F)|U\longrightarrow(E+F)|U\) is defined by a matrix \(\begin{pmatrix}a&b\\ c&d\end{pmatrix}\), we have
\[
  \Tr_{(E+F)|U}(u)=\Tr_{E|U}(a)+\Tr_{F|U}(d).
\]

\sgaStableTarget{sga6.e01.para.8-1-2}8.1.2. Let \(E\) be a strictly perfect complex, \(U\in\Ob \underline{S}\), and \(u:E|U\longrightarrow E|U\) a homomorphism of complexes, with components \(u^i:E^i|U\to E^i|U\). We then have
\[
  \Tr_{E|U}(u)=\sum(-1)^i\Tr_{E^i|U}(u^i).
\]
Indeed, since the question is local, we can suppose that \(u\) is defined by a section of \(\underline{H}^0(E^\vee\otimes E)\), and even that this section is of the form \(\coprod_i\sum_j x_{ij}^\vee\otimes x_{ij}\), where \(x_{ij}\) (resp.\ \(x_{ij}^\vee\)) is a section of \(E^i\) (resp.\ \(E^{i\vee}\)). Making the arrow (7.5) \(E^\vee\otimes E\to A\) explicit, we find
\[
  \Tr_{E|U}(u)=\sum(-1)^i\langle x_{ij},x_{ij}^\vee\rangle
  =\sum(-1)^i\Tr_{E^i|U}(u^i),
\]
q.e.d.

\sgaStableTarget{sga6.e01.para.8-2}8.2. Let \(E,F\in\Ob\Dpx(S)_{\underline{\mathrm{parf}}}\). We have seen (7.8) that there exists a canonical homomorphism
\[
  \gamma:\mathbb{R}\underline{\Hom}(E,F)\otL\mathbb{R}\underline{\Hom}(F,E)\longrightarrow A.
\]
Let \(U\in\Ob \underline{S}\), and take \(u\in H^0(U,\mathbb{R}\underline{\Hom}(E,F))\) and \(v\in H^0(U,\mathbb{R}\underline{\Hom}(F,E))\).
%%%%%%%%%%  scan idx 169  |  folio 156  %%%%%%%%%%
\SGAPrintPage{156}%
We can interpret \(u\) (resp.\ \(v\)) as a homomorphism, in \(\Dpx(U)\), from \(A|U\) to \(\mathbb{R}\underline{\Hom}(E,F)|U\) (resp.\ \(\mathbb{R}\underline{\Hom}(F,E)|U\)), and form
\[
  u\otL v:A|U\longrightarrow(\mathbb{R}\underline{\Hom}(E,F)\otL\mathbb{R}\underline{\Hom}(F,E))|U.
\]
The map \((u,v)\mapsto\langle u,v\rangle=\gamma\circ(u\otL v)\) defines an \(H^0(U,A)\)-linear map,
\[
  \Hom_{\Dpx(U)}(E|U,F|U)\otimes_{H^0(U,A)}\Hom_{\Dpx(U)}(F|U,E|U)
  \longrightarrow H^0(U,A),
\]
which we shall call the \uline{cup-product}\SGAAuthorFootnote{raisedstar}{the same remark as for the definition of \(\Tr_{E|U}\)}.

The relation between trace and cup-product is expressed in

\sgaStableTarget{sga6.e01.prop.8-3}\textbf{Proposition 8.3.} \uline{In the situation of} (8.2), \uline{for} \(u\in\Hom(E|U,F|U)\) \uline{and} \(v\in\Hom(F|U,E|U)\), \uline{we have}
\[
  \Tr_{E|U}(vu)=\Tr_{F|U}(uv)=\langle u,v\rangle.
\]
\uline{In particular}, \uline{when} \(E=F\), \uline{we have}
\[
  \Tr_{E|U}(u)=\langle u,\Id\rangle=\langle\Id,u\rangle.
\]

\uline{Proof}. For simplicity, we shall suppose that \(U=S\). We therefore have two arrows \(u:A\longrightarrow E^\vee\otL F\) and \(v:A\longrightarrow F^\vee\otL E\), hence
\[
  u\otL v:A\longrightarrow E^\vee\otL F\otL F^\vee\otL E.
\]
On the other hand, we have \(vu:A\to E^\vee\otL E\) and \(uv:A\to F^\vee\otL F\). The proposition follows from the commutativity of the following diagrams (whose verification is ``purely mechanical''):
%%%%%%%%%%  scan idx 170  |  folio 157  %%%%%%%%%%
\SGAPrintPage{157}%
\[
\begin{tikzcd}[column sep=large,row sep=large]
 A \arrow[r,"u\otL v"] \arrow[dr,"vu"'] & E^\vee\otL F\otL F^\vee\otL E \arrow[d,"\Id\otL\tau_F"] \\
 & E^\vee\otL E
\end{tikzcd}
\qquad
\begin{tikzcd}[column sep=large,row sep=large]
 A \arrow[r,"u\otL v"] \arrow[dr,"uv"'] & E^\vee\otL F\otL F^\vee\otL E \arrow[d,"\tau_E\otL\Id"] \\
 & F^\vee\otL F
\end{tikzcd}
\]
\[
\begin{tikzcd}[column sep=
large,row sep=large]
E^\vee\otL F\otL F^\vee\otL E \arrow[r,"\Id\otL\tau_E"] \arrow[d,"\Id\otL\tau_F"'] & F^\vee\otL F \arrow[d,"\tau_F"]\\
E^\vee\otL E \arrow[r,"\tau_E"'] & A .
\end{tikzcd}
\]

\sgaStableTarget{sga6.e01.cor.8-4}\textbf{Corollary 8.4.} \(E,F,U\) \uline{being given as in} (8.2), \uline{let} \(u\in\Hom(E|U,E|U)\), \(s:E|U\to F|U\) \uline{be an isomorphism}, \uline{so that} \(sus^{-1}\in\Hom(F|U,F|U)\). \uline{We then have}
\[
  \Tr_{F|U}(sus^{-1})=\Tr_{E|U}(u).
\]

\uline{Proof}. This is immediate.

\sgaStableTarget{sga6.e01.rem.8-5}\textbf{Remark 8.5.} When \(\underline{C}\) is a fibered \(\underline{S}\)-category, denote by \(E(\underline{C})\) the fibered \(\underline{S}\)-category whose fiber over \(U\in\Ob \underline{S}\) is the category with objects the pairs \((L,u)\), where \(L\in\Ob\underline{C}_U\) and \(u\in\operatorname{End}_{\underline{S}}(L)\), and with arrows from \((L,u)\) to \((M,v)\) the arrows \(f:L\to M\) such that \(fu=vf\). By (8.4), we can regard the trace homomorphism as an \(\underline{S}\)-functor from \(E(\Dpx(\underline{S})_{\underline{\parf}})^{\mathrm{is}}\) to \(\underline A\) (where the exponent \(\mathrm{is}\) means that the only arrows considered are the isomorphisms, and \(\underline A\) is the \(\underline{S}\)-category with discrete fibers defined by \(A\) (cf.\ (6.6.a))). Since every perfect complex is locally of finite tor-dimension, the forgetful \(\underline{S}\)-functor
\[
  E(\Dpx(\underline{S})_{\underline{\parf}})^{\mathrm{is}}
  \longrightarrow
  E(\Dpx(\underline{S})_{\parf})^{\mathrm{is}}
\]
induces an isomorphism on the associated stacks. Consequently, the trace
%%%%%%%%%%  scan idx 171  |  folio 158  %%%%%%%%%%
\SGAPrintPage{158}%
homomorphism extends uniquely to an \(\underline{S}\)-functor from \(E(\Dpx(\underline{S})_{\parf})^{\mathrm{is}}\) to \(\underline A\).

For \(E\in\Ob\Parf(U)\) (\(U\in\Ob \underline{S}\)), we again denote by
\[
  \Tr_E:\Hom(E,E)\longrightarrow H^0(U,A)
\]
the corresponding homomorphism, which is a homomorphism of \(H^0(U,A)\)-modules. We see similarly that the definition of the cup-product (8.2), as well as assertions (8.3) and (8.4), generalize to the case of perfect complexes (not necessarily of finite tor-dimension).

\sgaStableTarget{sga6.e01.prop.8-6}\textbf{Proposition 8.6.} \uline{Let} \(U\in\Ob \underline{S}\), \(E,F\in\Parf(U)\), \(u\in\Hom(E,E)\), \(v\in\Hom(F,F)\), \uline{so that} \(u\otL v\in\Hom(E\otL F,E\otL F)\). \uline{We have}
\[
  \Tr_{E\otL F}(u\otL v)=\Tr_E(u)\,\Tr_F(v).
\]

\uline{Proof}. Since the question is local, by (8.4) and (4.10) we can suppose that \(u\) and \(v\) are arrows of \(K^b(\underline{C}_0)\). Using (8.1.2), we reduce to the case where \(u\) and \(v\) are arrows of \(\underline{C}_0\). We conclude by (8.1.1).

%%%%%%%%%%  scan idx 172  |  folio 159  %%%%%%%%%%
\SGAPrintPage{159}%
\sgaStableTarget{sga6.e01.bibliography}\section*{BIBLIOGRAPHY}
\begin{description}[leftmargin=2.5cm,style=nextline]
\item[{[H]}]
Hartshorne, R. Residues and Duality, Lecture Notes no.\ 20, Springer, 1966.
\item[{[V]}] Verdier, J-L. Catégories dérivées, I.H.E.S., 1963.
\item[{[G]}] Giraud, J. Cohomologie non abélienne, to appear in North Holl. Pub. Comp.
\item[{SGA 4}] Cohomologie étale des schémas, by M. Artin, A. Grothendieck, and J-L. Verdier, to appear in North Holl. Pub. Comp.
\item[{SGA 5}] Cohomologie \(\ell\)-adique et fonctions L, by A. Grothendieck, to appear in North Holl. Pub. Comp.
\item[{[1]}] Hakim, M. Topos annelés et schémas relatifs, thesis (in preparation).
\item[{[2]}] Verdier, J-L. Les catégories dérivées des catégories abéliennes, to appear in North Holl. Pub. Comp.
\item[{[MULT]}] Cartan, H. and Eilenberg, S. Homological Algebra, Princeton, 1956.
\end{description}

\providecommand{\cS}{\mathcal S}
\providecommand{\cCz}{\mathcal C_0}
\providecommand{\cCone}{\mathcal C_1}
\providecommand{\cA}{\mathcal A}
\providecommand{\cB}{\mathcal B}
\providecommand{\cL}{\mathcal L}
\providecommand{\cI}{\mathcal I}
\providecommand{\cJ}{\mathcal J}
\providecommand{\Db}{\mathrm D}
\providecommand{\Kb}{\mathrm K}
\providecommand{\Dminus}{\mathrm D^{-}}
\providecommand{\Kminus}{\mathrm K^{-}}
\providecommand{\Kbd}{\mathrm K^{b}}
\providecommand{\Mod}{\operatorname{Mod}}
\providecommand{\Lib}{\operatorname{Lib}}
\providecommand{\Loclib}{\operatorname{Loclib}}
\providecommand{\Qcoh}{\operatorname{Qcoh}}
\providecommand{\Coh}{\operatorname{Coh}}
\providecommand{\Ext}{\operatorname{Ext}}
\providecommand{\Hom}{\operatorname{Hom}}
\renewcommand{\Ob}{\operatorname{ob}}
\providecommand{\im}{\operatorname{Im}}
\providecommand{\coker}{\operatorname{Coker}}
\providecommand{\coh}{\operatorname{coh}}
\providecommand{\parf}{\operatorname{parf}}
\providecommand{\amp}{\operatorname{amp}}
\providecommand{\RHom}{\operatorname{RHom}}

%%%%%%%%%%  scan idx 173  |  folio 160  %%%%%%%%%%
\SGAPrintPage{160}%
\sgaStableTarget{sga6.e02.expose}\section*{\underline{Exposé II}\\[3pt] \underline{\underline{EXISTENCE OF GLOBAL RESOLUTIONS}}}
\begin{center}
\underline{by L. ILLUSIE}
\end{center}

\sgaStableTarget{sga6.e02.sec.1}\section*{1. General globalization criteria}

\sgaStableTarget{sga6.e02.sec.1.0}1.0. In this section, we are given a site $\underline{S}$, a flat abelian $\underline{S}$-category $\underline{C}$ (I 1.1), and a strictly full additive sub-$\underline{S}$-category $\underline{C}_0$ (I 1.1). We assume that $\underline{C}_0$ is locally quasi-liftable in $\underline{C}$ (I 1.2) and locally quasi-stable under kernels of epimorphisms (I 1.2).

\sgaStableTarget{sga6.e02.prop.1-1}\textbf{Proposition 1.1.} \uline{Let} $S\in\Ob \underline{S}$, \uline{and suppose that}:
\begin{enumerate}[label=(\roman*)]
\item $\underline{C}_{0S}$ \uline{is quasi-liftable in} $\underline{C}_S$ (I 1.2) ($\underline{C}_S$ \uline{being regarded as a category fibered over a punctual site});
\item \uline{every object of} $\underline{C}_S$ \uline{which is of} $\underline{C}_0$\uline{-finite type} (I 1.2) (\uline{resp.\ of} $\underline{C}_0$\uline{-finite presentation} (I 2.8)) \uline{is of} $\underline{C}_{0S}$\uline{-finite type}.
\end{enumerate}
\uline{Then}:
\begin{enumerate}[label=\alph*)]
\item $\underline{C}_{0S}$ \uline{is quasi-stable under kernels of epimorphisms} (I 1.2) ($\underline{C}_{0S}$ \uline{being regarded as a category fibered over a punctual site}).
\item \uline{Let} $E\in\Ob \mathrm D^{-}(\underline{C}_S)$ \uline{and} $n\in\mathbb Z$. \uline{If} $E$ \uline{is} $n$-$\underline{C}_0$\uline{-pseudo-coherent} (I 2.3), $E$ \uline{is} $n$-$\underline{C}_{0S}$\uline{-pseudo-coherent} (\uline{resp.} $(n+1)$-$\underline{C}_{0S}$\uline{-pseudo-coherent}). \uline{In other words}, \uline{we have}
\[
\mathrm D^{-}(\underline{C}_S)_{n-\underline{C}_0\text{-coh}}
=
\mathrm D^{-}(\underline{C}_S)_{n-\underline{C}_{0S}\text{-coh}}
\]
(\uline{resp.}
\[
\mathrm D^{-}(\underline{C}_S)_{n-\underline{C}_0\text{-coh}}
\subset
\mathrm D^{-}(\underline{C}_S)_{(n+1)-\underline{C}_{0S}\text{-coh}}
\subset
\mathrm D^{-}(\underline{C}_S)_{(n+1)-\underline{C}_0\text{-coh}},
\]
\uline{inclusions of strictly full subcategories of} $\mathrm D(\underline{C}_S)$, \uline{the inclusion on the right being trivial, moreover}).
\item \uline{For} $E\in\Ob\mathrm D^{-}(\underline{C}_S)$ \uline{to be pseudo-coherent relative to} $\underline{C}_0$ (I 2.3), \uline{it is necessary and sufficient that it be so relative to} $\underline{C}_{0S}$; \uline{in other words}, \uline{we have}
\[
\mathrm D^{-}(\underline{C}_S)_{\underline{C}_0\text{-coh}}
=
\mathrm D^{-}(\underline{C}_S)_{\underline{C}_{0S}\text{-coh}}.
\]
\uline{Moreover}, \uline{the canonical functor}
\[
\mathrm K^{-}(\underline{C}_{0S})/\mathrm K^{-,\varnothing}(\underline{C}_{0S})
\longrightarrow \mathrm D(\underline{C}_S),
\]
\uline{where} $\mathrm K^{-,\varnothing}(\underline{C}_{0S})$ \uline{denotes the thick subcategory of} $\mathrm K^{-}(\underline{C}_{0S})$ \uline{formed by}
%%%%%%%%%%  scan idx 174  |  folio 161  %%%%%%%%%%
\SGAPrintPage{161}%
\uline{acyclic complexes}, \uline{is fully faithful and has as essential image} $\mathrm D^{-}(\underline{C}_S)_{\underline{C}_0\text{-coh}}$. If $\underline{C}_{0S}$ \uline{is liftable in} $\underline{C}_S$ (I 1.2) (\uline{i.e.\ if the objects of} $\underline{C}_{0S}$ \uline{are projective}), \uline{the category} $\mathrm K^{-,\varnothing}(\underline{C}_{0S})$ \uline{is reduced to zero objects}, \uline{and consequently the functor}
\[
\mathrm K^{-}(\underline{C}_{0S})\longrightarrow \mathrm D(\underline{C}_S)
\]
\uline{is fully faithful} (\uline{and has as essential image} $\mathrm D^{-}(\underline{C}_S)_{\underline{C}_0\text{-coh}}$).
\end{enumerate}

\uline{Proof.} a) Let
\[
0\longrightarrow L\longrightarrow L^0\longrightarrow\ldots\longrightarrow L^m\longrightarrow0
\]
be an exact sequence in $\underline{C}_S$, with $L^i\in\Ob\underline{C}_{0S}$ for every $i$; we have to show that $L$ is of $\underline{C}_{0S}$-finite type. Now, by (I 2.2), $L$ is $\underline{C}_0$-pseudo-coherent, hence (I 2.9) of $\underline{C}_0$-finite presentation, and the assertion follows from (ii).

b) Let us show, by descending induction on $i$, that $E$ is $i$-$\underline{C}_{0S}$-pseudo-coherent for $i\ge n$ (resp. $i\ge n+1$). We may already assume that $E$ is strictly $(n+1)$-$\underline{C}_{0S}$-pseudo-coherent (resp. $(n+2)$-$\underline{C}_{0S}$-pseudo-coherent) (I 2.1), and we have to prove that $E$ is $n$-$\underline{C}_{0S}$-pseudo-coherent (resp. $(n+1)$-$\underline{C}_{0S}$-pseudo-coherent). We therefore have an exact sequence
\[
0\longrightarrow E'\longrightarrow E\longrightarrow E''\longrightarrow0,
\]
where $E'$ is strictly $\underline{C}_{0S}$-perfect (I 2.1) and $E''{}^i=0$ for $i>n$ (resp. $i>n+1$). Replacing $E$ by $E''$, we may therefore assume (I 2.6) that $E^i=0$ for $i>n$ (resp. $i>n+1$). Then, by (I 2.10 a)), $H^n(E)$ (resp. $H^{n+1}(E)$) is of $\underline{C}_0$-finite type (resp. of $\underline{C}_0$-finite presentation). Thus, by (ii), $H^n(E)$ (resp. $H^{n+1}(E)$) is of $\underline{C}_{0S}$-finite type. Applying (I 2.10 b)) again, we deduce that $E$ is $n$-$\underline{C}_{0S}$-pseudo-coherent (resp. $(n+1)$-$\underline{C}_{0S}$-pseudo-coherent), q.e.d. (N.B.\ hypothesis (i) and assertion a) were used implicitly in the references to (I 2.6 and 2.10)).

c) The first assertion follows from b), the second from (I 2.7); the last is evident.

\sgaStableTarget{sga6.e02.prop.1-2}\textbf{Proposition 1.2.} \uline{Suppose that} $\underline{C}_0$ \uline{is locally liftable in} $\underline{C}$ (I 1.2) \uline{and} \uline{locally stable under kernels of epimorphisms} (I 1.2), \uline{and make the}
%%%%%%%%%%  scan idx 175  |  folio 162  %%%%%%%%%%
\SGAPrintPage{162}%
\uline{hypotheses} (1.1 (i) \uline{and} (ii)
\uline{in the resp.\ case). Denote by} $\underline{C}_{1S}$ \uline{the strictly full subcategory of} $\underline{C}_S$ \uline{formed by objects which are locally in} $\underline{C}_0$\uline{. Then:}
\begin{enumerate}[label=\alph*)]
\item $\underline{C}_{1S}$ \uline{is stable under kernels of epimorphisms} (I 1.2) \uline{and quasi-liftable in} $\underline{C}_S$ (I 1.2) \uline{(as a category fibered over a punctual site).}
\item \uline{Let} $E\in\Ob\mathrm D(\underline{C}_S)$ \uline{and} $[m,n]$ \uline{be an interval in} $\mathbb Z$\uline{. For} $E$ \uline{to have} $\underline{C}_0$\uline{-perfect amplitude contained in} $[m,n]$\uline{, it is necessary and sufficient that} $E$ \uline{be isomorphic, in} $\mathrm D(\underline{C}_S)$\uline{, to a complex} $F$ \uline{such that} $F^i=0$ \uline{for} $i\notin[m,n]$, $F^i\in\Ob\underline{C}_{0S}$ \uline{for} $i>m$\uline{, and} $F^m\in\Ob\underline{C}_{1S}$.
\item \uline{The canonical functor}
\[
\mathrm K^b(\underline{C}_{1S})/\mathrm K^{b,\varnothing}(\underline{C}_{1S})
\longrightarrow \mathrm D(\underline{C}_S)
\]
\uline{is fully faithful and has as essential image the full subcategory of} $\mathrm D(\underline{C}_S)$ \uline{formed by complexes of finite} $\underline{C}_0$\uline{-perfect amplitude.}
\item \uline{If} $\underline{C}_{1S}$ \uline{is not only quasi-liftable but liftable in} $\underline{C}_S$\uline{, one may replace in} b) \uline{the expression \textquotedblleft}$E$ \uline{is isomorphic to} $F$ \uline{in} $\mathrm D(\underline{C}_S)$\uline{\textquotedblright{} by \textquotedblleft{}there exists a quasi-isomorphism} $F\to E$\uline{\textquotedblright, and in} c) \uline{the category} $\mathrm K^{b,\varnothing}(\underline{C}_{1S})$ \uline{is reduced to zero objects.}
\end{enumerate}

\uline{Proof.} a) It is clear that $\underline{C}_{1S}$ is stable under kernels of epimorphisms. On the other hand, the objects of $\underline{C}_{1S}$ are $\underline{C}_0$-pseudo-coherent, hence $\underline{C}_{0S}$-pseudo-coherent by (1.1 c)); the second assertion therefore follows from (I 2.14) (where we take the punctual site for $\underline{S}$).

b) It suffices to prove necessity. We already know, by (1.1 c)), that $E$ is $\underline{C}_{0S}$-pseudo-coherent. Since $E$ is acyclic in degree $>n$, we may therefore assume (I 2.10), after replacing $E$ by an isomorphic complex in $\mathrm D^{-}(\underline{C}_S)$, that $E$ is strictly $(m+1)$-$\underline{C}_{0S}$-pseudo-coherent and that $E^i=0$ for $i>n$. Since $E$ is acyclic in degree $<m$, we may further assume, after replacing $E$ by its truncation
\[
0\longrightarrow E^m/B^m\longrightarrow E^{m+1}\longrightarrow\ldots\longrightarrow E^n\longrightarrow0,
\]
that $E^i=0$ for $i<m$. Then, by (I 4.16), $E^m$ is locally an object of $\underline{C}_0$, hence an object of $\underline{C}_{1S}$, which proves the assertion.

%%%%%%%%%%  scan idx 176  |  folio 163  %%%%%%%%%%
\SGAPrintPage{163}%
c)
In view of (1.2 a)), the assertion follows from (1.2 b)) and (I 4.12.2).

d) The fact that $\mathrm K^{b,\varnothing}(\underline{C}_{1S})$ is reduced to zero objects is evident. The other assertion follows from (I 4.6) (applied over the punctual site).

\sgaStableTarget{sga6.e02.sec.2}\section*{2. Application to certain categories of Modules}
\sgaStableTarget{sga6.e02.sec.2.0}\subsection*{2.0. Notation}
Let $(\underline{S},A)$ be a ringed topos. For $U\in\Ob \underline{S}$, we write $\Mod(U)$ for the category of left $A_U$-Modules, and $\Lib(U)$ (resp. $\Loclib(U)$) for the full subcategory of $\Mod(U)$ formed by free $A_U$-Modules (resp. Modules locally direct factors of free Modules) of finite type.

\sgaStableTarget{sga6.e02.sec.2.1}\subsection*{2.1. Lifting lemmas}

\sgaStableTarget{sga6.e02.lem.2-1-1}\textbf{Lemma 2.1.1} (cf. EGA $0_{\mathrm I}$ 5.2.3). \uline{Let} $(\underline{S},A)$ \uline{be a ringed topos, and} $S\in\Ob \underline{S}$\uline{. We assume that} $S$ \uline{is} quasi-compact\uline{, i.e.\ that every covering family} $(S_i\to S)$ \uline{is dominated by a finite covering family.}
\begin{enumerate}[label=\alph*)]
\item \uline{Let} $F$ \uline{be a left} $A_S$\uline{-Module which is the inductive limit, over a filtered category} $\Lambda$\uline{, of} $A_S$\uline{-Modules} $F_\lambda$\uline{, and let}
\[
u=\varinjlim u_\lambda:F\longrightarrow G
\]
\uline{be an epimorphism, where} $G$ \uline{is an} $A_S$\uline{-Module of finite type. Then there exists} $\lambda\in\Lambda$ \uline{such that} $u_\lambda$ \uline{is an epimorphism.}
\item \uline{Let} $F$ \uline{be a left} $A_S$\uline{-Module} generated by its sections \uline{(i.e.\ such that there exists an epimorphism} $A_S^{(I)}\to F$ \uline{for a suitable index set} $I$\uline{). Then, if} $F$ \uline{is of finite type, there exists an epimorphism} $A_S^p\to F$ \uline{for} $p\in\mathbb N$.
\end{enumerate}
The proof of a) is analogous to that of (loc.\ cit.) and is left as an exercise to the reader, and b) follows immediately from a).

\sgaStableTarget{sga6.e02.lem.2-1-2}\textbf{Lemma 2.1.2.} \uline{Let} $(\underline{S},A)$ \uline{be a ringed topos,} $\underline{C}$ \uline{the fibered} $\underline{S}$\uline{-category of left} $A$\uline{-Modules,} $\underline{C}_0$ \uline{the fibered subcategory of free} $A$\uline{-Modules of finite type} (I 1.3 d)), $S\in\Ob \underline{S}$, $\underline{C}'_S$ \uline{a strictly full additive subcategory of} $\underline{C}_S$ \uline{containing} $\underline{C}_{0S}$ \uline{and stable under kernels, cokernels and extensions. We assume that}
%%%%%%%%%%  scan idx 177  |  folio 164  %%%%%%%%%%
\SGAPrintPage{164}%
\[
H^1(S,F)=0\quad
\text{\uline{for every} }F\in\Ob\underline{C}'_S.
\]
\uline{Then:}
\begin{enumerate}[label=\alph*)]
\item \uline{Every exact sequence in} $\underline{C}'_S$:
\[
0\longrightarrow E\longrightarrow F\longrightarrow G\longrightarrow0,
\]
\uline{such that} $\underline{\Hom}(G,E)
\in\Ob\underline{C}'_S$ \uline{and such that} $G$ \uline{is locally a direct factor of a free Module of finite type, is split. The category} $\underline{C}_{0S}$ \uline{is liftable in} $\underline{C}'_S$ (I 1.2).
\item \uline{Denote by} $\Mod(A(S))$ \uline{the category of} $A(S)$\uline{-modules. The functor}
\[
\Gamma_S:\underline{C}'_S\longrightarrow\Mod(A(S))
\]
\uline{is exact. Moreover:}
\begin{enumerate}[label=(\roman*)]
\item \uline{If every object} $F$ \uline{of} $\underline{C}'_S$ \uline{admits a finite global presentation}
\[
A_S^p\longrightarrow A_S^q\longrightarrow F\longrightarrow0,
\]
\uline{the functor} $\Gamma_S:\underline{C}'_S\longrightarrow\Mod(A(S))$ \uline{is fully faithful and has as essential image the full subcategory of} $\Mod(A(S))$ \uline{formed by} $A(S)$\uline{-modules of finite presentation;}
\item \uline{If} $\underline{C}'_S$ \uline{contains the free} $A_S$\uline{-Modules, if every object} $F$ \uline{of} $\underline{C}'_S$ \uline{admits a global presentation}
\[
A_S^{(I)}\longrightarrow A_S^{(J)}\longrightarrow F\longrightarrow0,
\]
\uline{and if the functor} $\Gamma_S$ \uline{commutes with filtered inductive limits}\SGAAuthorFootnote{1}{a hypothesis satisfied when $\underline{S}$ is locally of finite type (SGA 4 VI 1.2) or when $\underline{S}$ is the topos of sheaves of sets on a compact space and $S$ the final object (T.F II 3.10.1).}\uline{, the functor} $\Gamma_S:\underline{C}'_S\longrightarrow\Mod(A(S))$ \uline{is an equivalence of categories.}
\end{enumerate}
\end{enumerate}

\uline{Proof.} a) Since $G$ is locally a direct factor of a free Module of finite type, the spectral sequence
\[
E_2^{pq}=H^p(S,\underline{\Ext}^q(G,E))\Longrightarrow\Ext^{*}(G,E)
\]
degenerates (I 1.3.1) and gives an isomorphism
\[
H^1(S,\underline{\Hom}(G,E))\simeq\Ext^1(G,E),
\]
%%%%%%%%%%  scan idx 178  |  folio 165  %%%%%%%%%%
\SGAPrintPage{165}%
whence the
first assertion. The second is a particular case, since, for $G\in\Ob\underline{C}_{0S}$, i.e.\ $G$ free and of finite type, and $E\in\Ob\underline{C}'_S$, we have $\underline{\Hom}(G,E)\in\Ob\underline{C}'_S$.

b) The first assertion is trivial. Let us now prove, in case (i) or (ii), the full faithfulness of $\Gamma_S:\underline{C}'_S\longrightarrow\Mod(A(S))$. Let $F,G\in\Ob\underline{C}'_S$, and let
\sgaStableTarget{sga6.e02.eq.unnamed}\[
A_S^{(I)}\longrightarrow A_S^{(J)}\longrightarrow F\longrightarrow0
\tag{$*$}
\]
be a global presentation of $F$ (with $I$ and $J$ finite in case (i)). Applying to $(*)$ the functor $\Gamma_S$, which is exact on $\underline{C}'_S$ and commutes, in case (ii), with filtered inductive limits, we obtain an exact sequence
\sgaStableTarget{sga6.e02.eq.unnamed.occ2}\[
A(S)^{(I)}\longrightarrow A(S)^{(J)}\longrightarrow\Gamma_S(F)\longrightarrow0.
\tag{$**$}
\]
On the other hand, applying to $(*)$ (resp. $(**)$) the left-exact functor $\operatorname{Hom}_{A_S}(.,G)$ (resp. $\operatorname{Hom}_{A(S)}(.,\Gamma_S(G))$), we obtain a commutative diagram with exact rows:
\[
\begin{tikzcd}[column sep=small]
0 \arrow[r] & \operatorname{Hom}_{A_S}(F,G) \arrow[r] \arrow[d] &
\operatorname{Hom}_{A_S}(A_S^{(J)},G) \arrow[r] \arrow[d] &
\operatorname{Hom}_{A_S}(A_S^{(I)},G) \arrow[d] \\
0 \arrow[r] & \operatorname{Hom}_{A(S)}(\Gamma_S(F),\Gamma_S(G)) \arrow[r] &
\operatorname{Hom}_{A(S)}(A(S)^{(J)},\Gamma_S(G)) \arrow[r] &
\operatorname{Hom}_{A(S)}(A(S)^{(I)},\Gamma_S(G)).
\end{tikzcd}
\]
The two vertical arrows on the right are visibly isomorphisms; the same is therefore true of the vertical arrow on the left, which proves the full faithfulness of $\Gamma_S:\underline{C}'_S\longrightarrow\Mod(A(S))$. The characterization of the essential image, in case (i) or (ii), is immediate and is left to the reader.

\sgaStableTarget{sga6.e02.lem.2-1-3}\textbf{Lemma 2.1.3.} \uline{Let} $A$ \uline{and} $B$ \uline{be abelian categories, and} $\phi:A\to B$ \uline{an exact functor.}
\begin{enumerate}[label=\alph*)]
\item \uline{For} $M\in\Ob A$\uline{, the following conditions are equivalent:}
\begin{enumerate}[label=(\roman*)]
\item \uline{for every} $L\in\Ob A$\uline{, the canonical arrow (induced by} $\phi$)
\[
\Ext^n_{A}(L,M)\longrightarrow\Ext^n_{B}(\phi(L),\phi(M))
\]
\uline{is an isomorphism for every} $n$;
\item \uline{for every} $L\in\Ob A$\uline{, the canonical arrow}
%%%%%%%%%%  scan idx 179  |  folio 166  %%%%%%%%%%
\SGAPrintPage{166}%
\[
\operatorname{Hom}_{A}(L,M)\longrightarrow\operatorname{Hom}_{B}(\phi(L),\phi(M))
\]
\uline{is an
isomorphism, and for every} $n\ge1$ \uline{and every} $x\in\Ext^n_{B}(\phi(L),\phi(M))$ \uline{there exists an epimorphism} $L'\to L$ \uline{in} $A$ \uline{which effaces} $x$.
\end{enumerate}
\item \uline{The following conditions are equivalent}:
\begin{enumerate}[label=(\roman*)]
\item \uline{for all} $L,M\in\Ob A$ \uline{and every} $n\in\mathbb Z$\uline{, the canonical arrow}
\[
\Ext^n_{A}(L,M)\longrightarrow\Ext^n_{B}(\phi(L),\phi(M))
\]
\uline{is an isomorphism};
\item \uline{the canonical functor} (\uline{induced by} $\phi$)
\[
\mathrm D^b(A)\longrightarrow\mathrm D^b(B)
\]
\uline{is fully faithful and has as essential image the full subcategory of} $\mathrm D^b(B)$ \uline{formed by complexes} $F$ \uline{such that} $H^i(F)$ \uline{is in the essential image of} $\phi$ \uline{for every} $i$.
\end{enumerate}
\end{enumerate}

\uline{Proof.} a) For the implication (i)$\Rightarrow$(ii), it suffices to show the following: given $n\ge1$, $L\in\Ob A$, and $y\in\Ext^n_{A}(L,M)$, there exists an epimorphism $L'\to L$ which effaces $y$. For this, one can use Yoneda's interpretation ([18] III 3.2): $y$ is the class of an exact sequence
\[
0\longrightarrow M\longrightarrow E_{n-1}\longrightarrow\ldots\longrightarrow E_0\longrightarrow L\longrightarrow0;
\]
the epimorphism $E_0\to L$ effaces $y$ ([18] III 3.2.5 and 3.2.8). Let us prove (ii)$\Rightarrow$(i). For brevity, put
\[
\Ext^n_{A}(.,M)=F^n,\qquad
\Ext^n_{B}(\phi(.),\phi(M))=G^n.
\]
The effaceability hypothesis means that, for $n\ge1$ and $L\in\Ob A$, we have
\[
G^n(L)=\varinjlim\operatorname{Coker}(G^{n-1}(L_0)\to G^{n-1}(L_1)),
\]
the inductive limit being taken over the set of exact sequences
%%%%%%%%%%  scan idx 180  |  folio 167  %%%%%%%%%%
\SGAPrintPage{167}%
$0\to L_1\to L_0\to L\to0$. From
what we saw above, we likewise have
\[
F^n(L)=\varinjlim\operatorname{Coker}(F^{n-1}(L_0)\to F^{n-1}(L_1)).
\]
since $\phi=\phi^0:F^0\to G^0$ is an isomorphism, it follows, by induction on $n$, that the morphism $\phi^n:F^n\to G^n$ induced by $\phi$ is an isomorphism for every $n$.

b) The equivalence (i)$\Leftrightarrow$(ii) follows from an immediate dévissage, left to the reader.

\sgaStableTarget{sga6.e02.rem.2-1-3-1}\textbf{Remarks 2.1.3.1.} There is of course a dual equivalence to a), which we leave the reader to state. On the other hand, when $A$ and $B$ have enough injectives, conditions (i) and (ii) of b) are equivalent to
\begin{enumerate}[label=(\roman*),start=3]
\item for $L\in\Ob\mathrm D^{-}(A)$ and $M\in\Ob\mathrm D^{+}(A)$, the canonical arrow
\[
\underline{\underline{R}}\!\operatorname{Hom}(L,M)\longrightarrow
\underline{\underline{R}}\!\operatorname{Hom}(\phi(L),\phi(M))
\]
is an isomorphism,
\end{enumerate}
as one sees immediately by a spectral-sequence argument.

\sgaStableTarget{sga6.e02.sec.2.2}\subsection*{2.2. Noetherian schemes and divisorial schemes}

\sgaStableTarget{sga6.e02.notation.2-2-1}\textbf{Notation 2.2.1.} Let $S$ be a scheme. We write $\Qcoh(S)$ (resp. $\Coh(S)$) for the (abelian) category of quasi-coherent (resp. coherent) $\underline{O}_S$-Modules. Let $[m,n]$ be an interval in $\mathbb Z$. We shall say that an object $E$ of $\mathrm D(\Qcoh(S))$ is $n$-pseudo-coherent (resp. pseudo-coherent, resp. of perfect amplitude contained in $[m,n]$, resp. perfect) when $E$ is $n$-pseudo-coherent (resp. pseudo-coherent, resp. of perfect amplitude contained in $[m,n]$, resp. perfect) relative to the category fibered over the Zariski site of $S$ defined by $U\mapsto\Loclib(U)$ ((I 2.3), (I 4.7) and (2.0)). We denote by $\mathrm D^{-}(\Qcoh(S))_{\coh}$ (resp. $\mathrm D(\Qcoh(S))_{\parf}$) the full triangulated subcategory of $\mathrm D(\Qcoh(S))$ formed by pseudo-coherent (resp. perfect) complexes.

\sgaStableTarget{sga6.e02.prop.2-2-2}\textbf{Proposition 2.2.2}\SGAAuthorFootnote{raisedstar}{This ancillary result will not be used in the remainder of the Seminar: the reader is advised to skip it on a first reading.}. \uline{Let} $S$ \uline{be a noetherian scheme. The canonical}
\uline{functor}
%%%%%%%%%%  scan idx 181  |  folio 168  %%%%%%%%%%
\SGAPrintPage{168}%
\[
\mathrm D^{-}(\Coh(S))\longrightarrow \mathrm D(\Qcoh(S))
\]
\uline{is fully faithful and has as essential image} $\mathrm D^{-}(\Qcoh(S))_{\coh}$.

\uline{Proof.} We apply (1.1), taking for $\underline{S}$ the Zariski site of $S$, and for $\underline{C}$ (resp. $\underline{C}_0$) the fibered category $U\mapsto\Qcoh(U)$ (resp. $\Coh(U)$). Since $\underline{O}_S$ is coherent, the objects of $\underline{C}_0$ are pseudo-coherent (I 3.5), so an object of $\mathrm D(\Qcoh(U))$ ($U$ = open subset of $S$) is pseudo-coherent (in the sense of (2.2.1)) precisely when it is so relative to $\underline{C}_0$ (I 2.14.1). It therefore suffices, by (1.1 c)), to show that conditions (i) and (ii) in the non-resp.\ case of (1.1) are satisfied. For (ii), this is trivial, since a quasi-coherent Module which is a quotient of a coherent Module is coherent, and the notion of coherence is local. For (i), we use (2.1.1 a)) and the fact that every quasi-coherent $\underline{O}_S$-Module is a filtered inductive limit of coherent $\underline{O}_S$-Modules (EGA I 9.4.9).

Admitting (3.5) and (3.7), which are independent of the present section, one easily deduces:

\sgaStableTarget{sga6.e02.cor.2-2-2-1}\textbf{Corollary 2.2.2.1.} \uline{Let} $S$ \uline{be a noetherian scheme (resp.\ noetherian and of finite dimension). The canonical functor}
\[
\mathrm D^b(\Coh(S))\longrightarrow\mathrm D(S)
\]
\[
(\text{resp. }\mathrm D^{-}(\Coh(S))\longrightarrow\mathrm D(S))
\]
\uline{is fully faithful and has as essential image} $\mathrm D^b(S)_{\coh}$ (resp. $\mathrm D^{-}(S)_{\coh}$) (I 2.15).

\sgaStableTarget{sga6.e02.prop.2-2-3}\textbf{Proposition 2.2.3} (cf. EGA II 4.5.2 and 4.5.5). \uline{Let} $S$ \uline{be a quasi-compact and quasi-separated scheme, and let} $(L_i)_{i\in I}$ \uline{be a family of invertible sheaves on} $S$\uline{. For every} $\underline{O}_S$\uline{-Module} $F$\uline{, and every pair} $(i,n)\in I\times\mathbb Z$\uline{, we denote by} $F_{(i,n)}$ \uline{the sub-}$\underline{O}_S$\uline{-Module of} $F\otimes L_i^{\otimes n}$ \uline{generated by the sections of} $F\otimes L_i^{\otimes n}$ \uline{over} $S$\uline{. The following conditions are equivalent}:
\begin{enumerate}
%%%%%%%%%%  scan idx 182  |  folio 169  %%%%%%%%%%
\item[(i)] \SGAPrintPage{169}%
\uline{The
open sets} $S_f$\uline{, for} $f\in\Gamma(S,L_i^{\otimes n})$\uline{,} $i\in I$ \uline{and} $n\in\mathbb N$\uline{, form a basis for the topology of} $S$.
\item[(i$'$)] \uline{Those of the} $S_f$ \uline{which are affine cover} $S$.
\item[(ii)] \uline{For every quasi-coherent} $\underline{O}_S$\uline{-Module} $F$ \uline{of finite type, there exists a family} $(r_i)_{i\in I}$ (resp. $(n_i)_{i\in I}$) \uline{of integers} $\ge0$ (resp. $>0$), \uline{and an epimorphism}
\[
\coprod_{i\in I}\underline{O}_S^{r_i}\otimes L_i^{\otimes-n_i}\longrightarrow F.
\]
\item[(iii)] \uline{For every} $F\in\Ob\Qcoh(S)$\uline{, we have}
\[
F=\sum_{(i,n)\in I\times\mathbb N^+}F_{(i,n)}\otimes L_i^{\otimes-n}
\]
(\uline{sum (not direct) of subsheaves of} $F$).
\item[(iii$'$)] \uline{The same condition as} (iii), \uline{but with} $F$ \uline{a quasi-coherent Ideal of} $\underline{O}_S$.
\end{enumerate}

\uline{Proof.} It consists of a simple paraphrase of (loc.\ cit.); we give it briefly for the reader's convenience.

(i)$\Rightarrow$(i$'$): For $x\in S$, choose an affine open subset $U$ containing $x$. An open subset $S_f$ such that $x\in S_f\subset U$, with $f\in\Gamma(S,L_i^{\otimes n})$ ($n\ge0$), is necessarily affine, as follows from

\sgaStableTarget{sga6.e02.lem.2-2-3-1}\textbf{Lemma 2.2.3.1.} (EGA II 5.5.8). \uline{Let} $X$ \uline{be a scheme,} $L$ \uline{an invertible} $\underline{O}_X$\uline{-Module,} $f\in\Gamma(X,L)$\uline{. The canonical immersion} $i:X_f\to X$ \uline{is affine}.

\uline{Proof.} This is trivial: the question is local on $X$, so we may assume that $X$ is affine and $L\simeq\underline{O}_X$, and we are done.

(i$'$)$\Rightarrow$(ii): Take a finite covering of $S$ by affine open subsets $S_{f_{ij}}$, where $f_{ij}\in\Gamma(S,L_i^{\otimes m_{ij}})$; put $S_i=\bigcup_j S_{f_{ij}}$. Replacing the $f_{ij}$ by suitable powers, we may assume that the $m_{ij}$ corresponding to the same index $i$ are equal to an integer $m_i$. The sheaf
%%%%%%%%%%  scan idx 183  |  folio 170  %%%%%%%%%%
\SGAPrintPage{170}%
$F\,|\,S_{f_{ij}}$ is
generated by finitely many sections $h_{ijk}$ over $S_{f_{ij}}$. There exists (EGA I 9.3.1 and IV 1.7.5) an integer $q_i\ge0$ such that $h_{ijk}\otimes f_{ij}^{\otimes q_i}$ extends to a section of $F\otimes L_i^{\otimes q_im_i}$ over $S$. In other words, $F\otimes L_i^{\otimes q_im_i}$ is generated over $S_i$ by finitely many of its sections over $S$. Proceeding as in the proof of (EGA II 4.5.5), we deduce that there exists an integer $n_i$, which we may take to be $>0$, such that, for every $n\ge n_i$, $F\otimes L_i^{\otimes n}$ is generated over $S_i$ by finitely many of its sections over $S$. In particular, there exist an integer $r_i\ge0$ and a morphism
\[
\underline{O}_S^{r_i}\otimes L_i^{\otimes-n_i}\longrightarrow F
\]
which induces an epimorphism over $S_i$. The implication (i$'$)$\Rightarrow$(ii) follows immediately.

(ii)$\Rightarrow$(iii): Since every quasi-coherent sheaf on $S$ is an inductive limit of its quasi-coherent subsheaves of finite type (EGA I 9.4.9 and IV 1.7.7), we may restrict ourselves to proving (iii) for $F$ of finite type. There then exist, for $i\in I$, morphisms
\[
u_i:\underline{O}_S^{r_i}\otimes L_i^{\otimes-n_i}\longrightarrow F
\]
such that $\sum u_i$ is an epimorphism. Now we have
\[
\begin{aligned}
\operatorname{Im}(u_i)&=\operatorname{Im}\bigl(u_i\otimes\operatorname{Id}(L_i^{\otimes n_i})\otimes\operatorname{Id}(L_i^{\otimes-n_i})\bigr)\\
&=\operatorname{Im}\bigl(u_i\otimes\operatorname{Id}(L_i^{\otimes n_i})\bigr)\otimes L_i^{\otimes-n_i}
\end{aligned}
\]
But $\operatorname{Im}(u_i\otimes\operatorname{Id}(L_i^{\otimes n_i}))$ is a subsheaf of $F_{(i,n_i)}$, whence the assertion.

(iii)$\Rightarrow$(iii$'$): This is trivial.

(iii$'$)$\Rightarrow$(i): Let $x\in S$, $U$ be an open subset containing $x$, and $F$ a
%%%%%%%%%%  scan idx 184  |  folio 171  %%%%%%%%%%
\SGAPrintPage{171}%
quasi-coherent Ideal
defining $S-U$. By (iii$'$), there exist $(i,n)\in I\times\mathbb N^+$ and a section $f$ of $F\otimes L_i^{\otimes n}$ over $S$ such that $f(x)\ne0$. We have $x\in S_f\subset U$, which proves (i).

\sgaStableTarget{sga6.e02.def.2-2-4}\textbf{Definition 2.2.4.} \uline{Let} $S$ \uline{be a quasi-compact and quasi-separated scheme, and let} $(L_i)_{i\in I}$ \uline{be a family of invertible sheaves on} $S$\uline{. We say that the} $L_i$ \uline{form an} ample family\SGAAuthorFootnote{raisedstar}{This notion, in form (i$'$), was introduced by M.~BORELLI and S.~L.~KLEIMAN; the terminology \textquotedblleft ample family\textquotedblright{} is proposed by the editor.} \uline{when the equivalent conditions of} (2.2.3) \uline{hold. (In particular, a family consisting of a single sheaf} $L$ \uline{is ample precisely when} $L$ \uline{is ample in the usual sense} (EGA II 4.5.3)).

\sgaStableTarget{sga6.e02.def.2-2-5}\textbf{Definition 2.2.5.} \uline{A scheme} $S$ \uline{is said to be} divisorial \uline{when} $S$ \uline{is quasi-compact and quasi-separated and there exists an ample family} (2.2.4) \uline{of invertible} $\underline{O}_S$\uline{-Modules.}

\sgaStableTarget{sga6.e02.lem.2-2-6}\textbf{Lemma 2.2.6.} \uline{Let} $X$ \uline{be a locally noetherian scheme. Let} $U$ \uline{be an open subset of} $X$ \uline{such that the inclusion} $U\to X$ \uline{is affine (for example} $U$ \uline{affine and} $X$ \uline{separated). Then} $X-U$ \uline{is purely of codimension} $0$ \uline{or} $1$ \uline{(i.e.\ at every maximal point} $x$ \uline{of} $X-U$ \uline{we have} $\dim(\underline{O}_{X,x})=0$ or $1$\uline{), hence purely of codimension} $1$ \uline{when} $U$ \uline{contains the maximal points of} $X$.

\uline{Proof.} Put $Y=X-U$. Let $x$ be a maximal point of $Y$, and make the base change $X'=\operatorname{Spec}(\underline{O}_{X,x})\to X$: we obtain a cartesian diagram:
\[
\begin{tikzcd}
U' \arrow[r] \arrow[d] & X' \arrow[d] & Y' \arrow[l] \arrow[d] \\
U \arrow[r] & X & Y \arrow[l]
\end{tikzcd},
\]
where $Y'=X'-U'$. Since $x$ is a maximal point of $Y$, $Y'$ is precisely the closed point of $X'$. On the other hand, since $X'$ is affine, $U'$ is an affine open subset of $X'$. Let us show that the dimension $n$ of $X'$ is equal to $0$ or $1$. Indeed, we know (SGA 2 V 5.7) that $H^n_{Y'}(X',\underline{O}_{X'})\ne0$. But, for $n\ge2$, we have
\[
H^n_{Y'}(X',\underline{O}_{X'})\simeq H^{n-1}(U',\underline{O}_{U'}),
\]
whereas
\[
H^{n-1}(U',\underline{O}_{U'})=0,
\]
%%%%%%%%%%  scan idx 185  |  folio 172  %%%%%%%%%%
\SGAPrintPage{172}%
since $U'$ is affine and $n-1\ge1$, whence a contradiction.

\sgaStableTarget{sga6.e02.rem.2-2-6-1}\textbf{Remark 2.2.6.1.} Without the separation hypothesis on $X$, (2.2.6) for an affine open subset $U$ of $X$ would be false, as shown by the example of the affine plane with the origin doubled (EGA I 5.5.11).

\sgaStableTarget{sga6.e02.prop.2-2-7}\textbf{Proposition 2.2.7.} \uline{Let} $S$ \uline{be a noetherian, separated, locally factorial scheme (i.e.\ one whose local rings are factorial). Then} $S$ \uline{is divisorial} (2.2.5).

In particular (EGA IV 21.11.1):

\sgaStableTarget{sga6.e02.cor.2-2-7-1}\textbf{Corollary 2.2.7.1.} \uline{Every noetherian regular separated scheme is divisorial.}

\uline{Proof} of 2.2.7. We may assume that $S$ is integral. Let $(S_i)_{i\in I}$ be a finite covering of $S$ by nonempty affine open subsets. By (2.2.6), $S-S_i$ is purely of codimension $1$, hence, by (EGA IV 21.7.4), is the support of a positive divisor. In other words, for each $i\in I$, there exist an invertible sheaf $L_i$ and a section $f_i\in\Gamma(S,L_i)$ such that $S_i=S_{f_i}$. The family $(L_i)$ is then ample (2.2.4) by (2.2.3 (i$'$)). This completes the proof.

\sgaStableTarget{sga6.e02.rem.2-2-7-2}\textbf{Remark 2.2.7.2.} As in (2.2.6), the separation hypothesis on $S$ is essential. Indeed, again take for $S$ the affine plane with the origin doubled. Recall that $S$ is obtained by gluing two copies of the affine plane $\underline{\underline{E}}^2_k$ ($k$, a field) by the identity on the complement $U$ of the origin $O$; we therefore have a canonical projection
\[
p:S\longrightarrow\underline{\underline{E}}^2_k,
\]
which induces an isomorphism over $U$, while $p^{-1}(O)$ consists of two points $O_1$ and $O_2$. An invertible sheaf on $S$ is obtained by gluing two invertible sheaves on $\underline{\underline{E}}^2_k$ by an isomorphism over $U$; since every invertible sheaf on $\underline{\underline{E}}^2_k$ is trivial, an invertible sheaf on $S$ is therefore defined by an element $g\in\Gamma(U,\underline{O}_U^*)$; but $g$ is necessarily a constant, since the restriction $\Gamma(\underline{\underline{E}}^2_k,\underline{O}_{\underline{\underline{E}}^2_k})\to\Gamma(U,\underline{O}_U)$ is an isomorphism, $O$ being of depth $2$ in the plane; thus every
%%%%%%%%%%  scan idx 186  |  folio 173  %%%%%%%%%%
\SGAPrintPage{173}%
invertible sheaf on $S$ is trivial. On the other hand, it is immediate that $p$ induces an isomorphism
\[
p^*:\Gamma(\mathbb E^2_k,\underline{O}_{\mathbb E^2_k})\longrightarrow\Gamma(S,\underline{O}_S).
\]
Consequently, every global section of $\underline{O}_S$ which vanishes at $O_1$ also vanishes at $O_2$, so the $S_f$, for $f\in\Gamma(S,\underline{O}_S)$, do not form a basis for the topology of $S$.

\sgaStableTarget{sga6.e02.thm.2-2-8}\textbf{Lemma 2.2.8.} \uline{Let} $S$ \uline{be a divisorial scheme} (2.2.5), $(L_i)_{i\in I}$ \uline{an ample family of invertible} $\underline{O}_S$\uline{-Modules. Denote by} $\underline{L}(S)$ \uline{the strictly full additive subcategory of} $\Qcoh(S)$ (2.2.1) \uline{generated by the} $L_i^{\otimes-n}$ \uline{for} $(i,n)\in I\times\mathbb N^+$\uline{. Let} $E\in\Ob\mathrm D(\Qcoh(S))$.
\begin{enumerate}[label=\alph*)]
\item \uline{The categories} $\underline{L}(S)$ \uline{and} $\Loclib(S)$ (2.0) \uline{are quasi-liftable in} $\Qcoh(S)$ (I 1.2) \uline{and quasi-stable under kernels of epimorphisms} (I 1.2).
\item \uline{Let} $n\in\mathbb Z$\uline{. For} $E$ \uline{to be} $n$\uline{-pseudo-coherent (resp.\ pseudo-coherent)} (2.2.1), \uline{it is necessary and sufficient that} $E$ \uline{be} $n$\uline{-pseudo-coherent (resp.\ pseudo-coherent) relative to} $\underline{L}(S)$ \uline{or} $\Loclib(S)$ (I 2.3)\uline{. The canonical functors}
\[
\mathrm K^{-}(\underline{L}(S))/\mathrm K^{-,\varnothing}(\underline{L}(S))\longrightarrow \mathrm D(\Qcoh(S))
\]
\uline{and}
\[
\mathrm K^{-}(\Loclib(S))/\mathrm K^{-,\varnothing}(\Loclib(S))\longrightarrow \mathrm D(\Qcoh(S))
\]
\uline{are fully faithful and have as essential image} $\mathrm D^{-}(\Qcoh(S))_{\coh}$ (2.2.1).
\item \uline{Let} $[m,n]$ \uline{be an interval in} $\mathbb Z$\uline{. For} $E$ \uline{to have perfect amplitude contained in} $[m,n]$ (2.2.1)\uline{, it is necessary and sufficient that} $E$ \uline{be isomorphic, in} $\mathrm D(\Qcoh(S))$\uline{, to a complex} $F$ \uline{such that} $F^i=0$ \uline{for} $i\notin[m,n]$, $F^i\in\Ob \underline{L}(S)$ \uline{for} $i>m$, \uline{with} $F^m$ \uline{locally free of finite type. The canonical functor}
\[
\mathrm K^b(\Loclib(S))/\mathrm K^{b,\varnothing}(\Loclib(S))\longrightarrow \mathrm D(\Qcoh(S))
\]
%%%%%%%%%%  scan idx 187  |  folio 174  %%%%%%%%%%
\SGAPrintPage{174}%
is
fully faithful and has as essential image $\mathrm D(\Qcoh(S))_{\parf}$ (2.2.1).
\end{enumerate}

\uline{Proof.} Let $F\to G$ be an epimorphism in $\Qcoh(S)$, where $G$ is of finite type. Since $F$ is an inductive limit of its quasi-coherent subsheaves of finite type (EGA I 9.4.9), there exists, by (2.1.1 a)), a quasi-coherent subsheaf $F'$ of $F$ of finite type such that the composite $F'\to F\to G$ is an epimorphism. By (2.2.3 (ii)), there exists an epimorphism $E\to F'$, with $E\in\Ob \underline{L}(S)$, so the composite $E\to F'\to F$, composed with $F\to G$, is an epimorphism, which proves that the category of quasi-coherent $\underline{O}_S$-Modules of finite type is quasi-liftable in $\Qcoh(S)$; a fortiori, the same holds for the category $\Loclib(S)$ or $\underline{L}(S)$. This proves the first assertion of a). For the second, we apply (1.1), taking for $\underline{S}$ the Zariski site of $S$, and for $\underline{C}$ (resp. $\underline{C}_0$) the fibered category $U\mapsto\Qcoh(U)$ (resp. $\underline{L}(U)$) ($U$ denotes a variable open subset of $S$, and $\underline{L}(U)$ the strictly full additive subcategory of $\Qcoh(U)$ generated by the $(L_i^{\otimes-n}\,|\,U)$). Since every object of $\underline{L}$ is pseudo-coherent, the notions of $n$-pseudo-coherence (resp. pseudo-coherence) (2.2.1) in the fibered category $\mathrm D(\Qcoh)$ coincide with the same notions relative to $\underline{L}$ (I 2.14). Condition (ii) (in the non-resp.\ case) of (1.1) being satisfied by (2.2.3 (ii)), we deduce, by (1.1 a)), the second part of a) concerning $\underline{L}(S)$; the assertion concerning $\Loclib(S)$ is proved similarly. Assertions b) follow immediately from (1.1 b) and c)), taking into account the fact that, when $E\in\Ob\mathrm D(\Qcoh(S))$ is $n$-pseudo-coherent, $E$ is an object of $\mathrm D^{-}(\Qcoh(S))$, since $S$ is quasi-compact. An object of the fibered category $\Qcoh$ is locally in $\underline{L}$ precisely when it is locally free of finite type: the category $\underline{L}$ is therefore locally liftable in $\Qcoh$ and locally stable under kernels of epimorphisms. Consequently, assertions c) follow from (1.2) applied with $\underline{C}=\Qcoh$, $\underline{C}_0=\underline{L}$, $\underline{C}_{1S}=\Loclib(S)$. This completes the proof.

\sgaStableTarget{sga6.e02.prop.2-2-9}\textbf{Proposition 2.2.9.} \uline{The hypotheses and notation are those of} (2.2.8)\uline{, but we further assume that} $S$ \uline{is separated or noetherian.}
\begin{enumerate}[label=\alph*)]
\item \uline{The canonical functors}
\[
\mathrm K^{-,b}(\underline{L}(S))/\mathrm K^{-,\varnothing}(\underline{L}(S))\longrightarrow\mathrm D(S)
\]
and
\[
\mathrm K^{-,b}(\Loclib(S))/\mathrm K^{-,\varnothing}(\Loclib(S))\longrightarrow\mathrm D(S)
\]
%%%%%%%%%%  scan idx 188  |  folio 175  %%%%%%%%%%
\SGAPrintPage{175}%
\uline{are fully faithful and have as essential image} $\mathrm D^b(S)_{\coh}$ (I 2.15)\uline{. Moreover, when} $S$ \uline{has \textquotedblleft finite cohomological dimension\textquotedblright{} in the sense of} (3.7)\SGAAuthorFootnote{1}{Section~3 is independent of the rest of the exposé.} \uline{(for example, when the underlying space of} $S$ \uline{is noetherian and of finite Krull dimension), the functors}
\[
\mathrm K^{-}(\underline{L}(S))/\mathrm K^{-,\varnothing}(\underline{L}(S))\longrightarrow
\mathrm K^{-}(\Loclib(S))/\mathrm K^{-,\varnothing}(\Loclib(S))\longrightarrow
\mathrm D^{-}(S)_{\coh}
\]
\uline{are equivalences of categories.}
\item \uline{The canonical functor}
\[
\mathrm K^b(\Loclib(S))/\mathrm K^{b,\varnothing}(\Loclib(S))\longrightarrow\mathrm D(S)
\]
\uline{is fully faithful and has as essential image} $\mathrm D(S)_{\parf}$ (I 4.9)\uline{. For} $E\in\Ob\mathrm D(S)$ \uline{to have perfect amplitude contained in} $[m,n]$ (I 4.8)\uline{, it is necessary and sufficient that} $E$ \uline{be isomorphic, in} $\mathrm D(S)$\uline{, to a complex} $F$ \uline{such that} $F^i=0$ \uline{for} $i\notin[m,n]$, $F^i\in\Ob \underline{L}(S)$ \uline{for} $i>m$, \uline{with} $F^m$ \uline{locally free of finite type.}
\end{enumerate}

\uline{Proof.} In view of (3.5) and (3.7)\SGAAuthorFootMark{1}, this is an immediate corollary of (2.2.8): the only point to verify, which follows trivially from (3.5) (resp. (3.7)), is that, for an object $E$ of $\mathrm D^b(\Qcoh(S))$ (resp. $\mathrm D(\Qcoh(S))$ when $S$ has finite cohomological dimension in the sense of (3.7)), the notions of pseudo-coherence and perfection of (2.2.1) coincide with the usual notions (I 2.15 and 4.9) when $E$ is regarded as an object of $\mathrm D(S)$.

\sgaStableTarget{sga6.e02.rem.2-2-10}\textbf{Remark 2.2.10.} a) When $S$ is an affine scheme with ring $A$, the family consisting of the sheaf $\underline{O}_S$ is ample. For this family, we have $\underline{L}(S)=\Lib(S)$. It is also clear that the categories $\mathrm K^{-,\varnothing}(\Lib(S))$ and $\mathrm K^{-,\varnothing}(\Loclib(S))$ are reduced to zero objects. We therefore conclude from (2.2.9) that, writing $\Mod(A)$ for the category of $A$-modules and $\mathrm D(A)=\mathrm D(\Mod(A))$, the functor $\underline{\underline{R}}\Gamma_S:\mathrm D^+(S)\longrightarrow\mathrm D^+(A)$ induces a fully faithful functor $\mathrm D^b(S)_{\coh}\longrightarrow\mathrm D(A)$ (resp. $\mathrm D(S)_{\parf}\longrightarrow\mathrm D(A)$) whose essential image is
%%%%%%%%%%  scan idx 189  |  folio 176  %%%%%%%%%%
\SGAPrintPage{176}%
the
triangulated subcategory $\mathrm D^b(A)_{\coh}$ (resp. $\mathrm D(A)_{\parf}$) formed by complexes with bounded cohomology and pseudo-coherent (resp. perfect) relative to the subcategory $\Lib(A)$ (resp. $\Loclib(A)$) of free (resp. projective) modules of finite type. We have commutative diagrams of equivalences of categories
\[
\begin{tikzcd}
\mathrm K^{-,b}(\Lib(S)) \arrow[r] \arrow[d,"\Gamma_S"] &
\mathrm K^{-,b}(\Loclib(S)) \arrow[r] \arrow[d,"\Gamma_S"] &
\mathrm D^b(S)_{\coh} \arrow[d,"\mathbb R\Gamma_S"] \\
\mathrm K^{-,b}(\Lib(A)) \arrow[r] &
\mathrm K^{-,b}(\Loclib(A)) \arrow[r] &
\mathrm D^b(A)_{\coh},
\end{tikzcd}
\]
\[
\begin{tikzcd}
\mathrm K^b(\Loclib(S)) \arrow[r] \arrow[d,"\Gamma_S"] &
\mathrm D(S)_{\parf} \arrow[d,"\mathbb R\Gamma_S"] \\
\mathrm K^b(\Loclib(A)) \arrow[r] &
\mathrm D(A)_{\parf},
\end{tikzcd}
\]
(where, in the upper diagram, one may replace $b$ by $-$ when $S$ has finite cohomological dimension in the sense of (3.7)).

b) Let $S$ be a scheme. It follows from a) that a complex of $\underline{O}_S$-Modules which is pseudo-coherent and has bounded cohomology is locally isomorphic (in $\mathrm D(S)$) to a strictly pseudo-coherent complex, which was by no means evident a priori (cf. (I 2.3.1)).

\sgaStableTarget{sga6.e02.sec.2.3}\subsection*{2.3. Soft sheaves}

\sgaStableTarget{sga6.e02.para.2-3-1}2.3.1. Recall some definitions and results from [T.F]. Let $X$ be a paracompact space. A sheaf of sets $F$ on $X$ is said to be \uline{soft} when, for every closed subset $Y$ of $X$, every section of $F$ over $Y$ extends to a section of $F$ over $X$. Soft abelian sheaves are acyclic, i.e.\ when $F$ is soft, we have $H^i(X,F)=0$ for $i>0$ (T.F. II 4.4.3).

Let $A$ be a sheaf of rings on $X$. When $F$ is a soft $A$-Module, $F$ is generated by its sections (cf. 2.1.1): indeed, the homomorphism $A_X^{(I)}\to F$ defined by $I=H^0(X,F)$ is clearly an epimorphism.
%%%%%%%%%%  scan idx 190  |  folio 177  %%%%%%%%%%
\SGAPrintPage{177}%
On the other hand, when $A$ is soft, every $A$-Module is soft (T.F II 3.7.1). As examples of soft sheaves of rings, we mention the sheaf of continuous functions (real or complex) on a paracompact space, and the sheaf of $C^r$ functions on a manifold of class $C^r$.

\sgaStableTarget{sga6.e02.prop.2-3-2}\textbf{Proposition 2.3.2.} \uline{Let} $S$ \uline{be a compact space, equipped with a soft sheaf of rings} $A$\uline{. We write} $\Mod(S)$ \uline{for the category of left} $A$\uline{-Modules,} $\mathrm D(S)=\mathrm D(\Mod(S))$.
\begin{enumerate}[label=\alph*)]
\item \uline{The category} $\Lib(S)$ \uline{(resp.} $\underline{\operatorname{Loclib}}(S)$) (2.0) \uline{is liftable in} $\Mod(S)$ \uline{and quasi-stable} (resp. stable) \uline{under kernels of epimorphisms} (I 1.2)\SGAAuthorFootnote{1}{This assertion holds more generally for $S$ paracompact.}.
\item \uline{Let} $E\in\Ob\mathrm D(S)$ \uline{and} $n\in\mathbb Z$\uline{. For} $E$ \uline{to be} $n$\uline{-pseudo-coherent (resp.\ pseudo-coherent), it is necessary and sufficient that} $E$ \uline{be so relative to} $\Lib(S)$ \uline{or to} $\Loclib(S)$\uline{. The canonical functors}
\[
\mathrm K^{-}(\Lib(S))\longrightarrow \mathrm D(S)
\]
\uline{and}
\[
\mathrm K^{-}(\Loclib(S))\longrightarrow \mathrm D(S)
\]
\uline{are fully faithful and have as essential image} $\mathrm D^{-}(S)_{\coh}$ (I 2.15).
\item \uline{Let} $E\in\Ob\mathrm D(S)$\uline{, and} $[m,n]$ \uline{be an interval in} $\mathbb Z$\uline{. For} $E$ \uline{to have perfect amplitude contained in} $[m,n]$ (I 4.8), \uline{it is necessary and sufficient that there exist a quasi-isomorphism} $F\to E$\uline{, where} $F$ \uline{is a complex such that} $F^i=0$ \uline{for} $i\notin[m,n]$, $F^i\in\Ob\Lib(S)$ \uline{for} $i>m$ \uline{and} $F^m\in\Ob\Loclib(S)$\uline{. The canonical functor}
\[
\mathrm K^b(\Loclib(S))\longrightarrow\mathrm D(S)
\]
\uline{is fully faithful and has as essential image} $\mathrm D(S)_{\parf}$ (I 4.9).
\item \uline{Denote by} $\Mod(A(S))$ \uline{the category of left} $A(S)$\uline{-modules, and by} $\mathrm D(A(S))$ \uline{the corresponding derived category. The functor}
\[
H^0(S,.):\Mod(S)\longrightarrow\Mod(A(S))
\]
\uline{is exact and an equivalence of categories. Moreover, it
defines}
%%%%%%%%%%  scan idx 191  |  folio 178  %%%%%%%%%%
\SGAPrintPage{178}%
\uline{equivalences of categories:}
\[
\begin{gathered}
\mathrm D(S)\longrightarrow\mathrm D(A(S))\\
\mathrm D^{-}(S)_{\coh}\longrightarrow\mathrm D^{-}(A(S))_{\coh}\\
\mathrm D(S)_{\parf}\longrightarrow\mathrm D(A(S))_{\parf},
\end{gathered}
\]
\uline{the notions of pseudo-coherence (resp.\ perfection) in} $\mathrm D(A(S))$ \uline{being understood relative to the category of free (resp.\ projective) modules of finite type. It also defines an equivalence from the full subcategory of} $\mathrm D(S)$ \uline{formed by} $n$\uline{-pseudo-coherent objects (resp.\ objects of perfect amplitude contained in} $[m,n]$\uline{) onto the analogous full subcategory of} $\mathrm D(A(S))$.
\end{enumerate}

\uline{Proof.} The first assertion of a) follows from (2.1.2 a)) applied by taking for $\underline{S}$ the topos of sheaves of sets on $S$, and $\underline{C}'_S=\Mod(S)$ (in view of (2.3.1)). The second follows from (1.1 a)) (resp. (1.2 a))) with $\underline{C}=\Mod$ (the fibered category of $A$-Modules on variable open subsets), $\underline{C}_0=\Lib$ (resp. $\Loclib$), condition (ii) of (1.1) being verified by (2.3.1) and (2.1.1.b)). For assertion b), we apply (1.1 b),c)) (with $\underline{C}$ and $\underline{C}_0$ as before). Assertion c) follows from (1.2 b),c),d)), where $\underline{C}=\Mod$ and $\underline{C}_0=\Lib$. The first assertion of d) follows, by (2.3.1), from (2.1.2 b) (ii)), where we take $\underline{C}'_S=\underline{C}_S=\Mod(S)$. The other assertions of d) follow from the preceding one, in view of b) and c), and the fact that the functor $\Gamma_S$ sends free Modules (resp. direct factors of free Modules) of finite type to free modules (resp. projective modules of finite type).

\sgaStableTarget{sga6.e02.sec.2.4}\subsection*{2.4. Stein spaces}

\sgaStableTarget{sga6.e02.def.2-4-1}\textbf{Definition 2.4.1.} \uline{Let} $\underline{S}$ \uline{be a locally ringed topos} (I 2.15.1 (ii))\uline{, with final object} $S$\uline{. We say that} $\underline{S}$ \uline{is} Stein \uline{when the structure sheaf} $\underline{O}_S$ \uline{is coherent} (I 3.1) \uline{and the following properties hold:}

\textquotedblleft\uline{theorem} A\textquotedblright{} \uline{: every coherent} $\underline{O}_S$\uline{-Module} (I 3.1) \uline{is generated by its sections,}

\textquotedblleft\uline{theorem} B\textquotedblright{} \uline{: for every coherent} $\underline{O}_S$\uline{-Module} $F$\uline{, we have} $H^i(S,F)=0$ \uline{for} $i>0$.

%%%%%%%%%%  scan idx 192  |  folio 179  %%%%%%%%%%
\SGAPrintPage{179}%
\sgaStableTarget{sga6.e02.rem.2-4-1-1}\textbf{Remark 2.4.1.1.} When $\underline{S}$ has \textquotedblleft enough closed points\textquotedblright\SGAAuthorFootnote{1}{i.e.\ enough fiber functors (SGA~4 IV 6.4.1) of the type $i_s:s\to\underline{S}$, where $i_s$ is the inclusion of a closed subtopos of $\underline{S}$.} $s$ such that the Ideal $\underline{m}_s$ defined by $s$\SGAAuthorFootnote{2}{i.e.\ formed by sections of $\underline{O}_S$ whose fiber at $s$ belongs to the maximal ideal $m_s$ of $\underline{O}_{S,s}$.} is coherent, the relation \textquotedblleft$H^1(S,F)=0$ for every coherent sheaf $F$\textquotedblright{} implies theorem A. Indeed, let $F$ be a coherent $\underline{O}_S$-Module and $s$ a closed point of $S$. We have the exact sequence
\[
0\longrightarrow\underline{m}_sF\longrightarrow F\longrightarrow F/\underline{m}_sF\longrightarrow0,
\]
where $F/\underline{m}_sF$ is a coherent sheaf concentrated at $s$. Applying the functor $H^0(S,.)$ and writing the cohomology exact sequence, we find that $H^0(S,F)\longrightarrow H^0(S,F/\underline{m}_sF)$ is an epimorphism. Since $H^0(S,F/\underline{m}_sF)\simeq F_s/m_sF_s$, where $m_s$ is the maximal ideal of $\underline{O}_{S,s}$, we conclude by Nakayama.

The editor does not know whether the preceding implication is true in general (when one assumes only that $\underline{O}_S$ is a coherent sheaf of local rings).

The editor also does not know whether the relation \textquotedblleft$H^1(S,F)=0$ for every coherent sheaf $F$\textquotedblright{} implies theorem B.

\sgaStableTarget{sga6.e02.ex.2-4-2}\textbf{Examples 2.4.2.} A complex analytic Stein space [3], a compact convex subset of $\mathbb C^n$ equipped with the sheaf induced by the sheaf of holomorphic functions [3], the spectrum of a noetherian ring, and an affine rigid analytic space [13] define Stein topoi.

\sgaStableTarget{sga6.e02.lem.2-4-3}\textbf{Lemma 2.4.3.} \uline{Let} $\underline{S}$ \uline{be a Stein topos} (2.4.1)\uline{. We assume that the final object} $S$ \uline{is quasi-compact} (2.1.1)\uline{, and we denote by} $\Coh(S)$ \uline{the abelian category of coherent} $\underline{O}_S$\uline{-Modules, i.e.\ of finite presentation. Then:}
\begin{enumerate}[label=\alph*)]
\item \uline{The category} $\Lib(S)$ \uline{(resp.} $\Loclib(S)$) (2.0) \uline{is liftable in} $\Coh(S)$ \uline{and quasi-stable (resp.\ stable) under kernels of epimorphisms} (I 1.2).
\item \uline{The canonical functors}
\[
\mathrm K^{-}(\Lib(S))\longrightarrow\mathrm K^{-}(\Loclib(S))\longrightarrow\mathrm D^{-}(\Coh(S))
\]
\uline{are equivalences of categories.}
%%%%%%%%%%  scan idx 193  |  folio 180  %%%%%%%%%%
\item \SGAPrintPage{180}%
\uline{Put} $A=\Gamma(S,\underline{O}_S)$\uline{, and denote by} $\Mod(A)$ \uline{the category of} $A$\uline{-modules, and by} $\mathrm D(A)$ \uline{the corresponding derived category. The functor}
\[
\Gamma_S:\Coh(S)\longrightarrow\Mod(A)
\]
\uline{is exact, fully faithful and has as essential image the strictly full additive subcategory of} $\Mod(A)$ \uline{formed by} $A$\uline{-modules of finite presentation. It induces a fully faithful functor}
\[
\mathrm D^{-}(\Coh(S))\longrightarrow\mathrm D(A),
\]
\uline{whose essential image is the full subcategory} $\mathrm D^{-}(A)_{\coh}$ \uline{formed by pseudo-coherent complexes (relative to the subcategory of free modules of finite type).}
\end{enumerate}

\uline{Proof.} For the first assertion of a), we apply (2.1.2 a)) with $\underline{C}'_S=\Coh(S)$, the hypothesis of (2.1.2) being satisfied by theorem B: with the notation of (2.1.2 a)), when $E$ is coherent and $G$ locally free of finite type, $\underline{\operatorname{Hom}}(G,E)$ is coherent, so the category $\Loclib(S)$ (and a fortiori $\Lib(S)$) is liftable in $\Coh(S)$. The second assertion of a) follows from (1.1 a)) (resp.(1.2 a)) ) applied with $\underline{C}_0=\Lib$ (resp. $\Loclib$)\SGAAuthorFootnote{1}{Lib and Loclib are locally stable under kernels of epimorphisms by (I 2.15.1).}, the resp.\ condition (ii) of (1.1) being satisfied by theorem A and (2.1.1 b)). Assertion b) follows from (1.1 c)) applied successively with $\underline{C}_0=\Lib$ and $\underline{C}_0=\Loclib$. The first assertion of c) follows from (2.1.2 b)(ii)) applied with $\underline{C}'_S=\Coh(S)$. To prove that the functor $\Gamma_S:\mathrm D^{-}(\Coh(S))\to\mathrm D(A)$ is fully faithful, it suffices, in view of b), to prove that the composite
\[
\mathrm K^{-}(\Lib(S))\longrightarrow\mathrm D^{-}(\Coh(S))\xrightarrow{\Gamma_S}\mathrm D(A)
\]
is so. But this factors as
\[
\mathrm K^{-}(\Lib(S))\xrightarrow{(*)}\mathrm K^{-}(\Lib(A))\xrightarrow{(**)}\mathrm D(A),
\]
%%%%%%%%%%  scan idx 194  |  folio 181  %%%%%%%%%%
\SGAPrintPage{181}%
where $\Lib(A)$
denotes the category of free $A$-modules of finite type and $(*)$ is induced by $\Gamma_S$. It is trivial that $(*)$ is an equivalence of categories. On the other hand, (I 2.7 b)) applied with $\underline{C}=\Mod(A)$, $\underline{C}_0=\Lib(A)$ (and taking account of the fact that $\mathrm K^{-,\varnothing}(\Lib(A))=0$) shows that $(**)$ is fully faithful and has as essential image $\mathrm D^{-}(A)_{\coh}$. This proves the last assertion of c) and completes the proof of (2.4.3).

\sgaStableTarget{sga6.e02.lem.2-4-4}\textbf{Lemma 2.4.4.} \uline{Under the hypotheses of} (2.4.3)\uline{, the canonical functor}
\[
\mathrm D^b(\Coh(S))\longrightarrow\mathrm D(S)
\]
\uline{is fully faithful and has as essential image} $\mathrm D^b(S)_{\coh}$ (I 2.15).

\uline{Proof} We apply (2.1.3) with $A=\Coh(S)$, $B=\Mod(S)$, $\phi:A\to B$ being the inclusion functor. Since this functor is fully faithful, it suffices to check, by (2.1.3 a)(ii)), that for fixed $M\in\Ob\Coh(S)$ and $n>0$ the functor $L\longmapsto \Ext^n_{\Mod(S)}(L,M)$ is coeffaceable in $L\in\Ob\Coh(S)$. Now, by theorem A and (2.1.1 b)), every coherent $\underline{O}_S$-Module is a quotient of a free Module of finite type, and, when $L$ is free of finite type, we have $\Ext^n_{\Mod(S)}(L,M)=0$ by theorem B, whence our assertion. It follows that the functor $\mathrm D^b(\Coh(S))\to\mathrm D(S)$ is fully faithful and has as essential image the full subcategory of $\mathrm D(S)$ formed by complexes with bounded and coherent cohomology, which is precisely $\mathrm D^b(S)_{\coh}$ since $\underline{O}_S$ is coherent (I 3.5. b)).

\sgaStableTarget{sga6.e02.prop.2-4-5}\textbf{Proposition 2.4.5.} \uline{With the hypotheses and notation of} (2.4.3)\uline{, we have:}
\begin{enumerate}[label=\alph*)]
\item \uline{The canonical functors}
\[
\mathrm K^{-,b}(\Lib(S))\longrightarrow\mathrm K^{-,b}(\Loclib(S))\longrightarrow\mathrm D^b(S)_{\coh}
\]
\uline{are equivalences of categories. Moreover, the functor} $\underline{\underline{R}}\Gamma_S:\mathrm D^+(S)\to\mathrm D^+(A)$ \uline{induces a fully faithful functor}
\[
\underline{\underline{R}}\Gamma_S:\mathrm D^b(S)_{\coh}\longrightarrow\mathrm D^+(A),
\]
%%%%%%%%%%  scan idx 195  |  folio 182  %%%%%%%%%%
\SGAPrintPage{182}%
\uline{whose essential image is} $\mathrm D^b(A)_{\coh}$ \uline{(the triangulated subcategory of} $\mathrm D(A)$ \uline{formed by complexes with bounded cohomology and pseudo-coherent relative to the category of free modules of finite type).}
\item \uline{Let} $E\in\Ob\mathrm D(S)$\uline{, and let} $[m,n]$ \uline{be an interval of} $\mathbb Z$\uline{. When} $E$ \uline{has perfect amplitude contained in} $[m,n]$ (I 4.8), $E$ \uline{is isomorphic to a complex} $F$ \uline{such that} $F^i=0$ \uline{for} $i\notin[m,n]$, $F^i\in\Ob\Lib(S)$ \uline{for} $i>m$ \uline{and} $F^m\in\Ob\Loclib(S)$.
\item \uline{The canonical functor}
\[
\mathrm K^b(\Loclib(S))\longrightarrow\mathrm D(S)_{\parf}\qquad(\text{I }4.9)
\]
\uline{is an equivalence of categories. Moreover, the functor} $\underline{\underline{R}}\Gamma_S:\mathrm D^+(S)\to\mathrm D^+(A)$ \uline{induces a fully faithful functor}
\[
\underline{\underline{R}}\Gamma_S:\mathrm D(S)_{\parf}\longrightarrow\mathrm D^+(A),
\]
\uline{whose essential image is} $\mathrm D(A)_{\parf}$ \uline{(the triangulated subcategory of} $\mathrm D(A)$ \uline{formed by complexes perfect relative to the category of projective modules of finite type).}
\end{enumerate}

\uline{Proof.} a) The functors
\[
\mathrm K^{-,b}(\Lib(S))\longrightarrow\mathrm K^{-,b}(\Loclib(S))\longrightarrow\mathrm D^b(\Coh(S))\xrightarrow{\Gamma_S}\mathrm D^b(A)_{\coh}
\]
are equivalences of categories by (2.4.3 b) and c)). Since the (fully faithful) functor $\Gamma_S:\mathrm D^b(\Coh(S))\to\mathrm D(A)$ factors as
\[
\mathrm D^b(\Coh(S))\longrightarrow\mathrm D^b(S)_{\coh}\xrightarrow{\underline{\underline{R}}\Gamma_S}\mathrm D(A),
\]
we conclude by (2.4.4).

b) Suppose that $\operatorname{parf.amp}(E)\subset[m,n]$. In particular, $E$ is pseudo-coherent; by a), after replacing $E$ by an isomorphic object, we may therefore assume that $E^i\in\Ob\Lib(S)$ for $i>m$ and $E^i=0$ for $i>n$.
%%%%%%%%%%  scan idx 196  |  folio 183  %%%%%%%%%%
\SGAPrintPage{183}%
After truncation, we may even assume that $E^i=0$ for $i<m$. By (I 4.16) applied with $\underline{C}=\Mod$, $\underline{C}_0=\Lib$, we then have $E^m\in\Ob\Loclib(S)$, whence the assertion.

c) Hypothesis (1.1 (i))(resp. (1.1 (ii) resp.)) being satisfied for $\underline{C}_S=\Coh(S)$, $\underline{C}_0=\Lib$ by (2.4.3 a)) (resp. (2.1.1 b)) and theorem A), we know from (1.2 c) and d)) that the functor
\[
\mathrm K^b(\Loclib(S))\longrightarrow\mathrm D^b(\Coh(S))
\]
is fully faithful. Thus, by (2.4.4), the same holds for the functor
\[
\mathrm K^b(\Loclib(S))\longrightarrow\mathrm D(S),
\]
whose essential image is $\mathrm D(S)_{\parf}$ by b). By an argument analogous to that in the proof of (2.4.3 c)), we see that the functor $\mathrm K^b(\Loclib(S))\to\mathrm D(A)$ induced by $\Gamma_S$ is fully faithful and has as essential image $\mathrm D(A)_{\parf}$, and we complete the proof as in a).

\sgaStableTarget{sga6.e02.rem.2-4-6}\textbf{Remark 2.4.6.} The editor does not know whether, in (2.4.4), one may replace $\mathrm D^b$ by $\mathrm D^{-}$, and consequently whether (2.4.5 a)) remains true when one replaces $\mathrm K^{-,b}$ by $\mathrm K^{-}$ and $\mathrm D^b$ by $\mathrm D^{-}$.

To finish, here is an interesting complement to (2.4.5). Let us first give a definition:

\sgaStableTarget{sga6.e02.def.2-4-7}\textbf{Definition 2.4.7.} \uline{A complex pseudo-analytic space is a ringed space locally isomorphic to a ringed space of the form} $(Y,\underline{O}_X|Y)$\uline{, where} $X$ \uline{is a complex analytic space and} $Y$ \uline{a closed subset of} $X$\uline{.

The structure sheaf of a complex pseudo-analytic space is therefore coherent.

A complex pseudo-analytic space is said to be} Stein \uline{when the corresponding topos is Stein} (2.4.1).

% Exposé II continuation macros
\providecommand{\Qcoh}{\operatorname{Qcoh}}
\providecommand{\Mod}{\operatorname{Mod}}
\providecommand{\Coh}{\operatorname{Coh}}
\renewcommand{\Ob}{\operatorname{ob}}
\providecommand{\Hom}{\operatorname{Hom}}
\providecommand{\Ext}{\operatorname{Ext}}
\providecommand{\Spec}{\operatorname{Spec}}
\providecommand{\qcoh}{\mathrm{qcoh}}
\providecommand{\parf}{\mathrm{parf}}
\providecommand{\coh}{\mathrm{coh}}
\providecommand{\Loclib}{\operatorname{Loclib}}
\providecommand{\Lib}{\operatorname{Lib}}
\providecommand{\im}{\operatorname{Im}}

%%%%%%%%%%  scan idx 197  |  folio 184  %%%%%%%%%%
\SGAPrintPage{184}%
\sgaStableTarget{sga6.e02.expose.occ2}%
\sgaStableTarget{sga6.e02.prop.2-4-8}\textbf{Proposition 2.4.8.}\SGAAuthorFootnote{1}{This result, together with its proof, was kindly communicated to the editor by A. Grothendieck; it was obtained independently by J. Frisch [5].}. \uline{Let} $S$ \uline{be a compact complex pseudo-analytic Stein space} (2.4.7)\uline{. We assume that, for every open subset} $U$ \uline{of} $S$ \uline{and every reduced closed subspace} $X$ \uline{of} $U$\uline{, the normalization of} $X$ \uline{is locally connected}\SGAAuthorFootnote{2}{We shall admit that the notions of local Analytic Geometry in [10] extend to complex pseudo-analytic spaces; the reduction to the usual case is indeed immediate each time.} \uline{(a condition satisfied, for example, when} $S$ \uline{is locally defined by real analytic equations or inequalities in a space} $\underline{\underline{R}}^n$ [15]\uline{). Then the ring} $A=\Gamma(S,\underline{O}_S)$ \uline{is a noetherian ring.}

\uline{Proof.} We have to show that, for an $A$-module $M$ of finite presentation, the set of submodules (or, equivalently, the set of submodules of finite type) of $M$ satisfies the maximal condition. Now, when $N$ is a submodule of finite type of $M$, $M/N$ is of finite presentation by (I 2.6 and 2.9), so $N$ is of finite presentation by the equivalence (2.4.3 c)). Using again the equivalence (2.4.3 c)) between the category of coherent $\underline{O}_S$-Modules and the category of $A$-modules of finite presentation, and the compactness of $S$, we are reduced to proving the following lemma:

\sgaStableTarget{sga6.e02.lem.2-4-8-1}\textbf{Lemma 2.4.8.1.} \uline{Let} $S$ \uline{be a complex pseudo-analytic space satisfying the local connectedness condition of} (2.4.8)\uline{. Let} $F$ \uline{be a coherent} $\underline{O}_S$\uline{-Module,} $(E_n)$ \uline{an increasing sequence of coherent sub-}$\underline{O}_S$\uline{-Modules of} $F$, $s$ \uline{a point of} $S$\uline{. Then there exists an open neighbourhood} $U$ \uline{of} $s$ \uline{such that the sequence} $(E_n|U)$ \uline{is stationary.}

\uline{Proof.} Since the ring $\underline{O}_{S,s}$ is noetherian, there exists an integer $n_0$ such that $(E_n)_s=(E_{n_0})_s$ for $n\geq n_0$. Replacing the sequence $(E_n)$ by the sequence $(E_n/E_{n_0})$ for $n\geq n_0$, we may therefore assume that $(E_n)_s=0$ for every $n$. On the other hand, $F_s=M$ is an $\underline{O}_s$-module of finite type, so there exists (Bourbaki Alg. com. chap. IV \S1 th.1) a finite filtration
\[
 M=M_0\supset M_1\supset\ldots\supset M_q=0
\]
such that $M_i/M_{i+1}\simeq \underline{O}_s/p_i$, where $p_i$ is a prime ideal of $\underline{O}_s$. Since the category of $\underline{O}_s$-modules of finite type is equivalent to the category of germs at $s$ of coherent $\underline{O}_S$-Modules, we may suppose, after
%%%%%%%%%%  scan idx 198  |  folio 185  %%%%%%%%%%
\SGAPrintPage{185}%
localizing in a neighbourhood of $s$, that the $M_i$ (resp. $p_i$) are the fibers at $s$ of coherent sub-$\underline{O}_S$-Modules $F_i$ (resp. coherent Ideals $\underline{p}_i$) of $F$ (resp. $\underline{O}_S$) such that
\[
 F=F_0\supset F_1\supset\ldots\supset F_q=0
\]
and $F_i/F_{i+1}\simeq\underline{O}_S/\underline{p}_i$. Proving that the sequence $(E_n)$ is stationary in a neighbourhood of $s$ is evidently equivalent to proving that the sequences induced on each graded piece of the filtration $(F_i)$ of $F$ are stationary. Thus we may assume that $F=\underline{O}_S$, and that $\underline{O}_{S,s}$ is an integral domain, the $E_n$ forming an increasing sequence of Ideals of $\underline{O}_S$ such that $(E_n)_s=0$. Since $\underline{O}_s$ is an integral domain, hence in particular reduced, Oka's theorem ([10] IV \S3) further allows us to assume, after localizing at $s$, that $S$ is \uline{reduced}. Let us show that there exists an open neighbourhood $U$ of $s$ such that $E_n|U=0$ for sufficiently large $n$. For this, consider the normalization $\widetilde S$ of $S$ ([10] IV \S4). We know (loc. cit.) that $\widetilde S$ is the analytic spectrum ([10] II \S3) of a finite $\underline{O}_S$-Algebra $\widetilde{\underline{O}}_S$ whose fiber at each $t\in S$ is the integral closure of $\underline{O}_{S,t}$ (cf. note $(^2)$ to (2.4.8)). We have $E_n\subset E_n\cdot\widetilde{\underline{O}}_S$, so it suffices to prove that $E_n\cdot\widetilde{\underline{O}}_S=0$ for sufficiently large $n$. Now we have ([10] II \S3 Prop.8)
\[
 E_n\cdot\widetilde{\underline{O}}_S = p_*(\widetilde{E_n\cdot\widetilde{\underline{O}}_S}),
\]
where $p:\widetilde S\longrightarrow S$ is the canonical projection, and $\widetilde{E_n\cdot\widetilde{\underline{O}}_S}$ the Ideal of $\underline{O}_{\widetilde S}$ canonically defined by $E_n\cdot\widetilde{\underline{O}}_S$. It therefore suffices to show that there exists an open neighbourhood $V$ of $s$ in $\widetilde S$ such that $\widetilde{E_n\cdot\widetilde{\underline{O}}_S}|V=0$ for sufficiently large $n$, which reduces us to the case where $S$ is \uline{normal}, as we shall now assume. Since $(E_n)_s=0$, the closed immersion $i_n:S_n\longrightarrow S$ defined by $E_n$ is \uline{étale} at $s$. But, since $S$ is normal, the set of points where $i_n$ is étale is both open and closed (see the lemma below), so its image in $S$ is a subset of $S$ that is both open and closed. Thus the restriction of $E_n$ to the connected component of $s$ in $S$ is zero; this connected component is an open set, since $S$ is locally connected, and we are done. It therefore remains to establish the following

%%%%%%%%%%  scan idx 199  |  folio 186  %%%%%%%%%%
\SGAPrintPage{186}%
\sgaStableTarget{sga6.e02.lem.2-4-8-2}\textbf{Lemma 2.4.8.2.} (cf. EGA IV 18.10.3) \uline{Let} $Y$ \uline{be a normal pseudo-analytic space}\SGAAuthorFootnote{1}{It suffices that the local rings of $Y$ be integral domains.}\uline{, and} $f:X\longrightarrow Y$ \uline{an unramified morphism. The set of points of} $X$ \uline{where} $f$ \uline{is étale is both open and closed.}

\uline{Proof.} It is immediate that the set of points where $f$ is étale is open (since saying that $f$ is étale at a point means that $f$ is an isomorphism in a neighbourhood of that point ([8] VI)). It remains to show that the set of points where $f$ is not étale is open. Since an unramified morphism is locally an immersion ([8] VI), we may assume that $f$ is a closed immersion, defined by a coherent Ideal $I$. Let $x$ be a point of $X$ such that $I_x\ne0$. Let $\underline{p}$ be a prime ideal of $\underline{O}_{Y,x}$ containing $I_x$. We have $\dim(\underline{O}_{Y,x})_{\underline{p}}>0$, since otherwise we would have $\underline{p}=0$, as $\underline{O}_{Y,x}$ is an integral domain. Thus $X$ has codimension $>0$ at $x$, hence codimension $>0$ in a neighbourhood of $x$ ([10] III A Cor 1 to Lemma 3), and we are done.

\sgaStableTarget{sga6.e02.sec.3}\section*{3. Complements on quasi-coherent sheaves on schemes\protect\SGAAuthorFootnote{2}{This section is essentially due to J.L. Verdier.}}

\sgaStableTarget{sga6.e02.para.3-0}3.0. In this section, we fix a \uline{quasi-compact and quasi-separated} scheme $S$. The canonical (fully faithful) inclusion functor
\[
 \phi:\Qcoh(S)\longrightarrow\Mod(S)\qquad\text{\SGAAuthorFootMark{3}}
\]\SGAAuthorFootText{3}{$\Mod(S)$ (resp. $\Qcoh(S)$) denotes the category of $\underline{O}_S$-Modules (resp. quasi-coherent $\underline{O}_S$-Modules).}
extends to a functor, still denoted by $\phi$, from $\mathrm D(\Qcoh(S))$ to $\mathrm D(S)=\mathrm D(\Mod(S))$, which we propose to study\SGAAuthorFootnote{4}{We shall omit $\phi$ in formulas where there is no risk of confusion.}.

\sgaStableTarget{sga6.e02.lem.3-1}\textbf{Lemma 3.1.} \uline{The category} $\Qcoh(S)$ \uline{is an abelian category satisfying} (AB 5) \uline{and for which the quasi-coherent sheaves of finite type form a \textquotedblleft small\textquotedblright{} family of generators. In particular} (Tohoku 1.10.1), $\Qcoh(S)$ \uline{has enough injectives.}

%%%%%%%%%%  scan idx 200  |  folio 187  %%%%%%%%%%
\SGAPrintPage{187}%
\uline{Proof.} The assertion concerning generators follows immediately from the fact that every quasi-coherent sheaf on $S$ is the filtered inductive limit of its quasi-coherent subsheaves of finite type (EGA I 9.4.9 and IV 1.7.7). The other assertions are evident.

\sgaStableTarget{sga6.e02.rem.3-1-1}\textbf{Remark 3.1.1.} The editor does not know whether $\Qcoh(S)$ still has enough injectives when one assumes only that $S$ is quasi-separated.

\sgaStableTarget{sga6.e02.lem.3-2}\textbf{Lemma 3.2.} \uline{The functor} $\phi:\Qcoh(S)\longrightarrow\Mod(S)$ \uline{has a right adjoint}
\[
 Q:\Mod(S)\longrightarrow\Qcoh(S),
\]
\uline{(called the} coherer).

\uline{Proof.} When $S$ is affine, the existence of $Q$ is immediate: take $Q=\widetilde{\Gamma}_S$ (EGA $\mathrm{Err}_{\mathrm I}$ 1.8.7). Now let $f:U\longrightarrow S$ be an affine open subset of $S$; it is clear that, for $E\in\Ob\Qcoh(S)$ and $F\in\Ob\Mod(U)$, we have a canonical isomorphism
\sgaStableTarget{sga6.e02.eq.1}\begin{equation}\tag{1}
 \Hom(E,f_*(F))\simeq \Hom(E,f_*Q(F)).
\end{equation}
Then let $(f_i:U_i\longrightarrow S)$ be a finite covering of $S$ by affine open subsets, and for each $(i,j)$ let $(U_{ijk}\longrightarrow U_{ij})$ be a finite covering of $U_{ij}=U_i\cap U_j$ by affine open subsets ($U_{ij}$ is quasi-compact since $S$ is quasi-separated). Let $F\in\Ob\Mod(S)$. We have an exact sequence
\sgaStableTarget{sga6.e02.eq.2}\begin{equation}\tag{2}
 0\longrightarrow F\longrightarrow \coprod_i f_{i*}(F|U_i)
 \longrightarrow \coprod_{ijk}f_{ijk*}(F|U_{ijk}),
\end{equation}
(where $f_{ijk}$ denotes the inclusion of $U_{ijk}$ into $S$). Note that, for an inclusion $g:V\longrightarrow U$ of affine open subsets of $S$, we have an evident natural arrow $Q(F|U)\longrightarrow g_*Q(F|V)$. Consequently, (2) gives an arrow
\sgaStableTarget{sga6.e02.eq.3}\begin{equation}\tag{3}
 \coprod_i f_{i*}Q(F|U_i)\longrightarrow
 \coprod_{ijk}f_{ijk*}Q(F|U_{ijk}),
\end{equation}
whose kernel we shall denote by $QF$. Let $E\in\Ob\Qcoh(S)$. Applying the functor $\Hom(E,.)$ to the sequence (2), we obtain, in view of (1), an
%%%%%%%%%%  scan idx 201  |  folio 188  %%%%%%%%%%
\SGAPrintPage{188}%
isomorphism functorial in $E$:
\[
 \Hom(E,F)\simeq\Hom(E,QF).
\]
This proves the existence of $Q$; uniqueness is evident, since it is an adjoint.

\sgaStableTarget{sga6.e02.lem.3-3}\textbf{Lemma 3.3.} \begin{enumerate}[label=\alph*)]
\item \uline{The coherer is an additive functor taking injectives to injectives.}
\item \uline{Write}
\[
 \underline{\underline{R}}Q:\mathrm D^+(S)\longrightarrow \mathrm D^+(\Qcoh(S))
\]
\uline{for the right derived functor of the coherer. For} $E\in\Ob\mathrm D(\Qcoh(S))$ and $F\in\Ob\mathrm D^+(S)$\uline{, there are canonical functorial isomorphisms}
\sgaStableTarget{sga6.e02.eq.i}\sgaStableTarget{sga6.e02.eq.ii}\begin{align}
 \underline{\underline{R}}\Hom_{\Mod}(E,F)&\simeq \underline{\underline{R}}\Hom_{\Qcoh}(E,\underline{\underline{R}}Q(F)) \qquad\text{\SGAAuthorFootMark{1}}\tag{i}\\
 \Hom_{\Mod}(E,F)&\simeq \Hom_{\Qcoh}(E,\underline{\underline{R}}Q(F)). \tag{ii}
\end{align}\SGAAuthorFootText{1}{The subscript $\Mod$ (resp. $\Qcoh$) is an abbreviation for $\mathrm D(S)$ (resp. $\mathrm D(\Qcoh(S))$).}
\end{enumerate}

\uline{Proof.} Assertion a) is an immediate consequence of the definition. For b), it suffices to establish the isomorphism (i), since (ii) will follow by passing to cohomology. For this, we may assume that $F$ has injective components; $QF$ then has injective components by a), and the isomorphism (i) comes simply from the definition of $Q$.

\sgaStableTarget{sga6.e02.lem.3-4}\textbf{Lemma 3.4.} \uline{Let} $f:X\longrightarrow Y$ \uline{be a morphism of quasi-compact and quasi-separated schemes. For} $E\in\Ob\Mod(X)$\uline{, there is a canonical functorial isomorphism}
\[
 Qf_*(E)\simeq f_*Q(E).
\]

%%%%%%%%%%  scan idx 202  |  folio 189  %%%%%%%%%%
\SGAPrintPage{189}%
\uline{Proof.} Since $f$ is quasi-compact and quasi-separated (EGA IV 1.2.3 and 1.2.4), $f_*$ takes quasi-coherent sheaves to quasi-coherent sheaves (III 1.4.10 and IV 1.7.21). The functors $Qf_*$ and $f_*Q$ from $\Mod(X)$ to $\Qcoh(Y)$ are both right adjoints to the functor $f^*:\Qcoh(Y)\longrightarrow\Mod(X)$, and hence canonically isomorphic.

\sgaStableTarget{sga6.e02.prop.3-5}\textbf{Proposition 3.5.} \uline{When} $S$ \uline{is noetherian, or quasi-compact and separated, the functor}
\begin{equation}\tag{3.0}
 \phi:\mathrm D^+(\Qcoh(S))\longrightarrow \mathrm D^+(S)
\end{equation}
\uline{is fully faithful and has as essential image the full subcategory} $\mathrm D^+(S)_{\qcoh}$ of $\mathrm D^+(S)$ \uline{formed by complexes with quasi-coherent cohomology.}

\uline{Proof.} We may restrict ourselves to proving the assertion in the case where $S$ is separated, the case where $S$ is noetherian being well known ([H] II 7.19). It is evidently equivalent to prove the following assertions:

\begin{enumerate}[label={(3.5.\arabic*)~:}]
\item \uline{for} $E\in\Ob\mathrm D^+(\Qcoh(S))$, \uline{the adjunction arrow}
\[
 E\longrightarrow \underline{\underline{R}}Q(E)
\]
(\uline{deduced from} (3.3 b) (ii)) \uline{is an isomorphism};
\item \uline{for} $E\in\Ob\mathrm D^+(S)_{\qcoh}$, \uline{the adjunction arrow}
\[
 \underline{\underline{R}}Q(E)\longrightarrow E
\]
(\uline{deduced from} (3.3 b)(ii)) \uline{is an isomorphism}.
\end{enumerate}

\uline{Proof of} (3.5.1). By an immediate dévissage\SGAAuthorFootnote{1}{Using, for example, the spectral sequence $R^pQ(H^q(E))\Longrightarrow R^{*}Q(E)$.}, we are reduced to proving that quasi-coherent sheaves are acyclic for the coherer, (i.e.\ that, for $E\in\Ob\Qcoh(S)$, we have $R^iQ(E)=0$ for $i>0$), and that, on the other hand, for $E\in\Ob\Qcoh(S)$, the adjunction arrow $E\longrightarrow QE$ is an isomorphism. When $S$ is affine, this follows from the formula $Q=\widetilde{\Gamma}_S$ (cf. proof of (3.2)). In the general case, cover $S$ by finitely many affine open subsets $U_i$. Since $S$ is separated, the inclusions $f_i:U_i\longrightarrow S$ are affine morphisms,
%%%%%%%%%%  scan idx 203  |  folio 190  %%%%%%%%%%
\SGAPrintPage{190}%
as are the inclusions $f_{i_0\ldots i_p}:U_{i_0\ldots i_p}=U_{i_0}\cap\ldots\cap U_{i_p}\longrightarrow S$. Let $E\in\Ob\Qcoh(S)$, and consider the resolution of $E$ defined by the covering $(U_i)$ (T.F. II 5.2):
\sgaStableTarget{sga6.e02.eq.unnamed.occ3}\begin{equation}\tag{*}
 0\longrightarrow E\longrightarrow \coprod f_{i*}f_i^*(E)\longrightarrow\ldots\longrightarrow
 \coprod f_{i_0\ldots i_p*}f_{i_0\ldots i_p}^*(E)\longrightarrow .
\end{equation}
It gives a spectral sequence
\sgaStableTarget{sga6.e02.eq.unnamed.occ4}\begin{equation}\tag{**}
 E_1^{pq}=\coprod_{i_0\ldots i_p}R^qQ\bigl(f_{i_0\ldots i_p*}f_{i_0\ldots i_p}^*(E)\bigr)
 \Longrightarrow R^{*}Q(E).
\end{equation}
Since $f_{i_0\ldots i_p}$ is affine, (3.4) gives:
\sgaStableTarget{sga6.e02.eq.unnamed.occ5}\begin{equation}\tag{***}
 R^qQ\bigl(f_{i_0\ldots i_p*}f_{i_0\ldots i_p}^*(E)\bigr)
 \simeq f_{i_0\ldots i_p*}R^qQ\bigl(f_{i_0\ldots i_p}^*(E)\bigr)
\end{equation}
for every $q$, whence in particular
\[
 E_1^{pq}=0\quad\text{for }q>0,
\]
by the affine case, already treated. From (**) we deduce isomorphisms
\[
 R^nQ(E)\simeq H^n(QE^\bullet)\quad\text{for }n\geq0,
\]
where $E^\bullet$ is the resolution of $E$ defined by the exact sequence (*). But (***) gives $QE^n\simeq E^n$ for $n\geq0$, whence finally
\[
 R^nQ(E)=0\quad\text{for }n>0,
\]
and $Q(E)\simeq E$, q.e.d.

\uline{Proof of} (3.5.2). The spectral sequence
\[
 R^pQ(H^q(E))\Longrightarrow R^{*}Q(E),
\]
%%%%%%%%%%  scan idx 204  |  folio 191  %%%%%%%%%%
\SGAPrintPage{191}%
immediately reduces us to the case where $E\in\Ob\Qcoh(S)$. Since the composite
\[
 E\longrightarrow \underline{\underline{R}}QE\longrightarrow E
\]
of the adjunction arrows is the identity, the assertion therefore follows from (3.5.1). This completes the proof of (3.5).

\sgaStableTarget{sga6.e02.rem.3-6}\textbf{Remark 3.6.} The conclusion of (3.5) would be false with the hypothesis \textquotedblleft separated or noetherian\textquotedblright{} omitted, as a neat counterexample of J.L. Verdier shows (see Appendix I to this exposé).

\sgaStableTarget{sga6.e02.prop.3-7}\textbf{Proposition 3.7.} \uline{Under the hypothesis of} (3.5)\uline{, suppose further that there exists an integer} $N$ \uline{such that, for every} $F\in\Ob\Mod(S)$ \uline{and every quasi-compact open subset} $U$ \uline{of} $S$,
\[
 H^i(U,F)=0\quad\text{for }i>N\qquad\text{\SGAAuthorFootMark{1}}\ .
\]
\SGAAuthorFootText{1}{By (Tohoku 3.6.5), this condition holds when the underlying space of $S$ is noetherian and has Krull dimension $\leq N$.}
\uline{Then:}
\begin{enumerate}[label=\alph*)]
\item \uline{The coherer has finite cohomological dimension.}
\item \uline{The functor}
\[
 \phi:\mathrm D(\Qcoh(S))\longrightarrow \mathrm D(S)
\]
\uline{is fully faithful and has as essential image the full subcategory} $\mathrm D(S)_{\qcoh}$ \uline{of} $\mathrm D(S)$ \uline{formed by complexes with quasi-coherent cohomology.}
\end{enumerate}

\uline{Proof of} a). When $S$ is affine, the formula $Q=\widetilde{\Gamma}_S$ (cf. proof of (3.2)) implies that $Q$ has cohomological dimension $\leq N$. We pass to the general case. Let us show that there exists an integer $a$ such that $R^qQ(E)=0$ for every $E\in\Ob\Mod(S)$ and every $q>a$. The hypothesis implies that every $E\in\Ob\Mod(S)$ admits a finite right resolution $E^\bullet$ of length $N$ such that the $E^i$ are acyclic, i.e.\ satisfy $H^q(U,E^i)=0$ for every open subset $U$ of $S$ and every $q>0$. It therefore suffices to prove that there exists an integer $b$ such that $R^qQ(E)=0$ for every acyclic $\underline{O}_S$-Module $E$ and every $q>b$. (We may then take $a=b+N$). Let $(U_i)$ be a finite covering of $S$ by $N'$ affine
%%%%%%%%%%  scan idx 205  |  folio 192  %%%%%%%%%%
\SGAPrintPage{192}%
open subsets. Let $E$ be an acyclic $\underline{O}_S$-Module, and consider the resolution $C^\bullet(E)$ of $E$ defined by the covering $(U_i)$ (T.F. II 5.2). There are two cases to examine:

1°) $S$ is \uline{separated}. The components of $C^\bullet(E)$ are then sums of terms of the form $f_*f^*(E)$, where $f:U\longrightarrow S$ is the inclusion of an affine open subset (hence an affine morphism, since $S$ is separated). We have
\[
 \underline{\underline{R}}(Qf_*)f^*E\simeq \underline{\underline{R}}Q\,\underline{\underline{R}}f_*f^*E\simeq \underline{\underline{R}}Q\,f_*f^*E,
\]
since $f^*E$ is acyclic. On the other hand, by (3.4), we have
\[
 \underline{\underline{R}}(Qf_*)f^*E\simeq \underline{\underline{R}}(f_*Q)f^*E\simeq f_*\underline{\underline{R}}Qf^*E,
\]
since $f$ is affine. We deduce that
\[
 R^nQ(f_*f^*E)\simeq f_*R^nQ(f^*E)
\quad\text{for every }n\geq0,
\]
whence, by the affine case,
\[
 R^nQ(f_*f^*E)=0\quad\text{for }n>N.
\]
Since $C^\bullet(E)$ has length $N'$, we finally have
\[
 R^nQ(E)=0\quad\text{for }n>N+N',
\]
and we are done.

2°) $S$ is \uline{noetherian}. The components of $C^\bullet(E)$ are sums of terms of the form $f_*f^*(E)$, where $f:U\longrightarrow S$ is the inclusion of a quasi-compact and \uline{separated} open subset. As in case 1°), we have
%%%%%%%%%%  scan idx 206  |  folio 193  %%%%%%%%%%
\SGAPrintPage{193}%
\begin{equation}\tag{*}
 \underline{\underline{R}}(Qf_*)f^*E\simeq \underline{\underline{R}}Q\,f_*f^*E\simeq \underline{\underline{R}}(f_*Q)f^*E.
\end{equation}
Since $U$ is noetherian, $\Qcoh(U)$ has enough objects that are injective in $\Mod(U)$ ([H] II 7.18), so the right derived functor of $f_*:\Qcoh(U)\longrightarrow\Qcoh(S)$ is \textquotedblleft induced\textquotedblright{} by $\underline{\underline{R}}f_*:\mathrm D^+(U)\longrightarrow\mathrm D^+(S)$ when we identify $\mathrm D^+(\Qcoh(U))$ (resp. $\mathrm D^+(\Qcoh(S))$) with a full subcategory of $\mathrm D^+(U)$ (resp. $\mathrm D^+(S)$) by (3.5). Consequently, we have
\[
 \underline{\underline{R}}(f_*Q)f^*E\simeq \underline{\underline{R}}f_*\underline{\underline{R}}Qf^*E.
\]
We deduce a spectral sequence
\[
 E_2^{pq}=R^pf_*R^qQf^*E\Longrightarrow R^{*}(f_*Q)f^*E.
\]
By case 1°), there exists an integer $n(U)$ depending only on $U$ such that $Q$ has cohomological dimension $\leq n(U)$ on $\Mod(U)$. Thus we have
\[
 E_2^{pq}=0\quad\text{for }q>n(U).
\]
But we have
\[
 E_2^{pq}=0\quad\text{for }p>N,
\]
because of the \textquotedblleft cohomological dimension\textquotedblright{} hypothesis on $S$. We therefore have
\[
 R^n(f_*Q)f^*E=0\quad\text{for }n>N+N'',
\]
where $N''$ is an integer greater than all the $n(U)$, with $U$ ranging over the set of intersections $U_{i_0}\cap\ldots\cap U_{i_p}$. Taking account of (*), we obtain
\[
 R^nQ(f_*f^*E)=0\quad\text{for }n>N+N'',
\]
whence finally
\[
 R^nQ(E)=0\quad\text{for }n>N+N'+N'',
\]
%%%%%%%%%%  scan idx 207  |  folio 194  %%%%%%%%%%
\SGAPrintPage{194}%
which completes the proof of a).

\uline{Proof of} b). The coherer, having finite cohomological dimension by a), admits ([H] I 5.3) a right derived functor
\[
 \underline{\underline{R}}Q:\mathrm D(S)\longrightarrow\mathrm D(\Qcoh(S))
\]
(whose restriction to $\mathrm D^+(S)$ coincides with the functor $\underline{\underline{R}}Q$ defined above). For $E\in\Ob\mathrm D(\Qcoh(S))$, there is a canonical functorial isomorphism
\sgaStableTarget{sga6.e02.eq.1.occ2}\begin{equation}\tag{1}
 E\xrightarrow{\sim}\underline{\underline{R}}Q(E),
\end{equation}
since, by (3.5.1), the $E^i$ are acyclic for $Q$ and satisfy $E^i\xrightarrow{\sim}Q(E^i)$. On the other hand, for $E\in\Ob\mathrm D(S)$, we have a canonical functorial homomorphism
\sgaStableTarget{sga6.e02.eq.2.occ2}\begin{equation}\tag{2}
 \underline{\underline{R}}Q(E)\longrightarrow E
\end{equation}
(coming, in the case where the $E^i$ are $Q$-acyclic, from the ordinary adjunction homomorphism between $\phi$ and $Q$). We check trivially that the homomorphisms (1) and (2) make $\underline{\underline{R}}Q:\mathrm D(S)\longrightarrow\mathrm D(\Qcoh(S))$ a right adjoint to $\phi:\mathrm D(\Qcoh(S))\longrightarrow\mathrm D(S)$. The fact that (1) is an isomorphism implies that $\phi$ is fully faithful. Finally, using the spectral sequence
\[
 R^pQ(H^q(E))\Longrightarrow R^{*}Q(E),
\]
(which converges since $Q$ has finite cohomological dimension), we see that (2) is an isomorphism for $E\in\Ob\mathrm D(S)_{\qcoh}$, so $\phi$ has as essential image $\mathrm D(S)_{\qcoh}$. This proves b) and completes the proof of (3.7).

%%%%%%%%%%  scan idx 208  |  folio 195  %%%%%%%%%%
\SGAPrintPage{195}%
\sgaStableTarget{sga6.e02.appendix-i-a-counterexample-of-verdier}\section*{APPENDIX I\\[3pt] A COUNTEREXAMPLE OF VERDIER\protect\SGAAuthorFootnote{1}{Apart from a few details of presentation, this appendix reproduces two letters from J.-L. Verdier (30/9/67 and 11/10/67) to the editor, in reply to a question from the latter.}}

\sgaStableTarget{sga6.e02.sec.0}\subsection*{0. Introduction}
For a scheme $S$, we denote by $\Mod(S)$ (resp. $\Qcoh(S)$) the category of $\underline{O}_S$-Modules (resp. quasi-coherent $\underline{O}_S$-Modules). The purpose of this appendix is to draw the reader's attention to certain pathological properties of the category $\mathrm D(\Qcoh(S))$, and in particular to prove the following facts:

\sgaStableTarget{sga6.e02.para.0-1}0.1. Let $X$ be an affine scheme with ring $A$, and let $I$ be an injective $A$-module. We know ([H] II 7.14) that, when $A$ is noetherian, the sheaf $\widetilde I$ is an injective $\underline{O}_X$-Module. When $A$ is not noetherian, $\widetilde I$ need not be injective, or even flasque, even when the underlying space of $X$ is noetherian and of finite Krull dimension. Moreover, the restriction of $\widetilde I$ to an open subset $U$ of $X$ is not in general an injective object of $\Qcoh(U)$.

\sgaStableTarget{sga6.e02.para.0-2}0.2. Let $f:X\longrightarrow Y$ be a morphism of quasi-compact and quasi-separated schemes. The functor $\underline{\underline{R}}f_*:\mathrm D^+(X)\longrightarrow\mathrm D^+(Y)$ does not in general induce on $\mathrm D^+(\Qcoh(X))$ (a full subcategory of $\mathrm D^+(X)$ by (II 3.5)) the right derived functor of $f_*:\Qcoh(X)\longrightarrow\Mod(Y)$. More precisely, an injective object of $\Qcoh(X)$ need not be acyclic for $f_*:\Mod(X)\longrightarrow\Mod(Y)$.

\sgaStableTarget{sga6.e02.para.0-3}0.3. Let $S$ be a quasi-compact and quasi-separated scheme. An object of $\Qcoh(S)$, even an injective one, need not be acyclic for the coherer (II 3.2). The functor $\phi:\mathrm D^b(\Qcoh(S))\longrightarrow\mathrm D^b(S)$ (II 3.0) is not in general fully faithful.

\sgaStableTarget{sga6.e02.sec.1.occ2}\subsection*{1. The counterexamples}
We first recall the following lemma (SGA 2 II 9):

\sgaStableTarget{sga6.e02.lem.1-1}\textbf{Lemma 1.1.} \uline{Let} $A$ \uline{be a ring,} $X=\Spec(A)$, $\underline{f}=(f_\alpha)$ \uline{a finite system of elements of} $A$, $Y$ \uline{the closed subset defined by} $\underline{f}(x)=0$\uline{. For} $i>0$\uline{, the following conditions are equivalent:}
\begin{enumerate}[label=(\roman*)]
\item \uline{for every injective} $A$\uline{-module} $I$\uline{, we have} $H^i_Y(X,\widetilde I)=0$;
\item \uline{the pro-object} $n\mapsto H_i(\underline{f}^n,A)$ \uline{is zero, i.e. for every} $n$\uline{, there exists} $n'\geq n$
%%%%%%%%%%  scan idx 209  |  folio 196  %%%%%%%%%%
\SGAPrintPage{196}%
\uline{such that}
\[
 H_i(\underline{f}^{n'},A)\longrightarrow H_i(\underline{f}^n,A)
\]
\uline{is zero}.
\end{enumerate}

\sgaStableTarget{sga6.e02.para.1-2}1.2 We shall use (1.1) to prove (0.1). Let $k$ be a field, $B=k[[x,y]]$, $\underline{f}$ the system $(x,y)$, $M=\coprod_{n>0}B/(\underline{f}^n)$, $A$ the $B$-algebra $D_B(M)$ (EGA $0_{\mathrm{IV}}$ 18.2.3). We have trivially
\[
 H_2^A(\underline{f}^n,A)\simeq H_2^B(\underline{f}^n,A).\qquad\text{\SGAAuthorFootMark{1}}
\]
\SGAAuthorFootText{1}{The superscript A or B indicates the ring over which the Koszul complex is computed.}
On the other hand, we have
\[
 H_2^B(\underline{f}^n,A)\simeq H_2^B(\underline{f}^n,B)\oplus H_2^B(\underline{f}^n,M)
\]
but $H_2^B(\underline{f}^n,B)=0$ since $\underline{f}$ is a regular sequence (EGA III 1.1.4 and $0_{\mathrm{IV}}$ 15.1.20), so we have
\[
 H_2^B(\underline{f}^n,A)\simeq H_2^B(\underline{f}^n,M).
\]
Finally, we have (EGA III 1.1.3.3)
\[
 H_2^B(\underline{f}^n,M)\simeq H_B^0(\underline{f}^n,M),
\]
where $H_B^0(\underline{f}^n,M)$ is canonically identified with the submodule ${}_{(\underline{f}^n)}M$ of $M$ annihilated by $(\underline{f}^n)$. Hence finally:
\sgaStableTarget{sga6.e02.eq.unnamed.occ6}\begin{equation}\tag{*}
 H_2^A(\underline{f}^n,A)\simeq {}_{(\underline{f}^n)}M.
\end{equation}
Since the transition morphism $H_2^A(\underline{f}^{n'},A)\longrightarrow H_2^A(\underline{f}^n,A)$ is induced by multiplication by $\underline{f}^{\,n'-n}$ (EGA III 1.1.6), we check immediately, using (*), that \uline{the pro-object} $n\mapsto H_2^A(\underline{f}^n,A)$ \uline{is not zero}. From (1.1) we then deduce that there exists an injective $A$-module $I$ such that $H^2_Y(X,\widetilde I)\ne0$, and thus such that $\widetilde I$ is not flasque (SGA 2 I 2.12). We note in passing that the projection $\Spec(A)\longrightarrow\Spec(B)$ is a homeomorphism,
%%%%%%%%%%  scan idx 210  |  folio 197  %%%%%%%%%%
\SGAPrintPage{197}%
so the underlying space of $X$ is noetherian of dimension $2$. Let $I$ be an injective $A$-module such that $H^2_Y(X,\widetilde I)\ne0$, and denote by $U$ the open complement of $Y$ in $X$. We have (SGA 2 II 4)
\[
 H^2_Y(X,I)\simeq H^1(U,I)\ne0.
\]
Now, since $U$ is separated, the functor $\mathrm D^+(\Qcoh(U))\to\mathrm D^+(U)$ is fully faithful (II 3.5), so $\widetilde I|U$ is not injective in $\Qcoh(U)$. We have therefore proved assertion (0.1).

Denote by $Z$ the quasi-compact and quasi-separated scheme obtained by gluing two copies of $X$ along $U$, and by $j_1,j_2$ the two canonical open immersions of $X$ into $Z$, so that we have a cartesian square
\[
\begin{tikzcd}
 U \arrow[r,"i"] \arrow[d,"i"'] & X \arrow[d,"j_1"]\\
 X \arrow[r,"j_2"] & Z .
\end{tikzcd}
\]
We have
\[
 j_2^*R^1j_{1*}(\widetilde I)\simeq R^1i_*(i^*\widetilde I),
\]
whence
\[
 H^0(X,R^1j_{1*}(\widetilde I))\simeq H^1(U,\widetilde I)\ne0,
\]
so $R^1j_{1*}(\widetilde I)\ne0$, which proves (0.2).

On the other hand, by (II 3.4 and 3.5), we have
\[
 \underline{\underline{R}}Q\,\underline{\underline{R}}j_{1*}(\widetilde I)\simeq j_{1*}(\widetilde I).
\]
We deduce a spectral sequence whose exact sequence of low-degree terms gives the isomorphism
\[
 R^1j_{1*}(\widetilde I)\simeq R^2Q(j_{1*}\widetilde I).
\]
Since the sheaf $j_{1*}\widetilde I$ is injective in $\Qcoh(Z)$, we have proved the first
%%%%%%%%%%  scan idx 211  |  folio 198  %%%%%%%%%%
\SGAPrintPage{198}%
part of (0.3).

Finally, let $E\in\Ob\Qcoh(Z)$. We have the duality isomorphism
\[
 \underline{\underline{R}}\Hom_{\Mod}(E,\underline{\underline{R}}j_{1*}(\widetilde I))
 \simeq \underline{\underline{R}}\Hom_{\Mod}(j_1^*E,\widetilde I),
\]
but, since $X$ is affine, we have (II 3.5)
\[
 \underline{\underline{R}}\Hom_{\Mod}(j_1^*E,\widetilde I)
 \simeq \underline{\underline{R}}\Hom_{\Qcoh}(j_1^*E,I)
 \simeq \Hom(j_1^*E,\widetilde I).
\]
We deduce
\[
 \underline{\underline{R}}\Hom_{\Mod}(E,\underline{\underline{R}}j_{1*}(\widetilde I))
 \simeq \Hom(j_1^*E,\widetilde I),
\]
whence a spectral sequence whose exact sequence of low-degree terms gives the isomorphism
\[
 \Ext^2_{\Mod}(E,j_{1*}\widetilde I)
 \simeq \Hom(E,R^1j_{1*}\widetilde I).
\]
Consequently, the identity morphism of $E=R^1j_{1*}(\widetilde I)$ gives a nonzero element of $\Ext^2_{\Mod}(E,j_{1*}\widetilde I)$, so the functor $\mathrm D^b(\Qcoh(Z))\longrightarrow \mathrm D^b(Z)$ is not fully faithful, which completes the proof of (0.3).

% Original English continuation definitions, retained in their original order.
\providecommand{\Diff}{\operatorname{Diff}}
\providecommand{\Hom}{\operatorname{Hom}}
\providecommand{\im}{\operatorname{Im}}
\providecommand{\ind}{\operatorname{ind}}
\providecommand{\Cone}{\operatorname{Cone}}
\providecommand{\Loclib}{\operatorname{Loclib}}
\providecommand{\ob}{\operatorname{ob}}
\providecommand{\Parf}{\operatorname{Parf}}
\providecommand{\Mod}{\operatorname{Mod}}
\providecommand{\Ob}{\operatorname{Ob}}
\providecommand{\Coker}{\operatorname{Coker}}
\providecommand{\Ker}{\operatorname{Ker}}
\providecommand{\cH}{\mathcal H}
\providecommand{\cO}{\mathcal O}
\providecommand{\cP}{\mathcal P}
\providecommand{\cE}{\mathcal E}
\providecommand{\cF}{\mathcal F}
\providecommand{\cG}{\mathcal G}
\providecommand{\cD}{\mathcal D}
\providecommand{\cT}{\mathcal T}
\providecommand{\cL}{\mathcal L}
\providecommand{\RHom}{\mathrm R\!\operatorname{Hom}}
\providecommand{\RGamma}{\mathrm R\Gamma}
\providecommand{\dlim}{\operatorname*{lim}}
\providecommand{\plim}{\operatorname*{lim}}

% Reviewed French continuation definitions, retained verbatim.
% Additions for Expose II, Appendix II
\providecommand{\cO}{\mathcal O}
\providecommand{\cE}{\mathscr E}
\providecommand{\cF}{\mathscr F}
\providecommand{\cG}{\mathscr G}
\providecommand{\cH}{\mathcal H}
\providecommand{\cL}{\mathscr L}
\providecommand{\cP}{\mathscr P}
\providecommand{\cM}{\mathscr M}
\providecommand{\cD}{\mathcal D}
\providecommand{\R}{\mathbb R}
\providecommand{\C}{\mathbb C}
\providecommand{\N}{\mathbb N}
\providecommand{\Zell}{\mathbb Z_\ell}
\providecommand{\Qell}{\mathbb Q_\ell}
\providecommand{\C}{\mathbb C}
\providecommand{\Ch}{\operatorname{Ch}}
\providecommand{\Sym}{\operatorname{Sym}}
\providecommand{\Td}{\operatorname{Td}}
\providecommand{\Diff}{\operatorname{Diff}}
\providecommand{\Homop}{\operatorname{Hom}}
\providecommand{\Spec}{\operatorname{Spec}}
\providecommand{\pr}{\operatorname{pr}}
\providecommand{\ind}{\operatorname{ind}}
\providecommand{\soft}{\operatorname{soft}}
\providecommand{\Ban}{\operatorname{Ban}}
\providecommand{\Lop}{\operatorname{L}}
\providecommand{\im}{\operatorname{Im}}
\providecommand{\Parf}{\operatorname{Parf}}
\renewcommand{\id}{\operatorname{id}}

\newpage
%%%%%%%%%%  scan idx 212  |  folio 199  %%%%%%%%%%
\SGAPrintPage{199}%
\sgaStableTarget{sga6.e02.appendix-ii-definition-of-the-analytic-index-of}\section*{Appendix II\\ Definition of the analytic index of a relative elliptic complex}

In this appendix, we propose to show that the notion of a perfect complex leads to a simple and natural definition of the \textquotedblleft analytic index\textquotedblright{} of a continuous family of elliptic differential operators, or more generally of a \textquotedblleft relative elliptic complex\textquotedblright. For the initiated reader, let us say that the definition given here has the advantage over the classical definitions (see for example [2][17]) of requiring no deformation of the elliptic complex under consideration.

Sections 1 and 2 consist of reminders about the classical notion of a \textquotedblleft mixed structure\textquotedblright\SGAAuthorFootnote{1}{The informed reader is advised to skip them.}. We give, of course, only the minimum of definitions and results required to state and prove, in \S3, the central theorem we have in view, which says, a little more precisely, that the direct image, under the projection (assumed proper) onto the base (assumed locally compact), of a relative elliptic complex is a perfect complex. When the base $S$ is compact, the definition of the analytic index as an element of $K(S)$ (\textquotedblleft Grothendieck group\textquotedblright{} of complex vector bundles of finite rank on $S$) then follows from the globalization sorite of (II 2.3.2).

Throughout what follows, for a topological space $S$, we shall denote by $\underline{O}_S$ the sheaf of continuous real-valued functions on $S$.

\sgaStableTarget{sga6.e02.sec.1.occ3}\subsection*{1. The sorite of mixed structures}

In this section, we fix a topological space $S$.

\sgaStableTarget{sga6.e02.sec.1.1}\subsubsection*{1.1. Definition of mixed varieties}
Let $V$ be an open subset of $\underline{\underline{R}}^n$. For $r\in\underline{\underline{N}}$, we denote by $C^r(V)$ the Fréchet space of real-valued functions of class $C^r$ on $V$. We put
\[
  C^\infty(V)=\varprojlim_r C^r(V).
\]
Let $U$ be an open subset of $S$. For $r\in\underline{\underline{N}}\cup\{\infty\}$, we denote by ${}^r\underline{O}_{U\times V}$ the sheaf of rings on $U\times V$ defined by
%%%%%%%%%%  scan idx 213  |  folio 200  %%%%%%%%%%
\SGAPrintPage{200}%
\[
  \Gamma(U'\times V',{}^r\underline{O}_{U'\times V'})=C(U',C^r(V')),
\]
where $U'$ (resp. $V'$) is an open subset of $U$ (resp. $V$), and $C(U',C^r(V'))$ denotes the space of continuous maps from $U'$ to $C^r(V')$. There is a canonical injection
\[
  \pr_1^{-1}(\underline{O}_U)\longrightarrow{}^r\underline{O}_{U\times V},
\]
which makes $(U\times V,{}^r\underline{O}_{U\times V})$ a ringed space over $U$.

A ringed space $X$ over $S$ is called a \uline{mixed variety of class} $C^r$ ($r\in\underline{\underline{N}}\cup\{\infty\}$) when $X$ is locally isomorphic (as a ringed space over $S$) to a space of the form $(U\times V,{}^r\underline{O}_{U\times V})$. (When $S$ consists of one point, a mixed variety of class $C^r$ over $S$ is a variety of class $C^r$ in the usual sense).

Let $f:X\to S$ be a mixed variety of class $C^r$. For $s\in S$, the fiber $X_s=f^{-1}(s)$, equipped with the sheaf $i_s^{-1}(\underline{O}_X)/I_s$, where $i_s:X_s\to X$ is the injection and $I_s$ is the Ideal of sections that \textquotedblleft vanish\textquotedblright{} on $X_s$, is a variety of class $C^r$.

\sgaStableTarget{sga6.e02.sec.1.2}\subsubsection*{1.2. \textquotedblleft Relative\textquotedblright{} differential operators}
It is easy to make an infinitesimal study of mixed varieties on the model of Algebraic Geometry (EGA IV 16, SGA 3 VII). We shall limit ourselves to a brief review of the notions we shall need.

\sgaStableTarget{sga6.e02.para.1-2-1}1.2.1. Let $f:X\to S$ be a mixed variety of class $C^\infty$ (1.1). For every $n\in\underline{\underline{N}}$, we define the sheaf $P^n_{X/S}$ of \uline{principal parts of order} $n$ on $X$. In two different ways, corresponding to the two projections $\pr_1$ and $\pr_2$ from $X\times_S X$ onto $X$, it is an $\underline{O}_X$-Algebra augmented to $\underline{O}_X$, locally free of finite type as an $\underline{O}_X$-Module. When we do not specify the $\underline{O}_X$-Module structure considered on $P^n_{X/S}$, it will be understood to be the one coming from $\pr_1$. There are canonical isomorphisms
%%%%%%%%%%  scan idx 214  |  folio 201  %%%%%%%%%%
\SGAPrintPage{201}%
\[
  \begin{aligned}
  P^0_{X/S}&\simeq\underline{O}_X,\\
  P^1_{X/S}&\simeq\underline{O}_X\oplus\Omega^1_{X/S},
  \end{aligned}
\]
where $\Omega^1_{X/S}$ is the sheaf of relative differentials of $X$ over $S$, or \uline{relative cotangent sheaf}. We denote by $T_{X/S}$ the dual of $\Omega^1_{X/S}$, called the \uline{relative tangent sheaf}. The sheaf $T_{X/S}$ (resp. $\Omega^1_{X/S}$) is the sheaf of sections of a vector bundle of finite rank on $X$, called the \uline{relative} \uline{tangent bundle} (resp. \uline{cotangent bundle}), which is a mixed $S$-variety over $X$.

The $f^{-1}(\underline{O}_S)$-linear homomorphism (corresponding to the $\underline{O}_X$-Algebra structure defined by $\pr_2$)
\[
  d^n_{X/S}:\underline{O}_X\longrightarrow P^n_{X/S},
\]
which associates to a section of $\underline{O}_X$ its principal part of order $n$, is called the \uline{universal differential operator of order} $n$ on $X$ relative to $S$.

\sgaStableTarget{sga6.e02.para.1-2-2}1.2.2. More generally, for a locally free $\underline{O}_X$-Module $E$ of finite type, we can consider the sheaf $P^n_{X/S}(E)=P^n_{X/S}\otimes_{\underline{O}_X}E$ of principal parts of order $n$ of $E$ (the tensor product being defined by means of the second $\underline{O}_X$-Module structure on $P^n_{X/S}$), and we have a universal differential operator
\[
  d^n_{X/S}(E):E\longrightarrow P^n_{X/S}(E).
\]

Let $E,F$ be locally free $\underline{O}_X$-Modules of finite type. By definition, a differential operator from $E$ to $F$ of order $\le n$ is an $f^{-1}(\underline{O}_S)$-linear homomorphism from $E$ to $F$ that factors through $d^n_{X/S}(E)$ as an $\underline{O}_X$-linear homomorphism from $P^n_{X/S}(E)$ to $F$ (the factorization is then unique). In other words, the operator $d^n_{X/S}(E)$ gives a bijection
\sgaStableTarget{sga6.e02.eq.1-2-2-1}\begin{equation}\tag{1.2.2.1}
  \Homop_{\underline{O}_X}\bigl(P^n_{X/S}(E),F\bigr)
  \xrightarrow{\sim}
  \Diff^n(E,F),
\end{equation}
where $\Diff^n(E,F)$ denotes the set of differential operators from $E$
%%%%%%%%%%  scan idx 215  |  folio 202  %%%%%%%%%%
\SGAPrintPage{202}%
to $F$ of order $\le n$.

\sgaStableTarget{sga6.e02.para.1-2-3}1.2.3. We shall also have to consider differential operators with complex coefficients. Denote by $\underline{O}_S\otimes\underline{\underline{C}}$ (resp. $\underline{O}_X\otimes\underline{\underline{C}}$) the complexification of $\underline{O}_S$ (resp. $\underline{O}_X$). Let $E,F$ be locally free $(\underline{O}_X\otimes\underline{\underline{C}})$-Modules of finite type. By a differential operator from $E$ to $F$ of order $\le n$, we shall mean an $f^{-1}(\underline{O}_S\otimes\underline{\underline{C}})$-linear homomorphism from $E$ to $F$ that factors through $d^n_{X/S}(E)$ as an $(\underline{O}_X\otimes\underline{\underline{C}})$-linear homomorphism from $P^n_{X/S}(E)$ to $F$. Denoting the set of these operators by $\Diff^n_{\underline{\underline{C}}}(E,F)$, the operator $d^n_{X/S}(E)$ therefore defines a bijection
\sgaStableTarget{sga6.e02.eq.1-2-3-1}\begin{equation}\tag{1.2.3.1}
  \Homop_{\underline{O}_X\otimes\underline{\underline{C}}}\bigl(P^n_{X/S}(E),F\bigr)
  \xrightarrow{\sim}
  \Diff^n_{\underline{\underline{C}}}(E,F).
\end{equation}

\sgaStableTarget{sga6.e02.para.1-2-4}1.2.4. For $n\ge1$, there is a canonical exact sequence
\[
  0\longrightarrow S^n(\Omega^1_{X/S})
  \longrightarrow P^n_{X/S}
  \longrightarrow P^{n-1}_{X/S}
  \longrightarrow 0.\text{\SGAAuthorFootMark{1}}
\]
\SGAAuthorFootText{1}{$S^n(\ )$ denotes the $n$th symmetric power.}%
Hence, tensoring over $\underline{O}_X$\SGAAuthorFootnote{2}{For the second $\underline{O}_X$-Module structure of $P^n_{X/S}$ and $P^{n-1}_{X/S}$.} with a locally free $\underline{O}_X$-Module (resp. $(\underline{O}_X\otimes\underline{\underline{C}})$-Module) $E$ of finite type, we obtain an exact sequence of $\underline{O}_X$ (resp. $(\underline{O}_X\otimes\underline{\underline{C}})$)-Modules:
\sgaStableTarget{sga6.e02.eq.1-2-4-1}\begin{equation}\tag{1.2.4.1}
  0\longrightarrow S^n(\Omega^1_{X/S})\otimes_{\underline{O}_X}E
  \longrightarrow P^n_{X/S}(E)
  \longrightarrow P^{n-1}_{X/S}(E)
  \longrightarrow 0.
\end{equation}
Let $F$ be a locally free $\underline{O}_X$ (resp. $(\underline{O}_X\otimes\underline{\underline{C}})$)-Module of finite type. Applying the functor $\Homop_{\underline{O}_X}(.,F)$ (resp. $\Homop_{\underline{O}_X\otimes\underline{\underline{C}}}(.,F)$) to (1.2.4.1), we deduce, taking account of the identifications (1.2.2.1) and (1.2.3.1), an exact sequence
\sgaStableTarget{sga6.e02.eq.1-2-4-2}\begin{equation}\tag{1.2.4.2}
\begin{gathered}
  0\longrightarrow \Diff^{n-1}(E,F)
  \longrightarrow \Diff^n(E,F)
  \xrightarrow{\sigma_n}
  \Homop_{\underline{O}_X}\bigl(S^n(\Omega^1_{X/S})\otimes_{\underline{O}_X}E,F\bigr)\\
  (\text{resp. } 0\longrightarrow \Diff^{n-1}_{\underline{\underline{C}}}(E,F)
  \longrightarrow \Diff^n_{\underline{\underline{C}}}(E,F)
  \longrightarrow
  \Homop_{\underline{O}_X\otimes\underline{\underline{C}}}\bigl(S^n(\Omega^1_{X/S})\otimes_{\underline{O}_X}E,F\bigr)).
\end{gathered}
\end{equation}
%%%%%%%%%%  scan idx 216  |  folio 203  %%%%%%%%%%
\SGAPrintPage{203}%
(When $X$ is paracompact and $\underline{O}_X$ is soft (cf. II 2.3.1), we can complete the exact sequence (1.2.4.2) by adding a zero on the right). The homomorphism $\sigma_n$ is called the \uline{symbol}; it associates to a differential operator from $E$ to $F$ of order $\le n$ a \textquotedblleft homogeneous polynomial map of degree $n$ from $\Omega^1_{X/S}$ to $\underline{\Homop}_{\underline{O}_X}(E,F)$ (resp. $\underline{\Homop}_{\underline{O}_X\otimes\underline{\underline{C}}}(E,F)$)\textquotedblright, which we may also regard, if we wish, as a homomorphism $p^*(E)\longrightarrow p^*(F)$ (where $p$ is the projection onto $X$ from the relative cotangent bundle (1.2.1)), satisfying a certain homogeneity condition.

Let $E,F,G$ be locally free $\underline{O}_X$-Modules of finite type, and $u\in\Diff^m(E,F)$, $v\in\Diff^n(F,G)$, whence $vu\in\Diff^{m+n}(E,G)$. We easily check the relation
\sgaStableTarget{sga6.e02.eq.1-2-4-3}\begin{equation}\tag{1.2.4.3}
  \sigma_{m+n}(vu)=\sigma_n(v)\sigma_m(u).
\end{equation}
(The same relation holds for differential operators with complex coefficients). In particular, when the composite of two differential operators is zero, the composite of their symbols is also zero.

\sgaStableTarget{sga6.e02.para.1-2-5}1.2.5. Let $E$ be a complex, with bounded degrees, of differential operators with complex coefficients on $X$ relative to $S$ (1.2.3), with $d^i:E^i\longrightarrow E^{i+1}$ of order $\le n_i$. By (1.2.4.3), we have
\[
  \sigma_{n_{i+1}}(d^{i+1})\,\sigma_{n_i}(d^i)=0.
\]
We call $E$ an \uline{elliptic complex} of order $\underline{n}=(n_i)$ when the complex $(p^*(E^i),\sigma_{n_i}(d^i))$ (where $p$ is the projection onto $X$ from the cotangent bundle) is acyclic outside the zero section.

When $S$ consists of one point, the preceding notion of an elliptic complex on $X$ coincides with the usual notion, studied in detail in [4] or [16].

\sgaStableTarget{sga6.e02.para.1-2-6}1.2.6. Let $s\in S$, and let $i_s$ be the inclusion of the fiber $X_s$ into $X$ (considered as a monomorphism of ringed spaces) cf.~(1.1). For $n\in\underline{\underline{N}}$, there is a canonical isomorphism
%%%%%%%%%%  scan idx 217  |  folio 204  %%%%%%%%%%
\SGAPrintPage{204}%
\[
  i_s^*(P^n_{X/S})\simeq P^n_{X_s},
\]
whence a restriction homomorphism
\[
  i_s^*:\Diff^n(E,F)
  \longrightarrow
  \Diff^n\bigl(i_s^*(E),i_s^*(F)\bigr),
\]
for locally free $\underline{O}_X$-Modules $E,F$ of finite type (and an analogous homomorphism for operators with complex coefficients). Let $E$ be a complex, with bounded degrees, of differential operators with complex coefficients on $X$; for $E$ to be elliptic (1.2.5), it is necessary and sufficient that, for every $s\in S$, $i_s^*(E)$ be an elliptic complex (verification left to the reader).

\sgaStableTarget{sga6.e02.sec.2.occ2}\subsection*{2. Rigidity lemmas}

\sgaStableTarget{sga6.e02.para.2-0}2.0. It is \textquotedblleft well known\textquotedblright{} that compact differentiable varieties are \textquotedblleft rigid\textquotedblright, i.e. admit no continuous deformation other than a trivial one. However, since this result will be essential to us in \S3, and references seem rather scarce, it may be useful to give precise statements and a few proofs.

We fix a topological space $S$. The mixed varieties considered will be assumed to be of class $C^r$, with $r\ge1$ (possibly $r=\infty$).

\sgaStableTarget{sga6.e02.para.2-1}\textbf{2.1. Notation.} Let $X,Y$ be mixed $S$-varieties. We denote by $\Homop(X,Y)$ the set of $S$-morphisms from $X$ to $Y$, and by $\underline{\Homop}(X,Y)$ the sheaf on $X$ of germs of $S$-morphisms from $X$ to $Y$ (defined by
\[
  \underline{\Homop}(X,Y)(U)=\Homop(U,Y)
\]
for $U$ open in $X$).

Let $s\in S$. We denote by $i_s$ the inclusion of $X_s$ into $X$. There is a canonical induction homomorphism
\sgaStableTarget{sga6.e02.eq.2-1-1}\begin{equation}\tag{2.1.1}
  i_s^{-1}\underline{\Homop}(X,Y)
  \longrightarrow
  \underline{\Homop}(X_s,Y_s),
\end{equation}
%%%%%%%%%%  scan idx 218  |  folio 205  %%%%%%%%%%
\SGAPrintPage{205}%
which is an epimorphism, as we see immediately. When $Y=S\times\underline{\underline{R}}$, the sheaf $\underline{\Homop}(X,Y)$ is canonically identified with $\underline{O}_X$, and the epimorphism (2.1.1) with the epimorphism
\[
  i_s^{-1}(\underline{O}_X)\longrightarrow i_s^*(\underline{O}_X)=\underline{O}_{X_s}
\]
considered in (1.1).

\sgaStableTarget{sga6.e02.prop.2-2}\textbf{Proposition}\SGAAuthorFootnote{1}{Communicated to the editor by A. Douady.} \textbf{2.2.} \uline{Let} $X,Y$ \uline{be mixed} $S$\uline{-varieties, and} $s\in S$\uline{. Suppose that the fiber} $X_s$ \uline{is paracompact and that} $s$ \uline{has a fundamental system of paracompact neighborhoods. Let} $A$ \uline{be a closed subset of} $X_s$\uline{, and} $f\in H^0(A,\underline{\Homop}(X_s,Y_s))$\uline{.
The subsheaf} $F$ \uline{of} $i_s^{-1}\underline{\Homop}(X,Y)$ \uline{formed by the germs \textquotedblleft that extend} $f$\uline{\textquotedblright{} (i.e. whose image under} (2.1.1) \uline{coincides with} $f$ \uline{over} $A$\uline{) is soft.}

\uline{Proof.} Let $X'$ be an open neighborhood of $x$ in $X$, isomorphic (as a mixed variety) to $S'\times U$, where $S'$ is a paracompact neighborhood of $s$ and $U$ a relatively compact open subset of $\underline{\underline{R}}^m$, and such that the image under $f$ of $X'_s\cap A$ is contained in an open subset of $Y$ isomorphic to $S'\times V$, where $V$ is a convex open subset of $\underline{\underline{R}}^n$ containing the origin. By (T.F. II 3.4.1), it suffices to show that every section of $F$ over a closed subset $A_0$ of $X_s$ contained in $X'_s$ extends to $X'_s$. Now $X'_s$ ($\simeq\{s\}\times U$) is relatively compact, so a closed subset of $X_s$ contained in $X'_s$ is compact. Identify $X'$ (resp. $Y'$) with $S'\times U$ (resp. $S'\times V$), put $A'=A\cap X'_s$, and let $A_0$ be a compact subset of $X'_s=\{s\}\times U$. Since $A'$ is closed in $\{s\}\times U$, which is paracompact, we have ([T.F.]~II 3.3.1)\SGAAuthorFootnote{1}{A closed subset of a paracompact space has a fundamental system of paracompact neighborhoods.}
\[
  H^0\bigl(A',\underline{\Homop}(X_s,Y_s)\bigr)
  =\varinjlim_{U'\supset A'}\Homop(U',Y'_s),
\]
over the open subsets $U'$ of $X'_s$ containing $A'$. Similarly, since $A_0$ is closed in $S'\times \overline U$, which is paracompact (a product of a paracompact space and a compact space), we have
\[
  H^0\bigl(A_0,i_s^{-1}\underline{\Homop}(X,Y)\bigr)
  =\mathop{\underrightarrow{\mathrm{Lim}}}\limits_{U_0\supset A_0}\Homop(U_0,Y'),
\]
over the open subsets $U_0$ of $X'$ containing $A_0$. The proposition will therefore be a
%%%%%%%%%%  scan idx 219  |  folio 206  %%%%%%%%%%
\SGAPrintPage{206}%
consequence of the following lemma:

\sgaStableTarget{sga6.e02.lem.2-2-1}\textbf{Lemma 2.2.1.} \uline{Let} $U$ \uline{be an open subset of} $\underline{\underline{R}}^m$, $V$ \uline{a convex open subset of} $\underline{\underline{R}}^n$ \uline{containing} $0$\uline{. Let} $A_0,A_1$ \uline{be closed subsets of} $\{s\}\times U$\uline{, let} $U_0$ (resp. $U_1$\uline{) be an open subset of} $S\times U$ (resp. $\{s\}\times U$\uline{) containing} $A_0$ (resp. $A_1$\uline{). Let} $f_1:U_1\to V$ \uline{be a morphism and let} $f_0:U_0\to S\times V$ \uline{be an} $S$\uline{-morphism inducing on} $U_0\cap(\{s\}\times U)$ \uline{a morphism defining the same germ as} $f_1$ \uline{in a neighborhood of} $A_1$\uline{. There then exists an open subset} $U'$ \uline{of} $S\times U$ \uline{containing} $s\times U$\uline{, and an} $S$\uline{-morphism} $\bar f:U'\longrightarrow S\times V$ \uline{such that}
\begin{itemize}
\item[(i)] \uline{the germ of} $\bar f$ \uline{along} $\{s\}\times U$ \uline{extends the germ defined by} $f_0$;
\item[(ii)] \uline{the morphism induced by} $\bar f$ \uline{on} $\{s\}\times U$ \uline{defines the same germ as} $f_1$ \uline{in a neighborhood of} $A_1$.
\end{itemize}

\uline{Proof.} Replacing $f_1$ by $af_1$, where $a$ is a $C^r$ function on $U$, equal to $1$ in a neighborhood of $A_1$ and with support in $U_1$, we may assume that $U_1=U$. Let $b$ be a $C^r$ function on $U$, equal to $1$ on $A_0$ and with support in $U_0\cap(\{s\}\times U)$. Then
\[
  U'=U_0\cup(S\times(U-\operatorname{Supp}b))
\]
is an open subset of $S\times U$ containing $\{s\}\times U$. The morphism $\bar f:U'\to S\times V$ defined by
\[
  \bar f(t,x)=\bigl(t,\ b(x)\pr_2f_0(t,x)+(1-b(x))f_1(x)\bigr)
\]
visibly satisfies conditions (i) and (ii).

% Author-claim audit ASTRA-II-AUTHOR-CLAIM-206 is recorded separately.
% The source statement and proof are retained in both English forms.
\sgaStableTarget{sga6.e02.cor.2-3}\textbf{Corollary 2.3.} \uline{Under the hypotheses of} (2.2)\uline{, the sheaves} $i_s^{-1}\underline{\Homop}(X,Y)$, $\underline{\Homop}(X_s,Y_s)$ \uline{are soft. In particular, the sheaves} $i_s^{-1}(\underline{O}_X)$, $\underline{O}_{X_s}$ \uline{are soft}\SGAAuthorFootnote{1}{For $\underline{O}_{X_s}$, this is \textquotedblleft very well known\textquotedblright. The editor does not know whether $\underline{O}_X$ is soft when the hypotheses of (2.2) hold for every $s\in S$ and $X$ is paracompact. Under those hypotheses, this would trivially imply the softness of $i_s^{-1}(\underline{O}_X)$ as well as (2.5) below.}\uline{. For every closed subset} $A$ \uline{of} $X_s$\uline{, the map}
\[
  H^0\bigl(X_s,i_s^{-1}\underline{\Homop}(X,Y)\bigr)
  \longrightarrow
  H^0\bigl(A,\underline{\Homop}(X_s,Y_s)\bigr)
\]
\uline{deduced from} (2.1.1) \uline{is surjective.}

%%%%%%%%%%  scan idx 220  |  folio 207  %%%%%%%%%%
\SGAPrintPage{207}%
\uline{Proof.} The softness of $i_s^{-1}\underline{\Homop}(X,Y)$ follows from (2.2) applied with $A=\varnothing$; taking $S=s$, we deduce that $\underline{\Homop}(X_s,Y_s)$ is soft. The last assertion follows from (2.2).

\sgaStableTarget{sga6.e02.cor.2-4}\textbf{Corollary 2.4.} \uline{Under the hypotheses of} (2.2)\uline{, let} $E$ \uline{be an} $\underline{O}_X$\uline{-Module. For every closed subset} $A$ \uline{of} $X_s$\uline{, the canonical homomorphism}
\[
  H^0\bigl(X_s,i_s^{-1}(E)\bigr)
  \longrightarrow
  H^0\bigl(A,i_s^*(E)\bigr),
\]
\uline{deduced from the epimorphism}
\[
  i_s^{-1}(E)\longrightarrow i_s^*(E)
  =i_s^{-1}(E)\otimes_{i_s^{-1}(\underline{O}_X)}\underline{O}_{X_s},
\]
\uline{is an epimorphism.}

\uline{Proof.} Since $i_s^{-1}(\underline{O}_X)$ is soft (2.3), the functor $H^0(A,.)$ is exact on the category of $i_s^{-1}(\underline{O}_X)$-Modules (T.F. II 3.7.1, 3.4.2, 4.4.3), so the homomorphism
\[
  H^0\bigl(A,i_s^{-1}(E)\bigr)
  \longrightarrow
  H^0\bigl(A,i_s^*(E)\bigr)
\]
is an epimorphism. Since $i_s^{-1}(E)$ is soft, the homomorphism
\[
  H^0\bigl(X_s,i_s^{-1}(E)\bigr)
  \longrightarrow
  H^0\bigl(A,i_s^{-1}(E)\bigr)
\]
is an epimorphism, and we are done.

\sgaStableTarget{sga6.e02.cor.2-5}\textbf{Corollary 2.5.} \uline{Let} $p:X\to S$ \uline{be a mixed variety, and} $s\in S$\uline{. Suppose that} $s$ \uline{has a fundamental system of paracompact neighborhoods} $S'$ \uline{such that the} $p^{-1}(S')$ \uline{form a fundamental system of paracompact neighborhoods of} $X_s$ \uline{(a condition satisfied, for example, when} $S$ \uline{is locally compact and} $p$ \uline{is a proper map of topological spaces). Then, for every} $\underline{O}_X$\uline{-Module} $E$\uline{, we have}
\[
  R^i p_*(E)_s=0 \quad \text{for } i>0.
\]

%%%%%%%%%%  scan idx 221  |  folio 208  %%%%%%%%%%
\SGAPrintPage{208}%
\uline{Proof.} By (T.F. II 4.17.1), we have
\[
  R^i p_*(E)_s=H^i\bigl(X_s,i_s^{-1}(E)\bigr).
\]
But, since $i_s^{-1}(\underline{O}_X)$ is soft (2.3), we have
\[
  H^i\bigl(X_s,i_s^{-1}(E)\bigr)=0\quad \text{for } i>0,
\]
whence the assertion.

\sgaStableTarget{sga6.e02.cor.2-6}\textbf{Corollary 2.6.} \uline{Let} $X,Y$ \uline{be mixed} $S$\uline{-varieties, and} $s\in S$\uline{. Suppose that} $X$ \uline{satisfies the hypotheses of} (2.5)\uline{. For every morphism} $f:X_s\to Y_s$\uline{, there exists an open neighborhood} $S'$ \uline{of} $s$ \uline{and a morphism} $\bar f:X|S'\longrightarrow Y|S'$ \uline{such that} $\bar f_s=f$.

\uline{Proof.} This follows immediately from (2.3) (with $A=X_s$) and the fact that, by the hypotheses on $X$, we have
\[
  H^0\bigl(X_s,i_s^{-1}\underline{\Homop}(X,Y)\bigr)
  =\operatorname*{Lim}_{S'\ni s}\Homop(X|S',Y)\qquad (S'\ \text{open})
\]
(T.F. II 3.3.1).

\sgaStableTarget{sga6.e02.prop.2-7}\textbf{Proposition 2.7.} \uline{Let} $X,Y$ \uline{be mixed} $S$\uline{-varieties,} $f:X\to Y$ \uline{a morphism,} $s\in S$, $x\in X_s$\uline{. If the tangent linear map of} $f_s:X_s\to Y_s$ \uline{is invertible at} $x$\uline{, then} $f$ \uline{is a local isomorphism at} $x$\uline{, i.e. there exists an open neighborhood} $U$ \uline{of} $x$ \uline{in} $X$ \uline{and an open subset} $V$ \uline{of} $Y$ \uline{such that} $f$ \uline{induces an isomorphism from} $U$ \uline{onto} $V$.

\uline{Proof.} This is an easy consequence of the implicit function theorem. The details are left to the reader.

\sgaStableTarget{sga6.e02.cor.2-8}\textbf{Corollary 2.8.} \uline{Let} $X,Y$ \uline{be mixed} $S$\uline{-varieties, and} $s\in S$\uline{. Suppose that} $s$ \uline{has a fundamental system of neighborhoods} $S'$ \uline{such that the} $X|S'$ \uline{(resp.} $Y|S'$\uline{) form a fundamental system of neighborhoods of} $X_s$ \uline{(resp.} $Y_s$\uline{) (a condition satisfied, for example, when} $S$ \uline{is locally compact and} $X$ \uline{(resp.} $Y$\uline{) is proper over} $S$\uline{). Let} $f:X\to Y$ \uline{be an} $S$\uline{-morphism such that} $f_s:X_s\to Y_s$ \uline{is an isomorphism. There then exists a neighborhood} $S'$
%%%%%%%%%%  scan idx 222  |  folio 209  %%%%%%%%%%
\SGAPrintPage{209}%
\uline{of} $s$ \uline{such that} $f|X|S':X|S'\longrightarrow Y|S'$
\uline{is an isomorphism.}

\uline{Proof.} By (2.7), there exists an open neighborhood $U$ of $X_s$ in $X$ such that $f|U$ is étale (i.e. a local isomorphism). We may assume that $U=X|S'$, where $S'$ is a neighborhood of $s$ in $S$; in other words, after restricting $S$, we may assume that $f$ is étale. We then easily check that the fiber product $Z=X\times_Y X$ is naturally equipped with the structure of a mixed $S$-variety (étale over $X$ for each projection), and that the diagonal morphism $\Delta:X\to Z$ is an open immersion\SGAAuthorFootnote{1}{Note the analogy with Algebraic Geometry (SGA 1 I 5.1).}. Moreover, it is immediate that $Z$ satisfies the same hypothesis as $X$ and $Y$, namely that the open subsets of the form $Z|S'$, for $S'$ an open neighborhood of $s$ in $S$, form a fundamental system of neighborhoods of $Z_s$. Since $\Delta_s:X_s\to Z_s$ is an isomorphism, it follows, after restricting $S$, that $\Delta$ is an isomorphism. Thus $f$ is an open immersion, hence an isomorphism after shrinking $S$ again, and we are done.

\sgaStableTarget{sga6.e02.cor.2-9}\textbf{Corollary 2.9.} \uline{Let} $X$ \uline{be a mixed} $S$\uline{-variety, and} $s\in S$\uline{. Suppose that the hypothesis of} (2.5) \uline{holds. There then exists a neighborhood} $S'$ \uline{of} $s$ \uline{in} $S$ \uline{such that} $X|S'$ \uline{is} $S'$\uline{-isomorphic to} $S'\times X_s$.

\uline{Proof.} Apply (2.6) and (2.8).

\sgaStableTarget{sga6.e02.cor.2-10}\textbf{Corollary 2.10.} \uline{Let} $s\in S$ \uline{and let} $V$ \uline{be a paracompact variety. Put} $X=S\times V$\uline{, and denote by} $i_s$ \uline{the injection of} $V$ \uline{into} $X$ \uline{defined by} $s$\uline{. Suppose that the hypothesis of} (2.5) \uline{holds. Let} $E$ \uline{be a locally free} $\underline{O}_X$ \uline{(resp.} $\underline{O}_X\otimes\underline{\underline{C}}$\uline{)-Module of finite type. There then exists a neighborhood} $S'$ \uline{of} $s$ \uline{in} $S$ \uline{such that} $E|X|S'$ \uline{is isomorphic to} $\pr_2^*i_s^*(E)$ ($\pr_2:S'\times V\to V$ \uline{being the second projection).}

\uline{Proof.} Put $F=\pr_2^*i_s^*(E)$. By (2.4), applied to the Module $\underline{\Homop}(E,F)$, the induction homomorphism
\[
  \Homop\bigl(i_s^{-1}(E),i_s^{-1}(F)\bigr)
  \longrightarrow
  \Homop\bigl(i_s^*(E),i_s^*(F)\bigr)
\]
%%%%%%%%%%  scan idx 223  |  folio 210  %%%%%%%%%%
\SGAPrintPage{210}%
is an epimorphism. On the other hand, we have (T.F. II 3.3.1):
\[
\begin{aligned}
  \Homop\bigl(i_s^{-1}(E),i_s^{-1}(F)\bigr)
  &\simeq H^0\bigl(X_s,i_s^{-1}\underline{\Homop}(E,F)\bigr)\\
  &=\varinjlim_{S'\ni s} H^0\bigl(X|S',\underline{\Homop}(E,F)\bigr)\qquad (S'\ \text{open})\\
  &\simeq \varinjlim_{S'\ni s}\Homop(E|X|S',F|X|S').
\end{aligned}
\]
Thus the canonical isomorphism $i_s^*(E)\simeq i_s^*(F)$ is induced by a homomorphism $f:E\to F$ (after shrinking $S$). Since $f_s$ is an isomorphism, $f$ is an isomorphism in a neighborhood of $X_s$ by Nakayama, hence over $X|S'$, where $S'$ is a suitable neighborhood of $s$, q.e.d.

\sgaStableTarget{sga6.e02.sec.3.occ2}\subsection*{3. The finiteness theorem}

\sgaStableTarget{sga6.e02.sec.3.0}\subsubsection*{3.0. Conventions and notation}
\sgaStableTarget{sga6.e02.para.3-0-1}3.0.1. Let $X$ and $Y$ be topological spaces. We denote by $C(X,Y)$ the set of continuous maps from $X$ to $Y$, and by $\underline C(X,Y)$ the sheaf on $X$ defined by $U\mapsto C(U,Y)$.

\sgaStableTarget{sga6.e02.para.3-0-2}3.0.2. Let $E,F$ be Banach spaces (real or complex). We denote by $\Lop(E,F)$ the Banach space of continuous linear maps from $E$ to $F$.

\sgaStableTarget{sga6.e02.para.3-0-3}3.0.3. The mixed varieties considered in this section will be assumed to be of class $C^\infty$. Let $X$ be a mixed $S$-variety; we shall often identify a locally free $\underline{O}_X$ (resp. $\underline{O}_X\otimes\underline{\underline{C}}$)-Module of finite type with the corresponding real (resp. complex) vector bundle.

\sgaStableTarget{sga6.e02.para.3-0-4}3.0.4. Let $X$ be a compact (ordinary) variety, and let $E$ be a vector bundle (real or complex) of finite rank on $X$. We denote by $C^\infty(X,E)$ the Fréchet space of sections (of class $C^\infty$) of $E$ on $X$. For $r\in\underline{\underline{Z}}$, we denote by $C^r_2(X,E)$ the Hilbert space of sections of $E$ that are $r$ times differentiable in $L^2$ ([16]~IX). For $s\ge r$, $C^s_2(X,E)$ maps into $C^r_2(X,E)$ by a continuous injective linear map with dense image, and we have
%%%%%%%%%%  scan idx 224  |  folio 211  %%%%%%%%%%
\SGAPrintPage{211}%
(loc.cit.)
\[
  C^\infty(X,E)=\varprojlim_r C^r_2(X,E).
\]

\sgaStableTarget{sga6.e02.para.3-0-5}3.0.5.
Let $S$ be a topological space. We shall need the notion of a \textquotedblleft Hilbert bundle\textquotedblright{} over $S$, for which we refer the reader to [12] or [14]. There is an additive functor from the category of Hilbert bundles (real resp. complex) over $S$ to that of $\underline{O}_S$ resp. $\underline{O}_S\otimes\underline{\underline{C}}$-Modules, associating to each Hilbert bundle $E$ the sheaf $\underline E$ of continuous sections of $E$.

\sgaStableTarget{sga6.e02.sec.3.1}\subsubsection*{3.1. Auxiliary technical results}
\sgaStableTarget{sga6.e02.lem.3-1-1}\textbf{Lemma 3.1.1.} \uline{Let} $X,Y$ \uline{be topological spaces, with} $X$ \uline{locally compact. Denote by} $\underline C'(X,Y)$ \uline{the presheaf defined on the basis of relatively compact open subsets of} $X$ \uline{by} $\underline C'(X,Y)(U)=C(\bar U,Y)$ (3.0.1)\uline{. The canonical restriction morphism}
\[
  \underline C'(X,Y)\longrightarrow \underline C(X,Y)
\]
\uline{induces an isomorphism on the associated sheaves.}

\uline{Proof.} This is immediate.

\sgaStableTarget{sga6.e02.lem.3-1-2}\textbf{Lemma 3.1.2.} \uline{Let} $E,F$ \uline{be locally convex topological vector spaces (over} $\underline{\underline{C}}$\uline{, to fix ideas), and} $u:E\to F$ \uline{a continuous linear map (resp. continuous with dense image). Let} $X$ \uline{be a compact space. Equip} $C(X,E)$, $C(X,F)$ \uline{with the topology of uniform convergence. Then the linear map}
\[
  u_X:C(X,E)\longrightarrow C(X,F)
\]
\uline{induced by} $u$ \uline{is continuous (resp. continuous with dense image).}

%%%%%%%%%%  scan idx 225  |  folio 212  %%%%%%%%%%
\SGAPrintPage{212}%
\uline{Proof.} It is clear that $u_X$ is continuous. On the other hand (Bourbaki, Intégration, chap. III \S~2, lemmas 1 and 2), the subspace $C(X,\underline{\underline{C}})\otimes_{\underline{\underline{C}}}E$ of $C(X,E)$ (formed by the continuous maps from $X$ to $E$ whose values lie in a finite-dimensional vector subspace) is dense; similarly, $C(X,\underline{\underline{C}})\otimes_{\underline{\underline{C}}}F$ is dense in $C(X,F)$. Suppose $u$ is continuous with dense image. Then
\[
  u_X(C(X,\underline{\underline{C}})\otimes_{\underline{\underline{C}}}E)=C(X,\underline{\underline{C}})\otimes_{\underline{\underline{C}}}u(E)
\]
is dense in $C(X,\underline{\underline{C}})\otimes_{\underline{\underline{C}}}F$, whence, taking closures, we deduce that $u_X(C(X,E))$ is dense in $C(X,F)$, q.e.d.

\sgaStableTarget{sga6.e02.para.3-1-3}3.1.3. Let $V$ be a compact variety, $E$ and $F$ complex vector bundles\SGAAuthorFootnote{1}{To fix ideas (analogous results hold in the real case).} of finite rank on $V$, and $n\in\underline{\underline{N}}$. The space $\Diff^n_{\underline{\underline{C}}}(E,F)$, identified by (1.2.3.1) with
\[
  \Homop(P^n(E),F)\simeq C^\infty(V,\underline{\Homop}(P^n(E),F)),
\]
is naturally equipped with the structure of a Fréchet space. Let $u\in\Diff^n_{\underline{\underline{C}}}(E,F)$; we easily check that the linear map from $C^\infty(V,E)$ to $C^\infty(V,F)$ defined by $u$ is continuous. Denoting by $\Lop(C^\infty(V,E),C^\infty(V,F))$ the space of continuous linear maps from $C^\infty(V,E)$ to $C^\infty(V,F)$, we therefore have a linear map
\[
  \phi:\Diff^n_{\underline{\underline{C}}}(E,F)
  \longrightarrow \Lop(C^\infty(V,E),C^\infty(V,F)).
\]
On the other hand, for $r\in\underline{\underline{Z}}$, $u$ acts by extension of scalars ([16]~IX Th.~12) on the $C^r_2$ sections of $E$ and defines a linear map
\[
  \phi_r:\Diff^n_{\underline{\underline{C}}}(E,F)
  \longrightarrow \Lop(C^r_2(V,E),C^{r-n}_2(V,F)).
\]
We readily check that:
\begin{enumerate}[label=(\roman*)]
\item the $\phi_r$ are \uline{continuous} (the target is a Banach space
%%%%%%%%%%  scan idx 226  |  folio 213  %%%%%%%%%%
\SGAPrintPage{213}%
(3.0.2));
\item $\phi=\varprojlim\phi_r$ (i.e.
for $u\in\Diff^n_{\underline{\underline{C}}}(E,F)$, $\phi(u)$ is the limit of the projective system of the $\phi_r(u)$).
\end{enumerate}
(The preceding assertions follow easily from the results of (loc.cit.): we reduce immediately to $n=0$; (ii) then follows directly from the definitions, while (i) follows (after localizing the problem by a partition of unity) from the formula giving the successive derivatives of a product).

\sgaStableTarget{sga6.e02.lem.3-1-4}\textbf{Lemma 3.1.4.} \uline{Let} $S$ \uline{be a topological space,} $V$ \uline{a compact variety,} $E$ \uline{and} $F$ \uline{complex vector bundles of finite rank on} $V$\uline{, and} $n\in\underline{\underline{N}}$\uline{. Denote by} $X$ \uline{the mixed} $S$\uline{-variety} $S\times V$\uline{. For} $u\in\Diff^n_{\underline{\underline{C}}}(\pr_2^*(E),\pr_2^*(F))$\uline{, the map} $s\mapsto i_s^*(u)$ \uline{from} $S$ \uline{to} $\Diff^n_{\underline{\underline{C}}}(E,F)$ \uline{is continuous. The linear map}
\[
  \Diff^n_{\underline{\underline{C}}}(\pr_2^*(E),\pr_2^*(F))
  \longrightarrow C(S,\Diff^n_{\underline{\underline{C}}}(E,F))
\]
\uline{thus defined is an isomorphism.}

\uline{Proof.} Using the identification (1.2.3.1), we may restrict to $n=0$. Putting $\underline{\Homop}(E,F)=G$, we are then reduced to showing that, for $u\in\Gamma(X,\pr_2^*(G))$, the map $s\mapsto u_s=i_s^*(u)$ from $S$ to $C^\infty(V,G)$ is continuous, and that the linear map
\[
  \Gamma(X,\pr_2^*(G))\longrightarrow C(S,C^\infty(V,G))
\]
so obtained is an isomorphism. But this follows immediately from the definition of mixed varieties (1.1).

\sgaStableTarget{sga6.e02.para.3-1-5}3.1.5. Under the hypotheses of (3.1.4), let
\[
  u\in\Diff^n(\pr_2^*(E),\pr_2^*(F)).\text{\SGAAuthorFootMark{1}}
\]
\SGAAuthorFootText{1}{We omit the subscript $\underline{\underline{C}}$ for brevity.}
%%%%%%%%%%  scan idx 227  |  folio 214  %%%%%%%%%%
\SGAPrintPage{214}%
Identifying $u$, by (3.1.4), with a continuous map from $S$ to $\Diff^n(E,F)$, we see that, for $U$ open in $S$, the morphism
\[
  \pr_{1*}(u)(U):\Gamma(X|U,E)\longrightarrow \Gamma(X|U,F)
\]
is canonically identified with the linear map
\[
  C(U,C^\infty(V,E))\longrightarrow C(U,C^\infty(V,F))
\]
deduced from
\begin{equation}\tag{3.1.3}
  \phi u:U\longrightarrow \Lop(C^\infty(V,E),C^\infty(V,F)).
\end{equation}
Suppose $S$ is locally compact. Then (3.1.1) the morphism
\[
  \pr_{1*}(u):\pr_{1*}\pr_2^*(E)
  \longrightarrow \pr_{1*}\pr_2^*(F)
\]
is the morphism of sheaves associated to the morphism of presheaves
\[
  u_\infty:\underline C'(S,C^\infty(V,E))
  \longrightarrow \underline C'(S,C^\infty(V,F))
\]
defined by $\phi u$. On the other hand, for $r\in\underline{\underline{Z}}$,
\[
  \phi_r u:S\longrightarrow \Lop(C^r_2(V,E),C^{r-n}_2(V,F))
\]
defines a morphism of presheaves
\[
  u_r:\underline C'(S,C^r_2(V,E))
  \longrightarrow \underline C'(S,C^{r-n}_2(V,F)).
\]
By (3.1.3 (ii)), the $u_r$ form a projective system whose limit is $u_\infty$. It will be convenient to interpret $\phi_r u$ as a morphism of trivial Hilbert bundles (3.0.5)
\[
  S\times C^r_2(V,E)\longrightarrow S\times C^{r-n}_2(V,F),
\]
%%%%%%%%%%  scan idx 228  |  folio 215  %%%%%%%%%%
\SGAPrintPage{215}%
inducing on the sheaves of sections the morphism of sheaves associated to $u_r$.

\sgaStableTarget{sga6.e02.sec.3.2}\subsubsection*{3.2. Definition of the index}
\sgaStableTarget{sga6.e02.thm.3-2-1}\textbf{Theorem 3.2.1.} \uline{Let} $S$ \uline{be a locally compact space,} $p:X\to S$ \uline{a proper mixed variety,} $E^\bullet$ \uline{an elliptic complex of order} $\underline{n}$ \uline{of complex vector bundles on} $X$ (1.2.5)\uline{. Then the complex} $\underline{\underline{R}}p_*(E^\bullet)$ \uline{(which, by} (2.5)\uline{, is isomorphic, to} $p_*(E^\bullet)$ \uline{in the derived category of that of} $\underline{O}_S\otimes\underline{\underline{C}}$\uline{-Modules) is perfect.}

\uline{Proof.} The question being local on $S$, we may assume (2.9) that $X=S\times V$, where $V$ is a compact variety, and $p=\pr_1$. Moreover, we may assume (2.10) that, for every $i$, $E^i=\pr_2^*(F^i)$, where $F^i$ is a complex vector bundle of finite rank on $V$. By (3.1.5), the complex $p_*(E)$ is associated to the complex of presheaves $\bigl(\underline C'(S,C^\infty(V,F^i)),d_\infty^i\bigr)$ (where the $d^i:E^i\longrightarrow E^{i+1}$ are the given differential operators), which we shall denote simply by $\underline C'(S,C^\infty(V,F^\bullet))$. Put, for each $i\in\underline{\underline{Z}}$,
\[
  m_i=\sum_{j<i} n_j
  \qquad (\text{we agree that } n_i=0 \text{ when } E^i=0).
\]
For each $r\in\underline{\underline{Z}}$, we have (3.1.5) a complex of presheaves $\bigl(\underline C'(S,C_2^{r-m_i}(V,F^i)),d_{r-m_i}^i\bigr)$, which we shall denote by $\underline C'(S,C_2^r(V,F^\bullet))$.
The $\underline C'(S,C_2^r(V,F^\bullet))$ form a projective system, and we have (3.1.5):
\sgaStableTarget{sga6.e02.eq.3-2-1-1}\begin{equation}\tag{3.2.1.1}
  \underline C'(S,C^\infty(V,F^\bullet))=\varprojlim_r \underline C'(S,C_2^r(V,F^\bullet)).
\end{equation}
Finally, we shall denote by $F_r^\bullet$ the complex of trivial Hilbert bundles over $S$ defined by $\phi_{r-m_i}d^i$, and by $\underline F_r^\bullet$ the corresponding complex of sheaves of sections, which is precisely (3.1.5) the complex of sheaves associated to $\underline C'(S,C_2^r(V,F^\bullet))$. With this notation, to prove the theorem it suffices to prove the following assertions:

(a) \uline{For} $r\in\underline{\underline{Z}}$\uline{, the canonical morphism}
%%%%%%%%%%  scan idx 229  |  folio 216  %%%%%%%%%%
\SGAPrintPage{216}%
\[
  \underline C'(S,C^\infty(V,F^\bullet))
  \longrightarrow \underline C'(S,C_2^r(V,F^\bullet))
\]
\uline{is a
quasi-isomorphism of complexes of presheaves, i.e.\ induces isomorphisms on the cohomology presheaves.}

(b) \uline{For} $r\in\underline{\underline{Z}}$\uline{, the complex} $\underline F_r^\bullet$ \uline{is perfect.}

\uline{Proof of} (b). The ellipticity of $E^\bullet$ implies ([16]~XI Th.6\SGAAuthorFootnote{1}{which readily generalizes to elliptic complexes.}) that, for every $s\in S$, the differential of $(F_r^\bullet)_s$ has closed image and the $H^i((F_r^\bullet)_s)$ are finite-dimensional vector spaces. In other words, to use the terminology of [11], $F_r^\bullet$ is a direct quasi-acyclic complex of Hilbert bundles. Consequently, by ([12]~2.3), $F_r^\bullet$ is locally homotopically equivalent to a bounded complex of vector bundles of finite rank. Since the \textquotedblleft sheaf of sections\textquotedblright{} functor is additive (3.0.5), $\underline F_r^\bullet$ is locally homotopically equivalent to a bounded complex of locally free sheaves of finite type, and hence is perfect.

For the proof of (a), we shall need the following intermediate assertion:

(a') \uline{For} $p\geq q$\uline{, the canonical morphism}
\[
  \underline F_p^\bullet\longrightarrow \underline F_q^\bullet
\]
\uline{is a quasi-isomorphism.}

\uline{Proof of} (a'). Denote by $M_{pq}^\bullet$ the complex of Hilbert bundles that is the cone of the morphism $F_p^\bullet\longrightarrow F_q^\bullet$. The ellipticity of $E$ implies that, for every $s\in S$, the morphism $(F_p^\bullet)_s\longrightarrow (F_q^\bullet)_s$ is a quasi-isomorphism ([16]~XI Cor.~2 of Th.~6). In other words, $M_{pq}^\bullet$ is a direct acyclic complex. It follows ([11]~2.3) that $M_{pq}^\bullet$ is locally homotopically trivial, and hence that $\underline M_{pq}^\bullet$, which is precisely the cone of $\underline F_p^\bullet\longrightarrow\underline F_q^\bullet$, is acyclic, q.e.d.

\uline{Proof of} (a). We must show that, for every relatively compact open subset $U$ of $S$, the morphism
%%%%%%%%%%  scan idx 230  |  folio 217  %%%%%%%%%%
\SGAPrintPage{217}%
\[
  C(\overline U,C^\infty(V,F^\bullet))
  \longrightarrow C(\overline U,C_2^r(V,F^\bullet))
\]
is a
quasi-isomorphism. Replacing $S$ by $\overline U$ (and $E^\bullet$ by the complex induced over $\overline U$), we are reduced to showing that
\[
  C(S,C^\infty(V,F^\bullet))
  \longrightarrow C(S,C_2^r(V,F^\bullet))
\]
is a quasi-isomorphism, with $S$ assumed \uline{compact}. By (3.2.1.1), we have
\[
  C(S,C^\infty(V,F^\bullet))=
  \varprojlim_r C(S,C_2^r(V,F^\bullet))
\]
For each $i$, $C(S,C_2^r(V,F^i))$ is a Banach space (for the topology of uniform convergence), and the transition morphisms
\[
  C(S,C_2^p(V,F^i))\longrightarrow C(S,C_2^q(V,F^i))
\]
are, by (3.0.4) and (3.1.2), continuous linear maps with dense image. On the other hand, applying to the quasi-isomorphisms (a') $\underline F_p^\bullet\longrightarrow\underline F_q^\bullet$ the functor $\Gamma_S$, which is exact since $\underline{O}_S$ is soft, we see that the transition morphisms
\[
  C(S,C_2^p(V,F^\bullet))\longrightarrow C(S,C_2^q(V,F^\bullet))
\]
are quasi-isomorphisms. Assertion (a) then follows from (EGA $0_{\mathrm{III}}$ 13.2.3, \textquotedblleft topological\textquotedblright{} variant). This completes the proof of (3.2.1).

\sgaStableTarget{sga6.e02.para.3-2-2}3.2.2. For a topological space $S$, we shall denote by $\Parf(S)$ the triangulated category of perfect complexes of $(\underline{O}_S\otimes\underline{\underline{C}})$-Modules, and by $K^\bullet(S)$ the corresponding \textquotedblleft Grothendieck group\textquotedblright{}\SGAAuthorFootnote{1}{see (I 6.3). For a general review of the functors $K$, see Exp.~IV.}.

Let $S$, $X$, $E^\bullet$ be as in (3.2.1); the \uline{analytic index} of $E^\bullet$ is the element
\[
  \operatorname{ind.an}(E^\bullet)=C\ell(\underline{\underline{R}}p_*(E^\bullet))\in K^\bullet(S).
\]
%%%%%%%%%%  scan idx 231  |  folio 218  %%%%%%%%%%
\SGAPrintPage{218}%
For \uline{compact} $S$, $\operatorname{ind.an}(E^\bullet)$ canonically defines an element in the usual group $K(S)$ of topologists. Indeed, in general, for every topological space $S$, denote by $\mathrm{Loclib}(S)$ the additive category of locally free $(\underline{O}_S\otimes\underline{\underline{C}})$-Modules of finite type, and by $K(S)$ the \textquotedblleft Grothendieck group\textquotedblright{} of the triangulated category $K^b(\mathrm{Loclib}(S))/K^{b,\varnothing}(\mathrm{Loclib}(S))$ (the latter is, moreover, identified with $K^b(\mathrm{Loclib}(S))$ when $S$ is paracompact (II 2.3.2 (a)), so that in this case (I 6.4) $K(S)$ is canonically isomorphic to the \textquotedblleft Grothendieck group\textquotedblright{} of the additive category $\mathrm{Loclib}(S)$). There is a natural homomorphism
\[
  K(S)\longrightarrow K^\bullet(S).
\]
But, by (II 2.3.2), this is an isomorphism for compact $S$, proving our assertion. More generally, $\operatorname{ind.an}(E^\bullet)$ defines an element in the completed group
\[
  \widehat K(S)=\varprojlim K(A),
\]
the projective limit being taken over the compact subsets $A$ of $S$.

\sgaStableTarget{sga6.e02.para.3-2-3}3.2.3. The writer does not know whether the finiteness theorem (3.2.1) remains valid when, instead of assuming $S$ locally compact, one assumes only that every point of $S$ has a fundamental system of paracompact neighborhoods. Moreover, the local compactness hypothesis on $S$ does not seem very natural. The space of \textquotedblleft universal local parameters\textquotedblright{} for elliptic complexes of fixed order is indeed an (infinite-dimensional) subspace of a Fréchet space. Let us make this precise. Let $V$ be a compact variety, $F=(F^i)_{i\in\underline{\underline{Z}}}$ a sequence of complex vector bundles of finite rank on $V$, all but finitely many zero, and $\underline n=(n_i)$ a sequence of integers $\geq 0$. Denote by $\operatorname{Comp.diff}^{\underline n}(F)$ the space of complexes of differential operators \textquotedblleft with base $F$\textquotedblright{} and of order $\underline n$, i.e.\ the closed subspace of $\prod\Diff^{n_i}(F^i,F^{i+1})$ defined by the equations $d^{i+1}d^i=0$. (the space $\prod\Diff^{n_i}(F^i,F^{i+1})$ is a Fréchet space). The subset
%%%%%%%%%%  scan idx 232  |  folio 219  %%%%%%%%%%
\SGAPrintPage{219}%
$\operatorname{Ell}^{\underline n}(F)$ of
$\operatorname{Comp.diff}^{\underline n}(F)$ formed by the elliptic complexes of order $\underline n$ is open; it is this space (which in general is not locally compact) that plays the role of a space of universal local parameters.

\sgaStableTarget{sga6.e02.para.3-2-4}3.2.4. Nevertheless, even should the generalization of (3.2.1) mentioned above prove impossible, we can always associate to $E^\bullet$, under the sole hypothesis that the fibration $p$ is proper and locally trivial, an element of $K^\bullet(S)$ deserving the name analytic index of $E^\bullet$. Indeed, suppose (for simplicity\SGAAuthorFootnote{1}{for the general case, use the notion of a \textquotedblleft weak Hilbert bundle\textquotedblright{} or continuous field of Hilbert spaces from [12].}) that we are in the \textquotedblleft local\textquotedblright{} situation where $X=S\times V$ ($V$, a compact variety), with $E^\bullet$ given by a continuous map from $S$ to $\operatorname{Ell}^{\underline n}(F)$ (3.2.3). Assertions (a') and (b) in the proof of (3.2.1) are in fact valid without any hypothesis on $S$: it is therefore reasonable to put, with the notation of (loc.cit.):
\[
  \operatorname{ind.an}(E^\bullet)=C\ell(\underline F_r^\bullet)\in K^\bullet(S),
\]
this element being independent of $r$.

\sgaStableTarget{sga6.e02.para.3-2-5}3.2.5. Suppose $S$ is paracompact, and denote by $H$ a complex Hilbert space of infinite Hilbert dimension, and by $I(H)$ the subspace of $L(H,H)$ formed by the operators \textquotedblleft with index\textquotedblright{}. Then the complexes $F_r^\bullet$ define [12] a canonical element in the group $[S,I(H)]$ of homotopy classes of continuous maps from $S$ to $I(H)$, but the writer does not know, for noncompact $S$ (even when $S$ is locally compact), the relations between the groups $[S,I(H)]$ and $K^\bullet(S)$ (is there a natural arrow from one to the other? are they isomorphic?), and therefore a fortiori between the two indices thus defined.

For a generalization, under the hypotheses of (3.2.1), of the Atiyah--Singer index theorem, we refer to the recent work of Shih [17], or to work in preparation by M. F. Atiyah and I. M. Singer.

%%%%%%%%%%  scan idx 233  |  folio 220  %%%%%%%%%%
\SGAPrintPage{220}%
\sgaStableTarget{sga6.e02.bibliography}\subsection*{BIBLIOGRAPHY}
\begin{enumerate}[label={[\arabic*]}]
\item
Artin, M., Grothendieck, A., and Verdier, J-L., Cohomologie étale des schémas, to appear in North Holl. Pub. Comp., (cited as SGA 4).
\item Atiyah, M. and Segal, G., secret papers.
\item Cartan, H. Funzioni e varietà complesse, Faisceaux analytiques cohérents, Course at the Centro Internazionale Matematico Estivo, 1964, Roma.
\item Cartan, H., and Schwartz, L., Théorème d'Atiyah--Singer, Seminar 1963-64, Institut Henri Poincaré, Paris.
\item Frisch, J., Fonctions analytiques sur un ensemble semi-analytique, note in CRAS Paris, t. 260, p. 2974-2976 (1965).
\item Godement, R., Théorie des Faisceaux, Hermann (1958), (cited as T.F).
\item Grothendieck, A., Cohomologie locale des faisceaux cohérents et théorèmes de Lefschetz locaux et globaux, Séminaire de Géométrie Algébrique de l'IHES 1962, to appear in North Holl. Pub. Comp., (cited as SGA 2).
\item Grothendieck, A., Fondements de la Géométrie Analytique, Séminaire Cartan 1959-60, Institut Henri Poincaré, Paris.
\item Hartshorne, R., Residues and Duality, Lecture Notes, Springer n\textsuperscript{o}~20, cited as [H].
\item Houzel, C., Géométrie Analytique Locale, Séminaire Cartan 1959-60, Institut Henri Poincaré, Paris.
\item Illusie, L., Complexes de fibrés et définitions du foncteur $K$, Séminaire Shih Weishu, IHES 1964-65.
\item Illusie, L., Complexes quasi-acycliques directs de fibrés banachiques, note in CRAS Paris, t. 260, p. 6499-6502.
\item Kiehl, R., Theorem A und Theorem B in der nichtarchimedischen Funktionentheorie, Invent. math. 2, p. 256-273 (1967).
%%%%%%%%%%  scan idx 234  |  folio 221  %%%%%%%%%%
\SGAPrintPage{221}%
\item
Lang, S., Introduction to differentiable manifolds,
\item \L{}ojasiewicz, S., Ann. della Scuola Norm. Sup. Pisa, serie III, 18, fasc. IV, 1964.
\item Palais, R., Lectures at the Princeton Seminar 1963-64 on the index theorem.
\item Shih, W., Fiber cobordism and the index of a family of elliptic differential operators,
\item Verdier, J.-L., Les catégories dérivées des catégories abéliennes, to appear in North Holl. Pub. Comp.
\end{enumerate}

% Original English continuation definitions, retained in their original order.
\providecommand{\Ob}{\operatorname{Ob}}
\providecommand{\qcoh}{\mathrm{qcoh}}
\providecommand{\coh}{\mathrm{coh}}
\providecommand{\torf}{\mathrm{torf}}
\providecommand{\toramp}{\operatorname{tor.amp}}
\providecommand{\amp}{\operatorname{amp}}
\providecommand{\Id}{\operatorname{Id}}
\providecommand{\pr}{\operatorname{pr}}
\providecommand{\coker}{\operatorname{coker}}
\providecommand{\im}{\operatorname{Im}}
\providecommand{\Spec}{\operatorname{Spec}}

% Reviewed French continuation definitions, retained verbatim.
% Exposé III additions
\renewcommand{\Ob}{\operatorname{ob}}
\providecommand{\qcoh}{\mathrm{qcoh}}
\providecommand{\coh}{\mathrm{coh}}
\providecommand{\torf}{\mathrm{torf}}
\providecommand{\toramp}{\operatorname{tor.amp}}
\providecommand{\amp}{\operatorname{amp}}
\providecommand{\Id}{\operatorname{Id}}
\providecommand{\pr}{\operatorname{pr}}
\providecommand{\coker}{\operatorname{coker}}
\providecommand{\im}{\operatorname{Im}}
\providecommand{\Spec}{\operatorname{Spec}}

%%%%%%%%%%  scan idx 235  |  folio 222  %%%%%%%%%%
\SGAPrintPage{222}
\sgaStableTarget{sga6.e03.expose}\section*{\underline{EXPOSÉ III}\\[3pt] \underline{\underline{RELATIVE FINITENESS CONDITIONS}}}
\begin{center}
\underline{by L. ILLUSIE}
\end{center}

\sgaStableTarget{sga6.e03.sec.1}\section*{1. Relative pseudo-coherence}

\sgaStableTarget{sga6.e03.thm.1-1}\textbf{Theorem 1.1.} \uline{Let} \(n\in\mathbb Z\). \uline{For each locally finite type morphism of schemes} \(f:X\longrightarrow Y\), \uline{there exists one and only one full subcategory} \(D(f)_{n\text{-}\coh}\) \uline{of} \(D(X)\)\SGAAuthorFootnote{1}{\(D(X)\) = derived category of the category of \(\underline{O}_X\)-Modules.} \uline{such that the following conditions are satisfied}:
\begin{enumerate}[label=(\roman*)]
\item \uline{If} \(E\in\operatorname{Ob} D(f)_{n\text{-}\coh}\), \uline{then} \(i^*(E)\in\operatorname{Ob} D(fi)_{n\text{-}\coh}\) \uline{for every open immersion} \(i:U\to X\).

\item \uline{If} \(E\in\ob D(X)\) \uline{and there exists an open covering} \((i_\alpha:U_\alpha\to X)_{\alpha\in A}\) \uline{such that} \(i_\alpha^*(E)\in\ob D(fi_\alpha)_{n\text{-}\coh}\) \uline{for every} \(\alpha\in A\), \uline{then} \(E\in\ob D(f)_{n\text{-}\coh}\).

\item \uline{If}
\[
\begin{tikzcd}
X \arrow[rr,"f"] \arrow[dr,"h"'] & & Y \arrow[dl,"g"]\\
& Z &
\end{tikzcd}
\]
\uline{is a commutative triangle of locally finite type morphisms}, \uline{where} \(f\) \uline{is a closed immersion and} \(g\) \uline{a smooth morphism}, \uline{then we have}
\[
  E\in\operatorname{Ob} D(h)_{n\text{-}\coh}
  \quad\Longleftrightarrow\quad
  f_*(E)\in D(Y)_{n\text{-}\coh}\qquad\mbox{\SGAAuthorFootMark{2}}\ .
\]
\SGAAuthorFootText{2}{\(D(Y)_{n\text{-}\coh}\) = full subcategory of \(D(Y)\) formed by the \(n\)-pseudo-coherent objects (I~2.15).}
\end{enumerate}

\uline{Proof}. \uline{Uniqueness} is evident. Indeed, let \(f:X\to Y\) be a locally finite type morphism. Choose an open covering \((i_\alpha:X_\alpha\to X)_{\alpha\in A}\) such that, for each \(\alpha\in A\), \(fi_\alpha\) factors as
\[
  X_\alpha\xrightarrow{i'_\alpha}Y_\alpha\xrightarrow{f'_\alpha}Y,
\]
where \(i'_\alpha\) is a closed immersion and \(f'_\alpha\) a smooth morphism. It follows at once from (i), (ii), (iii) that the objects of \(D(f)_{n\text{-}\coh}\) are the objects \(E\) of \(D(X)\) such that
\[
  i'_{\alpha *}i_\alpha^*(E)\in \ob D(Y_\alpha)_{n\text{-}\coh};
\]
thus \(D(f)_{n\text{-}\coh}\)
%%%%%%%%%%  scan idx 236  |  folio 223  %%%%%%%%%%
\SGAPrintPage{223}
is indeed determined (since it is a full subcategory of \(D(X)\)). To prove existence, we shall need a few lemmas.

\sgaStableTarget{sga6.e03.lem.1-1-1}\textbf{Lemma 1.1.1.} \uline{Let} \(i:X\to Y\) \uline{be a closed immersion of schemes such that} \(i_*(\underline{O}_X)\) \uline{is pseudo-coherent}, \uline{and let} \(E\in\ob D(X)\). \uline{For} \(E\) \uline{to be \(n\)-pseudo-coherent}, \uline{it is necessary and sufficient that} \(i_*(E)\) \uline{be so}.

\uline{Proof}. The ``necessity'' follows from the general criterion for stability of pseudo-coherence under application of a functor (I~2.13). Conversely, suppose that \(i_*(E)\) is \(n\)-pseudo-coherent, and prove that the same is true of \(E\). Since \(i_*(E)\) has locally bounded-above cohomology, the same is true of \(E\); thus, the question being local, we may suppose that \(E\) is \(m\)-pseudo-coherent for sufficiently large \(m\). Let us show, by descending induction on \(m\), that \(E\) is \(m\)-pseudo-coherent for every \(m\geq n\). We may suppose it already proved that \(E\) is \((n+1)\)-pseudo-coherent. After localizing, we therefore have a distinguished triangle
\[
\begin{tikzcd}
& E'' \arrow[dl,"|" marking] & \\
E' \arrow[rr] & & E \arrow[ul]
\end{tikzcd}
\]
where \(E'\) is perfect and \(E''^i=0\) for \(i>n\). By (I~2.6), we are reduced to showing that \(E''\) is \(n\)-pseudo-coherent. Now, applying \(i_*\) to the preceding triangle, we obtain a distinguished triangle
\[
\begin{tikzcd}
& i_*(E'') \arrow[dl,"|" marking] & \\
i_*(E') \arrow[rr] & & i_*(E) \arrow[ul]
\end{tikzcd}\ .
\]
By the first part of the proof, \(i_*(E')\) is pseudo-coherent, so by (I~2.6), \(i_*(E'')\) is \(n\)-pseudo-coherent. Since \(i_*(E'')\) is zero in degree \(>n\), this means, by (I~2.10 b)), that \(H^n(i_*(E''))\) is of finite type. Since \(H^n(i_*(E''))\simeq i_*H^n(E'')\), we are reduced to showing, by (I~2.10 b)), that if \(M\) is an \(\underline{O}_X\)-Module such that \(i_*(M)\) is of finite type, then \(M\) is of finite type, which is immediate (and independent of the hypothesis on \(i_*(\underline{O}_X)\)).

%%%%%%%%%%  scan idx 237  |  folio 224  %%%%%%%%%%
\SGAPrintPage{224}
\sgaStableTarget{sga6.e03.rem.1-1-2}\textbf{Remark 1.1.2.} The preceding lemma applies in particular when \(i\) is a \uline{regular immersion} (VII~1.4). Indeed, by means of the Koszul complex, we have locally a finite left resolution of \(i_*(\underline{O}_X)\) by free \(\underline{O}_Y\)-Modules of finite type, so \(i_*(\underline{O}_X)\) is perfect (I 4.8 and 4.14) (and therefore a fortiori pseudo-coherent).

\sgaStableTarget{sga6.e03.lem.1-1-3}\textbf{Lemma 1.1.3.} \uline{Let there be a Cartesian square of schemes}
\[
\begin{tikzcd}
X \arrow[d,"i"'] & X' \arrow[l,"f"'] \arrow[d,"i'"]\\
Y & Y'{,} \arrow[l,"g"']
\end{tikzcd}
\]
\uline{where} \(i\) \uline{is a closed immersion}, \uline{and let} \(E\) \uline{be a complex of} \(\underline{O}_X\)-\uline{Modules}. \uline{Then} the \uline{canonical base-change arrow} (EGA~III 1.4.15.4)
\[
  g^*i_*(E)\longrightarrow i'_*f^*(E)
\]
\uline{is an isomorphism}.

(ii) \uline{Put} \(E'=f^*(E)\). \uline{If} \(g\) \uline{is flat and} \(i_*(E)\) \(n\)-\uline{pseudo-coherent}, \uline{then} \(i'_*(E')\) \uline{is} \(n\)-\uline{pseudo-coherent}. \uline{Conversely}, \uline{if} \(g\) \uline{is faithfully flat and} \(i'_*(E')\) \(n\)-\uline{pseudo-coherent}, \uline{then} \(i_*(E)\) \uline{is} \(n\)-\uline{pseudo-coherent}.

\uline{Proof}. The assertion is checked trivially on the stalks.

%%%%%%%%%%  scan idx 238  |  folio 225  %%%%%%%%%%
\SGAPrintPage{225}
\sgaStableTarget{sga6.e03.lem.1-1-4}\textbf{Lemma 1.1.4.}
\uline{Let there be a commutative diagram of schemes}
\[
\begin{tikzcd}
X' \arrow[dr,"f'"'] & X \arrow[l,hook',"i'"'] \arrow[r,hook,"i''"] & X'' \arrow[dl,"f''"]\\
& Y{,} &
\end{tikzcd}
\]
\uline{where} \(i'\) (\uline{resp.} \(i''\)) \uline{is a closed immersion}, \uline{and} \(f'\) (\uline{resp.} \(f''\)) \uline{a smooth morphism}, \uline{and let} \(E\in\ob D(X)\). \uline{We then have}:
\[
  i'_*(E)\ n\text{-\uline{pseudo-coherent}}
  \quad\Longleftrightarrow\quad
  i''_*(E)\ n\text{-\uline{pseudo-coherent}}.
\]

\uline{Proof}. Following a device due to Lichtenbaum, we send \(X\) diagonally into \(X'\times_YX''\), and are at once reduced to proving the assertion when one of the morphisms \(f'\), \(f''\) is the identity. Changing notation, we thus have a commutative triangle
\begin{equation}\tag{$*$}
\begin{tikzcd}
& Y \arrow[d,"f"]\\
X \arrow[r,hook,"j"'] \arrow[ur,hook,"i"] & Z{,}
\end{tikzcd}
\end{equation}
where \(i\) and \(j\) are closed immersions and \(f\) a smooth morphism, and we must show that ``\(i_*(E)\) \(n\)-pseudo-coherent'' is equivalent to ``\(j_*(E)\) \(n\)-pseudo-coherent''. Since the question is local in \(Y\) and in \(Z\) near a point of \(X\), we may suppose that \(X,Y,Z\) are affine and that we have a factorization
\[
\begin{tikzcd}
Y \arrow[r,hook,"i'"] \arrow[d,"f"'] & Z[T_1,\ldots,T_r]\arrow[dl,"f'"]\\
Z{,}&
\end{tikzcd}
\]
where \(i'\) is a closed immersion, and the oblique arrow the canonical projection. Then (VII 1.10) \(i'\) is a regular immersion, so (1.1.1) an object \(F\) of \(D(Y)\) is \(n\)-pseudo-coherent if and only if \(i'_*(F)\) is so.

%%%%%%%%%%  scan idx 239  |  folio 226  %%%%%%%%%%
\SGAPrintPage{226}
We may therefore suppose that \(Y=Z[T_1,\ldots,T_r]\), with \(f\) the canonical projection. The diagram \((*)\) then corresponds to a diagram of rings
\[
\begin{tikzcd}
& B=C[T_1,\ldots,T_r] \arrow[dl,"u"'] & \\
A & C \arrow[l,"v"] \arrow[u,"\varepsilon"'] & {,}
\end{tikzcd}
\]
where \(u\) and \(v\) are epimorphisms, and \(\varepsilon\) the canonical inclusion. For \(1\leq i\leq r\), put \(u(T_i)=x_i\), and let \(z_i\in C\) satisfy \(v(z_i)=x_i\). The retraction \(\sigma:B\longrightarrow C\) of \(\varepsilon\) defined by \(\sigma(T_i)=z_i\) satisfies \(v\sigma=u\). In geometric language, it corresponds to a section \(s\) of \(f\) such that \(i=sj\). By (VII 1.10), \(s\) is a regular immersion, so the assertion follows from (1.1.1).

\uline{Proof of} (1.1)(\uline{end}). Let \(f:X\to Y\) be a locally finite type morphism. There exists an open covering \((X_\alpha)\) of \(X\) such that each \(X_\alpha\) embeds by a closed immersion into a scheme \(Y_\alpha\) smooth over \(Y\). Putting \(\coprod X_\alpha=X'\), \(\coprod Y_\alpha=Y'\), we thus have a commutative diagram
\begin{equation}\tag{$a'$}
\begin{tikzcd}
X' \arrow[r,"i'"] \arrow[d,"p'"'] & Y' \arrow[ddl,"f'"]\\
X \arrow[d,"f"']\\
Y
\end{tikzcd}
\end{equation}
where \(p'\) is a morphism that is the sum of a surjective family of open immersions (we shall say \uline{special} for short), \(i'\) a closed immersion, and \(f'\) a smooth morphism. We shall define \(D(f)_{n\text{-}\coh}\) as the full subcategory of \(D(X)\) formed by the objects \(E\) such that \(i'_*p'^*(E)\) is \(n\)-pseudo-coherent. Let us first show that this definition makes sense, i.e. that \(D(f)_{n\text{-}\coh}\) is independent of the choice of \(i'\), \(p'\), \(f'\). Thus, let
%%%%%%%%%%  scan idx 240  |  folio 227  %%%%%%%%%%
\SGAPrintPage{227}
\begin{equation}\tag{$a''$}
\begin{tikzcd}
X'' \arrow[r,"i''"] \arrow[d,"p''"'] & Y'' \arrow[ddl,"f''"]\\
X \arrow[d,"f"']\\
Y
\end{tikzcd}
\end{equation}
be a commutative diagram analogous to \((a')\) (\(p''\) a special morphism, \(i''\) a closed immersion, \(f''\) a smooth morphism). From \((a')\) and \((a'')\), we deduce a commutative diagram
\[
\begin{tikzcd}
X'\times_X X'' \arrow[rr,hook] \arrow[dr] \arrow[dd] & & Y'\times_Y Y'' \arrow[dr] \arrow[dd] & \\
& X' \arrow[rr,hook,"i'"] \arrow[dd] & & Y' \arrow[dd,"f'"]\\
X'' \arrow[rr,hook,"i''"] \arrow[dr] & & Y'' \arrow[dr,"f''"'] & \\
& X \arrow[rr,"f"] & & Y,
\end{tikzcd}
\]
where the horizontal arrows, except for \(f\), are closed immersions, and the arrows of the left (resp. right) Cartesian square are special morphisms (resp. smooth morphisms). We are therefore reduced to proving the following assertion:

(A): \uline{Let}
\[
\begin{tikzcd}
X_1 \arrow[r,hook,"i_1"] \arrow[d,"p"'] & Y_1 \arrow[d,"q"]\\
X_2 \arrow[r,hook,"i_2"] & Y_2
\end{tikzcd}
\]
\uline{be a commutative square of schemes}, \uline{where} \(p\) \uline{is a special morphism}, \(q\) \uline{a smooth morphism}, \(i_1\) \uline{and} \(i_2\) \uline{closed immersions}, \uline{and let} \(E\in\operatorname{ob} D(X_2)\). \uline{Then}, \uline{for} \(i_{2*}(E)\) \uline{to be} \(n\)-\uline{pseudo-coherent}, \uline{it is necessary and sufficient that} \(i_{1*}p^*(E)\) \uline{be so}.

%%%%%%%%%%  scan idx 241  |  folio 228  %%%%%%%%%%
\SGAPrintPage{228}
\uline{Proof of} (A). By definition, \(p\) is the sum of a family of open subsets \(U_\alpha\) of \(X_2\) that cover \(X_2\). Each \(U_\alpha\) is induced by an open subset \(V_\alpha\) of \(Y_2\), and the \(V_\alpha\) cover an open neighborhood \(X'_2\) of \(X_2\) in \(Y_2\). We therefore have a commutative diagram with a Cartesian square
\[
\begin{tikzcd}
X_1 \arrow[r,hook,"i'_1"] \arrow[d,"p"'] & X'_1 \arrow[d,"p'"]\\
X_2 \arrow[r,hook,"i'_2"] \arrow[rr,bend right,"i_2"'] & X'_2 \arrow[r,hook,"j"] & Y_2,
\end{tikzcd}
\]
where \(i'_1\) and \(i'_2\) are closed immersions, \(p\) and \(p'\) special morphisms, and \(j\) an open immersion. The condition ``\(i_{2*}(E)\) \(n\)-pseudo-coherent'' is clearly equivalent to ``\(i'_{2*}(E)\) \(n\)-pseudo-coherent''. On the other hand, it follows at once from (1.1.3) that the condition ``\(i'_{2*}(E)\) \(n\)-pseudo-coherent'' is equivalent to ``\(i'_{1*}p^*(E)\) \(n\)-pseudo-coherent''. Applying (1.1.4) to the factorizations \((jp')i'_1\) and \(qi_1\) of \(i_2p\), and to \(p^*(E)\in\Ob D(X_1)\), we finally conclude that ``\(i_{2*}(E)\) \(n\)-pseudo-coherent'' is equivalent to ``\(i_{1*}p^*(E)\) \(n\)-pseudo-coherent'', q.e.d.

We have thus shown that the category \(D(f)_{n\text{-}\coh}\) defined above by means of the diagram \((a')\) in fact depends only on \(f\). Moreover, it is trivial from the definition that it does satisfy conditions (i) to (iii) of (1.1). This therefore completes the proof of (1.1).

\sgaStableTarget{sga6.e03.def.1-2}\textbf{Definition 1.2.} \uline{Let} \(f:X\to Y\) \uline{be a locally finite type morphism}, \(E\in\Ob D(X)\), \(n\in\mathbb Z\). \uline{We shall say that} \(E\) \uline{is} \(n\)-pseudo-coherent relative to \(f\) \uline{if} \(E\in\Ob D(f)_{n\text{-}\coh}\) (cf. (1.1)). \uline{We shall say that} \(E\) \uline{is} pseudo-coherent (\uline{or equivalently} \((-\infty)\)-pseudo-coherent) relative to \(f\) \uline{if} \(E\) \uline{is} \(n\)-\uline{pseudo-coherent relative to} \(f\) \uline{for every} \(n\in\mathbb Z\). \uline{We shall write} \(D(f)_{\mathrm{coh}}\) \uline{for the full subcategory}
%%%%%%%%%%  scan idx 242  |  folio 229  %%%%%%%%%%
\SGAPrintPage{229}
\uline{of} \(D(X)\) \uline{formed by the objects pseudo-coherent relative to} \(f\). \uline{We shall say that} \(f\) \uline{is} \(n\)-\uline{pseudo-coherent} (\uline{resp.} \uline{pseudo-coherent}) \uline{if} \(\underline{O}_X\) \uline{is} \(n\)-\uline{pseudo-coherent} (\uline{resp.} \uline{pseudo-coherent}) \uline{relative to} \(f\). \uline{We shall sometimes say}, \uline{when this cannot lead to confusion}, ``\(n\)-\uline{pseudo-coherent} (\uline{resp.} \uline{pseudo-coherent}) \uline{relative to} \(Y\)'' \uline{instead of} ``\uline{relative to} \(f\)'', \uline{and shall also write} \(D(X)_{Y\text{-}n\text{-}\coh}\) (\uline{resp.} \(D(X)_{Y\text{-}\coh}\)) \uline{instead of} \(D(f)_{n\text{-}\coh}\) (\uline{resp.} \(D(f)_{\mathrm{coh}}\)). \uline{Finally, for} \(*=+,-,b\), \uline{we shall put}
\[
\begin{gathered}
  D^*(f)_{n\text{-}\coh}=D^*(X)\cap D(f)_{n\text{-}\coh},\\
  D^*(f)_{\mathrm{coh}}=D^*(X)\cap D(f)_{\mathrm{coh}}.
\end{gathered}
\]

\sgaStableTarget{sga6.e03.rem.1-2-1}\textbf{Remarks 1.2.1.}
\begin{enumerate}[label=\alph*)]
\item The categories \(D(f)_{n\text{-}\coh}\), \(D(f)_{\mathrm{coh}}\) are not in general contained in \(D^{-}(X)\). This is so, however, when \(X\) is quasi-compact, since an object of \(D(X)\) with locally bounded-above cohomology then has bounded-above cohomology.

\item The objects of \(D(f)_{\mathrm{coh}}\) have quasi-coherent cohomology. Indeed, let \(E\in\Ob D(f)_{\mathrm{coh}}\); the question being local, we may suppose that \(f\) is a closed immersion. Then \(f_*(E)\) is pseudo-coherent, hence has quasi-coherent cohomology, and the result follows from the fact that an \(\underline{O}_X\)-Module \(M\) is quasi-coherent if and only if \(f_*(M)\) is so.

\item When speaking of an \(\underline{O}_X\)-Module, we shall use the terminology ``of finite presentation rel. to \(f\)'' in preference to ``\((-1)\)-pseudo-coherent rel. to \(f\)''. For an \(\underline{O}_X\)-Module \(E\) to be of finite presentation rel. to \(f\), it is necessary and sufficient, in view of (I 2.9), that there exist locally a closed immersion \(i:X\to X'\), with \(X'\) smooth over \(Y\), such that \(i_*(E)\) is of finite presentation (in the ordinary sense). In particular, saying that \(f\) is
%%%%%%%%%%  scan idx 243  |  folio 230  %%%%%%%%%%
\SGAPrintPage{230}
\((-1)\)-pseudo-coherent is equivalent to saying that \(f\) is locally of finite presentation in the sense of (EGA IV 1.4).
\end{enumerate}

\sgaStableTarget{sga6.e03.ex.1-3}\textbf{Examples 1.3.} A smooth morphism and a regular immersion (1.1.2) are pseudo-coherent morphisms. The same holds for any locally finite type morphism \(f:X\to Y\) when \(Y\) is locally noetherian.

\sgaStableTarget{sga6.e03.prop.1-4}\textbf{Proposition 1.4.} \uline{Let} \(f:X\to Y\) \uline{be a locally finite type morphism}, \uline{and let}
\[
\begin{tikzcd}
& G \arrow[dl] & \\
E \arrow[rr] && F \arrow[ul]
\end{tikzcd}
\]
\uline{be a distinguished triangle in} \(D(X)\). \uline{We then have the following table}, \uline{where} \(n\in\mathbb Z\cup\{-\infty\}\) \uline{indicates a degree of pseudo-coherence relative to} \(f\), \uline{and each row represents an implication whose conclusion is underlined}:
\[
\begin{array}{ccc}
E&F&G\\
n+1&n&\underline n\\
n&\underline n&n\\
\underline n&n&n-1.
\end{array}
\]
\uline{In particular}, \(D(f)_{\mathrm{coh}}\) \uline{is a triangulated subcategory of} \(D(X)\).

\uline{Proof}. By properties (i) to (iii) of (1.1), we may suppose that \(f\) is a closed immersion. We then have a distinguished triangle
\[
\begin{tikzcd}
& f_*(G) \arrow[dl] & \\
f_*(E) \arrow[rr] && f_*(F) \arrow[ul]
\end{tikzcd},
\]
and the assertion follows from (I 2.6).

\sgaStableTarget{sga6.e03.def.1-5}\textbf{Definition 1.5.} \uline{Let} \(X\) \uline{and} \(Y\) \uline{be two} \(S\)-\uline{schemes}. \uline{We shall say that} \(X\) \uline{and} \(Y\) \uline{are tor-independent over} \(S\) \uline{if we have}
\[
  \underline{\Tor}_i^S(\underline{O}_X,\underline{O}_Y)=0\quad\text{\uline{for} }i\ne 0,
\]
\uline{where the} \(\underline{\Tor}_i^S\) \uline{are the local hyper-tor defined in} (EGA III 6.5).
%%%%%%%%%%  scan idx 244  |  folio 231  %%%%%%%%%%
\SGAPrintPage{231}
This is a condition of a local nature; when \(X,Y,S\) are affine with respective rings \(B,C,A\), it means that the canonical arrow
\[
  B\otimes_A^{L}C\longrightarrow B\otimes_A C
\]
is an isomorphism. It is therefore satisfied in particular when \(X\) or \(Y\) is flat over \(S\).

\sgaStableTarget{sga6.e03.lem.1-5-1}\textbf{Lemma 1.5.1.} \uline{Let there be Cartesian squares of schemes}
\[
\begin{tikzcd}
X \arrow[d] & X' \arrow[l] \arrow[d] & X'' \arrow[l] \arrow[d]\\
Y & Y' \arrow[l] & Y'' \arrow[l].
\end{tikzcd}
\]
If \(X\) and \(Y'\) are tor-independent over \(Y\), and if \(X'\) and \(Y''\) are tor-independent over \(Y'\), then \(X\) and \(Y''\) are tor-independent over \(Y\).

\uline{Proof}. An easy exercise on associativity of the tensor product.

\sgaStableTarget{sga6.e03.notation.1-6}\textbf{Notation 1.6.} Let \(X\) and \(Y\) be two \(S\)-schemes, \(E\in\operatorname{ob} D^{-}(X)\), \(F\in\operatorname{ob} D^{-}(Y)\). We shall put
\[
  E\otimes_S^{L}F
  =
  L\pr_1^*(E)\otimes^{L}L\pr_2^*(F),
\]
where \(\pr_1,\pr_2\) are the two projections from \(X\times_SY\).

\sgaStableTarget{sga6.e03.lem.1-7}\textbf{Lemma 1.7.} \uline{Let} \(S\) \uline{be a scheme}, \uline{and, for} \(i=1,2\), \(f_i:X_i\longrightarrow Y_i\) \uline{be a closed} \(S\)-\uline{immersion}. \uline{Suppose that} \(X_1\) \uline{and} \(X_2\), \uline{as well as} \(Y_1\) \uline{and} \(Y_2\), \uline{are tor-independent over} \(S\). \uline{Then}, \uline{for} \(E_i\in\operatorname{ob} D^{-}(X_i)_{\mathrm{qcoh}}\)\SGAAuthorFootnote{1}{\(D(X_i)_{\mathrm{qcoh}}\) denotes the full subcategory of \(D(X_i)\) formed by complexes with quasi-coherent cohomology.}, \uline{there exists a canonical isomorphism}
\[
  f_{1*}(E_1)\otimes_S^{L}f_{2*}(E_2)
  \xrightarrow{\sim}
  (f_1\times_Sf_2)_*(E_1\otimes_S^{L}E_2).
\]

\uline{Proof}. Put \(f_1\times_Sf_2=f\). By adjunction, one easily defines a canonical arrow
%%%%%%%%%%  scan idx 245  |  folio 232  %%%%%%%%%%
\SGAPrintPage{232}
\[
  f_{1*}(E_1)\otimes_S^{L}f_{2*}(E_2)
  \longrightarrow
  f_*(E_1\otimes_S^{L}E_2).
\]
Let us show that it is an isomorphism. The question being local, we may suppose that \(S,X_i,Y_i\) are affine with respective rings \(C,A_i,B_i\). By dévissage, we reduce to \(E_i=\underline{O}_{X_i}\). We must then check that
\[
  A_1\otimes_{B_1}^{L}B\otimes_{B_2}^{L}A_2
  \longrightarrow A
\]
is an isomorphism (we have put \(A=A_1\otimes_C A_2\), \(B=B_1\otimes_C B_2\)). But, taking account of the tor-independence hypotheses, this follows at once from associativity of the tensor product.

\sgaStableTarget{sga6.e03.prop.1-8}\textbf{Proposition 1.8.} \uline{Let} \(S\) \uline{be a scheme}, \uline{and}, \uline{for} \(i=1,2\), \(f_i:X_i\longrightarrow Y_i\) \uline{be a locally finite type} \(S\)-\uline{morphism}. \uline{Suppose} \(X_1\) \uline{and} \(X_2\), \uline{as well as} \(Y_1\) \uline{and} \(Y_2\), \uline{tor-independent over} \(S\). \uline{Let} \([a_i,b_i]\) \uline{be an interval of} \(\mathbb Z\), \uline{and} \(E_i\in\operatorname{ob} D^{-}(X_i)_{\mathrm{qcoh}}\) \uline{a complex that is} \(a_i\)-\uline{pseudo-coherent rel.} \uline{to} \(f_i\) \uline{and acyclic in degree} \(>b_i\). \uline{Then} \(E_1\otimes_S^{L}E_2\) \uline{is} \(a\)-\uline{pseudo-coherent rel.} \uline{to} \(f\) \uline{and acyclic in degree} \(>b\), \uline{with} \(f=f_1\times_Sf_2\), \(a=\operatorname{Sup}(a_1+b_2,a_2+b_1)\), \(b=b_1+b_2\).

\uline{Proof}. The question being local on \(X_1\times_SX_2\), we may suppose that \(f_i\) factors as
\[
  X_i\xrightarrow{f'_i}Y'_i\xrightarrow{p_i}Y_i,
\]
where \(f'_i\) is a closed immersion and \(p_i\) a smooth morphism. Putting \(f'=f'_1\times_Sf'_2\), we must show that \(f'_*(E_1\otimes_S^{L}E_2)\) is \(a\)-pseudo-coherent and acyclic in degree \(>b\). By (1.5.1), \(Y'_1\) and \(Y'_2\) are tor-independent over \(S\): after changing notation, we may therefore suppose that \(f_i=f'_i\), in other words that \(f_i\) is a closed immersion. By hypothesis, \(f_{i*}(E_i)\) is \(a_i\)-pseudo-coherent and acyclic in degree \(>b_i\), so (I 2.16) \(f_{1*}(E_1)\otimes_S^{L}f_{2*}(E_2)\) is \(a\)-pseudo-coherent and acyclic in degree \(>b\). We conclude by (1.7).

%%%%%%%%%%  scan idx 246  |  folio 233  %%%%%%%%%%
\SGAPrintPage{233}
\sgaStableTarget{sga6.e03.cor.1-9}\textbf{Corollary 1.9.} \uline{Under the hypotheses of} (1.8), \uline{the functor}
\[
  \otimes_S^{L}:D^{-}(X_1)\times D^{-}(X_2)\longrightarrow D^{-}(X)
\]
\uline{induces a functor}
\[
  \otimes_S^{L}:D^{-}(f_1)_{\mathrm{coh}}\times D^{-}(f_2)_{\mathrm{coh}}
  \longrightarrow D^{-}(f)_{\mathrm{coh}}\qquad(\text{cf. }(1.2)).
\]
\uline{In particular}, \uline{if} \(f_1\) \uline{and} \(f_2\) \uline{are pseudo-coherent} (1.2), \uline{the same holds for} \(f\).

\sgaStableTarget{sga6.e03.cor.1-10}\textbf{Corollary 1.10.} \uline{Let}
\[
\begin{tikzcd}
X \arrow[d,"f"'] & X' \arrow[l,"g"'] \arrow[d,"f'"]\\
Y & Y' \arrow[l,"h"']
\end{tikzcd}
\]
\uline{be a Cartesian square}, \uline{with} \(f\) \uline{locally of finite type}. \uline{Then}:
\begin{enumerate}[label=\alph*)]
\item \uline{Suppose} \(X\) \uline{and} \(Y'\) \uline{tor-independent over} \(Y\) (1.5). \uline{For} \(n\in\mathbb Z\cup\{-\infty\}\), \uline{the functor} \(Lg^*:D^{-}(X)\to D^{-}(X')\) \uline{induces a functor}
\[
  Lg^*:D^{-}(f)_{n\text{-}\coh}\longrightarrow D^{-}(f')_{n\text{-}\coh}.
\]
\uline{In particular}, \uline{if} \(f\) \uline{is} \(n\)-\uline{pseudo-coherent} (1.2), \uline{the same holds for} \(f'\).

\item \uline{Suppose} \(h\) \uline{faithfully flat}, \uline{let} \(n\in\mathbb Z\cup\{-\infty\}\) \uline{and} \(E\in\operatorname{ob} D^{-}(X)\) \uline{be a complex with quasi-coherent cohomology}. \uline{For} \(E\) \uline{to be} \(n\)-\uline{pseudo-coherent rel.} \uline{to} \(f\), \uline{it is necessary and sufficient that} \(g^*(E)\) \uline{be} \(n\)-\uline{pseudo-coherent} \uline{relative to} \(f'\). \uline{In particular}, \uline{for} \(f\) \uline{to be} \(n\)-\uline{pseudo-coherent}, \uline{it is necessary and sufficient that} \(f'\) \uline{be so}.
\end{enumerate}

\uline{Proof}. Taking account of (1.2.1 b)), all the assertions of (1.9) and (1.10) follow trivially from (1.8), except the ``sufficiency'' part of (1.10 b)). For this, the question being local on \(X\), we may suppose that \(f\) factors as a closed immersion followed by a smooth morphism. By (1.1.3), we are then reduced to the case where \(f\) (resp. \(f'\)) is the identity of \(X\) (resp. \(X'\)). The assertion then follows, by a standard dévissage (cf. I 2.16.2), from the fact that a quasi-coherent sheaf \(L\) on \(X\) is of finite type if and only if \(g^*L\) is of finite type.

%%%%%%%%%%  scan idx 247  |  folio 234  %%%%%%%%%%
\SGAPrintPage{234}
\sgaStableTarget{sga6.e03.prop.1-11}\textbf{Proposition 1.11.} \uline{Let}
\[
\begin{tikzcd}
X \arrow[r,"f"] \arrow[dr,"h"'] & Y \arrow[d,"g"]\\
& Z
\end{tikzcd}
\]
\uline{be a commutative triangle of locally finite type morphisms}, \uline{let} \(E\in\operatorname{ob} D(X)\), \(F\in\operatorname{ob} D(Y)\), \uline{and} \([a,b]\), \([c,d]\) \uline{be intervals of} \(\mathbb Z\).
\begin{enumerate}[label=\alph*)]
\item \uline{If} \(E\) \uline{is} \(a\)-\uline{pseudo-coherent rel.} \uline{to} \(f\) \uline{and acyclic in degree} \(>b\) \uline{and if} \(F\) \uline{is} \(c\)-\uline{pseudo-coherent rel.} \uline{to} \(g\) \uline{and acyclic in degree} \(>d\), \uline{then} \(E\otimes_Y^{L}F\) (cf. (1.6)) \uline{is} \(m\)-\uline{pseudo-coherent rel.} \uline{to} \(h\) \uline{and acyclic in degree} \(>n\), \uline{where} \(m=\operatorname{Sup}(a+d,b+c)\), \uline{and} \(n=b+d\).

\item \uline{If} \(E\) \uline{is} \(a\)-\uline{pseudo-coherent rel.} \uline{to} \(h\) \uline{and acyclic in degree} \(>b\), \uline{and if} \(g\) \uline{is} \((a+1-b)\)-\uline{pseudo-coherent}, \uline{then} \(E\) \uline{is} \(a\)-\uline{pseudo-coherent rel.} \uline{to} \(f\).
\end{enumerate}

\uline{Proof}. The question being local on \(X\), we may suppose that we have a commutative diagram with a Cartesian square
\[
\begin{tikzcd}[column sep=large,row sep=large]
& & E^{r+s}_Z \arrow[dr] & \\
& E^r_Y \arrow[ur,hook,"j"] \arrow[d,"p"'] & & E^s_Z \arrow[d,"q"]\\
X \arrow[ur,hook,"i"] \arrow[r,"f"] & Y \arrow[urr,hook,"k"'] \arrow[rr,"g"] & & Z{,}
\end{tikzcd}
\]
where \(i,j,k\) are closed immersions, and \(p,q\) the canonical projections of the standard affine spaces. If \(F\) is \(c\)-pseudo-coherent rel. to \(Z\), then \(k_*(F)\) is \(c\)-pseudo-coherent by (1.1 (iii)), and \(j_*(p^*(F))\) is \(c\)-pseudo-coherent by (1.1.3). Similarly, if \(g\) is \(x\)-pseudo-coherent, then \(j_*(\underline{O}_{E^r_Y})\) is \(x\)-pseudo-coherent. By the definition of relative pseudo-coherence (1.2), we are therefore reduced to proving a) and b) when \(f\) and \(g\) are closed immersions.

%%%%%%%%%%  scan idx 248  |  folio 235  %%%%%%%%%%
\SGAPrintPage{235}
\uline{Proof of} a). By (1.1 (iii)), we must show the following: suppose \(f_*(E)\) is \(a\)-pseudo-coherent and acyclic in degree \(>b\), and \(g_*(F)\) is \(c\)-pseudo-coherent and acyclic in degree \(>d\); then \(h_*(E\otimes_Y^{L}F)\) is \(m\)-pseudo-coherent and acyclic in degree \(>n\). Now we have a canonical isomorphism
\[
  f_*(E)\otimes^{L}F\xrightarrow{\sim} f_*(E\otimes_Y^{L}F)
\]
(a special case of the isomorphism in (1.7)). Since the functor \(g_*(.\otimes^{L}F)\) takes locally free sheaves of finite type to \(c\)-pseudo-coherent complexes, and complexes acyclic in degree \(>0\) to complexes acyclic in degree \(>d\), the assertion therefore follows from (I 2.13).

\uline{Proof of} b). We must prove the following assertion: suppose \(h_*(E)\) is \(a\)-pseudo-coherent and acyclic in degree \(>b\), and \(g_*(\underline{O}_Y)\) is \((a+1-b)\)-pseudo-coherent; then \(f_*(E)\) is \(a\)-pseudo-coherent. Since the assertion concerns \(f_*(E)\), we may suppose \(X=Y\) and \(f=\Id_X\). We then proceed as in the proof of (1.1.1), proving by descending induction that \(E\) is pseudo-coherent in degree \(\ge a\). We suppose it already proved that \(E\) is \((a+1)\)-pseudo-coherent, whence (I 2.10), after localizing, a distinguished triangle
\[
\begin{tikzcd}
& E'' \arrow[dl,hook] & \\
E' \arrow[rr] & & E \arrow[ul],
\end{tikzcd}
\]
where \(E'^i\) is locally free for \(i\in[a+1,b]\), \(E'^i=0\) for \(i\notin[a+1,b]\), and \(E''^i=0\) for \(i>a\). The hypothesis that \(g_*(\underline{O}_Y)\) is \((a+1-b)\)-pseudo-coherent then implies (I 2.11 a)) that \(g_*(E')\) is \((a+1)\)-pseudo-coherent. Since \(g_*(E)\) is \(a\)-pseudo-coherent, it follows from (I 2.6) that \(g_*(E'')\) is \(a\)-pseudo-coherent. We conclude by (I 2.10 b)) and the fact that an \(\underline{O}_Y\)-Module \(M\) is of finite type if and only if \(g_*(M)\) is so.

The following corollaries are deduced trivially from (1.11):

%%%%%%%%%%  scan idx 249  |  folio 236  %%%%%%%%%%
\SGAPrintPage{236}
\sgaStableTarget{sga6.e03.cor.1-12}\textbf{Corollary 1.12.} \uline{In the situation of} (1.11), \uline{if} \(g\) \uline{is pseudo-coherent}, \uline{we have}, \uline{for} \(r\in\mathbb Z\cup\{-\infty\}\),
\[
  D^{-}(f)_{r\text{-}\coh}=D^{-}(h)_{r\text{-}\coh};
\]
\uline{in particular} (\uline{taking} \(X=Y\)), \uline{the notions of} \(r\)-\uline{pseudo-coherence over} \(Y\), ``\uline{in the absolute sense}'' \uline{or rel.} \uline{to} \(g\), \uline{coincide}.

\sgaStableTarget{sga6.e03.cor.1-13}\textbf{Corollary 1.13.} \uline{In the situation of} (1.11), \uline{if} \(f\) \uline{is pseudo-coherent}, \uline{the functor} \(\underline{O}_X\otimes_Y^{\mathrm L}\cdot=Lf^*:D^{-}(Y)\longrightarrow D^{-}(X)\) \uline{induces}, \uline{for} \(r\in\mathbb Z\cup\{-\infty\}\), \uline{a functor}
\[
  Lf^*:D^{-}(g)_{r\text{-}\coh}\longrightarrow D^{-}(h)_{r\text{-}\coh}.
\]

\sgaStableTarget{sga6.e03.cor.1-14}\textbf{Corollary 1.14.} \uline{In the situation of} (1.11), \uline{if} \(f\) \uline{is} \(r\)-\uline{pseudo-coherent} \uline{and} \(g\) \uline{is} \(s\)-\uline{pseudo-coherent} (\(r,s\in\mathbb Z\cup\{-\infty\}\)), \uline{then} \(h\) \uline{is} \(\operatorname{Sup}(r,s)\)-\uline{pseudo-coherent}. \uline{If} \(h\) \uline{is} \(r\)-\uline{pseudo-coherent} \uline{and} \(g\) \uline{is} \((r+1)\)-\uline{pseudo-coherent}, \(f\) \uline{is} \(r\)-\uline{pseudo-coherent}.

\sgaStableTarget{sga6.e03.rem.1-14-1}\textbf{Remark 1.14.1.} In particular, for \(r=s=-1\), we recover (cf. (1.2.1 c))) assertions (ii) and (v) of (EGA IV 1.4.3).

\sgaStableTarget{sga6.e03.sec.2}\section*{2. The finiteness theorem.}

\sgaStableTarget{sga6.e03.conj.2-1}\textbf{Conjecture 2.1.} \uline{Let} \(S\) \uline{be a scheme}, \(f:X\to Y\) \uline{be a} proper \(S\)-\uline{morphism of schemes locally of finite type over} \(S\), \(E\in\Ob D^+(X)\). \uline{Then, if} \(E\) \uline{is pseudo-coherent rel.} \uline{to} \(S\) (1.2), \uline{the same holds for} \(Rf_*(E)\).

\sgaStableTarget{sga6.e03.thm.2-2}\textbf{Theorem 2.2.} \uline{The preceding conjecture is true in the following cases}:
\begin{enumerate}[label=(\roman*)]
\item \(S\) \uline{is locally noetherian};
\item \(f\) \uline{is projective}.
\end{enumerate}

%%%%%%%%%%  scan idx 250  |  folio 237  %%%%%%%%%%
\SGAPrintPage{237}
\uline{Proof in case} (i). Since \(S\) is locally noetherian, the projection of \(X\) (resp. \(Y\)) onto \(S\) is a pseudo-coherent morphism (1.3), and the notions of pseudo-coherence over \(X\) (resp. \(Y\)) and relative to \(S\) coincide (1.12). On the other hand, since \(X\) (resp. \(Y\)) is locally noetherian, hence \(\underline{O}_X\) (resp. \(\underline{O}_Y\)) coherent, an object \(M\) of \(D(X)\) (resp. \(D(Y)\)) is pseudo-coherent if and only if (I 3.5) the \(H^i(M)\) are coherent and \(M\) is locally in \(D^{-}\). The theorem therefore follows from the following one:

\sgaStableTarget{sga6.e03.thm.2-2-1}\textbf{Theorem 2.2.1.} \uline{Let} \(f:X\to Y\) \uline{be a proper morphism of locally noetherian schemes}, \uline{and let} \(E\in\Ob D^+(X)\). \uline{If the} \(H^i(E)\) \uline{are coherent}, \uline{the same holds for the} \(R^if_*(E)\). \uline{If, moreover,} \(E\) \uline{is locally in} \(D^{-}\), \uline{the same holds for} \(Rf_*(E)\).

\uline{Proof}. The first assertion follows from the spectral sequence
\begin{equation}\tag{$*$}
  R^pf_*(H^q(E)) \Longrightarrow R^*f_*(E)
\end{equation}
and the finiteness theorem (EGA III 3.2.1). For the second, we may suppose \(Y\) affine; \(X\) is then quasi-compact and separated. Suppose it covered by \(r\) affine open subsets. We know (EGA III 1.4.12) that, under these conditions, \(R^if_*(M)=0\) for every quasi-coherent \(\underline{O}_X\)-Module \(M\) and every \(i>r\). If \(E\) is locally in \(D^{-}\), \(E\) is in \(D^{-}\), since \(X\) is quasi-compact; say \(H^i(E)=0\) for \(i>N\). From the spectral sequence (*) we then deduce that \(R^if_*(E)=0\) for \(i>N+r\), whence the assertion.

\uline{Proof in case} (ii). The question being local on \(Y\), we may suppose \(S\) and \(Y\) affine. We then have a commutative diagram with a Cartesian square
\[
\begin{tikzcd}[column sep=large,row sep=large]
& & P^r_{Y'} \arrow[dr,"p'"] & \\
& P^r_Y \arrow[ur,hook,"j"] \arrow[d,"p"'] & & Y' \arrow[d,"g'"]\\
X \arrow[ur,hook,"i"] \arrow[r,"f"] & Y \arrow[urr,hook,"j'"] \arrow[rr,"g"] & & S{,}
\end{tikzcd}
\]
%%%%%%%%%%  scan idx 251  |  folio 238  %%%%%%%%%%
\SGAPrintPage{238}
where \(i,j,j'\) are closed immersions, \(g'\) a smooth affine morphism, and \(p,p'\) the projections of the standard projective spaces. Suppose \(E\) pseudo-coherent rel. to \(S\). By (1.1 (iii)), this means that \((ji)_*(E)\) is pseudo-coherent. We must show that \(Rf_*(E)\) is pseudo-coherent rel. to \(S\), i.e. (1.1 (iii)) that \(j'_*Rf_*(E)\) is pseudo-coherent. But we have
\[
  j'_*Rf_*(E)\simeq Rp'_*(ji)_*(E),
\]
so the assertion will follow from the following more precise result:

\sgaStableTarget{sga6.e03.thm.2-2-2}\textbf{Theorem 2.2.2.} \uline{Let} \(Y\) \uline{be an affine scheme}, \(f:X=P^r_Y\longrightarrow Y\) \uline{be the projection of a standard projective space}, \(E\in\Ob D^b(X)_{\mathrm{coh}}\), \(a\in\mathbb Z\). \uline{Then} \(Rf_*(E)\) \uline{is pseudo-coherent}. \uline{Moreover}, \uline{if} \(E\) \uline{is acyclic in degree} \(>a\), \(Rf_*(E)\) \uline{is acyclic in degree} \(>a+r\), \uline{and there exists an integer} \(N\) \uline{such that}, \uline{for every} \(n\geq N\), \(Rf_*(E(n))\) (\uline{where} \(E(n)=E\otimes\underline{O}_X(n)\)) \uline{is acyclic in degree} \(>a\).

\uline{Proof}. Since \(Y\) is affine, \(X\) is quasi-compact and separated, so (II 2.2.9 a)) \(E\) is isomorphic, in \(D(X)\), to a complex \(F\) bounded above in degree such that, for every \(i\), \(F^i\in\Ob \underline{L}(X)\), i.e. (loc. cit.) is a finite direct sum of sheaves of the type \(\underline{O}_X(-n)\) \((n>0)\). Moreover, if \(E\) is acyclic in degree \(>a\), we may choose \(F\) so that \(F^i=0\) for \(i>a\): this follows formally from the preceding assertion, taking account of (I 2.10 a) and II 2.2.8 a)). Replacing \(E\) by \(F\), we may therefore suppose that \(E^i\in\Ob \underline{L}(X)\) for every \(i\), and \(E^i=0\) for \(i>a\). Since \(X\) is covered by \((r+1)\) affine open subsets, we have, by (EGA III 1.4.10) and the hypercohomology spectral sequence, \(R^qf_*(E)=0\) for \(q>a+r\). Let \(t\) be an integer \(\leq a\), and put \(E'=(0\to E^t\to\ldots\to E^a\to0)\). We have an exact sequence
\[
  0\longrightarrow E'\longrightarrow E\longrightarrow E''\longrightarrow0,
\]
%%%%%%%%%%  scan idx 252  |  folio 239  %%%%%%%%%%
\SGAPrintPage{239}
where \(E''\) satisfies the same conditions as \(E\), but with \(a\) replaced by \(t-1\). Applying \(Rf_*\), we obtain a distinguished triangle
\begin{equation}\tag{$*$}
\begin{tikzcd}
& Rf_*(E'') \arrow[dl,hook] & \\
Rf_*(E') \arrow[rr] & & Rf_*(E) \arrow[ul] & .
\end{tikzcd}
\end{equation}
Since \(t\) can be chosen arbitrarily far in the negative direction, it suffices, for the first assertion, to show that \(Rf_*(E)\) is \((t+r)\)-pseudo-coherent. Now \(Rf_*(E'')\) is acyclic in degree \(\geq t+r\), hence \((t+r)\)-pseudo-coherent. Thus, by (I 2.6), the result will follow once we prove that \(Rf_*(E')\) is perfect. For this, we reduce by dévissage (I 4.10) to the case where \(E'\) is concentrated in degree \(0\) and equal to \(\underline{O}_X(-n)\), \(n>0\). Then, by (EGA III 2.1.16), \(R^rf_*(\underline{O}_X(-n))\) is locally free of finite type, and we have
\[
  Rf_*(\underline{O}_X(-n))\simeq R^rf_*(\underline{O}_X(-n))[-r].
\]
Thus \(Rf_*(\underline{O}_X(-n))\) is perfect, and we deduce the first assertion of (2.2.2). We have already seen that \(Rf_*(E)\) is acyclic in degree \(>a+r\). It remains to prove the last assertion. For \(t<-r+a\), \(Rf_*(E'')\) is acyclic in degree \(>a\), so, taking account of the triangle \((*)\), we are reduced to proving the result for \(E'\) in place of \(E\). Since \(E'\) is bounded in degree and its components are objects of \(\underline{L}(X)\), there exists an integer \(N\) such that, for \(n\geq N\), the components of \(E'(n)\) are finite direct sums of sheaves of the type \(\underline{O}_X(p)\), with \(p\geq0\). Then, by (EGA III 2.1.16) and the hypercohomology spectral sequence, \(Rf_*(E'(n))\) is acyclic in degree \(>a\) for every \(n\geq N\), which completes the proof of (2.2.2), and hence of (2.2).

We have a ``classical'' complement to (2.2 (ii)), which we shall give for later reference:

%%%%%%%%%%  scan idx 253  |  folio 240  %%%%%%%%%%
\SGAPrintPage{240}
\sgaStableTarget{sga6.e03.cor.2-3}\textbf{Corollary 2.3} (EGA III 2.2.1). \uline{Let} \(S\) \uline{be a quasi-compact scheme}, \(f:X\to Y\) \uline{be a projective} \(S\)-\uline{morphism of schemes of finite type over} \(S\), \(L=i^*\underline{O}_P(1)\) \uline{be the invertible sheaf defined by an embedding} \(i:X\to P=P(F)\), \uline{where} \(F\) \uline{is} \uline{a quasi-coherent} \(\underline{O}_Y\)-\uline{Module of finite type}. \uline{For} \(M\in\Ob D(X)\) \uline{and} \(n\in\mathbb Z\), \uline{put} \(M(n)=M\otimes^{\mathrm L} L^{\otimes n}\). \uline{Then}, \uline{if} \(E\in\Ob D^+(X)\) \uline{is pseudo-coherent rel.} \uline{to} \(S\) \uline{and} \uline{acyclic in degree} \(>a\) \((a\in\mathbb Z)\), \uline{there exists an integer} \(N\) \uline{such that}, \uline{for every} \(n\geq N\), \(Rf_*E(n)\) \uline{satisfies the same conditions}.

\uline{Proof}. We reduce trivially, as in the proof of (2.2 (ii)), to the situation of (2.2.2): \(Y=S\), \(f:X=P^r_Y\to Y\) the canonical projection, \(L=\underline{O}_X(1)\) the canonical sheaf, \(E\in\Ob D^b(X)\) a pseudo-coherent complex acyclic in degree \(>a\), and the result follows by (2.2.2).

\sgaStableTarget{sga6.e03.cor.2-4}\textbf{Corollary 2.4.} \uline{Let} \(S\) \uline{be a quasi-compact scheme}, \uline{and} \(f:X\to Y\) \uline{be a proper} \(S\)-\uline{morphism of schemes of finite type over} \(S\). \uline{If} \(S\) \uline{is noetherian or} \(f\) \uline{is projective}, \uline{the functor} \(Rf_*:D^+(X)\longrightarrow D^+(Y)\) \uline{induces a functor}
\[
  Rf_*:D^b(X)_{S\text{-}\coh}\longrightarrow D^b(Y)_{S\text{-}\coh}
\]

\uline{Proof}. A consequence of (2.2) and the fact that there is then (EGA III 1.4.12 and IV 1.7.21) an integer \(N\) such that \(R^if_*(E)=0\) for every \(i>N\) and every quasi-coherent \(\underline{O}_X\)-Module \(E\).

\sgaStableTarget{sga6.e03.cor.2-5}\textbf{Corollary 2.5.} \uline{Let} \(f:X\to Y\) \uline{be a proper and pseudo-coherent morphism} (1.2), \uline{and} \(E\in\Ob D^+(X)\). \uline{Suppose} \(Y\) \uline{locally noetherian or} \(f\) \uline{projective}. \uline{Then}, \uline{if} \(E\) \uline{is pseudo-coherent} (\uline{in the absolute sense}), \uline{the same holds for} \(Rf_*(E)\). \uline{If} \(Y\) \uline{is quasi-compact}, \uline{the functor} \(Rf_*\) \uline{induces a functor}
\[
  Rf_*:D^b(X)_{\mathrm{coh}}\longrightarrow D^b(Y)_{\mathrm{coh}}.
\]

%%%%%%%%%%  scan idx 254  |  folio 241  %%%%%%%%%%
\SGAPrintPage{241}%
\uline{Proof}. A consequence of (2.2), (2.4) and (1.12).

\sgaStableTarget{sga6.e03.schol.2-6}\textbf{Scholium 2.6.} To prove conjecture (2.1) in the general case, it will no doubt be necessary to introduce new techniques. One may hope that these will one day help to prove the analogous conjecture in Analytic Geometry, for a proper morphism of analytic spaces over a locally ringed topos in suitably topologized \(\mathbb C\)-algebras, a conjecture that is at present ``proved'' only over a complex analytic space (Grauert's theorem [1]).\SGAAuthorFootnote{1}{(added in April 1971) see an article in preparation by R. Kiehl and J.L. Verdier, in which finiteness theorem 1 is proved by different, simpler techniques and under more general hypotheses.}

\sgaStableTarget{sga6.e03.sec.3}\section*{3. Relative finite Tor-dimension}

\sgaStableTarget{sga6.e03.def.3-1}\textbf{Definition 3.1.} \uline{Let} \(f:(X,\underline{O}_X)\longrightarrow(Y,\underline{O}_Y)\) \uline{be a morphism of ringed topoi}, \(E\in\Ob D^{-}(X)\), \([m,n]\) \uline{an interval of} \(\mathbb Z\). \uline{We say that} \(E\) \uline{is} of flat amplitude relative to \(f\) (or \(Y\), when no confusion is possible) contained in \([m,n]\), \uline{and write}
\[
  \toramp_f(E)\subset[m,n]\qquad(\text{or }\toramp_Y(E)\subset[m,n])
\]
\uline{if} \(E\), \uline{considered as a complex of} \(f^{-1}(\underline{O}_Y)\)-\uline{Modules}, \uline{is of flat amplitude contained} \uline{in} \([m,n]\) (I 5.2). \uline{We say that} \(f\) \uline{is of flat amplitude contained} \uline{in} \([-a,0]\) \((a\geq0)\), \uline{or of} tor-dimension \(\leq a\), \uline{if} \(\underline{O}_X\) \uline{is of flat amplitude} \uline{rel.} \uline{to} \(f\), \uline{contained in} \([-a,0]\). \uline{We say that} \(E\) \uline{is} of finite flat amplitude rel. to \(f\) \uline{if there is an interval} \([m,n]\) \uline{such that} \(\toramp_f(E)\subset[m,n]\).

\sgaStableTarget{sga6.e03.notation.3-2}\textbf{Notation 3.2.} We denote by \(D(f)_{\mathrm{torf}}\), or \(D(X)_{f\text{-}\torf}\), or \(D(X)_{Y\text{-}\torf}\) when no confusion is possible, the full subcategory of \(D^{-}(X)\) formed by the complexes of finite flat amplitude relative to \(f\). By (I 5.3), this is a triangulated subcategory of \(D^b(X)\).

%%%%%%%%%%  scan idx 255  |  folio 242  %%%%%%%%%%
\SGAPrintPage{242}%
\sgaStableTarget{sga6.e03.prop.3-3}\textbf{Proposition 3.3.} \uline{With the data of} (3.1), \uline{suppose that} \(X\) \uline{has enough points} (SGA 4 IV 6.4). \uline{The following conditions are then equivalent}:
\begin{enumerate}[label=(\roman*)]
\item \(\toramp_f(E)\subset[m,n]\);
\item \uline{for every point} \(x\) \uline{of} \(X\), \uline{we have}
\[
  \toramp_y(E_x)\subset[m,n],
\]
\uline{where} \(y\) \uline{is the image point of} \(x\) \uline{under} \(f\), \uline{ringed by} \(\underline{O}_{Y,y}\);
\item \uline{for every object} \(V\) \uline{of} \(Y\) \uline{and every} \(\underline{O}_V\)-\uline{Module} \(M\), \uline{we have}
\[
  H^i\left((E|f^{-1}(V))\otimes_V^{\mathrm L}M\right)=0
  \quad\text{\uline{for} }i\notin[m,n],
\]
\uline{where} \((E|f^{-1}(V))\otimes_V^{\mathrm L}M=(E|f^{-1}(V))\otimes^{\mathrm L}\mathrm Lf^*(M)\);
\item \uline{the same condition as} (iii), \uline{but with} \(M\) \uline{the cokernel of a homomorphism} \(\underline{O}_V^q\longrightarrow\underline{O}_V^p\), \uline{where} \(p,q\in\mathbb N\).
\end{enumerate}

\uline{Proof}. The equivalence (i)\(\Longleftrightarrow\)(ii) follows from (I 5.6.2), and the implications (ii)\(\Longrightarrow\)(iii)\(\Longrightarrow\)(iv) are trivial. Let us prove (iv)\(\Longrightarrow\)(ii). Let \(x\) be a point of \(X\), \(y=f(x)\); we must show that
\begin{equation}\tag{$*$}
  H^i(E_x\otimes_{\underline{O}_{Y,y}}^{\mathrm L}M)=0
\end{equation}
for every \(\underline{O}_{Y,y}\)-module \(M\) and every \(i\notin[m,n]\). Since every \(\underline{O}_{Y,y}\)-module is a filtered inductive limit of modules of finite presentation, we may restrict ourselves to checking the preceding relation for \(M\) of finite presentation. Now, when \(M\) is an \(\underline{O}_{Y,y}\)-module of finite presentation, the same argument as for ordinary ringed spaces (EGA \(0_{\mathrm I}\) 5.3.9) shows that there is a \(y\)-pointed object \(V\) of \(Y\) and an \(\underline{O}_V\)-Module \(P\), the cokernel of a homomorphism \(\underline{O}_V^q\longrightarrow\underline{O}_V^p\) \((p,q\in\mathbb N)\), such that \(P_y\simeq M\). We then have
\[
  \bigl((E|f^{-1}(V))\otimes_V^{\mathrm L}P\bigr)_x
  \simeq
  E_x\otimes_{\underline{O}_{Y,y}}^{\mathrm L}M,
\]
%%%%%%%%%%  scan idx 256  |  folio 243  %%%%%%%%%%
\SGAPrintPage{243}%
by the commutation of tensor product with passage to stalks, and hypothesis (iv) implies (*), q.e.d.

\sgaStableTarget{sga6.e03.rem.3-3-1}\textbf{Remark 3.3.1.} The writer does not know whether the equivalences (i)\(\Longleftrightarrow\)(iii)\(\Longleftrightarrow\)(iv) hold without the hypothesis that \(X\) has enough points.

\sgaStableTarget{sga6.e03.cor.3-4}\textbf{Corollary 3.4.} \uline{Let}
\[
\begin{tikzcd}[column sep=large]
X \arrow[rr,"f"] \arrow[dr,"h"'] & & Y \arrow[dl,"g"]\\
& Z &
\end{tikzcd}
\]
\uline{be a commutative triangle of ringed topoi}, \uline{where we suppose that} \(X\) \uline{has enough points}, \uline{and let} \(E\in\Ob D(X)\), \(F\in\Ob D(Y)\), \([a,b]\), \([c,d]\) \uline{be intervals of} \(\mathbb Z\). \uline{Then we have}:
\[
  \toramp_f(E)\subset[a,b]\quad\text{and}\quad\toramp_g(F)\subset[c,d]
\]
\[
  \Longrightarrow\toramp_h(E\otimes_Y^{\mathrm L}F)\subset[a+c,b+d].
\]

\uline{Proof}. By (3.3), we must show that the relations on the left imply
\[
  H^i\left(\bigl((E\otimes_Y^{\mathrm L}F)|h^{-1}(Z')\bigr)
  \otimes_{Z'}^{\mathrm L}M\right)=0
\]
for every \(i\notin[a+c,b+d]\) and every \(\underline{O}_{Z'}\)-Module \(M\), where \(Z'\) is an object of \(Z\). Now, putting \(Y'=g^{-1}(Z')\), \(X'=f^{-1}(Y')=h^{-1}(Z')\), we have
\[
  \bigl((E\otimes_Y^{\mathrm L}F)|X'\bigr)
  \otimes_{Z'}^{\mathrm L}M
  \simeq
  (E|X')\otimes_{Y'}^{\mathrm L}
  \bigl((F|Y')\otimes_{Z'}^{\mathrm L}M\bigr),
\]
and the assertion follows at once from the hypotheses.

In particular:

\sgaStableTarget{sga6.e03.cor.3-4-1}\textbf{Corollary 3.4.1.} \uline{Under the hypotheses of} (3.4),
\begin{enumerate}[label=\alph*)]
\item \uline{if} \(g\) \uline{is of tor-dimension} \(\leq n\), \uline{we have}
\[
  \toramp_f(E)\subset[a,b]\Longrightarrow\toramp_h(E)\subset[a-n,b];
\]
%%%%%%%%%%  scan idx 257  |  folio 244  %%%%%%%%%%
\SGAPrintPage{244}%
\item \uline{if} \(f\) \uline{is of tor-dimension} \(\leq m\), \uline{we have}
\[
  \toramp_g(F)\subset[c,d]\Longrightarrow\toramp_h(\mathrm Lf^*(F))\subset[c-m,d].
\]
\end{enumerate}

\sgaStableTarget{sga6.e03.rem.3-4-2}\textbf{Remark 3.4.2.} Let \(f:X\to Y\) be a morphism of finite tor-dimension and \(E\in\Ob D(X)\). Then, by (3.4.1 a)), when \(E\) is of finite flat amplitude (in the absolute sense), \(E\) is of finite flat amplitude rel. to \(f\). But the converse is false, even when \(f\) is flat, as the example of a singularity of a scheme over a field shows.

\sgaStableTarget{sga6.e03.prop.3-5}\textbf{Proposition 3.5.} \uline{Let} \(S\) \uline{be a scheme}, \uline{and}, \uline{for} \(i=1,2\), \(f_i:X_i\longrightarrow Y_i\) \uline{an} \(S\)-\uline{morphism}, \(E_i\in\Ob D^{-}(X_i)\), \([a_i,b_i]\) \uline{an interval of} \(\mathbb Z\). \uline{Suppose} \(X_1\) \uline{and} \(X_2\), \uline{as well as} \(Y_1\) \uline{and} \(Y_2\), \uline{are tor-independent over} \(S\) (1.5). \uline{Put} \(a=a_1+a_2\), \(b=b_1+b_2\), \(f=f_1\times_Sf_2:X\to Y\). \uline{Then}:
\begin{enumerate}[label=\alph*)]
\item \uline{If we have} \(\toramp_{f_1}(E_1)\subset[a_1,b_1]\) \uline{and if} \(E_2\) \uline{is acyclic outside} \([a_2,b_2]\), \(E_1\otimes_S^{\mathrm L}E_2\) (\uline{cf.} (1.6)) \uline{is acyclic outside} \([a,b]\).
\item \uline{If we have} \(\toramp_{f_i}(E_i)\subset[a_i,b_i]\) \uline{for} \(i=1,2\), \uline{we have}
\[
  \toramp_f(E_1\otimes_S^{\mathrm L}E_2)\subset[a,b].
\]
\end{enumerate}

\uline{Proof}. By (3.3), it suffices to make the verification ``on the stalks''. Let \(s\) be a point of \(S\), \(x\) a point of \(X\) over \(s\), \(y=f(x)\), \(x_i\) (resp. \(y_i\)) the canonical projection of \(x\) (resp. \(y\)) onto \(X_i\) (resp. \(Y_i\)). Put
\[
  (E_i)_{x_i}=M_i,\quad \underline{O}_{X_i,x_i}=A_i,\quad
  \underline{O}_{Y_i,y_i}=B_i,\quad \underline{O}_{S,s}=C,
\]
\[
  \underline{O}_{X,x}=A=A_1\otimes_C A_2,\qquad
  \underline{O}_{Y,y}=B=B_1\otimes_C B_2.
\]
We have
\begin{equation}\tag{$*$}
  (E_1\otimes_S^{\mathrm L}E_2)_x
  \simeq
  M_1\otimes_{A_1}^{\mathrm L}A\otimes_{A_2}^{\mathrm L}M_2.
\end{equation}
The tor-independence hypotheses imply, on the other hand:
\[
  A\simeq A_1\otimes_{B_1}^{\mathrm L}B\otimes_{B_2}^{\mathrm L}A_2.
\]
Consequently, the right-hand side of (*) is canonically isomorphic, in the
%%%%%%%%%%  scan idx 258  |  folio 245  %%%%%%%%%%
\SGAPrintPage{245}%
derived category \(D(B)\) of \(B\)-modules, to \(M_1\otimes_{B_1}^{\mathrm L}B\otimes_{B_2}^{\mathrm L}M_2\). The relations \(\toramp_{f_i}(E_i)\subset[a_i,b_i]\) \((i=1,2)\) imply \(\toramp_{B_i}(M_i)\subset[a_i,b_i]\), hence (I 5.6) \(\toramp_B(M_1\otimes_{B_1}^{\mathrm L}B\otimes_{B_2}^{\mathrm L}M_2)\subset[a,b]\), giving assertion b). On the other hand, when we have \(\toramp_{f_1}(E_1)\subset[a_1,b_1]\), \(M_1\) is isomorphic to a complex of flat \(B_1\)-modules, zero outside \([a_1,b_1]\), and consequently when \(E_2\) (hence \(M_2\)) is acyclic outside \([a_2,b_2]\), \(M_1\otimes_{B_1}^{\mathrm L}B\otimes_{B_2}^{\mathrm L}M_2\) is acyclic outside \([a,b]\), giving assertion a), which completes the proof.

\sgaStableTarget{sga6.e03.cor.3-5-1}\textbf{Corollary 3.5.1.} \uline{Under the hypotheses of} (3.5), \uline{if} \(f_1\) \uline{and} \(f_2\) \uline{are of finite tor-dimension} (3.1), \uline{the same holds for} \(f\).

\sgaStableTarget{sga6.e03.cor.3-5-2}\textbf{Corollary 3.5.2.} \uline{Let}
\[
\begin{tikzcd}
X \arrow[d,"f"'] & X' \arrow[l,"g"'] \arrow[d,"f'"]\\
Y & Y' \arrow[l,"h"']
\end{tikzcd}
\]
\uline{be a Cartesian square of schemes}, \uline{where} \(X\) \uline{and} \(Y'\) \uline{are supposed tor-independent over} \(Y\), \uline{and let} \(E\in\Ob D(X)\). \uline{Then we have}
\[
  \toramp_f(E)\subset[a,b]\Longrightarrow
  \toramp_{f'}(\mathrm Lg^*(E))\subset[a,b],
\]
\uline{and consequently the functor} \(\mathrm Lg^*:D^{-}(X)\to D^{-}(X')\) \uline{induces a functor}
\[
  \mathrm Lg^*:D(f)_{\mathrm{torf}}\longrightarrow D(f')_{\mathrm{torf}} .
\]


% Original English continuation definitions, retained in their original order.
\providecommand{\Ob}{\operatorname{Ob}}
\providecommand{\D}{\mathrm D}
\providecommand{\Ho}{\mathrm H}
\providecommand{\RHom}{\mathrm R\!\operatorname{Hom}}
\providecommand{\parfamp}{\operatorname{parf.amp}}
\providecommand{\tamp}{\operatorname{tor.amp}}
\providecommand{\Qcoh}{\operatorname{Qcoh}}
\providecommand{\Mod}{\operatorname{Mod}}
\providecommand{\Coh}{\operatorname{Coh}}
\providecommand{\Lib}{\operatorname{Lib}}
\providecommand{\Loclib}{\operatorname{Loclib}}
\providecommand{\qcoh}{\mathrm{qcoh}}
\providecommand{\coh}{\mathrm{coh}}
\providecommand{\parf}{\mathrm{parf}}
\providecommand{\torf}{\mathrm{torf}}
\providecommand{\id}{\operatorname{id}}
\providecommand{\Coker}{\operatorname{Coker}}
\providecommand{\rank}{\operatorname{rank}}
\providecommand{\rg}{\operatorname{rg}}
\providecommand{\Hom}{\operatorname{Hom}}
\providecommand{\Ext}{\operatorname{Ext}}
\providecommand{\Spec}{\operatorname{Spec}}
\providecommand{\Cone}{\operatorname{Cone}}
\providecommand{\pt}{\operatorname{Pt}}
\providecommand{\ind}{\operatorname{ind}}
\hbadness=10000
\hfuzz=20pt
\emergencystretch=3em

% Reviewed French continuation definitions, retained verbatim.
% Additions for Exposé III, pages 253--280
\renewcommand{\Ob}{\operatorname{ob}}
\providecommand{\D}{\mathrm D}
\providecommand{\Ho}{\mathrm H}
\providecommand{\RHom}{\mathrm R\!\operatorname{Hom}}
\providecommand{\parfamp}{\operatorname{parf.amp}}
\providecommand{\tamp}{\operatorname{tor.amp}}
\providecommand{\Qcoh}{\operatorname{Qcoh}}
\providecommand{\Mod}{\operatorname{Mod}}
\providecommand{\Coh}{\operatorname{Coh}}
\providecommand{\Lib}{\operatorname{Lib}}
\providecommand{\Loclib}{\operatorname{Loclib}}
\providecommand{\qcoh}{\mathrm{qcoh}}
\providecommand{\coh}{\mathrm{coh}}
\providecommand{\parf}{\mathrm{parf}}
\providecommand{\torf}{\mathrm{torf}}
\renewcommand{\id}{\operatorname{id}}
\providecommand{\uparf}{\underline{\mathrm{parf}}\rule[-1ex]{0pt}{0pt}}
\providecommand{\Coker}{\operatorname{Coker}}
\providecommand{\rank}{\operatorname{rank}}
\providecommand{\rg}{\operatorname{rg}}
\providecommand{\Hom}{\operatorname{Hom}}
\providecommand{\Ext}{\operatorname{Ext}}
\providecommand{\Spec}{\operatorname{Spec}}
\providecommand{\Cone}{\operatorname{Cone}}
\providecommand{\pt}{\operatorname{Pt}}
\providecommand{\ind}{\operatorname{ind}}
\hbadness=10000
\hfuzz=20pt
\emergencystretch=3em

%%%%%%%%%%  scan idx 259  |  folio 246  %%%%%%%%%%
\SGAPrintPage{246}%
\sgaStableTarget{sga6.e03.prop.3-6}\textbf{Proposition 3.6.} \uline{Let}
\[
\begin{tikzcd}
X \arrow[r,hook,"i"] \arrow[d,"f"'] & X' \arrow[dl,"f'"]\\
Y &
\end{tikzcd}
\]
\uline{be a commutative triangle of schemes}, \uline{where} \(i\) \uline{is a closed immersion and} \(f'\) \uline{a smooth morphism}, \uline{and let} \(E\in\Ob D(X)\). \uline{For} \(E\) \uline{to be locally of finite flat amplitude rel.} \uline{to} \(f\), \uline{it is necessary and sufficient that} \(i_*(E)\) \uline{be locally of finite flat amplitude} (\uline{in the absolute sense}). \uline{More precisely}, \uline{if} \(f'\) \uline{is of relative dimension} \(\leq d\), \uline{we have}
\begin{align*}
\text{(i)}\quad &\tamp(i_*(E))\subset[a,b] \Longrightarrow \tamp_f(E)\subset[a,b],\\
\text{(ii)}\quad &\tamp_f(E)\subset[a,b] \Longrightarrow \tamp(i_*(E))\subset[a-d,b].
\end{align*}

%%%%%%%%%%  scan idx 260  |  folio 247  %%%%%%%%%%
\SGAPrintPage{247}%
\uline{Proof}. For an interval \([m,n]\) of \(\mathbb Z\), we evidently have
\[
 \tamp_f(E)\subset[m,n]
 \Longleftrightarrow
 \tamp_{f'}(i_*(E))\subset[m,n],
\]
so that we may suppose \(X=X'\), \(i=\Id_X\), \(f=f'\). Relation (i) follows trivially from (3.4.1 a)). Let us prove (ii). We have a Cartesian square
\[
\begin{tikzcd}
X \arrow[d,"f"'] & X\times_Y X \arrow[l,"\operatorname{pr}_2"'] \arrow[d,"\operatorname{pr}_1"]\\
Y & X \arrow[l,"f"']
\end{tikzcd},
\]
and the diagonal morphism \(s:X\longrightarrow X\times_YX\) is a section of \(\operatorname{pr}_1\) and \(\operatorname{pr}_2\). Suppose \(f\) is smooth of relative dimension \(d\). Then (VII 1.10) \(s\) is a regular immersion of codimension \(\leq d\), hence is of tor-dimension \(\leq d\) (as is seen using the Koszul complex). By (3.5.2), we have
\[
 \tamp_f(E)\subset[a,b]
 \Longrightarrow
 \tamp_{\operatorname{pr}_1}(\operatorname{pr}_2^*(E))\subset[a,b],
\]
and by (3.4.1 b)) we have
\[
 \tamp_{\operatorname{pr}_1}(\operatorname{pr}_2^*(E))\subset[a,b]
 \Longrightarrow
 \tamp(\mathrm Ls^*\operatorname{pr}_2^*(E))\subset[a-d,b],
\]
giving implication (ii), since \(\mathrm Ls^*\operatorname{pr}_2^*\simeq\Id\).

\sgaStableTarget{sga6.e03.rem.3-6-1}\textbf{Remark 3.6.1.} Statement (3.6) would be false with only the assumption that \(f'\) is flat; see (3.4.2).

\sgaStableTarget{sga6.e03.prop.3-7}\textbf{Proposition 3.7.} (``\uline{Projection formula}'')\SGAAuthorFootnote{1}{cf. (H II 5.6).}. \uline{Let} \(S\) \uline{be a quasi-compact scheme}, \(f:X\to Y\) \uline{a quasi-compact and quasi-separated} \(S\)-\uline{morphism of quasi-compact schemes over} \(S\), \(E\in\Ob D^b(X)_{\qcoh}\), \(M\in\Ob D^b(S)_{\qcoh}\)\SGAAuthorFootnote{2}{The subscript ``qcoh'' indicates the full subcategory formed by complexes with quasi-coherent cohomology.}. \uline{If} \(M\) \uline{is} \uline{of finite flat amplitude}, \uline{or if} \(E\) \uline{is of finite flat amplitude rel.} \uline{to} \(S\), \uline{there}
%%%%%%%%%%  scan idx 261  |  folio 248  %%%%%%%%%%
\SGAPrintPage{248}%
\uline{is a canonical functorial isomorphism in} \(D(Y)_{\qcoh}\)
\[
 Rf_*(E)\otimes_S^{\mathrm L}M
 \xrightarrow{\sim}
 Rf_*(E\otimes_S^{\mathrm L}M)\qquad(\text{cf. }(1.6)).
\]

\uline{Proof}. First note that there is (EGA III 1.4.12 and IV 1.7.21) an integer \(N\) such that \(R^if_*(F)=0\) for every \(i>N\) and every quasi-coherent \(\underline{O}_X\)-Module \(F\). Consequently the functor \(Rf_*\) sends \(D^b(X)_{\qcoh}\) into \(D^b(Y)\) (and even \(D^b(Y)_{\qcoh}\) (EGA III 1.4.10 and IV 1.7.21)). Suppose \(E\) (resp. \(M\)) is of finite flat amplitude rel. to \(S\) (resp. finite flat amplitude). The adjunction arrow \(Lf^*Rf_*(E)\longrightarrow E\) defines, by tensoring with \(M\), a canonical arrow of \(D^-(X)\)
\[
 Lf^*(Rf_*(E)\otimes_S^{\mathrm L}M)
 \longrightarrow E\otimes_S^{\mathrm L}M,
\]
where \(E\otimes_S^{\mathrm L}M\in\Ob D^b(X)\). By adjunction we deduce a canonical arrow of \(D^-(Y)\)
\sgaStableTarget{sga6.e03.eq.unnamed}\[
 Rf_*(E)\otimes_S^{\mathrm L}M
 \longrightarrow Rf_*(E\otimes_S^{\mathrm L}M).\tag{*}
\]
Its target and source moreover have quasi-coherent cohomology, since \(Rf_*\) sends \(D^b(X)_{\qcoh}\) into \(D^b(Y)_{\qcoh}\) (EGA III 1.4.10 and IV 1.7.21), and \(\otimes_S^{\mathrm L}\) sends \(D^-(X)_{\qcoh}\times D^-(S)_{\qcoh}\) into \(D^-(X)_{\qcoh}\) (likewise with \(Y\) in place of \(X\)). Let us prove that \((*)\) is an isomorphism. The question being local on \(Y\), we may suppose \(S\) and \(Y\) affine. By dévissage, in the first case we immediately reduce to the case where \(M\) is an ordinary quasi-coherent \(\underline{O}_S\)-Module. Noting that the functors \(Rf_*(E)\otimes_S^{\mathrm L}.\) and \(Rf_*(E\otimes_S^{\mathrm L}.)\) are ``way-out to the left'' on the category \(D^b(S)_{\qcoh}\) ([H] I 7), we then reduce, by taking a left resolution of \(M\) by free Modules and truncating, to the case where \(M\) is free. Since the \(R^if_*\) commute with inductive limits (\(f\) being quasi-compact and quasi-separated), we may even suppose \(M\) free of finite type, and then \(M=\underline{O}_S\). The arrow \((*)\) is then tautologically an isomorphism. In the second case, by dévissage we reduce to \(M=P\), where \(P\) is a flat \(A\)-module
%%%%%%%%%%  scan idx 262  |  folio 249  %%%%%%%%%%
\SGAPrintPage{249}%
\((S=\Spec(A))\). Putting \(A'=D_A(P)\) (EGA \(0_{\mathrm{IV}}\) 18.2.3), \(S'=\Spec(A')\), everything reduces to proving that \(Rf_*(E)\) commutes with the flat base change \(S'\to S\), i.e. that the canonical arrow
\[
 q^*Rf_*(E)\longrightarrow Rf'_*p^*(E)
\]
is an isomorphism (we have put \(f_{(S')}=f'\) and denoted by \(p\) (resp. \(q\)) the canonical projection \(X_{(S')}\longrightarrow X\) (resp. \(Y_{(S')}\longrightarrow Y\))). But this follows from (EGA III 1.4.15 and IV 1.7.21). This completes the proof of (3.7).

\sgaStableTarget{sga6.e03.cor.3-7-1}\textbf{Corollary 3.7.1.} \uline{Let} \(S\) \uline{be a scheme}, \(f:X\to Y\) \uline{a quasi-compact and quasi-separated} \(S\)-\uline{morphism of quasi-compact schemes over} \(S\), \(E\in\Ob D^b(X)_{\qcoh}\). \uline{If} \(E\) \uline{is locally of finite flat amplitude rel.} \uline{to} \(S\), \uline{the same holds for} \(Rf_*(E)\). \uline{More precisely}, \uline{if} \(S\) \uline{is quasi-compact}, \uline{and if} \(N\) \uline{is an integer such that} \(R^if_*(F)=0\) \uline{for every} \(i>N\) \uline{and every quasi-coherent} \(\underline{O}_X\)-\uline{Module} \(F\) (EGA III 1.4.12 \uline{and} IV 1.7.21), \uline{then we have}
\[
 \tamp_S(E)\subset[a,b]
 \Longrightarrow
 \tamp_S(Rf_*(E))
 \subset[a,b+N].
\]

\uline{Proof}. We may evidently restrict ourselves to proving the last assertion. By (3.3 (i)\(\Longleftrightarrow\)(iv)), and taking into account that \(Rf_*\) ``commutes with restriction to an open subset of \(S\)'', we must prove that the hypothesis \(\tamp_S(E)\subset[a,b]\) implies
\[
 H^i(Rf_*(E)\otimes_S^{\mathrm L}M)=0
\]
for every \(i\notin[a,b+N]\) and every quasi-coherent \(\underline{O}_S\)-Module \(M\). Now the preceding hypothesis means that \(H^i(E\otimes_S^{\mathrm L}M)=0\) for every \(i\notin[a,b]\) and every quasi-coherent \(\underline{O}_S\)-Module \(M\), and (3.7) gives the result.

%%%%%%%%%%  scan idx 263  |  folio 250  %%%%%%%%%%
\SGAPrintPage{250}%
\sgaStableTarget{sga6.e03.cor.3-7-2}\textbf{Corollary 3.7.2.} \uline{Let} \(f:X\to Y\) \uline{be a quasi-compact}, \uline{quasi-separated morphism}, \uline{locally of finite tor-dimension} (3.1), \uline{and let} \(E\in\Ob D^b(X)_{\qcoh}\). \uline{If} \(E\) \uline{is locally of finite flat amplitude} (\uline{in the absolute sense}), \uline{the same holds for} \(Rf_*(E)\).

\uline{Proof}. Combine (3.4.1 a)) and (3.7.1).

\sgaStableTarget{sga6.e03.sec.4}\section*{4. Relative perfection}

\sgaStableTarget{sga6.e03.def.4-1}\textbf{Definition 4.1.} \uline{Let} \(f:X\to Y\) \uline{be a locally finite type morphism of schemes}, \(E\in\Ob D(X)\), \([a,b]\) \uline{an interval of} \(\mathbb Z\). \uline{We say that} \(E\) \uline{is} of perfect amplitude \uline{relative to} \(f\) (\uline{or} \(Y\) \uline{when there is no ambiguity}), contained in \([a,b]\), \uline{and write}
\[
 \parfamp_f(E)\subset[a,b]\qquad(\text{or }\parfamp_Y(E)\subset[a,b])
\]
\uline{if}, \uline{relative to} \(f\), \(E\) \uline{is pseudo-coherent} (1.2) \uline{and of flat amplitude contained in} \([a,b]\) (3.1). \uline{We say that} \(f\) is of perfect amplitude contained in \([-n,0]\), \uline{or of} perfect dimension \(\leq n\) \((n\geq0)\), \uline{if} \(\underline{O}_X\) \uline{is of perfect amplitude contained in} \([-n,0]\). \uline{We say that} \(E\) \uline{is} of finite perfect amplitude relative to \(f\) \uline{if there is an interval} \([a,b]\) \uline{such that} \(\parfamp_f(E)\subset[a,b]\). \uline{We say that} \(E\) \uline{is} perfect rel. to \(f\) \uline{if} \(E\) \uline{is locally of finite perfect amplitude rel.} \uline{to} \(f\). \uline{We say that} \(f\) is perfect \uline{if} \(f\) \uline{is} \uline{locally of finite perfect amplitude}.

\sgaStableTarget{sga6.e03.ex.4-1-1}\textbf{Example 4.1.1.} A regular immersion of codimension \(d\) (VII 1.4) is of perfect dimension \(\leq d\) (as is seen from the Koszul complex). A locally complete-intersection morphism (i.e. locally a composite
%%%%%%%%%%  scan idx 264  |  folio 251  %%%%%%%%%%
\SGAPrintPage{251}%
of a regular immersion and a smooth morphism) is a perfect morphism.

\sgaStableTarget{sga6.e03.notation.4-2}\textbf{Notation 4.2.} Let \(f:X\to Y\) be a locally finite type morphism of schemes. We denote by \(D(f)_{\parf}\), or \(D(X)_{f\text{-}\parf}\), or \(D(X)_{Y\text{-}\parf}\) when no confusion is to be feared (resp. \(D(f)_{\uparf}\), \(D(X)_{f\text{-}\uparf}\), \(D(X)_{Y\text{-}\uparf}\)) the full subcategory of \(D(X)\) formed by the complexes perfect (resp. of finite perfect amplitude) rel. to \(f\). By (I 5.3) and (1.4), \(D(f)_{\parf}\) and \(D(f)_{\uparf}\) are triangulated subcategories of \(D(X)\). When \(X\) is quasi-compact, we have \(D(f)_{\parf}=D(f)_{\uparf}\).

\sgaStableTarget{sga6.e03.prop.4-3}\textbf{Proposition 4.3} (cf. I 5.8). \uline{Let} \(f:X\to Y\) \uline{be a locally finite type morphism of schemes}, \(E\in\Ob D(X)\), \([a,b]\) \uline{an interval of} \(\mathbb Z\), \(x\) \uline{a point of} \(X\), \(y=f(x)\). \uline{The following conditions are equivalent}:
\begin{enumerate}[label=(\roman*)]
\item \(\parfamp_f(E)\subset[a,b]\) \uline{in a neighborhood of} \(x\):
\item \(E\) \uline{is}, \uline{in a neighborhood of} \(x\), \((a-1)\)-\uline{pseudo-coherent rel.} \uline{to} \(f\) \uline{and acyclic outside} \([a,b]\), \uline{and we have} \(\tamp_{\underline{O}_{Y,y}}(E_x)\subset[a,b]\).
\end{enumerate}

\uline{Proof}. The implication (i) \(\Longrightarrow\) (ii) is trivial; let us prove (ii) \(\Longrightarrow\) (i). The question being local in a neighborhood of \(x\), we immediately reduce, by the definitions, to the case where \(f\) is smooth. Then \(E\) is \((a-1)\)-pseudo-coherent in a neighborhood of \(x\); hence, after localizing, we may suppose \(E\) strictly \((a-1)\)-pseudo-coherent (I 2.1) and zero in degree \(>b\), and after truncating (i.e. replacing \(E\) by the complex \((0\to E^a/B^a\to E^{a+1}\to\ldots E^b\to 0)\)), we may suppose that \(E^i\) is free of finite type for \(i>a\), \(E^a\) of finite presentation, and \(E^i=0\) for \(i\notin[a,b]\). By (I 5.1.1), \(E^a_x\) is flat rel. to \(\underline{O}_{Y,y}\), hence (EGA IV 11.3.1) \(E^a\) is \(f\)-flat in a neighborhood of \(x\). Write \(E^a=F\) for brevity; we must show that \(F\) is pseudo-coherent (in a neighborhood of \(x\)). After embedding \(X\) in a standard affine space over \(Y\), we may
%%%%%%%%%%  scan idx 265  |  folio 252  %%%%%%%%%%
\SGAPrintPage{252}%
suppose that \(X=Y[t_1,\dots,t_r]\), with \(Y\) affine and \(f\) the canonical projection. Then \(Y\) is a projective limit of noetherian affine schemes \(Y_\alpha\). Denote by \(f_\alpha:X_\alpha=Y_\alpha[t_1,\dots,t_r]\to Y_\alpha\) the canonical projection. There is (EGA IV 8.5.2 and 11.2.6) an index \(\alpha\) and an \(\underline{O}_{X_\alpha}\)-Module \(F_\alpha\) of finite presentation and \(f_\alpha\)-flat such that \(g_\alpha^*(F_\alpha)\simeq F\), where \(g_\alpha:X\to X_\alpha\) is the canonical morphism. Since \(X_\alpha\) is noetherian, \(F_\alpha\) is coherent, hence pseudo-coherent rel. to \(f_\alpha\). Since, on the other hand, \(F_\alpha\) is \(f_\alpha\)-flat and \(f_\alpha\) is flat, we have \(g_\alpha^*(F_\alpha)\simeq F_\alpha\otimes^{\mathrm L}_{Y_\alpha}\underline{O}_Y\), and it follows from (1.10 a)) that \(F\) is pseudo-coherent rel. to \(Y\), hence pseudo-coherent, q.e.d.

\sgaStableTarget{sga6.e03.cor.4-3-1}\textbf{Corollary 4.3.1.} \uline{Let} \(f:X\to Y\) \uline{be a locally finite type morphism}, \uline{and} \(E\) \uline{an} \(\underline{O}_X\)-\uline{Module}. \uline{For} \(\parfamp_f(E)\subset[0,0]\) \uline{to hold, it is necessary and sufficient that} \(E\) \uline{be}, \uline{rel.} \uline{to} \(f\), \uline{flat and of finite presentation} (1.2.1 c)). \uline{In particular}, \uline{for} \(f\) \uline{to have perfect dimension zero} (4.1), \uline{it is necessary and sufficient that} \(f\) \uline{be flat and locally of finite presentation}.

\uline{Proof}. Apply (4.3 (ii) \(\Longrightarrow\) (i)) with \(a=b=0\).

\sgaStableTarget{sga6.e03.cor.4-3-2}\textbf{Corollary 4.3.2.} \uline{Let} \(f:X\to Y\) \uline{be a locally finite type morphism}, \uline{and} \(E\in\Ob D(X)\) \uline{a complex of} \(\underline{O}_X\)-\uline{Modules flat and of finite presentation} \uline{rel.} \uline{to} \(f\) (1.2.1 c)), \uline{and} \([a,b]\) \uline{an interval of} \(\mathbb Z\). \uline{If} \(E^i=0\) \uline{for} \(i\notin[a,b]\), \uline{we have} \(\parfamp_f(E)\subset[a,b]\). \uline{If} \(E\) \uline{is bounded above in degree}, \(E\) \uline{is pseudo-coherent rel.} \uline{to} \(f\).

\uline{Proof}. A consequence of (4.3.1) and (I 4.15 and 2.11).

\sgaStableTarget{sga6.e03.prop.4-4}\textbf{Proposition 4.4.} \uline{Let there be a commutative triangle of schemes}
\[
\begin{tikzcd}
X \arrow[r,hook,"i"] \arrow[d,"f"'] & Y \arrow[dl,"g"]\\
S &
\end{tikzcd},
\]
%%%%%%%%%%  scan idx 266  |  folio 253  %%%%%%%%%%
\SGAPrintPage{253}%
\uline{where} \(i\) \uline{is a closed immersion and} \(g\) \uline{a smooth morphism}, \uline{and let} \(E\in\Ob D(X)\). \uline{For} \(E\) \uline{to be perfect rel.} \uline{to} \(f\), \uline{it is necessary and sufficient that} \(i_*(E)\) \uline{be perfect} (\uline{in the absolute sense} (I 4.8)). \uline{More precisely}, \uline{suppose} \(g\) \uline{is of relative dimension} \(\leq d\), \uline{and let} \([a,b]\) \uline{be an interval of} \(\mathbb Z\). \uline{Then we have}:
\begin{enumerate}[label=(\roman*)]
\item \(\parfamp(i_*(E))\subset[a,b]\Longrightarrow\parfamp_f(E)\subset[a,b]\);
\item \(\parfamp_f(E)\subset[a,b]\Longrightarrow\parfamp(i_*(E))\subset[a-d,b]\).
\end{enumerate}

\uline{Proof}. An immediate consequence of (3.6).

\sgaStableTarget{sga6.e03.prop.4-5}\textbf{Proposition 4.5.} \uline{Let}
\[
\begin{tikzcd}
X \arrow[rr,"f"] \arrow[dr,"h"'] && Y \arrow[dl,"g"]\\
& S &
\end{tikzcd}
\]
\uline{be a commutative triangle of locally finite type morphisms of schemes}, \uline{and let} \(E\in\Ob D^-(X)\), \(F\in\Ob D^-(Y)\), \([a,b]\), \([c,d]\) \uline{be intervals of} \(\mathbb Z\). \uline{Then we have}:
\[
 \parfamp_f(E)\subset[a,b]\text{ and }\parfamp_g(F)\subset[c,d]
 \Longrightarrow
 \parfamp_h(E\otimes_Y^{\mathrm L}F)
 \subset[a+c,b+d].
\]
\uline{Consequently}, \uline{if} \(E\) \uline{is perfect rel.} \uline{to} \(f\) \uline{and} \(F\) \uline{perfect rel.} \uline{to} \(g\), \(E\otimes_Y^{\mathrm L}F\) \uline{is perfect rel.} \uline{to} \(h\).

\uline{Proof}. Combine (1.11) and (3.4).

\sgaStableTarget{sga6.e03.cor.4-5-1}\textbf{Corollary 4.5.1.} \uline{Under the hypotheses of} (4.5):
\begin{enumerate}[label=\alph*)]
\item \uline{if} \(g\) \uline{has perfect dimension} \(\leq n\), \uline{we have}
\[
 \parfamp_f(E)
 \subset[a,b]
 \Longrightarrow
 \parfamp_h(E)
 \subset[a-n,b];
\]
\uline{consequently}, \uline{if} \(g\) \uline{is perfect}, \(E\) \uline{perfect rel.} \uline{to} \(f\) \uline{implies} \(E\) \uline{perfect rel.} \uline{to} \(h\);
\item \uline{if} \(f\) \uline{has perfect dimension} \(\leq m\), \uline{we have}
\[
 \parfamp_g(F)
 \subset[c,d]
 \Longrightarrow
 \parfamp_h(\mathrm Lf^*(F))
 \subset[c-m,d];
\]
%%%%%%%%%%  scan idx 267  |  folio 254  %%%%%%%%%%
\SGAPrintPage{254}%
\uline{consequently}, \uline{if} \(f\) \uline{is perfect}, \(F\) \uline{perfect rel.} \uline{to} \(g\) \uline{implies} \(\mathrm Lf^*(F)\) \uline{perfect rel.} \uline{to} \(h\).
\end{enumerate}

\sgaStableTarget{sga6.e03.rem.4-5-2}\textbf{Remark 4.5.2.} With the preceding notation, when \(g\) is perfect, \(E\) perfect rel. to \(h\) does not necessarily imply \(E\) perfect rel. to \(f\), as the ever-recurring example of a singular point of a scheme over a field shows. We have the following criterion:

\sgaStableTarget{sga6.e03.prop.4-6}\textbf{Proposition 4.6} (\uline{criterion of perfection by fibers})\SGAAuthorFootnote{1}{Analogous to the ``criterion of flatness by fibers'' (EGA IV 11.3.10).}. \uline{Let}
\[
\begin{tikzcd}
X \arrow[rr,"f"] \arrow[dr,"h"'] && Y \arrow[dl,"g"]\\
& S &
\end{tikzcd}
\]
\uline{be a commutative triangle of morphisms of schemes}. \uline{We suppose} \(f\), \(h\) \uline{locally of finite type}, \uline{and} \(g\) \uline{locally of finite presentation and} flat. \uline{Let} \(x\in X\), \(y=f(x)\), \(s=h(x)=g(y)\). \uline{Finally}, \uline{let} \(E\in\Ob D^-(X)\), \uline{and let} \([a,b]\) \uline{be an interval of} \(\mathbb Z\). \uline{We denote by} \(f_s:X_s\to Y_s\) \uline{the fiber of} \(f\) \uline{at} \(s\), \uline{i.e.} \uline{the morphism obtained from} \(f\) \uline{by the base change} \(\Spec(k(s))\to S\). \uline{Then the following conditions are equivalent}:
\begin{enumerate}[label=(\roman*)]
\item \(\parfamp_h(E)\subset[a,b]\) \uline{in a neighborhood of} \(x\), \uline{and}
\[
 \tamp_{\underline{O}_{Y_s,y}}(E_x\otimes_{\underline{O}_{Y,y}}^{\mathrm L}\underline{O}_{Y_s,y})
 \subset[a,b];
\]
\item \(\parfamp_f(E)\subset[a,b]\) \uline{in a neighborhood of} \(x\).
\end{enumerate}

\uline{Proof}. (ii) \(\Longrightarrow\) (i). We already know (4.3.1) that \(g\) has perfect dimension zero. Thus (4.5.1 a)), condition (ii) implies \(\parfamp_h(E)\subset[a,b]\) in a neighborhood of \(x\). It also implies (by 4.3), \(\tamp_{\underline{O}_{Y,y}}(E_x)\subset[a,b]\), hence (I 5.1 (i)), \(\tamp_{\underline{O}_{Y_s,y}}(E_x\otimes_{\underline{O}_{Y,y}}^{\mathrm L}\underline{O}_{Y_s,y})\subset[a,b]\), which
%%%%%%%%%%  scan idx 268  |  folio 255  %%%%%%%%%%
\SGAPrintPage{255}%
completes the proof of (ii) \(\Longrightarrow\) (i).

(i) \(\Longrightarrow\) (ii). The question being local in a neighborhood of \(x\), we may suppose that \(f\) factors as \(X\overset{i}{\hookrightarrow}X'\overset{f'}{\longrightarrow}Y\), where \(i\) is a closed immersion and \(f'\) a smooth morphism. It is immediate that conditions (i) and (ii) are unchanged when we replace \(E\) by \(i_*(E)\), \(f\) by \(f'\) (and \(h\) by \(gf'\)), so that we may suppose \(f\) smooth. Since \(g\) is perfect (4.3.1), \(h\) is perfect, hence (1.12) \(E\) is pseudo-coherent. After localizing, by (II 2.2.10 b)) and (I 2.10 a)), we may suppose that \(E\) is strictly pseudo-coherent (I 2.1) and zero in degree \(>b\), and that the truncation arrow
\begin{equation}
 E\longrightarrow(0\longrightarrow E^a/B^a\longrightarrow E^{a+1}\longrightarrow\ldots\longrightarrow E^b\longrightarrow0)\tag*{$t)$}
\end{equation}
is a quasi-isomorphism (hypothesis (i) indeed implies that \(E\) is, in a neighborhood of \(x\), acyclic outside \([a,b]\)). We must prove that \(E^a/B^a\) is \(f\)-flat in a neighborhood of \(x\). Denote by \(i:X_s\longrightarrow X\) the canonical arrow. We have:
\begin{equation}
 E_x\otimes_{\underline{O}_{Y,y}}^{\mathrm L}\underline{O}_{Y_s,y}\simeq i^*(E)_x.\notag
\end{equation}
Since \(h\) is flat, hypothesis (i) implies, by (3.5.2), \(\tamp_{h_s}(i^*(E))\subset[a,b]\) in a neighborhood of \(x\). Thus, after localizing, we may suppose that the truncation arrow
\begin{equation}
 i^*(E)\longrightarrow(0\longrightarrow i^*(E^a)/B^a(i^*(E))\longrightarrow\ldots\longrightarrow i^*(E^b)\longrightarrow0)\tag*{$i^*(t)$}
\end{equation}
is a quasi-isomorphism. The hypothesis \(\tamp_{\underline{O}_{Y_s,y}}i^*(E)_x\subset[a,b]\) then implies, by (I 5.1.1), that \(i^*(E^a/B^a)_x\) is a flat \(\underline{O}_{Y_s,y}\)-module. On the other hand, since the arrow \((t)\) is a quasi-isomorphism, the hypothesis \(\tamp_h(E)\subset[a,b]\) implies, again by (I 5.1.1),
%%%%%%%%%%  scan idx 269  |  folio 256  %%%%%%%%%%
\SGAPrintPage{256}%
that \(E^a/B^a\) is \(h\)-flat. By the criterion of flatness by fibers (EGA IV 11.3.10), we conclude that \(E^a/B^a\) is \(f\)-flat in a neighborhood of \(x\), q.e.d.

\sgaStableTarget{sga6.e03.cor.4-6-1}\textbf{Corollary 4.6.1.} \uline{Let} \(f:X\to S\) \uline{be a flat morphism locally of finite presentation}, \(x\in X\), \(s=f(x)\), \(E\in\Ob D^-(X)\), \([a,b]\) \uline{an interval of} \(\mathbb Z\). \uline{Denote by} \(i_s:X_s=X\times_S\Spec(k(s))\to X\) \uline{the canonical projection}, \uline{and put} \(E_{X_s}=\mathrm Li_s^*(E)\). \uline{Then we have}
\[
\begin{gathered}
 \parfamp_f(E)\subset[a,b]\text{ \uline{in a neighborhood of} }x\text{ \uline{and} }\tamp(E_{X_s})_x\subset[a,b]\\
 \Longleftrightarrow \parfamp(E)\subset[a,b]\text{ \uline{in a neighborhood of} }x.
\end{gathered}
\]
\uline{Consequently}, \(E\) \uline{is perfect if and only if} \(E\) \uline{is perfect rel.} \uline{to} \(f\) \uline{and} \(E_{X_s}\) \uline{is perfect for every} \(s\in S\).

\uline{Proof}. Apply (4.6) with \(X=Y\), \(f=\Id_X\).

\sgaStableTarget{sga6.e03.rem.4-6-2}\textbf{Remark 4.6.2.} From (4.6.1) and (I 5.10) it follows that, when \(f\) is locally of finite presentation, flat and has regular schemes as fibers, \(E\) perfect rel. to \(f\) is equivalent to \(E\) perfect (this equivalence had already been obtained (4.4) for \(f\) smooth).

\sgaStableTarget{sga6.e03.prop.4-7}\textbf{Proposition 4.7.} \uline{Let} \(S\) \uline{be a scheme}, \uline{and for} \(i=1,2\), \(f_i:X_i\to Y_i\) \uline{an} \(S\)-\uline{morphism locally of finite type}, \(E_i\in\Ob D(X_i)\), \([a_i,b_i]\) \uline{an interval of} \(\mathbb Z\). \uline{Suppose} \(X_1\) \uline{and} \(X_2\), \uline{as well as} \(Y_1\) \uline{and} \(Y_2\), \uline{are tor-independent over} \(S\) (1.5). \uline{Put} \(a=a_1+a_2\), \(b=b_1+b_2\), \(f=f_1\times_Sf_2\). \uline{Then we have}:
\begin{enumerate}[label=\alph*)]
\item \uline{If} \(\parfamp_{f_1}(E_1)\subset[a_1,b_1]\), \uline{and if} \(E_2\) \uline{is pseudo-coherent rel.} \uline{to} \(f_2\) \uline{and acyclic in degree outside} \([a_2,b_2]\), \(E_1\otimes_S^{\mathrm L}E_2\) (\uline{cf.} (1.6)) \uline{is pseudo-coherent rel.} \uline{to} \(f\) \uline{and acyclic in degree outside} \([a,b]\).
%%%%%%%%%%  scan idx 270  |  folio 257  %%%%%%%%%%
\SGAPrintPage{257}%
\item \(\parfamp_{f_i}(E_i)\subset[a_i,b_i]\) \uline{for} \(i=1,2\)
\[
 \Longrightarrow\parfamp_f(E_1\otimes_S^{\mathrm L}E_2)\subset[a,b].
\]
\end{enumerate}

\uline{Proof}. Combine (1.9) and (3.5).

\sgaStableTarget{sga6.e03.cor.4-7-1}\textbf{Corollary 4.7.1.} \uline{Under the hypotheses of} (4.7), \uline{if}, \uline{for} \(i=1,2\), \(f_i\) \uline{has perfect dimension} \(\leq n_i\), \(f\) \uline{has perfect dimension} \(\leq n_1+n_2\). \uline{Consequently}, \uline{if} \(f_1\) \uline{and} \(f_2\) \uline{are perfect}, \uline{the same holds for} \(f\).

\sgaStableTarget{sga6.e03.cor.4-7-2}\textbf{Corollary 4.7.2.} \uline{Let}
\[
\begin{tikzcd}
X \arrow[d,"f"'] & X' \arrow[l,"g"'] \arrow[d,"f'"]\\
Y & Y' \arrow[l,"h"']
\end{tikzcd}
\]
\uline{be a Cartesian square of schemes}, \uline{where} \(f\) \uline{is locally of finite type and} \(X\) \uline{and} \(Y'\) \uline{are supposed tor-independent over} \(Y\), \uline{and let} \(E\in\Ob D(X)\), \uline{and} \([a,b]\) \uline{be an interval of} \(\mathbb Z\). \uline{We have}:
\[
 \parfamp_f(E)\subset[a,b]\Longrightarrow\parfamp_{f'}(\mathrm Lg^*(E))\subset[a,b],
\]
\uline{and consequently the functor} \(\mathrm Lg^*:D^-(X)\to D^-(X')\) \uline{sends} \(D(f)_{\parf}\) (resp. \(D(f)_{\uparf}\)) into \(D(f')_{\parf}\) (resp. \(D(f')_{\uparf}\)) (4.2).

\sgaStableTarget{sga6.e03.prop.4-8}\textbf{Proposition 4.8.} \uline{Let} \(S\) \uline{be a scheme}, \(f:X\to Y\) \uline{a proper} \(S\)-\uline{morphism of} \(S\)-\uline{schemes locally of finite type}, \(E\in\Ob D^b(X)\). \uline{Then}, \uline{subject to the validity of} (2.1) \uline{for} \(f\) (\uline{hence in any case when} \(S\) \uline{is locally noetherian} \uline{or when} \(f\) \uline{is projective}):
\begin{enumerate}[label=\alph*)]
\item \uline{If} \(E\) \uline{is perfect rel.} \uline{to} \(S\), \uline{the same holds for} \(Rf_*(E)\).
\item \uline{Suppose} \(S\) \uline{quasi-compact}, \(X\) \uline{and} \(Y\) \uline{of finite type over} \(S\); \uline{let} \([a,b]\) \uline{be an interval of} \(\mathbb Z\), \uline{and} \(N\) \uline{an integer such that} \(R^if_*(M)=0\) \uline{for every} \(i>N\) \uline{and}
%%%%%%%%%%  scan idx 271  |  folio 258  %%%%%%%%%%
\SGAPrintPage{258}%
\uline{every} \(\underline{O}_X\)-\uline{Module} \(M\) \uline{that is quasi-coherent} (\uline{such an integer exists by} (EGA III 1.4.12 \uline{and} IV 1.7.21))\,; \uline{we then have}
\[
 \parfamp_S(E)\subset[a,b]\Longrightarrow\parfamp_S(Rf_*(E))\subset[a,b+N].
\]
\end{enumerate}
\uline{Consequently}, \uline{the functor} \(Rf_*\) \uline{sends} \(D(X)_{S\text{-}\uparf}\) \uline{into} \(D(Y)_{S\text{-}\uparf}\)\SGAAuthorFootnote{1}{moreover, \(D(X)_{S\text{-}\uparf}=D(X)_{S\text{-}\parf}\) (and likewise with \(Y\)), since \(X\) and \(Y\) are quasi-compact (cf. (4.2)).}.

\uline{Proof}. Assertion a) follows from b), and b) is obtained by combining (2.1) and (3.7.1).

\sgaStableTarget{sga6.e03.cor.4-8-1}\textbf{Corollary 4.8.1.} \uline{Let} \(f:X\to Y\) \uline{be a proper and perfect morphism}, \uline{and} \(E\in\Ob D^+(X)\). \uline{Then}, \uline{subject to the validity of} (2.1) \uline{for} \(f\), \uline{if} \(E\) \uline{is perfect}, \uline{the same holds for} \(Rf_*(E)\)\,; \uline{moreover}, \uline{if} \(Y\) \uline{is quasi-compact}, \uline{the functor} \(Rf_*\) \uline{sends} \(D(X)_{\uparf}\) \uline{into} \(D(Y)_{\uparf}\)\SGAAuthorFootnote{2}{We have \(D(X)_{\uparf}=D(X)_{\parf}\) (and likewise with \(Y\)).}.

\uline{Proof}. Indeed, when \(E\) is perfect, \(E\) is perfect rel. to \(f\) by (4.5.1 a)).

\sgaStableTarget{sga6.e03.prop.4-9}\textbf{Proposition 4.9}\SGAAuthorFootnote{3}{This result and its corollaries will not be used in the remainder of this Seminar, but play an important role in the Lefschetz--Verdier formula for coherent sheaves; see SGA 5 Exp.\ III.}. \uline{Let} \(S\) \uline{be a noetherian scheme}, \(X\) \uline{and} \(Y\) \uline{be} \(S\)-\uline{schemes of finite type}, \(f:X\to Y\) \uline{a compactifiable} \(S\)-\uline{morphism} (\uline{i.e.} \uline{a composite of an open immersion and a proper morphism}), \(E\in\Ob D(X)\), \(F\in\Ob D^+(Y)\). \uline{If} \(E\) \uline{is perfect rel.} \uline{to} \(f\), \uline{and} \(F\) \uline{perfect rel.} \uline{to} \(S\), \(\mathrm R\underline{\Hom}(E,f^!(F))\) \uline{is perfect rel.} \uline{to} \(S\)\SGAAuthorFootnote{4}{For the definition of the functor \(f^!\), see [H] App.\ n\textsuperscript{o}~4.}.

\uline{Proof}. The question being local on \(X\), we may suppose that \(f\) factors as \(X\xrightarrow{i}X'\xrightarrow{f'}Y\), where \(i\) is a closed immersion and \(f'\) a smooth morphism. When \(E\) is perfect rel. to \(f\), \(E\) is perfect rel. to \(i\) (4.4), and when \(F\) is perfect rel. to \(S\), \(f'^!(F)\) is perfect rel. to \(S\) by (4.5.1 b)), since \(f'^!(F)\) is locally isomorphic to \(f'^*(F)\) up to a shift. Since \(f^!\simeq i^!f'^!\), we therefore reduce to proving the assertion when \(f\) is a closed immersion. Thus suppose \(E\) perfect rel. to \(f\), i.e. \(f_*(E)\) perfect, and \(F\) perfect rel. to \(S\). We must show that \(f_*\mathrm R\underline{\Hom}(E,f^!F)\) is
%%%%%%%%%%  scan idx 272  |  folio 259  %%%%%%%%%%
\SGAPrintPage{259}%
perfect rel. to \(S\). But by the duality theorem ([H] III 6.7), we have
\[
 f_*\mathrm R\underline{\Hom}(E,f^!F)\simeq\mathrm R\underline{\Hom}(f_*(E),F).
\]
Since \(f_*(E)\) is perfect, we have, by (I 7.7),
\[
 \mathrm R\underline{\Hom}(f_*(E),F)\simeq(f_*(E))^{\vee}\otimes^{\mathrm L}F,
\]
and (4.5) (where \(X=Y\) and \(f=\Id_X\)) gives the result.

\sgaStableTarget{sga6.e03.cor.4-9-1}\textbf{Corollary 4.9.1.} \uline{Under the hypotheses of} (4.9), \uline{if} \(f\) \uline{is perfect}, \(f^!\) \uline{sends} \(D(Y)_{S\text{-}\parf}\) \uline{into} \(D(X)_{S\text{-}\parf}\).

\uline{Proof}. Take \(E=\underline{O}_X\).

\sgaStableTarget{sga6.e03.cor.4-9-2}\textbf{Corollary 4.9.2.} \uline{Let} \(S\) \uline{be a noetherian scheme}, \(f:X\to S\) \uline{a compactifiable morphism}. \uline{Put} \(f^!(\underline{O}_S)=K_X\), \(\mathrm R\underline{\Hom}(.,K_X)=\underline{D}_X\). \uline{For} \(E\in\Ob D(f)_{\parf}\), \uline{we have} \(\underline{D}_X(E)\in\Ob D(f)_{\parf}\), \uline{and the canonical arrow} (I 7.2)
\[
 E\longrightarrow \underline{D}_X\underline{D}_X(E)
\]
\uline{is an isomorphism}.

\uline{Proof}. For the first assertion, apply (4.9) with \(Y=S\), \(F=\underline{O}_S\). For the second, we easily reduce to the case where \(f\) is a closed immersion. It then suffices to prove that the arrow obtained from the canonical arrow by applying \(f_*\) is an isomorphism. By global duality ([H] III 6.7), we have
\[
 f_*\underline{D}_X\simeq\underline{D}_Sf_*.
\]
We therefore reduce to showing that the canonical arrow
\[
 f_*(E)\longrightarrow \underline{D}_S\underline{D}_Sf_*(E)
\]
is an isomorphism for \(E\) perfect rel. to \(f\), i.e. \(f_*(E)\) perfect, which follows from ordinary biduality for perfect complexes (I 7.2).

\sgaStableTarget{sga6.e03.sec.5}\section*{5. Applications: exchange and semi-continuity theorems\SGAAuthorFootnote{1}{The results of this section will not be used in the remainder of the Seminar.}}

\sgaStableTarget{sga6.e03.lem.5-1}\textbf{Lemma 5.1.} \uline{Let} \((X,\underline{O}_X)\) \uline{be a ringed topos} (\uline{commutative}), \uline{denote by} \(\Mod(X)\) \uline{the category of} \(\underline{O}_X\)-\uline{Modules}, \(D(X)=D(\Mod(X))\). \uline{Let} \(E\in\Ob D^-(X)\), \uline{and} \(p\in\mathbb Z\). \uline{Consider the following conditions}:
%%%%%%%%%%  scan idx 273  |  folio 260  %%%%%%%%%%
\SGAPrintPage{260}%
\begin{enumerate}[label=\alph*)]
\item \uline{The functor} \(M\mapsto H^p(E\otimes^{\mathrm L}M)\) \uline{from} \(\Mod(X)\) \uline{to} \(\Mod(X)\) \uline{is right exact}.
\item[a')] \uline{The canonical arrow}
\[
 H^p(E)\otimes M\longrightarrow H^p(E\otimes^{\mathrm L}M)
\]
\uline{is an isomorphism for every} \(M\in\Ob\Mod(X)\).
\item[b)] \uline{The functor} \(M\mapsto H^{p+1}(E\otimes^{\mathrm L}M)\) \uline{from} \(\Mod(X)\) \uline{to} \(\Mod(X)\) \uline{is left exact}.
\item[b')] \uline{The canonical arrow defined below}
\begin{equation}
 H^{p+1}(E\otimes^{\mathrm L}M)\longrightarrow\underline{\Hom}(H^{-p-1}(E^{\vee}),M)\tag{$*$}
\end{equation}
\uline{is an isomorphism for every} \(M\in\Ob\Mod(X)\).
\end{enumerate}
\uline{Then we have}:
\[
 \mathrm a)\Longleftrightarrow\mathrm a')\Longleftrightarrow\mathrm b)\Longleftarrow\mathrm b').
\]
\uline{Moreover}, \uline{if} \(E\) \uline{is pseudo-coherent}, \uline{the conditions} a), a'), b), b') \uline{are equivalent}.

\uline{Definition of} \((*)\): The canonical arrow (I 7.5) \(E^{\vee}\otimes^{\mathrm L}E\longrightarrow\underline{O}_X\), where \(E^{\vee}=\mathrm R\underline{\Hom}(E,\underline{O}_X)\), gives, by tensoring with \(M\in\Ob\Mod(X)\), an arrow \(E^{\vee}\otimes^{\mathrm L}E\otimes^{\mathrm L}M\longrightarrow M\); hence, by the adjunction formula dear to Cartan (I 7.4), a canonical arrow
\begin{equation}
 E\otimes^{\mathrm L}M\longrightarrow\mathrm R\underline{\Hom}(E^{\vee},M).\tag{$**$}
\end{equation}
Composing the arrow induced by \((**)\) on \(H^{p+1}\) with the canonical arrow
\begin{equation}
 \underline{\Ext}^{p+1}(E^{\vee},M)\longrightarrow\underline{\Hom}(H^{-p-1}(E^{\vee}),M),\tag{$***$}
\end{equation}
we obtain the arrow \((*)\).

\uline{Proof}. The implications
\[
 \mathrm a')\Longrightarrow\mathrm a)\Longleftrightarrow\mathrm b)\Longleftarrow\mathrm b')
\]
are trivial. Let us prove \(\mathrm a)\Longrightarrow\mathrm a')\). Let \(M\) be an \(\underline{O}_X\)-Module. There is an exact sequence \(M^1\longrightarrow M^0\longrightarrow M\longrightarrow0\), where \(M^0\) and \(M^1\) are direct sums of sheaves
%%%%%%%%%%  scan idx 274  |  folio 261  %%%%%%%%%%
\SGAPrintPage{261}%
of the type \(\underline{O}_{U,X}\), where \(U\) is an object of \(X\). Using the right exactness and commutation with direct sums of the functors \(H^p(E)\otimes\,.\,\) and \(H^p(E\otimes^{\mathrm L}\,.\,)\), we reduce to proving that the arrow of a') is an isomorphism for \(M=\underline{O}_{U,X}\), \(U\in\Ob X\). Now \(H^p(E)\otimes\underline{O}_{U,X}\) (resp. \(H^p(E\otimes^{\mathrm L}\underline{O}_{U,X})\)) identifies canonically (SGA 4 IV 2.13) with \(i_!(H^p(E)|U)\) (resp. \(H^p(i_!(E|U))\)), where \(i_!\) is extension by \(0\) outside \(U\), and the result follows because \(i_!\) is exact.

Suppose \(E\) pseudo-coherent, and let us prove \(\mathrm b)\Longrightarrow\mathrm b')\). It clearly suffices to show that \((*)\) is an isomorphism for \(M\) injective. Since the arrow \((***)\) is an isomorphism for \(M\) injective, the assertion will therefore follow from the following one:

\((5.1.1)\): \uline{the arrow} \((**)\) \uline{is an isomorphism if} \(E\) \uline{is pseudo-coherent and} \(M\) \uline{injective}.

\uline{Proof of} \((5.1.1)\). Denote by \(T\) (resp. \(T'\)) the functor from \(D^-(X)\) to \(D(X)\) defined by \(F\mapsto F\otimes^{\mathrm L}M\) (resp. \(\mathrm R\underline{\Hom}(F^{\vee},M)\)). Let \(n\in\mathbb Z\); if \(H^i(F)=0\) for \(i>n\), then \(H^iT(F)=0\) for \(i>n\), \(H^i(F^{\vee})=0\) for \(i<-n\), and \(H^iT'(F)=0\) for \(i>n\) if \(M\) is injective. Since, on the other hand, the arrow \((**)\) commutes with restriction to an open, and an injective restricted to an open is still injective, we may, after localizing and truncating, suppose \(E\) strictly perfect (I 2.1). By dévissage we then reduce to \(E=\underline{O}_X\), and the result follows.

\sgaStableTarget{sga6.e03.rem.5-2}\textbf{Remarks 5.2.} (i) Condition a) (resp. a'), b), b')) is local in nature: if \((X_\alpha)\) is an open covering of the final object of \(X\), and if, for each \(\alpha\), condition a) (resp\ldots) holds for \(X_\alpha\) and \(E\,|\,X_\alpha\), then a) (resp\ldots) holds.

(ii) If \(x\) is a point of \(X\), denote by \(\mathrm a_x)\) (resp. \(\mathrm a'_x)\), etc.) the condi%
%%%%%%%%%%  scan idx 275  |  folio 262  %%%%%%%%%%
\SGAPrintPage{262}%
tion a) (resp. a'), etc.) relative to \(E_x\) and \(\Mod(\underline{O}_{X,x})\). Then, if \(X\) has enough points, condition a) (resp. a'), b)) is equivalent to the condition ``for every point \(x\) of \(X\), \(\mathrm a_x)\) (resp. \(\mathrm a'_x)\), \(\mathrm b_x))\)'', and is equivalent to the condition ``for every point \(x\) of \(X\), \(\mathrm b'_x)\)'' if \(E\) is pseudo-coherent; moreover, condition a') is equivalent to the analogous condition in which ``for every \(M\in\Ob\Mod(X)\)'' is replaced by ``for every \(\underline{O}_X\)-Module that is the cokernel of a homomorphism \(\underline{O}_{U,X}^p\longrightarrow\underline{O}_{U,X}^q\), where \(U\) is an object of \(X\)''.

(iii) The criterion (EGA III 7.4.2) extends word for word to our situation: if \(E\) has flat components, condition a) simply means that \(\Coker(E^p\longrightarrow E^{p+1})\) is flat, hence locally a direct factor of a free Module of finite type when \(E^p\) and \(E^{p+1}\) are also supposed of finite presentation (I 5.8.3). In particular, we deduce that, if \(E\) is perfect, then \(E\) satisfies condition a) relative to the integer \(p\) if and only if \(E^{\vee}\) satisfies condition a) relative to the integer \(-p-1\).

(iv) Let \(X\) be a \uline{quasi-separated} scheme, \(\underline{O}_X\) the structure sheaf, \(E\in\Ob D^-(X)\) a complex with locally bounded and quasi-coherent cohomology. Denote by \(\mathrm a_{\qcoh})\) (resp. \(\mathrm a'_{\qcoh})\), etc.) the condition analogous to a) (resp. a'), etc.) in which \(\Mod(X)\) is replaced by the category \(\Qcoh(X)\) of quasi-coherent \(\underline{O}_X\)-Modules. Then we have
\[
 \mathrm a)\Longleftrightarrow\mathrm a_{\qcoh})\Longleftrightarrow\mathrm a'_{\qcoh})\Longleftrightarrow\mathrm b_{\qcoh})\Longleftarrow\mathrm b'_{\qcoh}),
\]
and the preceding conditions are equivalent if \(E\) is pseudo-coherent.

Indeed, since \(X\) is quasi-separated, every quasi-coherent sheaf on an affine open of \(X\) is induced by a quasi-coherent sheaf on \(X\) (EGA I 9.4.2 and IV 1.2.7). It follows that condition \(\mathrm a_{\qcoh})\) (resp. \(\mathrm a'_{\qcoh})\), \(\mathrm b_{\qcoh})\), \(\mathrm b'_{\qcoh}))\) holds if and only if it holds over every affine open. To prove the implications above, we may therefore suppose \(X\) affine. Then (II 3.5) \(E\) is isomorphic in \(D(X)\) to a complex of quasi-coherent Modules, and the implications
\[
 \mathrm a_{\qcoh})\Longleftrightarrow\mathrm a'_{\qcoh})\Longleftrightarrow\mathrm b_{\qcoh})\Longleftarrow\mathrm b'_{\qcoh})
\]
%%%%%%%%%%  scan idx 276  |  folio 263  %%%%%%%%%%
\SGAPrintPage{263}%
follow from (5.1) applied to the ``one-point'' topos ringed by \(\Gamma(X,\underline{O}_X)\); the same holds, by (II 2.2.10 a)), for the implication \(\mathrm b_{\qcoh})\Longrightarrow\mathrm b'_{\qcoh})\) when \(E\) is pseudo-coherent. Finally, the equivalence \(\mathrm a)\Longleftrightarrow\mathrm a_{\qcoh})\), or, equivalently, \(\mathrm a')\Longleftrightarrow\mathrm a'_{\qcoh})\) follows from the last remark of (ii) (and the fact that quasi-coherent sheaves ``extend'', since \(X\) is quasi-separated).

\sgaStableTarget{sga6.e03.prop.5-3}\textbf{Proposition 5.3.} (\uline{exchange property}) (EGA III 7.7.5 II). \uline{Let} \(f:X\to Y\) \uline{be a morphism of quasi-compact and quasi-separated schemes}, \(E\in\Ob D^b(X)\) \uline{a complex with quasi-coherent cohomology and finite flat amplitude rel.} \uline{to} \(f\), \uline{and} \(p\in\mathbb Z\). \uline{Consider the following conditions}:
\begin{enumerate}[label=\alph*)]
\item[a)] \uline{The functor} \(M\mapsto R^pf_*(E\otimes_Y^{\mathrm L}M)\) \uline{from} \(\Qcoh(Y)\)\SGAAuthorFootnote{1}{\(\Qcoh(Y)\) = category of quasi-coherent \(\underline{O}_Y\)-Modules.} \uline{to} \(\Qcoh(Y)\)\SGAAuthorFootnote{2}{by (3.7).} \uline{is right exact}.
\item[a')] \uline{The canonical arrow}
\[
 R^pf_*(E)\otimes M\longrightarrow R^pf_*(E\otimes_Y^{\mathrm L}M)
\]
(\uline{the composite of the canonical arrow} \(R^pf_*(E)\otimes M\longrightarrow H^p(Rf_*(E)\otimes^{\mathrm L}M)\) \uline{and the isomorphism} (3.7) \(H^p(Rf_*(E)\otimes^{\mathrm L}M)\xrightarrow{\sim}R^pf_*(E\otimes_Y^{\mathrm L}M)\)) \uline{is an isomorphism for every} \(M\in\Ob\Qcoh(Y)\).
\item[b)] \uline{The functor} \(M\mapsto R^{p+1}f_*(E\otimes_Y^{\mathrm L}M)\) \uline{from} \(\Qcoh(Y)\) \uline{to} \(\Qcoh(Y)\) \uline{is left exact}.
\item[b')] \uline{The canonical arrow}
\[
 R^{p+1}f_*(E\otimes_Y^{\mathrm L}M)\longrightarrow\underline{\mathrm{Hom}}(H^{-p-1}(Rf_*(E)^{\vee}),M)
\]
(\uline{the composite of the isomorphism} (3.7) \(R^{p+1}f_*(E\otimes_Y^{\mathrm L}M)\simeq H^{p+1}(Rf_*(E)\otimes^{\mathrm L}M)\) \uline{and the canonical arrow} (5.1 (\(*\))) \(H^{p+1}(Rf_*(E)\otimes^{\mathrm L}M)\longrightarrow\underline{\mathrm{Hom}}(H^{-p-1}(Rf_*(E)^{\vee}),M)\)) \uline{is an isomorphism for every} \(M\in\Ob\Qcoh(Y)\).
%%%%%%%%%%  scan idx 277  |  folio 264  %%%%%%%%%%
\SGAPrintPage{264}%
\item[c)]
\uline{For every base change} \(Y'\to Y\), \uline{the canonical arrow}
\[
 R^pf_*(E)\otimes_Y\underline{O}_{Y'}\longrightarrow H^p(Rf_*(E)\otimes_Y^{\mathrm L}\underline{O}_{Y'})
\]
\uline{is an isomorphism}.
\end{enumerate}
\uline{Then we have}:
\[
 \mathrm a)\Longleftrightarrow\mathrm a')\Longleftrightarrow\mathrm b)\Longleftrightarrow\mathrm c)\Longleftarrow\mathrm b').
\]
\uline{Moreover}, \uline{if} \(f\) \uline{is proper}, \(E\) \uline{pseudo-coherent rel.} \uline{to} \(f\) (\uline{hence perfect rel.} \uline{to} \(f\)), \uline{and subject to the validity of} (2.1) \uline{for} \(f\) (\uline{hence in particular if} \(f\) \uline{is projective or if} \(Y\) \uline{is noetherian}), \uline{the preceding conditions are equivalent}.

\uline{Proof}. The implications \(\mathrm a)\Longleftrightarrow\mathrm a')\Longleftrightarrow\mathrm b)\Longleftarrow\mathrm b')\) follow at once from (3.7), (5.1) and (5.2 (iv)), as does the implication \(\mathrm b')\Longleftarrow\mathrm b)\) (for \(f\) proper and \(E\) pseudo-coherent rel. to \(f\)) by (2.1). It remains to prove that \(\mathrm a)\Longleftrightarrow\mathrm c)\). This will follow from the first assertion of the following lemma:

\sgaStableTarget{sga6.e03.lem.5-3-1}\textbf{Lemma 5.3.1.} \uline{Let} \(X\) \uline{be a quasi-separated scheme}, \uline{and let} \(E\in\Ob D^{-}(X)\) \uline{be a complex with locally bounded and quasi-coherent cohomology}. \uline{Then condition} (5.1 a)) \uline{is equivalent to the condition}
\begin{enumerate}[label=\alph*)]
\item[c)] \uline{For every scheme} \(X'\) \uline{over} \(X\), \uline{the canonical arrow}
\[
 H^p(E)\otimes_X\underline{O}_{X'}\longrightarrow H^p(E\otimes_X^{\mathrm L}\underline{O}_{X'})
\]
\uline{is an isomorphism}.
\end{enumerate}
\uline{Moreover}, \uline{if} \(E\) \uline{is pseudo-coherent}, \uline{condition} (5.1 a)) \uline{implies that}, \uline{for every morphism} \(X'\to X\), \uline{the canonical arrow}
\sgaStableTarget{sga6.e03.eq.5-3-1-1}\begin{equation}
 H^{-p-1}(E^{\vee})\otimes_X\underline{O}_{X'}\longrightarrow H^{-p-1}((E\otimes_X^{\mathrm L}\underline{O}_{X'})^{\vee}),\tag{5.3.1.1}
\end{equation}
%%%%%%%%%%  scan idx 278  |  folio 265  %%%%%%%%%%
\SGAPrintPage{265}%
\uline{the composite of} \(H^{-p-1}(E^{\vee})\otimes_X\underline{O}_{X'}\longrightarrow H^{-p-1}(E^{\vee}\otimes_X^{\mathrm L}\underline{O}_{X'})\) \uline{and the canonical arrow} (I 7.1.2) \(H^{-p-1}(E^{\vee}\otimes_X^{\mathrm L}\underline{O}_{X'})\longrightarrow H^{-p-1}((E\otimes_X^{\mathrm L}\underline{O}_{X'})^{\vee})\), \uline{is an isomorphism}.

\uline{Proof}. \(\mathrm a)\Longrightarrow\mathrm c)\). The question being local on \(X'\), we may suppose \(X\) (resp. \(X'\)) affine with ring \(A\) (resp. \(A'\)). By (II 3.5), we may also suppose that \(E=\widetilde F\), where \(F\) is a complex of \(A\)-modules. By hypothesis (condition a')), the arrow
\[
 H^p(F)\otimes_A A'\longrightarrow H^p(F\otimes_A^{\mathrm L}A')
\]
is an isomorphism, and the result follows since \(H^p(E)\otimes_X\underline{O}_{X'}\) (resp. \(H^p(E\otimes_X^{\mathrm L}\underline{O}_{X'})\)) identifies canonically with the sheaf on \(X'\) defined by the \(A'\)-module \(H^p(F)\otimes_A A'\) (resp. \(H^p(F\otimes_A^{\mathrm L}A')\)).

\(\mathrm c)\Longrightarrow\mathrm a)\). In view of (5.2 (iv)), it suffices to prove \(\mathrm c)\Longrightarrow\mathrm a'_{\qcoh})\). The question being local (since \(X\) is quasi-separated), we may suppose \(X\) affine with ring \(A\), and \(E=\widetilde F\), where \(F\) is a complex of \(A\)-modules. Let \(M=\widetilde N\) be a quasi-coherent \(\underline{O}_X\)-Module. We must show that the arrow
\[
 H^p(F)\otimes_A N\longrightarrow H^p(F\otimes_A^{\mathrm L}N)
\]
is an isomorphism. But, putting \(A'=D_A(N)\) (EGA \(\mathrm O_{\mathrm{IV}}\) 18.2.3), \(X'=\Spec(A')\), this follows from c) applied to the canonical morphism \(X'\to X\).

Let us prove the last assertion. The question being local on \(X'\), we may, by (II 2.2.10 b)), suppose \(E\) strictly pseudo-coherent. Moreover, since the assertion involves only a part of \(E\) of bounded length, we may, after truncating, suppose \(E\) strictly perfect. Then \((E\otimes_X^{\mathrm L}\underline{O}_{X'})^{\vee}\) identifies canonically (I 7.1.2) with \(E^{\vee}\otimes_X^{\mathrm L}\underline{O}_{X'}\), and the arrow (5.3.1.1) may be rewritten as the canonical arrow
\sgaStableTarget{sga6.e03.eq.unnamed.occ2}\begin{equation}
 H^{-p-1}(E^{\vee})\otimes_X\underline{O}_{X'}\longrightarrow H^{-p-1}(E^{\vee}\otimes_X^{\mathrm L}\underline{O}_{X'}).\tag{$*$}
\end{equation}
%%%%%%%%%%  scan idx 279  |  folio 266  %%%%%%%%%%
\SGAPrintPage{266}%
If \(E\) satisfies condition (5.1 a)) for the integer \(p\), \(E^{\vee}\) satisfies the same condition for the integer \(-p-1\) by (5.2 (iii)), and the arrow \((*)\) is an isomorphism by the first assertion of (5.3.1). This completes the proof of (5.3.1), and hence of (5.3).

\sgaStableTarget{sga6.e03.rem.5-3-2}\textbf{Remark 5.3.2.} Proposition (5.3) applies in particular when \(E\) is a complex of bounded degrees of quasi-coherent \(\underline{O}_X\)-Modules that are flat and of finite presentation rel. to \(f\): indeed, \(E\) then has finite perfect amplitude rel. to \(f\) by (4.3.2).

%%%%%%%%%%  scan idx 280  |  folio 267  %%%%%%%%%%
\SGAPrintPage{267}%
\sgaStableTarget{sga6.e03.def.5-4}\textbf{Definition 5.4.} \uline{Let} \(X\) \uline{be a topos}, \uline{and denote by} \(\pt(X)\) \uline{the set of its points} (SGA 4 IV 6.1). \uline{Let} \(M\) \uline{be an ordered set}, \(f:\pt(X)\to M\) \uline{a map}, \uline{and} \(x\in\pt(X)\). \uline{We say that} \(f\) \uline{is upper semi-continuous} (\uline{resp.} \uline{lower semi-continuous}, \uline{resp.} \uline{locally constant}) \uline{at} \(x\) \uline{if there exists} \uline{an} \(x\)-\uline{pointed object} \(U\) \uline{of} \(X\) \uline{such that for every point} \(y\) \uline{of} \(U\) \uline{we have} \(f(y)\leq f(x)\) (\uline{resp.} \(f(y)\geq f(x)\), \uline{resp.} \(f(y)=f(x)\)). \uline{We say that} \(f\) \uline{is upper semi-continuous} (\uline{resp.} \uline{lower semi-continuous}, \uline{resp.} \uline{locally constant}) \uline{if it is so at every point}\SGAAuthorFootnote{1}{When \(X\) comes from a topological space, we recover the usual notions.}.

\sgaStableTarget{sga6.e03.ex.5-4-1}\textbf{Examples 5.4.1.} a) Denote by \(M_X\) the constant sheaf with value \(M\) on \(X\), and let \(s\in H^0(X,M)\). The map \(x\mapsto s(x)=x^*(s)\) from \(\pt(X)\) to \(M\) is locally constant. In particular, let \(E\) be a perfect complex on \(X\), supposed locally ringed: the rank of \(E\) (I 6.11) defines a locally constant function on \(\pt(X)\) with values in \(\mathbb Z\).

b) Suppose \(X\) locally ringed, and let \(P\) be an \(\underline{O}_X\)-Module of finite type. For \(x\in\pt(X)\), denote by \(\underline{O}_{X,x}\) the local stalk ring of \(\underline{O}_X\) at \(x\), and by \(k(x)\) its residue field. Then, for \(x\in\pt(X)\), \(P_x\otimes_{\underline{O}_{X,x}}k(x)\) has finite rank over \(k(x)\), and the function \(x\mapsto\operatorname{rg}_{k(x)}(P\otimes_{\underline{O}_{X,x}}k(x))\) is upper semi-continuous. Indeed, if, for \(x\in\pt(X)\), \(P_x\otimes_{\underline{O}_{X,x}}k(x)\) has rank \(n\), there is an epimorphism \(\underline{O}_{X,x}^n\to P_x\) by Nakayama, and, since \(P\) is of finite type, there is an \(x\)-pointed object \(U\) of \(X\) and an epimorphism \(\underline{O}_U^n\to P\,|\,U\), giving the assertion.

\sgaStableTarget{sga6.e03.lem.5-5}\textbf{Lemma 5.5.} \uline{Let} \((X,\underline{O}_X)\) \uline{be a locally ringed topos}, \(E\in\Ob D^{-}(X)\), \uline{and} \(p\in\mathbb Z\). \uline{Then we have}, \uline{with the notation of} (5.4.1 b)):
\begin{enumerate}[label=(\roman*)]
%%%%%%%%%%  scan idx 281  |  folio 268  %%%%%%%%%%
\item\SGAPrintPage{268}%
\uline{If} \(E\) \uline{is perfect}, \uline{the} \(H^i(E_x\otimes_{\underline{O}_{X,x}}^{\mathrm L}k(x))\), \uline{for} \(x\in\pt(X)\), \uline{are vector spaces of finite rank over} \(k(x)\), \uline{all but finitely many zero}, \uline{and the function}
\[
 x\longmapsto\sum(-1)^i\operatorname{rg}_{k(x)}H^i(E_x\otimes_{\underline{O}_{X,x}}^{\mathrm L}k(x))
\]
\uline{is locally constant on} \(\pt(X)\).
\item \uline{If} \(E\) \uline{is} \((p-1)\)-\uline{pseudo-coherent}, \(H^p(E_x\otimes_{\underline{O}_{X,x}}^{\mathrm L}k(x))\) \uline{is}, \uline{for} \(x\in\pt(X)\), \uline{a vector space of finite rank over} \(k(x)\), \uline{and the function}
\[
 x\longmapsto\operatorname{rg}_{k(x)}H^p(E_x\otimes_{\underline{O}_{X,x}}^{\mathrm L}k(x))
\]
\uline{is upper semi-continuous on} \(\pt(X)\). \uline{If moreover the functor} \(M\mapsto H^p(E\otimes^{\mathrm L}M)\) \uline{from} \(\Mod(X)\) \uline{to} \(\Mod(X)\) \uline{is exact}, \(H^p(E)\) \uline{is locally free of finite type and the preceding function is locally constant}.
\item \uline{Suppose} \(E\) \(p\)-\uline{pseudo-coherent} (\uline{resp.} \((p-1)\)-\uline{pseudo-coherent}), \uline{and let} \(x\in\pt(X)\). \uline{If the functor} \(M\mapsto H^p(E_x\otimes_{\underline{O}_{X,x}}^{\mathrm L}M)\) \uline{is right exact} (\uline{resp.} \uline{left exact}) \uline{on the category} \(\Mod(x)\) \uline{of} \(\underline{O}_{X,x}\)-\uline{Modules}, \uline{there is an} \(x\)-\uline{pointed open} \(U\) \uline{of} \(X\) \uline{such that the functor} \(M\mapsto H^p(E\,|\,U\otimes^{\mathrm L}M)\) \uline{is right exact} (\uline{resp.} \uline{left exact}) \uline{on} \(\Mod(U)\).
\end{enumerate}

\uline{Proof}. (i) By (I 4.19.1), \(E_x\otimes_{\underline{O}_{X,x}}^{\mathrm L}k(x)\) is a perfect complex of \(k(x)\)-modules, so (I 5.9 and 3.5) the \(H^i(E_x\otimes_{\underline{O}_{X,x}}^{\mathrm L}k(x))\) have finite rank over \(k(x)\) and all but finitely many are zero. We have
\begin{equation}
 \operatorname{rg}(E_x\otimes_{\underline{O}_{X,x}}^{\mathrm L}k(x))=\sum(-1)^i\operatorname{rg}_{k(x)}H^i(E_x\otimes_{\underline{O}_{X,x}}^{\mathrm L}k(x)),\qquad\mbox{\SGAAuthorFootMark{1}}\tag{5.5.1}
\end{equation}
\SGAAuthorFootText{1}{\(\operatorname{rg}(\ )\) denotes the rank of a perfect complex (I 6.11).}%
as is seen by an immediate dévissage using the additivity of the rank function. In view of (5.4.1 a)), assertion (i) therefore follows from
%%%%%%%%%%  scan idx 282  |  folio 269  %%%%%%%%%%
\SGAPrintPage{269}%
the compatibility of rank with inverse images (I 6.12).

(ii) After truncating and shifting, we may, for the first assertion, suppose \(E\) strictly perfect, \(p=0\), and \(E^i=0\) for \(i\notin[-1,+1]\). For \(x\in\pt(X)\), put \(E_x\otimes_{\underline{O}_{X,x}}^{\mathrm L}k(x)=E_{k(x)}\). By (5.5.1) we have
\[
\begin{gathered}
 \operatorname{rg}_{k(x)}H^0(E_{k(x)})=\operatorname{rg}_{k(x)}H^{-1}(E_{k(x)})+\operatorname{rg}_{k(x)}H^1(E_{k(x)})\\
 +\operatorname{rg}(E_x),
\end{gathered}
\]
so it suffices to prove that the function \(x\mapsto\operatorname{rg}_{k(x)}H^{-1}(E_{k(x)})\) (resp. \(\operatorname{rg}_{k(x)}H^1(E_{k(x)})\)) is upper semi-continuous (since the function \(x\mapsto\operatorname{rg}(E_x)=\operatorname{rg}(E_{k(x)})\) is locally constant). For the function \(x\mapsto\operatorname{rg}_{k(x)}H^1(E_{k(x)})\), this is an immediate consequence of (5.4.1 b))\SGAAuthorFootnote{*}{taking into account that \(H^1(E_{k(x)})\simeq H^1(E)_{k(x)}\).}. For the function \(x\mapsto\operatorname{rg}_{k(x)}H^{-1}(E_{k(x)})\), we may suppose \(E^1=0\), i.e. replace \(E\) by \(E'=(0\to E^{-1}\to E^0\to0)\). We already know that the function \(x\mapsto\operatorname{rg}_{k(x)}H^0(E'_{k(x)})\) is upper semi-continuous, so the same holds for the function \(x\mapsto\operatorname{rg}_{k(x)}H^{-1}(E'_{k(x)})=\operatorname{rg}_{k(x)}H^{-1}(E_{k(x)})\), since the function \(x\mapsto\operatorname{rg}_{k(x)}H^0(E'_{k(x)})-\operatorname{rg}_{k(x)}H^{-1}(E'_{k(x)})=\operatorname{rg}(E'_x)\) is locally constant. Let us now prove the second assertion. After truncating, we may suppose \(E\) strictly perfect and \(E^i=0\) for \(i<p-1\). The \(\underline{O}_X\)-Module \(Z'^i(E)=\Coker(E^{i-1}\to E^i)\) is therefore of finite presentation. We have an exact sequence:
\begin{equation}
 0\to H^p(E)\to Z'^p(E)\to E^{p+1}\to Z'^{p+1}(E)\to0.\tag{$*$}
\end{equation}
Since the functor \(M\mapsto H^p(E\otimes^{\mathrm L}M)\) is right and left exact, \(Z'^{p+1}(E)\) and \(Z'^p(E)\) are flat by (5.2 (iii)), hence (I 5.8.3) locally free of finite type. The exactness of (*) then implies that \(H^p(E)\) is locally free of finite type, and consequently that the function \(x\mapsto\operatorname{rg}_{k(x)}H^p(E_{k(x)})=\operatorname{rg}(H^p(E)_x)\)
%%%%%%%%%%  scan idx 283  |  folio 270  %%%%%%%%%%
\SGAPrintPage{270}%
is locally constant, which completes the proof of (ii).

(iii) We may suppose \(E\) strictly perfect and \(E^i=0\) for \(i<p\) (resp. \(p-1\)). Saying that the functor \(M\mapsto H^p(E_x\otimes_{\underline{O}_{X,x}}^{\mathrm L}M)\) is right exact (resp. left exact) on \(\Mod(x)\) is equivalent, by (5.2 (iii)), to saying that \(Z'^{p+1}(E_x)\) (resp. \(Z'^p(E_x)\)) (with the notation of (ii)) is flat, hence free of finite type (since it is of finite presentation). Since \(Z'^{p+1}(E_x)\simeq Z'^{p+1}(E)_x\) (resp. \(Z'^p(E_x)\simeq Z'^p(E)_x\)), the usual argument (cf. EGA \(\mathrm O_I\) 5.2.7) shows that there is an \(x\)-pointed object \(U\) of \(X\) such that \(Z'^{p+1}(E\,|\,U)\) (resp. \(Z'^p(E\,|\,U)\)) is free of finite type, and the result follows.

\sgaStableTarget{sga6.e03.cor.5-6}\textbf{Corollary 5.6} (EGA III 7.6.9). \uline{Let} \(X\) \uline{be a scheme}, \(p\in\mathbb Z\), and \(E\in\Ob D^{-}(X)\) \uline{a} \((p-1)\)-\uline{pseudo-coherent complex}. \uline{Then}:
\begin{enumerate}[label=(\roman*)]
\item \uline{For} \(x\in X\), \(H^p(E\otimes_X^{\mathrm L}k(x))\) \uline{has finite rank over} \(k(x)\) \uline{and the function}
\[
 x\longmapsto\operatorname{rg}_{k(x)}H^p(E\otimes_X^{\mathrm L}k(x))
\]
\uline{is constructible and upper semi-continuous on} \(X\).
\item \uline{Suppose} \(X\) \uline{quasi-separated and} \(E\) \uline{with quasi-coherent and locally bounded cohomology}. \uline{If the functor} \(M\mapsto H^p(E\otimes^{\mathrm L}M)\) \uline{is exact on the category of quasi-coherent} \(\underline{O}_X\)-\uline{Modules}, \uline{the function} \(x\mapsto\operatorname{rg}_{k(x)}H^p(E\otimes_X^{\mathrm L}k(x))\) \uline{is locally constant and} \(H^p(E)\) \uline{is locally free of finite type}. \uline{Conversely}, \uline{if} \(X\) \uline{is locally noetherian and reduced and the preceding function is locally constant}, \uline{the functor} \(M\mapsto H^p(E\otimes^{\mathrm L}M)\) \uline{is exact on the category of quasi-coherent} \(\underline{O}_X\)-\uline{Modules}.
\end{enumerate}

\uline{Proof}. If \(X\) is integral with generic point \(x\), the functor \(M\mapsto H^p(E_x\otimes_{\underline{O}_{X,x}}^{\mathrm L}M)\) is exact on the category of \(\underline{O}_{X,x}\)-modules (since \(\underline{O}_{X,x}=k(x)\)), so (5.5 (iii)) there is an open neighborhood \(U\) of \(x\) such that the functor \(M\mapsto H^p(E\,|\,U\otimes^{\mathrm L}M)\) is exact on \(\Mod(U)\), and, on \(U\), the function
%%%%%%%%%%  scan idx 284  |  folio 271  %%%%%%%%%%
\SGAPrintPage{271}%
\(y\mapsto\operatorname{rg}_{k(y)}H^p(E\otimes_X^{\mathrm L}k(y))\) is locally constant by (5.5 (ii)). The constructibility assertion of (i) follows at once. The other assertions follow from (5.5 (ii) and (iii)) and (5.2 (iv)), except for the converse of (ii), for which it suffices, after reducing to \(E\) strictly perfect (by localization and truncation), to apply (EGA III 7.6.2 and 7.6.5) to the functor \(T_{-p}=H^p(E\otimes\,.\,)\).

\sgaStableTarget{sga6.e03.prop.5-7}\textbf{Proposition 5.7.} (EGA III 7.7.5 \uline{and} 7.9.4). \uline{Let} \(f:X\to Y\) \uline{be a proper morphism of quasi-compact schemes}, \(E\in\Ob D^b(f)_{\uparf}\) (4.2), \uline{and} \(p\in\mathbb Z\). \uline{Then}, \uline{subject to the validity of} (2.1), \uline{hence for example if} \(f\) \uline{is projective or if} \(Y\) \uline{is noetherian}:
\begin{enumerate}[label=(\roman*)]
\item \uline{For} \(y\in Y\), \uline{the} \(H^i(Rf_*(E)\otimes_Y^{\mathrm L}k(y))\) \uline{are vector spaces of finite rank over} \(k(y)\), \uline{all but finitely many zero}, \uline{and the function}
\[
 y\longmapsto\sum(-1)^i\operatorname{rg}_{k(y)}H^i(Rf_*(E)\otimes_Y^{\mathrm L}k(y))
\]
\uline{is locally constant on} \(Y\).
\item \uline{The function}
\[
 y\longmapsto\operatorname{rg}_{k(y)}H^p(Rf_*(E)\otimes_Y^{\mathrm L}k(y))
\]
\uline{is constructible and upper semi-continuous on} \(Y\).
\item \uline{Suppose} \(Y\) \uline{quasi-separated}. \uline{Then}, \uline{if the functor} \(M\mapsto R^pf_*(E\otimes_Y^{\mathrm L}M)\) \uline{is exact on the category of quasi-coherent} \(\underline{O}_Y\)-\uline{Modules}, \(R^pf_*(E)\) \uline{is a locally free} \(\underline{O}_Y\)-\uline{Module of finite type}, \uline{and the function of} (ii) \uline{is locally constant}. \uline{Conversely}, \uline{if} \(Y\) \uline{is noetherian and reduced}, \uline{and if the function of} (ii) \uline{is locally constant}, \uline{the functor} \(M\mapsto R^pf_*(E\otimes_Y^{\mathrm L}M)\) \uline{is exact on the category of quasi-coherent} \(\underline{O}_Y\)-\uline{Modules}.
\end{enumerate}

%%%%%%%%%%  scan idx 285  |  folio 272  %%%%%%%%%%
\SGAPrintPage{272}%
\uline{Proof}. By (4.8), \(Rf_*(E)\) is a perfect complex.

Assertion (i) (resp. (ii)) therefore follows from (5.5 (i)) (resp. (5.6 (i)) ). Finally, if \(M\) is a quasi-coherent \(\underline{O}_Y\)-Module, the canonical arrow
\[
 Rf_*(E)\otimes_Y^{\mathrm L}M\longrightarrow Rf_*(E\otimes_Y^{\mathrm L}M)
\]
is an isomorphism by (3.7), so assertion (iii) follows from (5.6 (ii)).

\sgaStableTarget{sga6.e03.rem.5-8}\textbf{Remarks 5.8.}
\begin{enumerate}[label=\alph*)]
\item In the situation of (5.7), denote by \(i_y:X_y\to X\) the inclusion of the fiber of \(X\) over a point \(y\) of \(Y\). One can show (cf. (EGA III 6) and (IV 3.1.0)) that, for \(f\) flat (resp. \(E\) bounded with quasi-coherent and \(Y\)-flat components), we have the Künneth formula:
\[
 H^i(Rf_*(E)\otimes_Y^{\mathrm L}k(y))\xrightarrow{\sim}H^i(X_y,E_y),
\]
where \(E_y=Li_y^*(E)\) (resp. \(i_y^*(E)\)).
%%%%%%%%%%  scan idx 286  |  folio 273  %%%%%%%%%%
\SGAPrintPage{273}%
\item One can develop for complex analytic spaces a body of results analogous to that of this section (exchange, semi-continuity, continuity properties)\SGAAuthorFootnote{1}{It was moreover in this setting that these properties were first formulated (see for example Grauert [1]), before being transposed to Algebraic Geometry in (EGA III 7).}. The essential ingredients are: 1) the finiteness theorem for proper morphisms (``Grauert's theorem''); 2) an analogue of the projection formula (3.7), which may be stated as follows: let \(f:X\to Y\) be a morphism of analytic spaces, where \(X\) is supposed of bounded dimension, and let \(E\in\Ob D^-(X)\), \(M\in\Ob D^-(Y)_{\coh}\) (``coh'' means ``with coherent cohomology''). The canonical arrow
\[
 Rf_*(E)\otimes^{\mathrm L}M\longrightarrow Rf_*(E\otimes_Y^{\mathrm L}M)
\]
is an isomorphism (the hypothesis on \(X\) ensures (T.F II 5.12) that \(f_*\) has finite cohomological dimension, hence that \(Rf_*\) is defined on \(D(X)\)).

Unlike statement 1), statement 2) is essentially trivial, and follows from a dévissage analogous to that of the proof of (3.7); we moreover have greater freedom since \(Rf_*\) is defined on \(D(X)\), and in particular no hypothesis on \(E\) is needed (apart from \(E\in\Ob D^-(X)\)).
\end{enumerate}

\sgaStableTarget{sga6.e03.bibliography}\subsection*{BIBLIOGRAPHY}
\begin{enumerate}[label={[\arabic*]}]
\item[{[1]}] Grauert, H., Ein Theorem der Analytischen Garbentheorie \ldots, Pub. IHES, n\textsuperscript{o}5.
\item[{[H]}] Hartshorne, R., Residues and Duality, Lecture Notes n\textsuperscript{o}20, Springer, 1966.
\item[{[T.F]}] Godement, R., Théorie des faisceaux, Hermann, 1958.
\end{enumerate}

% SGA_NATIVE_CONTINUATION_273
\providecommand{\Ob}{\operatorname{Ob}}
\providecommand{\Mod}{\operatorname{Mod}}
\providecommand{\Coh}{\operatorname{Coh}}
\providecommand{\Loclib}{\operatorname{Loclib}}
\providecommand{\coh}{\mathrm{coh}}
\providecommand{\parf}{\mathrm{parf}}
\providecommand{\naif}{\mathrm{naif}}
\providecommand{\loc}{\mathrm{loc}}
\providecommand{\fin}{\mathrm{fin}}
\providecommand{\clA}{\operatorname{cl}}
\providecommand{\Rg}{\operatorname{rg}}
\providecommand{\varinjlim}{\varinjlim}
\providecommand{\varprojlim}{\varprojlim}
\providecommand{\id}{\operatorname{id}}
\providecommand{\Th}{\operatorname{Th}}
\providecommand{\Hom}{\operatorname{Hom}}
\providecommand{\Dcat}{\mathrm D}
\providecommand{\Kcat}{\mathrm K}
\providecommand{\Kdot}{K_{\bullet}}
\providecommand{\Kprime}{K'}
\providecommand{\naif}{\mathrm{na\ddot{\i}f}}


% SGA_REVIEWED_FRENCH_CONTINUATION_273
% Exposé IV macros
\renewcommand{\Ob}{\operatorname{ob}}
\providecommand{\Mod}{\operatorname{Mod}}
\providecommand{\Coh}{\operatorname{Coh}}
\providecommand{\Loclib}{\operatorname{Loclib}}
\providecommand{\coh}{\mathrm{coh}}
\providecommand{\parf}{\mathrm{parf}}
\providecommand{\naif}{\text{naïf}}
\providecommand{\loc}{\mathrm{loc}}
\providecommand{\fin}{\mathrm{fin}}
\providecommand{\clA}{\operatorname{cl}}
\providecommand{\Rg}{\operatorname{rg}}
\providecommand{\varinjlim}{\varinjlim}
\providecommand{\varprojlim}{\varprojlim}
\renewcommand{\id}{\operatorname{id}}
\providecommand{\Th}{\operatorname{Th}}
\providecommand{\Hom}{\operatorname{Hom}}
\providecommand{\Dcat}{\mathrm D}
\providecommand{\Kcat}{\mathrm K}
\providecommand{\Kdot}{K_{\bullet}}
\providecommand{\Kprime}{K'}
\providecommand{\naif}{\text{naïf}}



% SGA_SOURCE_MACRO_RESTORATION_273: native provide must not shadow the reviewed accented label.
\renewcommand{\naif}{\text{naïf}}

%%%%%%%%%%  scan idx 287  |  folio 274  %%%%%%%%%%
\SGAPrintPage{274}%
\sgaStableTarget{sga6.e04.expose}\section*{\underline{EXPOSÉ IV}\\[3pt] \underline{\underline{GROTHENDIECK GROUPS OF RINGED TOPOI}}}
% The reviewed Exposé IV source renews this local abbreviation with one rule.
\renewcommand{\parf}{\underline{\mathrm{parf}}}
\begin{center}
\underline{by L. ILLUSIE}
\end{center}

\sgaStableTarget{sga6.e04.sec.1}\section*{1. Reminders and generalities on Grothendieck groups}

\sgaStableTarget{sga6.e04.para.1-1}1.1. Let \(C\) be a \uline{triangulated} category. Recall (I 6.3 and SGA 5 VIII 2) that a map \(f\) from \(\Ob C\) to an abelian group \(G\) is called \uline{additive} if \(f(E)=f(E')+f(E'')\) for every distinguished triangle \(E'\longrightarrow E\longrightarrow E''\longrightarrow E'[1]\). The additive maps from \(\Ob C\) to \(G\) form an abelian group, depending functorially on \(G\). The functor thus obtained is represented by an abelian group \(k(C)\) and a universal additive map \(\cl_C:\Ob C\longrightarrow k(C)\) (denoted by \(\cl\) when there is no risk of confusion).

The group \(k(C)\) depends functorially on \(C\) with respect to exact functors. If \(C,C'\) are triangulated categories, two isomorphic exact functors from \(C\) to \(C'\) induce the same homomorphism \(k(C)\to k(C')\); in particular, if \(u:C\to C'\) is an exact functor that is an equivalence of categories, \(k(u):k(C)\to k(C')\) is an isomorphism.

\sgaStableTarget{sga6.e04.lem.1-2}\textbf{Lemma 1.2.} \uline{Let} \(C\) \uline{be a triangulated category, and let} \(L\in\Ob C\).
\begin{enumerate}[label=\alph*)]
\item \uline{For every} \(n\in\mathbb Z\)\uline{, we have}
\[
 \cl(L[n])=(-1)^n\cl(L).
\]
\uline{If} \(L',L''\) \uline{are objects of} \(C\) \uline{such that} \(L\simeq L'\oplus L''\)\uline{, then} \(\cl(L)=\cl(L')+\cl(L'')\).
\item \uline{Suppose there exists} \(n\in\mathbb Z\) \uline{such that} \(\coprod_{i\geq0} L[2ni]\) \uline{is representable. Then} \(\cl(L)=0\).
\end{enumerate}

\uline{Proof.} Assertion a) follows directly from the definitions (cf. SGA 5 VIII 2). For b), it suffices to observe that we have
%%%%%%%%%%  scan idx 288  |  folio 275  %%%%%%%%%%
\SGAPrintPage{275}%
\[
 \coprod_{i\geq0} L[2ni]\simeq
 L\oplus\left(\coprod_{i\geq0}L[2ni]\right)[2n],
\]
and to apply a).

\sgaStableTarget{sga6.e04.lem.1-3}\textbf{Lemma 1.3.} \uline{Let} \(C\) \uline{be a triangulated category, and} \(A\) \uline{a thick subcategory} \(([\mathrm V]\ \mathrm I\ 2.1)\)\uline{. The inclusion and quotient functors define an exact sequence}
\[
 k(A)\longrightarrow k(C)\longrightarrow k(C/A)\longrightarrow 0.
\]

\uline{Proof.} See (SGA 5 VIII 3.1).

\sgaStableTarget{sga6.e04.lem.1-4}\textbf{Lemma 1.4.} \uline{Let} \(A\) \uline{be an additive (resp. abelian) category, and} \(f\) \uline{an additive function} (1.1) \uline{on} \(\Ob K^b(A)\) \uline{(resp.} \(\Ob D^b(A)\)\uline{). For} \(E\in\Ob K^b(A)\) \uline{(resp.} \(D^b(A)\)\uline{) we have}
\[
 f(E)=\sum(-1)^i f(E^i)
\]
\uline{(resp.}
\[
 f(E)=\sum(-1)^i f(H^i(E))
\]).

\uline{Proof.} Left as an exercise to the reader.

\sgaStableTarget{sga6.e04.para.1-5}1.5. Let \(C\) be an \uline{abelian} category and \(A\) a full subcategory of \(C\), stable under finite sums. A map \(f\) from \(\Ob A\) to an abelian group \(G\) is called \uline{additive} if \(\sum(-1)^i f(E^i)=0\) for every acyclic complex \(E\) of \(C\), with bounded degrees and components belonging to \(A\). When \(A\) is stable under kernels of epimorphisms (I 1.2), this amounts to saying that \(f(L)=f(L')+f(L'')\) for every exact sequence \(0\longrightarrow L'\longrightarrow L\longrightarrow L''\longrightarrow0\) of \(C\) whose terms belong to \(A\). The additive maps from \(\Ob A\) to \(G\) form an abelian group depending functorially on \(G\). Denote by \(K^{b,\emptyset}(A)\) the thick subcategory of \(K^b(A)\) formed by the acyclic complexes. It follows easily from (1.3) and (1.4) that the map
\sgaStableTarget{sga6.e04.eq.1-5-1}\begin{equation}\tag{1.5.1}
 \cl_A:\Ob A\longrightarrow k\bigl(K^b(A)/K^{b,\emptyset}(A)\bigr)
\end{equation}
%%%%%%%%%%  scan idx 289  |  folio 276  %%%%%%%%%%
\SGAPrintPage{276}%
induced by \(\cl_{K^b(A)/K^{b,\emptyset}(A)}\) (1.1) is a universal additive map, i.e. the pair \(\left(k\bigl(K^b(A)/K^{b,\emptyset}(A)\bigr),\cl_A\right)\) represents the functor assigning to each abelian group \(G\) the abelian group of additive maps from \(\Ob A\) to \(G\). We shall put
\sgaStableTarget{sga6.e04.eq.1-5-2}\begin{equation}\tag{1.5.2}
 k\bigl(K^b(A)/K^{b,\emptyset}(A)\bigr)=k(A)
\end{equation}
(we shall write \(k(A,C)\) instead of \(k(A)\) when we wish to avoid any ambiguity about the ambient abelian category)\SGAAuthorFootnote{1}{The notation \(k(A)\) was introduced in (I 6.3) to denote the group \(k(K^b(A))\), which is not in general isomorphic to the group defined by (1.5.2). We therefore change convention here: in the present exposé, the notation \(k(A)\) will denote exclusively the group defined by (1.5.2).}.
In particular, for \(A=C\), we have
\sgaStableTarget{sga6.e04.eq.1-5-3}\begin{equation}\tag{1.5.3}
 k(D^b(C))=k(C).
\end{equation}

\sgaStableTarget{sga6.e04.lem.1-6}\textbf{Lemma 1.6.} \uline{Let} \(C\) \uline{be an abelian category, and} \(A\) \uline{a thick subcategory (i.e. a full subcategory stable under kernels, cokernels, and extensions). Denote by} \(D_A(C)\) \uline{the full subcategory of} \(D(C)\) \uline{formed by the objects} \(E\) \uline{such that} \(H^i(E)\in\Ob A\) \uline{for every} \(i\)\uline{. Then} \(D_A(C)\) \uline{is a triangulated subcategory} \(([\mathrm V]\ \mathrm I\ 1.2.3)\) \uline{of} \(D(C)\)\uline{, and the canonical functor} \(D^b(A)\longrightarrow D_A^b(C)=D^b(C)\cap D_A(C)\) \uline{induces an isomorphism}
\[
 k(A)\xrightarrow{\sim} k(D_A^b(C)),
\]
\uline{where} \(k(A)\) \uline{is defined by} (1.5.2).

\uline{Proof.} The first assertion is an immediate consequence of the long exact cohomology sequence. For the second, consider the map \(f:\Ob D_A^b(C)\longrightarrow k(A)\) defined by
\[
 f(E)=\sum(-1)^i\cl_A(H^i(E)),
\]
where \(\cl_A\) is the map (1.5.1). By the long exact cohomology sequence, \(f\) is additive (1.1), hence defines a homomorphism
\[
 \varphi:k(D_A^b(C))\longrightarrow k(A).
\]
%%%%%%%%%%  scan idx 290  |  folio 277  %%%%%%%%%%
\SGAPrintPage{277}%
It follows at once from (1.4) that \(\varphi\) is inverse to the canonical homomorphism \(k(A)\longrightarrow k(D_A^b(C))\).

\sgaStableTarget{sga6.e04.rem.1-6-1}\textbf{Remark 1.6.1.} In the situation of (1.6), the canonical functor \(D^b(A)\longrightarrow D_A^b(C)\) is not in general an equivalence of categories, as shown by the example where \(S\) is a scheme, \(C\) the category of \(\underline{O}_S\)-Modules, and \(A\) the category of quasi-coherent \(\underline{O}_S\)-Modules (II Ap. I 0.3). See, however, the criterion (II 2.1.3).

\sgaStableTarget{sga6.e04.def.1-7}\textbf{Definition 1.7.} \uline{Let} \(C\) \uline{be a triangulated (resp. abelian) category, and} \(C'\) \uline{a subset of} \(\Ob C\)\uline{. We shall say that} \(C'\) \uline{is an} exact \uline{subset (cf.} EGA III 3.1.1\uline{) if, for every distinguished triangle} \(E'\to E\to E''\to E'[1]\) \uline{(resp. every exact sequence} \(0\to E'\to E\to E''\to0\)\uline{) such that two of the terms} \(E',E,E''\) \uline{belong to} \(C'\)\uline{, the third also belongs to it. We shall say that} \(C'\) \uline{is} dense \uline{in} \(\Ob C\) \uline{if} \(\Ob C\) \uline{is the smallest exact subset of} \(\Ob C\) \uline{containing} \(C'\).

\sgaStableTarget{sga6.e04.lem.1-8}\textbf{Lemma 1.8.} \uline{Let} \(C\) \uline{be a triangulated (resp. abelian) category, and} \(C'\) \uline{a dense subset of} \(\Ob C\)\uline{. Then the elements} \(\cl_C(E)\)\uline{, for} \(E\in C'\)\uline{, generate} \(k(C)\).

\uline{Proof.} Put \(C'=C_0\), and, for \(i\geq1\), define \(C_i\) inductively as the set of objects of \(C\) appearing in a distinguished triangle (resp. a short exact sequence) two of whose terms belong to \(C_{i-1}\). The \(C_i\) form an increasing filtration of \(\Ob C\), which is exhaustive by hypothesis. For \(i\geq0\), denote by \(k(C)_i\) the subgroup of \(k(C)\) generated by the elements \(\cl_C(E)\) for \(E\in C_i\). It is immediate from the definitions that, for every \(i\geq0\), \(k(C)_i=k(C)_{i+1}\), hence \(k(C)_i=k(C)_0\). Now let \(x\in k(C)\);
%%%%%%%%%%  scan idx 291  |  folio 278  %%%%%%%%%%
\SGAPrintPage{278}%
as is easily seen, we can write \(x=\cl(E')-\cl(E'')\) for \(E',E''\in\Ob C\) (in the first case, we can even arrange that \(E''=0\)). Since the filtration \((C_i)\) is exhaustive, there is an integer \(i\geq0\) such that \(E'\) and \(E''\) belong to \(C_i\); then \(x\in k(C)_i=k(C)_0\), as required.

\sgaStableTarget{sga6.e04.exercise.1-9}\textbf{Exercise 1.9.} Let \(C\) be an abelian category, and \(A\) a full subcategory stable under finite sums, subobjects, and quotients, such that \(\Ob A\) is dense in \(\Ob C\). Then the homomorphism \(k(A,C)\longrightarrow k(C)\) defined by the additive function \(E\mapsto\cl_C(E)\) on \(\Ob A\) (cf. (1.5)) is an isomorphism.

\uline{Hint}: Put \(A=A_0\), and define \(A_i\), for \(i\geq1\), as the full subcategory of \(C\) formed by the objects of \(C\) that are extensions of two objects of \(C_{i-1}\). Prove that \(A_i\) satisfies the same hypotheses as \(A\), then show that the canonical homomorphism \(\varinjlim k(A_i,C)\longrightarrow k(C)\) is an isomorphism (the desired conclusion follows at once).

\sgaStableTarget{sga6.e04.sec.2}\section*{2. The functors \(K_{\bullet}\) and \(K^{\bullet}\) of a ringed topos}

\sgaStableTarget{sga6.e04.para.2-1}2.1. Let \((X,A)\) be a ringed topos. Recall the following notation:

\(D^b({}_A X)_{\coh}\) = category of pseudo-coherent complexes with bounded cohomology of left \(A\)-Modules (I 2.15).

\(D({}_A X)_{\uparf}\) = category of complexes of left \(A\)-Modules of finite perfect amplitude (I 5.8.4).

\(\Loclib({}_A X)\) = category of left \(A\)-Modules locally direct summands of free Modules of finite type (II 2.0).

\(\Coh({}_A X)\) = category of coherent left \(A\)-Modules (relative to the category of free \(A\)-Modules of finite type) (I 3.1).

%%%%%%%%%%  scan idx 292  |  folio 279  %%%%%%%%%%
\SGAPrintPage{279}%
The last two categories are full subcategories of the category \(\Mod({}_A X)\) of left \(A\)-Modules, and the first two are full triangulated subcategories of \(D({}_A X)\) (\(= D(\Mod({}_A X))\)).

\sgaStableTarget{sga6.e04.def.2-2}\textbf{Definition 2.2.} \uline{Let} \((X,A)\) \uline{be a ringed topos. Unless otherwise stated, we shall put}
\[
 \begin{gathered}
  K_\bullet({}_A X)=k(D^b({}_A X)_{\coh}),\\
  K^{\bullet}({}_A X)=k(D({}_A X)_{\uparf}),
 \end{gathered}
\]
\uline{where} \(k(\ )\) \uline{is defined in} (1.1)\uline{, and}
\[
 \begin{gathered}
  K_\bullet({}_A X)_{\naif}=k(\Coh({}_A X)),\\
  K^{\bullet}({}_A X)_{\naif}=k(\Loclib({}_A X)),
 \end{gathered}
\]
\uline{where} \(k(\ )\) \uline{is defined by} (1.5.2)\uline{, the ambient abelian category being} \(\Mod({}_A X)\)\SGAAuthorFootnote{1}{We shall allow ourselves to omit the index \(A\) when no confusion can result. Sometimes we shall also drop the index ``naïf'', but shall indicate this explicitly.}\uline{. We shall say that} \(K_\bullet({}_A X)\) \uline{(resp.} \(K^{\bullet}({}_A X)\)\uline{) is the Grothendieck group of pseudo-coherent (resp. perfect) complexes of left} \(A\)\uline{-Modules, and that} \(K_\bullet({}_A X)_{\naif}\) \uline{(resp.} \(K^{\bullet}({}_A X)_{\naif}\)\uline{) is the Grothendieck group of coherent (resp. locally free) left} \(A\)\uline{-Modules.}

\sgaStableTarget{sga6.e04.rem.2-3}\textbf{Remarks 2.3.}
\begin{enumerate}[label=\alph*)]
\item One can define variants of \(K_\bullet(\ )\) and \(K^{\bullet}(\ )\) by putting
\[
 K_\bullet^{\loc}({}_A X)=k(D^{\loc b}({}_A X)_{\coh})
\]
(where ``\(\loc b\)'' means ``with locally bounded cohomology''), and
\[
 K^{\bullet}_{\loc}({}_A X)=k(D({}_A X)_{\mathrm{parf}}).
\]
These groups coincide with those defined in (2.2) when the final object of \(X\) is quasi-compact, but in general are less suited to intersection calculations.
%%%%%%%%%%  scan idx 293  |  folio 280  %%%%%%%%%%
\item\SGAPrintPage{280}%
The group \(K_\bullet({}_A X)_{\naif}\) is practically of interest only when \(A\) is coherent (otherwise there may well be very few coherent sheaves \ldots).
\item Considering the \(k\) of triangulated categories that are ``too large'' is of no interest. Thus, from (1.2 b)) one easily deduces the relations
\[
 k(D({}_A X))=k(D^-({}_A X))=k(D^+({}_A X))=k(D^b({}_A X))=k(D^-({}_A X)_{\coh})=0
\]
(for the last equality, observe that \(L\in\Ob D^-({}_A X)_{\coh}\) implies \(\coprod_{i\geq0}L[2i]\in\Ob D^-({}_A X)_{\coh}\)).
\end{enumerate}

\sgaStableTarget{sga6.e04.para.2-4}2.4. \uline{Case of a coherent sheaf of rings}. Let \((X,A)\) be a ringed topos, with \(A\) supposed \uline{coherent} as a left \(A\)-Module (I 3.1 and 3.6). Every coherent left \(A\)-Module is a pseudo-coherent complex (I 3.5), so we have a fully faithful functor
\sgaStableTarget{sga6.e04.eq.2-4-1}\begin{equation}\tag{2.4.1}
 \Coh({}_A X)\hookrightarrow D^b({}_A X)_{\coh}.
\end{equation}
Now (I 3.4) \(\Coh({}_A X)\) is a thick subcategory of \(\Mod({}_A X)\), and, by (I 3.5), \(D^b({}_A X)_{\coh}\) is the full subcategory of \(D^b({}_A X)\) formed by the objects \(L\) such that \(H^i(L)\in\Ob\Coh({}_A X)\) for every \(i\). Consequently, it follows from (1.6) that \uline{the homomorphism}
\sgaStableTarget{sga6.e04.eq.2-4-2}\begin{equation}\tag{2.4.2}
 K_\bullet({}_A X)_{\naif}\longrightarrow K_\bullet({}_A X)
\end{equation}
\uline{defined by} (2.4.1) \uline{is an isomorphism}.

\sgaStableTarget{sga6.e04.para.2-5}2.5. \uline{The homomorphism} \(K^{\bullet}\to K_\bullet\). Let \((X,A)\) be a ringed topos. The inclusion functor \(D({}_A X)_{\uparf}\hookrightarrow D^b({}_A X)_{\coh}\) defines a homomorphism
\sgaStableTarget{sga6.e04.eq.2-5-1}\begin{equation}\tag{2.5.1}
 K^{\bullet}({}_A X)\longrightarrow K_\bullet({}_A X).
\end{equation}
This is in general neither injective nor surjective. However, make the following hypotheses:
\begin{enumerate}[label=(\roman*)]
%%%%%%%%%%  scan idx 294  |  folio 281  %%%%%%%%%%
\item\SGAPrintPage{281}%
the final object of \(X\) is quasi-compact;
\item \(X\) has enough points, and, for every point \(x\) of \(X\), every left \(A_x\)-module has finite tor-dimension.
\end{enumerate}
Then, by (I 5.10), \(D({}_A X)_{\parf}=D^b({}_A X)_{\coh}\), and the homomorphism (2.5.1) is an isomorphism.

Hypotheses (i) and (ii) are satisfied, for example, by a compact space ringed by regular local rings.

\sgaStableTarget{sga6.e04.para.2-6}2.6. \uline{Pairings between} \(K^{\bullet}\) \uline{and} \(K^{\bullet}\) (\uline{noncommutative case}). Let
\[
 f:(X',A')\longrightarrow (X,A)
\]
be a morphism of ringed topoi. Denote by \(\Mod({}_{A'} X'_A)\) the category of bi-Modules, left over \(A'\) and right over \(f^{-1}(A)\), and by \(\Loclib_{A'}({}_{A'} X'_A)\) the full subcategory formed by the objects belonging to \(\Loclib({}_{A'} X')\). By analogy with (2.2), we shall put
\[
 K^{\bullet}_{\naif}({}_{A'} X'_A)_{A'}=k(\Loclib_{A'}({}_{A'} X'_A)).
\]
The tensor product
\[
 \otimes_A:\Mod({}_{A'} X'_A)\times\Mod({}_A X)\longrightarrow \Mod({}_{A'} X')
\]
\[
 (E',E)\longmapsto E'\otimes_{f^{-1}(A)}f^{-1}(E)
\]
induces a functor
\sgaStableTarget{sga6.e04.eq.2-6-1}\begin{equation}\tag{2.6.1}
 \Loclib_{A'}({}_{A'} X'_A)\times\Loclib({}_A X)\longrightarrow\Loclib({}_{A'} X').
\end{equation}
The map \((E',E)\mapsto\cl(E'\otimes_A E)\) deduced from (2.6.1) is biadditive in the sense of (1.5), and therefore defines a homomorphism
\sgaStableTarget{sga6.e04.eq.2-6-2}\begin{equation}\tag{2.6.2}
 K^{\bullet}_{\naif}({}_{A'} X'_A)_{A'}\otimes K^{\bullet}_{\naif}({}_A X)\longrightarrow K^{\bullet}_{\naif}({}_{A'} X').
\end{equation}
Denote by \(D({}_{A'} X'_A)\) the derived category of \(\Mod({}_{A'} X'_A)\), and by \(D({}_{A'} X'_A)_{A'\text{-}\parf}\) the full subcategory formed by the complexes of finite perfect amplitude rel. to \(A'\) (i.e. as objects of \(D({}_{A'} X')\)), and put, as in (2.2),
\[
 K^{\bullet}({}_{A'} X'_A)_{A'}=k(D({}_{A'} X'_A)_{A'\text{-}\parf}).
\]
%%%%%%%%%%  scan idx 295  |  folio 282  %%%%%%%%%%
\SGAPrintPage{282}%
The inclusion functors plainly define homomorphisms
\sgaStableTarget{sga6.e04.eq.2-6-3-1}\sgaStableTarget{sga6.e04.eq.2-6-3-2}\sgaStableTarget{sga6.e04.eq.2-6-3-3}\begin{align}
 K^{\bullet}_{\naif}({}_{A'} X'_A)_{A'}&\longrightarrow K^{\bullet}({}_{A'} X'_A)_{A'}, \tag{2.6.3.1}\\
 K^{\bullet}_{\naif}({}_A X)&\longrightarrow K^{\bullet}({}_A X), \tag{2.6.3.2}\\
 K^{\bullet}_{\naif}({}_{A'} X')&\longrightarrow K^{\bullet}({}_{A'} X'). \tag{2.6.3.3}
\end{align}
On the other hand, we know (I 4.18, 4.19) that the derived functor of \(\otimes_A\),
\[
 \otimes_A^{\mathrm L}:D^-({}_{A'} X'_A)\times D^-({}_A X)\longrightarrow D^-({}_{A'} X'),
\]
induces an (exact) functor
\sgaStableTarget{sga6.e04.eq.2-6-4}\begin{equation}\tag{2.6.4}
 D({}_{A'} X'_A)_{A'\text{-}\parf}\times D({}_A X)_{\parf}
 \longrightarrow D({}_{A'} X')_{\parf}.
\end{equation}
The latter defines a homomorphism
\sgaStableTarget{sga6.e04.eq.2-6-5}\begin{equation}\tag{2.6.5}
 K^{\bullet}({}_{A'} X'_A)_{A'}\otimes K^{\bullet}({}_A X)\longrightarrow K^{\bullet}({}_{A'} X').
\end{equation}
Since the functor (2.6.1) is induced by (2.6.4), the homomorphisms (2.6.2) and (2.6.5) are compatible with the homomorphisms (2.6.3.1), (2.6.3.2), and (2.6.3.3); in other words, we have a commutative diagram
\sgaStableTarget{sga6.e04.eq.2-6-6}\[
\begin{tikzcd}[column sep=tiny]
K^{\bullet}_{\naif}({}_{A'} X'_A)_{A'}\otimes K^{\bullet}_{\naif}({}_A X)\arrow[r]\arrow[d]
&K^{\bullet}_{\naif}({}_{A'} X')\arrow[d]\\
K^{\bullet}({}_{A'} X'_A)_{A'}\otimes K^{\bullet}({}_A X)\arrow[r]
&K^{\bullet}({}_{A'} X'),
\end{tikzcd}
\tag{2.6.6}
\]
where the upper and lower horizontal arrows are respectively (2.6.2) and (2.6.5), and the vertical arrows are defined by (2.6.3.1), (2.6.3.2), and (2.6.3.3).

%%%%%%%%%%  scan idx 296  |  folio 283  %%%%%%%%%%
\SGAPrintPage{283}%
\sgaStableTarget{sga6.e04.para.2-7}2.7. \uline{Contravariant law and ring structure on} \(K^{\bullet}\) We shall spell out some important special cases of the pairings (2.6) (their direct study is moreover immediate).
\begin{enumerate}[label=\alph*)]
\item Let \(f:(X',A')\longrightarrow (X,A)\) be a morphism of ringed topoi. The functor \(Lf^*:D^-({}_A X)\longrightarrow D^-({}_{A'} X')\) sends (I 4.19.1) \(D({}_A X)_{\parf}\) to \(D({}_{A'} X')_{\parf}\) (resp. \(\Loclib({}_A X)\) to \(\Loclib({}_{A'} X')\)), and defines a homomorphism
\sgaStableTarget{sga6.e04.eq.2-7-1}\begin{equation}\tag{2.7.1}
 f^*:K^{\bullet}({}_A X)\longrightarrow K^{\bullet}({}_{A'} X')
\end{equation}
(resp.
\sgaStableTarget{sga6.e04.eq.2-7-2}\begin{equation}\tag{2.7.2}
 f^*:K^{\bullet}({}_A X)_{\naif}\longrightarrow K^{\bullet}({}_{A'} X')_{\naif}
\end{equation}
), which can also be interpreted as multiplication by \(\cl(A')\) in (2.6.5) (resp. (2.6.2)). We have a commutative diagram (a special case of (2.6.6))
\[
\begin{tikzcd}
K^{\bullet}({}_A X)_{\naif}\arrow[r,"f^*"]\arrow[d]
&K^{\bullet}({}_{A'} X')_{\naif}\arrow[d]\\
K^{\bullet}({}_A X)\arrow[r,"f^*"]
&K^{\bullet}({}_{A'} X'),
\end{tikzcd}
\]
where the left and right vertical arrows are respectively (2.6.3.2) and (2.6.3.3).

Let \(g:(X,A)\to(X'',A'')\) be a morphism of ringed topoi. The transitivity of inverse-image functors implies that the homomorphisms of type (2.7.1) (resp. 2.7.2) satisfy the relation
\sgaStableTarget{sga6.e04.eq.2-7-3}\begin{equation}\tag{2.7.3}
 (gf)^*=f^*g^*.
\end{equation}
Thus \(K^{\bullet}({}_A X)\) (resp. \(K^{\bullet}({}_A X)_{\naif}\)) \uline{is a contravariant functor of} \((X,A)\).

\item Let \((X,\underline{O}_X)\) be a commutatively ringed topos. By the associativity and commutativity of the tensor product, \uline{the homomorphism} (2.6.5) (resp. (2.6.2)) (where we put \(A=\underline{O}_X\), \((X',A')=(X,A)\), \(f=\Id\)) \uline{endows} \(K^{\bullet}(X)\) (resp. \(K^{\bullet}(X)_{\naif}\)) \uline{with the structure of a commutative ring with unit}
%%%%%%%%%%  scan idx 297  |  folio 284  %%%%%%%%%%
\SGAPrintPage{284}%
\((1=\cl(\underline{O}_X))\), \uline{and the canonical homomorphism} (2.6.3.2)
\[
 K^{\bullet}(X)_{\naif}\longrightarrow K^{\bullet}(X)
\]
\uline{is a homomorphism of unital rings}.

Let \(f:(X,\underline{O}_X)\to(Y,\underline{O}_Y)\) be a morphism of topoi ringed by commutative rings. Then compatibility of the tensor product with inverse images shows that \uline{the homomorphism}
\[
 f^*:K^{\bullet}(Y)\longrightarrow K^{\bullet}(X)
\]
(resp.
\[
 f^*:K^{\bullet}(Y)_{\naif}\longrightarrow K^{\bullet}(X)_{\naif}
\]
) \uline{is a homomorphism of unital rings}.
\end{enumerate}

\sgaStableTarget{sga6.e04.para.2-8}2.8. \uline{Augmentation of} \(K^{\bullet}\) \uline{to} \(H^0(\ ,\mathbb Z)\). Let \((X,\underline{O}_X)\) be a locally ringed topos (I 2.15.1). Recall that we have defined (I 6.11 to 6.14) a homomorphism of (unital) rings
\[
 i:H^0(X,\mathbb Z)\longrightarrow K^{\bullet}(X)_{\naif},
\]
thus giving \(K^{\bullet}(X)_{\naif}\) (and consequently also \(K^{\bullet}(X)\) via the canonical homomorphism \(K^{\bullet}(X)_{\naif}\to K^{\bullet}(X)\)) the structure of an algebra (with unit) over \(H^0(X,\mathbb Z)\). Moreover (loc.cit.), the algebra \(K^{\bullet}(X)_{\naif}\) (resp. \(K^{\bullet}(X)\)) is augmented to \(H^0(X,\mathbb Z)\) by the \uline{rank}, and the canonical homomorphism \(K^{\bullet}(X)_{\naif}\to K^{\bullet}(X)\) is a homomorphism of augmented algebras.

Let \(f:(X,\underline{O}_X)\to(Y,\underline{O}_Y)\) be a morphism of locally ringed topoi. The homomorphisms \(f^*:K^{\bullet}(Y)\longrightarrow K^{\bullet}(X)\) and \(f^*:K^{\bullet}(Y)_{\naif}\longrightarrow K^{\bullet}(X)_{\naif}\) are compatible (I 6.12, 6.13) with the preceding augmented algebra structures; this is expressed by the commutativity of the diagram (where \(\Rg\) denotes rank):
%%%%%%%%%%  scan idx 298  |  folio 285  %%%%%%%%%%
\SGAPrintPage{285}%
\[
\begin{tikzcd}[column sep=tiny]
H^0(Y,\mathbb Z)\arrow[r,"i"]\arrow[d]
&K^{\bullet}(Y)_{\naif}\arrow[r]\arrow[d]
&K^{\bullet}(Y)\arrow[r,"\Rg"]\arrow[d]
&H^0(Y,\mathbb Z)\arrow[d]\\
H^0(X,\mathbb Z)\arrow[r,"i"]
&K^{\bullet}(X)_{\naif}\arrow[r]
&K^{\bullet}(X)\arrow[r,"\Rg"]
&H^0(X,\mathbb Z).
\end{tikzcd}
\]

\textsuperscript{*}We shall see later in the Seminar (V 3.9.1 and VI) that \(K^{\bullet}(X)_{\naif}\) can be given a richer structure than that of an algebra augmented to \(H^0(X,\mathbb Z)\), namely that of a \(\lambda\)-algebra augmented to \(H^0(X,\mathbb Z)\) regarded as a binomial ring (V 2.8), this structure being compatible with inverse-image homomorphisms. Deligne [1] recently proved that \(K^{\bullet}(X)\) can be given an analogous structure (at least when \(X\) has characteristic \(p\geq0\)) and that the canonical homomorphism \(K^{\bullet}(X)_{\naif}\longrightarrow K^{\bullet}(X)\) is a homomorphism of augmented \(\lambda\)-algebras, thus answering a conjecture of Grothendieck (XIV 1) in the affirmative.\textsubscript{*}

\sgaStableTarget{sga6.e04.para.2-9}2.9. \uline{Comparison of} \(K^{\bullet}\) \uline{and} \(K^{\bullet}_{\naif}\). Let \((X,A)\) be a ringed topos. In general there is no reason for the canonical homomorphism
\[
 K^{\bullet}(X)_{\naif}\longrightarrow K^{\bullet}(X)
\]
to be injective or surjective. However, the ``criteria for the existence of global resolutions'' of Exposé II show that it is an isomorphism in the following cases:
\begin{enumerate}[label=(\roman*)]
\item \(X\) is the Zariski topos of a divisorial scheme, separated or noetherian (and \(A=\underline{O}_X\)) (II 2.2.9);
\item \(X\) is the topos defined by a compact space and \(A\) is a soft sheaf (II 2.3.2);
\item \(X\) is a Stein topos whose final object is quasi-compact and \(A=\underline{O}_X\) (II 2.4.5).
\end{enumerate}

%%%%%%%%%%  scan idx 299  |  folio 286  %%%%%%%%%%
\SGAPrintPage{286}%
\sgaStableTarget{sga6.e04.para.2-10}2.10. \uline{Pairing between} \(K^{\bullet}\) \uline{and} \(K_{\bullet}\). Let \((X,\underline{O}_X)\) be a commutatively ringed topos. We know (I 5.8.4 b)) that there is a pairing
\[
 \otimes_A^{\mathrm L}:D(X)_{\uparf}\times D^b(X)_{\coh}
 \longrightarrow D^b(X)_{\coh}.
\]
This defines a homomorphism
\sgaStableTarget{sga6.e04.eq.2-10-1}\begin{equation}\tag{2.10.1}
 K^{\bullet}(X)\otimes K_{\bullet}(X)\longrightarrow K_{\bullet}(X).
\end{equation}
Thanks to the associativity of the tensor product, we see that (2.10.1) \uline{endows} \(K_{\bullet}(X)\) \uline{with the structure of a} \(K^{\bullet}(X)\)-\uline{module}. More generally, if \(f:(X,\underline{O}_X)\longrightarrow(Y,\underline{O}_Y)\) is a morphism of commutatively ringed topoi, we can regard \(K_{\bullet}(X)\) as a \(K^{\bullet}(Y)\)-module via the homomorphism \(f^*:K^{\bullet}(Y)\longrightarrow K^{\bullet}(X)\).

\sgaStableTarget{sga6.e04.para.2-11}2.11. \uline{Covariant law on} \(K_{\bullet}\) \uline{and projection formula}. Let \((X,A)\) be a ringed topos. Unlike \(K^{\bullet}({}_A X)\), the group \(K_\bullet({}_A X)\) has no natural variance with respect to \((X,A)\). However, in Algebraic Geometry and Analytic Geometry, \(K_{\bullet}(X)\) tends to be \uline{covariant} in \(X\) for \uline{proper} morphisms. Let us make this precise:

\sgaStableTarget{sga6.e04.para.2-11-1}2.11.1. Let \(f:X\to Y\) be a proper and pseudo-coherent morphism (III 1.2) of quasi-compact schemes. Then, subject to the validity of the finiteness conjecture for \(f\) (III 2.1), hence in particular (III 2.2) when \(Y\) is noetherian (in which case the pseudo-coherence hypothesis on \(f\) is automatically satisfied (III 1.3)) or when \(f\) is projective, the functor \(Rf_*\) induces, by (III 2.6), a functor
\[
 Rf_*:D^b(X)_{\coh}\longrightarrow D^b(Y)_{\coh}.
\]
This defines a homomorphism
\sgaStableTarget{sga6.e04.eq.2-11-1-1}\begin{equation}\tag{2.11.1.1}
 f_*:K_{\bullet}(X)\longrightarrow K_{\bullet}(Y)
\end{equation}
%%%%%%%%%%  scan idx 300  |  folio 287  %%%%%%%%%%
\SGAPrintPage{287}%
This homomorphism is \(K^{\bullet}(Y)\)-linear (cf. (2.10)); in other words, for \(y\in K^{\bullet}(Y)\) and \(x\in K_{\bullet}(X)\), we have the so-called \uline{projection formula}
\sgaStableTarget{sga6.e04.eq.2-11-1-2}\begin{equation}\tag{2.11.1.2}
 f_*(yx)=y f_*(x).
\end{equation}
Indeed, if \(x=\cl(E)\) for \(E\in\Ob D^b(X)_{\coh}\) and \(y=\cl(M)\) for \(M\in\Ob D(Y)_{\uparf}\), we have (III 3.7)
\[
 Rf_*(E)\otimes^{\mathrm L}M\simeq
 Rf_*(E\otimes_Y^{\mathrm L}M).
\]
Finally, if \(g:Y\to Z\) is a morphism satisfying the same hypotheses as \(f\) and for which the finiteness conjecture (III 2.1) holds, the transitivity of direct-image functors gives, for the homomorphisms (2.11.1.1), the formula
\sgaStableTarget{sga6.e04.eq.2-11-1-3}\begin{equation}\tag{2.11.1.3}
 g_*f_*=(gf)_*.
\end{equation}

\sgaStableTarget{sga6.e04.para.2-11-2}2.11.2. Let \(f:X\to Y\) be a proper morphism of finite-dimensional analytic spaces. Then, by Grauert's finiteness theorem [2], \(Rf_*\) sends \(D^b(X)_{\coh}\) to \(D^b(Y)_{\coh}\), and therefore defines a homomorphism
\sgaStableTarget{sga6.e04.eq.2-11-2-1}\begin{equation}\tag{2.11.2.1}
 f_*:K_{\bullet}(X)\longrightarrow K_{\bullet}(Y).
\end{equation}
Like its algebraic analogue (2.11.1.1), this is \(K^{\bullet}(Y)\)-linear, i.e. formula (2.11.1.2) is valid, as follows from the projection formula at the level of derived categories (III 5.8 b)). We likewise have the transitivity formula (2.11.1.3).

\sgaStableTarget{sga6.e04.para.2-11-3}2.11.3. Once we leave the algebraic or analytic setting, this covariance of \(K_{\bullet}\) for proper morphisms disappears. Thus, let \(f:X\to Y\) be a proper morphism of topological spaces ringed by the sheaf of complex continuous functions. Then, for \(E\in\Ob D^b(X)_{\coh}\), in general we do not have \(Rf_*(E)\in\Ob D^b(Y)_{\coh}\), as is already evident for \(E=\underline{O}_X\) and \(Y\) reduced to a point. In fact, definition (2.2) of \(K_{\bullet}\)
%%%%%%%%%%  scan idx 301  |  folio 288  %%%%%%%%%%
\SGAPrintPage{288}%
appears especially suited to the needs of Algebraic Geometry and Analytic Geometry, and a good definition of the \(K_{\bullet}\) of a topological space remains to be found. It seems moreover that we shall have to leave the setting of \(\underline{O}_X\)-linear complexes and define, for a morphism \(f:X\to Y\) (of topological spaces, resp. of mixed varieties (II App. II), etc.), a suitable class of complexes of topological vector sheaves on \(X\), which should include, for example, when \(f\) is the structural morphism of a mixed variety, the elliptic complexes on \(X\) relative to \(Y\). This is suggested by the finiteness theorem (II App. II), as well as by a similar theorem (where \(Y\) is supposed reduced to a point) proved by Grothendieck in [3].

\sgaStableTarget{sga6.e04.para.2-12}2.12. \uline{Unusual variances of} \(K^{\bullet}\) \uline{and} \(K_{\bullet}\). Let \(f:(X,A)\longrightarrow(Y,B)\) be a morphism of finite tor-dimension of ringed topoi (III 3.1). The functor \(Lf^*:D^-({}_B Y)\longrightarrow D^-({}_A X)\) sends \(D^b({}_B Y)\) to \(D^b({}_A X)\), hence (I 2.16.1) \(D^b({}_B Y)_{\coh}\) to \(D^b({}_A X)_{\coh}\), and therefore defines a homomorphism
\sgaStableTarget{sga6.e04.eq.2-12-1}\begin{equation}\tag{2.12.1}
 f^*:K_\bullet({}_B Y)\longrightarrow K_\bullet({}_A X).
\end{equation}
If \(A\) and \(B\) are commutative, commutation of inverse image with the tensor product shows that the homomorphism (2.12.1) is compatible with the module structures of \(K_{\bullet}(Y)\) and \(K_{\bullet}(X)\) over \(K^{\bullet}(Y)\) and \(K^{\bullet}(X)\) respectively (2.10); in other words, for \(a\in K^{\bullet}(Y)\) and \(y\in K_{\bullet}(Y)\), we have the formula
\sgaStableTarget{sga6.e04.eq.2-12-2}\begin{equation}\tag{2.12.2}
 f^*(ay)=a f^*(y).
\end{equation}
If \(g:(Y,B)\to(Z,C)\) is a morphism of finite tor-dimension of ringed topoi, transitivity of inverse-image functors implies that the homomorphisms (2.12.1) satisfy the formula:
\sgaStableTarget{sga6.e04.eq.2-12-1-1}\begin{equation}\tag{2.12.1.1}
 f^*g^*=(gf)^*.
\end{equation}

%%%%%%%%%%  scan idx 302  |  folio 289  %%%%%%%%%%
\SGAPrintPage{289}%
Suppose now that we are in situation (2.11.1) (resp. (2.11.2)), with \(f\) of finite tor-dimension. Then (III 4.8.1 and 5.8 b)) the functor \(Rf_*\) sends \(D(X)_{\uparf}\) to \(D(Y)_{\uparf}\), and therefore defines a homomorphism
\sgaStableTarget{sga6.e04.eq.2-12-3}\begin{equation}\tag{2.12.3}
 f_*:K^{\bullet}(X)\longrightarrow K^{\bullet}(Y).
\end{equation}
This is not in general a ring homomorphism. However, for \(y\in K_{\bullet}(Y)\) (resp. \(K^{\bullet}(Y)\)) and \(x\in K^{\bullet}(X)\), we have the formula
\sgaStableTarget{sga6.e04.eq.2-12-4}\begin{equation}\tag{2.12.4}
 f_*(f^*(y)x)=y f_*(x),
\end{equation}
where \(f^*\) is the homomorphism (2.12.1) (resp. (2.7.1)). The proof is analogous to that of (2.11.1.2). If \(g:Y\to Z\) is a morphism satisfying the same hypotheses as \(f\), the homomorphisms (2.12.3) satisfy the transitivity formula:
\sgaStableTarget{sga6.e04.eq.2-12-3-1}\begin{equation}\tag{2.12.3.1}
 g_*f_*=(gf)_*.
\end{equation}

The homomorphism (2.12.3) is one of the essential ingredients of the Riemann--Roch theorem, in both Algebraic Geometry and Analytic Geometry (for indications in the analytic case, see (Exposé 0, 5)).

Suppose that \(X\) and \(Y\) are separated and divisorial schemes (II 2.2.5) (as is the case, for example, when \(X\) and \(Y\) possess ample invertible sheaves). Then (2.9 (i)) \(K^{\bullet}(X)\) (resp. \(K^{\bullet}(Y)\)) is identified with \(K^{\bullet}(X)_{\naif}\) (resp. \(K^{\bullet}(Y)_{\naif}\)), and (2.12.3) therefore defines a homomorphism
\sgaStableTarget{sga6.e04.eq.2-12-3.occ2}\begin{equation}\tag{2.12.3'}
 f_*:K^{\bullet}(X)_{\naif}\longrightarrow K^{\bullet}(Y)_{\naif}.
\end{equation}
As already explained in (Exposé 0, 3.1), the definition of (2.12.3') was by no means evident a priori. Indeed, from the standpoint of our Seminar, it is one of the main motivations for the sorites of Exposé I through Exposé IV.

%%%%%%%%%%  scan idx 303  |  folio 290  %%%%%%%%%%
\SGAPrintPage{290}%
\sgaStableTarget{sga6.e04.sec.3}\section*{3. Supplements on the Grothendieck groups of schemes\protect\SGAAuthorFootnote{1}{The results of this section will not be used in Exposés V through VIII.}}

\sgaStableTarget{sga6.e04.prop.3-1-0}\textbf{Proposition 3.1.0.} Let
\[
\begin{tikzcd}
X \arrow[d,"f"'] & X' \arrow[l,"g"']\arrow[d,"f'"]\\
Y & Y' \arrow[l,"h"']
\end{tikzcd}
\]
be a Cartesian square of schemes, where \(Y\) is quasi-compact and \(f\) quasi-compact and quasi-separated, and let \(E\in\Ob D^b(X)\) be a complex with quasi-coherent cohomology. Suppose \(X\) and \(Y'\) tor-independent over \(Y\) (III 1.5). Suppose moreover either \(h\) of finite tor-dimension or \(E\) of finite flat amplitude rel. to \(f\). Under these conditions, the base-change arrow (cf. SGA 4 XVII 4)
\[
 Lh^*\,Rf_*E\longrightarrow Rf'_*\,Lg^*E
\]
is defined and is an isomorphism.

\uline{Proof}: To define the arrow, we use the trivial duality formula (SGA 4 XVII 2.3) for \(h\), which says that there is a canonical functorial isomorphism
\begin{equation}\tag{$*$}
 \Hom(Lh^*L,M)\xrightarrow{\sim}\Hom(L,Rh_*M)
\end{equation}
for \(L\in\Ob D^-(Y)\), \(M\in\Ob D^+(Y')\). Using (III 3.5.2) and (III 3.7.1), we easily verify that \(Lg^*E\in\Ob D^b(X')\) and \(Rf_*E\in\Ob D^-(Y)\), so, applying (*) with \(L=Rf_*E\) and \(M=Rf'_*Lg^*E\), we are reduced to defining an arrow
\[
 Rf_*E\longrightarrow Rf_*Rg_*Lg^*E.
\]
For this we take the image under \(Rf_*\) of the adjunction arrow
\[
 E\longrightarrow Rg_*Lg^*E,
\]
which comes from the trivial duality formula for \(g\). The base-change arrow
%%%%%%%%%%  scan idx 304  |  folio 291  %%%%%%%%%%
\SGAPrintPage{291}%
is therefore defined; it remains to show that it is an isomorphism. For this, we may first suppose \(Y\) and \(Y'\) affine. Then, by the Leray spectral sequence for an open covering of \(X\), we reduce to the case where \(X\) and \(X'\) are affine. We therefore have a cocartesian square of rings
\[
\begin{tikzcd}
B\arrow[r] & B'\\
A\arrow[u]\arrow[r] & A'\arrow[u],
\end{tikzcd}
\]
and we may suppose that \(E\) is defined by a bounded-above complex of \(B\)-modules. The tor-independence of \(A'\) and \(B\) over \(A\), which is written
\[
 B\otimes_A^{\mathrm L}A'\xrightarrow{\sim}B',
\]
then implies that in \(D(A')\) we have
\[
 E\otimes_A^{\mathrm L}A'\xrightarrow{\sim}E\otimes_B^{\mathrm L}B',
\]
whence the proposition.

Translating (3.1.0) into the language of Grothendieck groups, we obtain:

\sgaStableTarget{sga6.e04.prop.3-1-1}\textbf{Proposition 3.1.1.} \uline{Let}
\[
\begin{tikzcd}
X \arrow[d,"f"'] & X' \arrow[l,"g"']\arrow[d,"f'"]\\
Y & Y' \arrow[l,"h"']
\end{tikzcd}
\]
%%%%%%%%%%  scan idx 305  |  folio 292  %%%%%%%%%%
\SGAPrintPage{292}%
\uline{be a Cartesian square of schemes}. \uline{Suppose} \(Y\) \uline{and} \(Y'\) \uline{quasi-compact}, \(X\) \uline{and} \(Y'\) \uline{tor-independent over} \(Y\) (III 1.5), \(f\) \uline{proper}, \uline{and pseudo-coherent}, \uline{and the finiteness conjecture} (III 2.1) \uline{satisfied for} \(f\) \uline{and} \(f'\) (\uline{for example} \(Y\) \uline{and} \(Y'\) \uline{noetherian}, \uline{or} \(f\) \uline{projective}). \uline{Then}, \uline{if} \(f\) \uline{is perfect} (\uline{resp}. \(h\) \uline{of finite tor-dimension}), \uline{we have}, \uline{for} \(x\in K^{\bullet}(X)\) (\uline{resp}. \(K_{\bullet}(X)\)),
\[
 h^* f_*(x)=f'_* g^*(x),
\]
\uline{where} \(h^*,g^*\) \uline{are defined by} (2.7.1) (\uline{resp}. (2.12.1)), \uline{and} \(f_*,f'_*\) \uline{by} (2.12.3) (\uline{resp}. (2.11.1.1)) (\uline{note that by} (III 4.7.1) \(f'\) \uline{is perfect if} \(f\) \uline{is}, \uline{and that by} (III 3.5.1) \(g\) \uline{has finite tor-dimension if} \(h\) \uline{does}).

\sgaStableTarget{sga6.e04.para.3-2}3.2. \uline{Passage to the limit for} \(K^{\bullet}_{\naif}\) \uline{and} \(K_{\bullet}\) Let \((X_i,f_{ij})_{i\in I}\) be a filtered projective system of quasi-compact and quasi-separated schemes, with affine transition morphisms (cf. EGA IV 8.2.2). Put \(X=\varprojlim X_i\), and denote by \(f_i:X\to X_i\) the canonical homomorphism.

\sgaStableTarget{sga6.e04.para.3-2-1}3.2.1. The homomorphisms (2.7.2) \(f_i^*:K^{\bullet}(X_i)_{\naif}\longrightarrow K^{\bullet}(X)_{\naif}\) define, by the transitivity formula (2.7.3), a homomorphism
\sgaStableTarget{sga6.e04.eq.3-2-1-1}\begin{equation}\tag{3.2.1.1}
 \varinjlim K^{\bullet}(X_i)_{\naif}\longrightarrow K^{\bullet}(X)_{\naif}.
\end{equation}
I claim that \uline{the homomorphism} (3.2.1.1) \uline{is an isomorphism}. Indeed, let \(E\) be a locally free \(\underline O_X\)-Module of finite type. By (EGA IV 8.5.2 and 8.5.5), there is an index \(i\) and a locally free \(\underline O_{X_i}\)-Module of finite type \(E_i\) such that \(E\simeq f_i^*(E_i)\). By (EGA IV 8.5.2.5), the image of \(\cl(E_i)\) in \(\varinjlim K^{\bullet}(X_i)_{\naif}\) is independent of the choice of \(E_i\); denote it by \(u(E)\). By (EGA IV 8.5.2 and 8.5.6), the map
\(u:\Ob\Loclib(X)\longrightarrow \varinjlim K^{\bullet}(X_i)_{\naif}\) thus defined is additive (1.5);
%%%%%%%%%%  scan idx 306  |  folio 293  %%%%%%%%%%
\SGAPrintPage{293}%
it therefore defines a homomorphism \(K^{\bullet}(X)_{\naif}\longrightarrow \varinjlim K^{\bullet}(X_i)_{\naif}\), which is immediately seen, by (EGA IV 8.5.2), to be inverse to (3.2.1.1), proving our assertion.

\sgaStableTarget{sga6.e04.para.3-2-2}3.2.2. In the same way, one defines a canonical homomorphism
\sgaStableTarget{sga6.e04.eq.3-2-2-1}\begin{equation}\tag{3.2.2.1}
 \varinjlim K^{\bullet}(X_i)\longrightarrow K^{\bullet}(X).
\end{equation}
Using Deligne's localization technique (SGA 4 XVII), which implies that the category of perfect complexes on \(X\) is equivalent to the category \(\mathop{\underrightarrow{\mathrm{Lim}}}\) of the analogous categories \(\Parf(X_i)\), one again finds that (3.2.2.1) is an isomorphism.

\sgaStableTarget{sga6.e04.para.3-2-3}3.2.3. Suppose the \(X_i\) and \(X\) \uline{noetherian}, and the morphisms \(f_{ij}\) \uline{flat} (the \(f_i\) then are as well). Then the homomorphisms (2.12.1)
\(f_i^*:K_{\bullet}(X_i)\longrightarrow K_{\bullet}(X)\) define an isomorphism
\begin{equation}\tag{3.2.3.1} \varinjlim K_{\bullet}(X_i)\xrightarrow{\sim}K_{\bullet}(X).
\end{equation}
The proof is analogous to that of (3.2.1), using the fact (2.4) that \(K_{\bullet}(X)\simeq K_{\bullet}(X)_{\naif}\) is computed in terms of coherent sheaves on \(X\).

Since the module structures of the \(K_{\bullet}(X_i)\) over the \(K^{\bullet}(X_i)_{\naif}\) are compatible with the inverse images \(f_{ij}^*\) by (2.7 a)) and (2.12.2), we find, by passing to the limit, that \(\varinjlim K_{\bullet}(X_i)\) has a module structure over \(\varinjlim K^{\bullet}(X_i)_{\naif}\); this is compatible, via the isomorphisms (3.2.1.1) and (3.2.3.1) with the module structure of \(K_{\bullet}(X)\) over \(K^{\bullet}(X)_{\naif}\). In other words, we have a commutative diagram, where the vertical arrows are defined by (3.2.1.1) and (3.2.3.1):
\[
\begin{tikzcd}[column sep=tiny]
\varinjlim K^{\bullet}(X_i)_{\naif}\otimes \varinjlim K_{\bullet}(X_i)
\arrow[r]\arrow[d,"\simeq"]
&\varinjlim K_{\bullet}(X_i)\arrow[d,"\simeq"]\\
K^{\bullet}(X)_{\naif}\otimes K_{\bullet}(X)\arrow[r]
&K_{\bullet}(X).
\end{tikzcd}
\]

%%%%%%%%%%  scan idx 307  |  folio 294  %%%%%%%%%%
\SGAPrintPage{294}%
\sgaStableTarget{sga6.e04.para.3-3}3.3. \uline{Relative Grothendieck groups}

\sgaStableTarget{sga6.e04.para.3-3-1}3.3.1. Let \(f:X\longrightarrow Y\) be a locally finite type morphism of schemes. We shall put:
\[
\begin{gathered}
 K^{\bullet}(f)=k(D(f)_{\parf}),\\
 K_{\bullet}(f)=k(D^b(f)_{\coh}),
\end{gathered}
\]
where \(D^b(f)_{\coh}\) denotes the category of complexes pseudo-coherent rel. to \(f\) and with bounded cohomology (III 1.2), and \(D(f)_{\parf}\) the category of complexes of finite perfect amplitude rel. to \(f\) (III 4.2). When no ambiguity is possible, we shall sometimes write \(K^{\bullet}(X/Y)\) (resp. \(K_{\bullet}(X/Y)\)) instead of \(K^{\bullet}(f)\) (resp. \(K_{\bullet}(f)\)).

The groups \(K^{\bullet}(f)\) and \(K_{\bullet}(f)\) do not possess the rich structure of their ``absolute'' analogues \(K^{\bullet}(X)\) and \(K_{\bullet}(X)\), since the tensor product does not in general define a homomorphism \(K^{\bullet}(f)\otimes K^{\bullet}(f)\longrightarrow K^{\bullet}(f)\) (resp. \(K^{\bullet}(f)\otimes K_{\bullet}(f)\longrightarrow K_{\bullet}(f)\)). On the other hand, they satisfy interesting variance properties, which we shall review.

\sgaStableTarget{sga6.e04.para.3-3-2}3.3.2. Let
\[
\begin{tikzcd}
X\arrow[rr,"f"]\arrow[dr,"h"'] &&Y\arrow[dl,"g"]\\
&S&
\end{tikzcd}
\]
be a commutative triangle of locally finite type morphisms of schemes.
\begin{enumerate}[label=3.3.2.\arabic*., leftmargin=4em]
\item If \(g\) is pseudo-coherent, we have \(K_{\bullet}(f)=K_{\bullet}(h)\) (III 1.12). If \(g\) is smooth and \(X\) quasi-compact, we have \(K^{\bullet}(f)=K^{\bullet}(h)\) (III 4.4.1).

\item By (III 1.11 and 4.5), the functor \(\otimes_Y^{\mathrm L}\) defines pairings
\[
 D(f)_{\parf}\times D^b(g)_{\coh}\longrightarrow D^b(h)_{\coh},
\]
\[
 D(f)_{\parf}\times D(g)_{\parf}\longrightarrow D(h)_{\parf},
\]
which induce homomorphisms
\[
 K^{\bullet}(f)\otimes K_{\bullet}(g)\longrightarrow K_{\bullet}(h),
\]
\[
 K^{\bullet}(f)\otimes K^{\bullet}(g)\longrightarrow K^{\bullet}(h).
\]
%%%%%%%%%%  scan idx 308  |  folio 295  %%%%%%%%%%
\SGAPrintPage{295}%
We deduce in particular (by putting \(X=Y\), \(f=\Id\)) that \(K_{\bullet}(g)\) is a \(K^{\bullet}(Y)\)-module. On the other hand, if \(f\) is perfect, multiplication by \(\cl(\underline O_X)\in K^{\bullet}(f)\) defines \(K^{\bullet}(Y)\)-linear homomorphisms
\[
 f^*:K_{\bullet}(Y/S)\longrightarrow K_{\bullet}(X/S),
\]
and
\[
 f^*:K^{\bullet}(Y/S)\longrightarrow K^{\bullet}(X/S).
\]

\item Suppose \(S\) quasi-compact, \(X\) and \(Y\) of finite type over \(S\), and \(f\) proper. Then, subject to the validity of the finiteness conjecture (III 2.1), hence in particular when \(Y\) is noetherian or \(f\) is projective, the functor \(Rf_*\) sends \(D^b(h)_{\coh}\) to \(D^b(g)_{\coh}\), and therefore defines a homomorphism
\[
 f_*:K_{\bullet}(X/S)\longrightarrow K_{\bullet}(Y/S).
\]
If moreover \(f\) is perfect, \(Rf_*\) sends \(D(h)_{\parf}\) to \(D(g)_{\parf}\) (III 4.8), and defines a homomorphism
\[
 f_*:K^{\bullet}(X/S)\longrightarrow K^{\bullet}(Y/S).
\]
From the projection formula (III 3.7) we deduce the formula
\[
 f_*(a x)=a f_*(x)
\]
for \(f\) proper, \(a\in K^{\bullet}(S)\), \(x\in K_{\bullet}(X/S)\) (resp. \(f\) proper and perfect, \(a\in K_{\bullet}(S)\), \(x\in K^{\bullet}(X/S)\)).

\item Suppose \(S\) noetherian, \(X\) and \(Y\) of finite type over \(S\), and \(f\) compactifiable and perfect. Then (III 4.9.1) the functor \(f^!\) sends \(D(g)_{\parf}\) to \(D(h)_{\parf}\), and defines a homomorphism
\[
 f^!:K^{\bullet}(Y/S)\longrightarrow K^{\bullet}(X/S).
\]
\end{enumerate}

\sgaStableTarget{sga6.e04.para.3-3-3}3.3.3. Under the hypotheses of (III 4.7), the functor \(\otimes_S^{\mathrm L}\) defines pairings
\[
 D(f_1)_{\parf}\times D^b(f_2)_{\coh}\longrightarrow D^b(f)_{\coh},
\]
\[
 D(f_1)_{\parf}\times D(f_2)_{\parf}\longrightarrow D(f)_{\parf},
\]
which induce homomorphisms
\[
 K^{\bullet}(f_1)\otimes K_{\bullet}(f_2)\longrightarrow K_{\bullet}(f),
\]
\[
 K^{\bullet}(f_1)\otimes K^{\bullet}(f_2)\longrightarrow K^{\bullet}(f).
\]
%%%%%%%%%%  scan idx 309  |  folio 296  %%%%%%%%%%
\SGAPrintPage{296}%
(Of course, for \(X_i=Y_i=S\), we recover the \(K^{\bullet}(S)\)-module structure of \(K_{\bullet}(S)\).)

If \(X\) is flat over \(S\), we deduce, for \(n\geq1\), homomorphisms
\[
 K^{\bullet}(X/S)^{\otimes n}\longrightarrow K^{\bullet}\bigl((X/S)^n\bigr),
\]
which should in a sense replace the absent ring structure on \(K^{\bullet}(X/S)\).

\sgaStableTarget{sga6.e04.bibliography}\subsection*{BIBLIOGRAPHY}
\begin{enumerate}[label={[\arabic*]}]
\item Deligne, P., papiers secrets.
\item Grauert, H., Ein Theorem der Analytischen Garbentheorie \dots, Pub. Math. IHES, n\textsuperscript{o}5 (1960).
\item Grothendieck, A., Théorèmes de finitude pour la cohomologie des faisceaux, Bull. Soc. Math. France, 84, 1956, p.1-7.
\item[{[V]}] Verdier, J.-L., Catégories dérivées, IHES (1964), miméographié.
\end{enumerate}


% SGA_NATIVE_CONTINUATION_296
\providecommand{\Chat}{\widehat C}
\providecommand{\Ghat}{\widehat G}
\providecommand{\Gg}{\widehat G_g}
\providecommand{\Ga}{\widehat G_a}
\providecommand{\Grhat}{\widehat{\operatorname{Gr}}}
\providecommand{\Abfil}{\mathrm{Ab\,fil}}
\providecommand{\Kalgfil}{K\text{-}\mathrm{Alg.fil.}}
\providecommand{\Hom}{\operatorname{Hom}}
\providecommand{\Sym}{\mathfrak S}
\providecommand{\id}{\operatorname{Id}}
\providecommand{\Log}{\operatorname{Log}}
\providecommand{\Toddc}{\operatorname{Todd}}
\providecommand{\Ob}{\operatorname{Ob}}
\providecommand{\cG}{\mathcal G}
\providecommand{\cCzero}{\mathcal C_0}
\providecommand{\Fil}{\operatorname{Fil}}



% SGA_REVIEWED_FRENCH_CONTINUATION_296
% Expose V macros
\providecommand{\Chat}{\widehat C}
\providecommand{\Ghat}{\widehat G}
\providecommand{\Gg}{\widehat G_g}
\providecommand{\Ga}{\widehat G_a}
\providecommand{\Grhat}{\widehat{\operatorname{Gr}}}
\providecommand{\Abfil}{\underline{\mathrm{Ab}}\,\underline{\mathrm{fil}}}
\providecommand{\Kalgfil}{K\text{-}\mathrm{Alg.fil.}}
\providecommand{\Hom}{\operatorname{Hom}}
\providecommand{\Sym}{\mathfrak S}
\renewcommand{\id}{\operatorname{Id}}
\providecommand{\Log}{\operatorname{Log}}
\providecommand{\Toddc}{\operatorname{Todd}}
\renewcommand{\Ob}{\operatorname{Ob}}
\providecommand{\cG}{\mathcal G}
\providecommand{\cCzero}{\mathcal C_0}
\providecommand{\Fil}{\operatorname{Fil}}



% SGA_SOURCE_MACRO_RESTORATION_296: native provide must not shadow the two reviewed word rules.
\renewcommand{\Abfil}{\underline{\mathrm{Ab}}\,\underline{\mathrm{fil}}}

%%%%%%%%%%  scan idx 310  |  folio 297  %%%%%%%%%%
\SGAPrintPage{297}%
% Exposé V: back to the volume's parf
\renewcommand{\parf}{\operatorname{parf}}
\sgaStableTarget{sga6.e05.expose}\section*{\underline{Exposé V}\\[3pt] \underline{\underline{GENERALITIES ON \(\lambda\)-RINGS}}}
\begin{center}
\underline{by P. BERTHELOT}
\end{center}

The theory of \(\lambda\)-rings is useful whenever one uses rings such as the ring of classes of vector bundles on a ringed site, the representation ring of a group, etc..., for which the exterior-power operations define maps \(\lambda^i\). We shall use it extensively in the remainder of the Seminar, so it is necessary first to develop its general formalism. This exposé is thus entirely independent of the rest of the seminar and can serve users of \(\lambda\)-rings outside algebraic geometry.

The first section essentially gives Hirzebruch's theory of ``multiplicative sequences of polynomials'' (cf. [2]); the functorial viewpoint adopted here has the advantage, besides its generality, of considerably reducing the calculations. The following sections will show how, to every augmented \(\lambda\)-ring, one can associate a ``Chern ring'' that permits the development of a theory of Chern classes. This is the theory that will enable us to state the Riemann--Roch theorem in the general form of \hyperlink{sga6.e08.expose}{Exp. VIII}.

\sgaStableTarget{sga6.e05.sec.1}\section*{1. Universal polynomials}

\sgaStableTarget{sga6.e05.para.1-1}1.1. Let \(K\) be a commutative ring with unit, fixed throughout this section. We denote by \(\mathcal C\) the category of filtered commutative \(K\)-algebras with unit, whose filtration is decreasing and such that the group of elements of filtration zero is the whole algebra. We define a functor on \(\mathcal C\), with values in the category of groups
%%%%%%%%%%  scan idx 311  |  folio 298  %%%%%%%%%%
\SGAPrintPage{298}%
that are abelian, as follows: for \(A\in\operatorname{Ob}(\mathcal C)\), we associate to \(A\) the set of formal series \(\psi\in A[[t]]\) satisfying the following conditions: if \(\psi=a_0+a_1t+\ldots+a_nt^n+\ldots\), then
\sgaStableTarget{sga6.e05.eq.1-1-1}\begin{equation}\tag{1.1.1}
  a_0=1;\qquad \forall i\geq 1,\quad a_i\in\Fil^i A .
\end{equation}
This set will be denoted by \(\Ghat_K(A)\), or by \(\Ghat(A)\) when no confusion is possible; it is clear that it is an abelian group under multiplication of formal series, and that this does define a functor
\[
  \Ghat_K:\underline{K\text{-Alg.fil.}}\longrightarrow \underline{\mathrm{Ab}}.
\]
For \(A\) endowed with the coarse filtration, \(\Ghat(A)\) is the group of formal series of augmentation \(1\), which will also be denoted by \(1+A[[t]]^+\).

\sgaStableTarget{sga6.e05.lem.1-2}\textbf{Lemma 1.2.} \uline{The functor} \(\Ghat_K\) \uline{is representable}.

Let \((X_i)_{i\in N,\ i\geq 1}\) be a family of indeterminates; we assign the weight \(i\) to \(X_i\), and filter the ring \(K[X_i]_{i\geq 1}\) by weight with respect to the \(X_i\). We then define a bijection:
\[
  \Ghat_K(A)\xrightarrow{\sim}
  \Hom_{K\text{-Alg.fil.}}\bigl(K[X_i]_{i\geq 1},A\bigr)
\]
by associating to every series \(\psi=1+a_1t+\ldots+a_nt^n+\ldots\in\Ghat_K(A)\) the homomorphism \(f_\psi\) determined by \(f_\psi(X_i)=a_i\). This homomorphism is compatible with the filtrations, and we thus obtain an isomorphism of functors.

Put \(\phi=1+X_1t+\ldots+X_nt^n+\ldots\); for every filtered \(K\)-algebra \(A\) and every series \(\psi\in\Ghat(A)\), there exists exactly one homomorphism
%%%%%%%%%%  scan idx 312  |  folio 299  %%%%%%%%%%
\SGAPrintPage{299}%
\(f_\psi\) from \(K[X_i]_{i\geq 1}\) to \(A\) such that \(\Ghat(f_\psi)(\phi)=\psi\).

\sgaStableTarget{sga6.e05.para.1-3}1.3. For every \(A\in\operatorname{Ob}(\mathcal C)\) we define a filtration on \(\Ghat(A)\) as follows: a series \(\psi=1+a_1t+\ldots\) will be said to be of filtration \(\geq k\) if \(a_i=0\) for \(1\leq i<k\). We denote the subgroup so obtained by \(F^k\Ghat(A)\); this is a functor on \(\mathcal C\), and from now on we shall regard \(\Ghat\) as a functor with values in the category \(\underline{\mathrm{Ab}}.\underline{\mathrm{fil}}.\) of filtered abelian groups by means of this filtration (with \(F^0\Ghat=\Ghat\)).

Let \(H\) be a functor on \(\mathcal C\) with values in \(\underline{\mathrm{Ab}}\,\underline{\mathrm{fil}}\); we shall assume that \(H\) satisfies the following two conditions:
\[
\begin{array}{ll}
\text{i)} & \forall A\in\operatorname{Ob}(\mathcal C),\ H(A)\text{ is separated and complete}.\\[0.3em]
\text{ii)} & \text{If }A\in\operatorname{Ob}(\mathcal C)\text{ is a graded }K\text{-algebra, endowed with the filtration}\\
& \text{associated with its grading, and if }G\text{ is a finite group acting on }A\\
& \text{by graded morphisms, then }H(A^G)\xrightarrow{\sim} H(A)^G\\
& \text{(the exponent }G\text{ denoting invariants under }G\text{)}
\end{array}
\]

In fact, we shall need these conditions only when \(A\) is a polynomial ring and \(G\) is a finite product of symmetric groups acting by permutation of the indeterminates.

\sgaStableTarget{sga6.e05.ex.unnamed}\textbf{Examples}: 1)~\(H(A)=\Ghat(A)\). \sgaStableTarget{sga6.e05.ex.1-3-2}2)~\(H(A)=1+A[[t]]^+\); this functor will be denoted by \(\Gg\) (\(g\), the initial of ``grossière''). 3)~\(H(A)=\Ga(A)\), where \(\Ga(A)\) denotes the group, under addition, of formal series such that \(a_0=0\), \(a_i\in\Fil^i A\) for \(i\geq 1\), endowed with the filtration analogous to that of \(\Ghat\).

\sgaStableTarget{sga6.e05.para.1-4}1.4. For every \(n\), we shall denote by \(f_n\) the endomorphism of \(K[X_1,\ldots,X_n]\)
%%%%%%%%%%  scan idx 313  |  folio 300  %%%%%%%%%%
\SGAPrintPage{300}%
defined by \(f_n(X_i)=X_i\) for \(i<n\), \(f_n(X_n)=0\). It factors through the inclusion of \(K[X_1,\ldots,X_{n-1}]\) in \(K[X_1,\ldots,X_n]\), which shows that the homomorphism from \(H(K[X_1,\ldots,X_{n-1}])\) to \(H(K[X_1,\ldots,X_n])\) is injective, and will allow us to identify the former with a subgroup of the latter; similarly, we shall regard \(H(K[X_1,\ldots,X_n])\) as contained in \(H(K[X_i]_{i\geq 1})\).

Notice that \(f_1\) is precisely the augmentation of \(K[X]\).

\sgaStableTarget{sga6.e05.para.1-5}1.5. We embed \(K[X_1,\ldots,X_n]\) in \(K[T_1,\ldots,T_n]\) by sending \(X_i\) to the \(i\)-th symmetric function in the \(T_j\); letting the symmetric group \(\mathfrak S_n\) act by permutation of the \(T_j\), and using the theorem on symmetric polynomials and (1.3.ii)), we see that the resulting homomorphism from \(H(K[X_1,\ldots,X_n])\) to \(H(K[T_1,\ldots,T_n])\) is injective, which will again allow us to identify the former with a subgroup of the latter:
\[
  H(K[X_1,\ldots,X_n])\simeq H(K[T_1,\ldots,T_n]^{\mathfrak S_n})\xrightarrow{\sim} H(K[T_1,\ldots,T_n])^{\mathfrak S_n}\subset H(K[T_1,\ldots,T_n])
\]
We denote by
\[
  u_i^n:K[X]\longrightarrow K[T_1,\ldots,T_n]
\]
the homomorphism defined by \(u_i^n(X)=T_i\), and put
\[
  u_H^n=\sum_{i=1}^n H(u_i^n):H(K[X])\longrightarrow H(K[T_1,\ldots,T_n]).
\]
The symmetric group permutes the morphisms \(H(u_i^n)\), and therefore leaves the morphism \(u_H^n\) invariant, which consequently takes its values in
%%%%%%%%%%  scan idx 314  |  folio 301  %%%%%%%%%%
\SGAPrintPage{301}%
\(H(K[X_1,\ldots,X_n])\) by (1.3.ii)).

Finally, observe that for \(H=\Ghat\), \(u_H^n(1+Xt)=\prod\limits_{i=1}^n(1+T_it)=1+X_1t+\ldots+X_nt^n\).

\sgaStableTarget{sga6.e05.thm.1-6}\textbf{Theorem 1.6.} \uline{The map from} \(\Hom_{\underline{\mathrm{Ab}}\,\underline{\mathrm{fil}}}(\Ghat_K,H)\) \uline{to} \(H(K[X])\) \uline{which}, \uline{to a morphism of functors}, \uline{assigns its value on the element} \(1+Xt\) \uline{of} \(\Ghat_K(K[X])\), \uline{is injective}; \uline{its image is the subgroup of} \(H(K[X])\) \uline{formed by the elements} \(\theta\) \uline{such that} (\uline{with the notation of} 1.4. \uline{for the} \(f_n\)):
\[
\begin{array}{ll}
\text{i)}& H(f_1)(\theta)=0.\\[0.3em]
\text{ii)}& \forall n\geq1,\quad
H(f_n)\bigl(u_H^n(\theta)\bigr)\equiv u_H^n(\theta)
\ \mathrm{mod.}\ F^nH(K[X_1,\ldots,X_n]).
\end{array}
\]

Let us first prove injectivity. By Lemma 1.2, it is enough to see that a morphism of functors \(\Phi\) vanishing on \(1+Xt\) vanishes on \(\varphi=1+X_1t+\ldots+X_nt^n+\ldots\). The commutative diagram:
\sgaStableTarget{sga6.e05.eq.1-6-1}\[
\begin{tikzcd}
\Ghat(K[X]) \arrow[r,"u^n_{\Ghat}"] \arrow[d,"\Phi"'] &
\Ghat(K[X_1,\ldots,X_n]) \arrow[d,"\Phi"]\\
H(K[X]) \arrow[r,"u_H^n"] &
H(K[X_1,\ldots,X_n])
\end{tikzcd}
\tag{1.6.1}
\]
shows that for every \(n\), \(\Phi(1+X_1t+\ldots+X_nt^n)=0\).

Now \(\varphi\equiv 1+X_1t+\ldots+X_nt^n\ \mathrm{mod.}\ F^{n+1}\Ghat\); hence \(\Phi(\varphi)\equiv \Phi(1+\ldots+X_nt^n)\ \mathrm{mod.}\ F^{n+1}H\), and consequently \(\Phi(\varphi)\equiv0\ \mathrm{mod.}\ F^{n+1}H\), for every \(n\); since \(H(K[X_i]_{i\geq 1})\) is separated and complete, \(\Phi(\varphi)=0\).

Now let \(\Phi\) be any element of \(\Hom(\Ghat,H)\), and put
%%%%%%%%%%  scan idx 315  |  folio 302  %%%%%%%%%%
\SGAPrintPage{302}%
\(\theta=\Phi(1+Xt)\);
let us show that \(\theta\) satisfies conditions i) and ii).

i): \(H(f_1)(\theta)=H(f_1)\circ\Phi(1+Xt)=\Phi\circ\Ghat(f_1)(1+Xt)=\Phi(1)=0\).

ii): \(u_H^n(\theta)=\Phi(1+\ldots+X_nt^n)\) by (1.6.1); now \(1+\ldots+X_nt^n\equiv 1+\ldots+X_{n-1}t^{n-1}\ \mathrm{mod.}\ F^n\Ghat\), and \(1+\ldots+X_{n-1}t^{n-1}=\Ghat(f_n)(1+\ldots+X_nt^n)\). Since \(\Phi\circ\Ghat(f_n)=H(f_n)\circ\Phi\), and since \(\Phi\) is a morphism of filtered functors, we obtain ii).

Finally, we must show that if \(\theta\in H(K[X])\) satisfies i) and ii), there exists \(\Phi\in\Hom(\Ghat,H)\) such that \(\theta=\Phi(1+Xt)\).

We first determine the value of \(\Phi\) on the element \(1+\ldots+X_nt^n\) of \(\Ghat(K[X_1,\ldots,X_n])\) by means of (1.6.1); thus, by definition, \(\Phi(1+\ldots+X_nt^n)=\sum\limits_{i=1}^n \Phi(1+T_it)\), the \(\Phi(1+T_it)\) being known by functoriality through the morphisms \(u_i^n\).

We then glue the \(\Phi(1+\ldots+X_nt^n)\) for varying \(n\) by means of the commutative diagram:
\sgaStableTarget{sga6.e05.eq.1-6-2}\[
\begin{tikzcd}
K[X_1,\ldots,X_n] \arrow[r,"f_n"] \arrow[d,"\sigma_n"'] &
K[X_1,\ldots,X_n] \arrow[d,"\sigma'_n"']\\
K[T_1,\ldots,T_n] \arrow[r,"f'_n"] &
K[T_1,\ldots,T_n],
\end{tikzcd}
\tag{1.6.2}
\]
where \(f'_n\) is determined by \(f'_n(T_i)=T_i\) if \(i<n\), \(f'_n(T_n)=0\), where \(\sigma_n\) sends \(X_i\) to the \(i\)-th symmetric function of the \(T_j\), and \(\sigma'_n\) sends \(X_i\) to the \(i\)-th symmetric function of \(T_1,\ldots,T_{n-1}\) for \(i<n\), and \(X_n\) to \(0\).

Indeed, this gives:
\[
\begin{aligned}
H(f_n)\bigl(\Phi(1+\ldots+X_nt^n)\bigr)
&=H(f'_n)\left(\sum_{i=1}^n \Phi(1+T_it)\right)
=\sum_{i=1}^{n-1}\Phi(1+T_it)\ \text{(by 1.6.i))}\\
&=\Phi(1+\ldots+X_{n-1}t^{n-1})
\end{aligned}
\]
%%%%%%%%%%  scan idx 316  |  folio 303  %%%%%%%%%%
\SGAPrintPage{303}%
(working in \(H(K[T_1,\ldots,T_{n-1}])\) by means of the identifications 1.4).

In view of 1.6.ii), \(\Phi(1+\ldots+X_nt^n)\equiv\Phi(1+\ldots+X_{n-1}t^{n-1})\ \mathrm{mod.}\ F^nH(K[X_i]_{i\geq1})\), and the \(\Phi(1+\ldots+X_nt^n)\) define an element \(\psi\) of \(H(K[X_i]_{i\geq1})\), since this is a separated complete group; we therefore put \(\Phi(\varphi)=\psi\).

Finally, let \(A\in\operatorname{Ob}(\mathcal C)\), and \(\chi\in\Ghat(A)\); we put \(\Phi(\chi)=H(f_\chi)(\psi)\), where \(f_\chi\) is the homomorphism from \(K[X_i]_{i\geq 1}\) to \(A\) defined in 1.2. Since we have already defined \(\Phi\) on the elements \(1+\ldots+X_nt^n=\varphi_n\), we must verify that the two definitions give the same result, that is:
\[
  H(f_{\varphi_n})(\psi)=\Phi(1+\ldots+X_nt^n).
\]
Now, for every \(p\), we have: \(\psi\equiv \Phi(1+\ldots+X_pt^p)\ \mathrm{mod.}\ F^{p+1}H(K[X_i]_{i\geq 1})\). Hence, for every \(p\geq n\); \(H(f_{\varphi_n})(\psi)\equiv H(f_{\varphi_n})\bigl(\Phi(1+\ldots+X_pt^p)\bigr)\ \mathrm{mod.}\ F^{p+1}H(K[X_i]_{i\geq 1})\). Using the identifications of 1.4 and the gluing of the \(\Phi(\varphi_m)\) for varying \(m\), we obtain: \(H(f_{\varphi_n})(\Phi(\varphi_p))=\Phi(\varphi_n)\).
Thus, finally:
\[
  \forall p\geq n,\quad
  H(f_{\varphi_n})(\psi)\equiv \Phi(1+\ldots+X_nt^n)
  \ \mathrm{mod.}\ F^pH(K[X_i]_{i\geq 1}).
\]
Since \(H\) is separated, we obtain the required equality.

We have thus defined a system of maps \(\Phi_A:\Ghat(A)\to H(A)\); to verify that they constitute a morphism of functors (with values in \(\underline{\mathrm{Ens}}\)), we must show that for every morphism \(h:A\to A'\), the diagram:
\[
\begin{tikzcd}
\Ghat(A) \arrow[r,"\Ghat(h)"] \arrow[d,"\Phi_A"'] &
\Ghat(A') \arrow[d,"\Phi_{A'}"]\\
H(A) \arrow[r,"H(h)"] & H(A')
\end{tikzcd}
\]
%%%%%%%%%%  scan idx 317  |  folio 304  %%%%%%%%%%
\SGAPrintPage{304}%
is commutative. Let \(\chi\in\Ghat(A)\); then: \(h\circ f_\chi=f_{\Ghat(h)(\chi)}\); hence:
\[
  H(h)(\Phi_A(\chi))=H(h)\circ H(f_\chi)(\psi)
  =H(f_{\Ghat(h)(\chi)})(\psi)=\Phi_{A'}(\Ghat(h)(\chi)).
\]

It therefore remains to show that the morphism of functors thus constructed is compatible with the structures of filtered groups.

Suppose that \(\chi\in F^n\Ghat(A)\). To say that \(\chi\) has filtration \(\geq n\) is equivalent to saying that \(f_\chi:K[X_i]_{i\geq 1}\to A\) factors through the endomorphism \(g_n\) of \(K[X_i]_{i\geq 1}\) defined by: \(g_n(X_i)=0\) if \(1\leq i\leq n-1\), and \(g_n(X_i)=X_i\) if \(i\geq n\).
To show that \(\Phi\) is compatible with the filtrations, it therefore suffices to show that \(H(g_n)(\Phi(\varphi))\in F^nH\). Now, by construction:
\[
  \Phi(\varphi)\equiv \Phi(1+\ldots+X_{n-1}t^{n-1})\ \mathrm{mod.}\ F^nH,
\]
hence:
\[
  H(g_n)(\Phi(\varphi))
  \equiv H(g_n)\bigl(\Phi(1+\ldots+X_{n-1}t^{n-1})\bigr)
  =0\ \mathrm{mod.}\ F^nH.
\]

To verify the additivity of \(\Phi\), it suffices to do this for the elements of finite degree of \(\Ghat(A)\), and then pass to the limit. Thus let \(\chi_1,\chi_2\) be two elements of degree \(n\) of \(\Ghat(A)\), and \(\chi\) their product. Consider the commutative diagram:
\[
\begin{tikzcd}[column sep=tiny]
K[X_1,\ldots,X_{2n}] \arrow[rr,"f_\chi"] \arrow[d,hook] && A\\
\bigl(K[T_1,\ldots,T_{2n}]\bigr)^{\mathfrak S_n\times\mathfrak S_n}
\arrow[urr,"f'_\chi"] && ,
\end{tikzcd}
\]
where the arrows are defined as follows: \(\mathfrak S_n\times\mathfrak S_n\) acts on \(K[T_1,\ldots,T_{2n}]\)
%%%%%%%%%%  scan idx 318  |  folio 305  %%%%%%%%%%
\SGAPrintPage{305}%
in such a way that the first factor permutes \(T_1,\ldots,T_n\) and the second \(T_{n+1},\ldots,T_{2n}\); \(K[X_1,\ldots,X_{2n}]\) is embedded in the ring of invariants by sending \(X_j\) to the \(j\)-th symmetric function of the \(T_i\); \(f'\) sends the \(i\)-th symmetric function of \(T_1,\ldots,T_n\) (resp. \(T_{n+1},\ldots,T_{2n}\)) to the \(i\)-th coefficient of \(\chi_1\) (resp. \(\chi_2\)). Then:
\[
\begin{aligned}
\Phi(\chi)
&=H(f_\chi)\left(\Phi\left(\prod_{i=1}^{2n}(1+T_it)\right)\right)
=H(f'_\chi)\left(\Phi\left(\prod_{i=1}^{2n}(1+T_it)\right)\right);\\
\Phi(\chi_1) &=H(f'_\chi)\left(\Phi\left(\prod_{i=1}^{n}(1+T_it)\right)\right)\quad(\text{resp. }\chi_2,\ i=n+1,\ldots,2n).
\end{aligned}
\]
Since \(H(f'_\chi)\) is a homomorphism of groups, we are reduced to the case where \(\chi=\prod_{i=1}^{2n}(1+T_it)\), and in this case the additivity of \(\Phi\) follows from the definition of the \(\Phi(1+X_1t+\ldots+X_k t^k)\).

\sgaStableTarget{sga6.e05.para.1-7}1.7. We can generalize the preceding considerations to a product of functors \(\Ghat\); the functor \((\Ghat)^k\) is representable by the ring
\(K[(X_i^{(j)})]\), \(j=1,\ldots,k\), and \(i\geq 1\), where \(X_i^{(j)}\) has weight \(i\). We denote by \(\varphi^{(j)}\) the series \(1+X_1^{(j)}t+\ldots+X_n^{(j)}t^n+\ldots\); the element \((\varphi^{(j)})_j\) of \(\bigl(\Ghat(K[(X_i^{(j)})])\bigr)^k\) has the universal property analogous to that of \(\varphi\) (cf. 1.2).

We can introduce the morphisms \(f_n^{(j)}\), \(u_H^{n,(j)}\) as before, by repeating the construction for each family \((X_i^{(j)})_{i\geq 1}\) separately. We obtain a result analogous to Theorem 1.6 by seeking the homomorphisms of functors from \((\Ghat)^k\) to \(H\) that are \(\mathbb Z\)-multilinear and compatible with the filtrations in each of the \(k\) arguments separately: they are determined by their value on the element
%%%%%%%%%%  scan idx 319  |  folio 306  %%%%%%%%%%
\SGAPrintPage{306}%
\((1+X^{(j)}t)_{j=1,\ldots,k}\) of \(\bigl(\Ghat(K[(X^{(j)})_j])\bigr)^k\); an element \(\theta\) of \(H(K[(X^{(j)})_j])\) determines a morphism of functors if and only if it satisfies:
\[
\begin{array}{ll}
\text{i)}& \forall j,\quad H(f_1^{(j)})(\theta)=0\\[0.3em]
\text{ii)}& \forall n,j,\quad
H(f_n^{(j)})\bigl(u_H^{n,(j)}(\theta)\bigr)
\equiv u_H^{n,(j)}(\theta)\ \mathrm{mod.}\ F^nH.
\end{array}
\]

The reader will readily verify that, apart from complications of notation, it suffices to copy the proof of 1.6.

\sgaStableTarget{sga6.e05.para.1-8}1.8. We denote by \(\mathcal C_0\) the category of \(K\)-algebras, and by
\[
  i:\mathcal C_0\longrightarrow \mathcal C,
\]
the functor that assigns to a \(K\)-algebra \(A\) the \(K\)-algebra \(A\) endowed with the coarse filtration; we denote by
\[
  g:\mathcal C\longrightarrow\mathcal C_0,
\]
the functor ``forget the filtration.'' These functors satisfy the following relations:
\[
  g\circ i=\operatorname{Id}_{\mathcal C_0};\qquad
  \Ghat\circ i\circ g=\Gg;\qquad
  \Gg\circ i\circ g=\Gg
\]
(cf. 1.3, ex. 2)

\sgaStableTarget{sga6.e05.lem.1-9}\textbf{Lemma 1.9.} \uline{Let} \(H:\mathcal C_0\to\underline{\mathrm{Ab}}\,\underline{\mathrm{fil}}\) \uline{be a functor on} \(\mathcal C_0\); \uline{then the canonical homomorphism}:
\[
  \Hom_{\underline{\mathrm{Ab}}\,\underline{\mathrm{fil}}}(\Ghat,H\circ g)
  \longrightarrow
  \Hom_{\underline{\mathrm{Ab}}\,\underline{\mathrm{fil}}}(\Ghat\circ i,H)
\]
%%%%%%%%%%  scan idx 320  |  folio 307  %%%%%%%%%%
\SGAPrintPage{307}%
\uline{is an isomorphism}.

To see that it is injective, we use the sequence of canonical homomorphisms:
\[
  \Hom(\Ghat,H\circ g)\longrightarrow
  \Hom(\Ghat\circ i,H)\longrightarrow
  \Hom(\Gg,H\circ g)\longrightarrow
  \Hom(\Ghat,H\circ g),
\]
the last arrow coming from the fact that \(\Ghat\) is a subfunctor of \(\Gg\). To verify that the composite homomorphism is the identity, we observe that a morphism of functors: \(\Ghat\longrightarrow H\circ g\) is determined by its value on the objects of \(\mathcal C\) whose filtration is coarse; indeed, let \(A\in\operatorname{Ob}(\mathcal C)\), and let \(A_g\) be the algebra \(A\) endowed with the coarse filtration; the canonical morphism \(A\to A_g\) defines a commutative diagram:
\[
\begin{tikzcd}
\Ghat(A) \arrow[r,"\Phi_A"] \arrow[d] &
H\circ g(A) \arrow[d]\\
\Ghat(A_g) \arrow[r,"\Phi_{A_g}"] &
H\circ g(A_g),
\end{tikzcd}
\]
showing that \(\Phi_A\) is the restriction of \(\Phi_{A_g}\) to \(\Ghat(A)\).

We next consider the composite homomorphism:
\[
  \Hom(\Ghat\circ i,H)\longrightarrow
  \Hom(\Gg,H\circ g)\longrightarrow
  \Hom(\Ghat,H\circ g)\longrightarrow
  \Hom(\Ghat\circ i,H),
\]
and it is clear that this is the identity, giving the asserted result.

\sgaStableTarget{sga6.e05.thm.1-10}\textbf{Theorem 1.10.} \uline{Let} \(H\) \uline{be a functor on the category} \(\mathcal C_0\) \uline{of} \(K\)-\uline{algebras}, \uline{with values in the category of separated complete filtered abelian groups}, \uline{and commuting with invariants under a finite group}. \uline{Then the map from} \(\Hom_{\underline{\mathrm{Ab}}\,\underline{\mathrm{fil}}}(\Ghat\circ i,H)\) \uline{to} \(H(K[X])\) \uline{which}, \uline{to a morphism of functors}, \uline{assigns}
%%%%%%%%%%  scan idx 321  |  folio 308  %%%%%%%%%%
\SGAPrintPage{308}%
\uline{its value on the element} \(1+Xt\) \uline{of} \(\Ghat\circ i(K[X])\), \uline{is injective}; \uline{its image is the subgroup of} \(H(K[X])\) \uline{formed by the elements} \(\theta\) \uline{such that} (\uline{with notation analogous to that of} 1.4 \uline{and} 1.5):
\[
\begin{array}{ll}
\text{i)}& H(f_1)(\theta)=0;\\[0.3em]
\text{ii)}& \forall n\geq1,\quad
H(f_n)\bigl(u_H^n(\theta)\bigr)\equiv u_H^n(\theta)
\ \mathrm{mod.}\ F^nH(K[X_1,\ldots,X_n]).
\end{array}
\]

The theorem follows immediately from Theorem 1.6 and Lemma 1.9, since \(H\circ g(K[X])\) does not depend on the filtration chosen on \(K[X]\).

The analogous theorem obtained by replacing \(\Ghat\) by \((\Ghat)^k\) is also true (cf. 1.7)

When there is no risk of confusion, we shall denote the functor \(\Ghat\circ i\) by \(\Ghat\); this is legitimate, since \(i\) identifies \(\mathcal C_0\) with a subcategory of \(\mathcal C\).

\sgaStableTarget{sga6.e05.para.1-11}1.11. We denote by \(\mathcal G\) the category of graded \(K\)-algebras, and by \(j\) the functor from \(\mathcal G\) to \(\mathcal C\) that assigns to a graded \(K\)-algebra \(B\) the \(K\)-algebra \(B\) endowed with the filtration canonically defined by its grading; the ``associated graded'' functor: \(\mathcal C\to\mathcal G\) will be denoted by \(\operatorname{Gr}\). There is an isomorphism of functors \(\varepsilon:\operatorname{Id}_{\mathcal G}\xrightarrow{\sim}\operatorname{Gr}\circ j\).

Let \(B\in\operatorname{Ob}(\mathcal G)\); we denote by \(1+\widehat B^+\) the multiplicative group of elements of \(\prod\limits_{i=0}^{\infty}B^i\) with augmentation \(1\). Endowing it with the filtration analogous to that of \(\Ghat\), we thus define a functor: \(\mathcal G\to\underline{\mathrm{Ab}}\,\underline{\mathrm{fil}}\).

Furthermore, for every \(A\in\operatorname{Ob}(\mathcal C)\), we define a homomorphism of filtered groups: \(\Ghat(A)\longrightarrow 1+\widehat{\operatorname{Gr}A}^{+}\), by sending the series \(1+\ldots+a_nt^n+\ldots\),
%%%%%%%%%%  scan idx 322  |  folio 309  %%%%%%%%%%
\SGAPrintPage{309}%
to the series \(1+\ldots+\bar a_n+\ldots\), where \(\bar a_n\) is the class of \(a_n\) in \(\operatorname{Gr}^nA\). Its kernel is the group \(\Ghat'(A)\) of formal series such that:
\[
  a_0=1;\qquad \forall i\geq1,\quad a_i\in\Fil^{i+1}A;
\]
we therefore have an isomorphism, functorial in \(A\):
\[
  \Ghat(A)/\Ghat'(A)\xrightarrow{\sim}1+\widehat{\operatorname{Gr}A}^{+}.
\]

\sgaStableTarget{sga6.e05.lem.1-12}\textbf{Lemma 1.12.} \uline{Let} \(H:\mathcal G\to\underline{\mathrm{Ab}}\,\underline{\mathrm{fil}}\) \uline{be a functor on} \(\mathcal G\); \uline{then the canonical homomorphism}:
\[
  \rho:\Hom_{\underline{\mathrm{Ab}}\,\underline{\mathrm{fil}}}(1+\widehat{\,\mathord{.}\,}^{+},H)
  \longrightarrow
  \Hom_{\underline{\mathrm{Ab}}\,\underline{\mathrm{fil}}}(\Ghat,H\circ\operatorname{Gr}),
\]
\uline{is an isomorphism}.

The homomorphism \(\rho\) is defined as the composite:
\[
  \Hom(1+\widehat{\,\mathord{.}\,}^{+},H)
  \longrightarrow
  \Hom(1+\widehat{\operatorname{Gr}\,\mathord{.}\,}^{+},H\circ\operatorname{Gr})
  \longrightarrow
  \Hom(\Ghat,H\circ\operatorname{Gr}),
\]
the last arrow being defined by the homomorphism \(\Ghat\longrightarrow1+\widehat{\operatorname{Gr}\,\mathord{.}\,}^{+}\)

We shall construct a morphism in the opposite direction:
\[
  \rho':\Hom(\Ghat,H\circ\operatorname{Gr})
  \longrightarrow
  \Hom(1+\widehat{\,\mathord{.}\,}^{+},H).
\]
For this, we shall show that every morphism of functors from \(\Ghat\) to \(H\circ\operatorname{Gr}\) factors through the canonical morphism \(\Ghat\longrightarrow1+\widehat{\operatorname{Gr}\,\mathord{.}\,}^{+}\), which, by 1.11, is equivalent to showing that it vanishes on \(\Ghat'\).

Let \(A\in\operatorname{Ob}(\mathcal C)\), \(\chi\in\Ghat'(A)\); we define a homomorphism of
%%%%%%%%%%  scan idx 323  |  folio 310  %%%%%%%%%%
\SGAPrintPage{310}%
filtered \(K\)-algebras \(K[X_i]_{i\geq1}\longrightarrow A\) by sending \(X_1\) to \(0\) and \(X_i\) to the \((i-1)\)-th coefficient of \(\chi\) for \(i>1\); using this homomorphism, we see that it suffices to verify that the morphism of functors under consideration vanishes on the element \(1+X_2t+\ldots+X_{n+1}t^n+\ldots\). We then use the endomorphism \(f\) of \(K[X_i]_{i\geq1}\) defined by \(f(X_i)=X_{i+1}\) for every \(i\), and the commutative diagram:
\[
\begin{tikzcd}
\Ghat(K[(X_i)]) \arrow[r] \arrow[d,"\Ghat(f)"'] &
H\circ\operatorname{Gr}(K[(X_i)]) \arrow[d,"H(\operatorname{Gr}(f))"]\\
\Ghat(K[(X_i)]) \arrow[r] &
H\circ\operatorname{Gr}(K[(X_i)]);
\end{tikzcd}
\]
now \(\operatorname{Gr}(f)=\operatorname{Gr}(f_1)\), where \(f_1\) is the augmentation homomorphism of \(K[X_i]_{i\geq1}\); since the latter annihilates \(\varphi\), the analogous diagram written for \(f_1\) shows that \(H(\operatorname{Gr}(f))\) annihilates the image of \(\varphi\) in \(H\circ\operatorname{Gr}(K[X_i]_{i\geq1})\), giving the result.

Passing to the quotient, we therefore associate to every morphism of functors: \(\Ghat\longrightarrow H\circ\operatorname{Gr}\) a morphism of functors: \(1+\widehat{\operatorname{Gr}\,\mathord{.}\,}^{+}\longrightarrow H\circ\operatorname{Gr}\); by composing with \(j\), and using the isomorphism \(\varepsilon\), we obtain a morphism: \(1+\widehat{\,\mathord{.}\,}^{+}\longrightarrow H\), which defines the required homomorphism \(\rho'\). It is immediate to verify that \(\rho\) and \(\rho'\) are inverses of one another.

\sgaStableTarget{sga6.e05.thm.1-13}\textbf{Theorem 1.13.} \uline{Let} \(H\) \uline{be a functor on} \(\mathcal G\), \uline{with values in the category of separated complete filtered abelian groups}, \uline{and commuting with invariants under a finite group}. \uline{Then the map from} \(\Hom_{\underline{\mathrm{Ab}}\,\underline{\mathrm{fil}}}(1+\widehat{\,\mathord{.}\,}^{+},H)\) \uline{to} \(H(K[X])\) \uline{which}, \uline{to a morphism of functors}, \uline{assigns its value on the element}
%%%%%%%%%%  scan idx 324  |  folio 311  %%%%%%%%%%
\SGAPrintPage{311}%
\(1+X\) of \(1+\widehat{K[X]}^+\) \uline{is injective}; \uline{its image is the subgroup of} \(H(K[X])\) \uline{formed by the elements} \(\theta\) \uline{such that} (\uline{notation of} 1.4,1.5):
\[
\begin{array}{ll}
\text{i)}& H(f_1)(\theta)=0;\\[0.3em]
\text{ii)}& \forall n\geq1,\quad
H(f_n)\bigl(u_H^n(\theta)\bigr)\equiv u_H^n(\theta)
\ \mathrm{mod.}\ F^nH(K[X_1,\ldots,X_n]).
\end{array}
\]

Immediate from 1.6 and 1.12.

There is an analogous statement on replacing \(\Ghat\) by \((\Ghat)^k\).

\sgaStableTarget{sga6.e05.rem.1-14}\textbf{Remark 1.14.} If \(H\) is a functor on \(\mathcal G\), with values in the category of abelian groups, commuting with invariants under a finite group, and satisfying the condition:
\[
  \forall B\in\operatorname{Ob}(\mathcal G),\qquad
  H(B)\xrightarrow{\sim}\mathop{\underleftarrow{\mathrm{Lim}}}\limits_n H(B/B^{\geq n})\quad(\text{where }B^{\geq n}=\prod_{i\geq n}B^i),
\]
then condition ii) of 1.13 is always satisfied when we take as the filtration on \(H\) the filtration by the kernels of the \(H(B)\longrightarrow H(B/B^{\geq n})\) (which is separated and complete). Indeed, let \(p_n\) be the canonical homomorphism
\[
  K[X_1,\ldots,X_n]\longrightarrow K[X_1,\ldots,X_n]/K[X_1,\ldots,X_n]^{\geq n};
\]
we have: \(p_n\circ f_n=p_n\), hence: \(H(p_n)\bigl[H(f_n)\circ u_H^n(\theta)-u_H^n(\theta)\bigr]=0\), that is:
\[
  H(f_n)\circ u_H^n(\theta)-u_H^n(\theta)
  \in F^nH(K[X_1,\ldots,X_n]).
\]

Thus, to specify a morphism of functors from \(1+\widehat{\,\mathord{.}\,}^{+}\) to \(H\), it suffices to specify an element of \(H(K[X])\) satisfying only condition i) of 1.13.

\sgaStableTarget{sga6.e05.para.1-15}1.15. Let \(H:\mathcal C\to\underline{\mathrm{Ab}}\,\underline{\mathrm{fil}}\) (resp. \(\mathcal C_0\to\underline{\mathrm{Ab}}\,\underline{\mathrm{fil}}\), \(\mathcal G\to\underline{\mathrm{Ab}}\,\underline{\mathrm{fil}}\)) be a functor satisfying the conditions of 1.3 (resp. 1.10, 1.13), and let \(\Phi:(\Ghat)^k\to H\) (resp. \((\Ghat\circ i)^k\to H\), \((1+\widehat{\,\mathord{.}\,}^{+})^k\to H\)) be a morphism of functors. We have seen
%%%%%%%%%%  scan idx 325  |  folio 312  %%%%%%%%%%
\SGAPrintPage{312}%
(1.7) that \(\Phi\) is determined by its value on the element \((\varphi^{(j)})_{j=1..k}\) of \((\Ghat(K[X_i^{(j)}]_{\substack{i\geq1\\ j=1..k}}))^k\) (resp.\ \ldots)..\ Suppose that \(H\) is a functor such as \(\Ghat,\Gg,\Ga,1+\widehat{\operatorname{Gr}\,\mathord{.}\,}^{+}\), etc..; the element \(\Phi((\varphi^{(j)})_{j})\) is a formal series with coefficients in \(K[X_i^{(j)}]_{i\geq1,\ j=1..k}\); the \(n\)-th coefficient of this series will be called the \uline{\(n\)-th universal polynomial of} \(\Phi\); knowing these polynomials is therefore equivalent to knowing \(\Phi\). By the proof of 1.6, these polynomials are themselves determined by the element \(\Phi(1+X^{(1)}t,\ldots,1+X^{(k)}t)\), the calculation following the method of the proof: we calculate the elements \(\Phi(1+X_1^{(1)}t+\ldots+X_n^{(1)}t^n,\ldots,1+X_1^{(k)}t+\ldots+X_n^{(k)}t^n)\) by decomposing \(1+X_1^{(i)}t+\ldots+X_n^{(i)}t^n\) into \(\prod_{j=1}^{n}(1+T_j^{(i)}t)\) and using the additivity of \(\Phi\); the polynomials thus obtained are symmetric in the \(T_j^{(i)}\) (\(i\) fixed and \(j\) varying) and can be expressed in terms of the \(X_j^{(i)}\); by gluing with the aid of the filtration of \(H\), we thus determine the universal polynomials of \(\Phi\). If \((\psi^{(1)},\ldots,\psi^{(k)})\in(\Ghat(A))^k\), the calculation of \(\Phi(\psi^{(1)},\ldots,\psi^{(k)})\) is performed by substituting the coefficients \(a_j^{(i)}\) of the \(\psi^{(i)}\) for the \(X_j^{(i)}\) in the universal polynomials.

Suppose there is an integer \(m\) such that the coefficient of \(t^n\) in \(\theta=\Phi(1+X^{(1)}t,\ldots,1+X^{(k)}t)\) is a polynomial whose isobaric components have weights ranging between \(n\) and \(mn\); then the \(k\)-th universal polynomial of \(\Phi\) has isobaric components whose weights range between \(k\) and \(km\): this follows immediately from the calculation described above and from the theorem on symmetric functions.

%%%%%%%%%%  scan idx 326  |  folio 313  %%%%%%%%%%
\SGAPrintPage{313}%
Suppose in particular that \(H=\Ghat\); then the \(n\)-th universal polynomial of \(\Phi\) has no components of weight \(<n\). Now consider the canonical morphism: \(H\longrightarrow1+\widehat{\operatorname{Gr}\,\mathord{.}\,}^{+}\); composing with \(\Phi\) gives a morphism: \((\Ghat)^k\longrightarrow1+\widehat{\operatorname{Gr}\,\mathord{.}\,}^{+}\), which, by the proof of 1.12, factors through \((1+\widehat{\operatorname{Gr}\,\mathord{.}\,}^{+})^k\) and defines a morphism \(\Phi:(1+\widehat{\,\mathord{.}\,}^{+})^k\longrightarrow1+\widehat{\,\mathord{.}\,}^{+}\); we immediately see that the \(n\)-th universal polynomial of \(\Phi\) is the same as that of \((\Ghat)^k\longrightarrow1+\widehat{\operatorname{Gr}\,\mathord{.}\,}^{+}\), that is, the isobaric component of weight \(n\) of the \(n\)-th universal polynomial of \(\Phi\).

The same remarks apply to the case \(H=\Ga\): we can define a morphism of functors \(H\longrightarrow\widehat{\operatorname{Gr}\,\mathord{.}\,}^{+}\) (with \(\widehat{\operatorname{Gr}\,\mathord{.}\,}^{+}=\widehat{\,\mathord{.}\,}^{+}\circ\operatorname{Gr}\), where \(\widehat{\,\mathord{.}\,}^{+}\) is the functor defined on \(\mathcal G\) by \(\widehat B^+=\prod_{i\geq1}B^i\)) analogously to 1.11; we then repeat the reasoning and conclusions of the preceding paragraph.

\sgaStableTarget{sga6.e05.para.1-16}1.16. Let us finish with an example. Suppose that \(K=Q\). We define an endomorphism of the functor \(1+\widehat{\,\mathord{.}\,}^{+}\), denoted by \(\Toddc\) and called the Todd character, by putting:
\[
  \Toddc(1+X)=\frac{X}{1-e^{-X}}.
\]
Conditions i) and ii) of Theorem 1.13 are verified immediately, showing that we have indeed defined an endomorphism of \(1+\widehat{\,\mathord{.}\,}^{+}\).

\sgaStableTarget{sga6.e05.sec.2}\section*{2. Definition of \(\lambda\)-rings; examples}

In this section, the base ring of Section 1 will be \(\mathbb Z\) (except in 2.9); we work in the category \(\mathcal C_0\) (1.8) and denote by \(\Ghat\) the restriction of \(\Ghat\) to \(\mathcal C_0\) (1.10).

%%%%%%%%%%  scan idx 327  |  folio 314  %%%%%%%%%%
\SGAPrintPage{314}%
\sgaStableTarget{sga6.e05.def.2-1}\textbf{Definition 2.1.}: \uline{By a} pre-\(\lambda\)-ring \uline{we mean a ring} \(K\), \uline{commutative with unit}, \uline{endowed with a homomorphism of abelian groups}:
\[
  \lambda_t:K\longrightarrow1+K[[t]]^+,
\]
\uline{such that for every} \(x\in K\), \(\lambda_t(x)=1+xt+\ldots\)

We put \(\lambda_t(x)=\sum_{i\geq0}^{\infty}\lambda^i(x)t^i\); giving a pre-\(\lambda\)-ring structure on a ring \(K\) is therefore equivalent to giving a family of maps \(\lambda^i\), \(i\in N\), from \(K\) to \(K\), such that:
\[
\begin{array}{ll}
\text{i)}& \forall x\in K,\quad \lambda^0(x)=1,\\[0.3em]
\text{ii)}& \forall x\in K,\quad \lambda^1(x)=x,\\[0.3em]
\text{iii)}& \forall x,y\in K,\ \forall n\in N,\quad
\lambda^n(x+y)=\displaystyle\sum_{p+q=n}\lambda^p(x)\mathbin{.}\lambda^q(y).
\end{array}
\]

Let \(K,K'\) be two pre-\(\lambda\)-rings, and \(f:K\to K'\) a ring homomorphism; we say that \(f\) is a homomorphism of pre-\(\lambda\)-rings, or a \(\lambda\)-\uline{homomorphism}, when the diagram:
\[
\begin{tikzcd}
K \arrow[r,"\lambda_t"] \arrow[d,"f"'] &
1+K[[t]]^+ \arrow[d,"\Ghat(f)"]\\
K' \arrow[r,"\lambda_t"] &
1+K'[[t]]^+,
\end{tikzcd}
\]
is commutative.

This is also equivalent to the following relations:
\[
  \forall x\in K,\ \forall i\in N,\qquad
  \lambda^i(f(x))=f(\lambda^i(x)).
\]

%%%%%%%%%%  scan idx 328  |  folio 315  %%%%%%%%%%
\SGAPrintPage{315}%
\sgaStableTarget{sga6.e05.ex.2-2}\textbf{Examples 2.2.}: a) We can give all the pre-\(\lambda\)-ring structures on \(\mathbb Z\). For this it suffices to specify \(\lambda_t(1)\), which will therefore be a formal series with integer coefficients, subject only to the condition that it begin with \(1+t\).

b) Let \((T,\underline{A})\) be a commutatively ringed topos, and let \(K\) be the ``Grothendieck group'' of \(\underline{A}\)-Modules locally isomorphic to a direct summand of a free \(\underline{A}\)-Module of finite type; we know that \(K\) can be endowed with a ring structure by means of the tensor product (Exp. IV), and we shall endow it with a pre-\(\lambda\)-ring structure as follows: for every \(A\)-Module \(\underline{E}\) that is locally a direct summand of a free Module of finite type, we denote by \(\lambda^i(E)\) the class in \(K\) of the \(i\)-th exterior power of \(\underline{E}\) (which is also locally a direct summand of a free Module of finite type); this defines a function on the category of \(\underline{A}\)-Modules that are locally direct summands of a free Module of finite type, with values in the abelian group \(1+K[[t]]^+\); we must then show that this function is additive, which will imply that it factors through \(K\), giving the required structure.

\sgaStableTarget{sga6.e05.lem.2-2-1}\textbf{Lemma 2.2.1.} Let \(0\longrightarrow \underline{E}' \xrightarrow{u} \underline{E} \xrightarrow{v} \underline{E}''\longrightarrow0\) \uline{be an exact sequence of} \(\underline{A}\)-\uline{Modules that is locally split} (\uline{for example such that} \(E''\) \uline{is locally a direct summand of a free Module of finite type}). \uline{There exists on} \(\wedge^n(\underline{E})\) \uline{a finite filtration}, \uline{and isomorphisms}:
\[
  \wedge^{n-p}(\underline{E}')\otimes_A\wedge^p(\underline{E}'')
  \xrightarrow{\sim}
  \operatorname{Gr}^p\wedge^n(\underline{E}).
\]

We take as \(F^p\wedge^n(\underline{E})\) the image of the canonical morphism:
\[
\begin{gathered}
  \wedge^{n-p}(\underline{E}')\otimes_A\wedge^p(\underline{E})\longrightarrow\wedge^n(\underline{E})
  \quad\text{for } 0\leq p\leq n\ ;\quad \text{if } p\geq n,
\end{gathered}
\]
%%%%%%%%%%  scan idx 329  |  folio 316  %%%%%%%%%%
\SGAPrintPage{316}%
\(F^p\wedge^n(\underline{E})=\wedge^n(\underline{E})\), and if \(p<0\), \(F^p\wedge^n(\underline{E})=0\). This does define a finite filtration on \(\wedge^n(\underline{E})\).

We define a homomorphism:
\[
  f_p:\wedge^{n-p}(\underline{E}')\otimes_A\wedge^p(\underline{E}'')
  \longrightarrow F^p\wedge^n(\underline{E})/F^{p-1}\wedge^n(\underline{E})
\]
by putting \(f_p(x'_1\wedge\ldots\wedge x'_{n-p}\otimes x''_1\wedge\ldots\wedge x''_p)=\operatorname{cl}_{\operatorname{Gr}^p}\bigl(u(x'_1)\wedge\ldots\wedge u(x'_{n-p})\wedge x_1\wedge\ldots\wedge x_p\bigr)\), which visibly does not depend on the choice of the \(x_i\) lifting the given sections \(x''_i\) of \(\underline{E}''\). When the exact sequence is locally split, \(f_p\) is an isomorphism, since locally \(\wedge^n(\underline{E})\simeq\bigoplus_{p+q=n}\wedge^p(\underline{E}')\otimes_A\wedge^q(\underline{E}'')\).

When \(\underline{E}',\underline{E},\underline{E}''\) are locally direct summands of free Modules of finite type, the \(\operatorname{Gr}^p\wedge^n(\underline{E})\) are locally direct summands of free Modules of finite type, and therefore so are the \(F^p\wedge^n(\underline{E})\); we may thus write in \(K\):
\(\lambda^n(E)=\sum_{p=0}^{n}\operatorname{cl}(\operatorname{Gr}^p\wedge^n(\underline{E}))=\sum_{p=0}^{n}\lambda^{n-p}(E')\mathbin{.}\lambda^p(E'')\), giving additivity.

Let \((T,\underline{A})\), \((T',\underline{A}')\) be two commutatively ringed topoi, and \(f:(T,\underline{A})\longrightarrow(T',\underline{A}')\) a morphism of ringed topoi. We know (Exp. IV, 2.7) that \(f\) defines a ring morphism on the corresponding \(K\), namely \(f^*:K^{\bullet}_{\naif}(T',\underline{A}')\longrightarrow K^{\bullet}_{\naif}(T,\underline{A})\). Since inverse image commutes with exterior powers, the corresponding morphism \(f^*\) on the \(K\) is a \(\lambda\)-homomorphism for the structure just defined.

Let us mention two special cases:
\begin{itemize}[label=-]
\item Let \(X\) be a scheme; \(K^{\bullet}_{\naif}(X)\) is thus endowed with a pre-\(\lambda\)-ring structure; this is the case that will be used in the rest of the seminar, and it may be regarded as a ``model'' for the general theory of
%%%%%%%%%%  scan idx 330  |  folio 317  %%%%%%%%%%
\SGAPrintPage{317}%
\(\lambda\)-rings.
\item
Let \(A\) be a ring, \(X=\operatorname{Spec}A\); also let \(G\) be a group, acting trivially on \(A\); taking for \((T,A)\) the topos of sheaves of sets on \(X\) on which \(G\) acts, the ring \(K\) thus obtained is the representation ring of \(G\) in the projective \(A\)-modules of finite type.
\end{itemize}

\sgaStableTarget{sga6.e05.para.2-3}2.3. Let \(K\) be a commutative ring with unit; we propose to endow the abelian group \(\Ghat(K)=1+K[[t]]^+\) with a pre-\(\lambda\)-ring structure, functorial in \(K\). For this, we shall apply the results of (1.10) to the functor \(\Ghat\).

We first define multiplication by specifying a \(\mathbb Z\)-bilinear morphism of functors: \(\Ghat\times\Ghat\longrightarrow\Ghat\), by means of the element \(1+XYt\) of \(\Ghat(\mathbb Z[X,Y])\); this element indeed satisfies 1.7.i), and to see that it satisfies 1.7.ii), we must show that in \(\prod_{i=1}^{n}(1+T_iYt)\), expressed in terms of the elementary symmetric functions \(X_j\) of the \(T_i\), \(X_n\) does not occur in the coefficients of \(t^k\) for \(k<n\), which is evident for reasons of degree.

The law thus defined will be denoted by the sign \(\circ\); we therefore have, by definition:
\sgaStableTarget{sga6.e05.eq.2-3-1}\begin{equation}\tag{2.3.1}
  (1+Xt)\circ(1+Yt)=1+XYt.
\end{equation}
It is distributive by construction, and associative and commutative by the diagrams:
\[
\begin{tikzcd}[column sep=tiny]
\Ghat\times\Ghat\times\Ghat \arrow[r] \arrow[d] &
\Ghat\times\Ghat \arrow[d]\\
\Ghat\times\Ghat \arrow[r] & \Ghat
\end{tikzcd}
\qquad
\begin{tikzcd}[column sep=tiny]
\Ghat\times\Ghat \arrow[r] \arrow[d] &
\Ghat\times\Ghat \arrow[d]\\
\Ghat \arrow[r,"\id"] & \Ghat,
\end{tikzcd}
\]
%%%%%%%%%%  scan idx 331  |  folio 318  %%%%%%%%%%
\SGAPrintPage{318}%
where the meaning of the arrows is evident, and whose commutativity, by 1.10, need only be verified on the elements \(1+Xt,1+Yt,1+Zt\) of \(\Ghat(\mathbb Z[X,Y,Z])\) (resp. \(1+Xt,1+Yt\) of \(\Ghat(\mathbb Z[X,Y])\)), a verification that is immediate.

By 1.8, the product of two series under this law is calculated by substituting their coefficients in the universal polynomials of the law. We denote the \(n\)-th universal polynomial by \(P_n\); it is a polynomial in \(\mathbb Z[X_i,Y_j]_{1\leq i,j\leq n}\), isobaric of weight \(n\) in each of the two groups of variables (see 1.15), and symmetric with respect to these two groups. As an example, let us give the value of \(P_2\):
\[
  P_2(X_1,X_2;Y_1,Y_2)=X_2Y_1^2+Y_2X_1^2-2X_2Y_2.
\]

We also have the following relation: let \(K\in\operatorname{Ob}(\mathcal C_0)\), \(\psi\in\Ghat(K)\), \(a\in K\), then
\sgaStableTarget{sga6.e05.eq.2-3-2}\begin{equation}\tag{2.3.2}
  \psi(t)\circ(1+at)=\psi(at).
\end{equation}
As usual, the standard functorial argument reduces it to the case \(K=\mathbb Z[X,Y]\), \(\psi=1+Xt\), \(a=Y\), where it is precisely 2.3.1.

Applying 2.3.2 to \(\Ghat(K)=1+K[[t]]^+\), and to \(a=1\), we see that the element \(1+t\) is the unit for this multiplication. To avoid confusion, we sometimes denote it by \(\underline{1}\), or by \(\underline{1}_\circ\).

The functor \(\Ghat\) can now be regarded as a functor from \(\mathcal C_0\) to the category of commutative rings with unit, and can therefore be composed with itself; to avoid confusion,
we
%%%%%%%%%%  scan idx 332  |  folio 319  %%%%%%%%%%
\SGAPrintPage{319}%
shall choose different indeterminates for each factor in the product of \(\Ghat\) with itself, and indicate the indeterminate as a superscript when necessary: \(\Ghat^t\), \(\Ghat^u\), etc. For \(A\in\operatorname{Ob}(\mathcal C_0)\), we therefore have \(\Ghat^u\circ\Ghat^t(A)=1+\{1+A[[t]]^+\}[[u]]^+\). To define the operations \(\lambda^i\) on \(1+K[[t]]^+\), we shall therefore specify a morphism of functors: \(\Ghat\longrightarrow\Ghat^u\circ\Ghat^t\).
For this, we must first define a filtration of the functor \(\Ghat^u\circ\Ghat^t\) that permits the gluing of 1.6.

We first observe that \(F^k\Ghat\) is an ideal for the law \(\circ\); for this, functoriality reduces us to the case \(A=\mathbb Z[X_i,Y]_{i\geq1}\) and the product \(\left(1+\sum_{i\geq k}X_i t^i\right)\circ(1+Yt)\), and the assertion follows from 2.3.2. We then take as \(F^n(\Ghat^u\circ\Ghat^t)\) the set of series of the form \(1+\psi_1u+\ldots+\psi_i u^i+\ldots\), \(\psi_i\in F^{[n/i]}\Ghat\), where \([n/i]\) denotes the integer part of \(n/i\).

The element defining the morphism \(\lambda_u\) is \(\underline{1}+\{1+Xt\}u\); the reader will verify as an exercise that the conditions of 1.6 are indeed satisfied (in fact, the filtration is chosen ad hoc).

Here again, the operations \(\lambda^i\) are calculated by substitution from \(\lambda_u(1+\ldots+X_nt^n+\ldots)\), and are therefore determined by universal polynomials. We denote by \(Q_{i,j}(X_1,\ldots,X_{ij})\) the coefficient of \(t^j\) in \(\lambda^i(1+\ldots+X_nt^n+\ldots)\); it is an isobaric polynomial of weight \(ij\) in the variables \(X_1,\ldots,X_{ij}\) (a verification analogous to that of 1.15). It is calculated by symmetric functions from the defining element:
\sgaStableTarget{sga6.e05.eq.2-3-3}\begin{equation}\tag{2.3.3}
  \lambda_u(1+Xt)=\underline{1}+\{1+Xt\}u.
\end{equation}

As an example, let us give the value of \(Q_{2,2}\):
%%%%%%%%%%  scan idx 333  |  folio 320  %%%%%%%%%%
\SGAPrintPage{320}%
\[
  Q_{2,2}(X_1,X_2,X_3,X_4)=X_1X_3-X_4.
\]

Finally, let us indicate the origin of formulas (2.3.1) and (2.3.3). We consider Example 2.2.b), where the operations \(\lambda^i\) are defined by the exterior powers of \(A\)-Modules. We want \(\lambda_t:K\longrightarrow\Ghat(K)\) to be a \(\lambda\)-homomorphism; we must therefore define laws on \(\Ghat\) such that for all \(A\)-Modules \(\underline{E},\underline{F}\), we have:
\[
  \lambda_t(EF)=\lambda_t(E)\circ\lambda_t(F)~;\qquad
  \lambda_u(\lambda_t(E))=\Ghat^u(\lambda_t)(\lambda_u(E))~;
\]
now relations (2.3.1) and (2.3.3) are precisely the preceding relations applied to the case where \(\underline{E}\) and \(\underline{F}\) are invertible modules.

We shall see in Exp. VI to what extent the desired property is satisfied.

\sgaStableTarget{sga6.e05.def.2-4}\textbf{Definition 2.4.} \uline{A pre}-\(\lambda\)-\uline{ring} \(K\) \uline{will be called a} \(\lambda\)-\uline{ring if the homomorphism of abelian groups} \(\lambda_t:K\longrightarrow1+K[[t]]^+\) \uline{is a} \(\lambda\)-\uline{homomorphism} \uline{for the pre}-\(\lambda\)-\uline{ring structure on the second member defined in} 2.3.

This therefore means that we have the relations: \(\forall x,y\in K\),
\sgaStableTarget{sga6.e05.eq.2-4-1}\begin{equation}\tag{2.4.1}
  \lambda_t(1)=1+t~;\quad
  \lambda_t(xy)=\lambda_t(x)\circ\lambda_t(y)~;\quad
  \lambda_u(\lambda_t(x))=\Ghat(\lambda_t)\bigl(\lambda_u(x)\bigr).
\end{equation}

In terms of the maps \(\lambda^i\), the last two relations are equivalent to: \(x,y\in K\), \(i,j\in\mathbb N\),
\sgaStableTarget{sga6.e05.eq.2-4-2}\sgaStableTarget{sga6.e05.eq.2-4-3}\begin{align}
  \lambda^i(xy)
  &=P_i(\lambda^1(x),\ldots,\lambda^i(x);
       \lambda^1(y),\ldots,\lambda^i(y)) \tag{2.4.2}\\
  \lambda^j(\lambda^i(x))
  &=Q_{i,j}(\lambda^1(x),\ldots,\lambda^{ij}(x)) \notag
\end{align}

%%%%%%%%%%  scan idx 334  |  folio 321  %%%%%%%%%%
\SGAPrintPage{321}%
We note that here we depart from the usual terminology (cf. [1]), where the name \(\lambda\)-ring was used for what we call a pre-\(\lambda\)-ring, and ``special'' \(\lambda\)-rings for our \(\lambda\)-rings; the reason is that only ``special'' \(\lambda\)-rings are usable in practice.

\sgaStableTarget{sga6.e05.ex.2-5}\textbf{Examples 2.5.}
\begin{enumerate}[label=\alph*)]
\item We have determined all the pre-\(\lambda\)-ring structures on \(\mathbb Z\); it follows from 2.4.1 that only the one defined by \(\lambda_t(1)=1+t\) can be a \(\lambda\)-ring structure. We then have \(\lambda_t(n)=(1+t)^n\), hence:
\[
  \forall n,n'\in\mathbb Z,\quad
  \lambda_t(n)\circ\lambda_t(n')=(1+t)^n\circ(1+t)^{n'}
  =(1+t)^{nn'}=\lambda_t(nn')~;
\]
\[
  \lambda_u((1+t)^n)
  =
  (\lambda_u(1+t))^n
  =
  (\underline{1}+\{1+t\}u)^n
  =
  \Ghat(\lambda_t)((1+u)^n)
  =
  \Ghat(\lambda_t)(\lambda_u(n)).
\]
Consequently, this is indeed a \(\lambda\)-ring structure. Denoting by \(\binom{X}{k}\) the polynomial with rational coefficients \(\frac{1}{k!}X(X-1)\ldots(X-k+1)\), we therefore have:
\sgaStableTarget{sga6.e05.eq.2-5-1}\begin{equation}\tag{2.5.1}
  \lambda^i(n)=\binom{n}{i}.
\end{equation}

\item For every ring \(K\), the pre-\(\lambda\)-ring \(1+K[[t]]^+\) is a \(\lambda\)-ring. First, \(\lambda_u(\underline{1})=\lambda_u(1+t)=\underline{1}+\{1+t\}u=\underline{1}+\underline{1}.u\). We verify the other relations by the commutativity of the diagrams:
\[
\begin{tikzcd}[column sep=tiny]
\Ghat^t\times\Ghat^t \arrow[r,"\lambda_u\times\lambda_u"] \arrow[d,"\circ"'] &
(\Ghat^u\circ\Ghat^t)\times(\Ghat^u\circ\Ghat^t) \arrow[d,"\circ"]\\
\Ghat^t \arrow[r,"\lambda_u"] &
\Ghat^u\circ\Ghat^t,
\end{tikzcd}
\qquad
\begin{tikzcd}[column sep=tiny]
\Ghat^t \arrow[r,"\lambda_v"] \arrow[d,"\lambda_u"'] &
\Ghat^v\circ\Ghat^t \arrow[d,"\Ghat^v(\lambda_u)"]\\
\Ghat^u\circ\Ghat^t \arrow[r,"\lambda_v"] &
\Ghat^v\circ\Ghat^u\circ\Ghat^t,
\end{tikzcd}
\]
%%%%%%%%%%  scan idx 335  |  folio 322  %%%%%%%%%%
\SGAPrintPage{322}%
whose commutativity is verified, by 1.10, on the elements \(1+Xt,1+Yt\). For the first,
\[
\begin{aligned}
\lambda_u(1+Xt)\circ\lambda_u(1+Yt)
&=(\underline{1}+\{1+Xt\}u)\circ(\underline{1}+\{1+Yt\}u)\\
&=\underline{1}+\{(1+Xt)\circ(1+Yt)\}u\\
&=\underline{1}+\{1+XYt\}u\\
&=\lambda_u((1+Xt)\circ(1+Yt)).
\end{aligned}
\]
For the second:
\[
  \lambda_v(\lambda_u(1+Xt))
  =
  \underline{1}+\{\underline{1}+\{1+Xt\}u\}v
  =
  \Ghat^v(\lambda_u)(\lambda_v(1+Xt)).
\]

\item The geometric examples of 2.2 will be treated later.
\end{enumerate}

\sgaStableTarget{sga6.e05.para.2-6}2.6. We can generalize Example 2.5.a) as follows: let \(K\) be a commutative ring with unit, such that:
\[
\begin{array}{ll}
\text{i)}& K\text{ is torsion-free},\\
\text{ii)}& \forall x\in K,\ \forall n\in\mathbb N,\quad \text{we have}\ \binom{x}{n}\in K
\end{array}
\]
We then put, for every \(x\in K\), and every \(n\in\mathbb N\), \(\lambda^n(x)=\binom{x}{n}\), that is:
\sgaStableTarget{sga6.e05.eq.2-6-1}\begin{equation}\tag{2.6.1}
  \lambda_t(x)=\sum_{n=0}^{\infty}\binom{x}{n}t^n
  =(1+t)^x=\exp(x\Log(1+t)),
\end{equation}
where the last identity is considered in \((K\otimes_{\mathbb Z}\mathbb Q)[[t]]\). By the multiplicativity of the exponential, this defines a pre-\(\lambda\)-ring structure on \(K\).

To see that it is a \(\lambda\)-ring structure, we must prove relations 2.4.2, and it suffices to prove them in \(K\otimes_{\mathbb Z}\mathbb Q\), by 2.6.i). Functoriality then reduces us to the case \(K=\mathbb Q[X,Y]\) and the elements \(X,Y\). In view of the definition of the \(\lambda^i\), the two members of the
%%%%%%%%%%  scan idx 336  |  folio 323  %%%%%%%%%%
\SGAPrintPage{323}%
equalities 2.4.2 are polynomials in \(X,Y\), with rational coefficients, and to prove their equality it therefore suffices to show that they take the same value at all integer values of \(X,Y\). But the relations obtained by substituting integer values for \(X,Y\) are those expressing that the pre-\(\lambda\)-ring \(Z\) (2.5.a)) is a \(\lambda\)-ring; they are therefore satisfied.

\sgaStableTarget{sga6.e05.def.2-7}\textbf{Definition 2.7.} \uline{By a} binomial ring \uline{we mean a ring satisfying conditions} 2.6.i) \uline{and} ii), \uline{endowed with the} \(\lambda\)-\uline{ring structure defined by} 2.6.1.

\sgaStableTarget{sga6.e05.ex.2-8}\textbf{Example 2.8.} Let \(X\) be a site, \(K\) a binomial ring, and \(\underline{K}\) the sheaf associated to the constant presheaf with value \(K\) on \(X\); then \(H^0(X,\underline{K})\) is a binomial ring.

Indeed, conditions 2.6 are preserved under passage to an inductive limit.

In particular, \(H^0(X,\underline{\mathbb Z})\) is a binomial ring.

\sgaStableTarget{sga6.e05.para.2-9}2.9. Let \(K\) be a binomial ring, and \(A\) a \(K\)-algebra. For every series \(\varphi\in\Ghat(A)\), and every \(a\in K\), we can define a series \(\varphi^a\) in the usual way: we write \(\varphi=1+\psi\), and put:
\sgaStableTarget{sga6.e05.eq.2-9-1}\begin{equation}\tag{2.9.1}
  \varphi^a=(1+\psi)^a
  =
  \sum_{n=0}^{\infty}\binom{a}{n}\psi^n
  \left(=\exp(a\Log\varphi)\right),
\end{equation}
the last equality being writable only modulo torsion. These formulas are obtained by substitution from the analogous formulas for \(\varphi=1+t\); the following relations, the classical properties of the exponential, therefore remain valid for \(a,b\in K\), \(\varphi,\psi\in\Ghat(A)\):
\sgaStableTarget{sga6.e05.eq.2-9-2}\begin{equation}\tag{2.9.2}
  \varphi^a\varphi^b=\varphi^{a+b};\qquad
  \varphi^{ab}=(\varphi^a)^b;\qquad
  \varphi^a\psi^a=(\varphi\psi)^a .
\end{equation}

% Native English continuation macros retained at their source boundary.
\providecommand{\Ghat}{\widehat G}
\providecommand{\Gg}{\widehat G_g}
\providecommand{\Gcirc}{\widehat G_{\circ}}
\providecommand{\Gnabla}{\widehat G_{\nabla}}
\providecommand{\Grhat}{\widehat{\operatorname{Gr}}}
\providecommand{\Log}{\operatorname{Log}}
\providecommand{\Fil}{\operatorname{Fil}}
\providecommand{\eps}{\varepsilon}
\providecommand{\id}{\operatorname{id}}
\providecommand{\Sym}{\operatorname{Sym}}
\providecommand{\Mon}{\operatorname{Mon}}

% Exposé V continuation macros
\providecommand{\Ghat}{\widehat G}
\providecommand{\Gg}{\widehat G_g}
\providecommand{\Gcirc}{\widehat G_{\circ}}
\providecommand{\Gnabla}{\widehat G_{\nabla}}
\providecommand{\Grhat}{\widehat{\operatorname{Gr}}}
\providecommand{\Log}{\operatorname{Log}}
\providecommand{\Fil}{\operatorname{Fil}}
\providecommand{\eps}{\varepsilon}
\renewcommand{\id}{\operatorname{id}}
\providecommand{\Sym}{\operatorname{Sym}}
\providecommand{\Mon}{\operatorname{Mon}}

%%%%%%%%%%  scan idx 337  |  folio 324  %%%%%%%%%%
\SGAPrintPage{324}%
\sgaStableTarget{sga6.e05.continuation-of-2-9}We still suppose that \(a\in K\). Multiplication by \(\lambda_t(a)\) under the law \(\circ\) defines a morphism of functors: \(\Ghat_K\to\Ghat_K\), characterized by its value on the element \(1+Xt\) of \(\Ghat(K[X])\). We therefore have:
\[
  \lambda_t(a)\circ(1+Xt)=\lambda_{Xt}(a)\qquad\text{by (2.3.2), hence:}
\]
\sgaStableTarget{sga6.e05.eq.2-9-3}\begin{equation}\tag{2.9.3}
  \lambda_t(a)\circ(1+Xt)
  =(1+Xt)^a
  =\exp(a\Log(1+Xt))
  =\sum_{n=0}^{\infty}\binom{a}{n}X^nt^n .
\end{equation}
We deduce:
\sgaStableTarget{sga6.e05.eq.2-9-4}\begin{equation}\tag{2.9.4}
  \lambda_t(a)\circ\varphi=\varphi^a
  \quad(=\exp(a\Log\varphi)) ;
\end{equation}
Indeed, the two members are additive functors equal on \(1+Xt\). From (2.9.3) and 1.15, it follows that the \(n\)-th universal polynomial defining multiplication by \(\lambda_t(a)\) is isobaric of weight \(n\); it will be denoted by \(R_{n,a}\).

Finally, note that if \(A\) is a \(K\)-\(\lambda\)-algebra, i.e. a \(\lambda\)-ring endowed with a \(\lambda\)-homomorphism: \(K\to A\), and if \(b\in A\), we have:
\sgaStableTarget{sga6.e05.eq.2-9-5}\begin{equation}\tag{2.9.5}
  \lambda_t(ab)=(\lambda_t(b))^a .
\end{equation}
This follows from \(\lambda_t(ab)=\lambda_t(a)\circ\lambda_t(b)\), and from (2.9.4).

\sgaStableTarget{sga6.e05.sec.3}\section*{3. The operations \(\gamma\)}

We still work in the category \(\mathcal C_0\).

\sgaStableTarget{sga6.e05.para.3-1}3.1. Let \(K\) be a \(\lambda\)-ring. From the operations \(\lambda^i\) of \(K\), we shall define new operations, denoted by \(\gamma^i\). They are intended to introduce a notion of Chern classes under conditions that will be specified later. For the moment, let us give a ``justification'' of the following definitions: if \(\underline{E}\) is a vector bundle of rank \(d\), the \(d\)-th Chern class of \(\underline{E}\) will have to be (whatever theory of Chern classes is considered) the ``alternating sum of the exterior powers
%%%%%%%%%%  scan idx 338  |  folio 325  %%%%%%%%%%
\SGAPrintPage{325}%
of \(\underline{E}\)'';
to define the \(n\)-th Chern class, we shall proceed analogously, bringing \(\underline{E}\) to rank \(n\) by considering \((\cl(\underline{E})-d)+n\), and we shall take the alternating sum of the exterior powers of the element thus obtained; the following proposition is to be ``interpreted'' by setting \(x=\cl(\underline{E})-d\):

\sgaStableTarget{sga6.e05.prop.3-2}\textbf{Proposition 3.2.} \uline{Let} \(K\) \uline{be a} \(\lambda\)-\uline{ring}.
\begin{enumerate}[label=\roman*)]
\item \(\forall x\in K,\ \forall n\in\mathbb N,\quad \lambda^n(x+n-1)=(-1)^n\sum_{i=0}^n(-1)^i\lambda^i(x+n).\)
\item \uline{Putting} \(\gamma^n(x)=\lambda^n(x+n-1)\), \uline{and} \(\gamma_t(x)=\sum_{n\ge0}\gamma^n(x)t^n\), \uline{we have}:
\[
  \gamma_t(x)=\lambda_{t/(1-t)}(x),\qquad\text{\uline{hence}}
\]
\[
  \lambda_t(x)=\gamma_{t/(1+t)}(x).
\]
\end{enumerate}

We therefore see that knowing the operations \(\gamma\) is equivalent to knowing the operations \(\lambda\).

\uline{Proof}.

i) We have
\[
  \lambda^n(x+n-1)=\sum_{i=0}^{n}\lambda^i(x+n)\lambda^{n-i}(-1).
\]
From \(\lambda_t(1)=1+t\), we obtain: \(\lambda_t(-1)=1/(1+t)\), hence: \(\lambda^k(-1)=(-1)^k\), which gives the required relation.

ii) \(\lambda_{t/(1-t)}(x)=\sum_{n\ge0}\lambda^n(x).t^n(1-t)^{-n}.\)
We expand \(t^n(1-t)^{-n}=t^n\sum_{k\ge0}\binom{-n}{k}(-1)^kt^k=\sum_{p\ge n}\binom{p-1}{p-n}t^p\). \(\lambda_{t/(1-t)}(x)=\sum_{n\ge0}\bigl[\lambda^n(x).\sum_{p\ge n}\binom{p-1}{p-n}t^p\bigr]\); the coefficient of \(t^p\) is therefore:
\[
  \sum_{n=0}\lambda^n(x).\binom{p-1}{p-n}
  =\sum_{n=0}\lambda^n(x)\lambda^{p-n}(p-1)
  =\lambda^p(x+p-1).
\]

%%%%%%%%%%  scan idx 339  |  folio 326  %%%%%%%%%%
\SGAPrintPage{326}%
\sgaStableTarget{sga6.e05.para.3-3}3.3.
Let \(K\), \(K'\) be two \(\lambda\)-rings, and \(f:K\to K'\) a ring homomorphism. By Definition 3.2.ii), \(f\) is a \(\lambda\)-homomorphism if and only if, for every \(i\in\mathbb N\), and every \(x\in K\), \(f(\gamma^i(x))=\gamma^i(f(x))\), or equivalently \(\gamma_t(f(x))=\Gg(f)\circ\gamma_t(x)\).

\sgaStableTarget{sga6.e05.para.3-4}3.4. From relation 3.2.ii), we deduce the relations:
\sgaStableTarget{sga6.e05.eq.3-4-1}\begin{equation}\tag{3.4.1}
  \gamma^0(x)=1,
  \qquad \gamma^1(x)=x,
  \qquad \gamma^n(x+y)=\sum_{p+q=n}\gamma^p(x)\gamma^q(y).
\end{equation}
They show that the homomorphism \(\gamma_t\) defines on \(K\) a new pre-\(\lambda\)-ring structure; we shall see that in general it is not a \(\lambda\)-ring structure. Denoting by \(\rho_t\) the substitution in the ring of formal series in \(t\) with coefficients in \(K\) defined by \(\rho_t(t)=t/(1-t)\), we have:
\[
  \gamma_t=\rho_t\circ\lambda_t.
\]
The automorphism \(\rho_t\) is functorial in \(K\), and defines an automorphism of the functor \(\Ghat\) considered as a functor to abelian groups; with this automorphism, by transport of structure, we can define on \(\Ghat\) a new structure as a functor to \(\lambda\)-rings, such that for every \(\lambda\)-ring \(K\), the homomorphism \(\gamma_t:K\to\Ghat(K)\) is a \(\lambda\)-homomorphism; we shall make explicit the multiplication and the operations \(\lambda^i\) of this second structure.

\sgaStableTarget{sga6.e05.para.3-5}3.5. The multiplication is defined by the homomorphism of functors:
\[
  \Ghat\times\Ghat\xrightarrow{(\rho^{-1},\rho^{-1})}
  \Ghat\times\Ghat\xrightarrow{\circ}\Ghat\xrightarrow{\rho}\Ghat ,
\]
where the sign \(\circ\) denotes the multiplication defined in 2.3. This new multiplication will be denoted by \(v\). As usual, it is defined by its value on \((1+Xt,1+Yt)\); we therefore have:
%%%%%%%%%%  scan idx 340  |  folio 327  %%%%%%%%%%
\SGAPrintPage{327}%
\sgaStableTarget{sga6.e05.eq.3-5-1}\begin{align}
 (1+Xt)\nabla(1+Yt)
 &=\rho\!\left(\rho^{-1}(1+Xt)\circ\rho^{-1}(1+Yt)\right) \notag\\
 &=\rho\!\left(\frac{1+(X+1)t}{1+t}\circ
                 \frac{1+(Y+1)t}{1+t}\right) \notag\\
 &=\rho\!\left(\frac{(1+(X+1)(Y+1)t)(1+t)}
                      {(1+(X+1)t)(1+(Y+1)t)}\right) \notag\\
 &=\frac{1+(X+Y+XY)t}{(1+Xt)(1+Yt)} . \tag{3.5.1}
\end{align}

If \(K\) is a \(\lambda\)-ring, then for all \(x,y\in K\) we have:
\sgaStableTarget{sga6.e05.eq.3-5-2}\begin{equation}\tag{3.5.2}
  \gamma_t(xy)=\gamma_t(x)\nabla\gamma_t(y).
\end{equation}

We denote by \(P'_n\) the \(n\)-th universal polynomial of the law \(\nabla\); it is a polynomial in \(\mathbb Z[X_i,Y_j]_{1\le i,j\le n}\) whose isobaric components have weights ranging from \(n\) to \(2n\) with respect to all the variables \(X_i,Y_j\), (1.15) and which is symmetric with respect to these two groups of variables. For \(x,y\in K\), we therefore have:
\sgaStableTarget{sga6.e05.eq.3-5-3}\begin{equation}\tag{3.5.3}
  n\in\mathbb N,\quad \gamma^n(xy)
  =P'_n(\gamma^1(x),\ldots,\gamma^n(x),
       \gamma^1(y),\ldots,\gamma^n(y)).
\end{equation}

Finally, note that the law \(\nabla\) has a unit element, namely \(\rho_t(1+t)\), that is, \(1/(1-t)\); it will sometimes be denoted by \(\underline{1}_\nabla\).

\sgaStableTarget{sga6.e05.para.3-6}3.6. The ``operations \(\lambda\)'' are defined by the commutative diagram:
\sgaStableTarget{sga6.e05.eq.3-6-1}\[
\begin{tikzcd}[column sep=tiny]
1+K[[t]]^+ \arrow[r,"\lambda_u"] \arrow[d,"\rho_t"'] &
1+\{1+K[[t]]^+\}[[u]]^+ \arrow[d,"\Ghat(\rho_t)"']\\
1+K[[t]]^+ \arrow[r,"\lambda'_u"] &
1+\{1+K[[t]]^+\}[[u]]^+ ,
\end{tikzcd}
\tag{3.6.1}
\]
i.e. by the composite morphism of functors:
\[
  \lambda' : \Ghat\xrightarrow{\rho_t^{-1}}\Ghat\xrightarrow{\lambda_u}\Ghat\circ\Ghat\xrightarrow{\Ghat(\rho_t)}\Ghat\circ\Ghat ,
\]
and will be denoted by \(\lambda'\). They are defined by their value on \(1+Xt\); we have:
%%%%%%%%%%  scan idx 341  |  folio 328  %%%%%%%%%%
\SGAPrintPage{328}%
\[
\begin{aligned}
\lambda'_u(1+Xt) &=\Ghat(\rho_t)(\lambda_u(\rho_t^{-1}(1+Xt)))\\
&=\Ghat(\rho_t)\left(\lambda_u\left(\frac{1+(X+1)t}{1+t}\right)\right)\\
&=\Ghat(\rho_t)\left(\frac{\underline{1}_o+\{1+(X+1)t\}u}{\underline{1}_o+\{1+t\}u}\right)\\
&=\underline{1}_\nabla+\frac{\{1+Xt\}u}{\underline{1}_\nabla+\underline{1}_\nabla.u}\\
&=\rho_u^{-1}(\underline{1}_\nabla+\{1+Xt\}u).
\end{aligned}
\]

Denoting by \(\gamma'\) the operations \(\gamma\) corresponding to the operations \(\lambda'\), we therefore have:
\sgaStableTarget{sga6.e05.eq.3-6-2}\begin{equation}\tag{3.6.2}
  \gamma'_u(1+Xt)=\underline{1}_\nabla+\{1+Xt\}u .
\end{equation}

We denote by \(Q'_{i,j}\) the coefficient of \(t^j\) in \(\gamma'^i(1+\ldots+X_nt^n+\ldots)\); the reader can verify that it is a polynomial in \(Z[X_1,\ldots,X_{ij}]\) whose isobaric components have weights ranging between \(j\) and \(ij\).

Let \(K\) be a \(\lambda\)-ring, \(x\in K\). From diagram (3.6.1), we deduce the relations:
\sgaStableTarget{sga6.e05.eq.3-6-3}\begin{equation}\tag{3.6.3}
  \forall i\in\mathbb N,\quad\gamma_t(\lambda^i(x))=\lambda'^i(\gamma_t(x))~;
  \qquad
  \gamma_t(\gamma^i(x))=\gamma'^i(\gamma_t(x)).
\end{equation}
\sgaStableTarget{sga6.e05.eq.3-6-4}\begin{equation}\tag{3.6.4}
  \gamma^j(\gamma^i(x))
  =Q'_{i,j}(\gamma^1(x),\ldots,\gamma^{ij}(x)).
\end{equation}

Finally, we observe that the pre-\(\lambda\)-ring structure just defined on \(\Ghat\) is a \(\lambda\)-ring structure, since it is obtained from the one defined in 2.3 by the isomorphism \(\rho\).

\sgaStableTarget{sga6.e05.para.3-7}3.7. To avoid confusion, the functor \(\Ghat\), endowed with the structure of a functor to \(\lambda\)-rings defined in §2 (resp. in §3), will be denoted by \(\Ghat_o\) (resp. \(\Gnabla\)). It follows from the preceding considerations that to specify a
%%%%%%%%%%  scan idx 342  |  folio 329  %%%%%%%%%%
\SGAPrintPage{329}%
\(\lambda\)-ring structure on a ring \(A\), it is equivalent to specify a ring homomorphism \(A\to\Ghat(A)_o\) such that diagram 3.7.1 is commutative, or a ring homomorphism \(A\to\Ghat(A)_\nabla\) such that diagram 3.7.2 is commutative; the first defines the operations \(\lambda\), the second the operations \(\gamma\).
\sgaStableTarget{sga6.e05.eq.3-7-1}\[
\begin{tikzcd}[column sep=tiny]
A \arrow[r,"\lambda_u"] \arrow[d,"\lambda_t"'] & \Ghat^u(A)_o \arrow[d,"\Ghat^u(\lambda_t)"]\\
\Ghat^t(A)_o \arrow[r,"\lambda_u"] & \Ghat^u(\Ghat^t(A)_o)_o
\end{tikzcd}
\tag{3.7.1}
\]
\sgaStableTarget{sga6.e05.eq.3-7-2}\[
\begin{tikzcd}[column sep=tiny]
A \arrow[r,"\gamma_u"] \arrow[d,"\gamma_t"'] & \Ghat^u(A)_\nabla \arrow[d,"\Ghat^u(\gamma_t)"]\\
\Ghat^t(A)_\nabla \arrow[r,"\gamma'_u"] & \Ghat^u(\Ghat^t(A)_\nabla)_\nabla.
\end{tikzcd}
\tag{3.7.2}
\]

\sgaStableTarget{sga6.e05.para.3-8}3.8. Let \(K\) be a binomial ring (2.7); from (2.6.1) we deduce, \(\forall x\in K\), the relations:
\sgaStableTarget{sga6.e05.eq.3-8-1}\begin{equation}\tag{3.8.1}
  \gamma_t(x)=(1/(1-t))^x
  =\exp(-x\Log(1-t))
\end{equation}

Let \(a\) be an element of \(K\); multiplication by \(\gamma_t(a)\) under the law \(\nabla\) is characterized by its value on \(1+Xt\), namely:
\[
  \gamma_t(a)\nabla(1+Xt)=\rho_t\left(\lambda_t(a)\circ\frac{1+(X+1)t}{1+t}\right)=\rho_t\left(\left(\frac{1+(X+1)t}{1+t}\right)^a\right)
\]
\sgaStableTarget{sga6.e05.eq.3-8-2}\begin{equation}\tag{3.8.2}
  \gamma_t(a)\nabla(1+Xt)
  =(1+Xt)^a
  =\exp(a\Log(1+Xt))
  =\sum_{n=0}^{\infty}\binom{a}{n}X^nt^n .
\end{equation}
In particular, we see that the \(n\)-th universal polynomial defining multiplication by \(\gamma_t(a)\) under the law \(\nabla\) is the polynomial \(R_{n,a}\) (2.9).

%%%%%%%%%%  scan idx 343  |  folio 330  %%%%%%%%%%
\SGAPrintPage{330}%
\sgaStableTarget{sga6.e05.def.3-9}\textbf{Definition 3.9.} \uline{Let} \(K\) \uline{be a} \(\lambda\)-\uline{ring}; \uline{by an} augmented \(K\)-\(\lambda\)-algebra \uline{we mean the data of a} \(K\)-\(\lambda\)-\uline{algebra} \(A\) \uline{and a} \(\lambda\)-\uline{homomorphism of} \(K\)-\uline{algebras} \(A\to K\); a homomorphism of augmented \(K\)-\(\lambda\)-algebras \uline{is a} \(\lambda\)-\uline{homomorphism compatible with the augmentations}.

\sgaStableTarget{sga6.e05.ex.3-9-1}\textbf{Example 3.9.1.}: Let \((X,A)\) be a commutatively locally ringed topos, and \(K^\bullet(X)\) the Grothendieck group of the category of sheaves of locally free \(A\)-modules of finite type on \(X\). We have seen that \(K^\bullet(X)\) can be endowed with a pre-\(\lambda\)-ring structure (2.2.b, since every direct summand of a free module of finite type is then locally free of finite type) and that \(H^0(X,\underline{\mathbb Z})\) is canonically endowed with a binomial ring structure. We define a \(\lambda\)-homomorphism from \(H^0(X,\underline{\mathbb Z})\) to \(K^\bullet(X)\) as follows: let \(a\) be a section of the constant sheaf \(\underline{\mathbb Z}\) on \(X\); it defines a partition of the final object of \(X\) into a sum of disjoint opens \(U_n\), the section \(a\) taking the value \(n\) on the open \(U_n\); let \(\underline{F}\) be the sheaf equal to \(A^n\) on \(U_n\) for \(n\ge0\), and to \(0\) on \(U_n\) for \(n\le0\), and let \(\underline{F}'\) be the sheaf equal to \(A^{-n}\) on \(U_n\) for \(n\le0\), and to \(0\) on \(U_n\) for \(n\ge0\); we associate to \(a\) the element \(F-F'\) of \(K^\bullet(X)\).

It is immediate to verify that this does give a \(\lambda\)-homomorphism from \(H^0(X,\underline{\mathbb Z})\) to \(K^\bullet(X)\); we can then define the augmentation from \(K^\bullet(X)\) to \(H^0(X,\underline{\mathbb Z})\) by assigning to every locally free sheaf of finite type its rank, which is a section of \(\underline{\mathbb Z}\); additivity of rank implies that this defines a function on \(K^\bullet(X)\) with values in \(H^0(X,\underline{\mathbb Z})\). We have thus endowed \(K^\bullet(X)\) with an augmented pre-\(\lambda\)-algebra structure. This is always the structure referred to below.

%%%%%%%%%%  scan idx 344  |  folio 331  %%%%%%%%%%
\SGAPrintPage{331}%
\sgaStableTarget{sga6.e05.para.3-10}3.10
Let \(A\) be an augmented \(K\)-\(\lambda\)-algebra. We define a ring filtration on \(A\) as follows: for every \(n\ge0\), we denote by \(\Fil^n A\) the ideal generated by the products:
\sgaStableTarget{sga6.e05.eq.3-10-1}\begin{equation}\tag{3.10.1}
  \gamma^{i_1}(x_1)\ldots\gamma^{i_k}(x_k),
  \qquad \text{with } i_1+\ldots+i_k\ge n,\ \text{and } \forall i,\ \epsilon(x_i)=0
\end{equation}
\((\epsilon:A\to K\) being the augmentation).

Suppose that \(A\) is generated as an abelian group by a family of elements \((x_\alpha)_{\alpha\in I}\); then \(\Fil^n A\) is generated as a \(K\)-module by the products:
\sgaStableTarget{sga6.e05.eq.3-10-2}\begin{equation}\tag{3.10.2}
  \gamma^{i_1}(x_{\alpha_1}-\epsilon(x_{\alpha_1}))\ldots
  \gamma^{i_k}(x_{\alpha_k}-\epsilon(x_{\alpha_k})),
  \qquad \text{with } i_1+\ldots+i_k\ge n,\ \text{and } \alpha_i\in I
\end{equation}
Indeed, if \(x\) is an element of \(A\) with zero augmentation and we write \(x=\sum_i x_{\alpha_i}\), we have \(\sum_i \epsilon(x_{\alpha_i})=0\), hence \(x=\sum_i x_{\alpha_i}-\epsilon(x_{\alpha_i})\); furthermore, if \(a\in A\), we can write \(a=(a-\epsilon(a))+\epsilon(a)=\gamma^1(a-\epsilon(a))+\epsilon(a)\); using relations (3.4.1), we therefore see that the elements of \(\Fil^n A\) are linear combinations with coefficients in \(K\) of the elements (3.10.2); since these belong to \(\Fil^n A\), we obtain the required result.

The filtration thus defined is a ring filtration; it is decreasing, and we have:
\[
  \Fil^0A=A~;
  \qquad \Fil^1A=\Ker\epsilon.
\]
Moreover, it is functorial in \(A\): every \(\lambda\)-homomorphism of augmented \(K\)-\(\lambda\)-algebras is compatible with the filtrations. When we regard an augmented \(\lambda\)-ring as endowed with a filtration, this will mean the filtration above, unless expressly stated otherwise.

%%%%%%%%%%  scan idx 345  |  folio 332  %%%%%%%%%%
\SGAPrintPage{332}%
\sgaStableTarget{sga6.e05.prop.3-11}\textbf{Proposition 3.11.} \uline{Let} \(A\) \uline{be a} \(\lambda\)-ring, \uline{and} \(y_i\), \(i=1,\dots,n\) \uline{elements of} \(A\) \uline{such that} \(\lambda^j(y_i)=0\) \uline{for} \(j\ge2\). \uline{Put} \(x=\sum_i y_i\); \uline{then}:
\begin{enumerate}[label=\roman*)]
\item \uline{for} \(k>n\), we have \(\gamma^k(x-n)=0\)
\item \uline{for} \(1\le k\le n\), \(\gamma^k(x-n)\) \uline{is the} \(k\)-th \uline{symmetric function of the} \(y_i-1\).
\end{enumerate}
Let \(j\ge2\). We then have \(\gamma^j(y_i-1)=\lambda^j(y_i+j-2)=\sum_{p+q=j}\lambda^p(y_i)\lambda^q(j-2)\). Since \(\lambda^p(y_i)=0\) for \(p\ge2\), and \(\lambda^q(j-2)=0\) for \(q>j-2\), \(\gamma^j(y_i-1)=0\). The proposition then follows from \(\gamma_t(x-n)=\prod_{i=1}^{n}\gamma_t(y_i-1)\).

\sgaStableTarget{sga6.e05.sec.4}\section*{4. \(\lambda\)-rings generated by generators and relations}

We still work over \(\mathcal C_0\).

\sgaStableTarget{sga6.e05.para.4-1}4.1. Let \(K\) be a \(\lambda\)-ring. Consider a family of indeterminates \((X_i)_{i\in I}\), and let \(A=K[X_i,\lambda^n X_i]_{i\in I,\ n\ge2}\); we wish to define on \(A\) a \(\lambda\)-ring structure extending that of \(K\).

To specify a homomorphism \(\lambda_t:A\to\Ghat(A)\), it suffices to specify the images of the \(X_i\), \(\lambda^n X_i\), which defines a ring homomorphism: \(A\to\Ghat(A)_o\); for it to define a \(\lambda\)-ring structure on \(A\), we must verify the commutativity of diagram (3.7.1). We therefore put:
\sgaStableTarget{sga6.e05.eq.4-1-1}\begin{equation}\tag{4.1.1}
  \lambda_t(X_i)=1+X_it+\sum_{n\ge2}\lambda^n X_i\,t^n~;
  \qquad
  \lambda_t(\lambda^n X_i)=\lambda^n(\lambda_t(X_i)).
\end{equation}
Since diagram (3.7.1) is a diagram of ring homomorphisms, it suffices to verify its commutativity on the \(X_i\), \(\lambda^n X_i\) (it is known by hypothesis on the elements of \(K\)); this follows readily from definitions (4.1.1) and the fact that \(\lambda_u\) and \(\Ghat^u(\lambda_t)\) are \(\lambda\)-homomorphisms.

%%%%%%%%%%  scan idx 346  |  folio 333  %%%%%%%%%%
\SGAPrintPage{333}%
The \(\lambda\)-ring \(A\) is characterized by the following universal property:

\sgaStableTarget{sga6.e05.prop.4-2}\textbf{Proposition 4.2.} \uline{For every} \(K\)-\(\lambda\)-\uline{algebra} \(B\), \uline{and every family} \((b_i)_{i\in I}\) \uline{of elements of} \(B\), \uline{there exists exactly one} \(K\)-\(\lambda\)-\uline{homomorphism} \(f\) \uline{from} \(A\) \uline{to} \(B\) \uline{such that for every} \(i\in I\), \(f(X_i)=b_i\).

For \(f\) to be a \(\lambda\)-homomorphism, we must have \(f(\lambda^j X_i)=\lambda^j(b_i)\) for all \(i,j\); this defines \(f\) uniquely. The homomorphism thus obtained is indeed a \(\lambda\)-homomorphism, since to verify that \(\lambda_t\circ f=\Ghat(f)\circ\lambda_t\) it suffices to verify it on a system of generators of \(A\) over \(K\), which follows from the definition of \(f\) and relations (2.4.2).

\sgaStableTarget{sga6.e05.para.4-3}4.3. We can make the same construction using the \(\gamma\) in place of the \(\lambda\). Let \(A'=K[X_i,\gamma^n X_i]_{i\in I,\ n\ge2}\); to define the \(\lambda\)-ring structure on \(A'\), we specify a ring homomorphism: \(A'\to\Ghat(A')_\nabla\), by:
\sgaStableTarget{sga6.e05.eq.4-3-1}\begin{equation}\tag{4.3.1}
  \gamma_t(X_i)=1+X_it+\sum_{n\ge2}\gamma^n X_i\mathbin{.}t^n~;
  \qquad
  \gamma_t(\gamma^n X_i)=\gamma'^n(\gamma_t(X_i)).
\end{equation}
We then verify the commutativity of (3.7.2), which shows that we do obtain a \(\lambda\)-ring structure.

The \(\lambda\)-ring \(A'\) also satisfies the universal property 4.2; it is therefore canonically isomorphic (as a \(\lambda\)-ring) to the \(\lambda\)-ring \(A\).

\sgaStableTarget{sga6.e05.def.4-4}\textbf{Definition 4.4.} \uline{The} \(\lambda\)-\uline{ring} \(A\) \uline{will be called the} free \(\lambda\)-ring generated over \(K\) by the family of generators \((X_i)_{i\in I}\).

\sgaStableTarget{sga6.e05.def.4-5}\textbf{Definition 4.5.} \uline{Let} \(A\) \uline{be a} \(\lambda\)-\uline{ring}, \(J\) \uline{an ideal of} \(A\). \uline{We say that} \(J\) \uline{is} \uline{a} \(\lambda\)-ideal of \(A\) \uline{if for every} \(x\in J\), \uline{and every} \(n\ge1\), \(\lambda^n(x)\in J\).

If \(f:A\to B\) is a \(\lambda\)-homomorphism, the kernel of \(f\) is a \(\lambda\)-ideal.

Let \(J\) be a \(\lambda\)-ideal; the composite homomorphism: \(A\xrightarrow{\ \lambda_t\ }\Ghat(A)\longrightarrow\Ghat(A/J)\) vanishes on \(J\) and therefore factors through \(A/J\), which is thus endowed
%%%%%%%%%%  scan idx 347  |  folio 334  %%%%%%%%%%
\SGAPrintPage{334}%
with a \(\lambda\)-ring structure, relations (2.4.2) passing to the quotient.

Every intersection of \(\lambda\)-ideals is a \(\lambda\)-ideal. In particular, there exists a smallest \(\lambda\)-ideal containing a given family of elements \((x_i)_{i\in I}\); by relations (2.4.2), it is the (ordinary) ideal generated by the \(\lambda^n(x_i)\), \((n\ge1)\).

Again let \(K\) be a \(\lambda\)-ring, \((X_i)_{i\in I}\) a family of indeterminates, \(A=K[(X_i),(\lambda^n X_i)]\), and let \((P_{i'})_{i'\in I'}\) be a family of elements of \(A\): let \(J\) be the \(\lambda\)-ideal generated by the \(P_{i'}\).

\sgaStableTarget{sga6.e05.def.4-6}\textbf{Definition 4.6.} \uline{The} \(\lambda\)-\uline{ring} \(A/J\) \uline{will be called the} \(\lambda\)-ring generated over \(K\) by the family of generators \((X_i)_{i\in I}\) subject to the relations \((P_{i'})_{i'\in I'}\).

We can also specify relations between the \(X_i\), \(\gamma^n X_i\) and make the same construction, observing that an ideal is a \(\lambda\)-ideal if and only if it is stable under the operations \(\gamma^n\), \(n\ge1\), as follows from relations 3.2.

Let \(B\) be a \(K\)-\(\lambda\)-algebra, \((b_i)_{i\in I}\) a family of elements of \(B\). Since the \(P_{i'}\) are polynomials in the indeterminates \(X_i\), \(\lambda^n X_i\), we can substitute the \(b_i\), \(\lambda^n(b_i)\) for these indeterminates; if the resulting element is zero, we say that the \(b_i\) satisfy the relations \(P_{i'}\). The \(K\)-\(\lambda\)-algebra \(A/J\) is then characterized by the following universal property:

\sgaStableTarget{sga6.e05.prop.4-7}\textbf{Proposition 4.7.} \uline{For every} \(K\)-\(\lambda\)-\uline{algebra} \(B\) \uline{and every family} \((b_i)_{i\in I}\) \uline{of elements of} \(B\), \uline{satisfying the relations} \((P_{i'})_{i'\in I'}\), \uline{there exists exactly one} \(K\)-\(\lambda\)-\uline{homomorphism} \(f':A/J\to B\) \uline{such that for every} \(i\), we have \(f'(X_i)=b_i\).

Indeed, the kernel of the homomorphism \(f\) defined in 4.2 contains the \(P_{i'}\), hence \(J\) since it is a \(\lambda\)-ideal; \(f\) therefore passes to the quotient and defines
%%%%%%%%%%  scan idx 348  |  folio 335  %%%%%%%%%%
\SGAPrintPage{335}%
\(f'\), which is plainly unique.

\sgaStableTarget{sga6.e05.prop.4-8}\textbf{Proposition 4.8.} \uline{Let} \(n\) \uline{be an integer} \(\ge0\); \uline{in the free} \(\lambda\)-\uline{ring generated} \uline{over} \(K\) \uline{by an indeterminate} \(X\), \uline{the ideal generated by the} \(\lambda^k X\), \(k>n\), \uline{is} \uline{a} \(\lambda\)-\uline{ideal}.

Put \(A=K[X,\lambda^k X]_{k\ge2}\), and let \(J\) be the ideal generated by the \(\lambda^k X\) for \(k>n\); put \(B=A/J\). We must show that the composite homomorphism: \(A\xrightarrow{\ \lambda_t\ }\Ghat(A)\longrightarrow\Ghat(B)\) vanishes on \(J\), i.e. that \(\lambda_t(\lambda^k X)\), \(k>n\), has image \(1\) in \(\Ghat(B)\). Now \(\lambda_t(\lambda^k X)=\lambda^k(\lambda_t(X))\), and \(\Ghat(A)\to\Ghat(B)\) is a \(\lambda\)-homomorphism; it therefore suffices to show that for every ring \(C\) and every series \(\varphi\in\Ghat(C)\) whose coefficients vanish in degrees \(>n\), \(\lambda^k(\varphi)=1\) when \(k>n\). But to prove this last property it suffices to prove it for \(C=K[X_1,\ldots,X_n]\), and \(\varphi=1+X_1t+\ldots+X_nt^n\). Embedding \(K[X_1,\ldots,X_n]\) in \(K[T_1,\ldots,T_n]\) by means of symmetric functions, we can write: \(1+X_1t+\ldots+X_nt^n=\prod_{i=1}^{n}(1+T_it)\), hence
\[
  \lambda_u(1+X_1t+\ldots+X_nt^n)=\prod_{i=1}^{n}\lambda_u(1+T_it)=\prod_{i=1}^{n}\bigl(\underline{1}+\{1+T_it\}u\bigr),
\]
which has degree \(n\) in \(u\).

\sgaStableTarget{sga6.e05.cor.4-9}\textbf{Corollary 4.9.} \uline{The} \(\lambda\)-\uline{ring generated over} \(K\) \uline{by the generator} \(X\) \uline{subject to the relations} \(\lambda^k(X)=0\) \uline{for} \(k>n\), \uline{is isomorphic to} \(K[X,\ldots,\lambda^n X]\).

4.9.1. We would have the same statements on replacing the \(\lambda\) by the \(\gamma\). We can also replace \(X\) by a family of indeterminates \((X_i)_{i\in I}\), subject to the relations \(\lambda^k(X_i)=0\) for \(k>n_i\), where the \(n_i\) are integers \(\ge2\); the resulting \(\lambda\)-ring is then \(K[X_i,\lambda^k X_i]_{i\in I,\ 2\le k\le n_i}\).

\sgaStableTarget{sga6.e05.para.4-10}4.10. Let \(K\) be a binomial ring, and \(A\) an augmented \(K\)-\(\lambda\)-algebra. We denote by \(B=A[X,\gamma^k X]_{k\ge2}\) the free \(\lambda\)-ring generated by \(X\) over \(A\);
%%%%%%%%%%  scan idx 349  |  folio 336  %%%%%%%%%%
\SGAPrintPage{336}%
we define a \(\lambda\)-homomorphism: \(B\to K\) extending the augmentation of \(A\) by sending \(X\) to \(0\), which endows \(B\) with an augmented \(K\)-\(\lambda\)-algebra structure. We propose to determine its filtration (3.10) in terms of that of \(A\).

We grade \(B\) by weight with respect to the \(\gamma^k X\), assigning \(\gamma^k X\) the weight \(k\); we denote by \(M_k\) the set of monomials of weight \(k\) in the indeterminates \(\gamma^i X\).

\sgaStableTarget{sga6.e05.lem.4-11}\textbf{Lemma 4.11.} \uline{Let} \(a\in A\); \uline{the coefficient of a monomial of weight} \(n\) \uline{with respect} \uline{to the} \(\gamma^j X\) \uline{in} \(\gamma^k(aX)=P'_k(\gamma^i(a),\gamma^j X)\) \uline{has filtration at least} \(k-n\) (\uline{in} \(A\))

The polynomial \(P'_k\) has no isobaric components of weight \(<k\) in the \(\gamma^i(a)\), \(\gamma^j X\) (3.5), so the lemma is true when \(a\) has zero augmentation.

If \(\epsilon\) denotes the augmentation homomorphism: \(A\to K\), we can write in the general case:
\[
  \gamma_t(aX)=\gamma_t((a-\epsilon(a))X)\mathbin{.}\gamma_t(\epsilon(a)X).
\]
Now \(\gamma_t(\epsilon(a)X)=(\gamma_t(X))^{\epsilon(a)}=\sum_{n=0}^{\infty}R_{n,\epsilon(a)}((\gamma^j X))\mathbin{.}t^n\), the \(R_{n,\epsilon(a)}\) being isobaric of weight \(n\) in the \(\gamma^j X\) (3.8); we deduce the lemma by taking the product.

\sgaStableTarget{sga6.e05.prop.4-12}\textbf{Proposition 4.12.} \sgaStableTarget{sga6.e05.eq.4-12-1}\(\Fil^k B=\Fil^k A+\Fil^{k-1}A\mathbin{.}M_1+\ldots+\Fil^0 A\mathbin{.}M_k+B\mathbin{.}M_{k+1}\).

It is clear that the second ideal is contained in the first; it therefore remains only to prove the reverse inclusion.

Since the filtration defined by the family of ideals on the right-hand side for varying \(k\) is a ring filtration, it suffices to prove the reverse inclusion for an element of the form \(\gamma^k(b)\), with \(\epsilon(b)=0\), these elements
%%%%%%%%%%  scan idx 350  |  folio 337  %%%%%%%%%%
\SGAPrintPage{337}%
generating the canonical filtration.

Writing \(b=\sum a_i m_i\), with \(a_i\in A\) and the \(m_i\) monomials in the \(\gamma^j X\), and using the additivity of \(\gamma_t\), we see that it suffices to give the proof for \(\gamma^k(am_i)\), with \(a\in A\) and \(m_i\in M_i\).

Using Lemma 4.11. and the \(A\)-\(\lambda\)-homomorphism from \(B\) to \(B\) which sends \(X\) to \(m_i\), we see that it suffices to give the proof for an element of the form \(\gamma^k(m_i)\).

Finally, since the polynomial \(P'_k\) has no isobaric components of weight \(<k\), we are reduced to the case \(\gamma^k(\gamma^j X)\), which itself follows from the fact that the polynomials \(Q'_{j,k}\) have no isobaric components of weight \(<k\) (3.6).

\sgaStableTarget{sga6.e05.rem.4-13}\textbf{Remarks 4.13.} i) The decomposition of \(\Fil^k B\) given in 4.12 is a direct-sum decomposition of \(A\)-modules. In other words, if \(b\in\Fil^k B\), and if we write \(b=\sum_{j\le k}a_{i,j}m_{i,j}+b'\), with \(a_{i,j}\in A\), \(m_{i,j}\in M_j\), and \(b'\in B\mathbin{.}M_{k+1}\), then \(a_{i,j}\in\Fil^{k-j}A\).

ii) Statement 4.12. remains true for every quotient of \(B\) by a \(\lambda\)-ideal contained in the augmentation ideal.

iii) The statement analogous to 4.12., obtained by replacing the indeterminate \(X\) by a family of indeterminates \((X_i)\), is also true.

\sgaStableTarget{sga6.e05.cor.4-14}\textbf{Corollary 4.14.} \sgaStableTarget{sga6.e05.eq.4-14-1}\(\Gr^k B\approx\Gr^k A\oplus\Gr^{k-1}A\mathbin{.}M_1\oplus\ldots\oplus\Gr^0 A\mathbin{.}M_k\).

We define an arrow from right to left by associating to each element \(\bar a\mathbin{.}m\), where \(\bar a\) is the class in \(\Gr^{k-i}A\) of an element \(a\) of \(\Fil^{k-i}A\) and \(m\in M_i\), the class in \(\Gr^k B\) of \(am\). We define an arrow in the opposite direction by decomposition according to Proposition 4.12 and Remark 4.13.i). It is clear that these two arrows are inverse to one another.

%%%%%%%%%%  scan idx 351  |  folio 338  %%%%%%%%%%
\SGAPrintPage{338}%
\sgaStableTarget{sga6.e05.cor.4-15}\textbf{Corollary 4.15.} \uline{Let} \(B'\) \uline{be the} \(\lambda\)-\uline{ring generated over} \(K\) \uline{by the indeterminate} \(X\) \uline{subject to the relations} \(\gamma^k(X)=0\) \uline{for} \(k>n\), \uline{i}.\uline{e}. \(B'=K[X,\ldots,\gamma^n X]\) (4.9). \uline{Then the} \(\lambda\)-\uline{ring filtration of} \(B'\) \uline{is identical to the weight filtration with respect to the} \(\gamma^j X\).

Taking account of Proposition 4.12, Remark 4.13.ii), and the fact that \(A=K\), we have \(\Fil^k B'=K\mathbin{.} M_k+B'\mathbin{.} M_{k+1}\), so \(\Fil^k B'\) is the ideal generated by the monomials of weight \(\ge k\) with respect to the \(\gamma^j X\).

\sgaStableTarget{sga6.e05.sec.5}\section*{5. The \(\lambda^p(N,x)\).}

This section will be used only in Exp.VII~4, for the study of regular immersions; it may therefore be omitted on first reading.

\sgaStableTarget{sga6.e05.para.5-1}5.1. Let \(d\) be an integer \(\ge1\). We denote by \(A\) the \(\mathbb Z\)-\(\lambda\)-algebra generated by the generators \(N\), \(x\) subject to the relations \(\lambda^i(N)=0\) for \(i>d\); we therefore have:
\[
  A=\mathbb Z[N,\ldots,\lambda^dN;
       x,\ldots,\lambda^jx,\ldots].
\]

We put \sgaStableTarget{sga6.e05.eq.5-1-1}\(\lambda_{-1}(N)=\sum_{i=0}^{d}(-1)^i\lambda^i(N)\).

\sgaStableTarget{sga6.e05.lem.5-2}\textbf{Lemma 5.2.} \uline{For every} \(p\ge1\), \(\lambda^p(x\mathbin{.}\lambda_{-1}(N))\) \uline{is divisible by} \(\lambda_{-1}(N)\).

We know (2.3) that \(\lambda^p(x\mathbin{.}\lambda_{-1}(N))=P_p(\lambda^i x,\lambda^j(\lambda_{-1}(N)))\) is an isobaric polynomial of weight \(p\) with respect to the \(\lambda^j(\lambda_{-1}(N))\); it therefore suffices to prove the lemma for the \(\lambda^p(\lambda_{-1}(N))\).

Let \(B\) be the \(\lambda\)-ring generated over \(\mathbb Z\) by the indeterminates \(N_1,\ldots,N_d\), subject to the relations \(\lambda^i(N_j)=0\) if \(i\ge2\). It is clear that \(\lambda^k(\sum_{i=1}^{d}N_i)=0\) for \(k>d\). We can therefore define a \(\lambda\)-homomorphism \(f\) from \(A\) to \(B\) by sending \(N\) to \(\sum_{i=1}^{d}N_i\); we then have \(\Ghat(f)(\lambda_t(N))=\prod_{i=1}^{d}\lambda_t(N_i)=\prod_{i=1}^{d}(1+N_it)\),
%%%%%%%%%%  scan idx 352  |  folio 339  %%%%%%%%%%
\SGAPrintPage{339}%
which shows that, for \(1\le k\le d\), \(\lambda^k N\) maps to the \(k\)-th symmetric function of the \(N_i\), and consequently that \(f\) is injective and identifies \(A\) with the ring of symmetric polynomials in the \(N_i\).

We then have \(\lambda_{-1}(N)=\prod_{i=1}^{d}(1-N_i)\), and \(\lambda_t(\lambda_{-1}(N))=\lambda_t\bigl(\prod_{i=1}^{d}(1-N_i)\bigr)\). The \(\lambda^p(\lambda_{-1}(N))\) are polynomials in the \(N_i\); if, for some index \(i\), we set \(N_i=1\), we see that \(\lambda_t(\lambda_{-1}(N))\) becomes \(\lambda_t(0)=1\), hence that the \(\lambda^p(\lambda_{-1}(N))\) vanish for \(N_i=1\), and therefore that they contain \(1-N_i\) as a factor. Thus they contain the product of the \(1-N_i\) as a factor, i.e. they are divisible by \(\lambda_{-1}(N)\) in \(B\), and hence also in \(A\), since the quotient of two symmetric polynomials is evidently symmetric.

\sgaStableTarget{sga6.e05.def.5-3}5.3. Since \(A\) is an integral domain, we define an element \(\lambda^p(N,x)\) unambiguously by putting (\(p\ge1\)):
\sgaStableTarget{sga6.e05.eq.5-3-1}\begin{equation}\tag{5.3.1}
  \lambda^p(N,x)\mathbin{.}\lambda_{-1}(N)=\lambda^p(x\mathbin{.}\lambda_{-1}(N)).
\end{equation}

The element \(\lambda^p(N,x)\) is therefore a polynomial in the \(\lambda^i(x)\), \(\lambda^j(N)\), isobaric of weight \(p\) with respect to the \(\lambda^i(x)\), since the same is true of \(\lambda^p(x\mathbin{.}\lambda_{-1}(N))\) by 2.3.

Let \(d'\) be an integer \(\ge d\), and let \(A'\) be the \(\mathbb Z\)-\(\lambda\)-algebra generated by the generators \(N'\) and \(x'\) subject to the relations \(\lambda^i(N')=0\) for \(i>d'\); by Proposition 4.7, there exists exactly one \(\lambda\)-homomorphism \(f:A'\longrightarrow A\), sending \(N'\) to \(N\) and \(x'\) to \(x\). Definition (5.3.1) implies that \(f(\lambda^p(N',x'))=\lambda^p(N,x)\) for every \(p\ge1\).

If \(B\) is a \(\lambda\)-ring and \(N''\) is an element of \(B\) such that \(\lambda^k(N'')=0\) for \(k>d\), we can, for every \(x''\in B\), define an element \(\lambda^p(N'',x'')\) by substitution, or by using the unique \(\lambda\)-homomorphism from \(A\) to
%%%%%%%%%%  scan idx 353  |  folio 340  %%%%%%%%%%
\SGAPrintPage{340}%
\(B\) which sends \(N\) to \(N''\), and \(x\) to \(x''\); by the preceding remark, the element obtained does not depend on the integer \(d\) such that \(\lambda^k(N'')=0\) for \(k>d\). Using the same homomorphism, we see that relation (5.3.1) is satisfied by \(N''\) and \(x''\).

\sgaStableTarget{sga6.e05.para.5-4}5.4. Consider the \(\lambda\)-ring generated over \(\mathbb Z\) by the indeterminates \(N,x,y\), subject to the relations \(\lambda^k(N)=0\) for \(k>d\). From (5.3.1) we obtain:
\begin{align*}
  \lambda^p(N,x+y)\mathbin{.}\lambda_{-1}(N)
  &=\lambda^p\bigl((x+y)\mathbin{.}\lambda_{-1}(N)\bigr)\\
  &=\sum_{i+j=p}\lambda^i(x\mathbin{.}\lambda_{-1}(N))\mathbin{.}\lambda^j(y\mathbin{.}\lambda_{-1}(N))\\
  &=\sum_{\substack{i+j=p\\ i,j\ne0}}\lambda^i(N,x)\mathbin{.}\lambda^j(N,y)\mathbin{.}(\lambda_{-1}(N))^2
    +(\lambda^p(N,x)+\lambda^p(N,y))\lambda_{-1}N
\end{align*}

Since the ring under consideration is an integral domain, we deduce the following relation, valid by substitution in every \(\lambda\)-ring, and for every element \(N\) such that \(\lambda^k(N)=0\) for \(k>d\):
\sgaStableTarget{sga6.e05.eq.5-4-1}\begin{equation}\tag{5.4.1}
  \lambda^p(N,x+y)=\sum_{\substack{i+j=p\\ i,j\ne0}}\lambda^i(N,x)\mathbin{.}\lambda^j(N,y)\mathbin{.}\lambda_{-1}(N)+\lambda^p(N,x)+\lambda^p(N,y).
\end{equation}

Now consider the \(\lambda\)-ring generated over \(\mathbb Z\) by indeterminates \(N,N',x\) subject to the relations \(\lambda^k(N)=0\) for \(k>d\), and \(\lambda^k(N')=0\) for \(k>d'\). We then have \(\lambda^k(N+N')=0\) for \(k>d+d'\); we can therefore define the \(\lambda^p(N+N',x)\). We obtain:
\[
  \lambda^p(N+N',x)\mathbin{.}\lambda_{-1}(N+N')=\lambda^p(x\mathbin{.}\lambda_{-1}(N+N')).
\]
By the additivity of \(\lambda_t\), we have: \(\lambda_{-1}(N+N')=\lambda_{-1}(N)\mathbin{.}\lambda_{-1}(N')\). We deduce:
\begin{align*}
  \lambda^p(N+N',x)\mathbin{.}\lambda_{-1}(N)\mathbin{.}\lambda_{-1}(N')
  &=\lambda^p(x\mathbin{.}\lambda_{-1}(N)\mathbin{.}\lambda_{-1}(N'))\\
  &=\lambda^p(N',x\mathbin{.}\lambda_{-1}(N))\mathbin{.}\lambda_{-1}(N').
\end{align*}

%%%%%%%%%%  scan idx 354  |  folio 341  %%%%%%%%%%
\SGAPrintPage{341}%
Since the ring under consideration is an integral domain, we deduce the following relation, valid by substitution in every \(\lambda\)-ring, and for all elements \(N,N'\) such that \(\lambda^k(N)=0\) for \(k>d\), \(\lambda^k(N')=0\) for \(k>d'\);
\sgaStableTarget{sga6.e05.eq.5-4-2}\begin{equation}\tag{5.4.2}
  \lambda^p(N+N',x)\mathbin{.}\lambda_{-1}(N)=\lambda^p(N',x\mathbin{.}\lambda_{-1}(N)).
\end{equation}

\sgaStableTarget{sga6.e05.para.5-5}5.5. Now let \(A\) be any \(\lambda\)-ring, and let \(N\) be an element of \(A\) such that \(\lambda^k(N)=0\) for \(k>d\). We shall associate to \(A\) a new \(\lambda\)-ring \(A_N\).

We first define a new multiplication on \(A\) by putting, for \(x,y\in A\):
\begin{equation}\tag{5.5.1}
  x\mathbin{.}_N y=x\mathbin{.} y\mathbin{.}\lambda_{-1}(N).
\end{equation}

The ring \(A\) with this multiplication is no longer unital; we therefore consider the ring \(\mathbb Z\times A\), obtained by adjoining a unit to \(A\). Its operations are thus defined by:
\begin{equation}\tag{5.5.2}
\begin{aligned}
  (n,x)+(n',x')&=(n+n',x+x')\\
  (n,x)\mathbin{.}(n',x')&=(nn',nx'+n'x+x\mathbin{.}_N x')
\end{aligned}
\end{equation}
We regard \(A\) as embedded in this ring by \(x\rightsquigarrow(0,x)\). We then define a pre-\(\lambda\)-ring structure by:
\sgaStableTarget{sga6.e05.eq.5-5-1}\begin{equation}\tag{5.5.3}
  \lambda_t(n,x)=(1+t)^n\mathbin{.}\sum_{p=0}^{\infty}\lambda^p(N,x)\mathbin{.} t^p.
\end{equation}
This is a pre-\(\lambda\)-ring structure by (5.4.1). We shall denote by \(A_N\) the pre-\(\lambda\)-ring thus obtained. To see that it is a \(\lambda\)-ring, it suffices by substitution to give the proof in the case where \(A\) is the \(\lambda\)-ring generated over \(\mathbb Z\) by the generators \(N,X,Y\), subject to the relations \(\lambda^k(N)=0\) for \(k>d\). We then define a homomorphism \(A_N\to A\) by
%%%%%%%%%%  scan idx 355  |  folio 342  %%%%%%%%%%
\SGAPrintPage{342}%
sending \((1,0)\) to \(1\), and \((0,x)\) to \(x\mathbin{.}\lambda_{-1}(N)\). It follows immediately from (5.5.1) and (5.3.1) that this is a \(\lambda\)-homomorphism, and, since \(A\) is an integral domain in the case considered, one immediately verifies that it is injective; since \(A\) is a \(\lambda\)-ring, we deduce the result.

In the \(\lambda\)-ring \(A_N\), we can define operations \(\gamma\) (3.2). For \(x\in A\), the value of \(\gamma^p\) calculated on the element \((0,x)\) of \(A_N\) is given by:
\[
  \gamma^p(0,x)=\sum_{i=0}^{p-1}\binom{p-1}{i}\lambda^{p-i}(N,x),
\]
it is therefore an element of \(A\); it is defined by the relation:
\sgaStableTarget{sga6.e05.eq.5-5-2}\begin{equation}\tag{5.5.4}
  \gamma^p(N,x)\mathbin{.}\lambda_{-1}(N)=\gamma^p(x\mathbin{.}\lambda_{-1}(N)),
\end{equation}
and the universal polynomials resulting from it.

\sgaStableTarget{sga6.e05.para.5-6}5.6. Let \(A\) be the \(\lambda\)-ring generated over \(\mathbb Z\) by the generators \(N,L,x\), subject to the relations \(\lambda^k(N)=0\) for \(k>d\), \(\lambda^k(L)=0\) for \(k>1\). We propose to calculate the \(\lambda^p(N+L,x)\) in terms of the \(\lambda^q(N,x)\).

By definition:
\[
  \lambda^p(N+L,x)\mathbin{.}\lambda_{-1}(N)\mathbin{.}\lambda_{-1}(L)
  =\lambda^p(x\mathbin{.}\lambda_{-1}(N)\mathbin{.}\lambda_{-1}(L)),
\]
that is:
\[
\begin{aligned}
1&+\lambda_{-1}(N)\mathbin{.}\lambda_{-1}(L)\mathbin{.}\sum_{k\ge1}\lambda^k(N+L,x)\mathbin{.} t^k \\
 &=\lambda_t(x\mathbin{.}\lambda_{-1}(N))\circ \lambda_t(1-L) \\
 &=\frac{\lambda_t(x\mathbin{.}\lambda_{-1}(N))}{\lambda_{Lt}(x\mathbin{.}\lambda_{-1}(N))}
   \qquad\left(\text{since }\lambda_t(1-L)=\frac{\lambda_t(1)}{\lambda_t(L)}=\frac{1+t}{1+Lt}\text{, and we apply }(2.3.2)\right) \\
 &=\frac{1+\ldots+\lambda^k(x\mathbin{.}\lambda_{-1}(N))t^k+\ldots}{1+\ldots+\lambda^k(x\mathbin{.}\lambda_{-1}(N))L^k t^k+\ldots} \\
 &=1+\frac{\sum_{k\ge1}\lambda^k(N,x)\mathbin{.}\lambda_{-1}(N)\mathbin{.}(1-L^k)t^k}{1+\sum_{k\ge1}\lambda^k(N,x)\mathbin{.}\lambda_{-1}(N)\mathbin{.} L^k t^k}
\end{aligned}
\]

%%%%%%%%%%  scan idx 356  |  folio 343  %%%%%%%%%%
\SGAPrintPage{343}%
Thus \(\lambda^p(N+L,x)\) is the coefficient of \(t^p\) in the series:
\sgaStableTarget{sga6.e05.eq.5-6-1}\[
  \frac{\displaystyle\sum_{k\ge1}\lambda^k(N,x)
        (1+\ldots+L^{k-1})t^k}
       {\displaystyle 1+\sum_{k\ge1}\lambda^k(N,x)
        \lambda_{-1}(N)L^k t^k}.
\]

\sgaStableTarget{sga6.e05.prop.5-7}\textbf{Proposition 5.7.} \uline{Let} \(K\) \uline{be a binomial ring}, \(A\) \uline{an augmented} \(K\)-\(\lambda\)-\uline{algebra}, \uline{and} \(N\) \uline{an element of} \(A\) \uline{such that} \(\lambda^k(N)=0\) \uline{for} \(k>d\). \uline{Then for every} \(x\in A\) \uline{and every} \(p\ge1\), \uline{we have}:
\sgaStableTarget{sga6.e05.eq.5-7-1}\begin{equation}\tag{5.7.1}
  \gamma^p(N,x)\in\Fil^{p-d}A.
\end{equation}

Let \(A_0\) be the \(K\)-\(\lambda\)-algebra generated by generators \(N'_0,X'_0\), subject to the relations \(\gamma^k(N'_0)=0\) for \(k>d\). We define the augmentation homomorphism \(A_0\xrightarrow{\epsilon}K\) by \(\epsilon(N'_0)=\epsilon(X'_0)=0\). We know that \(A_0\) is the ring of polynomials with coefficients in \(K\) in the indeterminates \(\gamma^i(N'_0)\), with \(1\le i\le d\), and \(\gamma^j(X'_0)\), with \(j\ge1\), and that the \(\lambda\)-filtration on \(A_0\) is identical to the weight filtration, assigning the weight \(i\) to \(\gamma^i(N'_0)\), \(\gamma^i(X'_0)\) (4.12 and 4.15).

Put \(N_0=N'_0+d\), \(X_0=X'_0+\epsilon(x)\). We shall first prove 5.7 for \(N_0\) and \(X_0\). By the preceding discussion, (5.7.1) is equivalent to:
\[
  \gamma^p(N_0,X_0)\mathbin{.}\gamma^d(N'_0)\in\Fil^pA_0.
\]
Now \(\gamma^d(N'_0)=(-1)^d\lambda_{-1}(N_0)\); this relation can therefore be written:
\[
  (-1)^d\gamma^p(X_0\gamma^d(N_0))\in\Fil^pA_0
\]
by (5.5.4). This relation does hold, since \(\gamma^d(N_0)\) has zero augmentation.

By 4.7, there exists a \(K\)-\(\lambda\)-homomorphism \(A_0\longrightarrow A\) compatible with the augmentations, sending \(N'_0\) to \(N-d\), and \(X'_0\) to \(x-\epsilon(x)\), hence \(N_0\) to \(N\) and \(X_0\)
%%%%%%%%%%  scan idx 357  |  folio 344  %%%%%%%%%%
\SGAPrintPage{344}%
to \(x\), whence the result, since the homomorphism is compatible with the filtrations.

\sgaStableTarget{sga6.e05.sec.6}\section*{6. Chern ring}

In this section, the base ring will be a binomial ring (2.7) \(K\), fixed once and for all. We shall use the category \(\mathcal C\) (1.1), and therefore return to the notation of Section 1, reserving the notation \(\widehat G\) for the functor defined on \(\mathcal C\) (1.1), and denoting by \(\widehat G\circ i\) its restriction to \(\mathcal C_0\) (1.8).

\sgaStableTarget{sga6.e05.para.6-1}6.1. In 3.5 and 3.6 we defined morphisms of functors:
\[
  (\widehat G\circ i)\times(\widehat G\circ i)\xrightarrow{\nabla}\widehat G\circ i,
  \qquad
  \widehat G^t\circ i\xrightarrow{\gamma'}(\widehat G^u\circ i)\circ(\widehat G^t\circ i)_\nabla,
  \qquad
  \widehat G\circ i\xrightarrow{a}\widehat G\circ i,
\]
the last being multiplication under the law \(\nabla\) by \(\gamma_t(a)\), with \(a\in K\).

By 1.9, we deduce morphisms:
\[
  \widehat G\times\widehat G\xrightarrow{\nabla}\widehat G_g,
  \qquad
  \widehat G^t\xrightarrow{\gamma'}\widehat G^u\circ i\circ\widehat G^t_{g\nabla},
  \qquad
  \widehat G\xrightarrow{a}\widehat G_g;
\]
and it follows from the isobaricity properties of the polynomials \(P'_n\), \(Q'_{i,j}\), \(R_{n,a}\) (3.5, 3.6, 2.9) that the image of the first and third functors is in \(\widehat G\), and that of the second in the subgroup of formal series in \(u\) whose coefficients in nonzero degree are in \(\widehat G^t\).

Composing the first and third with the canonical morphism: \(\widehat G\longrightarrow 1+\widehat{\operatorname{Gr}\,\mathord{.}\,}^{+}\) (1.11), we obtain, by 1.12, two morphisms of functors on \(\mathcal G\):
%%%%%%%%%%  scan idx 358  |  folio 345  %%%%%%%%%%
\SGAPrintPage{345}%
\[
  (1+\widehat{\mathord{.}}{}^{+})\times(1+\widehat{\mathord{.}}{}^{+})\xrightarrow{\ \ast\ }1+\widehat{\mathord{.}}{}^{+},
\]
\[
  1+\widehat{\mathord{.}}{}^{+}\xrightarrow{\ a\ }1+\widehat{\mathord{.}}{}^{+},
\]
defining on the functor \(1+\widehat{\mathord{.}}{}^{+}\) a multiplication and operations of \(K\). By 1.13, the multiplication is defined by its value on the element \((1+X,1+Y)\) of \((1+\widehat{K[X,Y]}{}^{+})^{2}\); it will be denoted by \(\ast\); we therefore have:
\begin{equation}\tag{6.1.1}
  (1+X)\ast(1+Y)=\frac{1+X+Y}{(1+X)(1+Y)}.
\end{equation}
The \(n\)-th universal polynomial of this law is the isobaric component of weight \(n\) of \(P'_n\) (1.15), and will be denoted by \(P''_n\).

The operations of \(K\) will again be denoted by exponentiation (3.8). The universal polynomials defining these operations are the isobaric components of lowest degree of the \(R_{n,a}\) (1.15); since these are isobaric, they are the \(R_{n,a}\) themselves.

Let \(A\in\operatorname{Ob}(\mathcal G)\); the ring \(1+\widehat A^{+}\) defined by the preceding laws is in general without a unit; we therefore extend this ring by adjoining the scalars of \(K\), introducing the new ring \(K\times(1+\widehat A^{+})\), endowed with the composition laws defined by: \(\forall\,a,a'\in K\), \(\varphi,\varphi'\in1+\widehat A^{+}\),
\[
  (a,\varphi)+(a',\varphi')=(a+a',\varphi\varphi')~;
\]
\begin{equation}\tag{6.1.2}
  (a,\varphi).(a',\varphi')=(aa',\varphi^{a'}\mathbin{.}{\varphi'}^{\,a}\mathbin{.}(\varphi\ast\varphi')).
\end{equation}

The ring thus constructed admits (1,1) as its unit, and \(1+\widehat A^{+}\) is identified with the ideal \((0)\times(1+\widehat A^{+})\). It is naturally endowed with a \(K\)-algebra structure such that:
\begin{equation}\tag{6.1.3}
  a.(a',\varphi)=(a,1).(a',\varphi)=(aa',\varphi^{a}).
\end{equation}

%%%%%%%%%%  scan idx 359  |  folio 346  %%%%%%%%%%
\SGAPrintPage{346}%
We endow it with a \(K\)-\(\lambda\)-algebra structure as follows: we observed above that the image of the functor \(\gamma'\colon\widehat G^t\longrightarrow\widehat G^u\circ i\circ\widehat G^t_g\) is in the subgroup of formal series in \(u\) whose coefficients in nonzero degree are in \(\widehat G^t\). Using the canonical morphism: \(\widehat G\longrightarrow 1+\widehat{\operatorname{Gr}\,\mathord{.}\,}^{+}\), which gives: \(\widehat G\longrightarrow K\times(1+\widehat{\operatorname{Gr}\,\mathord{.}\,}^{+})\), we finally deduce a morphism:
\[
  \widehat G^t\xrightarrow{\gamma'}\widehat G^u\circ i\circ(K\times(1+\widehat{\operatorname{Gr}\,\mathord{.}\,}^{+})),
\]
which, by 1.12, defines a morphism of functors on \(\mathcal G\):
\[
  1+\widehat{\mathord{.}}{}^{+}\xrightarrow{\gamma'}\widehat G^u\circ i\circ(K\times(1+\widehat{\mathord{.}}{}^{+})).
\]
The universal polynomials defining this morphism are the coefficients of the \({\gamma'}^{i}(1+X_1+\ldots+X_n+\ldots)\); the \(j\)-th coefficient is the isobaric component of weight \(j\) of \(Q'_{i,j}\) (a proof analogous to that of 1.15), and will be denoted by \(Q''_{i,j}\).

Finally, we extend the morphism thus constructed to \(K\times(1+\widehat{\mathord{.}}{}^{+})\) naturally, using the \(\lambda\)-structure of \(K\). We therefore obtain a pre-\(\lambda\)-ring structure on \(K\times(1+\widehat{\mathord{.}}{}^{+})\), defined by the corresponding operations \(\gamma\), which satisfy:
\begin{equation}\tag{6.1.4}
  \gamma_t(a,\varphi)=\gamma_t(a).\Bigl(1+\sum_{p\ge1}{\gamma'}^{p}(\varphi)t^{p}\Bigr),
\end{equation}
and which will be denoted by \(\gamma\) instead of \(\gamma'\), since no confusion is possible.

To verify that this is a \(\lambda\)-ring structure, we verify the commutativity of (3.7.2); it is known for the elements of \(K\); for the elements of \(1+\widehat{\mathord{.}}{}^{+}\), it follows from that of (3.7.2) for \(\widehat G\) by passage to the quotient. The result follows at once.

\providecommand{\Ghat}{\widehat G}
\providecommand{\Aqhat}{\widehat A_{\mathbb Q}}
\providecommand{\Ch}{\operatorname{Ch}}
\providecommand{\Log}{\operatorname{Log}}
\providecommand{\id}{\operatorname{id}}
\providecommand{\gr}{\operatorname{gr}}
\providecommand{\Hom}{\operatorname{Hom}}
\providecommand{\epsi}{\varepsilon}
\providecommand{\clgr}{\operatorname{cl}}
\providecommand{\Fil}{\operatorname{Fil}}
\providecommand{\Gr}{\operatorname{Gr}}
\providecommand{\Sym}{\operatorname{Sym}}

% Local macros for Exposé V, §§6--7
\providecommand{\Ghat}{\widehat G}
\providecommand{\Aqhat}{\widehat A_{\mathbb Q}}
\providecommand{\Ch}{\operatorname{Ch}}
\providecommand{\Log}{\operatorname{Log}}
\renewcommand{\id}{\operatorname{id}}
\providecommand{\gr}{\operatorname{gr}}
\providecommand{\Hom}{\operatorname{Hom}}
\providecommand{\epsi}{\varepsilon}
\providecommand{\clgr}{\operatorname{cl}}
\providecommand{\Fil}{\operatorname{Fil}}
\providecommand{\Gr}{\operatorname{Gr}}
\providecommand{\Sym}{\operatorname{Sym}}

%%%%%%%%%%  scan idx 360  |  folio 347  %%%%%%%%%%
\SGAPrintPage{347}%
\sgaStableTarget{sga6.e05.expose.occ2}\sgaStableTarget{sga6.e05.def.6-2}\textbf{Definition 6.2.} \uline{The functor defined in} 6.1 \uline{on the category of} \uline{graded} \(K\)-\uline{algebras}, \uline{with values in the category of} \(K\)-\(\lambda\)-\uline{algebras}, \uline{is called the} \uline{Chern ring}; \uline{its value on} \(A\in\operatorname{Ob}(\mathcal G)\) \uline{is denoted by} \(\Ch(A)\) (\uline{or}, \uline{if} \uline{there is a risk of confusion}, \(\Ch_K(A)\)).

From formulas (6.1.1), (6.1.2), (6.1.4), and (3.6.2), we deduce the relations:
\sgaStableTarget{sga6.e05.eq.6-2-1}\begin{equation}\tag{6.2.1}
\begin{gathered}
 (1,1+X).(1,1+Y)=(1,1+X+Y)\\
 \lambda^i(1,1+X)=0\quad\text{for }i\ge2,
\end{gathered}
\end{equation}
relations which served to define the Chern ring in [1]. Conversely, from these formulas one can recover relations (6.1.1) and (6.1.4).

Let us indicate why the Chern ring is introduced. We place ourselves in the situation of Example 3.9, and denote by \(A\) the associated graded ring (for the filtration defined in 3.10) of the \(K^\bullet\) of a commutatively locally ringed topos (which we suppose to be a \(\lambda\)-ring for simplicity); \(K\) is the augmentation ring. We shall see below that, for every element of an augmented \(\lambda\)-ring, one can define Chern classes with values in the associated graded ring (6.7). To every locally free sheaf of finite type, we can therefore associate an element of \(\Ch(A)\) whose component on the factor \(K\) is the rank of the sheaf, and whose component on the factor \(1+\widehat A^+\) is the sum of the Chern classes. Formulas (6.2.1) are then the usual formulas giving the Chern classes and the rank of the tensor product of two invertible sheaves and of the exterior powers of an invertible sheaf (that is, the map just defined from \(K^\bullet\) to \(\Ch(A)\) is a \(\lambda\)-homomorphism).

%%%%%%%%%%  scan idx 361  |  folio 348  %%%%%%%%%%
\SGAPrintPage{348}%
\sgaStableTarget{sga6.e05.para.6-3}6.3.
For every graded \(K\)-algebra \(A\), we denote by \(A_{\mathbb Q}\) the \(K\)-algebra \(A\otimes_{\mathbb Z}\mathbb Q\), graded in the evident way. By 1.13, we define a functorial homomorphism: \(1+\widehat A^+\xrightarrow{\ \ch\ }\widehat A_{\mathbb Q}^{+}\) by specifying its value on \(1+X\), namely
\sgaStableTarget{sga6.e05.eq.6-3-1}\begin{equation}\tag{6.3.1}
  \ch(1+X)=\exp(X)-1.
\end{equation}

This homomorphism is compatible with the multiplication defined in 6.1 on \(1+\widehat A^+\): by 1.13, it suffices to show this for the elements \(1+X,1+Y\) of \(1+\widehat{K[X,Y]}{}^{+}\); now we have:
\[
\begin{aligned}
 \ch\bigl((1+X)*(1+Y)\bigr)
 &= \ch\!\left(\frac{1+X+Y}{(1+X)(1+Y)}\right)\\
 &= \ch(1+X+Y)-\ch(1+X)-\ch(1+Y)\\
 &= \exp(X+Y)-\exp(X)-\exp(Y)+1\\
 &= \ch(1+X).\ch(1+Y).
\end{aligned}
\]

Let \(\eta\) denote the endomorphism (functorial in \(A\)) of \(\widehat A_{\mathbb Q}^{+}\) defined by:
\[
 \eta\!\left(\sum_{k\ge1} a_k t^k\right)
 =
 \sum_{k\ge1}\frac{(-1)^{k-1}}{(k-1)!}.a_k t^k.
\]
Then, \(\forall\varphi\in1+\widehat A^+\),
\sgaStableTarget{sga6.e05.eq.6-3-2}\begin{equation}\tag{6.3.2}
  \ch(\varphi)=\eta(\Log\varphi).
\end{equation}
Indeed, \(\ch\) and \(\eta\circ\Log\) are two functorial homomorphisms of abelian groups, and for them to be equal it suffices that they agree on \(1+X\), which is evident.

We deduce that \(\ch\) commutes with the operations of \(K\): indeed, here again it suffices to show this for \(1+X\), and we have:
\[
\begin{aligned}
 \forall a\in K,\quad\ch((1+X)^a)
 &=\ch(\exp(a\Log(1+X)))\\
 &=\eta(a\Log(1+X))\\
 &=a.\eta(\Log(1+X))\\
 &=a.\ch(1+X).
\end{aligned}
\]

%%%%%%%%%%  scan idx 362  |  folio 349  %%%%%%%%%%
\SGAPrintPage{349}%
We can therefore extend \(\ch\) in the evident way to a homomorphism of \(K\)-algebras: \(K\times(1+\widehat A^+)\longrightarrow K\oplus\widehat A_{\mathbb Q}^{+}\). We therefore have:
\sgaStableTarget{sga6.e05.eq.6-3-3}\begin{equation}\tag{6.3.3}
  \ch(a,\varphi)=a+\eta(\Log\varphi).
\end{equation}

\sgaStableTarget{sga6.e05.def.6-4}\textbf{Definition 6.4.} \uline{The homomorphism} \(\ch:\Ch A\longrightarrow K\oplus\widehat A_{\mathbb Q}^{+}\) \uline{just} \uline{defined is called the Chern character}.

If \(A\) is of characteristic \(0\) (i.e. is a \(\mathbb Q\)-algebra), the Chern character is an isomorphism: indeed, let \(\psi\in K\oplus\widehat A^+\); we can write \(\psi=a+\psi_1\), where \(\psi_1\) has no term of degree zero; we then define an inverse morphism of \(\ch\) by \(\psi\rightsquigarrow(a,\exp(\eta^{-1}(\psi_1)))\).

\sgaStableTarget{sga6.e05.prop.6-5}\textbf{Proposition 6.5.} \uline{Let} \(A\) \uline{be a} \uline{graded} \(K\)-\uline{algebra}; \uline{on} \(\Ch(A)\) \uline{consider the filtration defined as follows} (\uline{degree filtration}):
\sgaStableTarget{sga6.e05.eq.6-5-1}\begin{equation}\tag{6.5.1}
  (\Ch(A))_0=\Ch(A)~;\qquad
  \forall n\ge1,\ (\Ch(A))_n=0\times 1+\widehat A^{\ge n},
\end{equation}
\uline{where} \(\widehat A^{\ge n}=\prod_{i\ge n}A^i\). \uline{We have the following congruences}:

\begin{enumerate}[label=\roman*)]
\item \uline{Let} \(n\ge1\), \uline{and} \(\varphi\in(\Ch A)_n\); \uline{put} \(\varphi=1+a_n+\ldots\); \uline{then}:
\sgaStableTarget{sga6.e05.eq.6-5-2}\begin{equation}\tag{6.5.2}
  \ch(\varphi)=\frac{(-1)^{n-1}}{(n-1)!}.a_n+\ldots~;
\end{equation}

\item \uline{Let} \(\varphi\in(\Ch A)_m\), \(\psi\in(\Ch A)_n\); \uline{put}
\(\varphi=1+a_m+\ldots\), \(\psi=1+b_n+\ldots\); \uline{then}:
\sgaStableTarget{sga6.e05.eq.6-5-3}\begin{equation}\tag{6.5.3}
  \varphi*\psi
  =
  1-\frac{(m+n-1)!}{(m-1)!(n-1)!}.a_m b_n+\ldots,
\end{equation}
\uline{and consequently}, \uline{the degree filtration is a ring filtration}.
\end{enumerate}

i) The Chern character is by construction a morphism of filtered functors; from the congruence \(\varphi\equiv1+a_n\ \mathrm{mod.}\ (\Ch A)_{n+1}\), we deduce
%%%%%%%%%%  scan idx 363  |  folio 350  %%%%%%%%%%
\SGAPrintPage{350}%
therefore \(\ch(\varphi)\equiv\ch(1+a_n)\ \mathrm{mod.}\ \widehat A_{\mathbb Q}^{\ge n+1}\); now \(\ch(1+a_n)=\eta(\Log(1+a_n))\)
\[
  =\frac{(-1)^{n-1}}{(n-1)!}\,a_n+\ldots
\]

ii) By substitution, it suffices to prove formula (6.5.3) for \(A=K[X_i;Y_j]_{i,j\in\mathbb N}\) and for \(\varphi=1+\sum_{k\ge m}X_k\), \(\psi=1+\sum_{k\ge n}Y_k\). Since \(K\) is torsion-free, we can embed this ring in its tensor product with \(\mathbb Q\), and thus see that \(\ch\) is injective. Put \(\varphi*\psi=1+Z_p+\ldots\), where \(Z_p\) is the first nonzero term, \(\neq1\), of \(\varphi*\psi\). The multiplicativity of \(\ch\) gives, by (6.5.2):
\[
  \frac{(-1)^{p-1}}{(p-1)!}.Z_p+\ldots
  =
  \left(\frac{(-1)^{m-1}}{(m-1)!}.X_m+\ldots\right)
  \left(\frac{(-1)^{n-1}}{(n-1)!}.Y_n+\ldots\right),
\]
whence the value of \(Z_p\) by identifying the terms of lowest degree.

\sgaStableTarget{sga6.e05.para.6-6}6.6. The \(\lambda\)-ring \(\Ch A\) is augmented over \(K\): indeed, we define a \(\lambda\)-homomorphism from \(\Ch A\) to \(K\) by \((a,\varphi)\rightsquigarrow a\). It is therefore endowed with a canonical filtration of an augmented \(\lambda\)-ring. We shall compare it with the one just defined.

\sgaStableTarget{sga6.e05.prop.6-6-1}\textbf{Proposition 6.6.1.} \uline{Let} \(x\in\Ch A\) \uline{be an element of the form} \(x=(n,1+\ldots+a_n)\). \uline{Then}:
\[
\begin{array}{ll}
\text{i)}& \lambda^i(x)=0\quad\text{\uline{for} }i>n~;\\[0.3em]
\text{ii)}& \lambda^n(x)=(1,1+a_1)\\[0.3em]
\text{iii)}& \displaystyle\sum_{i=0}^n(-1)^i\lambda^i(x)
       =(0,1-(n-1)!a_n+\ldots),
\end{array}
\]
\uline{the unwritten terms} \uline{being of degree} \(>n\).

By functoriality, it suffices to give the proof in the case
%%%%%%%%%%  scan idx 364  |  folio 351  %%%%%%%%%%
\SGAPrintPage{351}%
where \(A\) is the ring \(K[X_1,\ldots,X_n]\), filtered by weight with respect to the \(X_i\), \(X_i\) having weight \(i\), and for the element \((n,1+X_1+\ldots+X_n)\).

We embed \(A\) in \(A'=K[T_1,\ldots,T_n]\) by the symmetric functions, and thus obtain an injective homomorphism: \(\Ch A\longrightarrow\Ch A'\); we then have:
\[
 (n,1+X_1+\ldots+X_n)
 =
 \left(n,\prod_{i=1}^n(1+T_i)\right)
 =
 \sum_{i=1}^n(1,1+T_i)
\]
We deduce, by (6.2.1):
\[
  \lambda_t(x)=\prod_{i=1}^n\bigl(1+(1,1+T_i)t\bigr),
\]
and obtain:
\[
  \lambda^k(x)=0\ \text{for }k>n,\quad\text{and}\quad
  \lambda^n(x)=\prod_{i=1}^n(1,1+T_i)
   =\left(1,1+\sum_{i=1}^nT_i\right)=(1,1+X_1)
\]
Finally,
\[
  \sum_{i=0}^n(-1)^i\lambda^i(x)=\lambda_{-1}(x)
  =
  \prod_{i=1}^n\bigl(1-(1,1+T_i)\bigr)
  =
  (-1)^n.\prod_{i=1}^n(0,1+T_i).
\]
Now, by (6.5.3), \(\prod_{i=1}^n(0,1+T_i)=\left(0,1+(-1)^{n-1}(n-1)!\prod_{i=1}^nT_i+\ldots\right)\)
that is: \(\lambda_{-1}(x)=(0,1-(n-1)!X_n+\ldots)\).

\sgaStableTarget{sga6.e05.cor.6-6-2}\textbf{Corollary 6.6.2.} \uline{Relation} 6.6.1.iii) \uline{holds for every element} \(X\) \uline{of augmentation} \(n\).

To prove this, we shall use the lemma:

\sgaStableTarget{sga6.e05.lem.6-6-3}\textbf{Lemma 6.6.3.} \uline{For every} \(n\), \((\Ch A)_n\) \uline{is a} \(\lambda\)-\uline{ideal of} \(\Ch A\).

We may suppose that \(n\ge1\). Let \(x=(0,\varphi)\in(\Ch A)_n\); we can write \(\gamma^i(0,\varphi)=(0,1+Q''_{i,1}+\ldots+Q''_{i,j}+\ldots)\) where the \(Q''_{i,j}\) are the isobaric components of weight \(j\) of the \(Q'_{i,j}\) (3.6). For every \(j\), the polynomials \(Q''_{i,j}\) are therefore polynomials in the first \(j\) coefficients of \(\varphi\). If \(j<n\), they are therefore zero, and consequently the \(\gamma^i(x)\) belong to \((\Ch A)_n\), and
%%%%%%%%%%  scan idx 365  |  folio 352  %%%%%%%%%%
\SGAPrintPage{352}%
\((\Ch A)_n\) is a \(\lambda\)-ideal.

Now let \(x=(n,1+a_1+\ldots+a_n+\ldots)\) be an element of augmentation \(n\) of \(\Ch A\); \(x\) is congruent \(\mathrm{mod.}\ (\Ch A)_{n+1}\) to \((n,1+\ldots+a_n)\), and consequently \(\lambda_{-1}(x)\) is congruent to \(\lambda_{-1}(n,1+\ldots+a_n)=(0,1-(n-1)!a_n+\ldots)\), whence:
\[
  \lambda_{-1}(x)=(0,1-(n-1)!a_n+\ldots).
\]

\sgaStableTarget{sga6.e05.cor.6-6-4}\textbf{Corollary 6.6.4.} \uline{Let} \(\epsilon\) \uline{be the augmentation homomorphism of} \(\Ch A\); \uline{then}, \uline{for every} \(x\in\Ch A\), \uline{we have}:
\[
  \gamma^n(x-\epsilon(x))
  =
  (0,1+(-1)^{n-1}(n-1)!a_n+\ldots)
\]

By definition, \(\gamma^n(x-\epsilon(x))=(-1)^n\sum_{i=0}^n(-1)^i\lambda^i(x+n-\epsilon(x))\). Since \(x+n-\epsilon(x)\) has augmentation \(n\), we can apply Corollary 6.6.2 to it, whence the result.

\sgaStableTarget{sga6.e05.cor.6-6-5}\textbf{Corollary 6.6.5.} \uline{With the notation of} 3.10, \(\Fil_n\Ch A\subset(\Ch A)_n\).

By the preceding corollary, \(\gamma^n(x-\epsilon(x))\in(\Ch A)_n\). The corollary then follows from the fact that the two filtrations are ring filtrations, and that the \(\gamma^n(x-\epsilon(x))\), for varying \(n\) and \(x\), generate the \(\lambda\)-ring filtration.

\sgaStableTarget{sga6.e05.cor.6-6-7}\textbf{Corollary 6.6.7.} \uline{Let} \(A\) \uline{be a} \uline{graded} \(K\)-\uline{algebra of characteristic zero}, \uline{with bounded degrees}; \uline{then the two filtrations defined on} \(\Ch A\) \uline{are} \uline{identical}.

Suppose that \(\Fil^pA=0\), and let \(x=(0,1+a_n+\ldots+a_p)\) be an element of \((\Ch A)_n\). Put \(x'=\left(0,1+\dfrac{(-1)^{n-1}}{(n-1)!}a_n+a_{n+1}+\ldots+a_p\right)\); we then have:
\(\gamma^n(x')=(0,1+a_n+\ldots+a'_p)\), since \(x'\) has zero augmentation. Consequently,
%%%%%%%%%%  scan idx 366  |  folio 353  %%%%%%%%%%
\SGAPrintPage{353}%
\(x-\gamma^n(x')\in(\Ch A)_{n+1}\). Iterating the operation \(p-n\) times, we reach \((\Ch A)_{p+1}\), which is zero, and have therefore written \(x\) as an element of \(\Fil^n\Ch A\). In view of (6.6.5), the two filtrations are therefore identical.

\sgaStableTarget{sga6.e05.para.6-7}6.7. We now suppose that \(A\) is a \(\lambda\)-ring augmented over \(K\), the augmentation homomorphism being denoted by \(\epsilon\). We therefore have on \(A\) the canonical filtration of an augmented \(\lambda\)-ring, and can thus associate to \(\Gr A\) its Chern ring.

If \(x\in A\), then \(\gamma_t(x-\epsilon(x))\in\Ghat(A)\) by definition of the filtration of an augmented \(\lambda\)-ring. We denote by \(c(x)\), and call the \uline{total Chern class of} \(x\), the image of \(\gamma_t(x-\epsilon(x))\) in \(\Ghat(A)/\Ghat'(A)=1+\widehat{\Gr A}^{+}\); we denote by \(c^i(x)\), and call the \(i\)-th Chern class of \(x\), the \(i\)-th coefficient of \(c(x)\); we therefore have:
\sgaStableTarget{sga6.e05.eq.6-7-1}\begin{equation}\tag{6.7.1}
  c^i(x)=\cl_{\Gr^iA}\!\left(\gamma^i(x-\epsilon(x))\right)~;\qquad
  c(x)=1+c^1(x)+\ldots
\end{equation}
We then put:
\sgaStableTarget{sga6.e05.eq.6-7-2}\begin{equation}\tag{6.7.2}
  \widetilde c(x)=(\epsilon(x),c(x)),
\end{equation}
and call \(\widetilde c(x)\) the completed Chern class of \(x\).

\sgaStableTarget{sga6.e05.prop.6-8}\textbf{Proposition 6.8.} \uline{The completed Chern class} \(\widetilde c:A\longrightarrow\Ch(\Gr A)\), \uline{is} \uline{a} \(K\)-\(\lambda\)-\uline{homomorphism of} \uline{augmented} \(K\)-\(\lambda\)-\uline{algebras}.

It is clear that the homomorphism thus defined is additive; to verify that it is multiplicative, we write:
\[
\begin{aligned}
\gamma_t(xy-\epsilon(xy))
&=\gamma_t\bigl(\epsilon(y)(x-\epsilon(x))\bigr)\mathbin{.}
  \gamma_t\bigl(\epsilon(x)(y-\epsilon(y))\bigr)\mathbin{.}
  \gamma_t\bigl((x-\epsilon(x))(y-\epsilon(y))\bigr)\\
&=\bigl(\gamma_t(x-\epsilon(x))\bigr)^{\epsilon(y)}\mathbin{.}
  \bigl(\gamma_t(y-\epsilon(y))\bigr)^{\epsilon(x)}\mathbin{.}
  \bigl(\gamma_t(x-\epsilon(x))\nabla\gamma_t(y-\epsilon(y))\bigr)
\end{aligned}
\]
%%%%%%%%%%  scan idx 367  |  folio 354  %%%%%%%%%%
\SGAPrintPage{354}%
whence the result by (6.1.2), after passage to the quotient.

If \(a\in K\), \(\widetilde c(a)=(a,1)\), so \(\widetilde c\) is a homomorphism of \(K\)-algebras; moreover, it is compatible with the augmentations.

Finally, it remains to see that it is a \(\lambda\)-homomorphism. We must therefore prove the commutativity of the diagram:
\[
\begin{tikzcd}[column sep=large,row sep=large]
A \arrow[r,"\gamma_t"] \arrow[d,"\widetilde c"'] &
\widehat G_g(A) \arrow[d,"\widehat G_g(\widetilde c)"]\\
\Ch(\Gr A) \arrow[r,"\gamma_t"] &
\widehat G_g(\Ch(\Gr A))
\end{tikzcd}
\]

By additivity, it suffices to prove this for the elements of \(K\), for which it is evident, and for the elements of zero augmentation of \(A\), for which it follows, by the definition of \(c\) and (6.1.4), from that of diagram (3.7.2).

\sgaStableTarget{sga6.e05.prop.6-9}\textbf{Proposition 6.9.} \uline{Let} \(A\) \uline{be an augmented} \(K\)-\(\lambda\)-\uline{algebra}, \(z\in\Fil^mA\); \uline{then}:
\[
  \gamma^m(z)\equiv(-1)^{m-1}(m-1)!z\ \mathrm{mod.}\ \Fil^{m+1}A,
\]
\uline{i}.\uline{e}.
\[
  c^m(z)=\cl_{\Gr^mA}\!\left((-1)^{m-1}(m-1)!z\right).
\]

We must see that these two elements give the same image in \(\Gr^mA\), that is, that the \(m\)-th Chern class of \(z\) is the image of \((-1)^{m-1}(m-1)!z\).

We can write \(z\) as a sum \(\sum a\mathbin{.}\gamma^{i_1}(x_1)\ldots\gamma^{i_k}(x_k)\), with \(a\in A\), \(\epsilon(x_i)=0\) for every \(i\), and \(i_1+\ldots+i_k\ge m\). We then have:
\[
  \widetilde c(z)=
  \sum \widetilde c(a)\mathbin{.}
        \widetilde c(\gamma^{i_1}(x_1))\ldots
        \widetilde c(\gamma^{i_k}(x_k))
\]
Since \(\widetilde c\) is a \(\lambda\)-homomorphism, we can write:
\[
  \widetilde c(\gamma^{i_j}(x_j))=\gamma^{i_j}(\widetilde c(x_j)).
\]
%%%%%%%%%%  scan idx 368  |  folio 355  %%%%%%%%%%
\SGAPrintPage{355}%
Taking account of the definition of \(c\), and of (6.6.4), we obtain:
\[
  \widetilde c(\gamma^{i_j}(x_j))
  =
  \left(0,1+(-1)^{i_j-1}(i_j-1)!\,\gamma^{i_j}(x_j)+\ldots\right)
\]
By (6.5.3), and the fact that \(\widetilde c(a)=(a,1)\), we finally obtain:
\[
\begin{aligned}
\widetilde c(z)
&=\sum
\left(0,\left(1+
(-1)^{i_1+\ldots+i_k-1}
(i_1+\ldots+i_k-1)!
\,\gamma^{i_1}(x_1)\ldots\gamma^{i_k}(x_k)
+\ldots\right)^a\right)\\
&=
\left(0,1+
\sum_{\substack{(x_1,\ldots,x_k)\\ i_1+\ldots+i_k=m}}
(-1)^{m-1}(m-1)!a\,\gamma^{i_1}(x_1)\ldots\gamma^{i_k}(x_k)
+\ldots\right)\\
&=
(0,1+(-1)^{m-1}(m-1)!\bar z+\ldots),
\end{aligned}
\]
whence the result.

\sgaStableTarget{sga6.e05.prop.6-10}\textbf{Proposition 6.10.} \uline{Let} \(K\) \uline{be a binomial ring}, \(R\) \uline{an augmented} \(K\)-\(\lambda\)-\uline{algebra}, \(N\) \uline{an element of} \(R\) \uline{such that} \(\lambda^k(N)=0\) \uline{for} \(k>d\), \uline{and} \(x\in\Fil^iR\). \uline{Then}:
\[
  \gamma^{d+i}(N,x)-(-1)^{d+i-1}(d+i-1)!x\in\Fil^{i+1}R.
\]

We can find elements \(x_{\alpha,j}\in R\), and \(a_\alpha\in K\) such that:
\[
  x=\sum_\alpha a_\alpha\,
       \gamma^{i_1}(x_{\alpha,1})\ldots
       \gamma^{i_k}(x_{\alpha,k}),
\]
with \(i_1+\ldots+i_k\ge i\), and \(\epsilon(x_{\alpha,j})=0\) for every \(\alpha\) and every \(j\). Let \(R_0\) be the \(K\)-\(\lambda\)-algebra generated by generators \(N'_0\), \(X_{\alpha,j}\) subject to the relations:
\[
  \gamma^k(N'_0)=0\quad\text{for }k>d.
\]
We define an augmentation homomorphism \(R_0\to K\) by \(\epsilon(N'_0)=\epsilon(X_{\alpha,j})=0\). By (4.9.1), \(R_0\) is the ring of polynomials with coefficients in \(K\) in the indeterminates \(\gamma^n(N'_0)\), with \(1\le n\le d\), and \(\gamma^p(X_{\alpha,j})\) with \(p\ge1\), and the \(\lambda\)-filtration of \(R_0\) is identical to its weight filtration, assigning the weight \(k\) to \(\gamma^k(N'_0)\) and \(\gamma^k(X_{\alpha,j})\) (4.12 and 4.15).

We put:
\[
  X_0=\sum_\alpha a_\alpha\,
       \gamma^{i_1}(X_{\alpha,1})\ldots
       \gamma^{i_k}(X_{\alpha,k}),
\]
%%%%%%%%%%  scan idx 369  |  folio 356  %%%%%%%%%%
\SGAPrintPage{356}%
and shall first prove 6.10 for \(X_0\) and \(N_0=N'_0+d\) (which do satisfy the hypotheses of 6.10). In view of the preceding remark on the filtration of \(R_0\), it suffices for this to show that:
\[
  \gamma^{d+i}(N_0,X_0)\gamma^d(N'_0)
  -
  (-1)^{d+i-1}(d+i-1)!X_0\gamma^d(N'_0)
  \in\Fil^{d+i+1}R_0 .
\]
Now we have \(\lambda_{-1}(N_0)=(-1)^d\gamma^d(N'_0)\); this relation can therefore be written:
\[
  (-1)^d\gamma^{d+i}(X_0\ (-1)^d\gamma^d(N'_0))
  -
  (-1)^{d+i-1}(d+i-1)!X_0\gamma^d(N'_0)
  \in\Fil^{d+i+1}R_0 .
\]
Put \(z=(-1)^dX_0\gamma^d(N'_0)\); then \(z\in\Fil^{d+i}R_0\), and we are reduced to showing:
\[
  \gamma^{d+i}(z)-(-1)^{d+i-1}(d+i-1)!z\in\Fil^{d+i+1}R_0 .
\]
Now this follows from 6.9.

To deduce 6.10 from this particular case, we consider the \(K\)-\(\lambda\)-homomorphism \(R_0\to R\) which sends \(N'_0\) to \(N-d\) and \(X_{\alpha,j}\) to \(x_{\alpha,j}\) for every \(\alpha\) and every \(j\) (4.7). This homomorphism is compatible with the \(\lambda\)-filtrations, and sends \(N_0\) to \(N\), and \(X_0\) to \(x\), whence the result.

\sgaStableTarget{sga6.e05.prop.6-11}\textbf{Proposition 6.11.} Suppose that \(K\) is of characteristic zero. Let \(n\ge0\) be an integer. Then the functor \(\Ch\) is an equivalence of categories from the category of graded \(K\)-algebras \(A\) such that \(A_0=K\), \(A_k=0\) for \(k>n\), to the category of augmented \(K\)-\(\lambda\)-algebras \(\mathord{\wedge}\) such that \(\Fil^{n+1}\mathord{\wedge}=0\); a quasi-inverse functor is given by \(\mathord{\wedge}\rightsquigarrow\Gr\mathord{\wedge}\).

Let \(A\) be a graded \(K\)-algebra such that \(A_k=0\) for \(k>n\), endowed with the filtration associated to its grading; then \((\Ch A)_{n+1}=0\), hence \(\Fil^{n+1}\Ch A=0\) (6.6.5). Conversely, if \(\mathord{\wedge}\) is an augmented \(K\)-\(\lambda\)-algebra such that \(\Fil^{n+1}\mathord{\wedge}=0\), then \(\Gr^0\mathord{\wedge}=K\), \(\Gr^k\mathord{\wedge}=0\) for \(k>n\).

We shall define isomorphisms of functors:
\[
  A\longrightarrow \Gr(\Ch A),
  \qquad
  \mathord{\wedge}\longrightarrow \Ch(\Gr\mathord{\wedge}).
\]

%%%%%%%%%%  scan idx 370  |  folio 357  %%%%%%%%%%
\SGAPrintPage{357}%
Since \(A\) is of finite grading, the graded rings associated with the \(\lambda\)-ring filtration and the degree filtration on \(\Ch A\) are identical (6.6.7). We then define \(u:A\longrightarrow \Gr_{\deg}(\Ch A)\) by the identity on \(A_0\) and by:
\[
  \forall\, i\ge1,\quad\forall\, a_i\in A_i,\qquad u(a_i)=\cl_{\Gr^i}\left(0,1+(-1)^{i-1}(i-1)!a_i\right),
\]
and extend by additivity. By relation (6.5.3), this is a homomorphism of graded rings, and it is easy to see that it is compatible with the operations of \(K\). The inverse morphism \(u^{-1}\) can be defined as follows: \(\forall\, i\ge1\), \(\forall\,\bar\varphi\in\Gr^i(\Ch A)\), the \(i\)-th coefficient \(a_i\) of \(\varphi\) does not depend on the way in which \(\bar\varphi\) has been lifted to \((\Ch A)_i\), and we put:
\[
  u^{-1}(\bar\varphi)=\frac{(-1)^{i-1}}{(i-1)!}a_i.
\]
The homomorphism \(u\) is therefore an isomorphism.

The isomorphism \(\mathord{\wedge}\to\Ch(\Gr\mathord{\wedge})\) is defined by the completed Chern class \(\widetilde c\). To verify that it does define an isomorphism, it suffices to verify that the homomorphism obtained by passing to the associated graded rings is one, since the filtrations are discrete. We therefore obtain:
\[
  \Gr\mathord{\wedge}
   \xrightarrow{\Gr(\widetilde c)}
  \Gr_{\lambda}(\Ch(\Gr\mathord{\wedge}))
   \simeq
  \Gr_{\deg}(\Ch(\Gr\mathord{\wedge}))
   \xrightarrow{u^{-1}}
  \Gr\mathord{\wedge} .
\]
It suffices to show that the composite is the identity. Since the homomorphisms are homomorphisms of graded \(K\)-algebras, it suffices, by the definition of the \(\lambda\)-ring filtration on \(\mathord{\wedge}\), to show that it is the identity on the \(\overline{\gamma^i(x)}\), with \(\epsilon(x)=0\). We then obtain:
\[
  \Gr(\widetilde c)\bigl(\overline{\gamma^i(x)}\bigr)
  =
  \cl_{\Gr^i}\bigl(\gamma^i(\widetilde c(x))\bigr).
\]
Now \(\gamma^i(\widetilde c(x))=(0,1+(-1)^{i-1}(i-1)!c^i(x)+\ldots)\), by (6.6.4); hence:
\[
  \cl_{\Gr^i}\bigl(\gamma^i(\widetilde c(x))\bigr)=u(c^i(x)),
\]
%%%%%%%%%%  scan idx 371  |  folio 358  %%%%%%%%%%
\SGAPrintPage{358}%
which shows that \(u^{-1}\circ\Gr(\widetilde c)=\Id\).

\sgaStableTarget{sga6.e05.cor.6-12}\textbf{Corollary 6.12.} \uline{If} \(K\) \uline{is of characteristic zero}, \uline{and} \(\mathord{\wedge}\) \uline{is a} \(K\)-\(\lambda\)-\uline{algebra} \uline{with discrete filtration}, \uline{then the completed Chern class}:
\[
  \widetilde c:\mathord{\wedge}\longrightarrow \Ch(\Gr\mathord{\wedge})
\]
\uline{is an isomorphism}.

The proof is contained in that of 6.11.

\sgaStableTarget{sga6.e05.def.6-13}\textbf{Definition 6.13.} \uline{Let} \(A\) \uline{be a graded} \(K\)-\uline{algebra}, \uline{and} \(\mathord{\wedge}\) \uline{an} \uline{augmented} \(K\)-\(\lambda\)-\uline{algebra}; \uline{we call a Chern theory for} \(\mathord{\wedge}\) \uline{with values in} \(A\) \uline{any} \(\lambda\)-\uline{homomorphism}: \(\mathord{\wedge}\longrightarrow \Ch(A)\), \uline{compatible with the augmentations}.

We define a functor on the category of graded \(K\)-algebras, with values in the category of sets, by associating with every graded \(K\)-algebra \(A\) the set of Chern theories for \(\mathord{\wedge}\) with values in \(A\), namely \(\Hom_\lambda(\mathord{\wedge},\Ch(A))\).

\sgaStableTarget{sga6.e05.prop.6-14}\textbf{Proposition 6.14.} \uline{The functor defined in} 6.13 \uline{is representable}. \uline{If} \(K\) \uline{is} \uline{of characteristic zero}, \uline{and} \(\mathord{\wedge}\) \uline{has discrete filtration}, \uline{it is represented} \uline{by the graded ring associated with} \(\mathord{\wedge}\).

Let us first prove the first assertion. We consider the graded \(K\)-algebra generated by generators \(c^i(x)\), with \(i\ge1\), \(x\in\Fil^1\mathord{\wedge}\), where \(c^i(x)\) has degree \(i\), and take its quotient by the following relations:
\[
\begin{array}{ll}
\text{i)}&
 c^k(x+y)=\displaystyle\sum_{i+j=k}c^i(x)c^j(y),
 \quad\text{for } k\ge1,\ x,y\in\Fil^1\mathord{\wedge}\ ;\\[0.8em]
\text{ii)}&
 c^k(xy)=P''_k(c^i(x),c^j(y))
 \quad(\text{cf. }6.1),
 \quad\text{for } k\ge1,\ x,y\in\Fil^1\mathord{\wedge}\ ;\\[0.4em]
\text{iii)}&
 c^k(ax)=R_{k,a}(c^i(x)) \quad(\text{cf. }6.1),
 \quad\text{for } k\ge1,\ a\in K,\ x\in\Fil^1\mathord{\wedge}\ ;\\[0.4em]
\text{iv)}&
 c^k(\gamma^j(x))=Q''_{j,k}(c^i(x))
 \quad(\text{cf. }6.1),
 \quad\text{for } j,k\ge1,\ x\in\Fil^1\mathord{\wedge},
\end{array}
\]

%%%%%%%%%%  scan idx 372  |  folio 359  %%%%%%%%%%
\SGAPrintPage{359}%
These relations being isobaric, the quotient is a graded \(K\)-algebra \(\Omega\). We define a morphism \(\mathord{\wedge}\xrightarrow{\widetilde c_0}\Ch(\Omega)\) by:
\[
  x\rightsquigarrow
  \bigl(\epsilon(x),1+c^1(x-\epsilon(x))+\ldots\bigr).
\]
The imposed relations show that this is a Chern theory for \(\mathord{\wedge}\) with values in \(\Omega\). If in addition we have a Chern theory \(f\), with values in a graded \(K\)-algebra \(A\), there is a unique homomorphism: \(\Omega\to A\), such that \(f\) factors through \(\Ch(\Omega)\); it is defined by sending, for every \(i\ge1\) and every \(x\in\Fil^1\mathord{\wedge}\), the element \(c^i(x)\) to the \(i\)-th coefficient of \(f(x)\) (which is well defined since \(f\) is a \(\lambda\)-homomorphism).

Now suppose that \(K\) is of characteristic zero and \(\mathord{\wedge}\) has discrete filtration; let \(n\) be such that \(\Fil^n\mathord{\wedge}=0\). Then the graded algebra \(\Omega\) is zero in degrees \(\ge n\). To prove this, we must show that for \(x_1,\ldots,x_p\in\Fil^1\mathord{\wedge}\) and \(k_1+\ldots+k_p\ge n\), we have:
\[
  c^{k_1}(x_1)\mathinner{....}c^{k_p}(x_p)=0.
\]
Now, by hypothesis, we have:
\[
  \gamma^{k_1}(x_1)\mathinner{....}\gamma^{k_p}(x_p)=0\ ;
\]
Applying the \(\lambda\)-homomorphism \(\widetilde c_0\), we obtain:
\[
  \gamma^{k_1}(\widetilde c_0(x_1))\mathinner{....}
  \gamma^{k_p}(\widetilde c_0(x_p))=(0,1).
\]
Now, by (6.6.4) and (6.5.3), this implies:
\[
  (k_1+\ldots+k_p)!\,
  c^{k_1}(x_1)\ldots c^{k_p}(x_p)=0.
\]
Since \(K\) is of characteristic zero, we have the desired result.

The universal property of \(\Omega\), applied to the canonical homomorphism \(\widetilde c\), gives a commutative diagram:
%%%%%%%%%%  scan idx 373  |  folio 360  %%%%%%%%%%
\SGAPrintPage{360}%
\[
\begin{tikzcd}[column sep=large,row sep=large]
\mathord{\wedge} \arrow[r,"\widetilde c"] \arrow[d,"\widetilde c_0"'] &
\Ch(\Gr\mathord{\wedge})\\
\Ch(\Omega) \arrow[ur] &
\end{tikzcd}
\]
By Proposition 6.11, the homomorphism: \(\Omega\to\Gr\mathord{\wedge}\) is an isomorphism if and only if the homomorphism: \(\Ch(\Omega)\to\Ch(\Gr\mathord{\wedge})\) is one. Since it is visibly surjective, the latter is also surjective, and it therefore suffices to show that \(\Ch(\Omega)\to\Ch(\Gr\mathord{\wedge})\) is injective. By functoriality of \(\Ch\), we obtain the diagram:
\[
\begin{tikzcd}[column sep=large,row sep=large]
\mathord{\wedge} \arrow[r,"\widetilde c_0"] \arrow[d,"\widetilde c"'] &
\Ch(\Omega) \arrow[d,"\widetilde c"]\\
\Ch(\Gr\mathord{\wedge}) \arrow[r,"\Ch(\widetilde c_0)"] &
\Ch\bigl(\Gr_\lambda(\Ch(\Omega))\bigr);
\end{tikzcd}
\]
since \(\Omega\) has bounded degrees, the right vertical arrow is an isomorphism (6.12), and using 6.11, we therefore obtain a morphism: \(\Gr\mathord{\wedge}\longrightarrow\Omega\), the composite of the two being the identity by the universal property of \(\Omega\).

\sgaStableTarget{sga6.e05.sec.7}\section*{APPENDIX: The Adams operations \(\psi^k\)}

We follow here a presentation due to J.P. Serre.

\sgaStableTarget{sga6.e05.para.7-1}7.1. Let \(K\) be a pre-\(\lambda\)-ring. For every \(x\in K\), we denote by \(\frac{d}{dt}(\lambda_t(x))\) the derivative with respect to \(t\) of \(\lambda_t(x)\). We then define elements of \(K\), denoted by \(\psi^k(x)\), by putting:
\sgaStableTarget{sga6.e05.eq.7-1-1}\begin{equation}\tag{7.1.1}
  \frac{d}{dt}(\lambda_t(x))/\lambda_t(x)
  =
  \sum_{k=1}^{\infty}(-1)^{k-1}\psi^k(x)t^{k-1}.
\end{equation}

The maps \(\psi^k:K\to K\) are additive; indeed, we have:
\[
\begin{aligned}
 \frac{d}{dt}(\lambda_t(x+y))/\lambda_t(x+y)
 &=
 \frac{d}{dt}(\lambda_t(x)\lambda_t(y))/(\lambda_t(x)\lambda_t(y))\\
 &=
 \frac{d}{dt}(\lambda_t(x))/\lambda_t(x)
 +\frac{d}{dt}(\lambda_t(y))/\lambda_t(y)\ ;
\end{aligned}
\]
%%%%%%%%%%  scan idx 374  |  folio 361  %%%%%%%%%%
\SGAPrintPage{361}%
and therefore, for every \(k\):
\[
  \psi^k(x+y)=\psi^k(x)+\psi^k(y).
\]

By (7.1.1), the \(\psi^k\) are isobaric polynomials of weight \(k\) in the \(\lambda^i\), with \(i=1,\ldots,k\). For example, let us give the values of \(\psi^1\) and \(\psi^2\):
\[
  \psi^1(x)=\lambda^1(x)=x,
  \qquad
  \psi^2(x)=x^2-2\lambda^2(x).
\]

\sgaStableTarget{sga6.e05.para.7-2}7.2. Let \(K\) be a ring of characteristic \(0\); suppose given additive maps \(\psi^k\) from \(K\) to \(K\), such that \(\psi^1=\operatorname{Id}_K\). We then define a pre-\(\lambda\)-structure on \(K\) by putting:
\sgaStableTarget{sga6.e05.eq.7-2-1}\begin{equation}\tag{7.2.1}
  \lambda_t(x)=
  \exp\left(\sum_{k=1}^{\infty}(-1)^{k-1}\psi^k(x)t^k/k\right).
\end{equation}
Multiplicativity follows from the additivity of the \(\psi^k\). There is therefore an equivalence between giving a pre-\(\lambda\)-structure on \(K\) and giving the corresponding operations \(\psi^k\).

\sgaStableTarget{sga6.e05.para.7-3}7.3. It is clear from formula (7.1.1) that the \(\psi^k\) commute with \(\lambda\)-homomorphisms. If \(K,K'\) are pre-\(\lambda\)-rings of characteristic zero, it is necessary and sufficient that a homomorphism \(f:K\to K'\) commute with the \(\psi^k\) for it to be a \(\lambda\)-homomorphism: this follows from 7.2, the \(\lambda^i\) being expressed as polynomials in the \(\psi^k\).

\sgaStableTarget{sga6.e05.ex.7-4}\textbf{Examples 7.4.} i) For \(K=\mathbb Z\), endowed with its \(\lambda\)-ring structure, we have:
\[
  \frac{d}{dt}(\lambda_t(n))/\lambda_t(n)
  =
  n(1+t)^{n-1}/(1+t)^n
  =
  n.\sum_{k=1}(-1)^{k-1}t^{k-1}.
\]
Consequently, the \(\psi^k\) are all equal to the identity map of \(\mathbb Z\).

ii) Let us work in the category \(\mathcal C_0\) of \(\mathbb Z\)-algebras (1.8). The
%%%%%%%%%%  scan idx 375  |  folio 362  %%%%%%%%%%
\SGAPrintPage{362}%
operations \(\psi^k\) define, by 7.3, additive endomorphisms of the functor \(\widehat G\) with values in \(\lambda\)-rings; by 1.10, they are characterized by their value on \(1+Xt\). We have:
\[
  \frac{d}{du}(\lambda_u(1+Xt))/\lambda_u(1+Xt)
  =
  \frac{1+Xt}{\underline{1}+\{1+Xt\}u}
  =
  \sum_{k=1}^{\infty}(-1)^{k-1}\{1+Xt\}^{\circ k}u^{k-1}.
\]
Consequently:
\sgaStableTarget{sga6.e05.eq.7-4-1}\begin{equation}\tag{7.4.1}
  \psi^k(1+Xt)=1+X^k t.
\end{equation}

\sgaStableTarget{sga6.e05.prop.7-5}\textbf{Proposition 7.5.} i) \uline{If} \(K\) \uline{is a} \(\lambda\)-\uline{ring}, \uline{the operations} \(\psi^k\) \uline{satisfy} \uline{the identities}:
\sgaStableTarget{sga6.e05.eq.7-5-1}\begin{equation}\tag{7.5.1}
  \psi^k(1)=1~;\qquad
  \psi^k(xy)=\psi^k(x)\psi^k(y)~;\qquad
  \psi^k\circ\psi^{k'}=\psi^{kk'}.
\end{equation}
ii) \uline{Conversely}, \uline{if the} \(\psi^k\) \uline{satisfy these identities}, \uline{and if} \(K\) \uline{is of} \uline{characteristic zero}, \(K\) \uline{is a} \(\lambda\)-\uline{ring}.

i) If \(K\) is a \(\lambda\)-ring, we have \(\lambda_t(1)=1+t\), whence: \(\frac{d}{dt}(\lambda_t(1))/\lambda_t(1)=1/(1+t)=\sum_{k=1}^{\infty}(-1)^{k-1}t^{k-1}\), and consequently:
\[
  \psi^k(1)=1.
\]

To prove the second of relations (7.5.1), we consider the diagram:
\sgaStableTarget{sga6.e05.eq.7-5-2}\begin{equation}\tag{7.5.2}
\begin{tikzcd}[column sep=large,row sep=large]
K\times K \arrow[r,"\lambda_t\times\lambda_t"] \arrow[d] &
\widehat G(K)\times\widehat G(K) \arrow[r,"\Phi^k\times\Phi^k"] \arrow[d,"\circ"'] &
K\times K \arrow[d]\\
K \arrow[r,"\lambda_t"] &
\widehat G(K) \arrow[r,"\Phi^k"] &
K ,
\end{tikzcd}
\end{equation}
where the vertical arrows are defined by multiplication, and the maps \(\Phi^k\) as follows: the ``logarithmic derivative'' map is
%%%%%%%%%%  scan idx 376  |  folio 363  %%%%%%%%%%
\SGAPrintPage{363}%
a group homomorphism from \(\widehat G(K)\) to \(K[[t]]\); the map which assigns to a series its \((k-1)\)-st coefficient, multiplied by \((-1)^{k-1}\), is an additive homomorphism from \(K[[t]]\) to \(K\); the morphism \(\Phi^k\) will be the composite of the two. It is clear that \(\Phi^k\circ\lambda_t=\psi^k\) by definition (7.1.1).

The second of formulas (7.5.1) therefore expresses the commutativity of the diagram written above. If we suppose that \(K\) is a \(\lambda\)-ring, the left square is commutative, and it suffices to verify the commutativity of the right one. Now this is a diagram of homomorphisms of functors on \(\mathcal C_0\), with source \(\widehat G\times\widehat G\); it therefore suffices, by 1.10, to verify commutativity for \((1+Xt,1+Yt)\); we then have:
\[
  \Phi^k(1+Xt)=X^k\qquad\bigl(\text{$(k-1)$-st coefficient of }(-1)^{k-1}X/(1+Xt)\bigr)~;\qquad
  \Phi^k(1+Yt)=Y^k~;
\]
\(\Phi^k\bigl((1+Xt)\circ(1+Yt)\bigr)=\Phi^k(1+XYt)=(XY)^k=\Phi^k(1+Xt)\Phi^k(1+Yt)\), which gives the result.

To prove the last one, we first observe that the diagram:
\sgaStableTarget{sga6.e05.eq.7-5-3}\begin{equation}\tag{7.5.3}
\begin{tikzcd}[column sep=large,row sep=large]
\widehat G^t(K) \arrow[r,"\lambda_u"] \arrow[d,"\Phi^{kk'}"'] &
\widehat G^u\circ\widehat G^t(K) \arrow[d,"\Phi^{k'}_u"]\\
K &
\widehat G^t(K), \arrow[l,"\Phi^k"']
\end{tikzcd}
\end{equation}
is commutative. As above, it indeed suffices to see this for \(1+Xt\), and we have:
\[
  \Phi^{k'}_u(\lambda_u(1+Xt))
  =
  \Phi^{k'}_u(1+\{1+Xt\}u)
  =
  (1+Xt)^{\circ k'}
  =
  1+X^{k'}t~;
\]
\[
  \Phi^k(1+X^{k'}t)=X^{kk'}=\Phi^{kk'}(1+Xt).
\]

Consequently, \(\psi^{kk'}=\Phi^k\circ\Phi^{k'}_u\circ\lambda_u\circ\lambda_t=\Phi^k\circ\psi^{k'}\circ\lambda_t\).

Moreover, \(\psi^k\circ\psi^{k'}=\Phi^k\circ\lambda_t\circ\Phi^{k'}\circ\lambda_t=\Phi^k\circ\lambda_t\circ\psi^{k'}\).

%%%%%%%%%%  scan idx 377  |  folio 364  %%%%%%%%%%
\SGAPrintPage{364}%
Now, if \(K\) is a \(\lambda\)-ring, the diagram:
\sgaStableTarget{sga6.e05.eq.7-5-4}\begin{equation}\tag{7.5.4}
\begin{tikzcd}[column sep=large,row sep=large]
K \arrow[r,"\lambda_t"] \arrow[d,"\psi^{k'}"'] &
\widehat G(K) \arrow[d,"\psi^{k'}"]\\
K \arrow[r,"\lambda_t"] &
\widehat G(K)
\end{tikzcd}
\end{equation}
is commutative by 7.3, since \(\lambda_t\) is a \(\lambda\)-homomorphism. The desired relations follow at once.

ii) If \(K\) is of characteristic zero, the ``logarithmic derivative'' map is an isomorphism from \(\widehat G(K)\) onto \(K[[t]]\), since one can integrate. Consequently, two elements of \(\widehat G(K)\) are equal if and only if, for every \(k\), they have the same image under \(\Phi^k\).

To show that \(\lambda_t\) is a ring homomorphism, we must show the commutativity of the left square of diagram (7.5.2), and by the preceding remark, it suffices to show that for every \(k\), diagram (7.5.2) is commutative, since the right square always is. Now these are precisely the relations \(\psi^k(xy)=\psi^k(x)\psi^k(y)\).

To show that \(\lambda_t\) is a \(\lambda\)-homomorphism, it suffices to show that diagram (7.5.4) is commutative for every \(k'\), by 7.3. Now, by the preceding remark, this is equivalent to showing for all \(k,k'\) that \(\Phi^k\circ\psi^{k'}\circ\lambda_t=\Phi^k\circ\lambda_t\circ\psi^{k'}\), that is, the relations \(\psi^{kk'}=\psi^k\circ\psi^{k'}\).

Finally, the relation \(\lambda_t(1)=1+t\) is obtained by integration from the relations \(\psi^k(1)=1\).

\sgaStableTarget{sga6.e05.bibliography}\subsection*{BIBLIOGRAPHY}
\begin{enumerate}[label={[\arabic*]}]
\item A. Grothendieck, La théorie des classes de Chern, Bull. Soc. math. France, 86, 1958, p. 137 à 154.
\item F. Hirzebruch, Neue topologische Methoden in der algebraischen Geometrie.
\end{enumerate}

% SGA_NATIVE_CONTINUATION_364
\providecommand{\Sym}{\operatorname{Sym}}
\providecommand{\Proj}{\operatorname{Proj}}
\providecommand{\Projz}{\operatorname{Proj}_0}
\providecommand{\Spec}{\operatorname{Spec}}
\providecommand{\Comp}{\operatorname{Comp}}
\providecommand{\Dplus}{D_{+}}
\providecommand{\id}{\operatorname{id}}
\providecommand{\cG}{\mathcal G}
\providecommand{\cA}{\mathcal A}
\providecommand{\cB}{\mathcal B}
\providecommand{\cL}{\mathscr L}
\providecommand{\cM}{\mathscr M}
\providecommand{\cN}{\mathscr N}
\providecommand{\cE}{\mathscr E}
\providecommand{\cK}{\mathscr K}
\providecommand{\cP}{\mathscr P}
\providecommand{\cQ}{\mathscr Q}
\providecommand{\OO}{\mathcal O}
\providecommand{\rk}{\operatorname{rk}}
\providecommand{\colim}{\operatorname*{colim}}

% SGA_REVIEWED_FRENCH_CONTINUATION_364
% Exposé VI projective-bundle macro additions
\providecommand{\Sym}{\operatorname{Sym}}
\providecommand{\Proj}{\operatorname{Proj}}
\providecommand{\Projz}{\operatorname{Proj}_0}
\providecommand{\Spec}{\operatorname{Spec}}
\providecommand{\Comp}{\operatorname{Comp}}
\providecommand{\Dplus}{D_{+}}
\renewcommand{\id}{\operatorname{id}}
\providecommand{\cG}{\mathcal G}
\providecommand{\cA}{\mathcal A}
\providecommand{\cB}{\mathcal B}
\providecommand{\cL}{\mathscr L}
\providecommand{\cM}{\mathscr M}
\providecommand{\cN}{\mathscr N}
\providecommand{\cE}{\mathscr E}
\providecommand{\cK}{\mathscr K}
\providecommand{\cP}{\mathscr P}
\providecommand{\cQ}{\mathscr Q}
\providecommand{\OO}{\mathcal O}
\providecommand{\rk}{\operatorname{rk}}
\providecommand{\colim}{\operatorname*{colim}}

%%%%%%%%%%  scan idx 378  |  folio 365  %%%%%%%%%%
\SGAPrintPage{365}%
\sgaStableTarget{sga6.e06.expose}\section*{\underline{Exposé VI}\\[3pt] \underline{\underline{THE \(K^{\bullet}\) OF A PROJECTIVE BUNDLE:}}\\[3pt] \underline{\underline{COMPUTATIONS AND CONSEQUENCES}}}
% Exposé VI: reviewed local symmetric-algebra and Proj notation.
\renewcommand{\Sym}{\mathbf{S}}
\renewcommand{\Projz}{\underline{\operatorname{Proj}}_0}
\begin{center}
\underline{by Pierre BERTHELOT}
\end{center}

\sgaStableTarget{sga6.e06.sec.1}\section*{1. Computation of the \(K^{\bullet}\) of a projective bundle: the case of locally free sheaves of finite type}

In this section, the notation \(K^{\bullet}(X)\) will denote the Grothendieck group of locally free sheaves of finite type on a scheme \(X\).

The aim of this section is to prove the following theorem:

\sgaStableTarget{sga6.e06.thm.1-1}\textbf{Theorem 1.1.} \uline{Let} \(S\) \uline{be a quasi-compact scheme}, \(\underline{E}\) \uline{a} \(\underline{O}_S\)\uline{-Module locally free of rank} \(r+1\), \(X=\mathbf P(\underline{E})\) \uline{the associated projective bundle}; \uline{the structural morphism}: \(X\xrightarrow{f}S\) \uline{defines a ring homomorphism}: \(K^{\bullet}(S) \xrightarrow{f^*} K^{\bullet}(X)\) \uline{endowing} \(K^{\bullet}(X)\) \uline{with the structure of a} \(K^{\bullet}(S)\)-\uline{algebra}. \uline{Let} \(\xi\in K^{\bullet}(X)\) \uline{be the class of the sheaf} \(\underline{O}_X(1)\). \uline{Then} \(K^{\bullet}(X)\) \uline{is a free} \(K^{\bullet}(S)\)-\uline{module with basis} \(1,\ldots,\xi^r\).

\uline{Plan of the proof}:

\(1^\circ\)) The powers of \(\xi\) generate \(K^{\bullet}(X)\) as a \(K^{\bullet}(S)\)-module.

To prove this, we shall establish the following points:

\begin{enumerate}[label=\alph*)]
\item If \(\underline{F}\) is a \(\underline{O}_X\)-Module locally free of finite type, there exists a graded \(\Sym(\underline{E})\)-Module \(\underline{M}\) of finite presentation, flat over \(\underline{O}_S\), such that \(\underline{F}=\Projz(\underline{M})\) (\(\Sym(\underline{E})\) being the symmetric algebra of \(\underline{E}\); for the notation \(\Projz\), see EGA II 8.12.1).

\item There exists a finite graded resolution of \(\underline{M}\) by graded \(\Sym(\underline{E})\)-Modules \(\underline{L}_i\) locally free of finite type: \(0\longrightarrow \underline{L}_{r+1}\longrightarrow\ldots\longrightarrow \underline{L}_0\longrightarrow \underline{M}\longrightarrow 0\).

\item The \(\underline{L}_i\) admit a finite filtration whose quotients are of the form \(\underline{N}_j\otimes_{\underline{O}_S}\Sym(\underline{E})(j)\), where the \(\underline{N}_j\) are locally free of finite type on \(S\).
\end{enumerate}

%%%%%%%%%%  scan idx 379  |  folio 366  %%%%%%%%%%
\SGAPrintPage{366}%
\(2^\circ\)) \(1,\ldots,\xi^{r+1}\) are linearly dependent over \(K^{\bullet}(S)\).

\(3^\circ\)) \(1,\ldots,\xi^r\) are linearly independent over \(K^{\bullet}(S)\).

\sgaStableTarget{sga6.e06.sec.1.2}1.2.
For convenience, we shall recall here a few classical definitions and properties of graded Modules.

Let \(S\) be a scheme, \(\underline{B}\) a graded \(\underline{O}_S\)-Algebra, \(\underline{M}\) a graded \(\underline{B}\)-Module, For every graded \(\underline{B}\)-Module \(\underline{N}\), we denote as usual by \(\underline{N}(n)\) the graded \(\underline{B}\)-Module defined by:
\[
  (\underline{N}(n))_k=\underline{N}_{n+k}\quad\text{for every }k.
\]
We say that \(\underline{M}\) is a locally free graded \(\underline{B}\)-Module (resp. locally free of finite type) if every point of \(S\) has an open neighborhood \(U\) such that \(\underline{M}|U\) is isomorphic, as a graded \(\underline{B}\)-Module, to a direct sum of Modules of the form \(\underline{B}(n)|U\) (resp. a finite direct sum).

We say that \(\underline{M}\) is a graded \(\underline{B}\)-Module of finite presentation if every point of \(S\) has an open neighborhood \(U\) such that there is an exact sequence \(\underline{P}\longrightarrow \underline{Q}\longrightarrow \underline{M}|U\longrightarrow 0\), where \(\underline{P}\) and \(\underline{Q}\) are free graded \((\underline{B}|U)\)-Modules of finite type, and the homomorphisms are compatible with the gradings.

1.2.1. For every integer \(n\), we can define a canonical homomorphism of graded \(\underline{B}\)-Modules:
\sgaStableTarget{sga6.e06.eq.1-2-1}\begin{equation}\tag{1.2.1}
  \underline{M}_n\otimes_{\underline{O}_S} \underline{B}(-n)\longrightarrow \underline{M},
\end{equation}
by sending \(x\otimes a\), where \(x\) is a section of \(\underline{M}_n\), and \(a\) a section of \((\underline{B}(-n))_k=\underline{B}_{k-n}\), to the element \(a.x\) of \(\underline{M}_k\).

1.2.2. We say that a graded \(\underline{B}\)-Module \(\underline{M}\) is free of finite type and homogeneous of degree \(n\) if it is isomorphic as a graded \(\underline{B}\)-Module to \((\underline{B}(-n))^k\) for some integer \(k\).

%%%%%%%%%%  scan idx 380  |  folio 367  %%%%%%%%%%
\SGAPrintPage{367}%
Let \(\underline{M}\) be a free graded \(\underline{B}\)-Module of finite type. By 1.2, we can find a graded isomorphism:
\sgaStableTarget{sga6.e06.eq.1-2-2}\begin{equation}\tag{1.2.2}
  \underline{M}\simeq \bigoplus_{n=n_0}^{n_1} \bigl(\underline{B}(-n)\bigr)^{k_n}.
\end{equation}

We thus obtain a decomposition of \(\underline{M}\) as a direct sum of free graded \(\underline{B}\)-Modules of finite type that are homogeneous, namely:
\[
  \underline{M}\simeq \bigoplus_{n=n_0}^{n_1}\underline{M}^{(n)},
\]
where \(\underline{M}^{(n)}\) is free homogeneous of degree \(n\) and of rank \(k_n\).

\sgaStableTarget{sga6.e06.prop.1-2}\textbf{Proposition 1.2.3.} \uline{The integers} \(k_n\) \uline{occurring in decomposition} (1.2.2) \uline{do not depend on the chosen decomposition}.

Suppose given two decompositions, hence a graded isomorphism:
\[
  \bigoplus_{n=n_0}^{n_1}\bigl(\underline{B}(-n)\bigr)^{k_n}\xrightarrow{\sim}
  \bigoplus_{n=n'_0}^{n'_1}\bigl(\underline{B}(-n)\bigr)^{k'_n}.
\]
The proof proceeds by induction on the number of nonzero free homogeneous direct summands occurring in the decomposition on the left. If this is \(0\), then \(\underline{M}=0\), and the proposition is evident. Suppose it proved for \(n_1-n_0=m\), and let us prove it for \(n_1-n_0=m+1\). The given isomorphism, being compatible with the gradings, gives an isomorphism on the homogeneous components of the two sides, and in particular an isomorphism of the homogeneous components of degree \(<n_0\), which shows that for \(n<n_0\), we have \(k'_n=0\), and on the homogeneous components of degree \(n_0\), which shows that they have the same rank, and consequently that \(k_{n_0}=k'_{n_0}\). The given isomorphism therefore induces an isomorphism of the free homogeneous direct summands of lowest degree in the two decompositions, and hence passes to the quotient by these direct summands; we then apply the induction hypothesis to the quotient Module, whence the result.

%%%%%%%%%%  scan idx 381  |  folio 368  %%%%%%%%%%
\SGAPrintPage{368}%
It follows in particular from 1.2.3 that for every free graded Module of finite type \(\underline{M}\), the number of elements of a homogeneous basis of \(\underline{M}\) having a given degree does not depend on the chosen basis. In particular, the integers \(n_0\) and \(n_1\) occurring in decomposition (1.2.2), defined by the conditions \(k_{n_0}\neq0\), \(k_{n_1}\neq0\), do not depend on the decomposition.

\sgaStableTarget{sga6.e06.prop.1-3}\textbf{Proposition 1.3.} \uline{Let} \(S\) \uline{be a quasi-compact scheme}, \(\underline{E}\) \uline{a} \(\underline{O}_S\)\uline{-Module locally free of finite type}, \(X\) \uline{the associated projective bundle}. \uline{Let} \(\underline{F}\) \uline{be a} \(\underline{O}_X\)\uline{-Module of finite presentation}, \uline{flat over} \(S\). \uline{There exists an integer} \(n_0\), \uline{such that for} \(n\geq n_0\), \uline{we have}:
\[
\begin{array}{ll}
\text{i)}& R^i f_*(\underline{F}(n))=0\quad\text{\uline{for} }i\geq 1;\\[2mm]
\text{ii)}& f_*(\underline{F}(n))\text{ \uline{is locally free of finite type on} }S;\\[2mm]
\text{iii)}& \underline{M}=\displaystyle\bigoplus_{n\geq n_0} f_*(\underline{F}(n))\text{ \uline{is a graded} }\Sym(\underline{E})\text{\uline{-Module of finite presentation} (1.2).}
\end{array}
\]

The desired properties are local on \(S\); since \(S\) is quasi-compact, we can, by taking the largest of the integers found, reduce to the case where \(S\) is affine.

First suppose that \(S\) is noetherian.

Property i) then follows from Serre's theorem.

Property ii) follows from EGA III 7.9.10, with \(n_0\) chosen so that property i) holds.

By EGA II 3.4.5, there exists a quasi-coherent graded \(\Sym(\underline{E})\)-Module \(\underline{M}'\), of finite type, such that \(\underline{F}=\Projz(\underline{M}')\). Moreover, \(\Gamma_*(\underline{F})=\Gamma_*(\Projz \underline{M}')=\displaystyle\bigoplus_{n\in\mathbb Z} f_*(\underline{F}(n))\) is TN-isomorphic to \(\underline{M}'\) by EGA III 2.3.1; thus there exists an integer \(n_1\) such that \(\displaystyle\bigoplus_{n\geq n_1} f_*(\underline{F}(n))\simeq \bigoplus_{n\geq n_1} \underline{M}'_n\),
%%%%%%%%%%  scan idx 382  |  folio 369  %%%%%%%%%%
\SGAPrintPage{369}%
and hence such that \(\bigoplus_{n\geq n_1} f_*(\underline{F}(n))\) is a \(\Sym(\underline{E})\)-Module of finite type, and therefore of finite presentation by the noetherian hypotheses; the result moreover holds for every integer \(n_1\), since the \(f_*(\underline{F}(n))\) are coherent by Serre's theorem.

The proposition is therefore proved in the case where \(S\) is noetherian.

Suppose now that we have a Cartesian square:
\[
\begin{tikzcd}[column sep=large,row sep=large]
X \arrow[r,"g"] \arrow[d,"f"'] & X_0 \arrow[d,"f_0"]\\
S \arrow[r,"g_0"] & S_0,
\end{tikzcd}
\]
with \(S_0\) noetherian, \(X_0=\mathbf P(\underline{E}_0)\); let \(\underline{F}_0\) be a \(\underline{O}_{X_0}\)-Module of finite presentation, flat over \(S_0\), and let \(\underline{F}=g^*(\underline{F}_0)\). By the preceding argument the proposition holds for \(\underline{F}_0\). By EGA III 7.7.5 and 7.8.5, the base-change theorem holds under the preceding hypotheses; we immediately deduce that the theorem holds for \(\underline{F}\).

To complete the proof, it therefore suffices to show that the general case (with \(S\) affine) can be obtained from the noetherian case by base change.

Since \(S\) is affine, \(S\) is a filtered projective limit of noetherian affine schemes (if \(S=\Spec(A)\), we write \(A\) as the inductive limit of its subalgebras of finite type over \(\mathbb Z\)). Thus let \(S=\varprojlim S_\lambda\). By EGA IV 8.5.2 and 8.5.5, there exists an index \(\lambda\) and a \(\underline{O}_{S_\lambda}\)-Module locally free of finite type \(\underline{E}_\lambda\) such that \(\underline{E}=u_\lambda^*(\underline{E}_\lambda)\) (denoting by \(u_\lambda\) the morphism \(S\to S_\lambda\), and by \(u_{\lambda\mu}\) the transition morphisms \(S_\mu\to S_\lambda\)).

If we put, for \(\mu\geq\lambda\), \(\underline{E}_\mu=u_{\lambda\mu}^*(\underline{E}_\lambda)\), and \(X_\mu=\mathbf P(\underline{E}_\mu)\), then for every \(\mu\), we have \(X=X_\mu\times_{S_\mu}S\), and \(X=\varprojlim X_\mu\). Consequently, there exists an index \(\lambda'\geq\lambda\)
%%%%%%%%%%  scan idx 383  |  folio 370  %%%%%%%%%%
\SGAPrintPage{370}%
and a sheaf \(\underline{F}_{\lambda'}\) of finite presentation on \(X_{\lambda'}\), such that \(\underline{F}=u'^*_{\lambda'}(\underline{F}_{\lambda'})\) (\(u'_{\lambda'}\) being the morphism \(X\to X_{\lambda'}\)). Since the \(X_\mu\) and the \(\underline{F}_\mu\) are of finite presentation, we can, by EGA IV 11.2.6, find an index \(\lambda''\) such that \(\underline{F}_{\lambda''}\) is flat over \(S_{\lambda''}\). All the hypotheses of 1.3 are therefore satisfied, which completes the proof.

\sgaStableTarget{sga6.e06.cor.1-4}\textbf{Corollary 1.4.} \uline{Under the hypotheses of} 1.3, \uline{there exists a graded} \(\Sym(\underline{E})\)\uline{-Module} \(\underline{M}\), \uline{of finite presentation}, \uline{and flat over} \(\underline{O}_S\), \uline{such that} \(\underline{F}=\Projz(\underline{M})\).

Let \(n_0\) be an integer such that the properties of 1.3 hold. If \(\underline{M}=\bigoplus_{n\geq n_0} f_*(\underline{F}(n))\), we have seen that \(\underline{M}\) is of finite presentation, and it is flat over \(S\), since the \(f_*(\underline{F}(n))\), for \(n\geq n_0\), are locally free on \(S\).

Moreover, \(\underline{F}=\Projz(\underline{M})\) by EGA II 3.4.

Let us turn to part b) of the proof.

\sgaStableTarget{sga6.e06.lem.1-5-1}\textbf{Lemma 1.5.1.} \uline{Let} \(A\) \uline{be a ring}, \(B\) \uline{a positively graded} \(A\)-\uline{algebra} \uline{such that} \(B_0=A\), \(M\) \uline{a graded} \(B\)-\uline{module}. \uline{We suppose that} \(M\) \uline{is projective of finite type}, \uline{as an ungraded} \(B\)-\uline{module}. \uline{Let} \(\widetilde B\) \uline{be the sheaf of algebras on} \(\Spec(A)\) \uline{defined by} \(B\), \(\widetilde M\) \uline{the} \(\widetilde B\)-\uline{Module associated with} \(M\); \uline{then} \(\widetilde M\) \uline{is, locally on} \(\Spec(A)\), \uline{free of finite type over} \(\widetilde B\), \uline{as a graded} \(\widetilde B\)-\uline{Module} (1.2).

Since \(M\) is of finite type, its degrees are bounded below. Moreover, \(M\otimes_B A\) is projective of finite type over \(A\), and consequently has only finitely many nonzero homogeneous components, themselves projective of finite type over \(A\); by localizing on \(\Spec(A)\), we can therefore suppose that \(M\otimes_B A\) is free of finite type over \(A\). Finally, since \(M\) is projective, \(\Tor_1^B(M,A)=0\). These properties imply that \(M\) is free of finite type over \(B\) as a graded module (see Bourbaki, Alg. ch. II §11 prop. 7), so \(\widetilde M\) is free of finite type over \(\widetilde B\), q.e.d.

%%%%%%%%%%  scan idx 384  |  folio 371  %%%%%%%%%%
\SGAPrintPage{371}%
\sgaStableTarget{sga6.e06.lem.1-5-2}\textbf{Lemma 1.5.2.} \uline{Let} \(A\) \uline{be a ring}, \(B\) \uline{an} \(A\)-\uline{algebra of finite presentation}, \(M\) \uline{a} \(B\)-\uline{module of finite presentation}. \uline{We suppose} \(B\) \uline{and} \(M\) \uline{flat over} \(A\), \uline{and} \uline{that for every residue field} \(k\) \uline{of} \(A\), \(M\otimes_A k\) \uline{has projective dimension} \(\leq n\) \uline{over} \(B\otimes_A k\) (\uline{which is the case if} \(B=A[t_0,\ldots,t_r]\)). \uline{Under these conditions}, \(M\) \uline{admits a resolution of length} \(n\) \uline{by projective} \(B\)-\uline{modules of finite type}:
\[
  0\longrightarrow L_n\longrightarrow\ldots\longrightarrow L_0\longrightarrow M\longrightarrow 0.
\]
This lemma follows immediately from EGA IV 12.3.2.

\sgaStableTarget{sga6.e06.prop.1-6}\textbf{Proposition 1.6.} \uline{Let} \(S\) \uline{be a quasi-compact scheme}, \(\underline{E}\) \uline{a} \(\underline{O}_S\)\uline{-Module locally free of rank} \(r+1\), \(\underline{B}\) \uline{its symmetric algebra}, \uline{with its usual grading}, \(\underline{M}\) \uline{a} \(\underline{B}\)\uline{-Module that is graded and of finite presentation}, \uline{flat over} \(S\). \uline{Then there exists a graded resolution}:
\[
  0\longrightarrow \underline{L}_{r+1}\longrightarrow\ldots\longrightarrow \underline{L}_0\longrightarrow \underline{M}\longrightarrow 0,
\]
\uline{where the} \(\underline{L}_i\) \uline{are} \(\underline{B}\)\uline{-Modules that are graded and locally free of finite type} (\uline{as graded} \(\underline{B}\)\uline{-Modules}).

Let us first show that for every graded \(\underline{B}\)-Module \(\underline{M}\) satisfying the conditions of the proposition, there exists a surjective graded homomorphism: \(\underline{L}\to \underline{M}\), where \(\underline{L}\) is a \(\underline{B}\)-Module locally free of finite type.

Since \(\underline{M}\) is of finite type, there exists an integer \(N\) such that \(\underline{M}\) is generated over \(\underline{B}\) by its components of degree \(\leq N\). We then put:
\[
  \underline{L}=\bigoplus_{i\leq N} \underline{M}_i\otimes_{\underline{O}_S}\underline{B}(-i);
\]
the homomorphism from \(\underline{L}\) to \(\underline{M}\) is defined using (1.2.1), and it is clear that it is surjective.

By 1.5.2, \(\underline{M}\) admits an (ungraded) resolution by \(\underline{B}\)-Modules locally free of finite type on every affine open set. Consequently, if
%%%%%%%%%%  scan idx 385  |  folio 372  %%%%%%%%%%
\SGAPrintPage{372}%
\(\underline{M}\) has a global partial resolution:
\[
  \underline{L}_i\longrightarrow \underline{L}_{i-1}\longrightarrow\ldots\longrightarrow \underline{L}_0\longrightarrow \underline{M}\longrightarrow 0,
\]
by \(\underline{B}\)-Modules locally free of finite type, the kernel \(\underline{Z}_i\) of \(\underline{L}_i\to \underline{L}_{i-1}\) is again a \(\underline{B}\)-Module of finite type (I. 2.11). It is even of finite presentation (by making the same argument with \(i+1\) instead of \(i\)). Finally, it is flat over \(S\). It therefore satisfies the same hypotheses as \(\underline{M}\), which allows us to construct step by step a global graded resolution of \(\underline{M}\) by graded \(\underline{B}\)-Modules locally free of finite type, by the initial remark. For \(i\geq r+1\), \(\underline{Z}_i\) is a projective \(\underline{B}\)-Module on every affine open set (1.5.2), and consequently is locally free as a graded \(\underline{B}\)-Module, by 1.5.1. We can then, in the graded resolution constructed, replace \(\underline{L}_{r+1}\) by \(\underline{Z}_r\), and the following terms by \(0\), whence the asserted resolution.

Let us finally prove point c) of the proof.

\sgaStableTarget{sga6.e06.prop.1-7}\textbf{Proposition 1.7.} \uline{Let} \(S\) \uline{be a quasi-compact scheme}, \(\underline{E}\) \uline{a} \(\underline{O}_S\)\uline{-Module locally free of finite type}, \(\underline{B}=\Sym(\underline{E})\), \uline{and} \(\underline{L}\) \uline{a graded} \(\underline{B}\)\uline{-Module}, \uline{locally free of finite type} (1.2). \uline{Then there exists a decreasing filtration of} \(\underline{L}\) \uline{by graded sub-}\(\underline{B}\)\uline{-Modules} \(\underline{L}^{(m)}\), \uline{such that}:
\[
\begin{array}{ll}
\text{i)}& \underline{L}^{(m)}=\underline{L}\quad\text{\uline{for} }m\leq m_0;\\[1mm]
\text{ii)}& \underline{L}^{(m)}=0\quad\text{\uline{for} }m\geq m_1;\\[1mm]
\text{iii)}& \text{\uline{for every} }m,\quad \operatorname{gr}^m(\underline{L})=\underline{L}^{(m)}/\underline{L}^{(m+1)}\\
&\quad \text{\uline{is isomorphic as a graded} }\underline{B}\text{\uline{-Module to a graded}}\\
&\quad \underline{B}\text{\uline{-Module of the form} }\underline{N}_m\otimes_{\underline{O}_S}\underline{B}(m),\\
&\quad \text{\uline{where} }\underline{N}_m\text{ \uline{is an} }\underline{O}_S\text{\uline{-Module locally free of finite type}.}
\end{array}
\]

Since \(S\) is quasi-compact, we can find integers \(n_0\) and \(n_1\) such that in a neighborhood of every point of \(S\), the elements of a homogeneous basis
%%%%%%%%%%  scan idx 386  |  folio 373  %%%%%%%%%%
\SGAPrintPage{373}%
of \(\underline{L}\) have degrees between \(n_0\) and \(n_1\). We then proceed by induction on \(N=n_1-n_0\).

When \(N=0\), i.e. \(n_1=n_0\), let \(n\) be the common value. The canonical homomorphism (1.2.1):
\[
  \underline{M}_n\otimes_{\underline{O}_S}\underline{B}(-n)\longrightarrow \underline{M}
\]
is then an isomorphism. Indeed, the verification is local, and by (1.2.3), we can then suppose that \(\underline{M}\simeq(\underline{B}(-n))^k\), and the result is trivial.

In the case \(N>0\), we again consider the homomorphism (1.2.1):
\[
  \underline{M}_{n_0}\otimes_{\underline{O}_S}\underline{B}(-n_0)\longrightarrow \underline{M}.
\]
We verify locally that this homomorphism is injective, and that its cokernel \(\underline{M}'\) is locally free of finite type, with local generators of degree between \(n_0+1\) and \(n_1\): indeed, in a neighborhood of every point of \(S\), \(\underline{M}_{n_0}\) is either zero or the nonzero homogeneous component of \(\underline{M}\) of lowest degree; using a decomposition of \(\underline{M}\) into direct summands of the form (1.2.2), we see in the latter case, using (1.2.3), that \(\underline{M}_{n_0}\otimes_{\underline{O}_S}\underline{B}(-n_0)\) is isomorphic to the direct summand \((\underline{B}(-n_0))^{k_{n_0}}\) of \(\underline{M}\), whence the asserted properties. We then apply the induction hypothesis to \(\underline{M}'\), whence the desired conclusion follows at once.

\sgaStableTarget{sga6.e06.sec.1.8}1.8. Returning to the proof of Theorem 1.1, let us show that \(K^{\bullet}(X)\) is generated as a \(K^{\bullet}(S)\)-module by the powers of \(\xi\).

If \(\underline{F}\) is a \(\underline{O}_X\)-Module locally free of finite type, there exists, by 1.4, a graded \(\underline{B}\)-Module of finite presentation \(\underline{M}\), flat over \(S\), such that \(\underline{F}=\Projz(\underline{M})\). By 1.6, \(\underline{M}\) admits a graded resolution by graded \(\underline{B}\)-Modules locally free of finite type. We deduce an exact sequence:
%%%%%%%%%%  scan idx 387  |  folio 374  %%%%%%%%%%
\SGAPrintPage{374}%
\[
  0\longrightarrow \Projz(\underline{L}_{r+1})\longrightarrow\ldots\longrightarrow
  \Projz(\underline{L}_0)\longrightarrow\Projz(\underline{M})\longrightarrow 0,
\]
and the following equality in \(K^{\bullet}(X)\):
\[
  \cl(\underline{F})=\sum_i(-1)^i\cl\bigl(\Projz(\underline{L}_i)\bigr).
\]
Using the filtration on the \(\underline{L}_i\) defined in 1.7, we find:
\[
  \cl(\Projz(\underline{L}_i))=\sum_j \cl\bigl(\Projz(\underline{N}_{ij}\otimes_{\underline{O}_S}\underline{B}(j))\bigr)
\]
Now \(\Projz(\underline{N}_{ij}\otimes_{\underline{O}_S}\underline{B}(j))=f^*(\underline{N}_{ij})(j)\); we therefore find:
\[
\begin{aligned}
  \cl(\Projz(\underline{L}_i))
  &=\sum_j \cl\bigl(f^*(\underline{N}_{ij})\bigr).\xi^j \\
  &=\sum_j f^*\bigl(\cl(\underline{N}_{ij})\bigr).\xi^j.
\end{aligned}
\]
We have thus finally expressed \(\cl(\underline{F})\) as a linear combination, with coefficients in \(K^{\bullet}(S)\), of the powers of \(\xi\).

\sgaStableTarget{sga6.e06.sec.1.9}1.9. To see that \(1,\xi,\ldots,\xi^{r+1}\) are linearly dependent over \(K^{\bullet}(S)\) we consider the canonical surjective homomorphism: \(f^*(\underline{E})\to\underline{O}_X(1)\), whose kernel \(\underline{F}\) is locally free of rank \(r\). We can therefore write:
\[
  \lambda^{r+1}\bigl(f^*(E)-\xi\bigr)=0,
  \qquad \text{with } E=\cl\underline{E},
\]
the operations \(\lambda^i\) on \(K^{\bullet}(X)\) being defined in V 2.2. We therefore obtain:
\[
  \sum_{i=0}^{r+1}\lambda^i(f^*(E))\lambda^{r+1-i}(-\xi)=0.
\]
Since \(\underline{O}_X(1)\) is an invertible sheaf, \(\lambda_t(\xi)=1+\xi t\), and consequently, \(\lambda_t(-\xi)=1/(1+\xi t)\), whence finally the relation:
\sgaStableTarget{sga6.e06.eq.1-9-1}\begin{equation}\tag{1.9.1}
  \sum_{i=0}^{r+1}(-1)^{r+1-i}\xi^{r+1-i}\lambda^i(E)=0.
\end{equation}

\sgaStableTarget{sga6.e06.sec.1.10}Let us finally show that \(1,\xi,\ldots,\xi^r\) are independent over \(K^{\bullet}(S)\)\SGAAuthorFootnote{*}{This proof is due to J.P. Jouanolou.}.

%%%%%%%%%%  scan idx 388  |  folio 375  %%%%%%%%%%
\SGAPrintPage{375}%
\sgaStableTarget{sga6.e06.lem.1-10}\textbf{Lemma 1.10.} \uline{With the notation and hypotheses of} 1.1, \uline{let} \(\underline{F}\) and \(\underline{G}\) \uline{be two locally free sheaves of finite type on} \(X\), \uline{such that}:
\[
  \cl_{K^{\bullet}(X)}(\underline{F})=\cl_{K^{\bullet}(X)}(\underline{G});
\]
\uline{then there exists an integer} \(n_0\), \uline{such that for} \(n\geq n_0\), \uline{we have}:
\[
\begin{array}{ll}
\text{i)}& f_*(\underline{F}(n)),\ f_*(\underline{G}(n))\text{ \uline{are locally free of finite type on} }S;\\[1mm]
\text{ii)}& \cl_{K^{\bullet}(S)}\bigl(f_*(\underline{F}(n))\bigr)=\cl_{K^{\bullet}(S)}\bigl(f_*(\underline{G}(n))\bigr).
\end{array}
\]

Let \(A\) be the free abelian group generated by the classes (up to isomorphism) of locally free sheaves of finite type on \(X\). To say that, in \(K^{\bullet}(X)\), \(\underline{F}\) and \(\underline{G}\) have the same class is to say, by the construction of \(K^{\bullet}\), that there exist exact sequences of \(\underline{O}_X\)-Modules locally free of finite type:
\[
  0\longrightarrow \underline{L}_i\longrightarrow \underline{M}_i\longrightarrow \underline{N}_i\longrightarrow 0,
  \qquad i=1,\ldots,k,
\]
and integers \(n_i\), such that in \(A\) we have the equality:
\[
  (\underline{F})-(\underline{G})=\sum_i n_i\bigl((\underline{L}_i)-(\underline{M}_i)+(\underline{N}_i)\bigr).
\]
For every \(n\), tensoring with \(\underline{O}_X(n)\) defines an automorphism of \(A\), whence:
\[
  (\underline{F}(n))-(\underline{G}(n))=\sum_i n_i\bigl((\underline{L}_i(n))-(\underline{M}_i(n))+(\underline{N}_i(n))\bigr),
\]
the sequences \(0\to \underline{L}_i(n)\to \underline{M}_i(n)\to \underline{N}_i(n)\to0\) remaining exact.

Let \(B\) be the free abelian group generated by the classes (up to isomorphism) of quasi-coherent sheaves on \(S\), and \(A'\subset B\) the subgroup generated by the classes of locally free sheaves of finite type on \(S\). Direct image defines a homomorphism \(A\to B\), and we therefore still have:
%%%%%%%%%%  scan idx 389  |  folio 376  %%%%%%%%%%
\SGAPrintPage{376}%
\[
\begin{aligned}
  \bigl(f_*(\underline{F}(n))\bigr)-\bigl(f_*(\underline{G}(n))\bigr)
  ={}&\sum_i n_i\bigl((f_*(\underline{L}_i(n)))-(f_*(\underline{M}_i(n)))\\
      &\hspace{35mm}+(f_*(\underline{N}_i(n)))\bigr).
\end{aligned}
\]

By 1.3, for sufficiently large \(n\), the following conditions will hold:
\[
  \forall k\geq1,\ \forall i,
  \quad R^k f_*(\underline{L}_i(n))=0;
\]
\[
  \forall i,
  \quad f_*(\underline{L}_i(n)),\ f_*(\underline{M}_i(n)),\ f_*(\underline{N}_i(n)),\ f_*(\underline{F}(n)),\ f_*(\underline{G}(n))
  \text{ are locally free of finite type.}
\]

It follows that for sufficiently large \(n\), these elements all belong to \(A'\) and satisfy the relation written above, and that the sequences:
\[
  0\longrightarrow f_*(\underline{L}_i(n))\longrightarrow f_*(\underline{M}_i(n))\longrightarrow f_*(\underline{N}_i(n))\longrightarrow0
\]
are exact. This implies that in \(K^{\bullet}(S)\), \(f_*(\underline{F}(n))\) and \(f_*(\underline{G}(n))\) have the same class.

1.11. Suppose given an equation of linear dependence over \(K^{\bullet}(S)\) among \(1,\xi,\ldots,\xi^r\):
\sgaStableTarget{sga6.e06.eq.1-11-1}\begin{equation}\tag{1.11.1}
  f^*(x_0)+f^*(x_1)\xi+\ldots+f^*(x_r)\xi^r=0.
\end{equation}

We can write each of the \(x_j\) as the difference of the classes of two sheaves, namely: \(x_j=\cl(\underline{E}_j)-\cl(\underline{E}'_j)\), where \(\underline{E}_j\) and \(\underline{E}'_j\) are locally free of finite type on \(S\). Put
\[
  \underline{F}=\bigoplus_{j=0}^r f^*(\underline{E}_j)(j),
  \qquad
  \underline{F}'=\bigoplus_{j=0}^r f^*(\underline{E}'_j)(j).
\]

Equality (1.11.1) then reads:
\[
  \cl_{K^{\bullet}(X)}(\underline{F})=\cl_{K^{\bullet}(X)}(\underline{F}').
\]

By Lemma 1.10, we obtain for sufficiently large \(n\):
\[
  \cl_{K^{\bullet}(S)}\bigl(f_*(\underline{F}(n))\bigr)=\cl_{K^{\bullet}(S)}\bigl(f_*(\underline{F}'(n))\bigr).
\]
%%%%%%%%%%  scan idx 390  |  folio 377  %%%%%%%%%%
\SGAPrintPage{377}%
Now
\[
\begin{aligned}
  f_*(\underline{F}(n))
  &=\bigoplus f_*\bigl(f^*(\underline{E}_j)\otimes_{\underline{O}_X}\underline{O}_X(n+j)\bigr)\\
  &=\bigoplus \underline{E}_j\otimes_{\underline{O}_S} f_*(\underline{O}_X(n+j)).
\end{aligned}
\]

Moreover, \(f_*(\underline{O}_X(n+j))=\Sym(\underline{E})^{n+j}\); we shall denote by \(\sigma_p\) the class in \(K^{\bullet}(S)\) of \(\Sym(\underline{E})^p\). Relation (1.11.1) therefore finally gives:
\sgaStableTarget{sga6.e06.eq.1-11-2}\begin{equation}\tag{1.11.2}
  \forall n\geq n_0,
  \qquad
  \sum_{j=0}^r x_j\sigma_{n+j}=0.
\end{equation}

We then consider the complex:
\[
  0\longrightarrow \bigwedge^{r+1}(\underline{E})\otimes_{\underline{O}_S}\Sym(\underline{E})
  \longrightarrow \bigwedge^r(\underline{E})\otimes_{\underline{O}_S}\Sym(\underline{E})
  \longrightarrow\ldots\longrightarrow
  \underline{E}\otimes_{\underline{O}_S}\Sym(\underline{E})
  \longrightarrow \Sym(\underline{E})\longrightarrow0,
\]
whose arrows are defined by:
\[
\begin{aligned}
&(x_1\wedge\ldots\wedge x_n)\otimes(y_1\otimes\ldots\otimes y_p)\rightsquigarrow{}\\
&\qquad \sum(-1)^{i-1}(x_1\wedge\ldots\wedge\widehat{x_i}\wedge\ldots\wedge x_n)
\otimes(x_i\otimes y_1\otimes\ldots\otimes y_p)
\end{aligned}
\]
If we regard \(\bigwedge^n(\underline{E})\) as a graded \(\underline{O}_S\)-Module concentrated in degree \(n\), this is a complex of graded \(\Sym(\underline{E})\)-Modules, the morphisms themselves being graded.

This complex is a resolution of \(\underline{O}_S\): on an open set where \(\underline{E}\) is free, \(\Sym(\underline{E})\) identifies with a polynomial Algebra \(\underline{O}_S[T_0,\ldots,T_r]\), and the complex with the Koszul complex defined by the \(\Sym(\underline{E})\)-regular sequence \(T_0,\ldots,T_r\); we know that the latter is a resolution of \(\Sym(\underline{E})/(T_0,\ldots,T_r)=\underline{O}_S\).

Writing that each homogeneous component of the complex of degree \(\geq1\) is acyclic, we therefore obtain relations in \(K^{\bullet}(S)\):
\sgaStableTarget{sga6.e06.eq.1-11-3}\begin{equation}\tag{1.11.3}
  \sum_{i=0}^{r+1}(-1)^i\lambda^i(E)\sigma_{p-i}=0
  \qquad\text{for every }p\geq1.
\end{equation}
(putting \(\sigma_k=0\) for \(k<0\)).

Suppose relation (1.11.2) holds for \(n\geq n_0\); we shall show that it still holds for \(n_0-1\), which will allow us to descend by induction to \(n_0=-r\) (we shall see later that this suffices to imply the vanishing of the \(x_j\)).

%%%%%%%%%%  scan idx 391  |  folio 378  %%%%%%%%%%
\SGAPrintPage{378}%
By (1.11.3), we can write
\begin{equation}\tag{1.11.4}
  \sum_{j=0}^r x_j\left(\sum_{i=0}^{r+1}(-1)^i\lambda^i(E)\sigma_{p(j)-i}\right)=0,
\end{equation}
with \(p(j)\geq1\). We take \(p(j)=n_0+r+j\), which is indeed \(\geq1\) if \(n_0>-r\). We then obtain:
\[
\begin{aligned}
\sum_{j=0}^r x_j\left(\sum_{i=0}^{r+1}(-1)^i\lambda^i(E)\sigma_{n_0+r+j-i}\right)
&=\sum_{i=0}^{r+1}(-1)^i\lambda^i(E)
   \left(\sum_{j=0}^r x_j\sigma_{n_0+r+j-i}\right).
\end{aligned}
\]

For \(i\neq r+1\), \(n_0+r-i\geq n_0\), so by (1.11.2):
\[
  \sum_{j=0}^r x_j\sigma_{n_0+r+j-i}=0.
\]

Consequently, relation (1.11.4) reduces to:
\[
  \lambda^{r+1}(E).\left(\sum_{j=0}^r x_j\sigma_{n_0-1+j}\right)=0.
\]

Now \(\underline{E}\) has rank \(r+1\), so \(\lambda^{r+1}(E)\) is the class of an invertible sheaf, and is therefore invertible in \(K^{\bullet}(S)\). Consequently, relation (1.11.4) reduces to relation (1.11.2) for \(n_0-1\).

Thus relation (1.11.2) holds for \(n\geq-r\); for \(n=-r\), it gives:
\[
  x_r\sigma_0=0.
\]
Since \(\sigma_0=1\), we obtain \(x_r=0\). We thus show by descending induction on \(j\) that the \(x_j\) are all zero, using relations (1.11.2) for \(n\) ranging from \(-r\) to \(0\). This shows that \(1,\ldots,\xi^r\) are independent over \(K^{\bullet}(S)\), and completes the proof of Theorem 1.1.

\sgaStableTarget{sga6.e06.rem.1-12}\textbf{Remark 1.12.} For every \(\underline{O}_S\)-Module locally free of finite type \(\underline{E}\), let \(\sigma_i(E)\) be the class in \(K^{\bullet}(S)\) of the \(i\)-th symmetric power of \(\underline{E}\); put \(\sigma_t(E)=1+\sum_{i\geq1}\sigma_i(E)t^i\); relations (1.11.3) can also be written:
%%%%%%%%%%  scan idx 392  |  folio 379  %%%%%%%%%%
\SGAPrintPage{379}%
\begin{equation}\tag{1.12.1}
  \lambda_t(E).\sigma_{-t}(E)=1.
\end{equation}

This relation itself shows the additivity of the map \(\sigma_t\), which can therefore be extended to a homomorphism from \(K^{\bullet}(S)\) to \(1+K^{\bullet}(S)[[t]]^+\):

\sgaStableTarget{sga6.e06.rem.1-13}\textbf{Remark 1.13.} We can take as a basis of \(K^{\bullet}(X)\) over \(K^{\bullet}(S)\) any sequence of powers of \(\xi\); \(\xi^k,\ldots,\xi^{k+r}\), since \(\xi\) is invertible, and such a sequence therefore consists of elements linearly independent over \(K^{\bullet}(S)\); moreover, the other powers of \(\xi\) can be expressed explicitly as linear combinations of these with coefficients in \(K^{\bullet}(S)\), by equality (1.9.1) and the invertibility of \(\xi\).

Put \(x=1-\xi^{-1}\) (\(x\) is interpreted as the ``class of a hyperplane section''). We can then take as a basis of \(K^{\bullet}(X)\) over \(K^{\bullet}(S)\) the elements \(1,x,\ldots,x^r\). To obtain the equation of dependence of \(1,\ldots,x^{r+1}\) we consider the exact sequence:
\[
  0\longrightarrow \underline{F}\longrightarrow f^*(\underline{E})\longrightarrow \underline{O}_X(1)\longrightarrow0,
\]
and the exact sequence obtained by passing to the duals:
\[
  0\longrightarrow \underline{O}_X(-1)\longrightarrow f^*(\check{\underline{E}})\longrightarrow \check{\underline{F}}\longrightarrow0.
\]
We then observe that \(\gamma^{r+1}(\check F-r)=0\), from which we deduce:
\[
  \gamma^{r+1}(\check F-r)=\gamma^{r+1}(\check E-r-1+x)=0,
\]
that is:
\sgaStableTarget{sga6.e06.eq.1-13-1}\begin{equation}\tag{1.13.1}
  \sum_{k=0}^{r+1}\gamma^{r+1-k}(\check E-r-1).x^k=0.
\end{equation}

We obtain an analogous result by introducing \(x'=\xi-1\).

Finally, note that if \(\underline{E}=(\underline{O}_S)^{r+1}\), and hence \(X=\mathbf P_S^r\), we have \(E=r+1\), and equality (1.13.1) reduces to:
%%%%%%%%%%  scan idx 393  |  folio 380  %%%%%%%%%%
\SGAPrintPage{380}%
\begin{equation}\tag{1.13.2}
  x^{r+1}=0.
\end{equation}

\textbf{Remark 1.14.} Using the notion of a ``relative scheme'' (cf [3]), we can still give a meaning to 1.1 by assuming that \(S\) is a locally ringed topos. If we assume that the final object of the topos is quasi-compact, the proof given in the case of schemes carries over without difficulty, and the theorem therefore remains valid.

\sgaStableTarget{sga6.e06.sec.2}\section*{2. Computation of the \(K^{\bullet}\) of a projective bundle: the case of perfect complexes}

The present section will not be used in the rest of the seminar.

In this section, the notation \(K^{\bullet}(X)\) will denote the Grothendieck group of perfect complexes on \(X\) (IV 2.2). We propose to prove the following theorem, analogous to 1.1:

\sgaStableTarget{sga6.e06.thm.2-1}\textbf{Theorem 2.1.} \uline{Let} \(S\) \uline{be a quasi-compact scheme}, \(\underline{E}\) \uline{a locally free} \(\underline{O}_S\)\uline{-Module of rank} \(r+1\), \(X=\mathbf P(\underline{E})\) \uline{the associated projective bundle}; \uline{let} \(\xi\) \uline{be the class} \uline{in} \(K^{\bullet}(X)\) \uline{of the complex equal to} \(\underline{O}_X(1)\) \uline{in degree zero}, \uline{and zero in the other degrees}. \uline{Then} \(K^{\bullet}(X)\), \uline{endowed with its natural} \(K^{\bullet}(S)\)-\uline{algebra structure, is} \uline{a free} \(K^{\bullet}(S)\)-\uline{module with basis} \(1,\xi,\ldots,\xi^r\).

\sgaStableTarget{sga6.e06.sec.2.2}2.2. Let us fix the notation: \(\underline{B}\) will denote the symmetric algebra of \(\underline{E}\); we shall write \(\mathcal A\) for the category of \(\underline{O}_X\)-Modules, \(\mathcal C\) for the category of graded \(\underline{B}\)-Modules with graded morphisms as morphisms, \(\mathcal C_0\) for the full subcategory of \(\mathcal C\) whose objects are the \(\underline{B}\)-Modules locally free of finite type as graded \(\underline{B}\)-Modules, and \(\mathcal C'\) for the category of \(\underline{O}_S\)-Modules. The categories \(\mathcal A\), \(\mathcal C\), and \(\mathcal C'\) are abelian.

For every \(n\), and every complex \(\underline{L}^\bullet\) in \(D^+(\mathcal A)\), we shall write \(\underline{L}^\bullet(n)\) for the complex \(\underline{L}^\bullet\otimes_{\underline{O}_X}\underline{O}_X(n)\).

%%%%%%%%%%  scan idx 394  |  folio 381  %%%%%%%%%%
\SGAPrintPage{381}%
\sgaStableTarget{sga6.e06.sec.2.3}2.3.
We shall need to apply the functor \(\Projz\) to arbitrary graded \(\underline{B}\)-Modules, and not only to quasi-coherent graded \(\underline{B}\)-Modules as in the construction of EGA II 3.2. We shall therefore indicate briefly how this construction can be carried out in the general case.

\sgaStableTarget{sga6.e06.para.2-3-1}2.3.1. First let \(S\) be a scheme, \(\underline{A}\) a quasi-coherent \(\underline{O}_S\)-algebra, \(Y=\Spec(\underline{A})\), and \(f:Y\to S\) the structural morphism. From the isomorphism:
\[
  \underline{A}\xrightarrow{\sim} f_*(\underline{O}_Y)
\]
we deduce a canonical homomorphism:
\sgaStableTarget{sga6.e06.eq.2-3-1-1}\begin{equation}\tag{2.3.1.1}
  f^{-1}(\underline{A})\longrightarrow \underline{O}_Y
\end{equation}
where \(f^{-1}\) denotes inverse image in the sense of sheaves of sets; in other words, we have a morphism of ringed spaces:
\sgaStableTarget{sga6.e06.eq.2-3-1-2}\begin{equation}\tag{2.3.1.2}
  \bar f:(Y,\underline{O}_Y)\longrightarrow (S,\underline{A})
\end{equation}
From it we deduce an ``inverse-image'' functor \(\bar f^{\,*}\) from the category of \(\underline{A}\)-Modules to the category of \(\underline{O}_Y\)-Modules, left adjoint to the ``direct-image'' functor \(\bar f_*\) from the category of \(\underline{O}_Y\)-Modules to the category of \(f_*(\underline{O}_Y)\)-Modules. We therefore have:
\sgaStableTarget{sga6.e06.eq.2-3-1-3}\begin{equation}\tag{2.3.1.3}
  \bar f^{\,*}(\underline{M})=f^{-1}(\underline{M})\otimes_{f^{-1}(\underline{A})}\underline{O}_Y
\end{equation}
for every \(\underline{A}\)-Module \(\underline{M}\).

Since the homomorphism (2.3.1.1) is flat (since on each stalk it corresponds to a localization), the functor \(\bar f^{\,*}\) thus defined is an exact functor (cf [4], p. 164).

If \(\underline{M}\) is a quasi-coherent \(\underline{A}\)-Module, this construction gives the Module \(\widetilde{\underline{M}}\) defined in EGA II 1.4.3. Indeed, the isomorphism:
%%%%%%%%%%  scan idx 395  |  folio 382  %%%%%%%%%%
\SGAPrintPage{382}%
\[
  \underline{M}\xrightarrow{\sim} f_*(\widetilde{\underline{M}})=\bar f_*(\widetilde{\underline{M}})
\]
gives, by adjunction, a homomorphism:
\[
  \bar f^{\,*}(\underline{M})\longrightarrow \widetilde{\underline{M}}
\]
and we check immediately on each stalk that it is an isomorphism.

\sgaStableTarget{sga6.e06.para.2-3-2}2.3.2. Now let \(\underline{B}\) be a quasi-coherent graded \(\underline{O}_S\)-Algebra, \(\underline{M}\) a graded \(\underline{B}\)-Module, and \(X=\Proj(\underline{B})\). For every affine open \(U\) of \(S\), and every section \(s\) of \(\underline{B}\) over \(U\), homogeneous of degree \(d>0\), we write \(\underline{M}_{(s),U}\) for the homogeneous component of degree \(0\) of the Module of fractions \(\underline{M}_s\); we therefore have:
\[
  \underline{M}_{(s),U}=
  \bigl((\underline{M}|U)\otimes_{(\underline{B}|U)}(\underline{B}|U)_s\bigr)_0 .
\]
If \(V\) is an affine open of \(U\), then \(\underline{M}_{(s),V}\) is the restriction to \(V\) of \(\underline{M}_{(s),U}\).

We consider the basis of open sets of \(X\) formed by the affine opens:
\[
  D_+(U,s)\simeq \Spec\bigl((\underline{B}|U)_{(s)}\bigr) .
\]
We write \(f\) for the structural morphism \(X\to S\), and \(f_{U,s}\) for the affine morphism:
\[
  f_{U,s}:D_+(U,s)\longrightarrow U .
\]
We then put:
\[
  \Projz(\underline{M})_{U,s}
  =
  \bar f_{U,s}^{\,*}\bigl(\underline{M}_{(s),U}\bigr),
\]
with the notation of 2.3.1. We shall show that the \(\Projz(\underline{M})_{U,s}\) glue as \(U\) and \(s\) vary, thereby defining an \(\underline{O}_X\)-Module denoted \(\Projz(\underline{M})\). The gluing for variable \(U\) and fixed \(s\) being evident, we are reduced (cf. EGA II 2.5.2) to showing that, if \(t\) is a section of \(\underline{B}\) over \(U\), homogeneous of degree \(e>0\), then the restriction of \(\Projz(\underline{M})_{U,s}\) to \(D_+(U,st)\) is canonically isomorphic to \(\Projz(\underline{M})_{U,st}\). We consider:
%%%%%%%%%%  scan idx 396  |  folio 383  %%%%%%%%%%
\SGAPrintPage{383}%
\[
\begin{tikzcd}[column sep=large,row sep=large]
D_+(st) \arrow[rr,hook,"i"] \arrow[dr,"f_{U,st}"'] & &
D_+(s) \arrow[dl,"f_{U,s}"]\\
& U &
\end{tikzcd}
\]
we have to show that the Modules \(i^*\circ\bar f_{U,s}^{\,*}\bigl(\underline{M}_{(s),U}\bigr)\) and \(\bar f_{U,st}^{\,*}\bigl(\underline{M}_{(st),U}\bigr)\)
are canonically isomorphic. On the one hand, we have the canonical isomorphisms:
\[
\begin{aligned}
i^*\circ\bar f_{U,s}^{\,*}\bigl(\underline{M}_{(s),U}\bigr)
&\simeq
i^*\Bigl(f_{U,s}^{-1}\bigl(\underline{M}_{(s),U}\bigr)
 \otimes_{f_{U,s}^{-1}((\underline{B}|U)_{(s)})}\underline{O}_{D_+(s)}\Bigr)\\
&\simeq
i^{-1}\Bigl(f_{U,s}^{-1}\bigl(\underline{M}_{(s),U}\bigr)
 \otimes_{f_{U,s}^{-1}((\underline{B}|U)_{(s)})}\underline{O}_{D_+(s)}\Bigr)
 \otimes_{i^{-1}(\underline{O}_{D_+(s)})}\underline{O}_{D_+(st)}\\
&\simeq
f_{U,st}^{-1}\bigl(\underline{M}_{(s),U}\bigr)
 \otimes_{f_{U,st}^{-1}((\underline{B}|U)_{(s)})}\underline{O}_{D_+(st)} .
\end{aligned}
\]
On the other hand:
\[
  \bar f_{U,st}^{\,*}\bigl(\underline{M}_{(st),U}\bigr)
  \simeq
  f_{U,st}^{-1}\bigl(\underline{M}_{(st),U}\bigr)
  \otimes_{f_{U,st}^{-1}((\underline{B}|U)_{(st)})}\underline{O}_{D_+(st)} .
\]
Finally, it follows from EGA II 2.2.2, translated into sheaf language, that we have the canonical isomorphisms:
\[
  \underline{M}_{(st),U}\simeq
  \bigl(\underline{M}_{(s),U}\bigr)_{t^d/s^e}
  \simeq
  \underline{M}_{(s),U}\otimes_{(\underline{B}|U)_{(s)}}(\underline{B}|U)_{(st)},
\]
whence, finally, the desired isomorphism.

It is clear that, if \(\underline{M}\) is quasi-coherent, the Module \(\Projz(\underline{M})\) thus constructed is indeed the one defined in EGA II.

\sgaStableTarget{sga6.e06.para.2-3-3}2.3.3. It follows from 2.3.1 and 2.3.2 that the functor \(\Projz\) is exact. Returning to the notation of 2.2, by passage to the derived categories we therefore deduce a functor still denoted \(\Projz\):
\[
  \Projz:D^+(\mathcal C)\longrightarrow D^+(\mathcal A).
\]

We also see from the preceding constructions that the functor
%%%%%%%%%%  scan idx 397  |  folio 384  %%%%%%%%%%
\SGAPrintPage{384}%
\(\Projz\) on the category of graded \(\underline{B}\)-Modules is compatible with tensor product; the same remains true of the functor obtained on the derived categories.

\sgaStableTarget{sga6.e06.sec.2.4}2.4. For every integer $k$, and every $\underline{O}_X$-Module $\underline{L}$, put:
\sgaStableTarget{sga6.e06.eq.2-4-1}\begin{equation}\tag{2.4.1}
  F_k(\underline{L})=\bigoplus_{n\geq k} f_*\bigl(\underline{L}(n)\bigr).
\end{equation}
The functor $F_k:\mathcal A\to\mathcal C$ thus defined is left exact. Since the category of $\underline{O}_X$-Modules has enough injectives, we can define a derived functor:
\sgaStableTarget{sga6.e06.eq.2-4-2}\begin{equation}\tag{2.4.2}
  \mathbb R F_k:D^+(\mathcal A)\longrightarrow D^+(\mathcal C).
\end{equation}

If $\underline{L}$ is a locally free $\underline{O}_X$-Module of finite type, there exists an integer $k_0$ such that, for $k\geq k_0$, $\underline{L}$ is $F_k$-acyclic. To see this, it suffices to show that, for $n$ sufficiently large, $\underline{L}$ is acyclic for the functors $\underline{L}\leadsto f_*\bigl(\underline{L}(n)\bigr)$. Since tensoring with $\underline{O}_X(n)$ is an exact functor, this follows immediately from the spectral sequence of composite functors and from 1.3.

2.5. We propose to construct, for every $\underline{O}_X$-Module $\underline{F}$, a homomorphism functorial in $\underline{F}$:
\sgaStableTarget{sga6.e06.eq.2-5-1}\begin{equation}\tag{2.5.1}
  \beta:\Projz\bigl(F_k(\underline{F})\bigr)\longrightarrow \underline{F}.
\end{equation}
For this, we define the restriction of $\beta$ to the $D_+(U,s)$ (2.3.2). By construction we have, on putting $\underline{M}=F_k(\underline{F})$:
\[
  \Projz(\underline{M})|D_+(U,s)
  =\bar f^{*}_{U,s}\bigl(\underline{M}_{(s),U}\bigr).
\]
To define a homomorphism:
\[
  \bar f^{*}_{U,s}\bigl(\underline{M}_{(s),U}\bigr)\longrightarrow \underline{F}|D_+(U,s),
\]
%%%%%%%%%%  scan idx 398  |  folio 385  %%%%%%%%%%
\SGAPrintPage{385}%
it suffices, by the definition of $\bar f^{*}_{U,s}$ as left adjoint to the direct-image functor, to define an $f_{U,s*}(\underline{O}_X)$-linear homomorphism
\[
  \underline{M}_{(s),U}\longrightarrow f_{U,s*}\bigl(\underline{F}|D_+(U,s)\bigr).
\]
For this, we must give, for every open $V$ of $U$, a homomorphism:
\begin{equation}\tag{2.5.2}
  \Gamma\bigl(V,\underline{M}\bigr)_{(s)}\longrightarrow \Gamma\bigl(D_+(V,s),\underline{F}\bigr)
\end{equation}
compatible with restrictions to the opens of $U$, and take the corresponding homomorphism on the associated sheaves. Homomorphism (2.5.2) is defined as follows: let $z/s^n\in\Gamma(V,\underline{M})_{(s)}$, where, $s$ being of degree $d$, $z\in\Gamma(V,\underline{M}_{nd})=\Gamma\bigl(D_+(V,s),\underline{F}(nd)\bigr)$; we then associate to $z/s^n$ the section $z.\alpha_{nd}(s^n)^{-1}$ of $\underline{F}$ over $D_+(V,s)$, writing $\alpha_d$ for the canonical homomorphism $\Gamma(V,\underline{B}_d)\longrightarrow \Gamma\bigl(f^{-1}(V),\underline{O}_X(d)\bigr)$ (cf. EGA II 2.6.3 and 2.6.4). We easily check the necessary compatibilities, whence, by gluing, homomorphism (2.5.1).

It is clear that, in the case where $\underline{F}$ is quasi-coherent, the homomorphism thus constructed is the one in EGA II (3.3.4.1).

The homomorphism of functors $\Projz\circ F_k\to \operatorname{Id}$ on the category of $\underline{O}_X$-Modules gives, by passage to the derived categories, a homomorphism of functors on $D^+(\mathcal A)$:
\sgaStableTarget{sga6.e06.eq.2-5-2}\begin{equation}\tag{2.5.3}
  \Projz\circ\mathbb R F_k\longrightarrow \operatorname{Id}_{D^+(\mathcal A)}.
\end{equation}

\sgaStableTarget{sga6.e06.prop.2-6}\textbf{Proposition 2.6.} \uline{Let} $S$ \uline{be a quasi}-\uline{compact scheme}, $\underline{E}$ \uline{a locally free} $\underline{O}_S$\uline{-Module of finite type}, $X=\mathbf P(\underline{E})$, \uline{and} $\underline{L}^{\bullet}$ \uline{a perfect complex on} $X$. \uline{Then} \uline{there exists an integer} $k_0$ \uline{such that, for} $k\geq k_0$, \uline{the following properties hold}:
%%%%%%%%%%  scan idx 399  |  folio 386  %%%%%%%%%%
\SGAPrintPage{386}%
\[
\begin{array}{ll}
\mathrm{i)}& \mathbb R F_k(\underline{L}^{\bullet})\text{ \uline{is} }\mathcal C_0
\text{\uline{-perfect}};\\[2mm]
\mathrm{ii)}& \text{\uline{The canonical morphism} }(2.5.3):
\Projz\circ\mathbb R F_k(\underline{L}^{\bullet})\longrightarrow \underline{L}^{\bullet}
\text{ \uline{is an isomorphism in} }D^+(\mathcal A).
\end{array}
\]

The notion of a $\mathcal C_0$-perfect complex considered here is local, that is, it is perfection relative to the category fibered over $S$ defined by $\mathcal C_0$. This notion of perfection does have a meaning, since $\mathcal C_0$ satisfies the axioms of (I 4.0): only stability under kernels of epimorphisms is not immediate; it follows from 1.5.1.

Since $S$ is quasi-compact, it suffices to check the two assertions of 2.6 locally on $S$, so that we can assume $S$ affine. The invertible sheaf $\underline{O}_X(1)$ is then ample, and global resolutions on $X$ are available (II 2.2.9). There therefore exists a strictly perfect complex $\underline{L}'{}^{\bullet}$ in $D_{\mathrm{parf}}(\mathcal A)$, and an isomorphism in $D_{\mathrm{parf}}(\mathcal A)$: $\underline{L}^{\bullet}\xleftarrow{\sim}\underline{L}'{}^{\bullet}$, which allows us to assume $\underline{L}^{\bullet}$ strictly perfect. We saw in 2.4 that, for $k$ sufficiently large, the $\underline{L}^i$, all but finitely many of which are zero, are $F_k$-acyclic; we consequently have:
\[
  \mathbb R F_k(\underline{L}^{\bullet})=F_k(\underline{L}^{\bullet})
\]
for $k$ sufficiently large. By 1.3 and 1.6, we can find $k_0$ such that, for $k\geq k_0$, the preceding condition is satisfied, and the $F_k(\underline{L}^i)$ are $\mathcal C_0$-perfect $\underline{B}$-Modules for every $i$. The complex $F_k(\underline{L}^{\bullet})$ is then $\mathcal C_0$-perfect, whence i).

For $k$ chosen as above, homomorphism (2.5.3) is then the homomorphism:
%%%%%%%%%%  scan idx 400  |  folio 387  %%%%%%%%%%
\SGAPrintPage{387}%
\[
  \Projz\bigl(F_k(\underline{L}^{\bullet})\bigr)\longrightarrow \underline{L}^{\bullet},
\]
and it is an isomorphism on each component, since $\underline{L}^{\bullet}$ was assumed strictly perfect, whence ii).

\sgaStableTarget{sga6.e06.sec.2.7}2.7. On $\mathcal C$, the functor which associates to a graded $\underline{B}$-Module its $n$-th component is an exact functor with values in $\mathcal C'$; it therefore defines a functor:
\sgaStableTarget{sga6.e06.eq.2-7-1}\begin{equation}\tag{2.7.1}
  \operatorname{Comp}_n:D^+(\mathcal C)\longrightarrow D^+(\mathcal C').
\end{equation}
For every $n$, we can consider the functor $\mathcal C'\to\mathcal C$ defined by tensoring with $\underline{B}(-n)$. Homomorphism (1.2.1) defines a morphism of functors on $\mathcal C$, whence a morphism for the corresponding functors between derived categories:
\sgaStableTarget{sga6.e06.eq.2-7-2}\begin{equation}\tag{2.7.2}
  \operatorname{Comp}_n\otimes_{\underline{O}_S} \underline{B}(-n)
  \longrightarrow \operatorname{Id}_{D^+(\mathcal C)}.
\end{equation}
Let $\underline{M}^{\bullet}$ be a $\mathcal C_0$-perfect complex. We shall show that there exists a sequence of distinguished triangles:
\[
\begin{array}{ccccccccccc}
& \underline{P}_1^{\bullet} &&&& \underline{P}_2^{\bullet} &&&& \underline{P}_k^{\bullet} & \\
\swarrow && \nwarrow && \swarrow && \nwarrow && \swarrow && \nwarrow \\
\underline{N}_1^{\bullet} & \longrightarrow & \underline{M}^{\bullet}, & \qquad &
\underline{N}_2^{\bullet} & \longrightarrow & \underline{P}_1^{\bullet}, & ...., &
\underline{N}_k^{\bullet} & \longrightarrow & \underline{P}_{k-1}^{\bullet},
\end{array}
\]
where $\underline{P}_k^{\bullet}$ is acyclic, and the $\underline{N}_i^{\bullet}$ are of the form \(\underline{K}_i^{\bullet}\otimes_{\underline{O}_S}\underline{B}(d)\), the $\underline{K}_i^{\bullet}$ being perfect complexes in the ordinary sense on $S$.

Since $\underline{M}^{\bullet}$ is $\mathcal C_0$-perfect, and $S$ is quasi-compact, there exists a finite covering of $S$, and an integer $d_0$ such that, on every open of this covering, $\underline{M}^{\bullet}$ is isomorphic to a bounded complex whose terms are locally free of finite type over $\underline{B}$, say $\underline{M}_U'{}^{\bullet}$, such that $\operatorname{Comp}_d(\underline{M}_U'{}^{\bullet})=0$ for $d<d_0$.

%%%%%%%%%%  scan idx 401  |  folio 388  %%%%%%%%%%
\SGAPrintPage{388}%
We then consider the canonical morphism (2.7.2):
\[
  \underline{N}_1^{\bullet}=\operatorname{Comp}_{d_0}(\underline{M}^{\bullet})\otimes_{\underline{O}_S}\underline{B}(-d_0)
  \longrightarrow \underline{M}^{\bullet};
\]
let $\underline{P}_1^{\bullet}$ be the cone of this morphism. It is again a $\mathcal C_0$-perfect complex: indeed, for every $n$, $\operatorname{Comp}_n(\underline{M}^{\bullet})$ is a perfect complex on $S$, since this holds when $\underline{M}^{\bullet}$ is bounded with terms locally free of finite type over $\underline{B}$; consequently, the complex $\operatorname{Comp}_n(\underline{M}^{\bullet})\otimes_{\underline{O}_S}\underline{B}(-d_0)$ is $\mathcal C_0$-perfect, and the same therefore holds for the cone (I 4.10).

Moreover, on every open $U$ of the covering we can write the commutative diagram:
\[
\begin{tikzcd}[column sep=large,row sep=large]
\operatorname{Comp}_{d_0}(\underline{M}^{\bullet})\otimes_{\underline{O}_S}\underline{B}(-d_0) \arrow[r] \arrow[d,dash,"\wr"] & \underline{M}^{\bullet} \arrow[d,dash,"\wr"]\\
\operatorname{Comp}_{d_0}(\underline{M}_U'{}^{\bullet})\otimes_{\underline{O}_S}\underline{B}(-d_0) \arrow[r] & \underline{M}_U'{}^{\bullet},
\end{tikzcd}
\]
which induces an isomorphism between the cones; the lower homomorphism is an isomorphism in degree $d_0$; therefore $\underline{P}_1^{\bullet}$ is isomorphic over $U$ to a bounded complex with terms locally free of finite type over $\underline{B}$, whose components $\operatorname{Comp}_d$ are zero for $d<d_0+1$. Consequently, $\underline{P}_1^{\bullet}$ has the same properties as $\underline{M}^{\bullet}$, with $d_0$ replaced by $d_0+1$. We can therefore iterate the process, and it suffices to see that, after finitely many operations, the resulting complex $\underline{P}_k^{\bullet}$ is acyclic. Here again the question is local, so it suffices to check this for $\underline{M}_U'{}^{\bullet}$. For each term $\underline{M}_U'{}^i$, the morphisms considered above are those defining the filtration of Proposition 1.7; this filtration is finite, so, starting from $\underline{M}_U'{}^{\bullet}$, we reach the zero complex after finitely many operations, that is, the corresponding complex $\underline{P}_k^{\bullet}$ is acyclic.

%%%%%%%%%%  scan idx 402  |  folio 389  %%%%%%%%%%
\SGAPrintPage{389}%
\sgaStableTarget{sga6.e06.sec.2.8}2.8.
Let $\underline{L}^{\bullet}$ be a perfect complex on $X$. There exists an integer $k$ such that $\underline{M}^{\bullet}=\mathbb R F_k(\underline{L}^{\bullet})$ is a $\mathcal C_0$-perfect complex, and $\underline{L}^{\bullet}$ is isomorphic to $\Projz\bigl(\mathbb R F_k(\underline{L}^{\bullet})\bigr)$. Applying 2.7 to the complex $\underline{M}^{\bullet}$, the sequence of distinguished triangles constructed gives, on applying the functor $\Projz$, a sequence of distinguished triangles in $D_{\mathrm{parf}}(\mathcal A)$ (it is clear that $\Projz$ transforms a $\mathcal C_0$-perfect complex into a perfect complex). Moreover, we have:
\[
  \Projz\bigl(\underline{K}^{\bullet}\otimes_{\underline{O}_S}\underline{B}(d)\bigr)
  = f^*(\underline{K}^{\bullet})(d)
\]
(note that $f^*=\mathbb L f^*$ by the flatness of $f$) for $\underline{K}^{\bullet}$ in $D^+(\mathcal C')$.

We therefore obtain in $K^{\bullet}(X)$ the equality:
\[
  \cl(\underline{L}^{\bullet})
  =\sum_i\cl\bigl(f^*(\underline{K}_i^{\bullet})\bigr).\xi^{d_i}
  =\sum_i f^*\bigl(\cl(\underline{K}_i^{\bullet})\bigr).\xi^{d_i},
\]
where the $\underline{K}_i^{\bullet}$ are the perfect complexes on $S$ defined in 2.7. It follows that $K^{\bullet}(X)$ is generated as a $K^{\bullet}(S)$-module by the powers of $\xi$.

To see that $1,\xi,\ldots,\xi^{r+1}$ are linearly dependent, we consider the homomorphism: $K^{\bullet}_{\naif}(X)\to K^{\bullet}_{\mathrm{parf}}(X)$ which associates to the class of a locally free sheaf of finite type $\underline{L}$ the class of the complex equal to $\underline{L}$ in degree $0$, and zero in the other degrees. Equation (1.9.1) then gives a relation of dependence in $K^{\bullet}_{\mathrm{parf}}(X)$.

To see that $1,\xi,\ldots,\xi^r$ are linearly independent, we can use the $K^{\bullet}(S)$-linear morphism $f_*:K^{\bullet}(X)\to K^{\bullet}(S)$ (IV 2.11.1). We have:
\[
  f_*(1)=1,
  \quad\text{since } f_*(\underline{O}_X)=\underline{O}_S,\text{ and }R^i(f_*(\underline{O}_X))=0\text{ for }i\geq1;
\]
\[
  f_*(\xi^n)=0\quad\text{for }-r\leq n<0,
  \quad\text{since }R^i\bigl(f_*(\underline{O}_X(n))\bigr)=0\text{ for }i\geq0.
\]
Given a relation of the form:
\[
  f^*(x_0)+f^*(x_1)\xi+\ldots+f^*(x_r)\xi^r=0,
\]
we multiply it by $\xi^{-r}$, and apply $f_*$; we then find $x_r=0$; we
%%%%%%%%%%  scan idx 403  |  folio 390  %%%%%%%%%%
\SGAPrintPage{390}%
% SGA_NATIVE_CONTINUATION_389_390
\providecommand{\Sym}{\operatorname{Sym}}
\providecommand{\GL}{\operatorname{GL}}
\providecommand{\Grass}{\operatorname{Grass}}
\providecommand{\gr}{\operatorname{gr}}
\providecommand{\ctop}{c_{\mathrm{top}}}
\providecommand{\Rep}{\operatorname{Rep}}
\providecommand{\RZ}{R_{\mathbb Z}}
\providecommand{\Dflag}{D}
\providecommand{\Gmn}{G}
complete the induction, noting that, if $0\leq s\leq r$, and if $x_j=0$ for $s<j\leq r$, we have:
\[
  \sum_{i=0}^{s}f^*(x_i)\xi^i=0,
\]
whence:
\[
  \sum_{i=0}^{s}f^*(x_i)\xi^{i+r-s}=0,
\]
whence $x_s=0$ by what precedes.

\sgaStableTarget{sga6.e06.sec.3}\section*{3. Consequences of the structure theorem for the \(K^{\bullet}\) of a projective bundle}

\sgaStableTarget{sga6.e06.prop.3-1}\textbf{Proposition 3.1.} \uline{Let} \(S\) \uline{be a quasi}-\uline{compact scheme}, \(\underline{E}_1,\ldots,\underline{E}_n\) \uline{locally free} \(\underline{O}_S\)\uline{-Modules of finite type}, \(\mathbf D(\underline{E}_i)\) \uline{the flag scheme of} \(\underline{E}_i\), \(X=\mathbf D(\underline{E}_1)\times_S\ldots\times_S \mathbf D(\underline{E}_n)\).
\uline{Then the canonical homomorphism}:
\[
  K^{\bullet}(S)\longrightarrow K^{\bullet}(X)
\]
\uline{is injective} (\(K^{\bullet}\) \uline{being defined}, \uline{either from locally free} \(\underline{O}_S\)\uline{-Modules of finite type}, \uline{or from perfect complexes}).

We check immediately, using the universal property of flag schemes, that \(X\) can be constructed as follows: we consider the sequence of \(S\)-schemes \(X_1,\ldots,X_n\), defined inductively by taking \(X_k\) to be the flag scheme of the inverse image of \(\underline{E}_k\) over \(X_{k-1}\); then \(X=X_n\). We then see at once that, to prove the proposition, it suffices to prove it in the case \(n=1\), the general case following by induction. We can therefore assume that \(X\) is the flag scheme of a locally free \(\underline{O}_S\)-Module of finite type \(\underline{E}\).

We know that \(X\) can then be constructed as follows: we put \(\underline{E}=\underline{F}_0\), \(S=X_0\), and define inductively schemes \(X_1,\ldots,X_n\) and
%%%%%%%%%%  scan idx 404  |  folio 391  %%%%%%%%%%
\SGAPrintPage{391}%
for every \(i\) a locally free \(\underline{O}_{X_i}\)-Module of finite type \(\underline{F}_i\) by putting:
\[
  X_k=\mathbf P(\underline{F}_{k-1});\qquad
  \underline{F}_k=\Ker\bigl(f_k^*(\underline{F}_{k-1})\longrightarrow \underline{O}_{X_k}(1)\bigr),
\]
where \(f_k\) is the structural morphism \(X_k\to X_{k-1}\).

We then see that it suffices to prove the proposition in the case where \(X\) is a projective bundle over \(S\), the case of a flag bundle following by induction. This case follows immediately from 1.1 (resp. 2.1).

\sgaStableTarget{sga6.e06.thm.3-2}\textbf{Theorem 3.2.} \uline{Let} \(S\) \uline{be a quasi-compact scheme}. \uline{The pre}-\(\lambda\)-\uline{ring} \(K^{\bullet}(S)\), \uline{defined from locally free} \(\underline{O}_S\)\uline{-Modules of finite type}, \uline{is a} \(\lambda\)-\uline{ring}~(V 2.4).

The homomorphism \(\lambda_t:K^{\bullet}(S)\longrightarrow 1+\widehat{K^{\bullet}(S)[[t]]}^{+}\) being additive, it suffices, to check that it is a \(\lambda\)-homomorphism, to make the verification for elements of the form \(\cl(\underline{E})\), where \(\underline{E}\) is a locally free \(\underline{O}_S\)-Module of finite type. We are therefore reduced to showing that, for every pair of elements \(E,F\) of \(K^{\bullet}(S)\), classes of locally free \(\underline{O}_S\)-Modules of finite type \(\underline{E},\underline{F}\), relations V 2.4.2 hold:
\[
  \lambda^i(EF)=P_i\bigl(\lambda^1(E),\ldots,\lambda^i(E);\lambda^1(F),\ldots,\lambda^i(F)\bigr);
\]
\[
  \lambda^j(\lambda^i(E))=Q_{i,j}\bigl(\lambda^1(E),\ldots,\lambda^{ij}(E)\bigr).
\]

Since \(S\) is a quasi-compact scheme, we can find a finite covering of \(S\) by disjoint opens such that, on each, \(\underline{E}\) and \(\underline{F}\) have constant rank. The \(K^{\bullet}\) of a finite direct sum of opens is the product of the \(K^{\bullet}\) of these opens. It therefore suffices to check relations V 2.4.2 in each of the \(K^{\bullet}\) of the opens of the covering. We can therefore assume \(\underline{E}\) and \(\underline{F}\) to have constant rank on \(S\).

%%%%%%%%%%  scan idx 405  |  folio 392  %%%%%%%%%%
\SGAPrintPage{392}%
Now let \(X=\mathbf D(\underline{E})\times_S \mathbf D(\underline{F})\). The \(\lambda\)-homomorphism \(K^{\bullet}(S)\to K^{\bullet}(X)\) is injective by 3.1. Consequently, to prove relations V 2.4.2, we can work in \(K^{\bullet}(X)\). Since in \(K^{\bullet}(X)\) the elements \(E\) and \(F\) decompose into sums of classes of invertible sheaves, it suffices, by additivity, to prove these same relations for elements \(\xi,\eta\), classes of invertible sheaves; they can also be written:
\[
  1+\xi\eta t=(1+\xi t)\circ(1+\eta t);
  \qquad
  \lambda_u(1+\xi t)=1+\{1+\xi t\}u,
\]
by putting them in the form V 2.4.1. They are then precisely relations V 2.3.1 and 2.3.3 defining the \(\lambda\)-ring structure of \(1+\widehat{K^{\bullet}(S)[[t]]}^{+}\).

Let us state the following more general result:

\sgaStableTarget{sga6.e06.thm.3-3}\textbf{Theorem 3.3.} \uline{Let} \((T,\underline{A})\) \uline{be a commutatively ringed topos}; \uline{write} \(K\) \uline{for the Grothendieck group of} \(\underline{A}\)-\uline{Modules locally free} of bounded rank on \(T\). \uline{Then}, \(K\), \uline{endowed with its canonical pre}-\(\lambda\)-\uline{structure}~(V 2.2 b)), \uline{is} \uline{a} \(\lambda\)-\uline{ring}.

To show that \(K\) is a \(\lambda\)-ring, we must prove relations V 2.4.1, and, by additivity, it suffices to prove them for elements \(E,F\) which are classes of locally free \(\underline{A}\)-Modules of bounded rank. Since the ranks of \(\underline{E}\) and \(\underline{F}\) are bounded, we can decompose \(T\) into a finite ``disjoint sum'' of opens over which \(\underline{E}\) and \(\underline{F}\) have constant rank. Using the commutation of \(K^{\bullet}\) with finite direct sums, and putting relations V 2.4.1 in the form V 2.4.2, we see that it suffices to prove relations V 2.4.2 in each of the \(K^{\bullet}\) of the opens considered, that is, it suffices to prove relations V 2.4.2 for elements \(E,F\), classes of \(\underline{A}\)-Modules \(\underline{E},\underline{F}\) locally free of constant rank.

%%%%%%%%%%  scan idx 406  |  folio 393  %%%%%%%%%%
\SGAPrintPage{393}%
Let \(\underline{E}\) be a locally free \(\underline{A}\)-Module of finite rank \(n\). There exists a torsor \(\underline{P}\) under \(\underline{\GL}(n,\underline{A})\) such that \(\underline{E}=\underline{P}\times_{\underline{\GL}(n,\underline{A})} (\underline{A})^n\), where \((\underline{A})^n\) is regarded as a \(\underline{\GL}(n,\underline{A})\)-Module by the standard representation. We can then define a \(\lambda\)-homomorphism from \(R_{\mathbb Z}(\mathbf{Gl}(n))\) to \(K\) by associating to every representation \(\rho\) of \(\mathbf{Gl}(n)\) in a free \(\mathbb Z\)-module of rank \(p\) the class of the locally free \(\underline{A}\)-Module
\(\underline{P}\times_{\underline{\GL}(n,\underline{A})} (\underline{A})^p\), where \((\underline{A})^p\) is regarded as a \(\underline{\GL}(n,\underline{A})\)-module by the representation \(\rho\). Under this homomorphism, the class of \(\underline{E}\) is therefore the image of the standard representation. To check a polynomial relation among the \(\lambda^i(E)\), it therefore suffices that this relation hold for the standard representation in \(R_{\mathbb Z}(\mathbf{Gl}(n))\).

We can argue similarly with two locally free \(\underline{A}\)-Modules \(\underline{E},\underline{F}\) of ranks \(n,p\); we can then find a torsor \(\underline{P}\) under \(\underline{\GL}(n,\underline{A})\times\underline{\GL}(p,\underline{A})\) such that \(\underline{E}=\underline{P}\times_{\underline{\GL}(n,\underline{A})\times\underline{\GL}(p,\underline{A})} (\underline{A})^n\), where \((\underline{A})^n\) is regarded as a \(\underline{\GL}(n,\underline{A})\times\underline{\GL}(p,\underline{A})\)-Module by letting \(\underline{\GL}(n,\underline{A})\) act by the standard representation and \(\underline{\GL}(p,\underline{A})\) by the trivial representation, and conversely for \(\underline{F}\). We then define a \(\lambda\)-homomorphism from \(R_{\mathbb Z}(\mathbf{Gl}(n)\times\mathbf{Gl}(p))\) to \(K\) in the same way as above, and are thus reduced to checking the relations among the \(\lambda^i(E)\), \(\lambda^j(F)\) for the corresponding elements of \(R_{\mathbb Z}(\mathbf{Gl}(n)\times\mathbf{Gl}(p))\).

We are therefore reduced to showing that, for every group scheme \(G\) over \(\Spec\mathbb Z\) which is a finite product of linear group schemes, the pre-\(\lambda\)-ring \(R_{\mathbb Z}(G)\) is a \(\lambda\)-ring. For this, we can refer to [5], where \(R_{\mathbb Z}(G)\) is computed explicitly. To avoid using this result, we can reason in a manner analogous to 3.2, using 1.14.

%%%%%%%%%%  scan idx 407  |  folio 394  %%%%%%%%%%
\SGAPrintPage{394}%
\sgaStableTarget{sga6.e06.sec.4}\section*{4. Computation of the \(K^{\bullet}\) of a flag bundle}

\sgaStableTarget{sga6.e06.sec.4.1}4.1.
Let \(S\) be a scheme, and let \(\underline{E}\) be an \(\underline{O}_S\)-Module locally free of rank \(n\). Let \(\pi=(p_1,\ldots,p_k)\) be a sequence of integers \(\geq1\) such that:
\[
  \sum_{i=1}^k p_i=n.
\]
Recall that the flag functor of type \(\pi\) of \(\underline{E}\) is the functor on the category of \(S\)-schemes, with values in the category of sets, which associates to an \(S\)-scheme \(X\xrightarrow{f} S\) the set of sequences of quotients of \(f^*(\underline{E})\):
\[
  (\underline{L}_{k-1},\ldots,\underline{L}_1),
\]
where, for every \(i\), \(\underline{L}_i\) is a quotient of \(\underline{L}_{i+1}\), locally free of rank \(p_1+\ldots+p_i\).

Such a functor is representable by an \(S\)-scheme \(\mathbf D_\pi(\underline{E})\), endowed with a canonical sequence of quotients of the inverse image of \(\underline{E}\), locally free of ranks \(p_1,p_1+p_2,\ldots\)

If \(\pi=(1,n-1)\), then \(\mathbf D_\pi(\underline{E})=\mathbf P(\underline{E})\).

If \(\pi=(r,n-r)\), then \(\mathbf D_\pi(\underline{E})=\mathbf G_{r,n}(\underline{E})\), the Grassmannian of type \((r,n-r)\) of \(\underline{E}\).

If all the \(p_i\) are equal to \(1\), one obtains the usual flag bundle of \(\underline{E}\), denoted by \(\mathbf D(\underline{E})\).

\sgaStableTarget{sga6.e06.lem.4-2}\textbf{Lemma 4.2.} \uline{Let} \(S\) \uline{be a scheme}, \(\underline{E}\) \uline{an} \(\underline{O}_S\)\uline{-Module locally free of rank} \(n\), \(\pi=(p_1,\ldots,p_i,\ldots,p_k)\), \uline{with} \(p_1+\ldots+p_k=n\); \uline{let} \(p',p''\) \uline{be two integers} \(\geq1\) \uline{such that} \(p'+p''=p_i\); \uline{put} \(\pi'=(p_1,\ldots,p_{i-1},p',p'',p_{i+1},\ldots,p_k)\), \(\pi''=(p',p'')\). \uline{Let} \(\underline{L}_1,\ldots,\underline{L}_{k-1}\) \uline{be the canonical quotients of the inverse image of} \(\underline{E}\) \uline{on} \(\mathbf D_\pi(\underline{E})\), \uline{and let} \(\underline{F}\) \uline{be the kernel of} \(\underline{L}_i\to \underline{L}_{i-1}\). \uline{Then there is a canonical} \(S\)-\uline{isomorphism}:
%%%%%%%%%%  scan idx 408  |  folio 395  %%%%%%%%%%
\SGAPrintPage{395}%
\[
  \mathbf D_{\pi'}(\underline{E})\simeq \mathbf D_{\pi''}(\underline{F}).
\]

The proof, immediate from the functorial definition of flag schemes, is left to the reader.

\sgaStableTarget{sga6.e06.lem.4-3}\textbf{Lemma 4.3.} \uline{Let} \(A\) \uline{be a commutative ring with unit}, \(n\) \uline{an integer} \(>0\), \((c_i)_{1\le i\le n}\) \uline{a family of elements of} \(A\); \uline{let} \(C\) \uline{be the} \(A\)-\uline{algebra generated by generators} \(\xi_i\), \(i=1,\ldots,n\), \uline{subject to the relations}:
\[
  \sigma_i(\xi_1,\ldots,\xi_n)=c_i,
  \qquad i=1,\ldots,n,
\]
\uline{where the} \(\sigma_i\) \uline{denote the elementary symmetric functions}. \uline{Let} \(p_1,\ldots,p_k\) \uline{be integers} \(\geq1\) \uline{with sum} \(n\), \uline{and for} \(1\le j\le k\),
\[
  \xi^{(j)}=(\xi_1^{(j)},\ldots,\xi_{p_j}^{(j)}),
\]
\uline{with}
\[
  \xi_i^{(j)}=\xi_{\sum_{\alpha<j}p_\alpha+i}.
\]
\uline{Put}:
\[
  c_i^{(j)}=\sigma_i(\xi_1^{(j)},\ldots,\xi_{p_j}^{(j)}),
  \qquad 1\le j\le k,
  \quad 1\le i\le p_j,
\]
\uline{and let} \(B\) \uline{be the subalgebra of} \(C\) \uline{generated by the} \(c_i^{(j)}\). \uline{Then}:
\begin{enumerate}[label=\textup{\roman*)}]
\item \uline{the elements of the form}:
\[
  (\xi^{(1)})^{\alpha(1)}\ldots (\xi^{(k)})^{\alpha(k)},
\]
\uline{where} \(\alpha(j)=(\alpha_1^{(j)},\ldots,\alpha_{p_j}^{(j)})\), and \(0\le \alpha_i^{(j)}\le p_j-i\), \uline{form a basis of} \(C\) \uline{considered as a} \(B\)-\uline{module}.
\item \uline{the generators} \(c_i^{(j)}\) \uline{of the} \(A\)-\uline{algebra} \(B\) \uline{satisfy the} \(n\) \uline{isobaric relations deduced from the relation}:
\[
  1+\sum_{i=1}^{n}c_i=\prod_{j=1}^{k}\left(1+\sum_{i=1}^{p_j}c_i^{(j)}\right)
\]
%%%%%%%%%%  scan idx 409  |  folio 396  %%%%%%%%%%
\SGAPrintPage{396}%
(\(c_i\) \uline{and} \(c_i^{(j)}\) \uline{being assigned weight} \(i\)), \uline{and these} \(n\) \uline{relations generate the ideal of relations among the} \(c_i^{(j)}\).
\end{enumerate}

\sgaStableTarget{sga6.e06.cor.4-4}\textbf{Corollary 4.4.} \uline{The elements of} \(C\) \uline{of the form}:
\[
  \xi_1^{\alpha_1}\ldots\xi_{n-1}^{\alpha_{n-1}},
  \qquad 0\le\alpha_i\le n-i,
\]
\uline{form a basis of} \(C\) \uline{considered as a module over} \(A\).

For the proof of the lemma and the corollary, we refer to [2], Exposé 4, pages 4.19, from which they are borrowed.

\sgaStableTarget{sga6.e06.sec.4.5}4.5.
Let \(S\) be a quasi-compact scheme, \(\underline{E}\) an \(\underline{O}_S\)-Module locally free of rank \(n\), and \(\pi=(p_1,\ldots,p_k)\) a sequence of integers \(\geq1\) with sum \(n\). Put \(X=\mathbf D_\pi(\underline{E})\); \(X\) is endowed with a canonical sequence of quotients \(\underline{L}_i\); put \(\underline{L}_0=0\), \(\underline{L}_k=f^*(\underline{E})\), where \(f\) is the structural morphism \(X\to S\). For \(1\le i\le k\), denote by \(\underline{F}_i\) the kernel of \(\underline{L}_i\to \underline{L}_{i-1}\). Let \(K^{\bullet}(X)\) be the Grothendieck group of locally free sheaves of finite type on \(X\) (resp. of perfect complexes). Denote by \(F_j\) the class of \(\underline{F}_j\), and by \(\lambda_i^{(j)}\) the class of \(\wedge^i(\underline{F}_j)\). In the case where \(K^{\bullet}(X)\) is defined in terms of perfect complexes, the series \(\lambda_t(E)\), \(\lambda_t(F_j)\) still have an evident meaning, where \(E\) is the class of \(\underline{E}\). Finally, note that:
\[
  \sum_{j=1}^{k}F_j=E.
\]

\sgaStableTarget{sga6.e06.thm.4-6}\textbf{Theorem 4.6.} \uline{With the hypotheses and notation of} 4.5, \uline{the} \(K^{\bullet}(S)\)-\uline{algebra} \(K^{\bullet}(X)\) \uline{is generated by the} \(\lambda_i^{(j)}\); \uline{the ideal of relations among the} \(\lambda_i^{(j)}\) \uline{is generated by the relations deduced from}:
\[
  \lambda_t(E)=\prod_{j=1}^{k}\lambda_t(F_j).
\]
%%%%%%%%%%  scan idx 410  |  folio 397  %%%%%%%%%%
\SGAPrintPage{397}%
In other words, \(K^{\bullet}(X)\) is the \(K^{\bullet}(S)\)-\(\lambda\)-algebra generated by the generators \(F_1,\ldots,F_k\) subject to the relations (cf. V 4.6):
\begin{enumerate}[label=\textup{\roman*)}]
\item \(\lambda^p(F_i)=0\) for \(p>p_i\);
\item \(\sum_i F_i=E\).
\end{enumerate}

We begin with the case where all the \(p_i\) are equal to \(1\); the theorem then reads:

\sgaStableTarget{sga6.e06.cor.4-7}\textbf{Corollary 4.7.} \uline{Let} \(S\) \uline{be a quasi}-\uline{compact scheme}, \(\underline{E}\) \uline{an} \(\underline{O}_S\)\uline{-Module locally free of rank} \(n\), \(X\) \uline{the flag bundle of} \(\underline{E}\), and \(\xi_1,\ldots,\xi_n\) \uline{the classes of the invertible sheaves} \(\underline{F}_1,\ldots,\underline{F}_n\) (\uline{notation of} 4.5). \uline{Then} \(K^{\bullet}(X)\) \uline{is the} \(K^{\bullet}(S)\)-\uline{algebra generated by the} \(\xi_i\), \uline{subject to the relations}:
\[
  \sigma_i(\xi_1,\ldots,\xi_n)=\lambda^i(E).
\]

We proceed by induction on \(n\). For \(n=1\), the theorem is trivial. Suppose it has been proved for \(n-1\); we then consider \(\mathbf P(\underline{E})\), and on \(\mathbf P(\underline{E})\) the exact sequence:
\[
  0\longrightarrow \underline{F}\longrightarrow g^*(\underline{E})\longrightarrow \underline{O}_{\mathbf P(\underline{E})}(1)\longrightarrow 0,
\]
where \(g\) is the structural morphism of \(\mathbf P(\underline{E})\). By Lemma 4.2, \(X\) is identified with the usual flag bundle \(\mathbf D(\underline{F})\) over \(\mathbf P(\underline{E})\), the canonical sheaves corresponding under this identification. The diagram:
\[
\begin{tikzcd}[column sep=large,row sep=large]
\mathbf D(\underline{E})=\mathbf D(\underline{F}) \arrow[r,"h"] \arrow[dr,"f"'] & \mathbf P(\underline{E}) \arrow[d,"g"]\\
& S,
\end{tikzcd}
\]
shows that \(K^{\bullet}(S)\to K^{\bullet}(X)\) factors through \(K^{\bullet}(\mathbf P(\underline{E}))\). By the induction hypothesis applied to \(\mathbf D(\underline{F})\), and Corollary 4.4, \(K^{\bullet}(X)\) has as a basis over \(K^{\bullet}(\mathbf P(\underline{E}))\) the elements of the form:
\[
  \xi_2^{\alpha_2}\ldots\xi_{n-1}^{\alpha_{n-1}},
  \qquad \text{with } 0\le\alpha_i\le n-i;
\]
moreover, \(K^{\bullet}(\mathbf P(\underline{E}))\) has as a basis over \(K^{\bullet}(S)\) the elements \(1,\xi_1,\ldots,\xi_1^{n-1}\); thus \(K^{\bullet}(X)\) has as a basis over \(K^{\bullet}(S)\) the elements:
%%%%%%%%%%  scan idx 411  |  folio 398  %%%%%%%%%%
\SGAPrintPage{398}%
\[
  \xi_1^{\alpha_1}\ldots\xi_{n-1}^{\alpha_{n-1}},
  \qquad
\text{with } 0\le\alpha_i\le n-i.
\]
Moreover, the additivity of \(\lambda_t\) implies that \(\prod_i(1+\xi_i t)=\lambda_t(E)\), and the relations of the corollary are satisfied.

If \(C\) is the \(K^{\bullet}(S)\)-algebra generated by generators \(\xi_i\) subject to the relations:
\[
  \sigma_i(\xi_1,\ldots,\xi_n)=\lambda^i(E),
\]
there is therefore a surjective homomorphism from \(C\) onto \(K^{\bullet}(X)\). By 4.4, this homomorphism sends a basis of \(C\) to a basis of \(K^{\bullet}(X)\), so it is an isomorphism, which completes the proof of the corollary.

Let us treat the general case of a flag bundle of type \(\pi\). The usual flag bundle \(\mathbf D(\underline{E})\) can be obtained from \(\mathbf D_\pi(\underline{E})\) by successively constructing the flag bundles of the \(\underline{F}_j\) (4.2). It follows that \(K^{\bullet}(\mathbf D(\underline{E}))\), considered as a module over \(K^{\bullet}(X)\), has a basis formed by the elements considered in 4.3.i), the \(\xi_i^{(j)}\) being the classes of the sheaves analogous to the \(\underline{F}_j\) occurring in the splitting of the inverse images of the \(\underline{F}_j\) on \(\mathbf D(\underline{E})\) (notation of 4.5 and 4.7). Thus \(K^{\bullet}(X)\) is a subring of \(K^{\bullet}(\mathbf D(\underline{E}))\); moreover, the additivity of \(\lambda_t\) implies that:
\[
  \lambda_i^{(j)}=\sigma_i(\xi_1^{(j)},\ldots,\xi_{p_j}^{(j)}),
\]
and \(K^{\bullet}(X)\) contains the sub-\(K^{\bullet}(S)\)-algebra of \(K^{\bullet}(\mathbf D(\underline{E}))\) generated by the \(\lambda_i^{(j)}\). Let \(B\) be this subalgebra; by 4.3.i), and what has just been said, there is in \(K^{\bullet}(\mathbf D(\underline{E}))\) a family of elements that forms a basis both for the \(B\)-module structure and for the \(K^{\bullet}(X)\)-module structure; consequently, \(K^{\bullet}(X)=B\), which shows that \(K^{\bullet}(X)\) is generated by the \(\lambda_i^{(j)}\). The ideal of relations among the \(\lambda_i^{(j)}\) is then given by 4.3.ii), taking
%%%%%%%%%%  scan idx 412  |  folio 399  %%%%%%%%%%
\SGAPrintPage{399}%
into account the structure of \(K^{\bullet}(\mathbf D(\underline{E}))\) over \(K^{\bullet}(S)\), determined by 4.7, and the fact that the \(\xi_i^{(j)}\) are the classes of the sheaves occurring in the splitting of the inverse image of \(\underline{E}\) on \(\mathbf D(\underline{E})\). This completes the proof of the first form of 4.6.

To deduce the second form, note that the elements \(F_i\) of the \(K^{\bullet}(S)\)-\(\lambda\)-algebra \(K^{\bullet}(X)\) satisfy relations i) and ii). If \(A\) is the \(K^{\bullet}(S)\)-\(\lambda\)-algebra generated by elements \(F_i'\) satisfying i) and ii), there is therefore a \(K^{\bullet}(S)\)-\(\lambda\)-homomorphism from \(A\) to \(K^{\bullet}(X)\) sending \(F_i'\) to \(F_i\) (V 4.7). By 4.6, this homomorphism is surjective. Since the \(\lambda\)-ideal generated by relations i) and ii) in \(K^{\bullet}(S)[F_1,\ldots,F_k]\) contains the usual ideal generated by the relations of 4.6, this homomorphism is also injective, hence is an isomorphism.

\sgaStableTarget{sga6.e06.cor.4-8}\textbf{Corollary 4.8.} \uline{With the hypotheses and notation of} 4.5, \uline{put}:
\[
  c_i^{(j)}=\gamma^i(F_j-p_j),
  \qquad 1\le i\le p_j;
\]
\uline{then the} \(K^{\bullet}(S)\)-\uline{algebra} \(K^{\bullet}(X)\) \uline{is generated by the} \(c_i^{(j)}\); \uline{the ideal of relations among the} \(c_i^{(j)}\) \uline{is generated by the relations deduced from}:
\[
  \gamma_t(E-n)=\prod_{j=1}^{k}\gamma_t(F_j-p_j).
\]

This is just another statement of 4.6. Indeed, the \(c_i^{(j)}\) can be expressed in terms of the \(\lambda_i^{(j)}\), and conversely. Moreover, the relation: \(\lambda_t(E)=\prod\lambda_t(F_j)\) is equivalent to the relation: \(\gamma_t(E-n)=\prod\gamma_t(F_j-p_j)\).

\sgaStableTarget{sga6.e06.sec.4.9}4.9.
Let \(S\) be a scheme; denote by \(\mathbf G_S^{n,r}\) the Grassmannian of type \((n,r-n)\) of \((\underline{O}_S)^r\), with \(r\ge n\). As \(r\) varies, the \(\mathbf G_S^{n,r}\) form an inductive system, the morphism \(\mathbf G_S^{n,r}\longrightarrow \mathbf G_S^{n,r+1}\) being defined by considering \((\underline{O}_{\mathbf G_S^{n,r}})^r\) as the
%%%%%%%%%%  scan idx 413  |  folio 400  %%%%%%%%%%
\SGAPrintPage{400}%
quotient of \((\underline{O}_{\mathbf G_S^{n,r}})^{r+1}\) by the \((r+1)\)-st factor. In this system of morphisms, the canonical quotients therefore correspond by inverse image. The inductive system thus defined will be denoted by \(\mathbf G_S^{n,\infty}\).

Its \(K^{\bullet}\) is defined by the formula:
\sgaStableTarget{sga6.e06.eq.4-9-1}\begin{equation}\tag{4.9.1}
  K^{\bullet}(\mathbf G_S^{n,\infty})
  = \varprojlim_r \bigl(K^{\bullet}(\mathbf G_S^{n,r})\bigr) .
\end{equation}

Let \(\underline{F}_1,\underline{F}_2\) be the canonical sheaves on \(\mathbf G_S^{n,r}\) (4.5), and \(F_1,F_2\) their class in \(K^{\bullet}(\mathbf G_S^{n,r})\). Put \(c_i=\gamma^i(F_1-n)\). As \(r\) varies, the \(c_i\) correspond under the transition morphisms, and therefore define elements, still denoted by \(c_i\), of \(K^{\bullet}(\mathbf G_S^{n,\infty})\).

\sgaStableTarget{sga6.e06.prop.4-10}\textbf{Proposition 4.10.} \uline{Let} \(S\) \uline{be a quasi}-\uline{compact scheme}. \uline{Then}, \uline{with the notation of} 4.9, \uline{there is an isomorphism}:
\[
  K^{\bullet}(\mathbf G_S^{n,\infty})\xrightarrow{\sim} K^{\bullet}(S)[[c_1,\ldots,c_n]] .
\]

By 4.8, \(K^{\bullet}(\mathbf G_S^{n,r})\) is the \(K^{\bullet}(S)\)-algebra generated by the \(c_i=\gamma^i(F_1-n)\), \(1\le i\le n\), and the \(\gamma^j(F_2-r+n)\), \(1\le j\le r-n\), subject to the relations deduced from:
\[
  \gamma_t(F_1-n).\gamma_t(F_2-r+n)
  = \gamma_t(r-r)=1,
\]
or again:
\[
  \gamma_t(F_2-r+n)=1/\gamma_t(F_1-n).
\]

This relation expresses the \(\gamma^j(F_2-r+n)\) in terms of the \(c_i\), and gives a series of isobaric relations among the \(c_i\) (\(c_i\) being assumed to have weight \(i\)) by writing that the coefficients of \(1/\gamma_t(F_1-n)\) are zero in degree \(>r-n\); moreover, these relations generate the ideal of relations among the \(c_i\). We may therefore write:
%%%%%%%%%%  scan idx 414  |  folio 401  %%%%%%%%%%
\SGAPrintPage{401}%
\[
  K^{\bullet}(\mathbf G_S^{n,r})=K^{\bullet}(S)[c_1,\ldots,c_n]/I_r,
\]
where \(I_r\) is a homogeneous ideal of \(K^{\bullet}(S)[c_1,\ldots,c_n]\), of degree \(>r-n\). It follows from definition (4.9.1) that \(K^{\bullet}(\mathbf G_S^{n,\infty})\) is the separated completion of \(K^{\bullet}(S)[c_1,\ldots,c_n]\) for the topology defined by the \(I_r\). This topology is finer than the topology defined by weight; we shall show that it is equivalent to it, which will imply the proposition.

Put \(A=K^{\bullet}(S)[c_1,\ldots,c_n]/I_r\), and let \(C\) be the \(A\)-algebra generated by generators \(x_1,\ldots,x_n\) subject to the relations:
\[
  \sigma_j(x_1,\ldots,x_n)=c_j,
\]
where the \(\sigma_j\) denote the elementary symmetric functions. By 4.4, the canonical homomorphism \(A\to C\) is injective. We obtain:
\[
  \gamma_t(F_1-n)
  =1+c_1t+\ldots+c_nt^n
  =\prod_{i=1}^n(1+x_it).
\]
Since \(1/\gamma_t(F_1-n)\) is a polynomial, so is \(1/(1+x_it)\) for every \(i\). Consequently, the \(x_i\) are nilpotent, and there is an integer \(N\) such that every polynomial of degree \(>N\) in the \(x_i\), with coefficients in \(K^{\bullet}(S)\), is zero. Thus every polynomial of \(A\), of weight \(>N\), is zero, i.e. every polynomial of \(K^{\bullet}(S)[c_1,\ldots,c_n]\) of weight \(>N\) belongs to \(I_r\). It follows that the two topologies under consideration are equivalent, which completes the proof.

\sgaStableTarget{sga6.e06.sec.5}\section*{5. Applications to projective bundles; study of \texorpdfstring{$f_*$}{f*}}

\sgaStableTarget{sga6.e06.sec.5.1}5.1.
In this section, $S$ is a quasi-compact scheme, $\underline{E}$ an $\underline{O}_S$-Module locally free of rank $r+1$, and $X$ the projective bundle associated to $\underline{E}$. On $X$, there is the canonical exact sequence:
\[
  0\longrightarrow \underline{F}\longrightarrow f^*(\underline{E})\longrightarrow \underline{O}_X(1)\longrightarrow 0,
\]
%%%%%%%%%%  scan idx 415  |  folio 402  %%%%%%%%%%
\SGAPrintPage{402}%
where $f$ is the structural morphism $X\longrightarrow S$.

The $K^{\bullet}$ considered will be those defined using locally free sheaves of finite type. We shall denote by $E$, $F$, $\xi$ the classes of $\underline{E}$, $\underline{F}$, $\underline{O}_X(1)$. Put $x=1-\xi^{-1}$, $x'=\xi-1=\xi x$, and denote by $\bar x$, $\bar x'$ the classes of $x$, $x'$ in $\Gr^1K^{\bullet}(X)$. Since $\xi-1$ and $x$ belong to $\Fil^1K^{\bullet}(X)$, we have $(\xi-1)x\in\Fil^2K^{\bullet}(X)$, i.e. $x\equiv \xi x=x'$ mod. $\Fil^2K^{\bullet}(X)$ i.e. $\bar x'= \bar x$

\sgaStableTarget{sga6.e06.sec.5.2}5.2.
If $\underline{G}$ is an $\underline{O}_X$-Module locally free of finite type, the $R^if_*(\underline{G})$ are in general not locally free of finite type on $S$, so that a homomorphism $f_*:K^{\bullet}(X)\to K^{\bullet}(S)$ cannot be defined as it is for the $K^{\bullet}$ defined using perfect complexes. Theorem 1.1 allows us to get around this difficulty as follows: if $y\in K^{\bullet}(X)$, there is a unique way to write:
\begin{equation}\tag{5.2.1}
  y=a_0+a_1\xi^{-1}+\ldots+a_r\xi^{-r},\qquad \text{with } a_i\in K^{\bullet}(S);
\end{equation}
we then put:
\begin{equation}\tag{5.2.2}
  f_*(y)=a_0 .
\end{equation}

Clearly the homomorphism thus defined is additive and satisfies the projection formula, i.e. is $K^{\bullet}(S)$-linear.

The following relations hold:
\begin{equation}\tag{5.2.3}
  f_*(1)=1;\qquad
  f_*(\xi^n)=0\quad\text{for }-r\le n<0;\qquad
  f_*(\xi^n)=\sigma_n,
  \quad 1\le n .
\end{equation}

To verify the last group of relations, we proceed by induction on $n$. Suppose that, for $p\le n-1$, $f_*(\xi^p)=\sigma_p$; relation (1.9.1), multiplied by a suitable power of $\xi$, gives:
\[
  \sum_{i=0}^{r+1}(-1)^i\lambda^i(E)\xi^{\,n-i}=0 .
\]
We can apply $f_*$ to it:
%%%%%%%%%%  scan idx 416  |  folio 403  %%%%%%%%%%
\SGAPrintPage{403}%
\[
  \sum_{i=0}^{r+1}(-1)^i\lambda^i(E)f_*(\xi^{\,n-i})=0 .
\]
The result follows immediately from relation (1.11.3) and the induction hypothesis.

\sgaStableTarget{sga6.e06.prop.5-3}\textbf{Proposition 5.3.} \uline{With the notation and hypotheses of} 5.1, \uline{for} \uline{every} $k\ge 0$ \uline{there is a direct-sum decomposition of} $K^{\bullet}(S)$-\uline{modules}:
\begin{equation}\notag
  \Fil^kK^{\bullet}(X)=\Fil^kK^{\bullet}(S)+\Fil^{k-1}K^{\bullet}(S).x+\ldots+\Fil^{k-r}K^{\bullet}(S).x^r,
\end{equation}
\uline{the filtration under consideration being the filtration of an augmented} $\lambda$-\uline{algebra defined} \uline{in} V 3.9.1 \uline{and} 3.10.

First note that $x$ has augmentation zero. Moreover:
\[
  \gamma_t(x)=\lambda_{t/(1-t)}(1-\xi^{-1})
  =\left(1+\frac{t}{1-t}\right)\left(1+\frac{\xi^{-1}t}{1-t}\right)^{-1}
  =1/(1-xt) .
\]
Thus $\gamma^i(x)=x^i$ for every $i$. Consequently, $K^{\bullet}(X)$ is a quotient of the $K^{\bullet}(S)$-$\lambda$-algebra generated by a generator $X$ subject to the relations $\gamma^i(X)=X^i$ for every $i$, and the ideal by which the quotient is taken is generated by relation (1.13.1). We can therefore apply V 4.12 and V 4.13.ii), and are reduced to showing that, for every $n\ge 0$:
\[
  x^n\in \Fil^nK^{\bullet}(S)+\Fil^{n-1}K^{\bullet}(S).x+\ldots+\Fil^{n-r}K^{\bullet}(S).x^r .
\]

This is evident for $n\le r$; we continue by induction. Relation (1.13.1) gives:
\[
  x^n=-\sum_{i=0}^{r}\gamma^{r+1-i}(\check E-r-1).x^{n-r-1+i},\qquad \text{with } n\ge r+1;
\]
applying the induction hypothesis to the $x^{n-r-1+i}$, we deduce the desired result, the sum being direct since $1,\ldots,x^r$ are independent over $K^{\bullet}(S)$.

%%%%%%%%%%  scan idx 417  |  folio 404  %%%%%%%%%%
\SGAPrintPage{404}%
\sgaStableTarget{sga6.e06.cor.5-4}\textbf{Corollary 5.4.} \uline{The filtration induced by} $K^{\bullet}(X)$ \uline{on} $K^{\bullet}(S)$ \uline{is identical to the filtration of} $K^{\bullet}(S)$.

\sgaStableTarget{sga6.e06.cor.5-5}\textbf{Corollary 5.5.} \uline{Let} $S$ \uline{be a quasi}-\uline{compact scheme}, $\underline{E}_1,\ldots,\underline{E}_n$ \uline{be} $\underline{O}_S$-\uline{Modules} \uline{locally free of finite type}, $\mathbf D(\underline{E}_i)$ \uline{the flag scheme of} $\underline{E}_i$, $X=\mathbf D(\underline{E}_1)\times_S\ldots\times_S \mathbf D(\underline{E}_n)$. \uline{Then the canonical homomorphism}:
\[
  \Gr K^{\bullet}(S)\longrightarrow \Gr K^{\bullet}(X)
\]
\uline{is injective}.

The argument of 3.1 reduces us to proving the same property when $X$ is a projective bundle over $S$. This case then follows immediately from 5.4.

\sgaStableTarget{sga6.e06.cor.5-6}\textbf{Corollary 5.6.} \uline{With the notation and hypotheses of} 5.1, $\Gr K^{\bullet}(X)$ \uline{is a free} $\Gr K^{\bullet}(S)$-\uline{module}, \uline{with basis} $1,\bar x,\ldots,\bar x^r$. \uline{For every} $k$, \uline{we have}:
\[
  \Gr^kK^{\bullet}(X)=\Gr^kK^{\bullet}(S)\oplus \Gr^{k-1}K^{\bullet}(S).\bar x\oplus\ldots\oplus \Gr^{k-r}K^{\bullet}(S).\bar x^r .
\]
\uline{Finally}, \uline{the elements} $1,\bar x,\ldots,\bar x^{r+1}$ \uline{are related by the relation}:
\[
  \sum_{i=0}^{r+1}c^{r+1-i}(\check E)\bar x^i=0,
\]
\uline{where} $c^k(\check E)$ \uline{is the} $k$-\uline{th Chern class of} $\check E$ (V 6.7).

An immediate consequence of 1.1, 5.3, and (1.13.1).

\sgaStableTarget{sga6.e06.cor.5-7}\textbf{Corollary 5.7.} \uline{With the notation and hypotheses of} 4.5, \uline{the} $\Gr K^{\bullet}(S)$-\uline{algebra} $\Gr K^{\bullet}(X)$ \uline{is generated by the} $\bar c_i^{(j)}$, \uline{where} $\bar c_i^{(j)}$ \uline{is the class in} $\Gr^iK^{\bullet}(X)$ \uline{of} $\gamma^i(F_j-p_j)$. \uline{The ideal of relations among the} $\bar c_i^{(j)}$ \uline{is generated by the relations deduced from}:
\[
  c(E-n)=\prod_{j=1}^{k}c(F_j-p_j),
\]
%%%%%%%%%%  scan idx 418  |  folio 405  %%%%%%%%%%
\SGAPrintPage{405}%
\uline{where} $c$ \uline{is the total Chern class} (V 6.7).

This result is proved by repeating the proof of 4.6 and using 5.6 (compare [2], Exp 4).

\sgaStableTarget{sga6.e06.cor.5-8}\textbf{Corollary 5.8.} \uline{With the notation and hypotheses of} 5.1, \uline{for} \uline{every} $k$ \uline{we have}:
\[
  f_*(\Fil^kK^{\bullet}(X))\subset \Fil^{k-r}K^{\bullet}(S).
\]

An immediate consequence of 5.3 and the projection formula.

\sgaStableTarget{sga6.e06.prop.5-9}\textbf{Proposition 5.9.} \uline{With the notation and hypotheses of} 5.1, \uline{for} \uline{every} $i\ge 1$ \uline{we have}:
\[
  f_*(\lambda^i(F))=0\qquad\bigl(\text{whence }f_*(\gamma^i(F))=0\bigr).
\]
\uline{In particular}, \uline{putting} $\lambda_{-1}(F)=\sum_{i=0}^{r}(-1)^i\lambda^i(F)$, \uline{we obtain}:
\[
  f_*(\lambda_{-1}(F))=1.
\]

We define a map
\[
  \bar f_*:1+K^{\bullet}(X)[[t]]^+\longrightarrow 1+K^{\bullet}(S)[[t]]^+
\]
by putting:
\[
  \bar f_*\left(\sum_{i\ge0}a_it^i\right)=\sum_{i\ge0}f_*(a_i)t^i .
\]
In general this is not a group homomorphism; nevertheless, the following relation holds for $\varphi\in 1+K^{\bullet}(X)[[t]]^+$, $\psi\in 1+K^{\bullet}(S)[[t]]^+$:
\[
  \bar f_*\bigl(\varphi.\hat G(f^*)(\psi)\bigr)=\psi.\bar f_*(\varphi).
\]
Indeed, it holds coefficient by coefficient and follows immediately from the projection formula.

From the relation $E=F+\xi$, we therefore deduce:
%%%%%%%%%%  scan idx 419  |  folio 406  %%%%%%%%%%
\SGAPrintPage{406}%
\[
  \bar f_*(\lambda_t(F))=
  \bar f_*(\lambda_t(E)\lambda_t(-\xi))=
  \lambda_t(E)\bar f_*(\lambda_t(-\xi));
\]
Now $\bar f_*(\lambda_t(-\xi))=\sum_{i\ge0}(-1)^if_*(\xi^i)t^i=\sigma_{-t}(E)$ by (5.2.3).
Using relation (1.12.1), we finally obtain:
\[
  \bar f_*(\lambda_t(F))=1;
\]
whence the desired relations

\sgaStableTarget{sga6.e06.prop.5-10}\textbf{Proposition 5.10.} \uline{With the notation and hypotheses of} 5.1, \uline{let} $z\in K^{\bullet}(X)$. \uline{For} $z(1-\xi)=0$ \uline{to hold, it is necessary and sufficient that there exist} $y\in K^{\bullet}(S)$ \uline{such that}:
\[
  \text{i)}\quad z=y.\lambda_{-1}(F);
  \qquad
  \text{ii)}\quad y.\lambda_{-1}(E)=0.
\]
From $E=F+\xi$, we deduce $\lambda_{-1}(E)=\lambda_{-1}(F).\lambda_{-1}(\xi)=\lambda_{-1}(F).(1-\xi)$. The condition is therefore sufficient.

Let $z$ be such that $z(1-\xi)=0$, that is, $zx'=0$. By Remark 1.13, there are unique $a_i\in K^{\bullet}(S)$ such that:
\[
  z=\sum_{i=0}^{r}a_i{x'}^i
\]
Moreover, the ${x'}^i$ satisfy a dependence equation:
\[
  \sum_{i=0}^{r}c_i{x'}^i+{x'}^{r+1}=0,
  \qquad \text{with } c_i=(-1)^{r+i+1}\gamma^{r+1-i}(E-r-1).
\]
Since $zx'=0$, we obtain:
\[
  a_r{x'}^{r+1}=-\sum_{i=1}^{r}a_{i-1}{x'}^i
  =-\sum_{i=0}^{r}a_rc_i{x'}^i .
\]
Using the independence of $1,\ldots,{x'}^r$, we find:
\[
  a_rc_0=0;
  \qquad \text{for } 0\le i<r,
  \quad a_i=a_rc_{i+1}.
\]
%%%%%%%%%%  scan idx 420  |  folio 407  %%%%%%%%%%
\SGAPrintPage{407}%
Thus:
\[
  z=a_r.\left(\sum_{i=0}^{r-1}c_{i+1}{x'}^i+{x'}^r\right).
\]
Now:
\[
\begin{aligned}
\sum_{i=0}^{r-1}c_{i+1}{x'}^i+{x'}^r
&=(-1)^r.\left(\sum_{i=0}^{r-1}\gamma^{r-i}(E-r-1).(-1)^i{x'}^i+(-1)^r{x'}^r\right)\\
&=(-1)^r.\left(\sum_{i=0}^{r}\gamma^{r-i}(E-r-1)\gamma^i(-x')\right)\\
&=(-1)^r\gamma^r(E-r-1-x')\\
&=(-1)^r\gamma^r(F-r)\\
&=(-1)^r\lambda^r(F-1)\\
&=\lambda_{-1}(F).
\end{aligned}
\]
Putting $a_r=y$, we therefore obtain:
\[
  z=y.\lambda_{-1}(F).
\]
Moreover, $y.c_0=0$; now $c_0=(-1)^{r+1}\gamma^{r+1}(E-r-1)=\lambda_{-1}(E)$; thus $y.\lambda_{-1}(E)=0$.

\sgaStableTarget{sga6.e06.prop.5-11}\textbf{Proposition 5.11.} \uline{With the notation and hypotheses of} 5.1, \uline{let} $z\in \Gr K^{\bullet}(X)$. \uline{For} $z.\bar x'=0$ \uline{to hold, it is necessary and sufficient that there exist} $y\in \Gr K^{\bullet}(S)$ \uline{such that}:
\[
  \text{i)}\quad z=y.c^r(F);
  \qquad
  \text{ii)}\quad y.c^{r+1}(E)=0.
\]
The proof is identical to that of 5.10, working in the associated graded rings and using 5.6.

\sgaStableTarget{sga6.e06.sec.6}\section*{6. Study of the filtration of \texorpdfstring{$K^{\bullet}(X)$}{K^{\bullet}(X)}, $X$ having an ample sheaf}

In this section, $X$ is a scheme having an ample sheaf $\underline{L}$. Recall that in this case the $K^{\bullet}$ defined in terms of locally free sheaves of finite type, or of perfect complexes, are equal (IV, 2.9). We denote by
%%%%%%%%%%  scan idx 421  |  folio 408  %%%%%%%%%%
\SGAPrintPage{408}%
$\varepsilon$ the augmentation homomorphism: $K^{\bullet}(X)\longrightarrow H^0(X,\underline{\mathbb Z})$.

\sgaStableTarget{sga6.e06.prop.6-1}\textbf{Proposition 6.1.} \uline{Let} $X$ \uline{be a quasi}-\uline{compact scheme having an ample sheaf} $\underline{L}$. \uline{Then} $\Fil^1K^{\bullet}(X)$ \uline{is a nil ideal}.

Since $\Fil^1K^{\bullet}(X)$ is generated by elements of the form $E-\varepsilon(E)$, where $E$ is the class of an $\underline{O}_X$-Module locally free of finite type $\underline{E}$, and $\varepsilon:K^{\bullet}(X)\to H^0(X,\underline{\mathbb Z})$ is the augmentation of $K^{\bullet}(X)$, it suffices to show that such an element is nilpotent. The quasi-compactness of $X$, and the commutation of $K^{\bullet}$ with finite direct sums, reduce us to the case where $\underline{E}$ has constant rank. Finally, passing to the flag bundle of $\underline{E}$, which gives an injection on the $K^{\bullet}$ (3.1) and preserves the hypothesis of the existence of an ample sheaf, we may suppose that $\underline{E}$ is invertible.

Let $\underline{E}^*$ be the dual of $\underline{E}$. For $n$ sufficiently large, $\underline{E}^*(n)=\underline{E}^*\otimes_{\underline{O}_X}\underline{L}^{\otimes n}$ is generated by its global sections; we can therefore find an integer $m$ and a surjective homomorphism: $(\underline{O}_X)^m\longrightarrow \underline{E}^*(n)$, hence (tensoring by $\underline{E}(-n)$) a surjective homomorphism: $(\underline{E}(-n))^m\longrightarrow\underline{O}_X$. The kernel of this homomorphism is locally free of rank $m-1$. Thus, denoting by $L$ the class of $\underline{L}$, we obtain:
\[
  \lambda^m(mEL^{-n}-1)=0 .
\]
Now $\lambda^m(mEL^{-n}-1)=(-1)^m\lambda_{-1}(mEL^{-n})=(-1)^m\lambda_{-1}(EL^{-n})^m=(-1)^m(1-EL^{-n})^m$. Consequently, for every invertible sheaf $\underline{E}$, the element $1-EL^{-n}$ is nilpotent for $n$ sufficiently large, or equivalently its product $L^n-E$ by $L^n$ is nilpotent.

We can then write:
\[
  E-1=(L^n-1)-(L^n-E)
\]
and, by what precedes, $L^n-E$ and $L^n-1$ are nilpotent on choosing $n$ sufficiently large. Thus $E-1$ is nilpotent.

%%%%%%%%%%  scan idx 422  |  folio 409  %%%%%%%%%%
\SGAPrintPage{409}%
\sgaStableTarget{sga6.e06.prop.6-2}\textbf{Corollary 6.2.} \uline{With the hypotheses of} 6.1, \uline{let} $R\subset K^{\bullet}(X)$ \uline{be a sub}-$\lambda$-\uline{ring generated by the classes of finitely many sheaves}. \uline{Then the} $\lambda$-\uline{ring filtration of} $R$ \uline{is discrete}, \uline{and} $R\cap\Fil^1K^{\bullet}(X)$ \uline{is nilpotent}. \uline{Moreover}, $R$ \uline{is a module of finite type over} $R\cap H^0(X,\underline{\mathbb Z})$ (V 3.9.1).

Since $X$ is quasi-compact, there is a finite covering of $X$ by disjoint open sets such that, on each open set of this covering, the sheaves whose classes generate $R$ all have constant rank. The standard argument then reduces us to the case where these sheaves have constant rank on $X$; $R$ is then augmented to $\mathbb Z$.

Let $\underline{E}_i$, $i=1,\ldots,k$, be the sheaves whose classes generate $R$, $E_i$ their classes, and $x_i=E_i-\varepsilon(E_i)$. We can write $R$ as a quotient of the $\lambda$-ring generated by the $x_i$ subject to the relations $\gamma^j(x_i)=0$ for $j>r_i=\varepsilon(E_i)$, that is, $\mathbb Z[x_i,\gamma^j(x_i)]_{\substack{1\le i\le k\\2\le j\le r_i}}$; consequently, $R$ is noetherian. But by Proposition 6.1, the $x_i$, $\gamma^j(x_i)$ are nilpotent. Thus the augmentation ideal of $R$ is nilpotent; moreover, it is precisely $R\cap\Fil^1K^{\bullet}(X)$.

Every element of $\Fil^n R$ can be written as a polynomial with no components of weight $<n$ in the $x_i$, $\gamma^j(x_i)$. Let $m$ be the largest integer such that there exists $i$ such that $\gamma^{m-1}(x_i)\ne0$; let $p$ be an integer such that the $p$-th powers of the $x_i$, $\gamma^j(x_i)$ are all zero. For $n>kp(1+\ldots+m)$, one sees at once that $\Fil^nR=0$; the filtration of $R$ is therefore discrete.

Since the filtration of $R$ is discrete, there are only finitely many nonzero monomials in the $x_i$, $\gamma^j(x_i)$. These monomials form a system of generators of $R$ as a $\mathbb Z$-module, which completes the proof.

%%%%%%%%%%  scan idx 423  |  folio 410  %%%%%%%%%%
\SGAPrintPage{410}%
\sgaStableTarget{sga6.e06.cor.6-3}\textbf{Corollary 6.3.} \uline{With the hypotheses of} 6.1, \uline{let} $A$ \uline{be a commutative ring}. \uline{For an element} $y$ \uline{of} $K^{\bullet}(X)\otimes_{\mathbb Z}A=K^{\bullet}(X)_A$ \uline{to be invertible}, \uline{it is necessary and sufficient that} $z=(\varepsilon\otimes\Id_A)(y)$ \uline{be invertible}. \uline{In particular}, \uline{if} $\underline{E}$ \uline{is a locally free sheaf of finite type on} $X$, \uline{everywhere} $\ne0$, \uline{then its class in} $K^{\bullet}(X)_{\mathbb Q}$ \uline{is invertible}.

The augmentation ideal of $K^{\bullet}(X)\otimes_{\mathbb Z}A$ is $\Fil^1K^{\bullet}(X)\otimes_{\mathbb Z}A$, and is therefore a nil ideal by 6.1. Thus, for an element of $K^{\bullet}(X)\otimes_{\mathbb Z}A$ to be invertible, it is necessary and sufficient that its augmentation be invertible.

\sgaStableTarget{sga6.e06.cor.6-4}\textbf{Corollary 6.4.} \uline{Let} $X$ \uline{be a quasi}-\uline{compact scheme}, \uline{having an ample sheaf}, \uline{and} $f:X'\to X$ \uline{a proper and perfect morphism}. \uline{Suppose in addition that one of the following two conditions is satisfied}:
\[
  \text{i) } f\text{ \uline{projective}},\qquad
  \text{ii) } X\text{ \uline{noetherian}},
\]
\uline{so that the} ``\uline{finiteness theorem}'' \uline{holds} (III 2.2). \uline{We are given a ring} $A$, \uline{and} \uline{consider the} ``\uline{direct image}'' \uline{homomorphism} (IV 2.11):
\[
  f_*:K^{\bullet}(X')\otimes_{\mathbb Z}A\longrightarrow K^{\bullet}(X)\otimes_{\mathbb Z}A.
\]
\uline{Finally suppose that there is an element} $x'\in K^{\bullet}(X')\otimes_{\mathbb Z}A$ \uline{such that} $\varepsilon(f_*(x'))\in H^0(X,\underline{A})$ \uline{is invertible}. \uline{Then the homomorphism}:
\[
  f^*:K^{\bullet}(X)\otimes_{\mathbb Z}A\longrightarrow K^{\bullet}(X')\otimes_{\mathbb Z}A
\]
\uline{is injective}.

We write the projection formula:
\[
  f_*(f^*(x)x')=x.f_*(x').
\]
Now $f_*(x')$ has invertible augmentation, so it is invertible (6.3). Consequently, the composite homomorphism:
%%%%%%%%%%  scan idx 424  |  folio 411  %%%%%%%%%%
\SGAPrintPage{411}%
\[
  x\leadsto f^*(x)\leadsto f^*(x)x'\leadsto f_*(f^*(x)x')
\]
is injective, so $f^*$ is also injective.

\sgaStableTarget{sga6.e06.sec.6.5}6.5. Let $X$ be a noetherian scheme of dimension $d$, having an ample sheaf $\underline{L}$. We endow $K^{\bullet}(X)$ with the filtration defined as follows (cf. [1]): $x$ belongs to $K^{\bullet}(X)_i$ if and only if, for every closed subset $Y$ of $X$, there is a perfect complex $\underline{M}^\bullet(Y)$ such that:
\[
  \text{i) } \cl(\underline{M}^\bullet(Y))=x;
\]
\[
  \text{ii) } \operatorname{Codim}(\Supp(H^*(\underline{M}^\bullet(Y)))\cap Y,Y)\ge i,
  \quad\text{with } H^*(\underline{M}^\bullet(Y))=\bigoplus_{i\in\mathbb Z}H^i(\underline{M}^\bullet(Y)).
\]

The $K^{\bullet}(X)_i$ thus defined are additive subgroups of $K^{\bullet}(X)$. Indeed, let $x,x'\in K^{\bullet}(X)_i$, and, for a given closed subset $Y$ in $X$, let $\underline{M}^\bullet$ and $\underline{M}'{}^\bullet$ be the corresponding complexes; then $x+x'$ is the class of the complex $\underline{M}^\bullet\oplus \underline{M}'{}^\bullet$, and
\[
\Supp(H^*(\underline{M}^\bullet\oplus \underline{M}'{}^\bullet))
 =\Supp(H^*(\underline{M}^\bullet)\oplus H^*(\underline{M}'{}^\bullet))
 =\Supp(H^*(\underline{M}^\bullet))\cup\Supp(H^*(\underline{M}'{}^\bullet));
\]
since the codimension of a finite union of closed subsets is the infimum of the codimensions of the closed subsets, we indeed have $x+x'\in K^{\bullet}(X)_i$.

\sgaStableTarget{sga6.e06.prop.6-6}\textbf{Proposition 6.6.} \uline{The filtration} 6.5 \uline{of} $K^{\bullet}(X)$ \uline{is a ring filtration} \uline{and} $K^{\bullet}(X)_{d+1}=0$.

Let $x\in K^{\bullet}(X)_i$, $x'\in K^{\bullet}(X)_j$, and let $Y$ be a closed subset of $X$. There is a complex $\underline{M}^\bullet$ such that $x=\cl(\underline{M}^\bullet)$, and $\codim(\Supp(H^*(\underline{M}^\bullet))\cap Y,Y)\ge i$; there is a complex $\underline{M}'{}^\bullet$ such that $x'=\cl(\underline{M}'{}^\bullet)$, and:
\[
\codim\bigl(\Supp(H^*(\underline{M}^\bullet))\cap\Supp(H^*(\underline{M}'{}^\bullet))\cap Y,
            \Supp(H^*(\underline{M}^\bullet))\cap Y\bigr)\ge j.
\]

We then have: $xx'=\cl(\underline{M}^\bullet\overset{\mathbb L}{\otimes}\underline{M}'{}^\bullet)$, and:
\[
\codim\bigl(\Supp(H^*(\underline{M}^\bullet))\cap\Supp(H^*(\underline{M}'{}^\bullet))\cap Y,Y\bigr)\ge i+j
\]
(EGA $0_{\mathrm{IV}}$ 14.2.2.3).
%%%%%%%%%%  scan idx 425  |  folio 412  %%%%%%%%%%
\SGAPrintPage{412}%
Thus, to complete the proof, it suffices to note that:
\[
\Supp\bigl(H^*(\underline{M}^\bullet\overset{\mathbb L}{\otimes}\underline{M}'{}^\bullet)\bigr)
 \subset \Supp(H^*(\underline{M}^\bullet))\cap\Supp(H^*(\underline{M}'{}^\bullet)).
\]

To show that $K^{\bullet}(X)_{d+1}$ is zero, we take $x\in K^{\bullet}(X)_{d+1}$, and $Y=X$ so that there is a complex $\underline{M}^\bullet$ such that $x=\cl(\underline{M}^\bullet)$, and $\codim_X(\Supp(H^*(\underline{M}^\bullet)))\ge d+1$; since $X$ has dimension $d$, this implies $\Supp(H^*(\underline{M}^\bullet))=\varnothing$, hence $\underline{M}^\bullet$ is acyclic, and $x=0$.

\sgaStableTarget{sga6.e06.lem.6-7}\textbf{Lemma 6.7.} \uline{Let} $X=\Spec(A)$ \uline{be an affine noetherian scheme}, $x_1,\ldots,x_k$ \uline{finitely many points of} $X$, and $\underline{E}$ \uline{an invertible} $\underline{O}_X$\uline{-Module}. \uline{Then} \uline{there is a section} $s\in\Gamma(X,\underline{E})$ \uline{such that, for every} $i$, $s_{x_i}$ \uline{is a basis of} $\underline{E}_{x_i}$.

Put $M=\Gamma(X,\underline{E})$. For every $i$, let $x'_i$ be a closed point of $X$ such that $x'_i\in\overline{\{x_i\}}$. The $x'_i$ correspond to maximal ideals $\underline{m}_i$ of $A$, which we may suppose all distinct, changing the index set so that each of the $\underline{m}_i$ occurs once and only once in the family $(\underline{m}_i)$. Since $\underline{E}$ is invertible, for every $i$ there is an isomorphism $M_{\underline{m}_i}\simeq A_{\underline{m}_i}$ from which we deduce an isomorphism $M/\underline{m}_iM\simeq A/\underline{m}_i$, then $\prod_i M/\underline{m}_iM\simeq\prod_i A/\underline{m}_i$, or again $A/\prod_i\underline{m}_i\simeq M/(\prod_i\underline{m}_i)M$. Lifting this last isomorphism, we obtain a homomorphism $A\to M$ such that the corresponding homomorphism $A_{\underline{m}_i}\to M_{\underline{m}_i}$ is, for every $i$, an isomorphism modulo $\underline{m}_i$, hence an isomorphism by Nakayama. The image of $1$ under this homomorphism therefore gives a section $s$ of $\underline{E}$ that is a basis of $\underline{E}_{x'_i}$ for every $i$. Since the set of points where $s$ is a basis of $\underline{E}$ is an open subset of $X$, $s$ is also a basis of $\underline{E}_{x_i}$ for $i=1,\ldots,k$.

\sgaStableTarget{sga6.e06.lem.6-8}\textbf{Lemma 6.8.} \uline{Let} $X$ \uline{be a noetherian scheme having an ample sheaf} $\underline{L}$, $\underline{E}$ \uline{an invertible} $\underline{O}_X$\uline{-Module}, and $x_1\ldots,x_k$ \uline{points of} $X$. \uline{Then there is an}
%%%%%%%%%%  scan idx 426  |  folio 413  %%%%%%%%%%
\SGAPrintPage{413}%
\uline{integer} $n$ \uline{such that, for every} $k\ge1$, \uline{there is a section} $s\in\Gamma(X,\underline{E}(kn))$ \uline{that is a basis of} $\underline{E}(kn)_{x_i}$ \uline{for every} $i$.

There is an integer $n_0$ and a section $f\in\Gamma(X,\underline{L}^{\otimes n_0})$ such that, for every $i$, $x_i\in X_f$, and such that $X_f$ is an affine open subset (EGA II 4.5.4). On $X_f$, there is a section $t\in\Gamma(X_f,\underline{E})$ satisfying the conditions of Lemma 6.7; moreover, there is an integer $n'_0$ such that the section $f^{\otimes n'_0}\otimes t\in\Gamma(X_f,\underline{E}(n_0n'_0))$ extends to the whole of $X$. Put $n=n_0n'_0$; then $f^{\otimes n'_0}\in\Gamma(X,\underline{L}^n)$; taking the tensor product of the section of $\underline{E}(n)$ obtained above and $f^{\otimes kn'_0}$ therefore gives a section of $\underline{E}((k+1)n)$ for every $k\ge0$. Since, for every $p$, $X_{f^{\otimes p}}=X_f$, and the section $t$ is a basis of $\underline{E}_{x_i}$ for every $i$, the sections of $\underline{E}(kn)$ thus constructed for every $k\ge1$ form a basis of $\underline{E}(kn)_{x_i}$ for every $i$.

\sgaStableTarget{sga6.e06.thm.6-9}\textbf{Theorem 6.9.} \uline{Let} $X$ \uline{be a noetherian scheme of dimension} $d$, \uline{having an ample sheaf} $\underline{L}$. \uline{Then} $\Fil^{d+1}K^{\bullet}(X)=0$, \uline{and a fortiori} $(\Fil^1K^{\bullet}(X))^{d+1}=0$.

(This last relation is due to J.P. Serre, in the general case considered here).

We shall first show that if $\underline{E},\underline{E}'$ are two invertible sheaves on $X$, with classes $E,E'$, then $E-E'\in K^{\bullet}(X)_1$.

Let $Y$ be a closed subset of $X$, $Y_1,\ldots,Y_k$ its irreducible components, and for every $i$, let $x_i$ be a point of $Y_i-\bigcup_{j\ne i}Y_j$. By Lemma 6.8, there are integers $n,n'$ such that, for any integers $k,k'\ge1$, there are homomorphisms from $\underline{O}_X$ to $\underline{E}(kn)$, $\underline{E}'(k'n')$ that are isomorphisms at the points $x_i$; we therefore obtain morphisms:
\[
  \underline{L}^{\otimes -kn}\xrightarrow{s_k} \underline{E},
  \qquad
  \underline{L}^{\otimes -k'n'}\xrightarrow{s'_{k'}} \underline{E}',
\]
%%%%%%%%%%  scan idx 427  |  folio 414  %%%%%%%%%%
\SGAPrintPage{414}%
that are isomorphisms at the points $x_i$. Now consider the complex:
\[
  \underline{N}^\bullet=\ldots\to0\to \underline{L}^{\otimes -nn'}
    \xrightarrow{(s'_n,0)} \underline{E}'\oplus \underline{L}^{\otimes -nn'}
    \xrightarrow{0+s_{n'}} \underline{E}\to0\to\ldots
\]
The class of $\underline{N}^\bullet$ is $E-E'$; moreover, $\underline{N}^\bullet$ is acyclic at the points $x_i$, so $\Supp(H^*(\underline{N}^\bullet))\cap Y$ contains no irreducible component of $Y$ and consequently has codimension $\ge1$ in $Y$; this shows that $E-E'$ belongs to $K^{\bullet}(X)_1$.

Now let $x$ be an element of $\Fil^{d+1}K^{\bullet}(X)$; we shall show that $x=0$. We may suppose that $x$ can be written:
\[
  x=\gamma^{i_1}(E_1-F_1)\ldots\gamma^{i_k}(E_k-F_k),
  \qquad\text{with } i_1+\ldots+i_k\ge d+1,
\]
the $E_i,F_i$ being the classes of locally free sheaves of finite type and, for every $i$, $\underline{E}_i$ and $\underline{F}_i$ having the same rank. Moreover, decomposing $X$ into a sum of finitely many open sets and arguing in each open set, we may suppose the $\underline{E}_i,\underline{F}_i$ to have constant rank. Let $X'$ be the fibre product over $X$ of the flag schemes of the $\underline{E}_i,\underline{F}_i$. In $K^{\bullet}(X')$, the images of the $E_i,F_i$ are sums of classes of invertible sheaves. Let $\xi,\eta$ be the classes of two invertible $\underline{O}_{X'}$-Modules; we have:
\[
  \gamma_t(\xi-\eta)=\gamma_t(\xi-1)/\gamma_t(\eta-1)
  =(1+(\xi-1)t)/(1+(\eta-1)t).
\]
Thus $\gamma^i(\xi-\eta)=(-1)^{i-1}(\xi-\eta)(\eta-1)^{i-1}$. Now $X'$ is noetherian and has an ample sheaf, so, by what precedes and 6.6, $\gamma^i(\xi-\eta)\in K^{\bullet}(X')_i$. Using the additivity of $\gamma_t$ and 6.6, we deduce that $f^*(\gamma^{i_j}(E_j-F_j))\in K^{\bullet}(X')_{i_j}$, and $f^*(x)\in K^{\bullet}(X')_{d+1}$, where $f$ is the canonical morphism $X'\to X$.

We shall deduce that $f^*(x)=0$. Let $r$ be the relative dimension of $X'$ over $X$. There is an element $x'$ of $K^{\bullet}(X')$, belonging to a basis of $K^{\bullet}(X')$ over $K^{\bullet}(X)$, such that:
%%%%%%%%%%  scan idx 428  |  folio 415  %%%%%%%%%%
\SGAPrintPage{415}%
\[
  x'=\prod_{i=1}^{r}(\xi_i-1),
\]
where the $\xi_i$ are the classes of invertible sheaves: to show this, the argument of 3.1 reduces us to the case of a projective bundle over $X$. This last case is an immediate consequence of 1.13. The element $x'$ under consideration belongs to $K^{\bullet}(X')_r$ by what precedes, and consequently $x'f^*(x)\in K^{\bullet}(X')_{d+r+1}$. Now $\dim X'=d+r$; thus, by 6.6, $x'f^*(x)=0$. Since $x'$ belongs to a basis of $K^{\bullet}(X')$ over $K^{\bullet}(X)$, this implies $f^*(x)=0$. Since $f^*$ is injective (3.1), we finally have $x=0$.

\sgaStableTarget{sga6.e06.problem.6-10}\textbf{Problem 6.10.} Under the hypotheses of 6.9, is it true that, for every $k$, we have: $\Fil^kK^{\bullet}(X)\subset K^{\bullet}(X)_k$?

\sgaStableTarget{sga6.e06.bibliography-for-expose-vi}\subsection*{BIBLIOGRAPHY}
\begin{enumerate}[label={\textup{[\arabic*]}}]
\item Bass H., K-Theory and stable algebra, Publ. Math. de l'IHES, n°22.
\item Chevalley C., Séminaire 1958 : Anneaux de Chow et applications.
\item Hakim M., Topos localement annelés et schémas relatifs, thesis (in preparation).
\item Hartshorne R., Residues and duality, Lecture Notes in Mathematics n°20, Springer-Verlag.
\item Serre J.-P., Le groupe de Grothendieck des schémas en groupes réductifs, Publications Mathématiques de l'IHES, n°34.
\end{enumerate}

% Native Expose VII continuation definitions, retained at their boundary.
\providecommand{\Kosz}{\mathrm K}
\providecommand{\Proj}{\operatorname{Proj}}
\providecommand{\Sym}{\operatorname{Sym}}
\providecommand{\gr}{\operatorname{gr}}
\providecommand{\Tor}{\operatorname{Tor}}
\providecommand{\Spec}{\operatorname{Spec}}
\providecommand{\Imf}{\operatorname{Im}}
\providecommand{\Ker}{\operatorname{Ker}}
\providecommand{\Id}{\operatorname{Id}}
\providecommand{\rad}{\operatorname{rad}}
\providecommand{\cO}{\mathcal O}
\providecommand{\cS}{\mathcal S}
\providecommand{\cJ}{\mathcal J}
\providecommand{\cN}{\mathcal N}
\providecommand{\bfP}{\mathbf P}
\providecommand{\bfA}{\mathbf A}

% Expose VII additions
\providecommand{\Kosz}{\mathrm K}
\providecommand{\Proj}{\operatorname{Proj}}
\providecommand{\Sym}{\operatorname{Sym}}
\providecommand{\gr}{\operatorname{gr}}
\providecommand{\Tor}{\operatorname{Tor}}
\providecommand{\Spec}{\operatorname{Spec}}
\providecommand{\Imf}{\operatorname{Im}}
\providecommand{\Ker}{\operatorname{Ker}}
\providecommand{\Id}{\operatorname{Id}}
\providecommand{\rad}{\operatorname{rad}}
\providecommand{\cO}{\mathcal O}
\providecommand{\cS}{\mathcal S}
\providecommand{\cJ}{\mathcal J}
\providecommand{\cN}{\mathcal N}
\providecommand{\bfP}{\mathbf P}
\providecommand{\bfA}{\mathbf A}

%%%%%%%%%%  scan idx 429  |  folio 416  %%%%%%%%%%
\SGAPrintPage{416}%
% Exposé VII: back to the volume's Sym and Proj
\renewcommand{\Sym}{\operatorname{Sym}}
\renewcommand{\Projz}{\operatorname{Proj}_0}
\sgaStableTarget{sga6.e07.expose}\section*{Exposé VII\\[3pt] REGULAR IMMERSIONS AND COMPUTATION OF \texorpdfstring{$K^{\bullet}$}{K0} OF A BLOWN-UP SCHEME}
\begin{center}
\uline{by Pierre BERTHELOT}
\end{center}

\sgaStableTarget{sga6.e07.sec.1}\section*{1. Generalities on regular immersions}

As the classical notion of regular immersion (EGA IV 16.9.2) proves inadequate in the non-noetherian case, we shall define a weaker notion, which nevertheless gives the desired properties and is equivalent to the usual definition in the noetherian case.

\sgaStableTarget{sga6.e07.def.1-1}\textbf{Definition 1.1.} \uline{Let} $A$ \uline{be a commutative ring}, $E$ \uline{a projective} $A$-\uline{module of finite type}, $u:E\longrightarrow A$ \uline{an} $A$-\uline{linear homomorphism}. \uline{We say that} $u$ \uline{is regular if the Koszul complex} $K_\bullet(u)$ \uline{defined by} $u$ \uline{is acyclic in dimension} $>0$, \uline{and hence is a resolution of} $B=A/u(E)$.

Recall that the Koszul complex associated with $u$ is defined as follows: the graded $A$-module $K_\bullet(u)$ is the exterior algebra $\Lambda(E)$, and the boundary operator is interior multiplication by $u$ regarded as an element of the dual $\check E$: it is an antiderivation of degree $-1$ of $\Lambda(E)$, equal to $u$ on $E$.

\sgaStableTarget{sga6.e07.lem.1-2}\textbf{Lemma 1.2.} \uline{Let} $u:E\longrightarrow A$ \uline{be a regular homomorphism}. \uline{Let} $A'$ \uline{be an} $A$-\uline{algebra}, $u':E'\longrightarrow A'$ \uline{the homomorphism deduced from} $u$ \uline{by base change}; \uline{suppose that} $u'$ \uline{is regular}. \uline{Put} $J=u(E)$, $J'=JA'=u(E')$. \uline{Then we have the following relations for every integer} $n\geq0$:
\sgaStableTarget{sga6.e07.eq.1-2-1}\sgaStableTarget{sga6.e07.eq.1-2-2}\sgaStableTarget{sga6.e07.eq.1-2-3}\sgaStableTarget{sga6.e07.eq.1-2-4}\begin{align}
 J^n/J^{n+1}\otimes_AA'&\xrightarrow{\sim}J'^n/J'^{n+1}, \tag{1.2.1}\\
 \Tor_i^A(J^n/J^{n+1},A')&=0\quad\text{\uline{for} }i\geq1; \tag{1.2.2}\\
 J^n\otimes_AA'&\xrightarrow{\sim}J'^n, \tag{1.2.3}\\
 \Tor_i^A(J^n,A')&=0\quad\text{\uline{for} }i\geq1. \tag{1.2.4}
\end{align}

a) First suppose that the $J^n/J^{n+1}$ are projective over $B=A/J$; since $J$ is of finite type, they are locally free of finite type over $B$.
%%%%%%%%%%  scan idx 430  |  folio 417  %%%%%%%%%%
\SGAPrintPage{417}%
The relations to be proved can be checked pointwise on $\Spec(A)$, so we are reduced to the case where $A$ is local, with $u(E)\subset\rad(A)$. The $J^n/J^{n+1}$ are then free over $B$, and the relations $\Tor_i^A(J^n/J^{n+1},A')=0$ for $i\geq1$ are equivalent to the relations $\Tor_i^A(B,A')=0$ for $i\geq1$. Now we can compute the $\Tor_i^A$ using the projective resolution $K_\bullet(u)$ of $B$; they are the homology groups of the complex $K_\bullet(u)\otimes_AA'=K_\bullet(u')$. Since $u'$ is assumed regular, this last complex is acyclic in degrees $>0$, whence the relations $\Tor_i^A(J^n/J^{n+1},A')=0$, for $i\geq1$.

We deduce the relations $\Tor_i^A(A/J^n,A')=0$ for $i\geq1$ by induction on $n$ using the exact Tor sequence associated with the exact sequence:
\[
 0\longrightarrow J^n/J^{n+1}\longrightarrow A/J^{n+1}\longrightarrow A/J^n\longrightarrow0.
\]
Since for every $k$ we have $J^kA'=J'^k$, this exact sequence and $\Tor_1^A(A/J^n,A')=0$ give the relations $J^n/J^{n+1}\otimes_AA'\xrightarrow{\sim}J'^n/J'^{n+1}$. From the exact sequence:
\[
 0\longrightarrow J^n\longrightarrow A\longrightarrow A/J^n\longrightarrow0,
\]
and $\Tor_1^A(A/J^n,A')=0$, we deduce the relations $J^n\otimes_AA'\simeq J'^n$. Finally, from the exact Tor sequence associated with the exact sequence:
\[
 0\longrightarrow J^{n+1}\longrightarrow J^n\longrightarrow J^n/J^{n+1}\longrightarrow0,
\]
we deduce the relations $\Tor_i^A(J^n,A')=0$ for $i\geq1$ by induction on $n$, the case $n=0$ being evident.

b) We shall now show that the $J^n/J^{n+1}$ are projective over $B=A/J$. Since the question is local on $\Spec(A)$, we may suppose $E$ free, say $E=(A)^d$, so that $u$ is defined by a sequence $f_1,\ldots,f_d$ of elements of $A$. Put $A_0=\mathbb Z[T_1,\ldots,T_d]$ and let $A_0\longrightarrow A$ be the homomorphism that sends
%%%%%%%%%%  scan idx 431  |  folio 418  %%%%%%%%%%
\SGAPrintPage{418}%
$T_i$ to $f_i$; denote by $u_0:E_0=(A_0)^d\longrightarrow A_0$ the homomorphism defined by the sequence $(T_1,\ldots,T_d)$, so that $u$ is deduced from $u_0$ by base change from $A_0$ to $A$. Since the sequence $(T_1,\ldots,T_d)$ is regular (in the classical sense), the homomorphism $u_0$ is regular, and the hypotheses of the lemma hold for the base change $A_0\longrightarrow A$. Now the $J_0^n/J_0^{n+1}$ are free of finite type over $\mathbb Z$, so we can apply the results of a). Consequently $J^n/J^{n+1}\xleftarrow{\sim}J_0^n/J_0^{n+1}\otimes_{A_0}A$, and $J^n/J^{n+1}$ is free of finite type.

\sgaStableTarget{sga6.e07.prop.1-3}\textbf{Proposition 1.3.} \uline{Let} $A$ \uline{be a ring}, $E$ \uline{a projective} $A$-\uline{module of finite type}, $u:E\longrightarrow A$ \uline{a regular homomorphism}, $J=u(E)$. \uline{Then the ideal} $J$ \uline{is quasi-regular} (EGA IV 16.9.7), \uline{that is} (EGA IV 16.9.8):
\begin{enumerate}[label=\roman*)]
\item $J$ \uline{is of finite type}.
\item $J/J^2$ \uline{is a projective} $A/J$-\uline{module of finite type};
\item \uline{the canonical homomorphism}
\end{enumerate}
\sgaStableTarget{sga6.e07.eq.1-3-1}\begin{equation}\tag{1.3.1}
 \Sym_{A/J}(J/J^2)\longrightarrow \gr_J(A)=\bigoplus_{n\geq0}J^n/J^{n+1},
\end{equation}
\uline{is an isomorphism}.

We saw in the proof of 1.2 b) that $J/J^2$ is projective of finite type. Resume the notation of 1.2 b) (since the question is local on $\Spec(A)$, we again suppose $E$ free of finite type). Then the canonical homomorphism (1.3.1) is deduced from the analogous homomorphism for $J_0$ by the base change $A_0\longrightarrow A$, by the relations 1.2. The latter is an isomorphism, so the same is true of (1.3.1).

\sgaStableTarget{sga6.e07.def.1-4}\textbf{Definition 1.4.} \uline{Let} $(X,\underline{A})$ \uline{be a commutatively locally ringed site}, $\underline{E}$ \uline{an} $\underline{A}$-\uline{Module locally free of finite type}, $u:\underline{E}\longrightarrow \underline{A}$ \uline{an} $\underline{A}$-\uline{linear homomorphism}. \uline{We say that} $u$ \uline{is} regular \uline{if the Koszul complex} $\underline{K}_\bullet(u)$ \uline{defined by} $u$ \uline{is acyclic in degrees} $>0$. \uline{Let} $\underline{J}$ \uline{be an Ideal of} $\underline{A}$; \uline{we say that} $\underline{J}$ \uline{is}
%%%%%%%%%%  scan idx 432  |  folio 419  %%%%%%%%%%
\SGAPrintPage{419}%
\uline{regular}, \uline{if}, \uline{locally on} $X$, \uline{there exist a free} $\underline{A}$-\uline{Module of finite type} $\underline{E}$, \uline{and a surjective regular homomorphism}: $\underline{E}\longrightarrow \underline{J}$.

\uline{Let} $i:Y\longrightarrow X$ \uline{be an immersion of preschemes}, \uline{and let} $U$ \uline{be an open subset of} $X$ \uline{such that} $i(Y)\subset U$ \uline{and such that} $i$ \uline{is a closed immersion of} $Y$ \uline{into} $U$. \uline{We say that} $i$ \uline{is regular if the ideal of} $\underline{O}_U$ \uline{defining} $i(Y)$ \uline{is regular} (\uline{a condition independent of} $U$). \uline{If} $y\in Y$, \uline{we say that} $i$ \uline{is regular at} $y$ \uline{if there exists a neighborhood} $V$ \uline{of} $i(y)$ \uline{such that the immersion} $i^{-1}(V)\longrightarrow V$ \uline{is regular}. \uline{For an immersion to be regular}, \uline{it is therefore necessary and sufficient that it be regular in a neighborhood of each point}.

\uline{This definition conflicts with} EGA IV 16.9.2, \uline{and replaces that definition}. However, in the noetherian case the two definitions are equivalent (EGA IV 19.5.1). An immersion that can locally be defined by a regular sequence is regular; a regular immersion is quasi-regular (1.3).

1.4.1. Let $X$ be a scheme, $\underline{J}$ an Ideal of finite type of $\underline{O}_X$, $x$ a point of $\Supp(\underline{O}_X/\underline{J})$. We shall say that a family of sections $f_1,\ldots,f_d$ of $\underline{J}$ on a neighborhood of $x$ forms a \uline{minimum system of generators of} $\underline{J}$ \uline{in a neighborhood of} $x$ if the $(f_i)_x$ form a family of generators of $\underline{J}_x$ of minimum cardinality, i.e. their images in $\underline{J}_x\otimes k(x)$ form a basis of this vector space over $k(x)$. Note that the $f_i$ generate $\underline{J}$ on a neighborhood of $x$.

1.4.2. To say that an immersion is regular is also to say that it is defined in some open subset of $X$ by an ideal of finite type $\underline{J}$, and that for every point $x\in\Supp(\underline{O}_X/\underline{J})$ and every minimum system of generators of $\underline{J}$ in a neighborhood of $x$, the Koszul complex associated with this sequence is acyclic in a neighborhood of $x$: indeed, let $f_1,\ldots,f_d$ be a minimum system of
%%%%%%%%%%  scan idx 433  |  folio 420  %%%%%%%%%%
\SGAPrintPage{420}%
generators of $\underline{J}$ in a neighborhood of $x$ such that the associated Koszul complex is acyclic (one exists when the immersion is assumed regular); let $g_1,\ldots,g_d$ be another minimum system of generators of $\underline{J}$; there exists a neighborhood of $x$ on which $\underline{J}$ can be generated either by the $f_i$ or by the $g_j$, and consequently there exists on this neighborhood a matrix with coefficients in $\underline{O}_X$ expressing the $g_j$ in terms of the $f_i$. Now the images of the $f_i$ (resp. $g_j$) in $\underline{J}/\underline{J}^2$ form a basis at the point $x$, hence in a neighborhood of $x$, and consequently this matrix is invertible in a neighborhood of $x$. It therefore defines an isomorphism from the Koszul complex defined by the $g_j$ onto the Koszul complex defined by the $f_i$, which shows that the Koszul complex defined by the $g_j$ is acyclic.

1.4.3. Assuming only that $\underline{J}$ is of finite type and $\underline{J}/\underline{J}^2$ is locally free of finite type, note that every minimum system of generators of $\underline{J}$ in a neighborhood of a point $y$ of $Y$ gives a basis of $\underline{J}/\underline{J}^2$ in a neighborhood of $y$. Conversely, given a basis of $\underline{J}/\underline{J}^2$ in a neighborhood of $y$, we can lift it in a neighborhood of $y$ to a family of sections of $\underline{J}$, which generate $\underline{J}_y$ by Nakayama, hence $\underline{J}$ in a neighborhood of $y$ since $\underline{J}$ is of finite type; moreover these generators have minimum cardinality, since they give a basis of $\underline{J}/\underline{J}^2$.

\sgaStableTarget{sga6.e07.prop.1-5}\textbf{Proposition 1.5.} (\uline{Compare} EGA IV 19.1.5 (ii)). \uline{Let} $i:Y\longrightarrow X$ \uline{be an immersion}. \uline{If} $i$ \uline{is regular}, \uline{the same is true of every immersion deduced from} $i$ \uline{by flat base change}. \uline{For} $i$ \uline{to be regular}, \uline{it is necessary and sufficient that it be so locally for the faithfully flat quasi-compact topology}.

Stability under flat base change is evident.

Conversely, let $i'$ be the immersion deduced from $i$ by a
%%%%%%%%%%  scan idx 434  |  folio 421  %%%%%%%%%%
\SGAPrintPage{421}%
faithfully flat quasi-compact base change. If $i'$ is regular, then the Ideal $\underline{J}$ defining $i$ is of finite type, and $\underline{J}/\underline{J}^2$ is locally free of finite type, since these properties descend under faithfully flat quasi-compact morphisms. If $f_1,\ldots,f_d$ is a sequence of sections of $\underline{J}$ lifting a basis of $\underline{J}/\underline{J}^2$, the Koszul complex associated with the inverse images of the $f_i$ is acyclic in degrees $>0$ by 1.4; since it is the inverse image of the Koszul complex associated with the $f_i$, the latter is acyclic in degrees $>0$, and $i$ is regular.

\sgaStableTarget{sga6.e07.def.1-6}\textbf{Definition 1.6.} \uline{Let} $i:Y\longrightarrow X$ \uline{be an immersion defined in an open subset} $U$ \uline{of} $X$ \uline{by an Ideal} $\underline{J}$ \uline{of} $\underline{O}_X$; \uline{the} $\underline{O}_Y$-\uline{Module} $\underline{J}/\underline{J}^2$ (\uline{independent of} $U$) \uline{is called the} conormal sheaf of $Y$ in $X$, \uline{and is denoted by} $\underline{N}_{X/Y}$. \uline{If} $i$ \uline{is regular}, $\underline{J}/\underline{J}^2$ \uline{is locally free of finite type}, \uline{and its rank is called the} codimension of $Y$ in $X$, \uline{or the} \uline{codimension of} $i$.

Given two closed immersions: $Z\xrightarrow{j}Y\xrightarrow{i}X$, the conormal sheaves form an exact sequence:
\sgaStableTarget{sga6.e07.eq.1-6-1}\begin{equation}\tag{1.6.1}
 j^*(\underline{N}_{X/Y})\longrightarrow \underline{N}_{X/Z}\longrightarrow \underline{N}_{Y/Z}\longrightarrow0.
\end{equation}
Indeed, let $\underline{I}$ be the ideal of $Y$ in $X$, $\underline{J}$ the ideal of $Z$ in $X$, $\underline{K}$ the ideal of $Z$ in $Y$. We have the exact sequence:
\[
0\longrightarrow \underline{I}\longrightarrow \underline{J}\longrightarrow \underline{K}\longrightarrow0;
\]
whence, by tensoring with $\underline{O}_X/\underline{I}=\underline{O}_Y$, the exact sequence:
\[
 \underline{I}/\underline{I}^2\longrightarrow \underline{J}\otimes_{\underline{O}_X}\underline{O}_Y\longrightarrow \underline{K}\otimes_{\underline{O}_X}\underline{O}_Y\longrightarrow0;
\]
tensoring with $\underline{O}_Z=\underline{O}_X/\underline{J}=\underline{O}_Y/\underline{K}$, we obtain the exact sequence:
\[
 j^*(\underline{I}/\underline{I}^2)\longrightarrow \underline{J}\otimes_{\underline{O}_X}\underline{O}_Z\longrightarrow \underline{K}\otimes_{\underline{O}_Y}\underline{O}_Z\longrightarrow0,
\]
which is the required sequence.

%%%%%%%%%%  scan idx 435  |  folio 422  %%%%%%%%%%
\SGAPrintPage{422}%
\sgaStableTarget{sga6.e07.prop.1-7}\textbf{Proposition 1.7.} (\uline{Compare} EGA IV 19.1.5 (iv)). \uline{Let} $i:Y\longrightarrow X$, \uline{and} $j:Z\longrightarrow Y$ \uline{be two immersions}:
\begin{enumerate}[label=\roman*)]
\item \uline{If} $i$ \uline{and} $j$ \uline{are regular}, $i\circ j$ \uline{is regular}, \uline{and the sequence}:
\[
0\longrightarrow j^*(\underline{N}_{X/Y})\longrightarrow \underline{N}_{X/Z}\longrightarrow \underline{N}_{Y/Z}\longrightarrow0
\]
\uline{is exact}; \uline{in particular}, $\codim(i\circ j)=\codim(i)+\codim(j)$.
\item \uline{If} $i\circ j$ \uline{and} $i$ \uline{are regular}, \uline{and if the sequence}:
\[
0\longrightarrow j^*(\underline{N}_{X/Y})\longrightarrow \underline{N}_{X/Z}\longrightarrow \underline{N}_{Y/Z}\longrightarrow0,
\]
\uline{is exact and locally split}, $j$ \uline{is regular};
\item \uline{If} $X$ \uline{is noetherian}, \uline{and if} $i\circ j$ \uline{and} $j$ \uline{are regular}, $i$ \uline{is regular at the points of} $Z$.
\end{enumerate}

Since the question is local on $X$, we may suppose $X$ affine and the immersions closed; let $X=\Spec(A)$, $Y=\Spec(B)$, $Z=\Spec(C)$, with $B=A/I$, $C=A/J=B/K$.

i) After localizing, we may suppose $I/I^2$ free over $B$, and $K/K^2$ free over $C$ (1.3). Let $f_1,\ldots,f_n$ be elements of $A$ lifting a basis of $I/I^2$, $g_1,\ldots,g_m$ elements of $A$ lifting a basis of $K/K^2$; then, after further localization, we may suppose that $f_1,\ldots,f_n$ generate $I$, that $f_1,\ldots,f_n,g_1,\ldots,g_m$ generate $J$, and that $\bar g_1,\ldots,\bar g_m$ generate $K$, where $\bar g_j$ is the image of $g_j$ in $B$ (1.4); moreover, $f_1,\ldots,f_n$ form a minimum system of generators of $I$, and $\bar g_1,\ldots,\bar g_m$ a minimum system of generators of $K$.

Consider the diagram:
\sgaStableTarget{sga6.e07.eq.1-7-1}\begin{equation}\tag{1.7.1}
\begin{tikzcd}[column sep=large,row sep=large]
A_0=\mathbb Z[T_1,\ldots,T_{n+m}] \arrow[r] \arrow[d] & A \arrow[d]\\
B_0=\mathbb Z[T_1,\ldots,T_m] \arrow[r] \arrow[d] & B \arrow[d]\\
C_0=\mathbb Z \arrow[r] & C,
\end{tikzcd}
\end{equation}
%%%%%%%%%%  scan idx 436  |  folio 423  %%%%%%%%%%
\SGAPrintPage{423}%
where the horizontal arrows are defined by sending $T_i$ to $f_i$ for $1\leq i\leq n$, and $T_{n+j}$ to $g_j$ for $1\leq j\leq m$, where the homomorphism $A_0\longrightarrow B_0$ is defined by sending $T_i$ to $0$ for $1\leq i\leq n$, and $T_{n+j}$ to $T_j$ for $1\leq j\leq m$, and the homomorphism $B_0\longrightarrow C_0$ by setting the $T_j$, $1\leq j\leq m$, to zero. We then have:
\[
 B=A\otimes_{A_0}B_0,
 \qquad
 C=A\otimes_{A_0}C_0=B\otimes_{B_0}C_0.
\]
The situation under consideration is therefore deduced from the situation $(A_0,B_0,C_0)$ by the base change $A_0\longrightarrow A$. The corresponding immersions $i_0$, $j_0$ and $i_0\circ j_0$ are regular; consequently, the homology groups of the Koszul complexes associated with the sequences $(f_1,\ldots,f_n)$, $(f_1,\ldots,f_n,g_1,\ldots,g_m)$ and $(\bar g_1,\ldots,\bar g_m)$ are respectively the $\Tor_i^{A_0}(B_0,A)$, $\Tor_i^{A_0}(C_0,A)$, $\Tor_i^{B_0}(C_0,B)$. Consequently:
\sgaStableTarget{sga6.e07.eq.1-7-2}\begin{equation}\tag{1.7.2}
\begin{aligned}
 i\ \text{regular} &\Longleftrightarrow \forall i\geq1,
 \Tor_i^{A_0}(B_0,A)=0;\\
 i\circ j\ \text{regular} &\Longleftrightarrow \forall i\geq1,
 \Tor_i^{A_0}(C_0,A)=0;\\
 j\ \text{regular} &\Longleftrightarrow \forall i\geq1,
 \Tor_i^{B_0}(C_0,B)=0.
\end{aligned}
\end{equation}

We then use the transitivity spectral sequence ([1], XVI 5)
\sgaStableTarget{sga6.e07.eq.1-7-3}\begin{equation}\tag{1.7.3}
E^2_{pq}=\Tor_p^{B_0}\bigl(C_0,\Tor_q^{A_0}(B_0,A)\bigr)
\Longrightarrow \Tor_n^{A_0}(C_0,A)..
\end{equation}

By the hypotheses of i), $E^2_{pq}=0$ for $p$ or $q\ne0$. It therefore degenerates, and $i\circ j$ is regular by (1.7.2).

By (1.6.1), the sequence:
\[
 j^*(\underline{N}_{X/Y})\longrightarrow \underline{N}_{X/Z}\longrightarrow \underline{N}_{Y/Z}\longrightarrow0
\]
is exact. Since $i\circ j$ is regular, the three conormal sheaves are locally free of finite type; moreover, their rank at a point is the
%%%%%%%%%%  scan idx 437  |  folio 424  %%%%%%%%%%
\SGAPrintPage{424}%
number of elements in a minimum system of generators of the corresponding ideal; by the preceding proof, we have
\[
 \rang(\underline{N}_{X/Z})=\rang(\underline{N}_{X/Y})+\rang(\underline{N}_{Y/Z}).
\]
It follows that the kernel of $j^*(\underline{N}_{X/Y})\longrightarrow \underline{N}_{X/Z}$ is locally free of rank zero, and hence that the sequence (1.6.1) is exact.

ii) We may suppose $J/J^2$ free over $C$, and $I/I^2$ free over $B$; moreover, the hypothesis on the sequence of conormal sheaves allows us to suppose that $I/I^2\otimes_A A/J$ is a direct summand of $J/J^2$. We can therefore find elements $f_1,\ldots,f_n$ of $A$ lifting a basis of $I/I^2$, and elements $g_1,\ldots,g_m$ of $A$ such that $f_1,\ldots,f_n,g_1,\ldots,g_m$ lift a basis of $J/J^2$, the images of $g_1,\ldots,g_m$ generating $K/K^2$; after further localization, we may suppose that $f_1,\ldots,f_n$ form a minimum system of generators of $I$, that $f_1,\ldots,f_n,g_1,\ldots,g_m$ form a minimum system of generators of $J$, and that $\bar g_1,\ldots,\bar g_m$ generate $K$. We can then construct the diagram (1.7.1) and use the spectral sequence (1.7.3).

By (1.7.2), since $i$ is regular, we have $E^2_{pq}=0$ for $q\geq1$; we therefore obtain edge-isomorphisms: $E^2_{p,0}\simeq\Tor_p^{A_0}(C_0,A)$. Since $i\circ j$ is regular, this implies by (1.7.2) that $E^2_{p,0}=0$ for $p\geq1$, and hence that $j$ is regular, again by (1.7.2).

iii) This last property is proved in EGA IV 19.1.5 (iv), since in the noetherian case the definition in EGA IV coincides with the one we adopt here.

\sgaStableTarget{sga6.e07.prop.1-8}\textbf{Proposition 1.8.} \uline{Let} \(X\) \uline{be a scheme}, \(\underline{E}\) \uline{an} \(\underline{O}_X\)-\uline{Module locally free of finite type}, \(u:\underline{E}\to\underline{O}_X\) \uline{a regular homomorphism}, \(\underline{J}=u(\underline{E})\). \uline{Denote by} \(Y\)
%%%%%%%%%%  scan idx 438  |  folio 425  %%%%%%%%%%
\SGAPrintPage{425}%
\uline{the closed subscheme of} \(X\) \uline{defined by} \(\underline{J}\), \(X'\) \uline{the scheme obtained by blowing up} \(Y\), \(P\) \uline{the projective bundle associated with} \(\underline{E}\), \(p\) \uline{the structural morphism} \(P\to X\). \uline{On} \(P\), \uline{consider the canonical exact sequence}:
\sgaStableTarget{sga6.e07.eq.1-8-1}\[
0\to \underline{F}\to p^*(\underline{E})\to\underline{O}_P(1)\to0,
\]
\uline{and the composite morphism}:
\[
 v:\underline{F}\longrightarrow p^*(\underline{E})\xrightarrow{p^*(u)}\underline{O}_P .
\]

\uline{Then}:
\begin{enumerate}[label=\roman*)]
\item \uline{The formation of} \(X'\) \uline{commutes with every base change} \(f:T\to X\) \uline{such that} \(f^*(u)\) \uline{is regular}.
\item \uline{The closed immersion} \(j:X'\to P\) \uline{defined by the surjective homomorphism} \(\mathbf{S}(\underline{E})\longrightarrow\bigoplus_{n\ge0}\underline{J}^n\) \uline{defined by} \(u\) \uline{in degree} \(1\) \uline{is a regular immersion}, \uline{its ideal being the image of the regular morphism} \(v\); \uline{in particular, if} \(\operatorname{codim}(i)\) \uline{is constant}, \(\operatorname{codim}(j)=\operatorname{codim}(i)-1\).
\item \uline{The} \(\mathbf{S}(\underline{E})\)-\uline{Module} \(\bigoplus_{n\ge0}\underline{J}^n\) \uline{is perfect as a graded} \(\mathbf{S}(\underline{E})\)-\uline{Module}, \uline{i}.\uline{e}. \uline{locally there exists a finite}, \uline{graded resolution of} \(\bigoplus_{n\ge0}\underline{J}^n\) \uline{by graded} \(\mathbf{S}(\underline{E})\)-\uline{Modules locally free of finite type} (VI 1.2).
\end{enumerate}

i) Since \(X'=\Proj(\bigoplus_{n\ge0}\underline{J}^n)\), the result is an immediate consequence of the relations (1.2.3) and of EGA II.

ii) Since the stated properties are local on \(X\), we may suppose \(X\) affine, and moreover \(\underline{E}\) free over \(X\). Put \(X=\Spec(A)\).

First we give the proof in the standard case \(A=\mathbb Z[T_1,\ldots,T_d]\), with \(E\) endowed with a basis \(e_1,\ldots,e_d\) such that \(u(e_i)=T_i\). Then \(P=\Proj\bigl(\mathbb Z[T_1,\ldots,T_d,X_1,\ldots,X_d]\bigr)\), the grading being taken with respect to the \(X_i\). To show that \(v\) is a regular morphism whose image is the Ideal of \(j\), it suffices to do so locally on \(P\), hence on the
%%%%%%%%%%  scan idx 439  |  folio 426  %%%%%%%%%%
\SGAPrintPage{426}%
\(D^+(X_i)\); let us work, for example, on \(D^+(X_1)\). Then \(p^*(\underline{E})\) is defined by the module \(\mathbb Z[T_1,\ldots,T_d,X_2/X_1,\ldots,X_d/X_1]\otimes E\), and \(\underline{O}_P(1)\) by \(\mathbb Z[T_1,\ldots,T_d,X_2/X_1,\ldots,X_d/X_1]\), the canonical morphism sending \(e_i\) to \(X_i/X_1\) for every \(i\). Then \(F\) has as a basis the elements \(e_i-e_1.X_i/X_1\) for \(i>1\).

We have \(v(e_i-e_1.X_i/X_1)=T_i-T_1.X_i/X_1\); put \(T'_i=T_i-T_1X_i/X_1\).
Since the sequence of the \(T_i\) is regular in \(\mathbb Z[T_1,\ldots,T_d,X_2/X_1,\ldots,X_d/X_1]\), the same is true of the sequence \(T_1,T'_2,\ldots,T'_d\) deduced from it by an invertible matrix, and hence also of the subsequence \(T'_2,\ldots,T'_d\) (we are in a noetherian case). Consequently, \(v\) is a regular morphism.

We must therefore show that the \(T'_i\) generate the Ideal of \(j\). In \(\bigoplus_{n\ge0}J^n\), for \(m\le p\), we shall denote by \(x^{(m)}\) the element \(x\) of \(J^p\), considered as having degree \(m\). Over \(D^+(X_1)\), \(X'=\Spec\left(\left(\bigoplus_{n\ge0}J^n\right)_{(T_1^{(1)})}\right)\), and \(j\) is defined by the morphism sending \(X_i\) to \(T_i^{(1)}\). Under this morphism, the image of \(T'_i\) is \(T_i^{(0)}-T_1^{(0)}.T_i^{(1)}/T_1^{(1)}=\bigl((T_iT_1)^{(1)}-(T_1T_i)^{(1)}\bigr)/T_1^{(1)}=0\). Thus the \(T'_i\) belong to the Ideal of \(j\).

Since every element of \(\mathbb Z[T_1,\ldots,T_d,X_2/X_1,\ldots,X_d/X_1]\) can be expressed as a polynomial in \(T_1,T'_2,\ldots,T'_d,X_2/X_1,\ldots,X_d/X_1\), to show that the \(T'_i\) generate the Ideal of \(j\) it suffices to show that no nonzero element of \(\mathbb Z[T_1,X_2/X_1,\ldots,X_d/X_1]\) belongs to the Ideal of \(j\), which amounts to showing that the homomorphism
\[
\mathbb Z[T_1,X_1,\ldots,X_d]\longrightarrow\bigoplus_{n\ge0}J^n
\]
is injective, since \(T_1^{(1)}\) is not a zero divisor in \(\bigoplus_{n\ge0}J^n\). Grading \(\mathbb Z[T_1,X_1,\ldots,X_d]\) by assigning degree \(0\) to \(T_1\) and degree \(1\) to the \(X_i\), the kernel of the homomorphism under consideration is a graded ideal. Now it is clear that
%%%%%%%%%%  scan idx 440  |  folio 427  %%%%%%%%%%
\SGAPrintPage{427}%
this kernel is contained in that of the composite homomorphism:
\[
\mathbb Z[T_1,X_1,\ldots,X_d]\longrightarrow
\bigoplus_{n\ge0}J^n\longrightarrow A,
\]
the last homomorphism being defined in degree \(n\) by the inclusion of \(J^n\) in \(A\). Now the kernel of the composite homomorphism is the principal ideal generated by \((X_1-T_1)\), which, for reasons of integrality, contains no nonzero homogeneous elements. Thus the first homomorphism is injective, completing the proof in this case.

Return to the general case where \(A\) is arbitrary; we may still suppose \(\underline{E}\) free over \(X\). Attach the index \(0\) to the special case just treated; then, by part i) of the proposition, the situation under consideration is obtained by the base change \(A_0\to A\) from the situation studied above, the homomorphism \(A_0\to A\) being defined by taking a basis of \(E\), its image \(f_1,\ldots,f_d\) in \(J\), and sending \(T_i\) to \(f_i\). Put \(B_0=\left(\bigoplus_{n\ge0}J_0^n\right)_{(T_1^{(1)})}\), \(B=\left(\bigoplus_{n\ge0}J^n\right)_{(f_1^{(1)})}\); then \(B=B_0\otimes_{A_0}A\); the Koszul complex defined by \(v:\underline{F}\to\underline{O}_P\) is deduced by base change from the Koszul complex defined by \(v_0\); since the latter is a projective resolution of \(B_0\) (working on a suitable open subset), the homology groups of \(K_\bullet(v)\) on this open subset are the \(\Tor_i^{A_0}(A,B_0)\). These are the degree \(0\) components of the \(\Tor_i^{A_0}\left(A,\left(\bigoplus_{n\ge0}J_0^n\right)_{T_1^{(1)}}\right)=\left(\bigoplus_{n\ge0}\Tor_i^{A_0}(A,J_0^n)\right)_{T_1^{(1)}}\); the \(\Tor_i^{A_0}(A,J_0^n)\) vanish for \(i\ge1\) by the relations 1.2, so \(K_\bullet(v)\) is a resolution of \(B\), which shows that \(j\) is regular.

iii) Since the property is local on \(X\), we may again suppose \(X\) affine and \(\underline{E}\) free of finite type. Resume the notation of ii).

First consider the standard case of ii). The ring \(A_0\) is then noetherian and of finite cohomological dimension, and the algebra \(\mathbf{S}(E_0)\)
%%%%%%%%%%  scan idx 441  |  folio 428  %%%%%%%%%%
\SGAPrintPage{428}%
likewise. We can then construct step by step a partial graded resolution of \(\bigoplus_{n\ge0}J_0^n\) by graded free \(\mathbf{S}(E_0)\)-modules of finite type, of length greater than the cohomological dimension of \(A_0\); the last kernel is then a graded \(\mathbf{S}(E_0)\)-module, projective of finite type as an ungraded \(\mathbf{S}(E_0)\)-module, hence locally free of finite type over \(\Spec(A_0)\) as a graded \(\mathbf{S}(\underline{E}_0)\)-module by VI 1.5.1. We obtain the announced finite resolution by replacing the last module of the partial resolution by the corresponding kernel.

Tensoring the resolution thus constructed with \(A\), we obtain a graded complex of graded \(\mathbf{S}(E)\)-modules locally free of finite type, and must show that it is a resolution of \(\bigoplus_{n\ge0}J^n\). The \(i\)-th homology Module of the resulting complex is a graded \(\mathbf{S}(E)\)-Module whose \(n\)-th component is \(\Tor_i^{A_0}(A,J_0^n)\), which vanishes for \(i\ge1\) by (1.2.4). Consequently, by (1.2.3), we have constructed a resolution of the required type.

\sgaStableTarget{sga6.e07.prop.1-9}\textbf{Proposition 1.9.} \uline{Let} \(i:Y\to X\) \uline{be a regular closed immersion}, \(X'\) \uline{the} \(X\)-\uline{scheme obtained by blowing up} \(Y\), \uline{and} \(f:X'\to X\) \uline{the structural morphism}. \uline{Then} \(i\) \uline{and} \(f\) \uline{are perfect morphisms}. (III 4.1.).

Since the property is local on \(X\), we may suppose given an \(\underline{O}_X\)-Module locally free of finite type \(\underline{E}\), and a regular homomorphism \(u:\underline{E}\to\underline{O}_X\) whose image is the ideal \(\underline{J}\) of \(Y\). The Koszul complex \(\underline{K}_\bullet(u)\) is then a finite resolution of \(\underline{O}_Y\) by \(\underline{O}_X\)-Modules locally free of finite type; consequently, \(\underline{O}_Y\) is a perfect \(\underline{O}_X\)-Module. It follows that every regular immersion is a perfect morphism. By 1.8, \(f\) can be written as the composite of a regular immersion \(X'\to\bfP(\underline{E})\) and the structural morphism \(\bfP(\underline{E})\to X\), which is smooth and hence perfect; it is therefore
%%%%%%%%%%  scan idx 442  |  folio 429  %%%%%%%%%%
\SGAPrintPage{429}%
the composite of two perfect morphisms, and consequently a perfect morphism.

\sgaStableTarget{sga6.e07.prop.1-10}\textbf{Proposition 1.10.} \uline{Let} \(X\), \(Y\) \uline{be two schemes smooth over} \(S\). \uline{Then every} \(S\)-\uline{immersion of} \(X\) \uline{into} \(Y\) \uline{is regular}. \uline{In particular}, \uline{every section of a smooth morphism is a regular immersion}.

The proposition follows immediately from EGA IV 17.12.1, since a regular immersion in the sense of EGA IV is regular in the sense of 1.4.

\sgaStableTarget{sga6.e07.sec.2}\section*{2. Computations on regular immersions}

\sgaStableTarget{sga6.e07.para.2-1}\subsection*{\normalfont 2.1.}
Let \(A\) be a ring, and \(B\) and \(C\) two \(A\)-algebras. We propose to endow \(\bigoplus_{i\ge0}\Tor_i^A(B,C)\) with a graded \(A\)-algebra structure, and hence to define homomorphisms:
\begin{equation}\tag{2.1.1}
 \Tor_i^A(B,C)\otimes_A\Tor_j^A(B,C)\longrightarrow \Tor_{i+j}^A(B,C).
\end{equation}

Let \(P_\bullet\) be a resolution of \(B\) by projective \(A\)-modules; the \(\Tor_i^A(B,C)\) are the homology groups of the complex \(P_\bullet\otimes_A C\). We know how to define canonical homomorphisms ([1], IV 6):
\begin{equation}\tag{2.1.2}
 H_i(P_\bullet\otimes_A C)\otimes_A H_j(P_\bullet\otimes_A C)
 \longrightarrow H_{i+j}\bigl((P_\bullet\otimes_A C)\otimes_A(P_\bullet\otimes_A C)\bigr).
\end{equation}

Furthermore, \(P_\bullet\otimes_A P_\bullet\) is a complex of projective \(A\)-modules, in degrees \(\geq0\), and \(H_0(P_\bullet\otimes_A P_\bullet)=B\otimes_A B\); moreover, \(P_\bullet\) is a resolution of \(B\). Consequently, the augmentation morphism \(B\otimes_A B\to B\) extends to a morphism of complexes: \(P_\bullet\otimes_A P_\bullet\to P_\bullet\) and therefore defines, using the homomorphism \(C\otimes_A C\to C\) deduced from the algebra structure of \(C\), homomorphisms:
\begin{equation}\tag{2.1.3}
 H_{i+j}\bigl((P_\bullet\otimes_A C)\otimes_A(P_\bullet\otimes_A C)\bigr)
 \longrightarrow H_{i+j}(P_\bullet\otimes_A C).
\end{equation}

%%%%%%%%%%  scan idx 443  |  folio 430  %%%%%%%%%%
\SGAPrintPage{430}%
Composing (2.1.2) and (2.1.3), we obtain the homomorphisms (2.1.1).

The algebra structure defined by the homomorphisms (2.1.1) is associative (cf. [1] XI 2). Furthermore, the \(A\)-module \(H_i(P_\bullet\otimes_A C)\) has a standard \(B\otimes_A C\)-module structure, with \(C\) acting in the evident way, and \(B\) by functoriality; likewise, the \(H_j\bigl((P_\bullet\otimes_A C)\otimes_A(P_\bullet\otimes_A C)\bigr)\) carry two \(B\otimes_A C\)-module structures. The homomorphisms (2.1.2) and (2.1.3) are clearly compatible with these two structures: for (2.1.2), by functoriality (working in the category of \(A\)-modules), and for (2.1.3), by the classical result on morphisms of an acyclic complex. We have therefore endowed \(\bigoplus_{i\ge0}\Tor_i^A(B,C)\) with an associative \(B\otimes_A C\)-algebra structure.

\subsection*{\normalfont 2.2.}
Suppose that \(B\) admits a resolution by projective \(A\)-modules \(P_\bullet\) such that \(P_\bullet=\bigoplus_{i\ge0}P_i\) is a graded \(A\)-algebra, with the differential satisfying:
\[
\begin{array}{ll}
\text{a)}& P_0\longrightarrow B\text{ is a ring homomorphism~;}\SGAAuthorFootMark{1}\\[0.4ex]
\text{b)}& \forall x\in P_i,\ \forall y\in P_j,\quad d(xy)=d(x)y+(-1)^i x\,d(y).
\end{array}
\]
\SGAAuthorFootText{1}{One can show that such a resolution always exists.}
Then the ``product'' maps define commutative diagrams for \(i,j\ge0\):
\sgaStableTarget{sga6.e07.eq.2-2-1}\begin{equation}\notag
\begin{tikzcd}[column sep=large,row sep=large]
P_i\otimes_A P_j \arrow[r] \arrow[d] & (P_{i-1}\otimes_A P_j)\oplus(P_i\otimes_A P_{j-1}) \arrow[d]\\
P_{i+j} \arrow[r] & P_{i+j-1},
\end{tikzcd}
\end{equation}
and consequently a morphism of complexes \(P_\bullet\otimes_A P_\bullet\to P_\bullet\) extending the augmentation homomorphism: \(B\otimes_A B\to B\). This shows that the algebra structure of \(\bigoplus_{i\ge0}\Tor_i^A(B,C)\) can be computed from that of \(P_\bullet\). In particular, if \(P_\bullet\) is an alternating \(A\)-algebra (i.e. such that for every
%%%%%%%%%%  scan idx 444  |  folio 431  %%%%%%%%%%
\SGAPrintPage{431}%
\(x\in P_\bullet\), we have \(x^2=0\)), then the same is true of \(\bigoplus_{i\ge0}\Tor_i^A(B,C)\).

This last remark applies when \(B\) is the quotient of \(A\) by a regular ideal, since we can then take for \(P_\bullet\) the Koszul complex associated with a regular homomorphism defining the ideal of \(B\).

\subsection*{\normalfont 2.3.}
The \(B\otimes_A C\)-algebra structure defined in 2.1 on \(\bigoplus_{i\ge0}\Tor_i^A(B,C)\) is functorial in \(B\) and \(C\). We can therefore globalize the preceding constructions: let \(X\) be a scheme, \(X'\), \(X''\) two \(X\)-schemes; we can then define the morphisms (2.1.1) locally on \(X'\times_X X''\); they glue by functoriality and define on \(\bigoplus_{i\ge0}\underline{\Tor}_i^{\underline{O}_X}(\underline{O}_{X'},\underline{O}_{X''})\) an associative graded \(\underline{O}_{X'\times_X X''}\)-Algebra structure.

\sgaStableTarget{sga6.e07.lem.2-4}\textbf{Lemma 2.4.} \uline{Let} \(X',X''\) \uline{be two} \(X\)-\uline{schemes}, \uline{and} \(\underline{M}\) \uline{a flat} \(\underline{O}_{X'}\)-\uline{Module}. \uline{Then for every} \(i\ge0\), \uline{we have a canonical isomorphism}:
\sgaStableTarget{sga6.e07.eq.2-4-1}\begin{equation}\notag
\underline{\Tor}_i^{\underline{O}_X}(\underline{M},\underline{O}_{X''})\simeq
\underline{M}\otimes_{\underline{O}_{X'}}\underline{\Tor}_i^{\underline{O}_X}(\underline{O}_{X'},\underline{O}_{X''}).
\end{equation}

It suffices to define a functorial isomorphism in the affine case; we can then define it globally by gluing in the general case.

Let \(X=\Spec(A)\), \(X'=\Spec(A')\), \(X''=\Spec(A'')\); \(\underline{M}\) is defined by a flat \(A'\)-module \(M\). Let \(P_\bullet\) be a projective resolution of \(A''\) over \(A\); the \(\Tor_i^A(A',A'')\) are the homology groups of \(A'\otimes_A P_\bullet\). Since \(M\) is flat over \(A'\), we have a canonical isomorphism: \(M\otimes_{A'}\Tor_i^A(A',A'')\xrightarrow{\sim}H_i\bigl(M\otimes_{A'}(A'\otimes_A P_\bullet)\bigr)\). Furthermore, the canonical isomorphism: \(M\otimes_{A'}(A'\otimes_A P_\bullet)\xrightarrow{\sim}M\otimes_A P_\bullet\) gives a canonical isomorphism: \(H_i\bigl(M\otimes_{A'}(A'\otimes_A P_\bullet)\bigr)\xrightarrow{\sim}\Tor_i^A(M,A'')\), whence the result.

\sgaStableTarget{sga6.e07.prop.2-5}\textbf{Proposition 2.5.} \uline{Let} \(X\) \uline{be a scheme}, \(Y\) \uline{and} \(Z\) \uline{two closed subschemes of} \(X\), \(T=Y\times_X Z\). \uline{Suppose that for every} \(x\in T\), \uline{there exists an open neighborhood}
%%%%%%%%%%  scan idx 445  |  folio 432  %%%%%%%%%%
\SGAPrintPage{432}%
\(U\) \uline{of} \(x\), \uline{and a sequence of sections} \(f_1,\ldots,f_p,g_1,\ldots,g_r,h_1,\ldots,h_s\) \uline{of} \(\underline{O}_U\) \uline{such that the Koszul complex associated with} \(f_1,\ldots,f_p,g_1,\ldots,g_r\) \uline{is a resolution of} \(\underline{O}_{Y\cap U}\), \uline{the Koszul complex associated with} \(f_1,\ldots,f_p,h_1,\ldots,h_s\) \uline{a resolution of} \(\underline{O}_{Z\cap U}\), \uline{and the Koszul complex associated with} \(f_1,\ldots,f_p,g_1,\ldots,g_r,h_1,\ldots,h_s\) \uline{a resolution of} \(\underline{O}_{T\cap U}\) (\uline{conditions satisfied in particular in the noetherian case when the sequence} \(f_1,\ldots,f_p,g_1,\ldots,g_r,h_1,\ldots,h_s\) \uline{is regular and} \(Y=V(f_1,\ldots,f_p,g_1,\ldots,g_r)\), \(Z=V(f_1,\ldots,f_p,h_1,\ldots,h_s)\)). \uline{Let} \(\underline{I},\underline{J}\) \uline{be the ideals of} \(Y,Z\). \uline{We then have the following isomorphisms}:
\[
\begin{array}{ll}
\text{i)}& \forall i\ge0,\quad
\underline{\Tor}_i^{\underline{O}_X}(\underline{O}_Y,\underline{O}_Z)\xrightarrow{\sim}
\bigwedge^i\underline{\Tor}_1^{\underline{O}_X}(\underline{O}_Y,\underline{O}_Z);\\[0.8ex]
\text{ii)}&
\underline{\Tor}_1^{\underline{O}_X}(\underline{O}_Y,\underline{O}_Z)\simeq
\underline{I}\cap \underline{J}/\underline{I}.\underline{J}\simeq
\Ker\bigl(i^*(\underline{N}_{Y/X})\times j^*(\underline{N}_{Z/X})\to \underline{N}_{T/X}\bigr),
\end{array}
\]
\uline{where we denote by} \(i\) \uline{the immersion of} \(T\) \uline{into} \(Y\) \uline{and by} \(j\) \uline{the immersion of} \(T\) \uline{into} \(Z\).

\uline{In particular}, \uline{for every regular closed immersion} \(Y\to X\), \uline{we have}
\[
\begin{array}{ll}
\text{i')}& \forall i\ge0,\quad
\underline{\Tor}_i^{\underline{O}_X}(\underline{O}_Y,\underline{O}_Y)\xrightarrow{\sim}
\bigwedge^i\underline{\Tor}_1^{\underline{O}_X}(\underline{O}_Y,\underline{O}_Y);\\[0.8ex]
\text{ii')}&
\underline{\Tor}_1^{\underline{O}_X}(\underline{O}_Y,\underline{O}_Y)\simeq \underline{N}_{Y/X}.
\end{array}
\]

i) We saw in 2.3 that \(\bigoplus_{i\ge0}\underline{\Tor}_i^{\underline{O}_X}(\underline{O}_Y,\underline{O}_Z)\) could be endowed with an \(\underline{O}_T\)-Algebra structure; by 2.2, this algebra is alternating, since the immersion of \(Y\) (or \(Z\)) into \(X\) is regular. There therefore exists one and only one algebra homomorphism:
\sgaStableTarget{sga6.e07.eq.2-5-1}\begin{equation}\tag{2.5.1}
\bigwedge \underline{\Tor}_1^{\underline{O}_X}(\underline{O}_Y,\underline{O}_Z)\longrightarrow
\bigoplus_{i\ge0}\underline{\Tor}_i^{\underline{O}_X}(\underline{O}_Y,\underline{O}_Z),
\end{equation}
extending the identity in degree \(1\).

To see that it is an isomorphism, we can check locally. Thus we work in a neighborhood of a point of \(T\) and consider sections \(f_1,\ldots,f_p,g_1,\ldots,g_r,h_1,\ldots,h_s\) of \(\underline{O}_X\) on this neighborhood, having the properties stated in 2.5. Put \(\underline{E}=(\underline{O}_X)^p\), \(\underline{E}'=(\underline{O}_X)^r\), and let \(u:\underline{E}\to\underline{O}_X\) be the homomorphism defined by the \(f_i\), \(v:\underline{E}'\to\underline{O}_X\) the homomorphism defined by the \(g_j\), and \(\underline{K}_\bullet(u)\) and \(\underline{K}_\bullet(v)\) the associated
%%%%%%%%%%  scan idx 446  |  folio 433  %%%%%%%%%%
\SGAPrintPage{433}%
Koszul complexes. The Koszul complex associated with the sequence \(f_1,\ldots,f_p,g_1,\ldots,g_r\) is isomorphic to the tensor-product complex \(\underline{K}_\bullet(u)\otimes_{\underline{O}_X}\underline{K}_\bullet(v)\). Since it is a resolution of \(\underline{O}_Y\) by free \(\underline{O}_X\)-Modules of finite type, the \(\underline{\Tor}_i^{\underline{O}_X}(\underline{O}_Y,\underline{O}_Z)\) are the homology sheaves of the complex \(\underline{K}_\bullet(u)\otimes_{\underline{O}_X}\underline{K}_\bullet(v)\otimes_{\underline{O}_X}\underline{O}_Z\). We can regard it as the tensor product of the complexes \(\underline{K}_\bullet(u)\) and \(\underline{K}_\bullet(v)\otimes_{\underline{O}_X}\underline{O}_Z\), and compute its homology using the spectral sequences associated with a bicomplex (EGA \(0_{\mathrm{III}}\) (11.3.5)), which abut to the homology of the associated simple complex.

Consider the spectral sequence such that \({}'E^2_{pq}=H^I_p\bigl(H^{II}_q(K_{\bullet\bullet})\bigr)\), in the notation of EGA \(0_{\mathrm{III}}\) (11.3.5), with the first index referring to the complex \(\underline{K}_\bullet(u)\) and the second to the complex \(\underline{K}_\bullet(v)\otimes_{\underline{O}_X}\underline{O}_Z\). Now the complex \(\underline{K}_\bullet(v)\otimes_{\underline{O}_X}\underline{O}_Z\) is the Koszul complex over \(\underline{O}_Z\) associated with the sequence of sections \(\overline g_1,\ldots,\overline g_r\), where the \(\overline g_i\) are the images of the \(g_i\) in \(\underline{O}_Z\), and this complex is a resolution of \(\underline{O}_T\) over \(\underline{O}_Z\) since the immersions \(Z\to X\), and \(T\to Z\to X\) are regular (cf. the proof of 1.7 ii). Since \(\bigwedge^p\underline{E}\) is free over \(\underline{O}_X\), the complex \(\bigwedge^p\underline{E}\otimes_{\underline{O}_X}(\underline{K}_\bullet(v)\otimes_{\underline{O}_X}\underline{O}_Z)\) is a resolution of \(\bigwedge^p\underline{E}\otimes_{\underline{O}_X}\underline{O}_T\). Consequently the spectral sequence degenerates, and the \(\underline{\Tor}_i^{\underline{O}_X}(\underline{O}_Y,\underline{O}_Z)\) are the homology sheaves of the complex \(\underline{K}_\bullet(u)\otimes_{\underline{O}_X}\underline{O}_T\). The differential of \(\underline{K}_\bullet(u)\) is expressed in terms of the \(f_i\), which annihilate \(\underline{O}_T\); consequently the differential of \(\underline{K}_\bullet(u)\otimes_{\underline{O}_X}\underline{O}_T\) is zero, and we therefore obtain the isomorphisms:
\sgaStableTarget{sga6.e07.eq.2-5-2}\begin{equation}\tag{2.5.2}
\underline{\Tor}_i^{\underline{O}_X}(\underline{O}_Y,\underline{O}_Z)
\xrightarrow{\sim}
\underline{O}_T\otimes_{\underline{O}_X}\bigwedge^i\underline{E}
\longrightarrow
\bigwedge^i(\underline{O}_T)^p.
\end{equation}

In particular, for \(i=1\) we find:
\[
\underline{\Tor}_1^{\underline{O}_X}(\underline{O}_Y,\underline{O}_Z)\xrightarrow{\sim}(\underline{O}_T)^p,
\]
whence for every \(i\):
\sgaStableTarget{sga6.e07.eq.2-5-3}\begin{equation}\tag{2.5.3}
\underline{\Tor}_i^{\underline{O}_X}(\underline{O}_Y,\underline{O}_Z)
\xrightarrow{\sim}
\bigwedge^i\underline{\Tor}_1^{\underline{O}_X}(\underline{O}_Y,\underline{O}_Z).
\end{equation}

%%%%%%%%%%  scan idx 447  |  folio 434  %%%%%%%%%%
\SGAPrintPage{434}%
Now it follows from 2.2 that the isomorphism:
\[
\bigwedge(\underline{O}_T)^p\xrightarrow{\sim}
\bigoplus_{i\ge0}\underline{\Tor}_i^{\underline{O}_X}(\underline{O}_Y,\underline{O}_Z)
\]
defined by the isomorphisms (2.5.2) is compatible with the algebra structures of the two sides; hence the same is true of the isomorphism:
\[
\bigwedge\underline{\Tor}_1^{\underline{O}_X}(\underline{O}_Y,\underline{O}_Z)\xrightarrow{\sim}
\bigoplus_{i\ge0}\underline{\Tor}_i^{\underline{O}_X}(\underline{O}_Y,\underline{O}_Z)
\]
defined by the isomorphisms (2.5.3). By uniqueness, the latter is therefore the restriction of the homomorphism (2.5.1) that was defined globally. The homomorphism (2.5.1) is thus an isomorphism.

ii) Consider the exact sequence:
\[
0\longrightarrow \underline{I}\longrightarrow \underline{O}_X\longrightarrow \underline{O}_Y\longrightarrow0,
\]
and the resulting exact \(\underline{\Tor}\) sequence:
\[
0\longrightarrow\underline{\Tor}_1^{\underline{O}_X}(\underline{O}_Y,\underline{O}_Z)
\longrightarrow \underline{I}\otimes_{\underline{O}_X}\underline{O}_Z
\longrightarrow \underline{O}_Z.
\]
It follows that \(\underline{\Tor}_1^{\underline{O}_X}(\underline{O}_Y,\underline{O}_Z)\) is the kernel of the homomorphism:
\[
\underline{I}\otimes_{\underline{O}_X}(\underline{O}_X/\underline{J})\longrightarrow \underline{O}_X/J,
\]
that is, \(\underline{I}\cap \underline{J}/\underline{I}.\underline{J}\).

To prove the last relation, we shall show that there is an exact sequence:
\[
\underline{0}\longrightarrow \underline{I}\cap \underline{J}/\underline{I}.\underline{J}
\longrightarrow i^*(\underline{N}_{Y/X})\times j^*(\underline{N}_{Z/X})
\longrightarrow \underline{N}_{T/X}\longrightarrow0.
\]
The homomorphisms \(i^*(\underline{N}_{Y/X})\to \underline{N}_{T/X}\) and \(j^*(\underline{N}_{Z/X})\to \underline{N}_{T/X}\) were defined in (1.6.1); here we take their difference. Furthermore, \(\underline{I}\cap \underline{J}\) maps naturally into \(\underline{I}/\underline{I}^2\), hence into \(i^*(\underline{I}/\underline{I}^2)\), and likewise into \(j^*(\underline{J}/\underline{J}^2)\), and these homomorphisms vanish on \(\underline{I}.\underline{J}\), since \(\underline{I},\underline{J}\subset \underline{K}\) (where \(\underline{K}\) denotes the ideal of \(T\)), defining the arrow on the left.

%%%%%%%%%%  scan idx 448  |  folio 435  %%%%%%%%%%
\SGAPrintPage{435}%
To check exactness of the sequence, we can work locally on \(T\) and use the sections defined in the statement of 2.5. We then have a basis of \(\underline{N}_{T/X}\) formed by the images of the \(f_i,g_j,h_k\), a basis of \(i^*(\underline{N}_{Y/X})\) formed by the images of the \(f_i,g_j\), and a basis of \(j^*(\underline{N}_{Z/X})\) formed by the images of the \(f_i,h_k\). It is then clear that \(i^*(\underline{N}_{Y/X})\times j^*(\underline{N}_{Z/X})\to \underline{N}_{T/X}\) is surjective, and that its kernel consists of the elements whose coordinates on the \((\overline g_j,0)\) and the \((0,\overline h_k)\) vanish, and whose coordinates on \((\overline f_i,0)\) and \((0,\overline f_i)\) are equal for every \(i\). Consequently the sequence is also exact in the middle. The three \(\underline{O}_T\)-Modules involved are locally free of finite type; since \(\underline{\Tor}_1^{\underline{O}_X}(\underline{O}_Y,\underline{O}_Z)\) has rank \(p\) by the proof of i), we see immediately that the sequence is exact.

\sgaStableTarget{sga6.e07.para.2-6}\subsection*{\normalfont 2.6.}
We now suppose that \(X\) is a scheme having an ample sheaf. Every closed subscheme \(Y\) of \(X\) then has an ample sheaf. We know that in this case ``global resolutions exist'', and that the \(K^{\bullet}\) formed respectively with locally free sheaves of finite type and with perfect complexes are identical (II, 2.2.9). If moreover the immersion \(i:Y\to X\) is regular, it is a perfect morphism, and we can define an additive homomorphism \(i_*:K^{\bullet}(Y)\longrightarrow K^{\bullet}(X)\) (IV 2.12) by associating with the class of a perfect complex \(\underline{L}^\bullet\) the class of the complex \(\mathbb Ri_*(\underline{L}^\bullet)=i_*(\underline{L}^\bullet)\) (\(i_*\) being an exact functor).

\textbf{Proposition 2.7.} \uline{Let} \(X\) \uline{be a scheme having an ample sheaf}, \(i:Y\to X\) \uline{a regular closed immersion}, \(\underline{N}\) \uline{the conormal sheaf of} \(Y\) \uline{in} \(X\). \uline{Then}, \uline{for every} \(x\in K^{\bullet}(Y)\), \uline{we have the relation}:
\sgaStableTarget{sga6.e07.eq.2-7-1}\begin{equation}\notag
i^*(i_*(x))=x.\lambda_{-1}(N),
\end{equation}
\uline{where} \(N\) \uline{is the class of} \(\underline{N}\) \uline{in} \(K^{\bullet}(Y)\), \uline{and} \(\lambda_{-1}(N)=\sum_{i\ge0}(-1)^i\lambda^i(N)\).

%%%%%%%%%%  scan idx 449  |  folio 436  %%%%%%%%%%
\SGAPrintPage{436}%
Since both sides of the relation to be proved are additive, we may suppose that \(x\) is the class of an \(\underline{O}_Y\)-Module locally free of finite type \(\underline{M}\). The element \(i^*(i_*(x))\) is then the class of the complex \(\mathbb L i^*(\mathbb Ri_*(\underline{M}))\), obtained by taking a projective resolution of the complex \(\mathbb Ri_*(\underline{M})=i_*(\underline{M})\) and tensoring it with \(\underline{O}_Y\). The homology sheaves of this complex are therefore the \(\underline{\Tor}_i^{\underline{O}_X}(\underline{M},\underline{O}_Y)\). By 2.4, we have the isomorphism:
\[
\underline{\Tor}_i^{\underline{O}_X}(\underline{M},\underline{O}_Y)
\xrightarrow{\sim}
\underline{M}\otimes_{\underline{O}_Y}\underline{\Tor}_i^{\underline{O}_X}(\underline{O}_Y,\underline{O}_Y).
\]
By 2.5, we have the isomorphism:
\[
\underline{\Tor}_i^{\underline{O}_X}(\underline{O}_Y,\underline{O}_Y)
\xrightarrow{\sim}
\bigwedge^i\underline{N}.
\]

The \(\underline{\Tor}_i^{\underline{O}_X}(\underline{M},\underline{O}_Y)\) are therefore \(\underline{O}_Y\)-Modules locally free of finite type, and the class in \(K^{\bullet}(Y)\) of the complex \(\mathbb L i^*(\mathbb Ri_*(\underline{M}))\) is the alternating sum of the classes of its homology sheaves, namely:
\[
i^*(i_*(x))=\sum_{i\ge0}(-1)^i
\operatorname{cl}\bigl(\underline{M}\otimes_{\underline{O}_Y}\underline{\Tor}_i^{\underline{O}_X}(\underline{O}_Y,\underline{O}_Y)\bigr)
=\sum_{i\ge0}(-1)^i x.\lambda^i(N).
\]

\sgaStableTarget{sga6.e07.cor.2-8}\textbf{Corollary 2.8.} \uline{Under the hypotheses of} 2.7, \uline{let} \(x,y\in K^{\bullet}(Y)\); \uline{we have the relation}:
\sgaStableTarget{sga6.e07.eq.2-8-1}\begin{equation}\notag
i_*(xy\lambda_{-1}(N))=i_*(x).i_*(y).
\end{equation}

By 2.7, \(y\lambda_{-1}(N)=i^*(i_*(y))\).

Applying the projection formula, we therefore find:
\[
i_*(xy\lambda_{-1}(N))
=i_*(x.i^*(i_*(y)))
=i_*(x).i_*(y).
\]

\sgaStableTarget{sga6.e07.sec.3}\section*{3. Computation of \texorpdfstring{$K^{\bullet}$}{K} of a blown-up scheme}

\sgaStableTarget{sga6.e07.para.3-1}3.1. In this section, \(i:Y\to X\) denotes a regular closed immersion. Denote by \(X'\) the scheme obtained by blowing up \(Y\), and put \(Y'=Y\times_X X'\).
%%%%%%%%%%  scan idx 450  |  folio 437  %%%%%%%%%%
\SGAPrintPage{437}%
We know (EGA II 8.1.8) that \(Y'\) is a closed subscheme of \(X'\) defined by an ideal of \(\underline{O}_{X'}\), canonically isomorphic to \(\underline{O}_{X'}(1)\); since the immersion \(j:Y'\to X'\) is defined by an invertible ideal, it is a regular closed immersion of codimension \(1\). Denote by \(\underline{J}\) the ideal of \(\underline{O}_X\) defining \(Y\), and by \(\underline{J}'\) the ideal of \(\underline{O}_{X'}\) defining \(Y'\); thus \(\underline{J}'\simeq \underline{J}.\underline{O}_{X'}\simeq\underline{O}_{X'}(1)\). Finally, denote by \(f:X'\to X\), and \(g:Y'\to Y\) the structural morphisms. We therefore have the commutative diagram:
\sgaStableTarget{sga6.e07.eq.3-1-1}\begin{equation}\tag{3.1.1}
\begin{tikzcd}[column sep=large,row sep=large]
Y' \arrow[r,"j"] \arrow[d,"g"'] & X' \arrow[d,"f"]\\
Y \arrow[r,"i"'] & X.
\end{tikzcd}
\end{equation}

By definition, \(X'=\Proj(\underline{S})\), where \(\underline{S}=\bigoplus_{n\ge0}\underline{J}^n\). Consequently, we have:
\[
Y'=\Proj\bigl(i^*(\underline{S})\bigr)
=\Proj\left(\left(\bigoplus_{n\ge0}\underline{J}^n\right)\otimes_{\underline{O}_X}(\underline{O}_X/\underline{J})\right)
=\Proj\left(\bigoplus_{n\ge0}\underline{J}^n/\underline{J}^{n+1}\right).
\]
Since \(i\) is regular, \(\bigoplus_{n\ge0}\underline{J}^n/\underline{J}^{n+1}\simeq \underline{\Sym}_{\underline{O}_Y}(\underline{J}/\underline{J}^2)\) (1.3). Denote by \(\underline{N}\) the conormal sheaf \(\underline{J}/\underline{J}^2\); thus:
\[
Y'=\bfP(\underline{N}).
\]

On \(Y'\), we can write the exact sequence:
\sgaStableTarget{sga6.e07.eq.3-1-2}\begin{equation}\tag{3.1.2}
0\to \underline{F}\to g^*(\underline{N})\to\underline{O}_{Y'}(1)\to0.
\end{equation}

Observe that \(\underline{O}_{Y'}(1)=j^*(\underline{O}_{X'}(1))=\underline{J}'/\underline{J}'^2\); it is therefore the conormal sheaf of the immersion \(j\), and we shall denote it by \(\underline{L}\).

\sgaStableTarget{sga6.e07.lem.3-2}\textbf{Lemma 3.2.} \uline{With the hypotheses and notation of} 3.1, \uline{we have the following isomorphisms}:
\sgaStableTarget{sga6.e07.eq.3-2-i}\sgaStableTarget{sga6.e07.eq.3-2-ii}\[
\begin{array}{ll}
\text{i)}&
\forall i\ge0,\quad
\underline{\Tor}_i^{\underline{O}_X}(\underline{O}_{X'},\underline{O}_Y)
\simeq
\bigwedge^i\underline{\Tor}_1^{\underline{O}_X}(\underline{O}_{X'},\underline{O}_Y);\\[4pt]
\text{ii)}&
\underline{\Tor}_1^{\underline{O}_X}(\underline{O}_{X'},\underline{O}_Y)\simeq \underline{F}.
\end{array}
\]

%%%%%%%%%%  scan idx 451  |  folio 438  %%%%%%%%%%
\SGAPrintPage{438}%
i) Since \(\underline{O}_Y\) can locally be resolved by a Koszul complex, the graded algebra \(\bigoplus_{i\ge0}\underline{\Tor}_i^{\underline{O}_X}(\underline{O}_{X'},\underline{O}_Y)\) is alternating (2.2). There therefore exists one and only one algebra homomorphism:
\sgaStableTarget{sga6.e07.eq.3-2-1}\begin{equation}\tag{3.2.1}
\bigwedge \underline{\Tor}_1^{\underline{O}_X}(\underline{O}_{X'},\underline{O}_Y)
\longrightarrow
\bigoplus_{i\ge0}\underline{\Tor}_i^{\underline{O}_X}(\underline{O}_{X'},\underline{O}_Y),
\end{equation}
extending the inclusion of \(\underline{\Tor}_1^{\underline{O}_X}(\underline{O}_{X'},\underline{O}_Y)\) in
\(\bigoplus_{i\ge0}\underline{\Tor}_i^{\underline{O}_X}(\underline{O}_{X'},\underline{O}_Y)\). It suffices to check locally that it is an isomorphism.

We can therefore work on an affine open subset \(U=\Spec(A)\) of \(X\), with \(\underline{J}\) defined by an ideal \(J\) of \(A\), and choose a regular homomorphism \(u:E\to A\) defining \(J\), with \(E=(A)^n\). Let \(f_1,\ldots,f_d\) be the images of a basis of \(E\).

In \(\bigoplus_{n\ge0}J^n\), for \(m\le p\), we shall denote by \(x^{(m)}\) the element \(x\) of \(J^p\) considered as having degree \(m\). We can restrict to the open subset \(V\) of \(X'\) where, for example, the element \(f_1^{(1)}\) is invertible.

The \(\underline{\Tor}_i^{\underline{O}_X}(\underline{O}_{X'},\underline{O}_Y)\) are the homology sheaves of the complex \(K_\bullet(u)\otimes_{\underline{O}_X}\underline{O}_{X'}\). To study this complex on \(V\), make the following change of basis in \(E'=E\otimes_A\Gamma(V,\underline{O}_{X'})\): if \(e_1,\ldots,e_d\) is the given basis of \(E\), put:
\[
e'_1=e_1;\qquad \text{for }i\ne1,\quad
e'_i=e_i-(f_i^{(1)}/f_1^{(1)})e_1.
\]
Put \(A'=\Gamma(V,\underline{O}_{X'})\); the differential of the complex \(K_\bullet(u)\otimes_A A'\) is given by
\[
d(e'_1)=f_1^{(0)};\qquad \text{for }i\ne1,\quad d(e'_i)=0.
\]
Taking as a basis of \(\bigwedge^k(E')\) the elements
\(e'_{i_1}\wedge\cdots\wedge e'_{i_k}\) with \(i_1<\cdots<i_k\), we therefore find:
\[
\begin{array}{ll}
d(a.e'_{i_1}\wedge\cdots\wedge e'_{i_k})=0
&\text{if }i_1\ne1;\\[3pt]
d(a.e'_{i_1}\wedge\cdots\wedge e'_{i_k})=
af_1^{(0)}.e'_{i_2}\wedge\cdots\wedge e'_{i_k}
&\text{if }i_1=1.
\end{array}
\]
%%%%%%%%%%  scan idx 452  |  folio 439  %%%%%%%%%%
\SGAPrintPage{439}%
Denoting by \(E''\) the submodule of \(E'\) generated by \(e'_2,\ldots,e'_d\), we therefore see that the image of \(d:\bigwedge^kE'\longrightarrow\bigwedge^{k-1}E'\) is \(f_1^{(0)}\bigwedge^{k-1}E''\). Furthermore, \(f_1^{(0)}\) is regular in \(A'\), since on \(V\) the invertible ideal \(\underline{J}'\) is defined by the equation \(f_1^{(0)}\). The kernel of \(d:\bigwedge^kE'\longrightarrow\bigwedge^{k-1}E'\) is therefore \(\bigwedge^kE''\). Consequently, the \(k\)-th homology group of \(K_\bullet(u)\otimes_AA'\) is \(\bigwedge^kE''/f_1^{(0)}\bigwedge^kE''\); we thus obtain on \(V\) isomorphisms:
\sgaStableTarget{sga6.e07.eq.3-2-2}\begin{equation}\tag{3.2.2}
\underline{\Tor}_i^{\underline{O}_X}(\underline{O}_{X'},\underline{O}_Y)
\xrightarrow{\sim}
(\underline{O}_{X'}/f_1^{(0)}\underline{O}_{X'})\otimes_{\underline{O}_{X'}}
\bigwedge^i(\underline{O}_{X'})^{d-1}
\xrightarrow{\sim}
\bigwedge^i(\underline{O}_{Y'})^{d-1}.
\end{equation}
In particular, for \(i=1\) we find:
\[
\underline{\Tor}_1^{\underline{O}_X}(\underline{O}_{X'},\underline{O}_Y)
\xrightarrow{\sim}
(\underline{O}_{Y'})^{d-1}.
\]
Consequently, for every \(i\) we have:
\sgaStableTarget{sga6.e07.eq.3-2-3}\begin{equation}\tag{3.2.3}
\underline{\Tor}_i^{\underline{O}_X}(\underline{O}_{X'},\underline{O}_Y)
\xrightarrow{\sim}
\bigwedge^i\underline{\Tor}_1^{\underline{O}_X}(\underline{O}_{X'},\underline{O}_Y)
\end{equation}

By 2.2, the isomorphism:
\[
\bigwedge(\underline{O}_{Y'})^{d-1}
\xrightarrow{\sim}
\bigoplus_{i\ge0}\underline{\Tor}_i^{\underline{O}_X}(\underline{O}_{X'},\underline{O}_Y)
\]
defined by the isomorphisms (3.2.2) is an algebra isomorphism. Hence the same is true of the isomorphism:
\[
\bigwedge\underline{\Tor}_1^{\underline{O}_X}(\underline{O}_{X'},\underline{O}_Y)
\xrightarrow{\sim}
\bigoplus_{i\ge0}\underline{\Tor}_i^{\underline{O}_X}(\underline{O}_{X'},\underline{O}_Y)
\]
defined by the isomorphisms (3.2.3). By uniqueness, it is then the restriction to \(V\) of the homomorphism (3.2.1), and the latter is therefore an isomorphism.

ii) Consider the exact sequence:
\[
0\longrightarrow \underline{J}\longrightarrow\underline{O}_X\longrightarrow\underline{O}_Y\longrightarrow0,
\]
%%%%%%%%%%  scan idx 453  |  folio 440  %%%%%%%%%%
\SGAPrintPage{440}%
and the exact \(\underline{\Tor}\) sequence obtained by tensoring with \(\underline{O}_{X'}\):
\[
0\longrightarrow
\Tor_1^{\underline{O}_X}(\underline{O}_{X'},\underline{O}_Y)
\longrightarrow
\underline{J}\otimes_{\underline{O}_X}\underline{O}_{X'}
\longrightarrow
\underline{O}_{X'}
\longrightarrow
\underline{O}_{Y'}
\longrightarrow0.
\]
The image of the morphism: \(\underline{J}\otimes_{\underline{O}_X}\underline{O}_{X'}\to\underline{O}_{X'}\) is
\(\underline{J}.\underline{O}_{X'}=\underline{O}_{X'}(1)\). Consequently, we obtain the exact sequence:
\[
0\longrightarrow
\underline{\Tor}_1^{\underline{O}_X}(\underline{O}_{X'},\underline{O}_Y)
\longrightarrow
\underline{J}\otimes_{\underline{O}_X}\underline{O}_{X'}
\longrightarrow
\underline{O}_{X'}(1)
\longrightarrow0.
\]

Then tensor it with \(\underline{O}_{Y'}\); since
\(\underline{\Tor}_1^{\underline{O}_X}(\underline{O}_{X'},\underline{O}_Y)\) is an \(\underline{O}_{Y'}\)-Module, and
\(\underline{O}_{X'}(1)\) is invertible, we obtain the exact sequence:
\sgaStableTarget{sga6.e07.eq.3-2-4}\[
0\longrightarrow
\underline{\Tor}_1^{\underline{O}_X}(\underline{O}_{X'},\underline{O}_Y)
\longrightarrow
g^*(\underline{N})
\longrightarrow
\underline{O}_{Y'}(1)
\longrightarrow0.
\]
One readily checks that the morphism \(g^*(\underline{N})\to\underline{O}_{Y'}(1)\) occurring in this sequence is the canonical morphism, so that this sequence is identified with (3.1.2). By definition, we therefore have:
\[
\underline{\Tor}_1^{\underline{O}_X}(\underline{O}_{X'},\underline{O}_Y)=\underline{F}.
\]

\sgaStableTarget{sga6.e07.para.3-3}3.3. We now suppose that \(X\) is a scheme having an ample sheaf. Then the same is true of \(X'\), which is projective over \(X\), and of \(Y\) and \(Y'\). Furthermore, \(i\) and \(j\) are perfect morphisms, since they are regular immersions; \(g\) is smooth and hence perfect, since \(Y'=\bfP(\underline{N})\); finally, \(f\) is perfect by 1.9. Since these morphisms are projective, we can define the additive homomorphisms (IV
\[
i_*:K^{\bullet}(Y)\to K^{\bullet}(X);\quad
j_*:K^{\bullet}(Y')\to K^{\bullet}(X');\quad
f_*:K^{\bullet}(X')\to K^{\bullet}(X);\quad
g_*:K^{\bullet}(Y')\to K^{\bullet}(Y).
\]

%%%%%%%%%%  scan idx 454  |  folio 441  %%%%%%%%%%
\SGAPrintPage{441}%
\sgaStableTarget{sga6.e07.prop.3-4}\textbf{Proposition 3.4.} \uline{With the hypotheses and notation of} 3.1 \uline{and} 3.3, \uline{for every} \(x\in K^{\bullet}(Y)\) \uline{we have the relation}:
\sgaStableTarget{sga6.e07.eq.3-4-1}\[
f^*(i_*(x))=j_*\bigl(g^*(x).\lambda_{-1}(F)\bigr).
\]

Since both sides of the relation to be proved are additive, we may suppose that \(x\) is the class of an \(\underline{O}_Y\)-Module locally free of finite type \(\underline{M}\). The element \(f^*(i_*(x))\) is then the class of the complex \(\mathbb Lf^*(\mathbb Ri_*(\underline{M}))\), obtained by taking a projective resolution of \(i_*(\underline{M})\) and applying the functor \(f^*\) to it. The homology sheaves of the resulting complex are the \(\underline{\Tor}_i^{\underline{O}_X}(\underline{M},\underline{O}_{X'})\). By 2.4 and 3.2, we have:
\[
\underline{\Tor}_i^{\underline{O}_X}(\underline{M},\underline{O}_{X'})
\xrightarrow{\sim}
\underline{M}\otimes_{\underline{O}_Y}\underline{\Tor}_i^{\underline{O}_X}(\underline{O}_{X'},\underline{O}_Y)
\xrightarrow{\sim}
\underline{M}\otimes_{\underline{O}_Y}\bigwedge^i\underline{F}.
\]

It follows that these homology sheaves are locally free of finite type over \(\underline{O}_{Y'}\), hence perfect over \(\underline{O}_{X'}\); consequently, the class in \(K^{\bullet}(X')\) of the complex \(\mathbb Lf^*(\mathbb Ri_*(\underline{M}))\) is the alternating sum of the classes of its homology sheaves, namely:
\[
f^*(i_*(x))
=\sum_{i\ge0}(-1)^i\operatorname{cl}
\bigl(j_*(\underline{M}\otimes_{\underline{O}_Y}\bigwedge^i\underline{F})\bigr)
=\sum_{i\ge0}(-1)^i
j_*\bigl(g^*(x)\lambda^i(F)\bigr).
\]

\sgaStableTarget{sga6.e07.lem.3-5}\textbf{Lemma 3.5.} \uline{With the hypotheses and notation of} 3.1, \uline{we have the following isomorphisms}:
\begin{enumerate}[label=\roman*)]
\item \uline{for every} \(n\ge0\), \uline{and every} \(i\ge1\), \(R^if_*\bigl(\underline{O}_{X'}(n)\bigr)=0\);
\item \uline{for every} \(n\ge0\), \(f_*\bigl(\underline{O}_{X'}(n)\bigr)=\underline{J}^n\).
\end{enumerate}

We may clearly suppose \(X\) affine, hence quasi-compact.

i) First recall the following lemma, which was proved in exp. III.

%%%%%%%%%%  scan idx 455  |  folio 442  %%%%%%%%%%
\SGAPrintPage{442}%
\textbf{Lemma.} \uline{Let} \(X'\xrightarrow{f}X\) \uline{be a projective and perfect morphism}, \uline{with} \(X\) \uline{quasi-compact}, \uline{and let} \(\underline{L}^\bullet\) \uline{be a perfect complex over} \(\underline{O}_{X'}\), \uline{acyclic in degrees} \(>0\). \uline{Then there exists an integer} \(n_0\) \uline{such that for} \(n\ge n_0\) \uline{the complex} \(\mathbb Rf_*\bigl(\underline{L}^\bullet\otimes_{\underline{O}_{X'}}\underline{O}_{X'}(n)\bigr)\) \uline{is acyclic in degrees} \(>0\).

Apply this lemma to the complex reduced to \(\underline{O}_{X'}\); there is therefore an integer \(n_0\) such that for \(n\ge n_0\), we have \(R^if_*(\underline{O}_{X'}(n))=0\) when \(i\ge1\). We then descend by induction to \(n=0\). For this, consider the exact sequence:
\[
0\longrightarrow
\underline{O}_{X'}(n+1)
\longrightarrow
\underline{O}_{X'}(n)
\longrightarrow
\underline{O}_{Y'}(n)
\longrightarrow0,
\]
and apply the functor \(f_*\) to it. We have \(\mathbb Rf_*(\underline{O}_{Y'}(n))=i_*\circ\mathbb Rg_*(\underline{O}_{Y'}(n))\). Since \(Y'=\bfP(\underline{N})\), \(\underline{O}_{Y'}(n)\) is \(g_*\)-acyclic for \(n\ge0\). Thus for every \(n\ge0\), and every \(i\ge1\):
\[
R^if_*(\underline{O}_{Y'}(n))=0,
\]
and consequently we have the exact sequences:
\[
R^if_*(\underline{O}_{X'}(n+1))
\longrightarrow
R^if_*(\underline{O}_{X'}(n))
\longrightarrow0,
\]
which complete the induction.

ii) We must show that the canonical morphism:
\[
\bigoplus_{n\ge0}\underline{J}^n
\longrightarrow
\bigoplus_{n\ge0}f_*\bigl(\underline{O}_{X'}(n)\bigr)
\]
is an isomorphism. Since the property is local, we may suppose that \(\underline{J}\) is defined by a sequence \(f_1,\ldots,f_d\) of sections of \(\underline{O}_X\) such that the corresponding homomorphism \((\underline{O}_X)^d\longrightarrow\underline{O}_X\) is regular. Put: \(\underline{S}=\bigoplus_{n\ge0}\underline{J}^n\); we define a surjective algebra homomorphism: \(\underline{O}_X[T_1,\ldots,T_d]\xrightarrow{v}\underline{S}\), by sending
%%%%%%%%%%  scan idx 456  |  folio 443  %%%%%%%%%%
\SGAPrintPage{443}%
\(T_i\) to \(f_i^{(1)}\). This homomorphism defines a closed immersion \(X'\longrightarrow\bfP_X^{d-1}\), regular by 1.8 ii).

By 1.8 iii), \(\underline{S}\) is perfect as a graded \(\underline{O}_X[T_1,\ldots,T_d]\)-Module. In particular, \(\underline{S}\) is finitely presented as a graded Module. Let therefore:
\[
\underline{L}''\longrightarrow \underline{L}'\longrightarrow \underline{S}\longrightarrow0
\]
be a finite graded presentation of \(\underline{S}\) over \(\underline{O}_X[T_1,\ldots,T_d]\). On \(\bfP_X^{d-1}\) we obtain the exact sequence:
\[
\underline{M}''\longrightarrow \underline{M}'\longrightarrow\underline{O}_{X'}\longrightarrow0,
\]
where \(\underline{M}''=\underline{\Proj}(\underline{L}'')\), \(\underline{M}'=\underline{\Proj}(\underline{L}')\), so these are locally free of finite type over \(\bfP_X^{d-1}\). We can therefore write, for every \(n\), the commutative diagram:
\[
\begin{tikzcd}[column sep=large,row sep=large]
\underline{L}''_n \arrow[r] \arrow[d,"\sim"'] &
\underline{L}'_n \arrow[r] \arrow[d,"\sim"'] &
\underline{J}^n \arrow[r] \arrow[d] &
0\\
p_*(\underline{M}''(n)) \arrow[r] &
p_*(\underline{M}'(n)) \arrow[r] &
p_*(\underline{O}_{X'}(n)) \arrow[r] &
0,
\end{tikzcd}
\]
where \(p\) is the structural morphism: \(\bfP_X^{d-1}\to X\), and \(\underline{L}''_n\), \(\underline{L}'_n\) are the degree \(n\) components of \(\underline{L}''\), \(\underline{L}'\). The upper row is exact by hypothesis; the first two vertical arrows are isomorphisms by the classical results for the direct images of \(\underline{O}_P(n)\), \(P\) being a projective bundle (EGA III, 2.1.15); finally, the last row is exact when \(n\) is sufficiently large, since the kernel of \(\underline{M}'\to\underline{O}_{X'}\) is a perfect \(\underline{O}_P\)-Module, and we can apply the lemma recalled above, which shows that this kernel, tensored with \(\underline{O}_P(n)\) for \(n\) sufficiently large, is \(p_*\)-acyclic; likewise for the kernel of \(\underline{M}''\to \underline{M}'\). It follows that for \(n\) sufficiently large, \(\underline{J}^n\longrightarrow p_*(\underline{O}_{X'}(n))\) is an isomorphism. Since \(p_*(\underline{O}_{X'}(n))=f_*(\underline{O}_{X'}(n))\), we therefore have
%%%%%%%%%%  scan idx 457  |  folio 444  %%%%%%%%%%
\SGAPrintPage{444}%
the result for \(n\ge n_0\).

It remains to descend to \(n=0\). We proceed by descending induction, using the commutative diagram:
\[
\begin{tikzcd}[column sep=large,row sep=large]
0 \arrow[r] \arrow[d] &
\underline{J}^{n+1} \arrow[r] \arrow[d,"\sim"'] &
\underline{J}^n \arrow[r] \arrow[d] &
\underline{J}^n/\underline{J}^{n+1} \arrow[r] \arrow[d,"\sim"'] &
0 \arrow[d]\\
0 \arrow[r] &
f_*(\underline{O}_{X'}(n+1)) \arrow[r] &
f_*(\underline{O}_{X'}(n)) \arrow[r] &
f_*(\underline{O}_{Y'}(n)) \arrow[r] &
0.
\end{tikzcd}
\]
The bottom row is exact because \(R^1f_*(\underline{O}_{X'}(n+1))=0\) for every \(n\ge0\), by i), already proved. Furthermore, we have:
\[
f_*(\underline{O}_{Y'}(n))
=i_*\bigl(g_*(\underline{O}_{Y'}(n))\bigr)
=i_*(\mathbf S(\underline{N})_n)
=\underline{J}^n/\underline{J}^{n+1}.
\]
The arrow on the right is therefore an isomorphism; the one on the left is an isomorphism by the induction hypothesis; hence the middle one is also an isomorphism.

\sgaStableTarget{sga6.e07.prop.3-6}\textbf{Proposition 3.6.} \uline{With the hypotheses and notation of} 3.1 \uline{and} 3.3, \uline{we have the relation}:
\sgaStableTarget{sga6.e07.eq.3-6-1}\[
f_*\circ f^*=\Id_{K^{\bullet}(X)}.
\]

Let \(x\in K^{\bullet}(X)\); by the projection formula, we can write:
\[
f_*(f^*(x))=f_*(1.f^*(x))=x.f_*(1).
\]
It therefore suffices to show that \(f_*(1)=1\).

Now 3.5 shows that the cohomology of \(\mathbb Rf_*(\underline{O}_{X'})\) vanishes in degrees \(>0\), and equals \(\underline{O}_X\) in degree \(0\). The proposition follows at once.

\sgaStableTarget{sga6.e07.thm.3-7}\textbf{Theorem 3.7.} \uline{Let} \(i:Y\to X\) \uline{be a regular closed immersion}, \(X'\) \uline{the scheme obtained by blowing up} \(Y\), \uline{and} \(Y'=X'\times_XY\), \uline{so that we have the commutative diagram}
%%%%%%%%%%  scan idx 458  |  folio 445  %%%%%%%%%%
\SGAPrintPage{445}%
\[
\begin{tikzcd}[column sep=large,row sep=large]
Y' \arrow[r,"j"] \arrow[d,"g"'] & X' \arrow[d,"f"]\\
Y \arrow[r,"i"] & X.
\end{tikzcd}
\]
\uline{Suppose that} \(X\) \uline{is quasi-compact and has an ample sheaf}. \uline{Then}, \uline{we have the exact sequence}:
\sgaStableTarget{sga6.e07.eq.3-7-0}\[
0\longrightarrow K^{\bullet}(Y)\xrightarrow{u}K^{\bullet}(Y')\times K^{\bullet}(X)
\xrightarrow{v}K^{\bullet}(X'),
\]
\uline{where we put}, \uline{for} \(y\in K^{\bullet}(Y)\):
\sgaStableTarget{sga6.e07.eq.3-7-u}\[
u(y)=\bigl(g^*(y).\lambda_{-1}(F),-i_*(y)\bigr),
\]
\uline{and for} \(y'\in K^{\bullet}(Y')\), \(x\in K^{\bullet}(X)\):
\sgaStableTarget{sga6.e07.eq.3-7-v}\[
v(y',x)=j_*(y')+f^*(x),
\]
\uline{with the notation of} 3.1 \uline{and} 3.3. \uline{Moreover}, \(u\) \uline{has a left inverse} \(u'\), \uline{defined by} \(u'(y',x)=g_*(y')\).

\uline{If in addition} \(X\) \uline{is noetherian}, \uline{and} \(Y\) \uline{regular}, \uline{then} \(v\) \uline{is surjective}, \uline{so we have the exact sequence}:
\sgaStableTarget{sga6.e07.eq.3-7-2}\[
0\longrightarrow K^{\bullet}(Y)\xrightarrow{u}K^{\bullet}(Y')\times K^{\bullet}(X)
\xrightarrow{v}K^{\bullet}(X')\longrightarrow0.
\]

First, \(u\) has the left inverse \(u'\), and is a fortiori injective. Indeed, we have:
\[
g_*\bigl(g^*(y).\lambda_{-1}(F)\bigr)
=y.g_*(\lambda_{-1}(F))
\]
by the projection formula. Furthermore, \(g_*(\lambda_{-1}(F))=1\) (VI 5.9), which gives the result.

Next, \(v\circ u=0\). Indeed, for \(y\in K^{\bullet}(Y)\) we have:
\[
v\circ u(y)
=v\bigl(g^*(y)\lambda_{-1}(F),-i_*(y)\bigr)
=j_*\bigl(g^*(y)\lambda_{-1}(F)\bigr)-f^*(i_*(y))=0,
\]
by 3.4.

It therefore remains to show that \(\Ker(v)\subset \Im(u)\). Let therefore \(x\in K^{\bullet}(X)\), \(y'\in K^{\bullet}(Y')\), be such that:
%%%%%%%%%%  scan idx 459  |  folio 446  %%%%%%%%%%
\SGAPrintPage{446}%
\[
j_*(y')+f^*(x)=0.
\]
We must find \(y\in K^{\bullet}(Y)\) such that:
\sgaStableTarget{sga6.e07.eq.3-7-1}\begin{equation}\tag{3.7.1}
x=-i_*(y);\qquad y'=g^*(y).\lambda_{-1}(F).
\end{equation}
Since \(u\) has a left inverse \(u'\), we necessarily have
\[
y=u'(y',x)=g_*(y').
\]
Thus we must show that \(y\) satisfies the relations (3.7.1).

To show that \(x=-i_*(y)\), start with the relation \(f^*(x)=-j_*(y')\) and apply \(f_*\) to it. We therefore obtain: \(f_*(f^*(x))=-f_*(j_*(y'))=-i_*(g_*(y'))=-i_*(y)\) by the definition of \(y\). By Proposition 3.6, however, \(f_*(f^*(x))=x\); hence the required relation.

To complete the proof, we must show that
\[
y'=g^*(g_*(y'))\lambda_{-1}(F).
\]
Put \(z=y'-g^*(g_*(y'))\lambda_{-1}(F)\). We obtain: \(g_*(z)=g_*(y')-g_*(y')\,g_*(\lambda_{-1}(F))=0\). On the other hand, by hypothesis we have: \(j_*(y')=-f^*(x)\), whence, by the preceding calculation, \(j_*(y')=f^*(i_*(y))\), and by Proposition 3.4: \(j_*(y')=j_*\bigl(g^*(y)\lambda_{-1}(F)\bigr)\); consequently, we also have \(j_*(z)=0\). Now \(\Ker(j_*)\cap\Ker(g_*)=0\). Indeed, if \(j_*(z)=0\), then \(j^*(j_*(z))=0\); but \(j^*(j_*(z))=z(1-L)\) (Prop. 2.7), and for \(z(1-L)=0\) to hold there must exist \(z'\in K^{\bullet}(Y)\) such that \(z=z'\lambda_{-1}(F)\) (VI,5.10); applying \(g_*\) gives \(z'=g_*(z)=0\), hence \(z=0\).

Suppose now that \(X\) is noetherian and \(Y\) regular. Then \(Y'\) is also regular, since it is a projective bundle over \(Y\), and for every point \(y'\) of \(Y'\), the local ring \(\underline{O}_{X',y'}\) of \(X'\) at \(y'\) is regular, since \(\underline{O}_{Y',y'}\) is the quotient of \(\underline{O}_{X',y'}\) by a regular element (see EGA \(0_{\mathrm{IV}}\) 17.1.8).

%%%%%%%%%%  scan idx 460  |  folio 447  %%%%%%%%%%
\SGAPrintPage{447}%
We know that for every regular scheme having an ample sheaf, the Grothendieck groups formed with locally free sheaves of finite type and with coherent sheaves are identical. In particular, this is the case for \(Y'\).

Let \(\underline{E}\) be a coherent \(\underline{O}_{X'}\)-Module supported on \(Y'\); then \(\underline{E}\) is a perfect \(\underline{O}_{X'}\)-Module. Indeed, it is perfect at the points of \(Y'\), since \(X'\) is regular at these points, and at the points of \(X'-Y'\) since it vanishes there.

If \(E\) is the class of \(\underline{E}\) in \(K^{\bullet}(X')\), then \(E\) belongs to the image of \(j_*\). Indeed, we can filter \(\underline{E}\) by the \(\underline{J}'^n.\underline{E}\). This filtration is finite, since, \(X'\) being noetherian, there is an integer \(n_0\) such that \(\underline{J}'^{\,n_0}.\underline{E}=0\); the quotients of this filtration are then \(\underline{O}_{Y'}\)-Modules, so their classes in \(K^{\bullet}(X')\) belong to the image of \(j_*\). Since \(E\) is the sum of the classes of these quotients, \(E\) also belongs to this image.

Let now \(x\in K^{\bullet}(X')\). We can find a perfect complex \(\underline{M}^\bullet\) whose class is \(x\). Consider the distinguished triangle:
\[
\begin{tikzcd}[column sep=large,row sep=large]
& \underline{N}^\bullet \arrow[dr] &\\
\underline{M}^\bullet \arrow[ur] && \mathbb Lf^*(\mathbb Rf_*(\underline{M}^\bullet)) \arrow[ll] ,
\end{tikzcd}
\]
where the horizontal arrow is the canonical homomorphism. This arrow is an isomorphism on the complement of \(Y'\), and consequently the cohomology sheaves of \(\underline{N}^\bullet\) are coherent sheaves supported on \(Y'\). By the preceding argument, these are perfect \(\underline{O}_{X'}\)-Modules, and their classes belong to the image of \(j_*\). Thus we have:
\[
x-f^*(f_*(x))\in\operatorname{Im}(j_*)\,;
\]
and consequently the homomorphism \(v\) is surjective.



% Macros for Exposé VII continuation
\providecommand{\Q}{\mathbb Q}
\providecommand{\bfP}{\mathbf P}
\providecommand{\gr}{\operatorname{gr}}
\providecommand{\Id}{\operatorname{Id}}
\providecommand{\Fil}{\operatorname{Fil}}
\providecommand{\Imf}{\operatorname{Im}}
\providecommand{\rank}{\operatorname{rank}}
\providecommand{\codim}{\operatorname{codim}}
\providecommand{\cJ}{\mathcal J}
\providecommand{\cI}{\mathcal I}
\providecommand{\cN}{\mathcal N}
\providecommand{\cE}{\mathcal E}
\providecommand{\cO}{\mathcal O}



%%%%%%%%%%  scan idx 461  |  folio 448  %%%%%%%%%%
\SGAPrintPage{448}%
\sgaStableTarget{sga6.e07.remark-3-8}\subsection*{Remark 3.8.}

When \(X\) is noetherian and \(Y\) regular, 3.7 determines the abelian group \(K^{\bullet}(X')\) entirely from the \(\lambda\)-rings \(K^{\bullet}(X)\), \(K^{\bullet}(Y)\), the maps \(i_*,i^*\) between these rings, and the element \(N\) of \(K^{\bullet}(Y)\): indeed, \(K^{\bullet}(Y')\) can be computed from \(K^{\bullet}(Y)\) and \(N\) by (VI 1.1). We shall see in the next section that the multiplicative structure of \(K^{\bullet}(X')\) is determined by the same data, as is the \(\lambda\)-structure after tensoring with \(\mathbb Q\) (this restriction is probably unnecessary: see 4.5).

\sgaStableTarget{sga6.e07.sec.4}\section*{4. Regular immersions and the filtration of \(K^{\bullet}\)}

Throughout this section we keep the notation and hypotheses of 3.1.

\sgaStableTarget{sga6.e07.lem.4-1}\textbf{Lemma 4.1.} \uline{Suppose that there exists on} \(Y\) \uline{a locally free sheaf of finite type} \(\underline{N}'\), \uline{such that}, \uline{if} \(N,N'\) \uline{are the classes in} \(K^{\bullet}(Y)\) \uline{of} \(\underline{N},\underline{N}'\), \uline{we have} \(N=N'+2\). \uline{Then the following congruence holds}:
\[
  \lambda_{-1}(F)\equiv 0 \ \mathrm{mod.}(1-L).
\]

Let \(d\) be the rank of \(\underline{N}\) (which we may suppose constant); the rank of \(\underline{F}\) is then \(d-1\). By definition:
\[
\begin{aligned}
\lambda_{-1}(F)
 &= \sum_{i\ge0}(-1)^i\lambda^i(F)\\
 &= (-1)^{d-1}\lambda^{d-1}(F-1)\\
 &= (-1)^{d-1}\lambda^{d-1}(F+L-1-L)\\
 &= (-1)^{d-1}\lambda^{d-1}(N-2+1-L).
\end{aligned}
\]
Now \(\lambda_t(1-L)=\lambda_t(1)/\lambda_t(L)=(1+t)/(1+Lt)\). Thus for every \(i\ge1\), \(\lambda^i(1-L)\) is divisible by \(1-L\). Consequently:
\[
\begin{aligned}
\lambda_{-1}(F)
  &\equiv (-1)^{d-1}\lambda^{d-1}(N-2) \qquad\mathrm{mod.}(1-L)\\
  &\equiv(-1)^{d-1}\lambda^{d-1}(N') \qquad\mathrm{mod.}(1-L).
\end{aligned}
\]
%%%%%%%%%%  scan idx 462  |  folio 449  %%%%%%%%%%
\SGAPrintPage{449}%
Now \(\underline{N}'\) has rank \(d-2\), so \(\lambda^{d-1}(N')=0\), which gives the result.

\sgaStableTarget{sga6.e07.sec.4.2}\subsection*{\normalfont 4.2.}

Let \(i:Y\to X\) be a regular immersion. Let \(r\) be an integer \(\ge1\), and \(s:X\longrightarrow \mathbf P_X^r\) the section defined by the homomorphism: \((\underline{O}_X)^{r+1}\to\underline{O}_X\), passage to the quotient by the subsheaf generated by the last \(r\) coordinates. If \(\mathbf P_X^r\) is defined by the graded \(\underline{O}_X\)-Algebra \(\underline{O}_X[T_0,\ldots,T_r]\), the ideal \(\underline{I}\) defining the section \(s\) is the ideal generated by the sections \(T_1,\ldots,T_r\), so \(\underline{I}/\underline{I}^2\) is isomorphic to \((\underline{O}_X[T_0])^r\), and hence the corresponding ideal of \(\underline{O}_P\) is isomorphic to \((\underline{O}_X)^r\).

Since the section \(s\) is clearly a regular immersion, the same holds for \(s\circ i\), and we can write the corresponding exact sequence of conormal sheaves (1.7). We thus obtain the following equality in \(K^{\bullet}(Y)\), writing \(\underline{N}\) for the conormal sheaf of \(i\), and \(\underline{N}_1\) for the conormal sheaf of \(s\circ i\):
\[
  N_1=N+r.
\]
Thus for every regular immersion \(i\), composing it with a section: \(X\to\mathbf P_X^2\), or: \(X\to\mathbf P_{\mathbf P_X^1}^{1}\), gives a regular immersion satisfying the property of Lemma 4.1.

\sgaStableTarget{sga6.e07.thm.4-3}\textbf{Theorem 4.3.} \uline{Let} \(i:Y\to X\) \uline{be a regular closed immersion}, \(N\) \uline{the class in} \(K^{\bullet}(Y)\) \uline{of the conormal sheaf of} \(Y\) \uline{in} \(X\). \uline{Suppose that} \(X\) \uline{is quasi-compact and has an ample sheaf}. \uline{Then}, \uline{with the notation of} 3.3, \uline{we have for every} \(y\in K^{\bullet}(Y)\), \uline{and every} \(p\ge1\):
\sgaStableTarget{sga6.e07.eq.4-3-1}\begin{equation}\tag{4.3.1}
 i_*\bigl(\lambda^p(N,y)\bigr)\equiv \lambda^p(i_*(y))
 \qquad\text{\uline{mod}. \uline{torsion}},
\end{equation}
\(\lambda^p(N,y)\) \uline{being the element defined in} (V 5.3).

i) First suppose that the hypothesis of 4.1 is satisfied; in this case we shall show that \(i_*\bigl(\lambda^p(N,y)\bigr)=\lambda^p(i_*(y))\). Since \(f_*\circ f^*=\Id\), this
%%%%%%%%%%  scan idx 463  |  folio 450  %%%%%%%%%%
\SGAPrintPage{450}%
relation is equivalent to the relation:
\[
  f^*\bigl(i_*(\lambda^p(N,y))\bigr)=
  f^*\bigl(\lambda^p(i_*(y))\bigr).
\]
Now we have:
\sgaStableTarget{sga6.e07.eq.3-4}\[
\begin{aligned}
 f^*\bigl(\lambda^p(i_*(y))\bigr)
 &=\lambda^p\bigl(f^*(i_*(y))\bigr)\\
 &=\lambda^p\bigl(j_*(g^*(y)\lambda_{-1}(F))\bigr)
\qquad(3.4)
\end{aligned}
\]
On the other hand:
\sgaStableTarget{sga6.e07.eq.3-4.occ2}\[
  f^*\bigl(i_*(\lambda^p(N,y))\bigr)
  =
  j_*\bigl(g^*(\lambda^p(N,y))\lambda_{-1}(F)\bigr)\ (3.4).
\]
Since \(g^*(N)=F+L\), we can apply formula (V (5.4.2)):
\[
  g^*(\lambda^p(N,y))\lambda_{-1}(F)
  =
  \lambda^p(g^*(N),g^*(y))\lambda_{-1}(F)
  =
  \lambda^p(L,g^*(y)\lambda_{-1}(F)).
\]
Putting: \(z=g^*(y)\lambda_{-1}(F)\), we are therefore reduced to proving the relation:
\[
  j_*\bigl(\lambda^p(L,z)\bigr)=\lambda^p(j_*(z)),
\]
that is, the analogous relation for the immersion \(j\) and the element \(z\). By hypothesis, \(\lambda_{-1}(F)\) is a multiple of \((1-L)\); we can find \(z'\) such that \(z=z'(1-L)\). Then \(z=j^*(u')\) with \(u'=j_*(z')\) (2.7); since \(L=j^*(J')\), we finally have to verify:
\[
  j_*\bigl(j^*(\lambda^p(J',u'))\bigr)=
  \lambda^p\bigl(j_*(j^*(u'))\bigr).
\]
The exact sequence
\[
  0\longrightarrow \underline{J}'\longrightarrow \underline{O}_{X'}
  \longrightarrow \underline{O}_{Y'}\longrightarrow 0
\]
gives the relation:
\sgaStableTarget{sga6.e07.eq.4-3-2}\begin{equation}\tag{4.3.2}
  j_*(1)=1-J'.
\end{equation}
Thus, applying the projection formula, the relation to be verified finally becomes:
%%%%%%%%%%  scan idx 464  |  folio 451  %%%%%%%%%%
\SGAPrintPage{451}%
\[
  \lambda^p(J',u')(1-J')=\lambda^p(u'(1-J')).
\]
Since \(1-J'=\lambda_{-1}(J')\), this is the formula defining \(\lambda^p(J',u')\).

ii) We turn to the general case. Compose \(i\) with a section: \(X\to\mathbf P_X^1\), then with a section: \(\mathbf P_X^1\to\mathbf P_{\mathbf P_X^1}^1\), so that the composite immersion satisfies the condition of 4.1. (see 4.2). Proposition 4.3 is then true for the composite immersion by the preceding argument; it therefore remains to ``descend'' to \(X\). We descend first to \(\mathbf P_X^1\), then to \(X\), so that the proof consists in applying the following lemma twice (due to J.P. Jouanolou):

\sgaStableTarget{sga6.e07.lem.unnamed}\textbf{Lemma}: \uline{Under the conditions of} 4.3, \uline{let} \(s:X\to\mathbf P_X^1\) \uline{be the section defined in} 4.2. \uline{If Proposition} 4.3 \uline{is true for} \(s\circ i\), \uline{it is true for} \(i\).

Supposing 4.3 true for \(s\circ i\), we have for every \(p\ge1\), and every \(y\in K^{\bullet}(Y)\):
\sgaStableTarget{sga6.e07.eq.4-3-3}\begin{equation}\tag{4.3.3}
 (s\circ i)_*\bigl(\lambda^p(N+1,y)\bigr)
 \equiv
 \lambda^p\bigl((s\circ i)_*(y)\bigr)
 \quad\text{mod.\ torsion}.
\end{equation}

The section \(s\) is defined by an ideal isomorphic to \(\underline{O}_{\mathbf P}(-1)\) (4.2); write \(\xi\) for its class in \(K^{\bullet}(\mathbf P_X^1)\). We then have the following relation for every \(x\in K^{\bullet}(X)\):
\[
  s_*(x)=p^*(x)(1-\xi),
\]
where \(p\) is the structural morphism \(\mathbf P_X^1\to X\). Indeed, since \(p\circ s=\Id_X\), we have \(x=s^*(p^*(x))\), whence \(s_*(x)=s_*(s^*(p^*(x)))=p^*(x)s_*(1)\); since \(s_*(1)=1-\xi\), we have the required relation:

Relation (4.3.3) then becomes:
\[
\begin{aligned}
(s\circ i)_*\bigl(\lambda^p(N+1,y)\bigr)
&\equiv \lambda^p\bigl(p^*(i_*(y))(1-\xi)\bigr)\quad\text{mod.\ torsion}\\
&\equiv \lambda^p(\xi,p^*(i_*(y)))(1-\xi)\quad\text{mod.\ torsion}.
\end{aligned}
\]
%%%%%%%%%%  scan idx 465  |  folio 452  %%%%%%%%%%
\SGAPrintPage{452}%
Now \(\lambda^p(\xi,p^*(i_*(y)))-\lambda^p(1,p^*(i_*(y)))\) is divisible by \(1-\xi\): to see this, we can work in the universal \(\lambda\)-ring used to define the \(\lambda^p(\xi,X)\) (V 5); then \(\lambda^p(\xi,X)-\lambda^p(1,X)\) is a polynomial in \(\xi\), vanishing for \(\xi=1\), and hence divisible by \(\xi-1\). Furthermore, the dependence equation of \(\xi\) over \(K^{\bullet}(X)\) is \((1-\xi)^2=0\) (VI (1.13.1)). We therefore obtain:
\[
\begin{aligned}
(s\circ i)_*\bigl(\lambda^p(N+1,y)\bigr)
&\equiv \lambda^p(1,p^*(i_*(y)))(1-\xi) \quad\text{mod.\ torsion}\\
&\equiv p^*(\lambda^p(1,i_*(y)))(1-\xi) \quad\text{mod.\ torsion}.
\end{aligned}
\]
Applying \(p_*\), we obtain:
\sgaStableTarget{sga6.e07.eq.4-3-4}\begin{equation}\tag{4.3.4}
  i_*\bigl(\lambda^p(N+1,y)\bigr)
  \equiv \lambda^p(1,i_*(y)) \quad\text{mod.\ torsion},
\end{equation}
since \(p_*(\xi)=0\).

By the formula in (V 5.6), applied with \(L=1\), \(\lambda^p(N+1,y)\) is the coefficient of \(t^p\) in the series:
\[
  \frac{\sum_{k\ge1}\lambda^k(N,y)kt^k}
       {1+\sum_{k\ge1}\lambda^k(N,y)\lambda_{-1}(N)t^k}.
\]

Similarly, the same formula applied with \(N=0\) and \(L=1\) shows that \(\lambda^p(1,y)\) is the coefficient of \(t^p\) in the series:
\[
  \frac{\sum_{k\ge1}\lambda^k(y)kt^k}
       {\sum_{k\ge0}\lambda^k(y)t^k}.
\]

The congruences (4.3.4) can therefore be written:
\[
i_*\left(
  \frac{\sum_{k\ge1}k\lambda^k(N,y)t^k}
       {1+\sum_{k\ge1}\lambda^k(N,y)\lambda_{-1}(N)t^k}
\right)
\equiv
  \frac{\sum_{k\ge1}k\lambda^k(i_*(y))t^k}
       {\sum_{k\ge0}\lambda^k(i_*(y))t^k}
       \quad\text{mod.\ torsion}
\]
%%%%%%%%%%  scan idx 466  |  folio 453  %%%%%%%%%%
\SGAPrintPage{453}%
Applying formula 2.8, we obtain:
\[
  \frac{\sum_{k\ge0}k\,i_*(\lambda^k(N,y))t^k}
       {1+\sum_{k\ge1} i_*(\lambda^k(N,y))t^k}
  \equiv
  \frac{\sum_{k\ge0}k\lambda^k(i_*(y))t^k}
       {\sum_{k\ge0}\lambda^k(i_*(y))t^k}
       \quad\text{mod.\ torsion}
\]
Putting:
\[
  u(t)=1+\sum_{k\ge1}i_*(\lambda^k(N,y))t^k,\qquad
  v(t)=\sum_{k\ge0}\lambda^k(i_*(y))t^k,
\]
we can also write this last relation as:
\[
  \frac{d}{dt}\!\left(\frac{u(t)}{v(t)}\right)\equiv0
  \quad\text{mod.\ torsion}.
\]
Since the constant term of \(u(t)/v(t)\) is \(1\), and we are writing congruences modulo torsion, we can integrate and write the congruence \(u(t)\equiv v(t)\). This gives the stated congruences.

\sgaStableTarget{sga6.e07.cor.4-4}\textbf{Corollary 4.4.} \uline{With the notation and hypotheses of} 4.3, \uline{we have for every} \(p\ge1\), \uline{and every} \(y\in K^{\bullet}(Y)\):
\[
  i_*\bigl(\gamma^p(N,y)\bigr)\equiv \gamma^p(i_*(y))
  \quad\text{mod.\ torsion}.
\]

The corollary follows immediately from 4.3, using the expression of the \(\gamma^p\) as linear combinations with integer coefficients of the \(\lambda^p\).

\sgaStableTarget{sga6.e07.rem.4-5}\textbf{Remark 4.5.} The proof of 4.3 itself suggests that (4.3.1) should hold not only modulo torsion, but also as an equality. In [RRR], the equality was proved in characteristic \(0\) by means of a geometric characterization of the \(\lambda^p(N,x)\). An extension of this method to all characteristics would probably be possible if one knew how to define exterior-power operations on perfect complexes (cf. XIV).

%%%%%%%%%%  scan idx 467  |  folio 454  %%%%%%%%%%
\SGAPrintPage{454}%
\sgaStableTarget{sga6.e07.thm.4-6}\textbf{Theorem 4.6.}
\uline{Let} \(X\) \uline{be a quasi-compact scheme having an ample sheaf}, \(d\) \uline{an integer} \(>0\) \uline{and} \(i:Y\to X\) \uline{a regular immersion of codimension} \(d\). \uline{Then}, \uline{for every} \(k\), \uline{we have}:
\[
 i_*\bigl((\Fil^kK^{\bullet}(Y))\otimes_{\mathbb Z}\mathbb Q\bigr)
 \subset
 (\Fil^{k+d}K^{\bullet}(X))\otimes_{\mathbb Z}\mathbb Q.
\]

First observe that \(i_*(K^{\bullet}(Y))\subset \Fil^1K^{\bullet}(X)\). Indeed, \(X-Y\) is dense in \(X\): otherwise, \(Y\) would contain an affine open subset, and on this open subset the ideal \(\underline{J}\) defining \(Y\) would be a nil ideal, hence a locally nilpotent ideal since it is of finite type. This is impossible, however, since the \(\underline{J}^n/\underline{J}^{n+1}\) are locally free \(\underline{O}_Y\)-Modules of rank \(\ne0\). The augmentation in \(K^{\bullet}(X)\) of an element of \(i_*(K^{\bullet}(Y))\) is a locally constant function on \(X\); since it vanishes on \(X-Y\), it vanishes on all of \(X\).

For every abelian group \(A\), we put \(A_{\mathbb Q}=A\otimes_{\mathbb Z}\mathbb Q\). Let \(x\in\Fil^kK^{\bullet}(Y)_{\mathbb Q}\); we can write \(x\) in the form:
\[
  x=\sum a\,\gamma^{i_1}(x_1)\cdots\gamma^{i_n}(x_n),
\]
with \(\epsilon(x_i)=0\) for every \(i\), \(a\in\mathbb Q\), and \(i_1+\cdots+i_n\ge k\). The \(x_i\) can be expressed as linear combinations with integer coefficients of the classes \(E_\alpha\) of finitely many locally free \(\underline{O}_Y\)-Modules of finite type \(\underline{E}_\alpha\). Let \(R\) be the sub-\(\mathbb Q\)-\(\lambda\)-algebra of \(K^{\bullet}(Y)_{\mathbb Q}\) generated by the \(E_\alpha\) and by \(N\). We then have:
\[
  x\in\Fil^kR.
\]
Furthermore, the \(\lambda\)-ring filtration of \(R\) is discrete (VI 6.2); let \(m\) satisfy \(\Fil^mR=0\). Finally, by decomposing \(Y\) into a finite sum of disjoint open subsets and carrying out the proof on each open subset (since \(K^{\bullet}\) commutes with finite direct sums), we may suppose that the \(\underline{E}_\alpha\) have constant rank \(r_\alpha\). The augmentation ring of \(R\) is then \(\mathbb Q\).

%%%%%%%%%%  scan idx 468  |  folio 455  %%%%%%%%%%
\SGAPrintPage{455}%
We shall show that for every \(i\), we have:
\[
  i_*(\Fil^iR)\subset \Fil^{i+d}(K^{\bullet}(X)_{\mathbb Q}).
\]
This will imply in particular:
\[
  i_*(x)\in\Fil^{k+d}K^{\bullet}(X)_{\mathbb Q},
\]
and hence prove the theorem.

The proof is by descending induction on \(i\), the case \(i=m\) being trivial. Let \(y\in\Fil^iR\), and put, for every \(j\):
\[
  a_j=(-1)^{j-1}(j-1)!.
\]
Then by V 6.10,
\[
  \gamma^{d+i}(N,y)-a_{d+i}y\in\Fil^{i+1}R.
\]

We can therefore apply the induction hypothesis, supposing that:
\[
  i_*(\Fil^{i+1}R)\subset \Fil^{i+d+1}K^{\bullet}(X)_{\mathbb Q}.
\]
Consequently, we have:
\[
  i_*\bigl(\gamma^{d+i}(N,y)-a_{d+i}y\bigr)
  \in\Fil^{i+d+1}K^{\bullet}(X)_{\mathbb Q}.
\]
By 4.4, this implies:
\[
  \gamma^{d+i}\bigl(i_*(y)\bigr)-a_{d+i}i_*(y)
  \in\Fil^{i+d+1}K^{\bullet}(X)_{\mathbb Q}.
\]
Since \(i_*(y)\) has zero augmentation, \(\gamma^{d+i}\bigl(i_*(y)\bigr)\in\Fil^{d+i}K^{\bullet}(X)_{\mathbb Q}\).
Since we are calculating modulo torsion, we therefore find:
\[
  i_*(y)\in\Fil^{i+d}K^{\bullet}(X)_{\mathbb Q}.
\]

\sgaStableTarget{sga6.e07.rem.4-6-1}\textbf{Remark 4.6.1.} We see that, even if 4.3 were true without tensoring with \(\mathbb Q\), the proof above works only modulo torsion. For a counterexample, see XIV 4.6.

%%%%%%%%%%  scan idx 469  |  folio 456  %%%%%%%%%%
\SGAPrintPage{456}%
\sgaStableTarget{sga6.e07.sec.4.7}\subsection*{\normalfont 4.7.}

We return to the situation of 3.7, where we had the commutative diagram:
\[
\begin{tikzcd}[column sep=large,row sep=large]
Y' \arrow[r,"j"] \arrow[d,"g"'] & X' \arrow[d,"f"]\\
Y \arrow[r,"i"] & X,
\end{tikzcd}
\]
where \(i\) is a regular immersion, \(X'\) the \(X\)-scheme obtained by blowing up \(Y\), and \(Y'=X'\times_XY\). Suppose that \(X\) is noetherian with an ample sheaf, and \(Y\) regular. We saw in 3.7 that we then have the exact sequence:
\[
0\longrightarrow K^{\bullet}(Y)\xrightarrow{u}
K^{\bullet}(Y')\times K^{\bullet}(X)\xrightarrow{v}
K^{\bullet}(X')\longrightarrow0.
\]
We shall show that this exact sequence describes the whole situation formed by the \(K^{\bullet}\) of \(X,Y,X',Y'\) from the data of the \(\lambda\)-rings \(K^{\bullet}(X)\), \(K^{\bullet}(Y)\), the maps \(i^*\) and \(i_*\), and the element \(N\) of \(K^{\bullet}(Y)\). (at least up to torsion, for the time being).

Let therefore \(K,K'\) be two \(\lambda\)-rings, \(i_*:K\to K'\) an additive homomorphism, \(i^*:K'\to K\) a \(\lambda\)-homomorphism, \(d\) an integer, and \(N\) an element of \(K\) such that \(\lambda^i(N)=0\) for \(i>d\) and \(\lambda^d(N)\) is invertible. Suppose moreover that the following properties hold:
\begin{enumerate}[label=\roman*)]
\item \(i_*\) is \(K'\)-linear (``projection formula'');
\item \(i^*\circ i_*\) is the homomorphism ``multiplication by \(\lambda_{-1}(N)\)'';
\item for every \(n\ge1\), and every \(y\in K\), \(i_*\bigl(\lambda^n(N,y)\bigr)=\lambda^n(i_*(y))\) (relation (4.3.1)).
\end{enumerate}

First we use VI.1.1 to construct a \(K\)-\(\lambda\)-algebra \(K_1\) which plays the role of \(K^{\bullet}(Y')\). Let \(L\) be an indeterminate, and \(K[L]\) the \(K\)-\(\lambda\)-algebra generated by the generator \(L\) subject to the relations \(\lambda^i(L)=0\) for \(i\ge2\) (V 4.9). Put:
\[
  P(L)=\lambda^d(N-L)=\sum_{i=0}^{d}(-1)^i\lambda^{d-i}(N)L^i.
\]
%%%%%%%%%%  scan idx 470  |  folio 457  %%%%%%%%%%
\SGAPrintPage{457}%
The ideal generated by \(P\) is a \(\lambda\)-ideal. To see this, first observe that the relations \(\lambda^i(N)=0\) for \(i>d\) imply that the \(\lambda^j(N-L)\) for \(j\ge d\) are multiples of \(\lambda^d(N-L)\). Working in \(1+\widehat{K[L]/(P)[[t]]}^+\), the image of \(\lambda_t(N-L)\) is a polynomial of degree \(<d\), and a universal argument then shows that its \(\lambda^i\) for \(i\ge d\) equal \(1\) (cf. V 4.8). Since in \(1+\widehat{K[L][[t]]}^+\) we have \(\lambda_t(\lambda^d(N-L))=\lambda^d(\lambda_t(N-L))\), we deduce that the \(\lambda^i(\lambda^d(N-L))\) for \(i\ge1\) are multiples of \(P\). By the relations V 2.4.2, this implies that the ideal generated by \(P\) is a \(\lambda\)-ideal.

Put then \(K_1=K[L]/(P)\). We denote the natural homomorphism \(K\to K_1\) by \(g^*\); it is a \(\lambda\)-homomorphism. Furthermore, \(K_1\) is a free \(K\)-module with basis the elements \(1,L,\ldots,L^{d-1}\). The relation \(\lambda^d(N)=\lambda^{d-1}(N-L).L\) and the invertibility hypothesis on \(\lambda^d(N)\) imply that \(L\) is invertible. Consequently, \(K_1\) has as a basis over \(K\) the elements \(1,L^{-1},\ldots,L^{-d+1}\); we then define a map \(g_*:K_1\to K\) by assigning to every element \(y'=y_0+y_1L^{-1}+\ldots+y_{d-1}L^{-d+1}\) of \(K_1\) the element \(y_0\) of \(K\). Clearly the map thus defined is an additive, \(K\)-linear homomorphism, satisfying \(g_*\circ g^*=\Id\).

Next we endow \(K_1\times K'\) with a \(\lambda\)-ring structure, putting:
\sgaStableTarget{sga6.e07.eq.4-7-1}\sgaStableTarget{sga6.e07.eq.4-7-2}\begin{align}
 (y'_1,x_1)*(y'_2,x_2)
   &=(y'_1y'_2(1-L)+x_1y'_2+x_2y'_1,\; x_1x_2); \tag{4.7.1}\\
 \lambda^n_*(y',x)
   &=\left(\sum_{p=1}^{n}\lambda^{\,n-p}(x)\lambda^p(L,y'),\;
     \lambda^n(x)\right). \tag{4.7.2}
\end{align}
In these formulas, \(K_1\) is regarded as a \(K'\)-algebra through the homomorphism \(g^*\circ i^*\).
Next, for \(y_1,y_2\in K\) and \(n\ge1\), put:
\sgaStableTarget{sga6.e07.eq.4-7-3}\sgaStableTarget{sga6.e07.eq.4-7-4}\begin{align}
 y_1*y_2&=-\,y_1y_2\lambda_{-1}(N); \tag{4.7.3}\\
 \lambda^n_*(y_1)&=-\,\lambda^n(N,-y_1); \tag{4.7.4}
\end{align}
%%%%%%%%%%  scan idx 471  |  folio 458  %%%%%%%%%%
\SGAPrintPage{458}%
and define a homomorphism:
\[
  u:K\longrightarrow K_1\times K'
\]
by putting:
\[
  u(y)=\bigl(g^*(y).\lambda_{-1}(F),-i_*(y)\bigr),
\]
with \(F=N-L\).

\sgaStableTarget{sga6.e07.prop.4-7-5}\textbf{Proposition 4.7.5.} \uline{The operations defined by} (4.7.1) \uline{and} (4.7.2) \uline{endow} \(K_1\times K'\) \uline{with a} \(\lambda\)-\uline{ring structure}, \uline{and} \(u\) \uline{is an injective homomorphism identifying} \(K\), \uline{equipped with the operations defined by} (4.7.3) \uline{and} (4.7.4), \uline{with a} \(\lambda\)-\uline{ideal of} \(K_1\times K'\).

It is immediate to verify that \(K_1\times K'\) is a \(\lambda\)-ring: one can refer, for example, to V 5.5, where an entirely analogous construction is given. The second part of the proposition is also proved without difficulty using the preceding relations, the relation \(i_*(y_1y_2\lambda_{-1}(N))=i_*(y_1).i_*(y_2)\) which follows immediately from the hypotheses, the relation \(g_*(\lambda_{-1}(F))=1\) (see VI 5.9), and the relations of V 5.

Put then:
\[
  K'_1=(K_1\times K')/K
\]
and denote by \(v\) the canonical homomorphism \(K_1\times K'\to K'_1\). Regard \(K'_1\) as a \(\lambda\)-ring through the quotient structure.

Next define additive homomorphisms:
\[
\begin{array}{llll}
j_*:K_1\to K'_1; & j^*:K'_1\to K_1; \\
f_*:K'_1\to K'; & f^*:K'\to K'_1 .
\end{array}
\]

The homomorphism \(j_*\) is the composite \(K_1\to K_1\times K'\to K'_1\). Likewise, the homomorphism \(f^*\) is the composite \(K'\to K_1\times K'\to K'_1\). We then have:
\[
  v=j_*+f^*.
\]
%%%%%%%%%%  scan idx 472  |  folio 459  %%%%%%%%%%
\SGAPrintPage{459}%
Next define a homomorphism \(K_1\times K'\to K_1\) by assigning to \((y',x)\) the element \(y'(1-L)+g^*\circ i^*(x)\); this homomorphism vanishes on \(K\), and consequently defines a homomorphism \(j^*:K'_1\longrightarrow K_1\). Finally, define a homomorphism \(K_1\times K'\to K'\) by assigning to \((y',x)\) the element \(i_*\circ g_*(y')+x\); this homomorphism vanishes on \(K\) and therefore defines a homomorphism \(f_*:K'_1\longrightarrow K'\).

\sgaStableTarget{sga6.e07.prop.4-7-6}\textbf{Proposition 4.7.6.}
\begin{enumerate}[label=\roman*)]
\item \(j^*\) \uline{and} \(f^*\) \uline{are} \(\lambda\)-\uline{homomorphisms};
\item \(j_*\) \uline{is} \(K'_1\)-\uline{linear and} \(f_*\) \uline{is} \(K'\)-\uline{linear};
\item \(j^*\circ j_*\) \uline{is the homomorphism} ``\uline{multiplication by} \(1-L\)'';
\item \uline{for every} \(y'\in K_1\), \uline{and every} \(p\ge1\), \(j_*\bigl(\lambda^p(L,y')\bigr)=\lambda^p(j_*(y'))\);
\item \(f_*\circ f^*=\Id\).
\end{enumerate}

We need only set the machinery in motion.

Returning to the geometric situation recalled above, we see that \(K^{\bullet}(X')\) is constructed from \(K^{\bullet}(Y)\), \(K^{\bullet}(X)\), the maps \(i^*\) and \(i_*\), and the element \(N\) of \(K^{\bullet}(Y)\) by the preceding method; the \(\lambda\)-operations on \(K^{\bullet}(X')\), however, agree with the \(\lambda\)-operations defined above on \(K'_1\) only up to torsion, as long as relation (4.3.1) is proved only modulo torsion.

\sgaStableTarget{sga6.e07.thm.4-8}\textbf{Theorem 4.8.} \uline{Let} \(i:Y\to X\) \uline{be a regular closed immersion of constant codimension} \(d\), \(X'\) \uline{the scheme obtained by blowing up} \(Y\), \uline{and} \(Y'=X'\times_XY\), \uline{so that we have the commutative diagram}:
\[
\begin{tikzcd}[column sep=large,row sep=large]
Y' \arrow[r,"j"] \arrow[d,"g"'] & X' \arrow[d,"f"]\\
Y \arrow[r,"i"] & X.
\end{tikzcd}
\]
\uline{Suppose that} \(X\) \uline{is quasi-compact and has an ample sheaf}. \uline{Then for every} \(k\) \uline{we have the exact sequence}:
%%%%%%%%%%  scan idx 473  |  folio 460  %%%%%%%%%%
\SGAPrintPage{460}%
\[
0\longrightarrow
\operatorname{Gr}^{k-d}K^{\bullet}(Y)_{\mathbb Q}
\xrightarrow{u^{\gr}}
\operatorname{Gr}^{k-1}K^{\bullet}(Y')_{\mathbb Q}
\times
\operatorname{Gr}^{k}K^{\bullet}(X)_{\mathbb Q}
\xrightarrow{v^{\gr}}
\operatorname{Gr}^{k}K^{\bullet}(X')_{\mathbb Q}
\]
\uline{where we put}, \uline{for} \(y\in \operatorname{Gr}^{k-d}K^{\bullet}(Y)_{\mathbb Q}\):
\[
  u^{\gr}(y)=
  \bigl(g^{\gr}(y).(-1)^{d-1}c^{d-1}(F),-i_{\gr}(y)\bigr),
\]
\uline{and for} \(y'\in \operatorname{Gr}^{k-1}K^{\bullet}(Y')_{\mathbb Q}\) \uline{and} \(x\in \operatorname{Gr}^{k}K^{\bullet}(X)_{\mathbb Q}\):
\[
  v^{\gr}(y',x)=j_{\gr}(y')+f^{\gr}(x),
\]
\uline{with the notation of} 3.1 \uline{and} 3.3, \uline{and} \(c^{d-1}(F)\) \uline{denoting the} \((d-1)\)-\uline{st Chern class of} \(F\) (V 6.7). \uline{Moreover}, \(u^{\gr}\) \uline{has a left inverse} \(u'^{\gr}\), \uline{defined by} \(u'^{\gr}(y',x)=g_{\gr}(y')\).

\uline{If in addition} \(X\) \uline{is noetherian and} \(Y\) \uline{regular}, \uline{we have the exact sequence}:
\[
0\longrightarrow
\operatorname{Gr}^{k-d}K^{\bullet}(Y)_{\mathbb Q}
\xrightarrow{u^{\gr}}
\operatorname{Gr}^{k-1}K^{\bullet}(Y')_{\mathbb Q}
\times
\operatorname{Gr}^{k}K^{\bullet}(X)_{\mathbb Q}
\xrightarrow{v^{\gr}}
\operatorname{Gr}^{k}K^{\bullet}(X')_{\mathbb Q}
\longrightarrow0
\]

\uline{Proof}. Observe that, when \(Y\) is regular, one could give a purely algebraic proof in terms of the theory of \(\lambda\)-rings, using the description of the situation in 4.7. We shall instead follow the proof of 3.7.

We know that the ``inverse-image'' homomorphisms \(f^*,g^*\), being \(\lambda\)-homomorphisms, are compatible with the \(\lambda\)-ring filtrations, and consequently define homomorphisms on the associated gradeds; hence the homomorphisms:
\[
  f^{\gr}:\operatorname{Gr}^kK^{\bullet}(X)\to\operatorname{Gr}^kK^{\bullet}(X');
  \qquad
  g^{\gr}:\operatorname{Gr}^{k-d}K^{\bullet}(Y)\to\operatorname{Gr}^{k-d}K^{\bullet}(Y'),
\]
and the analogous homomorphisms obtained after tensoring with \(\mathbb Q\). Furthermore, 4.6 shows that after tensoring with \(\mathbb Q\), the homomorphisms \(i_*,j_*\) pass to the associated gradeds, and therefore define:
\[
  i_{\gr}:\operatorname{Gr}^{k-d}K^{\bullet}(Y)_{\mathbb Q}
  \to \operatorname{Gr}^{k}K^{\bullet}(X)_{\mathbb Q};
  \qquad
  j_{\gr}:\operatorname{Gr}^{k-1}K^{\bullet}(Y')_{\mathbb Q}
  \to \operatorname{Gr}^{k}K^{\bullet}(X')_{\mathbb Q}.
\]
%%%%%%%%%%  scan idx 474  |  folio 461  %%%%%%%%%%
\SGAPrintPage{461}%
Finally, observe that \(\lambda_{-1}(F)=(-1)^{d-1}\gamma^{d-1}(F)\), so the morphisms \(u^{\gr}\) and \(v^{\gr}\) are obtained by passage to the quotient from the morphisms \(u\) and \(v\) of 3.7. Thus we have \(v^{\gr}\circ u^{\gr}=0\).

By VI 5.8, \(g_*(\Fil^{k-1}K^{\bullet}(Y'))\subset \Fil^{k-d}K^{\bullet}(Y)\) for every \(k\), since \(Y'=\mathbf P(\underline{N})\), with \(\underline{N}\) of rank \(d\). We therefore obtain a homomorphism:
\[
  g_{\gr}:\operatorname{Gr}^{k-1}K^{\bullet}(Y')
  \to\operatorname{Gr}^{k-d}K^{\bullet}(Y),
\]
and an analogous homomorphism after tensoring with \(\mathbb Q\). The homomorphism \(u'^{\gr}\) is therefore obtained from \(u'\) (3.7) by passage to the associated graded. Since \(u'\) is a left inverse of \(u\), \(u'^{\gr}\) is a left inverse of \(u^{\gr}\), proving at the same time the injectivity of \(u^{\gr}\).

To show that \(\Ker(v^{\gr})=\Im(u^{\gr})\), we must show that for \(y'\in\Fil^{k-1}K^{\bullet}(Y')_{\mathbb Q}\), and \(x\in\Fil^kK^{\bullet}(X)_{\mathbb Q}\), satisfying:
\sgaStableTarget{sga6.e07.eq.4-6-1}\begin{equation}\tag{4.8.1}
  j_*(y')+f^*(x)\in\Fil^{k+1}K^{\bullet}(X')_{\mathbb Q},
\end{equation}
there exists \(y\in\Fil^{k-d}K^{\bullet}(Y)_{\mathbb Q}\) such that:
\[
\begin{aligned}
\text{i)}\quad &x+i_*(y)\in\Fil^{k+1}K^{\bullet}(X)_{\mathbb Q},\\
\text{ii)}\quad &y'-g^*(y)\lambda_{-1}(F)\in\Fil^{k}K^{\bullet}(Y')_{\mathbb Q}.
\end{aligned}
\]

Put:
\[
  y=g_*(y').
\]
Then by VI 5.8, \(y\in\Fil^{k-d}K^{\bullet}(Y)_{\mathbb Q}\). Next observe that \(f_*\) preserves the filtration (using results of Exposé VIII, which causes no difficulty, since 4.8 is not used to prove this fact). Indeed, it follows from 1.8 ii) that \(f\) is a complete-intersection morphism of virtual relative dimension zero (VIII 4.2). Since \(f\) is projective, the result follows from VIII 3.2. Applying \(f_*\) to the relation
%%%%%%%%%%  scan idx 475  |  folio 462  %%%%%%%%%%
\SGAPrintPage{462}%
(4.8.1), we obtain:
\[
  f_*(j_*(y'))+f_*(f^*(x))\in\Fil^{k+1}K^{\bullet}(X)_{\mathbb Q},
\]
or, by 3.6:
\[
  i_*(y)+x\in\Fil^{k+1}K^{\bullet}(X)_{\mathbb Q}.
\]

To prove relation ii), we must show that:
\[
  y'-g^*(g_*(y'))\lambda_{-1}(F)\in\Fil^kK^{\bullet}(Y')_{\mathbb Q}.
\]
Put \(z=y'-g^*(g_*(y'))\lambda_{-1}(F)\). We have:
\[
\begin{aligned}
j_*(z)
  &=j_*(y')-j_*\bigl(g^*(g_*(y'))\lambda_{-1}(F)\bigr)\\
  &=j_*(y')-f^*(i_*(y))\qquad\text{by 3.4.}
\end{aligned}
\]
Now relation i) gives:
\[
  f^*(i_*(y))+f^*(x)\in\Fil^{k+1}K^{\bullet}(X')_{\mathbb Q}.
\]
Taking (4.8.1) into account, we therefore obtain:
\[
  j_*(z)\in\Fil^{k+1}K^{\bullet}(X')_{\mathbb Q},
\]
whence:
\[
  j^*(j_*(z))=z(1-L)\in\Fil^{k+1}K^{\bullet}(Y')_{\mathbb Q}.
\]

Let \(\bar z\) be the image of \(z\) in \(\operatorname{Gr}^{k-1}K^{\bullet}(Y')_{\mathbb Q}\), and \(\overline{(1-L)}\) that of \((1-L)\) in \(\operatorname{Gr}^{1}K^{\bullet}(Y')_{\mathbb Q}\). We therefore have: \(\overline{z(1-L)}=0\). We can then apply (VI 5.11) (which remains valid after tensoring with \(\mathbb Q\), since its proof uses only the structure of \(K^{\bullet}(Y')\) as a free \(K^{\bullet}(Y)\)-algebra), and hence find \(z'\in K^{\bullet}(Y)_{\mathbb Q}\) such that:
\[
  z-g^*(z')\lambda_{-1}(F)\in\Fil^kK^{\bullet}(Y')_{\mathbb Q}.
\]
Applying \(g_*\), and using \(g_*(z)=0\), we find:
\[
  z'\in\Fil^{k+1-d}K^{\bullet}(Y)_{\mathbb Q},
\]
whence, since \(\lambda_{-1}(F)\in\Fil^{d-1}K^{\bullet}(Y')_{\mathbb Q}\):
\[
  z\in\Fil^kK^{\bullet}(Y')_{\mathbb Q},
\]
which completes the proof.

%%%%%%%%%%  scan idx 476  |  folio 463  %%%%%%%%%%
\SGAPrintPage{463}%
Suppose now that \(X\) is noetherian and \(Y\) regular, which implies that \(v\) is surjective (3.7). We saw in the proof of 3.7 that for every \(x'\in K^{\bullet}(X')\) there is \(y'\in K^{\bullet}(Y')\) such that:
\[
  x'=f^*(f_*(x'))+j_*(y').
\]
The same property clearly remains true after tensoring with \(\mathbb Q\). Consequently, for \(x'\in\Fil^kK^{\bullet}(X')_{\mathbb Q}\) there is \((y',x)\in K^{\bullet}(Y')_{\mathbb Q}\times K^{\bullet}(X)_{\mathbb Q}\) such that:
\[
  x'=v(y',x),\qquad x\in\Fil^kK^{\bullet}(X)_{\mathbb Q}.
\]
Indeed, it suffices to take \(x=f_*(x')\).

Filter the ring \(K^{\bullet}(Y')_{\mathbb Q}\times K^{\bullet}(X)_{\mathbb Q}\) defined in 4.7 by putting:
\[
  \Fil^k\bigl(K^{\bullet}(Y')_{\mathbb Q}\times K^{\bullet}(X)_{\mathbb Q}\bigr)
  =
  \Fil^{k-1}K^{\bullet}(Y')_{\mathbb Q}\times \Fil^kK^{\bullet}(X)_{\mathbb Q}.
\]
It follows immediately from (4.7.1) that this is a ring filtration. Thus, to show that for every \(k\) the homomorphism:
\[
  \Fil^k\bigl(K^{\bullet}(Y')_{\mathbb Q}\times K^{\bullet}(X)_{\mathbb Q}\bigr)
  \xrightarrow{v}
  \Fil^kK^{\bullet}(X')_{\mathbb Q}
\]
is surjective, it suffices to show that for \(x'= \gamma^k(x'_1)\), with \(\epsilon(x'_1)=0\), there is \((y',x)\in\Fil^k(K^{\bullet}(Y')_{\mathbb Q}\times K^{\bullet}(X)_{\mathbb Q})\) such that \(x'=v(y',x)\).

By the preceding argument, there is \((y'_1,x_1)\in K^{\bullet}(Y')_{\mathbb Q}\times K^{\bullet}(X)_{\mathbb Q}\) such that \(x'_1=v(y'_1,x_1)\) and \(x_1\in\Fil^1K^{\bullet}(X)_{\mathbb Q}\). By 4.7, we then have:
\[
  x'=v\bigl(\gamma^k_*(y'_1,x_1)\bigr).
\]
Now:
\[
  \gamma^k_*(y'_1,x_1)
  =
  \left(\sum_{p=1}^{k}\gamma^{k-p}(x_1)\gamma^p(L,y'_1),\;
  \gamma^k(x_1)\right)
\]
(an immediate consequence of (4.7.2)). Since \(x_1\) has zero augmentation, \(\gamma^k(x_1)\in\Fil^kK^{\bullet}(X)_{\mathbb Q}\), and, using V 5.7,
\[
  \sum_{p=1}^{k}\gamma^{k-p}(x_1)\gamma^p(L,y'_1)
  \in\Fil^{k-1}K^{\bullet}(Y')_{\mathbb Q},
\]
which gives the result.

%%%%%%%%%%  scan idx 477  |  folio 464  %%%%%%%%%%
\SGAPrintPage{464}%
\sgaStableTarget{sga6.e07.rem.4-9}\textbf{Remark 4.9.} When \(X\) is noetherian and \(Y\) regular, proceeding as in 4.7 we can describe the graded ring structure of \(\operatorname{Gr}^*K^{\bullet}(X')_{\mathbb Q}\) in terms of the graded rings \(\operatorname{Gr}^*K^{\bullet}(X)_{\mathbb Q}\), \(\operatorname{Gr}^*K^{\bullet}(Y)_{\mathbb Q}\), the additive homomorphisms \(i^{\gr}\) and \(i_{\gr}\) induced by \(i^*\) and \(i_*\), and the element \(c^{\bullet}(N)\) of \(\operatorname{Gr}^*K^{\bullet}(Y)\). Likewise, the same data describe the homomorphisms \(f^{\gr},f_{\gr},j^{\gr},j_{\gr}\) defined from \(f^*,f_*,j^*,j_*\) by passage to the associated gradeds.

\sgaStableTarget{sga6.e07.sec.4.10}\subsection*{\normalfont 4.10.}

Let \(X\) be a \uline{regular} scheme having an ample sheaf. We know (IV\qquad) that the Grothendieck groups formed with locally free \(\underline{O}_X\)-Modules of finite type and with coherent \(\underline{O}_X\)-Modules are isomorphic. We use this fact to define a ``topological'' filtration on \(K^{\bullet}(X)\), taking \(\Fil^i_{\mathrm{top}}K^{\bullet}(X)\) to be the subgroup generated by the classes of coherent sheaves whose support has codimension \(\ge i\); it is easy to show that \(\Fil^i_{\mathrm{top}}K^{\bullet}(X)\) is generated by the classes of the sheaves \(\underline{O}_Z\), where \(Z\) ranges over the integral closed subschemes of \(X\) of codimension \(\ge i\)(X\qquad).

\sgaStableTarget{sga6.e07.prop.4-11}\textbf{Proposition 4.11.} \uline{Let} \(X\) \uline{be a regular scheme having an ample sheaf}. \uline{Then} \uline{for every} \(k\ge0\):
\[
  \Fil_{\mathrm{top}}^kK^{\bullet}(X)_{\mathbb Q}
  \subset
  \Fil_\lambda^kK^{\bullet}(X)_{\mathbb Q}.
\]
\((\mathrm{N.B.}\) The reverse inclusion, valid without tensoring with \(\mathbb Q\), will be proved in X\qquad).

By 4.10, it suffices to show that for an integral closed subscheme \(Z\) of \(X\), of codimension \(k\), we have:
\[
  \operatorname{cl}(\underline{O}_Z)\in\Fil^kK^{\bullet}(X)_{\mathbb Q}.
\]

For this, we proceed by noetherian induction on the integral closed subschemes of \(X\). Suppose therefore that the property is proved for the integral closed subschemes of \(Z\), and let us prove it for \(Z\). Let \(\underline{J}\) be the Ideal of \(\underline{O}_X\)
%%%%%%%%%%  scan idx 478  |  folio 465  %%%%%%%%%%
\SGAPrintPage{465}%
defining \(Z\), and let \(z\) be the generic point of \(Z\). Since \(Z\) is integral, \(\underline{J}_z\) is the maximal ideal of \(\underline{O}_{X,z}\), and since \(X\) is regular, \(\underline{J}_z\) can be generated by an \(\underline{O}_{X,z}\)-regular sequence. We can therefore find a sequence of sections of \(\underline{J}\) in a neighborhood of \(z\) forming an \(\underline{O}_X\)-regular sequence. Consequently, there is an open neighborhood \(U\) of \(z\) such that \(Z'=X-U\) is contained in \(Z\), and \(\underline{J}|U\) is a regular ideal of \(\underline{O}_U\). The immersion \(Z\cap U\xrightarrow{i}U\) is then a regular immersion, with codimension \(\dim(\underline{O}_{X,z})=\codim(Z,X)\) (EGA IV 5.1.2). Furthermore, we have the exact sequence (IV\qquad):
\[
  K^{\bullet}(Z')\xrightarrow{j_*}K^{\bullet}(X)
  \xrightarrow{u^*}K^{\bullet}(U)\to0,
\]
where \(j\) is the immersion \(Z'\to X\) and \(u\) the immersion \(U\to X\); tensoring with \(\mathbb Q\) gives an analogous sequence. Now \(u^*(\underline{O}_Z)=i_*(\underline{O}_{Z\cap U})\), so by 4.6 the image \(u^*(1_Z)\) of the class of \(\underline{O}_Z\) under \(u^*\) belongs to \(\Fil^kK^{\bullet}(U)_{\mathbb Q}\). Since \(u^*\) is a surjective \(\lambda\)-homomorphism, it induces a \uline{surjection} \(\Fil^kK^{\bullet}(X)\longrightarrow \Fil^kK^{\bullet}(U)\), so \(u^*(1_Z)=u^*(x)\) with \(x\in\Fil^kK^{\bullet}(X)\); there is therefore \(y\in j_*(K^{\bullet}(Z'))\) such that:
\[
  \operatorname{cl}(\underline{O}_Z)-y\in\Fil^kK^{\bullet}(X)_{\mathbb Q}.
\]
Now \(Z'\subset Z\), and since \(z\notin Z'\), \(\codim(Z',X)>k\). Consequently, \(y\) can be expressed as a combination of classes of sheaves \(\underline{O}_{Z''_\alpha}\), where the \(Z''_\alpha\) are integral closed subschemes of \(Z\), distinct from \(Z\). By the induction hypothesis, \(y\in\Fil^{k+1}K^{\bullet}(X)_{\mathbb Q}\), which completes the proof.

\sgaStableTarget{sga6.e07.bibliography}\subsection*{BIBLIOGRAPHY}
\begin{enumerate}[label={[\arabic*]}]
\item Cartan--Eilenberg, Homological Algebra.
\end{enumerate}



% Expose VIII macro additions
\providecommand{\etale}{\'etale}
\providecommand{\Ch}{\operatorname{Ch}}


%%%%%%%%%%  scan idx 479  |  folio 466  %%%%%%%%%%
\SGAPrintPage{466}%
\sgaStableTarget{sga6.e08.expose}\section*{\underline{EXPOSÉ VIII}\\[3pt] \underline{\underline{THE RIEMANN--ROCH THEOREM}}}
\begin{center}
\underline{by Pierre BERTHELOT}
\end{center}

\sgaStableTarget{sga6.e08.sec.1}\subsection*{1. Complete-intersection morphisms}

\sgaStableTarget{sga6.e08.def.1-1}\textbf{Definition 1.1.} \uline{Let} $f:X\to Y$ \uline{be a morphism of schemes}. \uline{One says that} $f$ \uline{is a complete-intersection morphism if for every point} $x$ \uline{of} $X$, \uline{there exist an open neighbourhood} $U$ \uline{of} $x$, \uline{a smooth} $Y$-\uline{scheme} $V$, \uline{such that the restriction of} $f$ \uline{to} $U$ \uline{factors as}:
\[
\begin{tikzcd}[column sep=large,row sep=large]
U \arrow[rr,"i"] \arrow[dr] && V \arrow[dl]\\
& Y &
\end{tikzcd}
\]
\uline{where} $i$ \uline{is a regular immersion} (VII 1.4).

After replacing $V$ by an open subset, one may suppose that $i$ is a regular closed immersion.

In this definition one may replace $V$ by a vector bundle $Y[T_1,\ldots,T_n]$. Indeed, since the notion is local, one may suppose that $V$ is affine, and that its image in $Y$ is contained in an affine open subset $W$ of $Y$. There then exist an integer $n$ and a closed $W$-immersion $V\longrightarrow W[T_1,\ldots,T_n]$, hence a $Y$-immersion $V\to Y[T_1,\ldots,T_n]$, which is regular by VII 1.10. Since the composite of two regular immersions is regular, one thus obtains a regular $Y$-immersion $U\longrightarrow Y[T_1,\ldots,T_n]$ factoring $f$.

\sgaStableTarget{sga6.e08.prop.1-2}\textbf{Proposition 1.2.} \uline{Let} $f:X\to Y$ \uline{be a complete-intersection morphism}. \uline{Then every} $Y$-\uline{immersion} $i:X\to Z$ \uline{of} $X$ \uline{into a smooth} $Y$-\uline{scheme is regular}.

We shall first prove the corollary:

\sgaStableTarget{sga6.e08.cor.1-3}\textbf{Corollary 1.3.} \uline{Consider the commutative diagram}:
\[
\begin{tikzcd}[column sep=large,row sep=large]
& V' \arrow[d,"p"]\\
X \arrow[ur,"j"] \arrow[r,"i"] & V ,
\end{tikzcd}
\]
%%%%%%%%%%  scan idx 480  |  folio 467  %%%%%%%%%%
\SGAPrintPage{467}%
\uline{where} $i$ \uline{and} $j$
\uline{are immersions}, \uline{and} $p$ \uline{is a smooth morphism}. \uline{Then for} $i$ \uline{to be a regular immersion}, \uline{it is necessary and sufficient that} $j$ \uline{be a regular immersion}.

i) Suppose $i$ regular. Put $X'=V'\times_V X$. The datum of $j$ defines a section $j'$ of $X'$ over $X$, such that $j=i'\circ j'$:
\[
\begin{tikzcd}[column sep=large,row sep=large]
X' \arrow[r,"i'"] \arrow[d,"p'"'] & V' \arrow[d,"p"]\\
X \arrow[r,"i"] \arrow[ur,"j"] \arrow[u,bend right=35,"j'"'] & V .
\end{tikzcd}
\]

Now $j'$ is a regular immersion, since it is a section of a smooth morphism (VII 1.10); $i'$ is a regular immersion since it is obtained from $i$ by a smooth, hence flat, base change (VII 1.5). Thus $j$ is a regular immersion (VII 1.7).

ii) Suppose $j$ regular, and $p$ étale. Then $j'$ is an open immersion (EGA IV 17.9.3), hence an isomorphism of $X$ onto $j'(X)$; put $U=j'(X)$.

Let $i''$ be the restriction of $i'$ to $U$, and let $U'$ be an open subset of $V'$ such that $U=U'\cap X'$. We then have: $i''=j\circ(j')^{-1}$, so $i''$ is a regular immersion of $U$ into $U'$. Now the morphism $U'\to V$, being smooth, is flat and of finite presentation, hence open. It is consequently a faithfully flat and quasi-compact morphism from $U'$ onto an open neighbourhood of $X$ in $V$. It follows (VII 1.5) that $i$ is a regular immersion.

iii) Suppose $j$ regular, and that $V'$ is a vector bundle over $V$. Since the property to be proved ($i$ regular) is local on $V$, one may suppose $V$ affine, $i$ closed, and that $V'=V[T_1,\ldots,T_n]$. Put $V=\Spec(A)$,
%%%%%%%%%%  scan idx 481  |  folio 468  %%%%%%%%%%
\SGAPrintPage{468}%
$X=\Spec(B)$; we have the commutative diagram:
\[
\begin{tikzcd}[column sep=large,row sep=large]
& A[T_1,\ldots,T_n] \arrow[dl,"v"']\\
B & A \arrow[l,"u"'] \arrow[u]
\end{tikzcd}
\]
where the morphisms $u$ and $v$ are surjective. Let $t_i$, for $i=1,\ldots,n$, be elements of $A$ such that for every $i$ we have
\[
u(t_i)=v(T_i).
\]
One defines a homomorphism $w:A[T_1,\ldots,T_n]\longrightarrow A$ by setting $w(T_i)=t_i$. This gives a closed immersion $s:V\to V'$ which is a section of $p$, and which satisfies $j=s\circ i$. Moreover, being a section of a smooth morphism, it is a regular immersion (VII 1.10). Knowing that $s$ and $s\circ i$ are regular, we shall deduce that $i$ is regular if we show that the sequence:
\[
0\longrightarrow i^*(\underline{K}/\underline{K}^2)\longrightarrow \underline{J}/\underline{J}^2\longrightarrow \underline{I}/\underline{I}^2\longrightarrow0
\]
is exact and locally split, where $\underline{I}$, $\underline{J}$ and $\underline{K}$ denote respectively the ideals defining $i$, $j$, and $s$ (VII 1.7). For this it suffices to show that $i^*(\underline{K}/\underline{K}^2)$ is a direct summand of $\underline{J}/\underline{J}^2$. Now one sees at once that $j'$ is obtained from $s$ by the base change $i':X'\to V'$. Furthermore, the immersions $j'$ and $s$ are regular (VII 1.10). Consequently (VII 1.2), the ideal defining $j'$ is $i'^*(\underline{K})$, and the conormal sheaf of $j'$ is $i^*(\underline{K}/\underline{K}^2)$. The exact sequence of conormal sheaves for the decomposition $j=i'\circ j'$ therefore gives, in particular, a homomorphism: $\underline{J}/\underline{J}^2\longrightarrow i^*(\underline{K}/\underline{K}^2)$. It is immediate to check that the composite: $i^*(\underline{K}/\underline{K}^2)\longrightarrow \underline{J}/\underline{J}^2\longrightarrow i^*(\underline{K}/\underline{K}^2)$ is the identity (from the affine point of view, this corresponds to the decomposition of $J$ as the direct sum of an ideal of $A$, and the ideal generated by the $T_i-t_i$). This indeed gives the required splitting.

%%%%%%%%%%  scan idx 482  |  folio 469  %%%%%%%%%%
\SGAPrintPage{469}%
iv)
Finally suppose $j$ regular and $p$ smooth. For every $x\in X$, one can find an open neighbourhood $W$ of $i(x)$, an open neighbourhood $W'$ of $j(x)$ over $W$, and a commutative diagram:
\[
\begin{tikzcd}[column sep=large,row sep=large]
W' \arrow[r,"q"] \arrow[d,"p"'] & W[T_1,\ldots,T_n] \arrow[dl,"p_0"]\\
W & ,
\end{tikzcd}
\]
where $p_0$ is the structural morphism, and $q$ an étale morphism (EGA IV 17.11.4). One may then suppose $V=W$, $V'=W'$. It is easy to see that $q\circ j$ is an immersion (introduce the fiber product $X[T_1,\ldots,T_n]$). Applying the results ii) and iii), one deduces that $q\circ j$ is regular, and then that $i$ is regular.

Let us now prove 1.2.

Since the property is local on $Z$, to prove it one may work in a neighbourhood of a point $x$ of $X$. One can then find an open neighbourhood $U$ of $x$ in $X$, a smooth $Y$-scheme $V$, and a regular $Y$-immersion $i':U\to V$. We then have the following commutative diagram:
\[
\begin{tikzcd}[column sep=large,row sep=large]
& U \arrow[dl,"i"'] \arrow[dr,"i'"] &\\
Z \arrow[dr] && V \arrow[dl]\\
& Y . &
\end{tikzcd}
\]

This gives a closed immersion $i'':U\to Z\times_Y V$, and the commutative diagram:
\[
\begin{tikzcd}[column sep=large,row sep=large]
& U \arrow[dl,"i"'] \arrow[d,"i''"'] \arrow[dr,"i'"] &\\
Z & Z\times_Y V \arrow[l,"p"'] \arrow[r,"p'"] & V.
\end{tikzcd}
\]
%%%%%%%%%%  scan idx 483  |  folio 470  %%%%%%%%%%
\SGAPrintPage{470}%
Since $i'$ is regular, and $p$ and $p'$ are smooth, it follows by 1.3 that $i''$ and $i$ are regular.

The following proposition will not be used in the rest of the Seminar:

\sgaStableTarget{sga6.e08.prop.1-4}\textbf{Proposition 1.4.} \uline{Let} $f:X\to Y$ \uline{be a flat morphism locally of finite presentation}. \uline{For} $f$ \uline{to be a complete-intersection morphism in the sense of} 1.1, \uline{it is necessary and sufficient that it be a complete-intersection morphism in the sense of} EGA IV 19.3.6.

Suppose that $f$ is a complete-intersection morphism in the sense of 1.1, and let $x\in X$. Then there exist an open neighbourhood $U$ of $x$, and a factorization:
\[
\begin{tikzcd}[column sep=large,row sep=large]
U \arrow[rr,"i"] \arrow[dr,"f"'] && V \arrow[dl,"g"]\\
& Y &
\end{tikzcd}
\]
where $i$ is a regular immersion, and $g$ a smooth morphism. Let $y=f(x)$; then the immersion $i_y:U_y\to V_y$ is regular at $x$. Indeed, after restricting $V$, one may suppose that there exist a free $\underline{O}_V$-Module of finite type $\underline{E}$, and a regular homomorphism $u:\underline{E}\to\underline{O}_V$ such that the Koszul complex associated to $u$ is a resolution of $\underline{O}_U$. Consider then the homomorphism $u_y$ induced on the fiber at $y$, and the associated Koszul complex $\underline{K}_\bullet(u_y)=\underline{K}_\bullet(u)_y$. Since $g$ is smooth, hence flat, $\underline{K}_\bullet(u)$ is a $Y$-flat resolution of $\underline{O}_U$. Thus the homology groups of $\underline{K}_\bullet(u)_y$ are the $\underline{\Tor}_i^{\underline{O}_Y}(\underline{O}_U,\underline{k}(y))$, where $\underline{k}(y)$ is the residue field of $\underline{O}_{Y,y}$. Since $f$ is a flat morphism, they vanish, and $i_y$ is regular at $x$. Furthermore, since $g$ is smooth, the local ring of $V_y$ at $x$ is regular. Since the local ring of $U_y$ at $x$ is the quotient of a regular ring by a regular ideal,
%%%%%%%%%%  scan idx 484  |  folio 471  %%%%%%%%%%
\SGAPrintPage{471}%
it is a complete-intersection ring (EGA IV 19.3.2), and therefore $f$ is a complete intersection at $x$ in the sense of EGA IV 19.3.6.

Conversely, suppose that $f$ is a complete-intersection morphism in the sense of EGA IV 19.3.6, and let $x\in X$. Let $U$ be an affine neighbourhood of $x$, over an affine neighbourhood $U'$ of $y=f(x)$. One can then find a smooth $U'$-scheme $V$, and a closed immersion $U\xrightarrow{i}V$ (for example by taking for $V$ a vector bundle over $U'$). We shall show that $i$ is regular. Now by EGA IV 19.3.7, $i$ is transversally regular at $x$, since $V$ is smooth over $Y$, hence $V_y$ is regular at $x$, and $U'$ is flat over $Y$. It then follows from EGA IV 19.2.4, a)$\Rightarrow$c), that $i$ is regular at $x$. Thus $f$ is a complete-intersection morphism in the sense of 1.1.

\sgaStableTarget{sga6.e08.prop.1-5}\textbf{Proposition 1.5.} \uline{Let} $f:X\to Y$, \uline{and} $g:Y\to Z$ \uline{be two locally complete-intersection morphisms}. \uline{Then} $g\circ f$ \uline{is a locally complete-intersection morphism}.

Let $x\in X$. One can find open neighbourhoods $U,V$ of $x,f(x)$, and regular immersions $i,j$ such that we have the factorizations:
\[
\begin{tikzcd}[column sep=large,row sep=large]
U \arrow[r,"i"] \arrow[dr,"f"'] & V[T_1,\ldots,T_n] \arrow[d] &\\
& V \arrow[r,"j"] \arrow[dr,"g"'] & Z[T'_1,\ldots,T'_m] \arrow[d]\\
& & Z .
\end{tikzcd}
\]
This diagram can be completed by the commutative square:
%%%%%%%%%%  scan idx 485  |  folio 472  %%%%%%%%%%
\SGAPrintPage{472}%
\[
\begin{tikzcd}[column sep=large,row sep=large]
V[T_1,\ldots,T_n] \arrow[r,"j'"] \arrow[d] & Z[T'_1,\ldots,T'_m,T_1,\ldots,T_n] \arrow[d]\\
V \arrow[r,"j"] & Z[T'_1,\ldots,T'_m],
\end{tikzcd}
\]
where $j'$ is obtained from $j$ by base change; since this base change is flat, $j'$ is a regular immersion, and consequently $j'\circ i$ is a regular immersion of $U$ into a scheme smooth over $Z$, giving the result.

\sgaStableTarget{sga6.e08.prop.1-6}\textbf{Proposition 1.6.} \uline{Let} $f:X\to Y$ \uline{be a morphism and} $g:Y'\to Y$ \uline{a flat base change}. \uline{If} $f$ \uline{is a complete intersection}, \uline{then the morphism} $f':X\times_Y Y'\to Y'$ \uline{obtained from} $f$ \uline{is one as well}. \uline{The converse is true if} $g$ \uline{is also quasi-compact and surjective}.

One can find a covering of $X$ by open subsets $U_i$ such that for every $i$, the restriction of $f$ to $U_i$ factors as $U_i\longrightarrow V_i\longrightarrow Y$, where $U_i\to V_i$ is a regular immersion, and $V_i\to Y$ a smooth morphism. The $U_i\times_Y Y'$ form an open covering of $X\times_Y Y'$, the immersions $U_i\times_Y Y'\longrightarrow V_i\times_Y Y'$ are regular by VII 1.5, and the morphisms $V_i\times_Y Y'\to Y'$ are smooth, giving the first result. For the converse, EGA IV 2.7.1 (iv) first implies that $f$ is locally of finite presentation, which easily reduces us, using 1.2, to the case where $f$ is a closed immersion, in which case the assertion is precisely VII 1.5.

\sgaStableTarget{sga6.e08.prop.1-7}\textbf{Proposition 1.7.} \uline{Let} $f:X\to Y$ \uline{be a complete-intersection morphism}. \uline{Then} $f$ \uline{is a perfect morphism}.

Since the property is local on $X$, one may work in a neighbourhood of a point $x$ of $X$, and suppose that there exist an open subset $U$ containing $x$, and a factorization:
%%%%%%%%%%  scan idx 486  |  folio 473  %%%%%%%%%%
\SGAPrintPage{473}%
\[
\begin{tikzcd}[column sep=large,row sep=large]
U \arrow[rr,"i"] \arrow[dr,"f"'] && V \arrow[dl,"g"]\\
& Y, &
\end{tikzcd}
\]
where $i$ is a
regular immersion, and $g$ a smooth morphism. Since $g$ is smooth, it is a perfect morphism (III 4.1). By VII 1.9, $i$ is also a perfect morphism. Since the composite of two perfect morphisms is perfect (III 4.5), the proposition follows.

\sgaStableTarget{sga6.e08.prop.1-8}\textbf{Proposition 1.8.} \uline{Let} $f:X\to Y$ \uline{be a complete-intersection morphism}, $x$ \uline{a point of} $X$, \uline{and consider a factorization}:
\[
\begin{tikzcd}[column sep=large,row sep=large]
U \arrow[rr,"i"] \arrow[dr,"f"'] && V \arrow[dl,"g"]\\
& Y &
\end{tikzcd}
\]
\uline{where} $U$ \uline{is an open neighbourhood of} $x$, $i$ \uline{a regular immersion}, \uline{and} $g$ \uline{a smooth morphism}. \uline{Denote by} $\underline{J}$ \uline{the ideal of} $\underline{O}_V$ \uline{defining} $U$. \uline{Then the integer}
\[
d_x=\operatorname{rank}\bigl((\underline{\Omega}^1_{V/Y})_{i(x)}\bigr)-\operatorname{rank}\bigl((\underline{J}/\underline{J}^2)_x\bigr)
\]
\uline{is independent of the chosen factorization of} $f$.

Suppose that the restriction of $f$ to $U$ factors through two smooth $Y$-schemes $V$ and $V'$, and put $V''=V\times_Y V'$. Then $V''$ is smooth over $Y$, and also over $V$ and $V'$. Moreover, the immersions $X\to V$ and $X\to V'$ define an immersion $X\to V''$, regular by 1.2. One may therefore reduce to comparing the integers $d_x$ associated to two factorizations of $f$ through smooth $Y$-schemes $V$ and $V'$ such that there exists a commutative diagram over $Y$:
\[
\begin{tikzcd}[column sep=large,row sep=large]
& V' \arrow[d,"p"]\\
X \arrow[ur,"j"] \arrow[r,"i"] & V
\end{tikzcd}
\]
%%%%%%%%%%  scan idx 487  |  folio 474  %%%%%%%%%%
\SGAPrintPage{474}%
where $p$ is a
smooth morphism. One may also suppose $\underline{\Omega}^1_{V/Y}$, $\underline{\Omega}^1_{V'/Y}$, $\underline{\Omega}^1_{V'/V}$ of constant rank, as well as the conormal sheaves of $i$ and $j$.

From the commutative diagram (in which all morphisms are smooth):
\[
\begin{tikzcd}[column sep=large,row sep=large]
V' \arrow[rr,"p"] \arrow[dr] && V \arrow[dl]\\
& Y &
\end{tikzcd}
\]
one deduces:
\sgaStableTarget{sga6.e08.eq.1-8-1}\begin{equation}\tag{1.8.1}
\operatorname{rank}(\underline{\Omega}^1_{V'/Y})=
\operatorname{rank}(\underline{\Omega}^1_{V/Y})+
\operatorname{rank}(\underline{\Omega}^1_{V'/V}).
\end{equation}

Let $X'=V'\times_V X$; then the rank of $\underline{\Omega}^1_{V'/V}$ is equal to that of $\underline{\Omega}^1_{X'/X}$. Furthermore, $j$ defines a section $X\xrightarrow{j'} X'$, which is a regular immersion of codimension equal to the rank of $\underline{\Omega}^1_{X'/X}$ (EGA IV 17.10.4). If $i'$ is the immersion obtained from $i$ by the base change $V'\to V$, we have $j=i'\circ j'$ and consequently (VII 1.7):
\sgaStableTarget{sga6.e08.eq.1-8-2}\begin{equation}\tag{1.8.2}
\operatorname{rank}(\underline{N}_{V'/X})=
\operatorname{rank}(\underline{N}_{V/X})+
\operatorname{rank}(\underline{\Omega}^1_{V'/V}).
\end{equation}

Comparing (1.8.1) and (1.8.2) gives the proposition.

\sgaStableTarget{sga6.e08.def.1-9}\textbf{Definition 1.9.} \uline{The integer} $d_x$ \uline{of Proposition} 1.8 \uline{is called the} virtual relative dimension \uline{of} $X$ \uline{over} $Y$ (\uline{or of} $f$) \uline{at the point} $x$.

The virtual relative dimension of $X$ over $Y$ is therefore a locally constant function on $X$, with integer values (positive or negative). If $f$ is a smooth morphism of relative dimension $n$, its virtual relative dimension is $n$. If $f$ is a regular immersion of codimension $n$, its virtual relative dimension is $-n$.

\sgaStableTarget{sga6.e08.prop.1-10}\textbf{Proposition 1.10.} \uline{Let} $f:X\to Y$ \uline{and} $g:Y\to Z$ \uline{be two complete-intersection morphisms}, \uline{and let} $x$ \uline{be a point of} $X$. \uline{Then}:
\[
\operatorname{dim.rel.virt.}_x(g\circ f)=
\operatorname{dim.rel.virt.}_x(f)+
\operatorname{dim.rel.virt.}_{f(x)}(g)
\]

%%%%%%%%%%  scan idx 488  |  folio 475  %%%%%%%%%%
\SGAPrintPage{475}%
This follows immediately from 1.9 and from the proof of 1.5.

\sgaStableTarget{sga6.e08.sec.2}\subsection*{2. Relative cotangent complex}

\sgaStableTarget{sga6.e08.para.2-1}2.1. We shall give here a definition of the cotangent complex which, without having all the generality possible (see Exp. 0 4.4), will suffice for the needs of the Riemann--Roch theorem; of course, the definition given in exp. 0 coincides in the case considered with the definition given here.

Let $f:X\to Y$ be a smoothable morphism, that is, one admitting a factorization:
\[
\begin{tikzcd}[column sep=large,row sep=large]
X \arrow[rr,"i"] \arrow[dr,"f"'] && V \arrow[dl,"g"]\\
& Y, &
\end{tikzcd}
\]
where $g$ is a smooth morphism, and $i$ an immersion which may be supposed closed. Let $\underline{J}$ be the ideal of $X$ in $V$. The canonical derivation:
\[
d:\underline{O}_V\longrightarrow\underline{\Omega}^1_{V/Y}
\]
defines a homomorphism:
\[
\underline{J}\longrightarrow\underline{O}_V\xrightarrow{d}\underline{\Omega}^1_{V/Y}\longrightarrow
\underline{\Omega}^1_{V/Y}\otimes_{\underline{O}_V}\underline{O}_X
\]
which vanishes on $\underline{J}^2$, and hence defines a homomorphism, still denoted $d$:
\sgaStableTarget{sga6.e08.eq.2-1-1}\begin{equation}\tag{2.1.1}
\underline{J}/\underline{J}^2\xrightarrow{d} i^*(\underline{\Omega}^1_{V/Y}).
\end{equation}
One checks easily that this homomorphism is $\underline{O}_X$-linear. One then defines a chain complex $\underline{L}^{X/Y}_\bullet$ by putting:
\[
\underline{L}^{X/Y}_0=i^*(\underline{\Omega}^1_{V/Y}),\qquad
\underline{L}^{X/Y}_1=\underline{J}/\underline{J}^2;
\qquad
\underline{L}^{X/Y}_i=0\quad\text{if }i\neq0,1,
\]
and taking as homomorphism $\underline{L}^{X/Y}_1\to \underline{L}^{X/Y}_0$ the homomorphism (2.1.1).

%%%%%%%%%%  scan idx 489  |  folio 476  %%%%%%%%%%
\SGAPrintPage{476}%
\sgaStableTarget{sga6.e08.prop.2-2}\textbf{Proposition 2.2.} \uline{The complex} $\underline{L}^{X/Y}_\bullet$ \uline{does not depend}, \uline{up to canonical isomorphism in the derived category} $D(X)$, \uline{on the chosen factorization through a smooth} $Y$-\uline{scheme}.

Let $V,V'$ be two smooth $Y$-schemes through which $f$ factors. Put: $V''=V\times_Y V'$; then $V''$ is smooth over $Y$, and also over $V$ and $V'$. Moreover, the immersions $X\to V$ and $X\to V'$ define an immersion $X\to V''$ which factors $f$. One may therefore reduce to comparing the complexes $\underline{L}^{X/Y}_\bullet$ associated to two factorizations through smooth $Y$-schemes $V$ and $V'$, such that there exists a commutative diagram over $Y$:
\[
\begin{tikzcd}[column sep=large,row sep=large]
& V' \arrow[d,"p"]\\
X \arrow[ur,"j"] \arrow[r,"i"] & V,
\end{tikzcd}
\]
where $i$ and $j$ are closed immersions, and $p$ a smooth morphism.

By functoriality, one then obtains homomorphisms:
\[
\underline{I}/\underline{I}^2\longrightarrow \underline{J}/\underline{J}^2;
\qquad
p^*(\underline{\Omega}^1_{V/Y})\longrightarrow\underline{\Omega}^1_{V'/Y},
\]
hence a diagram:
\sgaStableTarget{sga6.e08.eq.2-2-1}\begin{equation}\tag{2.2.1}
\begin{tikzcd}[column sep=large,row sep=large]
\underline{I}/\underline{I}^2 \arrow[r] \arrow[d,"d"'] & \underline{J}/\underline{J}^2 \arrow[d,"d"]\\
i^*(\underline{\Omega}^1_{V/Y}) \arrow[r] & j^*(\underline{\Omega}^1_{V'/Y}).
\end{tikzcd}
\end{equation}

Since $V$ and $V'$ are smooth over $Y$, one obtains the exact sequence:
\[
0\longrightarrow p^*(\underline{\Omega}^1_{V/Y})\longrightarrow\underline{\Omega}^1_{V'/Y}
\longrightarrow\underline{\Omega}^1_{V'/V}\longrightarrow0.
\]
Since the Modules of differentials considered are locally free over $\underline{O}_{V'}$, applying the functor $j^*$ gives the exact sequence:
\sgaStableTarget{sga6.e08.eq.2-2-2}\begin{equation}\tag{2.2.2}
0\longrightarrow i^*(\underline{\Omega}^1_{V/Y})\longrightarrow j^*(\underline{\Omega}^1_{V'/Y})
\longrightarrow j^*(\underline{\Omega}^1_{V'/V})\longrightarrow0.
\end{equation}

%%%%%%%%%%  scan idx 490  |  folio 477  %%%%%%%%%%
\SGAPrintPage{477}%
Furthermore,
$j$ defines a section $j'$ of $X'=V'\times_V X$:
\[
\begin{tikzcd}[column sep=large,row sep=large]
X' \arrow[r,"i'"] \arrow[d,"p'"'] & V' \arrow[d,"p"]\\
X \arrow[r,"i"] \arrow[ur,"j"] \arrow[u,"j'"',bend right] & V .
\end{tikzcd}
\]
Let $\underline{I},\underline{J},\underline{K}$ be the ideals defining $i,j,j'$; we have: $j=i'\circ j'$, and the immersion $j'$ is regular; moreover, $i'$ is defined by the ideal $p^*(\underline{I})$ (since $p$ is flat), and its conormal sheaf is therefore $i'^*(p^*(\underline{I}))=p'^*(\underline{I}/\underline{I}^2)$. One then obtains the exact sequence:
\sgaStableTarget{sga6.e08.eq.2-2-3}\begin{equation}\tag{2.2.3}
0\longrightarrow \underline{I}/\underline{I}^2\longrightarrow \underline{J}/\underline{J}^2\longrightarrow \underline{K}/\underline{K}^2\longrightarrow0.
\end{equation}
We shall show:

a) that the diagram (2.2.1) is commutative; this will thus give a morphism of complexes $\theta$ from the complex $\underline{L}^{X/Y}_\bullet$ defined by $V$ to the complex $\underline{L}^{X/Y}_\bullet$ defined by $V'$;

b) that the homomorphism: $\underline{K}/\underline{K}^2\longrightarrow j^*(\underline{\Omega}^1_{V'/V})$ defined by (2.2.2), (2.2.3) and (2.2.1) is an isomorphism; this will imply that the mapping cylinder of $\theta$ is acyclic, hence that $\theta$ is a quasi-isomorphism, and therefore defines an isomorphism in the derived category.

The properties a) and b) are local on $X$. One may therefore suppose $X,Y,V$ and $V'$ affine. Put: $Y=\Spec(A)$, $V=\Spec(B)$, $V'=\Spec(B')$, $X=\Spec(C)$, and keep the same letters to denote the corresponding homomorphisms and ideals. We have the diagram:
%%%%%%%%%%  scan idx 491  |  folio 478  %%%%%%%%%%
\SGAPrintPage{478}%
\[
\begin{tikzcd}[column sep=large,row sep=large]
B\otimes_A B \arrow[r,"p\otimes p"] \arrow[d,"\Delta"'] \arrow[dr,phantom,"(1)"] & B'\otimes_A B' \arrow[d,"\Delta"]\\
B \arrow[r,"p"] \arrow[dr,"i"'] & B' \arrow[d,"j"'] \arrow[dr,"i'","(2)"' near start] & B'\otimes_B B' \arrow[l,"\Delta"'] \arrow[d]\\
{} & C \arrow[r,shift right=.7ex,"p'"'] & B'\otimes_B C = B'/IB'. \arrow[l,shift right=.7ex,"j'"']
\end{tikzcd}
\]

The homomorphism $\underline{I}/\underline{I}^2\to\underline{J}/\underline{J}^2$ is obtained from the homomorphism: $IB'\to J$ and coincides with the one obtained from $p:I\to J$; the homomorphism: $i^*(\underline{\Omega}^1_{V/Y})\longrightarrow j^*(\underline{\Omega}^1_{V'/Y})$ is obtained from the commutative square (1). Property a) is therefore evident.

The homomorphism: $\underline{K}/\underline{K}^2\to j^*(\underline{\Omega}^1_{V'/V})$ considered in b) can be defined as follows: let $\overline x\in K/K^2$, and let $x$ be an element of $J$ lifting $\overline x$ (one exists since $K=J/IB'$). To $\overline x$ one associates the image of $x$ under the composite homomorphism: $J\xrightarrow{d}\Omega_{B'/A}\otimes_{B'}C\longrightarrow\Omega^1_{B'/B}\otimes_{B'}C$. Furthermore, let $H$ be the augmentation ideal of: $B'\otimes_B B'\longrightarrow B'$; the commutative square (2) gives a homomorphism: $H\to K$, hence a homomorphism: $H/H^2\otimes_{B'}C\longrightarrow K/K^2$. Now $H/H^2=\Omega^1_{B'/B}$; consequently, it suffices to check that the homomorphisms constructed are inverse to each other, which is done without difficulty.

\sgaStableTarget{sga6.e08.def.2-3}\textbf{Definition 2.3.} \uline{Let} $f:X\to Y$ \uline{be a smoothable morphism}. \uline{The complex} $\underline{L}^{X/Y}_\bullet$, \uline{defined up to canonical isomorphism in the derived category} $D(X)$, \uline{is called the} relative cotangent complex \uline{of} $X$ \uline{over} $Y$ (\uline{or of} $f$).

\sgaStableTarget{sga6.e08.prop.2-4}\textbf{Proposition 2.4.} \uline{Let} $f:X\to Y$ \uline{be a complete-intersection morphism} (1.1) \uline{which is smoothable}. \uline{Then} $\underline{L}^{X/Y}_\bullet$ \uline{is perfect}. \uline{More precisely}, \uline{for every}
%%%%%%%%%%  scan idx 492  |  folio 479  %%%%%%%%%%
\SGAPrintPage{479}%
\uline{factorization of} $f$ \uline{as the composite of an immersion and a smooth morphism}, \uline{the complex} $\underline{L}^{X/Y}_\bullet$ \uline{defined in} 2.1 \uline{is strictly perfect}.

Let $f=g\circ i$, where $g$ is a smooth morphism, and $i$ an immersion. By 1.2, $i$ is a regular immersion. If $\underline{J}$ is the Ideal defining $i$, then $\underline{J}/\underline{J}^2$ is a locally free $\underline{O}_X$-Module of finite type. Similarly, since $g:V\to Y$ is smooth, $i^*(\underline{\Omega}^1_{V/Y})$ is locally free of finite type, giving 2.4.

\sgaStableTarget{sga6.e08.cor.2-5}\textbf{Corollary 2.5.} \uline{Let} $f:X\to Y$ \uline{be a smoothable complete-intersection morphism}, \uline{and} $f=g\circ i$ \uline{a factorization of} $f$ \uline{through a smooth} $Y$-\uline{scheme} $V$. \uline{Let} $\underline{J}$ \uline{be the ideal of} $\underline{O}_V$ \uline{defining} $i$. \uline{Then}, \uline{if} $K^{\bullet}(X)$ \uline{denotes the Grothendieck group of locally free} $\underline{O}_X$-\uline{Modules of finite type} (\uline{resp}. \uline{of perfect complexes on} $X$), \uline{the element} $T_f\in K^{\bullet}(X)$ \uline{defined by}:
\[
T_f=\operatorname{cl}(i^*(\underline{\Omega}^1_{V/Y}))-\operatorname{cl}(\underline{J}/\underline{J}^2)
\]
\[
(\text{resp. }T_f=\operatorname{cl}(\underline{L}^{X/Y}_\bullet))
\]
\uline{does not depend on the chosen factorization of} $f$. \uline{The value at a point} $x$ \uline{of} $X$ \uline{of the augmentation of} $T_f$ \uline{is precisely the virtual relative dimension of} $X$ \uline{over} $Y$ \uline{at the point} $x$ (1.9).

The element $T_f$ is well defined in both cases considered by 2.4, and its class in $K^{\bullet}(X)$ is independent of the factorization, by 2.2 in the second case, and by the proof of 2.2 in the first case. The last assertion of 2.5 is trivial.

\sgaStableTarget{sga6.e08.prop.2-6}\textbf{Proposition 2.6.} \uline{Let} $f:X\to Y$, \uline{and} $g:Y\to Z$ \uline{be two smoothable complete-intersection morphisms}, \uline{and suppose in addition that the following condition is satisfied}:

%%%%%%%%%%  scan idx 493  |  folio 480  %%%%%%%%%%
\SGAPrintPage{480}%
(G) \uline{There
exists a smooth} $Y$-\uline{scheme} $V$, \uline{factoring} $f$, \uline{such that for every closed immersion}: $Y\to V'$ \uline{one can find a smooth} $V'$-\uline{scheme} $V''$, \uline{with} $V\simeq V''\times_{V'}Y$ (\uline{this condition is satisfied in particular if} $f$ \uline{is quasi-projective and if} $Y$ \uline{has an ample sheaf}, \uline{by taking for} $V$ \uline{a} $\mathbf{P}^r_Y$).

\uline{Then} $g\circ f$ \uline{is a complete intersection}, \uline{is smoothable}, \uline{and we have a distinguished triangle}:
\[
\begin{tikzcd}[column sep=large,row sep=large]
& \underline{L}^{X/Y}_\bullet \arrow[dl] \\
\mathbb Lf^*(\underline{L}^{Y/Z}_\bullet) \arrow[rr] && \underline{L}^{X/Z}_\bullet . \arrow[ul]
\end{tikzcd}
\]

Let $V'$ be a smooth $Z$-scheme factoring $g$. Consider the commutative diagram:
\[
\begin{tikzcd}[column sep=large,row sep=large]
X \arrow[r,"i"] \arrow[dr,"f"'] & V \arrow[r,"j'"] \arrow[d,"p"] & V'' \arrow[d,"p''"]\\
& Y \arrow[r,"j"] \arrow[dr,"g"'] & V' \arrow[d,"p'"]\\
& & Z .
\end{tikzcd}
\]

Then the immersion $j'$ is regular, since it is obtained from $j$ by a flat base change. Consequently $j'\circ i$ is a regular immersion factoring $g\circ f$, which is therefore smoothable and a complete intersection. Moreover, we have the exact sequence of conormal sheaves:
\[
0\longrightarrow i^*(\underline{J}'/\underline{J}'^2)\longrightarrow \underline{K}/\underline{K}^2\longrightarrow \underline{I}/\underline{I}^2\longrightarrow0,
\]
where $\underline{I},\underline{J},\underline{J}',\underline{K}$ denote the sheaves defining respectively the immersions $i,j,j',j'\circ i$. Now, since $p''$ is flat, $\underline{J}'=p''^*(\underline{J})$. The sequence can therefore be written:
\sgaStableTarget{sga6.e08.eq.2-6-1}\begin{equation}\tag{2.6.1}
0\longrightarrow f^*(\underline{J}/\underline{J}^2)\longrightarrow \underline{K}/\underline{K}^2\longrightarrow \underline{I}/\underline{I}^2\longrightarrow0.
\end{equation}

Furthermore, since $V'$ and $V''$ are smooth over $Z$, we have the exact sequence:
\[
0\longrightarrow p''^*(\underline{\Omega}^1_{V'/Z})\longrightarrow\underline{\Omega}^1_{V''/Z}\longrightarrow\underline{\Omega}^1_{V''/V'}\longrightarrow0.
\]
%%%%%%%%%%  scan idx 494  |  folio 481  %%%%%%%%%%
\SGAPrintPage{481}%
Applying
$i^*\circ j'^*$ gives the exact sequence:
\[
0\longrightarrow (p''\circ j'\circ i)^*(\underline{\Omega}^1_{V'/Z})\longrightarrow (j'\circ i)^*(\underline{\Omega}^1_{V''/Z})\longrightarrow (j'\circ i)^*(\underline{\Omega}^1_{V''/V'})\longrightarrow0.
\]
Since $j'^*(\underline{\Omega}^1_{V''/V'})=\underline{\Omega}^1_{V/Y}$ by hypothesis (G), we obtain:
\sgaStableTarget{sga6.e08.eq.2-6-2}\begin{equation}\tag{2.6.2}
0\longrightarrow f^*(j^*(\underline{\Omega}^1_{V'/Z}))\longrightarrow (j'\circ i)^*(\underline{\Omega}^1_{V''/Z})\longrightarrow i^*(\underline{\Omega}^1_{V/Y})\longrightarrow0.
\end{equation}

Combining (2.6.1) and (2.6.2), we are reduced to checking that the diagram:
\[
\begin{tikzcd}[column sep=large,row sep=large]
0 \arrow[r] & f^*(\underline{J}/\underline{J}^2) \arrow[r] \arrow[d,"f^*(d)"'] & \underline{K}/\underline{K}^2 \arrow[r] \arrow[d,"d"'] & \underline{I}/\underline{I}^2 \arrow[r] \arrow[d,"d"'] & 0\\
0 \arrow[r] & f^*(j^*(\underline{\Omega}^1_{V'/Z})) \arrow[r] & (j'\circ i)^*(\underline{\Omega}^1_{V''/Z}) \arrow[r] & i^*(\underline{\Omega}^1_{V/Y}) \arrow[r] & 0
\end{tikzcd}
\]
is commutative, which is done locally without difficulty by returning to the definitions of the morphisms involved.

From the resulting exact sequence of complexes, one deduces the announced distinguished triangle in the derived category.

\sgaStableTarget{sga6.e08.cor.2-7}\textbf{Corollary 2.7.} \uline{With the notation and hypotheses of} 2.6, \uline{we have in} $K^{\bullet}(X)$ \uline{the following equality}:
\[
T_{g\circ f}=T_f+f^*(T_g).
\]

\sgaStableTarget{sga6.e08.sec.3}\subsection*{3. Riemann--Roch theorem: statement}

\sgaStableTarget{sga6.e08.para.3-1}3.1. In section 3, $Y$ is a quasi-compact scheme having an ample sheaf.

Let $f:X\to Y$ be a projective complete-intersection morphism (1.1); suppose in addition that the virtual relative dimension of $f$ (1.9)
%%%%%%%%%%  scan idx 495  |  folio 482  %%%%%%%%%%
\SGAPrintPage{482}%
is constant on $X$. Since $f$ is projective, there exists a closed immersion $X\to P$, where $P$ is a projective bundle over $Y$, and since $Y$ has an ample sheaf, one may suppose that $P=\mathbf{P}^r_Y$. We thus have a factorization:
\[
\begin{tikzcd}[column sep=large,row sep=large]
X \arrow[rr,"i"] \arrow[dr,"f"'] && \mathbf{P}^r_Y \arrow[dl,"g"]\\
& Y. &
\end{tikzcd}
\]
Since $f$ is a complete intersection, the immersion $i$ is regular (1.2).

By 1.7, $f$ is a perfect morphism. Moreover, $Y$ has an ample sheaf, and consequently so do $X$ and $P$. One can then define an additive homomorphism:
\[
f_*:K^{\bullet}(X)\longrightarrow K^{\bullet}(Y)
\]
(IV 2.12), the $K^{\bullet}$ defined in terms of locally free sheaves of finite type, or of perfect complexes, being identical. One can also define additive homomorphisms:
\[
i_*:K^{\bullet}(X)\longrightarrow K^{\bullet}(P);
\qquad
 g_*:K^{\bullet}(P)\longrightarrow K^{\bullet}(Y),
\]
and we have: $f_*=g_*\circ i_*$.

Let $d$ be the virtual relative dimension of $X$ over $Y$. The codimension of $i$ is then $r-d$. By VII 4.6, the homomorphism $i_*$ gives, after tensoring with $\mathbb Q$, inclusions:
\[
i_*\bigl(\Fil^kK^{\bullet}(X)_{\mathbb Q}\bigr)
\subset
\Fil^{k+r-d}K^{\bullet}(P)_{\mathbb Q}.
\]
By VI 5.8, the homomorphism $g_*$ gives inclusions:
\[
g_*\bigl(\Fil^{k'}K^{\bullet}(P)\bigr)
\subset
\Fil^{k'-r}K^{\bullet}(Y).
\]
Composing the two, we thus obtain:

%%%%%%%%%%  scan idx 496  |  folio 483  %%%%%%%%%%
\SGAPrintPage{483}%
\sgaStableTarget{sga6.e08.prop.3-2}\textbf{Proposition 3.2.} \uline{Let} $f:X\to Y$ \uline{be a projective morphism}, \uline{a complete intersection}, \uline{whose virtual relative dimension} $d$ \uline{is constant}. \uline{Suppose that} $Y$ \uline{is quasi-compact}, \uline{and has an ample sheaf}. \uline{Then for every} $k\in\mathbb Z$ \uline{we have}:
\[
f_*\bigl(\Fil^kK^{\bullet}(X)_{\mathbb Q}\bigr)
\subset
\Fil^{k-d}K^{\bullet}(Y)_{\mathbb Q}.
\]

We thus obtain, in particular, an additive homomorphism:
\[
f_{\mathrm{gr}}:\Gr K^{\bullet}(X)_{\mathbb Q}\longrightarrow
\Gr K^{\bullet}(Y)_{\mathbb Q}
\]
sending $\Gr^kK^{\bullet}(X)_{\mathbb Q}$ into $\Gr^{k-d}K^{\bullet}(Y)_{\mathbb Q}$.

\sgaStableTarget{sga6.e08.para.3-3}3.3. Recall that in (V 6.1) we defined a functor $A\longmapsto \operatorname{Ch}(A)$ on the category of graded $K$-algebras, with values in the category of augmented $K$-$\lambda$-algebras, where $K$ is a binomial ring (V 2.7). If $A$ is a graded $K$-algebra, we have:
\[
\operatorname{Ch}(A)=K\times(1+\widehat A^+).
\]

For every augmented $K$-$\lambda$-algebra $\Lambda$, we defined a $K$-$\lambda$-homomorphism, called the ``completed Chern class'' (V 6.7):
\[
\widetilde c=(\varepsilon,c):\Lambda\longrightarrow\operatorname{Ch}(\Gr_{\lambda}\Lambda).
\]

Finally, we defined a ring homomorphism, functorial in $A$, called the Chern character (V 6.4):
\[
\ch:\operatorname{Ch}(A)\longrightarrow K\oplus\widehat A^+_{\mathbb Q},
\]
and a group endomorphism, called the Todd character (V 1.16):
\[
\Todd:1+\widehat A^+\longrightarrow1+\widehat A^+.
\]

%%%%%%%%%%  scan idx 497  |  folio 484  %%%%%%%%%%
\SGAPrintPage{484}%
We shall apply these definitions to the augmented $\lambda$-rings $K^{\bullet}(X)$ and $K^{\bullet}(Y)$, the augmentation rings $H^0(X,\underline{\mathbb Z})$ and $H^0(Y,\underline{\mathbb Z})$ being binomial (V 2.8). To simplify notation, we shall put, for every scheme $Z$ and every $z\in K^{\bullet}(Z)$:
\[
\ch(z)=\ch(\widetilde c(z));
\qquad
\Todd(z)=\Todd(c(z)).
\]
For every graded algebra $A$, we shall regard the group $1+\widehat A^+$ as contained in $\widehat A$. Finally, given a homomorphism $f:A\to B$ which is additive and compatible with the gradings up to a shift, we shall still denote by $f$ the homomorphism $\widehat A\to\widehat B$ obtained from it.

\sgaStableTarget{sga6.e08.para.3-4}3.4. Let $X$ be a scheme. For every locally free $\underline{O}_X$-Module of finite type $\underline{E}$, denote by $\check E$ the class in $K^{\bullet}(X)$ of the dual $\underline{\check E}$ of $\underline{E}$. It is clear that the function associating $\check E$ to $\underline{E}$ is additive on the category of locally free $\underline{O}_X$-Modules of finite type, and therefore defines an endomorphism of $K^{\bullet}(X)$, which we shall denote by $x\rightsquigarrow\check x$.

\sgaStableTarget{sga6.e08.prop.3-5}\textbf{Proposition 3.5.} \uline{Let} $X$ \uline{be a quasi-compact scheme}. \uline{Then for every} $x\in K^{\bullet}(X)$ \uline{and every} $i\geq0$, \uline{we have}:
\[
c^i(\check x)=(-1)^i c^i(x).
\]

The relation to be proved may also be written:
\[
c_t(\check x)=c_{-t}(x),
\]
where we put: $c_t(x)=\sum_{i\geq0}c^i(x)t^i$.

By additivity, one may therefore suppose that $x$ is the class of a locally free $\underline{O}_X$-Module of finite type $\underline{E}$. Moreover, after replacing $X$ by the flag scheme of $\underline{E}$, which gives an injection on the
%%%%%%%%%%  scan idx 498  |  folio 485  %%%%%%%%%%
\SGAPrintPage{485}%
$\Gr K^{\bullet}$ (VI 5.5), one
may suppose that $\underline{E}$ is an invertible sheaf. We then have: $\underline{\check E}=\underline{E}^{-1}$; hence $\check x=x^{-1}$.

The $c^i(\check x)$, $c^j(x)$ vanish for $i,j\geq2$. It therefore suffices to show that $c^1(x^{-1})=-c^1(x)$. Now $c^1(x)$ is the class in $\Gr^1K^{\bullet}(X)$ of $x-1$, and $c^1(x^{-1})$ that of $x^{-1}-1$. The equality:
\[
x^{-1}-1=1-x+(x-1)(1-x^{-1})
\]
gives:
\[
x^{-1}-1\equiv1-x\qquad\mathrm{mod.}\ \Fil^2K^{\bullet}(X);
\]
the result follows.

\sgaStableTarget{sga6.e08.thm.3-6}\textbf{Theorem 3.6. (Riemann--Roch).} \uline{Let} $f:X\to Y$ \uline{be a projective complete-intersection morphism} (1.1), \uline{whose virtual relative dimension} $d$ \uline{is constant} (1.9). \uline{Suppose that} $Y$ \uline{is quasi-compact}, \uline{and has an ample sheaf}. \uline{With the notation of} 3.3, \uline{we then have for every} $x\in K^{\bullet}(X)$ \uline{the relation}:
\[
\ch(f_*(x))=f_{\mathrm{gr}}\bigl(\ch(x)\cdot\Todd(\check T_f)\bigr),
\]
$\check T_f$ \uline{denoting the dual element of the class of the relative cotangent complex} $\underline{L}^{X/Y}_\bullet$ (2.5 \uline{and} 2.3).

In other words, the diagram
\[
\begin{tikzcd}[column sep=huge,row sep=large]
K^{\bullet}(X) \arrow[r,"\ch\circ\widetilde c"] \arrow[d,"f_*"'] & \Gr K^{\bullet}(X)_{\mathbb Q} \arrow[d,"x\mapsto f_{\mathrm{gr}}(x\cdot\Todd(\check T_f))"]\\
K^{\bullet}(Y) \arrow[r,"\ch\circ\widetilde c"] & \Gr K^{\bullet}(Y)_{\mathbb Q}
\end{tikzcd}
\]
is commutative.

%%%%%%%%%%  scan idx 499  |  folio 486  %%%%%%%%%%
\SGAPrintPage{486}%
\sgaStableTarget{sga6.e08.continuation-of-the-proof-of-theorem-3-6}The proof of
Theorem 3.6 will occupy the rest of this exposé. Let us give some preliminaries.

\sgaStableTarget{sga6.e08.prop.3-7}\textbf{Proposition 3.7.} \uline{For arbitrary} $x\in \Gr K^{\bullet}(X)_{\mathbb Q}$, $y\in \Gr K^{\bullet}(Y)_{\mathbb Q}$, \uline{we have}:
\[
 f_{\mathrm{gr}}\bigl(x\cdot f^{\mathrm{gr}}(y)\bigr)
 =
 f_{\mathrm{gr}}(x)\,y .
\]
By additivity, one may suppose \(x\) and \(y\) homogeneous; say $x\in \Gr^iK^{\bullet}(X)_{\mathbb Q}$, $y\in \Gr^jK^{\bullet}(Y)_{\mathbb Q}$.
Let $x'\in \Fil^iK^{\bullet}(X)_{\mathbb Q}$, $y'\in \Fil^jK^{\bullet}(Y)_{\mathbb Q}$ be elements lifting \(x\) and \(y\). Then the projection formula gives:
\[
 f_*\bigl(x'\cdot f^*(y')\bigr)=f_*(x')\,y'.
\]
By 3.2, this equality is written in
\(\Fil^{i+j-d}K^{\bullet}(Y)_{\mathbb Q}\). It gives an equality in
\(\Gr^{i+j-d}K^{\bullet}(Y)_{\mathbb Q}\), and one checks at once that this is indeed the required equality.

\sgaStableTarget{sga6.e08.lem.3-8}\textbf{Lemma 3.8.} \uline{Let} \(f:X\to Y\), \uline{and} \(g:Y\to Z\) \uline{be two morphisms satisfying the hypotheses of} 3.6. \uline{Then} \(g\circ f\) \uline{also satisfies them}. \uline{If} \(x\) \uline{is an element of} \(K^{\bullet}(X)\), \uline{and if} 3.6 \uline{holds for} \(f\) \uline{and the element} \(x\), \uline{and for} \(g\) \uline{and the element} \(f_*(x)\), \uline{then it holds for} \(g\circ f\) \uline{and the element} \(x\).

By hypothesis, we have:
\sgaStableTarget{sga6.e08.eq.3-8-1}\begin{equation}\tag{3.8.1}
 \ch(f_*(x))=
 f_{\mathrm{gr}}\bigl(\ch(x)\cdot\Todd(\check T_f)\bigr),
\end{equation}
and:
\sgaStableTarget{sga6.e08.eq.3-8-2}\begin{equation}\tag{3.8.2}
 \ch(g_*(f_*(x)))=
 g_{\mathrm{gr}}\bigl(\ch(f_*(x))\cdot\Todd(\check T_g)\bigr).
\end{equation}

Combining (3.8.1) and (3.8.2), we obtain:
\[
\ch(g_*(f_*(x)))=
g_{\mathrm{gr}}\!\left(
 f_{\mathrm{gr}}\bigl(\ch(x)\Todd(\check T_f)\bigr)\cdot \Todd(\check T_g)
\right).
\]
%%%%%%%%%%  scan idx 500  |  folio 487  %%%%%%%%%%
\SGAPrintPage{487}%
Applying
3.7, we find:
\[
\ch(g_*(f_*(x)))=
g_{\mathrm{gr}}\!\left(
 f_{\mathrm{gr}}\bigl(\ch(x)\cdot\Todd(\check T_f)\cdot
 f^{\mathrm{gr}}(\Todd(\check T_g))\bigr)
\right).
\]
The functoriality of \(c\) and of \(\Todd\) implies:
\[
 f^{\mathrm{gr}}(\Todd(\check T_g))=\Todd(f^*(\check T_g)).
\]
By 2.7, we have:
\[
 \check T_{g\circ f}=\check T_f+f^*(\check T_g).
\]
The multiplicativity of the Todd character therefore implies:
\[
 \Todd(\check T_f)\cdot f^{\mathrm{gr}}(\Todd(\check T_g))
 =
 \Todd(\check T_{g\circ f}).
\]
Thus finally we obtain:
\[
 \ch(g_*(f_*(x)))=
 g_{\mathrm{gr}}\bigl(f_{\mathrm{gr}}\bigl(\ch(x)\cdot\Todd(\check T_{g\circ f})\bigr)\bigr),
\]
which is the announced result.

\sgaStableTarget{sga6.e08.para.3-9}3.9. Applying 3.8 to the factorization of \(f\) given in 3.1, one sees that, to prove the Riemann--Roch theorem, it suffices to prove it in the case of a regular immersion of constant codimension, and in the case of the structural morphism: $\mathbf{P}^r_Y\longrightarrow Y$.

Furthermore, one may suppose that the immersion \(i\) satisfies the condition of VII 4.1. Indeed, let \(s\) be the section of \(\mathbf{P}^2_{\mathbf{P}^r_Y}\) introduced in VII 4.2. Then \(s\circ i\) is a regular immersion satisfying the condition of VII 4.1 (VII 4.2), and by 3.8 it suffices to prove 3.6 for \(s\circ i\), \(\mathbf{P}^2_{\mathbf{P}^r_Y}\longrightarrow\mathbf{P}^r_Y\) and \(\mathbf{P}^r_Y\longrightarrow Y\).

We shall therefore prove 3.6 by proving the following two special cases:
\begin{enumerate}[label=\roman*)]
\item \(f\) is a regular immersion of constant codimension, satisfying the condition of VII 4.1;
%%%%%%%%%%  scan idx 501  |  folio 488  %%%%%%%%%%
\SGAPrintPage{488}%
\item \(f\) is the
structural morphism \(\mathbf{P}^r_Y\to Y\).
\end{enumerate}

\sgaStableTarget{sga6.e08.sec.4}\section*{4. Riemann--Roch theorem: case of a regular closed immersion}

\sgaStableTarget{sga6.e08.para.4-1}4.1. Let \(i:Y\to X\) be a regular closed immersion. The Riemann--Roch theorem for \(i\) is essentially relation VII 4.4, which makes it possible to compute the Chern classes of \(i_*(x)\) from the \(\gamma^p(N,x)\). More precisely, if one introduces the \(\lambda\)-ring \(K^{\bullet}(Y)_{\mathbb Q N}\) defined in V 5.5, then the homomorphism from \(K^{\bullet}(Y)_{\mathbb Q N}\) to \(K^{\bullet}(X)_{\mathbb Q}\) obtained from \(i_*\) is a \(\lambda\)-homomorphism by VII 2.8 and VII 4.3. The Riemann--Roch theorem then reduces to an analogous formula concerning the (additive) homomorphism \(K^{\bullet}(Y)\longrightarrow K^{\bullet}(Y)_{\mathbb Q N}\) defined by \(y\rightsquigarrow(0,y)\), that is, to a formal relation valid for every \(\lambda\)-ring. We shall not develop this proof here, in order to avoid still more formalism on \(\lambda\)-rings (cf. \(\RRR\) I 1.5), and shall return to the method of [1], which used blow-up schemes (a method already used in the proof of VII 4.3, and which is also that of the end of \(\RRR\)).

As indicated in 3.9, one may restrict to the case of a regular immersion of constant codimension \(d\) satisfying the condition of VII 4.1, the case of a general regular immersion then following from 3.8. In this section we return to the notation of VII 3.1. We thus have the cartesian square:
\[
\begin{tikzcd}[column sep=large,row sep=large]
Y' \arrow[r,"j"] \arrow[d,"g"'] & X' \arrow[d,"f"]\\
Y \arrow[r,"i"] & X ,
\end{tikzcd}
\]
%%%%%%%%%%  scan idx 502  |  folio 489  %%%%%%%%%%
\SGAPrintPage{489}%
where \(X'\) is the \(X\)-scheme
obtained by blowing up \(Y\). We suppose in addition that \(X\) is quasi-compact and has an ample sheaf.

\sgaStableTarget{sga6.e08.prop.4-2}\textbf{Proposition 4.2.} \uline{The morphism} \(f\) \uline{is projective}, \uline{a complete intersection}, \uline{and of virtual relative dimension zero}. \uline{We also have the relation}:
\[
 f_{\mathrm{gr}}\circ f^{\mathrm{gr}}=\operatorname{Id}_{\Gr K^{\bullet}(X)_{\mathbb Q}} .
\]

The first assertion follows immediately from (VII 1.8).

The second relation follows from the relation:
\[
 f_*\circ f^*=\operatorname{Id}_{K^{\bullet}(X)},
\]
proved in (VII 3.6).

\sgaStableTarget{sga6.e08.prop.4-3}\textbf{Proposition 4.3.} \uline{For every} \(y\in \Gr K^{\bullet}(Y)_{\mathbb Q}\), \uline{we have the relation}:
\[
 f^{\mathrm{gr}}\circ i_{\mathrm{gr}}(y)
 =
 j_{\mathrm{gr}}\bigl(g^{\mathrm{gr}}(y)\cdot c^{d-1}(\check F)\bigr),
\]
\uline{where} \(c^{d-1}\) \uline{denotes the} \((d-1)\)-\uline{th Chern class}.

This relation is contained in VII 4.8, since it is precisely \(v^{\mathrm{gr}}\circ u^{\mathrm{gr}}=0\).

\sgaStableTarget{sga6.e08.prop.4-4}\textbf{Proposition 4.4.} \uline{Let} \(X\) \uline{be a quasi-compact scheme}. \uline{If} \(\underline{E}\) \uline{is a locally free sheaf of finite type on} \(X\), \uline{of class} \(E\) \uline{in} \(K^{\bullet}(X)\), \uline{and of rank} \(n\), \uline{we have}:
\[
 \ch(\lambda_{-1}(E))=c^n(\check E)(\Todd(\check E))^{-1}.
\]

The two sides of the relation to be proved are multiplicative: this is clear for \(\ch\circ\lambda_{-1}\), and for \(\Todd\); it also holds for \(c^n\), which is the class of \((-1)^{\varepsilon(E)}\lambda_{-1}(E)\). Consequently, by passing to the flag bundle of \(\underline{E}\), one may suppose that \(E\) is a sum of classes of invertible sheaves, and reduce to proving the formula for an invertible sheaf. Let
%%%%%%%%%%  scan idx 503  |  folio 490  %%%%%%%%%%
\SGAPrintPage{490}%
\(\xi\) be the class of such a sheaf. We then have:
\[
 \ch(\lambda_{-1}(\xi))=\ch(1-\xi)=1-\ch(\xi).
\]
Since \(\widetilde c(\xi)=(1,1+c^1(\xi))\), the definition of \(\ch\) gives:
\[
 \ch(\lambda_{-1}(\xi))=1-\exp(c^1(\xi)).
\]

Furthermore, \(\xi\) has augmentation \(1\); the Chern class occurring in formula 4.4 for \(\xi\) is therefore \(c^1\). Taking into account the definition of \(\Todd\), and 3.5, the relation to be proved is written:
\[
 1-\exp(c^1(\xi))
 =
 -c^1(\xi)\cdot\left(\frac{-c^1(\xi)}{1-\exp(c^1(\xi))}\right)^{-1},
\]
and is therefore verified.

\sgaStableTarget{sga6.e08.rem.4-4-1}\textbf{Remark 4.4.1.} It is also possible to deduce 4.4 from a corresponding purely algebraic statement, cf. \(\RRR\) I 1.3.

\sgaStableTarget{sga6.e08.para.4-5}4.5. We shall now show that, to prove the Riemann--Roch theorem in the case considered in this section, one may reduce to the case where \(Y\) is a divisor of \(X\), and where the element to which the formula applies is the inverse image of an element of \(K^{\bullet}(X)\).

By 4.2, the formula to be proved:
\[
 \ch(i_*(x))=
 i_{\mathrm{gr}}\bigl(\ch(x)\cdot\Todd(\check T_i)\bigr)
\]
is equivalent to the formula:
\[
 f^{\mathrm{gr}}\bigl(\ch(i_*(x))\bigr)
 =
 f^{\mathrm{gr}}\circ i_{\mathrm{gr}}\bigl(\ch(x)\cdot\Todd(\check T_i)\bigr).
\]
Taking into account the functoriality of \(\ch\), and 4.3, it may also be written:
\[
 \ch\bigl(f^*(i_*(x))\bigr)
 =
 j_{\mathrm{gr}}\bigl(g^{\mathrm{gr}}\bigl(\ch(x)\cdot\Todd(\check T_i)\bigr)\cdot c^{d-1}(\check F)\bigr),
\]
or, by VII 3.4:
\[
 \ch\bigl(j_*(g^*(x)\cdot\lambda_{-1}(F))\bigr)
 =
 j_{\mathrm{gr}}\bigl(\ch(g^*(x))\cdot\Todd(g^*(\check T_i))\cdot c^{d-1}(\check F)\bigr).
\]

%%%%%%%%%%  scan idx 504  |  folio 491  %%%%%%%%%%
\SGAPrintPage{491}%
Now by
definition we have:
\[
 T_i=-\cl(\underline{J}/\underline{J}^2)=-N;\qquad
 T_j=-\cl(\underline{J}'/\underline{J}'^2)=-L.
\]
Furthermore, the exact sequence:
\[
0\longrightarrow \underline{F}\longrightarrow g^*(\underline{N})\longrightarrow \underline{O}_{Y'}(1)\longrightarrow0
\]
gives the relation:
\[
 g^*(T_i)=-(F+L),
\]
from which we deduce:
\[
 \Todd(g^*(\check T_i))=\Todd(\check T_j)\cdot(\Todd(\check F))^{-1}.
\]

Taking 4.4 into account, the relation to be proved is therefore written:
\[
 \ch\bigl(j_*(g^*(x)\cdot\lambda_{-1}(F))\bigr)
 =
 j_{\mathrm{gr}}\bigl(\ch(g^*(x)\cdot\lambda_{-1}(F))\cdot\Todd(\check T_j)\bigr).
\]

Putting \(z=g^*(x)\cdot\lambda_{-1}(F)\), this relation is precisely the Riemann--Roch theorem applied to the regular immersion \(j\) of codimension \(1\), and to the element \(z\) of \(K^{\bullet}(Y')\).

Moreover, since we supposed that \(i\) satisfies the hypothesis of (VII 4.1), \(\lambda_{-1}(F)\) is divisible by \((1-L)\). One can therefore find \(z'\) such that \(z=z'(1-L)\), or equivalently \(z=j^*(j_*(z'))\) (exp. VII 2.7). We are thus reduced to the case where \(Y\) is a divisor, and to an element \(x\) such that there exists \(y\), with \(x=i^*(y)\).

\sgaStableTarget{sga6.e08.para.4-6}4.6. We must therefore show the relation:
\[
 \ch\bigl(i_*(i^*(y))\bigr)
 =
 i_{\mathrm{gr}}\bigl(\ch(i^*(y))\cdot\Todd(\check T_i)\bigr).
\]
Now the projection formula gives:
\[
 i_*(i^*(y))=y\cdot i_*(1),
\]
and from the exact sequence: \(0\longrightarrow \underline{J}\longrightarrow\underline{O}_X\longrightarrow\underline{O}_Y\longrightarrow0\), where \(\underline{J}\) is invertible here, we obtain:
\[
 i_*(1)=1-J,
\]
hence
%%%%%%%%%%  scan idx 505  |  folio 492  %%%%%%%%%%
\SGAPrintPage{492}%
\[
 \ch(i_*(i^*(y)))=\ch(y)\cdot\ch(1-J).
\]
Since \(\underline{J}\) is invertible, we have \(\widetilde c(J)=(1,1+c^1(J))\), hence \(\ch(J)=\exp(c^1(J))\).
Consequently, we have:
\[
 \ch(i_*(i^*(y)))=\ch(y)\cdot(1-\exp(c^1(J))).
\]

Furthermore, the functoriality of \(\ch\), and the projection formula 3.7 imply:
\[
 i_{\mathrm{gr}}\bigl(\ch(i^*(y))\cdot\Todd(\check T_i)\bigr)
 =
 \ch(y)\cdot i_{\mathrm{gr}}(\Todd(\check T_i)).
\]
Now \(T_i=-\cl(\underline{J}/\underline{J}^2)=-i^*(J)\). Thus:
\[
 i_{\mathrm{gr}}(\Todd(\check T_i))
 =
 i_{\mathrm{gr}}\bigl(i^{\mathrm{gr}}(\Todd(-\check J))\bigr)
 =
 i_{\mathrm{gr}}(1)\cdot\Todd(-\check J).
\]

Since \(i_*(1)=1-J\), \(i_{\mathrm{gr}}(1)\) is the image in \(\Gr^1K^{\bullet}(X)\) of \(1-J\), and is therefore \(-c^1(J)\). Furthermore, we have:
\[
 \Todd(-\check J)=(\Todd(\check J))^{-1}
 =
 \frac{1-\exp(c^1(J))}{-c^1(J)}.
\]

Thus finally we obtain:
\[
 i_{\mathrm{gr}}\bigl(\ch(i^*(y))\cdot\Todd(\check T_i)\bigr)
 =
 \ch(y)\cdot(1-\exp(c^1(J))),
\]
which proves the required relation.

\sgaStableTarget{sga6.e08.sec.5}\section*{5. Riemann--Roch theorem: case of the structural morphism of a projective bundle}

\sgaStableTarget{sga6.e08.para.5-1}5.1. By 3.9, to prove the Riemann--Roch theorem, it remains to prove it for the structural morphism \(f:\mathbf{P}^r_Y\longrightarrow Y\), where \(Y\) is a quasi-compact scheme having an ample sheaf. Put \(P=\mathbf{P}^r_Y\), \(\xi=\cl(\underline{O}_P(1))\), and \(\underline{E}=(\underline{O}_Y)^{r+1}\).

%%%%%%%%%%  scan idx 506  |  folio 493  %%%%%%%%%%
\SGAPrintPage{493}%
\sgaStableTarget{sga6.e08.prop.5-2}\textbf{Proposition 5.2.} \uline{We have on} \(P\) \uline{the exact sequence}:
\[
0\longrightarrow\underline{\Omega}^1_{P/Y}\longrightarrow f^*(\underline{E})(-1)\longrightarrow\underline{O}_P\longrightarrow0,
\]
\uline{i}.\uline{e}. \uline{we have an isomorphism}
\[
\underline{\Omega}^1_{P/Y}\simeq \underline{F}(-1),\quad\text{where}\quad \underline{F}=\Ker\bigl(f^*(\underline{E})\longrightarrow\underline{O}_P(1)\bigr).
\]

Consider the vector bundle over \(Y\) defined by \(\underline{E}\), and denote by \(Z\) the complement of the zero section of this bundle. We have the commutative diagram:
\[
\begin{tikzcd}[column sep=large,row sep=large]
P \arrow[dr,"f"'] && Z \arrow[ll,"h"'] \arrow[dl,"g"]\\
& Y. &
\end{tikzcd}
\]
The morphism \(h\) identifies \(Z\) with the affine \(P\)-scheme defined by the sheaf of algebras \(\underline{B}=\bigoplus_{n\in\mathbb Z}\underline{O}_P(n)\); \(\underline{B}\) is a graded \(\underline{O}_P\)-Algebra.

Since \(f,g,h\) are smooth morphisms, one may write the exact sequence of sheaves of relative differentials; since \(Z\) is affine over \(P\), one may write this sequence in terms of \(\underline{B}\)-Modules; one thus obtains:
\sgaStableTarget{sga6.e08.eq.5-2-1}\begin{equation}\tag{5.2.1}
0\longrightarrow
 \underline{\Omega}^1_{P/Y}\otimes_{\underline{O}_P}\underline{B}
\longrightarrow
 \underline{\Omega}^1_{\underline{B}/Y}
\longrightarrow
 \underline{\Omega}^1_{\underline{B}/\underline{O}_P}
\longrightarrow0 .
\end{equation}

Now this exact sequence is a graded exact sequence of graded \(\underline{B}\)-Modules. Indeed, it follows from the commutative diagram:
\[
\begin{tikzcd}[column sep=large,row sep=large]
\mathbf{S}(\underline{E}) \arrow[rr] && \displaystyle\bigoplus_{n\in\mathbb Z}\mathbf{S}(\underline{E})(n)\simeq \mathbf{S}(\underline{E})[T,T^{-1}]\\
& \underline{O}_Y \arrow[ul] \arrow[ur] &
\end{tikzcd}
\]
where the arrows are algebra homomorphisms compatible with the grading of \(\mathbf{S}(\underline{E})[T,T^{-1}]\) (grading defined by degree in the indeterminate \(T\)), the two other algebras being supposed concentrated
%%%%%%%%%%  scan idx 507  |  folio 494  %%%%%%%%%%
\SGAPrintPage{494}%
in degree \(0\).

The module of differentials \(\underline{\Omega}^1_{\underline{B}/\underline{O}_P}\) is isomorphic to \(\underline{B}(-1)\otimes_{\underline{O}_P}\underline{O}_P(1)\), since one can define an \(\mathbf{S}(\underline{E})\)-derivation \(d\):
\[
\mathbf{S}(\underline{E})[T,T^{-1}]
\longrightarrow
\mathbf{S}(\underline{E})[T,T^{-1}]\otimes_{\mathbf{S}(\underline{E})}\mathbf{S}(\underline{E})\cdot dT
\]
so as to satisfy the universal property defining \(\underline{\Omega}^1_{\underline{B}/\underline{O}_P}\), by setting \(d(aT^n)=naT^{n-1}\otimes dT\); moreover, \(d\) is compatible with the gradings if the factor \(\mathbf{S}(\underline{E})\cdot dT\) is assigned degree \(1\).

Let \(V\) be the vector bundle over \(Y\) defined by \(\underline{E}\). Then \(\underline{\Omega}^1_{Z/Y}\) is the restriction to \(Z\) of \(\underline{\Omega}^1_{V/Y}\). Now \(V\) is the affine \(Y\)-scheme defined by \(\mathbf{S}(\underline{E})\), so \(\underline{\Omega}^1_{V/Y}\) is defined by \(\mathbf{S}(\underline{E})\otimes_{\underline{O}_Y}\underline{E}\) (since \(\underline{E}\) is locally free of finite type), hence \(\underline{\Omega}^1_{V/Y}\) is equal to \(g^*(\underline{E})\), and consequently \(\underline{\Omega}^1_{\underline{B}/Y}=f^*(\underline{E})\otimes_{\underline{O}_P}\underline{B}\). To see its grading, one looks locally, and sees at once that \(f^*(\underline{E})\) must be regarded as being of degree \(1\).

Taking the degree \(0\) terms in the exact sequence (5.2.1) gives the announced exact sequence.

\sgaStableTarget{sga6.e08.rem.5-2-1}\textbf{Remark 5.2.1.} Using the general results on the differential study of Hilbert functors ([2], 5), one directly obtains an isomorphism \(\underline{\Omega}^1_{P/Y}\xrightarrow{\sim}\underline{F}(-1)\) from the given functorial characterization of \(g^*(\underline{\Omega}^1_{P/Y})\) for every \(Y\)-morphism \(g:Y'\to P\).

\sgaStableTarget{sga6.e08.para.5-3}5.3. By 5.2, we have:
\[
 T_f=\cl(\underline{\Omega}^1_{P/Y})=f^*(E)\cdot\xi^{-1}-1.
\]
Consequently:
\[
 \widetilde c(T_f)=\widetilde c(E)\cdot\widetilde c(\xi^{-1})-\widetilde c(1).
\]

%%%%%%%%%%  scan idx 508  |  folio 495  %%%%%%%%%%
\SGAPrintPage{495}%
Since \(\underline{E}=(\underline{O}_Y)^{r+1}\), we have \(E=r+1\), and consequently:
\[
 \widetilde c(E)=(r+1,1).
\]
Put:
\[
 x=\cl_{\Gr^1K^{\bullet}(P)}(\xi-1)
 =
 \cl_{\Gr^1K^{\bullet}(P)}(1-\xi^{-1})
 =
 c^1(\xi).
\]
Since \(\xi^{-1}=\check \xi\), we have:
\[
 \widetilde c(\xi^{-1})=(1,1-x).
\]
By V (6.2.1), this gives:
\[
 \widetilde c(T_f)=(r,(1-x)^{r+1}).
\]
Taking 3.5 into account, we finally obtain:
\[
 \widetilde c(\check T_f)=(r,(1+x)^{r+1}).
\]

By the definition of \(\Todd\), we therefore have:
\sgaStableTarget{sga6.e08.eq.5-3-1}\begin{equation}\tag{5.3.1}
 \Todd(\check T_f)=
 \left(\frac{x}{1-\exp(-x)}\right)^{r+1}.
\end{equation}

\sgaStableTarget{sga6.e08.para.5-4}5.4. Using the basis of \(K^{\bullet}(P)\) over \(K^{\bullet}(Y)\) formed by the elements \(\xi^k\), with \(-r\leq k\leq0\), one sees by additivity that it suffices to prove the theorem for elements of the form \(f^*(a)\cdot\xi^k\). We then have:
\[
 \ch\bigl(f_*(f^*(a)\xi^k)\bigr)
 =
 \ch(a f_*(\xi^k))
 =
 \ch(a)\cdot\ch(f_*(\xi^k)).
\]
On the other hand:
\[
 f_{\mathrm{gr}}\bigl(\ch(f^*(a)\xi^k)\Todd(\check T_f)\bigr)
 =
 f_{\mathrm{gr}}\bigl(f^{\mathrm{gr}}(\ch(a))\ch(\xi^k)\Todd(\check T_f)\bigr)
 =
 \ch(a)\cdot f_{\mathrm{gr}}\bigl(\ch(\xi^k)\Todd(\check T_f)\bigr).
\]

One thus sees that it suffices to prove the theorem for the \(\xi^k\), with \(-r\leq k\leq0\). Now by (VI (5.2.3)), we have
\[
 f_*(\xi^k)=0\quad\text{for }-r\leq k<0~;\qquad f_*(1)=1.
\]
Consequently:
\sgaStableTarget{sga6.e08.eq.5-4-1}\begin{equation}\tag{5.4.1}
 \ch(f_*(\xi^k))=0\quad\text{if }-r\leq k<0~;\qquad
 \ch(f_*(1))=1.
\end{equation}

%%%%%%%%%%  scan idx 509  |  folio 496  %%%%%%%%%%
\SGAPrintPage{496}%
Let us now calculate the second side of the Riemann--Roch formula. For every \(k\), we have:
\[
 \widetilde c(\xi^k)=(\widetilde c(\xi))^k=(1,1+x)^k=(1,1+kx)
\]
by (V (6.2.1)). Consequently:
\sgaStableTarget{sga6.e08.eq.5-4-2}\begin{equation}\tag{5.4.2}
 \ch(\xi^k)=\exp(kx).
\end{equation}
The second side of the Riemann--Roch formula is therefore:
\sgaStableTarget{sga6.e08.eq.5-4-3}\begin{equation}\tag{5.4.3}
 f_{\mathrm{gr}}\left(\exp(kx)\cdot
 \frac{x^{r+1}}{(1-\exp(-x))^{r+1}}\right)
\end{equation}
Since \(x\) is the class in \(\Gr^1K^{\bullet}(P)\) of \(1-\xi^{-1}\), \(f_{\mathrm{gr}}(x^n)\) is, for \(0\leq n\leq r\), the class in
\(\Gr^{n-r}K^{\bullet}(Y)\) of \(f_*((1-\xi^{-1})^n)=1\). We thus have:
\sgaStableTarget{sga6.e08.eq.5-4-4}\begin{equation}\tag{5.4.4}
 f_{\mathrm{gr}}(x^n)=0\quad\text{if }0\leq n<r~;\qquad
 f_{\mathrm{gr}}(x^r)=1.
\end{equation}
Finally recall that the elements \(1,x,\ldots,x^{r+1}\) are linearly dependent over \(\Gr K^{\bullet}(Y)\) (VI 5.6), and satisfy the equation:
\sgaStableTarget{sga6.e08.eq.5-4-5}\begin{equation}\tag{5.4.5}
 x^{r+1}=0,
\end{equation}
and that \(1,\ldots,x^r\) are linearly independent over \(\Gr K^{\bullet}(Y)\).

One can then write the series \(\varphi_k=\exp(kx)\cdot\frac{x^{r+1}}{(1-\exp(-x))^{r+1}}\) uniquely as a linear combination, with coefficients in
\(\widehat{\Gr K^{\bullet}(Y)_{\mathbb Q}}\), of the elements \(1,\ldots,x^r\). The image of \(\varphi_k\) under \(f_{\mathrm{gr}}\) is then, by (5.4.4) and 3.7, its coefficient of \(x^r\). Taking (5.4.1) into account, the Riemann--Roch theorem for \(f\) is therefore equivalent to:
\sgaStableTarget{sga6.e08.eq.rr}\[
\left\{
\begin{array}{ll}
\text{for }-r\leq k<0,\text{ the coefficient of }x^r\text{ in }\varphi_k\text{ is }0~;\\
\text{for }k=0,\text{ the coefficient of }x^r\text{ in }\varphi_k\text{ is }1.
\end{array}
\right.
\tag{RR}
\]
This coefficient is the residue of the differential form
\[
 \frac{\exp(kx)\cdot dx}{(1-\exp(-x))^{r+1}}.
\]
%%%%%%%%%%  scan idx 510  |  folio 497  %%%%%%%%%%
\SGAPrintPage{497}%
Putting
\(y=1-\exp(-x)\), we are reduced to calculating the residue of:
\[
 \frac{(1-y)^{-k-1}\cdot dy}{y^{r+1}},
\]
that is, the coefficient of \(y^r\) in \((1-y)^{-k-1}\). For \(-r\leq k<0\), it is zero, and for \(k=0\), it is \(1\).

This completes the proof of the Riemann--Roch theorem.

\sgaStableTarget{sga6.e08.bibliography}\subsection*{BIBLIOGRAPHY}
\begin{enumerate}[label={[\arabic*]}]
\item A. Borel and J.P. Serre, Le théorème de Riemann--Roch, Bull. Soc. Math. de France, t. 86, 1958.
\item A. Grothendieck, Techniques de construction et théorèmes d'existence en Géométrie Algébrique, Sém. Bourbaki, 1960-61.
\end{enumerate}

%%%%%%%%%%  scan idx 511  |  folio 498  %%%%%%%%%%
\SGAPrintPage{498}%
\sgaStableTarget{sga6.e09.expose}\section*{\underline{EXPOSÉ IX}\\[3pt] \underline{\underline{SOME COMPUTATIONS OF GROUPS $K$}}}
\begin{center}
\underline{by P. BERTHELOT}
\end{center}

This exposé corresponds to no oral exposé, and will not be used in the rest of the seminar. It contains a number of computations of groups $K_\bullet$ for various fibrations: vector bundles, projective bundles, torsors under a linear group, etc\ldots{} Throughout this exposé, the base scheme will be supposed noetherian, and the notation $K_\bullet(X)$ will denote the Grothendieck group of coherent sheaves on a scheme $X$. Recall that if $X$ is noetherian, the homomorphism from the $K$ of the category of coherent $\underline{O}_X$-Modules to the $K$ of the category of pseudo-coherent complexes on $X$ is an isomorphism (IV 2.4).

The results presented here are analogous to those obtained for the Chow ring, which can be found in Grothendieck's exposés in the Chevalley Seminar 1958. The question also arises of determining to what extent statements of the same type can be obtained for the graded group associated with $K_\bullet$ for the filtration defined by the dimension of the support (cf. X). It does not seem that the answer is known, even for Proposition 1.1 (homotopy exact sequence), which is one of the properties most used in the computations that follow.

\sgaStableTarget{sga6.e09.sec.1}\subsection*{1. Vector bundles}

\sgaStableTarget{sga6.e09.prop.1-1}\textbf{Proposition 1.1. (Homotopy exact sequence).} \uline{Let} $X$ \uline{be a noetherian scheme}, $X'$ \uline{a closed subscheme of} $X$, $U$ \uline{the open complement of} $X'$ \uline{in} $X$, $i$ \uline{the immersion of} $X'$ \uline{into} $X$, \uline{and} $j$ \uline{the immersion of} $U$ \uline{into} $X$. \uline{Then the sequence}:
%%%%%%%%%%  scan idx 512  |  folio 499  %%%%%%%%%%
\SGAPrintPage{499}%
\[
K_\bullet(X')\xrightarrow{i_*}K_\bullet(X)\xrightarrow{j^*}K_\bullet(U)\longrightarrow 0
\]
\uline{is
exact}.

If $\underline{F}'$ is a coherent $\underline{O}_{X'}$-Module, $i_*(\underline{F}')$ is an $\underline{O}_X$-Module whose support is contained in $X'$; consequently its restriction to $U$ is zero, and $j^*\circ i_*(\cl(\underline{F}'))=0$, whence, since the elements of the form $\cl(\underline{F}')$ generate $K_\bullet(X')$, $j^*\circ i_*=0$. This gives a homomorphism
\[
 \rho:K_\bullet(X)/\operatorname{Im}K_\bullet(X')\longrightarrow K_\bullet(U).
\]
We shall define an inverse homomorphism to $\rho$. Let $\underline{G}$ be a coherent $\underline{O}_U$-Module. Since $X$ is noetherian, there exists a coherent $\underline{O}_X$-Module $\underline{F}$ whose restriction to $U$ is isomorphic to $\underline{G}$. If $\underline{F}'$ is a second coherent $\underline{O}_X$-Module whose restriction to $U$ is isomorphic to $\underline{G}$, then we have in $K_\bullet(X)$:
\[
 \cl(\underline{F})-\cl(\underline{F}')\in \operatorname{Im}K_\bullet(X').
\]
To see this, regard $\underline{G}$ as a subsheaf of $j^*(\underline{F}\oplus \underline{F}')$ by means of the diagonal map; $\underline{G}$ then extends to a subsheaf $\underline{F}''$ of $\underline{F}\oplus \underline{F}'$, and we are reduced to showing that
\[
 \cl(\underline{F})-\cl(\underline{F}'')\in \operatorname{Im}K_\bullet(X'),
\]
and the analogous assertion for $\underline{F}'$. Let then $\underline{H}$ and $\underline{K}$ be the kernel and cokernel of the composite homomorphism:
\[
\underline{F}''\longrightarrow \underline{F}\oplus \underline{F}'\longrightarrow \underline{F} ;
\]
this homomorphism reduces to the identity of $\underline{G}$ over $U$, and therefore $\underline{H}$ and $\underline{K}$ have support contained in $X'$. If $\underline{I}$ is the ideal of $X'$ in $X$, it follows that there exist integers $m$ and $n$ such that $\underline{I}^m\cdot \underline{H}=0$, $\underline{I}^n\cdot \underline{K}=0$ (EGA I 9.3.4). One may therefore express $\underline{H}$ and $\underline{K}$ as successive extensions of Modules
%%%%%%%%%%  scan idx 513  |  folio 500  %%%%%%%%%%
\SGAPrintPage{500}%
annihilated by $\underline{I}$, hence of $\underline{O}_{X'}$-Modules, and we have:
\[
\cl(\underline{H})\in \operatorname{Im}K_\bullet(X'),\qquad
\cl(\underline{K})\in \operatorname{Im}K_\bullet(X'),
\]
giving the result. We thus define a function on the category of coherent $\underline{O}_U$-Modules with values in $K_\bullet(X)/\operatorname{Im}K_\bullet(X')$ by assigning to every coherent $\underline{O}_U$-Module the class of one of its extensions to $X$.

This function is additive. Indeed, let:
\[
0\longrightarrow \underline{G}'\longrightarrow \underline{G}\longrightarrow \underline{G}''\longrightarrow 0
\]
be an exact sequence of coherent $\underline{O}_U$-Modules. Let $\underline{F}$ be an extension of $\underline{G}$ to $X$; then there exists a sub-Module $\underline{F}'$ of $\underline{F}$ extending $\underline{G}'$, and $\underline{F}/\underline{F}'$ extends $\underline{G}''$. It therefore factors through $K_\bullet(U)$, defining the required homomorphism $K_\bullet(U)\to K_\bullet(X)/\operatorname{Im}K_\bullet(X')$. One checks immediately that it is inverse to $\rho$.

\sgaStableTarget{sga6.e09.prop.1-2}\textbf{Proposition 1.2.} \uline{Let} $X,S$ \uline{be two noetherian schemes}, \uline{and} $f$ \uline{a} flat \uline{morphism from} $X$ \uline{to} $S$. \uline{Suppose given a family of elements} $(\xi_i)_{i\in I}$ \uline{of} $K^{\bullet}(X)$, \uline{and consider} \uline{the homomorphism}:
\[
\varphi:(K_\bullet(S))^{(I)}\longrightarrow K_\bullet(X)
\]
\uline{defined by}:
\[
\varphi((x_i)_{i\in I})=\sum_{i\in I}\xi_i\cdot f^*(x_i),
\]
\uline{where} $f^*$ \uline{is the} ``\uline{inverse-image}'' \uline{homomorphism from} $K_\bullet(S)$ \uline{to} $K_\bullet(X)$, \uline{defined by} IV 2.12, \uline{and where} $K_\bullet(X)$ \uline{is regarded as a} $K^{\bullet}(X)$\uline{-module by} IV 2.10. \uline{Suppose that for every point} $s\in S$, \uline{the images of the} $\xi_i$ \uline{in} $K_\bullet(X_s)$ \uline{under} \uline{the homomorphism}:
\[
K^{\bullet}(X)\longrightarrow K^{\bullet}(X_s)\longrightarrow K_\bullet(X_s)
\]
\uline{generate} $K_\bullet(X_s)$ \uline{as an abelian group}. \uline{Then} $\varphi$ \uline{is surjective}.

%%%%%%%%%%  scan idx 514  |  folio 501  %%%%%%%%%%
\SGAPrintPage{501}%
Let $S'$ be a closed subscheme of $S$, and $U$ the open complement of $S'$ in $S$. Denote by $X_U$ the restriction of $X$ over $U$, and by $X_{S'}$ the restriction of $X$ over $S'$. Finally, still denote by $\xi_i$ the images of the $\xi_i$ in $K^{\bullet}(X_{S'})$ and $K^{\bullet}(X_U)$, and denote by:
\[
\varphi_{S'}:K_\bullet(S')^{(I)}\to K_\bullet(X_{S'}),\qquad
\varphi_U:K_\bullet(U)^{(I)}\to K_\bullet(X_U)
\]
the homomorphisms defined in the same way as $\varphi$.

Consider the diagram
\sgaStableTarget{sga6.e09.eq.1-2-1}\[
\begin{tikzcd}[column sep=tiny]
K_\bullet(S')^{(I)} \arrow[r,"i_*"] \arrow[d,"\varphi_{S'}"'] & K_\bullet(S)^{(I)} \arrow[r,"j^*"] \arrow[d,"\varphi"] & K_\bullet(U)^{(I)} \arrow[r] \arrow[d,"\varphi_U"] & 0\\
K_\bullet(X_{S'}) \arrow[r,"u_*"] & K_\bullet(X) \arrow[r,"v^*"] & K_\bullet(X_U) \arrow[r] & 0
\end{tikzcd}\tag{1.2.1}
\]
where $i$ and $j$ denote the immersions of $S'$ and $U$ into $S$, and $u$ and $v$ the immersions of $X_{S'}$ and $X_U$ into $X$. Let us show that if $\varphi_{S'}$ and $\varphi_U$ are surjective, the same then holds for $\varphi$.

The two horizontal rows of the diagram are exact by 1.1. One then sees that to show that $\varphi$ is surjective, it suffices to show that the two squares commute.

Let $(x'_i)_{i\in I}\in K_\bullet(S')^{(I)}$; we have:
\[
\varphi\circ i_*((x'_i))=\varphi((i_*(x'_i)))=\sum_i \xi_i\cdot f^*(i_*(x'_i));
\]
\[
u_*\circ\varphi_{S'}((x'_i))=u_*\left(\sum_i \xi_i\cdot f^*_{S'}(x'_i)\right)=\sum_i \xi_i\cdot u_*(f^*_{S'}(x'_i)),
\]
the last equality following from the projection formula IV 2.11.1.2. The equality $\varphi\circ i_*=u_*\circ\varphi_{S'}$ therefore follows from the equality $f^*\circ i_*=u_*\circ f^*_{S'}$, which holds by IV 3.1.1 since $i$ is
proper and $f$ flat.

%%%%%%%%%%  scan idx 515  |  folio 502  %%%%%%%%%%
\SGAPrintPage{502}%
Let $(x_i)_{i\in I}\in K_\bullet(S)^{(I)}$; we have:
\[
\varphi_U\circ j^*((x_i))=\sum_i \xi_i\cdot f_U^*\circ j^*(x_i);
\]
\[
v^*\circ\varphi((x_i))=v^*\left(\sum_i \xi_i\cdot f^*(x_i)\right)=\sum_i \xi_i\cdot v^*\circ f^*(x_i)=\sum_i \xi_i\cdot f_U^*\circ j^*(x_i)
\]
the middle equality following from IV 2.12.2.

This result allows us to prove 1.2 by noetherian induction on the set of closed subschemes of $S$. If we suppose that for every closed subscheme $S'$ of $S$ distinct from $S$, $\varphi_{S'}$ is surjective, then the preceding argument shows that we may reduce to proving 1.2 for an irreducible affine base scheme, after replacing $S$ by an irreducible affine open subset. Moreover, the diagram
\[
\begin{tikzcd}
K_\bullet(S_{\red})^{(I)} \arrow[r] \arrow[d,"\varphi_{S_{\red}}"'] & K_\bullet(S)^{(I)} \arrow[d,"\varphi"]\\
K_\bullet(X_{S_{\red}}) \arrow[r] & K_\bullet(X),
\end{tikzcd}
\]
where $K_\bullet(X_{S_{\red}})\to K_\bullet(X)$ is surjective by 1.1, shows that we may also suppose $S$ reduced.

We therefore work over an integral affine base scheme. Put $S=\Spec(A)$, and denote by $K$ the field of fractions of $A$; if $s$ is the generic point of $S$, then $\underline{O}_{S,s}=K$.

Let $y\in K_\bullet(X)$. By hypothesis, $K_\bullet(X_s)$ is generated by the images of the $\xi_i$; if $y_s$ is the image of $y$ in $K_\bullet(X_s)$, we can find integers $n_i$ such that:
\[
y_s=\sum_i\xi_i\cdot n_i.
\]
%%%%%%%%%%  scan idx 516  |  folio 503  %%%%%%%%%%
\SGAPrintPage{503}%
Furthermore,
we have $K=\varinjlim_f A_f$, where $f$ runs through the set of nonzero elements of $A$. Consequently, $\Spec(K)=\varprojlim_f\Spec(A_f)$, and:
\[
X_s=X\times_S\Spec(K)=\varprojlim_f X\times_S\Spec(A_f)
\]
by EGA IV 8.2.3 and 8.2.5. Since the schemes considered are noetherian, and the transition morphisms affine and flat, we have by IV 3.2.3:
\[
K_\bullet(X_s)=\varinjlim_f K_\bullet(X\times_S\Spec(A_f)).
\]
Since the ``inverse-image'' homomorphisms on the $K_\bullet$ are linear with respect to the $K^{\bullet}$ (IV (2.12.2)), there exists a nonempty open subset $U=\Spec(A_f)$ of $S$ such that in $K_\bullet(X_U)$ we have:
\[
y_U=\sum_i\xi_i\cdot n_i.
\]
Let $S'$ be a closed subscheme of $S$ with support $S-U$. Then $\varphi_{S'}$ is surjective by the induction hypothesis, and consequently diagram (1.2.1) shows that:
\[
y-\sum_i\xi_i\cdot n_i\in \operatorname{Im}(\varphi),
\]
so, since $n_i=f^*(n_i)$, $\varphi$ is surjective.

\sgaStableTarget{sga6.e09.cor.1-3}\textbf{Corollary 1.3.} \uline{Let} $f:X\to S$ \uline{be a flat morphism of noetherian schemes}. \uline{Suppose that for every point} $s$ \uline{of} $S$, \uline{the homomorphism}:
\[
K^{\bullet}(X)\longrightarrow K_\bullet(X_s)
\]
\uline{is surjective}. \uline{Then the homomorphism obtained from} $f^*$:
\[
K_\bullet(S)\otimes_{K^{\bullet}(S)}K^{\bullet}(X)\longrightarrow K_\bullet(X)
\]
\uline{is surjective}.

%%%%%%%%%%  scan idx 517  |  folio 504  %%%%%%%%%%
\SGAPrintPage{504}%
Take
for the family $(\xi_i)_{i\in I}$ the family of all elements of $K^{\bullet}(X)$. The hypotheses of 1.2 are then satisfied, and consequently the homomorphism:
\[
\varphi:\bigl(K_\bullet(S)\bigr)^{(K^{\bullet}(X))}\longrightarrow K_\bullet(X)
\]
is surjective. By the definition of $\varphi$, this is equivalent to saying that the homomorphism:
\[
K_\bullet(S)\otimes_{K^{\bullet}(S)}K^{\bullet}(X)\longrightarrow K_\bullet(X)
\]
is surjective.

\sgaStableTarget{sga6.e09.cor.1-4}\textbf{Corollary 1.4.} \uline{Let} $f:X\to S$ \uline{be a flat morphism of noetherian schemes}. \uline{Suppose that for every point} $\underline{s}$ \uline{of} $S$, \uline{the canonical homomorphism} $\mathbb Z\to K_\bullet(X_s)$ \uline{is an isomorphism}. \uline{Then the homomorphism}:
\[
f^*:K_\bullet(S)\longrightarrow K_\bullet(X)
\]
\uline{is surjective}. \uline{If in addition} $f$ \uline{has a section} $g$ \uline{of finite Tor-dimension}, \uline{then} $f^*$ \uline{is an isomorphism}.

Take for the family $(\xi_i)_{i\in I}$ the family consisting of the single element $1\in K^{\bullet}(X)$; we then have $\varphi=f^*$, and, the conditions of 1.2 being satisfied, $f^*$ is surjective.

If $f$ has a section $g$ of finite Tor-dimension, one can define a homomorphism $g^*:K_\bullet(X)\to K_\bullet(S)$ (IV 2.12). Since we have $g^*\circ f^*=(f\circ g)^*=\operatorname{Id}$ by IV 2.12.1.1, $f^*$ is injective, hence an isomorphism.

\sgaStableTarget{sga6.e09.rem.1-5}\textbf{Remark 1.5.} For $g$ to have finite Tor-dimension, it is necessary and sufficient that the fibers of $X$ be regular at the points of $g(S)$, i.e.\ that $f$ be regular at the points of $g(S)$ (or smooth at these points if we suppose $f$
%%%%%%%%%%  scan idx 518  |  folio 505  %%%%%%%%%%
\SGAPrintPage{505}%
of finite presentation).

Since these properties are local on $X$ and $S$, one may suppose them affine; let $S=\Spec(A)$, $X=\Spec(B)$. We then have:
\[
\operatorname{tor.dim}_B(A)=\operatorname{dim.proj}_B(A)=\sup_{s\in S}\operatorname{dim.proj}_{B\otimes_A k(s)}(k(s))
\]
by EGA IV 12.3.2, where the hypothesis of finite presentation on $B$ can be replaced by noetherian hypotheses. Now $k(s)$ is isomorphic to the residue field of $\underline{O}_{X_s,g(s)}$, and by EGA $0_{\mathrm{IV}}$ 17.2.5 we have:
\[
\operatorname{dim.proj}_{B\otimes_A k(s)}(k(s))=
\operatorname{dim.proj}_{\underline{O}_{X_s,g(s)}}(k(s)).
\]
It then follows from EGA $0_{\mathrm{IV}}$ 17.2.11 that the latter is finite if and only if $\underline{O}_{X_s,g(s)}$ has finite cohomological dimension, that is, is regular.

\sgaStableTarget{sga6.e09.prop.1-6}\textbf{Proposition 1.6.} \uline{Let} $S$ \uline{be a noetherian scheme}, $\underline{E}$ \uline{a locally free} $\underline{O}_S$\uline{-Module of finite type}, $X$ \uline{the vector bundle associated to} $\underline{E}$, \uline{and} $f$ \uline{the structural morphism of} $X$. \uline{Then} \(f^*:K_\bullet(S)\longrightarrow K_\bullet(X)\) \uline{is an isomorphism}.

Since $f$ is smooth, and has a section, namely the zero section, to prove 1.6 it suffices to show that for every $s\in S$, $K_\bullet(X_s)\simeq \mathbb Z$ (1.4). We are thus reduced to the case where the base scheme is a field.

In this case, $\underline{E}$ is a direct sum of submodules of rank $1$, which allows us to decompose $f$ into a sequence of fibrations of rank $1$, and thus reduce to the case where $X$ has rank $1$, namely $X=\Spec(k[T])$. The proposition then follows from the next lemma:

%%%%%%%%%%  scan idx 519  |  folio 506  %%%%%%%%%%
\SGAPrintPage{506}%
\sgaStableTarget{sga6.e09.lem.1-7}\textbf{Lemma 1.7.} \uline{Let} $A$ \uline{be a principal ideal domain}, $X=\Spec(A)$. \uline{Then the canonical homomorphism}:
\[
\mathbb Z\longrightarrow K_\bullet(X)
\]
\uline{is an isomorphism}.

Since $X$ is regular, the homomorphism $K^{\bullet}(X)\to K_\bullet(X)$ is an isomorphism. Thus $K_\bullet(X)$ is generated by the classes of locally free $\underline{O}_X$-Modules of finite type, i.e.\ by the classes of projective $A$-modules of finite type. Since $A$ is a principal ideal domain, every projective module of finite type is free, giving surjectivity. Injectivity follows from the fact that $K^{\bullet}(X)$ is $\mathbb Z$-augmented.

\sgaStableTarget{sga6.e09.prop.1-8}\textbf{Proposition 1.8.} \uline{Let} $S$ \uline{be a noetherian scheme}, $\underline{E}$ \uline{a locally free} $\underline{O}_S$\uline{-Module of rank} $r$, $X$ \uline{the vector bundle associated to} $\underline{E}$, $f$ \uline{the structural morphism of} $X$, \uline{and} $g$ \uline{the zero section of} $X$. \uline{Then the composite homomorphism}:
\[
K_\bullet(S)\xrightarrow{g_*}K_\bullet(X)\xrightarrow{g^*}K_\bullet(S)
\]
(\uline{where} $g^*$ \uline{is defined by} IV 2.12.1 \uline{and} 1.5) \uline{is multiplication by the element} $\lambda_{-1}(E)$ \uline{of} $K^{\bullet}(S)$, \uline{denoting by} $E$ \uline{the class of} $\underline{E}$ \uline{in} $K^{\bullet}(S)$, \uline{and putting}
\[
\lambda_{-1}(E)=\sum_{i=0}^r(-1)^i\lambda^i(E).
\]

This statement is the analogue, for $K_\bullet$, of VII 2.7.

Let $\underline{F}$ be a coherent $\underline{O}_S$-Module, of class $F$ in $K_\bullet(S)$. Then $g_*(F)$ is the class in $K_\bullet(X)$ of the complex $\mathbb Rg_*(\underline{F})=g_*(\underline{F})$, hence of $\underline{F}$ regarded as an $\underline{O}_X$-Module by means of $g$. To compute $g^*(g_*(F))$, one must therefore take the class of \(\mathbb Lg^*(\underline{F})=\underline{F}\otimes^{\mathbb L}_{\underline{O}_X}\underline{O}_S\).

%%%%%%%%%%  scan idx 520  |  folio 507  %%%%%%%%%%
\SGAPrintPage{507}%
For this, consider the complex:
\[
0\longrightarrow \wedge^r(\underline{E})\otimes_{\underline{O}_S}\mathbb S(\underline{E})\longrightarrow
\wedge^{r-1}(\underline{E})\otimes_{\underline{O}_S}\mathbb S(\underline{E})\longrightarrow\ldots\longrightarrow
\underline{E}\otimes_{\underline{O}_S}\mathbb S(\underline{E})\longrightarrow \mathbb S(\underline{E})\longrightarrow 0
\]
where $\mathbb S(\underline{E})$ is the symmetric algebra of $\underline{E}$. This is a resolution of $\underline{O}_S$ over $\mathbb S(E)$ (VI 1.11). The associated complex on $X$, denoted $\underline{L}_\bullet$, has $k$-th term $\wedge^k(f^*(E))$ and is therefore a resolution of $\underline{O}_S$ by locally free $\underline{O}_X$-Modules of finite type; thus we have \(\underline{L}_\bullet\otimes_{\underline{O}_X}\underline{F}=\underline{O}_S\otimes^{\mathbb L}_{\underline{O}_X}\underline{F}\). Consequently:
\[
\begin{aligned}
g^*(g_*(F))
&=\cl(\underline{L}_\bullet\otimes_{\underline{O}_X}\underline{F})
=\sum_{i=0}^r(-1)^i\cl(\underline{L}_i\otimes_{\underline{O}_X}\underline{F})\\
&=\sum_{i=0}^r(-1)^i\cl(\wedge^i(\underline{E})\otimes_{\underline{O}_S}\underline{F})
=F\cdot\lambda_{-1}(E).
\end{aligned}
\]
The proposition follows by additivity.

\sgaStableTarget{sga6.e09.sec.2}\subsection*{2. Principal bundles under split tori}

\sgaStableTarget{sga6.e09.para.2-1}2.1. Let $S$ be a noetherian scheme, and $T$ a torus over $S$. Suppose $T$ split; we thus have an isomorphism $T\simeq D_S(M)$, where $M$ is a free abelian group of rank $r$, or equivalently $T\simeq \underline{G}_{m\,S}^r$. Let $P$ be a torsor under $T$. Then $P$ is affine over $S$, defined by a graded $\underline{O}_S$-Algebra $\underline{A}$ of type $M$; putting $\underline{A}=\bigoplus_{m\in M}\underline{A}_m$, we have that $\underline{A}_m$ is an invertible $\underline{O}_S$-Module for every $m$, and, for $m,m'\in M$, the homomorphism $\underline{A}_m\otimes_{\underline{O}_S}\underline{A}_{m'}\longrightarrow \underline{A}_{m+m'}$ defined by multiplication in $\underline{A}$ is an isomorphism (SGA 3 VIII 4.1). We thus define a multiplicative homomorphism $\eta:M\longrightarrow K^{\bullet}(S)$ by associating to $m\in M$ the class of $\underline{A}_m$ in $K^{\bullet}(S)$.
%%%%%%%%%%  scan idx 521  |  folio 508  %%%%%%%%%%
\SGAPrintPage{508}%

\sgaStableTarget{sga6.e09.prop.2-2}\textbf{Proposition 2.2.} \uline{With the notation and hypotheses of} 2.1, \uline{let} $f:P\to S$ \uline{be the structural morphism of} $P$. \uline{Then the homomorphism}:
\[
f^*:K_\bullet(S)\longrightarrow K_\bullet(P)
\]
\uline{is surjective}, \uline{and its kernel is the subgroup of} $K_\bullet(S)$ \uline{generated by the subgroups} $(\eta(m)-1)K_\bullet(S)$ \uline{for} $m\in M$, \uline{or}, \uline{equivalently}, \uline{for} $m$ \uline{running through a system of generators of} $M$.

First suppose $r=1$, i.e.\ $T=G_{m\,S}$. Then $M=\mathbb Z$, and for every $m\in M$ one has:
\[
\underline{A}_m\simeq \underline{A}_1^{\otimes m}.
\]
Consequently, putting $\underline{A}_1=\underline{L}$, one has $\underline{A}=\bigoplus_{n\in\mathbb Z}\underline{L}^{\otimes n}$, which shows that $P$ identifies with the complement of the zero section in the bundle $X=\mathbb V(\underline{L})$. Considering $S$ as a closed subscheme of $X$ by the zero section, one has, by 1.1, the exact sequence:
\[
K_\bullet(S)\xrightarrow{g_*}K_\bullet(X)\longrightarrow K_\bullet(P)\longrightarrow 0
\]
where $g:S\to X$ is the zero section of $X$. Moreover, $f^*$ factors as \(K_\bullet(S)\xrightarrow{f'^*}K_\bullet(X)\longrightarrow K_\bullet(P)\), where $f'$ is the structural morphism of $X$. By 1.6, $f'^*$ is an isomorphism, inverse to $g^*$. Thus one has the exact sequence:
\[
K_\bullet(S)\xrightarrow{g^*\circ g_*}K_\bullet(S)\xrightarrow{f^*}K_\bullet(P)\longrightarrow 0,
\]
from which one obtains, by 1.8:
\[
K_\bullet(P)\simeq K_\bullet(S)/\lambda_{-1}(L)K_\bullet(S)=K_\bullet(S)/(1-L)K_\bullet(S)
\]
where $L$ is the class in $K^{\bullet}(S)$ of $\underline{L}$. Since for every $n\in\mathbb Z$, $L^n-1$ is a multiple of $1-L$, this proves 2.2 when $P$ has rank $1$.


%%%%%%%%%%  scan idx 522  |  folio 509  %%%%%%%%%%
\SGAPrintPage{509}%

Continue the proof by induction on the rank $r$ of $P$. Suppose the result proved for torsors under split tori of rank $<r$. Write $M=M_1\oplus M_2$, with $M_1\simeq\mathbb Z^{r-1}$ and $M_2\simeq\mathbb Z$; put $\underline{A}'=\bigoplus_{m\in M_1}\underline{A}_m$ and $\underline{A}''=\bigoplus_{m\in M_2}\underline{A}_m$, so that $\underline{A}=\underline{A}'\otimes_{\underline{O}_S}\underline{A}''$. Let $P_1=\Spec(\underline{A}')$; then $P_1$ is a torsor under $D_S(M_1)$, and $P$ is a torsor over $P_1$ with group $\underline{G}_{m\,P_1}$ and algebra the inverse image of $\underline{A}''$ on $P_1$. By the preceding argument, $K_\bullet(P)$ is the quotient of $K_\bullet(P_1)$ by the subgroup generated by the $(\eta(m_2)-1)K_\bullet(P_1)$ for $m_2\in M_2$. By the induction hypothesis, $K_\bullet(P_1)$ is the quotient of $K_\bullet(S)$ by the subgroup generated by the $(\eta(m_1)-1)K_\bullet(S)$ for $m_1\in M_1$. Taking account of the relation
\[
\eta(m_1+m_2)-1=\eta(m_1)\eta(m_2)-1
=\eta(m_1)(\eta(m_2)-1)+\eta(m_1)-1,
\]
$K_\bullet(P)$ is therefore the quotient of $K_\bullet(S)$ by the subgroup generated by the $(\eta(m)-1)K_\bullet(S)$ for $m\in M$, or equivalently for $m$ running through a system of generators of $M$.

\sgaStableTarget{sga6.e09.cor.2-3}\textbf{Corollary 2.3.} \uline{With the notation and hypotheses of} 2.1, \uline{suppose that} $S$ \uline{is regular}; \uline{then one has a canonical isomorphism}:
\[
K^{\bullet}(P)\simeq K^{\bullet}(S)/\bigl((\eta(m)-1)_{m\in M}\bigr)
\]
\uline{where} $((\eta(m)-1)_{m\in M})$ \uline{is the ideal of} $K^{\bullet}(S)$ \uline{generated by the} $\eta(m)-1$ \uline{for} $m$ \uline{running through a system of generators of} $M$.

If $S$ is regular, the canonical homomorphism $K^{\bullet}(S)\to K_\bullet(S)$ is an isomorphism (IV 2.5). Likewise, $P$ is then regular, so $K^{\bullet}(P)\to K_\bullet(P)$ is an isomorphism. The corollary therefore follows immediately from 2.2.


%%%%%%%%%%  scan idx 523  |  folio 510  %%%%%%%%%%
\SGAPrintPage{510}%

\sgaStableTarget{sga6.e09.sec.3}\subsection*{3. Projective bundles and flag bundles}

\sgaStableTarget{sga6.e09.prop.3-1}\textbf{Proposition 3.1.} \uline{Let} $S$ \uline{be a noetherian scheme}, $\underline{E}$ \uline{a locally free} $\underline{O}_S$\uline{-Module of rank} $n$, $\pi=(p_1,\ldots,p_k)$ \uline{a sequence of integers such that} $\sum_i p_i=n$, \uline{and} $X=\mathbb D_\pi(\underline{E})$ \uline{the scheme of flags of type} $\pi$ \uline{in} $\underline{E}$. \uline{Then the canonical homomorphism}:
\[
K_\bullet(S)\otimes_{K^{\bullet}(S)}K^{\bullet}(X)\longrightarrow K_\bullet(X)
\]
\uline{is an isomorphism}.

Resume the notation of VI 4.5. Thus $\underline{L}_i$ denotes the canonical quotients of the inverse image of $\underline{E}$ on $X$, $\underline{F}_i$ the kernel of $\underline{L}_i\to \underline{L}_{i-1}$ ($i=1,\ldots,k$), and $\lambda_i^{(j)}$ the class in $K^{\bullet}(X)$ of $\wedge^i(\underline{F}_j)$. Then the $\lambda_i^{(j)}$ generate $K^{\bullet}(X)$ as a $K^{\bullet}(S)$-algebra (VI 4.6). Likewise, the images of the $\lambda_i^{(j)}$ in $K^{\bullet}(X_s)$ generate $K^{\bullet}(X_s)$ as an algebra over $K^{\bullet}(k(s))=\mathbb Z$ for every point $s$ of $S$. Since $X_s$ is a flag scheme over a field, it is a regular scheme, so the homomorphism $K^{\bullet}(X_s)\to K_\bullet(X_s)$ is an isomorphism, and thus the homomorphism $K^{\bullet}(X)\to K_\bullet(X_s)$ is surjective. By 1.3, this implies that the homomorphism:
\[
K_\bullet(S)\otimes_{K^{\bullet}(S)}K^{\bullet}(X)\longrightarrow K_\bullet(X)
\]
is surjective.

To prove injectivity, one may first reduce to the case of the complete flag bundle $\mathbb D(E)$. Indeed, suppose that
\[
K_\bullet(S)\otimes_{K^{\bullet}(S)}K^{\bullet}(\mathbb D(\underline{E}))\longrightarrow K_\bullet(\mathbb D(\underline{E}))
\]
is injective. Since $K_\bullet(S)\to K_\bullet(\mathbb D(\underline{E}))$ factors through $K_\bullet(X)$, this
%%%%%%%%%%  scan idx 524  |  folio 511  %%%%%%%%%%
\SGAPrintPage{511}%
homomorphism factors as:
\[
K_\bullet(S)\otimes_{K^{\bullet}(S)}K^{\bullet}(X)\otimes_{K^{\bullet}(X)}K^{\bullet}(\mathbb D)
\longrightarrow K_\bullet(X)\otimes_{K^{\bullet}(X)}K^{\bullet}(\mathbb D)
\longrightarrow K_\bullet(\mathbb D)
\]
and the first homomorphism is injective. Since, by VI 4.3, 4.6 and 4.7, $K^{\bullet}(\mathbb D)$ is free over $K^{\bullet}(X)$, this implies the injectivity of:
\[
K_\bullet(S)\otimes_{K^{\bullet}(S)}K^{\bullet}(X)\to K_\bullet(X).
\]

Thus let $X=\mathbb D(\underline{E})$. Proceed by induction on the rank of $\underline{E}$. The case $n=1$ being trivial, suppose the property proved for $n-1$. Introduce $P=\mathbb P(\underline{E})$ and the kernel $\underline{F}$ of the canonical homomorphism $p^*(\underline{E})\to \underline{O}_P(1)$, where $p$ is the structural morphism of $P$. Then $X$ identifies with the scheme of flags of $\underline{F}$ over $P$, and the induction hypothesis shows that it is enough to prove 3.1 for $P$. Taking VI 1.1 into account, the result to be proved can therefore be stated as follows:

\sgaStableTarget{sga6.e09.cor.3-2}\textbf{Corollary 3.2.} \uline{Let} $S$ \uline{be a noetherian scheme}, $\underline{E}$ \uline{a locally free} $\underline{O}_S$\uline{-Module of rank} $n$, $X=\mathbb P(\underline{E})$ \uline{the associated projective bundle}, $f$ \uline{the structural morphism of} $X$, \uline{and} $\xi$ \uline{the class in} $K^{\bullet}(X)$ \uline{of the sheaf} $\underline{O}_X(1)$. \uline{Then the homomorphism}:
\[
\varphi:(K_\bullet(S))^n\longrightarrow K_\bullet(X)
\]
\uline{defined by}:
\[
\varphi((x_i))=\sum_{i=0}^{n-1}\xi^{-i}f^*(x_i)
\]
\uline{is an isomorphism}.

We have already seen that this homomorphism is surjective, since the $\xi^i$, $i=0,\ldots,-n+1$, form a basis of $K^{\bullet}(X)$ over $K^{\bullet}(S)$; to see that it is injective, use the homomorphism $f_*$, defined because $f$ is projective.
%%%%%%%%%%  scan idx 525  |  folio 512  %%%%%%%%%%
\SGAPrintPage{512}%
Then, by IV (2.12.4), one has:
\[
f_*\left(\sum_{i=0}^{n-1}\xi^{-i}f^*(x_i)\right)
=\sum_{i=0}^{n-1}x_i f_*(\xi^{-i})=x_0,
\]
the last equality following from the fact that for $-n+1\le i<0$ one has $f_*(\xi^i)=0$ (VI 2.8). Suppose that $\varphi((x_i))=0$. Then the preceding argument shows that $x_0=0$; applying the same reasoning to the elements $\xi^i\varphi((x_i))$ for $i=1,\ldots,n-1$, one deduces by induction that the $x_i$ are all zero, whence injectivity.

\sgaStableTarget{sga6.e09.para.3-3}3.3. Let $S$ still be a noetherian scheme, and $\underline{E}$ a locally free $\underline{O}_S$-Module of rank $n$. Denote by $V=\mathbb V(\underline{E})$ the vector bundle (contravariant in $E$!) associated with $\underline{E}$ (SGA 3 I 4.6), regarded as a group scheme over $S$. Let $P$ be a torsor under $V$; we wish to compute $K_\bullet(P)$.

Recall how one can embed $P$ in a projective bundle over $S$. Torsors under $V$ correspond to extensions $\underline{F}$ of $\underline{E}$ by $\underline{O}_S$. Indeed, given an extension:
\sgaStableTarget{sga6.e09.eq.3-3-1}\[
0\longrightarrow \underline{O}_S\xrightarrow{s}\underline{F}\xrightarrow{u}\underline{E}\longrightarrow 0,
\tag{3.3.1}
\]
one obtains a torsor under $V$ by considering the homomorphism $\mathbb V(\underline{F})\to \mathbb V(\underline{O}_S)$ defined by $s$, the section of $\mathbb V(\underline{O}_S)=\Spec(\underline{O}_S[T])$ defined by $T=1$, and taking $P=\mathbb V(\underline{F})\times_{\mathbb V(\underline{O}_S)}S$, the actions of $V$ on $P$ being induced by those of $V$ on $\mathbb V(\underline{F})$.

Conversely, let $P$ be a torsor under $V$. Denote by $\underline P$ the sheaf of germs of homomorphisms from $S$ to $P$; it is a principal homogeneous sheaf under the sheaf of germs of homomorphisms from $S$ to $V$, which identifies with the dual $\check{\underline{E}}$ of $\underline{E}$. In the classical way this yields an extension of $\underline{O}_S$ by $\check{\underline{E}}$, whence, by passing to duals, the sequence (3.3.1).


%%%%%%%%%%  scan idx 526  |  folio 513  %%%%%%%%%%
\SGAPrintPage{513}%

Thus let $P$ be a torsor under $V$, corresponding to an extension $\underline{F}$ of $\underline{E}$ by $\underline{O}_S$. The algebra $\underline{A}(P)$ of $P$ over $\underline{O}_S$ is $\mathbb S(\underline{F})\otimes_{\mathbb S(\underline{O}_S)}\underline{O}_S$, that is, $\mathbb S(\underline{F})/(s-1)\mathbb S(\underline{F})$, which identifies with the ring of homogeneous fractions of degree $0$, $\mathbb S(\underline{F})_{(s)}$. Consequently $P$ identifies with the open $D_+(s)$ of $\mathbb P(\underline{F})$. The complement of $P$ in $\mathbb P(\underline{F})$ is then $\operatorname{Proj}(\mathbb S(\underline{F})/s\cdot \mathbb S(\underline{F}))$, which identifies, by (3.3.1), with $\mathbb P(\underline{E})$. Thus one has the immersions:
\[
\begin{tikzcd}[column sep=large,row sep=small]
\mathbb P(\underline{E}) \arrow[r,hook] \arrow[dr] & \mathbb P(\underline{F}) \arrow[d] & P \arrow[l,hook'] \arrow[dl]\\
& S &
\end{tikzcd}
\]
where $\mathbb P(\underline{E})$ is closed in $\mathbb P(\underline{F})$, and $P$ is open and complementary to $\mathbb P(\underline{E})$.

\sgaStableTarget{sga6.e09.prop.3-4}\textbf{Proposition 3.4.} \uline{Let} $S$ \uline{be a noetherian scheme}, $\underline{E}$ \uline{a locally free} $\underline{O}_S$\uline{-Module of rank} $n$, $V=\mathbb V(\underline{E})$ \uline{the associated vector bundle}, $P$ \uline{a torsor under} $V$, \uline{and} $f$ \uline{the structural morphism of} $P$. \uline{Then the canonical homomorphism}:
\[
f^*:K_\bullet(S)\longrightarrow K_\bullet(P)
\]
\uline{is an isomorphism}.

With the identifications of 3.3, one has the exact sequence:
\sgaStableTarget{sga6.e09.eq.3-4-1}\[
K_\bullet(\mathbb P(\underline{E}))\xrightarrow{i_*}K_\bullet(\mathbb P(\underline{F}))\longrightarrow K_\bullet(P)\longrightarrow 0,
\tag{3.4.1}
\]
where $i$ denotes the closed immersion $\mathbb P(\underline{E})\to \mathbb P(\underline{F})$ defined by $u$. By 3.2, one has $K_\bullet(\mathbb P(\underline{E}))=K_\bullet(S)\otimes_{K^{\bullet}(S)}K^{\bullet}(\mathbb P(\underline{E}))$, $K_\bullet(\mathbb P(\underline{F}))=K_\bullet(S)\otimes_{K^{\bullet}(S)}K^{\bullet}(\mathbb P(\underline{F}))$.
Since $\mathbb P(\underline{E})$ and $\mathbb P(\underline{F})$ are smooth, $i$ is a regular immersion, and one therefore has a homomorphism $i_*:K^{\bullet}(\mathbb P(\underline{E}))\to K^{\bullet}(\mathbb P(\underline{F}))$; the homomorphism between the $K_\bullet$ is then obtained by tensoring the homomorphism between the $K^{\bullet}$ by $K_\bullet(S)$: to see this, it is enough to show that for $a\in K_\bullet(S)$ and
%%%%%%%%%%  scan idx 527  |  folio 514  %%%%%%%%%%
\SGAPrintPage{514}%
$b\in K^{\bullet}(\mathbb P(\underline{E}))$, one has
\[
i_*(b\cdot h^*(a))=i_*(b)\cdot g^*(a),
\]
where $g$ and $h$ are the structural morphisms of $\mathbb P(\underline{F})$ and $\mathbb P(\underline{E})$. Since $h=g\circ i$, one has:
\[
i_*(b\cdot h^*(a))=i_*(b\cdot i^*g^*(a))=g^*(a)\cdot i_*(b)
\]
by IV 2.12.4.

By (3.4.1), one therefore has:
\[
K_\bullet(P)\simeq K_\bullet(S)\otimes_{K^{\bullet}(S)}
\bigl(K^{\bullet}(\mathbb P(\underline{F}))/\operatorname{Im}K^{\bullet}(\mathbb P(\underline{E}))\bigr).
\]
Let $\xi_F$ and $\xi_E$ be the classes of the canonical invertible sheaves on $\mathbb P(\underline{F})$ and $\mathbb P(\underline{E})$. Then, by VI 1.1, $K^{\bullet}(\mathbb P(\underline{F}))$ (resp. $K^{\bullet}(\mathbb P(\underline{E}))$) has as a basis over $K^{\bullet}(S)$ the elements $\xi_F^k$ for $k=0,\ldots,-n$ (resp. $\xi_E^k$ for $k=0,\ldots,-n+1$). The projection formula gives:
\[
i_*(\xi_E^k)=i_*(i^*(\xi_F^k))=\xi_F^k i_*(1).
\]
If $\underline{I}$ is the ideal of $\mathbb P(\underline{E})$ in $\mathbb P(\underline{F})$, $\underline{I}$ is an invertible ideal, and one has $i_*(1)=1-I$. Moreover, $\underline{I}$ is defined by the ideal $J$ of $\mathbb S(\underline{F})$ that is the kernel of $\mathbb S(\underline{F})\to \mathbb S(\underline{E})$, and is generated by a sub-Module of $\underline{F}$ isomorphic to $\underline{O}_S$; consequently, $\underline{I}\simeq \underline{O}_{\mathbb P(\underline{F})}(-1)$. The image of $K^{\bullet}(\mathbb P(\underline{E}))$ is therefore the $K^{\bullet}(S)$-submodule generated by the $\xi_F^k-\xi_F^{k-1}$ for $k=0,\ldots,-n+1$. The homomorphism $K^{\bullet}(S)\longrightarrow K^{\bullet}(\mathbb P(\underline{F}))/\operatorname{Im}(K^{\bullet}(\mathbb P(\underline{E})))$ is thus an isomorphism, and consequently the same holds for the homomorphism $K_\bullet(S)\to K_\bullet(P)$.

\sgaStableTarget{sga6.e09.sec.4}\subsection*{4. Principal bundles under the groups $\mathbf{Gl}(n)_S$}

\sgaStableTarget{sga6.e09.prop.4-1}\textbf{Proposition 4.1.} \uline{Let} $S$ \uline{be a noetherian scheme}, $G$ \uline{the} $S$-\uline{group} $\mathbf{Gl}(n)_S$, $P$ \uline{a torsor under} $G$, \uline{and} $\underline{E}$ \uline{the locally free} $\underline{O}_S$\uline{-Module of rank} $n$ \uline{corresponding to} $P$. \uline{Denote by} $E$ \uline{the class of} $\underline{E}$ \uline{in} $K^{\bullet}(S)$, \uline{and by} $J$ \uline{the} $\lambda$-\uline{ideal generated}
%%%%%%%%%%  scan idx 528  |  folio 515  %%%%%%%%%%
\SGAPrintPage{515}%
\uline{by} $(E-n)$. \uline{If} $f$ \uline{is the structural morphism of} $P$, \uline{then} $f^*$, \uline{defined because} $f$ \uline{is flat}, \uline{is surjective and defines an isomorphism}:
\[
K_\bullet(S)/J\cdot K_\bullet(S)\xrightarrow{\sim}K_\bullet(P).
\]

Let $\underline{E}_0=(\underline{O}_S)^n$, let $e_1,\ldots,e_n$ be the evident basis of $\underline{E}_0$, and let $\underline{E}_0^{(i)}$ be the sub-Module of $\underline{E}_0$ generated by $e_1,\ldots,e_i$. The subfunctor of $G$ that assigns to every $S$-scheme $S'$ the subgroup of automorphisms of $\underline{E}_{0S'}$ leaving invariant the submodules $\underline{E}_{0S'}^{(i)}$ for every $i$ is representable by a group subscheme $B$ of $G$, which is a Borel subgroup of $G$. The quotient sheaf $\underline{G}/\underline{B}$ (for the fpqc topology) is representable by the scheme of flags of $\underline{E}_0$, which will be denoted by $X_0$. Let $\underline{L}_0^{(i)}$ be the canonical quotients of the inverse image of $\underline{E}_0$ on $X_0$, let $\underline{G}_0^{(i)}$ be the kernel of $\underline{E}_{0X_0}\to \underline{L}_0^{(i)}$, and let $\underline{F}_0^{(i)}=\underline{G}_0^{(i-1)}/\underline{G}_0^{(i)}\simeq \Ker(\underline{L}_0^{(i)}\to \underline{L}_0^{(i-1)})$.

Moreover, $B$ is the semidirect product of its unipotent radical $U$ and its maximal subtorus $T$ leaving the $e_i$ invariant. Let $Y_0=G/U$. Then $Y_0$ is a torsor with base $X_0$ under the group $T_{X_0}$. If $S'$ is an $S$-scheme, giving an $S'$-point of $Y_0$ is equivalent to giving, over an $S''$ faithfully flat and quasi-compact over $S'$, an automorphism of $\underline{E}_{0S''}$ modulo the unipotents of $U_{S''}$, that is, a flag of $\underline{E}_{0S''}$ whose submodules are free and, for every $i$, a class of elements $f_i^\alpha$ free over $\underline{O}_{S''}$ such that $f_i^\alpha-f_i^\beta\in \underline{G}_{0S''}^{(i)}$; equivalently, it is giving a flag of $\underline{E}_{0S''}$ and, for every $i$, a basis of $\underline{F}_{0S''}^{(i)}$, these data being subject to the gluing condition on $S''\times_{S'}S''$; by fpqc descent, this amounts to the same data on $S'$. One then sees immediately that $Y_0$ is the affine $X_0$-scheme defined by the algebra:
%%%%%%%%%%  scan idx 529  |  folio 516  %%%%%%%%%%
\SGAPrintPage{516}%
\[
\underline{A}_0=\coprod_{k_i\in\mathbb Z}{\underline{F}_0^{(1)}}^{\otimes k_1}\otimes {\underline{F}_0^{(2)}}^{\otimes k_2}\otimes\cdots\otimes{\underline{F}_0^{(n)}}^{\otimes k_n}.
\]
The group $G$ acts on $\underline{E}_0$, on $X_0$ equipped with the sheaves $\underline{L}_0^{(i)}$, and on $Y_0$. Taking the product of these objects with $P$ under $G$, one obtains the sheaf $\underline{E}$ associated with $P$, the scheme $X=\mathbb D(\underline{E})=P/B$ of flags of $\underline{E}$, and the scheme $Y=P/U$. We shall denote the canonical sheaves on $X$ by the same letters as their counterparts on $X_0$. Thus one has the diagram:
\[
\begin{array}{ccccc}
&& Y=P/U &&\\[-0.2em]
& \nearrow && \searrow &\\[-0.2em]
P &&&& X=P/B=\mathbb D(\underline{E})\\[-0.2em]
& \searrow && \swarrow &\\[-0.2em]
&& S &&
\end{array}.
\]
Observe that $P$ is a torsor with base $Y$ and structural group $U_Y$, and that $Y$ is a torsor with base $X$ and structural group $T_X$, defined by the $\underline{O}_X$-Algebra:
\[
\underline{A}=\coprod_{k_i\in\mathbb Z}{\underline{F}^{(1)}}^{\otimes k_1}\otimes\cdots\otimes{\underline{F}^{(n)}}^{\otimes k_n}.
\]

The homomorphism $K_\bullet(S)\to K_\bullet(P)$ therefore factors as:
\[
K_\bullet(S)\longrightarrow K_\bullet(X)\longrightarrow K_\bullet(Y)\longrightarrow K_\bullet(P).
\]
By 3.1, one has $K_\bullet(X)=K_\bullet(S)\otimes_{K^{\bullet}(S)}K^{\bullet}(X)$. Let $\xi_i$ be the class of $\underline{F}^{(i)}$ in $K^{\bullet}(X)$. Then, by 2.2, $K_\bullet(Y)$ is the quotient of $K_\bullet(X)$ by the submodule generated by the $(\xi_i-1)K_\bullet(X)$. Thus, if $I$ is the ideal of $K^{\bullet}(X)$ generated by the $\xi_i-1$, one has:
\[
K_\bullet(Y)=K_\bullet(S)\otimes_{K^{\bullet}(S)}(K^{\bullet}(X)/I).
\]
%%%%%%%%%%  scan idx 530  |  folio 517  %%%%%%%%%%
\SGAPrintPage{517}%
Now, by VI 4.8, $K^{\bullet}(X)$ is the $K^{\bullet}(S)$-algebra generated by the $\xi_i-1$, subject to the relations deduced from:
\[
\gamma_t(E-n)=\prod_{i=1}^n\gamma_t(\xi_i-1).
\]
Thus $K^{\bullet}(X)/I=K^{\bullet}(S)/J$, where $J$ is the ideal generated by the $\gamma^j(E-n)$ for $j=1,\ldots,n$, hence the $\lambda$-ideal generated by $E-n$. Consequently, $K_\bullet(Y)\simeq K_\bullet(S)/J K_\bullet(S)$. The proof is completed by the lemma:

\sgaStableTarget{sga6.e09.lem.4-2}\textbf{Lemma 4.2.} \uline{Let} $S$ \uline{be a noetherian scheme}, $B$ \uline{a parabolic subgroup of a reductive} $S$\uline{-group}, \uline{and} $P$ \uline{a torsor under the unipotent radical} $U$ \uline{of} $B$; \uline{then one has an isomorphism}:
\[
K_\bullet(S)\xrightarrow{\sim}K_\bullet(P).
\]
Indeed, there is a sequence of subgroup schemes of $U$:
\[
U_0=U\supset U_1\supset U_2\supset\cdots\supset U_p=e,
\]
closed in $B$, finite because $B$ is noetherian, and such that, for every $i$, $U_i/U_{i+1}$ is isomorphic to an $S$-vector group (SGA 3 XXVI 2.1). Consequently, $P\to S$ factors through a sequence of torsors under vector groups, whence the result by 3.4.

\sgaStableTarget{sga6.e09.cor.4-3}\textbf{Corollary 4.3.} \uline{Let} $S$ \uline{be a noetherian scheme}, $G=\mathrm{Gl}(n)_S$. \uline{Then the canonical homomorphism} $K_\bullet(S)\longrightarrow K_\bullet(G)$ \uline{is an isomorphism}.

This follows immediately from 4.1 applied to $P=G$, since then $\underline{E}=\underline{O}_S^n$.


%%%%%%%%%%  scan idx 531  |  folio 518  %%%%%%%%%%
\SGAPrintPage{518}%

\sgaStableTarget{sga6.e09.cor.4-4}\textbf{Corollary 4.4.} \uline{Under the hypotheses of} 4.1, \uline{suppose} $S$ \uline{regular}. \uline{Then one has an isomorphism}:
\[
K^{\bullet}(S)/J\xrightarrow{\sim}K^{\bullet}(P).
\]
If $S$ is regular, $G$ is also regular, and so is $P$. Thus one has $K^{\bullet}(S)\simeq K_\bullet(S)$, $K^{\bullet}(P)\simeq K_\bullet(P)$, whence the result by 4.1.

\clearpage

\providecommand{\Filt}{}\renewcommand{\Filt}{\operatorname{Filt}}
\providecommand{\Coh}{}\renewcommand{\Coh}{\operatorname{Coh}}
\providecommand{\Supp}{}\renewcommand{\Supp}{\operatorname{Supp}}
\providecommand{\longg}{}\renewcommand{\longg}{\operatorname{length}}
\providecommand{\cl}{}\renewcommand{\cl}{\operatorname{cl}}
\providecommand{\Ker}{}\renewcommand{\Ker}{\operatorname{Ker}}
\providecommand{\codim}{}\renewcommand{\codim}{\operatorname{codim}}
\providecommand{\degtr}{}\renewcommand{\degtr}{\operatorname{deg.tr}}
\providecommand{\Pic}{}\renewcommand{\Pic}{\operatorname{Pic}}
\providecommand{\Spec}{}\renewcommand{\Spec}{\operatorname{Spec}}
\providecommand{\Gr}{}\renewcommand{\Gr}{\operatorname{Gr}}
\providecommand{\Gro}{}\renewcommand{\Gro}{\operatorname{Gr}^{\circ}}
\providecommand{\Grb}{}\renewcommand{\Grb}{\operatorname{Gr}_{\bullet}}
\providecommand{\fpqc}{}\renewcommand{\fpqc}{\mathrm{fpqc}}
\providecommand{\Sch}{}\renewcommand{\Sch}{\operatorname{Sch}}
\providecommand{\Ob}{}\renewcommand{\Ob}{\operatorname{Ob}}
\providecommand{\Tor}{}\renewcommand{\Tor}{\operatorname{Tor}}
\providecommand{\red}{}\renewcommand{\red}{\mathrm{red}}
\providecommand{\Mod}{}\renewcommand{\Mod}{\operatorname{Mod}}
\providecommand{\Aut}{}\renewcommand{\Aut}{\operatorname{Aut}}
\providecommand{\GL}{}\renewcommand{\GL}{\operatorname{GL}}
\providecommand{\SL}{}\renewcommand{\SL}{\operatorname{SL}}
\providecommand{\Gm}{}\renewcommand{\Gm}{\mathbb G_{\mathrm m}}
\providecommand{\Der}{}\renewcommand{\Der}{\operatorname{Der}}
\providecommand{\id}{}\renewcommand{\id}{\operatorname{id}}
\providecommand{\Ch}{}\renewcommand{\Ch}{\operatorname{Ch}}
\providecommand{\cO}{}\renewcommand{\cO}{\mathcal O}
\providecommand{\C}{}\renewcommand{\C}{\mathbb C}
\providecommand{\Q}{}\renewcommand{\Q}{\mathbb Q}
\providecommand{\Z}{}\renewcommand{\Z}{\mathbb Z}
\providecommand{\Zell}{}\renewcommand{\Zell}{\mathbf Z_\ell}
\providecommand{\Qell}{}\renewcommand{\Qell}{\mathbf Q_\ell}
\providecommand{\ch}{}\renewcommand{\ch}{\operatorname{ch}}
\providecommand{\Todd}{}\renewcommand{\Todd}{\operatorname{Todd}}
\providecommand{\Gtop}{}\renewcommand{\Gtop}{\operatorname{Gr}_{\mathrm{top}}}
\providecommand{\Fil}{}\renewcommand{\Fil}{\operatorname{Fil}}
\providecommand{\Qcoh}{}\renewcommand{\Qcoh}{\operatorname{Qcoh}}
\providecommand{\Parf}{}\renewcommand{\Parf}{\operatorname{Parf}}
\providecommand{\Lib}{}\renewcommand{\Lib}{\operatorname{Lib}}
\providecommand{\Loclib}{}\renewcommand{\Loclib}{\operatorname{Loclib}}
\providecommand{\Proj}{}\renewcommand{\Proj}{\operatorname{Proj}}
\providecommand{\Hom}{}\renewcommand{\Hom}{\operatorname{Hom}}
\providecommand{\Diff}{}\renewcommand{\Diff}{\operatorname{Diff}}
\providecommand{\Dcat}{}\renewcommand{\Dcat}{\mathrm D}
\providecommand{\Tr}{}\renewcommand{\Tr}{\operatorname{Tr}}
\providecommand{\toramp}{}\renewcommand{\toramp}{\operatorname{tor.amp}}
\providecommand{\parfamp}{}\renewcommand{\parfamp}{\operatorname{par.amp}}
\providecommand{\coh}{}\renewcommand{\coh}{\operatorname{coh}}
\providecommand{\qcoh}{}\renewcommand{\qcoh}{\operatorname{qcoh}}
\providecommand{\parf}{}\renewcommand{\parf}{\operatorname{parf}}
\providecommand{\torf}{}\renewcommand{\torf}{\operatorname{torf}}
\providecommand{\amp}{}\renewcommand{\amp}{\operatorname{amp}}
\providecommand{\num}{}\renewcommand{\num}{\operatorname{num}}
\providecommand{\naif}{}\renewcommand{\naif}{\operatorname{naif}}
\providecommand{\alg}{}\renewcommand{\alg}{\operatorname{alg}}

\hbadness=1000
\hfuzz=0.1pt
\allowdisplaybreaks


% Exposé X local macros
\providecommand{\Filt}{\operatorname{Filt}}
\providecommand{\Coh}{\operatorname{Coh}}
\renewcommand{\Ob}{\operatorname{Ob}}
\providecommand{\longg}{\operatorname{long}}
\providecommand{\degtr}{\operatorname{deg.tr}}
\providecommand{\dimdot}{\operatorname{dim.}}
\providecommand{\Frac}{\operatorname{Frac}}
\providecommand{\fpqc}{\mathrm{fpqc}}
\providecommand{\Sch}{\operatorname{Sch}}
\providecommand{\Picf}{\operatorname{Pic}}
\providecommand{\Ens}{\operatorname{Ens}}
\providecommand{\im}{\operatorname{Im}}

%%%%%%%%%%  scan idx 532  |  folio 519  %%%%%%%%%%
\SGAPrintPage{519}%

\sgaStableTarget{sga6.e10.expose}\section*{\underline{EXPOSÉ X}\\[3pt] \underline{\underline{FORMALISM OF INTERSECTIONS ON PROPER ALGEBRAIC SCHEMES}}}
\begin{center}
\underline{by O. Jussila}\\
(with an appendix by A. Grothendieck)
\end{center}

\sgaStableTarget{sga6.e10.sec.1}\section*{1. Compatibility of the filtrations with the composition law}
\[
K^{\bullet}(X)\times K_\bullet(X)\longrightarrow K_\bullet(X).
\]

\sgaStableTarget{sga6.e10.sec.1.1}\subsection*{\normalfont 1.1.}
First consider a noetherian scheme $X$ and the group $K_\bullet(X)$ of classes of coherent $\underline{O}_X$-modules. For every coherent $\underline{O}_X$-module $\underline{F}$, we shall denote its class in $K_\bullet(X)$ by $\cl_\bullet(\underline{F})$.

\sgaStableTarget{sga6.e10.def.1-1-1}\textbf{Definition 1.1.1.} \uline{For every integer} $n\in\mathbb Z$, \uline{let} $\Filt_n(X)$ \uline{denote the subgroup of} $K_\bullet(X)$ \uline{generated by the classes} $\cl_\bullet(\underline{F})$ \uline{of coherent} $\underline{O}_X$\uline{-modules} $\underline{F}$ \uline{such that}
\[
\dim(\Supp(\underline{F}))\leq n.
\]
This defines an increasing filtration $(\Filt_n(X))_{n\in\mathbb Z}$ of $K_\bullet(X)$ such that $\Filt_n(X)=0$ for $n<0$.

\sgaStableTarget{sga6.e10.prop.1-1-2}\textbf{Proposition 1.1.2.} \uline{Let} $X$ \uline{be a noetherian scheme,} $\underline{F}$ \uline{a coherent} $\underline{O}_X$\uline{-module such that} $\dim(\Supp(\underline{F}))\leq n$\uline{, and} $X(n)$ ($n\in\mathbb Z$\uline{) the subset} $\{x\in X\,;\,\dim(\overline{\{x\}})=n\}$ \uline{of} $X$\uline{. For every} $x\in X$, \uline{let} $Z(x)$ \uline{denote the integral closed subscheme associated with} $\overline{\{x\}}$\uline{. Then the length of the} $\underline{O}_{X,x}$\uline{-module} $\underline{F}_x$ \uline{is finite for every} $x\in X(n)$ \uline{and nonzero for only finitely many of the} $x\in X(n)$\uline{, and one has}
%%%%%%%%%%  scan idx 533  |  folio 520  %%%%%%%%%%
\SGAPrintPage{520}%
\sgaStableTarget{sga6.e10.eq.1-1-3}\begin{equation}\tag{1.1.3}
\cl_\bullet(\underline{F})=\sum_{x\in X(n)}\longg(\underline{F}_x)\cl_\bullet(\underline{O}_{Z(x)})+w
\end{equation}
\uline{where} $w\in\Filt_{n-1}(X)$.

By immediate induction on $n$, one obtains:

\sgaStableTarget{sga6.e10.cor.1-1-4}\textbf{Corollary 1.1.4.} \uline{For every noetherian scheme} $X$ \uline{and every integer} $n$, $\Filt_n(X)$ \uline{is the subgroup of} $K_\bullet(X)$ \uline{generated by the classes} $\cl_\bullet(\underline{O}_Z)$\uline{, where} $Z$ \uline{runs through the set of integral closed subschemes of dimension} $\leq n$ \uline{in} $X$.

\uline{Proof of} 1.1.2: Let $X'=\Supp(\underline{F})$. Since $X'$ is a noetherian space of dimension $\leq n$, the elements of $X'\cap X(n)$ are maximal in $X'$, hence finite in number, and $\longg(\underline{F}_x)$ is finite for every $x\in X'\cap X(n)$. This proves the first assertion.

Let ${K}=\Coh(X)$ be the category of coherent $\underline{O}_X$-modules, and let ${K}'$ be the subset of $\operatorname{Ob}({K})$ consisting of the coherent $\underline{O}_X$-modules concentrated on $X'$ that satisfy condition (1.1.3). Clearly ${K}'$ contains the sheaves $\underline{O}_{Z(x)}$ with $x\in X'$. Thus it is enough to show that ${K}'$ is an ``exact'' subset of $\operatorname{Ob}({K})$, since the dévissage lemma (EGA III 3.1.2) then gives $\underline{F}\in{K}'$.

Thus let
\sgaStableTarget{sga6.e10.eq.1-1-5}\begin{equation}\tag{1.1.5}
0\longrightarrow \underline{G}'\longrightarrow G\longrightarrow \underline{G}''\longrightarrow0
\end{equation}
be an exact sequence of $\underline{O}_X$-modules with two of its terms in ${K}'$. Then the sheaves $\underline{G}'$, $\underline{G}$ and $\underline{G}''$ are all concentrated on $X'$, so the lengths $\ell'_x$, $\ell_x$ and $\ell''_x$ of the $\underline{O}_{X,x}$-modules $\underline{G}'_x$, $\underline{G}_x$ and $\underline{G}''_x$ are finite
%%%%%%%%%%  scan idx 534  |  folio 521  %%%%%%%%%%
\SGAPrintPage{521}%
for every $x\in X(n)$, and one has
\sgaStableTarget{sga6.e10.eq.1-1-6}\begin{equation}\tag{1.1.6}
\ell_x=\ell'_x+\ell''_x\qquad (x\in X(n)).
\end{equation}
Assertion (1.1.3) for \uline{all} the terms of (1.1.5) then follows immediately from the hypothesis and the equation
\sgaStableTarget{sga6.e10.eq.1-1-7}\begin{equation}\tag{1.1.7}
\cl_\bullet(\underline{G})=\cl_\bullet(\underline{G}')+\cl_\bullet(\underline{G}'').
\end{equation}

It is known that $K_\bullet$ is a covariant functor on the category of locally noetherian schemes with respect to proper morphisms. The following proposition asserts that our filtration is functorial:

\sgaStableTarget{sga6.e10.prop.1-1-8}\textbf{Proposition 1.1.8.} \uline{Let} $f:X'\to X$ \uline{be a proper morphism of locally noetherian schemes, and let}
\[
f_*:K_\bullet(X')\longrightarrow K_\bullet(X)
\]
\uline{be the morphism induced by} $f$\uline{. Then}
\sgaStableTarget{sga6.e10.eq.1-1-9}\begin{equation}\tag{1.1.9}
f_*(\Filt_n(X'))\subset\Filt_n(X)
\end{equation}
\uline{for every integer} $n$.

\uline{Proof}: If $\underline{F}$ is a coherent $\underline{O}_{X'}$-module, then by definition (IV 2.11) one has:
\[
f_*(\cl_\bullet(\underline{F}))=\sum_{i\in\mathbb N}(-1)^i\cl_\bullet(R^if_*(\underline{F})).
\]
Now the $\underline{O}_X$-modules $R^if_*(\underline{F})$ are concentrated on $f(\Supp(\underline{F}))$, and by (EGA IV 5.4.1 (ii)) one has
\[
\dim(\Supp(\underline{F}))\geq\dim(f(\Supp(\underline{F}))),
\]
whence the assertion.


%%%%%%%%%%  scan idx 535  |  folio 522  %%%%%%%%%%
\SGAPrintPage{522}%

\sgaStableTarget{sga6.e10.sec.1.2}\subsection*{\normalfont 1.2.}
Now let $X$ be an arbitrary scheme, and consider the group $K^{\bullet}(X)$ of classes of locally free $\underline{O}_X$-modules (henceforth understood to be of finite type). It is an augmented pre-$\lambda$-algebra over the ring $H^0(X,\mathbb Z)$ (V 3.9.1). The map $f\rightsquigarrow f\cdot1\in K^{\bullet}(X)$ allows us henceforth to identify $H^0(X,\mathbb Z)$ with a subring of $K^{\bullet}(X)$. If $\underline{E}$ is a locally free $\underline{O}_X$-module, its class in $K^{\bullet}(X)$ will be denoted by $\cl^{\bullet}(\underline{E})$. Define the operators $\lambda^k$ ($k\in\mathbb N$) by
\sgaStableTarget{sga6.e10.eq.1-2-1}\begin{equation}\tag{1.2.1}
\lambda^k(\cl^{\bullet}(\underline{E}))=\cl^{\bullet}(\Lambda_{\underline{O}}^k\underline{E})
\end{equation}
and by
\sgaStableTarget{sga6.e10.eq.1-2-2}\begin{equation}\tag{1.2.2}
\lambda^k(x+y)=\sum_{\substack{i+j=k\\ i,j\in\mathbb N}}\lambda^i(x)\lambda^j(y).
\end{equation}
(V 2.2). Recall that in the $\lambda$-filtration $(\Filt^i(X))_{i\in\mathbb Z}$ of $K^{\bullet}(X)$, one has $\Filt^i(X)=K^{\bullet}(X)$ for $i\leq0$, and that for $i\geq1$, $\Filt^i(X)$ is the ideal generated by the products
\sgaStableTarget{sga6.e10.eq.1-2-3}\begin{equation}\tag{1.2.3}
P=\gamma^{i_1}(x_1)\cdots\gamma^{i_k}(x_k)
\end{equation}
with $1\leq i_1,\ldots,i_k$; $i_1+\cdots+i_k\geq i$ and $x_1,\ldots,x_k\in I=\Ker(\varepsilon)$ ($\varepsilon$ being the augmentation homomorphism, cf. Exposé V 3.10). It was proved (V 3.10) that $\Filt^i(X)$ is even the $H^0(X,\mathbb Z)$-submodule (hence the abelian subgroup, when $X$ is connected) generated by the products (1.2.3), where one may moreover suppose that the elements $x_p$ ($p=1,\ldots,k$) have the form
\sgaStableTarget{sga6.e10.eq.1-2-4}\begin{equation}\tag{1.2.4}
x_p=\cl^{\bullet}(\underline{E}_p)-\underline{n}_p,
\end{equation}
with $\underline{E}_p$ a locally free $\underline{O}_X$-module of rank $\underline{n}_p\in H^0(X,\mathbb Z)$.


%%%%%%%%%%  scan idx 536  |  folio 523  %%%%%%%%%%
\SGAPrintPage{523}%

\sgaStableTarget{sga6.e10.sec.1.3}\subsection*{\normalfont 1.3.}
Recall (IV 2.10) that the bifunctor
\[
(\underline{E},\underline{F})\rightsquigarrow \underline{E}\otimes_{\underline{O}_X}\underline{F},
\]
where $\underline{E}$ is a locally free $\underline{O}_X$-module and $\underline{F}$ a coherent $\underline{O}_X$-module, induces a composition law
\sgaStableTarget{sga6.e10.eq.1-3-1}\begin{equation}\tag{1.3.1}
K^{\bullet}(X)\times K_\bullet(X)\longrightarrow K_\bullet(X)
\end{equation}
that makes $K_\bullet(X)$ a unital $K^{\bullet}(X)$-module. The key result of the present exposé is the following

\sgaStableTarget{sga6.e10.thm.1-3-2}\textbf{Theorem 1.3.2.} \uline{Let} $X$ \uline{be a noetherian scheme. For all} $i,j\in\mathbb Z$, \uline{the composition law} (1.3.1) \uline{induces by restriction a composition law}
\sgaStableTarget{sga6.e10.eq.1-3-3}\begin{equation}\tag{1.3.3}
\Filt^i(X)\times\Filt_j(X)\longrightarrow\Filt_{j-i}(X).
\end{equation}

\sgaStableTarget{sga6.e10.cor.1-3-3}\textbf{Corollary 1.3.3.} \uline{If} $X$ \uline{is a noetherian scheme of dimension} $\leq j$\uline{, then the canonical homomorphism}
\[
\Theta_X:K^{\bullet}(X)\longrightarrow K_\bullet(X)
\]
\uline{defined by} $\Theta_X(x)=x\cdot\cl_\bullet(\underline{O}_X)$ ($x\in K^{\bullet}(X)$\uline{) induces for every} $i\in\mathbb Z$ \uline{a homomorphism}
\[
\Filt^i(X)\longrightarrow\Filt_{j-i}(X).
\]

\uline{Proof}: First let us show that (1.3.3), for $i+j\leq k\in\mathbb Z$ and for \uline{every} integral closed subscheme of dimension $\leq j$, implies (1.3.2) for $i+j\leq k$. By (1.1.4), it is enough to show that for every $x\in\Filt^i(X)$ and for \uline{every} integral closed subscheme $Z$ of dimension $\leq j$ one has
%%%%%%%%%%  scan idx 537  |  folio 524  %%%%%%%%%%
\SGAPrintPage{524}%
\[
x\cl_\bullet(\underline{O}_Z)\in\Filt_{j-i}(X).
\]
To avoid any confusion, denote by $1_Z$ the class of $\underline{O}_Z$ in $K_\bullet(Z)$. If $\alpha:Z\to X$ is the canonical immersion, then the projection formula (IV 2.11.1.2) gives
\[
x\cl_\bullet(\underline{O}_Z)=x\alpha_*(1_Z)=\alpha_*(\alpha^*(x)1_Z).
\]
One has $\alpha^*(x)\in\Filt^i(Z)$, so by (1.1.8) and (1.3.3) for $X=Z$, $i+j\leq k$, one obtains
\[
\alpha_*(\alpha^*(x)1_Z)\in\Filt_{j-i}(X),
\]
whence the assertion.

Theorem 1.3.2 is trivial for $i+j\leq0$. Suppose by induction that it has been proved for $i+j<k$ and for every noetherian scheme $X$. By the preceding remark, it remains to prove that for every \uline{integral} noetherian scheme $X$ of dimension $j$ and every product (1.2.3) satisfying (1.2.4), one has
\sgaStableTarget{sga6.e10.eq.1-3-4}\begin{equation}\tag{1.3.4}
P\cl_\bullet(\underline{O}_X)\in\Filt_{j-i}(X).
\end{equation}
If $k>1$ in (1.2.3), the induction hypothesis and the associativity of the law (1.3.1) give the result. We may therefore suppose ($X$ being connected) that
\sgaStableTarget{sga6.e10.eq.1-3-5}\begin{equation}\tag{1.3.5}
x=P=\gamma^m(\cl^{\bullet}(\underline{E})-n),
\end{equation}
where $m\geq i$ and $\underline{E}$ is a locally free $\underline{O}_X$-module of rank $n>0$.

Distinguish two cases:

I. Suppose that one can find invertible $\underline{O}_X$-modules $\underline{L}_1,\ldots,\underline{L}_n$ such that
\[
\cl^{\bullet}(\underline{E})=\cl^{\bullet}(\underline{L}_1)+\cdots+\cl^{\bullet}(\underline{L}_n)\in K^{\bullet}(X).
\]
%%%%%%%%%%  scan idx 538  |  folio 525  %%%%%%%%%%
\SGAPrintPage{525}%
It follows (V) that $x$ is the $m$-th elementary symmetric function in the $\cl^{\bullet}(\underline{L}_r)-1$ ($r=1,\ldots,n$), hence is a homogeneous polynomial of degree $m$ with integer coefficients in the variables $(\cl^{\bullet}(\underline{L}_r)-1)$ ($r=1,\ldots,n$). Applying the induction hypothesis and the associativity of the law (1.3.1) again, we are thus reduced to proving that, for every invertible $\underline{O}_X$-module $\underline{L}$, one has
\sgaStableTarget{sga6.e10.eq.1-3-6}\begin{equation}\tag{1.3.6}
(\cl^{\bullet}(\underline{L})-1)\cl_\bullet(\underline{O}_X)\in\Filt_{j-1}(X)
\end{equation}
Since the scheme $X$ is integral, $\underline{L}$ and $\underline{O}_X$ are isomorphic to submodules of the constant $\underline{O}_X$-module $\underline{R}_X$ of rational functions on $X$. The sum $\underline{O}_X$-module $\underline{F}=\underline{O}_X+\underline{L}$ is coherent, and one obtains
\begin{align*}
(\cl^{\bullet}(\underline{L})-1)\cl_\bullet(\underline{O}_X)
&=\cl_\bullet(\underline{L})-\cl_\bullet(\underline{O}_X)\\
&=(\cl_\bullet(\underline{F})-\cl_\bullet(\underline{O}_X))-(\cl_\bullet(\underline{F})-\cl_\bullet(\underline{L}))\\
&=\cl_\bullet(\underline{F}/\underline{O}_X)-\cl_\bullet(\underline{F}/\underline{L}).
\end{align*}
But the fibers at the generic point of the sheaves $\underline{F}/\underline{O}_X$ and $\underline{F}/\underline{L}$ are zero, so their supports in $X$ are nowhere dense ($X$ being irreducible), hence of dimension $\leq j-1$. Assertion (1.3.6) follows immediately.

II. It remains to prove the general case. Introduce the flag bundle
\[
g:D\longrightarrow X
\]
of $\underline{E}$ over $X$, on which one can find invertible $\underline{O}_D$-modules $\underline{L}_1,\ldots,\underline{L}_n$ such that
\sgaStableTarget{sga6.e10.eq.1-3-7}\begin{equation}\tag{1.3.7}
g^*(\cl^{\bullet}(\underline{E}))=\cl^{\bullet}(L_1)+\cdots+\cl^{\bullet}(\underline{L}_n)\in K^{\bullet}(D).
\end{equation}
Let $e$ be the generic point of $X$, and $\underline{k}(e)$ the residue field of $e$. Then one can find a rational point $u$ of the generic fiber $D_e$ of $D$
%%%%%%%%%%  scan idx 539  |  folio 526  %%%%%%%%%%
\SGAPrintPage{526}%
over $\underline{k}(e)$, since these points are in bijective correspondence with the flags of type $(1,\ldots,1)$ ($n$ terms) of the vector space $\underline{E}_e$ of dimension $n>0$ over $\underline{k}(e)$. On the other hand, these rational points also correspond to the germs of sections of $g$ in a neighborhood of $e$, so one can find such a section $s:U\to D$ over an open subset $U$ of $X$. Let $X'=\overline{s(U)}$ be the closure of $s(U)$ with its reduced induced structure, and let $f:X'\to X$ be the restriction of $g$ to $X'$. The morphism $f$ is projective and, since it is an isomorphism over $s(U)$, it is birational.
Thus $\dim(X')\leq\dim(X)$ (EGA IV 5.6.6.1). On the other hand, $f(X')=g(X')$ is closed and contains $U$, so $f$ is surjective and $\dim(X')\geq\dim(X)$ (EGA IV 5.4.1(ii)).

The hypotheses of case I are therefore satisfied for $X'$ and $x'=f^*(x)\in\Filt^i(X')$, and one obtains
\sgaStableTarget{sga6.e10.eq.1-3-7.occ2}\begin{equation}\tag{1.3.7'}
f^*(x)\cl_\bullet(\underline{O}_{X'})\in\Filt_{j-i}(X').
\end{equation}
Taking account of the projection formula and (1.1.8), one concludes that:
\sgaStableTarget{sga6.e10.eq.1-3-8}\begin{equation}\tag{1.3.8}
x f_*(\cl_\bullet(\underline{O}_{X'}))\in\Filt_{j-i}(X).
\end{equation}

On the other hand, the canonical homomorphism
\[
\underline{O}_X\longrightarrow f_*(\underline{O}_{X'})
\]
is an isomorphism over the dense open subset $U$, and the $\underline{O}_X$-modules $R^if_*(\underline{O}_{X'})$ ($i\geq1$) are concentrated on the nowhere dense closed subset $X-U$. It follows that
\sgaStableTarget{sga6.e10.eq.1-3-9}\begin{equation}\tag{1.3.9}
f_*(\cl_\bullet(\underline{O}_{X'}))=\cl_\bullet(\underline{O}_X)+w,
%%%%%%%%%%  scan idx 540  |  folio 527  %%%%%%%%%%
\end{equation}
\SGAPrintPage{527}%
where $w\in\Filt_{j-1}(X)$. By the induction hypothesis, one has $x\cdot w\in\Filt_{j-i-1}(X)$, whence, taking account of formulas (1.3.8) and (1.3.9), the assertion
\[
x\cl_\bullet(\underline{O}_X)\in\Filt_{j-i}(X),\qquad\text{Q.E.D.}
\]

\sgaStableTarget{sga6.e10.rem.1-4}\textbf{Remark 1.4.} Let $X$ be a locally noetherian scheme. Observe that, in this section, we have used \uline{only} the following properties of the ``dimension function'' $D(x)=D_Y(x)=\dim(\overline{\{x\}})$ on proper $X$-schemes $Y$:
\begin{enumerate}[label=(\roman*)]
\item For every proper $X$-scheme $Y$ and every $x\in Y$, one has $D_Y(x)\in\mathbb Z$.
\item If $y\ne x$ is a specialization of $x$, one has
\[
D(y)<D(x).
\]
\item For every $X$-morphism $f:Y_1\to Y_2$ of proper $X$-schemes and every $y\in Y_1$, one has
\[
D_{Y_2}(f(y))\leq D_{Y_1}(y).
\]
\item Equality holds in (iii) when $f$ is an immersion, or when $Y_1$ is irreducible, $f$ is birational and surjective, and $y$ is the generic point of $Y_1$.
\end{enumerate}

\uline{This means that all the results of this section, in particular} (1.3.2) \uline{and} (1.3.3)\uline{, remain valid unchanged when} ``\uline{dimension}'' \uline{is replaced by a function} $D$ \uline{satisfying conditions} (i) to (iv) \uline{above and the filtration of} $K_\bullet(X)$ \uline{is defined using such a function} $D$.


%%%%%%%%%%  scan idx 541  |  folio 528  %%%%%%%%%%
\SGAPrintPage{528}%

\sgaStableTarget{sga6.e10.ex.1-5}\textbf{Example 1.5.} Let $S$ be a locally noetherian universally catenary scheme (for example, a subscheme of a Cohen--Macaulay scheme). Put
\sgaStableTarget{sga6.e10.eq.1-5-1}\begin{equation}\tag{1.5.1}
D_S(x)=n-\codim(\overline{\{x\}},S)=n-\dimdot\underline{O}_{S,x}\qquad(x\in S),
\end{equation}
and, for every proper $S$-scheme $s_Y:Y\to S$ and $y\in Y$, also put:
\sgaStableTarget{sga6.e10.eq.1-5-2}\begin{equation}\tag{1.5.2}
D_Y(y)=D_S(s_Y(y))+\degtr_{k(s_Y(y))}k(y).
\end{equation}

If $f:Y_1\to Y_2$ is an $S$-morphism of proper $S$-schemes, the additivity of transcendence degrees and (1.5.2) immediately give
\sgaStableTarget{sga6.e10.eq.1-5-3}\begin{equation}\tag{1.5.3}
D_{Y_1}(y)=D_{Y_2}(f(y))+\degtr_{k(f(y))}k(y),\qquad(y\in Y_1).
\end{equation}
Let us show that the functions $D_Y$ so defined satisfy conditions (i)--(iv) of 1.4. By (1.5.3), condition (ii) is the only nontrivial one. By (iv), one may clearly suppose that $Y=\overline{\{x\}}$. Then factor the structural morphism $s_Y$ of $Y$ through its image $s_Y(Y)=Z\hookrightarrow S$, say $s_Y=g\circ f$:
\[
\begin{tikzcd}
Y \arrow[r,"f"] & Z \arrow[r,hook,"g"] & S,
\end{tikzcd}
\]
where $Y$ and $Z$ are irreducible, $f$ is proper and surjective, and $g$ is a closed immersion. Put $x'=s_Y(x)$, $y'=s_Y(y)$. Then (1.5.2) and (1.5.1) give:
\[
D_Y(x)-D_Y(y)=\dimdot\underline{O}_{S,y'}-\dimdot\underline{O}_{S,x'}+\degtr_{k(x')}k(x)-\degtr_{k(y')}k(y).
\]
Since $S$ is catenary, one has
%%%%%%%%%%  scan idx 542  |  folio 529  %%%%%%%%%%
\SGAPrintPage{529}%
\[
\dimdot\underline{O}_{S,y'}-\dimdot\underline{O}_{S,x'}=\dimdot\underline{O}_{Z,y'}-\dimdot\underline{O}_{Z,x'}
\]
(EGA $0_{IV}$ 14.3.2.1), whence
\sgaStableTarget{sga6.e10.eq.1-5-4}\begin{equation}\tag{1.5.4}
D_Y(x)-D_Y(y)=
[\dimdot\underline{O}_{Z,y'}-\degtr_{k(y')}k(y)]-[\dimdot\underline{O}_{Z,x'}-\degtr_{\underline{k}(x')}\underline{k}(x)].
\end{equation}
Since $Z$ is universally catenary, (EGA IV 5.6.5.1.) applies and gives
\[
C_1=\dimdot\underline{O}_{Y,y}-e,\qquad C_2=\dimdot\underline{O}_{Y,x}-e,
\]
where $C_1$ and $C_2$ are the two terms on the right in (1.5.4). Then (1.5.4) immediately gives
\[
D_Y(x)-D_Y(y)=\dimdot\underline{O}_{Y,y}-\dimdot\underline{O}_{Y,x}>0,\qquad \text{Q.E.D.}
\]

Thus (1.3.2) and (1.3.3) remain valid for a noetherian universally catenary scheme $X$ of dimension $n$ when the filtration of $K_\bullet(X)$ is defined by means of (1.5.1) according to the \uline{codimension} of supports. Observe that this filtration coincides with that of (1.1.1) when, in addition, $X$ is \uline{biequidimensional}.

\sgaStableTarget{sga6.e10.sec.1.6}\subsection*{\normalfont 1.6.}
One may ask whether it is always possible to define, on a noetherian scheme $X$ of finite dimension, a function $D$ having property (i) of 1.4 and also the following property, stronger than (ii):
\begin{quote}
(ii bis) If $x\in X$ and $y\ne x$ is an immediate specialization of $x$, then $D(x)=D(y)+1$.
\end{quote}


%%%%%%%%%%  scan idx 543  |  folio 530  %%%%%%%%%%
\SGAPrintPage{530}%

For this it is clearly necessary that $X$ be catenary, and sufficient that $X$ be biequidimensional (EGA $0_{IV}$ 14.3.3). Nevertheless, there exist noetherian schemes of finite dimension, catenary and equidimensional, having no such functions $D$. To see this, one can modify the example (EGA IV 5.6.11) as follows: take two copies $Y_1$ and $Y_2$ of the spectrum of the ring $E$ of loc. cit., and denote by $y_i$ and $z_i$ ($i=1,2$) the closed points of $Y_i$ ($i=1,2$) corresponding to the maximal ideals $\underline{m}$ and $\underline{m}'$ of $E$, of heights $2$ and $1$ respectively. Recall that the residue fields of $\underline{m}$ and $\underline{m}'$ are isomorphic. Thus one can construct a scheme $X$ by gluing $Y_1$ and $Y_2$, identifying $y_1$ with $z_2$ and $y_2$ with $z_1$. One can then identify $Y_1$ and $Y_2$ with the two irreducible components of $X$, and then, inside $X$, $Y_1\cap Y_2=\{y_1,y_2\}$ has dimension zero. Thus $X$ is a noetherian, catenary and equidimensional scheme (of dimension $2$). If (ii bis) held on $X$, one would have
\[
D(y_1)=D(z_1)-1=D(y_2)-1
\]
and, on the other hand,
\[
D(y_1)=D(z_2)=D(y_2)+1,
\]
giving a contradiction.


%%%%%%%%%%  scan idx 544  |  folio 531  %%%%%%%%%%
\SGAPrintPage{531}%

\sgaStableTarget{sga6.e10.sec.2}\section*{2. Snapper polynomials}

\sgaStableTarget{sga6.e10.sec.2.1}\subsection*{\normalfont 2.1.}
First recall the definition and some elementary properties of polynomial maps:

\sgaStableTarget{sga6.e10.def.2-1-1}\textbf{Definition 2.1.1.} \uline{Let} $f$ \uline{be a map from an abelian group} $A$ \uline{to an abelian group} $B$\uline{. We say that} $f$ \uline{is a} polynomial map \uline{when, for every sequence} $x_1,\ldots,x_r$ ($r\geq1$\uline{) of elements of} $A$\uline{, there is a family} $(y(x_1,\ldots,x_r;k_1,\ldots,k_r))_{k_1,\ldots,k_r\geq0}$ \uline{of elements of} $B$\uline{, almost all zero, satisfying the following condition: For every sequence} $n_1,\ldots,n_r$ \uline{of integers, one has}
\[
f(n_1x_1+\cdots+n_rx_r)=
\sum_{k_1,\ldots,k_r\geq0}(n_1^{k_1}\cdots n_r^{k_r})
y(x_1,\ldots,x_r;k_1,\ldots,k_r).
\]
\uline{We say that} $f$ \uline{has degree} $\leq m$ \uline{when, for all} $r\in\mathbb Z$ \uline{and} $x_1,\ldots,x_r\in A$\uline{, one has}
\[
y(x_1,\ldots,x_r;k_1,\ldots,k_r)=0
\]
\uline{for} $k_1+\cdots+k_r>m$.

\sgaStableTarget{sga6.e10.lem.2-1-2}\textbf{Lemma 2.1.2.} \uline{Every homomorphism of abelian groups is a polynomial map of degree} $\leq1$\uline{. If} $h=f\circ g$\uline{, where} $f$ \uline{is a polynomial map of degree} $p$ \uline{and} $g$ \uline{a polynomial map of degree} $q$\uline{, then} $h$ \uline{is a polynomial map of degree} $\leq pq$.

The proof is trivial.

\sgaStableTarget{sga6.e10.lem.2-1-3}\textbf{Lemma 2.1.3.} \uline{Let}
\[
0\longrightarrow G'\longrightarrow G\xrightarrow{j}G''\longrightarrow0
\]
\uline{be an exact sequence of abelian groups such that} $G'$ \uline{is divisible by}
%%%%%%%%%%  scan idx 545  |  folio 532  %%%%%%%%%%
\SGAPrintPage{532}%
\uline{a prime number} $p$ \uline{and} $G''$ \uline{is a torsion group. Then every polynomial map} $f:G\to\mathbb Q$ \uline{that factors through} $\mathbb Z$ \uline{is constant.}

\uline{Proof}: We must show that for every $x\in G$ one has $f(x)=f(0)$.

a) First suppose that $x$ is a torsion element, say $mx=0$ ($m\in Z$, $m\ne0$). Then the polynomial $P_x(n)=f(nx)-f(0)$ has infinitely many roots, namely all integers divisible by $m$. Thus the polynomial $P_x$ is identically zero, and one has $f(x)=f(0)$.

b) Now suppose that $x$ is an element of infinite order in $G$ divisible by every power $p^k$ ($k\in\mathbb Z$). Consider the polynomials $P_k$ defined by
\[
P_k(n)=f(nx/p^k)=\sum_{i\in\mathbb N}n^iy(x/p^k;i)\in\mathbb Z
\qquad(n\in\mathbb Z;\ y(x/p^k,i)\in\mathbb Q).
\]
For every $n\in\mathbb Z$ one has the identity
\[
P_k(n)=P_{k+1}(pn)
\]
i.e.
\[
\sum_{i\in\mathbb N}n^iy(x/p^k;i)=\sum_{i\in\mathbb N}p^in^iy(x/p^{k+1};i),
\]
from which comparison of coefficients gives
\[
y(x/p^k;i)=p^iy(x/p^{k+1};i)\qquad(i\in\mathbb N).
\]
%%%%%%%%%%  scan idx 546  |  folio 533  %%%%%%%%%%
\SGAPrintPage{533}%
But every polynomial $P:Q\to\mathbb Q$ is continuous for the topology induced by $\mathbb R$. Thus for sufficiently large $k$ one has
\[
|P_x(n/p^k)-P_x(0)|<1.
\]
By definition, however, $P_x(n/p^k)$ is an integer for every $n,k\in\mathbb Z$. It follows that $P_x(n/p^k)-P_x(0)=0$ for sufficiently large $k$. Thus the polynomial $P_x$ is constant, and in particular one obtains $P_x(1)=P_x(0)$, i.e. $f(x)=f(0)$.

The two cases a) and b) above show that the restriction of $f$ to $G'$ is constant. We shall now show that $f$ is constant on each coset in $G$ modulo $G'$. For every $b\in G$, denote by $g_b:G'\to\underline{Q}$ the map defined by
\[
g_b(y)=f(b+y)\qquad(y\in G').
\]
The map $g_b$ is a polynomial map from $G'$ to $\mathbb Q$ that factors through $\mathbb Z$, hence is constant by the preceding argument, whence the assertion that $f$ is constant on the cosets in $G$ modulo $G'$. It follows that one can write $f=h\circ j$, where $h:G''\to\mathbb Q$ is a polynomial map that factors through $\mathbb Z$. But by hypothesis $G''$ is a torsion group, so $h$ is constant. Thus $f$ is constant, Q.E.D.

\sgaStableTarget{sga6.e10.lem.2-1-4}\textbf{Lemma 2.1.4.} \uline{Let} $A$ \uline{be a unital commutative ring and} $I$ \uline{a nilpotent ideal of} $A$\uline{. Let} $P$ \uline{be the multiplicative subgroup of} $A^*$ \uline{consisting of the invertible elements of} $A$ \uline{of the form} $x=1+y$ \uline{with} $y\in I$\uline{, and let} $i:P\to A$ \uline{be the natural injection. Then the composite map}
%%%%%%%%%%  scan idx 547  |  folio 534  %%%%%%%%%%
\SGAPrintPage{534}%
\[
\alpha:P\xrightarrow{i}A\xrightarrow{j}A\otimes_{\mathbb Z}\mathbb Q
\]
\uline{(where} $j(x)=x\otimes1$\uline{) is a polynomial map. If} $I^{m+1}=0$\uline{, then} $\alpha$ \uline{has degree} $\leq m$.

\uline{Proof}: Let $x_1,\ldots,x_r$ be a sequence of elements of $P$ with $x_i=1+y_i$ ($i=1,\ldots,r$). Then the binomial formula gives, for every $n_1,\ldots,n_r\in Z$:
\sgaStableTarget{sga6.e10.eq.2-1-5}\begin{equation}\tag{2.1.5}
x_1^{n_1}\cdots x_r^{n_r}=
\sum_{k_1,\ldots,k_r\geq0}\binom{n_1}{k_1}\cdots\binom{n_r}{k_r}y_1^{k_1}\cdots y_r^{k_r},
\end{equation}
where the sum on the right is finite because the elements $y_1,\ldots,y_r$ are nilpotent. But the coefficients $\binom{n_1}{k_1}\cdots\binom{n_r}{k_r}$ are polynomials with rational coefficients in the variables $n_1,\ldots,n_r$, of degree $k_1+\cdots+k_r$, so $\alpha$ is indeed a polynomial map. If $I^{m+1}=0$, then $y_1^{k_1}\cdots y_r^{k_r}=0$ for $k_1+\cdots+k_r>m$. Thus $\alpha$ has degree $\leq m$, Q.E.D.

\sgaStableTarget{sga6.e10.sec.2.2}\subsection*{\normalfont 2.2.}
Now consider a noetherian scheme $X$ and an element $y$ of $\Filt_j(X)$ (1.1.1). Associating to every $x\in K^{\bullet}(X)$ the element $xy$ of $K_\bullet(X)$ gives a homomorphism of $K^{\bullet}(X)$-modules
\[
\Theta_y:K^{\bullet}(X)\longrightarrow K_\bullet(X)
\]
which, taking account of (1.3.2), factors through the canonical epimorphism
\[
j:K^{\bullet}(X)\longrightarrow K^{\bullet}(X)/\Filt^{j+1}(X)=A.
\]
%%%%%%%%%%  scan idx 548  |  folio 535  %%%%%%%%%%
\SGAPrintPage{535}%
Let $\beta_y:A\to K_\bullet(X)$ be the homomorphism so defined. Denoting by $I$ the ideal $\Filt^1(X)/\Filt^{j+1}(X)$ of $A$, one has $I^{j+1}=0$. Keeping the notation of Lemma 2.1.4, one obtains a homomorphism of abelian groups
\[
\gamma:\Pic(X)\longrightarrow P
\]
that associates to the class in $\Pic(X)$ of every invertible $\underline{O}_X$-module $\underline{L}$ the element $j(\cl^{\bullet}(\underline{L}))$ of $P$.

Lemmas 2.1.2 and 2.1.4 immediately give:

\sgaStableTarget{sga6.e10.prop.2-2-1}\textbf{Proposition 2.2.1.} \uline{For every noetherian scheme} $X$ \uline{and every} $y\in\Filt_j(X)$ ($j\in\mathbb Z$\uline{), the composite map}
\[
\Pic(X)\longrightarrow K_\bullet(X)\otimes_{\mathbb Z}\mathbb Q,
\]
\uline{defined by}
\[
\cl_P(\underline{L})\rightsquigarrow (\cl^{\bullet}(\underline{L})y)\otimes1
\]
\uline{is a polynomial map of degree} $\leq j$.

\sgaStableTarget{sga6.e10.sec.2.3}\subsection*{\normalfont 2.3.}
If $X$ is a proper scheme over a field $k$ and $g:X\to\Spec(k)$ is the structural morphism, then the homomorphism
\[
g_*:K_\bullet(X)\longrightarrow K_\bullet(\Spec(k))\cong\mathbb Z
\]
associates to the class $\cl_\bullet(\underline{F})$ of every coherent $\underline{O}_X$-module $\underline{F}$ the Euler--Poincaré characteristic $\chi(\underline{F})\in\mathbb Z$ of $\underline{F}$ on $X$. From 2.2.1 one immediately concludes:


%%%%%%%%%%  scan idx 549  |  folio 536  %%%%%%%%%%
\SGAPrintPage{536}%

\sgaStableTarget{sga6.e10.prop.2-3-1}\textbf{Proposition 2.3.1.} \uline{For every proper scheme} $X$ \uline{over} $k$\uline{, and every} $y\in\Filt_j(X)$\uline{, the map}
\[
s_y:\Pic(X)\longrightarrow\mathbb Z\hookrightarrow\mathbb Q
\]
\uline{defined by}
\[
\cl_P(\underline{L})\rightsquigarrow g_*(\cl^{\bullet}(\underline{L})y)=\chi(\cl^{\bullet}(\underline{L})y)
\]
\uline{is a polynomial map of degree} $\leq j$.

When $\underline{L}_1,\ldots,\underline{L}_n$ are invertible $\underline{O}_X$-modules and $\underline{F}$ is a coherent $\underline{O}_X$-module such that $\dim(\Supp(\underline{F}))=j$, then 2.3.1. implies in particular that the Euler--Poincaré characteristic
\sgaStableTarget{sga6.e10.eq.2-3-2}\begin{equation}\tag{2.3.2}
\chi(\underline{L}_1^{\otimes m_1}\otimes\cdots\otimes \underline{L}_n^{\otimes m_n}\otimes \underline{F})
\end{equation}
is a numerical polynomial in the variables $m_1,\ldots,m_n\in\mathbb Z$, of degree $\leq\dimdot\operatorname{supp.}F$, sometimes called the \uline{Snapper polynomial} for $\underline{F}$ and the sequence $\underline{L}_1,\ldots,\underline{L}_n$.

\sgaStableTarget{sga6.e10.sec.2.4}\subsection*{\normalfont 2.4.}
Recall [2] that the Picard functor $\text{\uwave{Pic}}_{X/k}$ of $X$ over $k$ is the sheaf of groups for the ($\fpqc$) topology over $k$ associated with the presheaf
\[
T\rightsquigarrow \Pic(X\times_kT)\qquad (T\in\operatorname{Ob}(\Sch/k)).
\]
When $k$ is algebraically closed, the canonical homomorphism
\[
\Pic(X)\longrightarrow\Pic_{X/k}(\Spec(k))
\]
is known to be an isomorphism.

% Local continuation macros
\providecommand{\Filt}{\operatorname{Fil}}
\providecommand{\fpqc}{\mathrm{fpqc}}
\providecommand{\Sch}{\operatorname{Sch}}
\renewcommand{\Ob}{\operatorname{ob}}
\providecommand{\Gm}{\mathbb G_m}
\providecommand{\Der}{\operatorname{Der}}
\providecommand{\SL}{\operatorname{SL}}
\providecommand{\Aut}{\operatorname{Aut}}
\providecommand{\Mod}{\operatorname{Mod}}
\providecommand{\Tors}{\operatorname{Tors}}
\providecommand{\Mat}{\operatorname{Mat}}
\providecommand{\len}{\operatorname{long}}
\renewcommand{\id}{\operatorname{id}}
\providecommand{\Spec}{\operatorname{Spec}}
\providecommand{\Picf}{\operatorname{Pic}}
\providecommand{\Ens}{\operatorname{Ens}}
\providecommand{\clp}{\operatorname{cl}_{P}}
\providecommand{\clo}{\operatorname{cl}^{\circ}}
\providecommand{\clb}{\operatorname{cl}_{\bullet}}
\providecommand{\Grb}{\operatorname{Gr}_{\bullet}}
\providecommand{\Gro}{\operatorname{Gr}^{\bullet}}
\providecommand{\longO}[1]
{\operatorname{length}_{#1}}

\renewcommand{\Gro}{\operatorname{Gr}^{\bullet}}

%%%%%%%%%%  scan idx 550  |  folio 537  %%%%%%%%%%
\SGAPrintPage{537}%

When \(X\) is proper over \(k\), it is known [4] that \(\text{\uwave{Pic}}_{X/k}\) is representable by a group scheme \(\underline{\Pic}_{X/k}\) locally of finite type over \(k\). The neutral component \(\underline{\Pic}^0_{X/k}\) of \(\underline{\Pic}_{X/k}\) is a connected commutative algebraic group over \(k\).

\sgaStableTarget{sga6.e10.def.2-4-1}\textbf{Definition 2.4.1.} \uline{Let} \(X\) \uline{be a proper scheme over a field} \(k\), \uline{and} \(\underline{L}\) \uline{an invertible} \(\underline{O}_X\)\uline{-module} (resp. \(x\) \uline{an element of} \(\Pic(X)\)\uline{). We say that} \(\underline{L}\) (resp. \(x\)\uline{) is} \(\tau\)\uline{-equivalent to zero when there is an integer} \(n\ne0\) \uline{such that} \(\underline{L}^{\otimes n}\) (resp. \(nx\)\uline{) is algebraically equivalent to zero. The elements of} \(\Pic(X)\) \uline{algebraically equivalent} (resp. \(\tau\)\uline{-equivalent) to zero form a subgroup} \(\Pic^0(X)\) (resp. \(\Pic^\tau(X)\)\uline{) of} \(\Pic(X)\).

\sgaStableTarget{sga6.e10.lem.2-4-3}\textbf{Lemma 2.4.3.} \uline{For every proper scheme} \(X\) \uline{over an algebraically closed field} \(k\), \uline{the group} \(\Pic^0(X)\) \uline{is divisible by every integer} \(n\) \uline{prime to the characteristic of} \(k\).

\uline{Proof}: Since \(k\) is algebraically closed, one has \(\Pic^0(X)\cong \underline{\Pic}^0_{X/k}(k)\). Since \(\underline{\Pic}^0_{X/k}\) is connected and of finite type, the \(n\)-th-power endomorphism is surjective when \(n\) is prime to the characteristic of \(k\) (SGA 3 XV 1.3 and VI 1.3.2). But \(k\) is algebraically closed, so the corresponding endomorphism of the group \(\underline{\Pic}^0_{X/k}(k)\) is surjective, Q.E.D.

From 2.4.3 one concludes that, in the exact sequence
\[
0\longrightarrow\Pic^0(X)\longrightarrow\Pic^\tau(X)\longrightarrow T\longrightarrow0,
\]
\(\Pic^0(X)\) is divisible by a prime number \(p\) and \(T\) is a torsion group, so one can apply 2.1.3 to the restriction of the polynomial map (2.3.1) to \(\Pic^\tau(X)\). One thus obtains:


%%%%%%%%%%  scan idx 551  |  folio 538  %%%%%%%%%%
\SGAPrintPage{538}%

\sgaStableTarget{sga6.e10.prop.2-4-4}\textbf{Proposition 2.4.4.} \uline{The polynomial map} \(s_y\) (2.3.1) \uline{is constant on the cosets modulo} \(\Pic^\tau(X)\).

\uline{Proof}: If \(K\) is an algebraically closed extension of \(k\), one obtains a commutative diagram
\[
\begin{tikzpicture}[>=stealth,baseline=(current bounding box.center)]
\node (PicX) at (0,2) {\(\Pic(X)\)};
\node (Kk) at (3.7,2) {\(K_\bullet(\Spec(k))\)};
\node (Z) at (5.35,1) {\(\mathbb Z\)};
\node (Q) at (6.6,1) {\(\mathbb Q\)};
\node (PicXK) at (0,0) {\(\Pic(X_K)\)};
\node (KK) at (3.7,0) {\(K_\bullet(\Spec(K))\)};
\draw[->] (PicX) -- node[above] {\(s_y\)} (Kk);
\draw[->] (PicX) -- (PicXK);
\draw[->] (PicXK) -- node[above] {\(\underline{s}_y\)} (KK);
\draw[->] (Kk) -- node[sloped,above] {\(\sim\)} (Z);
\draw[->] (KK) -- node[sloped,above] {\(\sim\)} (Z);
\draw[right hook->] (Z) -- (Q);
\end{tikzpicture},
\]
where \(\underline{s}_y\) is the polynomial map on \(\Pic(X_K)\) corresponding to \(s_y\). Thus one may suppose that \(k\) is algebraically closed. The assertion then follows immediately from 2.1.3 and 2.4.3.

\sgaStableTarget{sga6.e10.rem.2-4-5}\textbf{Remark 2.4.5.} In Exp. XIII this result will be proved without using the representability of the functor \(\Pic_{X/k}\).

\sgaStableTarget{sga6.e10.sec.3}\section*{3. Projection formulas for the associated graded groups}

\sgaStableTarget{sga6.e10.sec.3.1}\subsection*{\normalfont 3.1.}
Consider a noetherian scheme \(X\). Denote by \(\Gro(X)\) the graded group associated with the \(\lambda\)-filtration of \(K^{\bullet}(X)\), and by \(\Grb(X)\) the graded group associated with the filtration 1.1.1 of \(K_\bullet(X)\). By 1.3.2, the composition law (1.3.1)
\[
 K^{\bullet}(X)\times K_\bullet(X)\longrightarrow K_\bullet(X)
\]
induces, on passing to the associated graded groups, composition laws
\sgaStableTarget{sga6.e10.eq.3-1-1}\begin{equation}\tag{3.1.1}
 \Gr^i(X)\times\Gr_j(X)\longrightarrow \Gr_{j-i}(X)\qquad(i,j\in\mathbb Z).
\end{equation}
%%%%%%%%%%  scan idx 552  |  folio 539  %%%%%%%%%%
\SGAPrintPage{539}%
\(\Grb(X)\) thus becomes a module over the ring \(\Gro(X)\) (and even a graded module over the graded ring \(\Gro(X)\), provided the sign of the grading of \(\Grb(X)\) is changed by putting \(\Gr^i_\bullet(X)=\Gr_{-i}(X)\) \((i\in\mathbb Z)\)).
Now consider a proper morphism
\[
 f:X'\longrightarrow X
\]
of noetherian schemes. Then \(f\) induces a homomorphism
\[
 f^*:\Gro(X)\longrightarrow\Gro(X')
\]
of graded rings (V), and, by 1.1.8, a homomorphism
\[
 f_*:\Grb(X')\longrightarrow\Grb(X)
\]
of graded abelian groups. Passing to the associated graded groups, the projection formula (IV 2.11) and (3.1.1) then give the graded projection formula
\sgaStableTarget{sga6.e10.eq.3-1-2}\begin{equation}\tag{3.1.2}
 x\cdot f_*(y)=f_*\bigl(f^*(x)\cdot y\bigr)
 \qquad (x\in\Gro(X),\ y\in\Grb(X')).
\end{equation}

\sgaStableTarget{sga6.e10.sec.3.2}\subsection*{\normalfont 3.2.}
Let \(f:X'\to X\) be a projective complete-intersection morphism (Exp. VIII) of virtual dimension \(r\) between noetherian schemes, with \(X\) carrying an ample invertible \(\underline{O}_X\)-module. In Exp. VIII it was proved that \(f\) induces a homomorphism
\sgaStableTarget{sga6.e10.eq.3-2-1}\begin{equation}\tag{3.2.1}
 f_*:\Filt^k(X')_{\mathbb Q}\longrightarrow\Filt^{k-r}(X)_{\mathbb Q}\qquad(k\in\mathbb Z)
\end{equation}
which moreover satisfies the projection formula (Exp. IV):
\sgaStableTarget{sga6.e10.eq.3-2-2}\begin{equation}\tag{3.2.2}
 f_*(f^*(x)\cdot y)=x\cdot f_*(y)
 \qquad (x\in K^{\bullet}(X)_{\mathbb Q},\ y\in K^{\bullet}(X')_{\mathbb Q}).
\end{equation}
%%%%%%%%%%  scan idx 553  |  folio 540  %%%%%%%%%%
\SGAPrintPage{540}%
Passing to the associated graded groups in (3.2.1) and (3.2.2), one obtains a homomorphism of degree \(-r\):
\sgaStableTarget{sga6.e10.eq.3-2-3}\begin{equation}\tag{3.2.3}
 f_*:\Gro(X')_{\mathbb Q}\longrightarrow\Gro(X)_{\mathbb Q}
\end{equation}
and, by (1.3.2), a graded projection formula
\[
 f_*\bigl(f^*(x)\cdot y\bigr)=x\cdot f_*(y)
 \qquad (x\in\Gro(X)_{\mathbb Q},\ y\in\Gro(X')_{\mathbb Q}).
\]

\sgaStableTarget{sga6.e10.sec.3.3}\subsection*{\normalfont 3.3.}
Finally, consider a morphism \(f:X'\to X\) of finite Tor-dimension between noetherian schemes. Then it is known (Exp. IV) that \(f\) induces a homomorphism
\[
 f^*:K_\bullet(X)\longrightarrow K_\bullet(X')
\]
such that
\[
 f^*(\cl_\bullet(\underline{F}))=
 \sum_i(-1)^i\cl_\bullet\bigl(\underline{\Tor}_i^{\underline{O}_X}(\underline{F},\underline{O}_{X'})\bigr)
\]
for every coherent \(\underline{O}_X\)-module \(\underline{F}\).

The \(\underline{O}_X\)-modules \(\underline{\Tor}_i^{\underline{O}_X}(\underline{F},\underline{O}_{X'})\) are concentrated on \(f^{-1}(\Supp(\underline{F}))\), and when the fibers of \(f\) are everywhere of dimension \(\leq d\), it is known (EGA IV 5.6.7) that
\[
 \dim(f^{-1}(\Supp(\underline{F})))\leq \dim(\Supp(\underline{F}))+d.
\]
It follows that \(f\) induces, for every \(k\in\mathbb Z\), a homomorphism
\[
 f^*:\Filt_k(X)\longrightarrow\Filt_{k+d}(X'),
\]
whence a homomorphism of degree \(d\) of graded abelian groups
%%%%%%%%%%  scan idx 554  |  folio 541  %%%%%%%%%%
\SGAPrintPage{541}%
\sgaStableTarget{sga6.e10.eq.3-3-1}\begin{equation}\tag{3.3.1}
 f^*:\Grb(X)\longrightarrow\Grb(X').
\end{equation}

\sgaStableTarget{sga6.e10.sec.3.4}\subsection*{\normalfont 3.4.}
Finally, suppose that the morphisms (3.2.3) and (3.3.1) are both defined for a morphism \(f:X'\to X\) such that \(r=d\). This is the case, for example, when \(X\) is quasi-compact with an ample \(\underline{O}_X\)-module and \(X'=P_X(\underline{E})\), where \(\underline{E}\) is a locally free \(\underline{O}_X\)-module of rank \(r+1\). Then, for every \(x\in\Filt^i(X')_{\mathbb Q}\) and \(y\in\Filt_j(X)_{\mathbb Q}\), the projection formula in the derived category gives the projection formula
\[
 f_*\bigl(x\cdot f^*(y)\bigr)=f_*(x)\cdot y\in\Filt_{j+r}(X)_{\mathbb Q}.
\]
When \(X\) is noetherian, 1.3.2 yields the corresponding graded formula.

\sgaStableTarget{sga6.e10.sec.4}\section*{4. Intersection numbers}

\sgaStableTarget{sga6.e10.sec.4.1}\subsection*{\normalfont 4.1.}
Consider a proper scheme \(X\) over a field \(k\). Composing the product (3.1.1)
\[
 \Gr^i(X)\times\Gr_i(X)\longrightarrow \Gr_0(X)=\Filt_0(X)
 \qquad (i\in\mathbb Z)
\]
with the Euler--Poincaré characteristic
\[
 \chi:\Filt_0(X)\longrightarrow\mathbb Z,
\]
gives a pairing called the ``\uline{intersection number}'':
\sgaStableTarget{sga6.e10.eq.4-1-1}\begin{equation}\tag{4.1.1}
 \Gr^i(X)\times\Gr_i(X)\longrightarrow \mathbb Z,
\end{equation}
%%%%%%%%%%  scan idx 555  |  folio 542  %%%%%%%%%%
\SGAPrintPage{542}%
which we shall denote by
\[
 (x,y)\rightsquigarrow \langle x,y\rangle\in\mathbb Z
 \qquad (i\in\mathbb Z,\ x\in\Gr^i(X),\ y\in\Gr_i(X)).
\]

\sgaStableTarget{sga6.e10.ex.4-2}\textbf{Example 4.2.} Suppose that the scheme \(X\) has dimension \(n\). Take an element \(y\in\Gr_n(X)\), for example \(y=\cl_\bullet(\underline{O}_Z)\), where \(Z\) is an integral subscheme of dimension \(n\) in \(X\). Also take a homogeneous polynomial of weight \(n\)
\[
 P\in\mathbb Z[T_1,\ldots,T_n],
\]
where \(T_i\) is assigned weight \(i\) \((i=1,\ldots,n)\), for example a monomial \(T_{i_1}\cdot T_{i_2}\cdots T_{i_k}\) with \(i_1+i_2+\cdots+i_k=n\). Substituting in \(P\), for each variable \(T_i\) \((i=1,\ldots,n)\), the Chern class
\[
 c_i(x)=\cl^{\bullet}(\gamma^i(x-\epsilon(x)))\in\Gr^i(X)
 \qquad(x\in K^{\bullet}(X)),
\]
gives an element
\[
 P(c_1(x),\ldots,c_n(x))\in\Gr^n(X).
\]
The integer
\[
 P_y(x)=\bigl\langle P(c_1(x),\ldots,c_n(x)),y\bigr\rangle.
\]
is called the \uline{Chern number of} \(x\) \uline{relative to} \(y\).
When \(y=\cl_\bullet(\underline{O}_X)\), the reference to \(y\) is omitted from this terminology.

\sgaStableTarget{sga6.e10.sec.4.3}\subsection*{\normalfont 4.3.}
Without any hypothesis on the dimension of \(X\), consider a sequence \(\underline{L}_1,\ldots,\underline{L}_d\) of invertible \(\underline{O}_X\)-modules and a coherent \(\underline{O}_X\)-module \(\underline{F}\) such that \(\dim(\Supp(\underline{F}))\leq d\). Put \(L_i=\cl^{\bullet}(\underline L_i)\) \((i=1,\ldots,d)\) and \(F=\cl_\bullet(\underline F)\).


%%%%%%%%%%  scan idx 556  |  folio 543  %%%%%%%%%%
\SGAPrintPage{543}%

\sgaStableTarget{sga6.e10.prop.4-3-1}\textbf{Proposition 4.3.1.} \uline{With the above notation, for every sequence} \(\alpha_1,\ldots,\alpha_d\) \uline{of nonnegative integers with} \(\sum_{i=1}^{d}\alpha_i=d\)\uline{, the coefficient of}
\[
 \binom{n_1}{\alpha_1}\cdots\binom{n_d}{\alpha_d}
\]
\uline{in the Snapper polynomial} \(\chi(L_1^{n_1}\cdots L_d^{n_d}F)\) \uline{is equal to}
\sgaStableTarget{sga6.e10.eq.4-3-2}\begin{equation}\tag{4.3.2}
 \bigl\langle c_1(L_1)^{\alpha_1}\cdots c_1(L_d)^{\alpha_d},\lvert\underline{F}\rvert\bigr\rangle,
\end{equation}
\uline{where} \(\lvert\underline{F}\rvert\) \uline{is the class of} \(\underline{F}\) \uline{in} \(\Gr_d(X)\)\uline{. Thus, in the ordinary expansion of this polynomial in monomials in the} \(n_i\) \((i=1,\ldots,d)\)\uline{, the coefficient of the monomial} \(n_1^{\alpha_1}\cdots n_d^{\alpha_d}\) \uline{is equal to}
\[
 (1/(\alpha_1! \cdots \alpha_d!))
 \bigl\langle c_1(L_1)^{\alpha_1}\cdots c_1(L_d)^{\alpha_d},\lvert\underline{F}\rvert\bigr\rangle.
\]

In particular, one immediately obtains the following

\sgaStableTarget{sga6.e10.cor.4-3-3}\textbf{Corollary 4.3.3.} \uline{The coefficient of the monomial} \(n_1n_2\cdots n_d\) \uline{in the Snapper polynomial} \(\chi(L_1^{n_1}\cdots L_d^{n_d}F)\) \uline{is equal to}
\sgaStableTarget{sga6.e10.eq.4-3-4}\begin{equation}\tag{4.3.4}
 \bigl\langle c_1(L_1)\cdots c_1(L_d),\lvert\underline{F}\rvert\bigr\rangle
 =
 (c_1(L_1)\cdots c_1(L_d),\lvert\underline{F}\rvert).
\end{equation}

\uline{Proof}: Put \(L_i=1+M_i\) for every \(i=1,\ldots,d\). The definitions immediately show that the integer (4.3.2) is equal to \(\chi(M_1^{\alpha_1}\cdots M_d^{\alpha_d}F)\). Substituting \(A=K^{\bullet}(X)/\Filt^{d+1}(X)\), \(I=\Filt^1(X)/\Filt^{d+1}(X)\) in (2.1.4), formula (2.1.5) gives (taking account of the fact that \(M_i\in\Filt^1(X)\) for every \(i=1,\ldots,d\)) the congruence
\sgaStableTarget{sga6.e10.eq.4-3-5}\begin{equation}\tag{4.3.5}
 L_1^{n_1}\cdots L_d^{n_d}
 \equiv
 \sum_{\substack{\alpha_1,\ldots,\alpha_d\geq0\\ \alpha_1+\cdots+\alpha_d\leq d}}
 \binom{n_1}{\alpha_1}\cdots\binom{n_d}{\alpha_d}
 M_1^{\alpha_1}\cdots M_d^{\alpha_d}
 \quad\mathrm{mod.}\Filt^{d+1}(X).
\end{equation}
Thus the coefficient of \(\binom{n_1}{\alpha_1}\cdots\binom{n_d}{\alpha_d}\) in the polynomial \(\chi(L_1^{n_1}\cdots L_d^{n_d}F)\) is equal to \(\chi(M_1^{\alpha_1}\cdots M_d^{\alpha_d}F)\), Q.E.D.

%%%%%%%%%%  scan idx 557  |  folio 544  %%%%%%%%%%
\SGAPrintPage{544}%

\sgaStableTarget{sga6.e10.sec.4.4}\subsection*{\normalfont 4.4.}
The integer (4.3.4) is sometimes called the \uline{intersection number} of the invertible \(\underline{O}_X\)-modules \(\underline{L}_1,\ldots,\underline{L}_d\) with \(\underline{F}\), and is denoted by \((\underline{L}_1,\ldots,\underline{L}_d;\underline{F})\) (cf. [3]). Expression (4.3.4) for this number immediately gives the following properties:

\begin{enumerate}[label=\arabic*)]
\item The integer (4.3.4) is a symmetric, \(d\)-linear function of the variables \(L_1,\ldots,L_d\in\Pic(X)\), zero when \(\dim(\Supp(\underline{F}))<d\), and is ``additive'' in \(\underline{F}\), i.e. additive in \(\cl_\bullet(\underline F)=F\).

\item Using 2.1.2, one immediately concludes (with the notation of 1.1.2):
\sgaStableTarget{sga6.e10.eq.4-4-1}\begin{equation}\tag{4.4.1}
 (\underline{L}_1,\ldots,\underline{L}_d;\underline{F})
 =
 \sum_{x\in X(d)}
 \operatorname{long}_{\underline{O}_{X,x}}(\underline{F}_x)
 (\underline{L}_1,\ldots,\underline{L}_d;\underline{O}_{Z(x)}).
\end{equation}

\item Let \(Y\) be a closed subscheme of \(X\) and \(\underline{G}\) a coherent \(\underline{O}_Y\)-module such that \(\underline{F}=f_*(\underline{G})\), where \(f:Y\to X\) is the canonical immersion. Then one obtains
\[
 (f^*\underline{L}_1,\ldots,f^*\underline{L}_d;\underline{G})=(\underline{L}_1,\ldots,\underline{L}_d;\underline{F}).
\]
More generally, the projection formula gives
\sgaStableTarget{sga6.e10.eq.4-4-2}\begin{equation}\tag{4.4.2}
 \chi(f^*(x)\cdot y)=\chi\bigl(f_*(f^*(x)\cdot y)\bigr)=\chi(x\cdot f_*(y))
 \qquad (x\in K^{\bullet}(X),\ y\in K_\bullet(Y))
\end{equation}
for every \(k\)-morphism \(f:Y\to X\) of proper \(k\)-schemes.

\item Suppose that, in the morphism \(f:Y\to X\) above, \(X\) and \(Y\) are irreducible of dimension \(\leq d\). If \(f\) has degree \(n\), then one has
%%%%%%%%%%  scan idx 558  |  folio 545  %%%%%%%%%%
\SGAPrintPage{545}%
\sgaStableTarget{sga6.e10.eq.4-4-3}\begin{equation}\tag{4.4.3}
 (f^*\underline{L}_1,\ldots,f^*\underline{L}_d;O_Y)
 =
 n(\underline{L}_1,\ldots,\underline{L}_d;\underline{O}_X).
\end{equation}
Recall the definition of the degree \(n\) of \(f\) [3]: it is zero except when \(\dim(X)=\dim(Y)=\dim(f(Y))\). In this latter case, put
\sgaStableTarget{sga6.e10.eq.4-4-4}\begin{equation}\tag{4.4.4}
 n=(\len_{\underline{O}_{X,x}}(\underline{O}_{Y,y}))/(\len_{\underline{O}_{X,x}}(\underline{O}_{X,x})),
\end{equation}
where \(x\) is the generic point of \(X\) and \(y\) the generic point of \(Y\).
\end{enumerate}

In the proof of (4.4.3), the case \(\dim(X)=\dim(Y)\) is the only nontrivial one. If \(\dim(f(Y))<d\), then one has
\sgaStableTarget{sga6.e10.eq.4-4-5}\begin{equation}\tag{4.4.5}
 R^if_*(\underline{O}_Y)=0
\end{equation}
for \(i>0\) in a neighborhood of \(x\). But condition (4.4.5) also holds when \(\dim(f(Y))=d\), since the generic fiber then has dimension zero (EGA III 4.2.2 and EGA IV 5.6.6). Then 1.1.2 gives the congruences
\[
 f_*(\cl_\bullet(\underline{O}_Y))
 \equiv \cl_\bullet(f_*(\underline{O}_Y)) \equiv
 \len_{\underline{O}_{X,x}}(\underline{O}_{Y,y})
 \cl_\bullet(\underline{O}_{X_{\mathrm{red}}})
\]
modulo \(\Filt_{d-1}(X)\). By (4.4.2), it follows that
\[
 (f^*\underline{L}_1,\ldots,f^*\underline{L}_d;\underline{O}_Y)
 =
 \len_{\underline{O}_{X,x}}(\underline{O}_{Y,y})
 (\underline{L}_1,\ldots,\underline{L}_d;\underline{O}_{X_{\mathrm{red}}}).
\]
On the other hand, \(2^\circ\) gives
\[
 (\underline{L}_1,\ldots,\underline{L}_d;\underline{O}_X)
 =
 \len_{\underline{O}_{X,x}}(\underline{O}_{X,x})
 (\underline{L}_1,\ldots,\underline{L}_d;\underline{O}_{X_{\mathrm{red}}}),
\]
whence the assertion.


%%%%%%%%%%  scan idx 559  |  folio 546  %%%%%%%%%%
\SGAPrintPage{546}%

\sgaStableTarget{sga6.e10.sec.4.5}\subsection*{\normalfont 4.5}
The equivalence relation defined by the intersection product \(\langle\, ,\,\rangle\) is called numerical equivalence:

\sgaStableTarget{sga6.e10.def.4-5-1}\textbf{Definition 4.5.1.} \uline{Two elements} \(x\) \uline{and} \(y\) \uline{of} \(\Gr^i(X)\) \((i\in Z)\) (resp. \uline{of} \(\Gr_i(X)\)\uline{) are said to be numerically equivalent when they define the same homomorphism} \(\Gr_i(X)\longrightarrow\mathbb Z\) (resp. \(\Gr^i(X)\longrightarrow\mathbb Z\)\uline{). We say that two invertible} \(\underline{O}_X\)\uline{-modules} \(\underline{L}_1\) and \(\underline{L}_2\) \uline{are numerically equivalent when their first Chern classes} \(x,y\in\Gr^1(X)\) \uline{are numerically equivalent.}

\sgaStableTarget{sga6.e10.prop.4-5-2}\textbf{Proposition 4.5.2.} \uline{An invertible} \(\underline{O}_X\)\uline{-module} \(\underline{L}\) \uline{is numerically equivalent to zero if and only if}
\sgaStableTarget{sga6.e10.eq.4-5-3}\begin{equation}\tag{4.5.3}
 \chi(\underline{L}^n\otimes_{\underline{O}_X}\underline{O}_Z)=\chi(\underline{O}_Z)
\end{equation}
\uline{for every} \(n\in\mathbb Z\) \uline{and every integral closed curve} \(Z\subset X\).

This follows from 1.1.2 and 4.3.1.

Denote by \(\Pic^N(X)\) the subgroup of \(\Pic(X)\) consisting of the classes of invertible \(\underline{O}_X\)-modules numerically equivalent to zero. From 2.4.4 one obtains

\sgaStableTarget{sga6.e10.cor.4-5-4}\textbf{Corollary 4.5.3.} \(\Pic^\tau(X)\subset \Pic^N(X)\).

In Exp. XIII it will be proved that one even has \(\Pic^N(X)=\Pic^\tau(X)\).

\sgaStableTarget{sga6.e10.sec.5}\section*{5. The isomorphism \(\Pic(X)\cong\Gr^1(X)\)}

\sgaStableTarget{sga6.e10.sec.5.1}\subsection*{\normalfont 5.1.}
Let \(X\) be an arbitrary ringed topos. Then one can define the ring \(K^{\bullet}(X)\) and the filtration \((\Filt^i(X))_{i\in\mathbb Z}\) exactly as for schemes. Likewise, define the group \(\Pic(X)\) as the group of isomorphism classes of invertible \(\underline{O}_X\)-modules. Denote by
%%%%%%%%%%  scan idx 560  |  folio 547  %%%%%%%%%%
\SGAPrintPage{547}%
\sgaStableTarget{sga6.e10.eq.5-1-1}\begin{equation}\tag{5.1.1}
 i:\Pic(X)\longrightarrow K^{\bullet}(X)
\end{equation}
the map associating to the class \(\cl_P(\underline{L})\in\Pic(X)\) of an invertible \(\underline{O}_X\)-module \(\underline{L}\) the element \(\cl^{\bullet}(\underline{L})\) of \(K^{\bullet}(X)\). This map is clearly a homomorphism from \(\Pic(X)\) to the multiplicative group \(K^{\bullet}(X)^*\) of invertible elements of \(K^{\bullet}(X)\). Henceforth, the group law of \(\Pic(X)\) will be written multiplicatively.

The map \(i\) has a left inverse defined by the ``det'' functor: when \(\underline{E}\) is a locally free \(\underline{O}_X\)-module of rank \(\underline{n}\), one has
\[
 \det(\underline{E})=\bigwedge^{\underline{n}} \underline{E}.
\]
If
\sgaStableTarget{sga6.e10.eq.5-1-2}\begin{equation}\tag{5.1.2}
0\longrightarrow \underline{E}'\longrightarrow \underline{E}\longrightarrow \underline{E}''\longrightarrow0
\end{equation}
is an exact sequence of locally free \(\underline{O}_X\)-modules, there is a canonical isomorphism
\[
 \det(\underline{E})\xrightarrow{\sim}(\det(\underline{E}'))\otimes_{\underline{O}_X}(\det(\underline{E}''))
\]
arising from (6.2.9) and the fact that (5.1.2) is locally split. Thus one obtains a homomorphism of abelian groups
\sgaStableTarget{sga6.e10.eq.5-1-3}\begin{equation}\tag{5.1.3}
 \det:K^{\bullet}(X)\longrightarrow\Pic(X)
\end{equation}
that is the required left inverse of \(i\). In particular, \(i\) is injective, and henceforth we shall identify \(\Pic(X)\) with its image in \(K^{\bullet}(X)^*\).

\sgaStableTarget{sga6.e10.sec.5.2}\subsection*{\normalfont 5.2.}
Define another map
\sgaStableTarget{sga6.e10.eq.5-2-1}\begin{equation}\tag{5.2.1}
 \delta:\Pic(X)\longrightarrow\Filt^1(X)
\end{equation}
%%%%%%%%%%  scan idx 561  |  folio 548  %%%%%%%%%%
\SGAPrintPage{548}%
by putting, for every \(L\in\Pic(X)\),
\[
 \delta(L)=1-L^{-1}.
\]

\sgaStableTarget{sga6.e10.rem.5-2-2}\textbf{Remarks 5.2.2.}:
\begin{enumerate}[label=\alph*)]
\item Suppose that \(X\) is the locally ringed topos of Zariski sheaves on a scheme \(S\). When \(D\) is a Cartier divisor \(\geq0\) on \(X\) and
\[
 L=\cl^{\bullet}(\underline{O}_S(D)),
\]
one has
\sgaStableTarget{sga6.e10.eq.5-2-3}\begin{equation}\tag{5.2.3}
 \theta_S(\delta(L))=\cl_\bullet(\underline{O}_D)
\end{equation}
where \(\theta_S:K^{\bullet}(X)\to K_\bullet(S)\) is the canonical homomorphism of 1.2.3. This follows immediately from the exact sequence
\[
0\longrightarrow\underline{O}_S(-D)\longrightarrow\underline{O}_S\longrightarrow\underline{O}_D\longrightarrow0.
\]

\item Under the hypotheses of 5.1 and 5.2, one has, for every \(L\in\Pic(X)\),
\sgaStableTarget{sga6.e10.eq.5-2-4}\begin{equation}\tag{5.2.4}
 \delta(L)\equiv L-1\quad\operatorname{mod.}\Filt^2(X),
\end{equation}
which justifies the notation \(c_1(L)\) introduced below. Indeed, one has
\[
 (L-1)-(1-L^{-1})=(L-1)(1-L^{-1})\in\Filt^2(X).
\]
\end{enumerate}

\sgaStableTarget{sga6.e10.sec.5.3}\subsection*{\normalfont 5.3.}
Composing \(\delta\) with the canonical epimorphism
\[
 \Filt^1(X)\longrightarrow\Gr^1(X)
\]
gives the ``first Chern class'' on \(\Pic(X)\):
\sgaStableTarget{sga6.e10.eq.5-3-1}\begin{equation}\tag{5.3.1}
 c_1:\Pic(X)\longrightarrow\Gr^1(X).
\end{equation}

%%%%%%%%%%  scan idx 562  |  folio 549  %%%%%%%%%%
\SGAPrintPage{549}%

\sgaStableTarget{sga6.e10.thm.5-3-2}\textbf{Theorem 5.3.2.} \uline{The map} \(c_1\) (5.3.1) \uline{is an injective homomorphism, and the homomorphism}
\[
 \det:\Filt^1(X)\longrightarrow\Pic(X)
\]
\uline{induces, on passing to the quotient, a homomorphism that is a left inverse of} \(c_1\):
\[
 \alpha:\Gr^1(X)\longrightarrow\Pic(X).
\]
\uline{If every section} \(\underline{n}\) \uline{of the constant sheaf} \(\mathbb Z_X\) \uline{is bounded (for example, if} \(X\) \uline{is ``quasi-compact'', or if the final object} \(\underline{e}\) \uline{of} \(X\) \uline{is a finite sum of connected objects), then} \(c_1\) \uline{is an isomorphism and} \(\alpha\) \uline{is its inverse.}

\uline{Proof}: \(c_1\) is a homomorphism, since for every \(L,L'\in\Pic(X)\) one has the equation
\[
 (1-(LL')^{-1})-(1-L^{-1})-(1-{L'}^{-1})
 =
 -(1-L^{-1})(1-{L'}^{-1})\in\Filt^2(X).
\]

For the moment, admit that
\[
 \det(\Filt^2(X))=\{1\}\in\Pic(X)
\]
i.e., that \(\alpha\) is well defined. One then has
\[
 \alpha\circ c_1=\id_{\Pic(X)},
\]
since, by 5.2.2 b), one has
\[
 \det(1-L^{-1})=(\det(L-1))=L
\]
for every \(L\in\Pic(X)\). Next one shows that
\[
 c_1\circ\alpha=\id_{\Gr^1(X)}
\]
when every section \(\underline{n}\) of \(\mathbb Z_X\) is bounded. Since \(\Filt^1(X)\) is the abelian group generated by the elements \(\cl^{\bullet}(\underline{E})-\underline{n}\), where \(\underline{E}\) is a
%%%%%%%%%%  scan idx 563  |  folio 550  %%%%%%%%%%
\SGAPrintPage{550}%
locally free \(\underline{O}_X\)-module of rank \(\underline{n}\), it suffices, taking account of 5.2.2 b), to show that
\sgaStableTarget{sga6.e10.eq.5-3-3}\begin{equation}\tag{5.3.3}
 \det(\cl^{\bullet}(\underline{E}))-\cl^{\bullet}(\underline{E})+\underline{n}-1\in\Filt^2(X).
\end{equation}
Since \(\underline{n}\) is bounded, the final object \(\underline{e}\) of \(X\) decomposes into a \uline{finite} disjoint direct sum \(\coprod e_i\) of subobjects in such a way that \(\underline{n}\) is constant on each subtopos \(X/\underline{e}_i\) \((i=1,\ldots,p)\). Since the functor \(Y\rightsquigarrow K^{\bullet}(X/Y)\) \((Y\in\operatorname{Ob}(X))\) takes \uline{finite} direct sums to finite direct products, one may suppose that \(\underline n\) is constant on \(X\), say \(\underline n=n\in\mathbb Z\). Putting \(E=\cl^{\bullet}(\underline E)\), one finds (cf. Exposé V):
\[
 \det(\underline{E})=\lambda^n(E)=\gamma^n((E-n)+1)
 =
 \sum_{0\leq i\leq n}\gamma^i(E-n),
\]
This proves (5.3.3), whose left-hand side is equal to
\[
 \sum_{2\leq i\leq n}\gamma^i(E-n).
\]

It remains to prove the following

\sgaStableTarget{sga6.e10.lem.5-3-4}\textbf{Lemma 5.3.4.} \uline{For every ringed topos} \(X\) \uline{and every} \(x\in\Filt^2(X)\)\uline{, one has}
\[
 \det(x)=1.
\]

\uline{Proof}: Let \(\epsilon:K^{\bullet}(X)\to\Gamma(\mathbb Z)\) be the augmentation homomorphism. Then (6.2.8) and (6.5.3) below give, by linearity, the formula
\sgaStableTarget{sga6.e10.eq.5-3-5}\begin{equation}\tag{5.3.5}
 \det(xy)=(\det(x))^{\epsilon(y)}(\det(y))^{\epsilon(x)}
 \qquad(x,y\in K^{\bullet}(X)).
\end{equation}
%%%%%%%%%%  scan idx 564  |  folio 551  %%%%%%%%%%
\SGAPrintPage{551}%
It therefore suffices to show that, for every \(k>1\), one has
\[
 \det(\gamma^k(E-\underline{n}))=1,
\]
where, as above, \(E=\cl^{\bullet}(\underline E)\), \(\underline{n}=\epsilon(E)\). Now one has
\[
 \gamma^k(x)=\lambda^k(x+k-1)
\]
for every \(x\in K^{\bullet}(X)\), so that
\[
 \det(\gamma^k(E-\underline{n}))
 =
 \det(\lambda^k(E-\underline{n}+k-1))
 =
 \det\left(\sum_{p+q=k}\binom{-\underline{n}+k-1}{p}\lambda^q(E)\right).
\]
But (6.2.10) and (6.5.3) give
\[
 \det(\lambda^q(E))=(\det(E))^{\binom{\underline{n}-1}{\underline{q}-1}}
\]
for every \(q>0\), whence
\[
 \det(\gamma^k(E-\underline{n}))=(\det(E))^{\underline{r}},
\]
with
\[
 \underline{r}=\sum_{\substack{p+q=k\\ q\geq1}}
 \binom{-\underline{n}+k-1}{p}\binom{\underline{n}-1}{q-1}.
\]
Now \(\underline{r}\) is a section of \(\mathbb Z\), everywhere equal to the coefficient of \(t^{k-1}\) in \((1+t)^{k-2}\), and is therefore zero. The result follows.

\sgaStableTarget{sga6.e10.sec.6}\section*{6. Appendix: Calculation of Determinants of Locally Free Sheaves}

\sgaStableTarget{sga6.e10.sec.6.1}\subsection*{\normalfont 6.1.}
First consider an arbitrary ringed topos \(X=({C},\underline{O})\). For every positive section \(\underline{n}\in\Gamma(\mathbb Z_X)\) of the constant sheaf \(\mathbb Z\) and every \(\underline{O}\)-module \(M\in\Mod_{\underline{O}}(X)\), it is clear how to define the \(\underline{O}\)-modules
%%%%%%%%%%  scan idx 565  |  folio 552  %%%%%%%%%%
\SGAPrintPage{552}%
\sgaStableTarget{sga6.e10.eq.6-1-1}\begin{equation}\tag{6.1.1}
 M^{\underline{n}},\qquad \bigotimes^{\underline{n}} M,\qquad \bigwedge^{\underline{n}} M,\quad\text{etc.}
\end{equation}
The functors \((\ )^{\underline{n}}\), \(\bigotimes\limits^{\underline{n}}\) and \(\bigwedge\limits^{\underline{n}}\) take free (resp. locally free) modules to free (resp. locally free) modules.

More generally, consider sections \(\underline{n}_1,\ldots,n_r\) of \(\mathbb Z_X\) and a multifunctor (covariant, for definiteness)
\[
 T:\underline{\Mod}^{n_1}(\underline{O})\times\cdots\times
 \underline{\Mod}^{n_r}(\underline{O})\longrightarrow\underline{\Mod}(\underline{O})
\]
in \(r\) variables, where \(\Mod_{\underline{O}}^{\underline{n}}(X)\) denotes the category fibered over \(X\) of locally free \(\underline{O}\)-modules of rank \(\underline{n}\). Put
\[
 M=T(\underline{O}^{\underline{n}_1},\ldots,\underline{O}^{\underline{n}_r})
\]
and
\[
 G=\operatorname{G\ell}(\underline{n}_1)\times\cdots\times\operatorname{G\ell}(\underline{n}_r),
 \quad\text{where }\operatorname{G\ell}(\underline{n}_i)=\underline{\Aut}_{\underline{O}}(\underline{O}^{n_i}).
\]
It follows immediately that \(T\) defines a representation
\sgaStableTarget{sga6.e10.eq.6-1-2}\begin{equation}\tag{6.1.2}
 \alpha_T:G\longrightarrow\underline{\Aut}_{\underline{O}}(M).
\end{equation}
It is clear that, when \(M\) is itself locally free of rank \(\underline{m}\), for every composite functor
\[
 T'T,\ \text{where } T':\underline{\Mod}^{\underline{m}}(\underline{O})\longrightarrow\underline{\Mod}(\underline{O}),
\]
one has
\sgaStableTarget{sga6.e10.eq.6-1-3}\begin{equation}\tag{6.1.3}
 \alpha_{T'T}=\alpha_{T'}\,\alpha_T.
\end{equation}

Conversely, given the representation \(\alpha=\alpha_T:G\to\underline{\Aut}_{\underline{O}}(M)\), one recovers \(T(\underline{E}_1,\ldots,\underline{E}_r)\) (for every sequence
%%%%%%%%%%  scan idx 566  |  folio 553  %%%%%%%%%%
\SGAPrintPage{553}%
\(\underline{E}_1,\ldots,\underline{E}_r\) of locally free \(\underline{O}\)-modules of ranks \(\underline{n}_1,\ldots,\underline{n}_r\)) by gluing the free \(\underline{O}\)-module \(T(\underline{O}^{\underline{n}_1},\ldots,\underline{O}^{\underline{n}_r})\) on local trivializing charts by means of the representation \(\alpha\). More precisely, the sequence \(\underline{E}_1,\ldots,\underline{E}_r\) corresponds to a right \(G\)-torsor \(P=P_{\underline{E}_1,\ldots,\underline{E}_r}\) such that there is a canonical isomorphism
\sgaStableTarget{sga6.e10.eq.6-1-3.occ2}\begin{equation}\tag{6.1.3}
 \underline{E}_i\xrightarrow{\sim}P\overset{G}{\wedge}\underline{O}^{\underline{n}_i}\qquad(i=1,\ldots,r),
\end{equation}
where \(\overset{G}{\wedge}\) denotes the contracted product, and \(G\) acts on \(\underline{O}^{\underline{n}_i}\) through the \(i\)-th projection \(G\to\Gl(\underline{n}_i)\); then every representation
\[
 \alpha:G\longrightarrow\underline{\Aut}_{\underline{O}}(M)
\]
makes \(G\) act on \(M\) by \(\underline{O}\)-automorphisms, and thus defines a multifunctor \(T=T_\alpha\) such that
\sgaStableTarget{sga6.e10.eq.6-1-4}\begin{equation}\tag{6.1.4}
 T_\alpha(\underline{E}_1,\ldots,\underline{E}_r)=P_{\underline{E}_1,\ldots,\underline{E}_r}\overset{G}{\wedge} M.
\end{equation}
Moreover, for every composite representation \(\alpha=\alpha_1\circ\alpha_2\), one finds
\sgaStableTarget{sga6.e10.eq.6-1-5}\begin{equation}\tag{6.1.5}
 T_\alpha=T_{\alpha_1}\circ T_{\alpha_2}.
\end{equation}

\sgaStableTarget{sga6.e10.sec.6.2}\subsection*{\normalfont 6.2.}
For every locally free \(\underline{O}\)-module \(\underline{E}\) of rank \(\underline{n}\), we shall denote by \(\det(\underline{E})\) the invertible \(\underline{O}\)-module \(\bigwedge\limits^{\underline{n}} \underline{E}\). Keeping the notation of 6.1, the functor \(\det(T(\underline{E}_1,\ldots,\underline{E}_r))\) is determined, up to canonical isomorphism, by the corresponding representation
\sgaStableTarget{sga6.e10.eq.6-2-1}\begin{equation}\tag{6.2.1}
 \det\circ\alpha_T:G\longrightarrow{G}_m=\Gl(1).
\end{equation}

Suppose now that \(X\) is the ringed topos of sheaves on the ringed site \((\Sch/S)\) endowed with the étale topology. Then one can calculate \(\det(T(\underline{E}_1,\ldots,\underline{E}_r))\), or equivalently the representation \(\det\circ\alpha_T\), by means of the following
%%%%%%%%%%  scan idx 567  |  folio 554  %%%%%%%%%%
\SGAPrintPage{554}%

\sgaStableTarget{sga6.e10.prop.6-2-2}\textbf{Proposition 6.2.2.} \uline{For every positive section} \(\underline{n}\) \uline{of} \(\mathbb Z\)\uline{, the étale sheaf of groups} \(\underline{\Der}(\Gl(\underline{n}))\) \uline{of commutators of} \(\Gl(\underline{n})\) \uline{is represented by}
\(\operatorname{Sl}(\underline{n})=\Ker(\det:\Gl(\underline{n})\to{G}_m)\).

\sgaStableTarget{sga6.e10.cor.6-2-3}\textbf{Corollary 6.2.3.} \uline{Every representation}
\sgaStableTarget{sga6.e10.eq.6-2-4}\begin{equation}\tag{6.2.4}
 \alpha:G=\Gl(\underline{n}_1)\times\cdots\times\Gl(\underline{n}_r)\longrightarrow{G}_m
\end{equation}
\uline{is of the form}
\sgaStableTarget{sga6.e10.eq.6-2-5}\begin{equation}\tag{6.2.5}
 \alpha(Y)(s)=(\det(s_1))^{\underline{m}_1}\cdots(\det(s_r))^{\underline{m}_r},
\end{equation}
\uline{where} \(Y\in\operatorname{Ob}(\Sch/S)\), \(s=(s_1,\ldots,s_r)\in G(Y)\) \uline{and} \(\underline{m}_1,\ldots,\underline{m}_r\) \uline{are sections of} \(\mathbb Z\) \uline{uniquely determined by} \(\alpha\).

\sgaStableTarget{sga6.e10.cor.6-2-6}\textbf{Corollary 6.2.6.} \uline{With the notation of} 6.1. \uline{one has}
\sgaStableTarget{sga6.e10.eq.6-2-6}\[
 \det\bigl(T_\alpha(\underline{E}_1,\ldots,\underline{E}_r)\bigr)
 =
 \det(\underline{E}_1)^{\otimes \underline{m}_1}\otimes\cdots\otimes
 \det(\underline{E}_r)^{\otimes \underline{m}_r},
\]
\uline{where} \(\underline{m}_1,\ldots,\underline{m}_r\) \uline{are sections of} \(\mathbb Z\) \uline{uniquely determined by} \(\alpha\).

\sgaStableTarget{sga6.e10.rem.6-2-7}\textbf{Remark 6.2.7.} One can determine the exponents \(\underline{m}_1,\ldots,\underline{m}_r\) in (6.2.5) and (6.2.6) by successively substituting in (6.2.5)
\(s_i=\lambda.\ \operatorname{id}.(\lambda\in{G}_m(Y))\), \(s_j=\operatorname{id}_{\Gl(\underline{n}_j)(Y)}\), for \(j\ne i\) \((i,j=1,\ldots,r)\), whence
\[
 \det(s_i)=\lambda^{\underline{n}_i}\quad\text{and}\quad
 \alpha(Y)(s)=\lambda^{\underline{n}_im_i}.
\]
For example, one thus deduces from (6.2.6)
\sgaStableTarget{sga6.e10.eq.6-2-8}\sgaStableTarget{sga6.e10.eq.6-2-9}\begin{align}
 \det(\underline{E}_1\otimes \underline{E}_2)&=(\det(\underline{E}_1))^{\otimes \underline{n}_2}\otimes(\det(\underline{E}_2))^{\otimes \underline{n}_1}, \tag{6.2.8}\\
 \det(\underline{E}_1\oplus \underline{E}_2)&=(\det(\underline{E}_1))\otimes(\det(\underline{E}_2)) \tag{6.2.9}
\end{align}
and
\sgaStableTarget{sga6.e10.eq.6-2-10}\begin{equation}\tag{6.2.10}
 \det\bigl(\bigwedge\limits^{\underline{m}} \underline{E}_1\bigr)=(\det(\underline{E}_1))^{\binom{\underline{n}_1-1}{\underline{m}-1}}
\end{equation}
etc.

%%%%%%%%%%  scan idx 568  |  folio 555  %%%%%%%%%%
\SGAPrintPage{555}%

\sgaStableTarget{sga6.e10.sec.6.3}\subsection*{\normalfont 6.3.}
Let us first show how 6.2.3 follows from 6.2.2. For this, it plainly suffices to prove (6.2.5) for the
\(s_{(i)}=(1,\ldots,s_i,\ldots,1)\) \((i=1,\ldots,r)\), so one may suppose \(r=1\), \(G=\operatorname{G\ell}(\underline{n}_1)\). By 6.2.2, \(\alpha\) factors through \(\operatorname{G\ell}(\underline{n}_1)/\operatorname{Sl}(\underline{n}_1)\), i.e., there is a commutative diagram
\[
\begin{tikzcd}[column sep=large,row sep=large]
\operatorname{G\ell}(\underline{n}_1) \arrow[rr,"\alpha"] \arrow[dr,"j"'] && {G}_m\\
& \operatorname{G\ell}(\underline{n}_1)/\operatorname{Sl}(\underline{n}_1) \arrow[ur,"\beta"'] &
\end{tikzcd}
\]
where \(j\) is the canonical epimorphism and \(\beta\) is a representation of \(\operatorname{G\ell}(\underline{n}_1)/\operatorname{Sl}(\underline{n}_1)\) in \({G}_m\). But the homomorphism
\[
 \det:\operatorname{G\ell}(\underline{n}_i)\longrightarrow{G}_m
\]
has a section, so that the monomorphism
\[
 \operatorname{G\ell}(\underline{n}_i)/\operatorname{S\ell}(\underline{n}_i)\longrightarrow{G}_m
\]
induced by det is an isomorphism. On the other hand, it is known (SGA 3 VIII 1.3) that every representation of \({G}_m\) in itself is of the form
\[
 s\rightsquigarrow s^{\underline{m}},
\]
where \(\underline{m}\) is a section of \(\mathbb Z\), whence the assertion.

Corollary 6.2.6 is a simple reformulation of 6.2.3. Note that 6.2.3 also follows immediately from SGA 3 XXII 6.2.1, but we wish to give an elementary proof here via 6.2.2.

%%%%%%%%%%  scan idx 569  |  folio 556  %%%%%%%%%%
\SGAPrintPage{556}%

\sgaStableTarget{sga6.e10.sec.6.4}\subsection*{6.4. Proof of 6.2.2.}
It suffices to show that, for every \(Y\in\operatorname{Ob}(\Sch/S)\), \(y\in Y\) and \(s\in\operatorname{S\ell}(\underline{n})(Y)\), one can find an ``étale neighborhood'' \(Y'\to Y\) of \(y\) such that the restriction of \(s\) to \(Y'\) is a product of commutators of \(\operatorname{G\ell}(\underline{n})(Y')\). Since the question is local, one may suppose that \(Y\) is affine with ring \(A\), and that \(\underline{n}\) is constant with value \(n\) on \(Y\). Let \(\underline{p}\) be the prime ideal of \(y\) and \((s_{ij})(i,j=1,\ldots,n)\) the matrix of \(s\). For every \(\lambda\in A\) and \(h, k=1,\ldots,n\); \(h\ne k\), we shall denote by \(c_{h,k,\lambda}\) the matrix \((c_{ij})\) \((i,j=1,\ldots,n)\) such that \(c_{hk}=\lambda\) and \(c_{ij}=\delta_{ij}\) for \((i,j)\ne(h,k)\) \((\delta_{ij}=\) the Kronecker symbol). For every matrix \(\underline{a}=(a_{ij})(i,j=1,\ldots,n)\), one obtains \(c_{h,k,\lambda}\underline{a}\) by adding to the \(k\)-th column of \(\underline{a}\) \(\lambda\) times the \(h\)-th column. Similarly, one obtains the matrix \(\underline{a}c_{h,k,\lambda}\) by adding to the \(h\)-th row of \(\underline{a}\) \(\lambda\) times the \(k\)-th row. In particular, one has
\[
 \det(c_{h,k,\lambda})=1,
\]
and one sees at once that
\sgaStableTarget{sga6.e10.eq.6-4-1}\begin{equation}\tag{6.3.1}
 c_{h,k,\lambda}c_{h,k,\mu}=c_{h,k,\lambda+\mu}
\end{equation}
for every \(\lambda,\mu\in A\), so that one obtains a representation
\[
 c_{h,k}:G_a\longrightarrow\operatorname{S\ell}(n)
 \qquad(h,k=1,\ldots,n\ ;\ h\ne k).
\]
Unless expressly stated otherwise, in what follows we shall suppose that the \(\lambda\) occurring in \(c_{h,k,\lambda}\) \((\lambda\in A\ ;\ h,k=1,\ldots,n\ ;\ h\ne k)\) is invertible.
The proof then proceeds in two steps. We shall first show that \(s\) is a product of matrices \(c_{h,k,\lambda}\) in a sufficiently small Zariski neighborhood of \(y\). We shall then show that the matrices \(c_{h,k,\lambda}\) are products of commutators of \(\operatorname{G\ell}(n)\) in a sufficiently small étale neighborhood of \(y\).

%%%%%%%%%%  scan idx 570  |  folio 557  %%%%%%%%%%
\SGAPrintPage{557}%

For the first step, one must transform \(s\) into the identity matrix \(\underline{\underline{1}}=(\delta_{ij})\) in a Zariski neighborhood of \(y\) by multiplying \(s\) on the left and on the right by suitable matrices \(c_{h,k,\lambda}\). Let us first show that one can thereby make \(s_{11}=1\). If \(s_{11}\ne1\) and \(s_{21}\in\underline{p}\), we seek a matrix \(c=c_{2,h,\lambda}\) such that the product \(s'=sc=(s'_{ij})\) has \(s'_{21}\notin\underline{p}\). Since \(\det(s)=1\notin\underline{p}\), there is at least one index \(i\) such that \(s_{i1}\notin\underline{p}\). Then one has \(s_{21}+s_{i1}\notin\underline{p}\), and \(c_{2,i,1}\) answers the question. One may therefore suppose \(s_{21}\notin\underline{p}\) and, after localizing, suppose that \(s_{21}\) is invertible. If \(s_{11}-1\in\underline{p}\), then in \(s'=(s'_{ij})=sc_{1,2,1}\) one has
\(s'_{11}-1=(s_{11}+s_{21})-1\notin\underline{p}\). After localizing, one may then suppose that \(s_{11}-1\) is invertible. Then \(\lambda=(1-s_{11})/s_{21}\) is invertible and one finds \(s'_{11}=1\) in \(s'=(s'_{ij})=sc_{1,2,\lambda}\).

Supposing \(s_{11}=1\), one can then put \(s\) in the form
\[
 s_{11}=1,\qquad s_{1i}=s_{i1}=0\quad(i=2,\ldots,n)
\]
in a Zariski neighborhood, of \(y\) by multiplying \(s\) on the left (resp. on the right) by matrices \(c_{1,i,\lambda}\) (resp. \(c_{i,1,\lambda}\)). Next one makes \(s_{22}=1\), etc. By induction, one can then transform \(s\) in a Zariski neighborhood \(U\) of \(y\) into a matrix \(s'=(s'_{ij})\) with \(s'_{ij}=\delta_{ij}\) for \(i<n\) or \(j<n\). But by the very construction one then has \(s'=c_1sc_2\) on \(\mathcal U\), where \(c_1\) and \(c_2\) are products of matrices \(c_{h,k,\lambda}\). Thus one has
\[
 \det(s')=\det(s)=1,
\]
so that \(s'_{nn}=1\), \(s'=\underline{\underline{1}}\), and the result follows.

We therefore suppose that \(s\) is already, on \(Y\), a product of matrices \(c_{h,k,\lambda}\).
It remains to see that these are products of commutators in an étale neighborhood of \(y\). First, one sees that any two matrices
%%%%%%%%%%  scan idx 571  |  folio 558  %%%%%%%%%%
\SGAPrintPage{558}%
\(c_{h,k,\lambda}\) and \(c_{h,k,\mu}\) are conjugate; in other words, there is an invertible matrix \(a=(a_{ij})\) \((i,j=1,\ldots,n)\) such that
\[
 ac_{h,k,\lambda}=c_{h,k,\mu}a.
\]
Indeed, the matrix \((a_{ij})\) with \(a_{hh}=1/\mu\), \(a_{kk}=1/\lambda\), and otherwise \(a_{ij}=\delta_{ij}\), answers the question (recall that \(\lambda\) and \(\mu\) are invertible by hypothesis and \(h\ne k\)). Thus, for any homomorphism \(\beta\) from \(\operatorname{G\ell}(n)(U)\) to a sheaf of abelian groups \(H\), one has
\[
 \beta(c_{h,k,\lambda})=\beta(c_{h,k,\mu}).
\]
If \(\lambda+\mu\) is invertible, this and (6.3.1) imply that \(\beta(c_{h,k,\lambda})=0\) for every \(h,\beta\), so that \(c_{h,k,\lambda}\) is a product of commutators. If the residue field \(k\) of \(y\) is not the prime field \(\underline{F}_2=\{0,1\}\), then for every invertible \(\lambda\in A\) one can find a Zariski neighborhood \(U=\Spec(B)\) of \(y\) and an invertible element \(\mu\in B\) such that \(\lambda+\mu\in B\) is invertible. It follows that \(c_{h,k,\lambda}\) is a product of commutators on \(U\).

If \(k=\underline{F}_2\), then \(V=\Spec(A[T]/(T^2+T+1))\) is an étale neighborhood of \(y\) whose residue field differs from \(\underline{F}_2\). The result follows.

\sgaStableTarget{sga6.e10.sec.6.5}\subsection*{6.5. Remarks}
\sgaStableTarget{sga6.e10.para.6-5-1}6.5.1. When none of the residue fields of \(S\) is isomorphic to \(\underline{F}_2\), the preceding proof shows that 6.2.2. remains valid for the ringed topos of \uline{Zariski} sheaves on \(S\). On the other hand, already for \(S=\Spec(\underline{F}_2)\), \(\operatorname{S\ell}(2,\underline{F}_2)=\operatorname{G\ell}(2,\underline{F}_2)\) differs from the commutator group of \(\operatorname{G\ell}(2,\underline{F}_2)\). The order of the latter is \(3\), whereas the order of \(\operatorname{S\ell}(2,\underline{F}_2)\) is \(6\) (cf. [1]).

\sgaStableTarget{sga6.e10.para.6-5-2}6.5.2. On an arbitrary locally ringed topos \(X\), it is not true in general that every representation \(G_m\to G_m\) is defined by a
%%%%%%%%%%  scan idx 572  |  folio 559  %%%%%%%%%%
\SGAPrintPage{559}%
power. For example, on the topos \((\Ens)\) ringed by the field \(\mathbb Q\), the representation \(G_m\to G_m\) defined by
\[
 r\rightsquigarrow|r|\qquad(r\in\mathbb Q^*)
\]
is not a power. In particular, 6.2.3 does not remain valid for an arbitrary locally ringed topos \(X\) and an arbitrary representation \(\alpha:G\to G_m\).

\sgaStableTarget{sga6.e10.para.6-5-3}6.5.3. On the other hand, 6.2.3 remains valid for the representations associated with formulas (6.2.8), (6.2.9) and (6.2.10) on an arbitrary ringed topos \(X=({C},\underline{O})\), because these representations are ``algebraic'' in the following sense:

Let \(\underline{e}\) be the final object of \(X\), put \(S=\Spec(\underline{O}(\underline{e}))\), and consider the functor
\[
 u:{C}\longrightarrow(\Sch/S)
\]
which assigns to each \(Y\in\operatorname{Ob}({C})\) the affine scheme \(\Spec(\underline{O}(Y))\) over \(S\). Then, for every positive integer \(n\), the presheaf \(\operatorname{G\ell}(n)_X\) on \({C}\) is the composite with \(u\) of the corresponding presheaf \(\operatorname{G\ell}(n)_S\) on \((\Sch/S)\). A representation on \(X\)
\sgaStableTarget{sga6.e10.eq.6-5-4}\begin{equation}\tag{6.5.4}
 \alpha:\operatorname{G\ell}(n)\longrightarrow G_m
\end{equation}
is called \uline{algebraic} if it arises via \(u\) from a representation
\sgaStableTarget{sga6.e10.eq.6-5-5}\begin{equation}\tag{6.5.5}
 \alpha_S:\operatorname{G\ell}(n)_S\longrightarrow G_{m,S}
\end{equation}
of presheaves of groups on \((\Sch/S)\), i.e., if one has
\[
 \alpha(s)=\alpha_S(s)
\]
for every \(Y\in\operatorname{Ob}({C})\), \(s\in\operatorname{G\ell}(n)(Y)=\operatorname{G\ell}(n)_S(u(Y))\).

The representation \(\alpha\) in (6.5.4) is, of course, completely determined by \(\alpha_S\). It follows immediately that (6.2.3) remains valid in an arbitrary ringed topos when \(\alpha\) is algebraic. In particular, formulas (6.2.8), (6.2.9) and (6.2.10) therefore remain valid on an arbitrary ringed topos.

%%%%%%%%%%  scan idx 573  |  folio 560  %%%%%%%%%%
\SGAPrintPage{560}%

\sgaStableTarget{sga6.e10.sec.7}\section*{7. Appendix: Specialization in Intersection Theory}
\begin{center}
\underline{by A. Grothendieck}
\end{center}

\sgaStableTarget{sga6.e10.lem.7-1}\textbf{Lemma 7.1.} \uline{Let} \(i:Y\to X\) \uline{be a regular closed immersion} (VII 1.4)\uline{, of bounded codimension, so that} \(i\) \uline{is a perfect morphism} (VII 1.9)\uline{, and one can therefore define the morphism} (IV 2.12)
\[
 i^*:K_\bullet(X)\longrightarrow K_\bullet(Y).
\]
\uline{Let, on the other hand,} \(\mathbb N\) \uline{be the conormal sheaf of} \(i\)\uline{, which is a locally free sheaf of finite type on} \(Y\) (VII 1.6)\uline{, so that one can form}
\[
 \lambda_{-1}(\mathbb N)\in K^{\bullet}(Y).
\]
\uline{Consider also the direct-image homomorphism} (IV 2.11)
\[
 i_*:K_\bullet(Y)\longrightarrow K_\bullet(X).
\]
\uline{Then one has the formula}
\[
 i^*i_*(x)=x\lambda_{-1}(\mathbb N)\qquad\text{for } x\in K_\bullet(Y),
\]
\uline{where the product on the right is taken with respect to the} \(K^{\bullet}(Y)\)\uline{-module structure on} \(K_\bullet(Y)\) \uline{defined in} (IV 2.10).

The proof is, essentially without change, the same as that of the same formula, written for \(K^{\bullet}(X)\), \(K^{\bullet}(Y)\), established in (VII 2.7).

\sgaStableTarget{sga6.e10.cor.7-2}\textbf{Corollary 7.2.} \uline{Let} \(f:X\to S\) \uline{be a flat morphism, with} \(S\) \uline{a trait, i.e., the spectrum of a discrete valuation ring} \(A\), \(s\) \uline{the closed point of} \(S\), \(t\) \uline{its generic point, and} \(i:X_s\to X\) \uline{the canonical inclusion. Then} \(i\) \uline{has finite Tor-dimension, so} (IV 2.12) \uline{one can define}
\[
 i^*:K_\bullet(X)\longrightarrow K_\bullet(X_s),
\]
%%%%%%%%%%  scan idx 574  |  folio 561  %%%%%%%%%%
\SGAPrintPage{561}%
\uline{on the other hand}, (IV 2.11) \uline{one has a homomorphism}
\[
 i_*:K_\bullet(X_s)\longrightarrow K_\bullet(X).
\]
\uline{Then}, \uline{the composite}
\[
 i^*i_*:K_\bullet(X_s)\longrightarrow K_\bullet(X)\longrightarrow K_\bullet(X_s)
\]
\uline{is zero}.

Indeed, \(i:X_s\to X\) is a regular immersion of codimension \(1\) whose conormal sheaf \(\mathbb N\) is \uline{free} of rank \(1\), since \(X_s\) can be described by a global equation \(f^*(g)=0\) (where \(g\) is a uniformizer of \(A\)). Thus one has \(\lambda_{-1}(\mathbb N)=0\), and the conclusion follows from 7.1.

Throughout what follows, we shall suppose that \(X\) is \uline{noetherian} and \uline{flat over} \(S\). Let us now apply the exact homotopy sequence (IX 1.1) to the following situation: \(X\), with the open subscheme \(X_t\) and complementary closed subscheme \(X_s\):
\[
K_\bullet(X_s)\xrightarrow{i_*}K_\bullet(X)\xrightarrow{j^*}K_\bullet(X_t)\longrightarrow0,
\]
where \(j:X_t\to X\) is the inclusion, which is an open immersion. Using 7.2, one obtains:

\sgaStableTarget{sga6.e10.prop.7-3}\textbf{Proposition 7.3.} \uline{With the notation of} 7.2, \(X\) \uline{being assumed noetherian}, \uline{there is a unique group homomorphism}
\sgaStableTarget{sga6.e10.eq.7-3-1}\begin{equation}\tag{7.3.1}
 \sigma_\bullet:K_\bullet(X_t)\longrightarrow K_\bullet(X_s)
\end{equation}
(\uline{called the} ``specialization homomorphism'') \uline{making the diagram commutative}
\sgaStableTarget{sga6.e10.eq.7-3-2}\begin{equation}\tag{7.3.2}
\begin{tikzcd}[column sep=large,row sep=large]
K_\bullet(X) \arrow[r,"j^*"] \arrow[d,"i^*"'] & K_\bullet(X_t) \arrow[dl,"\sigma_\bullet"]\\
K_\bullet(X_s) & {.}
\end{tikzcd}
\end{equation}


% Local macro additions for Exposé X continuation
\providecommand{\Kalg}{K_{\mathrm{alg}}}
\providecommand{\Gralg}{\operatorname{Gr}_{\mathrm{alg}}}
\providecommand{\Grnum}{\operatorname{Gr}_{\mathrm{num}}}
\providecommand{\Knum}{K_{\mathrm{num}}}
\providecommand{\Cl}{\operatorname{Cl}}
\providecommand{\Ch}{\operatorname{Ch}}
\providecommand{\Nm}{\operatorname{N}}
\providecommand{\colim}{\varinjlim}
\providecommand{\Hens}{\operatorname{Hens}}
\providecommand{\Spec}{\operatorname{Spec}}
\providecommand{\coker}{\operatorname{coker}}


%%%%%%%%%%  scan idx 575  |  folio 562  %%%%%%%%%%
\SGAPrintPage{562}%

\noindent \uline{When the two sides of} (7.3.1) \uline{are regarded as modules over} \(K^{\bullet}(X)\) \uline{by means of the ring homomorphisms} \(K^{\bullet}(X)\xrightarrow{j^*}K^{\bullet}(X_t)\) \uline{and} \(K^{\bullet}(X)\xrightarrow{i^*}K^{\bullet}(X_s)\) (IV 2.7 b) \uline{and the module structures on} \(K_\bullet(X_t)\) \uline{and} \(K_\bullet(X_s)\) \uline{over} \(K^{\bullet}(X_t)\) \uline{and} \(K^{\bullet}(X_s)\) \uline{respectively} (IV 2.10), \uline{the homomorphism} \(\sigma_\bullet\) \uline{is a homomorphism of} \(K^{\bullet}(X)\)-\uline{modules}.

The last assertion follows from the fact that \(i^*:K_\bullet(X)\to K_\bullet(X_s)\) is a homomorphism of \(K^{\bullet}(X)\)-modules, as is the epimorphism \(j^*:K_\bullet(X)\longrightarrow K_\bullet(X_t)\), by (IV 2.12.2).

We shall review other properties of the specialization morphism.

\sgaStableTarget{sga6.e10.sec.7.4}\subsection*{7.4. Compatibility with filtrations.}
Endow the \(K_\bullet\) with the filtration \(\operatorname{Fil}_\bullet K_\bullet\) by the dimension of supports (1.1.1). Then, when \(f\) is of finite type, more generally when \(X\) is catenary, \uline{the specialization homomorphism is compatible with these filtrations}, and consequently induces a homomorphism (also called the specialization homomorphism):
\sgaStableTarget{sga6.e10.eq.7-4-1}\begin{equation}\tag{7.4.1}
  \sigma_\bullet:\Gr_\bullet(X_t)\longrightarrow \Gr_\bullet(X_s).
\end{equation}

To see this, by 1.1.4 one is reduced to proving that, when \(Y_t\) is an integral subscheme of dimension \(d\) of \(X_t\), \(\sigma_\bullet(\cl_\bullet(Y_t))\) has filtration \(\leq d\). Now let \(Y\) be the schematic closure of \(Y_t\) in \(X\), so that \(Y\) is flat over \(S\) (EGA IV 2.8.1). It follows readily that \(i^*(\cl_\bullet(Y))=\cl_\bullet(Y_s)\), where the two \(\cl_\bullet\) are taken in \(K_\bullet(X)\) and \(K_\bullet(X_s)\) respectively, whence \(\sigma_\bullet(\cl_\bullet(Y_t))=\cl_\bullet(Y_s)\). One checks readily (using the fact that \(X\) is catenary, hence \(Y\) is catenary) that \(\dim Y_s\leq \dim Y_t\); when \(f\) is of finite type, this follows at once, by reduction to the case of integral \(X\), from EGA IV 7.1.13. This gives the asserted conclusion.

%%%%%%%%%%  scan idx 576  |  folio 563  %%%%%%%%%%
\SGAPrintPage{563}%

\sgaStableTarget{sga6.e10.sec.7.5}\subsection*{7.5. ``Conservation of number''.}
Suppose now that \(f\) is proper, so that \(X_t\) and \(X_s\) are proper over \(k(t)\) and \(k(s)\) respectively, and one has the homomorphisms of 2.3:
\[
  \chi_t:K_\bullet(X_t)\longrightarrow \mathbb Z,\qquad
  \chi_s:K_\bullet(X_s)\longrightarrow \mathbb Z .
\]
Then one has the relation
\sgaStableTarget{sga6.e10.eq.7-5-1}\begin{equation}\tag{7.5.1}
  \chi_s(\sigma_\bullet(x))=\chi_t(x)\qquad\text{for } x\in K_\bullet(X_t)
\end{equation}
(the ``conservation-of-number theorem''); in other words, the diagram
\sgaStableTarget{sga6.e10.eq.7-5-2}\begin{equation}\tag{7.5.2}
\begin{tikzcd}[column sep=large,row sep=large]
K_\bullet(X_t) \arrow[r,"\sigma_\bullet"] \arrow[dr,"\chi_t"'] &
K_\bullet(X_s) \arrow[d,"\chi_s"]\\
& \mathbb Z .
\end{tikzcd}
\end{equation}
commutes. To prove (7.5.1), one is again reduced to the case where \(x=\cl_\bullet(Y_t)\), with \(Y_t\) an integral closed subscheme of \(X_t\), and, with the notation of the proof of 7.4, one is reduced to proving that
\[
  \chi(Y_s)=\chi(Y_t),
\]
which is well known (EGA III 7.9.4, or Exposé III 5.7).

\sgaStableTarget{sga6.e10.sec.7.6}\subsection*{\normalfont 7.6.}
It is immediate that
\sgaStableTarget{sga6.e10.eq.7-6-1}\begin{equation}\tag{7.6.1}
  \sigma_\bullet\bigl(\cl_\bullet(\underline{O}_{X_t})\bigr)=\cl_\bullet(\underline{O}_{X_s}),
\end{equation}
which, taking account of the \(K^{\bullet}(X)\)-linearity of \(\sigma_\bullet\), implies that the following diagram is commutative:
\sgaStableTarget{sga6.e10.eq.7-6-2}\begin{equation}\tag{7.6.2}
\begin{tikzcd}[column sep=large,row sep=large]
& K^{\bullet}(X) \arrow[dl,"j^*"'] \arrow[dr,"i^*"] &\\
K^{\bullet}(X_t) \arrow[d] && K^{\bullet}(X_s) \arrow[d]\\
K_\bullet(X_t) \arrow[rr,"\sigma_\bullet"] && K_\bullet(X_s),
\end{tikzcd}
\end{equation}
%%%%%%%%%%  scan idx 577  |  folio 564  %%%%%%%%%%
\SGAPrintPage{564}%
where the vertical arrows are the canonical homomorphisms \(K^{\bullet}\to K_\bullet\) (IV 2.5.1).

\sgaStableTarget{sga6.e10.sec.7.7}\subsection*{7.7. Specialization for \(K^\bullet\) and \(\Gr^\bullet\).}
When \(X_t\) and \(X_s\) are regular (for example, when \(f\) is smooth), the two preceding vertical arrows are isomorphisms (IV 2.5), and the specialization homomorphism can then be interpreted as a homomorphism
\sgaStableTarget{sga6.e10.eq.7-7-1}\begin{equation}\tag{7.7.1}
  \sigma^\bullet:K^\bullet(X_t)\longrightarrow K^\bullet(X_s),
\end{equation}
which can also be described as the unique homomorphism making the following diagram commutative
\sgaStableTarget{sga6.e10.eq.7-7-2}\begin{equation}\tag{7.7.2}
\begin{tikzcd}[column sep=large,row sep=large]
K^\bullet(X) \arrow[r,"j^*"] \arrow[d,"i^*"'] & K^\bullet(X_t) \arrow[dl,"\sigma^\bullet"]\\
K^\bullet(X_s) & {.}
\end{tikzcd}
\end{equation}
Since, in this last diagram, \(i^*\) and \(j^*\) are ring homomorphisms (IV 2.7 b), with \(j^*\) surjective, it follows that \uline{the specialization homomorphism} (7.7.1) \uline{is a ring homomorphism}. For the same reason, \uline{if} \(X\) \uline{is separated}, \uline{the homomorphism} (7.7.1) \uline{is even a homomorphism of} \(\lambda\)-\uline{rings}. To make this precise, one need only note that a regular separated noetherian scheme is divisorial (II 2.2.7.1), so its \(K^\bullet\) is identified with its ``naive'' \(K^\bullet\) (IV 2.9 (ii)), which is naturally endowed with a \(\lambda\)-ring structure (VI 3.2). On the other hand, the regularity of \(X_t\) and \(X_s\) implies that of \(X\) (EGA \(0_{\mathrm{IV}}\) 17.1.8), and the homomorphisms \(i^*:K^\bullet(X)\to K^\bullet(X_s)\) and \(j^*:K^\bullet(X)\to K^\bullet(X_t)\) are \(\lambda\)-homomorphisms, so the same is true of (7.7.1), which is obtained from \(i^*\) by passage to the quotient.

%%%%%%%%%%  scan idx 578  |  folio 565  %%%%%%%%%%
\SGAPrintPage{565}%

Note also that, since \(X\) is regular and hence normal, its connected components coincide with its irreducible components, which dominate \(S\), so that the homomorphism \(H^0(X,\mathbb Z)\to H^0(X_t,Z)\) induced by the inclusion \(X_t\to X\) is bijective. Consequently the homomorphism \(H^0(X,\mathbb Z)\to H^0(X_s,\mathbb Z)\) induced by the inclusion \(X_s\to X\) can also be interpreted as a homomorphism
\sgaStableTarget{sga6.e10.eq.7-7-3}\begin{equation}\tag{7.7.3}
  H^0(X_t,\mathbb Z)\longrightarrow H^0(X_s,\mathbb Z),
\end{equation}
which again deserves the name ``specialization homomorphism''. One then checks immediately from the definitions that the homomorphism (7.7.1) is compatible with the \uline{augmented algebra} structures over \(H^0(X_t,\mathbb Z)\) and \(H^0(X_s,\mathbb Z)\) respectively on the two sides, defined in (IV 2.8), taking account of the preceding homomorphism (7.7.3).

Still supposing that \(X\) is separated, the preceding discussion implies in particular that (7.7.1) induces a homomorphism of graded rings
\sgaStableTarget{sga6.e10.eq.7-7-4}\begin{equation}\tag{7.7.4}
  \sigma^\bullet:\Gr^\bullet(X_t)\longrightarrow\Gr^\bullet(X_s),
\end{equation}
and, using the notion of Chern class of (V 6), one obtains a commutative diagram
\[
\begin{tikzcd}[column sep=large,row sep=large]
K^\bullet(X_t) \arrow[r,"c_{X_t}"] \arrow[d] &
\Ch\,\Gr^\bullet(X_t) \arrow[d]\\
K^\bullet(X_s) \arrow[r,"c_{X_s}"] &
\Ch\,\Gr^\bullet(X_s).
\end{tikzcd}
\]

\sgaStableTarget{sga6.e10.sec.7.8}\subsection*{7.8. Specialization for \(\Pic\).}
Retain the hypotheses of 7.7: \(X\) separated, \(X_t\) and \(X_s\) regular. Taking account of the canonical isomorphism (5.3.2):
\[
  \Gr^1\simeq \Pic ,
\]
the homomorphism (7.7.4) therefore gives a canonical homomorphism
\sgaStableTarget{sga6.e10.eq.7-8-1}\begin{equation}\tag{7.8.1}
  \Pic(X_t)\longrightarrow \Pic(X_s).
\end{equation}
%%%%%%%%%%  scan idx 579  |  folio 566  %%%%%%%%%%
\SGAPrintPage{566}%
An equivalent description of this last homomorphism is as follows: the hypothesis that \(X\) is regular implies that these local rings are factorial (by the Auslander--Buchsbaum theorem, EGA IV 21.11.1), and hence that every invertible sheaf \(\underline{L}_t\) on \(X_t\) extends to an invertible sheaf \(\underline{L}\) on \(X\) (EGA IV 21.6.11). One then checks at once that the image under (7.8.1) of \(\cl(\underline{L}_t)\in\Pic(X_t)\) is \(\cl(\underline{L}_s)\in\Pic(X_s)\), where, of course, one puts \(\underline{L}_s=\underline{L}|X_s\).

In fact, using only the hypothesis that the local rings of \(X\) are factorial and that \(X_s\) is locally integral (without hypotheses of regularity or separatedness), one shows that \(\cl(\underline{L}_s)\) does not depend on the chosen extension \(\underline{L}\) of \(\underline{L}_t\) to an invertible Module on \(X\), so that one obtains a homomorphism (7.8.1) directly in this way. To verify the asserted independence, one is reduced to proving that, if \(\underline{L}_t\simeq\underline{O}_{X_t}\), then \(\underline{L}_s\simeq\underline{O}_{X_s}\). Now the hypothesis means that \(\underline{L}\simeq\underline{O}_X(D)\), with \(D\) a divisor on \(X\) having support \(\subset X_s\) (cf. EGA IV 21.2.11), and one reduces at once to the case where \(X_s\) is \uline{integral}, hence \(D=nX_s\); but then \(D=\operatorname{div}(g^n)\) (\(g\) a uniformizer of \(A\)), so \(\underline{O}_X(D)\simeq\underline{O}_X\) and a fortiori \(\underline{L}_s\simeq\underline{O}_{X_s}\).

\sgaStableTarget{sga6.e10.sec.7.9}\subsection*{7.9. Specialization and numerical equivalence.}
Suppose that \(f\) is proper, and \(X_s\) and \(X_t\) are regular, and consider the two specialization maps (7.4.1) and (7.7.4):
\sgaStableTarget{sga6.e10.eq.7-9-1}\begin{equation}\tag{7.9.1}
\begin{gathered}
  \sigma_\bullet:\Gr_\bullet(X_t)\longrightarrow\Gr_\bullet(X_s)\\
  \sigma^\bullet:\Gr^\bullet(X_t)\longrightarrow\Gr^\bullet(X_s).
\end{gathered}
\end{equation}
Since the homomorphism
\sgaStableTarget{sga6.e10.eq.7-9-2}\begin{equation}\tag{7.9.2}
  \sigma:K^\bullet(X_t)\longrightarrow K^\bullet(X_s)
\end{equation}
%%%%%%%%%%  scan idx 580  |  folio 567  %%%%%%%%%%
\SGAPrintPage{567}%
from which they are obtained is compatible with the ring structures, using (7.5.1) one finds that
\sgaStableTarget{sga6.e10.eq.7-9-3}\[
  \langle\sigma_\bullet(x_t),\sigma^\bullet(y_t)\rangle
  =
  \langle x_t,y_t\rangle
  \qquad\text{for }
  x_t\in\Gr_\bullet(X_t)\quad y_t\in\Gr^\bullet(X_t),
\tag{7.9.3}
\]
with the pairings \(\langle x,y\rangle\) of 4.1 between \(\Gr_\bullet\) and \(\Gr^\bullet\). It follows formally that \uline{if} \(x\) \uline{is an element of} \(\Gr_\bullet(X_t)\) (\uline{resp}. \uline{of} \(\Gr^\bullet(X_t)\)) \uline{such that} \(\sigma_\bullet(x)\) (\uline{resp}. \(\sigma^\bullet(x)\)) \uline{is numerically equivalent to zero} (4.5.1), \uline{then the same is true of} \(x\). The converse is true when (7.9.2) (or, equivalently, either of the homomorphisms (7.9.1) obtained from it by passage to the associated graded groups) is an epimorphism mod.torsion, - a circumstance that occurs only rather exceptionally.

\sgaStableTarget{sga6.e10.para.7-9-4}7.9.4. It is plausible, at least when \(f\) is proper and \uline{smooth}, that the converse in question always holds, i.e., that \(\sigma_\bullet(x)\) (resp. \(\sigma^\bullet(x)\)) is numerically equivalent to zero as soon as the same is true of \(x\). This would imply that the homomorphisms (7.9.1) induce \uline{monomorphisms} on the corresponding quotient groups by numerical equivalence, \(\Gr^{\mathrm{num}}_\bullet\) and \(\Gr_{\mathrm{num}}^\bullet\). At least the preceding discussion implies that the group \(\Gr_\bullet^{\mathrm{num}}(X_t)\) is isomorphic to a quotient of a subgroup of \(\Gr_\bullet^{\mathrm{num}}(X_s)\) (and in particular has rank at most that of \(\Gr_\bullet^{\mathrm{num}}(X_s)\)), and likewise for \(\Gr_{\mathrm{num}}^\bullet\), and also for \(K^\bullet_{\mathrm{num}}\) itself.

\sgaStableTarget{sga6.e10.sec.7.10}\subsection*{7.10. Compatibility with direct and inverse images}
Let
\[
  g:X\longrightarrow Y
\]
be a morphism of noetherian schemes flat over \(S\). Denote by \(\sigma_X\) (resp. \(\sigma_Y\)) the specialization homomorphism for \(X\) (resp. \(Y\)), and by \(i_X,j_X\) (resp. \(i_Y,j_Y\))
%%%%%%%%%%  scan idx 581  |  folio 568  %%%%%%%%%%
\SGAPrintPage{568}%
the morphisms \(i,j\) for \(X\) (resp. \(Y\)) from the special and generic fibers to the ambient scheme. Denote by
\[
  g_t:X_t\to Y_t,\qquad g_s:X_s\to Y_s
\]
the induced morphisms, which therefore give rise to the Cartesian squares
\sgaStableTarget{sga6.e10.eq.7-10-0}\begin{equation}\tag{7.10.0}
\begin{tikzcd}[column sep=large,row sep=large]
X_t \arrow[r,"g_t"] \arrow[d,"j_X"'] & Y_t \arrow[d,"j_Y"'] &
X_s \arrow[r,"g_s"] \arrow[d,"i_X"'] & Y_s \arrow[d,"i_Y"']\\
X \arrow[r,"g"] & Y &
X \arrow[r,"g"] & Y .
\end{tikzcd}
\end{equation}

First suppose that \(g\) is \uline{proper}. I claim that the following diagram then commutes:
\sgaStableTarget{sga6.e10.eq.7-10-1}\begin{equation}\tag{7.10.1}
\begin{tikzcd}[column sep=large,row sep=large]
K_\bullet(X_t) \arrow[r,"g_{t*}"] \arrow[d,"\sigma_X"'] &
K_\bullet(Y_t) \arrow[d,"\sigma_Y"]\\
K_\bullet(X_s) \arrow[r,"g_{s*}"] &
K_\bullet(Y_s)
\end{tikzcd}
\end{equation}
This result generalizes 7.5 (since, for \(X=S\), the arrows of (7.3.2) are identified with the identity map of \(\mathbb Z\)). Note that, when \(X\) and \(Y\) satisfy the conditions of 7.7, i.e., have regular fibers, \(g_t\) and \(g_s\) automatically have finite Tor-dimension, and the diagram can be interpreted as a commutative diagram with \(K^\bullet\), which we leave the reader to rewrite.

To prove the commutativity of (7.10.1), one sees at once, using the definitions of \(\sigma_X\) and \(\sigma_Y\), that one is reduced to proving the commutativity of the two front squares of the prism
\sgaStableTarget{sga6.e10.eq.unnamed}\[
\begin{tikzcd}[column sep=large,row sep=large]
K_\bullet(X_t) \arrow[rr,"g_{t*}"] \arrow[dd] && K_\bullet(Y_t) \arrow[dd]\\
& K_\bullet(X) \arrow[ul,"j_X^*" description] \arrow[dl,"i_X^*" description] \arrow[rr,"g_*"] &&
K_\bullet(Y) \arrow[ul,"j_Y^*" description] \arrow[dl,"i_Y^*" description]\\
K_\bullet(X_s) \arrow[rr,"g_{s*}"] && K_\bullet(Y_s).
\end{tikzcd}
\tag{*}
\]
%%%%%%%%%%  scan idx 582  |  folio 569  %%%%%%%%%%
\SGAPrintPage{569}%

Now these two squares are the ``base-change'' squares obtained from the two squares (7.10.0), so that one can apply (IV 3). One need only verify that these squares correspond to pairs \((X,Y_t)\) and \((X,Y_s)\) of schemes over \(Y\) that are Tor-independent over \(Y\). This is trivial for the first pair, since \(Y_t\to Y\) is flat, and also immediate for the second, since the immersion \(Y_s\to Y\) is regular of codimension 1 and remains so after the base change \(X\to Y\).

Now no longer suppose that \(g\) is proper, but suppose instead that it has finite Tor-dimension. I claim that the same is then true of \(g_t\) and \(g_s\). This is trivial for the first, and for the second follows from the fact (just noted) that \(X\) and \(Y_s\) are Tor-independent over \(Y\). One can therefore consider the square of homomorphisms
\sgaStableTarget{sga6.e10.eq.7-10-2}\begin{equation}\tag{7.10.2}
\begin{tikzcd}[column sep=large,row sep=large]
K_\bullet(Y_t) \arrow[r,"g_t^*"] \arrow[d,"\sigma_Y"'] &
K_\bullet(X_t) \arrow[d,"\sigma_X"]\\
K_\bullet(Y_s) \arrow[r,"g_s^*"] &
K_\bullet(X_s),
\end{tikzcd}
\end{equation}
where the horizontal arrows are defined in (IV 2.12). Note that, when \(X\) and \(Y\) satisfy the conditions of 7.7, i.e., have regular fibers, \(g\) automatically has finite Tor-dimension (since \(Y\) is regular), and the preceding diagram can be interpreted as a commutative diagram of rings \(K^\bullet\).

To prove the commutativity of (7.10.2), it again suffices to prove that of the front squares of the prism analogous to \((*)\), but with the direct-image homomorphisms \(g_{t*}\), \(g_*\), \(g_{s*}\) replaced by the inverse-image homomorphisms \(g_t^*\), \(g^*\), \(g_s^*\). Now the commutativity of these squares follows at 
%%%%%%%%%%  scan idx 583  |  folio 570  %%%%%%%%%%
\SGAPrintPage{570}%
once from that of the corresponding squares (7.10.0), taking account of the transitivity of inverse images (IV 2.12).

\sgaStableTarget{sga6.e10.rem.7-10-3}\textbf{Remark 7.10.3.} When \(g\) is proper, the commutativity of (7.10.1) immediately implies that of the following diagram, obtained from it by passage to the associated graded groups:
\sgaStableTarget{sga6.e10.eq.7-10-1bis}\begin{equation}\tag{7.10.1 bis}
\begin{tikzcd}[column sep=large,row sep=large]
\Gr_\bullet(X_t) \arrow[r,"g_{t*}"] \arrow[d,"\sigma_{X\bullet}"'] &
\Gr_\bullet(Y_t) \arrow[d,"\sigma_{Y\bullet}"]\\
\Gr_\bullet(X_s) \arrow[r,"g_{s*}"] &
\Gr_\bullet(Y_s).
\end{tikzcd}
\end{equation}
Likewise, when \(X\) and \(Y\) have regular fibers and are separated, the commutativity of (7.10.2) implies that of
\sgaStableTarget{sga6.e10.eq.7-10-2bis}\begin{equation}\tag{7.10.2 bis}
\begin{tikzcd}[column sep=large,row sep=large]
\Gr^\bullet(Y_t) \arrow[r,"g_t^*"] \arrow[d,"\sigma_Y^\bullet"'] &
\Gr^\bullet(X_t) \arrow[d,"\sigma_X^\bullet"]\\
\Gr^\bullet(Y_s) \arrow[r,"g_s^*"] &
\Gr^\bullet(X_s).
\end{tikzcd}
\end{equation}
When \(X\) and \(Y\) are moreover quasi-projective and smooth over \(S\), one also knows that \(g_t^*\) and \(g_s^*\) are compatible with the filtrations by codimension of supports, and one obtains a commutative diagram like (7.10.2 bis), but with the \(\Gr^\bullet\) replaced by the \(\Gr^\bullet_{\mathrm{top}}\) (cf. Exposé 0, 2.1). One can also note that, when \(X\) and \(Y\) are quasi-projective over \(S\) with regular fibers and separated, \(f_t\) and \(f_s\) are automatically complete-intersection morphisms (VIII 1.1), and one obtains a commutative diagram analogous to (7.10.1 bis), with the \(\Gr_\bullet\) replaced by the \(\Gr^\bullet(\ )_{\mathbb Q}\), taking account of the compatibility result (VIII 3.2).

%%%%%%%%%%  scan idx 584  |  folio 571  %%%%%%%%%%
\SGAPrintPage{571}%

\sgaStableTarget{sga6.e10.para.7-11}7.11. For a scheme \(T\) over an algebraically closed field \(k\), one can define the subgroup of \(K^\bullet(T)\) of elements \uline{algebraically equivalent to zero} as the subgroup of elements of the form
\sgaStableTarget{sga6.e10.eq.7-11-1}\begin{equation}\tag{7.11.1}
  E(x_1)-E(x_0),
\end{equation}
where \(x_0\) and \(x_1\) are two \(k\)-rational points of a proper, regular, connected algebraic curve \(C\) over \(k\), \(E\) is an element of \(K^\bullet(T\times_k C)\), and, for every point \(x\) of \(C\) rational over \(k\), one puts
\[
  E(x)=i_x^*(E)\in K^\bullet(T),\qquad\text{where }
  i_x:T\longrightarrow T\times_k C
\]
is the morphism \(t\mapsto(t,x)\) defined by \(x\). The fact that the elements of the form (7.11.1) form a subgroup of \(K^\bullet(T)\) follows by a standard argument, using a product \(C\times_k C'\) and its pair of points \((x_1,x'_1)\), \((x_0,x'_0)\), and the fact that through two points of an irreducible quasi-projective algebraic scheme over \(k\) (here \(C\times_k C'\)) one can always pass an irreducible algebraic curve. (A reader unwilling to admit this fact may define algebraic equivalence by the group generated by the elements of the form (7.11.1).) The elements algebraically equivalent to zero form an ideal of \(K^\bullet(T)\), as one sees at once, which is moreover stable under the operations \(\lambda^i\) when these are defined, i.e., when one works with the ``naive'' \(K^\bullet\); this ideal is moreover contained in the augmentation ideal of \(K^\bullet(T)\), as one verifies trivially. One can therefore define a quotient ring
\sgaStableTarget{sga6.e10.eq.7-11-2}\begin{equation}\tag{7.11.2}
  K^\bullet_{\mathrm{alg}}(T),
\end{equation}
which is even a \(\lambda\)-ring when one works with naive \(K^\bullet\).

In the same way, one can define a subgroup of elements algebraically equivalent to 0 in \(K_\bullet(T)\), taking \(E\) above in \(K_\bullet(T\times_k C)\), and
%%%%%%%%%%  scan idx 585  |  folio 572  %%%%%%%%%%
\SGAPrintPage{572}%
noting that \(i_x:T\to T\times_k C\) is a \uline{perfect} morphism, since it is obtained by the flat base change \(T\times_k C\to C\) from the morphism \(j_x:x\to C\), which is perfect because \(C\) is regular. One thus obtains a quotient group
\sgaStableTarget{sga6.e10.eq.7-11-3}\begin{equation}\tag{7.11.3}
  K_\bullet^{\mathrm{alg}}(T)
\end{equation}
of \(K_\bullet(T)\). This group is moreover naturally a module over \(K^\bullet_{\mathrm{alg}}(T)\), since one sees at once from the definitions that the product of an element of \(K^\bullet(T)\) with an element of \(K_\bullet(T)\) is algebraically equivalent to zero when one of the two factors is, using (IV 2.12.2) for \(i_x\).

One defines (when working with naive \(K^\bullet\)) the ideal of \(\Gr^\bullet(T)\) of elements algebraically equivalent to zero as the ideal generated by the \(c^i(x)\), where \(i\geq1\) and \(x\in K^\bullet(T)\) is algebraically equivalent to zero; the quotient graded ring
\sgaStableTarget{sga6.e10.eq.7-11-4}\begin{equation}\tag{7.11.4}
  \Gr^\bullet_{\mathrm{alg}}(T)
\end{equation}
is then precisely the Chern ring associated with the \(\lambda\)-ring \(K^\bullet_{\mathrm{alg}}(T)\).

Finally, one defines the graded subgroup of \(\Gr_\bullet(T)\) of elements algebraically equivalent to zero (\(T\) being assumed locally noetherian) as the subgroup whose component of dimension \(i\) is the image in \(\Gr_i(T)\) of the subgroup of \(\Fil_iK_\bullet(T)\) consisting of the elements algebraically equivalent to zero in \(K_\bullet(T)\). One finds at once that the quotient graded group
\sgaStableTarget{sga6.e10.eq.7-11-5}\begin{equation}\tag{7.11.5}
  \Gr_\bullet^{\mathrm{alg}}(T)
\end{equation}
is a module over the ring \(\Gr^\bullet_{\mathrm{alg}}(T)\), by passage to the quotient from the module structure of \(\Gr_\bullet(T)\) over \(\Gr^\bullet(T)\) (3.1).

\sgaStableTarget{sga6.e10.para.7-11-6}7.11.6. The preceding definitions would still make sense without assuming that \(k\) is algebraically closed, taking \(C\) above to be \uline{geometrically} connected over \(k\). It is not clear, however, that these definitions would be very useful; in particular, it does not appear that, in the case of the group \(\Pic(T)\), one
%%%%%%%%%%  scan idx 586  |  folio 573  %%%%%%%%%%
\SGAPrintPage{573}%
recovers the classical ``geometric'' notion of algebraic equivalence considered in 2.4.1. We therefore prefer to say that an element of \(K^\bullet(T)\) is algebraically equivalent to zero when its image in \(K^\bullet(T\otimes_k\bar k)\) is so, where \(\bar k\) is an algebraic closure of \(k\), and to proceed in the same way for the definition of algebraic equivalence in the groups \(K_\bullet(T)\), \(\Gr^\bullet(T)\) and \(\Gr_\bullet(T)\).

\sgaStableTarget{sga6.e10.para.7-12}7.12. With these definitions in place, I claim that \uline{the specialization homomorphism} (7.3.1) \uline{carries elements algebraically equivalent to zero to elements algebraically equivalent to zero}, and consequently defines, by passage to the quotient, a homomorphism
\sgaStableTarget{sga6.e10.eq.7-12-1}\begin{equation}\tag{7.12.1}
  K_\bullet^{\mathrm{alg}}(X_t)\longrightarrow K_\bullet^{\mathrm{alg}}(X_s).
\end{equation}

To see this, one reduces immediately to the case where \(A\) is complete and \(k(s)\) is algebraically closed (using a straightforward compatibility made explicit in 7.15). It suffices to see that \(\sigma_\bullet\) carries an element of the form (7.11.1) of \(K_\bullet(X_t)\) to an element algebraically equivalent to zero in \(K_\bullet(X_s)\). Here \(T=X_t\), and we shall accordingly write \(C_t\), \(E_t\) in place of \(C\), \(E\) above. One knows that \(C_t\) is projective over \(k(t)\), so (EGA IV 2.8.5) there is a projective flat scheme \(C\) over \(S\) whose generic fiber is \(C_t\). Using Abhyankar's resolution of singularities of arithmetic surfaces [6], one sees that one can even choose \(C\) to be a regular scheme. (N.B. Flatness is preserved under resolution, by EGA IV 14.3.8.) Since \(C\) is proper over \(S\), the points \(x_{0t}\) and \(x_{1t}\) of \(C_t\) rational over \(k(t)\) extend to sections \(x_0\), \(x_1\) of \(C\) over \(S\), and, since \(C\) is regular, it is well known that \(C\) is smooth over \(S\) at the points of \(x_0(S)\), \(x_1(S)\) (EGA IV, 17.12.1 c) \(\Rightarrow\) a) and 19.1.1). As a %%%%%%%%%%  scan idx 587  |  folio 574  %%%%%%%%%%
\SGAPrintPage{574}%
consequence, the immersions \(x_i:S\to C\) (\(i=0,1\)) are regular and remain regular immersions \(x_i:s\to C_s\) under the base change \(s\to S\). Note that, by the exact homotopy sequence (IX 1.1), \(E_t\in K_\bullet(X_t\times_{k(t)}C_t)\) comes from an element \(E\in K_\bullet(X\times_SC)\), so the element \(E_t(x_{1t})-E_t(x_{0t})\) of \(K_\bullet(X_t)\) comes from the element \(E(x_1)-E(x_0)\) of \(K_\bullet(X\times_SC)\), and consequently its image in \(K_\bullet(X_s)\) under (7.3.1) is precisely
\sgaStableTarget{sga6.e10.eq.unnamed.occ2}\[
  i^*(E(x_1)-E(x_0)).
  \tag{*}
\]
Of course, for a section \(x\) of \(C\) over \(S\), one puts
\[
  E(x)=i_x^*(E),\qquad\text{where } i_x:X\longrightarrow X\times_SC
\]
is defined as above, and can be regarded as obtained from \(x\) by the flat base change \(X\times_SC\to C\); it is therefore a regular immersion, a fortiori a perfect morphism, which gives a meaning to the expression \(E(x)\). For the same reason, the expression \(E_s(x_s)\) is defined, and, by transitivity of inverse images (IV 2.12) in the following commutative diagram of perfect morphisms:
\[
\begin{tikzcd}[column sep=large,row sep=large]
X \arrow[r,"i_x"] & X\times_SC\\
X_s \arrow[u,"i"] \arrow[r,"i_{x_s}"] & X_s\times_{k(s)}C_s \arrow[u]
\end{tikzcd}
\;,
\]
one finds that the left-hand side of \((*)\) is equal to
\sgaStableTarget{sga6.e10.eq.unnamed.occ3}\[
  E_s(x_{1s})-E_s(x_{0s}).
  \tag{**}
\]
Since \(k(s)\) is algebraically closed, the reduced irreducible components of \(C_s\) are geometrically integral over \(k(s)\) and their intersection points are rational over \(k(s)\). One concludes readily that the expression \((**)\) is algebraically equivalent to zero in \(K_\bullet(X_s)\), using the connectedness
%%%%%%%%%%  scan idx 588  |  folio 575  %%%%%%%%%%
\SGAPrintPage{575}%
of \(C_s\), which follows from the geometric connectedness of \(C_t\).

From the homomorphism (7.12.1), using 7.4, one obtains a homomorphism
\sgaStableTarget{sga6.e10.eq.7-12-2}\begin{equation}\tag{7.12.2}
  \Gr_\bullet^{\mathrm{alg}}(X_t)\longrightarrow \Gr_\bullet^{\mathrm{alg}}(X_s).
\end{equation}

When \(X_s\) and \(X_t\) are moreover assumed regular, (7.12.1) can also be interpreted as a ring homomorphism
\sgaStableTarget{sga6.e10.eq.7-12-3}\begin{equation}\tag{7.12.3}
  K^\bullet_{\mathrm{alg}}(X_t)\longrightarrow K^\bullet_{\mathrm{alg}}(X_s),
\end{equation}
which, when \(X\) is moreover separated, is a homomorphism of \(\lambda\)-rings, and consequently defines a homomorphism
\sgaStableTarget{sga6.e10.eq.7-12-4}\begin{equation}\tag{7.12.4}
  \Gr^\bullet_{\mathrm{alg}}(X_t)\longrightarrow \Gr^\bullet_{\mathrm{alg}}(X_s).
\end{equation}

\sgaStableTarget{sga6.e10.para.7-12-5}7.12.5. Let \(T\) be a smooth proper algebraic scheme over a field \(k\). One checks easily that an element \(x\) of \(K^\bullet(T)\), for which there is an integer \(n\geq1\) with \(nx\) algebraically equivalent to zero, is numerically equivalent to zero. It is plausible that the converse is always true, at least when \(T\) is projective. In this question, one can plainly suppose that \(k\) is algebraically closed, and replace \(K^\bullet(T)\) by \(\Gr^\bullet(T)\) or \(\Gr_\bullet(T)\) without changing the conjecture. One also obtains an equivalent conjecture by working with the Chow ring \(A(T)\) (cf. XIV 4.1), obtaining the following more or less classical conjecture: if \(Z\) is an algebraic cycle numerically equivalent to zero on \(T\), there is an integer \(n\geq1\) with \(nZ\) algebraically equivalent to zero. In the case where this conjecture is correct for \(X_t\), the compatibility of the homomorphism (7.3.1) with algebraic equivalence would imply that this specialization homomorphism (when \(f:X\to S\) is a projective smooth morphism) is also compatible with numerical equivalence, as conjectured in 7.9.

%%%%%%%%%%  scan idx 589  |  folio 576  %%%%%%%%%%
\SGAPrintPage{576}%

\sgaStableTarget{sga6.e10.sec.7.13}\subsection*{7.13. Relations with \(\ell\)-adic cohomology}
Recall (SGA 5 VII) that, for every locally free Module of finite type \(\underline{E}\) on a scheme \(T\), and every prime number distinct from the residue characteristics of \(T\), one defines \(\ell\)-adic Chern classes (for \(i\geq 0\))
\[
  c^i(E)\in H^{2i}(T,\mathbb Z_\ell(i)),
\]
satisfying an additivity formula
\[
  c(E)=c(E')c(E'')
\]
when \(E\) is an extension of a locally free module \(E'\) by a locally free module \(E''\). From this additivity one concludes that the total Chern class \(c(E)=\sum_i c^i(E)\) arises from a homomorphism
\sgaStableTarget{sga6.e10.eq.7-13-0}\begin{equation}\tag{7.13.0}
  c:K^\bullet(T)\longrightarrow
  \widehat G\bigl(H^{2*}(T,\mathbb Z_\ell(*))\bigr)
  =
  1+\prod_{i\geq 1}H^{2i}(T,\mathbb Z_\ell(i)),
\end{equation}
(where the target is a group under multiplication, as explained in V 1, and where \(K^\bullet(T)\) is the ``\uline{naive}'' \(K^\bullet\)). One obtains canonical maps
\sgaStableTarget{sga6.e10.eq.7-13-1}\begin{equation}\tag{7.13.1}
  c^i:K^\bullet(T)\longrightarrow H^{2i}(T,\mathbb Z_\ell(i)).
\end{equation}

Now let \(X\) be a \uline{proper} flat scheme over \(S\), with \uline{regular fibers}, and suppose that \(S\) is \uline{henselian}. Assume that \(\ell\neq\) the residue characteristic of \(S\). One can then form the diagram
\sgaStableTarget{sga6.e10.eq.7-13-2}\begin{equation}\tag{7.13.2}
\begin{tikzcd}[column sep=large,row sep=large]
K^\bullet(X_t) \arrow[r,"c^i"] \arrow[d,"\sigma"'] &
H^{2i}(X_t,\mathbb Z_\ell(i))\\
K^\bullet(X_s) \arrow[r,"c^i"] & H^{2i}(X_s,\mathbb Z_\ell(i)) \arrow[u,"\sigma^{-1}"']
\end{tikzcd}
\;,
\end{equation}
where the horizontal arrows are the ``Chern class'' maps for \(X_t\) and \(X_s\) respectively, and where \(\sigma^{-1}\) is the composite homomorphism
\[
  H(X_s)\xleftarrow{\substack{i^*\\\sim}}H(X)\xrightarrow{j^*}H(X_t);
\]
%%%%%%%%%%  scan idx 590  |  folio 577  %%%%%%%%%%
\SGAPrintPage{577}%
the arrows displayed here are restriction homomorphisms, and one uses the fact that \(i^*\) is an isomorphism by the proper base change theorem (SGA 4 XII 5.5). I claim that \uline{the square} (7.13.2) \uline{is commutative}. The proof is essentially the same as that establishing the commutativity of (7.10.2); here it suffices to use the compatibility of the formation of Chern classes with inverse images for the morphisms \(j:X_t\to X\) and \(i:X_s\to X\).

\sgaStableTarget{sga6.e10.para.7-13-3}7.13.3. Note that the square (7.13.2) arises by passage to the projective limit from the analogous squares in which cohomology with coefficients in the \(\ell\)-adic sheaf \(\mathbb Z_\ell(i)\) is replaced by cohomology with coefficients in the torsion sheaves \(\mu_n^{\otimes i}\) (where \(n\) runs through the powers of \(\ell\)). Thus the commutativity of (7.13.2) merely expresses the commutativity of the squares involving Chern classes mod. \(n\), for every power \(n\) of \(\ell\). Choose an algebraic closure \(\bar K\) of \(K=k(t)\), and hence an algebraic closure \(\bar k\) of \(k=k(s)\) (namely the residue field of the integral closure of \(V\) in \(\bar K\)). For every finite subextension \(K_\alpha\) of \(\bar K\), one can consider the analogous commutative diagram (for fixed \(n\)) in the situation obtained from \(f:X\to S\) by the base change \(S_\alpha\to S\), with \(S_\alpha\) the normalization of \(S\) in \(K_\alpha\). For this, note that, by Krull--Akizuki and since \(S\) is henselian, \(S_\alpha\) is indeed a trait; moreover, one must suppose here that the fibers of \(X\) over \(S\) are geometrically regular, i.e., that \(X\) is \uline{smooth} over \(S\), so that this hypothesis is preserved by base change. As \(\alpha\) varies, one thus obtains an inductive system of commutative squares. Passing to the inductive limit and using (SGA 4 VII 5.8) and (IV 3.2), one obtains a commutative square (cf. 7.16):
%%%%%%%%%%  scan idx 591  |  folio 578  %%%%%%%%%%
\SGAPrintPage{578}%
\[
\begin{tikzcd}[column sep=large,row sep=large]
K^\bullet(X_{\bar t}) \arrow[r,"c^i"] \arrow[d,"\sigma"'] &
H^{2i}(X_{\bar t},\mu_n^{\otimes i})\\
K^\bullet(X_{\bar s}) \arrow[r,"c^i"] &
H^{2i}(X_{\bar s},\mu_n^{\otimes i}) \arrow[u,"\sigma^{-1}"]
\end{tikzcd}
\;,
\]
where \(X_{\bar t}\), \(X_{\bar s}\) are the \uline{geometric} fibers of \(X\). Moreover, one knows that the cohomological specialization homomorphism \(\sigma^{-1}\) is now an isomorphism (SGA 4 XVI 2.2), so the commutativity of the preceding square can also be expressed as the commutativity of the square obtained by reversing the preceding arrow. Passing to the projective limit over \(n\) in the commutative squares thus obtained, one finds a commutative square
\sgaStableTarget{sga6.e10.eq.7-13-4}\begin{equation}\tag{7.13.4}
\begin{tikzcd}[column sep=large,row sep=large]
K^\bullet(X_{\bar t}) \arrow[r] \arrow[d,"\sigma"'] &
H^{2i}(X_{\bar t},\mathbb Z_\ell(i)) \arrow[d,"\wr"']\\
K^\bullet(X_{\bar s}) \arrow[r] &
H^{2i}(X_{\bar s},\mathbb Z_\ell(i)).
\end{tikzcd}
\end{equation}

\sgaStableTarget{sga6.e10.para.7-13-5}7.13.5. Another way to express this commutativity is to say that, when \(x_t\) is an element of \(K^\bullet(X_t)\) and \(x_s\) is its image in \(K^\bullet(X_s)\), \(c^i(x_{\bar t})\) and \(c^i(x_{\bar s})\) are the values, at the geometric points \(\bar t\) and \(\bar s\) respectively, of a \uline{section} of \(R^{2i}f_*(\mathbb Z_\ell(i))\). This section is moreover readily interpreted by taking an element \(x\in K^\bullet(X)\) inducing \(x_t\), whence a \(c^i(x)\in H^{2i}(X,\mathbb Z_\ell(i))\), which defines the required section of \(R^{2i}f_*(\mathbb Z_\ell(i))\).

\sgaStableTarget{sga6.e10.para.7-13-6}7.13.6. Let \(T\) be a scheme. Then, by passage to the associated graded groups, the homomorphism (7.13.0) defines homomorphisms
\[
  \psi_i:\Gr^i(T)\longrightarrow H^{2i}(T,\mathbb Z_\ell(i)),
\]
the image under \(\psi_i\) of an element \(x\) of \(\Filt^iK^\bullet(T)\) being, by definition, \(c^i(x)\). (The fact that (7.13.0) is compatible with the \(\lambda\)-filtration on
%%%%%%%%%%  scan idx 592  |  folio 579  %%%%%%%%%%
\SGAPrintPage{579}%
\(K^\bullet(T)\) and the degree filtration on the second member follows readily from (V 6.6.4) and from the fact that the \(c^i\) define a \uline{Chern theory} of \(K^\bullet(T)\) in the sense of (V 6.13), when \(K^\bullet(T)\) is regarded as a \(\lambda\)-algebra augmented over the binomial ring \(H^0(T,\mathbb Z)\)). One thereby obtains maps
\sgaStableTarget{sga6.e10.eq.7-13-7}\begin{equation}\tag{7.13.7}
  \operatorname{c\ell}:\Gr^i(T)\longrightarrow H^{2i}(T,\mathbb Q_\ell(i))
  \stackrel{\mathrm{dfn}}{=}
  H^{2i}(T,\mathbb Z_\ell(i))\otimes_{\mathbb Z_\ell}\mathbb Q_\ell,
\end{equation}
by the formula
\sgaStableTarget{sga6.e10.eq.7-13-8}\begin{equation}\tag{7.13.8}
  \operatorname{c\ell}(x^{(i)})=
  (1/(-1)^{i-1}(i-1)!)\,\psi_i(x^{(i)}),
\end{equation}
compare (XIV 5.1 and 5.2). With this understood, the commutativity of diagrams (7.13.2) and (7.13.4) trivially implies that of the analogous diagrams
\sgaStableTarget{sga6.e10.eq.7-13-9}\begin{equation}\tag{7.13.9}
\begin{tikzcd}[column sep=large,row sep=large]
\Gr^i(X_t)_{\mathbb Q} \arrow[r,"\operatorname{c\ell}"] \arrow[d,"\sigma_{\mathbb Q}"'] &
H^{2i}(X_t,\mathbb Q_\ell(i))\\
\Gr^i(X_s)_{\mathbb Q} \arrow[r,"\operatorname{c\ell}"] &
H^{2i}(X_s,\mathbb Q_\ell(i)) \arrow[u,"\sigma^{-1}"]
\end{tikzcd}
\;,
\end{equation}
\sgaStableTarget{sga6.e10.eq.7-13-10}\begin{equation}\tag{7.13.10}
\begin{tikzcd}[column sep=large,row sep=large]
\Gr^i(X_{\bar t})_{\mathbb Q} \arrow[r,"\operatorname{c\ell}"] \arrow[d] &
H^{2i}(X_{\bar t},\mathbb Q_\ell(i)) \arrow[d,"\wr"]\\
\Gr^i(X_{\bar s})_{\mathbb Q} \arrow[r,"\operatorname{c\ell}"] &
H^{2i}(X_{\bar s},\mathbb Q_\ell(i))
\end{tikzcd}
\;.
\end{equation}
\sgaStableTarget{sga6.e10.rem.7-13-11}\textbf{Remark 7.13.11.} The commutativity of the last two diagrams, which can be expressed by saying that specialization of algebraic cycles is compatible with formation of the associated \(\ell\)-adic cohomology classes, is not very satisfactory in the form stated, since we have had to replace the coefficients \(\mathbb Z_\ell(i)\) by \(\mathbb Q_\ell(i)\), that is, to neglect torsion. To obtain a statement that does not neglect torsion, it seems that one
%%%%%%%%%%  scan idx 593  |  folio 580  %%%%%%%%%%
\SGAPrintPage{580}%
will be forced to work with the Chow ring, since it is on that ring that the homomorphism \(\operatorname{c\ell}\) is naturally defined as a homomorphism with values in \(H^{2*}(-,\mathbb Z_\ell(*))\); compare (XIV 4.1). Unfortunately, no theory of specialization for the Chow ring has yet been developed (cf. 7.14.) that would allow one to write commutative squares analogous to (7.13.9) and (7.13.10) in terms of Chow rings and cohomology with coefficients \(\mathbb Z_\ell(*)\). All one can say when \(f:X\to S\) is smooth and proper is that there is a commutative square
\sgaStableTarget{sga6.e10.eq.7-13-12}\begin{equation}\tag{7.13.12}
\begin{tikzcd}[column sep=large,row sep=large]
Z^i(X_t) \arrow[r,"\operatorname{c\ell}"] \arrow[d,"\sigma"'] &
H^{2i}(X_t,\mathbb Z_\ell(i))\\
Z^i(X_s) \arrow[r,"\operatorname{c\ell}"] &
H^{2i}(X_s,\mathbb Z_\ell(i)) \arrow[u]
\end{tikzcd}
\;,
\end{equation}
analogous to (7.13.9), and a corresponding square analogous to (7.13.10), which we leave to the reader to write down, where \(Z^i\) denotes the group of algebraic cycles (modulo no relation at all) of codimension \(i\), where the horizontal arrows are the ``associated cohomology class'' homomorphisms defined in (SGA 5 IV), and where \(\sigma\) denotes the specialization homomorphism for algebraic cycles, associating with a cycle \(Z_t\) the cycle \(Z_s\), where \(Z\) is the schematic closure of \(Z_t\) in \(X\) (cf. 7.14). This commutativity follows from the following fact, established in (SGA 5 IV): if \(Z\) is a closed subscheme of \(X\) of codimension \(\geq i\) on every fiber, one can associate with it a cohomology class \(\operatorname{c\ell}(Z)\in H^{2i}(X,\mathbb Z_\ell(i))\) whose formation commutes with every base change and, in particular, such that its restriction to \(X_t\) (resp. \(X_s\)) equals \(\operatorname{c\ell}(Z_t)\) (resp. \(\operatorname{c\ell}(Z_s)\)). Of course, when one knows that the homomorphism \(\sigma\) in (7.13.12) carries cycles rationally equivalent to zero to cycles rationally equivalent to zero (\(X\) being projective over \(S\)),
%%%%%%%%%%  scan idx 594  |  folio 581  %%%%%%%%%%
\SGAPrintPage{581}%
one can replace the \(Z^i\) by the \(A^i\) in this diagram, and obtain a diagram of the desired kind in terms of the Chow rings of \(X_t\) and \(X_s\).

\sgaStableTarget{sga6.e10.rem.7-13-13}\textbf{Remark 7.13.13.} Let \(T\) be a smooth projective scheme over an algebraically closed field. It follows trivially from the formalism developed in (SGA 5 IV) that, if \(Z\) is a cycle of codimension \(i\) on \(T\) whose image in \(H^{2i}(T,\mathbb Q_\ell(i))\) is zero, then \(Z\) is numerically equivalent to zero. One of the standard conjectures in the theory of algebraic cycles asserts that the converse is also true. Moreover, by (VII 4.11), it is immediate that one obtains an equivalent conjecture by replacing the group of algebraic cycles with the group \(\Gr^\bullet(T)_{\mathbb Q}\), using definition (7.13.8), or again with the group \(K^\bullet(T)\), by stating that numerical equivalence to zero of an element of \(K^\bullet(T)\) is equivalent to vanishing of its augmentation and of its Chern classes in the \(H^{2i}(T,\mathbb Q_\ell(i))\).

With this understood, it follows immediately from the commutativity of (7.13.10) that, if \(X_{\bar t}\) satisfies the preceding conjecture, then the specialization homomorphism \(\Gr^\bullet(X_{\bar t})\longrightarrow \Gr^\bullet(X_{\bar s})\) (and, a fortiori, \(\Gr^\bullet(X_t)\longrightarrow \Gr^\bullet(X_s)\)) is compatible with numerical equivalence, as conjectured in 7.9. Note also that the conjecture invoked for this purpose (numerical equivalence implies \(\mathbb Q_\ell\)-cohomological equivalence) is weaker than the one considered in 7.12.5.

\sgaStableTarget{sga6.e10.sec.7.14}\subsection*{7.14. Specialization homomorphism for groups of cycles and cycle classes.}
If \(T\) is a noetherian scheme, denote by \(Z_i(T)\) (resp. \(Z^i(T)\)) the group of algebraic cycles on \(T\) purely of dimension \(i\) (resp. purely of codimension \(i\)). When \(f:X\to S\) is of finite type, one then defines
%%%%%%%%%%  scan idx 595  |  folio 582  %%%%%%%%%%
\SGAPrintPage{582}%
a specialization homomorphism:
\sgaStableTarget{sga6.e10.eq.7-14-1}\begin{equation}\tag{7.14.1}
  \sigma_i:Z_i(X_t)\longrightarrow Z_i(X_s),
\end{equation}
by defining, for every prime cycle \(Z_t\) of \(X_t\) purely of dimension \(i\), identified with an integral subscheme of \(X_t\), the element \(\sigma_i(Z_t)\) of \(Z_i(X_s)\) as the cycle associated with the closed subscheme \(Z_s\) of \(X_s\), where \(Z\) is the schematic closure of \(Z_t\) in \(X\) (which is a closed subscheme of \(X\) flat over \(S\) (EGA IV 2.8), with generic fiber equal to the original \(Z_t\), and satisfying \(\dim Z_s=\dim Z_t\) (EGA IV 7.1.13)). When \(X_s\) is locally equidimensional (i.e., all its irreducible components through a given point have the same dimension), one immediately verifies that the specialization homomorphism just defined induces homomorphisms
\sgaStableTarget{sga6.e10.eq.7-14-1bis}\begin{equation}\tag{7.14.1 bis}
  \sigma^i:Z^i(X_t)\longrightarrow Z^i(X_s).
\end{equation}

Now suppose that \(f\) is quasi-projective, so that the notion of rational equivalence of cycles is available on \(X_t\) and on \(X_s\) (cf. [5], where the theory is developed over an algebraically closed field). It seems very plausible that homomorphism (7.14.1) carries cycles rationally equivalent to zero to cycles rationally equivalent to zero, and thus induces a corresponding homomorphism on the ``Chow'' groups of cycle classes
\sgaStableTarget{sga6.e10.eq.7-14-2}\begin{equation}\tag{7.14.2}
  \sigma_i:A_i(X_t)\longrightarrow A_i(X_s),
\end{equation}
or again, when \(X_s\) is locally equidimensional:
\sgaStableTarget{sga6.e10.eq.7-14-2bis}\begin{equation}\tag{7.14.2 bis}
  \sigma^i:A^i(X_t)\longrightarrow A^i(X_s).
\end{equation}

One might try to prove this by a method paralleling that of 7.12, but in which the recourse to resolution of singularities and the reduction
%%%%%%%%%%  scan idx 596  |  folio 583  %%%%%%%%%%
\SGAPrintPage{583}%
to an algebraically closed residue field would become unnecessary, since \(C=\mathbb P^1_S\) is available here as a ``parameter curve'' along which to deform cycles. A difficulty nevertheless arises, for with the notation of 7.12, \(E\) must now denote a cycle on \(P^1_S\times_SX\) such that \(E_t(x_{1t})-E_t(x_{0t})\) ``is defined,'' but it by no means follows that the cycle \(E_s(x_{1s})-E_s(x_{0s})\) is also defined. As was to be expected, one is thus reduced to formulating a statement of the type of Chow's ``moving lemma'' in order to define (7.14.2).

Observe, however, that the existence of homomorphism (7.4.1) shows that (7.14.1) is compatible with the equivalence relation on cycles defined by the surjective homomorphism \(Z_\bullet(T)\longrightarrow \Gr_\bullet(T)\), which is only slightly coarser than linear equivalence; thus, when \(Z\) is quasi-projective and smooth over a field, we point out elsewhere that it differs from linear equivalence only up to torsion (XIV 4.1, 4.2). In this case one therefore obtains a homomorphism analogous to (7.14.2), or better to (7.14.2 bis), but with both members tensored by \(\mathbb Q\).

The most natural way to define a homomorphism (7.14.2) directly would seem to be to proceed as in 7.3. This would, however, require a theory of Chow groups \(A_i\) (or of the Chow ring in the regular case) for schemes more general than schemes quasi-projective over a field, including \(X\), and a contravariant law for these groups that would allow one to define the restriction homomorphism to \(X_s\), \(A_i(X)\longrightarrow A_i(X_s)\). Compare XIV 8. Once such a sufficiently general theory of Chow groups and rings is available, it should be possible to transcribe into Chow theory all the results enumerated above for the specialization homomorphism in \(K\)-theory.

%%%%%%%%%%  scan idx 597  |  folio 584  %%%%%%%%%%
\SGAPrintPage{584}%
\sgaStableTarget{sga6.e10.sec.7.15}\subsection*{7.15. Varying \(S\).}
Let \(S'\) be a second trait, and let
\[
  S'\longrightarrow S
\]
be a morphism of traits, i.e., a dominant local morphism (i.e., carrying the closed point to the closed point and the generic point to the generic point). Denote by \(s'\) (resp. \(t'\)) the closed point (resp. the generic point) of \(S'\), and set \(X'=X\times_SS'\). Suppose that \(X'\) is noetherian (e.g., \(X\) of finite type over \(S\)). There is a commutative diagram
\[
\begin{tikzcd}
X & X_t \arrow[l,hook,"j_X"'] & & \\
& & X' \arrow[ull] & X'_{t'} \arrow[l,hook,"i'_{X'}"] \arrow[ull] \\
X_s \arrow[uu,hook,"i_X"] & & & \\
& & X'_{s'} \arrow[uu,hook,"j_{X'}"'] \arrow[ull] &
\end{tikzcd}
\;,
\]
whose oblique arrows are flat morphisms. One can therefore take the inverse-image homomorphisms for \(K_\bullet\) along all the arrows of the diagram (IV 2.12). Applying definition 7.3, one finds that the following diagram
\sgaStableTarget{sga6.e10.eq.7-15-1}\begin{equation}\tag{7.15.1}
\begin{tikzcd}[column sep=large,row sep=large]
K_\bullet(X_t) \arrow[r,"\sigma_X"] \arrow[d] &
K_\bullet(X_s) \arrow[d]\\
K_\bullet(X'_{t'}) \arrow[r,"\sigma_{X'}"] &
K_\bullet(X'_{s'})
\end{tikzcd}
\end{equation}
is commutative, where the vertical arrows are the inverse-image homomorphisms. One immediately obtains an analogous commutative diagram with \(\Gr_\bullet\) in place of \(K_\bullet\); and when \(X_s\) and \(X_t\) are regular, one can also write \(K^\bullet\) and \(\Gr^\bullet\) in place of \(K_\bullet\) and \(\Gr_\bullet\).


%%%%%%%%%%  scan idx 598  |  folio 585  %%%%%%%%%%
\SGAPrintPage{585}%
\sgaStableTarget{sga6.e10.sec.7.16}\subsection*{7.16. Passage to geometric fibers.}
Let \(\bar K\) be an algebraic closure of \(K=k(t)\), the inductive limit of its finite subextensions \(K_i\), and let \(\bar A=\lim \bar A_i\) be the normalization of \(A\) in \(\bar K\), the inductive limit of its normalizations \(\bar A_i\) in the \(K_i\). Choose a maximal ideal \(\bar{\underline{m}}\) of \(\bar A\), thus inducing on each \(\bar A_i\) a maximal ideal \(\underline{m}_i\), and set \(B_i=\bar A_{i\,\underline{m}_i}\), which by Krull--Akizuki is again a discrete valuation ring, whose spectrum we denote by \(S_i\) and whose points we denote by \(s_i\) and \(t_i\). The residue field \(\bar k=\bar A/\bar{\underline{m}}\) is an algebraic closure of \(k=k(s)\), and is the inductive limit of the residue fields \(k_i=\bar A_i/\underline{m}_i=k(s_i)\) of the \(B_i\). Set \(\bar t=\Spec(\bar K)\), \(\bar s=\Spec(\bar k)\). \uline{Now suppose} \(X\) \uline{of finite type over} \(S\); then \(X_{\bar t}\) and \(X_{\bar s}\) are of finite type over \(\bar t\) and \(\bar s\), respectively, and hence \uline{noetherian}. By IV 3.2.3, we have
\[
  K_\bullet(X_{\bar t})\xleftarrow{\sim}\varinjlim_i K_\bullet(X_{t_i}),
  \qquad
  K_\bullet(X_{\bar s})\xleftarrow{\sim}\varinjlim_i K_\bullet(X_{s_i}).
\]
On the other hand, by 7.15 there is a homomorphism from the inductive system of the \(K_\bullet(X_{t_i})\) to the inductive system of the \(K_\bullet(X_{s_i})\). Passing to the limit, one obtains a canonical homomorphism:
\sgaStableTarget{sga6.e10.eq.7-16-1}\begin{equation}\tag{7.16.1}
  K_\bullet(X_{\bar t})\longrightarrow K_\bullet(X_{\bar s}),
\end{equation}
again called, as expected, the \uline{specialization homomorphism}. It induces a homomorphism \(\Gr_\bullet(X_{\bar t})\longrightarrow \Gr_\bullet(X_{\bar s})\) by 7.4, using the fact (whose verification is left to the reader) that the functor \(\Gr_\bullet\) (just like \(K_\bullet\)) commutes with the limits occurring here. Similarly, when \(f\) is smooth, one can interpret (7.16.1) as a homomorphism on \(K^\bullet\), which is moreover a ring homomorphism, and even a homomorphism of \(\lambda\)-rings when \(X\) is separated; in this case it therefore induces a homomorphism on \(\Gr^\bullet\).


%%%%%%%%%%  scan idx 599  |  folio 586  %%%%%%%%%%
\SGAPrintPage{586}%
\sgaStableTarget{sga6.e10.para.7-16-2}7.16.2.
We shall not review here, in the case of geometric fibers, the extension of all the properties of specialization homomorphisms proved in the preceding sections for the ``arithmetic'' fibers; they all follow trivially from the latter by passage to the limit. Let us mention only the following interesting consequence of 7.9.4 when \(f\) is smooth and proper: \(\Gr_\bullet^{\mathrm{num}}(X_{\bar t})\) is isomorphic to a quotient of a subgroup of \(\Gr_\bullet^{\mathrm{num}}(X_{\bar s})\); in particular, its rank is bounded above by that of the latter group. This is therefore a \uline{semicontinuity} result for the rank of the group of cycles modulo numerical equivalence.

\sgaStableTarget{sga6.e10.sec.7.17}\subsection*{7.17. Case of an arbitrary base.}
Now let
\[
  f:X\longrightarrow S
\]
be a flat and quasi-compact morphism of noetherian preschemes, and let \(s,t\) be two points of \(S\) such that \(s\) is a specialization of \(t\). We wish to relate the groups \(K_\bullet\), etc., for the two fibers \(X_t\) and \(X_s\), and for corresponding geometric fibers. For this, note (EGA II 7.1.9) that there exists a morphism
\[
  g:S'\longrightarrow S,
\]
where \(S'\) is a trait, and where \(g\) carries the closed point \(s'\) to \(s\) and the generic point \(t'\) to \(t\); one may even assume that \(k(t)\xrightarrow{\sim}k(t')\). Suppose that \(X'=X\times_S S'\) is noetherian, as is the case when \(f\) is of finite type or \(S'\) is essentially of finite type over \(S\) (EGA IV 1.3.8) (which one may assume when \(S\) itself is universally Japanese (EGA \(0_{\mathrm{IV}}\) 23.1.1)). Then the specialization homomorphism for \(X'\) defines a homomorphism
\sgaStableTarget{sga6.e10.eq.7-17-1}\begin{equation}\tag{7.17.1}
  \sigma_\bullet:K_\bullet(X_t)\xrightarrow{\sim}K_\bullet(X_{t'})
  \longrightarrow K_\bullet(X_{s'}).
\end{equation}
%%%%%%%%%%  scan idx 600  |  folio 587  %%%%%%%%%%
\SGAPrintPage{587}%
One thus obtains a specialization homomorphism on \(K_\bullet(X_t)\), but with values not in \(K_\bullet(X_s)\), but in \(K_\bullet(X_{s'})\), where \(X_{s'}\) is obtained from \(X_s\) by an extension \(k(s)\to k(s')\) of the base fields (which may be assumed to be of finite type when \(S\) is universally Japanese). Passing to geometric fibers as explained in 7.16, one likewise obtains a homomorphism
\sgaStableTarget{sga6.e10.eq.7-17-2}\begin{equation}\tag{7.17.2}
  \sigma_\bullet:K_\bullet(X_{\bar t})\longrightarrow K_\bullet(X_{\overline{s'}});
\end{equation}
so here again the ``geometric'' specialization homomorphism takes its values not in \(K_\bullet(X_{\bar s})\), but in \(K_\bullet(X_{\overline{s'}})\), where \(X_{\overline{s'}}\) is obtained from \(X_{\bar s}\) by the extension of algebraically closed fields \(k(\bar s)\to k(\overline{s'})\).

\sgaStableTarget{sga6.e10.para.7-17-3}7.17.3. In general, the homomorphism \(K_\bullet(X_{\bar s})\to K_\bullet(X_{\overline{s'}})\) is not an isomorphism (because of the possible presence of ``continuous'' components in the group \(K_\bullet(X_{\bar s})\)). It is, however, easy to verify that it is an isomorphism modulo algebraic equivalence (7.11), and also modulo numerical equivalence (4.5.1) when \(f\) is smooth and proper. Consequently, using 7.12, one finds that (7.17.2) induces a specialization homomorphism
\sgaStableTarget{sga6.e10.eq.7-17-3-1}\begin{equation}\tag{7.17.3.1}
  K_\bullet^{\mathrm{alg}}(X_{\bar t})\longrightarrow
  K_\bullet^{\mathrm{alg}}(X_{\bar s});
\end{equation}
and when \(f\) is proper and smooth, provided that specialization homomorphism (7.17.2) is compatible with numerical equivalence (cf. 7.9.4), one likewise obtains a homomorphism
\sgaStableTarget{sga6.e10.eq.7-17-3-2}\begin{equation}\tag{7.17.3.2}
  K_\bullet^{\mathrm{num}}(X_{\bar t})\longrightarrow
  K_\bullet^{\mathrm{num}}(X_{\bar s}).
\end{equation}
We leave to the reader the task of developing (if needed) the corresponding constructions for \(\Gr_\bullet\), \(K^\bullet\), and \(\Gr^\bullet\), and the properties of these specialization homomorphisms that follow trivially from the properties developed in the preceding sections. Let us mention only
%%%%%%%%%%  scan idx 601  |  folio 588  %%%%%%%%%%
\SGAPrintPage{588}%
that it follows from 7.16.2 that \(\Gr_\bullet^{\mathrm{num}}(X_{\bar t})\) is isomorphic to a quotient of a subgroup of \(\Gr_\bullet^{\mathrm{num}}(X_{\bar s})\) (\(f\) proper and smooth); in particular, its rank is again bounded above by the rank of the latter.

\sgaStableTarget{sga6.e10.para.7-17-4}7.17.4. When \(S\) is regular at \(s\) and \(t\) is maximal, one can construct \(S'\) by blowing up \(s\) in the localization of \(S\) at \(s\), and localizing at the generic point of the fiber of \(s\) in the blow-up. One then obtains for \(k(s')\) a \uline{pure extension} of \(k(s)\). Now it is not difficult to show that, when \(k'\) is a pure extension of a field \(k\), for every noetherian scheme \(T\) over \(k\) the canonical homomorphism \(K_\bullet(T)\to K_\bullet(T\otimes_k k')\) is an isomorphism, by representing \(k'\) as the inductive limit of the coordinate rings of affine open subsets of schemes of the form \(\Spec k[T_1,\ldots,T_n]\), and using (IX 1.6). Thus one concludes that (7.17.1) can be regarded as a canonical homomorphism
\sgaStableTarget{sga6.e10.eq.7-17-4-1}\begin{equation}\tag{7.17.4.1}
  K_\bullet(X_t)\longrightarrow K_\bullet(X_s).
\end{equation}
But note that one cannot deduce from this a homomorphism \(K_\bullet(X_{\bar t})\to K_\bullet(X_{\bar s})\), for when one passes to a finite extension \(K_i\) of \(K=k(t)\) and normalizes \(S\) in \(K_i\), the regularity hypothesis on \(S\) is lost.

Another way to define a homomorphism (7.17.4.1) is to reduce to the case where \(S\) is local with closed point \(s\) (which is trivial), and to consider an increasing sequence of \uline{regular} closed subschemes \(S_i\) of \(S\), with generic points \(t_i\), such that \(\dim S_i=i\) (\(0\leq i\leq n=\dim S\)). Since for every \(i\), \(0\leq i\leq n-1\), \(\underline{O}_{S_{i+1},t_i}\) is a discrete valuation ring with residue field \(k(t_i)\) and fraction field \(k(t_{i+1})\), there is a specialization homomorphism
\[
  K_\bullet(X_{t_{i+1}})\longrightarrow K_\bullet(X_{t_i}),
\]
%%%%%%%%%%  scan idx 602  |  folio 589  %%%%%%%%%%
\SGAPrintPage{589}%
and the composite of these homomorphisms for \(0\leq i\leq n-1\) defines a homomorphism (7.17.4.1), which one can show (by proceeding as in 7.3) to be the same as the one defined above by means of a blow-up. The advantage of the second method, however, is that here one even has homomorphisms of the form (7.16.1):
\[
  K_\bullet(X_{\bar t_{i+1}})\longrightarrow K_\bullet(X_{\bar t_i}),
\]
(with suitable coherent choices of the algebraic closures \(k(\bar t_i)\) of the \(k(t_i)\)), which allow one to define by composition a homomorphism
\sgaStableTarget{sga6.e10.eq.7-17-4-2}\begin{equation}\tag{7.17.4.2}
  K_\bullet(X_{\bar t})\longrightarrow K_\bullet(X_{\bar s}),
\end{equation}
that probably depends essentially on the choice of the \(S_i\), that is, of the \(t_i\).

\sgaStableTarget{sga6.e10.para.7-17-5}7.17.5. The construction of a specialization homomorphism (7.17.4.2) can be generalized to the case in which one no longer assumes that \(S\) is regular at \(s\) and that \(t\) is a maximal point, provided that resolution of singularities is available for the integral closed subscheme \(T\) of \(S\) with generic point \(t\), or merely for its localization at \(s\); this hypothesis holds in particular when \(\underline{O}_{S,s}\) is excellent of characteristic zero, by Hironaka. Indeed, one can then find a morphism \(S'\to \Spec(\underline{O}_{T,s})\) with \(S'\) regular, and, taking a closed point of the fiber \(S'_s\), apply the preceding construction to \(X'=X\times_S S'\) over \(S'\), the generic point \(t'\) of \(S'\) and its closed point \(s'\). This yields a homomorphism of the form (7.17.4.2), since \(k(s')\) is a finite extension of \(k(s)\) and therefore has the same algebraic closure.

\sgaStableTarget{sga6.e10.para.7-17-6}7.17.6. In all the constructions considered in section 7.17, the construction of a specialization homomorphism has been reduced to the
%%%%%%%%%%  scan idx 603  |  folio 590  %%%%%%%%%%
\SGAPrintPage{590}%
case where the base is a trait. It follows that the properties of the specialization homomorphism reviewed in the preceding sections, in the case where the base is a trait, extend automatically to the case of a more general base. Of course, the methods of definition considered apply equally well to \(\Gr_\bullet\), \(K^\bullet\) (when \(f\) is smooth), and \(\Gr^\bullet\) (when, in addition, \(f\) is separated).

\sgaStableTarget{sga6.e10.para.7-17-7}7.17.7. Finally, note that the definition of a homomorphism (7.17.3.1) or (7.17.3.2), or of a homomorphism (7.17.4.2), depends essentially on certain arbitrary choices (a suitable trait over \(S\) in 7.17.3, a sequence of \((S_i)\) in 7.17.4, and, in addition, a resolution of singularities in 7.17.5), and likewise for the variants for \(\Gr_\bullet\), etc. Obvious examples can be constructed with \(s\), say, an ordinary double point of an algebraic curve and \(X\) a projective smooth scheme over \(S\) (for example, an abelian scheme of relative dimension 2), in which the specialization homomorphism on \(\Gr^1=\Pic\) and even on \(\Gr^1_{\mathrm{alg}}=\mathrm{NS}\) (the Néron--Sévéri group) actually depends on the choice of the trait \(S'\) chosen over \(S\) or, what amounts to the same thing here, on the choice of a point over \(s\) in the normalization of \(S\).

\sgaStableTarget{sga6.e10.sec.7.18}\subsection*{7.18. Appendix to the Appendix: extension of the base field for certain functors with values in commutative groups}
We spell out here a technical point raised in 7.17.4.

\sgaStableTarget{sga6.e10.para.7-18-1}7.18.1. Let \(k\) be a field, \(C\) the category of integral \(k\)-algebras and flat homomorphisms of such algebras, and \(F\) a functor on \(C\) with values in the category \((\mathrm{Ens})\) of sets. Suppose that \(F\) commutes with filtered inductive limits. Let \(K\) be an extension of \(k\); we seek
%%%%%%%%%%  scan idx 604  |  folio 591  %%%%%%%%%%
\SGAPrintPage{591}%
criteria for \(F(k)\longrightarrow F(K)\) to be injective (resp. surjective).

\sgaStableTarget{sga6.e10.para.7-18-2}7.18.2. Suppose that we are in one of the following two cases:
\begin{enumerate}[label=\alph*)]
\item The field \(k\) is separably closed, and \(K/k\) is a separable extension.
\item The field \(k\) is infinite, and \(K\) is a pure extension of it.
\end{enumerate}
Then the map \(F(k)\to F(K)\) is injective.

Indeed, using the fact that \(K\) is the filtered inductive limit of subalgebras of finite type over \(k\), with flat transition morphisms, and the commutation of \(F\) with filtered inductive limits, one is reduced in case a) to proving that, if \(A\) is an integral algebra of finite type over \(k\) whose fraction field \(K\) is separable over \(k\), then \(F(k)\to F(A)\) is injective. Now the hypotheses on \(k\) and on \(K\) imply that there exists a \(k\)-homomorphism \(A\to k\); the composite \(F(k)\to F(A)\to F(k)\) is then the identity, whence the conclusion. In case b), passage to the limit similarly reduces the question to proving that, if \(A\) is an integral algebra of finite type over \(k\) whose fraction field is a pure extension of \(k\), then \(F(k)\to F(A)\) is injective. Since \(k\) is infinite, the hypothesis again implies that there exists a \(k\)-homomorphism \(A\to k\), and one concludes as above.

\sgaStableTarget{sga6.e10.para.7-18-3}7.18.3. Now suppose that \(F\) is a functor with values in the category \((\mathrm{Ab})\) of abelian groups, and suppose that for every finite extension \(k'\) of \(k\) a norm homomorphism
\[
  N_{k'/k}:F(k')\longrightarrow F(k)
\]
has been defined such that the composite
\[
  F(k)\longrightarrow F(k')\xrightarrow{N_{k'/k}}F(k)
\]
is multiplication by \([k':k]\). Under these conditions:
\begin{enumerate}[label=(\roman*)]
%%%%%%%%%%  scan idx 605  |  folio 592  %%%%%%%%%%
\item
\SGAPrintPage{592}%
For every extension \(K\) of \(k\), the kernel of \(F(k)\to F(K)\) is a torsion group.
\item For every pure extension \(K\) of \(k\), the homomorphism \(F(k)\to F(K)\) is injective. (NB. Note that \(k\) is not necessarily assumed infinite).
\end{enumerate}

In case (i), one is reduced to proving that every element \(x\) in the kernel of \(F(k)\to F(A)\), where \(A\) is an integral algebra of finite type over \(k\), is a torsion element. Using a \(k\)-homomorphism \(A\to k'\), where \(k'\) is a finite \(k\)-algebra, one is reduced to the case where \(A=k'\) is a finite extension of \(k\). One then has \([k':k]x=0\) by the property postulated above for the norm homomorphism. This proves that \(x\) is a torsion element. One even has \(x=0\) when the greatest common divisor of the degrees of the finite extensions \(k'\) of \(k\) that are quotients of \(A\) equals 1. This is the case when the fraction field \(K\) of \(A\) is a pure extension of \(k\), as is readily checked (only the case where \(k\) is finite requires a new verification). This observation then gives the conclusion of (ii).

\sgaStableTarget{sga6.e10.para.7-18-4}7.18.4. Suppose that, for every nonempty affine open subset \(U\) of the affine space \(\Spec k[T_1,\ldots,T_n]\), with \(A\) the ring of \(U\), the homomorphism \(F(k)\to F(A)\) is surjective (compare IX 1.6). Then, using the commutation of \(F\) with filtered inductive limits, one finds that for every pure extension \(K\) of \(k\), \(F(k)\to F(K)\) is surjective. If one also assumes that \(k\) is infinite, or that norm homomorphisms can be defined as in 7.18.3, one concludes that this homomorphism is even bijective.

\sgaStableTarget{sga6.e10.para.7-18-5}7.18.5. As an example, consider a scheme \(X\) of finite type over \(k\), and set
\[
  F(A)=K_\bullet(X\otimes_k A),
\]
%%%%%%%%%%  scan idx 606  |  folio 593  %%%%%%%%%%
\SGAPrintPage{593}%
which is a functor on the category \(C\) (by (IV 2.12)). By (IV 3.2.3), this functor commutes with filtered inductive limits. Moreover, when \(k'\) is a finite extension of \(k\), the projection \(p:X\otimes_k k'\to X\) is a finite, flat morphism of finite presentation, and therefore defines a direct-image homomorphism \(p=N_{k'/k}:K_\bullet(X\otimes_k k')\longrightarrow K_\bullet(X)\). One immediately verifies that the condition on the norm homomorphisms considered in 7.18.3 does hold. Finally, if \(A\) is the coordinate algebra of an open subset \(U\) of the affine space \(\Spec k[T_1,\ldots,T_n]\), then \(F(k)\to F(A)\) is identified with the map \(K_\bullet(X)\longrightarrow K_\bullet(X\times_k U)\), which is surjective for \(X\) noetherian, by IX 1.6. Using the results proved above, one then obtains:
\begin{enumerate}[label=(\roman*)]
\item If \(K/k\) is a pure extension of \(k\), then \(K_\bullet(X)\longrightarrow K_\bullet(X_K)\) is injective, and is even bijective when \(X\) is noetherian.
\item If \(k\) is separably closed and \(K/k\) is a separable extension, then \(K_\bullet(X)\longrightarrow K_\bullet(X_K)\) is injective.
\item For every extension \(K\) of \(k\), the kernel of \(K_\bullet(X)\longrightarrow K_\bullet(X_K)\) is a torsion group.
\end{enumerate}
These results remain valid when one assumes only that \(X\) is noetherian, provided one restricts to extensions \(K\) of \(k\) that are of finite type (so that \(X_K\) is still noetherian). Moreover, essentially the same results hold for \(K^\bullet\) in place of the functor \(K_\bullet\), with the sole difference that in (i) one cannot assert bijectivity in general, for lack of the homotopy theorem for \(K^\bullet\) in place of \(K_\bullet\); to conclude bijectivity, one would have to assume that \(X\) is noetherian and regular, which implies that the groups \(K^\bullet\) and \(K_\bullet\) are the same for \(X\) and for \(X_K\), respectively.

%%%%%%%%%%  scan idx 607  |  folio 594  %%%%%%%%%%
\SGAPrintPage{594}%
One again obtains essentially the same results for \(\Gr^\bullet\) in place of \(K_\bullet\) or \(K^\bullet\), with the same reservation concerning the bijectivity assertion in (i). In the case where one restricts to \(\Gr^1=\Pic\), this bijectivity assertion is valid when \(X\) is noetherian and \uline{normal}, for in this case, for every open subset \(U\) of an affine space over \(k\), one sees immediately that \(\Pic(X)\longrightarrow \Pic(X\times_k U)\) is surjective (EGA \(\mathrm{Err}_{\mathrm{IV}}\) 53, 21.4.13).

\sgaStableTarget{sga6.e10.bibliography}\subsection*{BIBLIOGRAPHY}
\begin{enumerate}[label={[\arabic*]}]
\item Dickson, E. Linear groups, Dover, New York, 1958.
\item Grothendieck, A. TDTE, V, VI; Séminaire Bourbaki 232, 236.
\item Kleiman, S. Towards a numerical theory of ampleness, Ann. of Math. 84 (1966), pp. 293-344.
\item Murre, J. On contravariant functors from the category of préschemes over a field into the category of abelian groups, Publications Mathématiques de l'IHES, 23, 1964.
\item Chevalley, C. Les classes d'équivalence rationnelles, I et II, Séminaire à l'ENS 1958, exposés 2 et 3.
\item S. Abhyankar, Resolution of singularities of arithmetical surfaces, in: Arithmetical Algebraic Geometry, Purdue University, Harper--Row, New York 1965.
\end{enumerate}

\providecommand{\Isom}{\operatorname{Isom}}
\providecommand{\Hom}{\operatorname{Hom}}
\providecommand{\Coker}{\operatorname{Coker}}
\providecommand{\Nil}{\operatorname{Nil}}
\providecommand{\im}{\operatorname{Im}}
\providecommand{\coker}{\operatorname{coker}}
\providecommand{\Supp}{\operatorname{Supp}}
\providecommand{\fppf}{\mathrm{fppf}}
\providecommand{\zar}{\mathrm{zar}}
\providecommand{\GL}{\operatorname{Gl}}
\providecommand{\Gr}{\operatorname{Gr}}
\providecommand{\Pic}{\operatorname{Pic}}
\providecommand{\Spec}{\operatorname{Spec}}
\providecommand{\V}{\operatorname{V}}
\providecommand{\Ass}{\operatorname{Ass}}

\providecommand{\Hom}{\operatorname{Hom}}
\providecommand{\Isom}{\operatorname{Isom}}
\providecommand{\Coker}{\operatorname{Coker}}
\providecommand{\Ass}{\operatorname{Ass}}
\providecommand{\sgacircled}[1]{\tikz[baseline=(char.base)]\node[shape=circle,draw,inner sep=.35pt](char){\mbox{\scriptsize #1}};}
\newpage
%%%%%%%%%%  scan idx 608  |  folio 595  %%%%%%%%%%
\SGAPrintPage{595}%
\sgaStableTarget{sga6.e12.expose}\section*{\underline{EXPOSÉ XII}\\[3pt] \underline{\underline{A RELATIVE REPRESENTABILITY THEOREM FOR}}\\[3pt] \underline{\underline{THE PICARD FUNCTOR}}}
\begin{center}
\underline{by M. Raynaud (written up by S. Kleiman)}
\end{center}

\sgaStableTarget{sga6.e12.warning}\subsection*{Warning}

The purpose of the two exposés XII, XIII is to give proofs of the nine finiteness theorems for the Picard functor stated by A. Grothendieck in the comments (pp. C 07-C 11) to his ``Fondements de la Géométrie Algébrique'' (cited as FGA). In the first exposé, the first finiteness theorem is proved in the following form, more precise than the one stated in FGA: let \(X\) and \(Y\) be two proper schemes over an integral noetherian base \(S\), and let \(f:X\longrightarrow Y\) be a surjective \(S\)-morphism. Then there is a nonempty open subset \(V\) of \(S\) over which the induced morphism \(f^*:\Pic_{Y/S}\longrightarrow \Pic_{X/S}\) is representable by a quasi-affine morphism of finite presentation. In the second exposé, the rest is proved.

The two exposés differ notably in their proof techniques. In the first, essential use is made of Oort's dévissage, non-flat descent, and the representation of unramified functors. In the second, one first reduces to the case where the schemes \(X/S\) that occur are projective and flat with geometrically integral fibers, or even, as needed, normal or smooth, always by applying the first finiteness theorem to a suitable morphism \(X'\to X\). One then uses the tool of \((b)\)-sheaves forged by G. Castelnuovo and sharpened and used by T. Matsusaka, D. Mumford, Y. Nakai and S. Kleiman. This tool makes it possible, in particular, to dispense with Chow coordinates, the Frobenius morphism and the Lefschetz theorems used in Grothendieck's initial proofs.

The present exposé XII was presented orally, following Grothendieck, by M. Raynaud, whom the writer wishes to thank for communicating his personal notes.


%%%%%%%%%%  scan idx 609  |  folio 596  %%%%%%%%%%
\SGAPrintPage{596}%
\sgaStableTarget{sga6.e12.sec.1}\section*{1. Statement of the main theorem and applications}

\sgaStableTarget{sga6.e12.thm.1-1}\textbf{Theorem 1.1.} \uline{Let} \(S\) \uline{be a noetherian scheme,} \(X\) \uline{and} \(Y\) \uline{two proper} \(S\)\uline{-schemes, and} \(f:X\to Y\) \uline{a morphism. Suppose that} \(S\) is integral \uline{and} \(f\) surjective. \uline{Then there is a nonempty open subset} \(V\) \uline{of} \(S\) \uline{such that the morphism}
\[
  f^*:\underline{\Pic}_{Y/S}|V\longrightarrow \underline{\Pic}_{X/S}|V
\]
\uline{between the Picard functors (localized for the} f.p.p.f \uline{topology) is representable by a quasi-affine morphism of finite presentation.}

\sgaStableTarget{sga6.e12.cor.1-2}\textbf{Corollary 1.2.} \uline{Let} \(X\) \uline{be a proper scheme over a noetherian scheme} \(S\)\uline{. Suppose that} \(S\) is integral. \uline{Then there is a nonempty open subset} \(V\) \uline{of} \(S\) \uline{such that the Picard functor} \(\Pic_{X/S}|V\) \uline{is representable by a} \(V\)\uline{-group scheme that is a sum of quasi-projective open subschemes of finite type, hence a fortiori separated over} \(V\).

Indeed, recall the following fact (FGA N\textsuperscript{o}232): let \(S\) be a scheme, and let \(X\) be a projective \(S\)-scheme, of finite presentation and integral over \(S\) (that is, flat with geometrically integral fibers). Then the Picard functor is represented by an \(S\)-group scheme \(\underline{\Pic}_{X/S}\), an increasing union of quasi-projective open subschemes of finite presentation.

Consideration of the ``Hilbert polynomials'' makes it possible here to decompose \(\underline{\Pic}_{X/S}\) as a sum of open subschemes \(\underline{\Pic}^q_{X/S}\), \(q\in\mathbb Q[t]\). In the next exposé it is proved, under the additional hypothesis that \(S\) is noetherian, that the \(\underline{\Pic}^q_{X/S}\) are of finite type over \(S\). It is then known that 1.2 holds in the case where \(X\) is in addition projective and integral over \(S\); we shall deduce the general case using the

\sgaStableTarget{sga6.e12.lem.1-3}\textbf{Lemma 1.3.} \uline{Let} \(X\) \uline{be of finite type over an integral noetherian} \(S\)\uline{. Then there are a nonempty open subset} \(V\) \uline{of} \(S\)\uline{, an integral scheme} \(S'\) \uline{and a finite and}
%%%%%%%%%%  scan idx 610  |  folio 597  %%%%%%%%%%
\SGAPrintPage{597}%
\uline{flat morphism} \(S'\to V\) \uline{such that every irreducible component} \(X'_i\) \uline{of} \(X_{S'}\)\uline{, endowed with the induced reduced structure, is integral over} \(S'\).

Indeed, there is a finite extension \(L\) of the function field \(k(S)\) such that every irreducible component of \((X_L)_{\text{red}}\) is geometrically integral (by EGA IV 4.6.8). Let \(S_1\) be an integral \(S\)-scheme such that \(k(S_1)=L\) and such that there are a nonempty open subset \(V_1\) of \(S\) and a finite flat morphism \(h:S_1\to V_1\). There is a nonempty open subset \(S_2\) of \(S_1\) such that every irreducible component of \(X_{S_2}\), endowed with the induced reduced structure, is integral over \(S_2\) (EGA IV 9.7.7). Take \(V=V_1-h(S_1-S_2)\) and \(S'=h^{-1}(V)\).

Let us now prove 1.2. Let \(h:S'\to V_1\) and \(X'_i\) be as in 1.3. For every \(i\), let \(f_i:Z_i\to X'_i\) be a Chow presentation with \(Z_i\) integral and \(f_i\) birational. After shrinking \(V\) and \(S'\), we may suppose that \(Z_i\) is integral over \(S'\) (EGA IV 9.7.7). Let \(f:Z=\coprod_i Z_i\longrightarrow X'=X_{S'}\) be the canonical morphism. Then \(\underline{\Pic}_{Z/S'}=\prod_i\underline{\Pic}_{Z_i/S'}\) is representable by a group scheme that is a sum of quasi-projective open subschemes of finite type.

Finally, 1.1 says that, after shrinking \(V\) and \(S'\), the morphism \(f^*:\underline{\Pic}_{X'/S'}\longrightarrow \underline{\Pic}_{Z/S'}\) is representable by a quasi-affine morphism. Consequently, \(\Pic_{X'/S'}\) is representable by a group scheme that is a sum of quasi-projective open subschemes of finite type; by finite flat descent (SGA 1 VIII 7.6 and EGA II 6.6.1), the same holds for \(\Pic_{X/S}|V\).

\sgaStableTarget{sga6.e12.lem.1-4}\textbf{Lemma 1.4.} \uline{Let} \(k\) \uline{be a field,} \(P\) \uline{and} \(Q\) \uline{two} \(k\)\uline{-group schemes locally of finite type, and let} \(f:P\to Q\) \uline{be a quasi-affine morphism. Then} \(f\) \uline{is affine.}

Indeed, let \(R=\Ker(f)\); \(R\) is a \(k\)-group scheme of finite type that is quasi-affine; by (SGA 3 \(\mathrm{VI}_{\mathrm B}\) 11.11), \(R\) is therefore affine.
%%%%%%%%%%  scan idx 611  |  folio 598  %%%%%%%%%%
\SGAPrintPage{598}%
Thus \(P\to P/R\), which makes \(P\) an \(R\)-torsor over \(P/R\), is affine. Moreover, it is known (SGA 3 1.4.2) that \(P/R\to Q\) is a closed immersion, hence affine, and hence the same holds for the composite \(P\to P/R\to Q\).

\sgaStableTarget{sga6.e12.cor.1-5}\textbf{Corollary 1.5.} \uline{Let} \(k\) \uline{be a field.}
\begin{enumerate}[label=\alph*)]
\item \uline{Let} \(X\) \uline{be a proper} \(k\)\uline{-scheme. Then} \(\underline{\Pic}_{X/k}\) \uline{is representable} (J. P. Murre).
\item \uline{Let} \(f:X\to Y\) \uline{be a surjective morphism between two proper} \(k\)\uline{-schemes. Then the induced morphism} \(f^*:\underline{\Pic}_{Y/k}\longrightarrow \underline{\Pic}_{X/k}\) \uline{is affine.}
\end{enumerate}

\sgaStableTarget{sga6.e12.cor.1-6}\textbf{Corollary 1.6.} \uline{Let} \(S\) \uline{be an integral noetherian scheme, and} \(f:X\to Y\) \uline{a surjective} \(S\)\uline{-morphism between two proper} \(S\)\uline{-schemes.}
\begin{enumerate}[label=\alph*)]
\item \uline{If} \(Y\) \uline{is normal over} \(S\) \uline{(that is, flat with geometrically normal fibers), then, after shrinking} \(S\), \(f^*:\underline{\Pic}_{Y/S}\to\underline{\Pic}_{X/S}\) \uline{is quasi-finite.}
\item \uline{If} \(\underline{\Pic}_{Y/S}\) \uline{is essentially proper over} \(S\) \uline{(that is, if it satisfies the valuative criterion of properness, for example, if} \(Y\) \uline{is smooth over} \(S\)\uline{), then, after shrinking} \(S\), \(f^*:\underline{\Pic}_{Y/S}\to\underline{\Pic}_{X/S}\) \uline{is finite (and in particular affine).}
\end{enumerate}

Indeed, in a), for every \(s\in S\), \(R_s=\Ker(f_s^*)\) is an affine algebraic group; its identity component is a proper algebraic group over \(k\) (FGA 236, Theorem 2.1 (iii)); \(R_s\) is therefore finite. In b), \(f^*\) is both proper and quasi-affine, hence finite.

\sgaStableTarget{sga6.e12.rem.1-7}\textbf{Remark 1.7.} It is not known whether 1.1 remains valid when ``quasi-affine'' is replaced there by affine, as 1.5 and 1.6 might suggest.


%%%%%%%%%%  scan idx 612  |  folio 599  %%%%%%%%%%
\SGAPrintPage{599}%
\sgaStableTarget{sga6.e12.sec.2}\section*{2. First reductions}

\sgaStableTarget{sga6.e12.para.2-1}2.1. Apply 1.3 to \(X/S\) and \(Y/S\): we obtain a nonempty open subset \(V\) of \(S\), an integral scheme \(S'\) and a finite flat morphism \(S'\to V\) such that every irreducible component \(X'_i\) of \(X_{S'}\) and \(Y'_j\) of \(Y_{S'}\), endowed with the induced reduced structure, is integral over \(S'\). By f.p.q.c. descent (EGA IV 2.7.1 (xiv) and SGA 1 VIII 7.9), we may replace \(S\) by \(S'\). It then follows from the following commutative diagram:
\[
\begin{tikzcd}
X_{\red} \arrow[r,"f_{\red}"] \arrow[d,"i"] & Y_{\red} \arrow[d,"j"]\\
X \arrow[r,"f"] & Y & ,
\end{tikzcd}
\]
implying the commutative diagram
\[
\begin{tikzcd}
\underline{\Pic}_{X_{\red}/S} & \underline{\Pic}_{Y_{\red}/S} \arrow[l]\\
\underline{\Pic}_{X/S} \arrow[u] & \underline{\Pic}_{Y/S} \arrow[l] \arrow[u] & ,
\end{tikzcd}
\]
and from trivial sorites on morphisms of functors representable by quasi-affine morphisms (following from EGA II 5.1.10 (ii) (v)) that it suffices to prove the following two assertions:
\begin{enumerate}[label=\Roman*.]
\item \uline{Theorem} 1.1 \uline{when} \(f\) \uline{is an immersion (necessarily nilpotent).}
\item \uline{Theorem} 1.1 \uline{when} \(X\) \uline{and} \(Y\) \uline{are reduced and every irreducible component} \(X_i\) \uline{of} \(X\) \uline{and} \(Y_i\) \uline{of} \(Y\)\uline{, endowed with the induced reduced structure, is integral over} \(S\).
\end{enumerate}

\sgaStableTarget{sga6.e12.para.2-2}2.2. Let us study assertion II. For every \(i\), let \(j=\varphi(i)\) be an index such that \(f(X_i)\subset Y_j\), and put \(\Sigma_j=\coprod_{\varphi(i)=j} X_i\). We then have the two diagrams
%%%%%%%%%%  scan idx 613  |  folio 600  %%%%%%%%%%
\SGAPrintPage{600}%
\[
\begin{tikzcd}
Y & \coprod_j Y_j \arrow[l]\\
X \arrow[u] & \coprod_j \Sigma_j \arrow[l] \arrow[u] & ,
\end{tikzcd}
\]
\[
\begin{tikzcd}
\underline{\Pic}_{Y/S} \arrow[r,"\sgacircled{1}"] \arrow[d,"\sgacircled{2}"'] &
\underline{\Pic}_{\coprod_jY_j/S}=\prod_j\underline{\Pic}_{Y_j/S} \arrow[d,"\sgacircled{4}"]\\
\underline{\Pic}_{X/S} \arrow[r,"\sgacircled{3}"] &
\underline{\Pic}_{\coprod_j\Sigma_j/S}=\prod_j\underline{\Pic}_{\Sigma_j/S}.
\end{tikzcd}
\]
And to establish II, it suffices to see that, after shrinking \(S\), \(\sgacircled{1}\), \(\sgacircled{3}\) and \(\sgacircled{4}\) are quasi-affine.

Consider \(\sgacircled{4}\) first. Replacing \(Y\) by \(Y_j\) and \(X\) by \(\Sigma_j\) reduces us to the case where \(Y\) is integral over \(S\). Since \(f\) is surjective, one of the \(X_i\), say \(X_0\), dominates \(Y\), and we have the two sequences
\[
  X_0\hookrightarrow \coprod X_i\longrightarrow Y,
\]
\[
  \Pic_{X_0/S}\longleftarrow \Pic_{\coprod X_i/S}\longleftarrow \Pic_{Y/S}.
\]
Thus, knowing that the \(\underline{\Pic}_{X_i/S}\), and therefore also \(\Pic_{\coprod X_i/S}\), are representable and that the composite morphism is quasi-affine, we conclude that \(\sgacircled{4}\) is quasi-affine. In view of the proof of 1.2, for \(\sgacircled{4}\) we are therefore reduced to proving II'. \uline{Theorem} 1.1 \uline{when} \(X\) \uline{and} \(Y\) \uline{are integral over} \(S\).

For \(\sgacircled{1}\) and \(\sgacircled{3}\), we must treat the case of a morphism \(\coprod X_i\to X\) where \(X\) is reduced, the \(X_i\) are the irreducible components of \(X\) and are integral over \(S\). Now we may suppose that \(X\) has connected fibers. Indeed, let \(\eta\) be the geometric point of \(S\). Then the irreducible components of \(X_\eta\) are the \((X_i)_\eta\) and, after shrinking \(S\), we may suppose that \(X_i\cap X_j=\varnothing\) if and only if
%%%%%%%%%%  scan idx 614  |  folio 601  %%%%%%%%%%
\SGAPrintPage{601}%
\((X_i)_\eta\cap(X_j)_\eta=\varnothing\). Let \(Y_k=\bigcup_{i\in I(k)}X_i\) be the connected components of \(X\); then \(f\) decomposes, \(f=\coprod_k f_k:\coprod_k(\coprod_{i\in I(k)}X_i)\longrightarrow \coprod_kY_k\), and the assertion follows. Finally, an immediate induction reduces us to the case of a morphism \(X_1\amalg X_2\to X\). After shrinking \(S\), \(\sgacircled{1}\) and \(\sgacircled{3}\) then follow, by (EGA III 7.8.6), from the

\sgaStableTarget{sga6.e12.prop.2-3}\textbf{Proposition 2.3.} \uline{Let} \(S\) \uline{be a scheme,} \(g:Y\to S\) \uline{a proper morphism, and} \(Y_1,Y_2\) \uline{two closed subschemes defined by the ideals} \(I_1,I_2\)\uline{. Suppose that}
\begin{enumerate}[label=\alph*)]
\item \(Y,Y_1,Y_2,Y_1\cap Y_2\) \uline{are faithfully flat and of finite presentation over} \(S\).
\item \(g_*{O}_Y={O}_S\) and \((g_i)_*{O}_{Y_i}={O}_S\) \uline{universally (where} \(g_i=g|Y_i\), \(i=1,2\)).
\item \(I_1\cap I_2=0\).
\end{enumerate}
\uline{Then} \(\underline{\Pic}_{Y/S}\longrightarrow \underline{\Pic}_{(Y_1\amalg Y_2)/S}\) \uline{is representable by an affine morphism of finite presentation.}

Indeed, first note that the hypotheses are stable under every base change. For (c), this follows from the two exact sequences:
\[
  0\longrightarrow I_1+I_2\longrightarrow{O}_Y
  \longrightarrow{O}_{Y_1\cap Y_2}\longrightarrow0
\]
and
\[
  0\longrightarrow I_1\cap I_2\longrightarrow I_1\oplus I_2
  \longrightarrow I_1+I_2\longrightarrow0~;
\]
the first shows that \(I_1+I_2\) is flat over \(S\); then the second shows that \(I_1\cap I_2\) is stable under base change.

The assertion is local for the f.p.p.f. topology. Making the base change \(Y_1\cap Y_2\to S\), we may therefore suppose that \(Y_1\cap Y_2/S\) has a section \(\alpha\), which defines sections \(\alpha_i\) of \(Y_i\) and \(\beta\) of \(Y\).

After a suitable base change, we are reduced to seeing that every point \(\mu:S\to\underline{\Pic}_{(Y_1\amalg Y_2)/S}\) has as its inverse image a subfunctor of
%%%%%%%%%%  scan idx 615  |  folio 602  %%%%%%%%%%
\SGAPrintPage{602}%
\(\underline{\Pic}_{Y/S}\) that is representable by an affine scheme over \(S\). Now \(\mu\) is defined by an invertible sheaf \(L_i\) on \(Y_i\) equipped with a rigidification \(u_i:\alpha_i^*L_i\xrightarrow{\sim}{O}_S\) for \(i=1,2\). Likewise, a point \(\lambda:S\to\underline{\Pic}_{Y/S}\) corresponds to an invertible sheaf \(L\) on \(Y\) equipped with a rigidification \(u:\alpha^*L\xrightarrow{\sim}{O}_S\). The image of \(\lambda\) in \(\Pic_{Y_1\amalg Y_2/S}\) is \(\mu\) if and only if there is an isomorphism of rigidified sheaves \(f_i^*(L)\simeq L_i\) where \(f_i:Y_i\hookrightarrow Y\) is the inclusion for \(i=1,2\). Observe that these isomorphisms define an isomorphism \(\varphi:L_1|Y_1\cap Y_2\xrightarrow{\sim}L_2|Y_1\cap Y_2\) such that \(u_1=u_2\circ\varphi\).

Conversely, given an isomorphism \(\varphi\), \(L_1|Y_1\cap Y_2\to L_2|Y_1\cap Y_2\), such that \(u_1=u_2\circ\varphi\), there is a rigidified invertible sheaf \(L\) on \(Y\) obtained by gluing along \(\varphi\); explicitly, \(L\) is the kernel of the canonical homomorphism \((f_1)_*L_1\oplus(f_2)_*L_2\longrightarrow(f|Y_1\cap Y_2)_*L_2\) defined by \(\varphi\). This \(L\) then corresponds to a point \(\lambda\) of \(\underline{\Pic}_{Y/S}\) mapping to \(\mu\).

Finally, let \(I\) be the functor defined, as \(T\) runs through the schemes over \(S\), by the sets of isomorphisms \(\varphi:(L_1)_T|(Y_1\cap Y_2)_T\xrightarrow{\sim}(L_2)_T|(Y_1\cap Y_2)_T\) such that \((u_1)_T=(u_2)_T\circ\varphi\). We have just established an isomorphism between the inverse image of \(\mu\) in \(\underline{\Pic}_{Y/S}\) and the functor \(I\). It therefore suffices to apply the following lemma (taking \(X=Y_1\cap Y_2\)):

\sgaStableTarget{sga6.e12.lem.2-4}\textbf{Lemma 2.4.} \uline{Let} \(X\) \uline{be a flat, proper scheme of finite presentation over} \(S\)\uline{. Let} \(L,M\) \uline{be two invertible sheaves on} \(X\), \uline{and} \(\underline{\Isom}(L,M)\) \uline{the functor of isomorphisms} \(L_T\to M_T\) \uline{(for} \(T/S\)\uline{). Then}
\begin{enumerate}[label=\roman*)]
\item \(\underline{\Isom}(L,M)\) \uline{is representable by an affine scheme over} \(S\) \uline{of finite presentation.}
\item \uline{If, moreover,} \(L\) \uline{and} \(M\) \uline{are equipped with rigidifications (relative to a}
%%%%%%%%%%  scan idx 616  |  folio 603  %%%%%%%%%%
\SGAPrintPage{603}%
\uline{given section of} \(X\)\uline{), the subfunctor of} \(\underline{\Isom}(L,M)\) \uline{formed by the isomorphisms compatible with these rigidifications is representable by a closed subscheme of} \(\underline{\Isom}(L,M)\).
\end{enumerate}

Indeed, ii) is clear. As for i), recall that the functors \(T\mapsto\Hom(L_T,M_T)\) and \(T\mapsto\Hom(M_T,L_T)\) are representable by vector bundles \(\mathbb V(E)\) and \(\mathbb V(F)\) (by EGA III 7.7.8). Then, for \(u\in\Hom(L_T,M_T)\) and \(v\in\Hom(M_T,L_T)\), the conditions \(u\circ v=1\) and \(v\circ u=1\) define a closed subscheme \(F\) of \(\mathbb V(E)\times_S\mathbb V(F)\), and it is clear that \(F\simeq I\).

\sgaStableTarget{sga6.e12.para.2-5}2.5. Finally, let us reduce assertion II' to the following assertion:

\(\mathrm{II}^*\). \uline{Theorem} 1.1 \uline{when} \(X\) \uline{and} \(Y\) \uline{are integral over} \(S\) \uline{and the sequence}
\sgaStableTarget{sga6.e12.eq.unnamed}\begin{equation}\tag{*}
  {O}_Y \longrightarrow f_*{O}_X \rightrightarrows g_*\bigl({O}_{X\times_YX}\bigr)
\end{equation}
\uline{(where} \(g:X\times_YX\to Y\) \uline{is the structural morphism) is exact.}

In fact, for this it suffices to prove that there is a factorization of \(f\)
\[
  X\xrightarrow{f_0}Y_0\xrightarrow{f_1}Y_1\longrightarrow\cdots\xrightarrow{f_\eta}Y_\eta=Y
\]
such that every \(f_i\) satisfies (*) and, after shrinking \(S\), every \(Y_i\) is integral over \(S\).

For this, put \(Y_0=\Spec(f_*(\underline{O}_X))\), with \(f_0\) arising from the Stein factorization of \(f\) (EGA III 4.3.3), so that \(f_{0*}(\underline{O}_X)\xleftarrow{\sim}\underline{O}_{Y_0}\), and hence \(f_0\) does satisfy (*). To factor \(Y_0\to Y\) in the same way, we use the

\sgaStableTarget{sga6.e12.lem.2-6}\textbf{Lemma 2.6.} \uline{Let} \(Y\) \uline{be a noetherian scheme, and} \(f:X\to Y\) \uline{a finite epimorphism (i.e. such that} \({O}_Y\to f_*({O}_X)\) \uline{is injective); then} \(f\) \uline{factors into a finite sequence of finite morphisms} \(f_i\) \(X\simeq Y_0\xrightarrow{f_1}Y_1\longrightarrow\cdots\longrightarrow Y_n=Y\) \uline{that are effective epimorphisms (i.e. such that the sequence}
%%%%%%%%%%  scan idx 617  |  folio 604  %%%%%%%%%%
\SGAPrintPage{604}%
\({O}_{Y_i}\longrightarrow (f_i)_*({O}_{Y_{i-1}})\rightrightarrows (g_i)_*\bigl({O}_{Y_{i-1}}\otimes_{{O}_{Y_i}}{O}_{Y_{i-1}}\bigr)\) \uline{is exact).}

Put \(A={O}_Y\), \(B_0=f_*({O}_X)\), and define \(B_i=\Ker\bigl(B_{i-1}\rightrightarrows B_{i-1}\otimes_A B_{i-1}\bigr)\) recursively. We verify at once that \(B_{i-1}\otimes_A B_{i-1}\xrightarrow{\sim}B_{i-1}\otimes_{B_i}B_{i-1}\). It remains to prove that, for \(i\) sufficiently large, \(B_i=A\). Indeed, consider the decreasing sequence of \(A\)-modules \(M_i=B_i/A\), and put \(S_i=\Supp(M_i)\). The decreasing sequence of closed subsets \(S_i\) of \(Y\) is stationary; let \(S\) be the stationary value. Let us show that \(S=\varnothing\). Otherwise, let \(s\) be a maximal point of \(S\). After replacing \(Y\) by \(\Spec{O}_{Y,s}\), we may suppose that \(Y\) is local with closed point \(s\). By the choice of \(s\), \(M_i\) is artinian for \(i\) sufficiently large, so the sequence \((M_i)\) is stationary. Thus there exists \(i\) such that \(B_i=B_{i+1}\). Consequently, \(\Spec(B_i)\to Y\) is a monomorphism, hence a closed immersion (Nakayama); since \(A\subset B_i\), we have \(A=B_i\) and \(s\notin S\), a contradiction.

\sgaStableTarget{sga6.e12.sec.3}\section*{3. Proof of I: Oort's dévissage}

\sgaStableTarget{sga6.e12.prop.3-1}\textbf{Proposition 3.1.} (Oort, Bull. soc. math. France, vol. 90-1962). \uline{Let} \(S\) \uline{be a locally noetherian scheme,} \(f:X\to S\) \uline{a proper morphism, and} \(X'\subset X\) \uline{a closed subscheme. Suppose that:}
\begin{enumerate}[label=\roman*)]
\item \(X'\) \uline{is defined by an ideal} \(I\) \uline{contained in the nilradical} \(N\) \uline{of} \(X\) \uline{and} \(N.I=0\).
\item \(X'\) \uline{is flat over} \(S\) \uline{and cohomologically flat over} \(S\) \uline{in dimension} \(0\) (EGA III 7.8.1).
\item \(I\) \uline{is flat over} \(S\) \uline{and cohomologically flat over} \(S\) \uline{in dimensions} \(0\) \uline{and} \(1\)\uline{, so that (by} EGA III 7.4.6\uline{) there are a coherent} \({O}_S\)\uline{-module} \(Q\) \uline{and an isomorphism, functorial in the coherent} \({O}_S\)\uline{-module} \(M\),
%%%%%%%%%%  scan idx 618  |  folio 605  %%%%%%%%%%
\SGAPrintPage{605}%
\[
  R^2f_*(I\otimes f^*M)\simeq \underline{\Hom}_{{O}_S}(Q,M).
\]
\item \uline{The cokernel} \(E\) \uline{of} \(f_*({O}_{X'})\to R^1f_*(I)\) \uline{is a locally free} \({O}_S\)\uline{-module.}
\item \uline{The homomorphism} \((f_*{O}_X)_{\red}\longrightarrow (f_*{O}_{X'})_{\red}\) \uline{is surjective.}
\end{enumerate}
\uline{Then the canonical morphism} \(\underline{\Pic}_{X/S}\longrightarrow \underline{\Pic}_{X'/S}\) \uline{is representable by an affine morphism of finite presentation; more precisely, there is a canonical exact sequence of} f.p.p.f. \uline{sheaves,}
\[
  0\longrightarrow \mathbb V(E^\vee)\longrightarrow \underline{\Pic}_{X/S}
  \longrightarrow\underline{\Pic}_{X'/S}\xrightarrow{u}\mathbb V(Q)
\]
\uline{and, for} \(P=\Ker(u)\)\uline{, the monomorphism} \(P\hookrightarrow\Pic_{X'/S}\) \uline{is representable by a closed immersion and} \(\Pic_{X/S}\) \uline{is a torsor over} \(P\) \uline{under the group} \(V(\check E)\) \uline{(that is, the sequence of} f.p.p.f. \uline{sheaves}
\[
  0\longrightarrow \mathbb V(\check E)\longrightarrow \Pic_{X/S}\longrightarrow P\longrightarrow0
\]
\uline{is exact).}

For an \(S\)-scheme \(T\), consider the following two diagrams with exact rows:
\[
\begin{tikzcd}
0 \arrow[r] & I_T \arrow[r] & {O}_{X_T} \arrow[r,dash] & {O}_{X'_T} \arrow[r] & 0\\
0 \arrow[r] & I_T \arrow[r,"\alpha",dash] \arrow[u,equal] & {O}_{X_T}^{*} \arrow[r] \arrow[u,hook] & {O}_{X'_T}^{*} \arrow[r] \arrow[u,hook] & 0
\end{tikzcd}
\qquad (\alpha(i)=1+i)
\]
\[
\begin{tikzcd}[column sep=tiny]
(f_T)_*({O}_{X'_T}) \arrow[r,"\delta_T"] &
R^1(f_T)_*(I_T) \arrow[d,equal]\\
(f_T)_*({O}_{X'_T}^{*}) \arrow[r,"\delta_T^*"] \arrow[u] &
R^1(f_T)_*(I_T) \arrow[r] &
R^1(f_T)_*({O}_{X_T}^{*}) \arrow[r] &
R^1(f_T)_*({O}_{X'_T}^{*}) \arrow[r] &
R^2(f_T)_*(I_T) .
\end{tikzcd}
\]

Observe carefully that the square on the left of this diagram \uline{is not necessarily commutative}. Nevertheless, we have:
\sgaStableTarget{sga6.e12.para.3-2}\[
  \textnormal{3.2.}\qquad \im\delta_T^*=\im\delta_T.
\]


%%%%%%%%%%  scan idx 619  |  folio 606  %%%%%%%%%%
\SGAPrintPage{606}%
Indeed, we may clearly suppose that \(S\) and \(T\) are affine. Take an affine covering \(\{V_i\}\) of \(X_T\). Let \(g\in\Gamma(X_T;{O}_{X'_T}^{*})\), and let \(g_i\) be a lift of \(g|V_i\) in \(\Gamma(V_i;{O}_{X_T}^{*})\). Then \(\delta_T^*(g)\) is the class of the \(1\)-cocycle \(\delta_{ij}^*=(g_ig_j^{-1}-1)\); \(\delta_T(g)\) is the class of \(\delta_{ij}=(g_i-g_j)=g_j\delta_{ij}^*=g\,\delta_{ij}^*\), so we have
\[
  \delta_T(g)=g\,\delta_T^*(g)
\]
taking into account the action of \(\Gamma(X_T,\underline{O}_{X'_T})\) on \(I_T\), hence on \(H^1(X_T,I_T)\).

By ii), we have \(\Gamma(X_T;{O}_{X'_T})=\Gamma(X;{O}_{X'})\otimes_S\Gamma(T,{O}_T)\), and hence \(g=\sum g_\beta\otimes t_\beta\) for \(g_\beta\in\Gamma(X,{O}_{X'})\) and \(t_\beta\in\Gamma(T,{O}_T)\). Now, by (v), \((f_*{O}_X)_{\red}\longrightarrow(f_*{O}_{X'})_{\red}\) is surjective; consequently \(g_\beta=v(h_\beta)+n_\beta\) for \(h_\beta\in\Gamma(X;{O}_X)\), \(v:\Gamma(X,{O}_X)\to
\Gamma(X,{O}_{X'})\) and \(n_\beta\in\operatorname{nilrad}(\Gamma(X,{O}_{X'}))\).

By i), \((n_\beta\otimes t_\beta)I_T=0\). Consequently, \(\delta_T(g)=h\,\delta_T^*(g)\), where \(h=\sum h_\beta\otimes t_\beta\in\Gamma(X_T,\underline{O}_{X_T})\), whence \(\delta_T^*(g)=h^{-1}\delta_T(g)=\delta_T(h^{-1}g)\) and hence \(\im\delta_T^*\subset\im\delta_T\).

Finally, let \(g=\sum g_\beta\otimes t_\beta\in
\Gamma(X_T,{O}_{X'_T})\). Let \(g_\beta=v(h_\beta)+n_\beta\), where \(n_\beta\in\operatorname{nilrad}(\Gamma(X,{O}_{X'}))\); we then have \(\delta_T(g)=\delta_T\bigl(\sum n_\beta\otimes t_\beta\bigr)=\delta_T\bigl(1+\sum n_\beta\otimes t_\beta\bigr)=\delta_T^*\bigl(1+\sum n_\beta\otimes t_\beta\bigr)\), whence \(\im\delta_T\subset\im\delta_T^*\).

By 3.2, iii) and iv), there is an exact sequence functorial in \(T\):
\[
  0\longrightarrow \mathbb V(E^\vee)(T)\longrightarrow
  R^1(f_T)_*({O}_{X_T}^{*})
  \longrightarrow
  R^1(f_T)_*({O}_{X'_T}^{*}) \longrightarrow \mathbb V(Q)(T),
\]
which, on passing to the associated f.p.p.f. sheaves, gives an exact sequence
\[
  0\longrightarrow \mathbb V(E^\vee)\longrightarrow\underline{\Pic}_{X/S}
  \longrightarrow\underline{\Pic}_{X'/S}\xrightarrow{u}\mathbb V(Q).
\]

For \(P=\Ker(u)\), the monomorphism \(P\hookrightarrow\underline{\Pic}_{X'/S}\) is obtained by base change from the unit section \(S\hookrightarrow\mathbb V(Q)\); since \(\mathbb V(Q)\) is representable and separated, it follows that \(P\hookrightarrow\Pic_{X'/S}\) is representable by a
%%%%%%%%%%  scan idx 620  |  folio 607  %%%%%%%%%%
\SGAPrintPage{607}%
closed immersion.

Finally, \(0\longrightarrow \mathbb V(E^\vee)\longrightarrow\underline{\Pic}_{X/S}\xrightarrow{w}P\longrightarrow0\) is a torsor for f.p.p.f.; a fortiori, by f.p.q.c. descent from a ``splitting'' of the torsor, \(w\) is representable by a faithfully flat affine morphism of finite presentation (see: Oort, ``Commutative group schemes'', Proposition (17.4)).

\sgaStableTarget{sga6.e12.cor.3-3}\textbf{Corollary 3.3.} \uline{Under the conditions of} 3.1\uline{, we have}
\begin{enumerate}[label=\roman*)]
\item \uline{If} \(\underline{\Pic}_{(\zar)(X'/S)}\to\underline{\Pic}_{(\fppf)(X'/S)}\) \uline{is an isomorphism, then} \(\Pic_{(\zar)(X/S)}\to\Pic_{(\fppf)(X/S)}\) \uline{is an isomorphism.}
\item \uline{Supposing that the} \(S\)\uline{-scheme} \(T\) \uline{is affine, if} \(\Pic(X'_T)\to\underline{\Pic}_{(\zar)(X'/S)}^{(T)}\) \uline{is surjective, then} \(\Pic(X_T)\to\underline{\Pic}_{(\zar)(X/S)}^{(T)}\) \uline{is surjective.}
\end{enumerate}

Indeed, (i) follows from the following diagram of sheaves for the Zariski topology
\[
\begin{tikzcd}[column sep=small]
0 \arrow[r] \arrow[d,"S"] & \mathbb V(E^\vee) \arrow[r] \arrow[d,"S"] &
  \underline{\Pic}_{(\zar)(X/S)} \arrow[r] \arrow[d] &
  \underline{\Pic}_{(\zar)(X'/S)} \arrow[r] \arrow[d,"S"] &
  \mathbb V(Q) \arrow[d,"S"]\\
0 \arrow[r] & \mathbb V(E^\vee) \arrow[r,dash] &
\underline{\Pic}_{(\fppf)(X/S)} \arrow[r] &
\underline{\Pic}_{(\fppf)(X'/S)} \arrow[r] &
\mathbb V(Q),
\end{tikzcd}
\]
by applying the five lemma to it.

(ii) follows from the following diagram with exact rows
\[
\begin{tikzcd}[column sep=small]
H^1(X_T,I_T) \arrow[r] \arrow[d] &
H^1(X_T,{O}_{X_T}^{*}) \arrow[r,dash] \arrow[d] &
H^1(X_T,{O}_{X'_T}^{*}) \arrow[r,dash] \arrow[d] &
H^2(X_T,I_T) \arrow[d]\\
H^0(T,R^1(f_T)_*(I_T)) \arrow[r] &
H^0(T,R^1(f_T)_*{O}_{X_T}^{*}) \arrow[r] &
H^0(T,R^1(f_T)_*{O}_{X'_T}^{*}) \arrow[r] &
H^0(T,R^2f_*I_T)
\end{tikzcd}
\]
because the outer vertical arrows are isomorphisms, \(T\) being affine.


%%%%%%%%%%  scan idx 621  |  folio 608  %%%%%%%%%%
\SGAPrintPage{608}%
\sgaStableTarget{sga6.e12.cor.3-4}\textbf{Corollary 3.4.} \uline{Let} \(k\) \uline{be a field, and} \(X\) \uline{a proper} \(k\)\uline{-scheme. Suppose that} \(\underline{\Pic}_{X_{\red}/k}\) \uline{possesses a universal invertible sheaf. (For example, by descent theory this holds when every connected component} \(X'_i\) \uline{of} \(X_{\red}\) \uline{has a rational point, for example when} \(k\) \uline{is algebraically closed, or more generally when the} g.c.d. \uline{of the degrees of the} \(0\)\uline{-cycles on every} \(X'_i\) \uline{is} \(1\)\uline{). Then, over every affine open subset} \(V\) \uline{of} \(\Pic_{X/k}\)\uline{, there is a universal invertible sheaf on} \(X\times_kV\).

Finally, we prove I in the following more precise form:

\sgaStableTarget{sga6.e12.prop.3-5}\textbf{Proposition 3.5.} \uline{Let} \(S\) \uline{be an integral noetherian scheme, and} \(f:X\hookrightarrow Y\) \uline{a nilpotent immersion of proper} \(S\)\uline{-schemes. Then there is a nonempty open subset} \(V\) \uline{of} \(S\) \uline{such that the morphism}
\[
  f^*:\Pic_{Y/S}|V\longrightarrow\Pic_{X/S}|V
\]
\uline{is representable by an affine morphism of finite presentation; in fact,} \(N=\Ker f^*\) \uline{is an affine group scheme with unipotent fibers (in the sense of} SGA 3 XVII).

Indeed, \(X\) is defined by a nilpotent ideal \(I\). Then \(I\) has a filtration by ideals \(I_0=I\supset I_1\supset\ldots\supset I_n=0\) such that \(NI_{i-1}\subset I_i\), where \(N\) is the nilradical of \({O}_Y\) (for example \(I_k=N^kI\)). Let \(X_i\) be the subscheme of \(Y\) defined by \(I_i\). By induction it suffices to study the immersions \(X_{i-1}\hookrightarrow X_i\). We may therefore suppose that \(N.I=0\); this is condition i) of 3.1, and we finish by reducing to the case where the other conditions (ii) to (v) also hold. By the generic flatness theorem (EGA IV 6.9.1), applied to \(\underline{O}_{X'}\), to the \(R^if_*(I)\) and to \(E\), and taking account of the Künneth relations EGA III 6.9.9, we may suppose, after shrinking \(S\), that conditions ii), iii) and iv) hold.


%%%%%%%%%%  scan idx 622  |  folio 609  %%%%%%%%%%
\SGAPrintPage{609}%
There is a finite extension \(K\) of \(k(S)\) such that \((f_K)_*({O}_{X'_K})\) has all its residue field extensions trivial. For an integral scheme \(S'\) finite and flat over an open subset \(V\) of \(S\), with function field \(k(S')=K\), there is an open subset \(V'\) of \(S'\) such that
\[
  \bigl[(f_{S'})_*({O}_{X'_{S'}})\bigr]_{\red}
  \simeq
  \bigl[(f_{S'})_*({O}_{X_{S'}})\bigr]_{\red}
\]
over \(V'\); after shrinking \(S\) and \(S'\), we may suppose that \(S=V\) and \(S'=V'\). Since the assertion to be proved is local for f.p.q.c., we may replace \(S\) by \(S'\). Finally, apply 3.1.

\sgaStableTarget{sga6.e12.rem.3-6}\textbf{3.6. Remark.} It follows in the same way, taking into account the exact sequences \(\underline{\Pic}_{X_i/S}\longrightarrow\underline{\Pic}_{X_{i-1}/S}\longrightarrow\mathbb V(Q_i)\), that, for every \(s\in S\),
\[
  \operatorname{coker}\bigl(\underline{\Pic}^0_{Y_s/s}\longrightarrow\underline{\Pic}^0_{X_s/s}\bigr)
\]
is a unipotent group.

\sgaStableTarget{sga6.e12.sec.4}\section*{4. Proof of II*: the most delicate part of the proof}

Recall assertion II*: Theorem 1.1 when \(X\) and \(Y\) are integral over \(S\) and the sequence (*) \({O}_Y\longrightarrow f_*{O}_X\rightrightarrows g_*({O}_{X\times_YX})\) is exact; applying the generic flatness theorem EGA IV 6.9.1 and EGA III 7.8.8 reduces us to proving the

\sgaStableTarget{sga6.e12.prop.4-1}\textbf{Proposition 4.1.} \uline{Let} \(S\) \uline{be an integral noetherian scheme,} \(\lambda:X\to S\) \uline{and} \(\mu:Y\to S\) \uline{proper morphisms of finite presentation,} \(f:X\to Y\) \uline{a surjective} \(S\)\uline{-morphism, and} \(g:X\times_YX\to Y\) \uline{the canonical morphism. Suppose that}
\begin{enumerate}[label=\arabic*)]
\item \(X,Y,X\times_YX\) \uline{are flat over} \(S\).
\item \(\lambda_*{O}_X={O}_S\), \(\mu_*{O}_Y={O}_S\) \uline{universally.}
%%%%%%%%%%  scan idx 623  |  folio 610  %%%%%%%%%%
\SGAPrintPage{610}%
\item
\uline{For} \(i,j\ge0\), \(R^if_*{O}_X\) \uline{and} \(R^jg_*{O}_{X\times_YX}\) \uline{are flat over} \(S\).
\item \uline{The sequence}
\sgaStableTarget{sga6.e12.eq.unnamed.occ2}\begin{equation}\tag{*}
  {O}_Y\longrightarrow f_*{O}_X
  \mathrel{\substack{\xrightarrow{\pi_1}\\[-0.7ex]\xrightarrow[\pi_2]{}}}
  g_*{O}_{X\times_YX}
\end{equation}
\uline{is exact and} \(\operatorname{coker}(\pi_1-\pi_2)\) \uline{is an} \({O}_Y\)\uline{-module flat over} \(S\).
\end{enumerate}
\uline{Then} \(f^*:\Pic_{Y/S}\longrightarrow\Pic_{X/S}\) \uline{is representable by a quasi-affine morphism of finite presentation.}

Indeed, first note that 1)-4) are stable under every base change \(T\to S\). Only the stability of 3) and 4) is not evident. But, by the Künneth formula (EGA III 6.9.9), 1) and 3) imply that the formation of the \(R^if_*({O}_X)\) and \(R^ig_*({O}_{X\times_YX})\) commutes with base changes \(T\to S\), which immediately implies the stability of conditions 3). Moreover, the formation of the cokernel of \(\pi_1-\pi_2:f_*{O}_X\longrightarrow g_*{O}_{X\times_YX}\) therefore also commutes with base changes; since it is flat over \(S\) (by 4)), \(\im(\pi_1-\pi_2)\) is flat over \(S\) and commutes with base changes. Consequently, \(\Ker(\pi_1-\pi_2)\) commutes with base changes \(T\to S\) and is therefore universally equal to \({O}_Y\).

The stability of the hypotheses implies two things: a. Since the assertion to be proved is local for f.p.p.f., by descent we may suppose that \(X\) has a section \(\alpha\), so \(Y\) has the section \(\beta=f\circ\alpha\). b. To establish the assertion it suffices to consider an \(S\)-point of \(\Pic_{X/S}\), that is, by a) and 2), an invertible sheaf \(L\) on \(X\) equipped with a rigidification \(u\) relative to \(\alpha\), and to show that the f.p.p.f. sheaf \(Z\), the inverse image of this point in \(\Pic_{Y/S}\), is representable by a quasi-affine \(S\)-scheme of finite presentation over \(S\).


%%%%%%%%%%  scan idx 624  |  folio 611  %%%%%%%%%%
\SGAPrintPage{611}%
\sgaStableTarget{sga6.e12.lem.4-2}\textbf{Lemma 4.2.} \uline{Under the conditions of} 4.1, \(f\) \uline{is a descent morphism for the category of locally free sheaves; that is, for every pair} \(L,M\) \uline{of locally free sheaves on} \(Y\)\uline{, the diagram}
\[
\begin{tikzcd}[column sep=tiny]
\Hom_Y(L,M) \arrow[r] &
\Hom_X(f^*L,f^*M) \arrow[r,shift left=.7ex] \arrow[r,shift right=.7ex] &
\Hom_{X\times_YX}(g^*L,g^*M)
\end{tikzcd}
\]
\uline{is exact.}

Indeed, for \(N=\underline{\Hom}_{{O}_Y}(L,M)\), we must show that the sequence
\[
  \Gamma(Y,N)\longrightarrow\Gamma(X,f^*N)
  \rightrightarrows
  \Gamma(X\times_YX,g^*N)
\]
is exact, or (by EGA \(0_1\) 5.4.10.1) that the sequence
\[
  \Gamma(Y,N)\longrightarrow\Gamma(Y,f_*{O}_X\otimes_{{O}_Y}N)
  \rightrightarrows
  \Gamma(Y,g_*{O}_{X\times_YX}\otimes_{{O}_Y}N)
\]
is exact. But the sequence (*) \({O}_Y\to f_*{O}_X\rightrightarrows g_*{O}_{X\times_YX}\) is exact by hypothesis; the lemma follows by applying to it the functor \(\Gamma(N\otimes -)\), which is left exact, \(N\) being flat over \(Y\).

Let \(L\) be an invertible sheaf on \(X\) equipped with a rigidification \(u\) relative to \(\alpha\); by 4.2 there is an equivalence between the category \(C\) of pairs \((M,\varphi)\) consisting of an invertible sheaf \(M\) on \(Y\) and an isomorphism \(\varphi:f^*M\xrightarrow{\sim}L\), and the category \(\mathrm{Deff}(S)\) of effective descent data on \(L\) relative to \(f\). But, for \(f^*M\simeq L\), since \(L\) is rigidified relative to \(\alpha\), \(M\) is canonically equipped with a rigidification relative to \(\beta\). Consequently, \(C\) is equivalent to the discrete category defined by the set of classes of \(\beta\)-rigidified invertible sheaves \((M,v)\) on \(Y\) such that \(f^*(M,v)\simeq(L,u)\). But this is precisely the set \(Z(S)\) of \(S\)-points of the inverse image in \(\Pic_{Y/S}\) of the \(S\)-point of \(\Pic_{X/S}\) defined by \((L,u)\). Moreover, this description of \(Z(S)\) is functorial in \(S\). Thus \(Z\) is identified with the functor \(\mathrm{Deff}\), and we may subsequently forget the rigidifications.


%%%%%%%%%%  scan idx 625  |  folio 612  %%%%%%%%%%
\providecommand{\Isom}{\operatorname{Isom}}
\providecommand{\Coker}{\operatorname{Coker}}
\providecommand{\Ker}{\operatorname{Ker}}
\providecommand{\Supp}{\operatorname{Supp}}
\providecommand{\Ass}{\operatorname{Ass}}

\SGAPrintPage{612}%

There is a canonical monomorphism from \(\mathrm{Deff}\) into the functor \(\mathrm{Rec}\) of gluing data on \(L\). By definition, \(\mathrm{Rec}\) is the functor \(\underline{\Isom}(p_1^*L,p_2^*L)\), where the \(p_i:X\times_YX\to X\) are the projections, which is represented by an affine scheme \(I\) of finite presentation over \(S\) by 2.4. It therefore remains to represent the monomorphism
\[
  i:\mathrm{Deff}\hookrightarrow \underline{\Isom}(p_1^*L,p_2^*L)
\]
by a morphism of finite presentation, which will necessarily be quasi-affine (by EGA IV 8.11.2).

Observe in passing that, if desired, we may restrict \(\underline{\Isom}(p_1^*L,p_2^*L)\) by imposing the cocycle condition, but this is not indispensable for what follows.

After the base change \(I\to S\), we are therefore reduced to the following problem: for an invertible sheaf \(L\) on \(X\) equipped with a gluing datum \(\varphi:p_1^*L\xrightarrow{\sim}p_2^*L\), represent the subfunctor of \(S\) ``that makes this datum effective''. But we have the

\sgaStableTarget{sga6.e12.lem.4-3}\textbf{Lemma 4.3.} \uline{Under the conditions of} 4.1\uline{, let} \(L\) \uline{be an invertible sheaf on} \(X\) \uline{equipped with a gluing datum}
\[
  \varphi:p_1^*L\xrightarrow{\sim}p_2^*L.
\]
\uline{For this datum to be effective, it is necessary and sufficient that the following five conditions hold.}
\begin{enumerate}[label=\arabic*)]
\item \uline{For every} \(i\ge0\), \(R^if_*(L)\) \uline{is flat over} \(S\).
\item \uline{For every} \(j\ge0\), \(R^jg_*(p_1^*L)\) \uline{is flat over} \(S\).
\item \uline{The cokernel of} \(\pi_1-\pi_2:f_*(L)\longrightarrow g_*(p_1^*L)\simeq g_*(p_2^*(L))\) \uline{is flat over} \(S\)\uline{, (where the last isomorphism is} \(g_*(\varphi)\)).
\item \uline{The kernel} \(N\) \uline{of} \(\pi_1-\pi_2:f_*(L)\to g_*(p_1^*L)\) \uline{is invertible on} \(Y\).
\item \uline{The composite morphism} \(f^*N\to f^*f_*L\to L\) \uline{is an isomorphism.}
\end{enumerate}
%%%%%%%%%%  scan idx 626  |  folio 613  %%%%%%%%%%
\SGAPrintPage{613}%
Indeed, these conditions are sufficient, since already 4) and 5) imply effectivity. Conversely, let \(M\) be invertible on \(Y\), with an isomorphism \(f^*M\simeq L\) compatible with the gluing datum. We may cover \(Y\) by open subsets \(V\) over which \(M\) is isomorphic to \({O}_Y\); hence the descent datum on \(L\) restricted over \(V=f^{-1}(V)\) is isomorphic to the trivial descent datum on \({O}_X|V\). Finally, conditions 1)-5) are local on \(Y\) and hold, by the hypotheses of 4.1, for \(L={O}_X\); they therefore hold for \(L\).

It therefore remains to represent each of conditions 1) to 5) successively by monomorphisms of finite presentation. For 1) and 2), we apply the

\sgaStableTarget{sga6.e12.lem.4-4}\textbf{Lemma 4.4.} \uline{Let} \(S\) \uline{be a scheme,} \(\lambda:X\to S\) \uline{and} \(\mu:Y\to S\) \uline{proper morphisms of finite presentation, and} \(f:X\to Y\) \uline{an} \(S\)\uline{-morphism. Let} \(F\) \uline{be a sheaf on} \(X\) \uline{of finite presentation and} flat \uline{over} \(S\) \uline{and, for an integer} \(n\)\uline{, let} \(S_n\) \uline{be the subfunctor of} \(S\) \uline{that makes} \(R^if_*(F)\) \uline{flat} \uline{over} \(S\) \uline{for every} \(i\ge n\) \uline{(more precisely,} \(S_n(T)=\{\varnothing\}\) \uline{when} \(R^i(f_T)_*(F_T)\) \uline{is} \uline{flat} \uline{for every} \(i\ge n\), \uline{and} \(S_n(T)=\varnothing\) \uline{otherwise). Then} \(S_n\hookrightarrow S\) \uline{is representable by a surjective monomorphism of finite presentation, and then, for} \(i\ge n\)\uline{, the} \(R^i(f_{S_n})_*(F_{S_n})\) \uline{commute with base changes} \(T\to S_n\)\uline{. Moreover,} \(S_n\) \uline{coincides with the functor} \(S'_n\) \uline{that makes the functors}
\[
  M\longmapsto R^i(f_V)_*(F_V\otimes\lambda_V^*M)\qquad \text{\uline{for} } i\ge n
\]
\uline{exact, where} \(M\) \uline{runs through the quasi-coherent} \({O}_V\)\uline{-modules for every affine open subset} \(V\) \uline{of} \(S\).

Indeed, the assertion is local on \(S\); we may therefore suppose that \(S\) is quasi-compact. The relative dimension of \(X/Y\) is then bounded, say by \(r\), and, for \(n>r\), \(R^nf_*(F_V\otimes\lambda_V^*M)=0\). The assertion therefore holds for \(n>r\);
%%%%%%%%%%  scan idx 627  |  folio 614  %%%%%%%%%%
\SGAPrintPage{614}%
proceeding by descending induction, suppose it is known for \(n+1\). After replacing \(S\) by \(S_{n+1}=S'_{n+1}\), we may suppose that \(R^if_*(F)\) is \(S\)-flat for \(i\ge n+1\). Since \(F\) is flat over \(S\), it then follows from EGA III 6.9.8 that \(R^nf_*(F)\) commutes with every base change. Consequently, the representability of \(S_n\) follows from the next lemma (see: J.P. Murre, Representation of unramified functors, Séminaire Bourbaki; May 1965, p.294-11, th.2):

\sgaStableTarget{sga6.e12.lem.4-5}\textbf{Lemma 4.5.} \uline{Let} \(X\to S\) \uline{be a proper morphism of finite presentation, and} \(F\) \uline{a sheaf of finite presentation on} \(X\)\uline{. Then the subfunctor} \(S_F\) \uline{of} \(S\) \uline{that makes} \(F\) \uline{flat over} \(S\) \uline{is representable by a surjective monomorphism of finite presentation (which is of course not an immersion in general).}

To represent condition 3) of 4.3, note that, by 1) and 2), \(f_*(L)\) and \(g_*(p_1^*L)\) commute with base changes \(I\to S\), and hence \(\operatorname{coker}(\pi_1-\pi_2)\) does also; then apply 4.5.

To represent 4), note that, by 3) and the fact that \(f_*(L)\) and \(g_*(p_1^*L)\) are flat and commute with base changes, \(N=\ker(\pi_1-\pi_2)\) also commutes; then apply the

\sgaStableTarget{sga6.e12.lem.4-6}\textbf{Lemma 4.6.} \uline{Let} \(\lambda:X\to S\) \uline{be a proper morphism of finite presentation, and} \(u:F\to G\) \uline{a homomorphism of} \({O}_X\)\uline{-Modules of finite presentation, with} \(G\) \uline{flat over} \(S\)\uline{. Then the subfunctor} \(S_1\) \uline{of} \(S\) \uline{that makes} \(u\) \uline{an isomorphism is representable by an open subscheme of} \(S\)\uline{, of finite presentation.}

As usual, we reduce to the case where \(S\) is noetherian, using EGA IV 8 and 11.2.6, which relieves us of the need to consider finite presentation. First let us represent the functor that makes \(u\) an epimorphism,
%%%%%%%%%%  scan idx 628  |  folio 615  %%%%%%%%%%
\SGAPrintPage{615}%
i.e. makes \(\operatorname{Coker}u\) zero. Since the formation of \(\operatorname{Coker}u\) commutes with base change, we see that it is represented by the open subset \(S-\lambda(\operatorname{supp}(\operatorname{Coker}u))\) of \(S\). We may therefore suppose that \(u\) is already an epimorphism. But, now using the hypothesis that \(G\) is \(S\)-flat, we find that the formation of \(\operatorname{Ker}u\) commutes with base change, and consequently the functor under consideration is represented by the open subset \(S-\lambda(\operatorname{supp}(\operatorname{Ker}u))\) of \(S\).

% Expose XIII local macros
\providecommand{\Ass}{\operatorname{Ass}}
\providecommand{\Hom}{\operatorname{Hom}}
\providecommand{\Ext}{\operatorname{Ext}}
\providecommand{\Proj}{\operatorname{Proj}}
\providecommand{\NS}{\operatorname{NS}}
\providecommand{\Tors}{\operatorname{Tors}}
\providecommand{\Picnum}{\operatorname{Pic}^{\mathrm{num}}}
\providecommand{\Piczero}{\operatorname{Pic}^{0}}
\providecommand{\Pictau}{\operatorname{Pic}^{\tau}}
\providecommand{\length}{\operatorname{length}}
\providecommand{\etale}{\text{\'etale}}

\newpage

%%%%%%%%%%  scan idx 629  |  folio 616  %%%%%%%%%%
\SGAPrintPage{616}%
\sgaStableTarget{sga6.e13.expose}\section*{\underline{EXPOS\'E XIII}\\[3pt] \underline{\underline{FINITENESS THEOREMS FOR THE PICARD FUNCTOR}}}
\begin{center}
\underline{by S. KLEIMAN}
\end{center}

\sgaStableTarget{sga6.e13.sec.1}\subsection*{1. The (b)-sheaves}
Let $k$ be an algebraically closed field and $X$ a projective $k$-scheme, endowed with an ample invertible module ${O}_X(1)$.

\sgaStableTarget{sga6.e13.def.1-1}\textbf{Definition 1.1 (Mumford).} \uline{Let} $m\in\mathbb Z$. \uline{A coherent sheaf} $F$ \uline{on} $X$ \uline{is said to be} $m$-regular \uline{if, at every point of the support of} $F$, ${O}_X(1)$ \uline{is generated by its global sections and if, for every} $q\ge1$, \uline{we have}
\[
 H^q(F(m-q))=0.
\]

\sgaStableTarget{sga6.e13.lem.1-2}\textbf{Lemma 1.2.} \uline{Let} \(0\longrightarrow F(-1)\longrightarrow F\longrightarrow G\longrightarrow 0\) \uline{be an exact sequence of coherent sheaves on} $X$\SGAAuthorFootnote{*}{Here, as below in analogous situations, it is understood that \(F(-1)\to F\) is defined by tensoring with a given section of \(\underline{O}_X(1)\).}. \uline{If} $F$ \uline{is} $m$\uline{-regular, then so is} $G$.

Indeed, we have the corresponding exact sequence:
\[
H^q(F(m-q))\longrightarrow H^q(G(m-q))\longrightarrow H^{q+1}(F(m-q-1)).
\]

\sgaStableTarget{sga6.e13.prop.1-3}\textbf{Proposition 1.3.} \uline{If a sheaf} $F$ \uline{on} $X$ \uline{is} $m$\uline{-regular, then, for every} $n\ge m$, \uline{we have}~:
\begin{enumerate}[label=(\roman*)]
\item $F$ \uline{is} $n$\uline{-regular}.
\item \(H^0(F(n))\otimes H^0({O}_X(1))\longrightarrow H^0(F(n+1))\) \uline{is surjective}.
\item $F(n)$ \uline{is generated by} \(H^0(F(n))\).
\end{enumerate}

%%%%%%%%%%  scan idx 630  |  folio 617  %%%%%%%%%%
\SGAPrintPage{617}%
The proof is by induction on the dimension \(s\) of the support of \(F\); when \(s=-1\), it is trivial. Take a section \(\sigma\in H^0({O}_X(1))\) which does not vanish at any point of the finite set \(\Ass(F)\); then \(\sigma\) defines an exact sequence \(0\longrightarrow F(-1)\longrightarrow F\longrightarrow G\longrightarrow 0\), and \(G\) is \(m\)-regular by 1.2., with support of dimension \(s-1\).

Consider the corresponding exact sequence:
\[
H^q(F(n-q-1))\longrightarrow H^q(F(n-q))\longrightarrow H^q(G(n-q)).
\]
By the induction hypothesis, we have \(H^q(G(n-q))=0\) for \(n\ge m\), so \(H^q(F(n-q-1))=0\) implies \(H^q(F(n-q))=0\); but \(H^q(F(m-q))=0\), which gives (i).

To establish (ii), we chase through the following commutative diagram:
\[
\begin{tikzcd}[column sep=1.8em,row sep=2.5em]
 & & & H^0(F(n))\otimes H^0({O}_X(1)) \arrow[rr] \arrow[d]
 & & H^0(G(n))\otimes H^0({O}_X(1)) \arrow[r] \arrow[d]
 & 0 \\
0 \arrow[rr]
& & H^0(F(n)) \arrow[ur] \arrow[r]
& H^0(F(n+1)) \arrow[rr]
& & H^0(G(n+1)) \arrow[d] \\
& & & & & 0
\end{tikzcd}
\]
where the oblique arrow is defined by \(\varphi\mapsto\varphi\otimes\sigma\), the horizontal rows are exact, the first by the \(n\)-regularity of \(F\) and by (i), already proved, and the second column is exact by the induction hypothesis.

Finally, (iii) follows from the following commutative diagram:
\[
\begin{tikzcd}[column sep=large,row sep=large]
H^0(F(n))_X\otimes H^0({O}_X(1))_X
  \arrow[r,"\alpha_n"]
  \arrow[d]
& H^0(F(n)(1))_X
  \arrow[d,"\beta_{n+1}"] \\
\bigl[H^0(F(n))_X\bigr](1)
  \arrow[r,"\beta_n\otimes\mathrm{id}"]
& F(n)(1)
\end{tikzcd}
\]
%%%%%%%%%%  scan idx 631  |  folio 618  %%%%%%%%%%
\SGAPrintPage{618}%
where, if \(M\) is a vector space over \(k\), we write \(M_X\) for \(M\otimes_k{O}_X\). For \(n\ge m\), \(\alpha_n\) is surjective by (ii); therefore, if \(\beta_{n+1}\) is surjective, so is \(\beta_n\). By Serre's theorem, \(\beta_{n+1}\) is surjective for \(n\gg0\); hence (iii).

\sgaStableTarget{sga6.e13.prop.1-4}\textbf{Proposition 1.4.} \uline{Let} \(0\to F(-1)\to F\to G\to0\) \uline{be an exact sequence on} \(X\). \uline{If} \(G\) \uline{is} \(m\)\uline{-regular, then we have}~:
\begin{enumerate}[label=(\roman*)]
\item \(H^q(F(n))=0\) \uline{for} \(q\ge2\) \uline{and} \(n\ge m-q\).
\item \(h^1(F(n-1))\ge h^1(F(n))\) \uline{for} \(n\ge m-1\).
\item \(H^1(F(n))=0\) \uline{for} \(n\ge (m-1)+h^1(F(m-1))\).
\end{enumerate}
\uline{In particular}, \(F\) \uline{is} \([m+h^1(F(m-1))]\)\uline{-regular}.

Indeed, in the exact sequence
\[
H^{q-1}(G(n))\longrightarrow H^q(F(n-1))\longrightarrow H^q(F(n))\longrightarrow H^q(G(n)),
\]
the two outer groups are \(0\) for \(q\ge2\) and \(n\ge m-(q-1)\) by 1.3 (i). Thus \(H^q(F(m-q))\xrightarrow{\sim}H^q(F(m-q+1))\xrightarrow{\sim}H^q(F(m-q+2))\xrightarrow{\sim}\ldots\), but, by Serre's theorem, \(H^q(F(n))=0\) for \(n\gg0\); hence (i).

For \(n\ge m-1\), consider the following exact sequence:
\[
0\to H^0(F(n-1))\to H^0(F(n))
  \xrightarrow{\alpha_n}H^0(G(n))\to H^1(F(n-1))
  \to H^1(F(n))\to0
\]
It implies (ii). Now suppose that \(\alpha_n\) is surjective and consider the following commutative diagram:
\[
\begin{tikzcd}[column sep=large,row sep=large]
H^0(F(n))\otimes H^0({O}_X(1))
  \arrow[r,"\alpha_n\otimes\mathrm{id}"]
  \arrow[d]
& H^0(G(n))\otimes H^0({O}_X(1))
  \arrow[r]
  \arrow[d,"\beta_n"]
& 0 \\
H^0(F(n+1))
  \arrow[r,"\alpha_{n+1}"]
& H^0(G(n+1)).
\end{tikzcd}
\]
%%%%%%%%%%  scan idx 632  |  folio 619  %%%%%%%%%%
\SGAPrintPage{619}%
If \(n\ge m\), then \(\beta_n\) is surjective by 1.3 (ii); therefore \(\alpha_{n+1}\) is surjective. Consequently, \(H^1(F(n-1))\xrightarrow{\sim}H^1(F(n))\xrightarrow{\sim}\ldots\xrightarrow{\sim}0\). On the other hand, if \(\alpha_n\) is not surjective, then \(h^1(F(n-1))>h^1(F(n))\); therefore \(n\mapsto h^1(F(n))\) is \uline{strictly} decreasing from \(h^1(F(m-1))\) until it becomes \(0\); hence (iii).

\sgaStableTarget{sga6.e13.def.1-5}\textbf{Definition 1.5.} \uline{Let} \(F\) \uline{be a coherent sheaf on} \(X\), \(r\) \uline{an integer} \(\ge\operatorname{dim{.}supp}(F)\) \uline{and} \((b)=(b_0,\ldots,b_r)\) \uline{a sequence of} \(r+1\) \uline{integers. Then} \(F\) \uline{is said to be a} \((b)\)-sheaf \uline{if it satisfies one of the following two equivalent sets of conditions}~:

\sgaStableTarget{sga6.e13.para.1-5-1}\((1.5.1)\)
\begin{enumerate}[label=(\roman*)]
\item \uline{At every point of the support of} \(F\), \({O}_X(1)\) \uline{is generated by} \(H^0({O}_X(1))\).
\item \(h^0(F(-1))\le b_0\).
\item (\uline{When} \(r\ge1\)), \uline{there exists a section} \(\sigma\in H^0({O}_X(1))\) \uline{defining an exact sequence} \(0\to F(-1)\to F\to G\to0\) \uline{such that} \(G\) \uline{is a} \((b_1,\ldots,b_r)\)\uline{-sheaf}.
\end{enumerate}

\sgaStableTarget{sga6.e13.para.1-5-2}\((1.5.2)\) \uline{There exists an} \(F\)\uline{-regular sequence of} \(r\) \uline{sections} \(\sigma_1,\ldots,\sigma_r\in H^0({O}_X(1))\) \uline{such that, for} \(i=0,\ldots,r\), \uline{if} \(F_i\) \uline{denotes the restriction of} \(F\) \uline{to the zero scheme of} \(\sigma_1,\ldots,\sigma_i\), \uline{then} \(h^0(F_i(-1))\le b_i\).

\sgaStableTarget{sga6.e13.prop.1-6}\textbf{Proposition 1.6.} \uline{Let} \(F\) \uline{be a} \((b)\)\uline{-sheaf on} \(X\). \uline{Then}~:
\begin{enumerate}[label=(\roman*)]
\item \uline{For every ``sufficiently general'' sequence} \(\sigma\) \uline{of} \(r\) \uline{sections} \(\sigma_1,\ldots,\sigma_r\in H^0({O}_X(1))\), \uline{if} \(F_{\sigma,i}\) \uline{denotes the restriction of} \(F\) \uline{to the zero scheme of} \(\sigma_1,\ldots,\sigma_i\), \uline{then} \(h^0(F_{\sigma,i}(-1))\le b_i\), \uline{where} \(0\le i\le r\).
\item \uline{Every subsheaf} \(G\) \uline{of} \(F\) \uline{is a} \((b)\)\uline{-sheaf}.
\end{enumerate}

%%%%%%%%%%  scan idx 633  |  folio 620  %%%%%%%%%%
\SGAPrintPage{620}%
Indeed, let \(S'\) be the affine space \(\mathbb A_k^N\) whose \(k\)-points correspond to the sequences \(\sigma\), and \(T\) the open subset of \(S'\) corresponding to the \(F\)-regular sequences. For fixed \(i\), the sheaves \(F_{\sigma,i}(-1)\) corresponding to the \(k\)-points of \(T\) are contained in a flat family over \(T\), and by hypothesis \(T\) is nonempty. Thus (i) follows from the upper semi-continuity property of the function $\sigma\longmapsto h^0(F_{\sigma,i}(-1))$ (EGA III, 7.7.5). Finally, since every ``sufficiently general'' sequence is \(G\)-regular and such that \(G_{\sigma,i}\) is a subsheaf of \(F_{\sigma,i}\) for every \(i\), (ii) follows from (i).

\sgaStableTarget{sga6.e13.lem.1-7}\textbf{Lemma 1.7.} \uline{Let} \(0\to F(-1)\to F\to G\to0\) \uline{be an exact sequence on} \(X\). \uline{If the Hilbert polynomial of} \(F\) \uline{is written in the form}
\[
 \chi(F(n))=\sum_{i=0}^{r}a_i\binom{n+i}{i},
\]
\uline{then}
\[
 \chi(G(n))=\sum_{i=0}^{r-1}a_{i+1}\binom{n+i}{i}.
\]
Indeed, \(\chi(G(n))=\chi(F(n))-\chi(F(n-1))=\sum_{i=0}^{r}a_i\left[\binom{n+i}{i}-\binom{n-1+i}{i}\right]\); therefore \(\chi(G(n))=\sum_{i=1}^{r}a_i\binom{n+i-1}{i-1}\), which gives the assertion.

\sgaStableTarget{sga6.e13.prop.1-8}\textbf{Proposition 1.8.} \uline{Let} \(F\) \uline{be a} \((b)\)\uline{-sheaf on} \(X\), \(s=\operatorname{dim{.}supp}(F)\) \uline{and} \(\chi(F(n))=\sum_{i=0}^{s}a_i\binom{n+i}{i}\) \uline{the Hilbert polynomial of} \(F\). \uline{Then}~:
\begin{enumerate}[label=(\roman*)]
\item \uline{For} \(n\ge -1\), \(h^0(F(n))\le \sum_{i=0}^{s}b_i\binom{n+i}{i}\).
\item \(a_s\le b_s\) \uline{and} \(F\) \uline{is also a} \((b_0,\ldots,b_{s-1},a_s)\)\uline{-sheaf}.
\end{enumerate}
Indeed, we argue by induction on \(s\). When \(s=0\), this is clear since \(a_0=h^0(F)=h^0(F(-1))\le b_0\). When \(s\ge1\), we have an exact sequence
\[
0\longrightarrow F(-1)\longrightarrow F\longrightarrow G\longrightarrow0
\]
%%%%%%%%%%  scan idx 634  |  folio 621  %%%%%%%%%%
\SGAPrintPage{621}%
with \(G\) a \((b_1,\ldots,b_r)\)-sheaf and \(\operatorname{dim{.}supp}(G)=s-1\). Then
\[
h^0(F(n))-h^0(F(n-1))\le h^0(G(n)),
\]
but
\[
h^0(G(n))\le\sum_{i=0}^{s-1}b_{i+1}\binom{n+i}{i}
\]
by induction on \(s\) and \(h^0(F(-1))\le b_0\); hence (i) by induction on \(n\ge -1\). Finally, by 1.7 and induction on \(s\), \(a_s\le b_s\) and \(G\) is a \((b_1,\ldots,b_{s-1},a_s)\)-sheaf; hence (ii).

\sgaStableTarget{sga6.e13.def.1-9}\textbf{Definition 1.9.} \uline{The following polynomials with rational coefficients} \(\ge0\), \uline{defined by induction, are called the} \((b)\)-polynomials~:
\[
\begin{cases}
P_{-1}=0\\
P_r(X_0,\ldots,X_r)=P_{r-1}(X_1,\ldots,X_r)
+\displaystyle\sum_{i=0}^{r}X_i
\binom{P_{r-1}(X_1,\ldots,X_r)-1+i}{i}.
\end{cases}
\]

\sgaStableTarget{sga6.e13.rem.1-10}\textbf{Remark 1.10.} For the \((b)\)-polynomials, one sees at once by induction on \(r\ge t\) that
\[
P_r(X_0,\ldots,X_t,0,\ldots,0)=P_t(X_0,\ldots,X_t).
\]

\sgaStableTarget{sga6.e13.thm.1-11}\textbf{Theorem 1.11.} (\uline{The main theorem on} \((b)\)\uline{-sheaves). Let} \(F\) \uline{be a} \((b)\)\uline{-sheaf on} \(X\) \uline{and} \(\chi(F(n))=\sum_{i=0}^{r}a_i\binom{n+i}{i}\) \uline{the Hilbert polynomial of} \(F\); \uline{let} \((c)=(c_0,\ldots,c_r)\) \uline{be a sequence of integers} \(\ge (b)-(a)\) \uline{and} \(m=P_r(c)\) \uline{the integer given by the} \(r\)\uline{-th} \((b)\)\uline{-polynomial} (1.9). \uline{Then} \(m\ge0\) \uline{and} \(F\) \uline{is} \(m\)\uline{-regular. In particular, if} \(s=\operatorname{dim{.}supp}(F)\), \uline{then} \(F\) \uline{is} \(P_{s-1}(c_0,\ldots,c_{s-1})\)\uline{-regular}.

%%%%%%%%%%  scan idx 635  |  folio 622  %%%%%%%%%%
\SGAPrintPage{622}%
Indeed, we argue by induction on \(r\). When \(r=0\), this is clear since \(m\ge0\) and \(F\) is \(n\)-regular for every \(n\). When \(r\ge1\), we have an exact sequence $0\longrightarrow F(-1)\longrightarrow F\longrightarrow G\longrightarrow0$ where \(G\) is a \((b_1,\ldots,b_r)\)-sheaf. By 1.7 and the induction hypothesis, we see that if \(n=P_{r-1}(c_1,\ldots,c_r)\), then \(n\ge0\) and \(G\) is \(n\)-regular. Thus, by 1.4, \(F\) is \([n+h^1(F(n-1))]\)-regular and \(h^q(F(n-1))=0\) for \(q\ge2\). Therefore, \(h^1(F(n-1))=h^0(F(n-1))-\chi(F(n-1))\), which is \(\le \sum_{i=0}^{r}(b_i-a_i)\binom{n-1+i}{i}\) by 1.8 (i), since \(b_i\ge0\). Finally, \(F\) is \(\left[n+\sum_{i=0}^{r}c_i\binom{n-1+i}{i}\right]\)-regular by 1.3 (i); this gives the first assertion. The second follows from it by 1.8 (ii) and 1.10.

\sgaStableTarget{sga6.e13.para.1-12}1.12. Let \(S\) be a noetherian scheme and \(X\) an \(S\)-scheme of finite type. Let \(\mathcal F\) be a \uline{family of classes of coherent sheaves on the fibers of} \(X/S\); that is, for every point \(s\in S\) and every extension \(K\) of \(k(s)\), coherent sheaves \(F_K\) on the algebraic scheme \(X_K\) are given, and we say that an \(F_K\) and an \(F'_{K'}\) belong to the same class if there exist \(k(s)\)-homomorphisms from \(K,K'\) into a common extension \(K''\) of \(k(s)\) such that \(F_{K''}(=F_K\otimes_K K'')\) and \(F'_{K''}\) are isomorphic on \(X_{K''}\). Then the family \(\mathcal F\) is said to be \uline{bounded by the coherent sheaf} \(F\) \uline{on} \(X_T=X\times_S T\), where \(T\) is of finite type over \(S\), if \(\mathcal F\) is contained in the family of classes of the sheaves \(F_{k(t)}\) with \(t\in T\); we say that \(\mathcal F\) is \uline{bounded} if there exist \(T,F\) as above.

Suppose in addition that \(X/S\) is projective and endowed with a (relatively) ample invertible sheaf \({O}_X(1)\). Then \(\mathcal F\) is said to be a \((b)\)-\uline{family}, for a sequence of integers \((b)=(b_0,\ldots,b_r)\), if every class in \(\mathcal F\) is represented by an \(F_K\) with \(K\) algebraically closed, which is a \((b)\)-sheaf.

%%%%%%%%%%  scan idx 636  |  folio 623  %%%%%%%%%%
\SGAPrintPage{623}%
\sgaStableTarget{sga6.e13.thm.1-13}\textbf{Theorem 1.13.}
\uline{Let} \(S\) \uline{be a noetherian scheme and} \(X\) \uline{a projective} \(S\)\uline{-scheme endowed with an ample sheaf} \({O}_X(1)\) \uline{whose induced sheaves} \({O}_{X_s}(1)\), \uline{with} \(s\in S\), \uline{are generated by their} \(H^0({O}_{X_s}(1))\); \uline{let} \(\mathcal F\) \uline{be a family of classes of coherent sheaves on the fibers of} \(X/S\). \uline{Then the following conditions are equivalent}:
\begin{enumerate}[label=(\roman*)]
\item \(\mathcal F\) \uline{is bounded. Moreover, if all the} \(F_K\in\mathcal F\) \uline{are locally free of rank} \(p\), \uline{then one may suppose that} \(\mathcal F\) \uline{is bounded by a sheaf} \(F\) \uline{on an} \(X_T\), \uline{with} \(F\) \uline{locally free of rank} \(p\).
\item \uline{The set of Hilbert polynomials} \(\chi(F_K(n))\) \uline{of the} \(F_K\in\mathcal F\) \uline{is finite, and there exists a sequence of integers} \((b)\) \uline{such that} \(\mathcal F\) \uline{is a} \((b)\)\uline{-family}.
\item \uline{The set of Hilbert polynomials} \(\chi(F_K(n))\) \uline{of the} \(F_K\in\mathcal F\) \uline{is finite, and there exists an integer} \(m\) \uline{such that all the} \(F_K\in\mathcal F\) \uline{are} \(m\)\uline{-regular}.
\item \uline{The set of Hilbert polynomials} \(\chi(F_K(n))\) \uline{of the} \(F_K\in\mathcal F\) \uline{is finite, and} \(\mathcal F\) \uline{is contained in the family of classes of quotients of sheaves of the form} \(E_K\), \uline{where} \(E\) \uline{is a coherent sheaf on an} \(X_T\). \uline{Moreover, one may suppose} \(T=S\) \uline{and} \(E\) \uline{of the form} \({O}_X(-m)^{\oplus M}\).
\item \(\mathcal F\) \uline{is contained in the family of classes of cokernels of homomorphisms of the form} \(E'_K\to E_K\), \uline{where} \(E'\) \uline{and} \(E\) \uline{are coherent sheaves on an} \(X_T\). \uline{Moreover, one may suppose} \(T=S\) \uline{and} \(E'\) \uline{and} \(E\) \uline{of the forms} \({O}_X(-m')^{\oplus M'}\) \uline{and} \({O}_X(-m)^{\oplus M}\).
\end{enumerate}
Indeed, suppose that \(\mathcal F\) is bounded by a sheaf \(F\) on an \(X_T\); after decomposing \(T\) into pieces, using the generic flatness theorem (EGA IV 6.9.1), one may suppose \(F\) flat over \(T\). Then,
%%%%%%%%%%  scan idx 637  |  folio 624  %%%%%%%%%%
\SGAPrintPage{624}%
the number of polynomials \(\chi(F(n))\) is at most equal to the number of connected components of \(T\) (EGA III 7.9.5), and it is a small lemma that if, for a \(t\in T\), \(F_t\) is locally free of rank \(p\), then so is \(F\) above an open neighborhood of \(t\) in \(T\) (see EGA IV 11.3.10, where one takes \(X=Y\), \(f=\mathrm{id}\)). Now it follows from the hypotheses and (EGA III 7.7.6) that, after decomposing \(T\) more finely, one can find a sequence \(\sigma\) of sections \(\sigma_1,\ldots,\sigma_r\in H^0({O}_X(1))\) which is \(F\)-regular. Thus (ii) follows from the upper semi-continuity property of the function \(t\longmapsto h^0(F_{t,i}(-1))\) where, for \(t\in T\) and \(0\le i\le r\), \(F_{t,i}\) denotes the restriction of \(F_t\) to the zero scheme of \(\sigma_1,\ldots,\sigma_i\) (EGA III 7.7.5).

The implication (ii) \(\Rightarrow\) (iii) follows immediately from 1.11. The implication (iii) \(\Rightarrow\) (iv) follows from 1.3 (iii) by taking \(M=\max(\chi(F_K(m)))\) and \(E={O}_X(-m)^{\oplus M}\).

Suppose that \(\mathcal F\) satisfies (iv); there exist exact sequences
\[
0\longrightarrow F'_K\longrightarrow E_K\longrightarrow F_K\longrightarrow0
\]
and the set of polynomials \(\chi(F_K(n))\) is finite. By hypothesis, the family of classes of the \(E_K\) is bounded; thus, by the implication (i) \(\Rightarrow\) (ii) already proved, the set of polynomials \(\chi(E_K(n))\) is finite and there exists a sequence of integers \((b)\) such that all the \(E_K\) (with \(K\) algebraically closed) are \((b)\)-sheaves. Consequently, the set of polynomials \(\chi(F'_K(n))\) is finite and, by 1.6, all the \(F'_K\) are \((b)\)-sheaves. Thus, applying to the family of classes of the \(F'_K\) the implication (ii) \(\Rightarrow\) (iv) already proved, we obtain (v).

Finally, suppose that \(\mathcal F\) satisfies (v); let us prove that \(\mathcal F\) is bounded. One may suppose \(E\) (resp. \(E'\)) of the form \({O}_X(-m)^{\oplus M}\). Indeed,
%%%%%%%%%%  scan idx 638  |  folio 625  %%%%%%%%%%
\SGAPrintPage{625}%
by (i) \(\Rightarrow\) (iv), take a surjection \(L={O}_{X_T}(-m)^{\oplus M}\longrightarrow E\). After decomposing \(T\) into pieces so that \(E\) becomes flat over \(T\), the formation of the exact sequence \(0\longrightarrow I\xrightarrow{u}L\longrightarrow E\longrightarrow0\) commutes with passage to the fibers. By (i) \(\Rightarrow\) (iii) and (i) \(\Rightarrow\) (iv), take \(m_1\) large enough that all the \(\operatorname{Ext}^1({O}_{X_K}(-m_1),I_K)=H^1(I_K(m_1))\) vanish and that there exist surjections \(L_1={O}_{X_T}(-m_1)^{\oplus M_1}\longrightarrow E'\) and \(\alpha:L_2={O}_{X_T}(-m_1)^{\oplus M_2}\longrightarrow I\). Then the maps \(\operatorname{Hom}(L_{1,K},L_K)\longrightarrow \operatorname{Hom}(L_{1,K},E_K)\) are surjective. Let \(\beta:E'_K\to E_K\) be a homomorphism and let \(\gamma\) be the composition \(L_{1,K}\to E'_K\to E_K\); then \(\gamma\) comes from a homomorphism \(\delta:L_{1,K}\to L_K\), and $(\delta,u_K\circ\alpha):L_{1,K}\oplus L_{2,K}\longrightarrow L_K$ and \(\beta:E'_K\to E_K\) have the same cokernel. Consequently, \(\mathcal F\) is contained in the family of classes of cokernels of homomorphisms of the form \({O}_{X_K}(-m_1)^{\oplus(M_1+M_2)}\longrightarrow {O}_{X_K}(-m)^{\oplus M}\).

Again after decomposing \(T\) into pieces, one may suppose in addition that \(X_T\) (hence \(\underline{\Hom}_{{O}_{X_T}}(E',E)\)) and all the \(R^q(f_T)_*\underline{\Hom}_{{O}_{X_T}}(E',E)\) (where \(f_T:X_T\to T\) is the structural morphism) are flat over \(T\), so that the formation of \(G=(f_T)_*\underline{\Hom}_{{O}_{X_T}}(E',E)\) commutes with arbitrary base change \(T'\to T\) (EGA III 6.9.9.2). Let $R=\mathbb V(G^*)$; then \(X_R\) canonically carries a ``universal'' exact sequence
\[
E'_R\longrightarrow E_R\longrightarrow F\longrightarrow0
\]
and \(F\) bounds \(\mathcal F\).

\sgaStableTarget{sga6.e13.sec.2}\subsection*{2. Several technical lemmas}
\sgaStableTarget{sga6.e13.para.2-1}2.1. Let \(k\) be a field and \(X\) a proper \(k\)-scheme. Let \(L_1,\ldots,L_s\) be invertible sheaves on \(X\), \(c_1(L_i)\) the Chern class of \(L_i\) and \(Y\) a closed subscheme of \(X\) of dimension \(\le s\). Recall that we denote by
%%%%%%%%%%  scan idx 639  |  folio 626  %%%%%%%%%%
\SGAPrintPage{626}%
\(\langle c_1(L_1)\ldots c_1(L_s)\cdot Y\rangle\) the intersection number, which can be calculated as the coefficient of the monomial \(n_1\ldots n_s\) in the Snapper polynomial \(\chi\bigl(L_1^{\otimes n_1}\otimes\ldots\otimes L_s^{\otimes n_s}\otimes{O}_Y\bigr)\) (X 4.3.1).

\sgaStableTarget{sga6.e13.para.2-2}2.2. From now on, we suppose that \(k\) is algebraically closed, that \(X\) is projective, endowed with a very ample invertible sheaf \(H={O}_X(1)\), and that \(X\) is integral, of dimension \(r\).

Let \(L\) be an invertible sheaf on \(X\). The integer \(\langle c_1(L)\cdot c_1(H)^{r-1}\rangle\) is called the \uline{degree} of \(L\) (relative to \(H\)) and is denoted by \(d(L)\).

In what follows we consider an invertible sheaf \(L\) and the Hilbert polynomials: \(\chi(L(n))=\sum_{i=0}^{r}a_i\binom{n+i}{i}\) and \(\chi({O}_X(n))=\sum_{i=0}^{r}e_i\binom{n+i}{i}\).

\sgaStableTarget{sga6.e13.lem.2-3}\textbf{Lemma 2.3.} \uline{We have}:
\begin{enumerate}[label=(\roman*)]
\item \(a_r=e_r=d(H)\).
\item (\uline{If} \(r\ge1\)), \(a_{r-1}-e_{r-1}=d(L)\).
\end{enumerate}

Indeed, we obtain (i) and (ii) at once by taking \(m=0,1\) in the following polynomial:
\[
\chi\bigl(L^{\otimes m}\otimes H^{\otimes n}\bigr)
=
\sum_{i+j\le r} f_{ij}\binom{m+i}{i}\binom{n+j}{j}.
\]

\sgaStableTarget{sga6.e13.lem.2-4}\textbf{Lemma 2.4.} \uline{If} \(d(L)\le0\) \uline{and} \(h^0(L)\ne0\), \uline{then} \(L\simeq{O}_X\); \uline{therefore} \(d(L)=0\) \uline{and} \(h^0(L)=1\).

Indeed, since \(X\) is integral, there exists a divisor \(Y\ge0\) such that \(L\simeq{O}_X(Y)\). If \(Y>0\), we would have \(0\ge\langle c_1(L)\cdot c_1(H)^{r-1}\rangle=\langle c_1(H)^{r-1}\cdot Y\rangle\ge1\).

\sgaStableTarget{sga6.e13.lem.2-5}\textbf{Lemma 2.5.} \uline{If} \(d(L)\le0\), \uline{then} \(L\) \uline{is a} \((0,\ldots,0,a_r)\)\uline{-sheaf (with} \(r\) \uline{zeros}).

%%%%%%%%%%  scan idx 640  |  folio 627  %%%%%%%%%%
\SGAPrintPage{627}%
Indeed, if \(r=0\), \(L\) is evidently an \(a_0\)-sheaf. If \(r\ge1\), by 2.4, \(h^0(L(-1))=0\) since \(d(L(-1))<0\). Let \(Y\) be a hyperplane section of \(X\) and let \(0\longrightarrow L(-1)\longrightarrow L\longrightarrow L_Y\longrightarrow0\) be an exact sequence defined by \(Y\). If \(r=1\), by 1.7, \(h^0(L_Y(-1))=a_1\); thus \(L\) is a \((0,a_1)\)-sheaf. If \(r\ge2\), one can take \(Y\) integral by Bertini's theorem and one may thus suppose by induction that \(L_Y\) is a \((0,\ldots,a_r)\)-sheaf; this gives the assertion.

\sgaStableTarget{sga6.e13.lem.2-6}\textbf{Lemma 2.7.} \uline{If} \(r=2\) \uline{and} \(d(L)\le0\), \uline{then}:
\begin{enumerate}[label=(\roman*)]
\item \uline{For every} \(n\), \(h^1(L(n))\le a_1(a_1+1)a_2-a_0\)
\item \(a_0\le a_1(a_1+1)\).
\end{enumerate}

Let \(Y\) be an integral hyperplane section of \(X\). By 1.7 and 2.5, \(L_Y\) is a \((0,a_2)\)-sheaf with Hilbert polynomial \(\chi(L_Y(n))=a_1+a_2\binom{n+1}{1}\). Thus, by 1.11, \(-a_1\ge0\) and \(L_Y\) is \([-a_1]\)-regular.

Consider the following exact sequence:
\[
H^0(L_Y(n))\longrightarrow H^1(L(n-1))
\longrightarrow H^1(L(n))\longrightarrow H^1(L_Y(n)).
\]
For \(n\ge -a_1-1\), we have \(h^1(L(-a_1-2))\ge h^1(L(n))\) since \(h^1(L_Y(n))=0\) by 1.3 (i); for \(n\le -1\), we have \(h^1(L(n-1))\le h^1(L(-1))\) since \(h^0(L_Y(n))=0\) by 2.4. Finally, for every \(n\), \(h^0(L_Y(n-1))\le h^0(L_Y(n))\) since there exists an injection \(L_Y(n-1)\longrightarrow L_Y(n)\), and \(h^0(L_Y(-a_1-1))=a_2\binom{-a_1}{1}+a_1\) since \(h^1(L_Y(-a_1-1))=0\); thus, for \(n\le -a_1-2\), we have \(h^1(L(n))-h^1(L(n+1))\le -a_2a_1+a_1\). As a consequence, we conclude that for every \(n\),
\sgaStableTarget{sga6.e13.eq.2-6-1}\[
h^1(L(n))\le h^1(L(-a_1-2))+(-a_1-1)(-a_2a_1+a_1).
\tag{2.7.1}
\]

%%%%%%%%%%  scan idx 641  |  folio 628  %%%%%%%%%%
\SGAPrintPage{628}%
Since \(L_Y\) is \([-a_1]\)-regular, \(h^2(L(-a_1-2))=0\) by 1.4; hence \(h^1(L(-a_1-2))=h^0(L(-a_1-2))-\chi(L(-a_1-2))\). Now \(h^0(L(-a_1-2))\le a_2\binom{-a_1}{2}\) by 1.8 (i), since \(L\) is a \((0,0,a_2)\)-sheaf by 2.5. Consequently,
\sgaStableTarget{sga6.e13.eq.2-6-2}\[
0\le h^1(L(-a_1-2))\le -a_1(-a_1-1)-a_0.
\tag{2.7.2}
\]
Finally, 2.7.2 gives (ii), and 2.7.1 and 2.7.2 give (i).

\sgaStableTarget{sga6.e13.lem.2-8}\textbf{Lemma 2.8.} \uline{We suppose that} \(r=2\) \uline{and} \(d(L)\le0\). \uline{Let} \(\delta=d(L)\) \uline{and} \(\beta=\delta(e_2-1)+e_1\). \uline{Then}:
\[
a_0\ge
\langle c_1(L)^2\rangle+2e_0+
\sum_{i=1}^{2}(2e_i-a_i)\binom{\delta-1+i}{i}
-\beta(\beta+1).
\]

Indeed, let \(\chi\bigl(L^{\otimes m}(\delta-1)\bigr)=\alpha_2m^2+\alpha_1m+\alpha_0\); then, by comparing polynomials, it follows that \(2\alpha_2=\langle c_1(L)^2\rangle\), that \(\alpha_0=\sum_{i=0}^{2}e_i\binom{\delta-1+i}{i}\) and that \(\alpha_2+\alpha_1+\alpha_0=\sum_{i=0}^{2}a_i\binom{\delta-1+i}{i}\). Let \(\chi\bigl(L^{-1}(\delta)(n)\bigr)=\sum_{i=0}^{2}\beta_i\binom{n+i}{i}\); then \(\beta_0=\alpha_2-\alpha_1+\alpha_0\) and, by 2.3, \(\beta_1=d(L^{-1}(\delta))+e_1\), which equals \(\beta\). Thus \(a_0=2\alpha_2+2\alpha_0-\sum_{i=1}^{2}a_i\binom{\delta-1+i}{i}-\beta_0\), and since \(d(L^{-1}(\delta))\le0\), we have \(\beta_0\le\beta(\beta+1)\) by 2.7 (ii); this gives the assertion.

\sgaStableTarget{sga6.e13.lem.2-9}\textbf{Lemma 2.9.} \uline{We suppose that} \(r\ge2\), \uline{that} \(d(L)\le0\) \uline{and that} \(X\) \uline{is normal. Then}:
\begin{enumerate}[label=(\roman*)]
\item \(h^1(L(-n))=0\) \uline{for} \(n\ge2+h_2(a_{r-2},a_{r-1},a_r)\),
\item \(h^1(L(n))\le h_r(a_{r-2},a_{r-1},a_r)\) \uline{for every} \(n\),
\item \(|a_i|\le g_{r-i}(a_{r-2},a_{r-1},a_r)\) \uline{for} \(i=0,\ldots,r\),
\end{enumerate}
\uline{where} \(h_2(X_0,X_1,X_2),h_3(X_0,X_1,X_2),\ldots\) \uline{and} \(g_0(X_0,X_1,X_2),g_1(X_0,X_1,X_2),\ldots\) \uline{are universal polynomials with coefficients in} \(\mathbb Q\), \uline{which depend on neither} \(k\) \uline{nor} \(X\).

%%%%%%%%%%  scan idx 642  |  folio 629  %%%%%%%%%%
\SGAPrintPage{629}%
Indeed, let us first treat the case \(r=2\). Put \(g_i(X_0,X_1,X_2)=X_{2-i}^2\) for \(i=0,1,2\) and \(h_2(X_0,X_1,X_2)=X_1(X_1+1)X_2-X_0\). Then (iii) is evident and (ii) is 2.7 (i).

Let \(X\) lie in \(P=\mathbb P_k^N\), and let \(\omega=\underline{\Ext}^{N-2}_{{O}_P}({O}_X,\Omega^N_{P/k})\). Then \(h^1(L(-n))=h^1(\omega\otimes L^{-1}(n))\) and \(\omega\) is invertible outside the finite set \(Z\) of singular points of \(X\) (see FGA 149 \S 6). Let \(Y\) be a normal hyperplane section of \(X\) which does not contain \(Z\). Then \(\omega_Y(1)\simeq\Omega^1_{Y/k}\) and we have an exact sequence:
\[
0\longrightarrow \omega\otimes L^{-1}(-1)
\longrightarrow \omega\otimes L^{-1}
\longrightarrow \Omega^1_{Y/k}\otimes L_Y^{-1}(-1)
\longrightarrow0.
\]
Now \(h^1(\Omega^1_{Y/k}\otimes L_Y^{-1}(n))=h^0(L_Y(-n))\) for every \(n\), and \(h^0(L_Y(-n))=0\) for \(n\ge1\) by 2.5. Thus \(\Omega^1_{Y/k}\otimes L_Y^{-1}(-1)\) is 3-regular and, by 1.4, \(\omega\otimes L^{-1}\) is \(\bigl[3+h^1(\omega\otimes L^{-1}(2))\bigr]\)-regular. Consequently, \(h^1(L(-n))=0\) for \(n\ge2+h^1(L(-2))\)\,; hence (i).

Now suppose that \(r\ge3\) and argue by induction on \(r\). Let \(Y\) be a normal hyperplane section of \(X\) and consider the exact sequence:
\sgaStableTarget{sga6.e13.eq.2-9-1}\[
H^0(L_Y(-n))\longrightarrow H^1(L(-n-1))
\longrightarrow H^1(L(-n))\longrightarrow H^1(L_Y(-n)).
\tag{2.9.1}
\]
By 1.7, \(\chi(L_Y(n))=\sum_{i=0}^{r-1}a_{i+1}\binom{n+i}{i}\). Thus, for \(n\ge2+h_2(a_{r-2},a_{r-1},a_r)\), the induction hypothesis gives \(h^1(L_Y(-n))=0\), and since \(n\ge2\) by (ii) for \(r=2\), 2.4 gives \(h^0(L_Y(-n))=0\)\,; hence \(h^1(L(-n-1))=h^1(L(-n))\). But, for \(n\gg0\), \(h^1(L(-n))=0\) (see FGA 149 \S 6)\,; hence (i).

As induction hypothesis, suppose that the polynomials \(g_i\) and \(h_i\) are defined for \(i\le r-1\). Then, for every \(n\), \(h^1(L(n))\le h^1(L(n-1))+h_{r-1}(a_{r-2},a_{r-1},a_r)\) by 2.9.1. Thus, by (i),
%%%%%%%%%%  scan idx 643  |  folio 630  %%%%%%%%%%
\SGAPrintPage{630}%
already proved, we conclude that for any \(m\ge -2\),
\sgaStableTarget{sga6.e13.eq.2-9-2}\[
h^1(L(n))\le
\bigl[2+h_2(a_{r-2},a_{r-1},a_r)+m\bigr]
h_{r-1}(a_{r-2},a_{r-1},a_r)
\tag{2.9.2}
\]
for every \(n\le m\).

By 2.5, \(L_Y\) is a \((0,\ldots,0,a_r)\)-sheaf. Put \(c_i=g_{r-i}(a_{r-2},a_{r-1},a_r)\) for \(i=1,\ldots,r\) and \(m=P_{r-2}(c_1,\ldots,c_{r-1})\). By 1.11, \(m\ge0\) and \(L_Y\) is \(m\)-regular. Thus, by 1.4 (ii), \(h^1(L(m-2))\ge h^1(L(n))\) for every \(n\ge m-1\). Therefore, if we define polynomials by \(M=P_{r-2}(g_{r-1},\ldots,g_1)\) and \(h_r=(h_2+M)h_{r-1}\), then we have (ii) by 2.9.2.

Finally, since \(L_Y\) is \(m\)-regular, \(h^q(L(m))=0\) for \(q\ge2\) by 1.4 (i)\,; hence \(\sum_{i=0}^{r}a_i\binom{m+i}{i}=h^0(L(m))-h^1(L(m))\). But \(h^0(L(m))\le a_r\binom{m+r}{r}\) by 1.8 (i). Thus, if we define \(g_r=\sum_{i=1}^{r}g_{r-i}\binom{M+i}{i}+a_r\binom{M+r}{r}+h_r\), we have (iii).

\sgaStableTarget{sga6.e13.rem.2-10}\textbf{Remark 2.10.} For fixed \(m\), consider an equality of polynomials in \(n\):
\[
\sum_{i=0}^{r}a_{m,i}\binom{n+i}{i}
=
\sum_{j=0}^{r}a_j\binom{m+n+j}{j}.
\]
Then one easily finds the formula
\[
a_{m,i}=\sum_{j=0}^{r-i}a_{j+i}\binom{m-1+j}{j}
\]
by induction on \(r\) so that \(a_{m,i}\) becomes \(a_{m,0}\), and by taking \(n=-1\).

Similarly, one finds that there exists a polynomial with rational coefficients \(P_i^{(m,r)}(X_r,\ldots,X_i)\) which is increasing in \(X_i\), such that in the following equality:
\[
\sum_{i=0}^{r}a_i^{(m)}\binom{n+i}{i}
=
\sum_{j=0}^{r}a_j\binom{mn+j}{j}
\]
%%%%%%%%%%  scan idx 644  |  folio 631  %%%%%%%%%%
\SGAPrintPage{631}%
we have: \(a_i^{(m)}=P_i^{(m,r)}(a_r,\ldots,a_i)\).

\sgaStableTarget{sga6.e13.lem.2-11}\textbf{Lemma 2.11.} \uline{Let} \(S\) \uline{be a noetherian scheme and} \(X\) \uline{a projective} \(S\)\uline{-scheme, endowed with a very ample sheaf} \(H={O}_X(1)\). \uline{We suppose that the fibers of} \(X/S\) \uline{are geometrically integral, of the same dimension} \(r\). \uline{Let} \(\mathcal L\) \uline{be a family of classes of invertible sheaves} \(L_K\) \uline{on the fibers of} \(X/S\). \uline{Then for} \(\mathcal L\) \uline{to be bounded, it is (necessary and) sufficient that, in the Hilbert polynomials} \(\chi(L_K(n))=\sum_{i=0}^{r}a_i\binom{n+i}{i}\) \uline{of the} \(L_K\in\mathcal L\), \uline{all the coefficients} \(a_i\) \uline{remain bounded. Moreover, if the fibers of} \(X/S\) \uline{are also geometrically normal, then it suffices that the three coefficients} \(a_r,a_{r-1},a_{r-2}\) \uline{remain bounded}.

Indeed, let \(\delta=\sup\{d(L_K)\}\). If \(\delta>0\), one may replace \(\mathcal L\) by the family of classes of the \(L_K(-\delta)\) by 2.3 and 2.10\,; one may thus suppose \(\delta\le0\). Then \(\mathcal L\) is a \((b)\)-family by 2.5, hence bounded if all the coefficients \(a_i\) remain bounded, by 1.13, and, when the fibers of \(X/S\) are geometrically normal, it suffices for this that the three coefficients \(a_r,a_{r-1},a_{r-2}\) remain bounded, by 2.9.

\sgaStableTarget{sga6.e13.sec.3}\subsection*{3. General finiteness theorems}
\sgaStableTarget{sga6.e13.lem.3-1}\textbf{Lemma 3.1.} \uline{Let} \(S\) \uline{be a scheme and} \(X\) \uline{a quasi-compact and quasi-separated} \(S\)\uline{-scheme\,; let} \(\underline{\Pic}_{X/S}\) \uline{be the Picard functor, (localized for the f.p.p.f. topology})\SGAAuthorFootnote{*}{Cf. SGA 3 IV 6.3, and FGA no.\ 232 for the definition of \(\underline{\Pic}_{X/S}\), where, as indicated, the fpqc topology should here be replaced by the fppf topology.}.
\uline{Then for every filtered inverse system} \((Z_\alpha)\) \uline{of quasi-compact and quasi-separated} \(S\)\uline{-schemes}
%%%%%%%%%%  scan idx 645  |  folio 632  %%%%%%%%%%
\SGAPrintPage{632}%
\uline{with affine transition morphisms, the canonical map}
\sgaStableTarget{sga6.e13.eq.3-1-1}\[
\varinjlim \underline{\Pic}_{X/S}(Z_\alpha)\longrightarrow
\underline{\Pic}_{X/S}\bigl(\varprojlim Z_\alpha\bigr)
\tag{3.1.1}
\]
\uline{is bijective}.

\uline{Consequently} (EGA IV 8.14.2), \uline{if the Picard functor is represented by an} \(S\)\uline{-scheme} \(\Pic_{X/S}\), \uline{then} \(\Pic_{X/S}\) \uline{is locally of finite presentation}.

Indeed, an element \(\lambda_\alpha\) of \(\underline{\Pic}_{X/S}(Z_\alpha)\) is represented by an invertible sheaf \(L'_\alpha\) on an \(X_{Z'_\alpha}\) where \(Z'_\alpha\to Z_\alpha\) is an f.p.p.f. morphism. Put \(Z=\varprojlim Z_\alpha\). The image \(\lambda\) of \(\lambda_\alpha\) in \(\underline{\Pic}_{X/S}(Z)\) is represented by the invertible sheaf \(L'\) on \(X_{Z'}\), the pullback of \(L'_\alpha\), where \(Z'=Z'_\alpha\times_{Z_\alpha}Z\).

If \(\lambda\) is the identity element of \(\underline{\Pic}_{X/S}(Z)\), there exists an f.p.p.f. morphism \(Z''\to Z'\) such that the pullback of \(L'\) to \(X_{Z''}\) is \(\simeq{O}_{X_{Z''}}\). By (EGA IV 8.8.2, 8.10.5 (vi), 11.2.6 (ii), 8.5.2 (i), and 8.5.2.4), there exist a \(\beta\geq\alpha\) and an f.p.p.f. morphism \(Z''_\beta\to Z'_\beta\), where \(Z'_\beta=Z'_\alpha\times_{Z_\alpha}Z_\beta\), such that the pullback of \(L'_\alpha\) to \(X_{Z''_\beta}\) is \(\simeq{O}_{X_{Z''_\beta}}\). The image of \(\lambda_\alpha\) is therefore already the identity element in \(\underline{\Pic}_{X/S}(Z_\beta)\). Consequently, 3.1.1 is injective.

The surjectivity of 3.1.1 is proved in an analogous way.

\sgaStableTarget{sga6.e13.prop.3-2}\textbf{Proposition 3.2.} \uline{Let} \(S\) \uline{be a noetherian scheme and} \(X\) \uline{a proper} \(S\)\uline{-scheme. We suppose that the Picard functor is represented by an} \(S\)\uline{-scheme} \(\Pic_{X/S}\). \uline{Then}~:
\begin{enumerate}[label=(\roman*)]
\item \(\Pic_{X/S}\) \uline{is locally of finite type over} \(S\).
%%%%%%%%%%  scan idx 646  |  folio 633  %%%%%%%%%%
\item\SGAPrintPage{633}%
\uline{The subsets} \(\Lambda\) \uline{of} \(\Pic_{X/S}\) \uline{correspond bijectively to the families} \(\mathscr L\) \uline{of classes of invertible sheaves on the fibers of} \(X/S\).
\item \uline{A family} \(\Lambda\) \uline{of} \(\Pic_{X/S}\) \uline{is quasi-compact if and only if the corresponding family} \(\mathscr L\) \uline{is bounded}.
\end{enumerate}

Indeed, (i) follows from (3.1)\,; (ii), from the following fact: if \(s\) is a point of \(S\) and \(K\) is an algebraically closed extension of \(k(s)\), then the \(K\)-points of \(\Pic_{X/S}\) correspond bijectively to the isomorphism classes of invertible sheaves on \(X_K\).

As for (iii), suppose that the subset \(\Lambda\) of \(\Pic_{X/S}\) is quasi-compact. Then \(\Lambda\) is contained in an open subset \(U\) of finite type, and the inclusion \(U\to\Pic_{X/S}\) corresponds to an invertible sheaf \(L\) on an \(X_T\), where \(T\to U\) is a surjective, (flat) morphism of finite type. Finally, \(T\to S\) is of finite type, so \(L\) bounds \(\mathscr L\).

Conversely, suppose that \(\mathscr L\) is bounded by a coherent sheaf \(L\) on an \(X_T\), with \(T\) of finite type over \(S\). One may suppose \(L\) invertible (1.13 (i)) so that \(L\) defines a morphism \(f:T\to\Pic_{X/S}\). The image of \(f\) contains \(\Lambda\) and is quasi-compact, hence noetherian by (i). Consequently, \(\Lambda\) is quasi-compact, and we have (iii).

By 3.2, the usual definition of morphisms of finite type between two Picard schemes generalizes to the following:

\sgaStableTarget{sga6.e13.def.3-3}\textbf{Definition 3.3.} \uline{Let} \(X\) \uline{and} \(Y\) \uline{be proper over noetherian} \(S\). \uline{A morphism between the Picard functors}
\[
\begin{tikzcd}
\underline{\Pic}_{Y/S}\arrow[r]&\underline{\Pic}_{X/S}
\end{tikzcd}
\]
%%%%%%%%%%  scan idx 647  |  folio 634  %%%%%%%%%%
\SGAPrintPage{634}%
\uline{is said to be} of finite type \uline{if every bounded family of classes of invertible sheaves on the fibers of} \(X/S\) \uline{has as its inverse image a bounded family of classes of sheaves on the fibers of} \(Y/S\).

\sgaStableTarget{sga6.e13.lem.3-4}\textbf{Lemma 3.4.} \uline{Let} \(X,Y,X',Y'\) \uline{be proper over noetherian} \(S\), \uline{and} \(T\to S\) \uline{surjective of finite type. Then}~:
\begin{enumerate}[label=(\roman*)]
\item \uline{For a morphism} \(\Phi:\underline{\Pic}_{Y/S}\to\underline{\Pic}_{X/S}\) \uline{to be of finite type, it is necessary and sufficient that} \(\Phi_T:\underline{\Pic}_{Y_T/T}\to\underline{\Pic}_{X_T/T}\) \uline{be of finite type}.
\item \uline{In a commutative diagram}
\[
\begin{tikzcd}
\Phi':\underline{\Pic}_{Y'/S}\arrow[r]&
\underline{\Pic}_{X'/S}\\
\Phi:\underline{\Pic}_{Y/S}\arrow[r]\arrow[u,"\Phi_Y"]&
\underline{\Pic}_{X/S}\arrow[u,"\Phi_X"]
\end{tikzcd}
\]
\uline{with} \(\Phi'\) \uline{and} \(\Phi_Y\) \uline{of finite type}, \(\Phi\) \uline{is of finite type}.
\end{enumerate}

Indeed, suppose that \(\Phi_T\) is of finite type. Let \(\mathscr L\) be a family on the fibers of \(X/S\), bounded by an \(L\) on an \(X_R\). The family \(\mathscr L_T=\mathscr L\times_S T\) on the fibers of \(X_T/T\), defined in the evident way, is bounded by the pullback of \(L\) to \(X_{R\times T}\). If \(\Phi_T^{-1}(\mathscr L_T)\) is bounded by an \(M\) on a \(Y_Q\), then \(M\) also bounds \(\Phi^{-1}(\mathscr L)\).

The converse of (i) and assertion (ii) are proved in an analogous way.

\sgaStableTarget{sga6.e13.thm.3-5}\textbf{Theorem 3.5.} \uline{Let} \(S\) \uline{be a noetherian scheme}, \(X\) \uline{and} \(Y\) \uline{two proper} \(S\)\uline{-schemes and} \(f:X\to Y\) \uline{a morphism. If} \(f\) \uline{is surjective, then the induced morphism between the Picard functors}
\[
f^*:\underline{\Pic}_{Y/S}\longrightarrow\underline{\Pic}_{X/S}
\]
%%%%%%%%%%  scan idx 648  |  folio 635  %%%%%%%%%%
\SGAPrintPage{635}%
\uline{is of finite type}.

Indeed, by 3.4 (i), one may evidently suppose \(S\) reduced and irreducible. Then there exists a nonempty open subset \(U\) of \(S\) such that \(\underline{\Pic}_{Y/S}\) and \(\underline{\Pic}_{X/S}\) are representable and the restriction of \(f^*\) is quasi-affine (XII 1.1), hence of finite type by 3.2 (i). Endowing \(T=S-U\) with the induced reduced structure, one may suppose that the morphism \(\underline{\Pic}_{Y_T/T}\longrightarrow\underline{\Pic}_{X_T/T}\) induced by \(f^*\) is of finite type by noetherian induction\,; this gives the assertion, by 3.4 (i) applied to \(U\amalg T\to S\).

\sgaStableTarget{sga6.e13.thm.3-6}\textbf{Theorem 3.6.} \uline{Let} \(S\) \uline{be a noetherian scheme}, \(X\) \uline{a proper} \(S\)\uline{-scheme and} \(\underline{\Pic}_{X/S}\) \uline{the Picard functor. Then, for every integer} \(n\ne0\), \uline{the} \(n\)\uline{-th power homomorphism}
\[
\Phi_n:\underline{\Pic}_{X/S}\longrightarrow\underline{\Pic}_{X/S}
\]
\uline{is a morphism of finite type}.

Indeed, let \(f:X'\to X\) be a Chow presentation: \(X'\) is projective over \(S\) and \(f\) is surjective. If the theorem is admitted for \(X'\), then the theorem for \(X\) follows by 3.5 and 3.4 (ii). One may therefore suppose \(X/S\) projective.

There exist a nonempty open subset \(U\) of \(S\), an integral scheme \(S'\) and a finite surjective morphism \(S'\to U\), such that every irreducible component \(X'_i\) of \(X_{S'}\), endowed with the induced reduced structure, is integral over \(S'\) (XII 1.3). Endowing \(T=S-U\) with the induced reduced structure, one may admit the theorem for \(X_T/T\) by noetherian induction. If it is admitted for all the \(X'_i/S'\), it follows for \(\coprod X'_i/S'\) by the
%%%%%%%%%%  scan idx 649  |  folio 636  %%%%%%%%%%
\SGAPrintPage{636}%
following easy lemma, hence for \(X_{S'}/S'\) by 3.5 and 3.4 (ii), and finally for \(X/S\) by 3.4 (i).

\sgaStableTarget{sga6.e13.lem.3-7}\textbf{Lemma 3.7.} \uline{Let} \(X=\coprod X_i\) \uline{be of finite type over noetherian} \(S\), \uline{and} \(\mathcal F\) \uline{a family of classes of coherent sheaves on the fibers of} \(X/S\). \uline{For} \(\mathcal F\) \uline{to be bounded, it is necessary and sufficient that for every} \(i\) \uline{the induced family} \(\mathcal F_i\) \uline{on the fibers of} \(X_i/S\) \uline{be bounded}.

In the proof of 3.6, we reduced to the case where \(X\) is projective and integral over \(S\)\,; endow \(X/S\) with a very ample sheaf \({O}_X(1)\). Let \(\mathcal L\) be a bounded family of classes of invertible sheaves on the fibers of \(X/S\), and \(L\) an invertible sheaf on an \(X_T\), with \(T\) of finite type over \(S\), which bounds \(\mathcal L\). After decomposing \(T\) into pieces, one may suppose \(L\) flat over \(T\), and after replacing \(T\) by each of its connected components, one may suppose \(T\) connected. Then the polynomial in \(q\) and \(m\), \(P(q,m)=\chi(L_t^{\otimes q}(m))\), does not depend on \(t\in T\) (EGA III 7.9.4).

Let \(\mathcal L_n=\Phi_n^{-1}(\mathcal L)\): \(\mathcal L_n\) is the family of classes of invertible sheaves \(M_K\) such that the class of \(M_K^{\otimes n}\) belongs to \(\mathcal L\). Let \(\chi(M_K^{\otimes p}(m))=\sum_{i=0}^{r}a_i(p)\binom{m+i}{i}\), where \(a_i(p)\) is a polynomial in \(p\) of degree \(\le r-i\). For variable \(q,m\), we have \(\chi(M_K^{\otimes qn}(m))=P(q,m)\), so each \(a_i(qn)\) is independent of the chosen \(M_K\)\,; consequently, each polynomial \(a_i(p)\) is also independent of \(M_K\). Therefore the family \(\mathcal L_n\) has a single Hilbert polynomial \(\chi(M_K(m))\), hence is bounded by 2.11.

\sgaStableTarget{sga6.e13.thm.3-8}\textbf{Theorem 3.8.} \uline{Let} \(S\) \uline{be a noetherian scheme}, \(Y\) \uline{a projective} \(S\)\uline{-scheme endowed with a relatively ample invertible sheaf} \({O}_Y(1)\), \(X\) \uline{the zero scheme of a section} \(\varphi\) \uline{of} \({O}_Y(1)\), \uline{and} \(f:X\to Y\) \uline{the inclusion morphism. We}
%%%%%%%%%%  scan idx 650  |  folio 637  %%%%%%%%%%
\SGAPrintPage{637}%
\uline{suppose that for every point} \(s\in S\) \uline{and every irreducible component} \(Y_s^i\) \uline{of} \(Y_s\), \uline{the irreducible components of} \(X_s\cap Y_s^i=V(\varphi|Y_s^i)\) \uline{have dimension} \(\ge2\), (\uline{which is the case if the irreducible components of the fibers of} \(Y\) \uline{have dimension} \(\ge3\)). \uline{Then the induced morphism between the Picard functors}
\[
f^*:\underline{\Pic}_{Y/S}\longrightarrow\underline{\Pic}_{X/S}
\]
\uline{is of finite type}.

\uline{Moreover, if} \(S\) \uline{is integral, there exists a nonempty open subset} \(U\) \uline{of} \(S\) \uline{such that the Picard functors restricted to} \(U\) \uline{are representable and the induced morphism}
\[
\underline{\Pic}_{Y/S}|U\longrightarrow\underline{\Pic}_{X/S}|U
\]
\uline{is quasi-affine (of finite type}).

Let us first prove the first assertion. As in the proof of 3.6, one sees by noetherian induction that one may replace \(S\) by any nonempty open subset, and one easily reduces to the case where \(S\) is integral and \(Y\) is integral over \(S\), always replacing \(X\) by \(X\times_YY'\) when \(Y\) is replaced by \(Y'\)\,; this done, one reduces in the same way to the case where \(Y\) is normal over \(S\), using the following lemma in place of XII 1.3.

\sgaStableTarget{sga6.e13.lem.3-9}\textbf{Lemma 3.9.} \uline{Let} \(S\) \uline{be an integral noetherian scheme and} \(X\) \uline{an} \(S\)\uline{-scheme, integral over} \(S\). \uline{Then there exist a nonempty open subset} \(U\) \uline{of} \(S\), \uline{an integral scheme} \(S'\) \uline{finite and faithfully flat over} \(U\), \uline{an} \(S'\)\uline{-scheme} \(X'\) \uline{normal and integral over} \(S'\), \uline{and a finite surjective} \(S'\)\uline{-morphism} \(X'\to X_{S'}\).

%%%%%%%%%%  scan idx 651  |  folio 638  %%%%%%%%%%
\SGAPrintPage{638}%
Indeed, let \(\eta\) be the generic point of \(S\) and \(K\) the algebraic closure of \(k(\eta)\). It is clear that \(\Spec(K)\) is the limit of the filtered inverse system \((S_\alpha)\) of integral affine schemes endowed with a finite faithfully flat morphism \(S_\alpha\to U_\alpha\), where \(U_\alpha\) is a nonempty open subset of \(S\). The normalization of \(X_K\) therefore comes from an \(S_\alpha\)-scheme \(X'\) for some \(\alpha\), and the finite surjective morphism \(X'_K\to X_K\) from a finite surjective \(S_\alpha\)-morphism \(X'\to X_{S_\alpha}\) (EGA IV 8.8.2, 8.10.5). Finally, since \(X'_K\) is normal and integral, there exists an open subset \(S'\) of \(S_\alpha\) such that \(X'_{S'}\) is normal and integral over \(S'\) (EGA IV 9.9.5).

\sgaStableTarget{sga6.e13.rem.3-10}\textbf{Remark 3.10.} Retaining the notation and conditions of 3.9, suppose that there exists a resolving morphism \(X'_K\to X_K\) (EGA IV 7.9.1), (for example if \(\dim(X_K)\le2\), by Abhyankar). Then a similar argument shows that there exist a finite faithfully flat morphism \(S'\to U\), a scheme \(X'\) smooth over \(S'\) and a proper surjective birational morphism \(X'\to X_{S'}\).

In the proof of the first part of 3.8, we reduced to the case where \(S\) is integral and \(Y\) is normal and integral over \(S\), so \(Y\) is integral. If \(\varphi=0\), then \(X=Y\) and the assertion is trivial. Otherwise, after restricting \(S\), one may suppose that for every \(s\in S\) we have \(\varphi_s\ne0\), so \(\varphi_s\) is \({O}_{Y_s}\)-regular and the dimension \(r\) of \(Y_s\) is \(\ge3\).

Now let \(\mathcal L\) be a bounded family of classes of invertible sheaves on the fibers of \(X/S\). Let \(M_K\) be an invertible sheaf on a geometric fiber \(Y_K\) of \(Y/S\), with Hilbert polynomial \(\chi(M_K(n))=\sum_{i=0}^{r}a_i\binom{n+i}{i}\); then the Hilbert polynomial of \(f^*M_K\) is \(\sum_{i=0}^{r-1}a_{i+1}\binom{n+i}{i}\). As the class of \(M_K\) ranges over \((f^*)^{-1}(\mathcal L)\), that is, the
%%%%%%%%%%  scan idx 652  |  folio 639  %%%%%%%%%%
\SGAPrintPage{639}%
class of \(f^*M_K\) belongs to \(\mathcal L\), the three coefficients \(a_r,a_{r-1},a_{r-2}\) remain bounded by 1.13.

Let \(H={O}_Y(m)\) be a multiple of \({O}_Y(1)\) which is very ample. Then, in the Hilbert polynomials \(\chi(M_K(mn))=\sum_{i=0}^{r}a_i^{(m)}\binom{n+i}{i}\) relative to \(H\) of the \(M_K\) in \((f^*)^{-1}(\mathcal L)\), the three coefficients \(a_r^{(m)},a_{r-1}^{(m)},a_{r-2}^{(m)}\) remain bounded by 2.10. Consequently, \((f^*)^{-1}(\mathcal L)\) is bounded by 2.11\,; this gives the first part of 3.8.

As for the second part of 3.8, the ``generic'' representability of the Picard functors was established in XII 1.2, as was their separatedness. By an argument similar to that of XII 2, using 3.9 and XII 1.1, one reduces to the case of \(Y\) normal over \(S\). Now, for the morphism \(\underline{\Pic}_{Y/S}\to\underline{\Pic}_{X/S}\) to be quasi-affine, it suffices that it be (separated and) quasi-finite (EGA IV 8.12.8), which can be checked on the geometric fibers, and thus follows from the following lemma:

\sgaStableTarget{sga6.e13.lem.3-11}\textbf{Lemma 3.11.} \uline{Let} \(k\) \uline{be an algebraically closed field}, \(Y\) \uline{a normal projective} \(k\)\uline{-scheme endowed with an ample sheaf} \({O}_Y(1)\), \uline{and} \(X\) \uline{the zero scheme of a section} \(\varphi\) \uline{of} \({O}_Y(1)\). \uline{We suppose that} \(Y\) \uline{has dimension} \(\ge3\) (resp. \(\ge2\)). \uline{Then the morphism} \(\Pic_{Y/k}\longrightarrow\Pic_{X/k}\) (resp. \uline{the morphism} \(\Pic^\tau_{Y/k}\to\Pic^\tau_{X/k}\) (\uline{see} 4.1)) \uline{is of finite type and has a kernel} \(N\) \uline{which is a finite group (in general nonreduced}).

Indeed, one may evidently suppose \(\varphi\ne0\)\,; thus \(\varphi\) is regular, \(Y\) being integral. The morphism under consideration is of finite type by the first part of 3.8 (resp. by 4.7 (iii)). Consequently, the kernel \(N\) is an algebraic group, hence contained in \(\Pic^\tau_{Y/k}\), and is therefore the same
%%%%%%%%%%  scan idx 653  |  folio 640  %%%%%%%%%%
\SGAPrintPage{640}%
in the two cases under consideration. Since \(\Pic^\tau_{Y/k}\) is proper, \(Y\) being normal (FGA 236-2.13), and since \(N\) is a subgroup, necessarily closed (SGA 3 VI\(_B\) 1.9.2), for \(N\) to be finite it suffices that it contain no subgroup of type \(\mu_\ell(=\mathbb Z/\ell\mathbb Z)\) where \((\ell,\operatorname{car}(k))=1\).

Consider Kummer theory, that is, the following commutative diagram, coming from the exact sequence
\[
0\longrightarrow\mu_\ell\longrightarrow\mathbb G_m
\xrightarrow{x\mapsto x^\ell}\mathbb G_m\longrightarrow0
\]
of étale sheaves on \(Y\) (resp. on \(X\)):
\[
\begin{tikzcd}[column sep=1.6em,row sep=1.2em]
& & H^1(Y,\mu_\ell)\arrow[r]\arrow[d]&H^1(Y,\mathbb G_m)\arrow[r,"\ell"]\arrow[d]&H^1(Y,\mathbb G_m)\arrow[d]\\
H^0(X,\mathbb G_m)\arrow[r,"\ell"]&H^0(X,\mathbb G_m)\arrow[r]&H^1(X,\mu_\ell)\arrow[r]&H^1(X,\mathbb G_m)\arrow[r,"\ell"]&H^1(X,\mathbb G_m).
\end{tikzcd}
\]
By Hilbert's ``Theorem 90'', we have \(H^1(Y,\mathbb G_m)=\Pic_{Y/k}(k)\) (resp. \(H^1(X,\mathbb G_m)=\Pic_{X/k}(k)\)). Since \(X\) is proper and integral, we have \(H^0(X,\mathbb G_m)=k^*\)\,; therefore \(H^1(X,\mu_\ell)\to H^1(X,\mathbb G_m)\) is injective.

It therefore suffices to see that \(H^1(Y,\mathbb Z/\ell\mathbb Z)\longrightarrow H^1(X,\mathbb Z/\ell\mathbb Z)\) is injective, which follows from \(\pi_1(X,\xi)\to\pi_1(Y,\xi)\) being surjective, or again (SGA 1 V 6.9) from every connected étale covering \(\widetilde Y\) of \(Y\) with group \(\mathbb Z/\ell\mathbb Z\) having a restriction \(\widetilde X=X\times_Y\widetilde Y\) which is also connected; this is a well-known consequence of Bertini's theorem and Zariski's connectedness theorem (SGA 1 X 2.10).

\sgaStableTarget{sga6.e13.rem.3-12}\textbf{Remark 3.12.} Retain the conditions of 3.11. Then:
\begin{enumerate}[label=(\roman*)]
\item The kernel \(N\) is a finite \uline{unipotent} group (SGA 3 XVII 1.1), since it contains no subgroup of multiplicative type \(\mu_n\) for any \(n\) (SGA 3 XVII 4.6.1 vi), which is proved in the same way, using instead
%%%%%%%%%%  scan idx 654  |  folio 641  %%%%%%%%%%
\SGAPrintPage{641}%
of Kummer theory the following fact, due to Grothendieck (SGA 1 XI, on the last page): if \(X\) is a proper scheme over a field \(k\), with \(H^0(X,{O}_X)=k\), then there exists an isomorphism of \(k\)-group schemes \(\underline{\Hom}_{\mathrm{gr}}(\mu_n,\underline{\Pic}_{X/k})\simeq \underline{H}^1(X,\mathbb Z/n\mathbb Z)\); consequently, if \(k\) is algebraically closed, we have \(\Hom_{\mathrm{gr}}(\mu_n,\underline{\Pic}_{X/k})\simeq H^1(X,\mathbb Z/n\mathbb Z)\).
\item If \(k\) has characteristic zero, then the kernel \(N\) is zero, since every finite unipotent group is then zero (SGA 3 XVII 1.7). One thus approximately recovers Kodaira's ``Vanishing theorem'' in the form given by Mumford [2].
\end{enumerate}

\sgaStableTarget{sga6.e13.thm.3-13}\textbf{Theorem 3.13.} \uline{Let} \(S\) \uline{be a noetherian scheme and} \(X\) \uline{a projective} \(S\)\uline{-scheme, endowed with a relatively ample invertible sheaf} \(H={O}_X(1)\). \uline{We suppose that the geometric fibers of} \(X/S\) \uline{are integral and of dimension} \(r\). \uline{Then, for a family} \(\mathcal L\) \uline{of classes of invertible sheaves on the fibers of} \(X/S\), \uline{the following conditions are equivalent}~:
\begin{enumerate}[label=(\roman*)]
\item \(\mathcal L\) \uline{is a bounded family}.
\item \uline{In the Hilbert polynomials} \(\chi(L_K(n))=\sum_{i=0}^{r}a_i\binom{n+i}{i}\) \uline{of the} \(L_K\in\mathcal L\), \uline{the coefficient} \(a_{r-1}\) \uline{remains bounded and the coefficient} \(a_{r-2}\) \uline{remains bounded below}.
\item \uline{As} \(L_K\) \uline{ranges over} \(\mathcal L\), \uline{the degree} \(d(L_K)=\langle c_1(L_K)\cdot c_1(H_K)^{r-1}\rangle\) \uline{remains bounded and the intersection number} \(\langle c_1(L_K)^2\cdot c_1(H_K)^{r-2}\rangle\) \uline{remains bounded below}.
\item \uline{There exists an integer} \(N\) \uline{such that we have}
\[
H^{-N}<\mathcal L<H^N
\]
\uline{in the sense of the order relation defined by ampleness\,; that is, for every} \(L_K\in\mathcal L\), \(H_K^{\otimes N}\otimes L_K\) \uline{and} \(H_K^{\otimes N}\otimes L_K^{-1}\) \uline{are ample}.
\end{enumerate}

%%%%%%%%%%  scan idx 655  |  folio 642  %%%%%%%%%%
\SGAPrintPage{642}%
Let us prove that (i) implies (ii), (iii), and (iv). Thus suppose that \(\mathcal L\) is bounded by an invertible sheaf on an \(X_T\). Then \(\mathcal L\) satisfies (iv) by EGA II 4.6.12. Suppose \(L\) flat over \(T\). Then the polynomial in \(m\) and \(n\), \(\chi(L_t^{\otimes m}(n))\), depends only on the connected component of \(T\) (EGA III 7.9.4). Consequently, \(\mathcal L\) satisfies (ii) and, taking 2.1 into account, (iii).

Next, note that the family of classes of the \(H_K\) (resp. \({O}_{X_K}\)) is bounded, hence satisfies (iii) (resp. (ii)).

Suppose that \(\mathcal L\) satisfies (iv). Then \(\langle [N\cdot c_1(H_K)+c_1(L_K)]\cdot c_1(H_K)^{r-1}\rangle>0\), hence \(d(L_K)>-Nd(H_K)\), and likewise \(Nd(H_K)>d(L_K)\). Moreover, \(\langle [N c_1(H_K)+c_1(L_K)]^2\cdot c_1(H_K)^{r-2}\rangle>0\), hence \(\langle c_1(L_K)^2\cdot c_1(H_K)^{r-2}\rangle>-2Nd(L_K)-N^2d(H_K)\). Thus \(\mathcal L\) satisfies (iii).

Finally, let us prove that (ii) (resp. (iii)) implies (i). By 2.10 and 2.3, one may replace \(H\) by a multiple\,; therefore suppose \(H\) very ample. Evidently one may decompose \(S\) into pieces\,; thus suppose that there exists a linear section \(Y\) of \(X\) which is integral over \(S\), with fibers of dimension \(2\). By 3.8, one may replace \(X,\mathcal L\) by \(Y,\mathcal L|Y\), (taking account of 1.7 (resp. the projection formula), which implies that condition (ii) (resp. (iii)) on \(X,\mathcal L\) is equivalent to (ii) (resp. (iii)) on \(Y,\mathcal L|Y\))\,; therefore suppose \(r=2\). If \(\delta=\sup\{d(L_K)\}\) is \(>0\), one may replace \(\mathcal L\) by the family of classes of the \(L_K(-\delta)\) by 2.3 and 2.10\,; thus suppose \(\delta\le0\).

If, in the Hilbert polynomials \(\sum_{i=0}^{2}a_i\binom{n+i}{i}\), \(a_1\) remains bounded and \(a_0\) remains bounded below, then the three coefficients remain bounded by 2.3 and 2.7\,; thus \(\mathcal L\) is bounded by 2.11\,; hence (ii) implies
%%%%%%%%%%  scan idx 656  |  folio 643  %%%%%%%%%%
\SGAPrintPage{643}%
(i). On the other hand, (iii) implies (ii) by 2.3 and 2.8.

\sgaStableTarget{sga6.e13.sec.4}\subsection*{4. Finiteness theorems for $\underline{\Pic}^\tau_{X/S}$}

\sgaStableTarget{sga6.e13.para.4-1}4.1. Let \(S\) be a scheme and \(G\) a commutative group functor on \(S\). We denote by \(G^0\) (resp. \(G^\tau\)) the group subfunctor whose \(T\)-points \(f:T\to G\) are those of \(G\) such that, for every algebraically closed field \(K\) and every \(K\)-point \(t\) of \(T\), there exist a connected \(K\)-scheme \(Z\), two \(K\)-points \(z,z_0\) of \(Z\) and a morphism \(g:Z\to G\) such that \(g(z)=f(t)\) and \(g(z_0)=0\) in \(G(K)\) (resp. there exists a nonzero integer \(n\) such that \(n.f(t)\in G^0(K)\)).

It is evident that the formation of \(G^0\) (resp. \(G^\tau\)) commutes with every base change, that a morphism of group functors \(G\to G'\) induces a morphism \(G^0\to G'^0\) (resp. \(G^\tau\to G'^\tau\)), and that if \(S\) is a field and \(G\) is represented by a group scheme locally of finite type, then \(G^0\) is represented by the connected component of the identity element, which is an open (and closed) subscheme, of finite type and geometrically irreducible (SGA 3 VI\(_A\) 2.4).

\sgaStableTarget{sga6.e13.lem.4-2}\textbf{Lemma 4.2.} \uline{Let} \(\Phi:G\to G'\) \uline{be a morphism between two commutative group functors on a scheme} \(S\). \uline{We suppose that, for every algebraically closed field} \(K\) \uline{and every morphism} \(\Spec(K)\to S\), \(G_K\) (resp. \(G'_K\)) \uline{is representable by a group scheme locally of finite type, and} \(\Phi_K\) \uline{is representable by a morphism of finite type, (that is, quasi-compact). Then}
\[
  \Phi^{-1}(G'^\tau)=G^\tau.
\]

%%%%%%%%%%  scan idx 657  |  folio 644  %%%%%%%%%%
\SGAPrintPage{644}%
Indeed, we have already noted the inclusion \(\Phi^{-1}(G'^\tau)\supset G^\tau\). For the reverse inclusion, note that \(g\in G(K)\) belongs to \(G^\tau(K)\) if and only if the sequence of \(ng\) (\(n\in\mathbb N\)) is quasi-compact\,; since \(\Phi_K\) is quasi-compact, it suffices for this that the sequence of \(\Phi(ng)=n\Phi(g)\) have the same property\,; this gives the assertion.

\sgaStableTarget{sga6.e13.lem.4-3}\textbf{Lemma 4.3.} \uline{Let} \(X\), \(Y\) \uline{be proper over a scheme} \(S\) \uline{and} \(f:X\to Y\) \uline{a surjective morphism} (resp. \uline{the inclusion morphism of the zero scheme of a section} \(\varphi\) \uline{of an ample sheaf} \({O}_Y(1)\), \uline{under the hypotheses of} 3.8). \uline{Then}
\[
  (f^*)^{-1}\underline{\Pic}^\tau_{X/S}=\underline{\Pic}^\tau_{Y/S}.
\]

Indeed, the assertion is a formal consequence of 4.2, (XII 1.2), 3.1, and 3.5 (resp. 3.8).

\sgaStableTarget{sga6.e13.para.4-4}4.4. Let \(k\) be an algebraically closed field, \(X\) a proper \(k\)-scheme, and \(L\) an invertible sheaf on \(X\). We say that \(L\) is \uline{algebraically equivalent} to \(0\) (resp. \(\tau\)-\uline{equivalent} to \(0\)) if there exist a connected \(k\)-scheme \(T\) of finite type, an invertible sheaf \(M\) on \(X_T=X\times_k T\), and two \(k\)-points \(t,t_0\) of \(T\) such that \(L\simeq M_t\) and \({O}_X\simeq M_{t_0}\) (resp. if there exists a nonzero integer \(n\) such that \(L^{\otimes n}\) is algebraically equivalent to \(0\)), or, what is easily seen to be equivalent, if the \(k\)-point \(\lambda\) of \(\underline{\Pic}_{X/k}\) corresponding to \(L\) belongs to \(\underline{\Pic}^0_{X/k}(k)\) (resp. to \(\underline{\Pic}^\tau_{X/k}(k)\)).

Recall also that \(L\) is said to be \uline{numerically equivalent} to \(0\) if, for every integral closed curve \(Y\) in \(X\), the intersection number \(\langle c_1(L)\cdot Y\rangle\) is \(0\) (see 2.1 or X 4.5).

%%%%%%%%%%  scan idx 658  |  folio 645  %%%%%%%%%%
\SGAPrintPage{645}%
\sgaStableTarget{sga6.e13.prop.4-5}\textbf{Proposition 4.5. (Functorial behavior).} \uline{Let} \(f:X'\to X\) \uline{be a morphism between two proper} \(k\)\uline{-schemes, and let} \(L\) \uline{be an invertible module on} \(X\). \uline{Then}~:
\begin{enumerate}[label=(\roman*)]
\item \uline{If} \(L\) \uline{is} \(\tau\)\uline{-equivalent to} \(0\) (resp. \uline{numerically equivalent to} \(0\)), \uline{then so is} \(f^*L\).
\item \uline{Conversely, when} \(f\) \uline{is surjective, if} \(f^*L\) \uline{is} \(\tau\)\uline{-equivalent to} \(0\) (resp. \uline{numerically equivalent to} \(0\)), \uline{then so is} \(L\).
\end{enumerate}

Indeed, assertion (i) (resp. (ii)) for \(\tau\)-equivalence follows at once from 4.1 (resp. 4.3). On the other hand, let \(Y'\) be an integral closed curve in \(X'\), and \(Y=f(Y')\) its image, endowed with the induced reduced structure. When \(\dim(Y)=1\), we have the projection formula (X 4.4.4)
\sgaStableTarget{sga6.e13.eq.4-5-1}\[
  \tag*{4.5.1.}
  \langle c_1(f^*L)\cdot Y'\rangle
  =[k(Y'):k(Y)]\,\langle c_1(L)\cdot Y\rangle
\]
where \(k(Y')\) (resp. \(k(Y)\)) is the function field of \(Y'\) (resp. \(Y\)), and when \(\dim(Y)=0\), we have \((f^*L)_{Y'}\simeq{O}_{Y'}\). Consequently, if \(L\) is numerically equivalent to \(0\), then \(\langle c_1(f^*L)\cdot Y'\rangle=0\)\,; hence (i).

Conversely, when \(f\) is surjective, given an integral closed curve \(Y\) in \(X\), one can find an integral closed curve \(Y'\) in \(X'\) such that \(Y=f(Y')\): let \(y\) be the generic point of \(Y\), \(y'\) a point of the fiber \(f^{-1}(y)\) algebraic over \(k(y)\), and \(Y'\) the closure of \(y'\). Then 4.5.1 shows that \(\langle c_1(f^*L)\cdot Y'\rangle=0\) implies \(\langle c_1(L)\cdot Y\rangle=0\)\,; hence (ii).

\sgaStableTarget{sga6.e13.thm.4-6}\textbf{Theorem 4.6. (The numerical characterization of \(\underline{\Pic}^\tau_{X/k}\)).} \uline{Let} \(k\) \uline{be an algebraically closed field}, \(X\) \uline{a proper} \(k\)\uline{-scheme, and} \(L\) \uline{an invertible sheaf on} \(X\). \uline{Then the following conditions are equivalent}~:
\begin{enumerate}[label=(\roman*)]
\item \(L\) \uline{is} \(\tau\)\uline{-equivalent to} \(0\).
\item \uline{For every coherent sheaf} \(F\) \uline{on} \(X\), \uline{we have}
%%%%%%%%%%  scan idx 659  |  folio 646  %%%%%%%%%%
\SGAPrintPage{646}%
\[
  \chi(F\otimes L)=\chi(F)
\]
\item
\(L\) \uline{is numerically equivalent to} \(0\).
\end{enumerate}
\uline{If} \(X/k\) \uline{is projective and equipped with an ample invertible module} \(H={O}_X(1)\), \uline{the preceding conditions are also equivalent to the following}:
\begin{enumerate}[label=(\roman*),start=4]
\item \uline{For every integer} \(m\) (\uline{positive or negative}), \(L^{\otimes m}(1)\) \uline{is ample}.
\item (\uline{If} \(X\) \uline{is irreducible). For every pair of integers} \(m,n\), \uline{we have}
\[
  \chi(L^{\otimes m}(n))=\chi({O}_X(n))
\]
\item (\uline{If} \(X\) \uline{is irreducible of dimension} \(r\ge2\)). \uline{We have}
\[
  \langle c_1(L)\cdot c_1(H)^{r-1}\rangle=0,\qquad
  \langle c_1(L)^2\cdot c_1(H)^{r-2}\rangle=0.
\]
\end{enumerate}

Let us prove that (i) implies (ii). Applying the d\'evissage lemma (EGA III 3.1) and 4.5 (i), we reduce to the case where \(X\) is integral and \(F={O}_X\). We must then prove that the polynomial \(\chi(L^{\otimes m})\) is constant. For this, we may replace \(L\) by a multiple, and thus assume that \(L\) is algebraically equivalent to \(0\). In this case, if \(L\) is deformed into \({O}\) by the algebraic family defined by an invertible sheaf \(M\) on some \(X\times_kT\) where \(T\) is connected, of finite type, the polynomial \(\chi(M_t^{\otimes m})\) does not depend on \(t\in T\) (EGA III 7.9.5); hence the implication.

Let \(Y\) be a closed integral curve in \(X\). Taking \(F={O}_Y\) in (ii), we find \(\chi({O}_Y)=\chi(L_Y)\). Thus the formula (``Riemann's theorem'') \(\chi(L_Y)=\langle c_1(L)\cdot Y\rangle+\chi({O}_Y)\) gives \(\langle c_1(L)\cdot Y\rangle=0\); hence (ii) implies (iii).

Let \(f:X'\to X\) be a Chow presentation: \(X'\) is projective and \(f\) is surjective. If we grant that (iii) implies (i) for \(X'\), then
%%%%%%%%%%  scan idx 660  |  folio 647  %%%%%%%%%%
\SGAPrintPage{647}%
by 4.5 the implication for \(X\) follows. Henceforth, then, suppose that \(X\) is projective, equipped with an ample sheaf \(H\).

Let \(X_i\) be the irreducible components of \(X\), equipped with the induced reduced structure. Applying 4.5 and EGA III 2.6.2 to the finite surjective morphism \(\coprod X_i\to X\) reduces the proof that (iii) implies (i) and (iv) to the case where \(X\) is integral.

If \(X\) is irreducible, it is clear that (ii) implies (v), that (iii) implies (vi), and, taking 2.1 into account, that (v) implies (vi), and that (vi) for \(X\) is equivalent to (vi) for \(X_{\red}\) (X 4.4.1).

It therefore remains to show that (vi) implies (i) and (iv), and that (iv) implies (i), assuming \(X\) integral. But then (vi) (resp. (iv)) implies that the family \(\mathcal L\) of classes of the \(L^{\otimes m}\) (\(m\in\mathbb N\)) satisfies 3.13 (iii) (resp. (iv)), and is therefore bounded. There thus exists an invertible sheaf \(L\) on some \(X_T\), with \(T\) of finite type over \(k\), which bounds \(\mathcal L\). Let \(p,q\) be two distinct integers such that \(L^{\otimes p}\) and \(L^{\otimes q}\) correspond to the same connected component of \(T\); then \(L^{\otimes(p-q)}\) is therefore algebraically equivalent to \(0\); hence (vi) (resp. (iv)) implies (i).

Finally, we likewise have that (vi) implies the existence of an integer \(N\) such that, for every integer \(m\), \(L^{\otimes m}\otimes H^{\otimes N}\) is ample. Then, for every \(m\), \((L^{\otimes m}\otimes H)^{\otimes N}\) is ample; hence (vi) implies (iv), which completes the proof.

\sgaStableTarget{sga6.e13.thm.4-7}\textbf{Theorem 4.7.} \uline{Let} \(S\) \uline{be a noetherian scheme, and} \(X\) \uline{a proper} \(S\)\uline{-scheme. Then}:
\begin{enumerate}[label=(\roman*)]
\item \uline{The inclusion morphism} \(\underline{\Pic}^\tau_{X/S}\to\underline{\Pic}_{X/S}\) \uline{is representable by an open immersion}.
\item
\uline{If} \(X\) \uline{is projective and integral over} \(S\), \uline{then} \(\underline{\Pic}^\tau_{X/S}\to\underline{\Pic}_{X/S}\)
%%%%%%%%%%  scan idx 661  |  folio 648  %%%%%%%%%%
\SGAPrintPage{648}%
\uline{is also representable by a closed immersion}.
\item \(\underline{\Pic}^\tau_{X/S}\) \uline{is of finite type over} \(S\), \uline{that is, the family of classes of invertible sheaves} \(\tau\)\uline{-equivalent to} \(0\) \uline{on the fibers of} \(X/S\) \uline{is bounded}.
\end{enumerate}

Indeed, in (i) (resp. (ii)), we must treat the case of a morphism \(T\to\underline{\Pic}_{X/S}\), where \(T\) is an arbitrary \(S\)-scheme, and show that the inverse image of \(\underline{\Pic}^\tau_{X/S}\) is representable by an open (resp. and closed) subscheme of \(T\). It suffices to take \(T=\Spec(B)\) over an affine open \(\Spec(A)\) of \(S\). Then \(T\) is the limit of the filtering inverse system of affine \(S\)-schemes \(T_\alpha=\Spec(B_\alpha)\), where \(B_\alpha\) runs through the finite-type \(A\)-subalgebras of \(B\). By 3.1, the morphism in question comes from a morphism \(T_\alpha\to\underline{\Pic}_{X/S}\). We may therefore suppose that \(T\) is of finite type over \(S\).

The morphism \(T\to\underline{\Pic}_{X/S}\) corresponds to an f.p.p.f. morphism \(R\to T\) and an invertible sheaf \(L\) on \(X_R\). By f.p.p.f. descent, we may replace \(T\) by \(R\). After the base change \(R\to S\), taking into account the compatibility of the formation of \(\underline{\Pic}^\tau_{X/S}\) with every base change (4.1), we are thus reduced in (i) and (ii) to studying the case of a section \(S\to\underline{\Pic}_{X/S}\), defined by an invertible sheaf \(L\) on \(X\).

We must prove that the set \(S_0\) of \(s\in S\) such that \(L_s\) is \(\tau\)-equivalent to \(0\) is open, and also closed under the conditions of (ii); since \(S\) is noetherian, this means three things: for an \(s\in S_0\),
\begin{enumerate}[label=(\alph*)]
\item every generization of \(s\) remains in \(S_0\).
\item every specialization of \(s\), near \(s\), remains in \(S_0\).
\item under the conditions of (ii), every specialization of \(s\) remains in \(S_0\).
\end{enumerate}

%%%%%%%%%%  scan idx 662  |  folio 649  %%%%%%%%%%
\SGAPrintPage{649}%
Let \(f:X'\to X\) be a Chow presentation: \(X'\) is projective over \(S\); and \(f\) is surjective. If we grant (a), (b), (c) and (iii) for \(X'\), then they follow for \(X\) by 4.3 and 3.5. We may therefore suppose that \(X\) is projective over \(S\).

Let \(H\) be ample for \(X/S\). If, for an \(s\in S\), \((L^{\otimes m}\otimes H)_s\) is ample, then \((L^{\otimes m}\otimes H)_t\) is ample for every generization \(t\) of \(s\) (EGA III 4.7.1). By the equivalence of 4.6 (i) and (iv), we have therefore established (a).

For (b), (c) and 4.7 (iii), we may suppose that \(S\) is integral, with \(s\) the generic point in (b) and (c) (For (b) and (c), this is clear, and for (iii), it follows from 3.4 (i) applied to the morphism \(\coprod S_i\to S\), where the \(S_i\) are the irreducible components of \(S\), equipped with the induced reduced structure).

There exist an integral scheme \(S'\) and an f.p.p.f. morphism \(S'\to S\) such that every irreducible component \(X'_i\) of \(X_{S'}\), equipped with the induced reduced structure, is integral over \(S'\) (XII 1.3). To prove (b) and (iii), it suffices to do so for \(X_{S'}/S'\), (in (b), by descent; in (iii), by noetherian induction and 3.4 (i)).

Finally, consider the canonical surjective morphism \(\coprod X'_i\to X_{S'}\). Then (b) and (iii) for \(X_{S'}/S'\) follow from (c) and (iii) for the \(X'_i/S'\), by 4.3, 3.5 and 3.7.

We have therefore reduced to proving (c) and (iii), assuming \(S\) integral with generic point \(s\), and \(X\) projective and integral over \(S\). Let \({O}_X(1)\) be a very ample sheaf for \(X/S\). Then, for every \(m,n\), \(\chi(L_t^{\otimes m}(n))\) and \(\chi({O}_{X_t}(n))\) do not depend on \(t\in S\) (EGA III 7.9.4). Thus (c) follows from 4.6 (i) \(\Longleftrightarrow\) (v), which also implies that the family of classes of invertible sheaves \(\tau\)-equivalent to \(0\) on the fibers of \(X/S\) has just
%%%%%%%%%%  scan idx 663  |  folio 650  %%%%%%%%%%
\SGAPrintPage{650}%
one Hilbert polynomial; thus (iii) follows from 2.11.

\sgaStableTarget{sga6.e13.rem.4-8}\textbf{Remark 4.8.} Let \(X\) be proper and integral over noetherian \(S\). Without a projectivity hypothesis on \(X/S\), it is unknown whether it remains true that \(\underline{\Pic}^\tau_{X/S}\to\underline{\Pic}_{X/S}\) is a \uline{closed} immersion, except when \(X\) is normal over \(S\) and \(\underline{\Pic}_{X/S}\) is \uline{representable} (FGA 236, 2.3).

% Added in pages 654--669
\providecommand{\NS}{\operatorname{NS}}
\providecommand{\Picard}{\operatorname{Pic}}
\providecommand{\Et}{\mathrm{et}}
\providecommand{\rk}{\operatorname{rk}}
\providecommand{\im}{\operatorname{Im}}
\providecommand{\HH}{\mathrm H}


\sgaStableTarget{sga6.e13.sec.5}\subsection*{5. Finiteness theorems of ``N\'eron-S\'everi'' type.}

\sgaStableTarget{sga6.e13.thm.5-1}\textbf{Theorem 5.1.} \uline{Let} \(X\) \uline{be proper over noetherian} \(S\); \uline{then the N\'eron-S\'everi groups} \(\underline{\Pic}_{X_K/K}(K)/\underline{\Pic}^0_{X_K/K}(K)\) \uline{of the geometric fibers} \(X_K/K\) \uline{of} \(X/S\) \uline{are finitely generated, and their ranks and the orders of their torsion subgroups} \(\underline{\Pic}^\tau_{X_K/K}(K)/\underline{\Pic}^0_{X_K/K}(K)\) \uline{remain bounded}.

Indeed, by noetherian induction, we may restrict ourselves to proving 5.1 with \(S\) replaced by an arbitrary nonempty open of \(S_{\red}\); thus we may suppose that \(\underline{\Pic}^\tau_{X/S}\) is representable by an \(S\)-scheme of finite type (XII 1.2 and 4.7 (iii)). Then the order of \(\underline{\Pic}^\tau_{X_K/K}(K)/\underline{\Pic}^0_{X_K/K}(K)\) is equal to the number of connected components of \((\Pic^\tau_{X/S})_K\), a number that remains the same over a nonempty open of \(S\) (EGA IV 9.7.9).

It remains to bound the rank of \(\underline{\Pic}_{X_K/K}(K)/\underline{\Pic}^\tau_{X_K/K}(K)\). As in the proof of 4.7 (iii), we readily reduce to the case where \(S\) is integral, and \(X\) is projective and integral over \(S\). Let \(r\) be the common dimension of the fibers of \(X/S\). If \(r\) is \(0\) (resp. \(1\)), then this rank is clearly \(0\) (resp. \(1\)).

If \(r\) is \(\ge2\), we shall reduce to the case where \(r=2\) and \(X\) is smooth over \(S\). After restricting \(S\), we can find a linear section \(Y\) of \(X\) which is
%%%%%%%%%%  scan idx 664  |  folio 651  %%%%%%%%%%
\SGAPrintPage{651}%
still integral over \(S\), with fibers of dimension \(2\); by 4.3, we may replace \(X\) by \(Y\). Finally, by 4.3 and 3.10 (using resolution of singularities for surfaces), after replacing \(S\) by a flat covering of a nonempty open of \(S\), we may replace \(X\) by a smooth \(S\)-scheme.

Now, by 4.6 (i) \(\Longleftrightarrow\) (iii) and the finiteness and specialization theorem for \'etale cohomology groups (SGA 4 XVI 2.2), 5.1 follows from the next proposition.

\sgaStableTarget{sga6.e13.prop.5-2}\textbf{Proposition 5.2.} \uline{Let} \(k\) \uline{be an algebraically closed field and} \(X\) \uline{a proper smooth} \(k\)\uline{-scheme. Then the group} \(C^p_{\rm num}(X)\) \uline{of algebraic cycles of codimension} \(p\) \uline{on} \(X\), \uline{modulo numerical equivalence, is free of rank} \(\le b^{2p}\), \uline{the Betti number}.

Indeed, if we forget the action of the roots of 1, \'etale cohomology \(H^*(X,\Q_\ell)=_{\rm dfn}\left[\varprojlim H^*(X,\Z/\ell^\nu\Z)\right]\otimes_{\Z_\ell}\Q_\ell\) where \((\ell,\operatorname{car}(k))=1\) becomes a Weil cohomology: it satisfies Poincar\'e duality and the K\"unneth formula; and it receives a functorial, nonzero homomorphism from the ring of algebraic cycles modulo rational equivalence \(C^*_{\rm rat}(X)\) (see X 7.13).

Let \(a_1,\ldots,a_m\in C^{r-p}_{\rm rat}(X)\), where \(r=\dim(X)\), be elements that map to a basis of the \(\Q_\ell\)-vector subspace of \(H^{2(r-p)}(X,\Q_\ell)\) generated by algebraic cycles. Consider the homomorphism
\[
   \alpha:C^p_{\rm rat}(X)\longrightarrow \Z^m
\]
defined by \(\alpha(b)=(\langle b.a_1\rangle,\ldots,\langle b,a_m\rangle)\), where \(\langle\ \rangle\) denotes the intersection number. Then one readily sees that the image of \(\alpha\) is \(C^p_{\rm num}(X)\).

%%%%%%%%%%  scan idx 665  |  folio 652  %%%%%%%%%%
\SGAPrintPage{652}%
\sgaStableTarget{sga6.e13.rem.5-3}\textbf{Remark 5.3.} Let \(k\) be an algebraically closed field, and \(X\) a proper smooth surface over \(k\). Then the second Betti number \(b^2(X)\) can be interpreted as \(c_2(X)-2+2b^1(X)\), where \(b^1(X)=2\dim(\Pic_{X/k})\), and \(c_2(X)\) = the second Chern class = the Euler-Poincar\'e characteristic = the intersection number of the diagonal of \(X\times X\) with itself; thus \(b^2(X)\) has an interpretation independent of cohomology. The inequality \(\rang(C^1_{\rm num}(X))\le b^2(X)\) in this form for \(\operatorname{car}(k)>0\) is due (before \'etale cohomology) to IGUSA, who used deep results on the structure of the tame fundamental group of an algebraic curve minus a few points.

\sgaStableTarget{sga6.e13.thm.5-4}\textbf{Theorem 5.4.} \uline{Let} \(k\) \uline{be an algebraically closed field, and} \(X\) \uline{a proper} \(k\)\uline{-scheme. Then there exist finitely many closed integral curves} \(Y_1,\ldots,Y_p\) \uline{in} \(X\) \uline{with the following property: a subset} \(\Lambda\) \uline{of the scheme} \(\Pic_{X/k}\) \uline{is quasi-compact if and only if the intersection numbers} \(\langle\lambda.Y_i\rangle\) \uline{for} \(\lambda\in\Lambda\) \uline{remain bounded. Moreover, we may take for} \(p\) \uline{the rank of the N\'eron-Severi group} \(\Pic_{X/k}(k)/\Pic^0_{X/k}(k)\).

Indeed, since \(\Pic^\tau_{X/k}\) is an open quasi-compact subgroup of \(\Pic_{X/k}\) (4.7), \(\Lambda\) is quasi-compact if and only if the corresponding set of points in the vector space \(\bigl[\Pic_{X/k}(k)/\Pic^\tau_{X/k}(k)\bigr]\otimes\Q\) is finite. But \(\tau\)-equivalence is equal to numerical equivalence (4.6 (i) \(\Longleftrightarrow\) (iii)), so we may choose among the curves in \(X\) curves that map to a basis of the dual space; hence the assertion.

%%%%%%%%%%  scan idx 666  |  folio 653  %%%%%%%%%%
\SGAPrintPage{653}%
\sgaStableTarget{sga6.e13.sec.6}\subsection*{6. Appendix: Study of $(b)$-sheaves on $P=\mathbb P^N_k$.}

\sgaStableTarget{sga6.e13.para.6-1}6.1. Let $k$ be an algebraically closed field, and let $P=\mathbb P^N_k$. Let $F$ be a coherent sheaf on $P$, with Hilbert polynomial
\[
   \chi(F(n))=\sum_{i=0}^{r} a_i \binom{n+i}{i},
\]
and let
\[
   (c)=(c_0,\ldots,c_r)
\]
be a sequence of integers $\ge (a)$. Put
\[
   a_{m,i}=\sum_{j=0}^{r-i}a_{j+i}\binom{m-1+j}{j}\quad\text{and}\quad c_{m,i}=\sum_{j=0}^{r-i}c_{j+i}\binom{m-1+j}{j}.
\]
Finally, for every integer $q\ge1$, denote the largest coherent subsheaf of $F$ with support of dimension $\le q-1$ by $N_q$, and put $F_q=F/N_q$.

\sgaStableTarget{sga6.e13.lem.6-2}\textbf{Lemma 6.2.} \uline{Under the conditions of} 6.1, \uline{consider an exact sequence}
\[
   0\longrightarrow I\longrightarrow {O}_P^{\oplus M}\longrightarrow F(m)\longrightarrow0
\]
\uline{with} $m\ge0$. \uline{Then} $I$ \uline{is} $p$\uline{-regular with} \(p=P_r(c_{m,0},\ldots,c_{m,r})\), \uline{where} $P_r$ \uline{is the} $r$\uline{-th} $(b)$\uline{-polynomial} (1.9).

Indeed, ${O}_P$ is clearly a $(0,\ldots,0,1)$-sheaf ($N+1$ terms), so $I$ is a $(0,\ldots,0,M)$-sheaf by 1.6. By 2.10, \(\chi(I(n))=M\binom{n+N}{N}-\sum_{i=0}^{r} a_{m,i}\binom{n+i}{i}\). Thus the assertion follows from 1.11 and 1.10.

\sgaStableTarget{sga6.e13.prop.6-3}\textbf{Proposition 6.3.} \uline{Under the conditions of} 6.1, \uline{suppose that} $F$ \uline{is a} $(b)$\uline{-sheaf. Then one can calculate, by means of universal polynomials in the} $a_i$, $b_i$ \uline{and} $N$, \uline{integers} $m, M, m_1$ \uline{and} $M_1$ \uline{such that there is an exact sequence}
\[
   {O}_P(-m_1)^{\oplus M_1}\longrightarrow
   {O}_P(-m)^{\oplus M}\longrightarrow F\longrightarrow0.
\]
\uline{Moreover}, $N$ \uline{occurs only in the calculation of} $M_1$.

%%%%%%%%%%  scan idx 667  |  folio 654  %%%%%%%%%%
\SGAPrintPage{654}%
Indeed, take $m=P_{r-1}(c_0,\ldots,c_{r-1})$ with $c_i=b_i-a_i$, and $M=\sum a_i\binom{m+i}{i}$, so that there is an exact sequence
\[
   0\longrightarrow I\longrightarrow {O}_P^{\oplus M}
   \longrightarrow F(m)\longrightarrow0
\]
by 1.11 and 1.3. Similarly, taking 6.2 into account, take $m_1=m+p$ with $p=P_r(a_{m,0},\ldots,a_{m,r})$, and $M_1=M\cdot\binom{p+N}{N}-\sum a_i\binom{m_1+i}{i}\quad(=\chi(I(p)))$.

\sgaStableTarget{sga6.e13.thm.6-4}\textbf{Theorem 6.4.} \uline{Under the conditions of} 6.1, \uline{suppose that} $F$ \uline{is a quotient of} ${O}_P(-m)^{\oplus M}$ \uline{with} $m\ge0$. \uline{Let} $b_i=P_{r-i}(c_{m,i},\ldots,c_{m,r})$ \uline{for} $i=0,\ldots,r$, \uline{where} $P_s$ \uline{is the} $s$\uline{-th} $(b)$\uline{-polynomial} (1.9); \uline{let} $b=b_1+m-1$, \uline{and let} $B=\sum_{i=1}^{r} a_i\binom{b+i}{i}$. \uline{Then}:
\begin{enumerate}[label=(\roman*)]
\item $F$ \uline{is} $b$\uline{-regular}.
\item $-B\le a_0=h^0(F(b))-B$
\item $F$ \uline{is a} $(b_0,\ldots,b_r)$\uline{-sheaf}.
\end{enumerate}

Indeed, take a section $\sigma$ of ${O}_P(1)$ such that we obtain a commutative diagram with exact rows and columns
\[
\begin{tikzcd}[column sep=large,row sep=large]
&0 \arrow[d]&0 \arrow[d]&0 \arrow[d]&\\
0 \arrow[r]&I(-1) \arrow[r] \arrow[d]&I \arrow[r] \arrow[d]&J \arrow[r] \arrow[d]&0\\
0 \arrow[r]&{O}_P(-1)^{\oplus M} \arrow[r] \arrow[d]&{O}_P^{\oplus M} \arrow[r] \arrow[d]&{O}_Y^{\oplus M} \arrow[r] \arrow[d]&0\\
0 \arrow[r]&F(m-1) \arrow[r] \arrow[d]&F(m) \arrow[r] \arrow[d]&G(m) \arrow[r] \arrow[d]&0\\
&0&0&0&.
\end{tikzcd}
\]
%%%%%%%%%%  scan idx 668  |  folio 655  %%%%%%%%%%
\SGAPrintPage{655}%
Since \(\chi(G(n))=\sum_{i=0}^{r-1} a_{i+1}\binom{n+i}{i}\) by 1.7, $J$ is $b_1$-regular by 6.2; thus \(H^q(I(p))=0\) for $q\ge2$ and $p\ge b_1-q$ by 1.4 (i). The exact sequence \(H^q({O}_P^{\oplus M}(p))\longrightarrow H^q(F(m+p))\longrightarrow H^{q+1}(I(p))\) then shows that $F(m)$ is $(b_1-1)$-regular; hence (i).

Assertion (i) implies that \(0\le h^0(F(b))=\chi(F(b))=\sum_{i=0}^{r}a_i\binom{b+i}{i}\); hence (ii). By 2.10 and 1.9,
\[
   \chi(F(b))=\sum_{i=0}^{r} a_{m,i}\binom{b_1-1+i}{i}\le b_1+\sum_{i=0}^{r} c_{m,i}\binom{b_1-1+i}{i}=b_0\,;
\]
thus, $h^0(F(-1))\le h^0(F(b))\le b_0$. By induction on $r$, we may suppose that $G$ is a $(b_1,\ldots,b_r)$-sheaf; hence (iii) by 1.6.2.

\sgaStableTarget{sga6.e13.lem.6-5}\textbf{Lemma 6.5.} \uline{Under the conditions} 6.1, \uline{suppose that} $F$ \uline{has no associated prime cycle of dimension} $0$.
\begin{enumerate}[label=(\roman*)]
\item \uline{If} $h^0(F)\ge1$, \uline{then} $h^0(F(-1))\le h^0(F)-1$.
\item $H^0(F(-n))=0$ \uline{for} $n\ge h^0(F)$.
\item \uline{Let an exact sequence} $0\to F(-1)\to F\to G\to0$ \uline{and an integer} $n_0\ge0$ \uline{satisfy} $H^0(G(-n_0))=0$; \uline{then} $H^0(F(-n_0))=0$ \uline{and} $h^0(F)\le n_0.h^0(G)$.
\end{enumerate}
Indeed, let $\sigma\in H^0(F)$, $\sigma\ne0$; then $\sigma$ defines a nonzero subsheaf $F'={O}_P.\sigma$ of $F$, so $s=\dim.\operatorname{supp}(F')$ is $\ge1$. Take a commutative diagram with exact rows and columns
\[
\begin{tikzcd}[column sep=large,row sep=large]
&0 \arrow[d]&0 \arrow[d]&0 \arrow[d]&\\
0 \arrow[r]&F'(-1) \arrow[r] \arrow[d]&F' \arrow[r] \arrow[d]&G' \arrow[r] \arrow[d]&0\\
0 \arrow[r]&F(-1) \arrow[r]&F \arrow[r]&G \arrow[r]&0 & .
\end{tikzcd}
\]
%%%%%%%%%%  scan idx 669  |  folio 656  %%%%%%%%%%
\SGAPrintPage{656}%
From $s\ge1$ follows $G'\ne0$; consequently, $F'$ is not contained in $F(-1)$; that is, the section $\sigma$ is not contained in the image of the injection $H^0(F(-1))\to H^0(F)$, whence (i).

Assertion (ii) follows immediately from (i), taking into account that $h^0(F)=0$ clearly implies $h^0(F(-1))=0$, since $h^0(F(-1))\le h^0(F)$. In (iii), we find $0\le h^0(F(n))-h^0(F(n-1))\le h^0(G(n))$. Thus, for $n\ge n_0$, $h^0(F(-n_0))=h^0(F(-n))$ since $h^0(G(-n_0))=0$; hence $h^0(F(-n_0))=0$ by (ii). Consequently, $h^0(F(p-n_0))\le h^0(G(1-n_0))+\ldots+h^0(G(p-n_0))\le ph^0(G(p-n_0))$ for $p\ge1$; whence (iii).

\sgaStableTarget{sga6.e13.prop.6-6}\textbf{Proposition 6.6.} \uline{Under the conditions of} 6.1, \uline{suppose that} $F$ \uline{is a} $(b_0,\ldots,b_r)$\uline{-sheaf. Then} $F_q$ \uline{is a} $\bigl((b_{q-1})^q,\ldots,(b_{q-1})^2,b_{q-1},b_q,\ldots,b_r\bigr)$\uline{-sheaf}.

Indeed, let $\sigma$ be a ``general'' section of ${O}_P(1)$; then $\sigma$ defines a commutative diagram with exact rows and columns
\[
\begin{tikzcd}[column sep=large,row sep=large]
&0 \arrow[d]&0 \arrow[d]&0 \arrow[d]&\\
0 \arrow[r]&N_q(-1) \arrow[r] \arrow[d]&N_q \arrow[r] \arrow[d]&G' \arrow[r] \arrow[d]&0\\
0 \arrow[r]&F(-1) \arrow[r] \arrow[d]&F \arrow[r] \arrow[d]&G \arrow[r] \arrow[d]&0\\
0 \arrow[r]&F_q(-1) \arrow[r] \arrow[d]&F_q \arrow[r] \arrow[d]&G'' \arrow[r] \arrow[d]&0\\
&0&0&0&
\end{tikzcd}
\]
where $G$ is a $(b_1,\ldots,b_r)$-sheaf (1.6 (i)) and where $G''=G_{q-1}$ (a question of depth).

%%%%%%%%%%  scan idx 670  |  folio 657  %%%%%%%%%%
\SGAPrintPage{657}%
Suppose that $q=1$. Then $\operatorname{dim{.}supp}(N_q)=0$; thus $H^1(N_q(-1))=0$ and $h^0(F_q(-1))\le h^0(F(-1))\le b_0$. Moreover, $G'=0$; hence $G\simeq G''$. Thus $F_1$ is indeed a $(b_0,\ldots,b_r)$-sheaf. In addition, 6.5 (ii) shows that $H^0(F_1(-b_0))=0$. If $q\ge2$, we may suppose that $G_{q-1}$ is a $((b_{q-1})^{q-1},\ldots,b_{q-1},\ldots,b_r)$-sheaf and that $H^0(G_{q-1}(-b_{q-1}))=0$, by induction on $q$. Then 6.5 shows that $H^0(F_q(-b_{q-1}))=0$ and $h^0(F_q(-1))\le (b_{q-1})^q$; whence the assertion.

\sgaStableTarget{sga6.e13.thm.6-7}\textbf{Theorem 6.7.} \uline{There exist two sequences} $\{A_i(X_0,\ldots,X_i;Y)\}$ \uline{and} $\{A_i^{(q)}(X_0,\ldots,X_q;Y)\}$ \uline{of polynomials with the following properties. Assume that the conditions of} 6.1 \uline{are satisfied and that} $F$ \uline{is a quotient of} ${O}_P(-m)^{\oplus M}$ \uline{with} $m\ge0$.
\begin{enumerate}[label=(\roman*)]
\item \uline{If} $F$ \uline{is a} $(b_0,\ldots,b_r)$\uline{-sheaf, then, for} $i=0,\ldots,r$,
\[
   |a_i|\le A_{r-i}(b_i,\ldots,b_r;m).
\]
\item \uline{If} $\chi(F_q(n))=\sum_{i=0}^{r} a_i^{(q)}\binom{n+i}{i}$ \uline{is the Hilbert polynomial of} $F_q$, \uline{then, for} $i=0,\ldots,q-1$,
\[
   |a_i^{(q)}|\le A_{r-i}^{(r-q+1)}(c_{q-1},\ldots,c_r;m).
\]
\end{enumerate}
(\uline{Notice that} $a_{q-1}^{(q)}\le a_{q-1}$ \uline{and} $a_q^{(q)}=a_q,\ldots,a_r^{(q)}=a_r$ \uline{since} $\operatorname{dim{.}supp}(N_q)\le q-1$).

Indeed, let an exact sequence $0\to F(-1)\to F\to G\to0$ be such that $G$ is a $(b_1,\ldots,b_r)$-sheaf; arguing by induction on $r$, we may suppose that $A_0,\ldots,A_{r-1}$ have been defined, taking 1.7 into account. Then 6.4 (ii) makes the definition of $A_r$ clear, taking 1.8 into account. Finally, (ii) follows formally from 6.4 (iii), 6.6, and assertion (i) (where we take $F_q$ for $F$).

%%%%%%%%%%  scan idx 671  |  folio 658  %%%%%%%%%%
\SGAPrintPage{658}%
\sgaStableTarget{sga6.e13.cor.6-8}\textbf{Corollary 6.8. (Grothendieck).} \uline{Let} $X$ \uline{be projective over noetherian} $S$, \uline{let} ${O}_X(1)$ \uline{be very ample for} $X/S$, \uline{and let} $\mathcal J$ \uline{be a family of classes of coherent sheaves on the fibers of} $X/S$. \uline{Assume}:
\begin{enumerate}[label=(\alph*)]
\item \uline{There exists a coherent sheaf} $E$ \uline{on} $X$ \uline{such that} $\mathcal J$ \uline{is contained in the family of classes of quotients of sheaves of the form} $E_K$.
\item[$(b_q)$] \uline{In the Hilbert polynomials} $\chi(F_K(n))$ \uline{of the} $F_K\in\mathcal J$, \uline{the coefficients in degrees} $\ge q-1$ \uline{remain bounded above}.
\end{enumerate}
\uline{Then the} $(F_K)_q$ (\uline{as} $F_K$ \uline{ranges over} $\mathcal J$) \uline{form a bounded family; moreover, the coefficients in} $\chi(F_K(n))$ \uline{in degrees} $\ge q-2$ \uline{remain bounded below}.

Indeed, we may clearly suppose that $X=\mathbb P^N_S$ and $E={O}_X(-m)^{\oplus M}$. Then the first assertion follows formally from 6.7 (ii) and 1.13; the second follows from 6.4 (ii), by induction using an exact sequence $0\to F(-1)\to F\to G\to0$.

\sgaStableTarget{sga6.e13.def.6-9}\textbf{Definition 6.9.} \uline{Let} $k$ \uline{be a field. A} special positive $k$-cycle of dimension $r$ \uline{is a projective} $k$\uline{-scheme} $X$, \uline{equipped with a very ample sheaf} ${O}_X(1)$, \uline{which is a union of closed subschemes} $X_j$ \uline{of dimension} $r$, \uline{with} $X_j$ \uline{obtained by a base change} $\operatorname{Spec}(k)\to\operatorname{Spec}(k_j)$ \uline{as the inverse image of an} integral $k_j$\uline{-scheme} $X'_j$ (\uline{equipped with} ${O}_{X'_j}(1)$). \uline{The} degree \uline{of} $X$ \uline{is the coefficient} $a_r$ \uline{in the Hilbert polynomial} $\chi({O}_X(n))=\sum a_i\binom{n+i}{i}$.

\sgaStableTarget{sga6.e13.lem.6-10}\textbf{Lemma 6.10.} \uline{Let} $k$ \uline{be an algebraically closed field and} $X$ \uline{a special positive} $k$\uline{-cycle of dimension} $r$ \uline{and degree} $d$. \uline{Then} ${O}_X$ \uline{is a} $(0,\ldots,0,d)$\uline{-sheaf}.

%%%%%%%%%%  scan idx 672  |  folio 659  %%%%%%%%%%
\SGAPrintPage{659}%
Indeed, let $X$ lie in $\mathbb P^N_k$. After replacing $k$ by the algebraic closure of a purely transcendental extension, which is allowed by 1.6 (i), we may suppose that there are $r(N+1)$ elements in $k$ which are algebraically independent over all the $k_j$. Then the $r$ successive ``generic'' hyperplane sections of $X$, defined by means of these elements, are still special positive $k$-cycles of degree $d$.

Arguing by induction on $r$, it therefore remains to note that, when $r=0$, we have $h^0({O}_X(-1))=d$, and, when $r\ge1$, we have $h^0({O}_X(-1))=0$. For this last point, we may suppose that $X$ itself comes from an integral scheme $X'$, by the injection ${O}_X\to \coprod{O}_{X_j}$; but then $h^0({O}_X(-1))=h^0({O}_{X'}(-1))$, which is $0$ since $H^0({O}_{X'})$ is a field.

\sgaStableTarget{sga6.e13.cor.6-11}\textbf{Corollary 6.11.} (i) \uline{Let} $r,d$ \uline{be two integers, let} $c_i=A_{r-i}(0,\ldots,0,d;0)$ \uline{for} $i=0,\ldots,r$ $(6.7\ \mathrm{(i)})$, \uline{and let} $p=P_r(c_0,\ldots,c_r)$ $(1.9)$. \uline{Let} $k$ \uline{be an algebraically closed field, and} $X$ \uline{a special positive} $k$\uline{-cycle of dimension} $\le r$ \uline{and degree} $\le d$, \uline{with Hilbert polynomial} $\chi({O}_X(n))=\sum_{i=0}^{r}a_i\binom{n+i}{i}$. \uline{Then, for} $i=0,\ldots,r$, \uline{we have}
\[
   |a_i|\le c_i.
\]
\uline{Moreover}, $X$ \uline{embeds as a subscheme of} $\mathbb P^N_k$ \uline{with} $N=d(r+1)-1$, \uline{defined by (at most}) $\binom{N+p}{N}$ \uline{equations of degree} $p$.

(ii) (Chow) \uline{Let} $S$ \uline{be a noetherian scheme. As} $K$ \uline{ranges over the algebraically closed extensions of the} $k(s)$, $s\in S$, \uline{the special positive} $K$\uline{-cycles of bounded dimension and degree form a ``bounded family'' (either with the} $X$ \uline{regarded as subschemes of a fixed} $\mathbb P^N_S$, \uline{or in an evident abstract sense}).

%%%%%%%%%%  scan idx 673  |  folio 660  %%%%%%%%%%
\SGAPrintPage{660}%
Indeed, in (i), the first assertion follows from 6.7 (i); the second, from 6.10 and 1.8 (i), and from 6.2 (applied to $0\to I\to{O}_P\to{O}_X\to0$, $P=\mathbb P^N_k$) and 1.3. Finally, (ii) follows from (i) and 1.13.

\sgaStableTarget{sga6.e13.rem.6-12}\textbf{Remark 6.12.} Let $k$ be an algebraically closed field, and $X$ a projective $k$-scheme of pure dimension $r$, with no embedded associated prime cycle. If we allow arbitrary nilpotent elements in ${O}_X$, we can no longer bound the coefficients in the Hilbert polynomial $\chi({O}_X(n))=\sum_{i=0}^{r}a_i\binom{n+i}{i}$ in terms of $a_r$ alone.

For example, let $Z$ be a smooth curve of degree $d$ in $\mathbb P^3_k$. For every $n\gg0$, there exists a smooth surface $Y$ of degree $n$ which contains $Z$. Then $Z$ is a Cartier divisor on $Y$; put ${O}_{X_n}={O}_Y/{O}_Y(-2Z)$. Finally, we have the exact sequence $0\to{O}_Z(-Z)\to{O}_{X_n}\to{O}_Z\to0$ and the formula $\langle Z^2\rangle_Y=2p_a(Z)-2-(n-4)d$; hence $a_0=\chi({O}_{X_n}(-1))\to\infty$ ($n\to\infty$), although $a_1=2d$.

\sgaStableTarget{sga6.e13.cor.6-13}\textbf{Corollary 6.13.} \uline{Under the conditions of} 6.1, \uline{consider an exact sequence}
\[
   {O}_P(-m_1)^{\oplus M_1}\longrightarrow
   {O}_P(-m)^{\oplus M}\xrightarrow{\alpha}F\longrightarrow0
\]
\uline{with} $m_1\ge m\ge0$. \uline{Then, for} $i=0,\ldots,N$,
\[
   |a_i|\le A_{N-i}(0,\ldots,0,M;m_1)+M\binom{m}{N-i}.
\]

Indeed, put $I=\operatorname{Ker}(\alpha)$; clearly, $I$ is a $(0,\ldots,0,M)$-sheaf, and $\chi(I(n))=\sum_{i=0}^{N}\left[M\binom{-m-1+N-i}{N-i}-a_i\right]\binom{n+i}{i}$; whence the assertion by 6.7 (i), and the formula $\binom{-m-1+Q}{Q}=(-1)^Q\binom{m}{Q}$.

%%%%%%%%%%  scan idx 674  |  folio 661  %%%%%%%%%%
\SGAPrintPage{661}%
\sgaStableTarget{sga6.e13.cor.6-14}\textbf{Corollary 6.14. (G. Hermann).} \uline{For every} $N\ge0$, \uline{there exists a polynomial} $R_N(m)$ \uline{such that, for every field} $k_0$ \uline{and every ideal} $I$ \uline{of the polynomial ring} $k_0[T_1,\ldots,T_N]$, \uline{generated by elements of degree} $\le m$, \uline{the radical} $\sqrt I$ \uline{of} $I$ \uline{is generated by elements of degree} $\le R_N(m)$.

Indeed, let $k$ be the algebraic closure of $k_0$. Introducing an auxiliary indeterminate, consider $Y$ (resp. $X$), the subscheme of $P=\mathbb P^N_k$ defined by the homogenized ideal of $I$ (resp. of $\sqrt I$). In the Hilbert polynomial $\chi({O}_Y(n))=\sum_{i=0}^{N}b_i\binom{n+i}{i}$, we have $|b_i|\le c_i$ with $c_i=A_{N-i}(0,\ldots,0,1;m)$ by 6.13.

Let $X=\bigcup X^q$ be the decomposition according to dimension. Then $\text{degree}(X^q)\le e_q$ with $e_q=P_{N-q}(c_q,\ldots,c_N)$. Indeed, after cutting $X$ (resp. $Y$) by a ``general'' linear space of codimension $q$, we may suppose that $q=0$. Then clearly, $\text{degree}(X^q)=h^0({O}_{X^q})$ and $h^0({O}_{X^q})\le h^0({O}_Y)$, and by 6.4 (iii), $h^0({O}_Y)\le e_q$.

Consequently, ${O}_{X^q}$ is a $(0,\ldots,0,e_q)$-sheaf by 6.10. It therefore follows from the injection ${O}_X\to \prod{O}_{X^q}$ that ${O}_X$ is an $(e_0,\ldots,e_N)$-sheaf. Thus, in the Hilbert polynomial $\chi({O}_X(n))=\sum_{i=0}^{N}a_i\binom{n+i}{i}$, we have $|a_i|\le f_i$ with $f_i=A_{N-i}(e_i,\ldots,e_N;0)$ by 6.7. Finally, the ideal $\sqrt I$ of $X$ is $R_N(m)$-regular with $R_N(m)=P_N(f_0,\ldots,f_N)$; whence the assertion by 1.3.

%%%%%%%%%%  scan idx 675  |  folio 662  %%%%%%%%%%
\SGAPrintPage{662}%
\sgaStableTarget{sga6.e13.sec.7}\subsection*{7. Appendix: The Hodge index theorem}

Here we recover and complete some earlier results by a direct argument.

\sgaStableTarget{sga6.e13.thm.7-1}\textbf{Theorem 7.1.} \uline{Let} $X$ \uline{be an irreducible surface, proper over an algebraically closed field. Suppose there exists an invertible sheaf} $H$ \uline{on} $X$ \uline{such that} $\langle c_1(H)^2\rangle>0$ (\uline{for example}, $H$ \uline{ample). Then an invertible sheaf} $L$ \uline{on} $X$ \uline{such that}
\[
   \langle c_1(L)\cdot c_1(H)\rangle=0
   \quad\hbox{\uline{and}}\quad
   \langle c_1(L)^2\rangle\geq 0,
\]
\uline{is numerically equivalent to} $0$. (\uline{Consequently, the intersection form on} $A^1_{\mathrm{num}}(X)_{\Q}$ ($=\bigl[\Pic(X)/\Pic^\tau(X)\bigr]\otimes\Q$, \uline{by} 4.6) \uline{is nondegenerate of type} $(1,\rho-1)$).

\sgaStableTarget{sga6.e13.para.7-1-1}7.1.1. First suppose: $X$ reduced (hence integral); $H$ ample; $\langle c_1(L)\cdot c_1(H)\rangle=0$; and $\langle c_1(L)^2\rangle>0$. Arguing by contradiction, take $m>0$ such that $H_1=L\otimes H^{\otimes m}$ is ample, and take $n$ such that, putting $N=L^{\otimes n}\otimes H^{-1}$, we have that $\langle c_1(N)\cdot c_1(H_1)\rangle=n\langle c_1(L)^2\rangle-m\langle c_1(H)^2\rangle$ is $>0$. Then we also have that $\langle c_1(N)\cdot c_1(H)\rangle=-\langle c_1(H)^2\rangle$ is $<0$, and that $\langle c_1(N)^2\rangle=n^2\langle c_1(L)^2\rangle+\langle c_1(H)^2\rangle$ is $>0$, which contradicts the following lemma.

\sgaStableTarget{sga6.e13.lem.7-1-2}\textbf{Lemma 7.1.2.} \uline{Let} $X$ \uline{be an integral surface, projective over an algebraically closed field. For an invertible sheaf} $N$ \uline{on} $X$ \uline{such that} $\langle c_1(N)^2\rangle>0$, \uline{the following conditions are equivalent}:
\begin{enumerate}[label=(\roman*)]
\item \uline{For every ample sheaf} $H$, \uline{we have} $\langle c_1(N)\cdot c_1(H)\rangle>0$.
\item[(i')] \uline{There exists an ample sheaf} $H$ \uline{such that we have} $\langle c_1(N)\cdot c_1(H)\rangle>0$.
\item \uline{There exists an integer} $n>0$ \uline{such that we have} $H^0(X,N^{\otimes n})\ne0$.
\end{enumerate}

%%%%%%%%%%  scan idx 676  |  folio 663  %%%%%%%%%%
\SGAPrintPage{663}%
Indeed, (ii) implies that there exists an effective divisor $D$ such that $N^{\otimes n}\simeq{O}_X(D)$, and $\langle c_1(N)^2\rangle>0$ implies that $D>0$; whence (i).

To see that (i') implies (ii), we may clearly suppose that $H={O}_X(1)$ is very ample; then there exists a closed integral, locally principal curve $Y$ such that $H={O}_X(Y)$. By Riemann's theorem, (i') implies that $H^1\bigl(N_Y^{\otimes n}(p)\bigr)=0$ for $p\geq0$ and $n\geq n_0$ with $n_0=2h^1({O}_Y)-1$. Since $H^2\bigl(N^{\otimes n}(p)\bigr)=0$ for $p\gg0$, the exact sequence
\[
   H^1\bigl(N_Y^{\otimes n}(p)\bigr)
   \longrightarrow H^2\bigl(N^{\otimes n}(p-1)\bigr)
   \longrightarrow H^2\bigl(N^{\otimes n}(p)\bigr)
   \longrightarrow0
\]
gives $H^2(N^{\otimes n})=0$ for $n\geq n_0$. Thus, for $n\gg0$, we have
\[
   h^0(N^{\otimes n})
   =\chi(N^{\otimes n})+h^1(N^{\otimes n})
   \geq\chi(N^{\otimes n})
   =1/2\langle c_1(N)^2\rangle n^2+\ldots>0,
\]
which completes the proof of the lemma.

\sgaStableTarget{sga6.e13.para.7-1-3}7.1.3. Now suppose in 7.1 that there exists an invertible sheaf $M$ on $X$ such that $\langle c_1(L)\cdot c_1(M)\rangle\ne0$, and argue by contradiction. Replace $X$ by the underlying reduced scheme of a Chow presentation of $X$, which is allowed by the projection formula X 4.4.4, and let $G$ be ample on $X$. Take $m>0$ such that $H_1=G^{\otimes m}\otimes M$ is ample, and such that $\langle c_1(L)\cdot c_1(H_1)\rangle\ne0$; take $p,q$ ($q\ne0$) such that, putting $L_1=L^{\otimes p}\otimes H^{\otimes q}$, we have that $\langle c_1(L_1)\cdot c_1(H_1)\rangle=p\langle c_1(L)\cdot c_1(H_1)\rangle+q\langle c_1(H)\cdot c_1(H_1)\rangle$ is $0$; then we also have that $\langle c_1(L_1)^2\rangle=p^2\langle c_1(L)^2\rangle+q^2\langle c_1(H)^2\rangle$ is $>0$; this returns us to case 7.1.1, with $L=L_1$ and $H=H_1$.

\sgaStableTarget{sga6.e13.para.7-1-4}7.1.4. Finally, consider the general case of 7.1. Still arguing by contradiction, suppose that $L$ is not numerically equivalent to $0$; in addition, replace $X$ by the normalization of $X_{\red}$, so that $X$ has
%%%%%%%%%%  scan idx 677  |  folio 664  %%%%%%%%%%
\SGAPrintPage{664}%
only finitely many singular points, a reduction justified by the projection formula X 4.4.4 and by 4.5 (ii). There therefore exists a closed integral curve $Y$ on $X$ such that $\langle c_1(L)\cdot Y\rangle\ne0$, and $Y$ is locally principal outside a finite set $Z$. Let $f:X'\to X$ be the blowup of the ideal $I$ of ${O}_X$. The ideal $M=I\cdot{O}_{X'}$ is invertible, and defines a subscheme of $X'$ which is the union of two subschemes $Y'$, $Y''$ such that $f(Y')=Y$ and $f(Y'')\subset Z$. By the projection formula X 4.4.4, we have that $\langle c_1(f^*L)\cdot c_1(M)\rangle=-\langle c_1(f^*L)\cdot Y'\rangle-\langle c_1(f^*L)\cdot Y''\rangle$ equals $-\langle c_1(L)\cdot Y\rangle-0\ne0$, and we may replace $X,L,H$ by $X',f^*L,f^*H$; we are therefore reduced to the preceding case, and the proof of 7.1 is complete.

\sgaStableTarget{sga6.e13.rem.7-2}\textbf{Remark 7.2.} It is unknown whether an irreducible proper surface (over an algebraically closed field), which possesses an invertible sheaf $H$ such that $\langle c_1(H)^2\rangle>0$, is necessarily projective. It is also unknown whether, for every irreducible proper surface, the intersection form on $A^1_{\mathrm{num}}(X)_{\Q}$ is nondegenerate; it follows from 7.1 that, when it is degenerate, it is necessarily negative (nondefinite).

\sgaStableTarget{sga6.e13.cor.7-3}\textbf{Corollary 7.3.} \uline{Under the conditions of} 7.1, \uline{let} $M$ \uline{be an invertible sheaf on} $X$ \uline{which is not numerically proportional to} $H$, \uline{and let} $d=\langle c_1(M),c_1(H)\rangle$ \uline{and} $h=\langle c_1(H)^2\rangle$. \uline{Then}, $\langle c_1(M)^2\rangle<d^2/h$.

Indeed, put $L=M^{\otimes h}\otimes H^{\otimes -d}$. Since $\langle c_1(L)\cdot c_1(H)\rangle=0$, by 7.1, $\langle c_1(L)^2\rangle<0$; whence the assertion.

%%%%%%%%%%  scan idx 678  |  folio 665  %%%%%%%%%%
\SGAPrintPage{665}%
\sgaStableTarget{sga6.e13.cor.7-4}\textbf{Corollary 7.4.} \uline{Let} $X$ \uline{be an irreducible scheme of dimension} $r\geq2$, \uline{projective over an algebraically closed field, and let} $H$ \uline{be an ample sheaf on} $X$. \uline{Then}:
\begin{enumerate}[label=(\roman*)]
\item \uline{For an invertible sheaf} $L$ \uline{on} $X$ \uline{to be numerically equivalent to} $0$, \uline{it is (necessary and) sufficient that}
\[
   \langle c_1(L)\cdot c_1(H)^{r-1}\rangle=0
   \quad\hbox{\uline{and}}\quad
   \langle c_1(L)^2\cdot c_1(H)^{r-2}\rangle\geq0.
\]
\item \uline{For every invertible sheaf} $M$ \uline{on} $X$ \uline{which is not numerically proportional to} $H$, \uline{we have}
\[
   \langle c_1(M)^2\cdot c_1(H)^{r-2}\rangle
   <
   \langle c_1(M)\cdot c_1(H)^{r-1}\rangle^2/\langle c_1(H)^r\rangle.
\]
\end{enumerate}

Indeed, (ii) follows immediately from (i), by the argument of 7.3. To establish (i), let $Y$ be a closed integral curve in $X$; we wish to prove that $\langle c_1(L)\cdot Y\rangle=0$. Embed $X$ in $\mathbb P^N_k$ by means of $H^{\otimes n}$ ($n>0$). Let $I=\bigoplus I_\nu$ be the homogeneous ideal of $Y$, and let $\nu_0$ be an integer such that $I_{\nu_0}$ generates $I_\nu$ for $\nu\geq\nu_0$. Then the hypersurfaces $Z$ of degree $\nu_0$ that contain $Y$ have no common points outside $Y$. Arguing by induction on $r=\dim(X)$, we find, using Bertini's theorem, $r-2$ hypersurfaces $Z_1,\ldots,Z_{r-2}$ of degree $\nu_0$ that contain $Y$, such that $X'=Z_1\cap\ldots\cap Z_{r-2}\cap X$ is irreducible. Since $\langle c_1(L)\cdot Y\rangle=\langle c_1(L_{X'})\cdot Y\rangle$, it suffices to prove that $L_{X'}$ is numerically equivalent to $0$ on $X'$. Now, by the projection formula, we have $\langle c_1(L_{X'})\cdot c_1(H_{X'})\rangle=0$ and $\langle c_1(L_{X'})^2\rangle\geq0$; thus the assertion follows from 7.1.

\sgaStableTarget{sga6.e13.cor.7-5}\textbf{Corollary 7.5.}
\uline{Let} $k$ \uline{be an algebraically closed field}, $Y$ \uline{a projective} $k$\uline{-scheme equipped with an ample sheaf} $H$, \uline{and} $X$ \uline{the zero scheme of a section} $\varphi$ \uline{of} $H$. \uline{Suppose that, for every irreducible component} $Y^i$ \uline{of} $Y$,
%%%%%%%%%%  scan idx 679  |  folio 666  %%%%%%%%%%
\SGAPrintPage{666}%
\uline{the irreducible components of} $X\cap Y^i=V(\varphi|Y^i)$ \uline{have dimension} $\geq2$. \uline{Then the morphism} $A^1_{\mathrm{num}}(Y)\longrightarrow A^1_{\mathrm{num}}(X)$ \uline{is injective}.

Indeed, since we may clearly suppose that $Y$ is integral, the assertion follows from 7.4 (i), by the projection formula.

\sgaStableTarget{sga6.e13.bibliography}\subsection*{BIBLIOGRAPHY}
\begin{enumerate}[label={[\arabic*]}]
\item Abhyankar S. : Resolution of singularities of arithmetical surfaces, in : Arithmetical Algebraic Geometry, Purdue University, Harper-Row, New York 1965.
\item Mumford D : Pathologies III, Amer. J. Math 1967 p. 94-104.
\item Grothendieck A. and Dieudonn\'e J. EGA, El\'ements de G\'eom\'etrie Alg\'ebrique, Pub. Math. Inst des Hautes Etudes Sci., 1960.
\item Grothendieck A. : FGA, Fondements de la g\'eom\'etrie alg\'ebrique, Extracts from the S\'eminaire Bourbaki 1957-1962.
\item Grothendieck A. et alii, SGA, S\'eminaire de g\'eom\'etrie alg\'ebrique, 1960-66, I.H.E.S.
\item Jussila O. : X, Formalisme des intersections sur les sch\'emas alg\'ebriques propres, Expos\'e X in this seminar.
\item Raynaud M. : XII, Un th\'eor\`eme de repr\'esentabilit\'e relative sur le foncteur de Picard, Expos\'e XII in this seminar.
\end{enumerate}



% Exposé XIV local notation
\providecommand{\Kb}{K^\bullet}
\providecommand{\Db}{D}
\providecommand{\Gtop}{\operatorname{Gr}_{\mathrm{top}}}
\providecommand{\Ftop}{\operatorname{Fil}_{\mathrm{top}}}


%%%%%%%%%%  scan idx 680  |  folio 667  %%%%%%%%%%
\SGAPrintPage{667}%
\sgaStableTarget{sga6.e14.expose}\section*{\underline{Exposé XIV}\\[3pt] \underline{\underline{OPEN PROBLEMS IN INTERSECTION THEORY}}}
\begin{center}
\underline{by A. Grothendieck}
\end{center}

In this exposé we review some of the problems connected with the questions encountered in this seminar that deserve to be studied. We leave aside the problems connected with Serre's \uline{local theory} of intersections [14], which we have not had occasion to use in the Seminar, and refer in particular to loc. cit. for the two fundamental conjectures that remain open in this direction (loc. cit. Chap. V \S~4). We also pass over the problems, very important for arithmetic applications, connected with the \uline{structure of groups of cycle classes}, especially on a smooth projective algebraic scheme; for a discussion of these problems we refer to S. KLEIMAN, Algebraic cycles and the Weil conjectures in Dix exposés sur la cohomologie des schémas (Masson-North Holland Publishing Cie). In no.~4 of this exposé we also give some supplements to the Seminar, concerning in particular the \uline{Chow ring}, which had not found a place in the preceding exposés.

\sgaStableTarget{sga6.e14.sec.1}\section*{1. Operations $\Lambda^i$ in the derived category $D(X)$\SGAAuthorFootMark{*}}
\SGAAuthorFootText{*}{cf. the note on page 23.}

Let $X$ be a locally ringed topos, and let $K^\bullet(X)$ be the ring of classes of perfect complexes on $X$ (Exp. IV). We would like to define on $K^\bullet(X)$ a $\lambda$-ring structure (Exp. V), functorial in $X$, such that the natural homomorphism
\[
 K^\bullet_{\mathrm{naif}}(X)\longrightarrow K^\bullet(X)
\]
(where the first member is the ring of classes of locally free Modules on $X$, endowed with the $\lambda$-structure defined in Exp. V) be a $\lambda$-homomorphism. For this purpose, one would need to define functors (generally nonadditive) of ``exterior power''
\[
 \Lambda^i:D^-(X)\longrightarrow D^-(X),
\]
taking pseudo-coherent (resp. perfect) complexes to pseudo-coherent (resp. perfect) complexes and having the formal properties needed, with respect to exact triangles, to define a $\lambda$-structure on $K^\bullet(X)$ by the formula
\sgaStableTarget{sga6.e14.eq.1-1}\begin{equation}\tag{1.1}
 \lambda^i(\mathrm{c}\ell^\bullet(L_\bullet))=\mathrm{c}\ell^\bullet(\Lambda^i L_\bullet).
\end{equation}

%%%%%%%%%%  scan idx 681  |  folio 668  %%%%%%%%%%
\SGAPrintPage{668}%
When $X$ is of characteristic zero, i.e. its structure sheaf is a sheaf of $Q$-algebras, the problem is solved simply by taking
\sgaStableTarget{sga6.e14.eq.1-2}\begin{equation}\tag{1.2}
 \Lambda^i L_\bullet=\bigl(\overset{i}{\otimes}{}^{\mathbb L}L_\bullet\bigr)^{\mathrm{alt}},
\end{equation}
where the $i$-th tensor power (in the sense of derived categories) of $L_\bullet$ is regarded as an object of $D^-(\underline{O}_i)$, with $\underline{O}_i$ the group algebra of the symmetric group $\mathfrak S_i$ with coefficients in $\underline{O}_X$, and where the superscript alt denotes passage to the ``alternating part under the action of $\mathfrak S_i$'' (an operation that makes sense by the characteristic-zero hypothesis). A formula of this type was already used in [8] (cited below as [RRR]) for the first proof of the Riemann--Roch theorem.

Without a hypothesis on the characteristic, when $L_\bullet$ is zero in degrees $>0$, Dold--Puppe [6] shows that $L_\bullet$ is canonically homotopic to the complex associated with a semi-simplicial Module $L'_\bullet$ canonically associated with $L_\bullet$, and a fairly natural candidate for $\Lambda^iL_\bullet$ (for $L_\bullet$ flat, say) would be the complex associated with the semi-simplicial Module $\Lambda^iL'_\bullet$ (whose component of degree $n$ is $\Lambda^i L^{\prime}_n$). But to extend the definition to complexes of $D^-(X)$ no longer assumed to be zero in degrees $>0$, one would have to study the behaviour of the preceding operation with respect to the translation operation $L_\bullet\mapsto L_\bullet[1]$ (or \uline{suspension}, in the terminology of [6]); unfortunately, there does not seem to be an isomorphism $\Lambda^i(L_\bullet[1])\simeq(\Lambda^i L_\bullet)[i]$ that would make the desired extension trivial.

\sgaStableTarget{sga6.e14.sec.2}\section*{2. The Riemann--Roch formula without projectivity hypotheses}

\sgaStableTarget{sga6.e14.para.2-1}2.1. Once a natural $\lambda$-structure on the $K^\bullet(X)$ is available, it is possible to formulate the Riemann--Roch theorem, in the form considered in Exp. VIII, for any proper complete-intersection morphism $f:X\to Y$ (and even, more generally, for a relative scheme $X$ over a locally ringed topos $Y$, with $X$ still subject to the corresponding relative conditions). The Todd class of $f$ is defined by means of the class in $K^\bullet(X)$ of the relative cotangent complex $T_f$ for $f$, constructed in Exp. 0 4.4; this complex is perfect here by the local hypothesis on $f$. It seems clear that a proof of the Riemann--Roch formula in this general case will require the introduction of essentially new ideas, and it is possible that these same ideas could provide the key
%%%%%%%%%%  scan idx 682  |  folio 669  %%%%%%%%%%
\SGAPrintPage{669}%
to a proof of the analogous formula in the complex analytic (or rigid analytic!) case. Recall that in the complex analytic case, only the case of a locally free sheaf $F$ on a compact nonsingular variety $X$ over a point $Y$ is currently known (by the Atiyah--Singer index formula); the case in which $F$ is merely assumed coherent still seems to lie beyond the techniques known so far.

\sgaStableTarget{sga6.e14.para.2-2}2.2. As an essential ingredient of the Riemann--Roch formula, let us also recall the \uline{finiteness problem} raised in Exp. III: for a proper morphism $f:X\to Y$ and a complex $L_\bullet$ on $X$ that is pseudo-coherent relative to $Y$, is $Rf_*(L_\bullet)$ pseudo-coherent? A general proof of this conjecture will probably use algebraic ideas that will make it possible to prove Grauert's finiteness ``theorem'' [7] under considerably more general conditions than those of loc. cit., starting directly with topological $\C$-algebra structures on base spaces (or topoi) $Y$ more general than analytic spaces (including the case of a locally compact space $Y$ with the sheaf of continuous complex functions, or a $C^m$ manifold with the sheaf of $C^m$ complex functions, or a real analytic space, etc.).

\sgaStableTarget{sga6.e14.para.2-3}2.3. Finally, in these questions one would need to replace the properness hypothesis on the morphism $f$ by the hypothesis that the cohomology sheaves $H^i(L_\bullet)$ of the complexes under consideration on $X$ have their supports proper over $Y$.

\sgaStableTarget{sga6.e14.sec.3}\section*{3. A Riemann--Roch formula ``without denominators'' for an immersion}

\sgaStableTarget{sga6.e14.para.3-1}3.1. Let $f:Y\to X$ be a regular closed immersion (Exp. VII), with conormal sheaf $\underline N$. For simplicity, suppose that $X$ admits an ample invertible sheaf. The Riemann--Roch formula for $y\in K^\bullet(Y)$:
\[
 \ch_X(f_*(y))=f_*\bigl(\ch_Y(y)\operatorname{Todd}(-\check{\underline N})\bigr),
\]
together with the general formula (Exp. VII):
\[
 f_*(y')\,f_*(y'')=f_*\bigl(y'y''c_d(\underline N)\bigr),\qquad y',y''\in \operatorname{Gr}(Y)_{\underline Q},
\]
(with $f$ assumed to have codimension $d$), makes it possible to express the Chern classes of $f_*(y)$ in $\operatorname{Gr}^\bullet(X)_\Q$ in the form
%%%%%%%%%%  scan idx 683  |  folio 670  %%%%%%%%%%
\SGAPrintPage{670}%
\sgaStableTarget{sga6.e14.eq.3-1}\begin{equation}\tag{3.1}
 c_i(f_*(y))=f_*\bigl(P_i((c_\alpha(y))_\alpha,(c_\beta(\underline N))_\beta)\bigr),
\end{equation}
where $P_i$ is a universal polynomial in the $c_\alpha(y)$, $c_\beta(\underline N)$, a polynomial whose coefficients are a priori rational. In fact, one finds [RRR] that these polynomials have \uline{integer} coefficients. This raises the question of the validity of a Riemann--Roch formula ``without denominators''. In fact, such a formula for $c_i(f_*(y))$ cannot in general be sought in $\operatorname{Gr}^\bullet(X)$ itself, since the homomorphism $f_*:K^\bullet(Y)\to K^\bullet(X)$ is not itself compatible with the $\lambda$-filtrations up to a shift by $d$ (4.6), but is so only up to torsion; thus one cannot define a homomorphism $f_*:\operatorname{Gr}^\bullet(Y)\longrightarrow \operatorname{Gr}^\bullet(X)$ giving the natural homomorphism $f_*:\operatorname{Gr}^\bullet(Y)_\Q\to\operatorname{Gr}^\bullet(X)_\Q$ and allowing formula (3.1) to be interpreted as a formula in $\operatorname{Gr}^\bullet(X)$, and not merely in $\operatorname{Gr}^\bullet(X)_\Q$. But when $X$ is assumed regular, so that $K^\bullet(X)\xrightarrow{\sim} K_\bullet(X)$ (IV 2.5), one can endow $K^\bullet(X)$ with the \uline{topological filtration} by the codimension of supports, inherited from the corresponding filtration of $K_\bullet(X)$. It is plausible that this filtration is compatible with the ring structure, a fact that can be proved at least [RRR] when $X$ is smooth over a field $k$, using the theory of the Chow ring [4]. When this is so, the associated graded
\[
 \Gtop^\bullet(X)
\]
is a \uline{graded ring}, and, also assuming $Y$ regular, we have done what is needed to obtain a natural homomorphism of degree $d$ of graded groups
\[
 f_*=f_*^{\mathrm{top}}:\Gtop^\bullet(Y)\longrightarrow \Gtop^\bullet(X).
\]
Since the $\lambda$-filtration of $K^\bullet(X)$ is finer than the topological filtration by codimension just considered (X 1.3.3), one obtains a canonical homomorphism of graded groups (and even of graded rings)
\[
 \operatorname{Gr}^\bullet(X)\longrightarrow \Gtop^\bullet(X),
\]
which is moreover an isomorphism mod torsion (Exp. VIII). This allows us, as in [RRR], to consider a theory of Chern classes on $X$ with values in $\Gtop^\bullet(X)$, and likewise on $Y$, and thus to interpret formula (3.1) as a formula in $\Gtop^\bullet(X)$. In this form, the formula is proved in [RRR] when $X$ is of \uline{characteristic zero}, at least when $X$ is smooth over a field (a hypothesis serving only to ensure that the topological filtration of $K^\bullet(X)$ is compatible with the ring structure). The characteristic-zero hypothesis is used essentially
%%%%%%%%%%  scan idx 684  |  folio 671  %%%%%%%%%%
\SGAPrintPage{671}%
to obtain formulas of the type (1.1), (1.2), and essentially the same proof should apply without a hypothesis on the characteristic, once the question raised in no.~1 has been answered affirmatively.

\sgaStableTarget{sga6.e14.para.3-2}3.2. When $X$ is smooth over a field, one also has the Chow ring
\[
 A(X),
\]
and a theory of Chern classes on $X$ with values in $A(X)$ [9], finer than the theory with values in $\Gtop^\bullet(X)$ just introduced---these latter Chern classes being the images of the Chow-style Chern classes under the canonical surjective homomorphism (cf. 4.1)
\[
 \phi:A(X)\longrightarrow \Gtop^\bullet(X).
\]
This homomorphism is an isomorphism mod. torsion, but not an isomorphism in general (4.7). With this understood, for smooth $X$ and $Y$ the inclusion $f:Y\to X$ again defines a homomorphism of degree $d$ of graded groups
\[
 f_*:A(Y)\longrightarrow A(X),
\]
which allows formula (3.1) to be interpreted as a formula in $A^i(X)$\SGAAuthorFootnote{*}{The formula, as well as the variant mentioned in 3.3, has since been proved by P. JOUANOLOU (Nov.\ 1969). Cf.}. This formula reduces to $0=0$ for $i<d$ and is therefore of interest only for $i\geq d$. Even for $i=d=2$ this formula has not been proved at present (cf. 4.3), nor has the less precise formula deduced from it by replacing rational equivalence of algebraic cycles with algebraic equivalence. Unlike the analogous formula in $\Gtop^\bullet(X)$, there seems to be no very convincing reason to suppose that the formula without denominators must hold in $A(X)$. In view of our almost total ignorance of the nature of the groups $A^i(X)$ for $i\geq 2$, a good way to test it seems to be to test the formulas that can be deduced from it in groups $A^1(Z)=\Pic(Z)$, for example by applying $g_*:A(X)\longrightarrow A(Z)$ to both members of (3.1), where $g:X\to Z$ is a proper morphism into a smooth quasi-projective $Z$ of dimension $\dim X-i+1$.

\sgaStableTarget{sga6.e14.para.3-3}3.3. Finally, to give a meaning to (3.1), one can also use the cohomological theory of Chern classes (SGA 5 VII), whenever \uline{Gysin homomorphisms}
\[
 f_*:H^p(Y,\mu_n^{\otimes i})\longrightarrow H^{p+2d}(X,\mu_n^{\otimes(i+d)})
\]
are available (cf. no.~6); this is the case at least when $Y$ and $X$ are both smooth over a scheme $S$ (by the theorem of relative cohomological purity
%%%%%%%%%%  scan idx 685  |  folio 672  %%%%%%%%%%
\SGAPrintPage{672}%
SGA 4 XVI 3.7), and, conjecturally, when $Y$ and $X$ are both regular (by the absolute cohomological purity theorem, proved at least for excellent $X$ of characteristic zero (SGA 4 XIX 3.2)). When $X$ and $Y$ are smooth over the field of complex numbers, one can also use the integral Chern class on the locally compact space $X(\C)$ of complex points of $X$. Formula (3.1) is established in this case by \uline{Atiyah} and \uline{Hirzebruch} [1], who work more generally with an immersion of complex analytic varieties (using the ``naive'' $K^\bullet$); it follows that the $\ell$-adic formula (3.1) holds whenever $Y$ and $X$ are smooth over an \uline{algebraically closed field of characteristic zero}, using the ``Lefschetz principle'' and the general results of SGA 4. This therefore makes the validity of the cohomological formula (3.1) extremely plausible; but observe that it has not been proved at present even when $X$ and $Y$ are smooth over a field $k$ of characteristic zero, when $k$ is not algebraically closed.

\sgaStableTarget{sga6.e14.para.3-4}3.4. One can also work with the cohomological theory of Chern classes in the De Rham or Hodge style [13]. The question of the validity of (3.1) in this setting arises only when $X$ is not of characteristic zero, since in any case (no.~6) the formula in question is known to hold modulo torsion. For smooth quasi-projective schemes over a field $k$, the De Rham or Hodge version of formula (3.1) would moreover follow from the cycle-class version of formula (3.1) mod. algebraic equivalence, so that the De Rham or Hodge cohomological version gives another accessible way to test the validity of the ``cycle-class'' version of (3.1).

\sgaStableTarget{sga6.e14.para.3-5}3.5. The proof [1] of the transcendental cohomological version of (3.1) is essentially transcendental in nature, using tubular neighbourhoods of $Y$ in $X$. This suggests that a proof of a formula of type (3.1) might be connected with the systematic consideration of groups $K^\bullet$ and cohomology groups \uline{with support in} $Y$, which might provide a sufficient technical means of ``localization around $Y$''.

\sgaStableTarget{sga6.e14.sec.4}\section*{4. Relations between $K^\bullet(X)$ and the Chow ring $A(X)$}

\sgaStableTarget{sga6.e14.para.4-1}4.1. For simplicity, suppose that $X$ is a smooth quasi-projective scheme over a field $k$ (cf no.~8 for a more general case), so that we
%%%%%%%%%%  scan idx 686  |  folio 673  %%%%%%%%%%
\SGAPrintPage{673}%
have the Chow ring $A(X)$ of cycle classes modulo rational equivalence [4]. As variants of this ring, which is moreover defined under substantially more general considerations, there are the rings $\operatorname{Gr}^\bullet(X)$ (the associated graded of the $\lambda$-filtration of $K^\bullet(X)$) and $\Gtop^\bullet(X)$ (the associated graded of the filtration of $K^\bullet(X)\simeq K_\bullet(X)$ by the codimension of supports). There are theories of Chern classes with values in each of these three rings (cf. [9] for $A(X)$, and this seminar for $\operatorname{Gr}^\bullet(X)$, which maps to $\Gtop^\bullet(X)$). There is also a cohomological theory ($\ell$-adic, say, with $\ell$ prime to the characteristic of $k$) of Chern classes (SGA 5 VII):
\[
 c_i(\underline{E})\in H^{2i}(X,\underline{Z}_\ell(i)),
\]
(where $\underline{E}$ is a vector bundle on $X$, or indeed any element of $K^\bullet(X)$). These cohomological theories are related by the canonical homomorphisms (compatible with the multiplicative structures)
\sgaStableTarget{sga6.e14.eq.4-1}\begin{equation}\tag{4.1}
\begin{tikzcd}[column sep=large,row sep=large]
 H^{2*}(X,\underline{Z}_\ell(*)) & {} \\
 A(X) \arrow[u,"\gamma"] \arrow[r,"\phi"] & \Gtop^\bullet(X) \\
 {} & \operatorname{Gr}^\bullet(X), \arrow[u,"j"]
\end{tikzcd}
\end{equation}
where we have set
\[
 H^{2*}(X,\underline{Z}_\ell(*))=\coprod_i H^{2i}(X,\underline{Z}_\ell(i)).
\]
The homomorphism $j$ comes from the fact that the $\lambda$-filtration of $K^\bullet(X)$ is finer than the topological filtration by the codimension of the support (Exp. X), and it is an isomorphism modulo torsion (Exp. VII 4.11). The homomorphism $\phi$ is defined, for example, in [RRR], and associates to the class of the cycle of codimension $i$ associated with a closed subscheme $Y$ of $X$ of codimension $i$ the element $\mathrm{c}\ell^\bullet(\underline{O}_Y)$ of $\operatorname{Filt}^i_{\mathrm{top}}(K^\bullet(X))$ mod $\operatorname{Filt}^{i+1}_{\mathrm{top}}(K^\bullet(X))$. It is a surjective homomorphism, which we shall also see (4.2) to be an isomorphism modulo torsion. Thus, \uline{modulo torsion the three rings} $A(X)$, $\Gtop^\bullet(X)$ \uline{and} $\operatorname{Gr}^\bullet(X)$ \uline{are canonically isomorphic}, so \uline{a posteriori} the Riemann--Roch formulas expressed using one or another of these rings (tensored with $\Q$) are equivalent. The graded ring $\operatorname{Gr}^\bullet(X)$, used particularly in this seminar as a substitute for the Chow ring, has the appreciable advantage of being defined for every locally ringed topos and, in particular, for every scheme
%%%%%%%%%%  scan idx 687  |  folio 674  %%%%%%%%%%
\SGAPrintPage{674}%
(regular or otherwise). It has, on the other hand, the disadvantage (4.6) relative to $A(X)$, $\Gtop^\bullet(X)$ that, as it stands, it is not covariant with respect to proper morphisms (even for smooth quasi-projective schemes over a field); it acquires this property only after tensoring with $\Q$. Another advantage of $A(X)$ over its competitors $\Gtop^\bullet$ and $\operatorname{Gr}^\bullet(X)$ is the existence of a homomorphism from $A(X)$ to $H^{2*}(X,\underline{Z}_\ell(*))$ that allows cohomological Chern classes to be interpreted as images of Chern classes with values in $A(X)$. In general this homomorphism does not factor through $\Gtop^\bullet(X)$ (4.7), and it is at least doubtful that there is in general a homomorphism from $\operatorname{Gr}^\bullet(X)$ to $H^{2*}(X,\underline{Z}_\ell(*))$ taking Chern classes to Chern classes (Cf. no.~5). Of course, these difficulties disappear after tensoring with $\underline{Q}$, and one obtains a homomorphism
\sgaStableTarget{sga6.e14.eq.4-2}\begin{equation}\tag{4.2}
 \Gtop^\bullet(X)\longrightarrow H^{2*}(X,\underline{Q}_\ell(*))=
 H^{2*}(X,\underline{Z}_\ell(*))\otimes_{\underline{Z}_\ell}\underline{Q}_\ell
\end{equation}
taking Chern classes to Chern classes.

\sgaStableTarget{sga6.e14.para.4-2}4.2. To establish that the surjective homomorphism $\phi$ of (4.1) is an isomorphism modulo torsion, one uses the Chern homomorphism (deduced from the theory of Chern classes [9])
\[
 K^\bullet(X)\longrightarrow \Ch(A(X))
\]
(where the second member is the Chern ring of $A(X)$ defined in Exp.~V), which is, as is easily seen, compatible with the topological filtration of $K^\bullet(X)$ and the natural filtration of $\Ch(A(X))$, and gives, on passing to the associated graded groups, a homomorphism of graded groups
\sgaStableTarget{sga6.e14.eq.4-3}\begin{equation}\tag{4.3}
 \psi:\Gtop^\bullet(X)\longrightarrow A(X).
\end{equation}
With this understood, we shall prove the formulas
\sgaStableTarget{sga6.e14.eq.4-4}\begin{equation}\tag{4.4}
\left\{
\begin{aligned}
 \psi\phi(x)&=(-1)^{i-1}(i-1)!\,x\quad \text{mod. torsion}\quad \text{for } x\in A^i(X)\\
 \phi\psi(x)&=(-1)^{i-1}(i-1)!\,x\quad \text{mod. torsion}\quad \text{for } x\in \Gtop^i(X),
\end{aligned}
\right.
\end{equation}
the first of which moreover implies the second, since $\phi$ is surjective. These formulas clearly imply that $\phi$ and $\psi$ are isomorphisms modulo torsion and that, up to a factor $(-1)^{i-1}(i-1)!$ in degree $i$, they are inverses of one another.

To prove the first formula (4.4), one can clearly suppose that $x$ is defined by an irreducible closed subscheme $Y$ of $X$ of codimension $i$, in which case the formula is equivalent to
%%%%%%%%%%  scan idx 688  |  folio 675  %%%%%%%%%%
\SGAPrintPage{675}%
\sgaStableTarget{sga6.e14.eq.4-5}\begin{equation}\tag{4.5}
 c_i(\underline{O}_Y)=(-1)^{i-1}(i-1)!\,\cl(Y)\quad \bmod
\text{ torsion},
\end{equation}
where in the first member $\underline{O}_Y$ is regarded as a coherent sheaf on $X$, thus defining an element of $K^\bullet(X)$, and in the second member $\mathrm{c}\ell(Y)$ denotes the class of $Y$ in $A^i(X)$. Since $A^i(X)$ does not change on removing from $X$ a closed subset of codimension $>i$, we are reduced to the case in which $Y$ is regular, hence smooth over the base field $k$ when the latter is perfect. But in this latter case, formula (4.5) is precisely the Riemann--Roch formula for the inclusion $Y\to X$ and the element $1$ of $K^\bullet(Y)$, written using the theory of Chern classes with values in the Chow ring $A(X)$. This formula is established in [2] under the conditions considered here, by a method very close to the one used in this seminar for the analogous formula using $\operatorname{Gr}^\bullet(X)$. This therefore establishes, at least for perfect $k$, formula (4.4) and consequently the fact that $\phi$ is an isomorphism mod. torsion (which in turn implies, a posteriori, that the Riemann--Roch formula ``with denominators'' for a proper morphism of smooth algebraic schemes over $k$ is essentially the same whether written using $A(X)$ as in [2] or using $\operatorname{Gr}^\bullet(X)$ as in this seminar!). The case with no hypothesis on $k$ reduces easily to the case in which $k$ is algebraically closed, using the fact that, for a finite extension $k'$ of degree $n$ of $k$, the kernel of the natural map $A(X)\xrightarrow{f^*}A(X\otimes_k k')$ is annihilated by $n$ (as is seen by using the trace homomorphism $f_*:A(X\otimes_k k')\to A(X)$, satisfying $f_*f^*=n\,\id$).\SGAAuthorFootText{*}{It is therefore established, thanks to P. JOUANOLOU, for $X$ smooth and quasi-projective over a field (cf. note p.5).}

\sgaStableTarget{sga6.e14.para.4-3}4.3. It seems quite plausible that formulas (4.4) hold as they stand and not only modulo torsion, or, equivalently, that formula 4.5. holds not only modulo torsion. This question is a very special case of the validity in $A(X)$ of the ``Riemann--Roch formula without denominators'' (3.1) for an immersion $Y\to X$ (namely, the case $i=d$). The validity of this special case may be regarded as quite plausible. Indeed, we already pointed out in no.~3 that, when $k$ is of characteristic zero, the Riemann--Roch formula without denominators holds at least in $\Gr^{\bullet}_{\mathrm{top}}(X)$; this implies in that case that both members of (4.5) have the same image in $\Gr^{\bullet}_{\mathrm{top}}(X)$, or equivalently that the second of formulas (4.4) holds; this result will extend of itself to arbitrary characteristic once
%%%%%%%%%%  scan idx 689  |  folio 676  %%%%%%%%%%
\SGAPrintPage{676}%
the question raised in no.~1 has been satisfactorily resolved. In support of formula (4.5) in $A(X)$ itself, observe that it is clear a priori that the restriction of the first member to $X-Y$ is zero; thus, with $Y$ assumed irreducible, the first member must a priori be an integer multiple of $\mathrm{c}\ell(Y)$:
\[
 c_i(Y)=n(Y)\mathrm{c}\ell(Y).
\]
This already proves the exact formula (4.5) when $\mathrm{c}\ell(Y)$ is not a torsion element of $A^i(X)$. It is immediate that this is so when $X$ is projective (and not merely quasi-projective) over the field $k$: note that for a polarization $\xi\in A^1(X)$ of $X$, one has $\deg(Y\cdot\xi^{\dim Y})>0$; this proves the desired formula in this case. Consequently, this formula also holds whenever $X$ is isomorphic to an open subscheme of a smooth projective scheme over $k$, which will be the case when a suitable solution to the resolution-of-singularities problem is available. This is so, in particular, when $k$ has characteristic zero, by Hironaka [11].

The validity of the exact formula (4.5), hence of (4.4), would imply that the kernel of $\phi^i:A^i(X)\longrightarrow \Gtop^i(X)$ is annihilated by $(i-1)!$, and in particular \uline{that} $\phi^2:A^2(X)\to\Gtop^2(X)$ \uline{is an isomorphism}. (That $\phi^1$ is an isomorphism already follows immediately from X 5.3.2.) Unless I am mistaken, the fact that $\phi^2$ is an isomorphism has not been proved at present, even in the special case $\dim X=3$. In any event, the writer was hasty in asserting the validity of (4.4) and (4.5) in [9, no.~4, 3\textsuperscript{o}].

\sgaStableTarget{sga6.e14.para.4-4}4.4. At present, the main advantage of $\Gtop^\bullet(X)$ (or $\operatorname{Gr}^\bullet(X)$) over the Chow ring $A(X)$ is that it is better suited to explicit computation of certain intersection formulas between cycles, in the case of ``excess'' intersections, i.e. of excessively large dimension. Indeed, defining the intersection class by means of Chow's lemma, which allows a cycle to be moved sufficiently within its class, does not readily lend itself to computation. A striking example is the following fundamental formula [Exp. VII 2.7] for a closed immersion $f:Y\to X$ of smooth quasi-projective schemes over $k$:
\sgaStableTarget{sga6.e14.eq.4-6}\begin{equation}\tag{4.6}
 f^*f_*(y)=y\,c_d(\underline{N}),
\end{equation}
a formula valid in $\Gtop^\bullet(Y)$, with $y$ in $\Gtop^\bullet(Y)$ and $\underline{N}$ the conormal sheaf
%%%%%%%%%%  scan idx 690  |  folio 677  %%%%%%%%%%
\SGAPrintPage{677}%
of $Y$ in $X$, assumed everywhere to have rank $d$. In loc. cit., this formula follows from a very simple computation of $\underline{\Tor}$ sheaves. By what we saw in 4.2, the preceding formula implies the same formula in $A(Y)$, but only \uline{modulo torsion}. It is quite extraordinary that formula (4.6) in $A(Y)$ itself has not been proved at present\SGAAuthorFootnote{*}{cf. the note on page 23.\newline (**)}, nor even the special case obtained by taking $y=1$, which gives the self-intersection of $Y$:
\sgaStableTarget{sga6.e14.eq.4-7}\begin{equation}\tag{4.7}
 f^*(\cl(Y))=c_d(\underline{N}),
\end{equation}
even when $k$ is the field of complex numbers!

An analogous problem arises when one attempts to compute the Chow ring of the variety $X'$ obtained by blowing up $X$ along $Y$, imitating the method used to compute $\operatorname{Gr}^\bullet(X')_\Q$ [Exp. VII 4.8], which also works (without tensoring with $\Q$) for the computation of $\Gtop^\bullet(X')$. The key formula for this determination (with the notation of loc. cit.) is
\sgaStableTarget{sga6.e14.eq.4-8}\begin{equation}\tag{4.8}
 f^*(i_*(y))=j_*\bigl(g^*(y)c_{d-1}(\check{\underline F})\bigr),
\end{equation}
a formula currently established in $\Gtop^\bullet(X')$ (for $y\in\Gtop^\bullet(Y)$) by a computation of $\underline{\Tor}$ (Exp. VII 3.4), whose validity in $A(X')$ (for $y\in A(Y)$) remains to be established.

Observe that even the $\ell$-adic analogues of formulas (4.6), (4.7), (4.8) have not been proved at present; when $y$ is the cohomology class of an algebraic cycle, (4.6) and (4.8) hold in any case mod. torsion.

\sgaStableTarget{sga6.e14.para.4-5}4.5. In this section and the following ones, we shall give examples in which, in diagram (4.1), the homomorphism $j$, respectively the homomorphism $\phi$, is not an isomorphism. For a given $X$, it clearly amounts to the same thing to say that
\[
 j:\operatorname{Gr}^\bullet(X)\longrightarrow \Gtop^\bullet(X)
\]
is an isomorphism, or a monomorphism, or an epimorphism. Moreover, the image of $j$ is the subring of $\Gtop^\bullet(X)$ generated by the Chern classes of locally free sheaves on $X$. When $K^\bullet(X)$ is generated as an algebra over $\Z$ by the classes of invertible sheaves on $X$, one sees immediately that the subring of $\Gtop^\bullet(X)$ under consideration is also the subring generated by $\Gtop^1(X)$, the component of degree $1$ of $\Gtop^\bullet(X)$.
%%%%%%%%%%  scan idx 691  |  folio 678  %%%%%%%%%%
\SGAPrintPage{678}%
Now take for $X$ a \uline{flag variety}
\[
 X=G/B,
\]
where $G$ is a simply connected semisimple algebraic group over the algebraically closed field $k$, and $B$ is a Borel subgroup [5]. When $k=\C$, the Bruhat cell decomposition of $X$ [5, Exp. 13] implies that the canonical homomorphism
\sgaStableTarget{sga6.e14.eq.unnamed}\[
 (*)\qquad K^\bullet(X)\longrightarrow K(X(\C))
\]
is an isomorphism, as is moreover the canonical homomorphism
\sgaStableTarget{sga6.e14.eq.unnamed.occ2}\[
 (**)\qquad A(X)\longrightarrow H^*(X(\C),\Z),
\]
where $X(\C)$ is the compact space of complex points of $X$, endowed with its usual topology induced by the topology of $\C$. On the other hand, L. HODGKIN [12] has proved that $K(X(\C))$ is indeed generated by the classes of invertible sheaves on $X$. The same therefore holds for $K^\bullet(X)$ itself, and it is moreover not difficult to deduce, by a specialization technique using split semisimple groups over $\Z$ (SGA 3), that the same result holds over every algebraically closed field $k$. Moreover, the cell decomposition of $X$ (or the bijectivity of $(**)$ when $k=\C$) readily implies that $A(X)$ is torsion-free, so by 4.2 the homomorphism $\phi:A(X)\longrightarrow\Gtop^\bullet(X)$ is an isomorphism. These facts show that $j$ is an isomorphism if and only if $A(X)$ is generated by its elements of degree $1$, or again, when $k=\C$, if and only if the integral cohomology of $X(\C)$ is generated by the elements of degree $2$, i.e. if the classifying space $B_{G(\C)}$ is torsion-free. It is also known [15, p. 217 th. B (1)$\Leftrightarrow$(2), and 2.5] that this is so if and only if $G$ is "torsion-free", or equivalently if $G$ is isomorphic to a product of groups $\operatorname{S\ell}(n)$ and $\operatorname{Sp}(m)$. Thus it suffices to take for $G$ a group $\operatorname{SO}(n)$ with $n\geq 7$, or an exceptional simple group, to obtain an example in which $j$ is not an isomorphism.

\sgaStableTarget{sga6.e14.para.4-6}4.6. Observe that, in this example, there is no doubt that the "right" ring is the torsion-free ring $A(X)\simeq\Gtop^\bullet(X)$, with the cells of the Bruhat decomposition as a basis, a ring also isomorphic to the integral cohomology of $X(\C)$ when $k=\C$; on the other hand, the ring $\operatorname{Gr}^\bullet(X)$ appears to have torsion of a rather pathological nature. This example therefore confirms the view that the invariant $\operatorname{Gr}^\bullet(X)$ is of little
%%%%%%%%%%  scan idx 692  |  folio 679  %%%%%%%%%%
\SGAPrintPage{679}%
interest except modulo torsion, i.e. after tensoring with $\Q$. At the same time, this example shows that only after tensoring with $\Q$ is $\operatorname{Gr}^\bullet(X)$ covariant with respect to proper morphisms, even when restricted to closed immersions of smooth quasi-projective schemes. Indeed, for $X$ such that $j$ is not an isomorphism, there is in $X$ an integral closed subscheme $Y$ such that $\mathrm{c}\ell^\bullet(Y)\in K^\bullet(X)$ does not lie in $\operatorname{Filt}^i(X)$, where $i=\operatorname{codim}(Y,X)$. We may take $Y$ of minimal dimension, so that $\operatorname{Filt}^{i+1}(X)=\operatorname{Filt}_{\mathrm{top}}^{i+1}(X)$. With $Z$ the singular locus of $Y$ and $X'=X-Z$, $Y'=Y-Z$, we conclude that $\mathrm{c}\ell^\bullet(Y')\in K^\bullet(X')$ does not lie in $\operatorname{Filt}^i(X')$ (since $K^\bullet(X)\to K^\bullet(X')$ is surjective and compatible with the augmentations, and its kernel is $\operatorname{Im}\bigl(K(Z)\longrightarrow K(X)\bigr)\subset\operatorname{Filt}_{\mathrm{top}}^{i+1}(X)=\operatorname{Filt}^{i+1}(X)$). Now consider the inclusion
\[
 f:Y'\longrightarrow X',
\]
which is a closed immersion of codimension $i$ of smooth quasi-projective schemes over $k$, \uline{such that}
\[
 f_*:K^\bullet(Y')\longrightarrow K^\bullet(X')
\]
\uline{is not compatible with the} $\lambda$-\uline{filtrations} (\uline{up to a shift by} $i$), since $f_*(1)$ does not lie in $\operatorname{Filt}^i(X')$.

\sgaStableTarget{sga6.e14.para.4-7}4.7. Keeping the notation of 4.5, observe now that (Exp. IX) shows that the cited theorem of HODGKIN is equivalent to the existence of an isomorphism
\[
 K^\bullet(G)\simeq\Z,\qquad\text{therefore}\qquad \Gtop^\bullet(G)\simeq\Z,
\]
the homomorphism $K^\bullet(G/B)\to K^\bullet(G)$ being a priori surjective, with kernel the ideal generated by the $\mathrm{c}\ell^\bullet(\underline{L})-1$, where $\underline{L}$ ranges over the invertible Modules on $X=G/B$ associated with the representations $B\to\underline{G}_m$ (which in fact give all the invertible modules on $X$ [5, Exp. 15]). On the other hand, an analogous argument shows that
\[
 A(G/B)\longrightarrow A(G)
\]
is also surjective and that its kernel is the ideal generated by the elements of $A^1(G/B)$ defined by the representations $B\to\underline{G}_m$, thus in fact by all of $A^1(G/B)$.\SGAAuthorFootnote{*}{Cf. IX 3.4, 2.2.} This shows that $A(G)$ reduces to $A^0(X)=\Z$ if and
%%%%%%%%%%  scan idx 693  |  folio 680  %%%%%%%%%%
\SGAPrintPage{680}%
only if "$G$ is torsion-free", i.e. $G$ is isomorphic to a product of groups $\operatorname{S\ell}(n)$ and $\operatorname{Sp}(m)$. Thus in this case, and in this case alone, the homomorphism
\[
 \phi:A(G)\longrightarrow\Gtop^\bullet(G)
\]
is an isomorphism. Thus, when $G$ has torsion, the underlying scheme provides an example of a nonbijective $\phi$.

\sgaStableTarget{sga6.e14.para.4-8}4.8. Similarly, one finds an example in which the homomorphism $\gamma$ of (4.1)
\sgaStableTarget{sga6.e14.eq.4-9}\begin{equation}\tag{4.9}
 \gamma:A(G)\longrightarrow H^{2*}(G,\underline{Z}_\ell(*))
\end{equation}
does not factor through $\Gtop^\bullet(G)$, or, equivalently here, an example in which this homomorphism is nonzero in degrees $>0$. Restricting ourselves (as suitable specialization theorems allow) to the case $k=\C$, we are reduced to studying the homomorphism
\sgaStableTarget{sga6.e14.eq.4-10}\begin{equation}\tag{4.10}
 A(G)\longrightarrow H^*(G(\C),\Z),
\end{equation}
which is moreover deduced by passage to the quotient from the canonical homomorphism
\[
 A(G/B)\simeq H^*(G/B(\C),\Z)\longrightarrow H^*(G(\C),\Z).
\]
It is tempting to conjecture that the homomorphism (4.10) is in fact injective, since the torsion primes of $G(\C)$ are known to be exactly those dividing the order of the finite group $A^+(G)=\coprod_{i>0}A^i(G)$. Were this so, as soon as $G$ had torsion the homomorphism (4.10) would not factor through $\phi$, and thus, for suitable $\ell$ (more precisely, for $\ell$ dividing the order of $A^+(G)$, i.e. $\ell$ a torsion coefficient of $G$), the same would hold for (4.9). One can at least assert that, for the group $G=\underline{G}_2$ over the field $\C$, the homomorphism (4.10) is nonzero in degree $>0$; more precisely, there is an element of $A^3(G)$ whose image in $H^6(G(\C),\Z)$ has order $2$; consequently, the homomorphism (4.9) does not factor through $\phi$ when $\ell=2$. The fact used concerning $\underline{G}_2$, which amounts to the existence of an element of $H^6(G/B(\C),\Z)$ whose image in $H^6(G(\C),\Z)$ is nonzero, is found in [3], using the familiar dictionary between the languages of complex semisimple algebraic groups and compact Lie groups; it was kindly communicated to me by J.P. SERRE.

%%%%%%%%%%  scan idx 694  |  folio 681  %%%%%%%%%%
\SGAPrintPage{681}%
\sgaStableTarget{sga6.e14.sec.5}\section*{5. Relations between \(\operatorname{Gr}^\bullet(X)\) and \(H^{2*}(X,\underline{Z}_\ell(*))\)}

\sgaStableTarget{sga6.e14.para.5-1}5.1. The example given in 4.8 does not rule out the possibility of finding, for an arbitrary scheme $X$, a ring homomorphism
\[
 g^\bullet:\operatorname{Gr}^\bullet(X)\longrightarrow H^{2*}(X,\underline{Z}_\ell(*)),
\]
such that, for every locally free Module $\underline{E}$ on $X$, $g^i$ takes the Chern class $c_i(\underline{E})$ of Exp. V to the cohomological Chern class of SGA 5 VII. Of course, $\ell$ here denotes a prime distinct from the residue characteristics of $X$. Since the ring $\operatorname{Gr}^\bullet(X)$ is generated by the $c^i(\underline{E})$, one sees in any case that $g^\bullet$ is uniquely determined by the requirement of compatibility with Chern classes.

\sgaStableTarget{sga6.e14.para.5-2}5.2. Using the Chern homomorphism
\[
 K^\bullet(X)\longrightarrow\Ch\bigl(H^{2*}(X,\Zell(*))\bigr),
\]
and passing to the associated graded groups for the filtration of the first member and the degree filtration of the second, respectively, one obtains a homomorphism of graded groups (but not of rings!)
\[
 \psi:\operatorname{Gr}^\bullet(X)\longrightarrow H^{2*}(X,\Zell(*)),
\]
satisfying the relation
\[
 \psi(c_i(\underline{E}))=(-1)^{i-1}(i-1)!\,c_i^\ell(\underline{E}),
\]
where $c_i^\ell(\underline{E})$ denotes the $\ell$-adic Chern class. By division starting from the $\psi^i$, one obtains a ring homomorphism:
\[
 g_\Q:\operatorname{Gr}^\bullet(X)\longrightarrow
 H^{2*}(X,\Qell(*))=
 H^{2*}(X,\Zell(*))\otimes_{Z_\ell}\Qell,
\]
taking Chern classes to Chern classes, such that the image of the first member is contained in that of $H^{2*}(X,\Zell(*))$; consequently, the question of the existence of $g^\bullet$ has an affirmative answer upon replacing $H^{2*}(X,\Zell(*))$ by its quotient by its torsion subgroup. The first doubtful case in defining the components $g^i$ of the hypothetical homomorphism $g$ is that of $g^3$; a priori, one succeeds only in defining $2g^3=\psi^3$. The answer to the existence question for $g^3$ would be negative if, for example, one could find three invertible sheaves $\underline{L}_1,\underline{L}_2,\underline{L}_3$ on $X$ satisfying the relation
\[
 (\underline{L}_1-1)(\underline{L}_2-1)(\underline{L}_3-1)=0\quad\text{in }K^\bullet(X),
\]
%%%%%%%%%%  scan idx 695  |  folio 682  %%%%%%%%%%
\SGAPrintPage{682}%
but such that
\[
 c_1(\underline{L}_1)c_1(\underline{L}_2)c_1(\underline{L}_3)\ne0
 \quad\text{in }H^6(X,\mu_2^{\otimes3}).
\]

\sgaStableTarget{sga6.e14.para.5-3}5.3. When $X$ is smooth and quasi-projective over a field $k$, one can likewise ask about the existence of a ring homomorphism
\[
 g:\operatorname{Gr}^\bullet(X)\longrightarrow A(X)
\]
with values in the Chow ring $A(X)$, taking Chern classes to Chern classes. One can repeat the preceding considerations, which show in particular how to construct $(-1)^{i-1}(i-1)!g^i$ directly in the absence of $g^i$ itself.

\sgaStableTarget{sga6.e14.para.5-4}5.4. To conclude this section, observe that essentially the same question arises when, in place of a scheme $X$, one considers a compact topological space $X$, say; one can again consider the associated graded $\operatorname{Gr}^\bullet(X)$ of $K^\bullet(X)$ endowed with its $\lambda$-filtration, and seek a ring homomorphism
\[
 \operatorname{Gr}^\bullet(X)\longrightarrow H^{2*}(X,\Z)
\]
taking Chern classes to Chern classes.

\sgaStableTarget{sga6.e14.sec.6}\section*{6. Cohomological Riemann--Roch theorem and Gysin homomorphism}

\sgaStableTarget{sga6.e14.para.6-1}6.1. For simplicity, we restrict ourselves here to schemes $X$ admitting an ample invertible Module; this implies that the naive $K^\bullet(X)$ and the other one can be identified. Let us begin with the Riemann--Roch formula ``with denominators'' for a proper complete-intersection morphism $f:X\to Y$. Since (for a prime $\ell$ prime to the residue characteristics of $Y$) one can define the $\ell$-adic Chern classes of a locally free Module $\underline{E}$,
\[
 c_i(\underline{E})\in H^{2i}(X,\Zell(i))
\]
(SGA 5 VII), corresponding to a homomorphism
\sgaStableTarget{sga6.e14.eq.6-1}\begin{equation}\tag{6.1}
 K^\bullet(X)\longrightarrow \Ch\bigl(H^{2*}(X,\Zell(*))\bigr),
\end{equation}
the question arises of giving a cohomological variant of the Riemann--Roch formula. The only ingredient missing to give the formula a meaning is a \uline{Gysin homomorphism}
%%%%%%%%%%  scan idx 696  |  folio 683  %%%%%%%%%%
\SGAPrintPage{683}%
\sgaStableTarget{sga6.e14.eq.6-2}\begin{equation}\tag{6.2}
 H^n(X,\underline{Q}_\ell(i))\longrightarrow H^{n-2d}(Y,\underline{Q}_\ell(i-d)),
\end{equation}
where $d$ is the virtual relative dimension of $f$. Assuming this homomorphism defined, and taking account of the Riemann--Roch formula in the form of Exp. VIII with values in $\operatorname{Gr}^\bullet(Y)_{\Q}$, one finds that the validity of the $\ell$-adic Riemann--Roch formula for the morphism $f$ is exactly equivalent to the commutativity of the square
\sgaStableTarget{sga6.e14.eq.6-3}\begin{equation}\tag{6.3}
\begin{tikzcd}[column sep=large,row sep=large]
 \operatorname{Gr}^\bullet(X)_{\Q} \arrow[r] \arrow[d] & H^{2*}(X,\underline{Q}_\ell(*)) \arrow[d] \\
 \operatorname{Gr}^\bullet(Y)_{\Q} \arrow[r] & H^{2*}(Y,\Qell(*)),
\end{tikzcd}
\end{equation}
where the horizontal arrows are defined from (6.1) by passing to associated graded groups and dividing so as to take Chern classes to Chern classes, and the vertical arrows are respectively the arrow of Exp. VIII and the Gysin homomorphism in cohomology (assumed given).

\sgaStableTarget{sga6.e14.para.6-2}6.2. Observe that the Gysin homomorphism is in any case defined when $Y$ is smooth over a field $k$ (SGA 5 IV), using the relative cohomological purity theorem over $k$; the same construction applies when $Y$ is excellent, regular and of characteristic zero, by the absolute cohomological purity theorem of SGA 4 XIX, and it will apply automatically to every excellent regular $Y$, without restriction on the characteristic, once the absolute cohomological purity theorem has been proved in full generality. Finally, the same construction applies, by the relative cohomological purity theorem (SGA 4 XVI 3.7), when $X$ and $Y$ are both smooth over the same scheme $S$ and $f$ is an $S$-morphism. One should conjecture that, in each of these cases in which the Gysin homomorphism is defined, diagram (6.3) is indeed commutative. The only case in which this assertion is currently verified is that in which $X$ and $Y$ are both smooth over a field $k$; it then follows from the commutativity of the diagram analogous to (6.3):
\sgaStableTarget{sga6.e14.eq.6-4}\begin{equation}\tag{6.4}
\begin{tikzcd}[column sep=large,row sep=large]
 A(X) \arrow[r] \arrow[d] & H^{2*}(X,\Zell(*)) \arrow[d] \\
 A(Y) \arrow[r] & H^{2*}(Y,\Zell(*)),
\end{tikzcd}
\end{equation}
%%%%%%%%%%  scan idx 697  |  folio 684  %%%%%%%%%%
\SGAPrintPage{684}%
involving the rings of cycle classes (mod. rational equivalence) $A(X)$ and $A(Y)$, whose commutativity is established in SGA 5 IV.

\sgaStableTarget{sga6.e14.para.6-3}6.3. Observe that, in each of the cases considered above in which one can define a Gysin homomorphism, one in fact obtains a construction for cohomology with coefficients in $\Zell(*)$, and not only in $\Qell(*)$. When $f$ is an \uline{immersion}, this allows the question of the validity of the Riemann--Roch formula ``without denominators'' to be posed, as explained in no.~3. When $X$ and $Y$ are both smooth over a field $k$, the commutativity of (6.4) shows that it would suffice to prove the formula without denominators in the Chow ring $A(Y)$ cf. (3.3).

\sgaStableTarget{sga6.e14.para.6-4}6.4. Without a regularity hypothesis on $Y$, even when $X$ and $Y$ are of finite type over the field $\C$, one no longer sees a plausible definition of a cohomological Gysin homomorphism (6.2), even for a regular immersion of codimension $1$. Can one nevertheless find one with reasonable properties, in particular making square (6.3) commute? The question is clearly connected with that of defining, for a complete intersection $X$ of codimension $i$ in $Y$, ``its cohomology class''
\[
 \gamma(X)\in H^{2i}(Y,\Zell(i)),
\]
so as to take intersections (without excess components) to cup products. This latter question has an affirmative answer at least modulo torsion, i.e. in $H^{2*}(Y,\Qell(*))$, by defining
\[
 \gamma(X)=1/(-1)^{i-1}(i-1)!\,c_i(\underline{O}_X).
\]
But I do not know whether $c_i(\underline{O}_X)$ is divisible by $(i-1)!$ in $H^{2i}(Y,\Zell(i))$. In particular, for a projective scheme $Y$ of dimension $3$ over the field of complex numbers, say, with $X$ reduced to a (singular) point of $Y$, is $c_3(\underline{O}_X)$ divisible by $2$ in $H^6(Y,\Z_2(3))$, and in particular does the Chern class vanish in $H^6(Y,\Z/2\Z)$? In the same direction, let us also mention the following question: let $Y$ be a projective (singular) scheme of dimension $4$, and $X$ a closed subscheme of dimension $1$ that is globally an intersection of three positive divisors $D_1,D_2,D_3$; consider the product of their Chern classes mod. $2$:
\[
 c_1(D_1)c_1(D_2)c_1(D_3)\in H^6(Y,\Z/2\Z);
\]
%%%%%%%%%%  scan idx 698  |  folio 685  %%%%%%%%%%
\SGAPrintPage{685}%
is this element of $H^6$ independent of the particular choice of the $D_i$ with intersection $X$?

\sgaStableTarget{sga6.e14.sec.7}\section*{7. Chern classes of perfect complexes}

\sgaStableTarget{sga6.e14.para.7-1}7.1. In seeking to develop a cohomological Riemann--Roch formula, with or without denominators, beyond the rather artificial hypothesis that the schemes under consideration admit an ample invertible Module, one is necessarily confronted with the question of defining, for a perfect complex of Modules $L_\bullet$ on a scheme $X$, Chern classes
\[
 c_i(L_\bullet)\in H^{2i}(X,\Zell(i)),
\]
generalizing the Chern classes of locally free Modules in SGA 5 VII. The same question also arises for general locally ringed topoi (compare [10]) and for perfect complexes on paracompact spaces ringed by continuous complex functions, working then with ordinary integral cohomology. For a compact space, the $K^\bullet(X)$ constructed with perfect complexes is known to be the same as the ``naive'' one constructed with locally free Modules (Exp. IV), and in this case one is trivially reduced to the classical definition. But even for a locally compact space the problem posed seems nontrivial, since the integral cohomology of such a space is not always the inverse limit of the cohomologies of its compact subsets.

\sgaStableTarget{sga6.e14.para.7-2}7.2. More precisely, for a closed subset $Y$ of $X$ such that $L_\bullet|(X-Y)$ is zero (in the sense of derived categories), i.e. such that the cohomology sheaves of $L_\bullet$ have supports in $Y$, we would like to define Chern classes \uline{with support in} $Y$
\[
 c_i^Y(L_\bullet)\in H_Y^{2i}(X,\Zell(i))\quad\text{resp.}\quad H_Y^{2i}(X,\Z),
\]
which would recover the ``ordinary'' (sic) Chern classes of 7.1 by means of the canonical homomorphism
\[
 H_Y^*(X,-)\longrightarrow H^*(X,-).
\]
%%%%%%%%%%  scan idx 699  |  folio 686  %%%%%%%%%%
\SGAPrintPage{686}%
Even when these reduce to the Chern classes of locally free Modules ($X$ a scheme with an ample invertible Module, or a compact space), these Chern classes ``with supports'' have not been constructed, and their construction will doubtless have to proceed together with that of the general Chern classes postulated in 7.1\SGAAuthorFootnote{*}{Added in October 1968. The problems raised here have been solved in the setting of \uline{Hodge cohomology} by L. ILLUSIE; cf. [13]}

It is probable that Chern classes with supports will have a role to play in proving the cohomological Riemann--Roch formula without denominators (cf. 3.5). On the other hand, they should allow an improved version of the Riemann--Roch formula (with or without denominators) by introducing families of supports into that formula.

\sgaStableTarget{sga6.e14.sec.8}\section*{8. Chow ring of regular schemes}

The problem is to define rational equivalence of cycles and to study the resulting ring of classes (in particular its functorial properties) for a regular noetherian scheme $X$ admitting an ample invertible Module, generalizing the theory known when $X$ is smooth over a field [4]. The key technical result will again be a ``moving lemma'', ensuring that a cycle can be varied within its class so as to intersect a given cycle in the correct dimension.

The problem considered is quite natural a priori, particularly since we have seen (4.5 to 4.8) that the substitutes $\operatorname{Gr}^\bullet(X)$, $\Gtop^\bullet(X)$ for the Chow ring $A(X)$ are hardly satisfactory beyond a theory ``modulo torsion''. When $X$ is obtained from a smooth projective algebraic scheme $Y$ over a field $k$ by taking the local scheme at the vertex of a projecting cone $C$ and removing the origin, one finds (admitting that $A(X)$ is indeed the direct limit of the Chow rings of the $U-(O)$, where $U$ ranges over the open neighbourhoods in $C$ of the vertex $O$) that $A(X)$ is canonically isomorphic to the quotient ring $A(Y)/\xi A(Y)$, where $\xi\in A^1(Y)$ is induced by the projective immersion of $Y$ under consideration. This shows that knowledge of the Chow ring $A(X)$ (of a nonclassical type) is equivalent to knowledge of $A(Y)/\xi A(Y)$, which might be called the ``primitive part'' of the Chow ring $A(Y)$ (in a sense suggested by the Lefschetz--Hodge cohomological theory of nonsingular projective algebraic varieties). This example shows that the ``local Chow rings'' of local rings with a singularity
%%%%%%%%%%  scan idx 700  |  folio 687  %%%%%%%%%%
\SGAPrintPage{687}%
that is isolated are geometric invariants that are by no means trivial, and will undoubtedly have their part to play in the ``arithmetic'' study of these local rings.

Of course, once one has a reasonable notion of a Chow ring for a regular noetherian scheme admitting an ample invertible Module, one should generalize to this case the relations established or conjectured, respectively, in n\textsuperscript{o} 4 between this ring and the ring $\Gtop^\bullet(X)$, and formulate in terms of this ring the question of the ``Riemann--Roch formula without denominators'' stated in n\textsuperscript{o} 3.

\sgaStableTarget{sga6.e14.bibliography-of-expose-xiv}\subsection*{BIBLIOGRAPHY}
\begin{enumerate}[label={[\arabic*]}]
\item M. F. Atiyah and F. Hirzebruch, The Riemann--Roch theorem for analytic embeddings, Topology, vol. 1, p. 151-166 (1962).
\item A. Borel and J.P. Serre, Le théorème de Riemann--Roch, Bull. Soc. Math. France, vol. 86, p. 97-136 (1958).
\item R. Bott and H. Samelson, Amer. Journ. of Math. vol 80, p. 1004 (1958).
\item C. Chevalley, Les classes d'équivalence rationnelle, I and II, Seminar at the ENS 1958, exposés 2 and 3.
\item C. Chevalley, Classification des groupes de Lie algébriques, Seminar at the ENS 1956/58.
\item A. Dold and D. Puppe, Homologie nicht-additiver Funktoren. Anwendungen, Annales de l'Inst. Fourier t. XI, p. 201-312 (1961).
\item H. Grauert, Ein Theorem der analytischen Garbentheorie, Pub. Math. n\textsuperscript{o} 5, p. 5-64 (1960).
\item A. Grothendieck, Classes de faisceaux et théorème de Riemann--Roch, mimeographed notes prepared by the Institute for Advanced Study, 1957 (work cited as [RRR] in the present exposé, and reproduced in our seminar as an Appendix to Chap. O).
\item A. Grothendieck, La théorie des classes de Chern, Bull. Soc. Math. France vol. 86, p. 137-154 (1958).
%%%%%%%%%%  scan idx 701  |  folio 688  %%%%%%%%%%
\SGAPrintPage{688}%
\item A.
Grothendieck, Classes de Chern et représentations linéaires des groupes discrets, in: Dix exposés sur la cohomologie des schémas, Masson--North Holland Pub. Cie
\item H. Hironaka, Resolution of singularities of an algebraic variety over a field of characteristic zero, Ann. of Math. 79, p. 109-326 (1964).
\item L. Hodgkin, thesis (to be published?).
\item L. Illusie, work in preparation.
\item J.P. Serre, Algèbre Locale, Multiplicités, Lecture Notes in Math. n\textsuperscript{o} 11 (1965) Springer.
\item A. Borel, Sous-groupes commutatifs et torsion des groupes de Lie compacts connexes, Tohoku Math. Journal, vol. 13, p. 216-240 (1961).
\item[{[16]}(*)] P. Deligne, Exp. XVIII in Théorie des topos et cohomologie étale des schémas (SGA 4), by M. Artin, A. Grothendieck and J. L. Verdier.
\item[{[17]}(*)] U. Manin, Correspondances, motifs et transformations monoïdales (in Russian). Mat. Sbornik, t 77 (119) n\textsuperscript{o} 4, 1968.
\end{enumerate}

\SGAAuthorFootText{*}{Added in February 1969}

%%%%%%%%%%  scan idx 702  |  folio 689  %%%%%%%%%%
\SGAPrintPage{689}%
\sgaStableTarget{sga6.e14.note-for-page-1}\subsection*{Note for page 1.}

The theory of the operations $\Lambda^i$ in $D^-(X)$, and of the operations $\lambda^i$ in $K^\bullet(X)$, was written down by P. Deligne in 1968. He proves that $K^\bullet(X)$ thereby becomes a pre-$\lambda$-ring, and, when $\underline{O}_X$ is a sheaf of Algebras over a field, that $K^\bullet(X)$ is even a $\lambda$-ring. It is plausible that a refinement of Deligne's method would make it possible to dispense with the restrictive hypothesis on $\underline{O}_X$.

\sgaStableTarget{sga6.e14.notes-for-page-11}\subsection*{Notes for page 11}

\noindent(*) (April 1969) D. MUMFORD points out to us a proof of (4.6) in the \uline{Chow ring} (without disregarding torsion), which he obtained in a seminar at Harvard in 1959! The reader will find the proof of this beautiful result in SGA 5 VII 9. This result obviously contains the corresponding cohomological result mentioned in note (**). However, the cohomological result is proved by J. P. JOUANOLOU under more general conditions (without a quasi-projectivity hypothesis, and in the relative smooth case over an arbitrary base).

\medskip

\noindent(**) (January 1969) J. P. JOUANOLOU has just proved these formulas for every $y$ without disregarding torsion, cf. SGA 5 VII \S{} 4.9. DELIGNE and MANIN 17 had previously observed independently, and by (easy) arguments of a rather different nature, that it was enough to establish formulas for ``algebraic'' $y$, and even for $y=1$ in 16, to conclude the same relation in the general case.

% Final-index helper macros
\providecommand{\Qcoh}{\operatorname{Qcoh}}
\providecommand{\Coh}{\operatorname{Coh}}
\providecommand{\Parf}{\operatorname{Parf}}
\providecommand{\Lib}{\operatorname{Lib}}
\providecommand{\Loclib}{\operatorname{Loclib}}
\providecommand{\Proj}{\operatorname{Proj}}
\providecommand{\Mod}{\operatorname{Mod}}
\renewcommand{\Ob}{\operatorname{Ob}}
\providecommand{\Hom}{\operatorname{Hom}}
\providecommand{\toramp}{\operatorname{tor.amp}}
\providecommand{\parfamp}{\operatorname{par.amp}}
\providecommand{\coh}{\operatorname{coh}}
\providecommand{\qcoh}{\operatorname{qcoh}}
\providecommand{\parf}{\operatorname{parf}}
\providecommand{\torf}{\operatorname{torf}}
\providecommand{\amp}{\operatorname{amp}}
\providecommand{\num}{\operatorname{num}}
\providecommand{\naif}{\operatorname{naif}}
\providecommand{\alg}{\operatorname{alg}}
\providecommand{\Pic}{\operatorname{Pic}}
\providecommand{\Gr}{\operatorname{Gr}}
\providecommand{\Fil}{\operatorname{Fil}}
\providecommand{\Ch}{\operatorname{Ch}}
\providecommand{\Todd}{\operatorname{Todd}}
\providecommand{\Diff}{\operatorname{Diff}}
\providecommand{\Dcat}{\mathrm D}
\providecommand{\Tr}{\operatorname{Tr}}
\providecommand{\Z}{\mathbb Z}
\providecommand{\Q}{\mathbb Q}

\newpage
%%%%%%%%%%  scan idx 703  |  folio 690  %%%%%%%%%%
\SGAPrintPage{690}%

%%%%%%%%%%  scan idx 704  |  folio 691  %%%%%%%%%%
\SGAPrintPage{691}%

\sgaStableTarget{sga6.back.terminological-index}\section*{Terminological index}
\begingroup\small
\begin{longtable}{@{}>{\raggedright\arraybackslash}p{0.64\linewidth}>{\raggedright\arraybackslash}p{0.14\linewidth}>{\raggedright\arraybackslash}p{0.12\linewidth}@{}}
\toprule
 & Exp. & §\\
\midrule
additive (map) & I & \hyperlink{sga6.e01.def.6-1}{6.1}\\
additive (map) & IV & \hyperlink{sga6.e04.para.1-5}{1.5}\\
augmented \(K-\lambda\)-algebra & V & \hyperlink{sga6.e05.def.3-9}{3.9}\\
algebraically equivalent to zero (in \(K^{\bullet}\)) & X App. & \hyperlink{sga6.e10.para.7-11}{7.11}\\
algebraically equivalent to zero (invertible sheaf) & XIII & \hyperlink{sga6.e13.para.4-4}{4.4}\\
of \(\underline{C}_0\)-perfect amplitude contained in \([a,b]\) (complex) & I & \hyperlink{sga6.e01.def.4-7}{4.7}\\
of finite \(\underline{C}_0\)-perfect amplitude (complex) & I & \hyperlink{sga6.e01.def.4-7}{4.7}\\
of \(\underline{C}_0\) perfect amplitude contained in \([a,b]\) (complex) & I & \hyperlink{sga6.e01.ex.4-8}{4.8}\\
of finite perfect amplitude (complex) & I & \hyperlink{sga6.e01.ex.4-8}{4.8}\\
of flat amplitude contained in \([a,b]\) (complex) & I & \hyperlink{sga6.e01.def.5-2}{5.2}\\
of finite flat amplitude (complex) & I & \hyperlink{sga6.e01.def.5-2}{5.2}\\
of perfect amplitude relative to \(f\) (or \(Y\)) contained in \([a,b]\) (complex) & III & \hyperlink{sga6.e03.def.4-1}{4.1}\\
of perfect amplitude contained in \([-n,0]\) (morphism) & III & \hyperlink{sga6.e03.def.4-1}{4.1}\\
of finite perfect amplitude relative to \(f\) (complex) & III & \hyperlink{sga6.e03.def.4-1}{4.1}\\
of flat amplitude relative to \(f\) (or \(Y\)) contained in \([a,b]\) (complex) & III & \hyperlink{sga6.e03.def.3-1}{3.1}\\
of flat amplitude contained in \([-n,0]\) (morphism) & III & \hyperlink{sga6.e03.def.3-1}{3.1}\\
of finite flat amplitude relative to \(f\) (complex) & III & \hyperlink{sga6.e03.def.3-1}{3.1}\\
projective amplitude & I & \hyperlink{sga6.e01.rem.4-12-1}{4.12.1}\\
\(\lambda\)-ring & V & \hyperlink{sga6.e05.def.2-4}{2.4}\\
\(\lambda\)-ring generated by generators and relations & V & \hyperlink{sga6.e05.def.4-6}{4.6}\\
free \(\lambda\)-ring generated by a family of generators & V & \hyperlink{sga6.e05.def.4-4}{4.4}\\
Chern ring & V & \hyperlink{sga6.e05.def.6-2}{6.2}\\
polynomial map & X & \hyperlink{sga6.e10.def.2-1-1}{2.1.1}\\
binomial (ring) & V & \hyperlink{sga6.e05.def.2-7}{2.7}\\
%%%%%%%%%%  scan idx 705  |  folio 692  %%%%%%%%%%
\SGAPrintPage{692}%
Chern character & V & \hyperlink{sga6.e05.def.6-4}{6.4}\\
Todd character & V & \hyperlink{sga6.e05.para.1-16}{1.16}\\
abelian \(\underline S\)-category & I & \hyperlink{sga6.e01.para.1-1-1}{1.1.1}\\
flat abelian \(\underline S\)-category & I & \hyperlink{sga6.e01.para.1-1-1}{1.1.1}\\
additive \(\underline S\)-category & I & \hyperlink{sga6.e01.para.1-1-1}{1.1.1}\\
triangulated \(\underline S\)-category & I & \hyperlink{sga6.e01.para.1-1-1}{1.1.1}\\
Chern class & V & \hyperlink{sga6.e05.para.6-7}{6.7}\\
completed Chern class & V & \hyperlink{sga6.e05.para.6-7}{6.7}\\
total Chern class & V & \hyperlink{sga6.e05.para.6-7}{6.7}\\
codimension (of a regular immersion) & VII & \hyperlink{sga6.e07.def.1-6}{1.6}\\
coherent (complex) & I & \hyperlink{sga6.e01.def.3-1}{3.1}\\
coherer & II & \hyperlink{sga6.e02.sec.3.2}{3.2}\\
relative cotangent complex & VIII & \hyperlink{sga6.e08.def.2-3}{2.3}\\
elliptic complex & II App. II & \hyperlink{sga6.e02.para.1-2-5}{1.2.5}\\
cup-product & I & \hyperlink{sga6.e01.para.8-2}{8.2}\\
special positive \(k\)-cycle & XIII & \hyperlink{sga6.e13.def.6-9}{6.9}\\
degree (of a polynomial map) & X & \hyperlink{sga6.e10.def.2-1-1}{2.1.1}\\
degree (of an invertible sheaf) & XIII & \hyperlink{sga6.e13.para.2-2}{2.2}\\
degree (of a morphism of schemes) & X & \hyperlink{sga6.e10.sec.4.4}{4.4}\\
dense (subset of an abelian or triangulated category) & IV & \hyperlink{sga6.e04.def.1-7}{1.7}\\
perfect dimension & I & \hyperlink{sga6.e01.def.4-14-1}{4.14.1}\\
virtual relative dimension & VIII & \hyperlink{sga6.e08.def.1-9}{1.9}\\
divisorial & II & \hyperlink{sga6.e02.def.2-2-5}{2.2.5}\\
dual & I & \hyperlink{sga6.e01.def.7-3}{7.3}\\
perfect duality & I & \hyperlink{sga6.e01.cor.7-8}{7.8}\\
numerical equivalence (in \(\operatorname{Gr}^{\bullet}(X)\) and \(\operatorname{Gr}_{\cdot}(X)\)) & X & \hyperlink{sga6.e10.sec.4.5}{4.5}\\
\(\tau\)-equivalence & X & \hyperlink{sga6.e10.def.2-4-1}{2.4.1}\\
\qquad'' & XIII & \hyperlink{sga6.e13.para.4-4}{4.4}\\
complex pseudo-analytic space & II & \hyperlink{sga6.e02.def.2-4-7}{2.4.7}\\
complex Stein pseudo-analytic space & II & \hyperlink{sga6.e02.def.2-4-7}{2.4.7}\\
exact (subset of an abelian or triangulated category) & IV & \hyperlink{sga6.e04.def.1-7}{1.7}\\
\((b)\)-sheaf & XIII & \hyperlink{sga6.e13.def.1-5}{1.5}\\
conormal sheaf & VII & \hyperlink{sga6.e07.def.1-6}{1.6}\\
relative cotangent sheaf & II App. II & \hyperlink{sga6.e02.para.1-2-1}{1.2.1}\\
%%%%%%%%%%  scan idx 706  |  folio 693  %%%%%%%%%%
\SGAPrintPage{693}%
relative tangent sheaf & II App. II & \hyperlink{sga6.e02.para.1-2-1}{1.2.1}\\
\((b)\)-family & XIII & \hyperlink{sga6.e13.para.1-12}{1.12}\\
ample family & II & \hyperlink{sga6.e02.def.2-2-4}{2.2.4}\\
family of classes of coherent sheaves & XIII & \hyperlink{sga6.e13.para.1-12}{1.12}\\
bounded family of classes of coherent sheaves & XIII & \hyperlink{sga6.e13.para.1-12}{1.12}\\
family of classes of coherent sheaves bounded by a coherent sheaf & XIII & \hyperlink{sga6.e13.para.1-12}{1.12}\\
relative cotangent bundle & II App. II & \hyperlink{sga6.e02.para.1-2-1}{1.2.1}\\
relative tangent bundle & II App. II & \hyperlink{sga6.e02.para.1-2-1}{1.2.1}\\
faithfully flat (morphism of ringed topoi) & I & \hyperlink{sga6.e01.exercise.2-16-2}{2.16.2}\\
filtration of an augmented \(\lambda\)-algebra & V & \hyperlink{sga6.e05.para.3-10}{3.10}\\
additive \(\underline S\)-functor & I & \hyperlink{sga6.e01.para.1-1-3}{1.1.3}\\
exact \(\underline S\)-functor & I & \hyperlink{sga6.e01.para.1-1-3}{1.1.3}\\
projection formula & IV & \hyperlink{sga6.e04.para.2-11}{2.11}\\
Grothendieck group & IV & \hyperlink{sga6.e04.def.2-2}{2.2}\\
\(\lambda\)-homomorphism & V & \hyperlink{sga6.e05.def.2-1}{2.1}\\
specialization homomorphism & X App. & \hyperlink{sga6.e10.prop.7-3}{7.3}\\
\qquad''\hspace{8em}'' & X App. & \hyperlink{sga6.e10.sec.7.16}{7.16}\\
\(\lambda\)-ideal & V & \hyperlink{sga6.e05.def.4-5}{4.5}\\
analytic index & II App. II & \hyperlink{sga6.e02.para.3-2-2}{3.2.2}\\
complete intersection (morphism) & VIII & \hyperlink{sga6.e08.def.1-1}{1.1}\\
\(n\)-isomorphism & I & \hyperlink{sga6.e01.def.1-4-3}{1.4.3}\\
homogeneous free (graded module) & VI & \hyperlink{sga6.e06.eq.1-2-2}{1.2.2}\\
locally stable under kernels of epimorphisms & I & \hyperlink{sga6.e01.sec.1.2}{1.2}\\
locally quasi-liftable & I & \hyperlink{sga6.e01.sec.1.2}{1.2}\\
locally quasi-stable under kernels of epimorphisms & I & \hyperlink{sga6.e01.sec.1.2}{1.2}\\
locally liftable & I & \hyperlink{sga6.e01.sec.1.2}{1.2}\\
Chern number & X & \hyperlink{sga6.e10.ex.4-2}{4.2}\\
intersection number & X & \hyperlink{sga6.e10.sec.4.1}{4.1}\\
intersection number (for a family of invertible sheaves) & X & \hyperlink{sga6.e10.sec.4.4}{4.4}\\
numerically equivalent to zero (invertible sheaf) & XIII & \hyperlink{sga6.e13.para.4-4}{4.4}\\
universal differential operator of order \(n\) & II App. II & \hyperlink{sga6.e02.para.1-2-1}{1.2.1}\\
%%%%%%%%%%  scan idx 707  |  folio 694  %%%%%%%%%%
\SGAPrintPage{694}%
Adams operations & V & \hyperlink{sga6.e05.para.7-1}{7.1}\\
\(\underline{C}_0\)-perfect (complex) & I & \hyperlink{sga6.e01.def.4-7}{4.7}\\
perfect (complex) & I & \hyperlink{sga6.e01.ex.4-8}{4.8}\\
perfect (morphism) & III & \hyperlink{sga6.e03.def.4-1}{4.1}\\
perfect relative to \(f\) (complex) & III & \hyperlink{sga6.e03.def.4-1}{4.1}\\
precoherent & I & \hyperlink{sga6.e01.def.3-1}{3.1}\\
of finite \(n\)-\(\underline{C}_0\)-presentation & I & \hyperlink{sga6.e01.def.2-8}{2.8}\\
of finite presentation & I & \hyperlink{sga6.e01.def.2-8}{2.8}\\
of finite \(\infty\)-\(\underline{C}_0\)-presentation & I & \hyperlink{sga6.e01.def.2-8}{2.8}\\
\(n\)-\(\underline{C}_0\)-pseudo-coherent (complex) & I & \hyperlink{sga6.e01.prop.2-2}{2.2}\\
\(\underline{C}_0\)-pseudo-coherent (complex) & I & \hyperlink{sga6.e01.prop.2-2}{2.2}\\
\(n\)-pseudo-coherent (complex) & I & \hyperlink{sga6.e01.ex.2-15}{2.15}\\
\(n\)-pseudo-coherent (morphism) & III & \hyperlink{sga6.e03.def.1-2}{1.2}\\
pseudo-coherent (complex) & I & \hyperlink{sga6.e01.ex.2-15}{2.15}\\
pseudo-coherent (morphism) & III & \hyperlink{sga6.e03.def.1-2}{1.2}\\
\(n\)-pseudo-coherent relative to \(f\) (complex) & III & \hyperlink{sga6.e03.def.1-2}{1.2}\\
pseudo-coherent relative to \(f\) (complex) & III & \hyperlink{sga6.e03.def.1-2}{1.2}\\
\((b)\)-polynomial & XIII & \hyperlink{sga6.e13.def.1-9}{1.9}\\
Snapper polynomial & X & \hyperlink{sga6.e10.prop.2-3-1}{2.3.1}\\
universal polynomial & V & \hyperlink{sga6.e05.para.1-15}{1.15}\\
pre-\(\lambda\)-ring & V & \hyperlink{sga6.e05.def.2-1}{2.1}\\
\(n\)-quasi-isomorphism & I & \hyperlink{sga6.e01.def.1-4-3}{1.4.3}\\
rank (of a perfect complex) & I & \hyperlink{sga6.e01.def.6-11}{6.11}\\
\(m\)-regular (sheaf) & XIII & \hyperlink{sga6.e13.def.1-1}{1.1}\\
regular (homomorphism) & VII & \hyperlink{sga6.e07.def.1-1}{1.1}\\
regular (ideal) & VII & \hyperlink{sga6.e07.def.1-4}{1.4}\\
regular (immersion) & VII & \hyperlink{sga6.e07.def.1-4}{1.4}\\
triangulated sub-\(\underline S\)-category & I & \hyperlink{sga6.e01.para.1-1-4}{1.1.4}\\
special (\(\lambda\)-ring) & V & \hyperlink{sga6.e05.def.2-4}{2.4}\\
special (morphism) & III & \hyperlink{sga6.e03.lem.1-1-4}{1.1.4}\\
Stein (locally ringed topos) & II & \hyperlink{sga6.e02.def.2-4-1}{2.4.1}\\
strictly \(\underline{C}_0\)-perfect & I & \hyperlink{sga6.e01.def.2-1}{2.1}\\
strictly \(\underline{C}_0\)-\(n\)-pseudo-coherent & I & \hyperlink{sga6.e01.def.2-1}{2.1}\\
strictly \(\underline{C}_0\)-pseudo-coherent & I & \hyperlink{sga6.e01.def.2-1}{2.1}\\
minimum system of generators & VII & \hyperlink{sga6.e07.def.1-4}{1.4.1}\\
%%%%%%%%%%  scan idx 708  |  folio 695  %%%%%%%%%%
\SGAPrintPage{695}%
of finite Tor-dimension & I & \hyperlink{sga6.e01.def.5-2}{5.2}\\
Tor-independence & III & \hyperlink{sga6.e03.def.1-5}{1.5}\\
trace & I & \hyperlink{sga6.e01.para.8-1}{8.1}\\
of finite \(\underline{C}_0\)-type & I & \hyperlink{sga6.e01.sec.1.2}{1.2}\\
mixed variety of class \(C^r\) & II App. II & \hyperlink{sga6.e02.sec.1.1}{1.1}\\
\bottomrule
\end{longtable}
\endgroup
\newpage
%%%%%%%%%%  scan idx 709  |  folio 696  %%%%%%%%%%
\SGAPrintPage{696}%
\sgaStableTarget{sga6.back.index-of-notations}\section*{Index of notations}
\begingroup\small
\begin{longtable}{@{}>{\raggedright\arraybackslash}p{0.50\linewidth}>{\raggedright\arraybackslash}p{0.17\linewidth}>{\raggedright\arraybackslash}p{0.15\linewidth}@{}}
\toprule
 & Exposé & §\\
\midrule
\endhead
\(A_N\) & V & \hyperlink{sga6.e05.para.5-5}{5.5}\\
\(\ch\) & V & \hyperlink{sga6.e05.para.6-3}{6.3}\\
\(\Ch(A),\ \Ch_K(A)\) & V & \hyperlink{sga6.e05.def.6-2}{6.2}\\
\((\Ch(A))_n\) & V & \hyperlink{sga6.e05.prop.6-5}{6.5}\\
\(\cl_C\) & IV & \hyperlink{sga6.e04.para.1-1}{1.1}\\
\(\Coh(S)\) & II & \hyperlink{sga6.e02.notation.2-2-1}{2.2.1}\\
\(\underline{C}_0\text{-}\parf\text{-}\amp(F)\) & I & \hyperlink{sga6.e01.def.4-7}{4.7}\\
\(C^p_{\num}(X)\) & XIII & \hyperlink{sga6.e13.prop.5-2}{5.2}\\
\(C^r(V),\ C^\infty(V)\) & II App. II & \hyperlink{sga6.e02.sec.1.1}{1.1}\\
\(c(x),\ c^i(x),\ \widetilde c(x)\) & V & \hyperlink{sga6.e05.para.6-7}{6.7}\\
\(\Dcat({}_{A}S),\ \Dcat({}_{A}\underline S)\) & I & \hyperlink{sga6.e01.ex.2-15}{2.15}\\
\(\Dcat({}_{A}\underline S)_{\coh},\ \Dcat({}_{A}\underline S)_{n\text{-}\coh}\) & I & \hyperlink{sga6.e01.ex.2-15}{2.15}\\
\(\Dcat({}_{A}\underline S)_{\parf}\) & I & \hyperlink{sga6.e01.notation.4-9}{4.9}\\
\(\Dcat^{-}({}_{A}\underline S)_{\torf},\ \Dcat^{-}({}_{A}X)_{\torf}\) & I & \hyperlink{sga6.e01.para.5-3}{5.3}\\
\(\Dcat(\underline{C})_{\underline{C}_0\text{-}\coh},\ \Dcat(\underline{C})_{\underline{C}_0\text{-}n\text{-}\coh}\) & I & \hyperlink{sga6.e01.schol.2-6}{2.6}\\
\(\Dcat(\underline{C})_{\underline{C}_0\text{-}\parf}\) & I & \hyperlink{sga6.e01.notation.4-9}{4.9}\\
\(\Dcat(f)_{\coh}\) & III & \hyperlink{sga6.e03.def.1-2}{1.2}\\
\(\Dcat(f)_{n\text{-}\coh}\) & III & \hyperlink{sga6.e03.def.1-2}{1.2}\\
\(\Dcat(f)_{\parf},\ \Dcat(f)_{\underline{\parf}}\) & III & \hyperlink{sga6.e03.notation.4-2}{4.2}\\
\(\Dcat(f)_{\torf}\) & III & \hyperlink{sga6.e03.notation.3-2}{3.2}\\
\(\Diff^n(E,F)\) & II App. II & \hyperlink{sga6.e02.para.1-2-2}{1.2.2}\\
\(d(L)\) & XIII & \hyperlink{sga6.e13.para.2-2}{2.2}\\
\(d^n_{X/S}\) & II App. II & \hyperlink{sga6.e02.para.1-2-1}{1.2.1}\\
\(d^n_{X/S}(E)\) & II App. II & \hyperlink{sga6.e02.para.1-2-2}{1.2.2}\\
%%%%%%%%%%  scan idx 710  |  folio 697  %%%%%%%%%%
\SGAPrintPage{697}%
\(\Dcat^{-}(\Qcoh(S))_{\coh},\ \Dcat(\Qcoh(S))_{\parf}\) & II & \hyperlink{sga6.e02.notation.2-2-1}{2.2.1}\\
\(\Dcat(S),\ \Dcat(\underline S),\ \Dcat(\underline S,A)\) & I & \hyperlink{sga6.e01.ex.2-15}{2.15}\\
\(\Dcat(X)_{f\text{-}\parf},\ \Dcat(X)_{f\text{-}\underline{\parf}}\) & III & \hyperlink{sga6.e03.notation.4-2}{4.2}\\
\(\Dcat(X)_{f\text{-}\torf}\) & III & \hyperlink{sga6.e03.notation.3-2}{3.2}\\
\(\Dcat(X)_{\qcoh}\) & III & \hyperlink{sga6.e03.lem.1-7}{1.7}\\
\(\Dcat(X)_{Y\text{-}\coh},\ \Dcat(X)_{Y\text{-}n\text{-}\coh}\) & III & \hyperlink{sga6.e03.def.1-2}{1.2}\\
\(\Dcat(X)_{Y\text{-}\parf},\ \Dcat(X)_{Y\text{-}\underline{\parf}}\) & III & \hyperlink{sga6.e03.notation.4-2}{4.2}\\
\(\Dcat(X)_{Y\text{-}\torf}\) & III & \hyperlink{sga6.e03.notation.3-2}{3.2}\\
\(\Dcat(Y)_{n\text{-}\coh}\) & III & \hyperlink{sga6.e03.thm.1-1}{1.1}\\
\(E^\vee\) & I & \hyperlink{sga6.e01.def.7-3}{7.3}\\
\(*\) & V & \hyperlink{sga6.e05.para.6-1}{6.1}\\
\(\lvert\underline F\rvert\) & X & \hyperlink{sga6.e10.prop.4-3-1}{4.3.1}\\
\(f^*\) (between the \(K^{\bullet}\)) & IV & \hyperlink{sga6.e04.para.2-7}{2.7}\\
\(f^*\) (between the \(K_{\cdot}\)) & IV & \hyperlink{sga6.e04.para.2-12}{2.12}\\
\(f_*\) (between the \(K_{\cdot}\)) & IV & \hyperlink{sga6.e04.para.2-11}{2.11}\\
\(f_*\) (between the \(K^{\bullet}\)) & IV & \hyperlink{sga6.e04.para.2-12}{2.12}\\
\(f_*\) (projective bundle case) & VI & \hyperlink{sga6.e06.sec.5.2}{5.2}\\
\(f_{\operatorname{gr}}\) & VIII & \hyperlink{sga6.e08.prop.3-2}{3.2}\\
\(\Filt_n(X)\) & X & \hyperlink{sga6.e10.def.1-1-1}{1.1.1}\\
\(F'_{>n},\ F'_{\leq n}\) & I & \hyperlink{sga6.e01.prop.2-12}{2.12}\\
\(\varphi^a\) & V & \hyperlink{sga6.e05.para.2-9}{2.9}\\
\(\Ghat,\ \Ghat_K\) & V & \hyperlink{sga6.e05.para.1-1}{1.1}\\
\(\Gcirc,\ \Gnabla\) & V & \hyperlink{sga6.e05.para.3-7}{3.7}\\
\(\Ga\) & V & \hyperlink{sga6.e05.para.1-3}{1.3.3}\\
\(\Gg\) & V & \hyperlink{sga6.e05.para.1-3}{1.3.2}\\
\(\operatorname{Gr}^{\bullet}(X),\ \operatorname{Gr}_{\cdot}(X)\) & X & \hyperlink{sga6.e10.sec.3.1}{3.1}\\
\(\operatorname{Gr}^{\bullet}_{\alg},\ \operatorname{Gr}_{\cdot}^{\alg}\) & X App. & \hyperlink{sga6.e10.para.7-11}{7.11}\\
%%%%%%%%%%  scan idx 711  |  folio 698  %%%%%%%%%%
\SGAPrintPage{698}%
\(\operatorname{Gr}^{\bullet}_{\num},\ \operatorname{Gr}_{\cdot}^{\num}\) & X App. & \hyperlink{sga6.e10.para.7-9-4}{7.9.4}\\
\(\gamma^n,\ \gamma_t\) & V & \hyperlink{sga6.e05.prop.3-2}{3.2}\\
\(\gamma^p(N,x)\) & V & \hyperlink{sga6.e05.para.5-5}{5.5}\\
\(\widehat G^t,\ \widehat G^u\) & V & \hyperlink{sga6.e05.para.2-3}{2.3}\\
\(\underline{\Hom}_{\underline{M}}(E,F)\) & I & \hyperlink{sga6.e01.notation.4-5}{4.5}\\
\(k(A)\) & IV & \hyperlink{sga6.e04.para.1-5}{1.5}\\
\(K^{\bullet}_{\alg},\ K_{\cdot}^{\alg}\) & X App. & \hyperlink{sga6.e10.para.7-11}{7.11}\\
\(K^{\bullet}({}_{A}X),\ K_{\cdot}({}_{A}X)\) & IV & \hyperlink{sga6.e04.def.2-2}{2.2}\\
\(K^{\bullet}({}_{A}X)_{\naif},\ K_{\cdot}({}_{A}X)_{\naif}\) & IV & \hyperlink{sga6.e04.def.2-2}{2.2}\\
\(k(C)\) & I & \hyperlink{sga6.e01.para.6-3}{6.3}\\
\(K^{\bullet}(f),\ K_{\cdot}(f)\) & IV & \hyperlink{sga6.e04.para.3-3-1}{3.3.1}\\
\(K^{\bullet}_{\num},\ K_{\cdot}^{\num}\) & X App. & \hyperlink{sga6.e10.para.7-9-4}{7.9.4}\\
\(K^{\bullet}(X),\ K_{\cdot}(X)\) & IV & \hyperlink{sga6.e04.def.2-2}{2.2}\\
\(K^{\bullet}(X)_{\naif},\ K_{\cdot}(X)_{\naif}\) & IV & \hyperlink{sga6.e04.def.2-2}{2.2}\\
\(K^{\bullet}(X/Y),\ K_{\cdot}(X/Y)\) & IV & \hyperlink{sga6.e04.para.3-3-1}{3.3.1}\\
\(\Lib(U)\) & II & \hyperlink{sga6.e02.sec.2.0}{2.0}\\
\(\operatorname{Loc\,lib}(U)\) & II & \hyperlink{sga6.e02.sec.2.0}{2.0}\\
\(\underline L_{\cdot}^{X/Y}\) & VIII & \hyperlink{sga6.e08.para.2-1}{2.1}\\
\(\lambda^i\) & V & \hyperlink{sga6.e05.def.2-1}{2.1}\\
\(\lambda_{-1}(N)\) & V & \hyperlink{sga6.e05.para.5-1}{5.1}\\
\(\lambda_{-1}(N)\) & VII & \hyperlink{sga6.e07.para.2-6}{2.7}\\
\(\lambda^p(N,x)\) & V & \hyperlink{sga6.e05.def.5-3}{5.3}\\
\(\lambda_t\) & V & \hyperlink{sga6.e05.def.2-1}{2.1}\\
\(\Mod(U)\) & II & \hyperlink{sga6.e02.sec.2.0}{2.0}\\
%%%%%%%%%%  scan idx 712  |  folio 699  %%%%%%%%%%
\SGAPrintPage{699}%
\(\operatorname{par.amp}_f(E),\ \operatorname{par.amp}_Y(E)\) & III & \hyperlink{sga6.e03.def.4-1}{4.1}\\
\(\Parf({}_{A}\underline S),\ \Parf(\underline{C},\underline{C}_0),\ \Parf(\underline S)\) & I & \hyperlink{sga6.e01.notation.4-9}{4.9}\\
\(\Pic^N(X)\) & X & \hyperlink{sga6.e10.prop.4-5-2}{4.5.2}\\
\(\Pic^\circ(X)\) & X & \hyperlink{sga6.e10.def.2-4-1}{2.4.1}\\
\(\underline{\Pic}^\circ_{X/k}\) & X & \hyperlink{sga6.e10.sec.2.4}{2.4}\\
\(\Pic^\tau(X)\) & X & \hyperlink{sga6.e10.def.2-4-1}{2.4.1}\\
\(\Pic(X)\) & X & \hyperlink{sga6.e10.sec.2.4}{2.4}\\
\(\underline{\Pic}_{X/k}\) & X & \hyperlink{sga6.e10.sec.2.4}{2.4}\\
\(P^n_{X/S}\) & II App. II & \hyperlink{sga6.e02.para.1-2-1}{1.2.1}\\
\(P^n_{X/S}(E)\) & II App. II & \hyperlink{sga6.e02.para.1-2-2}{1.2.2}\\
\(\underline{\operatorname{Proj}}_0\) & VI & \hyperlink{sga6.e06.para.2-3-2}{2.3.2}\\
\(\psi^k(x)\) & V & \hyperlink{sga6.e05.para.7-1}{7.1}\\
\(\Qcoh(S)\) & II & \hyperlink{sga6.e02.notation.2-2-1}{2.2.1}\\
\({}^r\underline{O}_{U\times V}\) & II App. II & \hyperlink{sga6.e02.sec.1.1}{1.1}\\
\(\rho_t\) & V & \hyperlink{sga6.e05.para.3-4}{3.4}\\
\(\circ\) & V & \hyperlink{sga6.e05.para.2-3}{2.3}\\
\(T_f\) & VIII & \hyperlink{sga6.e08.cor.2-5}{2.5}\\
\(\Todd\) & V & \hyperlink{sga6.e05.para.1-16}{1.16}\\
\(\toramp(F)\) & I & \hyperlink{sga6.e01.def.5-2}{5.2}\\
\(\toramp_f(E),\ \toramp_Y(E)\) & III & \hyperlink{sga6.e03.def.3-1}{3.1}\\
\(\Tr_{E|U}\) & I & \hyperlink{sga6.e01.para.8-1}{8.1}\\
\(T_{X/S}\) & II App. II & \hyperlink{sga6.e02.para.1-2-1}{1.2.1}\\
\(\tau\) & I & \hyperlink{sga6.e01.para.8-1}{8.1}\\
\(\nabla\) & V & \hyperlink{sga6.e05.para.3-5}{3.5}\\
%%%%%%%%%%  scan idx 713  |  folio 700  %%%%%%%%%%
\SGAPrintPage{700}%
\(1+A[[t]]^+\) & V & \hyperlink{sga6.e05.para.1-1}{1.1}\\
\(1+\widehat B^+\) & V & \hyperlink{sga6.e05.para.1-11}{1.11}\\
\(X(n)\) & X & \hyperlink{sga6.e10.prop.1-1-2}{1.1.2}\\
\(x\cdot_N y\) & V & \hyperlink{sga6.e05.para.5-5}{5.5}\\
\(\langle x,y\rangle\) & X & \hyperlink{sga6.e10.sec.4.1}{4.1}\\
\(Z(x)\) & X & \hyperlink{sga6.e10.prop.1-1-2}{1.1.2}\\
\bottomrule
\end{longtable}
\endgroup

%%%%%%%%%%  scan idx 714  |  folio rear-715  %%%%%%%%%%
\SGAPrintPage{rear-715}%
% Blank page: inspected in full; no printed text.

%%%%%%%%%%  scan idx 715  |  folio rear-716  %%%%%%%%%%
\SGAPrintPage{rear-716}%
% Blank page: inspected in full; no printed text.

%%%%%%%%%%  scan idx 716  |  folio rear-717  %%%%%%%%%%
\SGAPrintPage{rear-717}%
\section*{Lecture Notes in Mathematics}
\textbf{Comprehensive leaflet on request}

% Left column, volumes38–73.
Vol. 38: R. Berger, R. Kiehl, E. Kunz und H.-J. Nastold, Differentialrechnung in der analytischen Geometrie IV, 134 Seiten. 1967 DM 12,--

Vol. 39: Séminaire de Probabilités I. II, 189 pages. 1967. DM 14,--

Vol. 40: J. Tits, Tabellen zu den einfachen Lie Gruppen und ihren Darstellungen. VI, 53 Seiten. 1967. DM 6.80

Vol. 41: A. Grothendieck, Local Cohomology. VI, 106 pages. 1967. DM 10,--

Vol. 42: J. F. Berglund and K. H. Hofmann, Compact Semitopological Semigroups and Weakly Almost Periodic Functions. VI, 160 pages. 1967. DM 12,--

Vol. 43: D. G. Quillen, Homotopical Algebra. VI, 157 pages. 1967. DM 14,--

Vol. 44: K. Urbanik, Lectures on Prediction Theory. IV, 50 pages. 1967. DM 5,80

Vol. 45: A. Wilansky, Topics in Functional Analysis. VI, 102 pages. 1967. DM 9,60

Vol. 46: P. E. Conner, Seminar on Periodic Maps. IV, 116 pages. 1967. DM 10,60

Vol. 47: Reports of the Midwest Category Seminar I. IV, 181 pages. 1967. DM 14,80

Vol. 48: G. de Rham, S. Maumary et M. A. Kervaire, Torsion et Type Simple d'Homotopie. IV, 101 pages. 1967. DM 9,60

Vol. 49: C. Faith, Lectures on Injective Modules and Quotient Rings. XVI, 140 pages. 1967. DM 12,80

Vol. 50: L. Zalcman, Analytic Capacity and Rational Approximation. VI, 155 pages. 1968. DM 13.20

Vol. 51: Séminaire de Probabilités II. IV, 199 pages. 1968. DM 14,--

Vol. 52: D. J. Simms, Lie Groups and Quantum Mechanics. IV, 90 pages. 1968. DM 8,--

Vol. 53: J. Cerf, Sur les difféomorphismes de la sphère de dimension trois \((\Gamma_4=0)\). XII, 133 pages. 1968. DM 12,--

Vol. 54: G. Shimura, Automorphic Functions and Number Theory. VI, 69 pages. 1968. DM 8,--

Vol. 55: D. Gromoll, W. Klingenberg und W. Meyer, Riemannsche Geometrie im Großen. VI, 287 Seiten. 1968. DM 20,--

Vol. 56: K. Floret und J. Wloka, Einführung in die Theorie der lokalkonvexen Räume. VIII, 194 Seiten. 1968. DM 16,--

Vol. 57: F. Hirzebruch und K. H. Mayer, O (n)-Mannigfaltigkeiten, exotische Sphären und Singularitäten. IV, 132 Seiten. 1968. DM 10,80

Vol. 58: Kuramochi Boundaries of Riemann Surfaces. IV, 102 pages. 1968. DM 9,60

Vol. 59: K. Jänich, Differenzierbare G-Mannigfaltigkeiten. VI, 89 Seiten. 1968. DM 8,--

Vol. 60: Seminar on Differential Equations and Dynamical Systems. Edited by G. S. Jones. VI, 106 pages. 1968. DM 9,60

Vol. 61: Reports of the Midwest Category Seminar II. IV, 91 pages. 1968. DM 9,60

Vol. 62: Harish-Chandra, Automorphic Forms on Semisimple Lie Groups X, 138 pages. 1968. DM 14,--

Vol. 63: F. Albrecht, Topics in Control Theory. IV, 65 pages. 1968. DM 6,80

Vol. 64: H. Berens, Interpolationsmethoden zur Behandlung von Approximationsprozessen auf Banachräumen. VI, 90 Seiten. 1968. DM 8,--

Vol. 65: D. Kölzow, Differentiation von Maßen. XII, 102 Seiten. 1968. DM 8,--

Vol. 66: D. Ferus, Totale Absolutkrümmung in Differentialgeometrie und -topologie. VI, 85 Seiten. 1968. DM 8,--

Vol. 67: F. Kamber and P. Tondeur, Flat Manifolds. IV, 53 pages. 1968. DM 5,80

Vol. 68: N. Boboc et P. Mustată, Espaces harmoniques associés aux opérateurs différentiels linéaires du second ordre de type elliptique. VI, 95 pages. 1968. DM 8,60

Vol. 69: Seminar über Potentialtheorie. Herausgegeben von H. Bauer. VI, 180 Seiten. 1968. DM 14,80

Vol. 70: Proceedings of the Summer School in Logic. Edited by M. H. Löb. IV, 331 pages. 1968. DM 20,--

Vol. 71: Séminaire Pierre Lelong (Analyse), Année 1967 -- 1968. VI, 190 pages. 1968. DM 14,

Vol. 72: The Syntax and Semantics of Infinitary Languages. Edited by J. Barwise. IV, 268 pages. 1968. DM 18,--

Vol. 73: P. E. Conner, Lectures on the Action of a Finite Group. IV, 123 pages. 1968. DM 10,--

% Right column, volumes74–110.
Vol. 74: A. Fröhlich, Formal Groups. IV, 140 pages. 1968. DM 12,--

Vol. 75: G. Lumer, Algèbres de fonctions et espaces de Hardy. VI, 80 pages. 1968. DM 8,--

Vol. 76: R. G. Swan, Algebraic K-Theory. IV, 262 pages. 1968. DM 18,--

Vol. 77: P.-A. Meyer, Processus de Markov: la frontière de Martin. IV, 123 pages. 1968. DM 10,--

Vol. 78: H. Herrlich, Topologische Reflexionen und Coreflexionen. XVI, 166 Seiten. 1968. DM 12,--

Vol. 79: A. Grothendieck, Catégories Cofibrées Additives et Complexe Cotangent Relatif. IV, 167 pages. 1968. DM 12,--

Vol. 80: Seminar on Triples and Categorical Homology Theory. Edited by B. Eckmann. IV, 398 pages. 1969. DM 20,--

Vol. 81: J.-P. Eckmann et M. Guenin, Méthodes Algébriques en Mécanique Statistique. VI, 131 pages. 1969. DM 12,--

Vol. 82: J. Wloka, Grundräume und verallgemeinerte Funktionen. VIII, 131 Seiten. 1969. DM 12,--

Vol. 83: O. Zariski, An Introduction to the Theory of Algebraic Surfaces. IV, 100 pages. 1969. DM 8,--

Vol. 84: H. Lüneburg, Transitive Erweiterungen endlicher Permutationsgruppen. IV, 119 Seiten. 1969. DM 10.--

Vol. 85: P. Cartier et D. Foata, Problèmes combinatoires de commutation et réarrangements. IV, 88 pages. 1969. DM 8,--

Vol. 86: Category Theory, Homology Theory and their Applications I. Edited by P. Hilton. VI, 216 pages. 1969. DM 16,--

Vol. 87: M. Tierney, Categorical Constructions in Stable Homotopy Theory. IV, 65 pages. 1969. DM 6,--

Vol. 88: Séminaire de Probabilités III. IV, 229 pages. 1969. DM 18,--

Vol. 89: Probability and Information Theory. Edited by M. Behara, K. Krickeberg and J. Wolfowitz. IV, 256 pages. 1969. DM 18,--

Vol. 90: N. P. Bhatia and O. Hajek, Local Semi-Dynamical Systems. II, 157 pages. 1969. DM 14,--

Vol. 91: N. N. Janenko, Die Zwischenschrittmethode zur Lösung mehrdimensionaler Probleme der mathematischen Physik. VIII, 194 Seiten. 1969. DM 16,80

Vol. 92: Category Theory, Homology Theory and their Applications II. Edited by P. Hilton. V, 308 pages. 1969. DM 20,--

Vol. 93: K. R. Parthasarathy, Multipliers on Locally Compact Groups. III, 54 pages. 1969. DM 5,60

Vol. 94: M. Machover and J. Hirschfeld, Lectures on Non-Standard Analysis. VI, 79 pages. 1969. DM 6,--

Vol. 95: A. S. Troelstra, Principles of Intuitionism. II, 111 pages. 1969. DM 10,--

Vol. 96: H.-B. Brinkmann und D. Puppe, Abelsche und exakte Kategorien, Korrespondenzen. V, 141 Seiten. 1969. DM 10,--

Vol. 97: S. O. Chase and M. E. Sweedler, Hopf Algebras and Galois theory. II, 133 pages. 1969. DM 10,--

Vol. 98: M. Heins, Hardy Classes on Riemann Surfaces. III, 106 pages. 1969. DM 10,--

Vol. 99: Category Theory, Homology Theory and their Applications III. Edited by P. Hilton. IV, 489 pages. 1969. DM 24,--

Vol. 100: M. Artin and B. Mazur, Etale Homotopy. II, 196 Seiten. 1969. DM 12,--

Vol. 101: G. P. Szegö et G. Treccani, Semigruppi di Trasformazioni Multivoche. VI, 177 pages. 1969. DM 14,--

Vol. 102: F. Stummel, Rand- und Eigenwertaufgaben in Sobolewschen Räumen. VIII, 386 Seiten. 1969. DM 20,--

Vol. 103: Lectures in Modern Analysis and Applications I. Edited by C. T. Taam. VII, 162 pages. 1969. DM 12,--

Vol. 104: G. H. Pimbley, Jr., Eigenfunction Branches of Nonlinear Operators and their Bifurcations. II, 128 pages. 1969. DM 10,--

Vol. 105: R. Larsen, The Multiplier Problem. VII, 284 pages. 1969. DM 18,--

Vol. 106: Reports of the Midwest Category Seminar III. Edited by S. Mac Lane. III, 247 pages. 1969. DM 16,--

Vol. 107: A. Peyerimhoff, Lectures on Summability. III, 111 pages. 1969. DM 8,--

Vol. 108: Algebraic K-Theory and its Geometric Applications. Edited by R. M. F. Moss and C. B. Thomas. IV, 86 pages. 1969. DM 6,--

Vol. 109: Conference on the Numerical Solution of Differential Equations. Edited by J. Ll. Morris. VI, 275 pages. 1969. DM 18,--

Vol. 110: The Many Facets of Graph Theory. Edited by G. Chartrand and S. F. Kapoor. VIII, 290 pages. 1969. DM 18,--

%%%%%%%%%%  scan idx 717  |  folio rear-718  %%%%%%%%%%
\SGAPrintPage{rear-718}%
% Left column, volumes111–147.
Vol. 111: K. H. Mayer, Relationen zwischen charakteristischen Zahlen. III, 99 Seiten. 1969. DM 8,--

Vol. 112: Colloquium on Methods of Optimization. Edited by N. N. Moiseev. IV, 293 pages. 1970. DM 18,--

Vol. 113: R. Wille, Kongruenzklassengeometrien. III, 99 Seiten. 1970. DM 8,--

Vol. 114: H. Jacquet and R. P. Langlands, Automorphic Forms on GL (2). VII, 548 pages. 1970. DM 24,--

Vol. 115: K. H. Roggenkamp and V. Huber-Dyson, Lattices over Orders I. XIX, 290 pages. 1970. DM 18,--

Vol. 116: Séminaire Pierre Lelong (Analyse) Année 1969. IV, 195 pages. 1970. DM 14,--

Vol. 117: Y. Meyer, Nombres de Pisot, Nombres de Salem et Analyse Harmonique. 63 pages. 1970. DM 6.--

Vol. 118: Proceedings of the 15th Scandinavian Congress, Oslo 1968. Edited by K. E. Aubert and W. Ljunggren. IV, 162 pages. 1970. DM 12,--

Vol. 119: M. Raynaud, Faisceaux amples sur les schémas en groupes et les espaces homogènes. III, 219 pages. 1970. DM 14,--

Vol. 120: D. Siefkes, Büchi's Monadic Second Order Successor Arithmetic. XII, 130 Seiten. 1970. DM 12,--

Vol. 121: H. S. Bear, Lectures on Gleason Parts. III, 47 pages. 1970. DM 6,--

Vol. 122: H. Zieschang, E. Vogt und H.-D. Coldewey, Flächen und ebene diskontinuierliche Gruppen. VIII, 203 Seiten. 1970. DM 16,--

Vol. 123: A. V. Jategaonkar, Left Principal Ideal Rings. VI, 145 pages. 1970. DM 12,--

Vol. 124: Séminaire de Probabilités IV. Edited by P. A. Meyer. IV, 282 pages. 1970. DM 20,--

Vol. 125: Symposium on Automatic Demonstration. V, 310 pages. 1970. DM 20,--

Vol. 126: P. Schapira, Théorie des Hyperfonctions. XI, 157 pages. 1970. DM 14,--

Vol. 127: I. Stewart, Lie Algebras. IV, 97 pages. 1970. DM 10,--

Vol. 128: M. Takesaki, Tomita's Theory of Modular Hilbert Algebras and its Applications. II, 123 pages. 1970. DM 10,--

Vol. 129: K. H. Hofmann, The Duality of Compact Semigroups and C*-Bigebras. XII, 142 pages. 1970. DM 14,--

Vol. 130: F. Lorenz, Quadratische Formen über Körpern. II, 77 Seiten. 1970. DM 8,--

Vol. 131: A Borel et al., Seminar on Algebraic Groups and Related Finite Groups. VII, 321 pages. 1970. DM 22,--

Vol. 132: Symposium on Optimization. III, 348 pages. 1970. DM 22,--

Vol. 133: F. Topsøe, Topology and Measure. XIV, 79 pages. 1970. DM 8,--

Vol. 134: L. Smith, Lectures on the Eilenberg-Moore Spectral Sequence. VII, 142 pages. 1970. DM 14,--

Vol. 135: W. Stoll, Value Distribution of Holomorphic Maps into Compact Complex Manifolds. II, 267 pages. 1970. DM 18,--

Vol. 136: M. Karoubi et al., Séminaire Heidelberg-Saarbrücken-Strasbourg sur la K-Théorie. IV, 264 pages. 1970. DM 18,--

Vol. 137: Reports of the Midwest Category Seminar IV. Edited by S. MacLane. III, 139 pages. 1970. DM 12,--

Vol. 138: D. Foata et M. Schützenberger, Théorie Géométrique des Polynômes Eulériens. V, 94 pages. 1970. DM 10,--

Vol. 139: A. Badrikian, Séminaire sur les Fonctions Aléatoires Linéaires et les Mesures Cylindriques. VII, 221 pages. 1970. DM 18,--

Vol. 140: Lectures in Modern Analysis and Applications II. Edited by C. T. Taam. VI, 119 pages. 1970. DM 10,--

Vol. 141: G. Jameson, Ordered Linear Spaces. XV, 194 pages. 1970. DM 16,--

Vol. 142: K. W. Roggenkamp, Lattices over Orders II. V, 388 pages. 1970. DM 22,--

Vol. 143: K. W. Gruenberg, Cohomological Topics in Group Theory. XIV, 275 pages. 1970. DM 20,--

Vol. 144: Seminar on Differential Equations and Dynamical Systems, II. Edited by J. A. Yorke. VIII, 268 pages. 1970. DM 20,--

Vol. 145: E. J. Dubuc, Kan Extensions in Enriched Category Theory. XVI, 173 pages. 1970. DM 16,--

Vol. 146: A. B. Altman and S. Kleiman, Introduction to Grothendieck Duality Theory. II, 192 pages. 1970. DM 18,--

Vol. 147: D. E. Dobbs, Cech Cohomological Dimensions for Commutative Rings. VI, 176 pages. 1970. DM 16,--

% Right column, volumes148–181.
Vol. 148: R. Azencott, Espaces de Poisson des Groupes Localement Compacts. IX, 141 pages. 1970. DM 14,--

Vol. 149: R. G. Swan and E. G. Evans, K-Theory of Finite Groups and Orders. IV, 237 pages. 1970. DM 20,--

Vol. 150: Heyer, Dualität lokalkompakter Gruppen. XIII, 372 Seiten. 1970. DM 20,--

Vol. 151: M. Demazure et A. Grothendieck, Schémas en Groupes I. (SGA 3). XV, 562 pages. 1970. DM 24,--

Vol. 152: M. Demazure et A. Grothendieck, Schémas en Groupes II. (SGA 3). IX, 654 pages. 1970. DM 24,--

Vol. 153: M. Demazure et A. Grothendieck, Schémas en Groupes III. (SGA 3). VIII, 529 pages. 1970. DM 24,--

Vol. 154: A. Lascoux et M. Berger, Variétés Kähleriennes Compactes. VII, 83 pages. 1970. DM 8,--

Vol. 155: Several Complex Variables I, Maryland 1970. Edited by J. Horváth. IV, 214 pages. 1970. DM 18,--

Vol. 156: R. Hartshorne, Ample Subvarieties of Algebraic Varieties. XIV, 256 pages. 1970. DM 20,--

Vol. 157: T. tom Dieck, K. H. Kamps und D. Puppe, Homotopietheorie. VI, 265 Seiten. 1970. DM 20,--

Vol. 158: T. G. Ostrom, Finite Translation Planes. IV. 112 pages. 1970. DM 10,--

Vol. 159: R. Ansorge und R. Hass. Konvergenz von Differenzenverfahren für lineare und nichtlineare Anfangswertaufgaben. VIII, 145 Seiten. 1970. DM 14,--

Vol. 160: L. Sucheston, Contributions to Ergodic Theory and Probability. VII, 277 pages. 1970. DM 20,--

Vol. 161: J. Stasheff, H-Spaces from a Homotopy Point of View. VI, 95 pages. 1970. DM 10,--

Vol. 162: Harish-Chandra and van Dijk, Harmonic Analysis on Reductive p-adic Groups. IV, 125 pages. 1970. DM 12,--

Vol. 163: P. Deligne, Equations Différentielles à Points Singuliers Reguliers. III, 133 pages. 1970. DM 12,--

Vol. 164: J. P. Ferrier, Seminaire sur les Algebres Complètes. II, 69 pages. 1970. DM 8,--

Vol. 165: J. M. Cohen, Stable Homotopy. V, 194 pages. 1970. DM 16,--

Vol. 166: A. J. Silberger, \(\mathrm{PGL}_2\) over the p-adics: its Representations, Spherical Functions, and Fourier Analysis. VII, 202 pages. 1970. DM 18,--

Vol. 167: Lavrentiev, Romanov and Vasiliev, Multidimensional Inverse Problems for Differential Equations. V, 59 pages. 1970. DM 10,--

Vol. 168: F. P. Peterson, The Steenrod Algebra and its Applications: A conference to Celebrate N. E. Steenrod's Sixtieth Birthday. VII, 317 pages. 1970. DM 22,--

Vol. 169: M. Raynaud, Anneaux Locaux Henséliens. V, 129 pages. 1970. DM 12,--

Vol. 170: Lectures in Modern Analysis and Applications III. Edited by C. T. Taam. VI, 213 pages. 1970. DM 18,--

Vol. 171: Set-Valued Mappings, Selections and Topological Properties of \(2^X\). Edited by W. M. Fleischman. X, 110 pages. 1970. DM 12,--

Vol. 172: Y.-T. Siu and G. Trautmann, Gap-Sheaves and Extension of Coherent Analytic Subsheaves. V, 172 pages. 1971. DM 16,--

Vol. 173: J. N. Mordeson and B. Vinograde, Structure of Arbitrary Purely Inseparable Extension Fields. IV, 138 pages. 1970. DM 14,--

Vol. 174: B. Iversen, Linear Determinants with Applications to the Picard Scheme of a Family of Algebraic Curves. VI, 69 pages. 1970. DM 8,--

Vol. 175: M. Brelot, On Topologies and Boundaries in Potential Theory. VI, 176 pages. 1971. DM 18,--

Vol. 176: H. Popp, Fundamentalgruppen algebraischer Mannigfaltigkeiten. IV, 154 Seiten. 1970. DM 16,--

Vol. 177: J. Lambek, Torsion Theories, Additive Semantics and Rings of Quotients. VI, 94 pages. 1971. DM 12,--

Vol. 178: Th. Bröcker und T. tom Dieck, Kobordismentheorie. XVI, 191 Seiten. 1970. DM 18,--

Vol. 179: Seminaire Bourbaki -- vol. 1968/69. Exposés 347-363. IV. 295 pages. 1971. DM 22,--

Vol. 180: Séminaire Bourbaki -- vol. 1969/70. Exposés 364-381. IV, 310 pages. 1971. DM 22,--

Vol. 181: F. DeMeyer and E. Ingraham, Separable Algebras over Commutative Rings. V, 157 pages. 1971. DM 16.--

%%%%%%%%%%  scan idx 718  |  folio rear-719  %%%%%%%%%%
\SGAPrintPage{rear-719}%
% Left column, volumes182–216.
Vol. 182: L. D. Baumert, Cyclic Difference Sets. VI, 166 pages. 1971. DM 16,--

Vol. 183: Analytic Theory of Differential Equations. Edited by P. F. Hsieh and A. W. J. Stoddart. VI, 225 pages. 1971. DM 20,--

Vol. 184: Symposium on Several Complex Variables, Park City, Utah, 1970. Edited by R. M. Brooks. V, 234 pages. 1971. DM 20,--

Vol. 185: Several Complex Variables II, Maryland 1970. Edited by J. Horváth. III, 287 pages. 1971. DM 24,--

Vol. 186: Recent Trends in Graph Theory. Edited by M. Capobianco/J. B. Frechen/M. Krolik. VI, 219 pages. 1971. DM 18.--

Vol. 187: H. S. Shapiro, Topics in Approximation Theory. VIII, 275 pages. 1971. DM 22,--

Vol. 188: Symposium on Semantics of Algorithmic Languages. Edited by E. Engeler. VI, 372 pages. 1971. DM 26,--

Vol. 189: A. Weil, Dirichlet Series and Automorphic Forms. V, 164 pages. 1971. DM 16,--

Vol. 190: Martingales. A Report on a Meeting at Oberwolfach, May 17-23, 1970. Edited by H. Dinges. V, 75 pages. 1971. DM 12,--

Vol. 191: Séminaire de Probabilités V. Edited by P. A. Meyer. IV, 372 pages. 1971. DM 26,--

Vol. 192: Proceedings of Liverpool Singularities -- Symposium I. Edited by C. T. C. Wall. V, 319 pages. 1971. DM 24,--

Vol. 193: Symposium on the Theory of Numerical Analysis. Edited by J. Ll. Morris. VI, 152 pages. 1971. DM 16,--

Vol. 194: M. Berger, P. Gauduchon et E. Mazet. Le Spectre d'une Variété Riemannienne. VII, 251 pages. 1971. DM 22,--

Vol. 195: Reports of the Midwest Category Seminar V. Edited by J. W. Gray and S. Mac Lane. III, 255 pages. 1971. DM 22,--

Vol. 196: H-spaces -- Neuchâtel (Suisse)- Août 1970. Edited by F. Sigrist, V, 156 pages. 1971. DM 16,--

Vol. 197: Manifolds -- Amsterdam 1970. Edited by N. H. Kuiper. V, 231 pages. 1971. DM 20,--

Vol. 198: M. Hervé, Analytic and Plurisubharmonic Functions in Finite and infinite Dimensional Spaces. VI, 90 pages. 1971. DM 16.--

Vol. 199: Ch. J. Mozzochi, On the Pointwise Convergence of Fourier Series. VII, 87 pages. 1971. DM 16,--

Vol. 200: U. Neri, Singular Integrals. VII, 272 pages. 1971. DM 22,--

Vol. 201: J. H. van Lint, Coding Theory. VII, 136 pages. 1971. DM 16,--

Vol. 202: J. Benedetto, Harmonic Analysis on Totally Disconnected Sets. VIII, 261 pages. 1971. DM 22,--

Vol. 203: D. Knutson, Algebraic Spaces. VI, 261 pages. 1971. DM 22,--

Vol. 204: A. Zygmund, Intégrales Singulières. IV, 53 pages. 1971. DM 12,--

Vol. 205: Séminaire Pierre Lelong (Analyse) Année 1970. VI, 243 pages. 1971. DM 20,--

Vol. 206: Symposium on Differential Equations and Dynamical Systems. Edited by D. Chillingworth. XI, 173 pages. 1971. DM 16,--

Vol. 207: L. Bernstein, The Jacobi-Perron Algorithm -- Its Theory and Application. IV, 161 pages. 1971. DM 16,--

Vol. 208: A. Grothendieck and J. P. Murre, The Tame Fundamental Group of a Formal Neighbourhood of a Divisor with Normal Crossings on a Scheme. VIII, 133 pages. 1971. DM 16,--

Vol. 209: Proceedings of Liverpool Singularities Symposium II. Edited by C. T. C. Wall. V, 280 pages. 1971. DM 22,--

Vol. 210: M. Eichler, Projective Varieties and Modular Forms. III, 118 pages. 1971. DM 16,--

Vol. 211: Théorie des Matroïdes. Edité par C. P. Bruter. III, 108 pages. 1971. DM 16,--

Vol. 212: B. Scarpellini, Proof Theory and Intuitionistic Systems. VII, 291 pages. 1971. DM 24,--

Vol. 213: H. Hogbe-Nlend, Théorie des Bornologies et Applications. V, 168 pages. 1971. DM 18,--

Vol. 214: M. Smorodinsky, Ergodic Theory, Entropy. V, 64 pages. 1971. DM 12,--

Vol. 215: P. Antonelli, D. Burghelea and P. J. Kahn, The Concordance-Homotopy Groups of Geometric Automorphism Groups. X, 140 pages. 1971. DM 16,--

Vol. 216: H. Maaß, Siegel's Modular Forms and Dirichlet Series. VII, 328 pages. 1971. DM 20,--

% Right column, volumes217–225.
Vol. 217: T. J. Jech, Lectures in Set Theory with Particular Emphasis on the Method of Forcing. V, 137 pages. 1971. DM 16,--

Vol. 218: C. P. Schnorr, Zufälligkeit und Wahrscheinlichkeit. IV, 212 Seiten 1971. DM 20,--

Vol. 219: N. L. Alling and N. Greenleaf, Foundations of the Theory of Klein Surfaces. IX, 117 pages. 1971. DM 16,--

Vol. 220: W. A. Coppel, Disconjugacy. V, 148 pages. 1971. DM 16,--

Vol. 221: P. Gabriel und F. Ulmer, Lokal präsentierbare Kategorien. V, 200 Seiten. 1971. DM 18,--

Vol. 222: C. Meghea, Compactification des Espaces Harmoniques. III, 108 pages. 1971. DM 16,--

Vol. 223: U. Felgner, Models of ZF-Set Theory. VI, 173 pages. 1971. DM 16,--

Vol. 224: Revêtements Etales et Groupe Fondamental. (SGA 1). Dirigé par A. Grothendieck XXII, 447 pages. 1971. DM 30,--

Vol. 225: Théorie des Intersections et Théorème de Riemann-Roch. (SGA 6). Dirigé par P. Berthelot, A. Grothendieck et L. Illusie. XII, 700 pages. 1971. DM 40,--

%%%%%%%%%%  scan idx 719  |  folio rear-720  %%%%%%%%%%
\SGAPrintPage{rear-720}%
ISBN 3-540-05647-5\par
ISBN 0-387-05647-5

\end{document}
